% Reader-facing source curated from the forbidden-pc-minor campaign.
% Edit this file directly; the historical extraction script preserves it.

\section{Basic operations and structural decompositions}
\label{sec:definitions}

\subsection{Definitions and elementary operations}
\label{sec:damp-basic-operations}

All graphs are finite partial cubes.  A cube embedding is
\emph{irredundant} if every coordinate occurs on an edge; its number
of coordinates is the \emph{isometric dimension}.  We identify each
graph with an irredundant isometric family
\(X\subseteq Q_E=\{0,1\}^E\).  For \(e\in E\) and
\(b\in\{0,1\}\), write
\[
 \rho_e^bX
 =\{x|_{E\setminus\{e\}}:x\in X,\ x_e=b\}
\]
for the restriction to one \(e\)-semicube, with the constant coordinate
removed, and write
\[
 \pi_eX=\rho_e^0X\cup\rho_e^1X
\]
for contraction of the complete \(e\)-coordinate class.  Graphically,
the contraction identifies both endpoints of every edge in that
\(\Theta\)-class.  A \emph{pc-minor} is obtained by a finite sequence of
nonempty convex restrictions and coordinate contractions
\cite[Section~2.1]{KnauerMarc20}; it is \emph{proper} if it is not
isomorphic to the original partial cube.

For \(A\subseteq E\), the \emph{reduction} along \(A\) is
\[
 X^A=\{z\in Q_{E\setminus A}:
             \{z\}\times Q_A\subseteq X\}.
\]
In particular,
\(X^{\{e\}}=\rho_e^0X\cap\rho_e^1X\).
Unlike a restriction or contraction, a reduction need not be a
pc-minor.

CCMW call ordinary coordinate projection a \emph{restriction} and the
full-fibre operation above a \emph{reduction}. We use reduction throughout.
For pc-minors we retain the graph terms \emph{convex restriction} and
\emph{coordinate contraction}: $\rho_e^b$ is a coordinate-face intersection,
whereas $\pi_e$ is CCMW's restriction. Older source labels and verifier
filenames using the historical name refer to this same operation $X^A$.

For \(S\subseteq E\), the family \(X\) \emph{shatters} \(S\) if its
projection onto \(S\) is all of \(Q_S\).  It \emph{strongly shatters}
\(S\) if \(X\) contains a full \(S\)-dimensional cube obtained by fixing
all coordinates outside \(S\).  The family is \emph{ample} if its
shattered and strongly shattered coordinate sets coincide
\cite{BandeltCDK06}.
The Sandwich theorem states that every finite cube family satisfies
\[
 |\operatorname{SSh}(X)|\le |X|\le |\operatorname{Sh}(X)|,
\]
where \(\operatorname{SSh}(X)\) and \(\operatorname{Sh}(X)\) denote the
strongly shattered and shattered sets, respectively.  Equality in either
inequality holds exactly when \(X\) is ample
\cite[Corollary~1 and Theorem~2]{BandeltCDK06}.  We refer to these
cardinality statements as the \emph{Sandwich inequalities} and to the
ample case as the \emph{Sandwich equality}.

A vertex is a \emph{corner} if it belongs to a unique maximal cube.  A
\emph{corner peeling}, written in deletion order as
\(x_1,\ldots,x_m\), is an ordering in which \(x_i\) is a corner of
\(X\setminus\{x_1,\ldots,x_{i-1}\}\) for every \(i\).  Its reversal is
the prefix-building convention used for ample and isometric orders in
\cite[Proposition~4.2]{ChalopiChepoiMoran22}.  We use \(\Damp\) for the
ample families admitting such an ordering.

For a vertex $x$ and $A\subseteq E$, write $x^A$ for the vertex
obtained by flipping the coordinates in $A$.  Put
\[
 I_X(x)=\{e\in E:x^{\{e\}}\in X\},\qquad
 \partial_e^bX=\{x\in X:x_e=b,\ x^{\{e\}}\notin X\}.
\]
Thus $I_X(x)$ is the set of edge directions incident with $x$,
and $\partial_e^bX$ is the signed $e$-boundary on side $b$.
Write $\operatorname{Cor}(X)$ for the set of corners of $X$.
If $x$ is a corner, $Q_X(x)$ denotes its unique maximal cube.
Finally,
\[
 Q_n^{--}=Q_n\setminus\{x,\bar x\}
\]
denotes an $n$-cube with one antipodal pair removed; its
isomorphism type is independent of the choice of $x$.

\subsection{Representation maps and boundary conditions}
\label{sec:damp-uso-boundaries}

The orientation formulation identifies which information a gluing argument
must preserve. An orientation $D$ of the one-inclusion graph of an ample
$H\subseteq Q_F$ has \emph{out-map}
\[
 r_D(x)=\{f\in F:x^{\{f\}}\in H,\ x\longrightarrow x^{\{f\}}\}.
\]
Following Chalopin--Chepoi--Moran--Warmuth (CCMW), it is a
\emph{unique-sink orientation} (USO) if both conditions hold:
\begin{enumerate}
 \item[(C1)] $\{x^J:J\subseteq r_D(x)\}\subseteq H$ for every $x\in H$;
 \item[(C2)] every contained cube has exactly one sink in its induced orientation.
\end{enumerate}
A \emph{representation map} is a bijection
$r:H\to\operatorname{Sh}(H)$ such that distinct $x,y$ differ on at least
one coordinate of $r(x)\cup r(y)$. CCMW's Theorems~6.1 and~6.8 identify
representation maps with the out-maps of USOs. Their Proposition~6.12
identifies $H\in\Damp$ with existence of an acyclic USO
\cite{ChalopiChepoiMoran22}. Every topological order of that orientation
is a corner deletion order. We use these established results throughout.

The two USO conditions must be kept together. On the path
$\{00,01,11\}$, directing both edges out of $01$ satisfies (C2), since
only vertices and edges are contained cubes, but violates (C1): its
outgoing square would require $10$. Directing the path
$00\to01\to11$ satisfies both and realizes sink $11$.

Write $\mathcal A(H)$ for the set of acyclic USOs of $H$. For an ample
family, forbiddenness is the condition $\mathcal A(H)=\varnothing$ while
$\mathcal A(M)\ne\varnothing$ for every nonempty proper pc-minor $M$.
The witnesses on different minors are chosen independently. Positivity
of a minor does not require that every representation map on it be acyclic.
Existence of a representation map on every ample class is not assumed.

For nonempty ample $T\subseteq H$, say that $H$ \emph{peels to $T$}
when corners outside $T$ can be deleted until exactly $T$ remains.
An \emph{inward boundary at $T$} means that every edge between
$H\setminus T$ and $T$ is directed into $T$.

\begin{proposition}[Relative peeling and extension of a representation map]
\label{prop:damp-relative-uso-extension}
Suppose a USO $D_T$ on $T$ is given. Then $H$ peels to $T$ if and only
if $D_T$ extends to a USO on $H$ with inward boundary at $T$ and an
acyclic induced orientation on $H\setminus T$. If $T\in\Damp$, relative
peeling is equivalent to existence of an acyclic USO on $H$ with inward
boundary at $T$. Every prescribed acyclic USO of $T$ then extends.
\end{proposition}
\begin{proof}
Reverse the relative peeling, adding each corner as a source. Old outgoing
cubes are unchanged and the new source has its corner cube. If a contained
cube $B$ contains the new source $v$ and has more than one vertex, its old
part $B\setminus\{v\}$ is a coordinate-face intersection of the previous
ample family. Its induced orientation is a USO by CCMW Lemma~6.9 and has
one sink: by Theorem~6.8 its representation map takes the value
$\varnothing$ exactly once. Adding the source preserves this sink, so
(C2) holds as well. The added vertices are acyclic and the boundary enters $T$.

Conversely, take a source of the remaining outside acyclic digraph. The
inward boundary makes it a source of the whole remaining USO. By CCMW
Proposition~A.2, deleting this singleton source component leaves an ample
class with the restricted representation map; the source is a corner by
(C1). Repeating deletes precisely the outside vertices. If $D_T$ is
acyclic, so is the whole extension, since a cycle cannot enter $T$ and leave.
\end{proof}

For $T=\{a\}$, this is existence of an acyclic USO with prescribed sink
$a$. For a nonpeelable retained target, the extension may have cycles inside
$T$. Thus relative peeling must not be replaced unconditionally by existence
of an acyclic USO on the whole family. The rank-two invariance theorem below
also covers targets with no supplied representation map and the empty target;
its deletion formulation preserves that full scope.

\begin{proposition}[A quadratic form of metric elimination]
\label{prop:damp-quadratic-orientation-score}
Let $H$ be a nonempty isometric cube family, let $f_2(H)$ be its number
of contained squares, and let $O$ be an acyclic orientation of its
one-inclusion graph. Write $d_O^+(v)$ for the outdegree of $v$ and set
\[
 \varepsilon_H(O)=\sum_{v\in H}\binom{d_O^+(v)}2-f_2(H).
\]
Then $\varepsilon_H(O)\ge0$. Equality holds exactly when a topological
order of $O$ is a corner peeling; in that case every topological order
is a corner peeling and $H$ is ample. In particular, for ample $H$,
\[
 H\in\Damp\quad\Longleftrightarrow\quad
 \min_{O\text{ acyclic}}\varepsilon_H(O)=0.
\]
The minimum over acyclic orientations in which $a$ is a sink is zero
exactly when $H$ peels to $a$.
\end{proposition}
\begin{proof}
An outgoing pair at $v$ consists of two edges $v\to u,v\to w$.
Either their other square vertex $z=u\mathbin\oplus v\mathbin\oplus w$
is absent, or the pair is a source at $v$ of a contained square.
Every acyclically oriented square has one or two sources. Consequently
$\varepsilon_H(O)$ is exactly the number of outgoing pairs with absent
other square vertex, plus the number of squares with two sources.
This proves nonnegativity.

Take a topological order, in deletion order. If the score is zero, any
two remaining neighbours $u,w$ of the next vertex $v$ have their other
square vertex $z$ in $H$. Moreover $z$ has not been deleted: otherwise
both $v$ and $z$ would precede $u,w$, giving two sources in that square.
Thus deleting $v$ preserves every distance, by replacing any two-edge
segment through $v$ by the surviving two-edge segment through $z$.
Every suffix is isometric. Proposition~\ref{prop:damp-yang-shellability-metric-elimination}
identifies this order with a corner peeling, and also proves ampleness
without assuming it initially. This reasoning applies to every topological
order. Conversely a corner order supplies a surviving $z$ for every outgoing
pair, so neither counted defect occurs and the score is zero.
Finally an acyclic orientation with sink $a$ has a topological order ending
at $a$; the preceding equivalence proves the rooted assertion.
\end{proof}

The proposition is a scalar reformulation of the existing metric-elimination
criterion, not a new characterization of its generable inputs. Its outgoing
pair count is the source version of the two-frame objective used in the
abstract-objective-function literature
\cite[Theorems~1 and~5]{JoswigKaibelKorner2002KSystems}.
The cited two-face theorem can be applied to each contained cube after
acyclicity and (C1) have been imposed. The metric argument above is stronger
for this purpose: zero score already implies (C1), so no separate
higher-dimensional outgoing-cube test is required. Acyclicity remains
essential; a directed square has score $-1$.

For the fixed rooted-negative $292$-vertex control, the stored order ending
at its root has score exactly one. Its only defect is the outgoing pair
\[
 1116\longrightarrow1628,\qquad1116\longrightarrow92,
 \qquad 604=1116\mathbin\oplus1628\mathbin\oplus92\notin H.
\]
Every contained square has one source. The existing exhaustive certificate
that no corner order ends at the root excludes score zero, so the rooted
minimum is exactly one. The earlier relaxed order had score $17$; this
was the value of that order, not a lower bound. The rooted-positive
$294$-vertex control has minimum zero, witnessed by its existing order.
\path{research/quadratic_orientation_score.py/json} replays both bounds,
the defect count, and the old rooted-negative certificate.
These examples distinguish the two rooted problems numerically; they do
not imply that every rooted factor has minimum one, that local score descent
finds a minimum, or that a score-one orientation alone certifies
nonpeelability. The lower bound here still uses the earlier negative proof.

\begin{remark}[The score can increase under last-fibre contraction]
\label{ex:damp-quadratic-score-contraction}
An order on a cube family induces an acyclic orientation by directing each
edge toward its later endpoint. Under contraction of coordinate $j$, order
the image vertices by the positions of their last preimages. This operation
does not preserve an upper bound of one on the score. Identify the vertices
of $Q_4$ with the integers $0,\ldots,15$ in binary, with coordinate $0$ the
least significant bit. The order
\[
 14,6,10,12,8,4,9,0,1,11,3,7,2,15,5,13
\]
has score one: among its $24$ squares, precisely $\{2,3,6,7\}$ has two
sources (namely $6$ and $3$). Contracting coordinate $0$, and retaining the
zero bit in the image labels, gives the last-fibre order
\[
 8,0,10,6,2,14,4,12.
\]
This is an order on $Q_3$. Two of its six squares have two sources:
$\{0,2,4,6\}$, with sources $0,6$, and $\{2,6,10,14\}$, with sources
$6,10$. The score is therefore two. Equivalently, the outgoing-pair totals
are $25$ and $8$, respectively. All other squares have one source, and
there are no missing diagonals in either full cube.

Both cubes admit corner orders ending at the displayed final vertex,
so their optimized rooted scores are both zero. Thus this example refutes
nonincrease for the transported order, not for the minimum over all orders.
The zero-score transport supplied by corner-peeling minor closure is
unaffected. The orders and the independent square-source counts are replayed
in \path{research/quadratic_score_transport.py/json}.
\end{remark}

\subsection{Closure and elementary constructions}
\label{sec:damp-closure-constructions}

A coordinate $e$ is \emph{peripheral} if, after orienting it,
$\rho_e^1X\subseteq\rho_e^0X$.  Equivalently,
\[
 X=(B\times\{0\})\cup(S\times\{1\}),\qquad S\subseteq B,
\]
where $B=\pi_eX$ is the contraction base and $S$ is the
duplicated expansion site.
The family is \emph{periphery-free} if it has no peripheral
coordinate, and \emph{Cartesian-prime} if it is not a Cartesian
product of two nontrivial cube families.  We write
\[
 \Obs(\Damp)=\{O:O\notin\Damp:
 \text{ every proper pc-minor of }O\text{ belongs to }\Damp\}.
\]
We use that $\Damp$ is pc-minor closed and that Cartesian products
satisfy $A\square B\in\Damp$ if and only if
$A,B\in\Damp$; the product statement is proved in
Proposition~\ref{prop:damp-exact-expansion-product-tests}.

\begin{proposition}[Pc-minor closure]
\label{prop:damp-pc-minor-closure}
The class $\Damp$ is closed under pc-minors.
\end{proposition}

\begin{proof}
Reverse a corner-peeling order to obtain an ample prefix order
\(x_1,\ldots,x_m\), by
\cite[Proposition~4.2]{ChalopiChepoiMoran22}, and put
\(X_i=\{x_1,\ldots,x_i\}\).  Intersecting these prefixes with a
coordinate semicube, omitting empty sets and repetitions, gives an ample
prefix order of the nonempty restriction: every retained set is an
intersection of an ample class with a cube.

For contraction along \(e\), list distinct projected concepts in order
of their first occurrence.  If the \(j\)-th new projected concept first
occurs at index \(i_j\), the first \(j\) concepts of the new order form
exactly
\[
 \pi_eX_{i_j}=\rho_e^0X_{i_j}\cup\rho_e^1X_{i_j}.
\]
This is the projection of an ample prefix and is ample by
\cite[Section~3.1]{ChalopiChepoiMoran22}.  It is not the reduction \(X_{i_j}^{\{e\}}\), which uses the intersection of the
sections.  Thus both operations give ample prefix orders.  Reversing
them gives corner peelings of the elementary pc-minors.  Nonempty convex
restrictions are intersections of coordinate semicubes, and restrictions
and contractions generate the pc-minor relation, proving the claim.
\end{proof}

\begin{remark}[Multi-connectivity and game strategy classes]
\label{rem:damp-maat-multiconnected}
Maat's recursive notion of a \emph{multi-connected} binary class
\cite[Definition~4.6]{Maat25Strategy} gives another description of
$\Damp$. In that definition the empty class is allowed; a nonempty class
must admit a one-vertex deletion to a multi-connected class, and in every
cube face both the class and its complement must induce connected graphs.
Here connectivity is a pairwise condition, so empty intersections are allowed.
For nonempty binary classes, multi-connectivity is equivalent to membership
in $\Damp$.

Indeed, every class in the recursive deletion chain is ample by
\cite[Lemma~4.10]{Maat25Strategy}. A one-vertex deletion between ample
classes deletes a corner \cite[Lemma~4.1]{ChalopiChepoiMoran22}, giving a
corner peeling. Conversely, every remainder of a corner peeling is ample.
An ample class and its complement have connected intersections with every
cube face. Induction along the peeling therefore verifies all the conditions
of multi-connectivity.

Consequently the hypercube-realizable classes of
\cite[Definition~4.2 and Lemma~4.8]{Maat25Strategy} belong to $\Damp$.
This includes the binary strategy classes for the games specified in
\cite[Definition~4.4 and Lemma~4.5]{Maat25Strategy}. These constructions
supply positive inputs, not forbidden obstructions; their pc-minors remain
positive by the closure theorem above. No such assertion is made for
arbitrary subfamilies. The known non-corner-peelable ample examples also
answer Question~5.1 of that preprint negatively, via this dictionary;
no new counterexample is needed.
\end{remark}

\begin{proposition}[Reduction closure]
\label{prop:damp-strong-projection-closure}
If \(X\in\Damp\), then for every \(A\subseteq E\), the reduction
\(X^A\) is either empty or belongs to \(\Damp\).
\end{proposition}

\begin{proof}
If \(X^A=\varnothing\), there is nothing to prove.  Otherwise, reverse a corner peeling of \(X\) to obtain an ample prefix order
\(x_1,\ldots,x_m\), and put \(X_i=\{x_1,\ldots,x_i\}\).
For \(z\in X^A\), define its completion time to be the largest index of a
concept in the full fibre \(\{z\}\times Q_A\).  Order the members of
\(X^A\) by increasing completion time.  Distinct full fibers are disjoint,
so their completion times are distinct.  At every completion time \(i\),
the resulting prefix is exactly \(X_i^A\).  Reductions of ample
families are ample \cite[Section~3.1]{ChalopiChepoiMoran22}, so this is an
ample prefix order of \(X^A\).  Reversing it gives a corner peeling.
\end{proof}

\begin{proposition}[Exact peripheral-expansion and product tests]
\label{prop:damp-exact-expansion-product-tests}
Let \(B,S\subseteq Q_F\) be nonempty ample families with \(S\subseteq
B\), and let \(e\) be a new coordinate.  Then
\[
 Y=(B\times\{0\})\cup(S\times\{1\})
 \quad\text{satisfies}\quad
 Y\in\Damp\ \Longleftrightarrow\ B,S\in\Damp.
 \tag{A4g1a}\label{eq:damp-exact-peripheral-expansion-test}
\]
If \(A\) and \(B\) are nonempty ample families on disjoint coordinate
sets, then
\[
 A\square B\in\Damp
 \quad\Longleftrightarrow\quad
 A,B\in\Damp.
 \tag{A4g1b}\label{eq:damp-exact-product-test}
\]
In particular, peripheral expansions and Cartesian products preserve
non-corner-peelability as soon as one section or factor is
non-corner-peelable.
\end{proposition}

\begin{proof}
A set of old coordinates is shattered by \(Y\) exactly when it is
shattered by \(B\), while a set containing \(e\) is shattered exactly
when its old-coordinate part is shattered by \(S\).  Hence
\[
 |\operatorname{Sh}(Y)|
 =|\operatorname{Sh}(B)|+|\operatorname{Sh}(S)|
 =|B|+|S|=|Y|,
\]
so \(Y\) is ample.  If \(Y\in\Damp\), its two signed
\(e\)-restrictions show, by pc-minor closure, that \(B,S\in\Damp\).
Conversely, take a prefix-isometric order of \(B\times\{0\}\), followed
by the lift to the \(1\)-side of a prefix-isometric order of \(S\).  The
old layer is isometric.  Two vertices in a partial new layer have a
geodesic supplied by the corresponding prefix of \(S\), while a geodesic
from the old layer to a new vertex \((s,1)\) may first reach \((s,0)\)
inside \(B\times\{0\}\) and then use the \(e\)-edge.  Thus every prefix
is isometric, proving
\eqref{eq:damp-exact-peripheral-expansion-test}.

The product of two ample families is ample.  If \(A\square B\in\Damp\),
fixing a vertex in either factor realizes the other factor as a face
restriction, so both factors lie in \(\Damp\).  Conversely, let
\(a_1,\ldots,a_p\) and \(b_1,\ldots,b_q\) be prefix-isometric orders of
\(A\) and \(B\), respectively, and order the product fiber by fiber:
\[
 (a_1,b_1),\ldots,(a_1,b_q),
 (a_2,b_1),\ldots,(a_2,b_q),\ldots,
 (a_p,b_1),\ldots,(a_p,b_q).
\]
A prefix ending in the \(j\)-th vertex of the \(i\)-th fiber is
\[
 \bigl(\{a_1,\ldots,a_{i-1}\}\square B\bigr)
 \cup
 \bigl(\{a_i\}\square\{b_1,\ldots,b_j\}\bigr).
\]
Two vertices in the last partial fiber have a geodesic there.  For a
vertex in a completed fiber and one in the partial fiber, first move
inside the completed \(B\)-fiber to the second vertex's
\(B\)-coordinate, and then follow an \(A\)-geodesic with that coordinate
fixed.  This is a product geodesic contained in the prefix.  Hence the
displayed order is prefix-isometric and proves
\eqref{eq:damp-exact-product-test}.
\end{proof}

\begin{lemma}[Cartesian primality]
\label{lem:prime}
Every member of $\Obs(\Damp)$ is Cartesian-prime.
\end{lemma}

\begin{proof}
If $O=A\square B$ with both factors nontrivial, then $A$ and $B$
are proper convex subgraphs of $O$ and hence belong to $\Damp$.
The product equivalence then gives $O\in\Damp$, a contradiction.
\end{proof}

\begin{proposition}[The nonample layer]
\label{prop:nonample-layer}
The nonample part of the forbidden basis is
\[
 \Obs(\Damp)\setminus\Ample
   =\{Q_n^{--}:n\ge3\}.
\]
\end{proposition}

\begin{proof}
The elementary restrictions of $Q_n^{--}$ are singly punctured cubes and
its elementary contractions are hypercubes.  These families admit corner
peelings, whereas $Q_n^{--}$ is not ample, so $Q_n^{--}$ is an obstruction.
Conversely, the forbidden-pc-minor characterization of ampleness
\cite[Corollary~7.4]{KnauerMarc20} gives such a pc-minor in every
nonample family.  Minimality therefore forces $O\cong Q_n^{--}$.
\end{proof}

\begin{proposition}[Protected missing cubes force cycles]
\label{prop:damp-protected-uso-cycle}
Let $H\subseteq Q_F$ be ample and let $S\subseteq H$ with $|S|\ge2$.
Suppose for each $x\in S$ there is $J_x\subseteq F$ such that
$x^{\{e\}}\in S$ for every $e\in J_x$, but $x^{J_x}\notin H$.
Every USO of $H$, if one exists, has a directed cycle contained in $S$.
In particular, $H\notin\Damp$.
\end{proposition}
\begin{proof}
Not all directions in $J_x$ can be outgoing at $x$, by (C1).
Thus every vertex of the finite digraph induced on $S$ has an incoming
neighbour in $S$. Following predecessors until a vertex repeats yields
a directed cycle. CCMW's acyclic-USO characterization gives the conclusion.
\end{proof}

This is also a first-deletion certificate: no vertex of $S$ can be the
first one removed while all its specified neighbours remain. The statement
quantifies over every USO but does not assert a fixed common cycle, existence
of any USO, or that $S$ is a source strongly connected component. CCMW
Proposition~A.2 deletes source components of a chosen representation map;
that map-dependent operation is not a pc-minor operation.

\subsection{One-concept additions and square deletions}
\label{sec:damp-local-extension-foundations}

We collect the elementary steps used in the core-and-attachment
construction.  The local test identifies which missing concepts may be
adjoined.  Over a forbidden family every allowed addition either
restores peelability or gives another forbidden family.  Additions
of support size at most two always take the second alternative,
and their inverse corner deletions preserve forbiddenness as well.

For words \(x,y\in Q_F\), let
\(\operatorname{Diff}(x,y)=\{f\in F:x_f\ne y_f\}\).
The shattered supports form a simplicial complex.  A \emph{minimal
nonface} is a support outside that complex whose every proper subset
belongs to it.  The unique new shattered support in the next lemma
is called the \emph{acquired support} of the addition.

We use the following standard facts about ample families: their
restrictions, coordinate contractions, reductions, and cube
complements are ample, and every nonempty ample family is isometric
\cite[Section~3.1 and Theorems~3.1--3.2]{ChalopiChepoiMoran22}.
The \emph{ample corner-deletion criterion} says that deleting a
concept from an ample family leaves an ample family exactly when
the concept is a corner \cite[Lemma~4.1(iii)]{ChalopiChepoiMoran22}.
A reduction to a retained subset \(T\) successively deletes corners outside
\(T\) until precisely \(T\) remains; the empty sequence is allowed.
Reduction to \(\{r\}\), followed by deletion of \(r\), is a
peeling with prescribed endpoint \(r\).

\begin{lemma}[Local test for an ample one-concept extension]
\label{lem:damp-one-concept-extension-test}
Let \(S\subseteq Q_F\) be a nonempty ample family, let
\(a\in Q_F\setminus S\), and put
\[
 I_S(a)=\{f\in F:a^{\{f\}}\in S\}.
\]
Then \(S\cup\{a\}\) is ample if and only if
\begin{align}
 a^U&\in S &&\text{for every nonempty }U\subseteq I_S(a),
 \label{eq:damp-one-extension-punctured-cube}\\
 I_S(a)&\notin\operatorname{Sh}(S).
 \label{eq:damp-one-extension-missing-support}
\end{align}
When these conditions hold, \(a\) is a corner of \(S\cup\{a\}\), its
unique maximal cube has support \(I_S(a)\), and
\begin{equation}
 \operatorname{Sh}(S\cup\{a\})
 =\operatorname{Sh}(S)\mathbin{\dot\cup}\{I_S(a)\}.
 \label{eq:damp-one-extension-shadow}
\end{equation}
Moreover, \(I_S(a)\) is a minimal nonface of the shattered-set complex
\(\operatorname{Sh}(S)\).
In fact, \eqref{eq:damp-one-extension-punctured-cube} holds for every
\(S,a\) in the hypotheses.  Thus the extension is ample exactly when
it is isometric, equivalently when
\begin{equation}
 \operatorname{Diff}(a,x)\cap I_S(a)\ne\varnothing
 \qquad\text{for every }x\in S.
 \label{eq:damp-one-extension-metric-test}
\end{equation}
\end{lemma}

\begin{proof}
Put \(A=S\cup\{a\}\) and \(I=I_S(a)\).
We first establish the asserted automatic punctured cube.  There is
nothing to show when \(I=\varnothing\).  Otherwise restrict \(S\)
to the face agreeing with \(a\) outside \(I\).  This is an ample
family containing every \(I\)-neighbor of \(a\), but not \(a\).
Its complement in that face is ample, hence connected, and contains
\(a\) with none of its neighbors.  The complement is therefore
\(\{a\}\), proving \eqref{eq:damp-one-extension-punctured-cube}.
All \(I\)-traces other than \(a|_I\) are consequently present in
\(S\).  Condition~\eqref{eq:damp-one-extension-missing-support}
is therefore equivalent to \eqref{eq:damp-one-extension-metric-test},
including when \(I=\varnothing\), when both are false.
A cube geodesic from \(a\) to \(x\in S\) must start along a
direction in this intersection.  Conversely, any such first edge
can be followed by a geodesic in the isometric family \(S\), giving
a path of Hamming length from \(a\) to \(x\).  Thus the displayed
condition is equivalent to isometry of \(A\).

Suppose first that \(A\) is ample.  Since \(S=A\setminus\{a\}\) is also ample, the ample
corner-deletion criterion makes \(a\) a corner of \(A\).  Its unique
maximal cube contains every edge of \(A\) incident with \(a\), and hence
has support \(I\).  Its other vertices are precisely the
\(a^U\) with \(\varnothing\ne U\subseteq I\), which proves
\eqref{eq:damp-one-extension-punctured-cube}.

Deleting \(a\) preserves every cube avoiding \(a\).  It also preserves a
cube on each proper subset \(T\subsetneq I\): if \(f\in I\setminus T\),
then the \(T\)-face through \(a^{\{f\}}\) lies in \(S\).  Thus deletion
can remove only \(I\) from the strongly shattered sets.  Since the two
ample families have sizes differing by one, it must remove \(I\), so
\(I\notin\operatorname{Sh}(S)\) and
\eqref{eq:damp-one-extension-missing-support} and
\eqref{eq:damp-one-extension-shadow} follow.

Conversely, suppose that
\eqref{eq:damp-one-extension-punctured-cube} and
\eqref{eq:damp-one-extension-missing-support} hold.  The family \(A\)
strongly shatters every member of \(\operatorname{Sh}(S)\), because
\(S\) is ample, and it strongly shatters \(I\) by
\eqref{eq:damp-one-extension-punctured-cube}.  Hence
\[
 |\operatorname{SSh}(A)|\ge
 |\operatorname{Sh}(S)|+1=|S|+1=|A|.
\]
The Sandwich inequality gives the reverse inequality; equality in its
lower bound characterizes ample families.  Therefore \(A\) is ample.
The displayed punctured cube contains all neighbors of \(a\), so it is
the unique maximal cube through \(a\), and the first part of the proof
gives \eqref{eq:damp-one-extension-shadow}.

Finally, \(I\notin\operatorname{Sh}(S)\), while the parallel face through
\(a^{\{f\}}\) used above shows that every proper subset
\(T\subsetneq I\) is strongly shattered by \(S\).  Thus \(I\) is a
minimal nonface of \(\operatorname{Sh}(S)\).
\end{proof}

\begin{definition}[Punctured-face parametrization]
\label{def:damp-punctured-face-atlas}
For a nonempty ample family \(S\subseteq Q_F\), define
\begin{equation}
\mathcal P(S)
 =\left\{(a,I):
 \begin{array}{l}
  a\in Q_F\setminus S,\quad I=I_S(a),\\
  a^U\in S\text{ for every }\varnothing\ne U\subseteq I,\quad
  I\notin\operatorname{Sh}(S)
 \end{array}\right\}.
 \label{eq:damp-punctured-face-atlas}
\end{equation}
Thus an element of \(\mathcal P(S)\) is an oriented cube face whose only
missing vertex is \(a\), and whose support is not shattered by \(S\).
The support alone is not the datum: different missing patterns on the
same coordinate set give different oriented punctured faces.
\end{definition}

\begin{proposition}[One-concept extensions of a forbidden family]
\label{prop:damp-one-concept-atlas}
Let \(S\in\Obs(\Damp)\cap\Ample\) be irredundantly embedded in \(Q_F\).
Then
\begin{equation}
(a,I)\longmapsto S\cup\{a\}
 \label{eq:damp-one-concept-atlas-bijection}
\end{equation}
is a bijection from \(\mathcal P(S)\) onto the ample one-concept
extensions of \(S\).  Every such extension belongs to exactly one of
\(\Damp\) and \(\Obs(\Damp)\cap\Ample\).  Consequently
\begin{equation}
\mathcal P(S)
 =\mathcal P_{\mathrm{peel}}(S)
  \mathbin{\dot\cup}\mathcal P_{\mathrm{obs}}(S),
 \label{eq:damp-punctured-face-dichotomy}
\end{equation}
where the first part consists of the punctured faces whose completed
family is corner-peelable, and every element of the second part completes
to another forbidden pc-minor.
\end{proposition}

\begin{proof}
Lemma~\ref{lem:damp-one-concept-extension-test} proves the asserted
bijection.  Let \(E=S\cup\{a\}\) be one of these ample extensions.  The
same lemma makes \(a\) a corner of \(E\).

Let \(\mu\) be an elementary restriction or contraction of \(E\).  Since
the embedding of \(S\) is irredundant, \(\mu(S)\) is a proper pc-minor of
\(S\), and hence belongs to \(\Damp\).  Applying \(\mu\) either erases the
new concept, identifies it with an old image, or retains one new concept:
\begin{equation}
\mu(E)=\mu(S)
 \quad\text{or}\quad
 \mu(E)=\mu(S)\cup\{a_\mu\}.
 \label{eq:damp-one-concept-minor-image}
\end{equation}
In the second case both displayed families are ample, because ampleness
is pc-minor closed.  The ample corner-deletion criterion makes \(a_\mu\)
a corner of \(\mu(E)\).  Peeling \(a_\mu\) first and then \(\mu(S)\)
shows that \(\mu(E)\in\Damp\).  Thus all elementary pc-minors of \(E\)
belong to \(\Damp\), and pc-minor closure gives the same conclusion for
every proper pc-minor.

If \(E\) is corner-peelable, it belongs to the first part of
\eqref{eq:damp-punctured-face-dichotomy}.  Otherwise the preceding
paragraph proves \(E\in\Obs(\Damp)\cap\Ample\), which is the second part.
The two alternatives are disjoint by definition.
\end{proof}

The next invariance concerns every retained target. For targets admitting a
representation map, Proposition~\ref{prop:damp-relative-uso-extension}
interprets it as preservation of inward-boundary extension possibilities.
Its deletion formulation also covers targets for which no map is supplied.

\begin{proposition}[Additions of support size at most two]
\label{prop:damp-rank-two-extension-invariance}
Let \(S\subseteq Q_F\) be nonempty and ample, and suppose that
\(H=S\cup\{a\}\) is an ample one-concept extension with acquired
support \(P=I_S(a)\).  If \(|P|\le2\), then for every
\(T\subseteq S\), including \(T=\varnothing\),
\begin{equation}
 H\text{ reduces to }T\quad\Longleftrightarrow\quad
 S\text{ reduces to }T.
 \label{eq:damp-rank-two-relative-invariance}
\end{equation}
Both reductions use only corner deletions of concepts outside \(T\);
a sequence of length zero is allowed.  In particular,
\[
 H\in\Damp\quad\Longleftrightarrow\quad S\in\Damp,
\]
and every old prescribed endpoint is preserved.  Thus a one-concept
addition that makes any previously impossible retained target reachable
must acquire a support of size at least three.
\end{proposition}

\begin{proof}
The added concept \(a\) is a corner of \(H\), by the ample
corner-deletion criterion.  Any sequence from \(S\) to \(T\) can
therefore be prefixed by deletion of \(a\).

For the converse, suppose that \(H\) reduces to \(T\) but \(S\)
does not, and choose a counterexample with \(|S|\) least.  One has
\(|S|\ge2\), since a singleton reduces to each of its subsets.
Let \(z\) be the first deleted concept in the chosen sequence from
\(H\).  It lies outside \(T\) and is not \(a\), since deleting
\(a\) first would leave a sequence from \(S\) to \(T\).

If \(z\) is also a corner of \(S\), put
\(S'=S\setminus\{z\}\) and \(H'=H\setminus\{z\}\).
Both are ample, \(S'\) is nonempty and contains \(T\), and
\(H'=S'\cup\{a\}\) reduces to \(T\).  Its acquired support is
\(I_{S'}(a)\subseteq P\), by
Lemma~\ref{lem:damp-one-concept-extension-test}.
Minimality gives a sequence from \(S'\) to \(T\).  Prefixing it
by \(z\) contradicts the choice of \(S\).

Suppose instead that \(z\) is not a corner of \(S\).  The unique
maximal cube through \(z\) in \(H\) must contain \(a\): if it
avoided \(a\), it would still contain every cube through \(z\)
in \(S\), making \(z\) a corner there.  Since \(a\) is also
a corner of \(H\), this cube is precisely the completed \(P\)-cube
through \(a\).  The concepts \(a,z\) cannot be adjacent.  Removing
an adjacent vertex from that cube would leave at \(z\) the full face
on its remaining incident directions, again making \(z\) a corner
of \(S\).

It follows that \(|P|=2\) and \(a,z\) are opposite vertices of
the completed square.  Each has exactly the other two square vertices
as neighbors.  Hence exchanging \(a,z\) and fixing every other
vertex is a graph automorphism of \(H\).  It fixes \(T\) pointwise,
because neither \(a\) nor \(z\) belongs to \(T\), and identifies
\(H\setminus\{z\}\) with \(S=H\setminus\{a\}\).
Applying it to the remaining deletion sequence gives a sequence from
\(S\) to \(T\), the final contradiction.

Taking \(T=\varnothing\) gives ordinary peelability.  Taking
\(T=\{r\}\) and then deleting \(r\) gives the assertion for
each prescribed endpoint \(r\in S\).
\end{proof}

The proof also transports a supplied deletion sequence without searching
for another one.  If its first concept is \(a\), remove that first
entry.  If the first concept \(z\) is a corner of both families,
retain it and apply the same procedure to the remaining sequence.
In the opposite-square case, exchange \(a,z\) in the remaining
sequence.  Each recursive step reduces the family size.

The bounded checker \path{verify_damp_relative_low_support.py}
independently computes reachable retained subsets for all \(336\)
eligible one-concept extensions in \(Q_3\).  It checks \(5{,}616\)
extension--target pairs and replays \(4{,}056\) transported deletion
orders, exercising all three cases of the proof.  It also checks
ordinary peelings and complete negative relative-deletion graphs for
fixed ample families of orders \(294,295,296\), with a retained
vertex and a retained incident edge.  The respective negative graphs
have \(1,2,4\) states for each target.  These tests, recorded in
\path{damp_relative_low_support_verification.json}, verify the stated
finite instances; the proposition is proved uniformly above.

\begin{corollary}[Descent at corners of degree at most two]
\label{cor:damp-low-degree-corner-descent}
Let \(H\subseteq Q_F\) be an irredundantly embedded ample forbidden
pc-minor, and let \(a\) be a corner of \(H\) of degree at most two.
Then \(S=H\setminus\{a\}\) is again an irredundantly embedded ample
forbidden pc-minor.  Consequently, if no corner deletion of \(H\)
leaves an ample forbidden pc-minor, every corner of \(H\) has degree
at least three.
\end{corollary}

\begin{proof}
The corner-deletion criterion makes \(S\) ample, and
Proposition~\ref{prop:damp-rank-two-extension-invariance} makes it
nonpeelable.  Its embedding is irredundant: if a coordinate became
constant, \(S\) would be a proper signed restriction of \(H\), which
must peel.  Put \(I=I_S(a)\), so \(|I|\le2\).

For an elementary map \(\mu\), the family \(\mu H\) peels, and
either \(\mu H=\mu S\) or \(\mu H\) is an ample one-concept
extension of the ample \(\mu S\).  In the latter case its acquired
support still has size at most two.  Indeed, a restriction in coordinate
\(e\) retaining \(a\) leaves exactly the incident directions
\(I\setminus\{e\}\).  A contraction in a direction \(e\in I\)
identifies \(a\) with its old neighbor.  If \(e\notin I\), contraction
retains precisely the shattered sets avoiding \(e\); since
\(\operatorname{Sh}(H)=\operatorname{Sh}(S)\mathbin{\dot\cup}\{I\}\),
the acquired support of the contracted extension is still \(I\).
Proposition~\ref{prop:damp-rank-two-extension-invariance} therefore
makes \(\mu S\) peelable as well; empty images are already peelable.
All elementary pc-minors of \(S\) peel, so \(S\) is forbidden.
The terminality assertion follows immediately.
\end{proof}

\subsection{Product attachments and relative obstructions}
\label{sec:damp-product-relative-obstructions}

A relative peeling of an ample family $A$ to a nonempty subfamily $K$
deletes only corners outside $K$ and leaves exactly $K$; write $A\searrow K$.
The next construction turns this question into ordinary peelability and
transfers minimal relative failures to forbidden pc-minors. No convexity or
gatedness of $K$ in $A$ is required.

We use the elementary incident cube criterion: a vertex is a corner exactly
when the cube spanned by all its incident directions is contained in the
family. Every incident edge extends to a maximal cube, which proves the
criterion. Intersecting that cube with a coordinate face also proves that
a peeling restricts to a peeling of every coordinate section.

Fix a finite ample family $R\subseteq Q_F$ containing $0_F$, which peels but
has no peeling ending at $0_F$. The construction uses only these two peeling
properties. A fixed example is the 289-concept family of Remark~\ref{ex:damp-terminal-cut-vertex-1158}, whose only corner is its root.
For nonempty ample families $K\subseteq A\subseteq Q_E$, with $E,F$ disjoint,
put
\begin{equation}\label{eq:damp-product-attachment}
 P_R(A,K)=(A\times\{0_F\})\cup(K\times R).
\end{equation}
Thus a copy of $R$ is attached over each vertex of $K$, and the copies share
all cubes supplied by $K$. This differs from attaching independent copies
at the individual vertices: the mixed cubes are essential.

\begin{theorem}[Testing a relative peeling]\label{thm:damp-relative-attachment-test}
The family $P_R(A,K)$ is ample, has $|A|+(|R|-1)|K|$ vertices, and satisfies
\[
 P_R(A,K)\text{ peels}\quad\Longleftrightarrow\quad
 K\text{ peels and }A\searrow K.
\]
In particular, a fixed choice of $R$ converts relative peeling to ordinary
peeling with a constant factor increase in the number of vertices.
\end{theorem}
\begin{proof}
Write $T\subseteq E$ and $I\subseteq F$ for the two parts of a shattered
support. If $I$ is empty, ampleness of $A$ supplies its cube. If $I$ is
nonempty, a nonzero trace on $I$ can occur only in $K\times R$.
Realizing every trace on $T$ at that fixed nonzero trace shows that $K$
shatters $T$. Projection onto $F$ also shows that $R$ shatters $I$.
Their witnessing cubes have a product cube in $K\times R$. This proves
ampleness; the cardinality follows by subtracting the overlap $K\times\{0\}$.

Suppose $P_R(A,K)$ peels. Fix $k\in K$. Its section obtained by fixing all
$E$ coordinates to $k$ is a copy of $R$, so the restricted peeling is a
peeling of $R$. At the deletion of $(k,0)$ there must be a remaining neighbor
in this copy outside its root: otherwise the root could be postponed to
last, contradicting the choice of $R$. Let $f\in F$ be its direction.
If a remaining $(a,0)$ with $a\in A\setminus K$ were adjacent to $(k,0)$,
its direction $e\in E$ and $f$ would span a missing mixed square, since
$(a,0^{\{f\}})$ is absent. The incident cube criterion excludes this.

Restrict the peeling to $A\times\{0\}$. Every $k\in K$ is therefore deleted
without a remaining neighbor outside $K$. Retain all these vertices and run
only the deletions in $A\setminus K$. A retained vertex introduces no new
neighbor at a later outside deletion, so the original corner cube and all
incident directions at that deletion are unchanged. This gives $A\searrow K$.
Finally choose any $r\in R\setminus\{0\}$. The section with $F$ coordinates
fixed to $r$ is $K$, so $K$ peels by restriction to this coordinate face.

Conversely, perform a relative peeling of $A$ to $K$. Vertices of
$A\setminus K$ have no neighbors with nonzero $F$ coordinates, so these
are also valid deletions in the attachment. What remains is $K\times R$.
Given peelings of $R$ and $K$, delete the $R$ layers in the first order,
using the second order within each layer. At each step the incident cube
is the product of the current corner cubes: the current layer contains
its remaining part of $K$, and all later layers still contain all of $K$.
This verifies every deletion by the incident cube criterion and finishes the proof.
\end{proof}

Call $(R,0)$ a \emph{minimal rooted failure} when $R$ peels, $R\not\searrow\{0\}$,
and every elementary root-preserving image $\rho_f^0R$ and $\pi_fR$ peels to
its root. Assume all coordinates of $A$ and $R$ vary. Set $P_R(A,\varnothing)=A$.

\begin{theorem}[Transfer of minimality]\label{thm:damp-relative-attachment-minimality}
Let $(R,0)$ be a minimal rooted failure and let $K\subseteq A$ be nonempty
ample peelable families. Then $P_R(A,K)$ is a forbidden pc-minor if and only if
\begin{enumerate}
\item $A\not\searrow K$;
\item for every elementary operation $\mu\in\{\rho_e^0,\rho_e^1,\pi_e\}$
      on a coordinate of $A$ with $\mu K\ne\varnothing$, one has
      $\mu A\searrow\mu K$.
\end{enumerate}
\end{theorem}
\begin{proof}
For an old coordinate $e\in E$, direct application of the operation gives
\[
 \mu P_R(A,K)=P_R(\mu A,\mu K).
\]
For an added coordinate $f\in F$, the three identities are
\begin{align*}
 \rho_f^0P_R(A,K)&=P_{\rho_f^0R}(A,K),&
 \pi_fP_R(A,K)&=P_{\pi_fR}(A,K),\\
 \rho_f^1P_R(A,K)&=K\times\rho_f^1R.
\end{align*}
The first two added-coordinate images peel: peel their rooted factor down
to the root, layer by layer over $K$ as in Theorem~\ref{thm:damp-relative-attachment-test}, leaving
$A$, and then peel $A$. The third peels by the product construction and
pc-minor closure of peelability for $R$.

An old-coordinate image with empty $\mu K$ is simply $\mu A$, which peels.
If $\mu K$ is nonempty, it is ample and peelable by pc-minor closure;
Theorem~\ref{thm:damp-relative-attachment-test} says that the image peels exactly when
$\mu A\searrow\mu K$. The same theorem says the original attachment is
nonpeelable exactly when $A\not\searrow K$. Since every active coordinate
operation produces a proper image, and elementary images suffice by
pc-minor closure, these statements prove both directions.
\end{proof}

\begin{corollary}[Terminal product attachments along an edge]
\label{cor:damp-terminal-edge-product}
Let $A,D,R$ satisfy the hypotheses and both conditions of the forbiddenness
criterion in Theorem~\ref{thm:damp-relative-attachment-minimality}, and suppose that
$D$ is an edge in direction $e$. Put $r=\pi_eD$ and $X=P_R(A,D)$.
Every nonempty coordinate reduction of $X$ belongs to $\Damp$ if and
only if $A^e$ peels to $r$.

There is a fixed certified choice of $A$ with $289$ vertices in twelve
coordinates and $D=\{0,2\}$ for which all these hypotheses hold.
Consequently every rooted factor $(R,0)$ gives a forbidden class of order
$2|R|+287$ in $12+\dim R$ coordinates, all of whose coordinate reductions
are positive. Each output has minimum degree at least three. If $0$ is
the only corner of $R$, the output is cornerless.
\end{corollary}
\begin{proof}
Use the disjoint coordinate blocks $E,F$ of $A,R$. For $f\in E\setminus
\{e\}$, the edge $D$ is constant in $f$, so
$X^f=A^f\times\{0_F\}$. For $f\in F$, one has $X^f=D\times R^f$.
These families are positive when nonempty by reduction closure and the
product theorem. In direction $e$, the reduction is the vertex union
\[
 X^e=(A^e\times\{0_F\})\cup(\{r\}\times R).
\]
Both pieces are positive, and $R$ cannot peel to its root. The
vertex-gluing criterion therefore makes this reduction positive exactly
when $A^e$ peels to $r$. Forbiddenness itself is the existing minimality
transfer theorem.

For the fixed input take the $289$-vertex rooted factor, normalized at its
unique corner $0$, stored in \path{research/root_cube_completion.json}.
Its ordinary order is replayed in \path{research/terminal_edge_product.py/json}.
It cannot peel to $D=\{0,2\}$, since its only initial corner lies in $D$.
The verifier reuses the input orders from
\path{research/rooted_square_assembly.json}: relative orders for all $22$ root-side restrictions
and contractions in coordinates other than $e=1$. In the shared coordinate,
both sections and the contraction have orders ending at the image of $D$.
These are all $25$ elementary operations with nonempty target; the other
eleven sections have empty target. The reduction $A^e$, of order $65$,
also has a newly replayed order ending at $0$. Every full order is independently checked
through outgoing cubes, distinct out-sets, and nonclashing. Thus the
minimality-transfer hypotheses and the terminality criterion hold.

The fixed $A$ has minimum degree three. A rooted factor $R$ is
two-connected by Theorem~\ref{thm:damp-cut-vertex-obstructions}, hence has
minimum degree at least two. Vertices outside the copy of $A$ have their
degree in $D\times R$, at least three, and vertices of $A$ lose no
neighbours. This proves the degree assertion. If $0$ is the only corner
of $R$, its nonroot missing-cube witnesses survive in each product layer;
the nonroot witnesses of $A$ survive as well. At $0$, an incident direction
of $A$ other than $e$ and an incident direction of $R$ span an absent
mixed square. Thus no vertex is a corner.
\end{proof}

Taking $R=A$ gives a cornerless terminal forbidden class of order $865$
and minimum degree four. This is the previously recorded cornerless proper
halfspace of the $1601$-vertex square-code control in
\path{research/cube_corner_compatibility.json}. The new verifier identifies
the families by moving bit $23$ to bit $1$, shifting bits $1,\ldots,22$
up one position, and then reflecting bit $1$. The new assertion is its
forbiddenness and positive reductions, not the existence of the support.
In this marked presentation the peripheral piece is $D\times R$: its
duplicated site is $R$ itself and cannot peel to the root. Consequently the
root-positive-site hypothesis in
Theorem~\ref{thm:damp-rooted-downset-completion} is not automatic even for
terminal forbidden classes with a peripheral piece. The present operation
reuses product attachment; it does not intrinsically generate rooted factors
or exhaust the terminal inputs.

This gives a uniform construction of forbidden minors from minimal relative
failures of pairs $(A,K)$. Its input admissibility still uses relative
peelings. It therefore transfers a classification problem; it does not
supply intrinsic admissibility for those pairs or exhaust all forbidden
minors by such attachments.


\begin{remark}[Verification scope]
The checker \path{verify_gated_relative_tests.py} replays an ordinary peeling
and all 24 elementary root-preserving orders of the fixed 289-concept factor,
and checks that its root is its only corner. It constructs specified
boundary-record orders, including a nonconvex attachment site, and verifies
the coordinate-minor identities on that attachment. A 577-concept positive
attachment has a replayed order; a 577-concept negative attachment is
cornerless. The record is \path{gated_relative_tests_verification.json}.
The universal statements above follow from their proofs; no exhaustive
classification of relative pairs is claimed.
\end{remark}

\subsection{Gluing at vertices and separating cubes}
\label{sec:damp-gluing-constructions}

\paragraph{Gluing at a vertex.}
A \emph{rooted family} is a pair \((A,a)\) with \(a\in A\).
We say that \(A\) \emph{peels to \(a\)} if it has a corner peeling
whose last concept is \(a\).  A \emph{root-preserving pc-minor}
is obtained by coordinate contractions, carrying the root to its image,
and convex restrictions that contain the current root.  Peeling to the
root is preserved by these operations: reverse the peeling to obtain
an ample prefix order starting at the root, restrict the prefixes, or
order projected concepts by their first occurrence.  The resulting
ample prefix order still starts at the image of the root.  Here an
ample prefix order means an ordering whose every nonempty initial
segment is ample; reversing such an order gives a corner peeling.

\begin{proposition}[Peeling after gluing at one vertex]
\label{prop:damp-vertex-gluing}
Let \(A_1,\ldots,A_m\) be nonempty ample families, each containing
\(0\), on pairwise disjoint coordinate sets.  Embed them in the
product cube by setting all other coordinates to zero, and put
\[
 X=A_1\vee\cdots\vee A_m=\bigcup_{i=1}^m A_i.
\]
Then \(X\) is ample.  It peels to \(0\) if and only if every
\(A_i\) peels to \(0\).  Moreover,
\begin{equation}
 X\in\Damp
 \quad\Longleftrightarrow\quad
 \left\{
 \begin{array}{l}
 A_i\in\Damp\text{ for every }i,\\
 \text{at most one }A_i\text{ fails to peel to }0.
 \end{array}\right.
 \label{eq:damp-vertex-gluing-test}
\end{equation}
\end{proposition}

\begin{proof}
A pair of coordinates from different blocks is not shattered, because
the trace \(11\) is absent.  The shadow of \(X\) is therefore
the union of the block shadows, with the empty support shared.  Hence
\[
 |\operatorname{Sh}(X)|
 =1+\sum_i\bigl(|\operatorname{Sh}(A_i)|-1\bigr)
 =1+\sum_i(|A_i|-1)=|X|,
\]
which proves ampleness.

For the peeling assertion, use Proposition~\ref{prop:damp-relative-uso-extension}
with retained singleton $\{0\}$. A USO on the union restricts to each
piece. Its outgoing directions at $0$ can lie in at most one exclusive
block, by (C1), since mixed squares are absent. Thus in every other piece
$0$ is a sink, and that piece peels to $0$. A sink at $0$ in the union
requires this in every piece.

Conversely, choose acyclic USOs with sink $0$ in all but a possible
exceptional piece, and any acyclic USO in that piece. At $0$, all outgoing
edges belong to the exceptional piece, so (C1) holds; elsewhere it holds
within the unique piece. Every cube lies in a piece, giving (C2). A simple
cycle cannot cross between pieces through their only common vertex.
The union is therefore acyclic. Choosing sink $0$ in every piece also
proves the prescribed-sink assertion.
\end{proof}

Recall that a \emph{cut vertex} disconnects a connected graph when
deleted.  Its \emph{blocks} are its maximal two-connected subgraphs
and its bridges, with each bridge regarded as a two-vertex block.

\begin{theorem}[Forbidden pc-minors with a cut vertex]
\label{thm:damp-cut-vertex-obstructions}
An ample family with a cut vertex belongs to \(\Obs(\Damp)\)
if and only if it is obtained by gluing two rooted ample families
\((A_1,a_1)\), \((A_2,a_2)\) at their roots, where each pair
satisfies all three conditions:
\begin{enumerate}
 \item \(A_i\in\Damp\);
 \item no acyclic USO of \(A_i\) has sink \(a_i\);
 \item every proper root-preserving pc-minor of \((A_i,a_i)\)
       admits an acyclic USO with sink at the image of the root.
\end{enumerate}
By Proposition~\ref{prop:damp-relative-uso-extension}, conditions 2 and 3
are exactly failure and proper-minor restoration of peeling to the root.
Each \(A_i\) has no cut vertex. The glued graph has exactly two
blocks and one cut vertex, so the two rooted factors are recovered
uniquely up to interchange.
\end{theorem}

\begin{proof}
First note how a cut vertex \(a\) of a partial cube separates its
coordinates.  No coordinate can differ from its value at \(a\) in two components
of the graph with \(a\) removed.  Otherwise choose vertices \(x,y\)
in those components whose bits in that coordinate differ from the bit
at \(a\).  Every path between them passes through \(a\), so its
length is at least
\(d_H(x,a)+d_H(a,y)>d_H(x,y)\), where \(d_H\) is Hamming
distance, contradicting isometry.
Thus, after translating \(a\) to zero, the induced families on each
component together with \(a\) occupy disjoint coordinate blocks.
They and any union of them are convex: a shortest path cannot leave
a component through \(a\) and return through the same vertex.
They therefore give the vertex-gluing presentation above.

Let \(X\in\Obs(\Damp)\) have a cut vertex.  Each rooted
component family is a proper convex restriction and belongs to
\(\Damp\).  By \eqref{eq:damp-vertex-gluing-test}, at least two
of them fail to peel to the cut vertex.  If there were any additional
component, the union of two failing ones would be a proper convex
nonpeelable restriction.  Hence there are exactly two component
families \(A_1,A_2\), and both satisfy conditions 1 and 2.

An elementary root-preserving restriction or contraction in \(A_1\)
produces a proper minor of \(X\) obtained by gluing the resulting
rooted family to the unchanged \(A_2\).  That minor peels, whereas
\(A_2\) does not peel to its root.  The vertex-gluing test forces
the changed family to peel to its root.  The same reasoning applies
to \(A_2\).  Since peeling to the root is preserved by subsequent
root-preserving pc-minors, condition 3 follows in full.

Conversely, glue two pairs satisfying the three conditions, using
disjoint coordinate sets and identifying their roots with \(0\).
The result is ample and nonpeelable by the vertex-gluing test.
For a coordinate of \(A_1\), contraction or restriction to the
semicube containing \(0\) replaces \(A_1\) by a proper
root-preserving minor.  Peel that minor to \(0\), then peel
\(A_2\).  Restriction to the opposite semicube discards \(A_2\)
entirely and leaves a pc-minor of the peelable \(A_1\).
All three elementary images therefore peel.  The argument for a
coordinate of \(A_2\) is identical, proving forbidden-minor status.

It remains to show that a rooted factor \((A,a)\) satisfying the
conditions has no cut vertex.  If its root is a cut vertex, the rooted
part of Proposition~\ref{prop:damp-vertex-gluing} forces one component
family to fail to peel to the root.  That family is a proper
root-preserving convex restriction, contrary to condition 3.
If another vertex \(b\ne a\) is a cut vertex, let \(U\) consist
of \(b\) and the component containing \(a\), and let \(V_j\)
be the other component families, rooted at \(b\).
The same deletion argument shows that \(A\) peels to \(a\)
exactly when \(U\) peels to \(a\) and every \(V_j\) peels
to \(b\): before deleting \(b\), all the other branches must
disappear, because the branch containing the retained \(a\)
still has an edge at \(b\).
The pair \((U,a)\) is a proper root-preserving convex restriction,
so it peels to its root.  Each \((V_j,b)\) is also a proper
root-preserving pc-minor of \((A,a)\): contract all coordinates
outside its block, which sends \(a\) to \(b\).
Condition 3 makes all these pairs peel to their roots, a contradiction.
Thus the factors have no cut vertices.  Condition 2 also excludes a singleton or a single edge, since either
peels to any chosen root.  Thus the factors are two-connected, and the
asserted block decomposition and uniqueness follow.
\end{proof}

The theorem completely describes the block decomposition of an ample
forbidden pc-minor with a cut vertex.  Its conditions still use peeling
and rooted minor minimality.  An intrinsic classification must therefore
also identify these rooted factors without those recognition tests;
the theorem is not such an admissibility criterion.

\begin{proposition}[Minimality under ample inclusion at a cut vertex]
\label{prop:damp-inclusion-minimal-cut-vertex}
Let $X$ be ample, negative for $\Damp$, and suppose every proper ample
subfamily of $X$ belongs to $\Damp$. Then $X$ is cornerless and every
coordinate reduction of $X$ belongs to $\Damp$.

If $r$ is a cut vertex, it separates $X$ into exactly two component
families $A_1,A_2$, including $r$ in each. Conversely, a vertex gluing
$X=A_1\vee_r A_2$ has the stated inclusion-minimality property if and
only if its rooted pieces satisfy
\begin{enumerate}
 \item $A_i\in\Damp$, but $A_i$ does not peel to $r$;
 \item every proper ample $B\subsetneq A_i$ peels relatively to
 $B\cap\{r\}$, where an empty target means an ordinary full peeling.
\end{enumerate}
The coordinate blocks and pieces are recovered from the marked cut
vertex as in Theorem~\ref{thm:damp-cut-vertex-obstructions}. No assertion
that these rooted pieces are minimal under pc-minors is made.
\end{proposition}
\begin{proof}
If $X$ had a corner, deleting it would leave a proper ample positive
subfamily, and its peeling could be prefixed by that deletion. This
contradicts negativity. A coordinate reduction $R$ has a copy
$R\times\{0\}$ inside the corresponding section of $X$. That copy
is ample and proper, so $R$ is positive.

Reuse the coordinate separation and vertex-gluing criterion from
Theorem~\ref{thm:damp-cut-vertex-obstructions} and
Proposition~\ref{prop:damp-vertex-gluing}. All component families at
$r$ are proper ample subfamilies, hence positive. At least two fail to
peel to $r$. The union of any two such pieces is ample and negative;
inclusion-minimality therefore forces exactly two pieces.
For a proper ample $B\subsetneq A_1$ containing $r$, the union
$B\vee_r A_2$ is a proper ample subfamily of $X$, hence positive.
Since $A_2$ fails to peel to $r$, the vertex-gluing criterion forces
$B$ to peel to $r$. If $r\notin B$, positivity follows directly from
$B\subsetneq X$. This proves condition 2, and the other piece is
handled identically.

Conversely, suppose the two conditions hold. Vertex gluing makes $X$
ample and negative. Let $Y\subsetneq X$ be ample and nonempty. If
$r\notin Y$, connectedness puts $Y$ in one piece, where condition 2
makes it positive. If $r\in Y$, put $Y_i=Y\cap A_i$. Each is an
ample coordinate restriction of $Y$. Both are positive by condition 1
or 2, and at least one is proper in its piece because $Y\ne X$.
That piece peels to $r$ by condition 2. The vertex-gluing criterion
then makes $Y$ positive. The empty subfamily is positive as well.
\end{proof}

Unlike the pc-minor criterion, condition 2 concerns arbitrary ample
subfamilies. It is not an intrinsic description of the rooted inputs.
The proposition explains what the stronger minimality requires at a
separator; it does not exclude separators or replace contraction tests.

\begin{remark}[Positive faces do not replace positive contractions]
\label{rem:damp-separated-roots-face-control}
Take two rooted factors $(A,0)$ and $(B,0)$ satisfying the three
conditions of Theorem~\ref{thm:damp-cut-vertex-obstructions}, embedded
on disjoint active coordinate sets and extended by zero in the other
block. For a new coordinate $e$, form
\[
 X=(A\times\{0\})\cup(B\times\{1\}).
\]
Thus the two roots are joined by a single $e$-edge. The family $X$ is
ample, every elementary pc-minor except $\pi_eX$ belongs to $\Damp$,
and $\pi_eX=A\vee_0 B$ is forbidden. In particular, all proper
coordinate faces are positive, but $X$ is not a forbidden pc-minor.

Indeed, the contraction $A\vee_0B$ is ample, the intersection of the
two sections is the root, and no edge joins their exclusive parts.
The ample component-expansion criterion therefore gives ampleness of
$X$. Its two $e$-faces are $A$ and $B$, hence positive; its
$e$-contraction is forbidden by the cut-vertex theorem. For a coordinate
$q$ of $A$, the face away from the root is a proper face of $A$ and
is positive. The face through the root and the contraction in $q$
replace $A$ by a proper root-preserving minor $A'$ which peels to its
root. Peel $A'$ to that root, delete the root as the pendant endpoint
of the remaining bridge, and then peel $B$. This gives positive orders
for both images. The same argument, with the factors interchanged,
handles every coordinate of $B$. All further proper coordinate-face
restrictions inherit positivity. Finally, the negative contraction
precludes both positivity and forbidden-minor minimality of $X$.
This is an application of the rooted gluing theorem; the extra edge
makes the distinction between face tests and contraction tests explicit.
\end{remark}

\begin{corollary}[Inclusion-minimal negative classes of orders $527$ and $528$]
\label{cor:damp-inclusion-minimal-527-528}
Let $(A,r)$ be the $264$-word rooted input with $r=750$ in
Proposition~\ref{prop:damp-rooted264-double527}. The $527$-word gluing
of two copies at their roots has every proper ample subfamily in $\Damp$.
The $528$-word class obtained by joining the two roots with a new edge
also has this property and is cornerless, but is not a forbidden pc-minor:
contracting the new edge gives the negative $527$-word class.
Both examples have all coordinate reductions in $\Damp$.
\end{corollary}
\begin{proof}
The companion ample-corners manuscript proves
\texttt{prop:root264-relative-ample-inclusion}: every proper ample
$B\subsetneq A$ peels to $B\cap\{r\}$. Its exact all-subfamily
formula omits the cornerlessness condition only at $r$, permits $r$
to be absent, and excludes subfamilies contained in $\{r\}$. The
independently checked refutation is in
\path{../ample-corners/evidence/root264_ample_inclusion/}.
The input $A$ is positive and its only corner is $r$, so it cannot
peel to $r$. Proposition~\ref{prop:damp-inclusion-minimal-cut-vertex}
proves the $527$ assertion.

To handle the bridge, attach a new leaf $s$ to the root of a copy of
$A$, obtaining $L$ rooted at $s$. Deleting $s$ and peeling $A$ shows
$L$ positive. A peeling retaining $s$ would have to peel $A$ to $r$:
otherwise deleting $r$ would disconnect a nonempty remainder from $s$,
contrary to ampleness. Thus $L$ fails to peel to $s$.
If a proper ample $Y\subsetneq L$ contains $s$ and another vertex,
connectedness forces it to contain $r$. Its old-coordinate section
$Y\cap A$ is proper in $A$, and peels to $r$; then delete $r$, leaving
$s$. If $s\notin Y$, either $Y=A$ or the input property makes $Y$
positive. The empty and singleton cases are immediate. Hence $L$
satisfies condition 2 at $s$. Gluing $L$ at $s$ to the root of the
other copy of $A$ gives the $528$-word class, so the same proposition
proves its inclusion-minimality and cornerlessness.

Its negative bridge contraction is exactly the construction of
Remark~\ref{rem:damp-separated-roots-face-control}. That excludes
pc-minor minimality. Positivity of all reductions follows from the
first part of Proposition~\ref{prop:damp-inclusion-minimal-cut-vertex}.
The companion verifier also directly checks both corner sets and the
bridge contraction. The old $527$ forbidden-minor certification is
reused, not reproved by this argument.
\end{proof}

Thus even cornerlessness, positivity of every proper ample subfamily,
and positivity of all reductions together do not imply forbidden-pc-minor
status. The $528$ example strengthens the earlier face-only warning by
adding the all-subfamily property; the bridge construction itself is reused.

\paragraph{Separating coordinate faces.}
A nonempty coordinate-face intersection \(D\) is \emph{separating} when deleting all its vertices
disconnects the graph.  In particular, a separating edge means that its
two endpoints form a vertex cut; it does not mean that the edge is a
bridge.  For \(D\subseteq A\), a peeling of \(A\) \emph{to \(D\)}
deletes corners outside \(D\) until precisely \(D\) remains.

\begin{proposition}[Ample assembly and USO compatibility at a coordinate face]
\label{prop:damp-uso-cube-gluing}
Let $\varnothing\ne D\subseteq Q_E$ and
$A_i\subseteq Q_{E\mathbin{\dot\cup}F_i}$ be ample, with pairwise
disjoint exclusive blocks $F_i$ and
$A_i\cap(Q_E\times\{0\}^{F_i})=D$.
Embed the pieces by putting zero in the other blocks, and set
$X=\bigcup_{i=1}^m A_i$. Then $X$ is ample if and only if
\begin{equation}
 \operatorname{Sh}(A_i)\cap\operatorname{Sh}(A_j)
 =\operatorname{Sh}(D)\qquad(i\ne j).
 \label{eq:damp-face-assembly-shadow-test}
\end{equation}
This condition is automatic when $D=Q_E$.
When it holds, USOs on the pieces glue to a USO of $X$ exactly when they agree on $D$
and, at each $d\in D$, at most one piece has an outgoing edge from $d$
into $A_i\setminus D$. The glued USO is acyclic exactly when its union
orientation is acyclic. For $|D|\le2$, acyclicity of the piece USOs already
implies acyclicity of their union.
\end{proposition}
\begin{proof}
The pieces intersect pairwise in $D$, so
$|X|=\sum_i|A_i|-(m-1)|D|$. Every contained cube belongs to a piece:
a cube varying coordinates in two exclusive blocks would contain a word
nonzero in both blocks, which is absent; a cube with a fixed nonzero
exclusive coordinate stays in that piece. Cubes with all exclusive
coordinates fixed to zero lie in $D$. Consequently
\[
 \operatorname{SSh}(X)=\bigcup_i\operatorname{Sh}(A_i),
\]
where $\operatorname{SSh}$ denotes the strongly shattered supports.
Each piece shadow contains $\operatorname{Sh}(D)$, of cardinality $|D|$.
Thus this union has cardinality $|X|$ exactly when the supports outside
$\operatorname{Sh}(D)$ occur in only one piece, which is
\eqref{eq:damp-face-assembly-shadow-test}. Equality in the lower Sandwich
inequality characterizes ampleness, proving both directions. For $D=Q_E$,
any support common to two piece shadows lies in $E$ and already belongs
to $\operatorname{Sh}(D)=2^E$, so the test is automatic.

Every contained cube belongs to a piece, so (C2)
holds for the union whenever the orientations agree. Outside $D$, all
neighbours lie in one piece. At $d\in D$, outgoing directions into two
exclusive blocks violate (C1) because their square is absent. If at most
one piece supplies outgoing exclusive directions, all outgoing directions
have their cube in that piece. This proves necessity and sufficiency.
Conversely, the pieces inherit USOs from the union by CCMW Lemma~6.9,
as they are coordinate-face intersections.

A simple directed cycle crossing a singleton intersection is impossible.
If $D=\{s,t\}$ is oriented $s\to t$, a simple directed cycle visiting
multiple pieces would contain a directed path $t\to s$ within one piece,
with interior outside $D$. This path and the common edge contradict
acyclicity in that piece. For larger faces the union acyclicity condition
is retained; agreement alone is not being asserted sufficient.
\end{proof}

\begin{remark}[Ample subpieces need not inherit any acyclic USO]
\label{rem:damp-nonconvex-peripheral-cover}
The coordinate-face hypothesis in the preceding gluing result cannot be
replaced by requiring positive ample pieces and overlap, even when every
cube of the union belongs to a piece. Here is a positive union for which
no acyclic USO restricts to a USO on one of its two proper pieces.

Let $A\subseteq Q_F$ be positive with unique corner $r$, and suppose
$r$ has two incident directions $i,j$. Put
$C=\{r,r^i,r^j\}$. The square vertex $r^{\{i,j\}}$ belongs to $A$
because $r$ is a corner, but is absent from $C$. Thus $C$ is an ample
three-vertex path and $r$ is not its corner. Every acyclic USO of $A$
has $r$ as a source: an acyclic orientation has a source, and every
source of a USO is a corner. Both edges at $r$ in $C$ therefore point
outward. Their required outgoing square is absent, so no such orientation
restricts to a USO on $C$.

Introduce fresh coordinates $g,e$, embed $A,C$ with $g=0$, and set
$L=C\cup\{r^g\}$. Define
\[
\begin{split}
 X&=(L\times\{0\})\cup(A\times\{1\}),\\
 P_0&=(L\times\{0\})\cup(C\times\{1\}),\qquad
 P_1=(C\times\{0\})\cup(A\times\{1\}),
\end{split}
\]
where the displayed layers are indexed by $e$. Both pieces are positive
peripheral expansions, and their intersection is $C\times Q_1$.
The union $X=P_0\cup P_1$ is obtained from positive $P_1$ by attaching
one pendant edge in the fresh direction $g$ at $(r,0)$; hence $X$ is
ample and positive. Both pieces are proper, and every cube of $X$ lies
in $P_1$ except for the new edge, which lies in $P_0$.
Nevertheless, an acyclic USO of $X$ restricts to an acyclic USO on its
coordinate face $e=1$, namely $A$. At $(r,1)$ its two outgoing edges
in directions $i,j$ remain in $P_0$, whereas their opposite square
vertex does not. The restriction to $P_0$ fails the outgoing-cube axiom.
This excludes every acyclic USO, independently of any choice of orders.

For the certified $264$-vertex $A$ of
Proposition~\ref{prop:damp-rooted264-double527}, take $r=750$,
$i=0$, and $j=4$. Then $C=\{750,751,766\}$ and the missing square
vertex is $767$. The two pieces have orders $7,267$, their overlap
has order $6$, and their positive union has order $268$.
The check \path{research/nonconvex_uso_restriction.py/json} verifies
these sets, ampleness, the unique corner, and explicit positive orders.
The universal orientation obstruction is the preceding source argument.
Thus testing compatibility of USOs on these nonconvex peripheral pieces
would incorrectly reject a positive assembly; it cannot be used as a
necessary terminal-input criterion. This does not affect the gluing theorem
for coordinate-face pieces or CCMW's restriction theorem.
\end{remark}

\begin{corollary}[Finite boundary data for acyclic gluing]
\label{cor:damp-finite-boundary-gluing}
Use the hypotheses and notation of
Proposition~\ref{prop:damp-uso-cube-gluing}, assume its shadow-intersection
condition, and put $k=|D|$.
For each piece $A_i$ and each total order $\sigma$ of $D$, let
$\mathcal S_i(\sigma)$ consist of the sets $S\subseteq D$ for which
there is an acyclic USO of $A_i$ with a topological order restricting to
$\sigma$ on $D$, and with
\[
 S=\{d\in D: d\text{ has an outgoing neighbour in }A_i\setminus D\}.
\]
Then $X$ is corner-peelable if and only if there are one order $\sigma$
and pairwise disjoint sets $S_i\in\mathcal S_i(\sigma)$.
One may replace every $\mathcal S_i(\sigma)$ by its inclusion-minimal
members without changing this criterion. Thus a piece has at most
$k!2^k$ possible boundary records before this reduction, independently
of the number of its interior vertices.
\end{corollary}
\begin{proof}
Restrict an acyclic USO of $X$ to its coordinate-face pieces and restrict
one global topological order to $D$. The preceding proposition gives
the disjointness of the resulting sets $S_i$, proving necessity.

Conversely, choose witnessing piece orientations and topological orders.
They orient each edge of $D$ from its earlier to its later endpoint in
$\sigma$, so they agree on $D$. Disjointness of the $S_i$ gives a USO
of $X$ by the preceding proposition. To see that it is acyclic, split
each witnessing order into blocks of outside vertices before the first
vertex of $\sigma$, between successive vertices of $\sigma$, and after
the last vertex. Concatenate the pieces' blocks within each interval,
and insert each boundary vertex once. The resulting order restricts to
the chosen order on each piece. Every edge belongs to a piece, so this
is a topological order of the union. CCMW Proposition~6.12 makes it a
corner peeling.

If disjoint records have been chosen, replace each $S_i$ by a minimal
member of $\mathcal S_i(\sigma)$ contained in it. Such a member exists
by finiteness, and disjointness persists. The converse is immediate.
There are $k!$ orders and $2^k$ subsets, giving the stated bound.
\end{proof}

This eliminates a test on the full union digraph from the compatibility
rule, once the attainable boundary records of the pieces are supplied.
It does not intrinsically characterize those records or generate their
realizing pieces. In particular, the finite bound is for a fixed separator;
it is not a bound on the primitive inputs to the forbidden-minor problem.
Unlike the original cube case, this statement also applies to an ample
coordinate-face intersection that is not a full cube. In particular, it
applies to all eight canonical faces in
Theorem~\ref{thm:damp-four-edge-face-types}: their recovered pieces satisfy
the shadow test because their union is ample. In forward assembly, the
shadow test must be imposed explicitly.

For example, let $D=\{0,1,3\}\subseteq Q_2$ and take two copies of
$Q_3\setminus\{2\}$, identifying their zero third-bit sections with $D$
and using different exclusive coordinates. Both pieces are ample, but their
shadows intersect in $2^{\{0,1\}}$, strictly larger than
$\operatorname{Sh}(D)=\{\varnothing,\{0\},\{1\}\}$. Their union has
$11$ vertices, $10$ strongly shattered supports and $12$ shattered supports,
so it is not ample. Replacing the second piece by $D\square Q_1$ gives
an ample union of order $10$. The direct checks in
\path{research/face_assembly_shadow.py/json} cover these controls and all
$1032$ assemblies obtained from nonempty ample $D\subseteq Q_2$ and two
ample pieces in three coordinates with zero exclusive section $D$.

\begin{proposition}[Transport of boundary records]
\label{prop:damp-boundary-record-transport}
Let $D=Q_E\times\{0\}^F\subseteq A\subseteq Q_{E\mathbin{\dot\cup}F}$
be ample, and let a corner order of $A$ realize a boundary record
$(\sigma,S)$ as in Corollary~\ref{cor:damp-finite-boundary-gluing}.
For $f\in F$, identify $D$ with its image after deleting coordinate $f$.
\begin{enumerate}
\item Restriction to $f=0$ realizes a record $(\sigma,T)$ with $T\subseteq S$.
\item Contraction in $f$, with the order transported by last-fibre positions,
realizes a record $(\tau,T)$ with $T\subseteq S$. Moreover, if
$x\in D\setminus S$ precedes $y\in D$ in $\sigma$, then $x$ precedes
$y$ in $\tau$.
\end{enumerate}
In particular, contraction preserves the boundary order whenever $S$ is empty.
\end{proposition}
\begin{proof}
Both transported orders are corner orders by pc-minor closure and its
order-level proof. Under restriction, the boundary vertices keep their
positions relative to one another, and every remaining outward edge was
already an outward edge with the same orientation. This proves the first
assertion.

For contraction, the representative of a double fibre is its later endpoint,
equivalently its sink. CCMW Proposition~6.14 gives the outgoing directions
of the image vertex from this representative, with $f$ erased. If a boundary
vertex $d$ has its outside mate $d^f$ as representative, the original edge
$d\to d^f$ puts $d$ in $S$. If its representative is $d$ itself, any outgoing
exclusive direction in the image was already such a direction at $d$.
In either case, an outward vertex of the contracted record belongs to $S$.

Write $p(v)$ for the position in the original order. The new position used
to order a boundary vertex $d$ is $\max\{p(v):\pi_f(v)=d\}$, which is at
least $p(d)$. For $d\notin S$, a present mate $d^f$ precedes $d$, so this
maximum is exactly $p(d)$. Thus $p(x)<p(y)$ with $x\notin S$ implies that
the new position of $x$ is smaller than that of $y$. This proves the
second assertion and its final consequence.
\end{proof}

The statement concerns a transported witness. It does not assert that the
original order $\sigma$ is unattainable in the minor when $\tau\ne\sigma$.
For example, on the square $A=\{0,1,2,3\}$ with boundary edge
$D=\{0,1\}$, the corner order $0,1,3,2$ has $\sigma=(0,1)$ and
$S=D$. Contracting bit $1$ gives the last-fibre boundary order $(1,0)$
and empty outward set, but the minor edge also admits $(0,1)$.
The transport argument does not preserve the order of every attainable
record, even if the minor orientation may be chosen anew, as the next
proposition shows. The verifier
\path{research/boundary_record_transport.py/json} checks the stated law,
not that stronger existential assertion. If the minor peels to $D$,
same-order attainability does hold: restrict
the original corner order to $D$ to obtain $\sigma$, and append $\sigma$
to a relative peeling of the minor to $D$. This realizes $(\sigma,\varnothing)$.
Consequently, a counterexample to same-order attainability must have a
minor that itself fails relative peeling to $D$; relative-minimal pieces
cannot supply such a counterexample in their exclusive proper minors.

\begin{proposition}[An attainable boundary order can disappear under contraction]
\label{prop:damp-boundary-attainability-failure}
Preservation of attainable records with the same boundary order and a
smaller outward set fails under exclusive-coordinate contraction, already
for an edge boundary. It fails even when only inclusion-minimal outward
sets are recorded for each boundary order.
\end{proposition}
\begin{proof}
Use the positive $289$-concept class $C\subseteq Q_{12}$ of
Remark~\ref{ex:damp-terminal-cut-vertex-1158}, translated so that its
unique corner is $0$. It contains $2$ and $4$, the neighbours of $0$
in bits $1$ and $2$. Put $D=\{0,2\}$ and, for a fresh coordinate $f$,
form the peripheral expansion
\[
 A=(D\times\{0\})\cup(C\times\{1\}).
\]
Both $A$ and $C$ are ample and corner-peelable. Deleting $(2,0)$,
then $(0,0)$, then following any corner order of the upper copy of $C$
realizes the lower-boundary order $\sigma=(2,0)$ with outward set $D$.

In fact this is the only attainable outward set for that boundary order.
Every corner of the upper layer would have to be a corner of $C$:
a corner cube in $A$ restricts to a corner cube in that layer.
The only candidate is $(0,1)$, which is not a corner because its
neighbours $(0,0)$ and $(4,1)$ have missing square vertex $(4,0)$.
Thus initially the only corners of $A$ are its two lower vertices.
After deleting $(2,0)$, the same argument still blocks $(0,1)$, and
the only corner is $(0,0)$. Any order realizing $\sigma$ must therefore
start with these two deletions. Both have their upper mates still present,
so both belong to the outward set. Hence $D$ is inclusion-minimal, indeed
the unique record for $\sigma$.

Contracting $f$ returns $C$ and preserves the marked edge $D$.
Every corner order of $C$ starts at its unique corner $0$, so no order
can restrict to $(2,0)$ on $D$, regardless of its outward set. By the
acyclic-USO correspondence, changing the minor orientation cannot repair
this failure. The saved class, orders, and forced first corners are
replayed in \path{research/boundary_attainability_failure.py/json}.
\end{proof}

This does not contradict minor closure of positive gluing. The lost record
has outward set all of $D$, so any other piece compatible with that record
must have empty outward set and hence peel to $D$. Such pieces can be
peeled to $D$ first, leaving the positive contracted piece $C$.
Consequently this lost record alone does not distinguish positive gluing
contexts. Inclusion-minimality at each fixed boundary order is insufficient
to make the raw record families monotone under minors.

\begin{theorem}[Finite gluing algebra and comparison by contexts]
\label{thm:damp-finite-gluing-algebra}
Fix a nonempty corner-peelable ample family $D\subseteq Q_E$.
A piece is an ample family $A\subseteq Q_{E\mathbin{\dot\cup}F}$
whose zero $F$-section is exactly $D$. Private coordinate sets are
renamed to be disjoint before gluing. Put
\[
 U_D=2^E\setminus\operatorname{Sh}(D),\qquad
 T_A=(\operatorname{Sh}(A)\cap2^E)\setminus\operatorname{Sh}(D),
\]
and let $\Sigma_D$ be the set of corner orders of $D$.
For each piece define its record family
\[
 \mathcal R(A)=\{(T_A,\sigma,S):\sigma\in\Sigma_D,
                      \ S\in\mathcal S_A(\sigma)\}.
\]
Here $\mathcal S_A(\sigma)$ is defined in
Corollary~\ref{cor:damp-finite-boundary-gluing}. In particular,
$\mathcal R(A)=\varnothing$ if $A$ is not corner-peelable.
Set
\[
 \Omega_D=2^{U_D}\times\Sigma_D\times2^D.
\]
For arbitrary $P,Q\subseteq\Omega_D$ define
\begin{equation}
 P*Q=\{(T\cup T',\sigma,S\cup S'):
 (T,\sigma,S)\in P,\ (T',\sigma,S')\in Q,
 T\cap T'=\varnothing,\ S\cap S'=\varnothing\}.
 \label{eq:damp-record-product}
\end{equation}
Then the following statements hold.
\begin{enumerate}
\item This product is associative and commutative, with zero
$\varnothing$ and identity
$I_D=\{(\varnothing,\sigma,\varnothing):\sigma\in\Sigma_D\}$.
For two pieces glued along $D$,
\[
 \mathcal R(A\cup_D B)=\mathcal R(A)*\mathcal R(B),
\]
where the left side is defined to be empty if the union is not ample.
Thus the union is ample and corner-peelable exactly when the product
is nonempty.
\item Let $\mathfrak R_D$ consist of the record families realized by
pieces, together with zero. It is a finite monoid. On it define
\[
 P\preceq Q\quad\Longleftrightarrow\quad
 \text{for every }R\in\mathfrak R_D,
       \ P*R\ne\varnothing\ \Longrightarrow\ Q*R\ne\varnothing.
\]
Mutual comparison is an equivalence relation compatible with product.
Its classes form a partially ordered commutative monoid with at most
$2^{|\Omega_D|}$ elements.
\item If $A'$ is obtained from $A$ by contractions of private
coordinates and restrictions of private coordinates to zero, and its
zero remaining-private-coordinate section is still exactly $D$, then
\[
 \mathcal R(A)\preceq\mathcal R(A').
\]
If $X=A\cup_D B$ is an ample forbidden pc-minor and the corresponding
operation on $X$ is proper, this comparison is strict modulo mutual
comparison.
\end{enumerate}
\end{theorem}
\begin{proof}
The first coordinate $T_A$ records exactly the supports that can
prevent ample gluing. Indeed, private coordinate sets are disjoint,
so the common shattered supports of two pieces lie in $2^E$.
Proposition~\ref{prop:damp-uso-cube-gluing} says that their union is
ample exactly when $T_A\cap T_B=\varnothing$. If it is ample, its
shattered supports are the union of the piece shadows, as shown in
that proposition's proof. Its first record coordinate is consequently
$T_A\cup T_B$.

An order on a piece restricts to a corner order on $D$, by coordinate
face restriction. In an acyclic USO of the union, restriction to the
two pieces supplies a common boundary order and disjoint outward sets;
the outward set of the union is their union. Conversely the proof of
Corollary~\ref{cor:damp-finite-boundary-gluing} combines any two records
with a common order and disjoint outward sets into an acyclic USO,
with exactly that union of outward sets. This proves the record identity
when the union is ample. If the shadows conflict, every record pair
is excluded by its first coordinates. If either piece is nonpeelable,
the union cannot be ample and peelable, since that piece is a coordinate
face of the union. Thus the identity and positivity test also cover
all zero cases.

In a product of three formal record families, a resulting record is
specified by one common order, three pairwise disjoint first-coordinate
sets, and three pairwise disjoint outward sets. These conditions and
the two unions are independent of parenthesization, proving associativity.
Commutativity, the zero, and the identity follow directly from
\eqref{eq:damp-record-product}. The identity is realized by $D$ itself.
The record identity shows that a product of realized families is either
realized by their ample union or is zero. Hence $\mathfrak R_D$ is a
submonoid of the finite monoid $2^{\Omega_D}$.

The relation $\preceq$ is reflexive and transitive. If $P\preceq Q$
and $H\in\mathfrak R_D$, apply its definition to $H*R$, for each
$R\in\mathfrak R_D$, to obtain $P*H\preceq Q*H$.
It follows that mutual comparison is a product congruence and the
induced relation on equivalence classes is a partial order. Finiteness
and the stated bound follow from the ambient monoid.

For the minor assertion, choose any piece $B$ in fresh private coordinates.
Every allowed operation on $A\cup_D B$ fixes $B$ and induces the specified
operation on $A$. Thus $A'\cup_D B$ is a pc-minor of $A\cup_D B$.
Ample corner-peelable families are closed under these minors, proving the
implication defining $\preceq$ for every realized context; the zero
context imposes no condition. Finally, if $X$ is forbidden and the
resulting minor is proper, then
\[
 \mathcal R(A)*\mathcal R(B)=\varnothing,\qquad
 \mathcal R(A')*\mathcal R(B)\ne\varnothing.
\]
The context $B$ therefore excludes the reverse comparison, proving
strictness.
\end{proof}

This comparison keeps the effect of a piece on every gluing context,
rather than individual boundary orders. It is monotone even in
Proposition~\ref{prop:damp-boundary-attainability-failure}, where a raw
order disappears. For a cube boundary $U_D$ is empty, so the first
record coordinate is unnecessary, and every private-coordinate operation
in item~3 preserves the marked face. For a noncube boundary the first
coordinate enforces the shadow condition and must be retained. Moreover,
a private contraction can enlarge the zero private-coordinate section;
item~3 applies only when the resulting section is still $D$. No comparison
between different marked faces is asserted.

The algebra is finite and explicit, but its realizable submonoid
$\mathfrak R_D$ has not been intrinsically characterized. Nor does the
finite number of comparison classes bound the sizes of pieces minimal
for a class: proper minors may immediately change class, however large
the original piece is. Thus the theorem gives a minor-monotone comparison
and a necessary strictness test for forbidden assemblies, not an exhaustive
piece generator. The new content beyond the boundary compatibility
criterion is its associative composition and its comparison under all
realizable contexts; no preservation of individual records is asserted.

\begin{remark}[A square already permits a cyclic compatible union]
\label{rem:damp-square-gluing-cycle}
The acyclicity condition in
Proposition~\ref{prop:damp-uso-cube-gluing} cannot be dropped for a square.
Use binary integer labels and let $D=\{0,1,2,3\}$. Orient a first
three-cube $A=\{0,\ldots,7\}$ by the corner order
\[
 4,0,1,5,7,6,2,3.
\]
Orient a second copy by the corner order $0,2,4,6,7,5,1,3$, then
replace its bit $2$ by a fresh bit $3$. The two orientations agree
on the diamond $0\to1,0\to2,1\to3,2\to3$ in $D$.
Their boundary vertices with outgoing exclusive edges are respectively
$\{1\}$ and $\{0,2\}$, which are disjoint. Their union is therefore
a USO, but contains the directed cycle
\[
 1\to5\to7\to6\to2\to10\to11\to9\to1.
\]
The first piece forces $1$ before $2$ in every topological order,
whereas the second forces $2$ before $1$. They cannot supply the
same boundary order in Corollary~\ref{cor:damp-finite-boundary-gluing}.
The union family is $Q_2\square P_3$ and is corner-peelable; the
failure concerns these prescribed orientations, not the support.
The verifier \path{research/cube_separator_cycle.py/json} replays both
piece orders and independently checks (C1), (C2), and the directed cycle
in the union. Thus the automatic acyclicity for singleton and edge
intersections is sharp at the next cube size.
\end{remark}

\begin{theorem}[Pieces at a separating cube]
\label{thm:damp-cube-separator-bound}
Let \(D=Q_E\), and let \(A_1,\ldots,A_m\) be ample families
containing \(D\), on coordinate sets \(E\mathbin{\dot\cup}F_i\),
where the sets \(F_i\) are pairwise disjoint.  Embed each family by
putting zero in the other blocks, and put \(X=\bigcup_i A_i\).
Then \(X\) is ample.  If \(X\) is corner-peelable, at most
\(|D|\) of the pieces \(A_i\) fail to peel to \(D\).

Now suppose \(X\) is an ample forbidden pc-minor and \(D\) is
a separating \(k\)-cube.  The components \(K_i\) of \(X\setminus D\)
give a presentation of the preceding form with
\(A_i=D\cup K_i\).  Each piece is a proper convex restriction,
belongs to \(\Damp\), and fails to peel to \(D\).  Moreover,
\begin{equation}
 2\le m\le |D|+1=2^k+1.
 \label{eq:damp-cube-separator-bound}
\end{equation}
The pieces and their exclusive coordinate sets are uniquely recovered
from the specified separating cube, up to permutation of the pieces.
\end{theorem}

\begin{proof}
Ampleness follows from Proposition~\ref{prop:damp-uso-cube-gluing}.
If $X\in\Damp$, choose an acyclic USO and restrict it to each piece.
A piece failing to peel to $D$ must have an outgoing edge from some
$d\in D$ to its exclusive part, by
Proposition~\ref{prop:damp-relative-uso-extension}. By the compatibility
law, different such pieces use disjoint nonempty sets of boundary vertices.
There are at most $|D|$ of them. This does not require $A_i\setminus D$
to be connected.

For the recovery assertion, normalize the separating cube by making
its fixed coordinates zero, and let \(E\) be its varying coordinates.
A coordinate \(f\notin E\) cannot take value one in two different
components of \(X\setminus D\).  A path between such concepts
must pass through \(D\), where \(f=0\), and hence must change
\(f\) twice, contrary to the existence of a cube geodesic.
Assign each coordinate outside \(E\) to its unique component.
The piece \(D\cup K_i\) is exactly the restriction setting all
other exclusive coordinates to zero.  The same holds for the union
of any selected pieces.  These are convex restrictions and have the
claimed coordinate form.

Each piece belongs to \(\Damp\) by proper-minor minimality.
If one peeled to \(D\), its outside concepts could first be removed
from \(X\), followed by a peeling of the proper convex restriction
formed by the other pieces.  This would peel \(X\), a contradiction.
Thus every piece fails to peel to \(D\).  If there were more than
\(|D|+1\) pieces, the union of any \(|D|+1\) of them would be
a proper convex restriction which cannot peel, by the first part.
This proves \eqref{eq:damp-cube-separator-bound}.  The component
construction also proves uniqueness for the marked cube.
\end{proof}

\begin{proposition}[Pieces at a coordinate face]
\label{prop:damp-face-separator-bound}
Let $X$ be ample and let $D$ be a nonempty coordinate-face intersection
of $X$. Normalize the fixed coordinates of $D$ to zero, and let $E$ be
its free coordinate set. The components $W_1,\ldots,W_m$ of $X\setminus D$
have disjoint exclusive coordinate blocks $F_i$ outside $E$.
Each $A_i=D\cup W_i$, and each union of selected such pieces, is a
coordinate restriction of $X$.
If $X\in\Damp$, at most $|D|$ of the $A_i$ fail to peel to $D$.
If $X$ is forbidden and $m\ge2$, every $A_i$ is positive, none peels
to $D$, and $2\le m\le |D|+1$.
\end{proposition}
\begin{proof}
This extends the recovery and counting argument of
Theorem~\ref{thm:damp-cube-separator-bound}; $D$ need not itself be a cube.
If a fixed coordinate $j$ takes value one in two outside components,
a geodesic between the corresponding vertices stays in the $j=1$
semicube and avoids $D$, contradicting their separation. Thus each active
fixed coordinate belongs to one component. Every outside vertex has a
nonzero fixed coordinate, and all its nonzero fixed coordinates belong
to its component. Setting the other blocks to zero consequently recovers
$D\cup W_i$, and likewise recovers any selected union of pieces.

Suppose $X$ is positive and restrict an acyclic USO to each of these
coordinate faces. If no edge at a vertex of $D$ points into $W_i$,
a topological order can put all of $W_i$ before $D$, proving relative
peeling. Thus each piece failing relative peeling must use a boundary
vertex with an outgoing exclusive edge. Two different pieces cannot use
the same vertex: the square in their two exclusive directions is absent,
contradicting (C1). This gives the bound $|D|$.

If $X$ is forbidden and $m\ge2$, every piece is a proper restriction
and hence positive. A piece peeling to $D$ could be removed first from
$X$, since its outside vertices have no neighbours in the other pieces.
The remaining proper restriction is positive, a contradiction. If
$m>|D|+1$, the union of any $|D|+1$ pieces is a proper positive restriction,
contradicting the preceding count. This proves the assertion.
\end{proof}

\begin{proposition}[Projections of pieces outside a coordinate face]
\label{prop:damp-face-piece-projections}
In Proposition~\ref{prop:damp-face-separator-bound}, assume $X\ne D$,
let $p$ project onto
the free coordinates $E$, put $B=p(X)$ and $C_i=p(A_i)$, and identify
$D$ with $p(D)$. Then $B$ and the $C_i$ are ample,
\[
 B=\bigcup_i C_i,\qquad C_i\cap C_j=D\quad(i\ne j),
\]
and no edge joins $C_i\setminus D$ to $C_j\setminus D$ for $i\ne j$.
Thus each $C_i\setminus D$ is a union of components of $B\setminus D$.
If every coordinate in $E$ varies on $D$, then $D$ is gated in $A_i$
if and only if $C_i=D$. Consequently the number of pieces in which $D$
is not gated is at most the number of components of $B\setminus D$.
\end{proposition}
\begin{proof}
Ampleness follows by coordinate contraction. Suppose distinct pieces have
projected vertices $u,v\notin D$ that are equal or adjacent. Choose
preimages in their outside components. A geodesic between them in $X$
has $E$-trace constantly $u$ if $u=v$, and otherwise only $u$ or $v$:
an isometric cube path never changes a coordinate on which its endpoints
agree. This geodesic therefore avoids $D$, contradicting that its endpoints
lie in different components of $X\setminus D$. This proves the intersection
and edge assertions, and hence the component statement.

Recall that a gate of $x$ in $D$ is a vertex $g\in D$ satisfying
$d(x,z)=d(x,g)+d(g,z)$ for every $z\in D$; $D$ is gated when each
vertex has a gate. If $p(x)\in D$, the Hamming-distance identity makes
$p(x)$ a gate of $x$. Thus $C_i=D$ implies gatedness. Conversely, let
$g$ be a gate of $x\in A_i$. If $g$ and $x$ differed in some coordinate
of $E$, choose $z\in D$ with the bit of $x$ there, using variation on
$D$. The route through $g$ would change this coordinate twice, contradicting
the gate equality in the isometric cube metric. Thus $g=p(x)$, proving
$p(x)\in D$ and the converse. Each nongated piece uses at least one of
the disjoint components already identified, giving the bound.
\end{proof}

For an edge, the bound leaves at most two pieces that can fail relative
peeling in a peelable union.  Whether those two pieces are compatible
depends on the order in which the endpoints are removed.

\begin{definition}[Acyclic USO boundary conditions at an edge]
\label{def:damp-edge-peeling-profile}
Let $D=\{u,v\}$ be an edge in direction $e$ of an ample family $A$.
Define five Boolean values, each quantifying over acyclic USOs of $A$
with out-map $r$:
\begin{itemize}
 \item $r_D(A)=1$ if some such map satisfies $r(u),r(v)\subseteq\{e\}$;
 \item $t_s(A)=1$, for $s\in D$, if some such map satisfies $r(s)=\varnothing$;
 \item $\ell_s(A)=1$, for $s\in D$, if some such map satisfies $r(s)=\{e\}$.
\end{itemize}
For two such families, write $\chi_D(A,B)=1$ if some ordering $(s,t)$
of the endpoints satisfies
\begin{equation}
 (t_t(A)=1\text{ and }\ell_s(B)=1)
 \quad\text{or}\quad
 (t_t(B)=1\text{ and }\ell_s(A)=1).
 \label{eq:damp-edge-endpoint-compatibility}
\end{equation}
Otherwise put $\chi_D(A,B)=0$.
\end{definition}

These retain the original deletion interpretations. By
Proposition~\ref{prop:damp-relative-uso-extension}, $r_D$ means peeling to
$D$, and $t_s$ means peeling to $s$. In a topological deletion order the
remaining neighbours of a deleted vertex are precisely its outgoing
neighbours. Thus $\ell_s$ means a complete peeling in which $s$ is the
first endpoint deleted and is then a leaf adjacent to the other endpoint.
Conversely each such peeling induces the required out-map value.


These five values record attainable boundary conditions on acyclic USOs.
The following theorem makes their gluing effect exact; intrinsic structural
tests for attainability remain open.

\begin{theorem}[Peeling after gluing along an edge]
\label{thm:damp-edge-gluing-profile}
Let \(A_1,\ldots,A_m\in\Damp\) have the coordinate-block
presentation of Theorem~\ref{thm:damp-cube-separator-bound}, with
common edge \(D\).  Put
\(I=\{i:r_D(A_i)=0\}\).  Then
\begin{equation}
 \bigcup_i A_i\in\Damp
 \quad\Longleftrightarrow\quad
 |I|\le1\quad\text{or}\quad
 \bigl(I=\{i,j\}\text{ and }\chi_D(A_i,A_j)=1\bigr).
 \label{eq:damp-edge-gluing-profile}
\end{equation}
\end{theorem}

\begin{proof}
Suppose the union has an acyclic USO, and orient its common edge $s\to t$.
Each piece in $I$ must have outgoing exclusive directions at some endpoint.
Proposition~\ref{prop:damp-uso-cube-gluing} permits at most one such piece
at each endpoint, so $|I|\le2$. If $I=\{i,j\}$, the two pieces must
use different endpoints. The piece using $s$ has no outgoing edge at $t$,
so realizes $t_t=1$. The piece using $t$ has just the common outgoing
edge at $s$, so realizes $\ell_s=1$. This is $\chi_D(A_i,A_j)=1$.

Conversely, every piece outside $I$ can extend either prescribed orientation
of $D$ with inward boundary, by
Proposition~\ref{prop:damp-relative-uso-extension}. If $|I|\le1$, choose
an acyclic USO of that possible exceptional piece and extend its edge
orientation to the others; if $I$ is empty, choose either edge orientation.
The compatibility law gives an acyclic USO of the union. If
$t_t(A_i)=\ell_s(A_j)=1$, the witnessing maps both orient $s\to t$.
Only $A_i$ can have outgoing exclusive directions at $s$, and only $A_j$
can have them at $t$. Extend the common edge orientation to all other
pieces with inward boundary. Again Proposition~\ref{prop:damp-uso-cube-gluing}
gives an acyclic USO, hence a peeling of the union.
\end{proof}

\begin{proposition}[Prescribed sinks after attachments along edges]
\label{prop:damp-edge-attachment-leaf-constraints}
Let $S$ be ample. Along each of a collection of distinct edges $E_i$ of
$S$, attach an ample family $P_i\in\Damp$, using private coordinate
sets disjoint from those of $S$ and of all other pieces. Each $P_i$
meets $S$ exactly in $E_i$, and fixes the other coordinates of $S$ at
their values on that edge. Write $Y=S\cup\bigcup_i P_i$, and let $T$
be a vertex or an edge of $S$.

Discard from the list of constraints every piece that peels to $E_i$.
For each remaining piece, let $H_i\subseteq E_i$ be its set of
attainable endpoints. Then $Y$ peels to $T$ if and only if $S$ has an
acyclic USO with out-map $r$, inward boundary at $T$, and, for each
remaining piece, an ordering $(u_i,v_i)$ of $E_i$ such that
\begin{equation}
 v_i\in H_i,\qquad u_i\notin T,\qquad r(u_i)=\{e_i\},
 \label{eq:damp-edge-attachment-leaf-constraints}
\end{equation}
where $e_i$ is the direction of $E_i$.
Thus a singleton $H_i$ specifies the outgoing edge at its other endpoint;
if $H_i=E_i$, either direction may realize this condition. An empty $H_i$
precludes peeling to $T$.
\end{proposition}
\begin{proof}
Successive edge gluings show that $Y$ is ample. The central family $S$
and every $P_i$ are coordinate restrictions of $Y$.
Suppose $Y$ peels to $T$, and take an acyclic USO with inward boundary
at $T$ using Proposition~\ref{prop:damp-relative-uso-extension}.
Its unique sink belongs to $T$. The restriction to $P_i$ cannot have
a sink in $P_i\setminus E_i$: such a vertex has no neighbours outside
$P_i$ and would be a global sink outside $T$. Its sink is therefore an
attainable endpoint $v_i$, and the common edge points from the other
endpoint $u_i$ to $v_i$.

If $P_i$ cannot peel to $E_i$, its out-map at $u_i$ has a private
outgoing direction. Otherwise its boundary would be inward at $E_i$,
since $v_i$ is a sink. There is no square mixing this private direction
with a direction of $S$ other than $e_i$. Condition (C1) consequently
forces the restricted out-map on $S$ to satisfy $r(u_i)=\{e_i\}$.
The private outgoing edge also gives $u_i\notin T$. This proves
necessity of all displayed conditions.

Conversely, choose the asserted map on $S$. For a nonretaining piece
choose an acyclic USO with sink $v_i$. Its only possible boundary vertex
with outgoing private directions is $u_i$, and its common edge agrees
with the central orientation. No two such pieces can have the same
$u_i$: the central map would then specify two different singleton
directions at that vertex, since the attachment edges are distinct.
For a piece retaining its attachment edge, extend the central edge
orientation with inward boundary. The outgoing directions are therefore
compatible at every attachment vertex. Successively apply the acyclic
edge-gluing law of Proposition~\ref{prop:damp-uso-cube-gluing}.
This gives an acyclic USO of $Y$; all outgoing private directions at
attachment vertices occur outside $T$. Its boundary at $T$ is inward,
so $Y$ peels to $T$.
\end{proof}

\begin{proposition}[Edge attachments to a fixed path expansion; computer-assisted]
\label{prop:damp-eleven-vertex-edge-attachments}
Encode four coordinates by weights $1,2,4,8$, and put
\[
 S=\{0,1,2,3,7,8,9,11,12,13,15\},\qquad T=\{0,8\}.
\]
Attach arbitrary positive ample pieces to distinct edges of $S$ whose
direction is not $8$, with the private-coordinate convention of
Proposition~\ref{prop:damp-edge-attachment-leaf-constraints}.
If the resulting family peels to $0$ and to $8$, it peels to $T$.
Its reduction in direction $8$ is still the path $0-1-3-7$.
\end{proposition}
\begin{proof}
We give the finite reduction and its exhaustive verification. For an
acyclic USO of $S$, record two kinds of conditions at each allowed edge
$\{u,v\}$: the directed condition $r(u)=\{e\}$, and the undirected
condition $r(u)=\{e\}$ or $r(v)=\{e\}$. In a record for retaining
$T$, only endpoints outside $T$ may realize either condition. Directed
conditions with tail in $T$ cannot hold in both prescribed-sink maps,
so they can be omitted throughout the comparison.

There are $16$ edges in $S$. The verifier
\path{research/edge_attachment_constraints.py/json} checks all $2^{16}$
orientations using the outgoing-cube condition, nonclashing, and
acyclicity. It obtains $784$ acyclic USOs, independently recovered by
enumerating corner deletions and identifying equal out-maps. Of these,
$69$ have sink $0$, $69$ have sink $8$, and $72$ have inward boundary
at $T$. The intersections of the two prescribed-sink records give
$79$ distinct sets of conditions. Every one is contained in the record
of a map retaining $T$; the report stores a witnessing map for each.

If both prescribed sinks are attainable in the assembled family,
the preceding proposition gives one central map for each sink, both
satisfying the attachment conditions. Each nonretaining piece contributes
a directed condition when it has one attainable endpoint and an
undirected condition when it has two. It cannot have none. These
conditions lie in the intersection of the two records, and hence in
a record retaining $T$. Apply the preceding proposition again to obtain
the required relative peeling. Pieces retaining their attachment edges
impose no conditions. Finally each attachment fixes the coordinate of
weight $8$, and its private coordinates are disjoint from the other
pieces. Thus it creates no additional doubled fibre in that direction.
\end{proof}

The central family in the last proposition is the expansion of the
punctured cube $\{0,1,2,3,4,5,7\}$ that places the complementary
components $\{2\}$ and $\{4,5\}$ on opposite sides of $0-1-3-7$.
It supplies the desired separating geometry, but even arbitrary positive
edge attachments cannot turn it into a both-endpoints retention failure.
This is an exclusion of that assembly class, not of attachments along
larger faces or of arbitrary contractions. The preceding general
proposition reduces large edge pieces to small central constraints;
it does not intrinsically generate all pieces with specified endpoint sets.

\begin{proposition}[Square attachments retaining a punctured boundary]
\label{prop:damp-square-source-attachment}
Let $S$ be ample and let $T$ be a vertex or edge of $S$. Attach ample
pieces $P_i$ along squares $Q_i$ of $S$, fixing the other coordinates
of $S$ and using pairwise disjoint private coordinate sets. Suppose an
acyclic USO $r$ of $S$ has inward boundary at $T$. For each $i$, choose
$u_i\in Q_i\setminus T$ so that $P_i$ peels to $Q_i\setminus\{u_i\}$
and $r(u_i)$ consists exactly of the two directions of $Q_i$.
If the chosen vertices $u_i$ are distinct, the assembled family peels to $T$.
\end{proposition}
\begin{proof}
The square $Q_i$ has source $u_i$ in the central orientation. Deleting
that source leaves its three-vertex path with an acyclic USO. Relative
extension lifts this orientation to $P_i$ with inward boundary at that
path. Both square edges at $u_i$ point away from it, so the orientations
agree on all of $Q_i$. Only $u_i$ can have outgoing private directions
among the four boundary vertices. These directions are compatible with
the central map because its outgoing directions at $u_i$ are exactly
those of $Q_i$. Distinct choices of $u_i$ prevent conflicts between
private directions of different pieces. All contained cubes lie in $S$
or in one piece, giving the remaining USO conditions.

We check acyclicity, which does not follow from arbitrary square gluing.
A directed excursion into the private vertices of $P_i$ starts at
$u_i$ and, if it returns to the boundary, ends at a different vertex
$v\in Q_i$. A square USO with source $u_i$ has a directed path from
$u_i$ to every other square vertex. Replace each private excursion of
a hypothetical directed cycle by such a central path. This gives a
directed closed walk in the acyclic central orientation, a contradiction.
A cycle contained in a piece is already excluded by the chosen acyclic
extension. Finally no chosen source belongs to $T$, so the boundary at
$T$ is inward in the union. Relative extension gives the claimed peeling.
\end{proof}

For the eleven-vertex family $S$ above, consider its three horizontal
squares, in the displayed local binary order,
\[
 Q_0=(0,1,2,3),\quad Q_1=(8,9,12,13),\quad
 Q_2=(9,11,13,15).
\]
Let $h_i\in\{0,1\}$ specify the opposite pair of local indices
$\{h_i,3\mathbin{\mathrm{xor}}h_i\}$. Suppose each attached piece
can retain its square with either member of that pair omitted.
For every pattern except $(h_0,h_1,h_2)=(0,0,1)$, the union peels to
$\{0,8\}$. The following central out-maps certify this assertion.
Entries are supports encoded as masks, in vertex order
$(0,1,2,3,7,8,9,11,12,13,15)$.
\[
\begin{array}{c|c|l}
(h_0,h_1,h_2)&(u_0,u_1,u_2)&r\\\hline
(0,0,0)&(3,13,15)&(8,9,2,3,12,0,1,10,4,5,6)\\
(0,1,0)&(3,12,15)&(8,9,2,3,12,0,1,10,5,4,6)\\
(0,1,1)&(3,12,13)&(8,9,2,3,12,0,1,10,5,6,4)\\
(1,0,0)&(2,13,15)&(8,9,3,10,12,0,1,2,4,5,6)\\
(1,0,1)&(2,13,11)&(8,9,3,10,12,0,1,6,4,5,2)\\
(1,1,0)&(2,12,15)&(8,9,3,10,12,0,1,2,5,4,6)\\
(1,1,1)&(2,12,13)&(8,9,3,10,12,0,1,2,5,6,4)
\end{array}
\]
Each row is checked by a corner order and the nonclashing condition in
\path{research/square_attachment_sources.py/json}.
The chosen vertices have out-maps $3,5,6$, respectively, so the
proposition applies. The square case of
Theorem~\ref{thm:damp-cube-boundary-retention} supplies concrete pieces
of order $1165$. In fact these concrete pieces retain all four punctured
boundaries, as that theorem also proves. Consequently its first displayed
central map, with sources $(3,13,15)$, applies to three such pieces in
\emph{every} alignment. All resulting $3494$-vertex assemblies retain
$\{0,8\}$, including the hole-pair pattern omitted from the table.

The omitted pattern does not admit these source choices. The target
excludes $u_0=0$ and $u_1=8$, forcing $u_0=3$, $u_1=13$.
Choosing $u_2=13$ would require both $r(13)=5$ and $r(13)=6$.
Choosing $u_2=11$ would require $r(11)=6$, orienting the edge $3,11$
toward $11$, whereas $r(3)=3$ orients it toward $3$. This proves
failure of the sufficient source condition only. It does not prove
failure of edge retention or attainability of either endpoint in the
assembled family under the weaker hypothesis of just two prescribed
punctured-boundary orders. The concrete cube-boundary pieces avoid this
limitation by their additional orders, so they cannot supply a
both-endpoints retention failure on this template.

\begin{corollary}[Six contexts determine every edge-gluing type]
\label{cor:damp-edge-context-algebra}
Fix an ordered edge $D=\{u,v\}$ and use the context comparison of
Theorem~\ref{thm:damp-finite-gluing-algebra}. A piece has one of the
following types:
\begin{itemize}
\item $\mathbf 0$ if it is not corner-peelable;
\item $\mathbf 1$ if it peels to $D$;
\item $[p]$, where $p=(t_u,t_v,\ell_u,\ell_v)\in\{0,1\}^4$,
if it is corner-peelable but does not peel to $D$.
\end{itemize}
Two pieces have the same context type if and only if they have the same
entry in this list. In particular, there are at most eighteen types;
this does not assert that all sixteen vectors $p$ are realized.
Among realized vector types,
\[
 [p]\preceq[q]\quad\Longleftrightarrow\quad p\le q
 \text{ coordinatewise},
 \qquad \mathbf0\prec[p]\prec\mathbf1.
\]
Write $Z=[(0,0,0,0)]$. The multiplication law is
\begin{equation}
 [p]*[q]=
 \begin{cases}
 Z,&p_{t_u}q_{\ell_v}\lor p_{t_v}q_{\ell_u}
       \lor q_{t_u}p_{\ell_v}\lor q_{t_v}p_{\ell_u}=1,\\
 \mathbf0,&\text{otherwise},
 \end{cases}
 \label{eq:damp-edge-type-product}
\end{equation}
with identity $\mathbf1$ and absorbing element $\mathbf0$.
Every product of three vector types is $\mathbf0$.

Six fixed positive pieces suffice to distinguish all the types and test
the comparison: the edge $D$, a piece of type $Z$, and four pieces realizing
the four unit vectors in $\{0,1\}^4$. These pieces are supplied by
Remark~\ref{rem:damp-two-edge-failures-positive} and its positive union,
together with endpoint reversals.
\end{corollary}
\begin{proof}
The edge-gluing theorem makes a piece retaining $D$ neutral for ordinary
positivity, and a nonpeelable piece makes every union nonpeelable by
coordinate-face closure. For two positive pieces failing retention,
that theorem is exactly the Boolean test in
\eqref{eq:damp-edge-type-product}.

Suppose such a pair has a positive union. Its union cannot retain $D$,
since relative peeling restricts to either piece. We show that all four
boundary values of the union vanish. If an acyclic USO made $u$ a sink,
its restrictions would make $u$ a sink in both pieces. Each piece must
have an outward edge at some boundary vertex, or it would retain $D$.
Both would therefore have outward edges at $v$, violating the gluing
condition. The same argument excludes a sink at $v$. If a map instead
had out-map $\{e\}$ at $u$, where $e$ is the edge direction, neither
piece could have an outward edge at $u$. Both would again require
outward edges at $v$, a contradiction. Interchanging endpoints excludes
the fourth boundary value. Thus the positive product has type $Z$.
The displayed multiplication law follows, and multiplication of $Z$
by any vector type gives $\mathbf0$, proving the three-factor assertion.

The two pieces in Remark~\ref{rem:damp-two-edge-failures-positive}
have vector types $(0,1,0,0)$ and $(0,0,1,0)$; reversing endpoints supplies
$(1,0,0,0)$ and $(0,0,0,1)$. Their known positive union has type $Z$
by the preceding paragraph. Gluing a vector-type piece to these four
unit-vector pieces tests, respectively, its four dual Boolean entries:
$t_u$ is tested by the $\ell_v$ unit, $t_v$ by the $\ell_u$ unit,
and conversely. Hence different vectors are distinguished by realized
contexts, and a failed coordinatewise comparison is likewise witnessed
by one of these four contexts.

Coordinatewise inclusion is sufficient for comparison against every
context, by the displayed Boolean formula and the two exceptional types.
The edge distinguishes $\mathbf0$ from every positive type. The realized
piece $Z$ distinguishes $\mathbf1$ from every vector type: it glues
positively to the former and negatively to the latter. These observations
also establish the strict comparisons and show that the six stated
contexts determine both equivalence and comparison for all pieces.
\end{proof}

Thus edge context comparison needs no unknown collection of test pieces.
The required sink and leaf values themselves still express attainability
of acyclic USO boundary conditions. The corollary neither supplies intrinsic
tests for those values nor classifies their minor-minimal realizing pieces.
It sharpens the finite gluing algebra at an edge using the existing profile
theorem and existing witnesses; it does not enlarge the known forbidden
families. The elementary algebra is checked independently in
\path{research/edge_context_algebra.py/json}, allowing all sixteen formal
vectors without claiming their realizability.

\begin{theorem}[Forbidden minors with a separating edge]
\label{thm:damp-edge-separator-obstructions}
Let \(m\ge2\), and let \(A_1,\ldots,A_m\) be irredundantly
embedded ample families whose
coordinate sets intersect pairwise in precisely the coordinate of
\(D=\{u,v\}\), and whose common edge is \(D\).  Normalize
\(u=0\), and suppose each \(A_i\setminus D\) is nonempty
and connected.  Their union \(X\) is an ample forbidden pc-minor
if and only if all the following conditions hold.
\begin{enumerate}
 \item \(m\in\{2,3\}\), every \(A_i\in\Damp\), and
       \(r_D(A_i)=0\) for every \(i\).
 \item If \(m=2\), then \(\chi_D(A_1,A_2)=0\).
       If \(m=3\), then \(\chi_D(A_i,A_j)=1\) for every
       pair of distinct indices.
 \item For each exclusive coordinate \(f\) of \(A_i\) and
       each \(\mu\in\{\rho_f^0,\pi_f\}\), the following holds:
       when \(m=2\), writing \(j\ne i\) for the other index,
       \begin{equation}
        r_D(\mu A_i)=1\quad\text{or}\quad
        \chi_D(\mu A_i,A_j)=1;
        \label{eq:damp-two-piece-exclusive-minor}
       \end{equation}
       when \(m=3\), one has \(r_D(\mu A_i)=1\).
 \item For each of the two restrictions and the contraction in
       the shared coordinate, at most one image \(\mu A_i\)
       fails to peel to the singleton \(\mu D\).
\end{enumerate}
Every ample forbidden pc-minor with a specified separating edge arises
this way.  Its pieces are recovered uniquely as that edge united with
each component of its complement, up to permutation.  The union is
two-connected if and only if all its pieces are two-connected.
\end{theorem}

\begin{proof}
The union is ample by Theorem~\ref{thm:damp-cube-separator-bound}.
Its components outside \(D\) are exactly the stated
\(A_i\setminus D\).

Suppose \(X\) is forbidden.  The cube-separator theorem gives
item~1.  For two pieces, nonpeelability and the exact gluing law give
\(\chi_D(A_1,A_2)=0\).  For three pieces, each two-piece union is
a proper convex restriction and peels; the same law gives all three
compatibilities in item~2.
An elementary exclusive-coordinate restriction to zero or contraction
changes only \(A_i\) and retains \(D\).  If there are two pieces,
the exact gluing law is precisely
\eqref{eq:damp-two-piece-exclusive-minor}.  If there are three, the
other two pieces still fail relative peeling; thus the changed piece
must peel to \(D\).  Finally, an elementary operation in the shared
coordinate changes the intersection to a singleton, and
Proposition~\ref{prop:damp-vertex-gluing} gives exactly item~4.

Conversely, items~1 and 2 make \(X\) nonpeelable by
Theorem~\ref{thm:damp-edge-gluing-profile}.  A restriction to one
in an exclusive coordinate discards every other piece and is a
pc-minor of the peelable \(A_i\).  The two other elementary operations
in that coordinate give a peelable union by item~3 and the edge-gluing
law; in the three-piece case the unchanged pair is compatible by
item~2.  The shared-coordinate images peel by item~4 and the
vertex-gluing law.  Every elementary pc-minor therefore peels, so
\(X\) is forbidden.

For an arbitrary obstruction with a separating edge, the cube-separator
theorem recovers the coordinate blocks and their convex pieces.  These
meet the presentation hypotheses, and the first implication proves their
conditions.  Their recovery is canonical for the specified edge.

If every piece is two-connected, deleting any vertex leaves connected
residual pieces sharing at least one endpoint of \(D\), so the
union is two-connected.  Conversely, if deleting a vertex disconnects
a piece, the surviving endpoint or edge \(D\) lies in one
component, and any other component has no path to the rest of
\(X\).  That vertex is then a cut vertex of \(X\).  This proves
the final equivalence.
\end{proof}

This is an exact decomposition of the separating-edge branch. Its five
boundary values are still peeling predicates. The three-piece case is
realized by Theorem~\ref{thm:damp-three-piece-witness} below, so the bound
of three pieces is sharp. The three-piece case
forces every exclusive elementary image to peel to the edge.  In the
two-piece case, compatibility with the unchanged piece is essential:
Remark~\ref{ex:damp-coupled-edge-minors-1162} gives a forbidden
two-connected graph three-core for which the changed factor still
fails relative peeling while the full proper image peels.

The preceding decomposition assumes that a separating edge has been
specified.  Shattered pairs detect exactly which coordinate directions
can contain such an edge, and coordinate contractions recover the edges
themselves.  For an irredundant partial cube \(H\subseteq Q_F\), write
\(\operatorname{Cr}(H)\) for its \emph{crossing graph}: its vertices are
\(F\), and distinct \(f,g\) are adjacent when
\(\{f,g\}\in\operatorname{Sh}(H)\).  This is the usual crossing
graph, since shattering a pair means that all four intersections of its
coordinate semicubes are nonempty.  We use the usual convention that a
two-connected graph has at least three vertices.

\begin{proposition}[Detecting and recovering separating edges]
\label{prop:damp-crossing-separator-recovery}
Let \(H\subseteq Q_F\) be an irredundant partial cube with
\(|F|\ge3\).  Then \(H\) has a cut vertex if and only if
\(\operatorname{Cr}(H)\) is disconnected.
If \(H\) is two-connected, the following assertions hold for each
\(e\in F\).
\begin{enumerate}
 \item The separating edges of direction \(e\) are exactly the fibres
 \(\pi_e^{-1}(z)\cap H\), where \(z\) ranges over the cut vertices
 of \(\pi_e H\).  Components outside an edge correspond to components
 outside its image by projection and inverse image.
 \item Such an edge exists if and only if
 \(\operatorname{Cr}(H)-e\) is disconnected.
 \item If \(\operatorname{Cr}(H)-e\) has exactly two components,
 with vertex sets \(F_1,F_2\), there is exactly one separating
 \(e\)-edge \(D\).  The two pieces obtained by adjoining \(D\)
 to the components of \(H\setminus D\) have exclusive coordinate
 sets \(F_1,F_2\), and are convex coordinate restrictions.
\end{enumerate}
Consequently an ample forbidden pc-minor has neither a cut vertex nor
a separating edge exactly when its crossing graph is two-connected.
\end{proposition}

\begin{proof}
The cut-vertex assertion is
\cite[Theorem~3.4]{klavzar2002partial}.  The dimension assumption excludes
the one-edge degeneracy, both for \(H\) and for its one-coordinate
contractions.  We derive the edge assertions by applying this theorem
after contraction and then lifting the resulting cut vertices.

Projection preserves shattering of every remaining coordinate pair, so
\[
 \operatorname{Cr}(\pi_e H)=\operatorname{Cr}(H)-e.
\]
Each fibre of \(\pi_e\) is a singleton or an \(e\)-edge, and is
therefore connected.  Moreover every edge of the projected graph lifts
to an edge between its two fibres.  Indeed, choose one preimage of each
endpoint.  They differ either in the projected edge direction alone or
in that direction and \(e\).  In the latter case a length-two geodesic
in \(H\) supplies the required edge.  The same argument applies after
deleting any whole fibre: all edges between surviving fibres still lift.
Thus the connected components of
\(H\setminus\pi_e^{-1}(z)\) correspond bijectively to those of
\((\pi_e H)\setminus\{z\}\).  A cut vertex \(z\) cannot have a
singleton fibre, since that singleton would be a cut vertex of \(H\).
This proves item~1.  The contraction is again an irredundant partial
cube, and the cited theorem now proves item~2.

For item~3, recall why the components of the crossing graph of a partial
cube are exactly the coordinate sets of its blocks, including bridge
blocks.  Here a block is a maximal two-connected subgraph or a bridge
with its endpoints.  Each block is convex: a path leaving a block can
return to it only through the same cut vertex.  A coordinate cannot
occur in two different blocks.  Otherwise choose a cut vertex separating
the two blocks and, in each of them, an endpoint of an edge of that
coordinate on the side opposite the cut vertex.  A geodesic between
these endpoints would have to change that coordinate twice.  Also, two
coordinates occurring in different blocks cannot cross: a cut vertex
separating the blocks prevents their simultaneously taking values
opposite its values.  Finally, the crossing graph of a two-connected
block is connected by the cited theorem, and that of a bridge block is
a singleton.  These facts identify the blocks with the crossing-graph
components.

It follows that \(\pi_e H\) has precisely two blocks, meeting at
one cut vertex \(z\).  Item~1 gives a unique separating \(e\)-edge
\(D\), and exactly two components outside it.  Normalize
\(D=\{0,\mathbf1_{\{e\}}\}\).  A coordinate \(f\ne e\) cannot
take value one in both components: a geodesic between such vertices
would pass through \(D\), where \(f=0\), and repeat direction
\(f\).  Assign \(f\) to its unique component.  Adjoining \(D\)
to a component is therefore exactly the restriction setting all
coordinates assigned to the other component to zero.  Its exclusive
coordinates are those of the corresponding quotient block, namely
\(F_1\) or \(F_2\).  This proves item~3.

An ample forbidden pc-minor has VC dimension at least three, since
ample classes of VC dimension at most two are corner-peelable
\cite[Proposition~4.11]{ChalopiChepoiMoran22}.  It therefore satisfies
the dimension assumption.  The cut-vertex assertion and item~2 say
that absence of both kinds of separator is equivalent to connectivity
of the crossing graph and absence of a cut vertex in it, as required.
\end{proof}

\begin{corollary}[Construction and recovery when a shadow splits in two]
\label{cor:damp-shadow-edge-factorization}
Let \(\Sigma\subseteq2^F\) be a simplicial complex containing every
singleton, and let its graph of two-element faces be connected.
Suppose deleting a vertex \(e\) from that graph leaves exactly two
components, with vertex sets \(F_1,F_2\).  Put
\[
 \Sigma_i=\Sigma\cap2^{F_i\cup\{e\}}\qquad(i=1,2).
\]
All ample realizations of \(\Sigma\) are obtained as follows: choose
ample realizations \(H_i\subseteq Q_{F_i\cup\{e\}}\) of
\(\Sigma_i\), choose an oriented \(e\)-edge in each, translate
both chosen edges to \(D=\{0,\mathbf1_{\{e\}}\}\), put zero in
the other exclusive coordinates, and take any cube translate of
\(H_1\cup H_2\).  Every output has a unique separating edge of direction \(e\), and
that edge recovers the two convex pieces, up to their order and a
simultaneous reversal of the chosen edge orientation.

Within this construction, the forbidden pc-minors are exactly the
pairs satisfying Theorem~\ref{thm:damp-edge-separator-obstructions}
with \(m=2\).
\end{corollary}

\begin{proof}
Translation preserves shadows, ampleness, and separators, so it suffices
to check the normalized union.  No face of \(\Sigma\) meets both
\(F_1\) and \(F_2\), since
such a face would contain a pair joining the two components.  Hence
\[
 \Sigma=\Sigma_1\cup\Sigma_2,
 \qquad \Sigma_1\cap\Sigma_2=\{\varnothing,\{e\}\}.
\]
The coordinate-block union has shadow \(\Sigma\) by the argument
of Theorem~\ref{thm:damp-cube-separator-bound}, and
\[
 |H_1\cup H_2|=|H_1|+|H_2|-2=|\Sigma|.
\]
It is therefore ample.  Its crossing graph is connected, so it is
two-connected.  Proposition~\ref{prop:damp-crossing-separator-recovery}
gives uniqueness of its separating \(e\)-edge.

We check that the chosen edge is this separating edge.  The graph of
\(\Sigma_i\) is connected: any path from a vertex of \(F_i\)
to \(e\) stays in \(F_i\) until its last step.  Thus \(H_i\)
is two-connected.  If \(|F_i|\ge2\), its crossing graph remains
connected after deletion of \(e\), so the proposition shows that
\(H_i\setminus D\) is connected.  If \(|F_i|=1\), connectivity
of the graph forces \(\Sigma_i\) to be the full two-coordinate
simplex, and \(H_i=Q_2\); deleting \(D\) leaves the opposite
edge.  In both cases the two nonempty sets \(H_i\setminus D\)
are exactly the components outside the chosen edge in the union.

Conversely, let \(H\) be any ample realization of \(\Sigma\).
The proposition recovers the unique separating \(e\)-edge and
its two convex pieces with exclusive coordinate sets \(F_1,F_2\).
Normalize this edge to \(D\).  Projection onto
\(F_i\cup\{e\}\) gives precisely the corresponding piece,
since the other piece projects into \(D\).  Projection preserves
shattering on these coordinates, so that piece has shadow
\(\Sigma_i\).  Convex coordinate restrictions preserve ampleness.
Choosing an endpoint of the recovered edge also recovers the final
translation.  Thus these are exactly the construction data.  The pieces meet all
presentation hypotheses of
Theorem~\ref{thm:damp-edge-separator-obstructions}, which proves
the last assertion.
\end{proof}

\begin{remark}[What the shadow determines]
\label{rem:damp-crossing-separator-scope}
The shadow determines separating-edge directions, but in general it
does not determine how many separating edges occur in one direction.
For example, let the first factor have coordinate \(e\), and assign
three distinct coordinates to the edges of either four-vertex tree.
The ample partial cubes \(Q_1\mathbin{\square}P_4\) and
\(Q_1\mathbin{\square}K_{1,3}\) have the same shadow: besides the
empty face and singletons, it contains exactly the three pairs
\(\{e,f\}\).  They have respectively two and one separating
\(e\)-edges, as deleting such an edge amounts to deleting a vertex
of the tree.  The contraction, rather than its shadow alone, recovers
their positions and multiplicities.

For the \(1162\)-concept example in
Remark~\ref{ex:damp-coupled-edge-minors-1162}, the crossing graph
has just one cut vertex, coordinate \(10\).  Deleting it gives
the sets \(\{0,\ldots,13\}\setminus\{10\}\) and
\(\{14,\ldots,25\}\).  The corresponding pieces have orders
\(871\) and \(293\).  The corollary therefore forces a unique
separating edge and these piece orders in every ample realization
of that shadow.  In the stored realization the edge is
\(\{0,1024\}\).  By contrast, the fixed \(291\)-concept family
\(G\) in Remark~\ref{rem:damp-noncorner-rooted-factor} has crossing
graph \(K_{12}\), so every ample realization of its shadow has
neither a cut vertex nor a separating edge.

These statements identify and recover the separator branch without
assuming a marked separator.  They do not decide the five peeling
properties of its factors, or generate the remaining forbidden cores
whose crossing graphs are two-connected.  The checker
\path{verify_damp_crossing_separators.py} verifies the fixed examples
and all full-dimensional ample families in \(Q_3,Q_4\), as well as
two nonample partial-cube controls; the general conclusions follow
from the proofs above.
\end{remark}


\begin{remark}[Two relative failures can have a peelable union]
\label{rem:damp-two-edge-failures-positive}
Use the translated \(289\)-concept family \(C\) from
Remark~\ref{ex:damp-terminal-cut-vertex-1158}, with edge
\(D=\{u,v\}=\{0,2\}\), and set
\[
 A=C,\qquad
 B=\{0,8192\}\cup\{2+2^{12}x:x\in C\}.
\]
The old coordinates of the second copy are disjoint from those of \(A\).
The four words \(0,2,8192,8194\) form a square attached along
an edge of this second copy.  Both pieces are ample and two-connected.
Their five properties, in the orientation \((u,v)\), are
\[
 \begin{array}{c|ccccc}
       &r_D&t_u&t_v&\ell_u&\ell_v\\ \hline
 A     &0&0&1&0&0\\
 B     &0&0&0&1&0
 \end{array}
\]
The first family has only corner \(u\), of degree three, and a
stored peeling ending at \(v\).  The second has initial corners
\(0,8192\).  Keeping both endpoints forces deletion of \(8192\)
and then stops.  Keeping either endpoint as the final concept also
fails; the complete negative state graphs have two and four states.
Deleting \(8192\), then the leaf \(u\), then peeling the remaining
copy of \(C\) gives \(\ell_u(B)=1\).

Their \(578\)-concept union peels explicitly: delete \(8192\),
peel \(A\) to \(v\) with its last \(v\) omitted, and then
peel the second copy of \(C\).  Thus the vertex-gluing rule
``at most one piece fails relative peeling'' is false for edges.
The endpoint compatibility in
\eqref{eq:damp-edge-endpoint-compatibility} is necessary.

The checker \path{verify_damp_edge_separator_profiles.py} validates
these profiles and this complete order, a cornerless \(576\)-concept
incompatible double, and a nonminimal \(865\)-concept three-piece
union.  It also recovers the pieces of orders \(289,871\) in
Remark~\ref{ex:damp-terminal-edge-amalgam-1158}.  Both recovered
pieces have the first displayed profile.  All \(48\) exclusive
elementary factor images peel to the shared edge, and one factor
peels to its root in each of the three shared-coordinate images.
These are the conditions of
Theorem~\ref{thm:damp-edge-separator-obstructions} checked on the
recovered pieces.  The source endpoint order is recorded in
\path{computations/round807_edge_endpoint_order.json}, and the full
report is \path{damp_edge_separator_profiles_verification.json}.
\end{remark}
\begin{proposition}[Transferring rooted failure to an edge boundary]
\label{prop:damp-edge-boundary-transfer}
Let $0\in C\subseteq A,B\subseteq Q_F$ be ample, with $A\cap B=C$,
$A\cup B$ ample and
$\operatorname{Sh}(A)\cap\operatorname{Sh}(B)=\operatorname{Sh}(C)$.
Suppose $C\in\Damp$, both $A$ and $B$ peel to $C$, and both admit
acyclic USOs with sink $0$. For a fresh coordinate $e$, set
\[
 L=(A\times\{0\})\cup(B\times\{1\}),\qquad
 D=\{u,v\}=\{(0,0),(0,1)\}.
\]
Then $L$ is ample. If $p=1$ when $C$ admits an acyclic USO with sink
$0$, and $p=0$ otherwise, its five boundary values are exactly
\[
 (r_D(L),t_u(L),t_v(L),\ell_u(L),\ell_v(L))=(p,1,1,p,p).
\]
In particular, when $p=0$, either endpoint is an attainable sink but
neither leaf exit nor retention of the edge is possible.
\end{proposition}
\begin{proof}
The supports avoiding $e$ are those shattered by $A\cup B$, and the
supports containing $e$ are $I\cup\{e\}$ with
$I\in\operatorname{Sh}(A)\cap\operatorname{Sh}(B)=\operatorname{Sh}(C)$.
Their number is $|A\cup B|+|C|=|L|$, proving ampleness.
To end at $u$, delete $B\setminus C$ in the upper layer by its relative
order, then the whole upper copy of $C$ by an ordinary core order, and
finish with an order of $A$ to $0$. The first deletions have no vertical
mates. During the core deletions the lower $A$ contains every needed mate
cube. Thus every deletion is a corner deletion. Interchanging the layers
ends at $v$.

If $p=1$, clear both wings and peel $C\times\{0,1\}$ down to $D$
using an order of $C$ ending at $0$. Either orientation of the remaining
edge gives both leaf exits. Conversely, record a core word whenever its
full two-point fibre first breaks during a peeling. While a full fibre
exists, every neighbour in the current reduction lifts to a neighbour of
the deleted vertex, whose mate is still present. Its corner cube projects
to a corner cube in that reduction. These events therefore form a corner
deletion sequence starting from $C$.

If the peeling retains $D$, this sequence ends at $\{0\}$ without
deleting $0$. If instead an endpoint $s$ is the first one deleted and
is then a leaf, stop immediately before that deletion. Its mate is still
present and it has no old-coordinate neighbour. Hence $0$ is isolated in
the current reduction. That reduction is ample, nonempty and connected,
so it is $\{0\}$. Again the recorded events peel $C$ to $0$. This proves
that each of $r_D,\ell_u,\ell_v$ implies $p=1$ and finishes all cases.
\end{proof}

This reuses the two-layer contraction and endpoint schedules of the paired
root construction (Proposition~\ref{prop:damp-paired-root-transfer} and its
\href{research/relative_wing_square.pdf}{relative-wing extension}). The
additional conclusion is the exact five-value boundary profile. It refutes
the implication that two attainable endpoint sinks force edge retention;
its obstruction is the prescribed sink of the reduction, not absence of
ordinary positive orders on the layers.

\begin{theorem}[Relative obstructions at cube boundaries of arbitrary dimension]
\label{thm:damp-cube-boundary-retention}
Take the ample inputs $C,A,B$ and root $0$ of
Proposition~\ref{prop:damp-edge-boundary-transfer}. Thus $C$ is positive,
$A$ and $B$ peel both to $C$ and to $0$, and their intersection and
shadow intersection are those of $C$. For $k\ge1$, let $E$ be a fresh
$k$-coordinate set and put
\[
 R_A=Q_E\setminus\{1^E\},\quad
 R_B=Q_E\setminus\{0^E\},\quad D_k=\{0\}\times Q_E,
\]
\[
 Y_k=(C\times Q_E)\cup((A\setminus C)\times R_A)
                    \cup((B\setminus C)\times R_B).
\]
Then $Y_k$ is ample and positive, and it peels to every nonempty proper
face of $D_k$. It also peels to each of the two punctured cubes
$\{0\}\times R_A$ and $\{0\}\times R_B$.
For $k=2$, it peels to every nonempty proper ample subfamily of $D_2$,
including all four punctured squares.
It peels to $D_k$ if and only if $C$ peels to $0$.
If $(C,0)$ is a rooted factor using every old coordinate, every proper
pc-minor of the pair
$(Y_k,D_k)$ with nonempty target peels to its target image, although
$Y_k$ does not peel to $D_k$.
Here a pc-minor of the pair applies the same coordinate operations to
both entries and discards operations whose target image is empty.

The certified inputs $|C|=289$, $|A|=291$, $|B|=290$ give, for every
$k\ge1$, such a relative obstruction with
\[
 |Y_k|=292\,2^k-3,\qquad \operatorname{idim}(Y_k)=12+k.
\]
Thus even the simultaneous ability to retain every proper face of a
cube does not imply the ability to retain that cube.
\end{theorem}
\begin{proof}
Write $U=A\cup B$. Inclusion and the shadow-intersection hypothesis give
$\operatorname{Sh}(U)=\operatorname{Sh}(A)\cup\operatorname{Sh}(B)$:
the union on the right has cardinality
$|A|+|B|-|C|=|U|=|\operatorname{Sh}(U)|$.
For an old support $I$ and a new support $J\subsetneq E$, the combined
support $I\cup J$ is shattered in $Y_k$ exactly when $I$ is shattered
in $U$. For sufficiency choose the wing whose shadow contains $I$;
its punctured $E$-cube contains a $J$-face. Necessity follows by projection.
If $J=E$, the layers $0^E$ and $1^E$ force $I$ to be shattered in both
$A$ and $B$, equivalently in $C$. Conversely $C\times Q_E$ contains
all the required cubes. Consequently
\[
 |\operatorname{Sh}(Y_k)|=(2^k-1)|U|+|C|=|Y_k|,
\]
so $Y_k$ is ample.

We describe the corner orders used below. For a relative order from $A$
to $C$, delete the copies of each successive old vertex by a corner
order of $R_A$. Do the same for $B$ and $R_B$. The punctured cubes are
positive, for example after a translation they are downward-closed.
There are no old-coordinate edges from $A\setminus C$ to
$B\setminus C$: such an edge would give a missing-diagonal square in
the ample two-layer family of
Proposition~\ref{prop:damp-edge-boundary-transfer}. At each deletion,
therefore, the old neighbour cube is the one in the relative wing order;
the new neighbour cube is the one in its row order. Their product is
present, with core neighbours having every new-coordinate mate.
This verifies every corner deletion and clears both wings to $C\times Q_E$.
That product is positive, proving ordinary positivity of $Y_k$.
If $C$ peels to $0$, the same procedure followed by its product relative
order peels $Y_k$ to $D_k$.

Conversely, suppose a peeling retains $D_k$. Record an old core vertex
when its complete $E$-fibre first loses a vertex. These first-loss events
are corner deletions in the current $E$-reduction: at a full fibre, every
remaining reduction neighbour has all its mates, and the corner cube of
the deleted vertex projects to the requisite reduction cube. Initially
that reduction is $C$, since neither punctured wing has a full fibre.
At the end it is $\{0\}$, whose fibre was never touched. The recorded
sequence would peel $C$ to $0$. This is the relative form of the
reduction-closure argument, proving the equivalence.

Let $F$ be a nonempty proper face of $Q_E$. It omits at least one of
$0^E,1^E$, so it lies in $R_M$ for one wing $M\in\{A,B\}$.
A punctured cube $Q_E\setminus\{v\}$ peels to any face avoiding $v$:
choose a coordinate fixed on the face opposite to $v$, clear the punctured
opposite semicube using its full mates, then peel the remaining full cube
to the desired face. A full cube also peels to any of its faces, by
successively clearing opposite semicubes.

Clear the other wing completely. For each vertex in a relative order of
$M$ to $C$, delete only its copies in $R_M\setminus F$, using the
punctured-cube relative order just described. The same product corner
check applies. The remainder is
\[
 (C\times Q_E)\cup((M\setminus C)\times F).
\]
For each vertex in an ordinary order of $C$, clear its copies outside
$F$ using a full-cube relative order to $F$. Old neighbours at these
copies come only from the surviving core; all needed mates in $F$
remain. The remainder is $M\times F$, which peels to $\{0\}\times F$
using the given root order of $M$. This proves retention of every proper
face of $D_k$, including all its vertices.

For the punctured target $\{0\}\times R_A$, clear the $B$-wing by
its relative order to $C$. Next delete the single copy $(v,1^E)$ of
each $v$ in an ordinary peeling of $C$. Every old neighbour cube at
that copy has all its required mates in the remaining rows, so these
are corner deletions. What remains is $A\times R_A$; the prescribed
root order of $A$ then retains $\{0\}\times R_A$. Interchanging the
two wings proves the assertion for $R_B$.

For $k=2$, index the rows by $00,01,10,11$. They are respectively
$A,U,U,B$, where $U=A\cup B$. Consider a puncture $q\in\{01,10\}$,
and let $q'$ be the other member of that pair. First clear both repair
wings \emph{only in row $q$}, using their relative orders to $C$.
An $A$-exclusive vertex in that row has its only remaining new-coordinate
neighbour in row $00$, where its entire old corner cube is present.
A $B$-exclusive vertex has the analogous mate in row $11$. There are
no old edges between the two exclusive wings. These are therefore valid
corner deletions. Row $q$ is now $C$. Delete it by an ordinary peeling
of $C$: every current old neighbour cube has all its mates in the other
three rows, so the product corner condition holds.

The remaining row graph is the path
\[
 (A\text{ at }00)\;--\;(U\text{ at }q')\;--\;(B\text{ at }11).
\]
Keep the entire middle row while peeling each outer row to its root.
Each outer corner has its full mate cube in $U$. Then peel the middle
row to its root, using first the relative order clearing $B\setminus C$
and then the root order of $A$. At this last stage, no nonroot middle
vertex has a neighbour in either surviving outer row, so its corner
test is unchanged. The retained family is precisely the boundary square
minus its vertex at $q$. Together with the two designated punctures,
this proves all four cases. Every nonempty proper ample subfamily of
a square is a vertex, an edge, or a three-vertex path; these cases are
now all covered.

Assume now that $(C,0)$ is a rooted factor using every old coordinate.
An old-coordinate operation
with nonempty target is a root-preserving restriction or contraction.
The resulting $C'$ peels to $0$. The corresponding wings $A',B'$ peel
to $C'$: relative corner orders pass to their coordinate images by
omitting repeated families, as in
Corollary~\ref{cor:damp-edge-boundary-double}. The same corollary's
no-diagonal argument gives $A'\cap B'=C'$. The image of $Y_k$ consequently
has the same three-part form; clear its wings, then peel $C'\times Q_E$
to the target cube. A restriction in a new coordinate inherits a peeling
to its target facet from the proper-face result just proved. Contraction
in any new coordinate sends both punctured rows onto the full smaller
cube, giving $U\times Q_{E\setminus\{e\}}$. The family $U$ peels to $0$:
clear $B\setminus C$ first, then use the root order of $A$.
Thus this contraction also peels to its target. These cover every elementary
operation with nonempty target; relative minor closure covers subsequent
operations. Rooted failure of $C$ excludes retention in the original pair.

The finite inputs are those already used for the $581$-vertex edge
transfer, with $A\setminus C=\{128,160\}$ and
$B\setminus C=\{290\}$. Substituting $|U|=292$ and $|C|=289$ gives
the stated sizes. At $k=2$, \path{research/cube_boundary_retention.py/json}
checks exact shadows, an ordinary corner order and its nonclashing out-map,
all four boundary-edge relative orders, and all $30$ elementary proper
relative-minor orders. Failure at the full square follows from the
rooted-core proof, not a negative search.
\end{proof}

For $k=1$ this is the existing edge transfer. The opposite punctured
rows also occur in the ordinary common-seed constructions, for example
Proposition~\ref{prop:damp-terminal-cut-seed-towers}. The new conclusion
is their relative-boundary behaviour with a positive rooted-critical core:
ordinary positivity and every proper-face retention coexist with a
minimal failure to retain the entire cube. This rules out any fixed
proper-face dimension cutoff as a universal test for cube retention.
It does not generate the rooted input factors or classify all relative
obstructions.

\begin{corollary}[An edge compatibility test and a forbidden double]
\label{cor:damp-edge-boundary-double}
Assume $p=0$ in Proposition~\ref{prop:damp-edge-boundary-transfer}.
For a positive ample piece $K$ on coordinates disjoint from $F$ except
for the common edge coordinate, $L\cup_D K$ is positive exactly when
\[
 r_D(K)=1\quad\text{or}\quad \ell_u(K)=1\quad\text{or}\quad\ell_v(K)=1.
\]
If moreover $(C,0)$ is a rooted factor and uses every coordinate of $F$,
then two copies of $L$, on disjoint old coordinate sets and with $D$
identified, form a forbidden pc-minor of order $2(|A|+|B|)-2$.
\end{corollary}
\begin{proof}
The compatibility test is Theorem~\ref{thm:damp-edge-gluing-profile}
with the displayed five values substituted. In particular, two copies
of $L$ have a nonpeelable union.

Consider an elementary old-coordinate operation retaining $D$ in one
copy. It replaces $C$ by a proper root-preserving minor $C'$ that peels
to $0$. Images of relative deletion chains, with repetitions omitted,
are relative deletion chains: every intermediate image is ample and
each step loses at most one word outside the retained image. Thus the
new sections $A',B'$ peel to $C'$. Their intersection is $C'$.
For face intersections this is immediate; under contraction a new common
word outside $C'$ would come from opposite vertices of an old-coordinate
and $e$ square with both intermediate vertices absent, contradicting
isometry. Consequently the changed copy peels to $D$, by clearing its
wings and then its core product. Its union with the unchanged positive
copy peels by the edge-gluing theorem. An old-coordinate face intersection
not containing $D$ discards the other copy and leaves a positive minor
of $L$.

For either face intersection in $e$, or contraction of $e$, both copies
have positive images with the common root as an attainable sink: take
the root-preserving image of a peeling of $L$ ending at the corresponding
endpoint. The vertex-gluing theorem makes their union positive. These
are all elementary pc-minors, so the double is forbidden.
\end{proof}

The reduction of the forbidden double in $e$ is $C\vee C$, already a
forbidden class by the rooted-factor theorem. Thus this double belongs to
the established inverse-construction branch above negative reductions; it
does not generate a previously uncovered reduction-terminal seed.

The corollary is an application of the existing separator classification,
not intrinsic generation of the rooted input $C$. For a stored example,
use the first control in \path{relative_wing_square.json}: $|C|=289$,
$A=C\cup\{128,160\}$ and $B=C\cup\{290\}$. The root $0$ is the
only corner of $C$, so $p=0$. The resulting $L$ has $581$ words and
profile $(0,1,1,0,0)$; its forbidden double has $1160$ words.
The input rooted-factor certificates are the previously established ones.
The checker \path{research/verify_edge_boundary_transfer.py} replays both
endpoint orders and the relative wing orders and verifies the unique-corner
certificate of the core. The profile and double's minimality follow from
the uniform proofs above. This does not settle existence of forbidden
three-piece edge unions: it shows why attainable sinks alone cannot exclude
them, and why the leaf conditions in their classification must be retained.

\begin{proposition}[Rooted factors inside one edge context type]
\label{prop:damp-rooted-prism-context}
Let $C\subseteq Q_F$ be ample and corner-peelable, with marked vertex
$a$, and use a fresh coordinate $e$ to form
\[
 P=C\times Q_{\{e\}},\qquad D=\{a\}\times Q_{\{e\}}.
\]
If $h=1$ when $C$ peels to $a$ and $h=0$ otherwise, then
\[
 (r_D(P),t_{(a,0)}(P),t_{(a,1)}(P),
       \ell_{(a,0)}(P),\ell_{(a,1)}(P))=(h,h,h,h,h).
\]
Consequently $(C,a)$ is a rooted factor if and only if $P$ has the
positive zero-vector context type $Z$ and every proper elementary
private-coordinate minor retaining $D$ peels to $D$.
Here $C$ is embedded on its active coordinates, and a private restriction
retaining $D$ fixes its coordinate to the bit of $a$.

Moreover, every rooted factor occurs in this way as a piece of an ample
forbidden pc-minor with a separating edge. One may glue $P$ to the fixed
$581$-vertex piece $L$ of Corollary~\ref{cor:damp-edge-boundary-double},
identifying its marked edge with $D$ and renaming all private coordinates.
The resulting forbidden family has $2|C|+579$ vertices.
\end{proposition}
\begin{proof}
Products of ample corner-peelable families are ample and corner-peelable.
If $C$ peels to $a$, deleting its other vertex fibres in that order
peels $P$ to $D$. Indeed, at each fibre deletion the old corner cube
is multiplied by the remaining part of its edge fibre. Relative peeling
to $D$ implies all four sink and leaf values, by finishing with the
appropriate order on the edge.

Conversely, a relative peeling to $D$ restricts to a peeling of either
$e$-section ending at $a$. A sink at $(a,i)$ restricts to a sink at
$a$ in that section. Likewise an out-map equal to $\{e\}$ at $(a,i)$
has empty out-map after restriction to its $e$-section, so gives an
acyclic USO of $C$ with sink $a$. Thus each of the five values implies
that $C$ peels to $a$, proving the displayed profile.

A private root-containing restriction or contraction replaces the prism
by $(\mu C)\times Q_{\{e\}}$ and carries $D$ to the edge over $\mu a$.
The profile identity says that the image retains that edge exactly when
$\mu C$ peels to $\mu a$. These are precisely the elementary rooted
minimality conditions, and ordinary positivity of $C$ is assumed.
This proves the equivalence and also recovers $C$ and its root from
either section of the marked prism.

For the last assertion translate $a$ to zero. The fixed $L$ is positive,
fails retention of its edge, attains both endpoints as sinks, and every
proper elementary private minor retaining its edge retains that edge,
by Corollary~\ref{cor:damp-edge-boundary-double} and its proof.
Its two shared-coordinate restrictions and its shared-coordinate
contraction peel to the image of an endpoint, by rooted minor closure.
The union in the statement is exactly the product attachment
$P_C(L,D)$ of Theorem~\ref{thm:damp-relative-attachment-minimality}.
All old-coordinate images of $(L,D)$ with nonempty image of $D$ therefore
peel relatively to that image, whereas $L$ itself fails retention.
That theorem proves forbiddenness. The attachment formula gives
$|L|+(|C|-1)|D|=581+2(|C|-1)$. Both pieces have vertices outside $D$
and no edges between their private parts, so the common edge is a
separator.
\end{proof}

Thus classifying minor-minimal realizers even of the single edge type $Z$
contains the entire rooted-factor input problem, already within marked
prisms that actually occur in forbidden assemblies. The proposition is
an application of the earlier product-attachment and edge-transfer
results, not a new mechanism for generating rooted factors. In particular,
settling which Boolean types are realizable would not by itself settle
their support generation.

\begin{lemma}[Transport of edge boundary conditions]
\label{lem:damp-edge-boundary-minors}
Let $D=\{u,v\}$ be an edge of an ample class $K$. A pc-minor operation
\emph{retains $D$} if it is a coordinate-face intersection containing both
endpoints or a contraction in a coordinate other than the direction of $D$.
Each of $r_D,t_u,t_v,\ell_u,\ell_v$, when true, remains true after
such an operation, with the endpoints replaced by their images.
\end{lemma}
\begin{proof}
Apply the operation to the ample residual families of a witnessing peeling.
A step loses at most one image word; omit repetitions. This gives a peeling
of the image. For a face intersection the resulting order is the filtered
order. For a contraction it records a word when its last lift is deleted.
Retained endpoints persist, proving transport of $r_D$ and the two sink
conditions.

For $\ell_s$, consider the moment when $s$ is deleted as a leaf adjacent
to the other endpoint $t$. A face intersection containing $D$ preserves
this leaf and its neighbour. Suppose instead coordinate $f$ is contracted,
where $f$ differs from the edge direction. The $f$-mate of $s$ is absent,
so this deletion removes the last lift of its image. No new neighbour of
that image can appear outside the direction of $D$: if a quotient neighbour
had only a lift in the opposite $f$-layer, that lift would have distance two
from $s$ with both intermediate vertices absent (its matching-layer lift
and the $f$-mate of $s$). This contradicts isometry of the current ample
family. Thus the image of $s$ is still a leaf. The image of $t$ remains,
so it is still the first endpoint deleted. The image order witnesses
$\ell_s$. Iterating proves the claim for sequences of these operations.
\end{proof}

\begin{theorem}[An exact existence criterion for the three-piece branch]
\label{thm:damp-three-piece-boundary-equivalence}
The following statements are equivalent.
\begin{enumerate}
 \item There exists an ample forbidden pc-minor with an edge whose deletion
       leaves exactly three connected components.
 \item There exists a positive ample class $K$ with an edge $D=\{u,v\}$
       such that
       \[
       r_D(K)=0,\qquad
       (t_u(K)\lor t_v(K))=(\ell_u(K)\lor\ell_v(K))=1.
       \]
\end{enumerate}
In item 2 the sink and leaf exit may be witnessed by different acyclic USOs,
and their endpoints may agree or differ. No minor-minimality hypothesis or
connectivity assumption on $K\setminus D$ is needed.
\end{theorem}
\begin{proof}
For necessity, the separator classification gives three positive pieces,
all failing edge retention and pairwise compatible. Each compatible pair
has a witness consisting of a sink value in one piece and a leaf-exit value
in the other. If no piece possessed both kinds of value, pieces with a sink
value and pieces with a leaf-exit value would form disjoint groups, and
compatibility could occur only across the groups. Three pieces cannot all
be pairwise compatible in this way. Thus one piece satisfies item 2.

For sufficiency, pass to a minor $K_0$ retaining $D$ that still fails to
peel to $D$ and is minimal among such minors. All the positive boundary
values in item 2 survive by Lemma~\ref{lem:damp-edge-boundary-minors}.
Every proper elementary operation in an exclusive coordinate retaining
$D$ now makes $K_0$ peel to $D$. Moreover $K_0\setminus D$ is connected.
Otherwise the canonical pieces at $D$ are proper convex minors. They all
peel to $D$ by minimality, and their outside deletion orders concatenate
to peel $K_0$ to $D$, a contradiction.

First suppose $t_u(K_0)=\ell_v(K_0)=1$, interchanging endpoint names if
necessary. Glue three copies of $K_0$ along $D$ on disjoint exclusive
coordinate sets. Each pair is compatible, while all three fail retention;
thus their union is nonpeelable. An exclusive elementary operation retaining
$D$ changes one piece to a relative-positive piece and leaves a compatible
pair. An exclusive face intersection not containing $D$ leaves a positive
minor of a single piece. For the face image through $u$, all three pieces
peel to the common root, by restricting the sink-$u$ orders. The same is
true for the face through $v$: immediately before the leaf deletion of $v$
in a witnessing order, its own coordinate shore is a singleton. Indeed
isometry forces every other residual vertex into the shore of its only
neighbour $u$. Filtering the preceding order to the $v$-shore peels it to
$v$. Finally, contraction of the shared direction sends the sink-$u$
orders to orders ending at the common root. The vertex-gluing law therefore
makes all three shared-coordinate images positive. This proves forbiddenness.

Otherwise item 2 supplies a same-endpoint pair, say
$t_u(K_0)=\ell_u(K_0)=1$. Use a fixed relative-critical piece $L$ with
profile $(0,1,1,0,0)$ from
Corollary~\ref{cor:damp-edge-boundary-double}, chosen with
$L\setminus D$ connected. Such a piece exists by the certified $581$-word
instance above. Glue $L$, $K_0$, and an endpoint-reversed copy of $K_0$.
Every pair is compatible: $L$ supplies either sink and the other pieces
supply their respective leaf exits; the two copies of $K_0$ supply a sink
and the opposite leaf exit to each other. All three fail retention, so
the union is negative. Each exclusive elementary image retaining $D$
peels to $D$ and leaves a compatible pair; an image discarding $D$ lies
in a single positive piece. In either shared-coordinate face image at
most one piece can fail at the common root: $L$ attains both endpoints,
and one of the two copies attains the chosen endpoint. In the contraction
all three images attain the common root. The vertex-gluing law proves
positivity of these images. Hence this union is forbidden as well.
In both constructions the three complements of $D$ are connected, occupy
disjoint exclusive blocks, and have no cross edges, so deleting $D$ leaves
exactly three components.
\end{proof}

\begin{corollary}[Repairing a three-piece obstruction by component reflections]
\label{cor:damp-three-piece-component-repair}
Let $X=A_1\cup_D A_2\cup_D A_3$ be an ample forbidden pc-minor
with the canonical three-piece presentation of
Theorem~\ref{thm:damp-edge-separator-obstructions}. Write $e$ for the
direction of $D$ and $r=\pi_eD$. There is an index $i$ such that,
on putting $U=\pi_eA_i$ and $C=A_i^e$, both $U$ and $C$ peel to $r$.
For either $b\in\{0,1\}$, replace $A_i$ by
\[
 P_b=(U\times\{b\})\cup(C\times\{1-b\}),
\]
with $e$ written as the last coordinate and all private coordinates of
the other pieces fixed as before. The resulting class $X_b$ is ample,
has the same shattered sets and cardinality as $X$, and belongs to $\Damp$.

This replacement is a sequence of reflections of components of
$(\pi_eX)\setminus X^e$, all lying in the selected piece. If that
piece has $k_0$ components labelled zero and $k_1$ labelled one, the
shorter of the two replacements uses $\min(k_0,k_1)$ reflections.
This is an upper bound for reaching a positive signing, not a claim that
each intermediate reflection preserves forbiddenness or that the bound
is optimal.
\end{corollary}
\begin{proof}
The necessity argument of
Theorem~\ref{thm:damp-three-piece-boundary-equivalence} applies to these
three pieces: at least one piece $A_i$ admits both an endpoint sink and
an endpoint leaf exit, possibly in different acyclic USOs. Contracting
the shared direction in an order ending at that sink gives a peeling
of $U$ ending at $r$, by the order transport in
Proposition~\ref{prop:damp-pc-minor-closure}.

The leaf exit gives a peeling of $C$ ending at $r$. To see this directly,
record a reduction vertex each time its full $e$-fibre first loses an
endpoint during the witnessing peeling. These events are corner deletions
in the current reduction: the deleted endpoint's corner cube contains
the mate of every neighbour whose fibre is still full. Immediately before
the specified leaf exit, the fibre over $r$ is full and $r$ has no
neighbour in the current reduction. That reduction is ample and connected,
so it consists of $r$ alone. This is the first-fibre argument from
Proposition~\ref{prop:damp-edge-boundary-transfer}, which does not require
the special repair-wing hypotheses for this implication.

The peripheral expansion $P_b$ peels to $D$. First use the rooted order
of $C$ to delete its copy in layer $1-b$, except for $r$. Every required
mate cube remains in the untouched copy of $U$. Then delete $U\setminus
\{r\}$ in layer $b$ by its rooted order. The only remaining vertex in
the other layer is the mate of $r$, so it introduces no neighbour at any
of these deleted vertices. Both stages consist of corner deletions.
In the full union the other pieces have no neighbours outside $D$ in
this piece, so the same relative order is valid. It leaves the unchanged
two-piece union, which is a positive proper convex minor of $X$ by
Theorem~\ref{thm:damp-edge-separator-obstructions}. Thus $X_b$ peels.

For the geometric assertion put $A=\pi_eX$ and $B=X^e$. The projected
pieces intersect only at $r\in B$ and have no cross edges. Hence every
component of $A\setminus B$ belongs to a single projected piece, and
the components in the selected one are precisely those of $U\setminus C$.
Making their labels equal to $b$ gives exactly the replacement above.
The component-expansion theorem
\ref{thm:damp-component-expansion-signs} preserves ampleness, the shadow,
and cardinality under each reflection. The numbers of changed labels for
$b=0$ and $b=1$ are $k_1$ and $k_0$, respectively.
\end{proof}

For the $1237$-vertex class of Theorem~\ref{thm:damp-three-piece-witness},
the shared coordinate has six complementary components, three in each
layer. Each of the two $330$-vertex pieces contributes two components of
opposite labels. Reflecting either component of either piece therefore
gives a positive class. Reflecting either of the two components belonging
to the $581$-vertex piece instead gives a nonpeelable class with a protected
set of $1154$ vertices. The certificates in
\path{research/component_reflection_repair.py/json} verify all six cases:
four full orders are checked through outgoing cubes and nonclashing, and
two protected sets have explicit missing-cube witnesses. Two fixed greedy
orders fail on all six cases, so their failures alone would not distinguish
the positive repairs. This is a local repair theorem for the entire
three-piece branch; intrinsic admissibility of reverse sign changes and
coverage outside that branch remain open.

The equivalence reduces existence to joint attainability on one positive
piece. The preceding $(0,1,1,0,0)$ construction realizes only the sink
conditions, and the older $(0,0,0,1,0)$ collar realizes only a leaf exit.
The next construction realizes both on the same piece and therefore settles
existence of the three-piece branch.

\begin{proposition}[Rooted factors from a repair in a proper face]
\label{prop:damp-repair-face-rooted-factor}
Let $(C,0)$ be a rooted factor on its active coordinate set $F$, and
suppose that $0$ is its only corner. Let $c\notin C$ be an ample
one-concept addition with acquired support $P$, such that $C\cup\{c\}$
peels to $0$. Choose a proper coordinate face $Q_J(0)$ containing $c$
and its entire acquired cube. Put
\[
 B=C\cap Q_J(0),\quad T=B\cup\{c\},\quad
 K=(T\times\{0\})\cup(C\times\{1\}),\quad
 D=\{u,v\}=\{(0,0),(0,1)\}.
\]
Then $K$ is ample and ordinarily peelable, with
\[
 r_D(K)=t_v(K)=\ell_v(K)=0,\qquad \ell_u(K)=1.
\]
Every proper root-preserving pc-minor of $(K,u)$ peels to its root.
Every proper old-coordinate elementary image retaining $D$ peels to $D$.
Consequently either $(K,u)$ is a rooted factor, or $K$ has the mixed
profile $(0,1,0,1,0)$.

If $Q_J(0)$ omits a neighbour of $0$ in $C$, then $(K,u)$ is a rooted
factor. This last conclusion requires no search for a peeling of $K$.
\end{proposition}
\begin{proof}
The punctured acquired cube lies in $B$, so $T$ is ample and
$\operatorname{Sh}(T)=\operatorname{Sh}(B)\mathbin{\dot\cup}\{P\}$.
Since $P\notin\operatorname{Sh}(C)$, the section shadows intersect
exactly in $\operatorname{Sh}(B)$. The section union is $C\cup\{c\}$;
shadow counting proves ampleness of $K$.

As a proper root-containing face of $C$, the class $B$ peels to $0$.
Delete $(c,0)$ and then the lower copy of $B\setminus\{0\}$ using
such an order. The untouched upper $C$ supplies every vertical mate
cube. Now $u$ is a leaf at $v$; delete it and peel $C$ ordinarily.
This proves ordinary positivity and $\ell_u=1$. The upper shore cannot
peel to its root, so neither edge retention nor a sink at $v$ is possible.
A first-endpoint leaf exit at $v$ would also peel its own shore to $v$:
just before that exit its shore is a singleton by isometry, and the
preceding order restricts to a peeling of that shore. Thus $\ell_v=0$.

Let $\mu$ be an old-coordinate operation retaining $D$. Both $\mu C$
and $\mu B$ peel to their roots, by rooted minimality and rooted minor
closure respectively. The lower difference $\mu T\setminus\mu B$ has
at most one word. If it survives, it has no upper mate: a contraction
identifying $c$ with an old word would identify it with a neighbour in
its acquired cube, which already belongs to $B$. Deleting this lower
word leaves the ample nested expansion with sections $\mu B\subseteq\mu C$,
so the word is a corner. Peel the lower $\mu B$ to its root and then
the upper $\mu C$ to its root, retaining $D$. This proves all the stated
old-coordinate tests. The root-containing face in the new coordinate is
$T$, which peels to $0$ by deleting $c$ and then peeling $B$. The
contraction is $C\cup\{c\}$, which peels to $0$ by hypothesis.
These are every elementary root-preserving operation. Closure under further
such operations proves the assertion about all proper rooted minors and
hence the dichotomy.

Suppose finally that a root neighbour $h$ of $C$ is omitted by the chosen
face. In a hypothetical peeling ending at $u$, consider its first deletion
from the upper $C$. All upper words are still present, so the old directions
of that corner would make its projection a corner of $C$. Its projection
must therefore be $0$. But $v$ still has neighbours $u$ and $(h,1)$,
whereas $(h,0)$ is absent. Its incident cube is incomplete. This excludes
the first upper deletion and proves that $u$ is unattainable as a sink.
\end{proof}

The first-deletion argument reuses the missing-root-neighbour mechanism of
Theorem~\ref{thm:damp-root-section-square-pivots}. The additional conclusion
is a one-coordinate construction that supplies every rooted minor test;
the earlier square construction used two new coordinates and three repairs.
For the $289$-concept core and repair $c=290$, choose
$J=\{1,5,6,7,8\}$. Then $|B|=18$, $|T|=19$, and $h=4$ is omitted.
The output is a $308$-concept rooted factor of dimension $13$.
The checker \path{research/repair_face_rooted_factor.py} verifies its
shadow, ordinary order, and all $26$ proper rooted elementary orders;
its negative certificate is the uniform argument above. Adding coordinate
$2$ to $J$ gives the face used in the next theorem. Containing every root
neighbour removes this particular obstruction; it is not asserted to be
sufficient for a sink in general.

\begin{remark}[Recovering a repaired-face presentation and its limits]
\label{rem:damp-repair-face-recovery}
The construction has a necessary condition visible without its marked new
coordinate. Normalize the root of a candidate $H$ to zero. For each edge
at the root, with direction $e$, project its root shore and opposite shore
to $T_e$ and $C_e$. A repaired-face presentation requires
\[
 |T_e\setminus C_e|=1,\qquad
 T_e\cap C_e=C_e\cap Q_J(0),
\]
for a proper face containing $T_e$. The singleton is the recovered repair;
$C_e$ is the recovered input. The least possible $J$ is the union of the
supports of the words of $T_e$. Any larger face realizing the same data
has the same intersection with $C_e$. The acquired-cube and rooted-input
hypotheses of Proposition~\ref{prop:damp-repair-face-rooted-factor} must
still be checked; the singleton test alone is not sufficient.

On the $308$- and $330$-concept examples this recovers exactly one coordinate,
namely the added coordinate, and recovers $c=290$ and the stated input $C$.
For the established $292$-concept noncorner-root factor, the root-edge
directions are $1,2,5,6,8$. The respective root-shore exclusive counts are
$75,35,59,97,90$. Thus it has no repaired-face presentation, even before
any peeling condition is tested. This does not make it a new irreducible
input: the established corner deletions $32,96,64$ recover the same
$289$-concept factor, so the earlier ample-addition construction already
covers it. These recovery checks are recorded in
\path{research/repair_face_recovery.json}.

Finally every output of the repaired-face operation has the nonroot corner
$(c,0)$: it has no vertical mate and its old neighbours span its acquired
cube. Therefore its output cannot directly be reused as the unique-corner
input required by this operation. These are limits of this particular
construction, not of a combination with ample additions or other proved
growth operations. No exhaustive recovery assertion follows from it.
\end{remark}

\begin{theorem}[A forbidden three-piece edge union, with finite certificates]
\label{thm:damp-three-piece-witness}
Let $C\subseteq Q_{\{0,\ldots,11\}}$ be the translated $289$-concept
rooted factor used in Proposition~\ref{prop:damp-edge-boundary-transfer}.
Use integer words with coordinate $i$ represented by bit $2^i$, and put
\[
 c=290,\qquad J=\{1,2,5,6,7,8\},\qquad
 T=(C\cup\{c\})\cap Q_J(0),\qquad B=C\cap Q_J(0).
\]
For a fresh coordinate $e$, define
\[
 K=(T\times\{0\})\cup(C\times\{1\}),\qquad
 D=\{u,v\}=\{(0,0),(0,1)\}.
\]
Then $|T|=41$, $|K|=330$, and $K$ is ample with exact profile
\[
 (r_D,t_u,t_v,\ell_u,\ell_v)=(0,1,0,1,0).
\]
Every proper elementary pc-minor operation in an old coordinate retaining
$D$ makes $K$ peel to $D$, and $K\setminus D$ is connected.

Glue $K$, an endpoint-reversed copy of $K$, and the $581$-concept piece
$L$ of Corollary~\ref{cor:damp-edge-boundary-double} along $D$, using
disjoint old coordinate sets. The result is an ample forbidden pc-minor
with $1237$ concepts and isometric dimension $37$. Deleting $D$ leaves
exactly three components. All its nonempty-coordinate reductions are
corner-peelable.
\end{theorem}
\begin{proof}
The input has acquired support $P=\{5,6,7,8\}$, and the chosen proper
face contains its entire acquired cube. Proposition~\ref{prop:damp-repair-face-rooted-factor}
therefore proves ampleness, ordinary positivity, $r_D=t_v=\ell_v=0$,
and $\ell_u=1$, as well as every exclusive relative-minor condition.
It remains to certify the sink and apply the separator construction.

The remaining assertion $t_u=1$ is certified by the explicit $330$-word
order \texttt{SINK\_ORDER} in
\path{research/three_piece_edge_witness.py}. The checker reconstructs $K$
from the stored input, checks that this list contains every word exactly
once and ends at $u$, and verifies every corner deletion. Independently,
it forms the outgoing coordinate sets of the resulting acyclic orientation,
checks that they biject onto $\operatorname{Sh}(K)$, and checks nonclashing
for every pair of distinct words. CCMW's representation-map characterization
therefore also certifies the acyclic USO with sink $u$. The source words,
both positive orders, and the computed checks are recorded in
\path{research/three_piece_edge_witness.json}. This finite replay, rather
than a general assertion about lifting endpoint repairs, supplies $t_u=1$.

The checker also replays all $24$ relative minor orders furnished by
the preceding proposition. If $K\setminus D$ were
disconnected, its proper convex pieces at $D$ would all peel to $D$ by
this minimality, so their orders would concatenate to peel $K$ to $D$.
Thus it is connected; this is also checked directly.

The three profiles in the proposed union are
$(0,1,0,1,0)$, $(0,0,1,0,1)$, and $(0,1,1,0,0)$.
The same-endpoint construction in the proof of
Theorem~\ref{thm:damp-three-piece-boundary-equivalence} applies without
further minimalization: the preceding relative-minor checks supply its
hypotheses for $K$, and Corollary~\ref{cor:damp-edge-boundary-double}
supplies them for $L$. It proves nonpeelability and positivity of every
proper pc-minor. The three connected complements of $D$ are its three
components. Counting the common edge gives
$330+330+581-4=1237$ concepts and $12+12+12+1=37$ coordinates.

Finally, an exclusive-coordinate reduction equals the reduction of the
corresponding positive piece and is positive by reduction closure. Reduction in $e$
is $B\vee B\vee C$ at the common root. Both copies of $B$ peel to their
roots, and $C$ is ordinarily peelable, so the vertex-gluing theorem makes
this reduction positive. Further reductions remain positive. This also
agrees with the previously proved separator descent theorem
(Theorem~\ref{thm:damp-separator-strong-descent}).
\end{proof}

\begin{proposition}[A rooted factor as a peripheral third piece]
\label{prop:damp-three-piece-peripheral-input}
Let $K$ and $L$ be the fixed $330$- and $581$-concept edge pieces of
Theorem~\ref{thm:damp-three-piece-witness}, with common ordered edge
$D=\{u,v\}$. Their union $Y=K\cup_D L$ is a $909$-concept rooted
factor at $u$. Let $(R,0)$ be any rooted factor on its active coordinate
set $F$, disjoint from the private coordinates of $K,L$, and choose a
nonempty ample $S\subseteq R$ that peels to zero. Put
\[
 P_{R,S}=(R\times\{0\})\cup(S\times\{1\}),
 \qquad u=(0,0),\quad v=(0,1).
\]
Then $X_{R,S}=Y\cup_D P_{R,S}$ is an ample forbidden pc-minor of
order $|R|+|S|+907$ and isometric dimension $|F|+25$.
Its three canonical pieces have exactly one piece whose sections in
the direction of $D$ are nested, namely $P_{R,S}$.
All nonempty-coordinate reductions are positive. The output is
two-connected exactly when $|S|>1$.
With the standard $289$-concept rooted factor and an incident edge as
$S$, this gives a two-connected $1198$-concept example in dimension $37$.
\end{proposition}
\begin{proof}
The pieces $K,L$ fail to retain $D$, but have a positive union:
$t_v(L)=\ell_u(K)=1$ supplies the required compatibility. By
Corollary~\ref{cor:damp-edge-context-algebra}, their union has all
four endpoint and leaf values zero. In particular it cannot peel to $u$.
We check rooted minor minimality at $u$.

A private-coordinate contraction or root-side restriction makes the
changed piece peel to $D$. Clear that piece to $D$ and then use the
sink-$u$ order of the unchanged piece; both $K$ and $L$ have such an
order. The root-side restriction in direction $D$ and the contraction
in that direction each give a vertex gluing of two root-positive
images, by transport of these same sink orders. Thus every elementary
root-preserving minor peels to its root, and subsequent rooted minors
do also. This proves the rooted-factor assertion for $Y$.
The argument uses only compatibility, a common attainable endpoint,
and private-coordinate edge-retention minimality of the two pieces.

Now apply Theorem~\ref{thm:damp-rooted-downset-completion} to the
rooted factors $(Y,u)$ and $(R,0)$, with site $S$ and edge $uv$ in $Y$.
Its output is precisely $X_{R,S}$, so that theorem supplies ampleness,
forbiddenness and the two-connectivity criterion. This application does
not require any further output-minor tests. In the present notation its
order is
\[
 |X_{R,S}|=330+581+(|R|+|S|)-4=|R|+|S|+907,
 \qquad \dim(X_{R,S})=12+12+|F|+1.
\]

The complements of $D$ in $K,L$ are connected by their established
contracts. The complement in $P_{R,S}$ is connected as well:
$R\setminus\{0\}$ is connected, and every upper nonzero vertex of
$S$ has its mate there. The three private coordinate blocks are
pairwise disjoint, so these are exactly the three components of
$X_{R,S}\setminus D$. The fixed pieces $K,L$ have nonempty exclusive
shores on both sides, whereas $P_{R,S}$ is peripheral. Its precise
profile is $(0,0,0,0,1)$: clear $S\setminus\{0\}$ upstairs and delete
$v$ as a leaf; either endpoint sink, retention, or a leaf exit first at
$u$ would require $R$ to peel to zero. Finally, the separator descent
theorem rules out a negative nonempty-coordinate reduction of a
three-piece forbidden amalgam, since that alternative has exactly two
canonical pieces.
\end{proof}

This realizes peripheral pieces even in the two-connected three-piece
branch. It imports the established rooted ample completion operation;
the additional input is the rooted factor $(K\cup_D L,u)$ recovered
from the two boundary profiles. For $S=\{0\}$ the output has $1197$
vertices and is simply the vertex gluing of that factor with $R$.
Taking an incident edge as $S$ removes the cut vertex and gives $1198$.
The record \path{research/three_piece_peripheral_input.py/json}
constructs and replays all $111$ elementary minor orders in each case,
and checks the canonical pieces and cut vertices. Neither an optimal
order nor an exhaustive generator of nonnested pieces is asserted.

\begin{corollary}[Trace partitions and minimality in the three-piece example]
\label{cor:damp-obstruction-trace-blocks}
Let $X$ be the $1237$-concept forbidden class in
Theorem~\ref{thm:damp-three-piece-witness}. Put its shared edge in
coordinate $0$, and the three private coordinate blocks in positions
$1,\ldots,12$, $13,\ldots,24$, and $25,\ldots,36$, respectively.
Let $A=\pi_0X$ and $C=X^0$, retaining a zero in the erased bit position.
Then $|A|=870$, $|C|=367$, and both are corner-peelable.
Order the components of $A\setminus C$ by their least integer word;
their sizes are $249,1,249,1,2,1$.
For $t\in\{0,1\}^6$, assign component $i$ to layer $t_i$, retaining
both layers over $C$, and denote the resulting ample lift by $X_t$.
Then
\[
 X_{001100}\in\Damp,\qquad
 X_{100000},X_{001000}\notin\Damp.
\]
Consequently the positive label set cannot be a union of equality spaces
of partitions having at most one nonsingleton block, even though $(A,C)$
is recovered from an ample forbidden class.
Nevertheless, forbiddenness within this fixed pair has the exact selector
\begin{equation}
 X_t\in\Obs(\Damp)
 \quad\Longleftrightarrow\quad
 t_0\ne t_1,\quad t_2\ne t_3,\quad
 t_4\ne t_5,\quad t_1\ne t_3.
 \label{eq:damp-six-component-forbidden-selector}
\end{equation}
Its four outputs are all isomorphic to $X$.
\end{corollary}
\begin{proof}
Positivity of $A$ follows from minor closure; positivity of $C$ is part
of Theorem~\ref{thm:damp-three-piece-witness}. The stated component
counts follow directly from its three pieces and are also checked by
\path{research/obstruction_trace_partition_counterexample.py}.

For $001100$, each of the first two pieces has constant component labels.
It is a peripheral expansion of the $290$-concept repaired class along
its $40$-concept face. Both classes peel to their common root, so first
peeling the smaller layer to that root and then the larger layer to the
root retains the shared edge. Remove these two pieces outside the edge.
The remaining piece has constant labels on its two components, hence is
a peripheral expansion of its positive $292$-concept contraction along
the positive $289$-concept reduction. It is positive. The cited replay
checks the resulting complete order on all $1237$ vertices.

For each of $100000$ and $001000$, the accompanying JSON record supplies
a set $S$ of $155$ vertices and, for each $x\in S$, two distinct
directions $i,j$ such that
\[
 x^i,x^j\in S,\qquad x^{\{i,j\}}\notin X_t.
\]
The replay checks these inclusions and exclusions against the reconstructed
family. In any putative peeling, the first deleted vertex of $S$ would
retain both neighbours but lack the required square vertex. It could not
be a corner. This proves both negative assertions without assumptions
about which peeling choices are made.

We prove the selector without classifying all positive labelings.
If $t_0=t_1$ or $t_2=t_3$, the corresponding $330$-concept piece
peels to the common edge by the same argument as above. The full union
is positive if and only if the union of the other two pieces is positive:
one direction is restriction closure, and the other prefixes the relative
peeling. That smaller union is a proper coordinate face. Thus the full
union is either positive or has a proper negative minor, and cannot be
forbidden.

Suppose now both pairs are mixed. The shore at the singleton-component
label $t_1$ of the first piece is its repaired face $T$; the other shore
is the $289$-concept class with unique corner zero. The analogous labels
for the second piece are $t_3$ and $1-t_3$. If $t_1=t_3$, restriction
to the other shore therefore contains two copies of the $289$-concept
class joined at zero. Neither copy peels to zero, since zero is its
unique initial corner. Proposition~\ref{prop:damp-vertex-gluing} makes
this proper section negative. Consequently $t_1\ne t_3$ is necessary.

If in addition $t_4=t_5$, the third piece is a peripheral expansion
whose thin shore, at $1-t_4$, is also the $289$-concept class. Since
$t_1\ne t_3$, exactly one of the first two pieces has that same
rooted-failing shore. This proper section again has two rooted failures
at its common vertex and is negative. Thus $t_4\ne t_5$ is necessary.

Conversely, the four labelings satisfying all four inequalities are
obtained from the original labeling $100101$ by complementing coordinate
zero and exchanging the first two private coordinate blocks. These are
cube isometries of the recovered pair, so all four outputs are forbidden
by Theorem~\ref{thm:damp-three-piece-witness}. The replay
\path{research/six_component_forbidden_selector.py} verifies the four
isometries, the relative orders, and the rooted-failing sections. It
classifies $48$ labels by a relative-positive piece, $12$ by a negative
proper section, and the remaining four by these isometries. This proves
\eqref{eq:damp-six-component-forbidden-selector}.

An equality space with one nonsingleton block $B$ accepts $001100$ only
if $B$ lies in $\{0,1,4,5\}$ or in $\{2,3\}$. In the first case it
also accepts the negative label $001000$; in the second it also accepts
$100000$. A discrete partition accepts every label. Therefore no such
space can both contain $001100$ and be contained in the positive label
set, which proves the claimed failure of a union representation.
\end{proof}

\begin{corollary}[Equality conditions after imposing proper-minor positivity]
\label{cor:damp-minor-domain-not-affine}
For the fixed pair $(A,C)$ in
Corollary~\ref{cor:damp-obstruction-trace-blocks}, let $M$ be the set of
labels $t$ for which every elementary proper pc-minor of $X_t$ belongs
to $\Damp$. There is no affine subspace $L\subseteq\mathbb F_2^6$ such
that, for every $t\in M$,
\[
 X_t\in\Damp\quad\Longleftrightarrow\quad t\in L.
\]
In particular, no single partition-equality condition suffices on this
domain. In contrast, on $M$ positivity has the exact description
\[
 X_t\in\Damp\quad\Longleftrightarrow\quad
 (t_0=t_1)\ \lor\ (t_2=t_3)\ \lor\ (t_4=t_5)\ \lor\ (t_1=t_3).
\]
This is a consequence for the existing fixed pair, not a universal
selector for other input pairs.
\end{corollary}
\begin{proof}
Three positive labelings of this pair are
\[
 p=010000,\qquad q=111010,\qquad s=110000.
\]
Complete corner orders for $p,q$ occur in the existing six-component
trace record. The order for $s$ is obtained from the recorded positive
labeling $001100$ by exchanging coordinate blocks $1,\ldots,12$ and
$13,\ldots,24$. This cube isometry preserves both $A$ and $C$ and
permutes component indices by $(0\ 2)(1\ 3)$.
The replay \path{research/minor_admissible_partition_counterexample.py}
checks this isometry and all three complete orders. Constant labelings
$0^6,1^6$ are positive peripheral expansions of $A$ along $C$.
All five labels lie in $M$ by minor closure.

If $L$ existed, $0^6\in L$ would make it a linear subspace. It would
contain $p,q,s,1^6$, hence their sum
\[
 p+q+s+1^6=100101.
\]
This is the original forbidden labeling. It lies in $M$ but its class
is not positive, a contradiction. For the second assertion, every
negative labeling in $M$ is forbidden by definition and minor closure.
Negating the exact forbiddenness selector
\eqref{eq:damp-six-component-forbidden-selector} therefore gives the
stated disjunction on $M$.
\end{proof}

The two negative labels used to refute a one-block union description
on the entire label space, $100000$ and $001000$, do not belong to $M$:
each of its recorded $155$-vertex protected sets lies entirely in the
proper section with marked coordinate $0$ equal to one. The same
missing-square certificates therefore prove that section negative.
Thus the full-domain failure and the disjunctive description on $M$
are consistent. Computing $M$ intrinsically remains a separate obligation;
the restricted formula does not remove the proper-minor tests.

The coordinate pair is inherited from a forbidden class, but its other
labelings need not be forbidden. Thus minimality of one labeling does
not justify simplifying the full positive-label set of the uncoloured
pair. The counterexample concerns a proposed restriction on trace
partitions, not the completeness of the partition-certificate theorem
or the validity of the existing three-piece construction.

The construction shows that the three-piece alternative in the separator
classification is necessary, including among forbidden classes with all
proper reductions positive. It is a new realization of that alternative,
using the earlier rooted factor and repair as inputs; it does not supply
an intrinsic generator of all rooted factors or all reduced forbidden bases.

\begin{theorem}[Rooted ample completions of two blocks]
\label{thm:damp-rooted-downset-completion}
Let \(F,G\) be disjoint finite coordinate sets, and let
\(Y\subseteq Q_F\) and \(C\subseteq Q_G\) be irredundantly
embedded ample families
containing \(0\), with \(|C|>1\).  Suppose \(e\in F\) and
\(v=\mathbf1_{\{e\}}\in Y\).  Let \(K\subseteq C\) be a
nonempty ample subfamily containing \(0\), supplied with a corner-peeling
order ending at \(0\).
Embed the families by setting all other coordinates to zero, and put
\begin{equation}
 B=C\cup(v+K),\qquad
 X=Y\cup C\cup(v+K),\qquad D=\{0,v\}.
 \label{eq:damp-rooted-downset-completion}
\end{equation}
Here \(v+x\) sets coordinate \(e\) to one and retains the
\(G\)-coordinates of \(x\).  Then \(X\) is ample and
\begin{equation}
 X\in\Obs(\Damp)
 \quad\Longleftrightarrow\quad
 (Y,0),(C,0)\text{ are rooted factors}.
 \label{eq:damp-downset-completion-transfer}
\end{equation}
Here the rooted-factor conditions are those of
Theorem~\ref{thm:damp-cut-vertex-obstructions}.
For fixed \(F,G,e\), the parameters are recovered uniquely:
\(Y=X\cap Q_F\), \(C=X\cap Q_G\), and
\(K=\{x\in Q_G:v+x\in X\}\).

When these equivalent obstruction conditions hold, \(X\) is
two-connected exactly when \(|K|>1\).  If \(Y,C\) have minimum
degree at least three and \(K\) is a cube of dimension at least
two containing \(0\), then \(X\) also has minimum degree at
least three and is its own graph three-core.  Its order is
\(\lvert Y\rvert+\lvert C\rvert+\lvert K\rvert-2\).
\end{theorem}

\begin{proof}
By hypothesis \(K\) is ample and peels to \(0\). The family \(B\) is the
peripheral expansion of \(C\) with site \(K\), so it is ample
by Proposition~\ref{prop:damp-exact-expansion-product-tests}.
The edge-amalgam shadow count in
Theorem~\ref{thm:damp-cube-separator-bound} makes \(X=Y\cup B\)
ample as well.  The recovery formulas and the order formula follow
from \(Y\cap B=D\) and \(Y\cap C=\{0\}\).

We first prove the ordinary peeling identity
\begin{equation}
 X\in\Damp
 \quad\Longleftrightarrow\quad
 Y,C\in\Damp\ \text{and at least one of }Y,C\text{ peels to }0.
 \label{eq:damp-downset-completion-peeling}
\end{equation}
The families \(Y,C\) are convex restrictions of \(X\), so
their ordinary peelability is necessary.  Delete the words
\(v+x\), \(x\in K\setminus\{0\}\), in the supplied corner-peeling
order of \(K\). Each is a corner: its current upper-layer corner cube
and the matching lower cube in the untouched \(C\) are present.  These deletions leave
\(Y\cup C\).  Thus if one factor peels to \(0\), the
vertex-gluing law gives a complete peeling of \(X\).

To prove the remaining implication, suppose both factors peel
ordinarily but neither peels to \(0\).  In the notation of
Definition~\ref{def:damp-edge-peeling-profile}, the ordered-edge
profile of \(B\) is exactly
\begin{equation}
 (r_D(B),t_0(B),t_v(B),\ell_0(B),\ell_v(B))=(0,0,0,0,1).
 \label{eq:damp-downset-completion-profile}
\end{equation}
Indeed, an order ending at either endpoint would contract in
coordinate \(e\) to an order of \(C\) ending at \(0\).
Peeling to \(D\) would also give such an endpoint order.
The upper-layer deletion order just given leaves \(v\) a leaf;
delete it and then peel \(C\), proving \(\ell_v(B)=1\).
If \(0\) could be the first deleted endpoint and then a leaf,
its only neighbor would be the retained \(v\).  Isometry forces
every other residual word onto the \(e=1\) side: a geodesic
from \(0\) must start with its sole incident edge.  Restricting
the preceding deletion sequence to \(e=0\) would therefore
peel \(C\) to \(0\), a contradiction.  This proves
\(\ell_0(B)=0\).  Since \(t_0(Y)=0\), also \(r_D(Y)=0\),
and \eqref{eq:damp-downset-completion-profile} gives
\(\chi_D(Y,B)=t_0(Y)=0\).  The exact edge-gluing law
proves that \(X\) cannot peel.  This establishes
\eqref{eq:damp-downset-completion-peeling}.

Suppose now that \(X\) is forbidden.  The restrictions \(Y,C\)
are proper, and the identity just proved shows that both are
ordinarily peelable and fail to peel to \(0\).  For a coordinate
\(f\in F\setminus\{e\}\), a root-preserving restriction or
contraction replaces \(Y\) by its image and leaves \(C,K\)
unchanged.  That proper image of \(X\) peels, so
\eqref{eq:damp-downset-completion-peeling} forces the image of
\(Y\) to peel to \(0\).  For \(f=e\), either root-preserving
operation yields the vertex union of its image of \(Y\) with
\(C\); the vertex-gluing law forces the same conclusion.
For a coordinate of \(G\), either root-preserving operation
replaces \(C,K\) by their images.  The image of \(K\) is
again ample and peels to the image of \(0\), by rooted minor transport,
so the peeling identity forces the
image of \(C\) to peel to \(0\).  If that image is a singleton
the conclusion holds directly.  Root-preserving minor closure now
gives all three rooted-factor conditions for both inputs.

Conversely, assume those rooted-factor conditions.  The peeling
identity makes \(X\) nonpeelable.  A root-preserving elementary
operation in \(F\setminus\{e\}\) makes the image of \(Y\)
peel to \(0\), so the same identity makes the full image peel.
A root-preserving operation in \(G\) makes the image of \(C\)
peel to \(0\), with the same consequence.  The root restriction
and contraction in \(e\) give vertex unions of \(C\) with
root-positive images of \(Y\), and hence peel.  A restriction
opposite the root in \(F\setminus\{e\}\) discards \(B\)
and leaves a minor of \(Y\).  One in \(G\) discards \(Y\)
and leaves a peripheral expansion of a minor of \(C\) along the
corresponding minor of \(K\), or just its nonempty base if the site is
empty. The site is positive by pc-minor closure. Both
peel by Proposition~\ref{prop:damp-exact-expansion-product-tests}.
Finally, the opposite restriction in \(e\) gives the vertex
union of \(Y^1_e\) and \(K\).  The first family peels
ordinarily and the second peels to their common \(0\), so their
union peels.  Thus every proper elementary minor of \(X\) peels.

Both rooted factors have no cut vertex by
Theorem~\ref{thm:damp-cut-vertex-obstructions}.  If \(K=\{0\}\),
then \(X=Y\cup C\) has the cut vertex \(0\).  If \(|K|>1\),
the upper copy of \(K\) is connected and has at least two
matching edges to the two-connected \(C\).  Deleting a lower
vertex leaves that upper copy attached through another matching
edge.  Deleting an upper vertex leaves every surviving upper
vertex attached to \(C\).  Hence \(B\) is two-connected,
and gluing \(Y\) and \(B\) along \(D\) preserves
two-connectivity.  Old vertices lose no neighbors.  If \(K\)
is an \(r\)-cube, every new vertex \(v+x\), \(x\ne0\),
has its \(r\) upper cube neighbors and its matching lower
neighbor.  For \(r\ge2\) this proves the asserted degree bound.
\end{proof}

The former down-set condition is a sufficient special case: decreasing
support size peels a down-set to $0$. The general parameter can instead
be constructed starting from $\{0\}$ by ample one-concept additions
inside $C$; reversing their order supplies the required rooted peeling.
Each addition has the punctured-cube test of
Lemma~\ref{lem:damp-one-concept-extension-test}.
Conversely every admissible $K$ has such a construction. With rooted-factor
inputs fixed, every intermediate site gives a forbidden completion; no
further negative or proper-minor test is required. Thus this
extension enlarges the allowed sites without introducing a negative
membership test for them. The two rooted-factor inputs remain unresolved
structural parameters of the global campaign.

A non-down-set instance is recorded in
\path{research/rooted_ample_completion.py/json}. Use two copies of the
certified $289$-vertex rooted factor and take $K$ to be its coordinate-zero
face in coordinate $0$, with the rooted order already supplied by the
factor's proper-minor certificate. This face is not downward closed. The
site has $99$ vertices and the output has $675$. The verifier checks the
site order and the input rooted-minor orders, marked
recovery, and a protected set in the output. Forbiddenness of the output
uses the uniform theorem, not a new exhaustive minor search.

\begin{remark}[Factor-minor compatibility is essential in a three-core]
\label{ex:damp-coupled-edge-minors-1162}
Let \(C\subseteq Q_{12}\) and \(Y_0\subseteq Q_{14}\) be,
respectively, the \(289\)-concept core and the \(870\)-concept
rooted factor in Remark~\ref{ex:damp-terminal-cut-vertex-1158}.
Use binary integer notation with coordinates numbered from zero,
and define
\[
 v=1024,\qquad Y=Y_0\cup\{v\},\qquad
 C'=2^{14}C,\qquad K'=2^{14}\{0,2,4,6\},
 \qquad B=C'\cup(v+K').
\]
The added word \(v\) has acquired support \(1027\) in \(Y_0\).
The resulting \((Y,0)\) is a rooted factor: delete \(v\)
and ordinarily peel \(Y_0\); its complete root-retaining
deletion graph has eight states and never reaches \(\{0\}\).
Its \(28\) proper elementary root-preserving images peel to
\(0\).  Their orders are obtained from the stored orders for
\(Y_0\) by first deleting the image of \(v\) when that image
is new.  All these orders are checked directly.

Theorem~\ref{thm:damp-rooted-downset-completion} now gives an ample
forbidden pc-minor
\[
 X=Y\cup B,\qquad |X|=1162,\qquad
 \dim X=26,\qquad \operatorname{VC}(X)=4.
\]
It is two-connected, has minimum degree three, and is its own
graph three-core.  Deleting the edge endpoints \(D=\{0,1024\}\)
leaves two components, whose canonical pieces are precisely
\(Y\) and \(B\), of orders \(871\) and \(293\).

Consider the restriction \(\mu=\rho_0^0\), exclusive to \(Y\).
Its factor image is
\[
 H=\mu Y=\{0,1024\}\cup
       \{4x+2:x\in C\cup\{128\}\},\qquad |H|=292.
\]
The complete corner-deletion graph retaining \(D\) has five
states and never reaches \(D\), so \(r_D(H)=0\).  Nevertheless
\(H\) peels to \(0\): delete \(1024\), follow a stored
root-ending order of \(C\cup\{128\}\) in the upper layer,
and finish at \(0\).  On the unchanged side,
\eqref{eq:damp-downset-completion-profile} gives
\(r_D(B)=0\) and \(\ell_v(B)=1\).  Thus
\[
 r_D(H)=r_D(B)=0,
 \qquad \chi_D(H,B)=1,
 \qquad \mu X=H\cup B\in\Damp.
\]
In particular, the compatibility alternative in
\eqref{eq:damp-two-piece-exclusive-minor} cannot be replaced by
relative positivity of the changed factor, even for two-connected
forbidden graph three-cores.

The checker \path{verify_damp_coupled_edge_minors.py} replays all
\(52\) rooted input minor orders and all \(78\) output
elementary-minor orders.  It checks ampleness through exact shadows,
the canonical pieces, the graph degrees and connectivity, and the
complete \(32\)-state negative deletion graph of \(X\).
The explicit positive order of \(H\cup B\) has \(583\)
concepts.  This example is not terminal: deleting the corner
\(99328=v+2^{14}\cdot6\) replaces the square site by the
down-set \(\{0,2,4\}\), so the theorem still gives forbiddenness.
The down-set parameter is
intrinsic, but the two rooted-factor inputs remain unclassified;
this theorem does not supply a global generative classification.
\end{remark}
\begin{remark}[A forbidden pc-minor with two blocks]
\label{ex:damp-cut-vertex-587}
Let \(P\) be the \(294\)-concept peelable family of
Remark~\ref{rem:damp-one-corner-peelable}, translated by its only
corner \(1894\) so that this corner becomes \(0\).
The fixed checker \path{verify_damp_cut_vertex_obstruction.py}
verifies peelings ending at \(0\) for all \(24\) elementary
root-preserving pc-minors of \(P\).  The pair \((P,0)\)
therefore satisfies the theorem: it peels ordinarily, while its unique
corner cannot be left last.  Consequently
\[
 X=(P\times\{0\}^{12})\cup(\{0\}^{12}\times P)
 \in\Obs(\Damp).
\]
It has \(587\) concepts, isometric dimension \(24\), VC dimension
three, no corners, and the unique cut vertex \(0\).  Deleting that
vertex leaves two components of order \(293\).
The checker independently constructs and validates complete peelings
of all \(72\) elementary pc-minors of \(X\); the report
\path{damp_cut_vertex_587_verification.json} records every order.
Both copies of \(P\) are gated in \(X\), with gate \(0\)
for vertices in the other copy: for each such vertex \(x\) and
every \(y\in P\), the equality \(d(x,y)=d(x,0)+d(0,y)\) holds.  Thus \(\Damp\) is not closed
under vertex gluing, even when both pieces are gated and peelable.
\end{remark}


\begin{remark}[A minimal rooted failure whose root is not a corner]
\label{rem:damp-noncorner-rooted-factor}
The root of a factor in Theorem~\ref{thm:damp-cut-vertex-obstructions}
need not be a corner.  Let \(G\) be the fixed cornerless
\(291\)-concept family recorded as \texttt{base} in
\path{damp_fixed_291_repair_creation.json}, and put
\[
 H=G\cup\{91\},\qquad a=29.
\]
Here integers denote their twelve-bit cube words in the stored embedding.
The family \(H\) is ample and peelable, with VC dimension four and
initial corners \(61,91,93,125\); in particular, \(a\) is not a
corner.  Deleting \(91\) first leaves the cornerless \(G\).
The complete corner calculation with \(29\) retained has eight states.
It shows that every complete peeling begins with one of
\[
 (61,125,93),\quad(93,125,61),\quad
 (125,61,93),\quad(125,93,61),
\]
and then deletes \(29\) fourth.  Thus \(H\) cannot peel to
\(29\), although all \(24\) elementary root-preserving pc-minors
do peel to the image of \(29\).

The same rooted minimality holds after deleting any one of
\(61,93,125\), after deleting \(61,125\), and after deleting
all three.  Each remaining family also has a verified ordinary peeling.
In particular,
\[
 B=H\setminus\{61,125,93\}
\]
is a \(289\)-concept ample peelable family of VC dimension three,
with \(29\) as its only corner, and \((B,29)\) satisfies the
three conditions of Theorem~\ref{thm:damp-cut-vertex-obstructions}.
After translating \(29\) to \(0\), the vertex doubles
\(H\vee H\) and \(B\vee B\) are therefore ample forbidden
pc-minors of orders \(583\) and \(577\), respectively, both in
dimension \(24\).  The first has VC dimension four and eight corners;
the second has VC dimension three and no corners.  Both have one cut
vertex.  Deleting \(61,125,93\) in the first block, then in the
second block, gives six obstruction-preserving corner deletions from
the first double to the second.  These labels refer to the original
words within the corresponding block before translation.

The independent checker
\path{verify_damp_noncorner_rooted_obstruction.py} validates all
\(144\) rooted elementary peelings of the six displayed factors.
For each endpoint double it also validates all \(72\) elementary
pc-minor peelings and a complete negative deletion certificate, with
\(64\) states for the \(583\)-concept family and one state for
the cornerless \(577\)-concept family.  The families, orders, and
corner-descent chain are recorded in
\path{damp_noncorner_rooted_obstruction_verification.json}.
\end{remark}

\begin{corollary}[Corner descent in the cut-vertex branch]
\label{cor:damp-rooted-block-corner-descent}
Let \(X=A_1\vee A_2\) satisfy
Theorem~\ref{thm:damp-cut-vertex-obstructions}, with common root
\(a\).  For \(x\in\operatorname{Cor}(A_i)\setminus\{a\}\),
the family \(X\setminus\{x\}\) is a forbidden pc-minor if and only
if \((A_i\setminus\{x\},a)\) still satisfies the three rooted
factor conditions in that theorem.

Call an ample forbidden pc-minor \emph{corner-descent terminal} if
none of its corner deletions is another forbidden pc-minor.  A cornered
terminal ample forbidden pc-minor with a cut vertex exists if and only
if there is a rooted factor satisfying those three conditions that has
a nonroot corner, but no nonroot corner deletion preserving the conditions.
\end{corollary}

\begin{proof}
The corners of \(X\) are exactly the nonroot corners of its two
factors: the root has edges in both coordinate blocks and no mixed
square.  Deleting a nonroot corner leaves an ample family, still with
the root as a cut vertex.  If the result is forbidden, the cut-vertex
theorem shows that the changed branch is its other rooted block and
satisfies the three conditions.  Conversely, if those conditions hold,
the theorem applies directly to the two remaining factors.  Notice that
failure to peel to the root is automatic in the changed factor:
otherwise prefixing its rooted peeling by \(x\) would peel the
original factor to \(a\).

If a cornered terminal \(X\) exists, at least one of its rooted
factors has a nonroot corner, and the first assertion excludes every
condition-preserving deletion in that factor.  In the other direction,
glue a factor with the stated properties to \((B,29)\) from
Remark~\ref{rem:damp-noncorner-rooted-factor}.  This second factor's
only corner is its root.  All corners of the union therefore come from
the first factor, and the first assertion makes every one of their
deletions fail forbidden-minor status.  The union is the required
cornered terminal obstruction.
\end{proof}

\subsection{Canonical three-cores and square attachments}
\label{sec:damp-canonical-attachments}

The cut-vertex decomposition
(Theorem~\ref{thm:damp-cut-vertex-obstructions}) rules out degree-two
cut vertices in an ample forbidden minor.  Together with periphery-freeness and
the small-corner descent proved in
Section~\ref{sec:damp-local-extension-foundations}, this yields
the unique graph three-core.  We then prove the inverse construction:
all square attachments are specified by coordinate pairs and targets
in an acyclic dependency digraph.  This entire construction and
recovery argument uses the elementary results already established
above; the unresolved input is the choice of the forbidden core.

\begin{proposition}[Every dismantlable-ample obstruction is periphery-free]
\label{prop:damp-periphery-free}
If $O\in\Obs(\Damp)$, then $O$ has no peripheral coordinate.  Consequently,
every member of
\[
 \Obs(\Damp)\setminus\{Q_n^{--}:n\ge3\}
\]
is ample, periphery-free, and Cartesian-prime.
\end{proposition}

\begin{proof}
Suppose that $e$ is peripheral.  After orienting $e$, write
\[
 O=(B\times\{0\})\cup(S\times\{1\}),\qquad S\subseteq B.
\]
The contraction base $B=O/e$ and duplicated site $S$ are proper pc-minors
of $O$, so both belong to $\Damp$.  The exact peripheral-expansion theorem
for $\Damp$ says that the displayed expansion belongs to $\Damp$ if and only
if both $B$ and $S$ do.  This contradicts $O\notin\Damp$.  The remaining
claims follow from Proposition~\ref{prop:nonample-layer} and
Lemma~\ref{lem:prime}.
\end{proof}

\begin{theorem}[The graph three-core of an ample forbidden minor]
\label{thm:damp-three-core}
Let \(O\) be an ample forbidden pc-minor for \(\Damp\).
Every degree-two vertex of its graph is a corner.  Repeatedly delete
vertices of degree less than three, in any order.  Every deletion is
a degree-two corner deletion, and the final family \(K\) is a nonempty
ample forbidden pc-minor.  It is the graph's \emph{three-core}: the
unique largest induced subgraph of minimum degree at least three.
In particular, \(K\) is independent of every deletion choice.

The families \(K\) and \(O\) have the same active coordinates,
the same shattered supports of size at least three, and the same VC
dimension.  Conversely, every ample forbidden minor is obtained from
its three-core by a sequence of ample one-concept square completions
on those coordinates.  Starting from any ample forbidden minor of
minimum degree at least three, every such sequence produces an ample
forbidden minor with that same three-core.
\end{theorem}

\begin{proof}
We first record a local fact for any ample family \(A\).  A vertex
of degree two is either a corner or a cut vertex.  Translate it to
\(0\), and let its neighbors be \(0^{\{p\}},0^{\{q\}}\).
Every geodesic from \(0\) starts in one of these directions and
never changes that coordinate again.  Thus \(0\) is the only
concept with trace \(00\) on \(\{p,q\}\).
If any concept realizes \(11\), this pair is shattered.
Ampleness makes it strongly shattered, so a full \(\{p,q\}\)-square
exists.  Its \(00\)-vertex must be \(0\), making \(0\) a corner.
If no concept realizes \(11\), the remaining concepts split into
the nonempty \(10\)- and \(01\)-sections.  Each section is
connected by isometry, and no edge joins the sections.  Hence \(0\)
is a cut vertex.

In \(O\), a cut vertex has degree at least four: by
Theorem~\ref{thm:damp-cut-vertex-obstructions}, it belongs to two
two-connected rooted factors, each contributing at least two edges.
Consequently every degree-two vertex of \(O\) is a corner.
There is no vertex of degree zero, since \(O\) is connected and
is not a singleton.  There is no leaf either: isometry would make
its edge coordinate peripheral, contrary to
Proposition~\ref{prop:damp-periphery-free}.

Corollary~\ref{cor:damp-low-degree-corner-descent} preserves ample
forbidden-minor status after every degree-two deletion.  The preceding
paragraph therefore applies at each step.  The process cannot reach
the empty family, since every intermediate family is forbidden.
It terminates at a nonempty \(K\) of minimum degree at least three.
If \(L\) is any induced subgraph of \(O\) of minimum degree at
least three, none of its vertices can be deleted: as long as \(L\)
remains, each has at least three surviving neighbors.  Thus \(L\)
is contained in every possible terminal \(K\).  This proves the
largest-subgraph property and uniqueness.

Deleting a degree-two corner removes precisely its two-coordinate
support from the shadow.  It removes no singleton support and no
support of size at least three.  Active coordinates are therefore
unchanged.  An ample forbidden family has VC dimension at least
three, since every ample family of VC dimension at most two is
corner-peelable \cite[Proposition~4.11]{ChalopiChepoiMoran22}.
Thus VC dimension is unchanged as well.  Reversing the deletions
gives the asserted square completions.

For the last assertion, Proposition~\ref{prop:damp-rank-two-extension-invariance}
preserves nonpeelability at each square addition, and
Proposition~\ref{prop:damp-one-concept-atlas} preserves forbidden-minor
status.  In reverse addition order, each new vertex has degree two
when it is removed.  The original family of minimum degree at least
three remains.  The uniqueness just proved identifies it as the
three-core of the final family.
\end{proof}

This is a canonical reduction of the full classification to ample
forbidden families of minimum degree at least three.  The inverse
steps have intrinsic tests: choose a missing concept and two
coordinates whose other three square vertices are present, and
require that their support is not shattered.  The one-concept
extension criterion in Lemma~\ref{lem:damp-one-concept-extension-test}
then gives ampleness.  No peeling decision is
needed for these steps.  Recognizing the possible three-cores
intrinsically remains open; a three-core need not be cornerless.

\begin{lemma}[A support pattern forcing corners]
\label{lem:damp-shadow-forced-corners}
Let \(H\subseteq Q_F\) be ample with every coordinate active, and let
\(\Gamma\) be the graph on \(F\) in which \(f,g\) are adjacent
exactly when \(\{f,g\}\in\operatorname{Sh}(H)\).
Suppose that
\[
 F=\{e\}\mathbin{\dot\cup}I\mathbin{\dot\cup}A
             \mathbin{\dot\cup}B,
 \qquad I,A,B\ne\varnothing,
 \qquad P=\{e\}\cup I\in\operatorname{Sh}(H),
\]
that the induced graphs \(\Gamma[A]\) and \(\Gamma[B]\) are
connected, and that
\begin{equation}
 \{a,i\}\notin\operatorname{Sh}(H)
       \quad(a\in A,\ i\in I),
 \qquad
 \{b,e\}\notin\operatorname{Sh}(H)
       \quad(b\in B).
 \label{eq:damp-shadow-forced-corner-pattern}
\end{equation}
Then the unique cube of \(H\) with support \(P\) contains at least
\(2^{|I|}-1\) corners of \(H\).
All hypotheses depend only on the shattered supports, so the same
conclusion holds for every ample family with the same shadow.
\end{lemma}

\begin{proof}
The missing pairs prevent any support \(P\cup\{f\}\), \(f\notin P\),
from being shattered.  Ampleness supplies a cube \(C\) of support
\(P\).  It is unique: two such cubes would differ on some coordinate
\(f\notin P\), and their projections would shatter \(P\cup\{f\}\).

We first locate every edge whose direction is outside \(P\).
If \(\{f,g\}\) is not shattered, all edges of direction \(f\)
have the same value of coordinate \(g\).  Otherwise, one such edge
at \(g=0\) and one at \(g=1\) would realize all four traces
on \(\{f,g\}\).  Every active direction has an edge, so these
fixed values are defined.

It follows that all edges of each direction \(a\in A\) have a
fixed trace \(\alpha_a\in Q_I\).  If \(a,a'\) are adjacent
in \(\Gamma[A]\), ampleness supplies an \(\{a,a'\}\)-square.
Its two edge directions have the same trace on \(I\), so
\(\alpha_a=\alpha_{a'}\).  Connectivity gives one common
\(\alpha\in Q_I\) for all directions in \(A\).
For each \(b\in B\), all edges of direction \(b\) have a
fixed value \(\beta_b\) of coordinate \(e\).
An \(\{b,b'\}\)-square gives \(\beta_b=\beta_{b'}\) whenever
\(b,b'\) are adjacent in \(\Gamma[B]\).  Connectivity therefore
gives one common value \(\beta\in\{0,1\}\).

Choose any vertex \(x\in C\) such that
\[
 x_e\ne\beta,\qquad x|_I\ne\alpha.
\]
It has no neighbor in a direction of \(A\), because such an edge
would have trace \(\alpha\) on \(I\).  It has no neighbor in a
direction of \(B\), because such an edge would have value \(\beta\)
on \(e\).  Its incident directions are thus exactly \(P\), and
all their combinations occur in \(C\).  Every cube through \(x\)
is contained in \(C\), making \(x\) a corner.  The cube \(C\)
contains exactly \(2^{|I|}-1\) vertices satisfying the two displayed
conditions, which proves the bound.
\end{proof}

\begin{remark}[Sharpness of the support bound]
\label{rem:damp-shadow-forced-corners-sharp}
The bound is attained on the distinguished cube for every
\(k=|I|\ge1\).  Take \(A=\{a\}\), \(B=\{b\}\), and
let \(H\) be the union of the coordinate cubes with supports
\(\{e\}\cup I\), \(\{e,a\}\), and \(I\cup\{b\}\), all
based at the zero word.  This is a down-set and hence ample.
On the first cube, an \(a\)-edge occurs exactly when the
\(I\)-trace is zero, and a \(b\)-edge occurs exactly when
\(x_e=0\).  Such vertices are not corners: their neighbor in the
outside direction and a neighbor in a direction of the first cube
lack the fourth vertex of their square.  All other vertices of the
first cube are corners, so their number is exactly \(2^k-1\).
\end{remark}

\begin{corollary}[A forbidden three-core whose shadow forces corners]
\label{cor:damp-three-core-shadow-forces-corners}
There is a two-connected ample forbidden pc-minor \(X\) of minimum
degree three such that every ample \(H\subseteq Q_{26}\) with
\(\operatorname{Sh}(H)=\operatorname{Sh}(X)\) has at least three
corners.  In particular, preserving the full shadow cannot transform
every ample forbidden three-core into a cornerless ample family.
This remains impossible even if intermediate families need not be
forbidden pc-minors.
\end{corollary}

\begin{proof}
Take the \(1162\)-concept three-core of
Remark~\ref{ex:damp-coupled-edge-minors-1162}, with its stored coordinate
numbering, and put
\[
 e=10,\qquad I=\{15,16\},\qquad
 A=\{0,\ldots,13\}\setminus\{10\},\qquad
 B=\{14,\ldots,25\}\setminus\{15,16\}.
\]
The duplicated square supplies the \(P=\{10,15,16\}\)-cube.
The shadow of the down-set completion has no pairs between \(A\)
and \(I\), and no pairs between \(B\) and \(\{10\}\).
The induced graph on \(A\) contains the star with center \(0\),
and the one on \(B\) contains the star with center \(14\).
These support assertions are checked from the complete shadow by
\path{verify_damp_shadow_forced_corners.py}.  They depend on the
shadow alone, so Lemma~\ref{lem:damp-shadow-forced-corners} applies
to every ample realization and gives \(2^2-1=3\) corners.

The verifier independently checks the shadow, both spanning stars,
and all missing pairs.  On the original realization, the three
forced corners are \(33792,66560,99328\); their unique cube has
support \(P\).  It also replays all \(78\) proper elementary-minor
peelings and the complete \(32\)-state nonpeelability certificate,
and checks six translated sharpness examples with \(|I|=1,2,3\).
The record is \path{damp_shadow_forced_corners_verification.json}.
The universal same-shadow conclusion is the symbolic lemma, not an
enumeration of the possible signs.
\end{proof}

\begin{theorem}[Relative peeling when only small supports are added]
\label{thm:damp-relative-small-supports}
Let \(\varnothing\ne B\subseteq A\subseteq Q_F\) be ample, and assume
\begin{equation}
 |I|\le2\qquad
 \text{for every }I\in\operatorname{Sh}(A)\setminus\operatorname{Sh}(B).
 \label{eq:damp-relative-low-support-gap}
\end{equation}
Then \(A\) can be reduced to \(B\) by deleting corners of degree
at most two, without deleting any concept of \(B\).
More precisely, starting from \(B\), repeatedly adding a concept of
\(A\) with the largest number of neighbors in the current family
always gives an ample family and eventually reaches \(A\).
Every addition acquires a support of size at most two.

For every \(T\subseteq B\), including the empty target, \(A\)
can be reduced to \(T\) by corner deletions if and only if \(B\)
can be so reduced.  In particular, \(A\in\Damp\) if and only
if \(B\in\Damp\), and for every \(b\in B\), peeling \(A\)
to \(b\) is possible if and only if peeling \(B\) to \(b\)
is possible.
If the active coordinates of \(A\) and \(B\) agree, every
displayed deletion has degree exactly two.
\end{theorem}

\begin{proof}
We adapt the maximum-neighbor construction in
\cite[Proposition~4.11]{ChalopiChepoiMoran22}.  Its dimension bound
is needed only for cubes not already contained in the starting family.

First, every cube of \(A\) of dimension at least three lies in
\(B\).  To see this, extend it to a maximal cube \(Q\) of
\(A\), with support \(I\).  The reduction \(A^I\)
has one vertex for each full \(I\)-cube and is connected, by
\cite[Theorem~3.2]{ChalopiChepoiMoran22}.  The vertex corresponding
to \(Q\) is isolated: an incident edge would enlarge \(Q\)
by one coordinate.  Hence \(Q\) is the only \(I\)-cube of
\(A\).  Condition~\eqref{eq:damp-relative-low-support-gap} puts
\(I\) in \(\operatorname{Sh}(B)\).  Ampleness supplies an
\(I\)-cube in \(B\), which must therefore be \(Q\).
The same argument shows that any maximal square of \(A\) is the
only square of its support in \(A\).

Suppose that \(B\subseteq S\subsetneq A\) is ample, and
choose \(x\in A\setminus S\) maximizing its number of
neighbors in \(S\).  Since \(A\) is connected and \(S\)
is nonempty, this maximum is positive.  Put
\[
 P=\{e\in F:x^{\{e\}}\in S\}.
\]
By \cite[Lemma~4.1(i)]{ChalopiChepoiMoran22}, the punctured
\(P\)-cube at \(x\) is contained in \(S\).  Thus \(|P|\le2\):
otherwise its completion in \(A\) would be a cube of dimension
at least three containing \(x\notin B\).

If \(|P|=2\), the completed square is maximal in \(A\),
since no larger cube can contain \(x\).  It is the unique square
of that support.  Hence \(P\notin\operatorname{Sh}(S)\):
strong shattering in the ample \(S\) would give a second such
square.  Lemma~\ref{lem:damp-one-concept-extension-test} now
makes \(S\cup\{x\}\) ample, with acquired support \(P\).

If \(|P|=1\), every concept of \(A\setminus S\) has at
most one neighbor in \(S\).  Therefore \(S\) is locally
convex in \(A\): every concept of \(A\) at distance one from
two concepts of \(S\) belongs to \(S\).  Connected local
convexity implies convexity in an ample family by
\cite[Lemma~4.10]{ChalopiChepoiMoran22}.
Write \(P=\{p\}\) and \(u=x^{\{p\}}\in S\).
If some \(z\in S\) had \(z_p=x_p\), then \(x\) would
lie on a shortest \(u,z\)-path in \(A\), contrary to convexity.
Thus \(p\) is constant on \(S\), so \(P\) is not shattered
there.  The one-concept extension test again gives ampleness.

These arguments apply after each addition.  They produce the claimed
growth order, whose reverse is the required corner-deletion prefix.
For a fixed \(T\subseteq B\),
Proposition~\ref{prop:damp-rank-two-extension-invariance} preserves
the existence of a reduction to \(T\) at every addition.  Iteration
proves the asserted equivalence for every retained target; the empty
target and singleton targets give the stated special cases.
If the active coordinate sets agree, every
intermediate \(S\) already realizes both values of each active
coordinate.  The case \(|P|=1\) is then impossible, proving the
last assertion.
\end{proof}

\begin{theorem}[A canonical partial order for square attachments]
\label{thm:damp-square-interval-poset}
Under the hypotheses of Theorem~\ref{thm:damp-relative-small-supports},
suppose also that \(A\) and \(B\) have the same active coordinates,
and put \(D=A\setminus B\).  There is a unique orientation of
the edges with at least one endpoint in \(D\) such that
\begin{enumerate}
 \item every edge between \(D\) and \(B\) points toward \(B\);
 \item every vertex of \(D\) has exactly two outgoing edges;
 \item the orientation induced on \(D\) is acyclic.
\end{enumerate}
Let \(P_{A,B}\) be the partial order on \(D\) in which
\(x<y\) means that a directed path runs from \(x\) to \(y\).
The relative corner-deletion orders from \(A\) to \(B\) are
exactly its linear extensions.  Moreover,
\begin{equation}
 C\text{ is ample}
 \quad\Longleftrightarrow\quad
 A\setminus C\text{ is an order ideal of }P_{A,B}
 \qquad(B\subseteq C\subseteq A).
 \label{eq:damp-square-interval-ideals}
\end{equation}
Thus the ample intermediate families form a distributive lattice
under union and intersection.

If \(I_x\) is the pair of coordinate directions of the two
outgoing edges at \(x\), then
\[
 x\longmapsto I_x
 \quad\text{is a bijection }D\longrightarrow
 \operatorname{Sh}(A)\setminus\operatorname{Sh}(B).
\]
In particular, the support removed with each concept is independent
of the deletion order.
\end{theorem}

\begin{proof}
Take a relative deletion order supplied by the preceding theorem.
Orient each edge toward its later endpoint, regarding all of \(B\)
as retained after \(D\).  Each deleted vertex has degree two,
so this gives all three properties.

Suppose two orientations had these properties.  Their boundary edges
agree.  At each vertex of \(D\), the number of internal outgoing
edges is fixed: it is two minus the number of its neighbors in \(B\).
On the edges where the orientations differ, orient according to the
first orientation.  Each vertex then has equal indegree and outdegree
in this difference graph.  If it were nonempty, following outgoing
edges would give a directed cycle, contradicting acyclicity.  Hence
the orientations agree.

Any relative peeling loses only supports in the shadow difference.
All singleton supports remain in \(B\), so every deletion loses
a pair and has degree two.  Every relative deletion order therefore
induces this orientation and is a
linear extension of \(P_{A,B}\).  Conversely, any two linear
extensions of a finite partial order are connected by swaps of
consecutive incomparable elements: move the first element of the
desired order to the front and continue inductively.  Here incomparable
elements are nonadjacent in the graph.  Consecutive nonadjacent corner
deletions commute.  Indeed, deleting the first vertex leaves the
second's neighbor set unchanged, so the second was already a corner.
Its unique maximal cube avoids the first vertex.  The first vertex's
unique maximal cube must likewise avoid the second, since otherwise
the two cubes would coincide.  Either deletion therefore preserves
the other's corner cube.  Starting with the known relative peeling,
the swaps realize every linear extension.

If \(C\) is ample with \(B\subseteq C\subseteq A\), both
nested pairs \(C\subseteq A\) and \(B\subseteq C\) satisfy
the preceding theorem and have the same active coordinates.
Concatenate the resulting relative peelings.  The removed set
\(A\setminus C\) is an initial segment of a linear extension,
and hence an order ideal.  Conversely, every order ideal is the
initial segment of a linear extension; its corner deletions leave
an ample intermediate family.  This proves
\eqref{eq:damp-square-interval-ideals}.  Order ideals are closed
under union and intersection, so the complementary intermediate
families form the asserted distributive lattice.

Finally, each corner deletion removes exactly the support of its
two outgoing directions from the shadow.  These supports are
distinct and exhaust the shadow difference.  The orientation is
unique, proving the claimed bijection and order independence.
\end{proof}

Apply these results with \(B=K\), an ample forbidden minor of
minimum degree at least three.  Any ample \(A\supseteq K\) on
the same active coordinates, with no new shattered support of size
at least three, is obtained by square additions and is therefore
forbidden.  Its graph three-core is exactly \(K\).
This gives an admissibility test for the extension without supplying
an addition order, and a canonical description of every ample
intermediate family.  Intrinsic admissibility of the possible cores
themselves remains open.

For the \(579\)-concept example over its \(577\)-concept core,
the partial order is the chain \(8198<8194\), giving exactly three
ample intermediate families.  For the nine-square completion, the
partial order is a nine-element antichain, giving \(512\) ample
intermediate families.  All these intermediate families are forbidden.

The checker \path{verify_damp_small_support_intervals.py} verifies
the greedy construction and the canonical partial orders on the
\(579\)- and \(586\)-concept examples over their \(577\)-concept
core, and on a small example of VC dimension three.
It checks every intermediate subset in those three intervals against
\eqref{eq:damp-square-interval-ideals}.  A further example with a new
coordinate checks the degree-one step.  Its report is
\path{damp_small_support_intervals_verification.json}.

The preceding criterion recognizes an extension when its upper family is
already given.  We next specify finite data that construct the new concepts
from the lower family.  The data record which three existing concepts
complete each square; the coordinates of the added concepts are derived
from these incidences.

\begin{definition}[Square attachment data]
\label{def:damp-square-attachment-data}
Fix a nonempty ample family \(K\subseteq Q_F\), and fix an order on \(F\).
Take a finite set \(D\) of abstract symbols disjoint from \(K\).
For each \(x\in D\), specify two coordinates \(p_x<q_x\) and three
targets
\[
 u_x,v_x,w_x\in K\mathbin{\dot\cup}D.
\]
Write \(I_x=\{p_x,q_x\}\).  The data are \emph{admissible} when the
following conditions hold.
\begin{enumerate}
 \item The digraph on \(D\) with arrows from \(x\) to each of its
 targets lying in \(D\) is acyclic.
 \item The pairs \(I_x\) are distinct, and none belongs to
 \(\operatorname{Sh}(K)\).
 \item Set \(\phi(k)=k\) for \(k\in K\) and define
 \(\phi(x)=\phi(u_x)^{\{p_x\}}\).  This recursion terminates by
 item~1.  The resulting map on \(K\mathbin{\dot\cup}D\) is
 injective and satisfies
 \begin{equation}
  \phi(v_x)=\phi(x)^{\{q_x\}},\qquad
  \phi(w_x)=\phi(x)^{\{p_x,q_x\}}
  \quad(x\in D).
  \label{eq:damp-square-attachment-consistency}
 \end{equation}
\end{enumerate}
Their output is \(A=K\cup\phi(D)\).
\end{definition}

All checks in this definition concern a finite labelled digraph, coordinate
flips, equality of derived words, and the shattered pairs of the given
base.  Neither words for the new concepts nor a total order are supplied;
there is no ampleness or peeling test on the output.  The value of
\(\phi(x)\) can be computed simply by following successive \(u\)-targets
to a base concept and applying the recorded coordinate flips.

\begin{theorem}[Construction and recovery from square attachment data]
\label{thm:damp-square-attachment-data}
Admissible square attachment data over \(K\) produce an ample family
with the same active coordinates as \(K\), and
\begin{equation}
 \operatorname{Sh}(A)=\operatorname{Sh}(K)
   \mathbin{\dot\cup}\{I_x:x\in D\}.
 \label{eq:damp-square-attachment-shadow}
\end{equation}
Every order of the symbols in which all targets precede their source
constructs \(A\) by ample square additions.  Conversely, every ample
\(A\supseteq K\) with the same active coordinates and with no new
shattered support of size at least three arises this way.  Its data
are unique up to renaming the symbols of \(D\).

The reachability order of the dependency digraph is exactly the partial
order \(P_{A,K}\) of Theorem~\ref{thm:damp-square-interval-poset},
under the identification by \(\phi\).  In particular,
\[
 |D|\le \binom{|F|}{2}
       -\bigl|\{I\in\operatorname{Sh}(K):|I|=2\}\bigr|.
\]
If \(K\) is an ample forbidden pc-minor of minimum graph degree at least
three, all outputs are ample forbidden pc-minors with graph three-core
\(K\).  Conversely, every ample forbidden pc-minor is obtained from its
canonical graph three-core by its uniquely recovered attachment data.
\end{theorem}

\begin{proof}
Choose an order with every target before its source, which exists by
acyclicity.  We prove the shadow formula throughout the construction.
Suppose the concepts preceding \(x\) have already been added and form
an ample family \(S\).  Injectivity makes \(\phi(x)\) a new concept,
and the three targets are precisely the other vertices of its
\(I_x\)-square.  By induction and the distinctness condition,
\(I_x\notin\operatorname{Sh}(S)\).

There cannot be an additional neighbor of \(\phi(x)\) in an outside
direction \(e\notin I_x\).  Indeed, the punctured-cube assertion of
\cite[Lemma~4.1(i)]{ChalopiChepoiMoran22} would put the full
\(I_x\)-square based at \(\phi(x)^{\{e\}}\) inside \(S\).
That would shatter \(I_x\), a contradiction.  Thus the incident support
of the new concept is exactly \(I_x\).  The one-concept criterion,
Lemma~\ref{lem:damp-one-concept-extension-test}, proves that the addition
is ample and acquires precisely \(I_x\).  Its two directions were already
active in the punctured square.  All other coordinate values already
occurred at its targets.  Induction proves the first assertions and
\eqref{eq:damp-square-attachment-shadow} for any allowed order.

For recovery, take an upper family satisfying the stated hypotheses.
Theorem~\ref{thm:damp-square-interval-poset} gives its unique acyclic
orientation toward \(K\), with two outgoing edges at each concept of
\(A\setminus K\).  Use those concepts as temporary symbols.  Let
\(p_x<q_x\) be their two outgoing directions and let \(u_x,v_x\)
be the corresponding targets.  Set \(w_x=x^{\{p_x,q_x\}}\).
In any relative peeling this is the opposite vertex of the square
deleted with \(x\), so it is present and is deleted later, or lies
in \(K\).  All three dependency arrows therefore point forward in
that peeling.  The pairs are distinct and absent from
\(\operatorname{Sh}(K)\), by the support--concept bijection of the
same theorem.  Finally, \(\phi(x)=x\) satisfies the recursion and
\eqref{eq:damp-square-attachment-consistency}.  The recovered data are
admissible and reconstruct \(A\).

Conversely, reversing any construction order from admissible data
gives a relative peeling.  At the deletion of \(x\), its two neighbors
are exactly its prescribed targets \(u_x,v_x\).  Hence the data induce
the unique orientation just used for recovery.  This determines their
coordinate pair and both edge targets; the opposite vertex then
determines \(w_x\).  Injectivity identifies its symbol uniquely.
This proves uniqueness up to renaming.

For the assertion about partial orders, the two edge-target arrows are
precisely the oriented graph edges incident to the new concepts.
Every relative peeling also retains \(w_x\) when it deletes \(x\).
Thus, when \(w_x\in D\), \(x\) precedes \(w_x\) in every linear
extension of \(P_{A,K}\).  In a finite partial order this means
\(x<w_x\): otherwise a linear extension with \(w_x\) before \(x\)
exists.  Adding the diagonal-target arrows therefore does not change
reachability.  The cardinality bound follows from the distinct pairs.
Finally, Theorem~\ref{thm:damp-three-core} supplies both forbidden
soundness for square additions over a forbidden core and the canonical
core reduction of every ample forbidden pc-minor.  The recovery just
proved applies to that reduction.
\end{proof}

This is an intrinsic construction of all the square attachments over a
given core.  The unresolved part of the full generative classification
is still intrinsic admissibility of the cores themselves, which may
have corners.

\subsection{Rooted constructions and endpoint dependence}
\label{sec:damp-rooted-constructions}

The boundary datum in this section is a prescribed sink. By
Proposition~\ref{prop:damp-relative-uso-extension}, every assertion that a
family peels to $a$ means that it admits an acyclic USO with sink $a$.
The constructions below transport this constraint, which ordinary positivity
does not guarantee. We retain deletion orders where they give explicit
certificates and preserve the stronger relative statements.

\begin{proposition}[Operations with a prescribed endpoint]
\label{prop:damp-rooted-operations}
The following statements hold for nonempty ample families.
\begin{enumerate}
 \item Let \(S\subseteq B\subseteq Q_F\) and
 \[
  Y=(B\times\{0\})\cup(S\times\{1\}).
 \]
 For \(b\in B\), the family \(Y\) peels to \((b,0)\) if and
 only if \(B\) peels to \(b\) and \(S\in\Damp\).
 For \(s\in S\), it peels to \((s,1)\) if and only if both
 \(B\) and \(S\) peel to \(s\).
 \item The product \(A\square B\) peels to \((a,b)\) if and
 only if \(A\) peels to \(a\) and \(B\) peels to \(b\).
 \item Suppose \(H=S\cup\{c\}\) is an ample one-concept
 extension with acquired support of size at most two.  For every
 \(r\in S\), the family \(H\) peels to \(r\) if and only if
 \(S\) peels to \(r\).
\end{enumerate}
In particular, an ample addition that makes peeling to an old concept
possible must acquire a support of size at least three.
\end{proposition}

\begin{proof}
For the necessity in item~1, contract the new coordinate and restrict
to the upper layer.  Contraction preserves the root and gives the
assertion for \(B\).  An upper root is also preserved by the
upper restriction; for a lower root, ordinary pc-minor closure gives
only the required ordinary peeling of \(S\).

For a lower endpoint, first delete the entire upper layer following
an ordinary peeling of \(S\), and then peel \(B\) to \(b\).
Each upper deletion is a corner deletion: its incident support consists
of its current directions in \(S\) and the new coordinate, and the
whole corresponding lower cube is present in \(B\).
For an upper endpoint \((s,1)\), first peel the upper layer to that
vertex, retaining it.  Then follow a peeling of \(B\) to \(s\),
retaining \((s,0)\).  The single remaining upper vertex affects only
the lower vertex \((s,0)\), so all other lower deletions remain
corners.  Finish by deleting \((s,0)\) and then \((s,1)\).
This proves sufficiency in both cases.

In item~2, root-preserving face restrictions prove necessity.
For sufficiency, use the product construction of
Proposition~\ref{prop:damp-exact-expansion-product-tests} with ample
prefix orders starting at \(a\) and \(b\).  Its first product
concept is \((a,b)\); reversing the order gives the required endpoint.

For item~3, apply
Proposition~\ref{prop:damp-rank-two-extension-invariance} with
retained target \(\{r\}\).  Reducing to this singleton and then
deleting it is exactly a peeling ending at \(r\).
\end{proof}

\begin{proposition}[Unrestricted sections at a degree-two corner]
\label{prop:damp-degree-two-section-shadows}
Let $A,D,C\subseteq Q_F$ be ample families containing zero, and put
\[
 U=A\cup D\cup C,\qquad
 Y=\{(0,00)\}\cup(A\times\{01\})
       \cup(D\times\{10\})\cup(C\times\{11\}).
\]
No containment or relative-peeling assumptions are imposed.
Then $Y$ is ample if and only if all four conditions hold:
\begin{align}
 \operatorname{Sh}(U)&=\operatorname{Sh}(A)\cup
                      \operatorname{Sh}(D)\cup\operatorname{Sh}(C),\notag\\
 \operatorname{Sh}(A\cup C)&=\operatorname{Sh}(A)\cup\operatorname{Sh}(C),\notag\\
 \operatorname{Sh}(D\cup C)&=\operatorname{Sh}(D)\cup\operatorname{Sh}(C),\notag\\
 \operatorname{Sh}(A)\cap\operatorname{Sh}(D)&\subseteq\operatorname{Sh}(C).
 \label{eq:damp-degree-two-section-shadows}
\end{align}
Every ample family with a marked degree-two corner has this presentation,
uniquely up to exchanging the corner's two directions and the corresponding
side sections, after translating the corner to zero. In particular, the
criterion includes every rooted factor whose distinguished root is such
a corner. It does not impose rooted failure or rooted minimality.
\end{proposition}
\begin{proof}
Write $\mathcal A=\operatorname{Sh}(A)$,
$\mathcal D=\operatorname{Sh}(D)$, and
$\mathcal C=\operatorname{Sh}(C)$. Sorting shattered supports according
to their two new coordinates gives the exact identity
\begin{equation}
 |\operatorname{Sh}(Y)|=
 |\operatorname{Sh}(U)|+
 |\mathcal A\cap\operatorname{Sh}(D\cup C)|+
 |\mathcal D\cap\operatorname{Sh}(A\cup C)|+1.
 \label{eq:damp-general-square-shadow-count}
\end{equation}
For example, a support using the first new coordinate and not the second
has its old part shattered by both projected shores $A$ and $D\cup C$.
A support using both new coordinates cannot contain an old coordinate,
because the $00$ section is a singleton. The support consisting of the
two new coordinates is shattered, since all sections contain zero.

Suppose first that $Y$ is ample. Its two coordinate faces with sections
$(A,C)$ and $(D,C)$ are ample. The first has shadow size
$|\operatorname{Sh}(A\cup C)|+|\mathcal A\cap\mathcal C|$
and cardinality $|A|+|C|=|\mathcal A|+|\mathcal C|$. Equality forces
$\operatorname{Sh}(A\cup C)=\mathcal A\cup\mathcal C$, since
shadows always contain the union of the individual shadows. The other
face gives the analogous equality for $D\cup C$.
With these two identities, subtracting
$|Y|=1+|\mathcal A|+|\mathcal D|+|\mathcal C|$ from
\eqref{eq:damp-general-square-shadow-count} yields
\begin{equation}
 |\operatorname{Sh}(Y)|-|Y|=
 |\operatorname{Sh}(U)|-|\mathcal A\cup\mathcal D\cup\mathcal C|
       +|(\mathcal A\cap\mathcal D)\setminus\mathcal C|.
 \label{eq:damp-general-square-shadow-excess}
\end{equation}
Both terms on the right are nonnegative. Their sum is zero, so both
vanish, proving the remaining two conditions.

Conversely, the middle two conditions give
\eqref{eq:damp-general-square-shadow-excess}, and the first and last
make its right side zero. Equality in the Sandwich bound therefore
makes $Y$ ample. Its marked vertex has exactly its two new neighbours,
and their full square is present, so it is a degree-two corner.

For recovery, translate a marked degree-two corner to zero and call its
incident directions the two new coordinates. In the coordinate face
fixing both to zero, the marked vertex is isolated, since it has no
remaining incident direction. This face is ample and hence connected,
so it consists of that vertex alone. The corner square supplies zero
in each of the other three sections. Those sections are ample by face
closure and are uniquely determined by the marked vertex and its two
directions. The criterion already proved applies to them. Exchanging
the two directions only exchanges $A$ and $D$.
\end{proof}

This gives support admissibility for the full marked degree-two-corner
class, including noncontained opposite sections. In the paired special
case below, the containment and shadow hypotheses imply these four
conditions. The unrestricted criterion does not make the three sections
relatively peelable to a common core or decide peeling to the marked root.
The verification \path{research/degree_two_section_shadows.py/json}
checks all triples of ample two-coordinate sections containing zero and
the archived $870$, $872$, and $879$ presentations. These are checks of
the shadow criterion, not a rooted-factor census.

\begin{proposition}[Supports required opposite a degree-two corner]
\label{prop:damp-square-opposite-support-bound}
Let $A,B,C\subseteq\{0,1\}^n$ contain zero, where $n\le7$, and suppose
\[
 Y=\{(0,00)\}\cup(A\times\{01\})\cup(B\times\{10\})
       \cup(C\times\{11\})
\]
is ample. Set $K=A\cap C$ and define
\[
 \mathcal G=\operatorname{Sh}(A)\setminus\operatorname{Sh}(K),
 \qquad
 \mathcal L=\{T\subseteq[n]:\text{ no }S\in\mathcal G
                                      \text{ satisfies }S\subseteq T\}.
\]
Then
\[
 |B\setminus C|\le |\mathcal L\setminus\operatorname{Sh}(A)|.
\]
If $(0,00)$ is the unique corner of $Y$, necessarily
\[
 |\mathcal L\setminus\operatorname{Sh}(A)|\ge9.
\]
The same conclusion holds after exchanging $A$ and $B$.
\end{proposition}
\begin{proof}
The coordinate faces with sections $(A,C)$ and $(B,C)$ are ample.
Their section--reduction identities give
\[
 \operatorname{Sh}(A\cap C)=\operatorname{Sh}(A)\cap\operatorname{Sh}(C),
 \qquad
 \operatorname{Sh}(B\cap C)=\operatorname{Sh}(B)\cap\operatorname{Sh}(C).
\]
For completeness, the first identity follows from shadow containment
and equality of cardinalities: the ample face and its projection imply
$\operatorname{Sh}(A\cup C)=\operatorname{Sh}(A)\cup\operatorname{Sh}(C)$;
then $|A\cap C|=|A|+|C|-|A\cup C|$ is the size of the shadow
intersection. The reduction $A\cap C$ is ample. The second identity
is proved in the same way.
By Proposition~\ref{prop:damp-degree-two-section-shadows},
$\operatorname{Sh}(A)\cap\operatorname{Sh}(B)
 \subseteq\operatorname{Sh}(C)$.
Thus $\operatorname{Sh}(B)$ contains no member of $\mathcal G$.
As a down-set it contains no superset of such a member either, so
$\operatorname{Sh}(B)\subseteq\mathcal L$. Moreover,
\[
 \operatorname{Sh}(B)\setminus\operatorname{Sh}(C)
       \subseteq\mathcal L\setminus\operatorname{Sh}(A).
\]
The second section--reduction identity and ampleness now give
\[
 |B\setminus C|
 =|\operatorname{Sh}(B)\setminus\operatorname{Sh}(B\cap C)|
 =|\operatorname{Sh}(B)\setminus\operatorname{Sh}(C)|,
\]
proving the size bound.

Suppose the right-hand side of the bound is at most eight.
If $B\setminus C$ is nonempty, the relative peeling theorem for at
most eight added vertices, proved in the argument of
Corollary~\ref{cor:damp-small-component-wing-inverse}, supplies a corner
of $B$ outside $B\cap C$. This corner is nonzero and has no neighbour
in either adjacent section $C$ or $\{0\}$. Its copy in $B\times\{10\}$
is therefore a corner of $Y$ distinct from $(0,00)$.
If $B\subseteq C$ and $|B|>1$, the prescribed-endpoint theorem
(Corollary~\ref{cor:damp-seven-coordinate-peeling}) supplies a peeling
of $B$ ending at zero, whose first vertex is a nonzero corner.
Its incident cube doubles into $C$, again giving a corner of $Y$.
Finally, if $B=\{0\}$, the vertex $(0,10)$ is a corner: its only
incident directions are the two new coordinates and their full square
is present. These cases exhaust the possibilities and prove the claim.
\end{proof}

\begin{remark}[A stalled pair excluded by the support bound]
\label{ex:damp-square-thirty-eight-pair}
Encode a vertex of $\{0,1\}^6$ by the integer whose binary expansion
lists its coordinates. Let
\begin{align*}
 K={}&\{0,1,3,4,5,8,9,10,11,12,13,20,21,28,29,40,42,44,46,52,53,58,60,62\},\\
 A={}&K\cup\{16,17,32,33,34,35,36,38,41,43,48,49,50,54\}.
\end{align*}
Direct trace and cube enumeration verifies that these classes are ample,
with shattered-support counts $(1,6,15,16)$ and $(1,6,15,2)$ for $A$
and $K$, respectively. The corners of $A$ are $3,10,29,53,58$, all
in $K$, so relative peeling cannot start. Nevertheless this pair cannot
occur as $(A,A\cap C)$ in a square with a unique degree-two corner.
Indeed, $\mathcal G$ consists of fourteen triples and meets the set of
triples of every four-element support. Hence $\mathcal L$ has no support
of size four or more, and
\[
 \mathcal L\setminus\operatorname{Sh}(A)=\{7,19,38,50\}
\]
in the same binary encoding of supports. The preceding proposition
applies with bound four. The finite checks, including a member of
$\mathcal G$ in every four-element support, are reproduced by
\path{scripts/verify_square_support_bound.py} in the ample-corners project.
This rules out this marked completion, not every use of the pair in
an obstruction construction.
\end{remark}

\begin{proposition}[Exact endpoint transfer through a square]
\label{prop:damp-paired-root-transfer}
Let \(C\subseteq A,D\subseteq Q_F\) be ample families containing
\(0\), with \(|A\setminus C|,|D\setminus C|\le1\) and
\(A\cap D=C\).  Suppose that \(A\cup D\) is ample and
\(\operatorname{Sh}(A)\cap\operatorname{Sh}(D)=\operatorname{Sh}(C)\).
For two fresh coordinates \(q,r\), put
\begin{equation}
 Y=\{(0,00)\}\cup(A\times\{01\})
       \cup(D\times\{10\})\cup(C\times\{11\}),
 \qquad a=(0,00).
 \label{eq:damp-paired-rooted-factor}
\end{equation}
Then \(Y\) is ample, \(a\) is a degree-two corner, and
\begin{align}
 Y\in\Damp&\quad\Longleftrightarrow\quad C\in\Damp,
 \label{eq:damp-paired-ordinary-transfer}\\
 Y\text{ peels to }a&\quad\Longleftrightarrow\quad
 C\text{ peels to }0.
 \label{eq:damp-paired-endpoint-transfer}
\end{align}
If \(C\) uses every coordinate of \(F\), then \((Y,a)\) is a
rooted factor of Theorem~\ref{thm:damp-cut-vertex-obstructions} if
and only if \((C,0)\) is a rooted factor and both \(A,D\) peel to
\(0\).  No corner condition on the root of \(C\) is required.

If $C\in\Damp$, the endpoints in all three other sections satisfy
\begin{align}
 Y\text{ peels to }(x,01)&\ \Longleftrightarrow\ A\text{ peels to }x
                                      &&(x\in A),\notag\\
 Y\text{ peels to }(x,10)&\ \Longleftrightarrow\ D\text{ peels to }x
                                      &&(x\in D),\notag\\
 Y\text{ peels to }(x,11)&\ \Longleftrightarrow\ C\text{ peels to }x
                                      &&(x\in C).
 \label{eq:damp-paired-all-endpoints}
\end{align}
Consequently, whenever $(Y,a)$ is a rooted factor, $a$ is its only
vertex $z$ for which $(Y,z)$ is a rooted factor. Moreover, its attainable
sink set is not isometric: it contains $(0,01),(0,10)$ but neither of
their common cube neighbours $(0,00),(0,11)$.

For \(|C|>1\), the corners of \(Y\) are \(a\), the at most two
concepts outside the three copies of \(C\), and the members of the following two sets:
\begin{equation}
 \begin{split}
 &\{(x,01):x\in\operatorname{Cor}(C)\setminus\{0\},
                 d_H(x,z)\ne1\text{ for all }z\in A\setminus C\},\\
 &\{(x,10):x\in\operatorname{Cor}(C)\setminus\{0\},
                 d_H(x,z)\ne1\text{ for all }z\in D\setminus C\}.
 \end{split}
 \label{eq:damp-paired-general-corners}
\end{equation}
\end{proposition}

\begin{proof}
We first compute the shadow and give the forward peeling constructions.
The converse for the prescribed endpoint uses the first deletion among
the three nonroot copies of old \(0\).  At that moment one of two
overlaps has already been peeled to \(0\).

A support using neither new coordinate is shattered exactly when it
belongs to \(\operatorname{Sh}(A\cup D)\).  The old-coordinate projections of
the two \(q\)-sections are \(A,D\); those of the two \(r\)-sections
are \(D,A\).  The \(00\)-fibre is a singleton.  Consequently
\begin{equation}
 \operatorname{Sh}(Y)=\operatorname{Sh}(A\cup D)
 \mathbin{\dot\cup}\{I\cup\{q\}:I\in\operatorname{Sh}(C)\}
 \mathbin{\dot\cup}\{I\cup\{r\}:I\in\operatorname{Sh}(C)\}
 \mathbin{\dot\cup}\{\{q,r\}\}.
 \label{eq:damp-paired-rooted-shadow}
\end{equation}
Its size is \(|A\cup D|+2|C|+1=|Y|\), proving ampleness.
The root has exactly its two new-coordinate neighbors and their square.
Each extra concept in the \(01\)- or \(10\)-layer has no neighbor
in another layer and is a corner of its ample one-concept extension.

Necessity in \eqref{eq:damp-paired-ordinary-transfer} follows by
restricting to \(q=r=1\).  For sufficiency, delete \(a\) and both
extra concepts that are present, then peel
\(C\times\{01,10,11\}\).  Given an ordinary peeling of \(C\),
delete each old word's copies in the order \(10,11,01\).
For the reverse implication in \eqref{eq:damp-paired-endpoint-transfer},
delete the extra concepts first, retain \(a\), and use a peeling of
\(C\) to \(0\).  Delete the copies of each nonzero old word in
that same order, then finish the square at old \(0\) in the order
\(11,01,10,00\).  Every deletion has its full incident product cube.

Suppose conversely that \(Y\) peels to \(a\), and stop just before
the first deletion of a vertex \((0,t)\) with \(t\in\{01,10,11\}\).
Until this time the full square at old \(0\) is retained.  The corner
\((0,t)\) has both new directions.  It cannot have an old-coordinate
neighbor: its corner cube would also require that neighbor in the
singleton \(00\)-fibre.  Its current old-coordinate section is ample
and connected, so that section is the singleton \(\{0\}\).

If \(t=01\) or \(11\), restrict the preceding deletion sequence to
\(r=1\) and strongly project \(q\).  Initially the result is
\(A\cap C=C\).  If \(t=10\), restrict to \(q=1\) and strongly
project \(r\), obtaining \(D\cap C=C\) initially.  In each case
the projected family just before the selected deletion is \(\{0\}\),
because one of the two sections has become a singleton.  Throughout
the preceding sequence the full two-point fibre of \(0\) remains.
Record an old word when its full two-point fibre first disappears,
omitting repeated projected families.  At such a deletion the deleted
vertex has an edge to its mate.  Each incident direction of its word
in the overlap lifts to an incident direction of that vertex; the
full corner cube together with its mate layer therefore projects to
a complete corner cube in the overlap.  Thus every recorded word is
a corner of the current overlap.  The sequence ends at \(\{0\}\)
without deleting \(0\).  Appending \(0\) proves
\eqref{eq:damp-paired-endpoint-transfer}.

Assume now that \(C\) uses every old coordinate.  A root-preserving
restriction or contraction \(\mu\) in an old coordinate gives the
same square construction with \(C'=\mu C\), \(A'=\mu A\),
and \(D'=\mu D\).  The side differences still have size at most one.
They cannot share a word outside \(C'\): its \(01\)- and \(10\)-copies
would have distance two with neither intermediate word present, contrary
to isometry of the ample minor.  The union of the three old sections
is ample by contraction, and the shadow count above forces
\(\operatorname{Sh}(A')\cap\operatorname{Sh}(D')=\operatorname{Sh}(C')\).
Thus the endpoint equivalence applies to every old-coordinate image,
even if \(C'\) has constant coordinates or is a singleton.

If \((Y,a)\) is a rooted factor, the two equivalences make \(C\)
ordinarily peelable and unable to peel to \(0\).  Every elementary
root-preserving image of \(C\) peels to \(0\), by applying the
endpoint equivalence to the corresponding proper image of \(Y\).
Hence \((C,0)\) is a rooted factor.  The \(q=0\) restriction of
\(Y\), followed by contraction of \(r\), gives \((A,0)\);
the \(r=0\) restriction followed by contraction of \(q\) gives
\((D,0)\).  Rooted minimality makes both endpoints possible.

Conversely, suppose that \((C,0)\) is a rooted factor and that
\(A,D\) peel to \(0\).  The equivalences give ordinary peelability
and failure at \(a\) for \(Y\), and all old-coordinate rooted
minors peel to their roots.  The \(q=0\) restriction is \(A\) with
a leaf attached at \(0\); follow a peeling of \(A\) to \(0\) and
delete the leaf root last.  The \(r=0\) restriction uses \(D\).
Contracting \(q\) gives \(D\) at \(r=0\) and \(A\) at \(r=1\).
Delete the upper extra concept, if present, then its copy of \(C\)
using an ordinary peeling; all its corner cubes lift into the lower
\(D\)-layer.  Finish with \(D\) to \(0\).  Contracting \(r\)
interchanges \(A,D\).  All elementary rooted minors peel, and rooted
pc-minor closure proves minimality.

We next prove \eqref{eq:damp-paired-all-endpoints}, assuming $C\in\Damp$.
Necessity follows by restricting a root-ending order to the indicated
coordinate section. For an endpoint $(x,01)$, first delete the extra
concept in the $10$ section, if present, and delete $a$. Both are corners
and they are nonadjacent, so these two deletions are valid in this order.
The $10$ and $11$ sections are now both $C$, with no $00$ section.
Delete the entire $10$ copy by an ordinary order of $C$, then the entire
$11$ copy by an ordinary order of $C$. At each step its old corner cube
has all its mates in the untouched adjacent layer: first the $11$ copy
of $C$, then the $01$ copy of $A\supseteq C$. No other layer supplies
a neighbour. These are corner deletions, leaving exactly $A$ at $01$.
Finish with an order ending at $x$. Interchanging the two side sections
gives the $10$ assertion. For an endpoint $(x,11)$, delete both possible
extras and $a$. The remaining family is $C\square\{01,10,11\}$.
The second factor is a three-vertex path and peels to its middle vertex
$11$. The product endpoint rule in
Proposition~\ref{prop:damp-rooted-operations} gives sufficiency.

If $(Y,a)$ is a rooted factor, the equivalence already proved makes both
side endpoints attainable and the core endpoint unattainable. The root
$a$ is unattainable as well. This proves the missing-common-neighbours
assertion. Finally, for any other unattainable $z$, the corresponding
equivalence in \eqref{eq:damp-paired-all-endpoints} locates its failure
in its own coordinate section. This is a proper root-preserving minor
of $Y$, so $(Y,z)$ is not a rooted factor. Attainable vertices cannot
be roots of rooted factors by definition. This proves uniqueness.

Finally suppose \(|C|>1\).  Every nonzero \((x,11)\) has both new
directions but lacks \((x,00)\), so is not a corner.  A vertex
\((x,01)\), with \(x\in C\setminus\{0\}\), has its \(q\)-mate.
Its corner cube would force \(x\) to be a corner of \(C\); if the
extra concept of \(A\) is adjacent to \(x\), that cube would also
require the absent copy of the extra concept at \(11\).
Conversely, when \(x\) is a corner of \(C\) and has no such extra
neighbor, its old corner cube times the \(q\)-edge is complete.
This proves the first line of \eqref{eq:damp-paired-general-corners};
interchanging \(q,r\) proves the second.  Each nonroot copy of old
\(0\) has both new directions and an old neighbor, since \(C\) is
nontrivial and connected.  Its cube would require that neighbor in
the \(00\)-fibre, excluding all three copies.  The root and extra
concepts were checked at the start of the proof.
\end{proof}

\begin{theorem}[Rooted transfer through the two section overlaps]
\label{thm:damp-noncontained-relative-square}
Use the presentation and the four shadow conditions of
Proposition~\ref{prop:damp-degree-two-section-shadows}. Put
\[
 R=A\cap C,\qquad S=D\cap C,
\]
and suppose that $R,S\in\Damp$ and
\[
 A\searrow R,\qquad C\searrow R,\qquad
 D\searrow S,\qquad C\searrow S.
\]
Here $H\searrow K$ denotes a corner peeling retaining $K$.
Then $Y\in\Damp$ and, for $a=(0,00)$,
\begin{equation}
 \begin{split}
 Y\text{ peels to }a\quad\Longleftrightarrow\quad
 &(R\text{ peels to }0\text{ and }D\text{ peels to }0)\\
 &\quad\text{or }(S\text{ peels to }0\text{ and }A\text{ peels to }0).
 \end{split}
 \label{eq:damp-two-overlap-root-transfer}
\end{equation}
Every endpoint in another section is attainable in $Y$ exactly when
it is attainable in that section alone.

Assume every coordinate of $F$ is active in $A\cup D\cup C$.
Then $(Y,a)$ is a rooted factor exactly when all the following hold:
\begin{enumerate}
\item $A$ and $D$ peel to zero;
\item neither $R$ nor $S$ peels to zero;
\item for every old coordinate $f\in F$ and each root-preserving
restriction or contraction $\mu$ in $f$, at least one of $\mu R,\mu S$
peels to the image of zero.
\end{enumerate}
The singleton $a$ is then the only critical root of $Y$. In particular,
if $(R,0)$ and $(S,0)$ are rooted factors each using every coordinate of
$F$, item~3 holds automatically. No containment of $C$ in either side
is required.
\end{theorem}
\begin{proof}
The shadow criterion makes $Y$ ample. Isometry implies
$A\cap D\subseteq C$: a nonzero word in $(A\cap D)\setminus C$
would give two vertices in opposite side sections with neither common
neighbour present. Similarly there are no edges between $A\setminus C$
and $C\setminus A$ in their union: such an edge would produce a
missing-intermediate diagonal square in the ample face with sections
$A,C$. The corresponding assertion holds for $D,C$. Thus relative orders
in the side sections also give
\[
 A\cup C\searrow C,\qquad D\cup C\searrow C.
\]
For example, while deleting $A\setminus R$ no neighbour in
$C\setminus A$ is introduced, so each corner cube remains valid.

All three sections are positive by the relative orders and positivity
of $R,S$. Delete $a$ first. To retain $A$, clear $D\setminus S$,
delete its remaining copy of $S$ with mates in $C$, clear
$C\setminus R$ in $11$, and delete that copy of $R$ with mates in $A$.
Interchanging the sides gives a relative order retaining $D$.
To retain $C$, clear $A\setminus R$ and $D\setminus S$, then delete
the side copies of $R$ and $S$, both with mates in $C$. In each
relative step the deleted word has no mate in the adjacent layer;
in each whole-layer step its corner cube has all its mates in an
untouched containing layer. These are corner orders. Appending a
root-ending order in the retained section proves ordinary positivity
and sufficiency for every non-singleton-section endpoint. Necessity
is coordinate-face restriction.

Suppose an order ends at $a$. The two root-containing faces with
sections $\{0\},A$ and $\{0\},D$ force both side endpoints to be
attainable. Stop just before its first deletion of one of
$(0,01),(0,10),(0,11)$. The root square is still present. Its
old-coordinate section must be a singleton, since an old neighbour
in its corner cube would require an absent old neighbour in the
$00$ section. If the first deletion is at $01$ or $11$, record the
first destruction of each edge fibre in the face with sections $A,C$;
if it is at $10$, use the face with sections $D,C$. As in
Proposition~\ref{prop:damp-paired-root-transfer}, a broken fibre's
corner cube projects to a corner cube in the remaining reduction.
The zero fibre survives until the stopping time. The resulting order
peels $R$ or $S$, respectively, to zero. This proves necessity in
\eqref{eq:damp-two-overlap-root-transfer}.

For sufficiency, suppose $R$ and $D$ peel to zero. Relatively clear
$A\setminus R$ in $01$, then peel its copy of $R$ to zero with mates
in $C$. Delete $(0,01)$, now a corner of its root square. Clear
$C\setminus S$ in $11$, and delete the remaining copy of $S$ by an
ordinary order with mates in $D$. The remaining family is $D$ with
a leaf $a$ attached at zero; peel $D$ to zero and finish at $a$.
The other alternative follows by interchanging the two sides.

Relative orders pass to coordinate images: their residual images are
ample and lose at most one word at each step; omitting repetitions
gives a relative order to the image target. For every old-coordinate
operation one also has
\[
 \mu R=\mu A\cap\mu C,\qquad
 \mu S=\mu D\cap\mu C.
\]
Restrictions are immediate. A new intersection under contraction would
give a missing-intermediate diagonal square in the original ample
adjacent face. Consequently all old-coordinate rooted images satisfy
the same hypotheses, including positivity of the image overlaps; their
shadow conditions hold because the images of $Y$ are ample.

If $(Y,a)$ is a rooted factor, its two new-coordinate root restrictions
force item~1. The endpoint equivalence then makes root failure exactly
item~2. Applying that equivalence to each old-coordinate rooted image,
and using endpoint minor closure for $A,D$, gives item~3. These images
are proper by the active-coordinate assumption.

Conversely, items~1--3 give positivity of every old-coordinate rooted
image by \eqref{eq:damp-two-overlap-root-transfer}. The two new-coordinate
root restrictions are positive leaves on $A,D$. One new contraction has
sections $A,D\cup C$ and root in $A$. Use
$D\cup C\searrow C\searrow R$, delete the remaining $R$ layer with
mates in $A$, and finish at zero in $A$. Here
$(D\cup C)\cap A=R$, so relative deletions outside $R$ have no mates
in the retained layer. The other contraction similarly uses
$A\cup C\searrow C\searrow S$ and finishes in $D$.
Thus all elementary rooted images are positive. Ordinary positivity
and root failure, already proved, give rooted minimality. Every other
unattainable vertex fails in its own proper section by the full endpoint
formula, so cannot be a critical root. The stated sufficient condition
for item~3 is immediate from the definition of the two rooted factors.
\end{proof}

If $C\subseteq D$, then $S=C$. Since $C\searrow R$, an endpoint
order of $R$ extends to $C$, including after every coordinate operation.
The endpoint condition reduces to attainability of zero in $A,C$.
If all old coordinates are active in $C$, the rooted-factor criterion
reduces to $(C,0)$ being a rooted factor and $A,D$ peeling to zero.
This is the one-containing-side rule, including the $872$ example;
the contained relative-wing construction is its special case $R=C$.

Both containments can fail. In the partial-tip construction with old
core $B$ and a tip $\gamma$ occurring only in the opposite section,
put $A=B\cup\{128\}$, $D=B\cup\{290\}$, and
$C=B\cup\{\gamma\}$. Then $R=S=B$. For the archived tips $64$ and
$604$, all four relative orders remove a single tip. The theorem covers
both $871$-concept rooted factors, including the case $\gamma=604$
where $C$ itself peels to zero. Thus the obstruction at the singleton
is controlled by the two overlaps, not by the opposite section alone.

The verification \path{research/two_overlap_square.py/json} checks these
two presentations, their clearing orders and the two new-coordinate
contraction orders. The earlier one-sided schedules remain in
\path{research/noncontained_relative_square.py/json}. Intrinsic selection
of admissible inputs within this construction class remains open. Universal
recovery of its four relative orders is false: Theorem~\ref{thm:damp-positive-overlap-rooted-factor}
gives a rooted factor with both overlaps root-positive and a side that
cannot retain its overlap. Thus this sound construction is not exhaustive.

\begin{corollary}[Small overlaps in a degree-two rooted factor]
\label{cor:damp-small-rooted-overlaps}
Let $(Y,a)$ be a rooted factor in the canonical degree-two presentation
with sections $A,D,C$, and put $R=A\cap C$, $S=D\cap C$.
Then $|R|,|S|\ge3$. If either overlap has three words, its one-inclusion
graph is a three-vertex path with zero at an endpoint.
Moreover, a rooted factor whose distinguished root is a degree-two
corner uses at least eight coordinates.
\end{corollary}
\begin{proof}
The two new-coordinate reductions are, respectively, a copy of $R$
with a leaf attached at zero and a copy of $S$ with such a leaf.
Their orders are $|R|+1$ and $|S|+1$. By
Theorem~\ref{thm:damp-star-reduction-positive}, every coordinate class
of a rooted factor has at least four edges, and a four-word reduction
must be a four-vertex path. This proves the size bounds. An ample
three-word family is a path; adjoining a leaf at its middle vertex
would give a star, whereas adjoining it at an endpoint gives the
required four-vertex path.

For the last assertion, suppose the whole family uses at most seven
coordinates. The three sections $A,D,C$ then use at most five, and
their nonempty overlaps $R,S$ are ample reductions of faces of $Y$.
Theorem~\ref{thm:damp-five-coordinate-relative-peeling} supplies
$A\searrow R$, $C\searrow R$, $D\searrow S$, and $C\searrow S$,
and also peelings of $A,D,R,S$ to zero. All hypotheses of
Theorem~\ref{thm:damp-noncontained-relative-square} therefore hold,
and its endpoint criterion makes $Y$ peel to its distinguished root.
This contradicts the rooted-factor condition.
\end{proof}

\begin{theorem}[From edge retention to a prescribed root]
\label{thm:damp-path-rooted-transfer}
Fix the $289$-vertex family $B$ and its path
$H=a-b-c-d=2-0-4-5$ from
Theorem~\ref{thm:damp-four-vertex-path-obstruction}.
Let $P$ be any ample corner-peelable family containing the edge
$D=\{c,d\}$, using fresh coordinates except for the direction of $cd$,
and suppose that $P$ peels to $c$. Embed it with $B\cap P=D$.
For a fresh coordinate $e$, put
\[
 K(P)=(B\times\{0\})\cup((H\cup P)\times\{1\}),
 \qquad r=(a,1).
\]
Then $K(P)$ is ample and corner-peelable, $r$ is a degree-two corner,
and
\begin{equation}
 K(P)\text{ peels to }r
 \quad\Longleftrightarrow\quad P\text{ peels to }D.
 \label{eq:damp-path-rooted-transfer}
\end{equation}
Moreover $(K(P),r)$ is a rooted factor if and only if $P$ fails to
peel to $D$ but every elementary restriction or contraction in a private
coordinate of $P$ that retains $D$ peels to its image of $D$.
The output has $|P|+291$ vertices and $\dim(P)+12$ active coordinates.
When these equivalent minimality conditions hold, gluing $K(P)$ at $r$
to the fixed rooted factor $(B,b)$ on disjoint coordinates gives an ample
forbidden pc-minor with $|P|+579$ vertices.
\end{theorem}
\begin{proof}
We list precisely the fixed properties of $B$ used here. Its only
geometric corner is $b$. It has an acyclic USO with sink $a$ and
out-map $\{h\}$ at $d$, where $h=0$ is the direction of $cd$.
For every $j\in\{3,\ldots,11\}$, both its $j=0$ restriction and its
$j$-contraction peel to $H$. The $18$ relative orders and the prescribed
USO are already recorded in the four-path construction;
\path{research/path_rooted_transfer.py/json} independently replays these
fixed hypotheses. No property of the old $330$-vertex input is assumed
for the variable piece $P$.

Write $f=1,g=2$ for the directions of $ab,bc$. The edge-gluing law
combines the prescribed map of $B$ with a map of $P$ having sink $c$:
at $d$ the first map has only the outgoing edge to $c$, and at $c$
the second has no outgoing edge. Thus $U=B\cup P$ is ample and has
an acyclic USO with sink $a$. The upper section $H\cup P$ is a vertex
amalgam of $P$ and the path $a-b-c$, hence ample and peelable to $a$.
Its shadow intersects that of $B$ in precisely
$\{\varnothing,\{f\},\{g\},\{h\}\}=\operatorname{Sh}(H)$.
The section shadow identity proves ampleness of $K(P)$.

For ordinary positivity, delete $(a,1)$ and then $(b,1)$, both as
square corners. The remaining family is the edge amalgam of $B\times\{0\}$
and
\[
 T=(P\times\{1\})\cup(D\times\{0\}).
\]
The contraction of $T$ is $P$ and its reduction is the convex edge $D$.
Theorem~\ref{thm:damp-convex-reduction-lift} gives an acyclic USO of $T$
with sink $(c,0)$. Together with the prescribed out-map $\{h\}$ at
$(d,0)$ in $B$, the edge-gluing law proves positivity of the remainder.
Also $(a,1)$ has exactly its two square neighbours, proving the corner
and degree assertion.

If $P\searrow D$, remove its private upper vertices in this relative
order, then delete $(d,1)$ as a square corner. The remaining expansion
has contraction $B$ and reduction $a-b-c$, a star. Since $B$ peels to $a$,
Theorem~\ref{thm:damp-star-reduction-positive} supplies the required
endpoint $r$.

For the converse, suppose an acyclic USO of $K(P)$ has sink $r$ and
that $P$ cannot retain $D$. Restrict the orientation to its lower face
$B$. Every source is a geometric corner, so $b$ is its unique source.
Every lower vertex, including $d$, is therefore reachable from $b$.
The lower edges from $b$ to $a,c$ both point outwards. The vertex
$(a+ c,1)=(6,1)$ is absent (here $b=0$), so (C1) forces the vertical
edge over $b$ to point from $(b,1)$ to $(b,0)$.
The sink $r$ forces $(b,1)\to(a,1)$. The absent upper square on $f,g$
then forces $(c,1)\to(b,1)$.

No edge of the upper copy of $P$ can point out of $(c,1)$. Any such
edge, together with the outgoing edge towards $(b,1)$, would require
an absent square: for direction $h$ its missing vertex is $(1,1)$;
for a private direction of $P$ the missing vertex has that direction
nonzero over $b$. Thus the restriction to $P$ has sink $c$.
If $d$ had no outgoing private edge in that restriction, both boundary
out-maps would be contained in $\{h\}$, giving $P\searrow D$ by
Proposition~\ref{prop:damp-relative-uso-extension}. Hence $d$ has an
outgoing private edge. Its missing lower mate forces the vertical edge
at $d$ to point from $(d,0)$ to $(d,1)$, again by (C1).
The acyclic USO of $P$ has unique sink $c$, so there is a directed path
from $(d,1)$ to $(c,1)$. Concatenating paths now gives a directed closed
walk
\[
 (b,0)\leadsto(d,0)\to(d,1)\leadsto(c,1)
 \to(b,1)\to(b,0),
\]
a contradiction. This proves \eqref{eq:damp-path-rooted-transfer}.

It remains to check exactly when all proper rooted minors are positive.
A private-coordinate operation on $P$ retaining $D$ changes $K(P)$
into $K(\mu P)$. The image still peels to $c$, by rooted minor closure.
The equivalence just proved makes these rooted images positive exactly
when all the stated relative images of $P$ retain $D$. This proves
necessity and handles all private coordinates in the converse.

For a coordinate $j\in\{3,\ldots,11\}$ of $B$, use its recorded
relative order to clear the lower section outside $H$. These are corners
of the whole image since they have no upper mates. Delete the remaining
lower path in order $d,c,b,a$: the first three are square corners and
the last is a leaf. The upper section $H\cup P$ then peels to $r$.
For each of the six root-preserving elementary operations in $f,g,h$,
the reduction is a nonempty star: contractions shorten $H$ to three
vertices; the root-containing restrictions leave respectively
$\{a\}$, $\{a,b\}$, and $\{a,b,c\}$.
Reduction commutes with these old-coordinate operations by the no-diagonal
square identity for ample expansions. Their contractions are the corresponding
root-preserving images of $U$, which already peels to $a$. The star endpoint
theorem therefore makes all six rooted images positive.
Finally the root-containing $e$-section is $H\cup P$, positive to $a$,
and the $e$-contraction is $U$, also positive to $a$.
These exhaust the elementary rooted minors. Rooted minor closure proves
sufficiency. The cardinality and dimension follow directly from the two
sections and their private coordinate sets. The fixed $(B,b)$ is the
established $289$-vertex rooted factor; Theorem~\ref{thm:damp-cut-vertex-obstructions}
therefore proves the final gluing assertion, with order
$(|P|+291)+289-1=|P|+579$.
\end{proof}

This is a uniform construction and exact recovery of the relative input
condition within the displayed family. The old $621$-vertex factor is its
instance $|P|=330$. The earlier $909$-vertex proof used the stronger
condition $\ell_d(P)=0$ to prohibit a first path deletion; the directed-cycle
argument above identifies the exact endpoint condition $r_D(P)=0$ instead.
The input $P$ and its boundary are recovered from the specified upper face.
Intrinsic generation of the minor-minimal edge-retention failures is still
required; the theorem does not assert that every rooted factor has this form.

\begin{corollary}[Iteration of minimal edge-retention failures]
\label{cor:damp-iterable-path-edge}
Assume that $P$ satisfies the minimality conditions of
Theorem~\ref{thm:damp-path-rooted-transfer}. Mark in $K(P)$ the edge
\[
 E_a=\{(a,0),(a,1)\}.
\]
Then $K(P)$ is again an ample positive family failing to retain its marked
edge, every private elementary minor retaining that edge does retain it,
and its endpoint $(a,0)$ is attainable. Its other endpoint $(a,1)$ is
unattainable. Both sections in the direction of $E_a$, and the reduction
in that direction, peel to their common projected endpoint $a$.
Thus even under private-coordinate relative minimality, these three
projected endpoint tests do not imply edge retention.

Starting from the certified $330$-vertex edge piece and repeatedly applying
this construction, always marking $E_a$ and taking $(a,0)$ as the attainable
endpoint for the next input, gives edge pieces $P_k$, $k\ge0$, with
\[
 |P_k|=330+291k,\qquad \dim(P_k)=13+12k.
\]
For each $k\ge0$, the vertex-gluing construction of the theorem gives an
ample forbidden pc-minor with $909+291k$ vertices in $37+12k$ dimensions.
\end{corollary}
\begin{proof}
The theorem makes $(a,1)$ unattainable, so retaining $E_a$ is impossible.
The ordinary-positivity proof there can end at $(a,0)$: after the two upper
square deletions, glue the map of $B$ with sink $a$ to the map of the other
piece with sink $(c,0)$. The only global sink of the glued map is $(a,0)$.
The lower section $B$ peels to $a$, the upper section $H\cup P$ peels to
$a$, and the reduction is the path $H$, which also peels to $a$.

We verify the stronger claim of retaining the entire edge in every private
minor, rather than merely attaining its upper endpoint. If a private
coordinate of $P$ is restricted to the boundary bit or contracted, the
changed piece $P'$ peels to $D$ by hypothesis. Clear its private upper
vertices, then delete $(d,1),(c,1),(b,1)$ as successive square corners.
What remains is lower $B$ and its upper mate at $a$. A peeling of $B$ to
$a$ deletes only vertices without vertical mates, leaving exactly $E_a$.

For operations in a coordinate of $B$, we use the following fixed finite
property, stronger than the $18$ orders needed in the preceding theorem:
every root-$a$-preserving elementary minor $B'$ of $B$ peels to the
corresponding path image $H'$. The $18$ previously recorded orders cover
coordinates $3,\ldots,11$. The six remaining orders are stored in
\path{research/path_rooted_transfer_extra_orders.json}; the same verifier
replays all $24$ orders. The target path is a singleton, an edge, or a
three- or four-vertex path with endpoint $a'$.

First clear the lower section to $H'$ using that relative order. The upper
section is $H'$ together with the image $P'$ when it survives. In a private
coordinate of $B$ outside $f,g,h$, the piece $P$ is unchanged. In directions
$f,g$, it is unchanged up to projection, or disappears in the root-containing
restriction. In direction $h$, it is replaced by a root-$c$-preserving image,
which still peels to $c'$. Whenever it survives, its intersection with $H'$
is the endpoint edge $c'd'$, or the singleton $c'$ after contraction or
restriction in $h$. The same assertion holds after the contractions in
$f,g$, with the shortened path.

Delete lower vertices of $H'$ successively from the endpoint opposite $a'$
towards $a'$, retaining $(a',0)$. Each deletion is a square corner since
both copies of the still relevant path edge are present. The remaining
upper section peels to $(a',1)$: clear $P'$ to $c'$ when present, then
clear the upper path towards $a'$. All deleted upper vertices lack lower
mates, since only $(a',0)$ remains. Thus these are valid deletions in the
whole family and leave the marked edge. If $H'$ is a singleton, the first
sequence is empty and the same conclusion applies. This covers every
private coordinate of the new marked edge; its own coordinate $e$ is
excluded from private operations. Relative minor closure extends the result
to every proper sequence retaining the edge.

Hence the output satisfies the same input conditions, with its lower
endpoint designated as $c$. Iteration is valid after renaming private
coordinates. Each step adds $291$ vertices and $12$ coordinates, giving
the displayed formulas. The forbidden-output formulas follow from the
preceding theorem's additions of $579$ vertices and $24$ coordinates.
\end{proof}

The stronger condition that \emph{both endpoints of the whole family}
be attainable is distinct from the three projected endpoint tests refuted
here: every output of this corollary has one unattainable endpoint. The
iteration is an explicit closed construction within this branch, not an
exhaustive description of all minimal edge-retention failures or all
forbidden pc-minors.

\begin{theorem}[Root-positive overlaps in a degree-two rooted factor]
\label{thm:damp-positive-overlap-rooted-factor}
There is a $621$-word rooted factor $(K,r)$ of dimension $25$ whose
root is a degree-two corner and whose canonical sections satisfy
\[
 |A|=331,\quad |D|=149,\quad |C|=140,\qquad
 R=A\cap C=\{0,4,5\},\quad |S|=|D\cap C|=65.
\]
All five families $A,D,C,R,S$ peel to zero. The three-word overlap
$R$ is a path rooted at its endpoint zero, attaining the lower bound
of Corollary~\ref{cor:damp-small-rooted-overlaps}. Nevertheless
$A$ does not peel relatively to $R$.
Thus neither root failure in both overlaps nor recovery of the four
relative peelings in Theorem~\ref{thm:damp-noncontained-relative-square}
is necessary for a degree-two rooted factor.
\end{theorem}
\begin{proof}
Use the concrete $289$-word family $B$, $330$-word edge piece $P$, and
their embeddings from Theorem~\ref{thm:damp-four-vertex-path-obstruction}.
In that construction
\[
 H=a-b-c-d=2-0-4-5,\qquad B\cap P=\{c,d\}.
\]
Let $e$ be its new coordinate and put
\[
 K=(B\times\{0\})\cup((H\cup P)\times\{1\}),
 \qquad r=(a,1).
\]
The $909$-word forbidden class constructed there is exactly $K\vee_r Q$,
where $Q$ is its additional $289$-word copy of $B$, attached only at $r$
and using a disjoint private coordinate block. Both pieces contain
vertices other than $r$, and no edge joins their exclusive parts, so
$r$ is a cut vertex. The converse in
Theorem~\ref{thm:damp-cut-vertex-obstructions} makes $(K,r)$ a rooted
factor. Its order is $909-(289-1)=621$, and deleting the twelve exclusive
coordinates of $Q$ leaves dimension $37-12=25$. In $H\cup P$, the
vertex $a$ has only neighbour $b$; its other neighbour in $K$ is its
lower mate. These two directions span the square on $a,b$ in the two
layers, so $r$ is a degree-two corner.

Write $f=1$ for the coordinate of $ab$ (bitmask $2$). Translate $r$
to zero and take $e,f$ as the two new directions of the canonical
presentation, with $A$ the upper section on the $b$ side. Suppressing
these fixed coordinates gives
\[
 A=\{b\}\cup P,\quad D=\pi_f(B_{f=1}),\quad
 C=\pi_f(B_{f=0}),\quad R=\{b,c,d\},\quad S=B^{\{f\}}.
\]
Here $b=0$ in the remaining coordinates. The stated orders of the
sections and overlaps follow from the explicit input sets; they are
checked in \path{research/positive_overlap_rooted_factor.py/json}.

The edge profile of $P$ in Theorem~\ref{thm:damp-three-piece-witness}
gives a peeling to $c$ and forbids a relative peeling to $\{c,d\}$.
Since $b$ is attached to $P$ as a leaf at $c$, first peeling $P$ to
$c$ and then deleting $c$ peels $A$ to $b$. A relative peeling of $A$
to $R$ would remove exactly the vertices of $P\setminus\{c,d\}$.
Their corner tests are unchanged by the retained leaf $b$, so it would
peel $P$ relatively to $\{c,d\}$, a contradiction.

The original $(B,b)$ is a rooted factor, hence its proper $f=0$
restriction peels to $b$, proving the assertion for $C$. The prescribed
core order in the four-edge construction peels $B$ to $a$, and its
$f=1$ restriction gives the assertion for $D$. The path $R$ peels to
its endpoint $b$. Finally $B$ has VC dimension three, so its nonempty
reduction $S$ has VC dimension at most two. The VC-two prescribed-endpoint
result recalled in Remark~\ref{rem:damp-one-corner-peelable} makes $S$
peel to zero. The verifier also checks its $65$-word shadow and an
explicit endpoint order.

For the finite certificates, restrict the archived proper-minor orders
of the $909$-word class to the changed $K$ piece. Its unchanged $Q$ piece
still fails at the common root. The vertex-gluing criterion therefore
forces every changed $K$ piece to peel to that root. This gives all
$50$ elementary rooted tests in dimension $25$: $48$ explicit endpoint
orders replay, and the two source contractions certified by the star
reduction theorem pass by the same vertex-gluing argument. An ordinary
order of $K$ is obtained by restriction from a proper face of the
$909$-word class. Thus the computational records agree with the
cut-vertex proof, rather than substituting a search for it.
\end{proof}

This is a recovered block of an established forbidden class, not a new
$909$-word obstruction. Its new role is decisive for recovery: even
ordinary positivity, all three section endpoints, both overlap endpoints,
and every proper rooted minor can coexist with failure at the singleton
root. The four-relative-order construction remains sound but is not
exhaustive. The failure here is the already established inability of the
edge piece $P$ to retain $cd$, now appearing as $A\not\searrow R$.

\begin{theorem}[An adjacent degree-three corner repairs a rooted factor]
\label{thm:damp-adjacent-rank-three-repair}
Let $(K,r)$ be the $621$-word rooted factor of
Theorem~\ref{thm:damp-positive-overlap-rooted-factor}, with
$r=(2,1)$ and new coordinate $e$. Put $v=(6,1)$ and
\[
 W=K\cup\{v\}.
\]
Then $W$ is ample and peels to $r$. The vertices $r,v$ are adjacent
corners of degree three, and deleting $v$ destroys peeling to $r$.
Thus degree three is the least possible degree of a nonroot corner
whose deletion can destroy a prescribed endpoint, even when that corner
is adjacent to a corner root and the resulting family is a rooted factor.

There is also a $621$-word positive partner $T=W\setminus\{(0,1)\}$
that peels to $r$ and has the same shadow as $K$. Together they realize
the two-square endpoint switch of
Theorem~\ref{thm:damp-two-square-endpoint-switch} over a $619$-word base.
\end{theorem}
\begin{proof}
Keep the concrete notation $B,P,H=2-0-4-5$ from the preceding theorem,
and put $f=1$, $g=2$, $h=0$ for the three old path directions.
The contraction of $W$ in $e$ is unchanged:
\[
 U=B\cup P,\qquad |U|=617.
\]
Its upper section is $P\cup L$, where $L=\{0,2,4,6\}$ is the
square on directions $f,g$. This is a vertex amalgam at $4$, using
disjoint exclusive coordinate blocks, so it is ample and has $333$ words.
The lower section is $B$. Their common shadow consists of exactly
\[
 \varnothing,\quad\{h\},\quad\{f\},\quad\{g\},\quad\{f,g\}.
\]
Indeed the directions of $P$ and $B$ intersect only in $h$, while
$B$ contains the full square $L$. The contraction $U$ is ample as a
minor of $K$. The two-section shadow identity therefore gives
$|\operatorname{Sh}(W)|=617+5=622=|W|$, proving ampleness.
Compared with $K$, its only new shattered support is $\{e,f,g\}$.

Crucially, its enlarged reduction is convex:
\[
 W^{\{e\}}=H\cup\{6\}=\{0,2,4,5,6\}
              =U\cap Q_{\{f,g,h\}}(0).
\]
The last equality is an explicit feature of the input $B$ and the
embedding of $P$, checked in the replay below. Since $(K,r)$ is a
rooted factor, $U$ peels to $2$. Apply the prescribed-endpoint part of
Theorem~\ref{thm:damp-convex-reduction-lift}. At the contraction's sink
$2$, no old edge is outgoing, so its vertical direction is free; direct
it towards the upper copy $r$. The resulting acyclic USO of $W$ has
sink $r$, and hence gives a peeling to $r$ by CCMW's correspondence.
The eight vertices on old words $0,2,4,6$ in the two layers form a
three-cube. Both $r$ and $v$ have exactly their three neighbours in
that cube, and are adjacent along $g$. They are therefore degree-three
corners. Deleting $v$ leaves the root-negative factor $K$.
The support-at-most-two endpoint invariance in
Proposition~\ref{prop:damp-rooted-operations} proves minimality of degree
three; it does not assert minimality of the example's order or dimension.

For the switch, write $z=(0,1)$ and $G=W\setminus\{r,v,z\}$.
The same cube makes $z$ a degree-three corner with support $\{e,f,g\}$.
The recorded endpoint order of $W$ can start with $z$; deleting that
first term gives a peeling of $T=W-z$ to $r$. Both corner deletions
$W-v$ and $W-z$ remove just the same triple support, so $K,T$ have the
same shadow. Moreover
\[
 |G|=619,\qquad K=G\cup\{z,r\},\qquad T=G\cup\{v,r\}.
\]
Deleting $r$ first from $W$ leaves $v,z$ as square corners; deleting
both makes $G$ ample. These are exactly the two square attachments in
the cited switch theorem, now with opposite endpoint outcomes.

The verifier \path{research/adjacent_621_repair.py/json} constructs the
lift from the archived contraction order, replays the $622$-word endpoint
order and the $621$-word positive partner order, and independently checks
nonclashing, the exact shadow and out-map bijection, acyclicity, the
convex face, and all three degree-three corners. The earlier
\path{research/adjacent_621_repair_probe.py/json} contains only failed
heuristics; those inconclusive searches are superseded by this positive
certificate.
\end{proof}

This closes existence of the adjacent rank-three repair and of an
endpoint-changing two-square switch. It does not characterize all such
repairs or make arbitrary degree-three corner deletions valid during a
generative descent. The effective ingredient here is the convexity of
the enlarged reduction, which allows a prescribed contraction sink to lift.

\begin{proposition}[Canonical saturation of the positive-overlap factor]
\label{prop:damp-saturation-621}
For the factor $K$ and repair $v$ of
Theorem~\ref{thm:damp-adjacent-rank-three-repair}, put
$F=\{0,\ldots,11\}\cup\{24,\ldots,36\}$, $J=\{1,2,36\}$,
and
\[
 A_J=\{x\in Q_F:x|_J\ne v|_J\}.
\]
Then $A_J$ peels to $K$. In particular, failure of the canonical
four-relative-order recovery for $K$ does not prevent saturation from
its one-word endpoint repair. Applying the relative-square construction
to $K$, $K\cup\{v\}$ and $A_J$, after translating its root to zero,
gives a rooted factor of order $29\,361\,372$.
\end{proposition}
\begin{proof}
Split the coordinates outside $J$ into the common coordinate $0$ and
private sets
\[
 U=\{3,\ldots,11\},\qquad V=\{24,\ldots,35\}.
\]
The explicit presentation of $K$ shows that every word has either its
$U$-block or its $V$-block zero. Set
\[
 Z=(Q_U\times\{0_V\})\cup(\{0_U\}\times Q_V),\qquad
 L=(Q_J\setminus\{v|_J\})\times Q_{\{0\}}\times Z.
\]
Thus $K\subseteq L\subseteq A_J$. The cube $Q_{U\cup V}$ peels to
$Z$. Indeed its complement of $Z$ is
$(Q_U\setminus\{0_U\})\times(Q_V\setminus\{0_V\})$, a product
of corner-peelable punctured cubes. Reversing a corner order of this
complement and taking complements yields an ample chain from the full
cube to $Z$, hence the required relative corner order. Taking a product
with a fixed ample family preserves relative peeling: for each corner
deleted in the first factor, delete its entire fibre in any corner order
of the second factor, when that factor is corner-peelable. Here the second
factor is a punctured three-cube times an edge, which is corner-peelable.
At each such deletion the corner cube is the product of the two current
corner cubes. Consequently $A_J\searrow L$.

The remaining finite order is recorded in
\path{research/saturation_621.json}. Its independent verifier
\path{research/saturation_621.py} reconstructs $L$ and checks every
subset of the incident directions at each deleted vertex, retaining all
of $K$. It verifies $63\,877$ corner deletions from
\[
 |L|=14(2^9+2^{12}-1)=64\,498
 \quad\text{to}\quad |K|=621.
\]
This proves $L\searrow K$ and hence the first assertion, without
materializing the $|A_J|=7\cdot2^{22}=29\,360\,128$ words of the wing.

For the construction consequence, $W=K\cup\{v\}$ peels to the root
by Theorem~\ref{thm:damp-adjacent-rank-three-repair}, and $W\searrow K$
by deleting $v$. The support $J$ is the unique support acquired from $K$;
thus $A_J\cap W=K$ and
$\operatorname{Sh}(A_J)\cap\operatorname{Sh}(W)=\operatorname{Sh}(K)$.
The union is ample by adjoining $v$ to the trace-avoidance wing: all its
neighbours are in directions $J$, their cube is complete after addition,
and $J$ was not shattered. A punctured cube admits every prescribed
endpoint other than its missing vertex, so $A_J$ peels to the root of $K$.
Theorem~\ref{thm:damp-noncontained-relative-square}, with both overlaps
equal to $K$, now gives a rooted factor. Its order is
$1+|A_J|+|W|+|K|=29\,361\,372$.
\end{proof}

This is a certified application on a new eligible input, not a proof
that every repaired rooted factor admits canonical saturation. The
compression removes only words supported on both private blocks;
its success does not establish a converse from arbitrary wing orders
to orders through $L$. The increasing-label search stalled and was
inconclusive; the stored decreasing-label order proves the positive result.

\begin{proposition}[Preserving a positive overlap under rooted minimization]
\label{prop:damp-positive-overlap-minimization}
The following existence statements are equivalent.
\begin{enumerate}
\item There is a rooted factor with a degree-two corner root whose
canonical sections $A,D,C$ have $A\cap C$ peelable to zero, after
possibly interchanging the two root directions.
\item There is an ample degree-two presentation
\[
 Y=\{(0,00)\}\cup(A\times\{01\})
      \cup(D\times\{10\})\cup(C\times\{11\}),\qquad a=(0,00),
\]
satisfying all four conditions:
\begin{enumerate}
\item $Y$ is ordinarily peelable;
\item $Y$ does not peel to $a$;
\item both contractions in the two new coordinates peel to the image of $a$;
\item $A\cap C$ peels to zero.
\end{enumerate}
No tests on old-coordinate minors are required in item~2.
\end{enumerate}
More precisely, every family in item~2 has a rooted-factor minor obtained
using only old-coordinate operations. Its marked root is still a
degree-two corner, and its corresponding overlap still peels to zero.
\end{proposition}
\begin{proof}
A rooted factor in item~1 satisfies item~2 by definition and the
canonical support presentation in
Proposition~\ref{prop:damp-degree-two-section-shadows}.
For the converse, first note that the two positive contractions also
force $A,D$ to peel to zero. Indeed, one contracted family has sections
$A,D\cup C$ with its marked root in the $A$ section; restriction of
its endpoint order gives an endpoint order of $A$. The other contraction
similarly gives one of $D$. The two root-containing new-coordinate
restrictions of $Y$ are therefore positive: each is a side section
with the root attached as a leaf at its zero word. Thus all four
elementary root-preserving operations in the two new coordinates give
positive endpoints.

Restrictions and contractions in distinct coordinates commute.
Consequently a root-preserving pc-minor using either new coordinate
peels to its root: move its first operation in that coordinate to the
front, use the preceding positivity, and apply endpoint minor closure
to the remaining operations. Every root-negative minor must therefore
retain both new coordinates.

Choose a minimal root-negative minor of $Y$. Equivalently, repeatedly
apply a root-negative elementary operation in an old active coordinate
until none is available. Such an operation exists whenever a proper
root-negative descendant exists, since positivity of its first image
would imply positivity of the descendant. Cardinality decreases, so
this procedure terminates. Ordinary peelability is preserved by pc-minors.
The final family is root-negative and has every proper root-preserving
minor positive, hence is a rooted factor.

It remains to check the extra structure, which is the point of retaining
the two new directions. Every old-coordinate operation $\mu$ preserves
the singleton $00$ section and the three other zero copies. The marked
root therefore remains a degree-two corner. Moreover,
\[
 \mu(A\cap C)=\mu A\cap\mu C.
\]
For restrictions this is immediate. For a contraction, any new
intersection not arising from the old intersection would give two
opposite vertices of a square in the ample face with sections $A,C$
with both intermediate vertices absent, contrary to isometry.
Iterating proves the identity for the whole minor sequence. The final
overlap is thus a root-preserving image of $A\cap C$, and peels to
zero by endpoint minor closure.
\end{proof}

The preceding theorem supplies such a counterexample directly. This
minimization reduction remains useful for other candidate inputs. It reuses
rooted minor minimization, as in
Proposition~\ref{prop:damp-adjacent-exchange-minor-reduction}, with a
different preserved property: the canonical positive overlap and both
marked root directions survive. It reduces the certificate burden of
a counterexample search, not the intrinsic selection problem for the
generator. The ten assemblies in
Remark~\ref{rem:damp-positive-overlap-square-wings} fail before this
reduction can apply: each violates ordinary positivity or root failure.

\begin{remark}[A positive-overlap test using nonrelative square wings]
\label{rem:damp-positive-overlap-square-wings}
The relative hypotheses in
Theorem~\ref{thm:damp-noncontained-relative-square} cannot be recovered
merely by citing positivity of the sections. An existing example in
\path{research/relative_wing_square.tex}, Section~6, makes this distinction
explicit. Let $B$ be the normalized $289$-word core and put
\[
 P=B\cup\{4096,4098\},\qquad Q=\{0,2,4096,4098\}.
\]
Then $P$ is ample and peels to both $2$ and $4098$, but has no corner
outside the square $Q$, so it cannot peel relatively to $Q$.
The first endpoint order deletes $4096,4098$ and then peels $B$ to $2$;
the second deletes $4096$, peels $B$ to $2$, and finishes at $4098$.
These orders and the absence of outside corners are recorded in
\path{research/root_cube_retention.py/json}.

Use two copies of $P$ as $A,D$, identifying their squares with a common
$C=Q_2$ and giving their eleven remaining coordinates disjoint supports.
In each copy send either $2$ or $4098$ to zero and choose either ordering
of the two square directions. There are four such pointed placements,
and ten unordered pairs of placements. Their degree-two presentations
have $587$ words and are ample: the union is an ample coordinate-face
amalgam and the two shadows intersect exactly in the square shadow,
so Proposition~\ref{prop:damp-degree-two-section-shadows} applies.
Both overlaps equal $Q_2$ and peel to zero; both sides peel to zero,
but neither side peels relatively to its overlap.

This construction nevertheless supplies no rooted factor. Explicit
orders show that three of the ten presentations peel to their singleton
root; complete corner-deletion graphs show that the other seven are
not ordinarily peelable. The records in
\path{research/positive_overlap_recovery_probe.py/json} replay the wing
orders, all three root orders, and both the ordinary and root-avoiding
negative graphs for the remaining presentations. This finite test does
not prove that a rooted factor must have two root-negative overlaps.
It shows why these nonrelative wings do not refute that assertion:
the root-negative assemblies already fail ordinary positivity.
\end{remark}

\begin{remark}[A positive overlap without rooted minimality]
\label{rem:damp-positive-overlap-nonminimal}
The same $291$-word family $P=B\cup\{4096,4098\}$ from
Remark~\ref{rem:damp-positive-overlap-square-wings}, now rooted at
$a=4096$, refutes a stronger shortcut: ordinary positivity, positive
side endpoints, and a positive overlap do not by themselves imply
peeling to the degree-two root.

Indeed, $P$ peels ordinarily by either displayed order above. The
root $4096$ is a degree-two corner with incident directions $1,12$.
In its canonical presentation, one side is the singleton projected
from $4098$, the other is the root-containing coordinate-$1$ face of
$B$, and their opposite section is the other coordinate-$1$ face of
$B$. Both side endpoints are attainable: the singleton trivially,
and the other by rooted minor closure for $(B,0)$. The overlap of
the singleton side and the opposite section is itself a singleton.
Nevertheless $P$ cannot peel to $4096$. Its only initial corners are
$4096,4098$; deleting the latter leaves $4096$ as the only corner.
Equivalently, the first deletion in the untouched copy of $B$ would
have to be its unique corner zero, but its upper mate $4096$ and
its old neighbour $4$ would require the missing word $4100$ in the
corner cube.

This is not a rooted factor: contraction of coordinate $12$ gives
$(B,0)$, which is already root-negative. The other marked contraction
is positive, as it gives the positive contraction of $B$ in coordinate
$1$ with the retained root attached as a leaf. The replay
\path{research/positive_overlap_nonminimal_control.py/json} checks the
ordinary and side endpoint orders, the two-state root-negative graph,
and both marked contractions. This reuses the peripheral family; its
additional role is to show why the contraction hypotheses in
Proposition~\ref{prop:damp-positive-overlap-minimization} cannot simply
be discarded. Theorem~\ref{thm:damp-positive-overlap-rooted-factor}
supplies the eligible rooted-factor example, with a three-word overlap
instead of this singleton.
\end{remark}

\begin{proposition}[An obstruction to filling the root section]
\label{prop:damp-root-section-filling-obstruction}
Let $K,A,D,C\subseteq Q_F$ contain zero, and suppose that
\[
 H=(K\times\{00\})\cup(A\times\{01\})
   \cup(D\times\{10\})\cup(C\times\{11\})
\]
is ample. Put $R=A\cap C$ and $S=D\cap C$.
Suppose that each $T\in\{R,S\}$ has a forced corner-deletion prefix
avoiding zero and leading to a family $T_*$ whose unique corner is zero.
Precisely, there are distinct nonzero words $x_1,\ldots,x_m$ such that
$x_i$ is the unique corner of $T\setminus\{x_1,\ldots,x_{i-1}\}$,
and $T_*=T\setminus\{x_1,\ldots,x_m\}$ has unique corner zero;
the prefix may be empty. Define the incident-coordinate sets
\[
 I_T=\{f\in F:e_f\in T_*\},\qquad
 I_K=\{f\in F:e_f\in K\},
\]
where $e_f$ is the word supported on the single coordinate $f$.
Then
\begin{equation}
 H\text{ peels to }(0,00)
 \quad\Longrightarrow\quad I_R\subseteq I_K
                         \ \text{or}\ I_S\subseteq I_K.
 \label{eq:damp-root-section-filling-obstruction}
\end{equation}
No relative peeling between the sections and their overlaps is assumed.

In particular, take an ample family $Y$ in the degree-two presentation
of Proposition~\ref{prop:damp-degree-two-section-shadows}. If its two
overlaps have such prefixes with $|I_R|,|I_S|\ge2$, no ample one-concept
addition adjacent to its singleton root can make peeling to that root
possible. This includes every adjacent rank-three corner addition.
\end{proposition}
\begin{proof}
Suppose an order ends at $(0,00)$, and consider its first deletion of
one of $(0,01),(0,10),(0,11)$. All four zero copies are present just
before that deletion. If the deleted copy is in $01$ or $11$, use the
coordinate face with sections $A,C$, whose reduction is $R$; if it is
in $10$, use the face with sections $D,C$, whose reduction is $S$.
Write $T$ for the selected reduction.

Record the first destruction of each full edge fibre in that face.
These events form a corner-deletion prefix of $T$: the residual face
is ample at every step, its reduction is ample, and a single deletion
changes the reduction by at most one word. Whenever a word is lost,
the resulting one-word ample deletion makes it a corner. This is the
first-fibre argument used in Proposition~\ref{prop:damp-paired-root-transfer}.
The fibre at zero is not destroyed before the selected deletion, and
is destroyed at that deletion because its other endpoint still survives.
Thus the recorded prefix avoids zero and must next delete zero.
The unique-corner hypotheses force this prefix to be exactly
$x_1,\ldots,x_m$: it cannot stop earlier, when zero is not a corner,
or continue further without deleting zero. The reduction immediately
before the selected deletion is therefore $T_*$.

For each $f\in I_T$, the full fibre at $e_f$ is still present.
Consequently the deleted zero copy has an old-coordinate neighbour
in direction $f$. It also has both of its new-coordinate neighbours,
since the zero square is intact. Its corner cube must contain every
combination of these three directions, in particular $(e_f,00)$.
This word therefore belongs to the original family $H$, so $f\in I_K$.
This proves \eqref{eq:damp-root-section-filling-obstruction} in both cases.

For the last assertion, a new word adjacent to the singleton root of
$Y$ cannot use either of its two already present new-coordinate
neighbours. It must be $(e_f,00)$ for an old coordinate $f$.
The enlarged root section is $K=\{0,e_f\}$, with $|I_K|=1$.
Neither of the two sets of size at least two can be contained in $I_K$,
so the enlarged family cannot peel to the root.
\end{proof}

For the $870$ factor and the two opposite-only $871$ tip extensions,
both overlaps are the $289$-word core, whose unique corner is zero
of degree three. For the $872$ factor one overlap is that core and
the other first deletes the forced tip $64$ to reach it. Hence the
proposition excludes adjacent endpoint repairs in all four examples.
The checker \path{research/root_section_fillings.py/json} also records
complete root-avoiding deletion graphs for their adjacent rank-three
corner additions. This is an application of the forced-prefix
obstruction, not a proof of the unrestricted adjacent-corner property
in Corollary~\ref{cor:damp-adjacent-corner-minimal-repair-test}.
Root failure alone does not supply the forced prefixes: after some
nonroot deletions a reduction can have zero as a corner while retaining
other vertices. The singleton-section proof cannot then conclude that
the enlarged root section supplies an endpoint order of the overlap.

\begin{remark}[Recovering the paired extensions]
\label{rem:damp-paired-root-recovery}
Let \(Z\) be an ample family with a corner root of degree two.
Translate the root to \(0\), and let \(q,r\) be its incident
coordinates.  The \(q=r=0\) section is \(\{0\}\): the root is
isolated in that section, which must be connected.  The other three
sections recover \(A,D,C\) in positions \(01,10,11\), respectively.
Suppose that \(C\subseteq A,D\), that both side differences have
size at most one, and that \(C\) uses every remaining coordinate.
These are intrinsic tests on the recovered sections.

Isometry forces \(A\cap D=C\), by the missing-intermediate-word
argument in the proof.  All three sections and their union are ample.
The exact shadow formula, with
\(\operatorname{Sh}(A)\cap\operatorname{Sh}(D)\) in the two
one-new-coordinate terms, and \(|\operatorname{Sh}(Z)|=|Z|\),
forces that intersection to equal \(\operatorname{Sh}(C)\).
Thus Proposition~\ref{prop:damp-paired-root-transfer} applies and
recovers all parameters uniquely, up to interchanging \(q,r\).
Within this type, rooted minimality is exactly rooted minimality of
\(C\) and the two repairing endpoint conditions.  This reduces the
admissibility problem to smaller inputs; it does not replace those
remaining recognition conditions by intrinsic ones, or assert that
every degree-two-root factor has this section pattern.

The root of the recovered core need not be a corner.  Take the
\(292\)-concept family \(H\), rooted at \(29\), from
Remark~\ref{rem:damp-noncorner-rooted-factor}, and translate its root
to zero to obtain \(C\).  The extensions by \(128\) and \(290\)
both peel to zero and satisfy the shadow hypotheses.  The resulting
\(Y\) is a rooted factor with \(879\) concepts, dimension \(14\),
VC dimension five, and eleven corners.  The replay
\path{verify_damp_paired_root_transfer.py} checks all \(24\) rooted
minor orders of \(C\), all \(28\) of \(Y\), and its complete
\(704\)-state negative certificate.  Four positive controls over
\(C=\{0,1,3,7\}\), with optional tips \(2,5\), test every one of
\(321\) reachable states before the first root-square deletion.
All \(231\) possible exits from these states induce a core peeling
ending at zero.  The data and the recovered corner sets are recorded in
\path{damp_paired_root_transfer_verification.json}.
\end{remark}

\begin{theorem}[Two endpoint repairs give a terminal cut-vertex obstruction]
\label{thm:damp-terminal-cut-vertex}
Let \((C,0)\), irredundantly embedded in \(Q_F\), satisfy the three
rooted-factor conditions of Theorem~\ref{thm:damp-cut-vertex-obstructions},
and suppose that \(0\) is its only corner.  Let
\(A=C\cup\{\alpha\}\), \(D=C\cup\{\delta\}\), where
\(\alpha,\delta\notin C\) are distinct.  Assume that \(A,D,A\cup D\)
are ample, that both \(A,D\) peel to \(0\), and that their shadows
intersect in \(\operatorname{Sh}(C)\).
Then the square construction \eqref{eq:damp-paired-rooted-factor}
is a rooted factor with corners exactly
\[
 a=(0,00),\qquad \alpha^*=(\alpha,01),\qquad\delta^*=(\delta,10).
\]
Neither nonroot corner deletion preserves rooted minimality.
Gluing \((Y,a)\) to a copy of \((C,0)\) on disjoint coordinates
gives a corner-descent-terminal ample forbidden pc-minor \(X=Y\vee C\),
with a cut vertex, exactly two corners, and order \(4|C|+2\).
It has no nontrivial replicated common-seed presentation with resolving tips.
\end{theorem}

\begin{proof}
Proposition~\ref{prop:damp-paired-root-transfer} gives the rooted-factor
conditions and, for every \(R\subseteq\{\alpha^*,\delta^*\}\),
\begin{equation}
 \operatorname{Cor}(Y\setminus R)
 =\{a\}\cup\bigl(\{\alpha^*,\delta^*\}\setminus R\bigr).
 \label{eq:damp-paired-rooted-corners}
\end{equation}
After deleting \(\alpha^*\), the same transfer criterion fails because
its \(01\)-section is now \(C\), which cannot peel to its root.
Explicitly, restrict to \(q=0\) and contract \(r\) to expose the
proper bad rooted minor \((C,0)\).  After deleting \(\delta^*\),
restrict to \(r=0\) and contract \(q\) to get the same witness.

The cut-vertex theorem makes \(X=Y\vee C\) ample and forbidden.
The root ceases to be a corner, so its corners are exactly the two
tips.  The preceding restriction and contraction, with the second
block retained, give the proper minor \(C\vee C\) of either corner
deletion.  This minor is nonpeelable by
Proposition~\ref{prop:damp-vertex-gluing}, proving terminality.
Every nontrivial replicated common-seed construction with resolving
tips is two-connected by
Proposition~\ref{prop:damp-replicated-two-connected}; the cut vertex
of \(X\) therefore excludes such a presentation.
\end{proof}


\begin{remark}[A fixed terminal example with 1158 concepts]
\label{ex:damp-terminal-cut-vertex-1158}
Let \(B=(G\cup\{91\})\setminus\{61,125,93\}\) be the
\(289\)-concept rooted factor from
Remark~\ref{rem:damp-noncorner-rooted-factor}.  Translate its root
\(29\) to zero, obtaining \(C=\{x\mathbin{\mathrm{xor}}29:x\in B\}\).
Use the two tips \(\alpha=128\) and \(\delta=290\), corresponding
to the original words \(157\) and \(319\).  Their extensions peel
to zero and acquire the distinct supports
\[
 \{1,2,7,8\}\quad\text{and}\quad\{5,6,7,8\},
 \qquad\text{with bitmasks }390,480.
\]
Their union is ample and its shadow is \(\operatorname{Sh}(C)\)
together with these two supports.  Thus the theorem applies.
The resulting \(Y\) has \(870\) concepts in dimension \(14\).
With new bits \(q,r\) in positions \(0,1\), its nonroot corners
have words \(514,1161\).  Gluing a second \(C\) in bit positions
\(14,\ldots,25\) produces \(X\) of order \(1158\), isometric
dimension \(26\), and VC dimension four.  Both corners have degree
four; the root has degree five.  The minimum degree is four, so \(X\)
is already its graph three-core.  Deleting the cut vertex leaves
components of orders \(869\) and \(288\).  Each corner deletion has
the cornerless \(577\)-concept double \(C\vee C\) as a proper minor.

The independent replay
\path{verify_damp_terminal_cut_vertex.py} checks the input shadows,
both repairing endpoint orders, all \(24\) rooted elementary orders
of \(C\), all \(28\) of \(Y\), and all \(78\) ordinary elementary
orders of \(X\).  It verifies the complete four-state negative
certificates for both \(Y\) with its root retained and \(X\), as
well as both terminality witnesses.  The report is
\path{damp_terminal_cut_vertex_verification.json}; the new endpoint
orders are stored in
\path{computations/round790_paired_endpoint_repairs.json}.

This example refutes universal preserving nonroot-corner deletion and
universal existence of a common-seed layer set for cornered terminal
obstructions.  Proposition~\ref{prop:damp-terminal-cut-seed-towers} below uses its
two resolving extensions to refute cornerlessness of the largest recovered
seed as well.  Intrinsic admissibility of the rooted factors and their
endpoint-repairing extensions remains open.
\end{remark}

\begin{lemma}[Rooted minimality along ample additions]
\label{lem:damp-rooted-ample-additions}
Let \((H,a)\) be a rooted factor, irredundantly embedded in \(Q_F\),
and let \(K=H\cup\{c\}\subseteq Q_F\) be ample, with \(c\notin H\).
Then either \(K\) peels to \(a\), or \((K,a)\) is a rooted factor.
More generally, suppose
\[
 H=H_0\subset H_1\subset\cdots\subset H_m\subseteq Q_F,
 \qquad |H_{i+1}\setminus H_i|=1,
\]
and every \(H_i\) is ample.  If \(H_m\) does not peel to \(a\),
then every \((H_i,a)\) is a rooted factor.
\end{lemma}

\begin{proof}
The new concept \(c\) is a corner of \(K\) by
Lemma~\ref{lem:damp-one-concept-extension-test}.  Deleting it and
following an ordinary peeling of \(H\) proves \(K\in\Damp\).
Let \(\mu\) be an elementary root-preserving restriction or contraction.
Every coordinate is active in \(H\), so \(\mu H\) is a proper
rooted minor and peels to \(\mu a\).  Both \(\mu H\) and
\(\mu K\) are ample; their difference has size at most one.
If the difference is nonempty, its concept is a corner of \(\mu K\)
by the same lemma and is distinct from \(\mu a\).  Delete it first,
then peel \(\mu H\) to \(\mu a\).  If the difference is empty,
use that peeling directly.  Thus all elementary rooted minors of
\(K\) peel to their roots.  Closure under further root-preserving
pc-minors gives the third rooted-factor condition of
Theorem~\ref{thm:damp-cut-vertex-obstructions}.  If \(K\) does not
peel to \(a\), all three conditions hold.

Peeling to \(a\) persists under an ample one-concept addition:
delete the new corner first.  Hence no \(H_i\) in the displayed
chain can peel to \(a\) when \(H_m\) cannot.  Apply the first
assertion successively.  The fixed coordinate set is essential to
this argument: it makes every elementary image of the initial factor
a proper minor.
\end{proof}

\begin{proposition}[A smaller rooted input and a second root]
\label{prop:damp-forced-prefix-rooted-inputs}
Let $C\subseteq Q_{12}$ be the $289$-concept rooted factor, normalized
at zero, used in Proposition~\ref{prop:damp-edge-boundary-transfer}.
In its integer encoding put $b=256$ and $H=C\setminus\{0\}$.
Then $|H|=288$ and all three pairs
\[
 (H,b),\qquad(C,0),\qquad(C,b)
\]
are rooted factors. The first two have corner roots of degree three;
the third has a noncorner root of degree four. The concepts attainable
as the final vertex of a peeling of $C$ are exactly $C\setminus\{0,b\}$.
These same $287$ concepts are the attainable final vertices of $H$.

Pairwise gluing of these three rooted inputs, allowing repetition,
gives six pairwise nonisomorphic ample forbidden pc-minors, all of
isometric dimension $24$. Their orders and numbers of corners are
\[
\begin{array}{c|r|r}
\text{rooted factors}&\text{order}&\text{corners}\\ \hline
(H,b),(H,b)&575&0\\
(H,b),(C,0)&576&0\\
(H,b),(C,b)&576&1\\
(C,0),(C,0)&577&0\\
(C,0),(C,b)&577&1\\
(C,b),(C,b)&577&2
\end{array}
\]
\end{proposition}
\begin{proof}
Direct corner checks give
\[
 \operatorname{Cor}(C)=\{0\},\qquad
 \operatorname{Cor}(H)=\{b\}.
\]
Thus every peeling of $C$ begins with $0,b$. The previously certified
ordinary order supplies an ordinary peeling of $H$ after its first
vertex is removed. Corner deletion preserves ampleness, and every
coordinate remains active in $H$.

The new rooted-minimality assertion is certified explicitly, rather than
inferred from the forced prefix. The replay
\path{research/forced_prefix_rooted_inputs.py} reconstructs $C$ from
\path{research/relative_wing_square.json}. For $H$ it checks an ordinary
peeling, its unique corner $b$, and complete orders ending at the root
image for all $24$ elementary root-preserving restrictions and contractions.
Rooted minor closure then covers every proper root-preserving pc-minor.
Since $b$ is the only initial corner and $|H|>1$, $H$ cannot peel to $b$.
Consequently $(H,b)$ is a rooted factor.

The pair $(C,0)$ is the established input. Since $C$ is an ample
one-concept extension of $H$, Lemma~\ref{lem:damp-rooted-ample-additions}
applies at the old root $b$. The forced prefix excludes a peeling of $C$
ending at $b$, so $(C,b)$ is also a rooted factor. The replay separately
checks its $24$ rooted elementary orders using the constructive proof
of that lemma. It also checks the three stated root degrees.

For each $a\in C\setminus\{0,b\}$ the accompanying JSON record supplies
a complete peeling of $C$ ending at $a$. All $287$ orders are replayed.
Their forced initial $0$ may be removed to obtain orders of $H$ with the
same endpoints. The unique corners exclude the remaining endpoints in
$C$ and $H$, proving both assertions about attainable final vertices.

Apply Theorem~\ref{thm:damp-cut-vertex-obstructions} to each unordered
pair of rooted inputs. It proves forbiddenness and dimension $24$.
The orders follow by identifying the two roots. A nonroot corner is
unaffected by vertex gluing, while the common root is not a corner,
since incident directions from different blocks lack their mixed square.
The rooted inputs have respectively zero, zero, and one nonroot corner,
which gives the displayed corner counts. The six order/corner-count
pairs are distinct, proving pairwise nonisomorphism. As an independent
finite check, the replay constructs and verifies complete orders for
all $72$ elementary pc-minors of each output, $432$ orders in total.
\end{proof}

The $577$-concept double of $(C,0)$ is reused here. The additional input
comes from deleting its first forced vertex and changing the distinguished
root, followed by a new rooted-minor certification. The proposition does
not assert that this operation preserves rooted minimality in general.
It therefore expands the collection of certified inputs and outputs,
without providing an exhaustive intrinsic rooted-factor generator.

\begin{proposition}[Rooted inputs of order $264$ and a forbidden double]
\label{prop:damp-rooted264-double527}
Let $X\subseteq Q_{12}$ be the recorded $267$-vertex class, in the integer
encoding of \path{evidence/order267_obstruction/record.json}, and put
\[
 A=\bigl(X\setminus\{236,492,1004,748\}\bigr)\cup\{220\},
 \qquad a=750.
\]
Then $(A,a)$ is a rooted factor of order $264$. Its only corner is $a$,
and all twelve coordinates are active. Every vertex of every nonempty
proper pc-minor of $A$ is attainable as a peeling endpoint. In particular,
no proper pc-minor of $A$ is a rooted factor at any choice of root. Two copies on disjoint coordinate
blocks, identified at $a$, give a cornerless ample forbidden pc-minor of
order $527$ and isometric dimension $24$.
More generally, the certified shrinking outputs supply $24$, $12$, and
$12$ distinct labelled rooted inputs of orders $266$, $265$, and $264$,
respectively. Any two may be glued at their roots to obtain a cornerless
ample forbidden pc-minor. No pairwise nonisomorphism assertion is made.
\end{proposition}
\begin{proof}
Start with
$H=(X\setminus\{236,492\})\cup\{220\}$.
The certificates in
\path{../ample-corners/evidence/shrinking_rooted_inputs/}
verify that $H$ is ample and has a full corner peeling, and that its
first three vertices are forced to be $1004,748,750$, in that order.
After deleting the first two, the remainder is $A$, whose unique corner
is $750$. Thus $A$ is positive, but cannot peel to $a$: every peeling
must delete $a$ first and $|A|>1$.

Rooted minimality is checked separately, not inferred from these forced
deletions. For each of the twelve coordinates the certificate gives
complete peelings ending at the root image for the contraction and the
section containing the root. These are all $24$ elementary
root-preserving pc-minors. Rooted minor closure therefore gives the
same assertion for every proper root-preserving pc-minor.
This proves the three hypotheses of
Theorem~\ref{thm:damp-cut-vertex-obstructions} for $(A,a)$.
For this particular $A$, the stronger endpoint assertion is independently
certified in \path{../ample-corners/evidence/root264_minor_endpoints/}:
all $5637$ endpoint choices across its $36$ elementary proper pc-minors
have explicit orders. The verifier reconstructs the minor families and
checks suffix cubes, distinct out-maps, and nonclashing. Any further
nonempty minor and any chosen vertex in it have a preimage vertex in one
of these elementary minors; rooted minor closure transports a peeling
ending at that preimage to the desired endpoint. Thus every proper minor
admits every endpoint, excluding any smaller rooted factor even after
changing the root.
The same replay checks all $48$ labelled inputs: ordinary peelability,
a unique corner root, twelve active coordinates, and all $1152$
elementary rooted-minor orders. These inputs are the distinct labelled
remainders along the forced initial deletions of the $24$ ample
$266$-vertex two-deletion/one-insertion outputs from $X$.

The cut-vertex theorem proves forbiddenness of every pairwise gluing.
Its order is the sum of the input orders minus one. Each nonroot vertex
retains its incident cubes and is not a corner. At the common root,
nonempty incident directions from both blocks have no mixed square,
so the root is not a corner either. All coordinates remain active.
For the $527$-vertex double, the independent verifier additionally
constructs and checks full ordinary corner orders for all $72$
elementary pc-minors; the explicit class and orders are in
\path{../ample-corners/evidence/shrinking_rooted_inputs/double527.json}.
\end{proof}

The smaller positive classes obtained by signed coordinate identifications
of $A$ need not provide further separator inputs. In particular, retaining
$x_1=x_{10}$ and suppressing coordinate $10$, or retaining $x_3=x_8$ and
suppressing coordinate $8$, gives ample classes of orders $107$ and $108$.
Each peels to every contained cube. The independent replay in
\path{../ample-corners/evidence/root264_identifications/cube_verification.json}
covers all $571$ and $579$ cubes, respectively, including vertices and edges.
Theorem~\ref{thm:damp-cube-separator-bound} therefore excludes either class
as a piece of an ample forbidden pc-minor cut along any separating cube.
This is a fixed-input exclusion using the existing recovery theorem,
not an intrinsic characterization of all admissible pieces.

The constructed double reduces the available cut-vertex upper bound from $575$ to $527$.
It does not improve the unrestricted upper bound $267$ or generate all
rooted factors. As in Proposition~\ref{prop:damp-forced-prefix-rooted-inputs},
changing the root after a forced deletion has been certified for the
specified inputs only.

\begin{proposition}[Attainable sinks need not form an ample family]
\label{prop:damp-disconnected-attainable-sinks}
For a corner-peelable ample family $A$, put
\[
 \operatorname{Sink}(A)=\{a\in A:A\text{ has a corner peeling ending at }a\}.
\]
There is a $658$-concept corner-peelable ample family $H$ of isometric
dimension $25$ such that the graph induced by $\operatorname{Sink}(H)$
is disconnected. In particular, $\operatorname{Sink}(H)$ is not ample.
For every positive integer $m$, a corner-peelable ample family exists
whose attainable-sink graph has at least $2^m$ connected components.
\end{proposition}
\begin{proof}
Take two copies $K_1,K_2$ of the $330$-concept piece from
Theorem~\ref{thm:damp-three-piece-witness}, sharing the edge
$D=\{u,v\}$ and with disjoint private coordinates. Reverse the
endpoints in the second copy. Their exact profiles are
\[
 (r_D,t_u,t_v,\ell_u,\ell_v)(K_1)=(0,1,0,1,0),\qquad
 (r_D,t_u,t_v,\ell_u,\ell_v)(K_2)=(0,0,1,0,1).
\]
Their union $H$ is ample by cube-face gluing, has $330+330-2=658$
concepts, and uses $12+12+1=25$ coordinates. Both compatibility
witnesses $t_u(K_1)\ell_v(K_2)$ and $t_v(K_2)\ell_u(K_1)$ hold,
so $H$ is corner-peelable.

The first witness constructs a peeling ending in $K_2\setminus D$.
Choose a peeling of $K_1$ ending at $u$ and a peeling of $K_2$ in
which $v$ is the first boundary vertex removed and is then a leaf
adjacent to $u$. Both orders restrict to $v,u$ on $D$. Merge their
outside blocks before, between, and after these two boundary vertices,
as in Corollary~\ref{cor:damp-finite-boundary-gluing}. The first
order has no vertex after $u$. The second must have a vertex after
$u$: otherwise it would end at $u$, contradicting $t_u(K_2)=0$.
Thus the merged order ends in $K_2\setminus D$.
Using the other compatibility witness interchanges the roles and gives
an attainable sink in $K_1\setminus D$.

Corollary~\ref{cor:damp-edge-context-algebra} shows that a positive
union of two pieces failing retention has all four boundary values zero.
Hence neither $u$ nor $v$ belongs to $\operatorname{Sink}(H)$.
There are no edges between $K_1\setminus D$ and $K_2\setminus D$.
The attainable-sink graph meets both sets and is therefore disconnected.
Every nonempty ample family has connected one-inclusion graph, proving
that this sink set is not ample. The replay
\path{research/disconnected_attainable_sinks.py/json} constructs and
checks the two complete orders; it does not enumerate all attainable sinks.

Finally, for positive ample families $A,B$ one has
\[
 \operatorname{Sink}(A\times B)
   =\operatorname{Sink}(A)\times\operatorname{Sink}(B).
\]
Necessity follows by restricting a root-ending order to the coordinate
sections through its final vertex. For sufficiency, take root-ending
orders in both factors and delete the $A$-fibres in the $B$ order,
using the $A$ order within each fibre. The incident cube at each
step is the product of the two relevant corner cubes; all later
fibres remain complete. This ends at the prescribed pair of roots.
The one-inclusion graph on the displayed sink set is the Cartesian
product of the two sink graphs. Its connected components are exactly
the products of their components, since paths project to paths and
paths within factors lift with the other coordinate fixed. Applying
this to $H^m$ gives at least $2^m$ components.
\end{proof}

The ample attainable-sink set in
Proposition~\ref{prop:damp-forced-prefix-rooted-inputs} is therefore a
property of that input, not a general recovery principle. The counterexample
uses two existing edge pieces and the existing gluing criterion; it supplies
no new forbidden-minor family. It excludes treating the attainable sinks of arbitrary positive ample
families as one ample residual class, or as a union with a universally
bounded number of connected components. The next proposition shows that rooted minimality does not restore
ampleness of the sink set, although it makes no claim about disconnectedness.

\begin{proposition}[Rooted minimality does not make the sink set ample]
\label{prop:damp-rooted-nonample-sink-set}
There is an ample rooted factor $(S,0)$ with $870$ concepts and isometric
dimension $14$ such that $\operatorname{Sink}(S)$ is not isometric and
therefore not ample. More precisely, in the integer encoding below,
$1,2\in\operatorname{Sink}(S)$ but $0,3\notin\operatorname{Sink}(S)$.
\end{proposition}
\begin{proof}
Use the established $870$-concept rooted factor from
Remark~\ref{ex:damp-terminal-cut-vertex-1158}. In its stored coordinates,
let $C$ be the normalized $289$-concept core, put
$A=C\cup\{128\}$ and $B=C\cup\{290\}$, and write
\[
 S=\{0\}\cup\{4x+2:x\in A\}
       \cup\{4x+1:x\in B\}\cup\{4x+3:x\in C\}.
\]
The pair $(S,0)$ is a rooted factor by the earlier construction.
In particular, $0$ cannot be an attainable sink. The coordinate face
with both low bits equal to one is the copy $\{4x+3:x\in C\}$.
An order ending at $3$ would restrict to an order of $C$ ending at zero,
which is impossible because zero is the only initial corner of $C$.
Thus $3$ is also unattainable.

Let $c=(0,c_2,\ldots,c_{289})$ be the stored ordinary order of $C$,
and let $a=(a_1,\ldots,a_{290})$ and $b=(b_1,\ldots,b_{290})$
be the stored orders of $A$ and $B$ ending at zero. The following
concatenated lists are complete corner orders of $S$:
\begin{align*}
 &(0,1,3,1161),\quad
 (4c_i+1)_{i=2}^{289},\quad
 (4c_i+3)_{i=2}^{289},\quad
 (4a_i+2)_{i=1}^{290};\\
 &(0,514,2,3),\quad
 (4c_i+2)_{i=2}^{289},\quad
 (4c_i+3)_{i=2}^{289},\quad
 (4b_i+1)_{i=1}^{290}.
\end{align*}
They end at $2$ and $1$, respectively. The prefixes remove the root,
the other two copies of core zero, and the unused repair. The next two
blocks clear copies of $C\setminus\{0\}$ while their corner cubes
have mates in an untouched adjacent layer. The last block uses the
repair kept in the final shore, which makes zero attainable there.

For completeness, the finite prefix assertions and every subsequent
corner deletion are replayed by
\path{research/rooted_nonample_sink_set.py}. It reconstructs these lists
from the earlier core and repair orders, checks the complete four-state
negative deletion graph protecting zero, and replays all $28$ rooted
elementary-minor orders. Hence the input's rooted minimality and the two
new endpoint assertions are checked independently of any greedy search.
The exact lists and source hashes are in the accompanying JSON record.

The two vertices $1$ and $2$ differ in their two low coordinates, so
their only common neighbours in the ambient cube are $0$ and $3$.
Both are absent from $\operatorname{Sink}(S)$. Its induced graph
therefore has no path of length two between two vertices at Hamming
distance two. It is not isometric and, since nonempty ample families
are isometric, is not ample.
\end{proof}

This strengthens the failure of the ample-sink-set principle to the actual
rooted-factor domain. It does not assert that the sink set of this factor
is disconnected, or classify all its attainable sinks. The two previously
unsuccessful greedy endpoint tests are resolved by retaining the repair
needed in the final shore; a greedy stall was not a negative certificate.

\begin{proposition}[Corner deletion followed by rooted minor descent]
\label{prop:damp-combined-rooted-descent}
Let $(H,a)$ be a rooted factor on its active coordinate set $F$.
Suppose $x\ne a$ is a corner and
$S=H\setminus\{x\}$ is ordinarily peelable. There is a root-preserving
pc-minor map $\mu$ for which $(C,b)=(\mu S,\mu a)$ is a rooted factor
with $|C|<|H|$.

If $(S,a)$ is already a rooted factor, take $\mu$ to be the identity.
Otherwise $\mu$ can be chosen proper, and
\[
 \mu H=C\cup\{d\},\qquad d=\mu x\notin C,
\]
is an ample one-concept addition that peels to $b$. Write $P$ for the
incident support of $x$ in $H$. In a normal form for $\mu$, let $R$ be
the restricted coordinates and $Q$ the contracted coordinates. Then
\[
 Q\cap P=\varnothing,\qquad x_i=a_i\ (i\in R),\qquad
 \text{the acquired support at }d\text{ is }P\setminus R,
 \qquad |P\setminus R|\ge3.
\]
Thus the repair retains a face of the original corner cube; it is not an
arbitrary new repairing addition.

Iterating these steps ends at a rooted factor $(K,r)$ in which every
complete peeling starts at $r$. Equivalently, $r$ is the unique source
in every acyclic USO of $K$.
\end{proposition}
\begin{proof}
The class $S$ is ample by corner deletion. It cannot peel to $a$, since
prefixing such an order by $x$ would peel $H$ to $a$. Choose a
root-preserving minor $(C,b)$ of $(S,a)$ minimal among those failing at
their roots. Ordinary pc-minor closure makes $C$ peelable; minimality
supplies every proper rooted minor condition. Hence it is a rooted factor,
and its size is at most $|S|<|H|$.

When $S$ is not already minimal, the chosen map uses at least one proper
coordinate operation. Applied to $H$ it gives a proper root-preserving
minor, which peels to $b$. Images of $H$ and $S$ differ by at most one
word. They cannot coincide, because $C$ fails at $b$. Their difference
is therefore exactly $d=\mu x$, which is distinct from $b\in C$.

Restrictions and contractions on distinct coordinates commute, so write
$\mu$ with disjoint sets $R,Q$ and with every restriction taking the
root's bit. A restriction discarding $x$ would identify the two images;
thus $x_i=a_i$ for every $i\in R$. Contracting a coordinate in $P$
would also identify them, since the neighbour $x^{\{i\}}$ lies in $S$.
Consequently $Q\cap P=\varnothing$.

First restrict in $R$. The corner cube at $x$ becomes its retained face
of support $P\setminus R$, and $x$ remains a corner. The shadow difference
between these two restricted ample families is exactly that support.
Now contract in $Q$. A contraction retains precisely the shattered
supports avoiding its coordinate; all members of $Q$ avoid $P\setminus R$.
The shadow difference therefore remains $\{P\setminus R\}$, and the
ample one-concept addition at $d$ has this acquired support. A support
of size at most two cannot change peeling to an old prescribed endpoint,
by Proposition~\ref{prop:damp-rooted-operations}. Since adding $d$ does
change it here, $|P\setminus R|\ge3$.

Whenever a peeling of the current rooted factor begins at a nonroot
vertex $x$, its suffix proves that deleting $x$ leaves an ordinarily
peelable class, so the preceding step applies. Size strictly decreases,
and iteration must stop. At a terminal factor no peeling begins away
from the root. CCMW's correspondence identifies peelings with topological
orders of acyclic USOs. A nonroot source in any such orientation can be
chosen first in a topological order; conversely the first word of a peeling
is a source. Thus the terminal property is exactly that the root is the
unique source in every acyclic USO.
\end{proof}

For a degree-three repair, the inequality in the proposition forbids
restrictions in its support, agreeing with the more specialized
Proposition~\ref{prop:damp-minimal-rank-three-endpoint-repair}.
The additional statement here is the combined descent for arbitrary
corner degree, including the retained-face description when restrictions
lower that degree. It does not assert that the terminal factor has only
one geometric corner: other corners could all leave nonpeelable classes
when deleted. Nor is this an intrinsic forward generator, since ordinary
peelability, rooted minimalization, and the inverse expansion data remain
to be selected structurally.

The two recent controls illustrate distinct steps. In the $292$-concept
factor, deleting $32$ leaves a rooted factor, as follows from its ample
addition chain above the $289$-concept factor and
Lemma~\ref{lem:damp-rooted-ample-additions}. For the $308$-concept factor,
delete the repair $(290,0)$ and contract the fresh coordinate. This recovers
the $289$-concept input; applying the same contraction before deletion
recovers its endpoint repair $C\cup\{290\}$, of support
$\{5,6,7,8\}$. Thus allowing the minor step handles a failure of direct
rooted minimality already present in the repaired-face construction.

\begin{corollary}[Joint descent for rooted factors and ordinary obstructions]
\label{cor:damp-joint-rooted-ordinary-descent}
Start with either a rooted factor $(H,a)$ or an ample forbidden pc-minor $H$.
The following procedure strictly decreases cardinality at every step and ends
at one of these two types:
\begin{enumerate}
 \item a rooted factor whose root is its only geometric corner;
 \item a cornerless ample forbidden pc-minor.
\end{enumerate}
At a rooted factor with a nonroot corner $x$, delete $x$. If the result is
ordinarily peelable, take a minimal root-preserving minor failing at its root,
as in Proposition~\ref{prop:damp-combined-rooted-descent}. If it is not
ordinarily peelable, forget the root and take an ordinary forbidden pc-minor
of the deletion. At an ordinary forbidden class with a corner, delete that
corner and again take an ordinary forbidden pc-minor of the deletion.
Once the procedure enters the ordinary branch it stays there.
\end{corollary}
\begin{proof}
At a rooted factor, deleting a nonroot corner preserves ampleness and failure
to peel to the root. If the deletion is ordinarily positive, the preceding
proposition gives a smaller rooted factor. Otherwise finite minor minimalization
gives an ordinary forbidden pc-minor; it is ample because ampleness is pc-minor
closed. At an ordinary nonpeelable ample family, a corner deletion remains
nonpeelable: any peeling of the remainder could be prefixed by that corner.
Thus ordinary minor minimalization is available at each later step as well.
Each step begins with a deletion and then takes a minor, so size strictly
decreases. In the rooted branch termination means that no nonroot corner is
present. Ordinary positivity supplies at least one corner, necessarily the
root. In the ordinary branch termination means cornerlessness. Both branches
retain their stated minimality conditions by construction.
\end{proof}

The ordinary part reuses the corner-deletion mechanism of
Theorem~\ref{thm:damp-prime-cornerless-seed-reduction}; taking a minimal
negative pc-minor after each deletion additionally keeps its states forbidden.
The purpose of the joint statement is to allow ordinary nonpeelability after
a rooted deletion, instead of treating it as a barrier. For example, the
$1158$-concept corner-descent-terminal class already constructed in
Theorem~\ref{thm:damp-terminal-cut-vertex} is not terminal for this procedure.
After either corner deletion, its recorded restriction and contraction give
the cornerless forbidden double $C\vee C$ of order $577$.
This is reuse of that theorem's witness, not a new example.

This is a backward reduction, not an intrinsic forward generator. It still
uses negative-minor selection; reversing a step requires admissible expansions
and corner additions that recover minimality. It proves neither that the
terminal input is unique nor that the two input types have been generated.
It does, however, remove the need to identify the source-based terminal
condition in Proposition~\ref{prop:damp-combined-rooted-descent} with a unique
geometric corner. That stronger source question can be considered separately.

\begin{corollary}[Lifting a repairing addition through one contraction]
\label{cor:damp-repair-lift-labels}
Let $\varnothing\ne D\subseteq C\subseteq Q_F$ be ample, let
$K_1,\ldots,K_k$ be the components of $C\setminus D$, and let
$d\notin C$ with $A=C\cup\{d\}$ ample. For a fresh coordinate $e$,
labels $t_i\in\{0,1\}$, and $b\in\{0,1\}$, put
\[
 S_t=(D\times\{0,1\})\cup\bigcup_i(K_i\times\{t_i\}),
 \qquad H_{t,b}=S_t\cup\{(d,b)\},
\]
and define
\[
 I(d)=\{i: d\text{ has a neighbour in }K_i\}.
\]
Then $H_{t,b}$ is ample if and only if $t_i=b$ for every $i\in I(d)$.
In that case $(d,b)$ is a corner, and the acquired support of its
addition is the acquired support of $d$ over $C$.

Suppose additionally that $A,D\in\Damp$. If an ample $H_{t,b}$ is
not in $\Damp$, then some $j\notin I(d)$ has $t_j=1-b$.
In particular, if $I(d)=\{1,\ldots,k\}$, no labeling gives an ample
forbidden pc-minor. This last conclusion requires positivity of $D$;
it does not discard the branch with a negative reduction.
\end{corollary}
\begin{proof}
The contraction of $H_{t,b}$ in $e$ is $A$ and its reduction is $D$.
In the graph induced by $A\setminus D$, adjoining $d$ merges precisely
the components $K_i$ with $i\in I(d)$ together with $d$; all other
components remain unchanged. Theorem~\ref{thm:damp-component-expansion-signs}
therefore says exactly that their labels must equal the label $b$ of $d$.
It also proves sufficiency. Both $S_t$ and $H_{t,b}$ are then ample,
so the one-concept extension criterion makes $(d,b)$ a corner.
There is no neighbour in direction $e$. Every old neighbour of $d$
is either in $D$, which has both lifts, or in a component indexed by
$I(d)$, which has label $b$. Thus its incident support is unchanged,
as asserted.

If there is no $j\notin I(d)$ of label $1-b$, all component labels
equal $b$. Consequently
\[
 H_{t,b}=(A\times\{b\})\cup(D\times\{1-b\}).
\]
This is a peripheral expansion of the positive base $A$ at the positive
site $D$, hence is positive by
Proposition~\ref{prop:damp-exact-expansion-product-tests}.
Taking the contrapositive proves the final assertions.
\end{proof}

This is an application of the existing component-label and peripheral
positivity theorems, rather than a new parametrization of ample classes.
It supplies a local constraint for reversing a contraction in the joint
descent: a repair must leave an untouched component in the opposite layer
if the repaired contraction and the reduction are positive. For example,
take $C=\{00,01,10\}$, $D=\{00\}$, and $d=11$. The two components
of $C\setminus D$ both meet $d$, so lifting the addition forces both
labels to agree and gives a positive peripheral expansion. Opposite
labels cannot even give an ample addition. This small example illustrates
the geometric constraint, not a forbidden-minor witness.
The condition is necessary, not sufficient, for forbiddenness: it supplies
neither failure of an acyclic USO nor positive USOs on the other proper
pc-minors. Those remain the substantive forward-admissibility problem.

\begin{remark}[Singleton reductions and a positive separated repair]
\label{rem:damp-singleton-repair-control}
The vertex-gluing criterion already settles the singleton-reduction case
of Corollary~\ref{cor:damp-repair-lift-labels}. More generally, any ample
expansion $X$ with contraction $A$ and reduction $\{a\}$ satisfies
\[
 X\in\Damp\quad\Longleftrightarrow\quad A\in\Damp.
\]
Indeed, the components of $A-a$, with $a$ restored, occupy disjoint
coordinate blocks, as proved in
Theorem~\ref{thm:damp-cut-vertex-obstructions}. In $X$, these pieces
attach to the two ends of the sole edge in the new direction, according
to their component labels. If $A$ peels, all its pieces peel and at most
one fails to peel to $a$, by Proposition~\ref{prop:damp-vertex-gluing}.
Peel every other piece to its attachment vertex. What remains is the
exceptional piece with a pendant edge, or just the edge if there is no
exceptional piece. Delete the pendant endpoint and peel the remaining
piece, or peel the edge. Conversely, positivity of $X$ implies positivity
of its contraction $A$. Thus no active singleton-reduction expansion is
an ample forbidden pc-minor.

This also gives a control inside the inverse-step problem, with a
negative starting contraction. Take two copies $(B_1,0),(B_2,0)$ of the
$289$-concept rooted factor in
Remark~\ref{ex:damp-terminal-cut-vertex-1158}, on disjoint coordinates.
Let $C=B_1\vee B_2$, let $D=\{0\}$, and add the repairing tip
$d=290$ in the first block. The family $C$ is the $577$-concept forbidden
double. The repaired family
$A=(B_1\cup\{d\})\vee B_2$ peels: peel the first block to zero,
then peel the second ordinarily. In the notation of
Corollary~\ref{cor:damp-repair-lift-labels}, give $B_1-0$ label zero,
$B_2-0$ label one, and lift $d$ at zero. These are precisely the two
components, since rooted factors have no cut vertex. The repair touches
only the first component, so the second is the required untouched
opposite-layer component. Nevertheless the $579$-concept output peels:
peel the repaired first block to its root, delete that root as a leaf
across the new edge, and peel the second block. All its proper pc-minors
therefore peel as well.

This reuses the known rooted factor and its certified endpoint repair;
it is not a new forbidden family. It shows why component separation alone
does not close forward admissibility, even with a forbidden input $C$ and
a positive repaired contraction. The following theorem strengthens this exclusion from singleton reductions
to all convex reductions and supplies a lift of the prescribed USO.
\end{remark}

\begin{theorem}[Lifting acyclic USOs through a convex reduction]
\label{thm:damp-convex-reduction-lift}
Let $X$ be ample, let $e$ be active, and write $A=\pi_eX$ and
$D=X^{\{e\}}$ for its contraction and reduction. If $D$ is convex
in $A$, every acyclic USO of $A$ lifts to an acyclic USO of $X$ whose
orientation on every edge in an old direction is the given orientation
on its projection. Consequently
\[
 X\in\Damp\quad\Longleftrightarrow\quad A\in\Damp.
\]
Moreover, for every $a\in X$,
\begin{equation}
 X\text{ peels to }a
 \quad\Longleftrightarrow\quad A\text{ peels to }\pi_ea.
 \label{eq:damp-convex-endpoint-transfer}
\end{equation}
In particular, neither an ample forbidden pc-minor nor a rooted factor
has a coordinate whose reduction is convex in its contraction. This
excludes every cubical reduction, in every dimension.
\end{theorem}
\begin{proof}
Fix an acyclic USO $\mathcal O$ of $A$. Write $K_i$ for the components
of $A\setminus D$ and $t_i$ for their lift labels, using
Theorem~\ref{thm:damp-component-expansion-signs}. Convexity in a partial
cube means that $D$ is obtained by fixing coordinates. Thus, for every
cube $Q\subseteq A$, the intersection $Q\cap D$ is a face, possibly
empty or all of $Q$. Its complement in $Q$, when nonempty, is connected
and hence lies in a single $K_i$.

Orient every old-direction edge of $X$ as its projection in $\mathcal O$.
At $d\in D$, the outgoing cube $Q_d$ of $\mathcal O$ is contained in
$A$ by (C1). If $d$ has any outgoing neighbour outside $D$, all such
neighbours belong to a single component $K_i$, by the preceding face
observation. Orient the vertical edge over $d$ towards layer $t_i$.
If there is no such neighbour, choose either direction independently.
We verify that this orientation is a USO.

For (C1), first take a vertex outside the doubled reduction. Its outgoing
cube in $A$ lifts wholly beside it: all vertices of that cube outside
$D$ belong to its own component, and all vertices in $D$ have both lifts.
Next take a lift of $d\in D$ with no outgoing vertical edge. Its outgoing
old directions are a subset of those of $Q_d$. If outgoing neighbours
outside $D$ occur in that layer, all of $Q_d\setminus D$ has that layer's
label, so the requisite cube lifts. If none occur, the outgoing directions
are along the face $Q_d\cap D$, which has both lifts. Finally, at the
source of the vertical edge, our choice ensures that no outgoing old
neighbour outside $D$ is present. Its old outgoing cube is a face contained
in $Q_d\cap D$. The product of this face with the vertical edge is
contained in $X$, proving (C1) there as well.

For (C2), a cube not using $e$ projects to a cube of $A$ with precisely
its given USO. A cube using $e$ is $Q\times\{0,1\}$ for a cube
$Q\subseteq D$. Let $z$ be the unique sink of $\mathcal O$ on $Q$.
Both lifts of every other vertex have an outgoing old edge within this
cube. Its only sink is therefore the lift of $z$ at the head of its
vertical edge. This proves (C2), regardless of the independent choices
of vertical directions.

A directed cycle containing an old-direction edge would project, after
omitting vertical steps, to a nonempty directed closed walk in the acyclic
orientation $\mathcal O$, which is impossible. A cycle using only vertical
edges is also impossible, since those edges form a matching. Thus the
lift is acyclic. CCMW's correspondence gives the forward positivity
implication; the reverse follows from closure under contraction.
If the given orientation has sink $\pi_ea$, this also gives the endpoint
assertion. When $\pi_ea\notin D$, its unique lift is $a$ and is a sink.
When $\pi_ea\in D$, it has no outgoing old neighbour, so its vertical
edge is free; direct it towards $a$. The lifted ample USO has $a$ as its
unique global sink, and is acyclic. The reverse endpoint implication is
root-preserving contraction closure. Thus
\eqref{eq:damp-convex-endpoint-transfer} follows from the existing lift,
without a separate lifting theorem.
For a forbidden $X$, the proper contraction $A$ is positive, so convexity
of $D$ would contradict the ordinary conclusion. For a rooted factor,
its proper contraction peels to the root image, contradicting the endpoint
conclusion. A contained cube is convex, proving the last assertion.
\end{proof}

\begin{proposition}[Retaining a face in one layer of a convex expansion]
\label{prop:damp-convex-relative-face}
Under the hypotheses of Theorem~\ref{thm:damp-convex-reduction-lift},
write $X_t=\{x:(x,t)\in X\}$ for its two projected sections.
Let $B$ be a nonempty coordinate-face intersection of $A$ with
$B\subseteq X_t$, and put $T=B\times\{t\}$ and $C=D\cap B$.
Then
\begin{equation}
 X\text{ peels to }T
 \quad\Longleftrightarrow\quad
 \bigl(A\text{ peels to }B\bigr)
 \text{ and }\bigl(C=\varnothing\text{ or }C\in\Damp\bigr).
 \label{eq:damp-convex-relative-face}
\end{equation}
In particular, if $B\in\Damp$, the second condition is automatic.
The target $B$ itself need not be corner-peelable in the general statement.
\end{proposition}
\begin{proof}
Suppose first that $X$ peels to $T$. Contracting $e$ in every residual
ample family gives a relative peeling of $A$ to $B$, after omitting
repeated images: each step removes at most one word, all residual images
are ample, and no word of $B$ is removed. Restrict the same residual
families to the coordinate face projecting to $B$ and having $e=1-t$.
Its initial family is $C\times\{1-t\}$ and its final family is empty.
Again each nontrivial step is a corner deletion, by the ample
corner-deletion criterion. Thus $C$ is empty or corner-peelable.

Conversely, lift a relative corner order from $A$ to $B$ by deleting
whole $e$-fibres successively. We justify this operation locally, without
assuming an orientation or a peeling of the retained target. At any stage,
the current reduction is the intersection of the current contraction with
the original coordinate face defining $D$, so remains convex. Let $x$ be
the next corner of the contraction and $Q$ its unique maximal corner cube.
Every neighbour of $x$ lies in $Q$. The intersection $Q\cap D$ is a face,
and its nonempty complement is connected, hence belongs to one component
of the complement of the reduction and has a single lift label $s$.

If $x\notin D$, its unique lift lies in layer $s$ and is a corner:
all of $Q$ lifts to that layer, and these are all its incident directions.
Delete it. If $x\in D$ and $Q\setminus D\ne\varnothing$, first delete
$(x,1-s)$. Its corner cube is $(Q\cap D)\times\{0,1\}$: it has the
vertical neighbour and precisely the old neighbours lying in $D$.
Afterwards $(x,s)$ is a corner with corner cube $Q\times\{s\}$;
delete it as well. If $Q\subseteq D$, the same two deletions work with
either choice of $s$. This also covers a singleton corner cube.
Thus every fibre deletion is a sequence of actual corner deletions.

After processing all words outside $B$, the remaining family is
\[
 (B\times\{t\})\ \cup\ (C\times\{1-t\}).
\]
If $C$ is empty we are finished. Otherwise delete its opposite-layer
copies in a corner order of $C$. At each step the corner cube of the
next vertex in the opposite layer is its current corner cube in $C$
multiplied by the vertical edge, since the entire matching layer $B$
remains. These are therefore corners of the whole residual family.
This leaves exactly $T$ and proves sufficiency.
Finally $C$ is a coordinate-face intersection of $B$; if $B\in\Damp$,
its nonempty pc-minor $C$ is in $\Damp$ as well.
\end{proof}

The extra condition in \eqref{eq:damp-convex-relative-face} is necessary
even for a full prism. For any nonpeelable ample $B$, put $X=B\times Q_1$
and retain $T=B\times\{0\}$. The contraction is $A=B$, which peels to
itself by the empty sequence, and its reduction is convex. But the opposite
layer is another copy of $B$, so $X$ cannot peel to $T$. For a positive
face target, the proposition instead permits every convex expansion to be
removed without changing relative attainability. It supplies a structural
simplification of these inputs, not a generation of all remaining inputs.

\begin{corollary}[A seven-vertex cube in every direction]
\label{cor:damp-every-direction-punctured-cube}
Let $X$ be an ample forbidden pc-minor or the underlying family of a
rooted factor. For every active coordinate $e$, there are two other coordinates $f,g$ and a coordinate three-face whose
intersection with $X$ has exactly seven vertices. In the two-dimensional
projection of this face along $e$, exactly three vertices have both lifts.
In particular, every nonempty single-coordinate reduction has at least
three vertices and is not a cube.
\end{corollary}
\begin{proof}
Put $A=\pi_eX$ and $D=X^{\{e\}}$. The preceding theorem makes $D$
nonconvex in $A$. Both are ample, and $D$ is nonempty and connected.
CCMW Lemma~4.10 characterizes convexity of a connected subclass by
containment of all ambient intervals between its vertices at distance two
\cite{ChalopiChepoiMoran22}. Since $D$ is isometric, it contains at least
one of the two possible intermediate vertices for any such pair. Failure
of this criterion therefore gives a square of $A$, in directions $f,g$,
with exactly three vertices in $D$. In $X$, those three vertices have
both lifts, while the fourth has exactly one. Fixing all other coordinates
gives the asserted seven-vertex intersection.
\end{proof}

The horizontal construction and its free-bit count already appear, for
representation maps without an acyclicity hypothesis, in the sibling
CompressionSchemes manuscript, in its Round~717 theorem on restricting
orientations and its coordinate-conditioning corollary, in
\path{drafts/round717_restricting_orientations.tex}.
The additions here are the projected-cycle proof of acyclicity, necessity
for preserving every horizontal edge, and the resulting constraints on
forbidden pc-minors. The construction and count are reused, not new results.

Unlike the earlier singleton argument, the lifting theorem handles every
convex reduction and gives an orientation extending prescribed horizontal
data. It uses the existing component-label parametrization and CCMW's
USO criterion; the local-convexity assertion used in the corollary is
CCMW's result. Nonconvex reductions require a different argument: at a vertical source,
the old outgoing neighbours inside the reduction can fail to span a cube
inside it. Thus an edge reduction does not require a new square-compatibility
theory here. The star theorem below handles the three-vertex path and
all larger stars without requiring horizontal preservation.
The punctured-cube condition is necessary, not a sufficient generator.

\begin{theorem}[All lifts preserving the horizontal orientation]
\label{thm:damp-horizontal-uso-lifts}
Let $X$ be ample, $e$ active, $A=\pi_eX$, and $D=X^{\{e\}}$.
Fix an acyclic USO $\mathcal O$ of $A$. There is an acyclic USO of $X$
whose every old-direction edge has the orientation of its projection in
$\mathcal O$ if and only if the induced orientation $\mathcal O|_D$
is a USO of $D$.

When this condition holds, let $s$ be the number of $d\in D$ with no
outgoing neighbour in $A\setminus D$. There are exactly $2^s$ such lifts.
At each remaining $d$, the vertical edge is forced towards the component
label of its outgoing neighbours outside $D$; at the $s$ specified vertices,
either vertical direction may be chosen independently.
\end{theorem}
\begin{proof}
Suppose first that a lift exists. At the source of the vertical edge over
$d$, all outgoing edges of $\mathcal O|_D$ are still outgoing. Condition
(C1) on the lift requires their whole cube, together with its vertical
translate, to lie in $X$. Their projected cube thus lies in $D$.
This proves (C1) for $\mathcal O|_D$. Condition (C2) holds because every
cube of $D$ is a cube of $A$ with the same orientation. Acyclicity is
inherited as well.

Conversely, assume $\mathcal O|_D$ is a USO. For any cube $Q$ of $A$,
$Q\cap D$ is ample, and its complement in $Q$ is ample. If nonempty,
that complement is connected. This is the cube-complement fact used in
Theorem~\ref{thm:damp-component-expansion-signs}; unlike the convex case,
the intersection need not be a face. Hence $Q\setminus D$ lies in a
single component of $A\setminus D$.
In particular, the outgoing neighbours outside $D$ of a fixed $d\in D$
all have one component label. Point the vertical edge towards that label,
or choose either direction when there are no such neighbours.

At a vertex outside the doubled reduction, its whole outgoing cube lifts
beside it by the preceding component observation. At the head of a
vertical edge, if there are outgoing neighbours outside $D$, their entire
outgoing cube lifts beside them by the same observation. If there are none,
the outgoing cube lies in $D$ by (C1) for $\mathcal O|_D$.
At the tail of a vertical edge, no outgoing external neighbour is present;
its old outgoing directions are exactly those of $\mathcal O|_D$.
Their cube lies in $D$, so its product with the vertical edge lies in $X$.
Thus (C1) holds everywhere. The proofs of (C2) and acyclicity from
Theorem~\ref{thm:damp-convex-reduction-lift} apply without convexity:
each doubled cube has its only sink over its base sink, and a directed
cycle would project to a directed closed walk in $\mathcal O$.

Finally, if an external outgoing neighbour has label $b$, pointing the
vertical edge away from $b$ would give its tail two outgoing directions
whose square lacks the neighbour's opposite lift. This violates (C1).
All other vertical directions were arbitrary in the construction, proving
both exhaustiveness and the count.
\end{proof}

This criterion concerns preservation of every horizontal edge direction.
It is more restrictive than allowing arbitrary acyclic USOs on $X$, or
prescribing only the representation maps induced by CCMW's minor transport.
It gives a direct special case of the general completion problem, with
independent vertical choices instead of a compatibility search. If one
acyclic USO of $A$ induces a USO on $D$, every component labeling of this
same pair gives a positive expansion. The theorem does not assert the
converse of this last implication.

\begin{corollary}[A forced orientation at a path reduction]
\label{cor:damp-path-reduction-forced-fork}
If an ample forbidden pc-minor $X$ has a reduction
$D=\{u,c,v\}$ inducing the path $u-c-v$, then every acyclic USO of
$A=\pi_eX$ directs both edges away from $c$. Equivalently, every peeling
of $A$ deletes $c$ before both $u$ and $v$.
\end{corollary}
\begin{proof}
The contraction $A$ is positive. On the three-vertex path, (C2) is automatic,
and the only failure of (C1) is directing both edges away from its middle
vertex. Any other induced orientation would lift by the theorem, making
$X$ positive. The peeling formulation follows from the CCMW correspondence.
\end{proof}

Failure for one prescribed orientation is not a negative certificate.
For example, take $A=\{0,1,2,3\}$, $D=\{0,1,2\}$, and orient the
square $A$ from $0$ towards $3$. The induced path orientation fails (C1),
so no lift preserves these horizontal directions. Nevertheless the
peripheral expansion $X=\{0,1,2,3,4,5,6\}$ peels in the order
$(5,6,4,3,2,1,0)$. Here the new bit has value $4$.
The two initial deletions are corners of squares, the third is a leaf,
and the remaining four vertices form a square. The universal quantifier
in Corollary~\ref{cor:damp-path-reduction-forced-fork} is therefore essential.

\begin{remark}[Positive expansions with no horizontally preserved USO]
\label{rem:damp-no-common-horizontal-control}
The existence condition in Theorem~\ref{thm:damp-horizontal-uso-lifts}
is not necessary for positivity even after varying the contraction USO.
Take the $289$-concept factor $A$ from
Remark~\ref{ex:damp-terminal-cut-vertex-1158}, with root zero, and put
$D=\{0,2,4\}$. The words $2,4$ are neighbours of zero, and zero is
the only geometric corner of $A$. Every acyclic USO of $A$ therefore
has zero as its unique source: sources are corners, and a finite acyclic
orientation has a source. Both path edges point out of zero, so no such
USO induces a USO on $D$.

Nevertheless every component expansion of $(A,D)$ peels. The fixed
connectivity check in \path{research/path_reduction_eligibility.py}
shows that $A\setminus D$ is connected. Thus the only two labelings give
\[
 X=(A\times\{0\})\cup(D\times\{1\})
 \quad\text{and its reflection in the new coordinate}.
\]
Each has $292$ concepts. In $X$, delete $(2,1)$ and $(4,1)$ as square
corners, then $(0,1)$ as a leaf, and follow the recorded peeling of $A$.
The verifier replays both complete orders. Thus neither positive expansion
has an acyclic USO preserving all horizontal directions of any acyclic
USO of its contraction. This strengthens the preceding fixed-orientation
example by quantifying over every contraction orientation; the input factor
and its corner certificate are reused.

The same bounded audit checks every three-vertex path in six stored
positive supports: the $289,292,308,330$-concept inputs and the two repairing
extensions of orders $291,290$. None of the $22316$ paths has a disconnected
complement. Hence these supports give only peripheral expansions over a
path, and cannot supply a forbidden example by this construction.
The result concerns these six inputs only. It neither excludes all path
reductions of forbidden classes nor extends the convex-reduction theorem
to arbitrary nonconvex reductions. The detailed inputs, counts, hashes,
and positive orders are stored in the verifier's JSON report.
\end{remark}

\begin{theorem}[Exact label criterion for a path reduction]
\label{thm:damp-path-reduction-reachability}
Let $A\in\Damp$ contain the ample path $D=\{u,c,v\}$ with edges
$uc,cv$. Let $K_1,\ldots,K_k$ be the components of $A\setminus D$,
and form the ample expansion
\[
 X_t=(D\times\{0,1\})\cup\bigcup_i(K_i\times\{t_i\}),
 \qquad t\in\{0,1\}^k.
\]
For an acyclic USO $\mathcal O$ of $A$, let $S_{\mathcal O}$ be the
vertices reachable from at least one vertex of $D$, including $D$ itself.
For $z\in\{u,v\}$, define the set of component indices
\[
 N_{\mathcal O}(z)=\{i: z\text{ has a neighbour in }
                K_i\cap S_{\mathcal O}\}.
\]
Then
\[
 X_t\in\Damp\quad\Longleftrightarrow\quad
 \exists\mathcal O\ \exists z\in\{u,v\}:
       |\{t_i:i\in N_{\mathcal O}(z)\}|\le1.
\]
Equivalently, the negative labelings are exactly the two-colourings for
which every member of the hyperedge family
\[
 \mathcal H(A,D)=\{N_{\mathcal O}(u),N_{\mathcal O}(v):
                  \mathcal O\text{ an acyclic USO of }A\}
\]
contains both colours. Empty or singleton hyperedges make such a colouring
impossible.
\end{theorem}
\begin{proof}
We use two consequences of the component description. An exclusive
vertex is a corner of the expansion exactly when its projection is a
corner of the contraction: all its old neighbours lift beside it, and
its corner cube does so by the cube-complement argument. At a doubled
vertex $z$, a lift in layer $b$ is a corner exactly when $z$ is a corner
of the current reduction and all its external neighbours have label $1-b$.
Indeed an external neighbour beside $z$ would require its absent opposite
lift in the corner cube containing the vertical direction; when none is
present, the corner cube is the product of the reduction corner cube
with that edge. These are the first-deletion rules from the
\href{research/component_label_partitions.pdf}{component-partition writeup}.

Fix $\mathcal O$ and an endpoint $z$ satisfying the displayed test.
The set $S_{\mathcal O}$ has no outgoing edge to its complement.
Choose a topological order deleting $A\setminus S_{\mathcal O}$ first.
It is a peeling of $A$ by CCMW, and its remaining family $S_{\mathcal O}$
is ample and positive. All deleted vertices are outside $D$, so the
exclusive-corner rule lifts this prefix to $X_t$. In the remainder, all
external neighbours of $z$ have one label $b$, or there are none. Delete
$(z,1-b)$, choosing either $b$ in the empty case. The doubled-corner rule
makes this a valid corner deletion. The contraction is still
$S_{\mathcal O}$, whereas the reduction is the edge $D\setminus\{z\}$.
This edge is convex in the contraction. Theorem~\ref{thm:damp-convex-reduction-lift}
therefore peels the remainder, proving sufficiency.

Conversely, choose a peeling of $X_t$ and consider its first deletion
of a lift of a vertex in $D$. Until then, the reduction is still $D$.
Its middle vertex $c$ is not a corner, so the doubled-corner rule makes
this first vertex a lift of an endpoint $z\in\{u,v\}$. All its external
neighbours in the remaining expansion have one label, or none remain.
The earlier exclusive deletions project to a valid corner-deletion prefix
of $A$. The remaining contraction $A'$ is positive, since the remaining
expansion is positive and positivity is contraction-closed. Extend this
prefix by a peeling of $A'$ and let $\mathcal O$ be its acyclic USO.
Every vertex deleted in the prefix precedes all of $D$; no directed path
from $D$ can reach it. Hence $S_{\mathcal O}\subseteq A'$.
The external neighbours of $z$ in $S_{\mathcal O}$ are therefore a subset
of the monochromatic set just identified. This gives the required test.
Negating the existential condition proves the hypergraph assertion.
\end{proof}

The general component-partition theorem permits several equality blocks
and successive changes of the reduction. Here each positive certificate
requires only one set of labels to agree: after the first reduction-vertex
deletion, completion is automatic. The criterion allows the two horizontal
copies of an edge to have different directions. In particular it accepts
the positive example of Remark~\ref{rem:damp-no-common-horizontal-control}:
there is only one complementary component, so every hyperedge has size
at most one. This is an exact classification of labelings over a fixed
positive contraction and path reduction, not an intrinsic generator of
all forbidden classes. Constructing $\mathcal H(A,D)$ still quantifies over
acyclic USOs of $A$, and forbidden-minor minimality is a further condition.

\begin{theorem}[Star reductions preserve every attainable endpoint]
\label{thm:damp-star-reduction-positive}
\label{thm:damp-path-reduction-positive}
Let $X$ be ample and let $e$ be active. Suppose the one-inclusion graph
of $D=X^{\{e\}}$ is a star, including a singleton or an edge. Then
\[
 X\in\Damp\quad\Longleftrightarrow\quad \pi_eX\in\Damp.
\]
More strongly, for every $a\in X$,
\begin{equation}
 X\text{ peels to }a
 \quad\Longleftrightarrow\quad
 \pi_eX\text{ peels to }\pi_ea.
 \label{eq:damp-star-endpoint-transfer}
\end{equation}
For every $z\in D$, the whole two-point fibre satisfies
\begin{equation}
 X\text{ peels to }\bigl(X\cap\pi_e^{-1}(z)\bigr)
 \quad\Longleftrightarrow\quad \pi_eX\text{ peels to }z.
 \label{eq:damp-star-fibre-transfer}
\end{equation}
Consequently, for any edge with an attainable endpoint that cannot be
retained by peeling, its coordinate class has at least four edges;
if it has exactly four, its reduction is a four-vertex path. No minimality
assumption is needed for this assertion.
Consequently no active coordinate reduction of an ample forbidden pc-minor
or a rooted factor is a star. In either class every coordinate edge class
has at least four edges; if it has exactly four, its reduction is a
four-vertex path.
\end{theorem}
\begin{proof}
Put $A=\pi_eX$ and assume $A$ positive. We induct on the number $r$ of
edges of the star $D$. For $r\le1$, the reduction is a singleton or an edge,
so Theorem~\ref{thm:damp-convex-reduction-lift} applies. Assume $r\ge2$.
We will remove one leaf lift, after clearing suitable exclusive vertices,
and then apply induction to the smaller star.

Choose a peeling of $A$ and lift its deletions up to, but not including,
its first vertex in $D$. These are exclusive corner deletions, by the
corner rules used in Theorem~\ref{thm:damp-path-reduction-reachability}.
Write $A'$ for the remaining positive contraction. Component labels below
are inherited from the original expansion.
If the next vertex is a leaf $z$, it is a corner of $A'$.
Its incident cube minus $D$ is connected when nonempty, by the
connected-exterior lemma used in
Theorem~\ref{thm:damp-horizontal-uso-lifts}. Thus all its external
neighbours have one component label. Delete the opposite lift of $z$
(or either lift if there is no external neighbour). It is a corner:
$z$ is a leaf of $D$ and this lift has no external neighbour.
The reduction becomes $D\setminus\{z\}$ and the contraction remains
$A'$. Induction completes the peeling.

It remains to treat the case where the centre $c$ is a corner of $A'$.
Translate $c$ to zero and write
\[
 D=\{0\}\cup\{z_i=0^{\{i\}}:i\in E\},\qquad |E|=r.
\]
All other coordinates are zero on $D$. Cornerhood puts the entire
$E$-cube at zero in $A'$. Its vertices of $E$-weight at least two form
a connected set: each can reach the all-one $E$-vertex by adding bits.
Let $K_*$ be their component in $A'\setminus D$.
Every neighbour of $c$ outside $D$ also belongs to $K_*$.
Indeed, for its direction $j\notin E$ and distinct $i,k\in E$, the
corner cube supplies the path
$0^{\{j\}},0^{\{j,i\}},0^{\{j,i,k\}},0^{\{i,k\}}$
entirely outside $D$.

Every vertex of $A'$ with $E$-weight at least two belongs to $K_*$:
join it by a geodesic to the vertex of the $E$-cube having the same
$E$-trace. The geodesic keeps that trace and avoids $D$.
Every vertex with zero $E$-trace, other than $c$, also belongs to $K_*$:
a geodesic from $c$ keeps that trace, starts at an external neighbour,
and cannot meet $D$ again. Hence each other component $K$ has only
unit $E$-traces. Connectedness forces its trace to be one fixed unit
vector, since two distinct unit vectors differ in two coordinates.
It meets the rest of the graph only at the corresponding leaf $z_i$.

These are one-vertex attachments on disjoint coordinate blocks.
To see this precisely, a coordinate outside $E$ cannot take value one
in two components of $A'\setminus D$. A geodesic between such vertices
would keep that coordinate one, whereas any path between the components
passes through $D$, where it is zero. For each component $K$, let $F_K$
be the set of outside-$E$ coordinates taking value one somewhere on $K$.
These sets are pairwise disjoint. Every vertex of $K\ne K_*$ has a
nonzero coordinate in $F_K$, since its unit trace alone is a vertex of $D$.
Thus the core $B=D\cup K_*$ is the coordinate-face intersection obtained
by setting all $F_K$, $K\ne K_*$, to zero. Each attached piece
$K\cup\{z_i\}$ is likewise obtained by fixing its $E$-trace and setting
all other outside-$E$ blocks to zero. All these pieces are positive.

Restrict an acyclic USO of $A'$ to the pieces. At most one leaf is a
sink in the induced USO of $B$: an ample USO has exactly one global
sink, since its out-map is a bijection onto its shattered sets and
only one vertex receives the empty set, by CCMW's representation-map
theorem. As $r\ge2$, choose a leaf $z$ with an outgoing edge in $B$.
No attached piece at $z$ can have an outgoing edge from $z$, because
that edge and the chosen outgoing core edge would span a missing mixed
square, violating (C1). Its restricted acyclic USO therefore has sink
$z$, so the piece peels to $z$ by CCMW Proposition~6.12.

Peel every attached piece at $z$ to $z$. Their other vertices have no
neighbours outside their piece, so these are valid deletions in the
whole contraction and lift to exclusive corner deletions in the expansion.
The resulting contraction is the coordinate-face intersection obtained by
fixing the deleted attachment blocks to zero; hence it is positive.
Every remaining external neighbour of $z$ lies in $K_*$ and has its
single label, or no such neighbour remains. Delete the opposite lift
of $z$, or either lift in the latter case. It is a corner by the same
leaf rule as above. The reduction is the smaller star
$D\setminus\{z\}$, so induction completes the proof of positivity.

We can keep an entire fibre, not just one endpoint. Suppose $z\in D$
and $A$ peels to $z$. We repeat the induction on the number of star edges,
now retaining both lifts of $z$. For a singleton or edge reduction, lift an
acyclic USO of $A$ with sink $z$ using
Theorem~\ref{thm:damp-convex-reduction-lift}. Neither lift of $z$ has an
outgoing old edge, so the boundary of its two-point fibre is inward.
Proposition~\ref{prop:damp-relative-uso-extension} gives relative peeling
to that edge. This argument in fact applies to every convex reduction.

For a larger star, choose the initial contraction order to end at $z$.
Its prefix before the first reduction vertex avoids $z$, and the remaining
contraction $A'$ still peels to $z$. If the first reduction vertex is a
leaf $w$, then $w\ne z$. The leaf step above removes a lift of $w$,
keeps both lifts of $z$, and leaves the contraction $A'$ unchanged.
Induction on the smaller star completes the relative peeling.

If the first reduction vertex is its centre $c$, then $z\ne c$, so $z$
is a leaf. Use the centre-case decomposition $B=D\cup K_*$ and its
one-vertex attachments proved above. Choose an acyclic USO of $A'$ with
sink $z$. Its restriction to the coordinate face $B$ also has sink $z$.
Choose a different leaf $w$, which exists since the star has at least two
edges. This $w$ has an outgoing edge in $B$. The same missing-square
argument makes every attachment at $w$ peel to $w$. Clear these attachments
as above. The remaining contraction is a coordinate face containing $z$,
so it still peels to $z$ by rooted minor closure. Now remove the leaf lift
of $w$ opposite the common exterior label. The reduction is a smaller
star still containing $z$, and neither lift of $z$ has been touched.
Induction gives the required relative peeling. The reverse implication
follows by contracting a relative peeling to the fibre. This proves
\eqref{eq:damp-star-fibre-transfer}.

The ordinary converse is contraction closure. To prove the stronger
endpoint assertion, necessity is endpoint minor closure. Conversely,
suppose $\pi_eX$ peels to $\pi_ea$. The ordinary assertion just proved
makes $X$ ordinarily positive. If $X$ failed at $a$, glue two copies of
$X$ at $a$, on disjoint coordinate sets, to form $Z=X\vee_a X$.
Proposition~\ref{prop:damp-vertex-gluing} makes $Z$ ample and ordinarily
negative, since both positive pieces fail at their common root.
Contract $e$ in the first copy. The resulting amalgam is ordinarily
positive: peel the contracted first copy to $\pi_ea$, then peel the
second copy ordinarily. On the other hand, the reduction $Z^{\{e\}}$
is exactly $X^{\{e\}}$, with the second copy's coordinates fixed at
zero. Indeed no vertex exclusive to the second copy has an $e$-mate.
This is a star, so the ordinary assertion applied to $Z$ makes it
positive, a contradiction. This proves
\eqref{eq:damp-star-endpoint-transfer} without choosing a prescribed
lift of the contraction's USO.

For the edge-retention consequence, an attainable endpoint makes its
projected vertex attainable in the contraction. A star reduction would
therefore allow retention of the entire edge by the fibre assertion.
A convex reduction also allows this, by its base-case proof above.
Every nonempty ample family of at most three vertices is a star, and an
ample family of four vertices is a square or a tree. The convex square
and the three-leaf star are excluded, leaving only the four-vertex path.
Corollary~\ref{cor:damp-iterable-path-edge} attains this bound: its marked
coordinate reduction is $H=a-b-c-d$, its lower endpoint is attainable,
and its edge cannot be retained.

For a forbidden class the proper contraction is positive, so its
reduction cannot be a star. For a rooted factor, its proper contraction
peels to the root image, so the endpoint assertion gives the same
exclusion. Active
reductions are nonempty, and every ample family of at most three vertices
is a star. An ample family with four vertices is either a square or a
tree: if it has no square, ampleness gives VC dimension at most one,
and its connected one-inclusion graph is a tree. The square is excluded
by convex-reduction lifting, and the three-leaf star by this theorem,
leaving only the four-vertex path. For a rooted factor $X$, the double
$X\vee_a X$ is a forbidden pc-minor by
Theorem~\ref{thm:damp-cut-vertex-obstructions}; each coordinate reduction
in either copy is unchanged. The same four-vertex restriction therefore
applies to $X$. Reduction vertices correspond bijectively to edges of
the coordinate class, proving the conclusions.
\end{proof}

The theorem extends the three-vertex path case to arbitrarily many leaves.
In particular the hypergraph in
Theorem~\ref{thm:damp-path-reduction-reachability} has no proper two-colouring
for a positive contraction. Its exact criterion remains valid, but cannot
produce a negative expansion. The proof uses a full cube at the first
central corner to separate all remaining attachments, and the unique
global sink of a CCMW USO to choose a leaf that can be cleared.

The doubling argument strengthens ordinary contraction positivity to
prescribed endpoints because gluing at the root leaves the selected
reduction unchanged. In particular, the attainable endpoint set of $X$
is the full inverse image of that of $\pi_eX$ whenever the reduction is
a star. The replay \path{research/star_endpoint_transfer.py/json} checks
this identity on the peripheral $291$-word family: its contraction has
unattainable endpoints $0,256$, and the family has exactly the three
unattainable endpoints $0,256,4096$. All $288$ positive endpoint orders
replay. This is a fixed control for the uniform proof, not its basis.

\begin{proposition}[Retaining a fibre near the centre of a tree reduction]
\label{prop:damp-tree-radius-two-fibre}
Let $X$ be ample and let its reduction $D=X^{\{e\}}$ in an active
coordinate have a tree as its one-inclusion graph. Fix $z\in D$ and
assume that every vertex of $D$ has graph distance at most two from $z$.
Then the following conditions are equivalent:
\begin{enumerate}
\item at least one of the two lifts of $z$ is an attainable peeling endpoint;
\item $X$ peels to its entire two-point fibre over $z$;
\item both lifts of $z$ are attainable endpoints.
\end{enumerate}
Consequently, if an edge with an attainable endpoint cannot be retained
and its coordinate reduction is a tree, the projected marked vertex has
eccentricity at least three in that tree. For a four-vertex-path reduction,
it must be an endpoint of the path. The bound three is attained by
Corollary~\ref{cor:damp-iterable-path-edge}.
\end{proposition}
\begin{proof}
Retaining the edge allows either order on its last two vertices, so
item 2 implies item 3, which implies item 1. We prove item 1 implies
item 2 by induction on $|D|$.
If $D$ is a star, contract a peeling to the specified endpoint and apply
\eqref{eq:damp-star-fibre-transfer}. This includes a singleton and an edge.

Suppose $D$ is not a star. Then $z$ is not a leaf: if $z$ were a leaf
and every vertex were within distance two of it, its neighbour would be
adjacent to every other vertex, making the tree a star.
Choose a peeling ending at one lift of $z$ and consider the first deletion
of a vertex lying over $D$. Until this step every fibre over $D$ is intact,
so the current reduction is still $D$. If the deleted vertex lies over $w$,
its corner cube projects to a corner cube at $w$ in $D$: every neighbour
of $w$ in the reduction has both lifts, and every subset of these directions
occurs in the corner cube. Hence $w$ is a corner of the tree $D$, and
therefore a leaf. In particular $w\ne z$.

Let $X'$ be the residual family just after this first fibre breaks. It is
ample, its reduction is $D\setminus\{w\}$, and the suffix of the chosen
order still peels to the same lift of $z$. Both lifts of $z$ remain.
Deleting a leaf from a tree preserves the distances among its remaining
vertices. The smaller reduction is thus a tree whose vertices are still
within distance two of $z$. Induction gives a relative peeling of $X'$
to the whole fibre over $z$. Prefixing the deleted vertices gives the
same relative peeling from $X$ and proves the equivalence.

The necessary eccentricity bound is the contrapositive. Only the endpoints
of a four-vertex path have eccentricity three. In the iterative path
construction, the marked vertex is just such an endpoint, its lower lift
is attainable, and its fibre cannot be retained. This proves sharpness.
\end{proof}

This implication uses an endpoint order of $X$ itself. It does not assert
that a contraction endpoint alone lifts through an arbitrary tree reduction.
For the smallest possible edge-retention failure with an attainable endpoint,
only fibres over the two ends of a four-vertex path remain to be considered.
Whether both lifts can be attainable in a failure of that particular shape
is not settled by this proposition.

\begin{proposition}[Attainable corners at an end of a path reduction]
\label{prop:damp-path-end-corner-retention}
Let $X$ be ample, let $D=X^{\{e\}}$ be the four-vertex path
$a-b-c-d$, and write $a_0,a_1$ for the two lifts of $a$.
If some $a_i$ is both a geometric corner of $X$ and an attainable
endpoint, then $X$ peels to $\{a_0,a_1\}$.

More generally, suppose both $a_0,a_1$ are attainable but their edge
cannot be retained. Put $A=\pi_eX$ and label each component of $A\setminus D$
by its unique lift layer. There exist disjoint sets $R_0,R_1\subseteq A\setminus D$,
with $R_i$ contained in components labelled $i$, such that
\begin{enumerate}
\item each $R_i$ has a corner-deletion order with all of $D$ retained;
      the two orders can also be performed successively;
\item $A_i=A\setminus R_i$ peels to $a$, for $i=0,1$;
\item $A_*=A\setminus(R_0\cup R_1)=A_0\cap A_1$ is ample,
      $a$ is a leaf of $A_*$ with neighbour $b$, and $A_*$ does not peel to $a$.
\end{enumerate}
In particular, neither lift is a geometric corner in this hypothetical
failure. The intersection in item 3 is not asserted to be corner-peelable,
and is not the original coordinate reduction $D$.
\end{proposition}
\begin{proof}
First observe that if a peeling to $a_i$ first breaks the fibre over $d$,
its prefix avoids both lifts of $a$ and leaves a three-vertex-star reduction.
Its suffix still ends at $a_i$. The star fibre theorem then retains the
whole fibre over $a$, contradicting any assumed failure of edge retention.
The first broken fibre cannot be over $b$ or $c$, since the projection
of such a corner would be a corner of $D$. Thus, under failure, every
peeling to $a_i$ first breaks the fibre over $a$ by deleting $a_{1-i}$.
All preceding deletions are exclusive vertices.

We will also use the following elementary consequence of star retention.
If $a$ is a leaf of a contraction $A'$ and the corresponding expansion
$X'$ still has reduction $D$, then
\[
 A'\text{ peels to }a\quad\Longrightarrow\quad
 X'\text{ peels to }\{a_0,a_1\}.
\]
Indeed the two lifts of $a$ form the end edge of a square over $ab$,
with no other incident edges. Removing both gives an ample family $Y$
with reduction $b-c-d$ and contraction $A'\setminus\{a\}$.
In a peeling of $A'$ ending at its leaf $a$, the neighbour $b$ is
penultimate: deleting it earlier would isolate $a$ while other vertices
remain, contrary to connectedness of nonempty ample residuals.
Thus $A'\setminus\{a\}$ peels to $b$, and the star theorem peels $Y$
to the fibre over $b$. Perform these deletions while keeping both lifts
of $a$: no deleted vertex is adjacent to them, and its corner cube
cannot contain either lift, since both have only their square neighbours.
The remaining square peels to the fibre over $a$.

We explain the independence of exclusive prefixes carefully. Every old
cube not contained in $D$ has all its vertices outside $D$ in one
component of $A\setminus D$: its intersection with $D$ is ample, so
its nonempty cube complement is connected. An exclusive corner test
therefore depends only on its own component and on the retained $D$.
Deletions in different components can be reordered, or omitted from a
prefix, without affecting these tests. The same holds in the expansion
by the exclusive-corner rule. Their projected steps are corner deletions
in the contraction.

Choose an endpoint order ending at $a_{1-i}$. From its prefix before
the first fibre breaks, retain just the deletions in components labelled
$i$, and call their projected set $R_i$. They remove every exterior
neighbour of $a$ with that label: otherwise the proposed deletion of
$a_i$, whose vertical mate is still present, would require a missing
mixed square. Reorder the original prefix to perform $R_i$ first and
its other components afterwards. The residual contraction after the
whole prefix peels to $a$, as the projection of the endpoint suffix.
Consequently $A_i=A\setminus R_i$ also peels to $a$.

If an attainable $a_j$ was initially a corner, it had no exterior
neighbour in its own layer. Apply the preceding construction to its
endpoint order, taking $R_{1-j}$. After this prefix, all exterior
neighbours of $a$ in both layers are absent, so $a$ is a leaf of
$A_{1-j}$, which peels to $a$. The leaf observation gives edge retention,
and prefixing $R_{1-j}$ gives it in $X$. This proves the first assertion.

For the last assertion, use the two endpoint orders to obtain $R_0,R_1$.
Their components are disjoint, so their orders can be concatenated.
The resulting contraction $A_*$ and expansion $X_*$ are ample, both
lifts of $a$ remain, and $a$ has just its neighbour $b$ in $A_*$.
If $A_*$ peeled to $a$, the leaf observation would retain the edge in
$X_*$ and hence in $X$, a contradiction. This proves all three items.
Finally both attainable lifts must be noncorners by the first assertion.
\end{proof}

This isolates a limitation of combining two endpoint orders. Their
componentwise exclusive prefixes can indeed be combined while preserving
ampleness, but the required endpoint need not survive in the intersection.
In the hypothetical failure it necessarily does not survive. No closure
of endpoint-positive families under such intersections has been established;
assuming it would leave precisely the missing step unproved.

\begin{proposition}[Contractions without short separating paths]
\label{prop:damp-path-separation-input-exclusion}
Let $B$ be ample. Suppose that $B\setminus V(P)$ is nonempty and
connected for every geodesic path $P$ in $B$ on at most four vertices,
including a single vertex. Form $A$ by attaching finite trees to vertices
of $B$, each tree meeting $B$ only at its attachment vertex and using
fresh coordinates, one per tree edge. Different trees use disjoint
coordinates. Let $X$ be an ample expansion with contraction $A$ and
reduction a four-vertex path $D$. For $z\in D$, if both lifts of $z$
are attainable endpoints of $X$, then $X$ peels to their edge.
\end{proposition}
\begin{proof}
For an internal vertex of $D$, this follows from
Proposition~\ref{prop:damp-tree-radius-two-fibre}. Assume that $z$ is an end.
Repeatedly remove a leaf outside $B\cup D$ from an attached tree.
Its unique lift in the expansion is also a leaf: it is outside the
reduction and has just its one horizontal neighbour. Deleting this lift
is a corner deletion and a coordinate restriction. Indeed the coordinate
of its remaining tree edge occurs on that leaf alone, after the earlier
branches have been removed. Both marked lifts survive, and their
attainability survives by root-preserving pc-minor closure. Denote the
remaining contraction and expansion by $A'$ and $X'$.

We claim that $A'\setminus D$ is nonempty and connected. If $D$ meets
$B$, its intersection with $B$ is a geodesic path on at most four
vertices: a simple path cannot enter and leave a tree through its
single attachment vertex. Every surviving tree vertex then belongs
to $D$. Thus $A'\setminus D=B\setminus(D\cap B)$, and the hypothesis
applies. If $D$ misses $B$, it lies in one attached tree. The surviving
vertices outside $B\cup D$ form precisely the path joining that tree's
attachment vertex to $D$, with its endpoint in $D$ omitted. Their union
with $B$ is connected and nonempty. This proves the claim.

All vertices of $A'\setminus D$ consequently have the same lift layer.
The other layer consists exactly of $D$. Its lift of the end $z$ is a
geometric corner: its two neighbours span the square over the end edge
of $D$. This lift is still attainable, so
Proposition~\ref{prop:damp-path-end-corner-retention} retains the marked
edge in $X'$. Prepending the deleted leaves gives the required peeling
of $X$.
\end{proof}

This excludes the standard paired repairs as a source of the missing
component geometry, even after tree attachments and changes of the marked
path. In the archived binary encoding let $C$ be the $289$-vertex core.
The four ample families
\[
 C,\qquad C\cup\{290\},\qquad C\cup\{128,160\},\qquad
 C\cup\{290,128,160\}
\]
satisfy the path-complement hypothesis. The exhaustive finite check in
\path{research/path_separation_inputs.py/json} verifies ampleness and
connectivity after every geodesic path on one through four vertices.
There are respectively $13368$, $13480$, $13561$, and $13675$ paths
on four vertices, counted up to reversal. The proposition proves the
unbounded tree-attachment consequence; the four input checks are
computational. This does not exclude other contractions or classify the
required component-separated intersections.

For comparison, the sibling CompressionSchemes manuscript already proves
representation-map lifting over \emph{every} tree reduction, in
Theorem~\texttt{round736-tree-overlap-lifting} of
\path{drafts/round736_tree_overlap_lifting.tex}. Its outgoing-cube labels
and quotient-tree sink select a lower root and give a representation map
for every component colouring. That theorem does not assert acyclicity
of the lifted orientation. The star theorem above proves corner-peeling
closure for its stated subfamily; it does not require a prescribed upper
map to survive. For arbitrary trees this positivity implication is false, as
Theorem~\ref{thm:damp-four-vertex-path-obstruction} below shows. A branch of a longer
tree can meet a complementary component at several vertices, so the
one-vertex attachment argument does not establish that extension.
The first possible reduction shape has four vertices and is a path;
the theorem below proves that it occurs in the ample forbidden basis.

\begin{remark}[Tree lifting does not preserve acyclicity for every choice]
\label{rem:damp-tree-lift-cycle}
The acyclicity qualification above is necessary even for a four-vertex path
inside a cube. Encode vertices by binary masks, with coordinate weights
$1,2,4$, and take
\[
 U=\{0,\ldots,7\},\qquad H=\{1,3,7,6\},\qquad
 (r(0),\ldots,r(7))=(7,6,1,4,3,0,5,2).
\]
Supports also use masks. The map $r$ is an acyclic USO: the order $(0,1,4,6,2,3,7,5)$ is a corner peeling and its outgoing
supports are precisely the displayed values. The reduction is the path
$1-3-7-6$. Its complement in $U$ is connected.

In the Round~736 construction the only vertex whose outgoing cube is
contained in $H$ is $3$. The labelled blocks are $\{1\}$ and $\{6,7\}$,
and the free singleton block is $\{3\}$; the quotient is directed toward
$\{6,7\}$. Choose lower root $6$, give the complementary component bit
zero, and choose the free vertical-source bit at $3$ to be zero.
The other three vertical-source bits are one. With new coordinate weight
$8$, the expansion and lifted map are
\[
\begin{array}{c|rrrrrrrrrrrr}
v&0&1&2&3&4&5&6&7&9&11&14&15\\ \hline
R(v)&7&6&1&12&3&0&5&2&10&4&8&9.
\end{array}
\]
This is a representation map by the cited tree-lifting theorem, but its
orientation contains the directed cycle
\[
 3\longrightarrow11\longrightarrow15\longrightarrow14
 \longrightarrow6\longrightarrow2\longrightarrow3.
\]
Each arrow follows by testing its coordinate in the displayed out-map.
Changing just the free bit at $3$ exchanges its two assignments:
$R(3)=4$ and $R(11)=12$. All vertical edges now point from the upper
copy of $H$ to the lower cube. The upper tree and lower cube are acyclic,
so this modified map is acyclic. Independently, the expansion peels in the
order $(9,11,15,14,0,1,4,6,2,3,7,5)$.
Thus the example refutes preservation for arbitrary certified choices;
it does not refute existence of an acyclic choice or tree-reduction closure.
The script \path{research/tree_lift_acyclicity.py} also checks the maps by
nonclashing, outgoing cubes, and unique sinks on every contained cube.
\end{remark}

\begin{proposition}[Canonical free vertical directions in a tree lift]
\label{prop:damp-tree-free-bit-normalization}
Let $\varnothing\ne H\subseteq U$ be ample, with $G(H)$ a tree, and
fix a representation map $r$ of $U$. Put
\[
 Q_h=Q(h,r(h)),\qquad Z=\{h\in H:Q_h\subseteq H\}.
\]
For $h\notin Z$, let $W_h$ be the component of $U\setminus H$
containing $Q_h\setminus H$. Contract the edges of $H$ whose endpoints
lie outside $Z$ and have the same label $W_h$, retaining every vertex
of $Z$ as a singleton. Let $K_*$ be the sink block of the quotient
oriented by $r$, as supplied by the Round~736 tree-lifting theorem.
Fix $q\in K_*$, root $H$ at $q$, and let $p(h)$ be the parent of
$h\ne q$. Write $s_q$ for the out-map of this rooted tree.

Fix component bits $t_W$ and use a fresh coordinate $x$ to form
\[
 X_t=(H\times\{0,1\})\cup
       \bigcup_W(W\times\{t_W\}).
\]
The cited theorem supplies a representation map $R_b$ for every choice
of source bits
\[
 b(h)=1-t_{W_h}\quad(h\notin Z),\qquad b(h)\in\{0,1\}\quad(h\in Z),
\]
by the assignments
\[
\begin{aligned}
 R_b(u,t_W)&=r(u) &&(u\in W),\\
 R_b(h,1-b(h))&=r(h) &&(h\in H),\\
 R_b(h,b(h))&=s_q(h)\cup\{x\} &&(h\in H).
\end{aligned}
\]
Define a canonical choice $b_*$ by retaining the forced bits, setting
$b_*(q)=0$ if $q\in Z$, and setting
\[
 b_*(h)=b_*(p(h))\qquad(h\in Z\setminus\{q\}),
\]
in increasing distance from $q$. Then
\[
 \text{some }R_b\text{ is acyclic}
 \quad\Longleftrightarrow\quad R_{b_*}\text{ is acyclic}.
\]
In particular, for fixed $r$ and $t$, existence of an acyclic map within
this construction requires at most $|K_*|$ cycle tests, one for each root;
there is no need to search the $2^{|Z|}$ free-bit assignments.
Moreover, moving a successful root along an outgoing $r$-edge inside
$K_*$ preserves success. Thus it suffices to test the sinks of
$r$ restricted to $K_*$.
\end{proposition}
\begin{proof}
The tree-lifting theorem states that $r$ and $s_q$ agree on every
uncontracted tree edge. For $h\in Z$, the outgoing cube of $r$ lies
in a tree, so its support is empty or a singleton. Every edge incident
to this singleton block is uncontracted. It follows that
\[
 r(h)=s_q(h)=\{\text{direction of }hp(h)\}\quad(h\in Z\setminus\{q\}),
\]
and $r(q)=s_q(q)=\varnothing$ when $q\in Z$.
Consequently changing a free bit changes only that vertical edge:
all horizontal directions stay fixed. At either lift of a free nonroot
vertex, the only outgoing horizontal edge goes to the corresponding
lift of its parent.

Consider a free $h\ne q$ whose bit differs from its parent's. Put
$a=b(p(h))$ and change $b(h)$ to $a$. Suppose this creates a directed
cycle in a previously acyclic orientation. That cycle must use the
new vertical edge, and its next edge is forced. Thus it contains
\[
 (h,a)\longrightarrow(h,1-a)\longrightarrow(p(h),1-a).
\]
Replace this segment by
\[
 (h,a)\longrightarrow(p(h),a)\longrightarrow(p(h),1-a).
\]
Both replacement edges existed before the change: the first is the
fixed outgoing horizontal edge at $h$, and the second is the parent's
vertical edge with source bit $a$. Every other edge of the cycle also
existed before the change. The replacement therefore gives a nonempty
directed closed walk in the old orientation, which contains a directed
cycle, a contradiction. Matching a free bit to its parent's thus preserves
acyclicity. Representation-map validity throughout is supplied by the
original tree-lifting theorem.

If $q\in Z$, its two lifts have no outgoing horizontal edge. Reversing
its vertical edge cannot create a cycle, since the new head has no
outgoing edge. Set that bit to zero first. Then process the remaining
free vertices in increasing distance from $q$, matching each bit to its
parent's already fixed bit. Each step preserves acyclicity and the final
assignment is $b_*$. Hence an acyclic $R_b$ implies an acyclic $R_{b_*}$.
The reverse implication is immediate, since $b_*$ is an allowed choice.
Repeating this single canonical test for each $q\in K_*$ proves the
first bound on the number of tests.

To prove the root reduction, the case $|K_*|=1$ is immediate. Otherwise
$K_*$ has one exterior label $W$, so its source bits all equal
$a=1-t_W$. The canonical bits do not depend on the root within $K_*$:
outside that subtree, paths toward any such root have the same initial
segment up to their first vertex in $K_*$, where the prescribed bit is
always $a$.
Suppose $q\to q'$ is an $r$-edge inside $K_*$ and move the root from
$q$ to $q'$. This reverses only the lower edge over $qq'$, from
$(q',a)\to(q,a)$ to $(q,a)\to(q',a)$. All other edges stay fixed.
If a cycle is created, it uses this new edge. The new lower root
$(q',a)$ has only its vertical edge outgoing, so the cycle contains
\[
 (q,a)\longrightarrow(q',a)\longrightarrow(q',1-a).
\]
Replace this segment by
\[
 (q,a)\longrightarrow(q,1-a)\longrightarrow(q',1-a).
\]
The latter edges existed before the root move: the first is vertical,
and the second is the upper edge oriented by $r$. The old orientation
would therefore contain a directed closed walk, a contradiction.
Starting with any successful root, follow outgoing $r$-edges in $K_*$.
Since its underlying graph is a finite tree, this ends at a sink, and
every root along the path remains successful. This proves the final claim.
\end{proof}

This normalization removes the vertical-choice part of the tree problem;
it does not prove that a successful root always exists. It strengthens
the earlier example's one-bit repair to a uniform operation on every tree
lift. Round~736 supplies the maps and their quotient-tree structure; the
additional argument here is that the indicated free-edge reversals cannot
create cycles. The earlier general orientation-gluing criterion tests
prescribed upper-layer assignments, whereas this proposition removes those
free assignments for the certified tree roots. An acyclic $r$ need not have any successful canonical root:
Theorem~\ref{thm:damp-four-vertex-path-obstruction} below gives a negative
expansion with positive contraction, and hence refutes that existence claim.

\begin{theorem}[A path reduction does not preserve corner peeling]
\label{thm:damp-four-vertex-path-obstruction}
There is an ample forbidden pc-minor $X$ for $\Damp$ with $909$ vertices
in $37$ active coordinates
and an active coordinate $e$ such that
\[
 X\notin\Damp,\qquad \pi_eX\in\Damp,\qquad
 X_{e=0},X_{e=1}\in\Damp,
\]
while $X^{\{e\}}$ is the four-vertex path $2-0-4-5$ in binary-mask
notation. In particular, expansions over tree reductions need not preserve
corner peeling. The class $X$ itself has a coordinate edge class
of exactly four edges, so the bound in
Theorem~\ref{thm:damp-star-reduction-positive} is sharp.
\end{theorem}
\begin{proof}
We first describe the construction in terms of three pieces and their
boundary conditions. The obstruction is the first attempted deletion from
the doubled path. We then specify the finite inputs and prove positivity of every elementary
pc-minor, grouping the checks by coordinate block.

\emph{The three inputs.}
Let $B$ be positive ample with unique corner $b$, containing the induced
path $H=a-b-c-d$. Denote its successive directions by $f,g,h$.
Assume $B$ has an acyclic USO with sink $a$ and out-map $r_B(d)=\{h\}$.
Let $(A,a)$ be positive ample but not peelable to $a$, using a fresh
coordinate block outside its attachment $a$. Let $P$ be positive ample,
meeting $B$ exactly in $\{c,d\}$ and using a fresh coordinate block
apart from the shared direction $h$. Require
\[
 t_c(P)=1,\qquad \ell_d(P)=0.
\]
The exclusive coordinate blocks of $A$ and $P$ are disjoint. Put
\[
 U=B\cup A\cup P,\qquad S=H\cup A\cup P,\qquad
 X=(B\times\{0\})\cup(S\times\{1\}).
\]
Thus $\pi_eX=U$ and $X^{\{e\}}=B\cap S=H$.

The coordinate-block gluing results make $U$ and $S$ ample.
They are also positive. For $U$, glue an arbitrary acyclic USO of $A$
to the prescribed map of $B$, which has sink $a$ at their shared vertex.
The resulting map still has out-map $\{h\}$ at $d$. A map of $P$ with
sink $c$ then glues along $cd$, by
Theorem~\ref{thm:damp-edge-gluing-profile}. For $S$, first peel $P$ to
$c$, delete $c$ and then $b$, and finish with any peeling of $A$.
The lower section $B$ is positive by hypothesis.

To check ampleness of $X$, use the section shadow identity. Within the
coordinates of $B$, the shadow of $S$ contains only the empty support and
$\{f\},\{g\},\{h\}$: the first two come from the path $a-b-c$, and
the third from the shared edge in $P$. All other coordinates of $A$ and
$P$ are exclusive. Consequently
\[
 \operatorname{Sh}(B)\cap\operatorname{Sh}(S)
   =\operatorname{Sh}(H)=\{\varnothing,\{f\},\{g\},\{h\}\}.
\]
The old-coordinate shadow of $X$ is $\operatorname{Sh}(U)$, and its
supports containing $e$ are obtained by adjoining $e$ to this intersection.
Since $U$ is ample, $|\operatorname{Sh}(X)|=|U|+4=|X|$, proving ampleness.

\emph{Why no peeling can start deleting the doubled path.}
Suppose $X$ has a complete peeling and consider its first deletion from
$H\times\{0,1\}$. No vertex of the lower copy of $B$ has previously
been deleted. Indeed, the first such deletion would restrict to a corner
deletion of $B$, whose only corner is $b$, already in the doubled path.
Both internal path vertices have noncorner projections in $H$, so neither
of their lifts can be the first doubled-path corner. Each endpoint $a,d$
has a neighbour in $B\setminus H$: otherwise its only neighbour in $B$
would be its path neighbour, making it a corner of $B$, contrary to
uniqueness of $b$. These neighbours are still present on the lower shore
and have no upper mates. They prevent the lower lifts of $a,d$ from being
corners while their vertical mates remain.

If $(a,1)$ is first, it can have no remaining neighbour in $A\setminus\{a\}$,
since the lower mate of such a neighbour is absent. Restrict the preceding
deletions to the upper copy of $A$. They are corner deletions retaining $a$;
the remaining class is ample and connected. With $a$ isolated, it is the
singleton $\{a\}$, contradicting the failure of peeling $A$ to $a$.
Finally, if $(d,1)$ is first, it has no remaining exclusive neighbour in
$P$. Restrict the complete peeling to the upper copy of $P$. It deletes
$d$ before $c$, and $d$ then has exactly the neighbour $c$. This witnesses
$\ell_d(P)=1$, the other contradiction. These exhaust all possible first
doubled-path deletions, proving $X\notin\Damp$.

\emph{Explicit realization.}
Take for $B$ the $289$-vertex class $C$ of
Theorem~\ref{thm:damp-three-piece-witness}, and put
\[
 (a,b,c,d)=(2,0,4,5),\qquad (f,g,h)=(2,4,1)
\]
as coordinate masks. Its only corner is zero. The certificate in
\path{research/four_vertex_path_obstruction.json} is a complete corner
order ending at $2$, in which $5$ is deleted with sole remaining neighbour
$4$. Take $A$ to be a disjoint-coordinate copy of $(C,0)$ attached at $2$.
Take $P$ to be the $330$-vertex piece of
Theorem~\ref{thm:damp-three-piece-witness}, identifying its lower endpoint
with $4$ and its upper endpoint with $5$. Its established profile
$(0,1,0,1,0)$ gives exactly $t_c(P)=1$ and $\ell_d(P)=0$.
Thus all hypotheses of the construction hold.

More explicitly, keep the twelve bits of $B$, embed $A$ by
$v\mapsto 2\mathbin{\oplus}(v\ll12)$, and embed $P$ by
\[
 v\longmapsto4\mathbin{\oplus}((v\bmod4096)\ll24)
                 \mathbin{\oplus}\lfloor v/4096\rfloor.
\]
Here $\oplus$ is bitwise exclusive-or and $v\ll k=2^kv$.
Use bit $2^{36}$ for the new coordinate. The resulting sizes are
$|U|=289+289+330-3=905$, $|S|=620$, and $|X|=289+620=909$.
The verifier replays the core, contraction and section orders and checks
the contraction map independently by support bijection and nonclashing.
It also checks the exact section shadows. Negativity uses the first-deletion
argument above and the previously proved profile, not an unsuccessful search.

\emph{Every proper minor is positive.}
It suffices to check contractions and the two sections of each active
coordinate, since $\Damp$ is pc-minor closed. We use four coordinate groups.
The new coordinate $e$ gives the three positive classes $U,B,S$ already
certified above.

For any of the three path directions $f,g,h$, the resulting reduction
along $e$ is a proper minor of $H$, hence a star or empty. To justify this,
reduction along $e$ commutes with operations on another coordinate $j$ in
an ample class. Restrictions commute immediately. For contraction, the
only nontrivial inclusion arises when the projected $e$-mates are represented
on opposite $j$-sides. These representatives have Hamming distance two;
isometry supplies an intermediate vertex of their square, and hence an
actual $e$-pair on one $j$-side. Thus
$(\pi_jX)^{\{e\}}=\pi_j(X^{\{e\}})$.
If $e$ remains active, its contraction is a positive minor of $U$, so
Theorem~\ref{thm:damp-star-reduction-positive} applies. If $e$ is constant,
the class is a minor of the corresponding positive section of $X$.
An empty reduction cannot occur with both sections nonempty: an ample
class is connected and would then have an $e$-edge.

Consider a coordinate exclusive to $A$. Its section away from the
attachment lies in $A$ and is positive. Contraction or the section retaining
$a$ replaces $A$ by a proper rooted minor $A'$, which peels to $a$ by the
rooted-factor property of $(C,0)$. Perform that relative peeling on the
upper shore. Each deleted vertex is outside the attachment and has no
neighbour in the rest of the class, so its corner property is preserved.
Then $(a,1)$ is a corner: its incident directions are $f,e$, and their
square is present. Delete it. The remaining reduction is $H\setminus\{a\}$,
a three-vertex path, and the contraction is $B\cup P$. This contraction
is positive by gluing the prescribed map of $B$, with leaf exit at $d$,
to the sink-$c$ map of $P$. The star theorem finishes the peeling.

For a coordinate exclusive to $P$, the section away from $cd$ is a positive
section of $P$. Contraction or the section retaining $cd$ replaces $P$ by
a piece $P'$ which peels to that edge. These are precisely the $24$
relative orders verified in Theorem~\ref{thm:damp-three-piece-witness}.
Peel its upper exclusive vertices, then delete $(d,1)$, whose two incident
directions $h,e$ span a present square. The remaining reduction is
$H\setminus\{d\}$; the contraction $B\cup A$ is positive by vertex gluing,
since $B$ has a map with sink $a$. Again the star theorem applies.

Finally, a coordinate exclusive to $B$ is one of bits $3,\ldots,11$.
Its section away from $H$ lies in $B$ and is positive. For each contraction
and each section retaining $H$, the resulting $B'$ admits a relative
peeling to $H$. The $18$ explicit orders are stored and replayed in
\path{research/four_vertex_path_minors.json}; their completed path orders
also pass independent out-map bijectivity and nonclashing checks.
Perform such a relative peeling on the lower shore. Its deleted vertices
have no vertical mates or neighbours in the upper shore, so these are also
corners of the full class. The remaining lower path can then be deleted
in the order $a,b,c,d$. At each step its endpoint and vertical edge span
a square, or only its vertical edge remains. The upper shore is now $S$,
whose positive order was already established.

This covers all $37$ coordinates and all $111$ elementary operations.
The replay additionally supplies complete orders for $109$ of these
operations; only contraction of bits $1$ and $2$ uses the star theorem
without an explicitly stored full order. The structural argument above
covers those two as well. Thus every proper pc-minor is positive, while
$X$ is negative. Its reduction $H$ has four vertices, so its $e$-class has
exactly four edges, proving sharpness with this explicit forbidden class.
\end{proof}

This also refutes successful-root existence for the prescribed tree lift,
rather than only the arbitrary-choice rule of
Remark~\ref{rem:damp-tree-lift-cycle}. For the certified contraction map
above, the exterior labels on $a,b,c,d$ are respectively the $A$ piece,
the $B$ component at $b$, the same $B$ component at $c$, and the $P$ piece.
The quotient sink block is the singleton $\{a\}$, and no vertex is free.
The unique certified lift is cyclic; the verifier records an explicit cycle
and checks its representation map. More fundamentally, no alternative
acyclic map can exist on this expansion, by the theorem. The normalization
proposition remains valid, but its remaining existence hypothesis is false.

The example does not imply that every four-edge coordinate class has the
same three-piece presentation. There is, however, a canonical bounded
separator for the entire four-edge branch. A coordinate face here means
its intersection with the class, and may be a proper subset of its
ambient cube.

\begin{theorem}[Eight face types for a four-edge coordinate class]
\label{thm:damp-four-edge-face-types}
Let $X$ be an ample forbidden pc-minor with a coordinate $e$ having exactly
four edges. Normalize its reduction to
$H=\{0,1,3,7\}\subseteq Q_F$, where $F$ consists of the three successive
path directions, and make all other coordinates of $H$ zero. Let $D$
be the coordinate-face intersection of $X$ obtained by fixing all
coordinates outside $F\cup\{e\}$ to zero.
Then $D$ separates $X$, and $9\le |D|\le12$.
Up to reversing the path and interchanging the $e$-sections, it is one
of the following eight types. The last column bounds the number of pieces
$A_i$ in which $D$ is not gated:
\[
 D=((H\cup L)\times\{0\})\cup((H\cup R)\times\{1\}).
\]
\begin{center}
\begin{tabular}{cccc}
\toprule
$L$ & $R$ & $|D|$ & nongated pieces, at most\\
\midrule
$\{2\}$ & $\varnothing$ & $9$ & $2$\\
$\{2,6\}$ & $\varnothing$ & $10$ & $1$\\
$\{2,5\}$ & $\varnothing$ & $10$ & $1$\\
$\{2\}$ & $\{5\}$ & $10$ & $1$\\
$\{2,4,6\}$ & $\varnothing$ & $11$ & $1$\\
$\{2,4,5\}$ & $\varnothing$ & $11$ & $1$\\
$\{2,6\}$ & $\{5\}$ & $11$ & $0$\\
$\{2,4,5,6\}$ & $\varnothing$ & $12$ & $0$\\
\bottomrule
\end{tabular}
\end{center}
Every component $W_i$ of $X\setminus D$ lies in one $e$-section.
Both sections contain such components. The pieces $A_i=D\cup W_i$ are
proper positive coordinate restrictions, use disjoint exclusive blocks,
and fail to peel to $D$. If $m_t$ is the number in section $t$, then
\[
 1\le m_t\le |D_{e=t}|,\qquad
 2\le m_0+m_1\le |D|\le12.
\]
The specified coordinate determines $D$ and its pieces uniquely, apart
from the displayed normalization and permutation of the pieces.
\end{theorem}
\begin{proof}
The star-reduction theorem makes the four-vertex reduction a path rather
than a star. A tree in a partial cube uses distinct coordinates on its
edges, giving the normalization above. Write $U=\pi_eX$ and let $K$
be its intersection with the $F$-face containing $H$. Then $K$ is ample,
$H\subseteq K\subseteq Q_F$, and $D$ has contraction $K$ and reduction
$H$. In particular $|D|=|K|+4\le12$.

We first record the elementary geometry of this three-dimensional face.
The complement of $H$ in $Q_F$ is the induced geodesic path
\[
 J=2-6-4-5.
\]
For an ample $T$ with $H\subseteq T\subseteq Q_F$, complement duality
says that $Q_F\setminus T$ is ample. If nonempty it is connected, so it
is an interval of $J$. Conversely every interval of $J$ is ample, since
it is a geodesic path; its cube complement is therefore ample.
There are exactly eleven possibilities, including the empty complement.
Moreover every such $T$ peels to $H$: enlarge its complementary interval
one endpoint at a time until it is all of $J$. If the complement is
initially empty, start at any vertex of $J$. Every intermediate complement,
and hence every remaining $T$, is ample, so the successive single deletions
are corner deletions.

Every vertex outside $D$ has a nonzero coordinate outside $F\cup\{e\}$.
No $e$-edge has an endpoint outside $D$, since all its projected pairs
belong to $H$. Thus each outside component stays in one section.
Suppose section $t$ has no outside vertices. It is then an ample
three-dimensional family containing $H$, and the preceding interval
argument peels it to $H$. These relative deletions also work in $X$:
the deleted vertices have no $e$-mates, so their full neighbourhoods are
unchanged. The remaining class is a peripheral expansion of the other
section at $H$. That section is positive by forbidden-minor minimality,
and $H$ is positive, so Proposition~\ref{prop:damp-exact-expansion-product-tests}
makes the remainder positive. This contradicts nonpeelability of $X$.
Both sections therefore have outside components, and $D$ separates $X$.

Apply Proposition~\ref{prop:damp-face-separator-bound}. It gives the
canonical disjoint blocks, positivity of the pieces, and their failures
to peel to $D$. For a component in section $t$, this failure is equivalent
to the failure of its section piece $D_{e=t}\cup W_i$ to peel to $D_{e=t}$:
each vertex deleted in a relative peeling is outside $D$, and has no
vertical mate, so the corner tests in the two families coincide throughout.
The whole section $X_{e=t}$ is positive. Applying the positive part of
that proposition to its coordinate face $D_{e=t}$ gives
$m_t\le |D_{e=t}|$, and summing gives the stronger bound $m_0+m_1\le|D|$.

If $K=H$, then $H$ is a coordinate face of $U$, hence convex. Since
$U$ is positive, Theorem~\ref{thm:damp-convex-reduction-lift} would make
$X$ positive. Therefore $K\ne H$ and $|D|\ge9$.
To list $D$, its omitted set $Q_F\setminus K$ is an interval of $J$,
or empty. The components of $K\setminus H$ are consequently the prefix
and suffix of $J$ left by that interval, omitting empty ones; when the
interval is empty, all of $J$ is one component. Each component has a
single $e$-label by the component-expansion theorem
(Theorem~\ref{thm:damp-component-expansion-signs}). Conversely that theorem
makes every such labeling ample.

There are $26$ nonprism labeled faces: $4$ with one added vertex, $8$
with two, $12$ with three, and $2$ with four. Reversing $J$ and swapping
the two labels leaves respectively one, three, three, and one types,
which are exactly the table. Thus the list is exhaustive; it is a list
of necessary face types, not a claim that every type occurs in a forbidden
class. The literal $81$ assignments of absent, lower, or upper status to
the four vertices of $J$ are also checked in
\path{research/four_edge_face_types.py/json}. Finally the face and graph
components were recovered without choices, proving uniqueness.

It remains to prove the last column. Project $X$ onto $F\cup\{e\}$,
obtaining an ample family $V$ containing $D$. Its $e$-reduction is still
$H$. Indeed, a new $e$-edge over $u\notin H$ would have preimages in $X$
with the same $F$-trace $u$ and opposite $e$-bits. A geodesic between
them keeps that $F$-trace and crosses an $e$-edge of $X$, whose trace
would have to belong to $H$, a contradiction. Put $K'=\pi_e V$.
The complements $I=Q_F\setminus K$ and $I'=Q_F\setminus K'$ are
intervals of $J$, possibly empty, with $I'\subseteq I$.

The vertices of $V\setminus D$ project bijectively onto $I\setminus I'$:
no extra lift can be added over $K\setminus H$, since that would enlarge
the $e$-reduction. Each component of $K'\setminus H$ has a single
$e$-label, so the induced graphs on $V\setminus D$ and $I\setminus I'$
are isomorphic. The difference of two nested intervals has at most two
components; when $|I|\le2$, it has at most one. Since $|I|=12-|D|$,
Proposition~\ref{prop:damp-face-piece-projections} gives the first six
bounds and the bound zero for $|D|=12$.
For the seventh row, $|I|=1$ and the two components of $K\setminus H$
already have different $e$-labels. If $I'$ were empty, $K'=Q_F$ and
$K'\setminus H=J$ would be connected and require a single label.
Thus $I'=I$, $V=D$, and every piece has $D$ gated.
\end{proof}

For the $909$-vertex class, this face is the first type and has nine
vertices. Its three outside components have sizes $284,288,328$.
Their projections outside the canonical face contain respectively three,
zero, and zero vertices. Thus precisely the first of these pieces fails
to have $D$ gated. The finite projection check in
\path{research/four_edge_face_projections.py/json} verifies this statement
and all containments between the $26$ labelled nonprism faces.
The bound and recovery apply to every forbidden class with a four-edge
coordinate, whereas that example establishes realization of only one type.
The pieces still carry relative-peeling conditions. These cannot in general
be replaced by separate attainable-root and retained-edge tests, even for
the smallest face type. Indeed, take $Y_2$ from
Theorem~\ref{thm:damp-cube-boundary-retention}, with its retained square,
and attach a leaf at one square vertex in a fresh coordinate. Call the
result $T$ and the square together with this leaf $K$. Then $K$ has the
five-vertex shape $H\cup\{2\}$ in the first row of the table. The family
$T$ is positive and peels to each vertex and edge of $K$, but not to $K$:
retaining $K$ would retain the square in $Y_2$. For a square edge, delete
the leaf and use its certified retention; for the new leaf edge, peel
$Y_2$ to its attaching root. Vertex retention follows from these orders.
The peripheral expansion $(T\times\{0\})\cup(H\times\{1\})$ is a
positive one-sided piece with common face
$(K\times\{0\})\cup(H\times\{1\})$. It fails relative peeling to
that face, and its outside part is connected for the displayed input.
Its proper exclusive-coordinate images do peel to the face, by the
relative-minor assertion of that theorem. This is an admissible local
piece, not a claim that it occurs in a forbidden assembly.
Generating their
attainable USO data intrinsically, and proving compatible forward assembly,
remain necessary before this decomposition becomes a generator of even
this branch. Forbidden classes whose coordinate classes all have more than
four edges are not covered by the eight-type assertion.

The separate source question is whether a rooted factor whose every peeling
starts at its root can have another geometric corner. Low-rank extension
invariance forces every such nonroot corner to have degree at least three.
Moreover a degree-three corner cannot be adjacent to another corner: deleting
both leaves an ample class whose two individual additions have the same
support of size two, and square-pair invariance makes that class positive.
Either single deletion is then positive, contradicting the required first
vertex. These observations apply the earlier invariance results; they do not
exclude higher-degree or nonadjacent additional corners.

The twelve stored positive one-concept extensions of the $291$-concept
obstruction do not supply a counterexample: each has complete peelings starting
at two different vertices. The explicit orders are replayed by
\path{research/unique_source_controls.py} and stored in its JSON report.
This is a test of that fixed family only; the general source question remains
open. It is not a prerequisite for the joint descent, whereas intrinsic forward
generation remains open.

\begin{theorem}[Endpoint transfer for a partially replicated tip]
\label{thm:damp-partial-tip-endpoint-transfer}
Let \((B,0)\) be a rooted factor using every coordinate of \(F\).
Let \(\alpha,\delta,\gamma\in Q_F\setminus B\) be distinct
ample one-concept additions with pairwise distinct acquired supports
\(P_\alpha,P_\delta,P_\gamma\).  Suppose that
\(A=B\cup\{\alpha\}\) and \(D=B\cup\{\delta\}\) both peel
to zero.  For two fresh coordinates \(q,r\), put
\[
 L=\{01,10,11\},\qquad
 S=\{(0,00)\}\cup(A\times\{01\})
       \cup(D\times\{10\})\cup(B\times\{11\}),
\]
and, for \(R\subseteq L\), define
\[
 S_R=S\cup(\{\gamma\}\times R),\qquad C=B\cup\{\gamma\}.
\]
With \(\operatorname{Sh}(\varnothing)=\varnothing\), one has
\begin{equation}
 \operatorname{Sh}(S_R)=\operatorname{Sh}(S)
   \mathbin{\dot\cup}\{P_\gamma\cup J:J\in\operatorname{Sh}(R)\}.
 \label{eq:damp-partial-tip-shadow}
\end{equation}
Thus \(S_R\) is ample exactly when \(R\ne\{01,10\}\).
For every such \(R\), the family \(S_R\) is ordinarily peelable,
all its proper root-preserving pc-minors peel to their roots, and
\begin{equation}
 S_R\text{ peels to }(0,00)
 \quad\Longleftrightarrow\quad
 |R|\ge2\ \text{ and }\ C\text{ peels to }0.
 \label{eq:damp-partial-tip-endpoint}
\end{equation}
Consequently \((S_R,(0,00))\) is a rooted factor precisely when
the right-hand side is false.  The root is always a degree-two corner.
\end{theorem}

\begin{proof}
The compatible-extension union formula
\eqref{eq:damp-replicated-repair-union-shadow} applies to every
subset of the three old-coordinate tips.  They are pairwise
nonadjacent by Lemma~\ref{lem:damp-compatible-repairs-do-not-collide}.
For old supports, adding a nonempty row acquires exactly \(P_\gamma\).
A support containing \(q\) and not \(r\) is shattered when its
old part is shattered by both projected \(q\)-sections.  The only
possible new old part is \(P_\gamma\), and it occurs on both
sides exactly when \(R\) shatters \(\{q\}\).  Interchanging
\(q,r\) gives the other one-coordinate term.  A support using both
new coordinates has empty old part, since the \(00\) section is
\(\{0\}\); it is already shattered by \(S\).  The row \(R\)
never shatters \(\{q,r\}\).  These observations prove
\eqref{eq:damp-partial-tip-shadow}, including the empty row.

The shadow of a singleton has size one; that of either edge of
\(L\) has size two; and both \(L\) and its diagonal pair
\(\{01,10\}\) have the three supports
\(\varnothing,\{q\},\{r\}\).  Since \(|S_R|=|S|+|R|\),
the Sandwich equality proves the ampleness assertion.  Each allowable
row can be built from the empty row by ample one-concept additions.
It can also be enlarged to \(L\) by such additions.  For an edge,
add either endpoint first; for \(L\), use a singleton, an incident
edge, and then its third vertex.

We prove the endpoint equivalence next.  If \(|R|\le1\), the
intersections of the \(11\) section with each adjacent section
are both exactly \(B\).  Suppose that a peeling ends at \((0,00)\),
and consider the first deletion among \((0,01),(0,10),(0,11)\).
The full square at old zero is still present.  As in the proof of
Proposition~\ref{prop:damp-paired-root-transfer}, the selected corner
cannot have any old-coordinate neighbor: its corner cube would require
that neighbor in the singleton \(00\) section.  Its current ample section is connected,
and is therefore \(\{0\}\).  Choose the overlap along its edge to
\(11\), or either edge if the selected vertex is \((0,11)\).
Initially this overlap is \(B\).  Recording the first destruction
of each full edge fibre gives corner deletions of the overlap; the
fibre of zero remains intact until the selected step.  This would
peel \(B\) to zero, a contradiction.  This argument uses the two
overlap equalities and does not require section containment.

The full-row family \(S_L\) is the paired construction over \(C\),
with sides \(C\cup\{\alpha\}\) and \(C\cup\{\delta\}\).
Proposition~\ref{prop:damp-paired-root-transfer} gives its endpoint
equivalence with \(C\).  If \(C\) fails at zero, no allowable
\(S_R\) can peel to its root: prefacing such a peeling by the
reverse of an ample addition chain from \(S_R\) to \(S_L\)
would give a peeling of \(S_L\) to its root.

It remains to prove sufficiency for an edge row when \(C\) peels
to zero.  Let \(R=\{10,11\}\).  Delete \((\delta,10)\), then
use a peeling of \(C\) to zero to delete all nonzero words in its
\(10\) copy.  Every corner cube lifts to the untouched \(11\)
copy of \(C\), and no deleted word has a \(00\) mate.
Delete \((\gamma,11)\); it is now a corner with no mate in
either adjacent section.  Next delete \((0,10)\), whose remaining
incident cube is the square at old zero.  The \(11\) section is
now \(B\).  Peel it by an ordinary peeling of \(B\); its corner
cubes lift into \(A\) at \(01\).  What remains is \(A\) at
\(01\) with the root attached at old zero.  Peel \(A\) to zero
and finish at \((0,00)\).  For \(R=\{01,11\}\), interchange
\(q,r\), \(A,D\), and \(\alpha,\delta\) throughout this
construction.  The required endpoint order is then that of \(D\).
This proves \eqref{eq:damp-partial-tip-endpoint} in every case.

The initial \(S\) is a rooted factor by
Proposition~\ref{prop:damp-paired-root-transfer}.  Every allowable
\(S_R\) lies above it in an ample one-concept chain on the same
active coordinates.  Reversing the additions proves ordinary
peelability.  If the endpoint test is false,
Lemma~\ref{lem:damp-rooted-ample-additions} proves rooted-factor
status; if it is true, root-preserving pc-minor closure supplies
the stated positive orders in all proper rooted minors.  Finally,
the root's two incident edges and their square are unchanged.
\end{proof}

\begin{remark}[Two individually ineffective additions restore the endpoint]
\label{rem:damp-cooperative-rooted-tip-pair}
Take the \(289\)-concept core \(B\) and tips \(\alpha=128\),
\(\delta=290\) from Remark~\ref{ex:damp-terminal-cut-vertex-1158}.
A third ample addition \(\gamma=604\), corresponding to the
original tip \(577\), acquires support
\(\{5,6,9,10\}\), with bitmask \(1632\), and peels to zero.
Its endpoint order is stored in
\path{computations/round793_third_endpoint_repair.json}.
For the resulting \(870\)-concept rooted factor \(S\), put
\[
 u=(\gamma,10),\qquad v=(\gamma,11).
\]
Both \(S\cup\{u\}\) and \(S\cup\{v\}\) are rooted factors
with \(871\) concepts.  Yet \(S\cup\{u,v\}\) peels to the
original root.  In the encoding \((x,qr)\mapsto4x+q+2r\),
the new words are \(2417,2419\).  Each individual addition acquires
support \(6528\), of size four; adding the second word after the
first acquires \(6530\), of size five.  The pair is adjacent,
and the two initial supports coincide.  This does not establish
cooperation for additions with distinct initial supports.

The checker \path{verify_damp_partial_tip_endpoints.py} tests all
eight occurrence rows for this repairing tip and for the nonrepairing
tip \(64\) from Remark~\ref{rem:damp-noncontainment-872}.
It checks the shadow formula, the two nonample diagonal rows, complete
negative graphs for the other negative cases, and explicit endpoint
orders for both edge rows and the full row at \(604\).
It also replays all \(28\) rooted elementary orders of \(S\)
and all \(84\) orders of its three \(871\)-concept extensions.
The report is \path{damp_partial_tip_endpoints_verification.json}.
Thus individual endpoint statuses do not suffice for simultaneous
ample additions, even over a rooted factor.  The theorem supplies
the exact occurrence-row condition in this construction; the intrinsic
criterion for whether \(B\cup\{\gamma\}\) repairs zero remains open.
\end{remark}


\begin{proposition}[The opposite section need not be contained in both sides]
\label{prop:damp-opposite-section-noncontainment}
Let \((B,0)\) be a rooted factor using every coordinate of \(F\).
Let \(x,\alpha,\delta\in Q_F\setminus B\) be distinct concepts
whose one-concept extensions are ample and acquire pairwise distinct
shattered supports \(P_x,P_\alpha,P_\delta\).  Suppose that
\(B\cup\{x\}\) does not peel to zero, whereas
\(B\cup\{\alpha\}\) and \(B\cup\{\delta\}\) do.
Then, on two fresh coordinates \(q,r\), the family
\begin{equation}
 \begin{split}
 Z={}&\{(0,00)\}
 \cup((B\cup\{\alpha\})\times\{01\})\\
 &\cup((B\cup\{x,\delta\})\times\{10\})
 \cup((B\cup\{x\})\times\{11\})
 \end{split}
 \label{eq:damp-noncontained-opposite-section}
\end{equation}
is a rooted factor with a degree-two corner root \((0,00)\).
Its opposite section \(C=B\cup\{x\}\) satisfies
\[
 C\setminus A=\{x\},\qquad C\subseteq D,\qquad
 A\cap D=B,\qquad A\setminus C=\{\alpha\},\qquad
 D\setminus C=\{\delta\},
\]
where \(A,D\) denote its \(01,10\) sections.
Deleting \((x,11)\) and then \((x,10)\) preserves rooted-factor
status and leaves the paired-extension construction over \(B\).
\end{proposition}

\begin{proof}
Apply Theorem~\ref{thm:damp-partial-tip-endpoint-transfer} with
\(\gamma=x\) and \(R=\{10,11\}\).  Since \(B\cup\{x\}\)
does not peel to zero, the theorem makes this family, its singleton-row
predecessor, and its empty-row predecessor rooted factors.  Reversing
the two row additions gives the stated corner deletions.  The section
differences follow directly from distinctness of the three tips and
their exclusion from \(B\).
\end{proof}

\begin{remark}[A noncontained section in a root-protected graph three-core]
\label{rem:damp-noncontainment-872}
Take the translated \(289\)-concept factor of
Remark~\ref{ex:damp-terminal-cut-vertex-1158} as \(B\), and use
\[
 x=64,\qquad\alpha=128,\qquad\delta=290,
 \qquad (P_x,P_\alpha,P_\delta)=(70,390,480)
\]
in bitmask notation.  The new support \(70\) is \(\{1,2,6\}\).
The family \(B\cup\{64\}\) has only the corner \(64\);
deleting it leaves \(B\), whose only corner is zero.  It therefore
cannot peel to zero.  The two repairing endpoints were verified in
that earlier remark, so Proposition~\ref{prop:damp-opposite-section-noncontainment}
applies.

The resulting \(Z\) has \(872\) concepts, dimension \(14\), and
VC dimension four.  Encode \((v,qr)\) by \(4v+q+2r\).
Its corners are precisely
\[
 0,\quad257,\quad259,\quad514,\quad1161.
\]
The root has degree two and the four other corners have degree four.
Every nonroot vertex has degree at least four.  Thus the graph is
unchanged by repeatedly deleting nonroot vertices of degree at most
two, with the root protected.  Nevertheless its \(11\) section
contains the word \(64\), absent from its \(01\) section.
Rooted minimality therefore does not force the containment required
in Remark~\ref{rem:damp-paired-root-recovery}, even after this graph
reduction and even when both side differences from the opposite
section are singletons.

The fixed replay \path{verify_damp_opposite_section_noncontainment.py}
checks all eight compatible tip unions and the ample chain of orders
\(870,871,872,873\).  Its successive shadow gains have bitmasks
\(280,282,281\).  Complete root-protected negative graphs have
\(4,8,16,28\) states, respectively.  It also checks ordinary peelings,
all \(24\) rooted elementary orders of \(B\), and all \(28\) of
\(Z\).  The data are stored in
\path{damp_opposite_section_noncontainment_verification.json}.
The two preserving deletions \(259,257\) return \(Z\) to the
known \(870\)-concept factor.  Thus this example refutes automatic
section containment but does not establish another terminal assembly
type. In fact the full sectionwise endpoint law
\eqref{eq:damp-paired-all-endpoints} still holds for this $Z$, despite
its failure of containment. To verify this, let $c$ be the stored ordinary
order of $B$. The following prefixes, followed by the indicated complete
copies of $c$, give relative orders:
\[
\begin{array}{c|c|c}
\text{prefix}&\text{layers cleared in order}&\text{retained family}\\ \hline
259,257,1161,0&4B+1,\ 4B+3&4A+2\\
259,514,0&4B+2,\ 4B+3&4D+1\\
259,257,514,1161,0&4B+1,\ 4B+2&4B+3
\end{array}
\]
Here $4F+t=\{4x+t:x\in F\}$. Each order is checked corner by corner
in \path{research/noncontained_endpoint_transfer.py/json}. The first
two orders prove sufficiency for every endpoint attainable in its side
section. For the opposite section $C=B\cup\{64\}$, every peeling
starts at its unique corner $64$, so its attainable endpoints are exactly
those of $B$. The third order therefore proves sufficiency there too.
Necessity follows in each case by coordinate-face restriction.

Every unattainable vertex other than zero consequently fails in its own
proper section; zero is the only critical root of $Z$. Thus the endpoint
law and critical-root uniqueness do not imply the paired presentation's
containment hypotheses. At the critical root, the two incident coordinates
fix its four sections, so choosing another presentation cannot repair this
failure. Intrinsic rooted-factor and endpoint-repair admissibility, and
the full forbidden-pc-minor classification, remain open.
\end{remark}

\begin{proposition}[An antipodal corner pair does not restore the endpoint]
\label{prop:damp-antipodal-tip-nonrepair}
Let \((B,a)\) be a rooted factor using every coordinate of \(F\),
and suppose that \(a\) is the only corner of \(B\).
Let \(H=B\cup\{c\}\subseteq Q_F\) be an ample one-concept
extension, with acquired support \(P\) and completed cube
\(Q=\{c^U:U\subseteq P\}\).  Define
\begin{equation}
 E=\{v\in Q\setminus\{a,c\}:
       v^{\{e\}}\notin H\text{ for every }e\in F\setminus P\}.
 \label{eq:damp-tip-available-old-corners}
\end{equation}
Then \(E=\operatorname{Cor}(H)\setminus\{a,c\}\), and every
member of \(E\) has Hamming distance at least two from \(c\).
Every peeling of \(H\) ending at \(a\) must start in \(E\).
If \(E=\varnothing\) or \(E=\{c^P\}\), then \((H,a)\)
is another rooted factor.  In either case, every ordinary peeling
of \(H\) deletes \(a\) among its first two concepts.
Consequently, an endpoint-restoring addition must have some \(v\in E\)
with \(2\le d_H(c,v)<|P|\).
\end{proposition}

\begin{proof}
The added concept is a corner whose unique maximal cube is \(Q\).
Any old concept that becomes a corner must have a corner cube
containing \(c\), since otherwise that cube already made it a
corner of \(B\).  Its cube is therefore \(Q\).
A vertex \(v\in Q\) is a corner exactly when it has no neighbor
outside \(Q\).  Indeed, all directions of \(P\) are incident
with \(v\).  An additional direction \(e\) in its corner cube
would force \(c^{\{e\}}\in H\), contrary to the incident support
of \(c\).  With no additional direction, \(Q\) is its full
incident cube.  Since \(a\) is the only old corner, this proves
the formula for \(E\).

If \(v\in E\) were adjacent to \(c\), deleting \(c\) would
leave \(v\) as a corner of the opposite face of \(Q\).  This
would be a corner of \(B\) different from \(a\), a contradiction.
A peeling ending at \(a\) cannot start at \(a\); it cannot start
at \(c\), since that leaves \(B\), which fails at \(a\).
Its first concept must therefore belong to \(E\).

If \(E\) is empty, this already excludes the endpoint.
Suppose now that \(E=\{d\}\), where \(d=c^P\).
The distance assertion gives \(|P|\ge2\).  The only possible
nonroot first deletions are \(c,d\).  Deleting \(c\) leaves
only the corner \(a\).  We claim that deleting \(d\) likewise
leaves no corner other than possibly \(a\).
The vertex \(c\) retains all directions of \(P\), since it
is not adjacent to \(d\), but their cube is missing \(d\).
Thus \(c\) is no longer a corner.  Any newly created corner
must be a neighbor of \(d\): deletion of a nonneighbor changes
no incident direction and cannot fill a missing cube.

All neighbors of \(d\) lie in \(Q\).  Write one different
from \(a\) as \(w=d^{\{e\}}\), with \(e\in P\).
It is neither \(c\) nor \(d\), and is not in \(E\), so it
has an outside neighbor \(w^{\{f\}}\in H\), with \(f\notin P\).
After deleting \(d\), it retains that neighbor and every direction
in \(P\setminus\{e\}\).  Because \(c,d\) are antipodes,
\(c=w^{P\setminus\{e\}}\).  A corner cube on these retained
directions would therefore contain \(c^{\{f\}}\), which is
absent from \(H\).  Hence \(w\) cannot be a corner.  This
proves the claim, including every possible new nonroot corner.

In both cases, a first deletion different from \(a\) forces
\(a\) to be the second deletion in any ordinary peeling.
The endpoint \(a\) is therefore impossible.  The family \(H\)
does peel ordinarily, by deleting \(c\) and then peeling \(B\).
Lemma~\ref{lem:damp-rooted-ample-additions} now makes \((H,a)\)
a rooted factor.  All tests in \eqref{eq:damp-tip-available-old-corners}
are incident-edge tests on the completed cube; they do not ask for
an endpoint order.
Finally, if \(H\) peels to \(a\), the set \(E\) is neither empty
nor the singleton antipode.  Its members all have distance at least two
from \(c\), and only the antipode has distance \(|P|\).  This proves
the last assertion.
\end{proof}

\begin{remark}[Creating an old nonroot corner is insufficient]
\label{rem:damp-rooted-antipodal-290}
Use the normalized \(289\)-concept rooted factor \(B\) from
Remark~\ref{ex:damp-terminal-cut-vertex-1158}.  Its only corner
is zero, of degree three.  The addition \(c=167\), corresponding
to the original word \(186\), is ample and acquires
\[
 P=\{0,1,6,8\},\qquad \text{with bitmask }323.
\]
The completed four-cube has exactly two vertices without outside
neighbors: \(c=167\) and its antipode \(d=484\).  The latter
belongs to \(B\) and becomes a new corner.  All other fourteen
cube vertices have outside neighbors, and zero lies outside this
cube.  Thus \(E=\{484\}\), and the proposition applies.
The \(290\)-concept extension \(H\) is a rooted factor in
dimension twelve and VC dimension four, with corners
\(0,167,484\).  Deleting either nonroot corner leaves zero as
the only corner.  In particular, exposing an old nonroot corner
is necessary but not sufficient to restore the endpoint.

Compare the restoring tip \(604\) from
Remark~\ref{rem:damp-cooperative-rooted-tip-pair}.  Both tips have
support size four, Hamming distance five from zero, and two
tip--root disagreements inside their acquired support.  A coordinate
permutation identifies their punctured cubes together with the marked
root and missing tip.  Their exterior attachments differ: the
restoring completion has three available old nonroot corners,
\(1084,1116,1148\), whereas the nonrestoring completion has only
the antipode \(484\).  Thus even the marked punctured cube does
not determine the endpoint without its exterior attachments.

The checker \path{verify_damp_rooted_clause_pivots.py} reconstructs
both fixed extensions, verifies all fourteen exterior-neighbor
witnesses and the marked-cube permutation, and replays the positive
endpoint order.  It checks the complete three-state root-protected
negative graph of \(H\), an ordinary peeling, and all \(24\)
rooted elementary minor orders of \(H\), as well as the \(24\)
core orders.  The report is
\path{damp_rooted_clause_pivots_verification.json}.
This gives an intrinsic sufficient condition for a nonrestoring
addition over a known unique-corner rooted factor.  It does not
characterize all endpoint repairs or settle intrinsic factor
admissibility or the full classification.
\end{remark}

\begin{theorem}[A single clause can change an arbitrarily remote endpoint]
\label{thm:damp-remote-rooted-clause-pivot}
For every \(j\ge0\), there are two ordinarily corner-peelable ample
families \(B_j,K_j\subseteq Q_{12+2j}\), both containing zero, with
the following properties.
\begin{enumerate}
 \item \((B_j,0)\) is a rooted factor, whereas \(K_j\) peels to zero.
       Every proper root-preserving pc-minor of either family peels to
       its root.
 \item Their shattered-set complexes coincide.  Their signed
       minimal-nonface presentations differ on exactly one four-element
       support, whose omitted traces are antipodal.
 \item They contain exactly the same concepts at Hamming distance at
       most \(j+4\) from zero.  The nearest difference has distance
       \(j+5\).
\end{enumerate}
Their common order, dimension, and VC dimension are
\begin{equation}
 |B_j|=|K_j|=\frac{581\cdot3^j-3}{2},\qquad
 \dim B_j=12+2j,\qquad \operatorname{VCdim}(B_j)=3+j.
 \label{eq:damp-remote-clause-tower-size}
\end{equation}
Their roots have degree three for \(j=0\), and degree two for
\(j\ge1\).  In particular, no fixed-radius neighborhood of the
root, even together with the entire shattered-set complex, determines
prescribed-root peelability or rooted-factor status among these families.
\end{theorem}

\begin{proof}
We give explicit recursive families.  First we show that the paired
square construction transports its own two endpoint repairs, so that
it can be iterated without further endpoint searches.

Let \(C\) be ordinarily peelable, and suppose that
\(A=C\cup\{\alpha\}\), \(D=C\cup\{\delta\}\) are compatible
ample additions with distinct acquired supports \(P_\alpha,P_\delta\),
both peeling to zero.  Write \(\mathcal S(C;\alpha,\delta)\) for
the square construction \eqref{eq:damp-paired-rooted-factor}.
Its two missing concepts \((\alpha,11)\), \((\delta,11)\)
are ample one-concept additions with acquired supports
\begin{equation}
 P_\alpha\cup\{q\},\qquad P_\delta\cup\{r\},
 \label{eq:damp-square-repair-propagation}
\end{equation}
respectively.  For the first, its old punctured cube is present at
\(11\), its \(q\)-mate and that entire cube are present at
\(01\), and it has no \(r\)-mate.  Its acquired support is
absent from \eqref{eq:damp-paired-rooted-shadow}, so the local
one-concept test proves ampleness.  Interchanging \(q,r\) and
\(\alpha,\delta\) proves the second assertion.  These two
supports are distinct, so the compatible-extension union formula
also applies to their joint addition.

Both additions repair the square root.  After adding \((\alpha,11)\),
peel the \(01\) copy of \(A\) down to zero, using its endpoint
order; every deleted corner cube lifts into the \(11\) copy of
\(A\).  Delete \((\alpha,11)\) and then \((0,01)\).
The remaining \(11\) section is \(C\); peel it ordinarily,
lifting its corner cubes into \(D\) at \(10\).
Finally peel \(D\) to zero and delete the root last.
For the addition \((\delta,11)\), interchange the two new
coordinates and the two side extensions in this sequence.
This proves the repair propagation independently of whether \(C\)
itself peels to zero.

Let \(B_0\) be the normalized \(289\)-concept factor used in
Remark~\ref{rem:damp-rooted-antipodal-290}, and set
\[
 K_0=(B_0\setminus\{1084\})\cup\{604\},\qquad
 \alpha_0=128,\qquad\delta_0=290.
\]
The verified peeling of \(B_0\cup\{604\}\) starts at \(1084\).
Its suffix peels \(K_0\) to zero.  That first corner has support
\(I_0=\{5,6,9,10\}\), with bitmask \(1632\), so its deletion
removes the newly acquired support and gives
\(\operatorname{Sh}(K_0)=\operatorname{Sh}(B_0)\).
By Proposition~\ref{prop:damp-signed-clause-exchange}, only the
omitted trace on \(I_0\) changes, from bitmask \(576\) to
\(1056\).  Both \(604\) and \(1084\) have weight five.
The two additions by \(\alpha_0,\delta_0\) are compatible
and repairing for both bases; these finite input conditions and
all four endpoint orders are checked in the replay below.

Inductively define
\begin{equation}
 \begin{split}
 B_{j+1}&=\mathcal S(B_j;\alpha_j,\delta_j),\qquad
 K_{j+1}=\mathcal S(K_j;\alpha_j,\delta_j),\\
 \alpha_{j+1}&=(\alpha_j,11),\qquad
 \delta_{j+1}=(\delta_j,11).
 \end{split}
 \label{eq:damp-remote-clause-tower}
\end{equation}
The propagation just proved maintains the two compatible repairs
on both sides.  Proposition~\ref{prop:damp-paired-root-transfer}
therefore makes every \((B_j,0)\) a rooted factor and every
\(K_j\) root-peelable.  The two cores and their repairs have
identical shadows at every step, by induction using
\eqref{eq:damp-paired-rooted-shadow} and
\eqref{eq:damp-square-repair-propagation}.  Ordinary peelability and
the assertions about proper rooted minors follow from the same
transfer theorem and root-preserving pc-minor closure.
The order recurrence is \(n_{j+1}=3n_j+3\), with \(n_0=289\).
Each step adds two active coordinates.  The initial VC dimension
is three and the initial repair supports have size four; the two
shadow formulas increase these quantities by one at each step.
This proves \eqref{eq:damp-remote-clause-tower-size} and the root degrees.

Put \(L=\{01,10,11\}\).  Since the repair tips are the same
for both families and the differing old words are nonzero, their
differences are exactly
\[
 B_j\setminus K_j=\{1084\}\times L^j,\qquad
 K_j\setminus B_j=\{604\}\times L^j.
\]
Each word of \(L\) has weight at least one, and weight one is
attained.  This proves the asserted agreement through radius \(j+4\)
and the first difference at radius \(j+5\).

The support \(I_0\), in the original twelve coordinates, remains
a minimal nonface.  For \(j\ge1\), the projection of \(B_j\)
onto those coordinates is \(B_0\cup\{\alpha_0,\delta_0\}\),
and the corresponding projection of \(K_j\) uses \(K_0\).
The two compatible repairs do not fill \(I_0\), so its omitted
traces remain \(576\) and \(1056\).
Every other minimal nonface \(M\) misses some coordinate \(e\in I_0\).
The \(e\)-neighbor of each differing word belongs to both families:
the original neighbors of \(604\) and \(1084\) lie in their
common punctured four-cube and are nonzero, so all their \(L^j\)
copies persist.  That neighbor has the same \(M\)-trace as the
differing word.  Thus the two \(M\)-projections coincide.
Only the omitted trace on \(I_0\) changes, proving item~2.

Given any fixed radius, choose \(j\) large enough that it is at
most \(j+4\).  The two inputs then have the same root neighborhood
and the same full shadow but opposite endpoint status.  This proves
the final assertion; no bound on the radius was assumed in advance.
\end{proof}

The replay \path{verify_damp_rooted_clause_pivots.py} checks the
first three levels, of orders \(289,870,2613\), including the
common shadows, all signed minimal-nonface traces, explicit propagated
repair orders, and the positive endpoint orders.  It checks all
\(24,28,32\) rooted elementary orders of the negative families.
The first differences occur at distances \(5,6,7\), respectively.
The uniform proof establishes all levels; the finite replay is stored
in \path{damp_rooted_clause_pivots_verification.json}.
This rules out fixed-radius root tests, even augmented by the shadow.
It does not exclude a global criterion using the signed clauses and
their attachment incidences, which remains the intrinsic admissibility
problem.











\begin{corollary}[A remote clause change can remove a forbidden pc-minor]
\label{cor:damp-remote-forbidden-clause-pivot}
For every \(j\ge0\), there are two two-connected ample families
\(X_j,Y_j\subseteq Q_{24+2j}\), containing zero, such that
\(X_j\in\Obs(\Damp)\) and \(Y_j\in\Damp\).
They have the same shattered-set complex, and their signed
minimal-nonface presentations differ on exactly one four-element
support, with antipodal omitted traces.  They agree on every concept
at Hamming distance at most \(j+4\) from zero, and their nearest
difference has distance \(j+5\).  Their common order and VC dimension are
\[
 |X_j|=|Y_j|=\frac{581\cdot3^j+575}{2},\qquad
 \operatorname{VCdim}(X_j)=\operatorname{VCdim}(Y_j)=3+j.
\]
Every proper pc-minor of either family belongs to \(\Damp\).
Thus, even within this two-connected class with all proper minors
peelable, a fixed-radius neighborhood of a marked vertex together
with the full shadow does not determine forbidden-minor membership.
\end{corollary}

\begin{proof}
Take \(B_j,K_j\) from Theorem~\ref{thm:damp-remote-rooted-clause-pivot}.
Let \(T\) be a copy of \(B_0\) on twelve fresh coordinates, with
root zero, and first form
\[
 X'_j=B_j\vee T,\qquad Y'_j=K_j\vee T.
\]
Theorem~\ref{thm:damp-cut-vertex-obstructions} gives
\(X'_j\in\Obs(\Damp)\).  Since \(K_j\) peels to zero and
\(T\) peels ordinarily, Proposition~\ref{prop:damp-vertex-gluing}
gives \(Y'_j\in\Damp\): peel \(K_j\) down to zero and then
peel \(T\) ordinarily.

Write \(F_j,G\) for the two coordinate blocks.  The shadow of each
union is the union of its two block shadows.  Since every coordinate
is active, its minimal nonfaces are exactly the block minimal nonfaces
and the pairs \(\{e,f\}\), where \(e\in F_j\), \(f\in G\).
Every such mixed pair omits \(11\).  Old projections are unchanged,
because the zero trace supplied by the other block was already present.
Consequently the only changed clause is the old four-element clause;
all mixed signed clauses coincide.  The concept differences are also
unchanged by gluing.  The common order at this stage is
\((581\cdot3^j+573)/2\).

Choose a common root direction \(e\) in \(B_j,K_j\) and a root
direction \(f\) in \(T\).  Such an \(e\) exists because the
two root neighborhoods coincide.  Complete the same square in both:
\[
 X_j=X'_j\cup\{0^{\{e,f\}}\},\qquad
 Y_j=Y'_j\cup\{0^{\{e,f\}}\}.
\]
Theorem~\ref{thm:damp-square-completed-rooted-blocks} makes
\(X_j\) a two-connected ample forbidden pc-minor.
The added concept in \(Y_j\) is a degree-two corner acquiring
\(\{e,f\}\).  Deleting it and peeling \(Y'_j\) proves
\(Y_j\in\Damp\).

For completeness, \(Y_j\) is also two-connected.  The finite input
\(K_0\) is two-connected, as checked in the replay below.  For
\(j\ge1\), the three copies of \(K_{j-1}\) in the paired square
induce its Cartesian product with the path \(\{01,11,10\}\).
A product of two nontrivial connected graphs is two-connected:
after deleting \((v,t)\), choose another second-coordinate vertex
\(t'\).  A surviving \((w,s)\) with \(w\ne v\) reaches the
intact \(t'\)-layer along its column.  If \(w=v\), then
\(s\ne t\), so the intact \(s\)-layer first reaches a different
column and then the \(t'\)-layer.  Thus the whole surviving graph
is connected.  Adding the square root, which has two distinct neighbors,
and each repairing tip, which has at least four neighbors in that
product, preserves two-connectivity.  Hence every \(K_j\) is
two-connected.  The graph of \(T\) is two-connected by the rooted-factor
theorem.  After deleting the common root in \(Y_j\), the two blocks
remain connected and the new square tip joins them.  Deleting any other
old vertex leaves the blocks connected through zero and the square tip
with a surviving neighbor.  Deleting the tip leaves \(Y'_j\) connected.
These cases prove two-connectivity of \(Y_j\).

The completion adds the same support \(\{e,f\}\) to both shadows
and removes precisely its common mixed clause.  No larger mixed support
can become a minimal nonface: any such support contains another mixed
pair, which is still a nonface.  The single changed four-element clause
therefore persists.  Both concept sets gain the same word, so the
distance assertions persist as well.  The shadow and order formulas
give the stated parameters.  All proper minors of \(X_j\) peel
by forbidden-minor status; those of \(Y_j\) peel by pc-minor closure.
\end{proof}

The replay \path{verify_damp_forbidden_clause_pivots.py} checks
\(j=0,1\), including the intermediate pairs of orders \(577,1158\)
and their two-connected completions of orders \(578,1159\).
It verifies both full signed presentations, all \(600\) elementary
minor orders across the eight families, four positive peeling orders,
the complete negative graphs, and two-connectivity of both completed
families in each pair.  Its report is
\path{damp_forbidden_clause_pivots_verification.json}.
The negative intermediate families are the previously constructed
double and terminal cut-vertex example.  The completions have a
preserving degree-two corner deletion, so no new terminal type is
asserted.  This gives opposite-status inputs for a global signed-clause
criterion; it does not supply that criterion or the full classification.


\begin{proposition}[Geometric realizations allow every final concept]
\label{prop:damp-realizable-prescribed-endpoint}
Call an ample family \(H\subseteq Q_F\) \emph{geometrically realizable}
if it is the set of chamber sign vectors of finitely many affine
hyperplanes, indexed by \(F\), on a nonempty open convex polyhedron.
This is realizability of its associated complex of oriented matroids.
If \(H\) is geometrically realizable, then it peels to every
prescribed \(a\in H\).  In particular, no rooted factor in
Theorem~\ref{thm:damp-cut-vertex-obstructions} is geometrically realizable.
\end{proposition}

\begin{proof}
We keep a point of the prescribed chamber while applying the geometric
corner sweep of Knauer--Marc.  The local step in their proof of
\cite[Proposition~5.5]{KnauerMarc23corners} is the following: translating
a generic bounding halfspace into a realizing polyhedron, through its
first nontrivial cell event, removes a COM corner and leaves a realizable
COM.  Here a COM corner may in general contain several topes.  For a
lopsided system it consists of a single ordinary corner, by
\cite[Proposition~5.3]{KnauerMarc23corners}.  Consequently, in the ample
case this step deletes one corner and leaves an ample geometric
realization.  We explain why the halfspace can always be chosen to
preserve the prescribed chamber; this is an additional conclusion from
the sweeping proof, rather than the statement of Proposition~5.5.

Let \(P\) realize \(H\), and fix a point \(p\in P\) with sign
vector \(a\).  We may make \(P\) bounded without changing its COM:
choose one witness for each realized sign vector, including those with
zero entries, and intersect with a sufficiently large open polytope
containing all these finitely many witnesses and \(p\).  Bounding
facets can be made generic while retaining these witnesses and keeping
the new polyhedron inside the old one: first move the facets a little
inward, and then perturb them by less than the resulting margin.
Thus no sign vector is added or lost.  This is also the generic-position
choice used in the cited sweep.

If \(H=\{a\}\), there is nothing to prove.  Otherwise choose
\(b\in H\setminus\{a\}\) and a coordinate on which \(a\) and
\(b\) differ.  Let \(\ell\) be the affine functional of that
hyperplane, oriented and scaled so that \(\ell(p)=1\); the chamber
of \(b\) lies in \(\ell<0\).  On the compact closure of \(P\),
choose a generic affine functional \(h\) with
\[
 \sup_{x\in\overline P}|h(x)-\ell(x)|<\tfrac14.
\]
Such choices are available arbitrarily close to \(\ell\), because
only finitely many nongeneric incidences must be avoided.  In particular,
\(h(p)>3/4\), whereas \(h(x)<1/4\) whenever
\(x\in P\) and \(\ell(x)\le0\).

Introduce the halfspace \(h>s\), initially with
\(s<\min_{\overline P}h\), so that it is redundant, and translate
its boundary by increasing \(s\) toward \(1/2\).  The point \(p\)
stays strictly inside throughout.  By the final position the chamber
of \(b\) has disappeared, so there is a first nontrivial event before
that position.  Apply the local sweeping step and stop just after this
event.  It deletes one corner different from \(a\), and the remaining
ample class is realized on a bounded open polyhedron still containing
\(p\).  Repeating this argument inductively on the number of concepts
peels \(H\) to \(a\).  A rooted factor fails to peel to its specified
root, so the last assertion follows by contraposition.
\end{proof}

Geometric realizability is inherited by coordinate faces: intersect the
realizing polyhedron with the open sign halfspaces of the fixed
coordinates, then forget their hyperplanes.  Therefore both families
\(X_j,Y_j\) of Corollary~\ref{cor:damp-remote-forbidden-clause-pivot},
and both intermediate families \(X'_j,Y'_j\) in its proof, fail
geometric realizability.  Indeed, fixing every coordinate of \(F_j\)
to zero recovers exactly the rooted factor \(T=B_0\); the added square
tip has a nonzero \(F_j\)-coordinate and is removed by this restriction.
In particular, the \(289\)-concept class \(B_0\), and the ordinarily
peelable classes \(Y'_j,Y_j\), show that ordinary peelability does not
imply geometric realizability.  The latter property cannot separate
these opposite-status pairs.

\begin{remark}[The affine oriented-matroid restriction is narrower]
\label{rem:damp-affine-om-shadow}
An affine oriented-matroid tope class is the projection of the
\(g=+\) topes of an oriented matroid onto its other coordinates,
where \(g\) is a nonloop.  Its shattered-set complex must be a matroid
independence complex.  To see this, let \(\mathcal T\) be the full
antipodal tope set and let \(M\) be the underlying matroid.  A set
\(S\) is shattered by \(\mathcal T\) exactly when it is independent
in \(M\): restriction to an independent set is a free oriented
matroid and has every sign vector as a tope; a dependent set contains
a signed circuit, whose sign pattern cannot be the restriction of a
tope by circuit--covector orthogonality.  This is the restriction form
of the rank--VC correspondence; see also
\cite[Claim~2.2]{Marc24compression}.  Antipodality gives, for the affine
class \(C\),
\[
 I\in\operatorname{Sh}(C)
 \quad\Longleftrightarrow\quad
 I\cup\{g\}\in\operatorname{Sh}(\mathcal T)
 \quad\Longleftrightarrow\quad
 I\in\operatorname{Ind}(M/g).
\]
The first equivalence holds because every pattern on \(I\) in the
positive half supplies every pattern in the negative half by negation.
Thus \(\operatorname{Sh}(C)\) satisfies augmentation: if
\(A,B\) are faces and \(|A|<|B|\), then some
\(e\in B\setminus A\) has \(A\cup\{e\}\) a face.

Already the common shadow of \(B_0,K_0\) violates this condition:
\(A=\{1,2\}\) and \(B=\{1,5,6\}\) are shattered, but neither
\(\{1,2,5\}\) nor \(\{1,2,6\}\) is shattered.  In fact their
common projection onto \(\{1,2,5,6\}\) has thirteen concepts and
peels to zero.  Thus failure of augmentation can occur in a positive
proper minor.  For every higher paired-square level, the fresh root
pair \(\{q,r\}\) is a maximal shattered set of size two, whereas
\(\{q,e,f\}\) is shattered for any shattered old pair \(\{e,f\}\).
This also violates augmentation.  Hence none of the fixed rooted pairs
is itself a full affine oriented-matroid tope class.

This restriction concerns the full positive half of one oriented
matroid.  It does not exclude fibers defined by several sign constraints
or a representation using auxiliary elements.  Likewise, the
vertex-shellings in
\cite[Theorems~1.1--1.2]{HochstattlerWilhelmi25shellings} order cocircuits
and the face lattices of tope cells; they do not assert an ample-class
corner peeling to a prescribed concept.  The mutation preservation
lemma in \cite[Lemma~21]{Wilhelmi25mutations} requires a Euclidean
uniform oriented-matroid program and a mutation containing the objective
element but excluding the distinguished affine element.  A signed-clause
exchange in an ample family does not by itself provide those data.
No such auxiliary encoding or converse to geometric peeling is claimed.
\end{remark}

The fixed replay \path{verify_damp_geometric_imports.py} checks the
unique corner and ordinary peeling of \(B_0\), recovers its coordinate
face in all eight ordinary families at \(j=0,1\), and replays the
four positive ordinary orders.  It checks fourteen augmentation
witnesses and explicit root-ending orders of the thirteen-concept
projection and the five-concept projection \(Q_2\vee K_2\).
Nonrealizability follows from
Proposition~\ref{prop:damp-realizable-prescribed-endpoint}; it is not
an output of a numerical realizability test.  The report is
\path{damp_geometric_imports_verification.json}.


\begin{proposition}[The ample completion of a single-clause exchange]
\label{prop:damp-single-clause-completion}
Let \(H_s,H_t\subseteq Q_F\) be ample families with common shadow
\(\Sigma\), whose signed presentations differ only at
\(I\in\mathcal N(\Sigma)\).  Write \(s\ne t\in Q_I\) for their
respective omitted traces.  Let \(R\) be the family satisfying all the
unchanged clauses, and put
\[
 \Sigma^+=\{J\subseteq F:N\nsubseteq J
                  \text{ for every }N\in\mathcal N(\Sigma)\setminus\{I\}\}.
\]
For \(u\in Q_I\), write
\(R_u=\{y\in Q_{F\setminus I}:(u,y)\in R\}\).  Then
\(R=H_s\cup H_t\) is ample with shadow \(\Sigma^+\), and
\[
 A:=R_s=R_t=\bigcap_{u\in Q_I}R_u
\]
is a nonempty ample family.  In particular, \(Q_I\times A\subseteq R\),
and the two differences are precisely
\[
 H_t\setminus H_s=\{s\}\times A,
 \qquad H_s\setminus H_t=\{t\}\times A.
\]
The new shadow faces are
\[
 \Sigma^+\setminus\Sigma
   =\{I\cup J:J\in\operatorname{Sh}(A)\}.
\]

All admissible choices for the omitted \(I\)-trace, with the other
clauses fixed, are recovered by the intrinsic slice condition
\[
 P=\{u\in Q_I:R_u=A\}.
\]
Conversely, every ample \(R\) and nonempty coordinate set \(I\)
with nonempty \(A=\bigcap_u R_u\) and at least two members of \(P\)
produce such a single-clause family: its common shadow is
\(\operatorname{Sh}(R)\setminus\{J:I\subseteq J\}\).
Thus the construction and the slice recovery describe all single-clause
exchanges with fixed shadow.
For \(u\in P\), the corresponding class is
\(H_u=R\setminus(\{u\}\times A)\).  Deleting just
\(\{u\}\times A\) by corner deletions is possible exactly when
\(A\in\Damp\): the corner orders of this slice are exactly the
corner orders of \(A\), with first coordinate \(u\).
If \(|I|\ge2\) and some \(H_u\) is peelable or is an ample forbidden
pc-minor, then \(A\in\Damp\) automatically.
\end{proposition}

\begin{proof}
A word satisfying the unchanged clauses cannot violate both distinct
\(I\)-clauses, so \(R=H_s\cup H_t\).  For every \(i\in I\), the
projections of \(H_s,H_t\) obtained by forgetting \(i\) coincide:
their shadows coincide, and the signed presentations of these projections
retain exactly the clauses supported away from \(i\), which are
unchanged.  This uses Theorem~\ref{thm:damp-fixed-shadow-signed-clutter}
on the ample projections.

Take \((t,y)\in H_s\setminus H_t\).  Equality of these projections
forces \((t^{\{i\}},y)\in H_t\) for every \(i\in I\).
The coordinate face \((H_t)_y\subseteq Q_I\) is ample, omits \(t\),
and contains every neighbor of \(t\).  In its signed presentation,
a clause excluding \(t\) must involve every coordinate of \(I\):
otherwise the neighbor obtained by changing an omitted coordinate would
violate that clause too.  Thus \(I\) is a minimal nonface of the shadow
of this face.  All its proper subsets are shattered, and ampleness gives
\((H_t)_y=Q_I\setminus\{t\}\).  In particular \((s,y)\in H_t\).
Interchanging the two classes now gives
\((H_s)_y=Q_I\setminus\{s\}\).  Consequently the two difference
slices have the same outside assignments \(A\), and \(R_y=Q_I\)
for every \(y\in A\).  Conversely any full \(I\)-cube in \(R\)
has its \(s\)-vertex in \(H_t\setminus H_s\), so its outside
assignment belongs to \(A\).  This proves the slice and intersection
formulas.  The family \(A\) is nonempty and ample, since it is the
\(s\)-coordinate face of \(H_t\).

The old strongly shattered faces \(\Sigma\), together with
\(I\cup J\) for \(J\in\operatorname{Sh}(A)\), are strongly
shattered in \(R\).  The latter assertion follows from
\(Q_I\times A\subseteq R\) and ampleness of \(A\).  These two
collections are disjoint and have total size
\(|\Sigma|+|A|=|R|\).  The lower Sandwich inequality therefore forces
\(R\) to be ample, with exactly these shadow faces.
To identify this shadow with \(\Sigma^+\), a support containing an
unchanged clause is not shattered.  On a support containing no unchanged
clause, the two projected families either are already full cubes, or
omit only the extensions of \(s\) and \(t\), respectively.  Their
union projects onto the full cube in either case.  This proves
\(\operatorname{Sh}(R)=\Sigma^+\).

For an arbitrary \(u\in Q_I\), imposing the \(I\)-clause with
omitted trace \(u\) gives \(R\setminus(\{u\}\times R_u)\).
Since \(A\subseteq R_u\), its cardinality is \(|\Sigma|\) exactly
when \(R_u=A\).  The fixed-shadow signing theorem makes this equality
equivalent to admissibility, proving the complete description of \(P\).

For the converse construction, let \(R,I,A,P\) have the stated
properties.  Strong coordinate reduction preserves ampleness, so \(A\)
is ample.  Moreover
\[
 \operatorname{Sh}(A)
 =\{J\subseteq F\setminus I:I\cup J\in\operatorname{Sh}(R)\}:
\]
ampleness identifies shattered supports with cube supports, and a cube
on either side is exactly the corresponding full product cube on the
other side.  For \(u\in P\), deleting \(\{u\}\times A\)
leaves every projection forgetting some \(i\in I\) unchanged, since
each deleted vertex has its \(i\)-neighbor in another slice of
\(Q_I\times A\).  It removes all shattered supports containing \(I\)
and retains every other shattered support.  The number of lost supports
is \(|\operatorname{Sh}(A)|=|A|\), exactly the number of deleted
vertices.  Hence \(H_u\) is ample with the asserted common shadow.
The set \(I\) is a minimal nonface of that shadow; every other minimal
nonface avoids some \(i\in I\), so its projection and omitted trace
are unchanged.  The presentations therefore differ only at \(I\).
Any two distinct deletions have union \(R\), completing recovery.

Suppose \(u\in P\), and only the vertices \(\{u\}\times D\)
have so far been deleted, where \(D\subseteq A\).  At a surviving
\((u,y)\), the incident directions are precisely all directions of
\(I\), together with the neighbor directions of \(y\) in
\(A\setminus D\).  They span a cube in the remaining family exactly
when the latter directions span a cube at \(y\) in \(A\setminus D\):
the \(u\)-slice imposes this condition, and every other slice contains
\(A\).  Thus the corner tests coincide at every step, proving the
relative peeling assertion.

Finally, \(P\) contains at least \(s,t\).  Given \(u\in P\),
choose \(v\in P\setminus\{u\}\).  The \(v\)-coordinate face of
\(H_u\) is \(A\).  When \(|I|\ge2\), this is a proper pc-minor,
since \(H_u\) contains \((Q_I\setminus\{u\})\times A\) and
therefore has at least \(3|A|\) concepts.  Either positive pc-minor
closure or forbidden-minor minimality gives \(A\in\Damp\).
\end{proof}

% Keep the statement and its table together when the page is nearly full.
\par
\ifdim\dimexpr\pagegoal-\pagetotal\relax<16\baselineskip
  \newpage
\fi
\begin{corollary}[All admissible signs in the fixed clause exchange]
\label{cor:damp-four-sign-clause-family}
Fix all clauses except the distinguished four-coordinate clause in
Theorem~\ref{thm:damp-remote-rooted-clause-pivot}, at any level \(j\).
There are exactly four admissible signs.  On the old coordinates
\((5,6,9,10)\), in that order, their endpoint behavior is
\[
\begin{array}{c|c|c}
 \text{omitted trace}&\text{deleted vertex in the base completion}
                    &\text{peels to zero}\\ \hline
 0110&604&\text{no}\\
 1001&1084&\text{yes}\\
 0101&1116&\text{yes}\\
 1101&1148&\text{yes}
\end{array}
\]
The outside family \(A\) in
Proposition~\ref{prop:damp-single-clause-completion} is a translate of
\(\{01,11,10\}^j\), so it has \(3^j\) concepts.  The common ample
completion has \((583\cdot3^j-3)/2\) concepts.
The same four choices exhaust the admissible signs for both ordinary
pairs in the proof of
Corollary~\ref{cor:damp-remote-forbidden-clause-pivot}, before and after
square completion.  In each case exactly the first sign gives a forbidden
pc-minor; the other three classes peel, and every proper pc-minor of all
four classes peels.  All four square-completed classes are two-connected.
\end{corollary}

\begin{proof}
At level zero the common completion is
\(R_0=B_0\cup\{604\}\).  The replay below checks that the only
\(I\)-slices equal to the full-cube outside family
\(A_0=\{28\}\), where \(I\) has bitmask \(1632\), are the four
listed slices.  Their traces have bitmasks \(576,1056,1088,1120\),
respectively.  Proposition~\ref{prop:damp-single-clause-completion}
therefore proves that the list is exhaustive.  Deleting \(604\)
returns the rooted factor \(B_0\).  The stored peeling of \(R_0\)
to zero begins \(1084,1148,1116\).  The other two prefixes
\(1116,1148,1084\) and \(1148,1084,1116\) also consist of corner
deletions and reach exactly the same remaining class.  Appending the
same suffix gives the three asserted positive endpoint orders.

In the paired-square construction, the outside family becomes
\(A_j=A_{j-1}\times\{01,11,10\}\), with the inherited fixed
coordinates.  Every old slice strictly larger than \(A_{j-1}\)
remains strictly larger after this product.  The additional root and
the two repairing tips have old \(I\)-traces \(0,0,32\), all outside
the four listed traces.  Thus they do not change the four minimal
slices, and no other slice can become minimal.  This proves exhaustion
at every level.  The order formula adds \(3^j\) to the order of
\(B_j\).

For each of the three positive choices the two repairing tips are still
corners: at the base level their completed cubes are disjoint from the
four exchanged vertices, and this separation persists in the paired
square.  Delete those tips, then lift a root-ending order of the old
class through the three-layer path, finishing with \(11,01,10,00\).
This is exactly the positive-order construction in the proof of
Theorem~\ref{thm:damp-remote-rooted-clause-pivot}, and proves the three
positive assertions by induction.  The first choice remains \(B_j\).

Gluing the fixed factor \(T\) adds concepts only in old \(I\)-trace
zero.  The common square tip also has old \(I\)-trace zero.  These
additions leave the four minimal slices unchanged, so the exhaustion
statement persists.  The gluing and square-completion theorems give the
first sign's forbidden status.  For each other sign, peel its factor
to zero and then \(T\) ordinarily, prefixing the square tip when it is
present.  Positive pc-minor closure gives all proper minors on these
three sides; the first side already has forbidden-minor status.
The three positive base factors are two-connected, as checked below.
The product-core and attachment argument in the proof of
Corollary~\ref{cor:damp-remote-forbidden-clause-pivot} applies to each
one and proves two-connectivity after square completion.
\end{proof}

The replay \path{verify_damp_single_clause_completions.py} checks the
three rooted levels and both ordinary stages at levels zero and one.
It recovers the common completion, its full shadow, the outside family,
and all admissible signs directly from the fixed clauses and slices.
It checks the three base endpoint prefixes and their propagated orders,
as well as the relative slice deletions.  Its report is
\path{damp_single_clause_completions_verification.json}.
The proposition replaces a general single-clause exchange by a recovered
ample completion and its unattached slices.  The corollary settles all
choices for these fixed exterior attachments; it does not characterize
endpoint behavior for arbitrary completions.  That remaining condition
must retain the other slices of \(R\), rather than only their common
intersection or the graph of the admissible traces.


% Keep the proposition statement together when the page is nearly full.
\par
\ifdim\dimexpr\pagegoal-\pagetotal\relax<30\baselineskip
  \newpage
\fi
\begin{proposition}[A cube layer can change adjacent endpoint choices]
\label{prop:damp-adjacent-clause-endpoint-switch}
Let \(S\subseteq Q_F\) be ample and peel to \(a\in S\).
Suppose that \(I\subseteq F\) is nonempty and that
\(C=Q_I\times\{y\}\subseteq S\), where
\(y\in Q_{F\setminus I}\).  For \(u\in Q_I\), put
\[
 S_u=\{z\in Q_{F\setminus I}:(u,z)\in S\},
 \qquad P_0=\{u:S_u=\{y\}\}.
\]
Introduce a new coordinate \(f\), and define
\[
 R=(S\times\{0\})\cup(C\times\{1\}),\qquad J=I\cup\{f\}.
\]
Then \(R\) is an ample single-clause completion on \(J\), with
outside family \(\{y\}\).  Its admissible omitted traces are exactly
\[
 P=(P_0\times\{0\})\cup(Q_I\times\{1\}).
\]
Write \(X_{u,b}=R\setminus\{(u,y,b)\}\) for the member
corresponding to \((u,b)\in P\).  Their common shadow is
\[
 \operatorname{Sh}(S)\ \cup
 \{T\cup\{f\}:T\subsetneq I\}.
\]
For every \(u\in Q_I\), the upper omission \(X_{u,1}\)
peels to \((a,0)\).  For \(u\in P_0\) with \((u,y)\ne a\),
\[
 X_{u,0}\text{ peels to }(a,0)
 \quad\Longleftrightarrow\quad
 S\setminus\{(u,y)\}\text{ peels to }a.
\]
After discarding any omission that removes \((a,0)\), the graph on
the remaining admissible traces, with edges at Hamming distance one,
is connected.  Nevertheless, whenever a nonroot corner deletion in
\(S\) destroys the endpoint and its corner cube is \(C\), two
adjacent retained choices in this graph have opposite endpoint behavior.
\end{proposition}

\begin{proof}
The family \(R\) is the peripheral expansion of \(S\) along
the cube \(C\), so it is ample by
Proposition~\ref{prop:damp-exact-expansion-product-tests}.  Its shadow is
\[
 \operatorname{Sh}(R)=\operatorname{Sh}(S)
       \ \cup\ \{T\cup\{f\}:T\subseteq I\}.
\]
Its \((u,0)\)-slice outside \(J\) is \(S_u\), whereas every
\((u,1)\)-slice is \(\{y\}\).  Their intersection is therefore
\(\{y\}\), and their minimal slices are precisely the displayed
set \(P\).  Since \(|Q_I|\ge2\),
Proposition~\ref{prop:damp-single-clause-completion} applies.  It gives
all the admissible choices, shows that their presentations differ only
at \(J\), and obtains their common shadow by removing \(J\)
from \(\operatorname{Sh}(R)\).

We use an explicit peeling of a punctured cube.  If \(v\in C\),
delete the concepts of \(C\setminus\{v\}\) in increasing Hamming
distance from \(v\), with arbitrary ordering at equal distance.
At each step, all neighbors closer to \(v\) have disappeared,
and all concepts farther from \(v\) remain.  The remaining incident
directions span the full cube increasing that distance, so every
deletion is a corner deletion.
Whenever a lower layer contains a remaining upper layer, an upper
corner has the same incident cube with the additional direction \(f\).
Consequently an ordinary peeling of the upper layer deletes it completely
while leaving the lower layer fixed.

In \(X_{u,1}\), the upper layer is \(C\setminus\{(u,y)\}\)
and the lower layer is \(S\).  Delete the upper punctured cube by
the preceding order, then peel \(S\) to \(a\).  This proves
every upper choice positive.
In \(X_{u,0}\), the upper concept \((u,y,1)\) is a corner:
its \(I\)-cube is present and its lower neighbor is absent.
Delete it first, then delete the remaining upper punctured cube while
retaining \(S\setminus\{(u,y)\}\) below.  This proves the
reverse implication of the endpoint equivalence.  For the forward
implication, restrict a peeling to the lower coordinate face.
Deleting an upper concept leaves that face unchanged; a deleted lower
corner restricts to a corner of the surviving face, since its incident
cube restricts to the cube on the remaining lower directions.
The restricted order ends at \(a\), as required.

The whole upper trace cube remains among the root-retaining choices.
Every remaining lower trace has its vertical neighbor in that cube,
so the graph is connected.  If \(c=(u,y)\ne a\) is a corner
of \(S\) with unique cube \(C\), then \(u\in P_0\).
Indeed, \(S_u\) is an ample coordinate face containing \(y\);
if it contained another concept, its connected graph would give a
neighbor of \(c\) outside \(C\).  Thus, if deleting \(c\)
destroys the endpoint, \(X_{u,0}\) is negative and \(X_{u,1}\)
is positive.  Their omitted traces differ only in \(f\), proving
the final assertion.
\end{proof}

For a fixed instance, take \(S=B_0\cup\{604\}\) from
Corollary~\ref{cor:damp-four-sign-clause-family}, with root zero,
\(I\) of bitmask \(1632\), and \(y=28\).
The four elements of \(P_0\) have bitmasks
\(576,1056,1088,1120\).  Use bit \(2^{12}=4096\) for \(f\).
Then \(R\) has \(306\) concepts, and its twenty admissible
omissions give \(305\)-concept classes of dimension thirteen and
VC dimension four.  Their graph is connected, and every omission
retains zero, since zero lies outside \(C\).
All sixteen upper choices peel to zero.  Among the four lower choices,
exactly the trace \(576\) fails, by the preceding corollary and the
endpoint equivalence.  In particular, the adjacent traces
\(576\) and \(4672=576+4096\), deleting concepts \(604\)
and \(4700\), respectively, have opposite endpoint behavior.

The replay \path{verify_damp_adjacent_clause_endpoint_switch.py}
checks all thirty-two possible signs, the twenty admissible
presentations with their common full shadow, the connected trace graph,
and nineteen explicit root-ending orders.  The negative member has
the proper rooted face \(B_0\), whose only corner is zero.
Its ordinary peeling is checked as well.  The report is
\path{damp_adjacent_clause_endpoint_switch_verification.json}.
Thus endpoint status need not be constant on a connected component of
the root-retaining admissible graph.  The negative member is not a
rooted factor, since its proper face \(B_0\) already fails at its root.
Theorem~\ref{thm:damp-root-section-square-pivots} below gives an
adjacent endpoint change whose negative member is itself a rooted factor.

\begin{proposition}[Reducing an adjacent endpoint change to a rooted factor]
\label{prop:damp-adjacent-exchange-minor-reduction}
Let \(H_s,H_t\subseteq Q_F\) be ample families with the same shadow,
differing only at a clause on \(I\subseteq F\), whose omitted traces
\(s,t\in Q_I\) differ in exactly one coordinate \(f\).
Let \(a\in H_s\cap H_t\), suppose that \(H_s\in\Damp\),
and suppose that \(H_t\) peels to \(a\) whereas \(H_s\) does not.
If the root-containing face
\[
 H_s^{a,f}:=\{x\in H_s:x_f=a_f\}
\]
peels to \(a\), then there is a common sequence \(\mu\) of
root-preserving pc-minor operations, none using \(f\), such that
\((\mu H_s,\mu a)\) is a rooted factor and \(\mu H_t\) peels
to \(\mu a\).  The two resulting families still have the same
shadow and differ at just one clause, whose omitted traces differ
exactly at \(f\).  Its support is a subset of \(I\), and its
outside family is a nonempty pc-minor of the outside family \(A\)
in Proposition~\ref{prop:damp-single-clause-completion}.

Consequently, the following two assertions are equivalent:
\begin{enumerate}
 \item no adjacent single-clause change retaining the root turns a
 rooted factor into a class that peels to its root;
 \item whenever \(H_s\in\Damp\) and \(H_t\) peels to a common
 root in such an exchange, \(H_s\) peels to that root if and only
 if its root-containing \(f\)-face does.
\end{enumerate}
Both assertions are false by
Theorem~\ref{thm:damp-root-section-square-pivots}.  An adjacent endpoint change
is impossible when \(|I|\le2\), and also when \(A\) is a singleton
and \(|I|\le3\).
\end{proposition}

\begin{proof}
Put \(R=H_s\cup H_t\).  The completion theorem gives
\[
 R_s=R_t=A,\qquad Q_I\times A\subseteq R,
 \qquad H_u=R\setminus(\{u\}\times A)\quad(u=s,t).
\]
We first describe how a common elementary root-preserving operation
acts on this presentation.  At every stage the image of \(H_t\)
peels to the image of \(a\), by root-preserving pc-minor closure.
Thus, while retaining endpoint failure on the other side, we cannot
use an operation that identifies the two families.

Contraction of a coordinate in \(I\) identifies the two images:
each removed concept has its neighbor in that coordinate in both
projections, so both images equal the corresponding projection of
\(R\).  Contraction of \(e\notin I\) instead preserves the
presentation, replacing \(R,A\) by \(\pi_eR,\pi_eA\).
Indeed, the two displayed minimal slices both project to \(\pi_eA\),
and every other projected slice contains \(\pi_eA\).
The converse part of the completion theorem gives the same-shadow
single-clause exchange on \(I\).

For a root-containing restriction in a coordinate \(e\notin I\),
restrict \(R\) and \(A\) to the root's value of \(e\).
If the resulting outside family is empty, the differences disappear
and the images coincide.  Otherwise the two minimal slices remain
equal to that restricted family, so the completion theorem again
preserves the adjacent exchange.
For \(e\in I\setminus\{f\}\), the two omitted traces agree
at \(e\).  If their common bit differs from \(a_e\), both
differences are discarded and the images coincide.  If it equals
\(a_e\), the restricted union has two minimal slices equal to
\(A\), now indexed by the distinct adjacent traces
\(s|_{I\setminus\{e\}},t|_{I\setminus\{e\}}\).
The completion theorem preserves the exchange, on support
\(I\setminus\{e\}\).

Only restriction in \(f\) could destroy this presentation without
identifying the two sides.  It cannot preserve endpoint failure here.
Any sequence of operations in the other coordinates commutes with
the root-containing \(f\)-restriction.  The resulting face of the
negative side is therefore a root-preserving minor of the initially
positive \(H_s^{a,f}\), and peels to its root.  This excludes an
\(f\)-restriction at every subsequent stage as well.

Now choose successive proper elementary root-preserving minors that
still fail at their roots, as long as any exist.  If a proper descendant
fails, its first elementary step must also fail, by root-preserving
pc-minor closure.  Apply the same operations to
the positive partner.  The preceding cases show inductively that
\(f\) survives and that the adjacent exchange is retained.  The
changed support can only lose coordinates, and the outside family
can only undergo restrictions and contractions, remaining nonempty.
The order of the negative family strictly decreases, so the process
terminates.  Its final family still peels ordinarily, by pc-minor
closure from \(H_s\); every proper root-preserving minor peels to
its root.  It is therefore a rooted factor.  This proves the reduction.

If assertion 1 holds, a failure of the reverse implication in
assertion 2 would give a rooted-factor counterexample by this reduction.
The forward implication in assertion 2 is face restriction closure.
Conversely, in a rooted factor the root-containing \(f\)-face is
proper and peels to its root, so assertion 2 would exclude a positive
adjacent partner.  The face is proper because the changed coordinate
is active.  Thus the assertions are equivalent; their equivalence
does not establish either one.

For the small-support statements, \(|I|=1\) gives disjoint members
and hence no common root.  If \(I=\{f,k\}\), the common root has
its \(k\)-bit opposite to \(s_k=t_k\).  The shared face
\[
 G=R\cap\{x:x_k=a_k\}
\]
contains the root.  Each member is a peripheral expansion of \(G\),
along one of its two copies of \(A\) distinguished by the \(f\)-bit.
Proposition~\ref{prop:damp-rooted-operations} makes either endpoint
possible exactly when \(G\) peels to \(a\) and \(A\in\Damp\).
Their endpoint statuses are therefore equal.

Finally suppose that \(A=\{y\}\) and \(|I|\le3\).
The two concepts \((s,y),(t,y)\) are adjacent corners of the same
full \(I\)-cube in \(R\).  Deleting either leaves the other a
corner on support \(I\setminus\{f\}\).  Hence their simultaneous
deletion leaves an ample family \(G\) containing \(a\), and
each of \(H_s,H_t\) is an ample one-concept extension of \(G\)
with acquired support of size at most two.  Item 3 of
Proposition~\ref{prop:damp-rooted-operations} makes their endpoint
statuses equal.
\end{proof}

To obtain a rooted-minimal counterexample, it therefore suffices to find
an ordinarily peelable counterexample whose root-containing \(f\)-face
peels to the root.  No assumptions on its other proper rooted minors
are required.  This reduction concerns the adjacency question; it does
not supply intrinsic admissibility for the full classification.

For the \(305\)-concept pair above, this hypothesis fails at exactly
the expected face.  The root-containing restriction in coordinate
\(12\) gives \(B_0\) on the negative side and \(B_0\cup\{604\}\)
on the positive side.  All other twenty-five root-preserving elementary
images of the negative member peel to zero.  Among all twenty-six
operations, ten identify the two sides, fifteen preserve an adjacent
single-clause exchange, and only that restriction loses the exchange
while retaining endpoint failure.
The replay \path{verify_damp_adjacent_exchange_minor_reduction.py}
checks these images, fifty-one endpoint orders, and a two-coordinate
control with a three-concept outside family.  Its report is
\path{damp_adjacent_exchange_minor_reduction_verification.json}.


\begin{theorem}[An enlarged root section permits an adjacent endpoint change]
\label{thm:damp-root-section-square-pivots}
Let \((B,0)\) be a rooted factor using every coordinate of \(F\),
and suppose that zero is its only corner.  Let
\(\alpha,\delta,\gamma\notin B\) be three distinct ample
one-concept additions, with pairwise distinct acquired supports
\(P_\alpha,P_\delta,P_\gamma\), such that all three enlarged
classes peel to zero.  Write
\[
 C_\gamma=\{\gamma^J:J\subseteq P_\gamma\}.
\]
Suppose that a coordinate face \(D\) of \(B\) contains zero and
\(C_\gamma\setminus\{\gamma\}\), but omits a neighbor
\(h=0^{\{e\}}\in B\) of zero.
On two fresh coordinates \(q,r\), define, for \(R\subseteq Q_2\),
\[
 \begin{split}
 Y_R={}&(D\times\{00\})
   \cup((B\cup\{\alpha\})\times\{01\})\\
 &\cup((B\cup\{\delta\})\times\{10\})
   \cup(B\times\{11\})\cup(\{\gamma\}\times R).
 \end{split}
\]
The family \(Y_R\) is ample exactly when \(R\) is ample,
equivalently when \(R\) is neither diagonal pair of the square.
For an ample row, with \(\operatorname{Sh}(\varnothing)=\varnothing\),
\[
 \operatorname{Sh}(Y_R)=\operatorname{Sh}(Y_\varnothing)
   \mathbin{\dot\cup}
   \{P_\gamma\cup J:J\in\operatorname{Sh}(R)\}.
\]
Every such \(Y_R\) is two-connected, peels ordinarily, and has all
proper root-preserving pc-minors peelable to their roots.  Its endpoint
condition is the intrinsic row test
\[
 Y_R\text{ peels to }(0,00)
 \quad\Longleftrightarrow\quad
 \{01,11\}\subseteq R\ \text{ or }\ \{10,11\}\subseteq R.
\]
Thus eight of the fourteen ample rows give rooted factors and six
give positive endpoints.

In particular, the four three-vertex rows \(R=Q_2\setminus\{z\}\)
give classes with the same shadow and a single changing clause on
\(P_\gamma\cup\{q,r\}\).  Their omitted traces have fixed old
part \(\gamma|_{P_\gamma}\) and new part \(z\).
Exactly the omission \(z=11\) gives a rooted factor; the other
three classes peel to the root.  Consequently an adjacent clause
change can restore the endpoint even when every proper rooted minor
of the negative member is positive.
\end{theorem}

\begin{proof}
We first construct each ample row above the known paired rooted factor.
For the negative rows, the two overlaps opposite the root section then
prevent any first deletion in a copy of the old core.

Let \(S\) be the paired construction with root section \(\{0\}\),
the same \(01,10,11\) sections, and no \(\gamma\)-copies.
Proposition~\ref{prop:damp-paired-root-transfer} makes \(S\) a rooted
factor.  The face \(D\) is proper, since \(h\notin D\), so it
peels to zero.  Reverse such a peeling to add its nonzero concepts
one by one in section \(00\).  If an addition to the current
old-coordinate face acquires support \(J\), the corresponding
addition in the full family acquires \(J\cup\{q,r\}\): its
old corner cube lifts through the other three copies of \(B\).
That support was not shattered before the addition, because its
\(00\)-section did not shatter \(J\).  The one-concept extension
criterion therefore proves ampleness at every step and reaches
\(Y_\varnothing\).

The punctured cube \(C_\gamma\setminus\{\gamma\}\) is now
present in every section.  Distinct acquired supports make
\(\alpha,\delta,\gamma\) pairwise nonadjacent, by
Lemma~\ref{lem:damp-compatible-repairs-do-not-collide}.
Thus a row addition at \(\gamma\), acquiring row support \(J\),
has precisely the combined incident support \(P_\gamma\cup J\).
Its punctured cube is present.  Moreover, among the old-coordinate
words in \(Y_\varnothing\), none realizes
\(\gamma|_{P_\gamma}\): that trace is absent from \(B\), and
neither of the other two additions acquires \(P_\gamma\).
Consequently the combined support is new exactly when \(J\) is
new in the row.  Every ample subfamily of a square has a corner
peeling, so reversing it constructs \(Y_R\) by ample one-concept
additions and proves the shadow formula.
For either diagonal row, its three shattered supports give three
new shattered supports joined to \(P_\gamma\), although only two
concepts have been added.  Thus the shadow has more supports than
the family has concepts, excluding ampleness.  This proves all the ampleness claims.

All additions took place on the active coordinate set of \(S\).
Lemma~\ref{lem:damp-rooted-ample-additions} therefore proves ordinary
peelability and positivity of all proper rooted minors for every
ample row.  It remains to decide the endpoint itself.

Suppose first that neither displayed edge is contained in \(R\).
The intersections of sections \(01,11\) and of sections \(10,11\)
are both exactly \(B\).  Consider a hypothetical peeling ending at
\((0,00)\), and its first deletion of a concept in
\(B\times\{01,10,11\}\).  All three old core copies are intact
up to this step.  Choose one of the two displayed edges incident to
the deleted concept's layer, and write \(b\) for its old word.
The corner has its mate on that edge and all the neighbors of \(b\)
in \(B\).  Its corner cube would therefore give the full cube on
these old directions in both sections.  Since their intersection is
still \(B\), the word \(b\) would be a corner of \(B\).
Hence \(b=0\).

All four copies of zero are still present: the original root is
protected, and the other three belong to the undeleted old core copies.
The candidate corner therefore has both fresh-coordinate directions
and the old \(e\)-direction toward \(h\), whose three nonroot
copies also remain.  Its corner cube would contain \((h,00)\).
That concept is absent, since \(h\notin D\) and
\(\gamma\notin B\).  This contradiction excludes the first
deletion in any nonroot old core copy, and hence excludes a peeling
to the root.  Ordinary peelability and proper-minor positivity already
proved above make \(Y_R\) a rooted factor.

Now suppose that \(R\) contains one of the displayed edges.
If \(00\in R\), then \(00\) is a corner of the row: the row
is either the full square or a three-vertex path with \(00\) as
an endpoint.  Its copy \((\gamma,00)\) is therefore a corner of
\(Y_R\), and may be deleted first.
Next delete \((D\setminus\{0\})\times\{00\}\) by the
chosen peeling of \(D\) to zero.  At each step its corner cube
lifts into the untouched copies of \(B\) in the other three layers.
The remaining family is the old partially replicated-tip construction
with row \(R\setminus\{00\}\), which is an edge or
\(\{01,10,11\}\).  Theorem~\ref{thm:damp-partial-tip-endpoint-transfer}
peels it to the root because \(B\cup\{\gamma\}\) does.
This proves the positive direction, and exhausts the fourteen ample rows.

The starting factor \(S\) is two-connected.  Each root-section
addition has its two fresh-coordinate neighbors in the preceding
family.  Each \(\gamma\)-addition has at least three old-coordinate
neighbors: a repairing support has size at least three, by
Proposition~\ref{prop:damp-rooted-operations}.  Adjoining a vertex
with at least two neighbors preserves two-connectivity, since deleting
an old vertex leaves at least one attachment and deleting the new
vertex leaves the old graph.  Thus every ample \(Y_R\) is
two-connected.

Finally take \(R=Q_2\setminus\{z\}\).  These four rows have
the same shadow.  Their common completion is \(Y_{Q_2}\), and
each differs from it by the single concept \((\gamma,z)\).
That concept has unique cube \(C_\gamma\times Q_2\).
Proposition~\ref{prop:damp-single-clause-completion} therefore gives
the common shadow and the asserted single clause on
\(P_\gamma\cup\{q,r\}\).  The row criterion makes exactly
\(z=11\) negative; its neighbors \(01,10\) are positive.
\end{proof}

For the fixed core \(B=B_0\), take
\(\alpha=128\), \(\delta=290\), \(\gamma=604\), and
\(P_\gamma\) of bitmask \(1632\).  Put
\[
 D=\{v\in B_0:v_e=0\text{ for every }
              e\notin\{2,3,4,5,6,9,10\}\}.
\]
This face has \(28\) concepts, contains the punctured repair cube,
and omits the old root neighbor \(h=2\).
The family \(Y_\varnothing\) has \(897\) concepts.
In the encoding \((v,qr)\mapsto4v+q+2r\), the four three-vertex
rows give the following \(900\)-concept classes:
\[
\begin{array}{c|c|c}
 \text{omitted row trace}&\text{deleted concept in }Y_{Q_2}
                         &\text{peels to zero}\\ \hline
 00&2416&\text{yes}\\
 10&2417&\text{yes}\\
 01&2418&\text{yes}\\
 11&2419&\text{no}
\end{array}
\]
They have dimension fourteen and VC dimension five.  The completion
has \(901\) concepts; the changed support has bitmask \(6531\),
and its outside family is \(\{112\}\).
The negative omitted trace \(2307\) is adjacent to the positive
trace \(2306\), differing in coordinate \(q\).
The negative member is a rooted factor, so its root-containing
\(q\)-face also peels to zero.  This refutes both equivalent
assertions in Proposition~\ref{prop:damp-adjacent-exchange-minor-reduction}.

The replay \path{verify_damp_root_section_square_pivots.py} checks
the sixteen rows, all ample addition chains, the eight two-overlap
obstruction certificates, six endpoint orders, and the four complete
signed presentations.  For the negative \(900\)-concept member it
replays all twenty-eight rooted elementary orders.
Its report is \path{damp_root_section_square_pivots_verification.json}.

\begin{corollary}[Adjacent signs can change forbidden-pc-minor status]
\label{cor:damp-forbidden-adjacent-sign-pivots}
There are four ample classes of order \(1188\) and dimension
twenty-six with a common shadow whose signed presentations vary on
a square of adjacent signs at a single clause.  Exactly one is a
forbidden pc-minor; the other three peel.  Every proper pc-minor of
all four peels.  The same conclusions hold for four two-connected
classes of order \(1189\), obtained by completing a common root square.
All eight classes have VC dimension five.
\end{corollary}

\begin{proof}
Glue each of the four \(900\)-concept classes above to a fixed copy
of \((B_0,0)\) on twelve fresh coordinates.  The negative choice
glues two rooted factors, so Theorem~\ref{thm:damp-cut-vertex-obstructions}
makes it forbidden.  Each other choice peels to the shared root and
then peels the fixed factor ordinarily.  The gluing adds the same
cross-coordinate clauses to all four presentations, so their single
varying clause and its four signs persist.

The first factor has the common root neighbor in coordinate \(0\);
the second has one in coordinate \(15\).  Add the common square tip
\(2^0+2^{15}=32769\).  This adds only the support \(\{0,15\}\)
to the shadow and leaves the old clause signs unchanged.
Theorem~\ref{thm:damp-square-completed-rooted-blocks} makes the
negative choice forbidden;
on a positive choice, delete the square tip first.
All root factors and all three positive partners are two-connected,
by the preceding theorem, so the square joins their two blocks after
removing the root and makes every completed graph two-connected.
Forbidden-minor minimality and positive pc-minor closure give all
the claimed proper-minor conclusions.
\end{proof}

The checker \path{verify_damp_forbidden_adjacent_sign_pivots.py}
replays the eight signed presentations, six ordinary positive orders,
and all \(156\) elementary orders of the two negative classes.
Its report is \path{damp_forbidden_adjacent_sign_pivots_verification.json}.
The negative root factor descends by corner deletions to the existing
paired factor, and its glued obstruction descends to the existing
\(1158\)-concept construction.  These examples refute adjacency and
component endpoint invariance under rooted minimality; they do not
establish another corner-descent terminal type or the full intrinsic
classification.


\begin{corollary}[Irreducibility and small corner deletions in rooted factors]
\label{cor:damp-rooted-irreducibility}
Let \((A,a)\) satisfy the three rooted-factor conditions of
Theorem~\ref{thm:damp-cut-vertex-obstructions}.  Then \(A\) is
periphery-free and Cartesian-prime.
Deleting a nonroot corner of degree at most two preserves all three
conditions.  If two nonroot corners are adjacent and have degree three,
deleting either or both also preserves the conditions.

Consequently, a rooted factor with a nonroot corner but no preserving
nonroot corner deletion has all its nonroot corners of degree at least
three.  A degree-three nonroot corner has no nonroot corner as a neighbor.
\end{corollary}

\begin{proof}
If \(A\) had a peripheral coordinate, write it as \(Y\) in
item~1 of Proposition~\ref{prop:damp-rooted-operations}.  For a lower
root, its contraction base is a proper root-preserving pc-minor and
peels to the projected root; the upper section peels ordinarily because
\(A\in\Damp\).  For an upper root, both the contraction base
and the upper section are proper root-preserving minors and peel to
their roots.  Either case makes \(A\) peel to its root, a
contradiction.  If \(A\) were a product of two nontrivial factors,
the two root-preserving factor restrictions would peel to their roots,
and item~2 would give the same contradiction.

For the deletion assertions, glue \((A,a)\) to the fixed rooted
factor of Remark~\ref{rem:damp-noncorner-rooted-factor}.  The result
is forbidden by Theorem~\ref{thm:damp-cut-vertex-obstructions}.
Every nonroot corner of \(A\) is a corner of the glued family with
the same degree.  Apply
Corollary~\ref{cor:damp-low-degree-corner-descent} for a corner of
degree at most two, and
Corollary~\ref{cor:damp-degree-three-corner-pair-descent} for an
adjacent degree-three pair.  The resulting forbidden families still
have a cut vertex, so their changed rooted factors satisfy all three
conditions by the cut-vertex theorem.  Finally, adjacent corners share
their unique maximal cube and hence have the same degree.  This proves
the last assertion as well.
\end{proof}

\begin{remark}[An adjacent corner can change the prescribed endpoint]
\label{rem:damp-rooted-rank-four-repair}
Let \(B\) be the \(289\)-concept factor rooted at \(29\) from
Remark~\ref{rem:damp-noncorner-rooted-factor}, and put
\[
 M=B\cup\{157\}.
\]
The addition is ample and acquires the support
\(\{1,2,7,8\}\), with bitmask \(390\) in the stored embedding.
Both \(29\) and \(157\) are corners of the completed four-cube;
they are adjacent in coordinate \(7\).
The \(290\)-concept family \(M\) has verified peelings ending
at either vertex.  In contrast, \(B\) cannot peel to \(29\).
Thus deleting a nonroot corner can destroy peeling to the root even
when the root is also a corner and is adjacent to the deleted corner.
This is a support-four change of endpoint behavior, whereas item~3 of
Proposition~\ref{prop:damp-rooted-operations} excludes such a change
for additions of support size at most two.

The checker \path{verify_damp_rooted_operations.py} replays both
endpoint orders and checks the full-cube addition.  It also validates
the peripheral and product constructions in the proof and checks
periphery-freeness and all \(2047\) nontrivial coordinate bipartitions
of each of the seven stored rooted factors.  The report is
\path{damp_rooted_operations_verification.json}.
Theorem~\ref{thm:damp-adjacent-rank-three-repair} gives the sharper
degree-three endpoint repair. This older degree-four repair does not itself
give a terminal cut-vertex obstruction.  Theorem~\ref{thm:damp-terminal-cut-vertex} combines it
with a second repair of the same core to construct one.
\end{remark}

\begin{proposition}[Proper images after deleting a degree-three corner]
\label{prop:damp-rooted-degree-three-images}
Let \((A,a)\) satisfy the rooted-factor conditions of
Theorem~\ref{thm:damp-cut-vertex-obstructions}, in an irredundant
embedding \(A\subseteq Q_F\).  Let \(c\ne a\) be a corner with
three-coordinate incident support \(P\), and put \(B=A\setminus\{c\}\).
For \(e\in F\), let \(\pi_e\) forget coordinate \(e\), and write
\[
 R_eC=\pi_e\{x\in C:x_e=a_e\},\qquad a_e'=\pi_ea.
\]
Then \(B\) is ample and cannot peel to \(a\).  If \(B\in\Damp\),
the pair \((B,a)\) is a rooted factor if and only if
\begin{equation}
 \begin{aligned}
  &\pi_eB\text{ peels to }a_e' &&(e\in F\setminus P),\\
  &R_eB\text{ peels to }a_e' &&(e\in F\setminus P,\ c_e=a_e).
 \end{aligned}
 \label{eq:damp-rooted-degree-three-exterior-tests}
\end{equation}
If one of these tests fails, the corresponding image of \(A\) is
obtained from the image of \(B\) by adding the nonroot concept
\(\pi_ec\), acquiring exactly \(P\); this addition makes peeling
to \(a_e'\) possible.

If \(a\) is itself a corner adjacent to \(c\), the assumption
\(B\in\Damp\) holds automatically.  In this case all root-containing
restrictions outside \(P\) occur in the displayed test.
\end{proposition}

\begin{proof}
Corner deletion preserves ampleness.  A peeling of \(B\) to \(a\),
prefixed by \(c\), would peel \(A\) to \(a\), which is impossible.
The punctured three-cube at \(c\) retains every direction in \(P\).
It also supplies a surviving concept with the same outside coordinates
as \(c\), so deleting \(c\) cannot make any outside coordinate
constant.  Thus \(B\) still uses every coordinate of \(F\).

Every elementary root-preserving image \(\mu A\) peels to its root.
We compare it with \(\mu B\).  If \(\mu=\pi_e\) and \(e\in P\),
the removed concept is identified with its neighbor \(c^{\{e\}}\)
in \(B\), so the two images coincide.  For \(\mu=R_e\) with
\(c_e\ne a_e\), the removed concept is discarded and again they
coincide.  If \(\mu=R_e\), \(e\in P\), and \(c_e=a_e\), the
image of \(A\) adds \(\pi_ec\) on the two-coordinate support
\(P\setminus\{e\}\).  Indeed, the incident cube at that concept
is precisely the retained face of the original corner cube.  Both
images are ample, and the root is an old concept of the smaller image.
Item~3 of Proposition~\ref{prop:damp-rooted-operations} therefore
makes \(\mu B\) peel to its root.

Only the operations in
\eqref{eq:damp-rooted-degree-three-exterior-tests} remain.  In an
outside contraction, \(c\) has no neighbor in the contracted direction,
so its image is a new concept.  Contraction retains exactly the shattered
supports avoiding \(e\); the identity
\(\operatorname{Sh}(A)=\operatorname{Sh}(B)\mathbin{\dot\cup}\{P\}\)
therefore preserves the support gain \(P\).  In an outside restriction
retaining \(c\), the entire corner cube survives and its incident
support is still \(P\).  The added concept differs from the image of
the root: equality under contraction would make \(a\) an outside
neighbor of \(c\), and a retained restriction is injective.
This proves the asserted description of every failed test.

If all displayed tests pass, every elementary root-preserving image
of \(B\) peels to its root.  The same is true for every proper
root-preserving pc-minor, by closure under further root-preserving
operations.  Together with \(B\in\Damp\) and the failure to peel
to \(a\), this gives exactly the three factor conditions.  Their
necessity gives the converse.

Finally suppose that \(a,c\) are adjacent corners.  They share their
unique maximal three-cube.  Write \(a=c^{\{f\}}\) and
\(I=P\setminus\{f\}\).  Deleting \(c\) leaves \(a\) as a
corner of its \(I\)-face, so \(G=A\setminus\{a,c\}\) is ample.
Each individual addition of \(a\) or \(c\) to \(G\) acquires
the same two-coordinate support \(I\).
Proposition~\ref{prop:damp-square-pair-invariance}, applied to
\(A=G\cup\{a,c\}\in\Damp\), gives \(G\in\Damp\).
Deleting the corner \(a\) first in \(B\), and then peeling
\(G\), gives \(B\in\Damp\).  The two adjacent concepts agree
outside \(P\), proving the final assertion about the tests.
\end{proof}

This proposition separates ordinary peelability from rooted minimality
when the deleted corner is adjacent to a corner root.
It does not show that its outside-coordinate tests always pass.
The fixed checks in \path{verify_damp_rooted_corner_images.py} validate
the image and support identities and replay stored endpoint orders for
three degree-three preserving deletions.  Their report is
\path{damp_rooted_corner_images_verification.json}.


\begin{proposition}[Minimal sources of rank-three endpoint repairs]
\label{prop:damp-minimal-rank-three-endpoint-repair}
Let \(S\subseteq Q_F\) be ample and ordinarily peelable, let
\(a\in S\), and suppose that \(H=S\cup\{c\}\) is an ample
one-concept extension with acquired support \(P\), where \(|P|=3\).
Assume that \(H\) peels to \(a\) and \(S\) does not.
There is a root-preserving pc-minor map \(\mu\), using only
coordinates in \(F\setminus P\), such that
\[
 (T,t)=(\mu S,\mu a)
\]
is a rooted factor of Theorem~\ref{thm:damp-cut-vertex-obstructions}.
Moreover, \(\mu H=T\cup\{\mu c\}\) is still an ample one-concept
extension with acquired support \(P\), and it peels to \(t\).
The whole punctured \(P\)-cube at the added concept is retained.

If \(a,c\) are adjacent corners of \(H\), they remain adjacent
corners of \(\mu H\), both of degree three.  In that case \(t\)
is a corner of \(T\) of degree two.
\end{proposition}

\begin{proof}
First consider an elementary root-preserving operation in a coordinate
\(p\in P\).  Contraction identifies \(c\) with its existing
neighbor, so the images of \(H\) and \(S\) coincide.  A
root-containing restriction that discards \(c\) also gives equal
images.  If the restriction retains \(c\), its image adds that
concept on the two-coordinate support \(P\setminus\{p\}\).
The root is an old concept of the smaller image.  The image of \(H\)
peels to the root, by root-preserving minor closure; item~3 of
Proposition~\ref{prop:damp-rooted-operations} gives the same conclusion
for the image of \(S\).  Thus every elementary operation in \(P\)
makes peeling to the root possible.

Choose a pc-minor \((T,t)\) of \((S,a)\), minimal among those
that cannot peel to their roots.  It is ample and ordinarily peelable
by ordinary pc-minor closure.  Minimality supplies the third rooted
factor condition.  Coordinate restrictions and contractions on distinct
coordinates commute: both simply retain the words with specified
restricted bits and then forget the contracted bits.  Every minor has
such a form with disjoint restricted and contracted coordinate sets.
If the map to \(T\) used a coordinate of \(P\), that operation
could therefore be performed first.  The preceding paragraph and
root-preserving minor closure would make \(T\) peel to \(t\),
a contradiction.  Hence only coordinates outside \(P\) are used.

Apply these operations simultaneously to \(S\) and \(H\).
At every intermediate stage the smaller image fails at its root,
since its descendant \(T\) does; the larger image peels to its root.
Consequently the images cannot coincide.  Their difference remains
one concept, so no restriction discards the added concept and no
contraction identifies it with an old image.  An outside restriction
retaining the concept retains its entire \(P\)-cube.  An outside
contraction preserves the shadow difference \(\{P\}\), because
contraction keeps precisely the shattered supports avoiding its
coordinate.  By the ample one-concept extension criterion, the
acquired support remains \(P\) at every stage.  This proves the
first assertion and the retained punctured-cube statement.

Suppose now that \(a,c\) are adjacent corners of \(H\).
They share their unique maximal cube, of support \(P\).
A root-containing restriction outside \(P\) retains this cube and
both corners.  An outside contraction also preserves either corner:
its fibre is a singleton, and projecting the ample family obtained
by deleting that corner leaves a one-concept difference of support
\(P\).  The extension criterion again makes its image a corner
of degree three.  Induction over the operations preserves both corners
and their edge, whose coordinate belongs to \(P\).
Deleting \(\mu c\) from \(\mu H\) leaves \(t\) in the
opposite face of that three-cube, with exactly two incident directions.
Thus \(t\) is a degree-two corner of \(T\).
\end{proof}

\begin{corollary}[Equivalent forms of adjacent degree-three endpoint repair]
\label{cor:damp-adjacent-corner-minimal-repair-test}
The following statements are equivalent.
\begin{enumerate}
 \item There is an ample family \(H\) with adjacent degree-three
 corners \(a,c\) such that \(H\) peels to \(a\), but
 \(H\setminus\{c\}\) does not peel to \(a\).
 \item There is a rooted factor \((T,t)\) whose root is a
 degree-two corner, and an ample one-concept addition adjacent to
 \(t\), whose acquired support has size three, that makes peeling to \(t\)
 possible.
\end{enumerate}
Both profiles occur in Theorem~\ref{thm:damp-adjacent-rank-three-repair}.
In particular, prescribed-endpoint invariance under deleting an adjacent
nonroot corner is false already in degree three.
\end{corollary}

\begin{proof}
For item~1, put \(G=H\setminus\{a,c\}\).
Deleting either corner leaves the other as a degree-two corner, so
\(G\) is ample and both individual additions to \(G\) have the
same two-coordinate support.  Proposition~\ref{prop:damp-square-pair-invariance}
makes \(G\) ordinarily peelable, and
Proposition~\ref{prop:damp-rank-two-extension-invariance} gives this
property to \(H\setminus\{c\}\).  The preceding proposition
therefore produces item~2.  Conversely, in item~2 the added concept
is a degree-three corner.  Its cube contains the neighboring root;
the root's two old incident directions and its new edge are exactly
the directions of that cube.  Thus both are degree-three corners,
giving item~1.

The existence assertion and the failure of endpoint invariance follow
from the certified $622$-to-$621$ example in
Theorem~\ref{thm:damp-adjacent-rank-three-repair}.
\end{proof}

This equivalence identifies rooted-factor sources for the failure, whose
existence is now established. Its factor conditions remain recognition
data, not intrinsic admissibility. The older small control below validates
the reduction but is not the counterexample.

The fixed checker \path{verify_damp_minimal_endpoint_repair.py}
validates the corner and support transport in all \(18\) outside
images of a \(71\)-concept example, obtained by \(219\) corner
deletions from the \(290\)-concept family.  Both this family and
its \(70\)-concept corner deletion peel to \(29\), so it is a
positive control, not a repair of endpoint failure.  The checker also
replays all \(48\) rooted elementary orders and the existing
peripheral example with a degree-two corner root that fails as an
endpoint.  That example has an incident-coordinate contraction already
failing at its root, so it is not a rooted factor.  The report is
\path{damp_minimal_endpoint_repair_verification.json}.


\begin{proposition}[A prescribed endpoint can be lost under reduction]
\label{prop:damp-rooted-strong-projection-failure}
Let \(B\) and \(M=B\cup\{157\}\) be the families in
Remark~\ref{rem:damp-rooted-rank-four-repair}, and put
\[
 Y=(M\times\{0\})\cup(B\times\{1\}),\qquad r=(29,0).
\]
Then \(Y\) is ample, has \(579\) concepts, dimension \(13\), and
VC dimension four.  It peels to \(r\), but its reduction
\[
 Y^{\{e\}}=\{x:(x,0),(x,1)\in Y\}=B
\]
in the new coordinate \(e\) cannot peel to the projected root \(29\).
The full fibre through the root is present.  The only corners of \(Y\)
are \((157,0)\) and \((29,1)\), and every peeling ending at \(r\)
starts by deleting \((29,1)\).
VC dimension four is the least possible for such a failure under a
nonempty reduction.
\end{proposition}

\begin{proof}
The peripheral expansion is ample.  Its shadow is
\[
 \operatorname{Sh}(Y)=\operatorname{Sh}(M)
 \mathbin{\dot\cup}
 \{I\cup\{e\}:I\in\operatorname{Sh}(B)\}.
\]
Since \(M,B\) have orders \(290,289\) and VC dimensions four and
three, respectively, this gives the asserted size and VC dimension.
Both use all twelve old coordinates, and the new coordinate is active.
By Proposition~\ref{prop:damp-rooted-operations}, \(Y\) peels to
\(r\): first peel the upper copy of \(B\) ordinarily, and then
peel \(M\) to \(29\).  Its reduction is exactly \(B\),
whose only corner is \(29\) and which has more than one concept.
That corner must be deleted first, so it cannot be an endpoint.

An upper vertex is a corner of \(Y\) exactly when its projection
is a corner of \(B\); thus the only upper corner is \((29,1)\).
The lower vertex \((157,0)\) has no upper mate and is a corner
because \(157\) is a corner of \(M\).  Any other lower corner
would have an upper mate, and its incident cube would force its
projection to be a corner of \(B\).  The only candidate is
\((29,0)\), but its neighbors \((157,0)\) and \((29,1)\)
have the missing opposite vertex \((157,1)\), excluding that corner.
This proves the two-corner assertion.

Deleting \((157,0)\) leaves \(B\square Q_1\).  This product
cannot peel to \(r\), since its root-containing lower restriction
is \(B\), which cannot peel to \(29\).  Hence a peeling ending
at \(r\) must start with \((29,1)\).

For the dimension threshold, if \(Z^D\) shatters \(I\), then
\(Z\) shatters \(I\cup D\): every trace on \(I\) in the
reduction comes from a full \(D\)-fibre.  Thus
\(\operatorname{VC}(Z^D)\le\operatorname{VC}(Z)-|D|\).
If \(Z\) is ample of VC dimension at most three and
\(D\ne\varnothing\), its nonempty reduction is ample
of VC dimension at most two.  Such a family peels to any prescribed
concept, by the VC-two endpoint result recalled in
Remark~\ref{rem:damp-one-corner-peelable}.  This excludes a smaller
VC dimension.  The checker
\path{verify_damp_rooted_strong_projection.py} independently checks
the shadow, the two corners, and explicit ordinary and endpoint orders;
its report is \path{damp_rooted_strong_projection_verification.json}.
\end{proof}

Ordinary corner peelings do pass to reductions by ordering
full fibres according to their first deleted member.  In this example
the fibre containing the retained root is destroyed first, so the
induced peeling starts at its projected root.  Thus the endpoint
condition is lost even though the ordinary projection argument applies.
The outside-coordinate tests in
Proposition~\ref{prop:damp-rooted-degree-three-images} use contractions
and root-containing restrictions; reduction cannot replace
their endpoint compatibility requirement.


\begin{proposition}[Peeling the section opposite a prescribed endpoint]
\label{prop:damp-rooted-section-obstructions}
Let
\[
 X=(A_0\times\{0\})\cup(A_1\times\{1\})
\]
be ample, with both sections nonempty, and put \(C=A_0\cap A_1\).
Suppose that \(A_1\) admits a \emph{relative corner peeling to \(C\)}:
a sequence of corner deletions removing exactly \(A_1\setminus C\)
and retaining all of \(C\).  For \(a\in A_0\), one then has
\begin{equation}
 X\text{ peels to }(a,0)
 \quad\Longleftrightarrow\quad
 A_0\text{ peels to }a\ \text{ and }\ C\in\Damp.
 \label{eq:damp-relative-section-endpoint-test}
\end{equation}

Consequently, in a rooted factor satisfying
Theorem~\ref{thm:damp-cut-vertex-obstructions}, the section opposite
the root in every active coordinate has no relative corner peeling to
the overlap.  Each such section has VC dimension at least three and
uses at least six coordinates.  It contains an ample family \(D\)
strictly containing the overlap \(C\), with every corner of \(D\)
lying in \(C\).  In particular, every rooted factor uses at least
seven coordinates.
\end{proposition}

\begin{proof}
The sections and their overlap are ample.  The overlap is nonempty
because an ample family is connected and an edge between the two
sections projects into \(C\).
If \(X\) peels to \((a,0)\), its root-containing lower restriction
peels to \(a\).  The overlap peels ordinarily: order its full
two-vertex fibres by their first deleted member in a peeling of \(X\).
Equivalently, reverse the peeling and order full fibres by completion
time; the resulting prefixes are reductions of ample prefixes
and hence are ample.  This establishes necessity without prescribing
an endpoint in \(C\).

Conversely, first carry out the given relative peeling in the upper
section of \(X\).  Every deleted vertex lies outside \(C\) and has
no lower mate, so its incident corner cube is unchanged by the presence
of the lower section.  The remaining family is
\((A_0\times\{0\})\cup(C\times\{1\})\).
The lower-root peripheral test in
Proposition~\ref{prop:damp-rooted-operations} now gives a peeling
to \((a,0)\), proving the equivalence.

Apply this to a rooted factor, orienting any active coordinate so that
the root lies in the lower section.  That section is a proper
root-preserving restriction and peels to its root.  The overlap belongs
to \(\Damp\), by the ordinary full-fibre argument above applied to
an ordinary peeling of the factor.  A relative peeling of the opposite
section would therefore contradict its failure to peel to the root.
The opposite section must have VC dimension at least three by
Lemma~\ref{lem:damp-relative-peeling-vc-two}, and it must use at least
six coordinates by
Theorem~\ref{thm:damp-five-coordinate-relative-peeling}.
Since the split coordinate is absent there, the whole factor uses
at least seven coordinates.

Finally, keep deleting corners of the opposite section outside \(C\)
for as long as possible.  The process cannot reach \(C\), since
that would give a relative peeling.  It stops at an ample family
\(D\supsetneq C\) all of whose corners lie in \(C\), as claimed.
No ordinary peeling of this terminal \(D\) is asserted.
\end{proof}

The pairs \(C\subsetneq D\) are upper-root-locked pairs in the
sense of Definition~\ref{def:damp-root-locked-relative-pair}; their
one-hole description is given by
Theorem~\ref{thm:damp-relative-one-hole-upper-lock}.
They give necessary local certificates in every opposite section.
Satisfying these conditions in every coordinate does not imply failure
at the root: Proposition~\ref{prop:damp-no-root-section-first-peeling}
gives endpoint-positive counterexamples.  A sufficient compatibility
condition remains open.  The checker
\path{verify_damp_rooted_section_obstructions.py} validates all
\(84\) section/overlap pairs of the seven stored rooted factors,
including their complete relative deletion graphs and ordinary
overlap orders.  Its report is
\path{damp_rooted_section_obstructions_verification.json}.


\begin{lemma}[Corner cubes and initial section obstructions]
\label{lem:damp-root-in-corner-cubes}
Let \(A\subseteq Q_F\) be ample and let \(a\in A\).  In each
active coordinate, orient the two sections so that the root is lower.
Every corner of the opposite section lies in its overlap with the
root-containing section, in every coordinate, if and only if
\begin{equation}
 a\in Q_{I_A(c)}(c)
 \qquad\text{for every corner }c\text{ of }A.
 \label{eq:damp-root-in-all-corner-cubes}
\end{equation}
Here \(I_A(c)\) is the incident support of the corner \(c\), and
\(Q_{I_A(c)}(c)\) is its unique maximal cube.
\end{lemma}

\begin{proof}
A corner of an opposite section outside the overlap has no mate in
the root-containing section.  Its lift \(c\) is therefore a corner
of \(A\) with the same incident support, and the split coordinate
belongs to \(\operatorname{Diff}(c,a)\setminus I_A(c)\).
Conversely, if a corner \(c\) has such a coordinate, its whole
corner cube lies in the opposite section and it has no mate across
that coordinate.  Its projection is a corner outside the overlap.
The absence of all these cases is exactly
\(\operatorname{Diff}(c,a)\subseteq I_A(c)\), which is the displayed
cube condition.
\end{proof}

\begin{proposition}[Endpoint peeling need not remove a coordinate section first]
\label{prop:damp-no-root-section-first-peeling}
There is an ample family \(P\) of VC dimension three and a concept
\(a\in P\) such that \(P\) peels to \(a\), but, for every
active coordinate \(e\), there is no relative corner peeling of
\(P\) to its root-containing section
\[
 P_e(a)=\{x\in P:x_e=a_e\}.
\]
Furthermore, in every coordinate the opposite section has no corner
outside its overlap with \(P_e(a)\).  The family is two-connected,
periphery-free, and Cartesian-prime.

VC dimension three is the least possible for the failure of this
section-first schedule in a nontrivial ample family.  The failure can
also occur when \(a\) is itself a corner, in VC dimension four.
\end{proposition}

\begin{proof}[Proof with fixed certificates]
Take the \(294\)-concept family \(P\) from
Remark~\ref{rem:damp-one-corner-peelable} and set \(a=1895\), in
its stored twelve-coordinate embedding.  A complete peeling ending
at \(1895\) is recorded in
\path{computations/round776_noncorner_roots.json}.
The same underlying family, rooted at its unique corner \(1894\),
is the rooted factor of Remark~\ref{ex:damp-cut-vertex-587};
Corollary~\ref{cor:damp-rooted-irreducibility} and
Theorem~\ref{thm:damp-cut-vertex-obstructions} give the three
asserted irreducibility properties.
The checker \path{verify_damp_rooted_section_counterexample.py}
replays the endpoint order and the relative certificates below.

The only corner of \(P\) is \(1894\), and \(1895\) is its
neighbor in coordinate \(0\).  For every \(e\ne0\), the unique
initial corner lies in \(P_e(a)\), so no relative peeling to that
section can start.  Coordinate \(0\) has exactly \(153\) concepts
opposite the odd root \(1895\).  Its complete relative deletion
certificate has \(13\) states.  The possible deletions follow the
chain
\[
 1894,\ 1638,\ 1634,\ 1632,\ 1600,\ 3648,
\]
with \(3942\) insertable anywhere after \(1894\).  The checker
verifies that these are exactly all allowed next deletions in every
state while the odd section remains fixed.  All six maximal orders
therefore remove the same seven concepts and then stop, leaving
\(146\) concepts of the opposite section.  This proves failure in
coordinate \(0\) as well.

The root \(1895\) lies in the unique corner cube at \(1894\).
Lemma~\ref{lem:damp-root-in-corner-cubes} proves the asserted initial
overlap condition in every coordinate.  Those opposite sections are
strictly larger than their overlaps, since \(P\) is periphery-free.
Thus all the local relative-obstruction conditions of
Proposition~\ref{prop:damp-rooted-section-obstructions} can hold in
an endpoint-positive family.

For the corner-root example use \(M=B\cup\{157\}\), rooted at
\(29\), from Remark~\ref{rem:damp-rooted-rank-four-repair}.
All five corners have the same maximal four-cube, which contains
\(29\).  Its twelve root-containing-section intervals have complete
negative certificates with \(34\) states in total.  The same checker
replays them, its endpoint peeling, and all \(24\) rooted elementary
minor orders of each positive example.  It also directly checks
two-connectivity, periphery-freeness, and Cartesian primality.
The report is \path{damp_rooted_section_counterexample_verification.json}.

Finally, if a nontrivial ample family has VC dimension at most two,
choose any active coordinate.  Its root-containing section is nonempty
and ample.  Lemma~\ref{lem:damp-relative-peeling-vc-two} gives a
relative peeling to that section, which then peels to the prescribed
root by the same lemma applied to a singleton.  Hence the failure
cannot occur below VC dimension three.
\end{proof}


The obstruction in an opposite section concerns which concepts may be
deleted while that section's overlap is retained.  A different possible
test uses every missing cube of the original family, in every rank.
Even this test loses information when a deletion removes a vertex of
a cube that was originally complete.

\begin{proposition}[Pruning with original missing cubes is incomplete]
\label{prop:damp-original-cube-pruning}
Let \(H\subseteq Q_F\) be finite and let \(A\subseteq H\) be a
set of protected concepts, possibly empty.  Starting with \(U=H\),
repeatedly delete any \(v\in U\setminus A\) for which
\begin{equation}
 Q_{I_U(v)}(v)\subseteq H,
 \qquad I_U(v)=\{e\in F:v^{\{e\}}\in U\}.
 \label{eq:damp-original-cube-pruning}
\end{equation}
Cube containment is tested in the original \(H\), not in the
current \(U\).  The final set, denoted \(K_A(H)\), is independent
of the choices.  It is the greatest set \(U\), with
\(A\subseteq U\subseteq H\), such that every \(v\in U\setminus A\)
has an original missing-cube witness
\begin{equation}
 P\subseteq I_U(v),\qquad v^P\notin H.
 \label{eq:damp-original-hole-witness}
\end{equation}
The witness may be chosen so that all \(v^J\), \(J\subsetneq P\),
belong to \(H\); thus it is a punctured cube already present in
the original family.

If \(H\) is ample and admits a corner peeling to \(A\) retaining
all of \(A\), then \(K_A(H)=A\).  The converse fails both for
\(A\) a singleton and for \(A=\varnothing\).  Specifically,
there is a \(292\)-concept rooted factor \((B,0)\) with
\(K_{\{0\}}(B)=\{0\}\), and a \(583\)-concept ample forbidden
pc-minor \(X\) with \(K_\varnothing(X)=\varnothing\).
Their dimensions are twelve and twenty-four, respectively, and both
have VC dimension four.
\end{proposition}

\begin{proof}
If a concept satisfies \eqref{eq:damp-original-cube-pruning}, it
continues to satisfy it after other concepts are deleted: its incident
support only shrinks.  More directly, every set satisfying
\eqref{eq:damp-original-hole-witness} survives any allowed deletion
sequence.  At the first deletion from such a set, all the witnessing
neighbors would still be present, contradicting
\eqref{eq:damp-original-cube-pruning}.  The set left when pruning
stops itself satisfies \eqref{eq:damp-original-hole-witness}.
It is therefore the greatest such set, proving uniqueness and
independence of the deletion choices.  Choose an inclusion-minimal
\(P\) in \eqref{eq:damp-original-hole-witness}; all its proper
subsets give vertices in \(H\), as asserted.

Suppose that a corner peeling retaining \(A\) exists and that
\(K_A(H)\setminus A\ne\varnothing\).  Just before its first
deletion of a member \(v\) of that difference, every witnessing
neighbor from \eqref{eq:damp-original-hole-witness} remains.  The
corner cube would contain \(v^P\), which was absent originally
and is still absent.  This is impossible, proving necessity.

For the counterexamples, normalize the family in
Remark~\ref{rem:damp-noncorner-rooted-factor} by putting
\[
 B=\{x\mathbin{\oplus}29:x\in G\cup\{91\}\},
 \qquad X=B\cup\{2^{12}x:x\in B\},
\]
where \(\oplus\) is bitwise exclusive-or and \(G\) is the fixed
\(291\)-concept seed used there.  That remark establishes the
rooted-factor and forbidden-minor assertions, with the stated dimensions.
The replay \path{verify_damp_original_cube_pruning.py} checks
pruning sequences of lengths \(291\) and \(583\), leaving
\(\{0\}\) and \(\varnothing\), respectively.  Uniqueness of
the final set gives the claimed values of \(K_A\).

The first discrepancy is explicit.  The corners of \(B\) are
\(32,64,70,96\).  The checked pruning begins \(32,64\).
After deleting \(32\), the incident support at \(64\) still
has bitmask \(102\).  Its full four-cube belongs to \(B\), so
the pruning rule permits deleting \(64\); but that cube's vertex
\(32\) has just disappeared.  Hence \(64\) is not a corner of
\(B\setminus\{32\}\).  The rule has forgotten a newly created
hole in an originally complete cube.  Already in the original square
\(\{0,32,64,96\}\), these first two deletions remove both common
neighbors of the opposite surviving concepts \(0,96\).  This is
exactly the full-square obstruction in the metric criterion of
Theorem~\ref{thm:damp-boolean-prefix-cutset}.

The same replay independently checks the complete negative deletion
graphs, with eight rooted states and sixty-four ordinary states,
an ordinary peeling of \(B\), and all \(24+72\) elementary
minor orders.  These establish failure of actual peeling independently
of the pruning computation.  Its report is
\path{damp_original_cube_pruning_verification.json}.
\end{proof}

Thus no protected-set certificate using only
\eqref{eq:damp-original-hole-witness}, even with all original
punctured cubes available, can detect these two failures.  The
sufficiency of a nontrivial protected set remains valid.  A complete
construction must also control deletions inside originally complete
cubes; the pruning rule does not supply intrinsic rooted-factor or
repair admissibility.


\begin{proposition}[Two-connectivity of replicated extensions]
\label{prop:damp-replicated-two-connected}
Let \(S\subseteq Q_F\) be connected with \(|S|\ge2\), let
\(E\ne\varnothing\) be disjoint from \(F\), and let
\(\mathcal A\subseteq Q_F\setminus S\).  Suppose every
\(a\in\mathcal A\) has at least two neighbors in \(S\).
For arbitrary rows \(R_a\subseteq Q_E\), the graph of
\[
 X=(S\times Q_E)\cup
       \bigcup_{a\in\mathcal A}(\{a\}\times R_a)
\]
is two-connected.  In particular, this holds when \(S\) is ample
and nonpeelable and every \(S\cup\{a\}\) is ample and peelable.
\end{proposition}

\begin{proof}
The core \(S\times Q_E\) is two-connected.  Indeed, after deleting
\((s,t)\), its layers indexed by \(Q_E\setminus\{t\}\)
form a connected graph: a nontrivial cube with one vertex removed is
connected.  Every surviving vertex in the \(t\)-layer has a neighbor
in one of these other layers.  The core has at least four vertices,
so this proves the assertion also when \(|E|=1\).

Every added vertex \((a,t)\) has at least two distinct neighbors in
this core.  After deleting any one vertex of \(X\), the surviving
core is connected, and every surviving added vertex still has a neighbor
in it.  The whole remaining graph is therefore connected.
For the last assertion, a resolving ample one-concept extension has
at least three neighbors in its nonpeelable base by
Proposition~\ref{prop:damp-rank-two-extension-invariance}.
\end{proof}

The terminal cut-vertex obstruction of
Theorem~\ref{thm:damp-terminal-cut-vertex} therefore has no nontrivial
replicated common-seed presentation.  Its two degree-four corners
admit no obstruction-preserving deletion.  Thus a separate terminal
assembly type is necessary; intrinsic rooted-factor admissibility
remains open.


\begin{theorem}[Square completion across rooted blocks]
\label{thm:damp-square-completed-rooted-blocks}
Let \((A_1,0)\) and \((A_2,0)\) be rooted factors satisfying
the three conditions of Theorem~\ref{thm:damp-cut-vertex-obstructions},
on disjoint coordinate sets \(F_1,F_2\).  Put
\[
 N_i=\{e\in F_i:0^{\{e\}}\in A_i\},\qquad X=A_1\vee A_2,
\]
and choose a nonempty set \(\Gamma\subseteq N_1\times N_2\).
Then
\begin{equation}
 X_\Gamma=X\cup\{0^{\{e,f\}}:(e,f)\in\Gamma\}
 \label{eq:damp-square-completed-blocks}
\end{equation}
is a two-connected ample forbidden pc-minor.  Its shadow is
\[
 \operatorname{Sh}(X_\Gamma)
 =\operatorname{Sh}(A_1)\cup\operatorname{Sh}(A_2)
   \cup\{\{e,f\}:(e,f)\in\Gamma\},
\]
with the empty support shared by the two block shadows.
Every added concept is a degree-two corner.  If each \(A_i\) has
its root as its only corner, these are all the corners of \(X_\Gamma\).

Conversely, let \(O\) be a two-connected ample forbidden pc-minor,
and let \(c\) be a degree-two corner such that \(O\setminus\{c\}\)
has a cut vertex.  Then \((O,c)\) has the displayed construction
with \(|\Gamma|=1\).  The two rooted factors and their chosen root
neighbors are recovered uniquely up to interchange.
\end{theorem}

\begin{proof}
The cut-vertex theorem makes \(X\) forbidden.  For \((e,f)\in\Gamma\),
the concept \(c_{ef}=0^{\{e,f\}}\) has exactly the two neighbors
\(0^{\{e\}}\) and \(0^{\{f\}}\) in the enlarged family.
Indeed, the old concepts use coordinates of only one block, and
distinct added concepts have even Hamming distance.  The opposite
vertex of their square is \(0\).  Until \(c_{ef}\) is added,
the trace \(11\) on \(\{e,f\}\) is absent: no other added
concept realizes it.  Thus each addition, in any order, is ample
with acquired support \(\{e,f\}\), by the one-concept extension
criterion.  This proves the shadow formula.

At every step ordinary nonpeelability is preserved by
Proposition~\ref{prop:damp-rank-two-extension-invariance}.
Proposition~\ref{prop:damp-one-concept-atlas} then makes each enlarged
family forbidden.  Induction proves forbidden-minor status for
\(X_\Gamma\).

Each \(A_i\) is two-connected.  After deleting \(0\), the two
remaining block graphs are connected, and any one added concept
joins them through its two neighbors.  Every other added concept
also has neighbors in these connected blocks.  If an old concept
other than \(0\) is deleted, the surviving old graph remains
connected through \(0\), and each added concept retains at least
one neighbor.  If an added concept is deleted, the connected old
graph \(X\) remains, with every other added concept attached to it.
These cases prove two-connectivity.

The added concepts have their displayed squares as unique maximal
cubes.  Suppose that each factor has only its root as a corner.
Then \(X\) is cornerless by Proposition~\ref{prop:damp-vertex-gluing}.
Any old concept becoming a corner would have a maximal cube containing
an added concept, hence that cube would be one of the displayed squares.
The old vertices of such a square are \(0,0^{\{e\}},0^{\{f\}}\).
The root has degree at least four in \(X\), since each factor is
two-connected.  Either other vertex has degree at least two in its
old block and gains a new neighbor, so has degree at least three.
None can be a corner of a square.  Therefore the new concepts are
exactly the corners.

For the converse, put \(S=O\setminus\{c\}\).
Corollary~\ref{cor:damp-low-degree-corner-descent} makes \(S\)
an ample forbidden pc-minor.  Let \(u,v\) be the neighbors of
\(c\), and let \(w\) be the opposite vertex of its square.
If \(r\) is a cut vertex of \(S\), the graph \(O\setminus\{r\}\)
can be connected only if \(r\notin\{u,v\}\) and \(u,v\)
lie in different components of \(S\setminus\{r\}\): the added
vertex \(c\) has no other neighbors.  Thus \(r\) separates
\(u\) from \(v\) in \(S\).  The path \(u,w,v\) lies in
\(S\), so \(r=w\).  The cut-vertex theorem recovers the two
rooted factors of \(S\) at \(w\); the neighbors \(u,v\)
belong to different blocks.  Translate \(w\) to zero.  Their edge
coordinates are in the two disjoint block coordinate sets, and
\(c=0^{\{e,f\}}\).  This is the construction with one selected
pair.  Its recovered root, blocks, and neighbors are determined by
\((O,c)\), proving uniqueness up to interchange.
\end{proof}

The theorem gives a construction and recovery rule for two-connected
obstructions whose deletion of a marked degree-two corner creates a
cut vertex.  The rooted factors still require intrinsic admissibility.
Every displayed completion has an obstruction-preserving corner descent
back to \(X\), so these examples are not corner-descent terminal.

\begin{remark}[A two-connected obstruction with one degree-two corner]
\label{ex:damp-square-completed-578}
Use the \(289\)-concept factor \((B,29)\) from
Remark~\ref{rem:damp-noncorner-rooted-factor}, translated to root zero,
and take two copies in coordinate blocks \(0,\ldots,11\) and
\(12,\ldots,23\).  The cornerless double \(X\) has \(577\)
concepts.  The root directions in either factor are \(1,2,8\).
Completing just the square on directions \(1,13\) adds the concept
\[
 2^1+2^{13}=8194.
\]
The result has \(578\) concepts, isometric dimension \(24\),
VC dimension three, and exactly one corner, of degree two.
It is two-connected, and deleting that corner returns the cut-vertex
obstruction \(X\).

More generally, any nonempty subset of the nine pairs between
\(\{1,2,8\}\) and \(\{13,14,20\}\) gives a two-connected
forbidden pc-minor.  A subset of size \(k\) gives \(577+k\)
concepts and exactly \(k\) degree-two corners.  This describes
\(511\) distinct families in the fixed labelled embedding; no
count of their isomorphism classes is asserted.

Degree two is the least possible corner degree in an ample forbidden
pc-minor.  Such a family is connected and is not a singleton.  If it
had a leaf, every geodesic from that leaf would begin with its unique
incident coordinate.  Isometry forces all other concepts to have the
opposite value in that coordinate, making the coordinate peripheral.
This contradicts Proposition~\ref{prop:damp-periphery-free}.

The checker \path{verify_damp_square_completed_blocks.py} verifies
the one-square and nine-square examples.  It checks the exact shadows
using block projections and mixed supports, two-connectivity, marked
recovery, all \(144\) elementary minor peelings, and complete
nonpeelability certificates with \(2\) and \(512\) states.
The report is \path{damp_square_completed_blocks_verification.json}.
\end{remark}


\subsection{Rooted reductions and attachment examples}
\label{sec:damp-rooted-attachment-examples}

\begin{proposition}[A three-core reduction retaining the root and repair cube]
\label{prop:damp-root-protected-three-core}
Let \((A,a)\) be a rooted factor of
Theorem~\ref{thm:damp-cut-vertex-obstructions}.  Repeatedly delete
nonroot vertices of degree less than three.  Each deletion is a
degree-two corner deletion preserving the rooted-factor conditions.
The final family is the unique largest induced subgraph containing
\(a\) in which every nonroot vertex has degree at least three.

If \(a\) is a degree-two corner and an adjacent ample support-three
addition restores peeling to \(a\), the same addition restores
peeling in this final factor.  Its whole punctured three-cube survives,
and the root is the final factor's only degree-two vertex.
Thus the equivalent repair profiles of
Corollary~\ref{cor:damp-adjacent-corner-minimal-repair-test}, if they
exist, can be realized with this additional degree condition.
\end{proposition}

\begin{proof}
A rooted factor is two-connected.  The local degree-two fact in the
preceding proof therefore makes every degree-two vertex a corner;
no vertex has smaller degree.  Deleting a nonroot such corner
preserves the factor conditions by
Corollary~\ref{cor:damp-rooted-irreducibility}.  Induction justifies
the entire process.  Its uniqueness follows by the same containment
argument as for the ordinary three-core, keeping the root throughout.

For the repair assertion, write \(H=A\cup\{c\}\), and let
\(C\) be the completed three-cube at the adjacent corners \(a,c\).
We claim that a nonroot degree-two vertex \(x\) of \(A\) lies
outside \(C\).  The only vertices of \(C\setminus\{c\}\)
with fewer than three neighbors inside this punctured cube are the
three neighbors of \(c\).  One is \(a\).  If \(x\) were one
of the other two, then \(c,x\) would be adjacent degree-three
corners of \(H\).
Put \(G=H\setminus\{c,x\}\).  Both individual additions of
\(c,x\) to \(G\) complete squares of the same support.
Attach the fixed rooted factor of
Remark~\ref{rem:damp-noncorner-rooted-factor} at \(a\), on fresh
coordinates.  The attachment makes ordinary peelability equivalent
to peeling to \(a\), by Proposition~\ref{prop:damp-vertex-gluing}.
It changes neither of the two square additions, since \(a\) is
distinct from their tips.  Proposition~\ref{prop:damp-square-pair-invariance}
therefore implies that \(G\) peels to \(a\), since \(H\) does.
Adding back the corner \(x\) would make \(A\) peel to \(a\),
a contradiction.  This proves the claim.

For a degree-two vertex \(x\) outside \(C\), the addition of
\(c\) gives it no new neighbor.  Its square in \(A\) is still
its unique maximal cube in \(H\), so it is a degree-two corner
there as well.  Deleting \(x\) preserves peeling of \(H\) to
\(a\), by item~3 of Proposition~\ref{prop:damp-rooted-operations}.
In \(A\), the deletion preserves the rooted-factor conditions,
and it leaves all of \(C\setminus\{c\}\) intact.  The support
of the addition remains the same three-coordinate support, since
its punctured cube remains and deleting a concept cannot shatter
that support.  Repeating this argument proves the assertion at
every stage and in the final factor.  The retained square at \(a\)
keeps its degree equal to two.
\end{proof}

For a concrete two-stage reduction, let \(X\) be the
\(577\)-concept double in Remark~\ref{ex:damp-square-completed-578}.
Add \(8194=2^1+2^{13}\) and then
\(8198=2^1+2^2+2^{13}\).  These complete squares on
\(\{1,13\}\) and \(\{2,13\}\), respectively.  The resulting
\(579\)-concept family is two-connected and has just one corner,
\(8198\), of degree two.  Deleting it exposes \(8194\) as the
next degree-two corner; deleting that concept gives the three-core
\(X\).  The first deletion alone still leaves a two-connected
family.  Thus the three-core reduction covers examples beyond the
marked single-deletion recovery in
Theorem~\ref{thm:damp-square-completed-rooted-blocks}.
The checker \path{verify_damp_three_core.py} verifies this example,
its \(72\) elementary minor orders, and the canonical reductions
of the earlier one-square and nine-square examples.  It also checks
the existing nonresolving extension of the \(291\)-concept family
by \(186\): this is a \(292\)-concept three-core with two
degree-four corners, so three-cores need not be cornerless.  Its report is
\path{damp_three_core_verification.json}.

The attachment targets cannot be replaced by the list of acquired pairs.
For example, use the \(577\)-concept core \(K\) from
Remark~\ref{ex:damp-square-completed-578}.  The following admissible data
give a four-element diamond; integers in the target columns denote base
concepts, and letters denote abstract new concepts.
\[
\begin{array}{c|c|rrr|r}
 x&I_x&u_x&v_x&w_x&\phi(x)\\ \hline
 a&\{1,13\}&8192&2&0&8194\\
 b&\{2,13\}&a&6&2&8198\\
 d&\{1,14\}&24576&a&8192&24578\\
 t&\{2,14\}&d&b&a&24582
\end{array}
\]
Its relative deletion order satisfies \(t<b<a\) and \(t<d<a\),
with \(b,d\) incomparable, so there are six ample intermediate families.
Using the same four pairs but taking
\((u_x,v_x,w_x)=(2^{q_x},2^{p_x},0)\) in every row instead gives the
four concepts \(8194,8196,16386,16388\).  Its partial order is an
antichain, giving sixteen ample intermediates.  Both outputs have
\(581\) concepts, VC dimension three, the same full shattered-set
complex, and the same labelled three-core.  They are not even graph
isomorphic: the first has one corner and the second has four.

The checker \path{verify_damp_square_attachment_data.py} constructs both
families from these abstract data, independently recovers their data
from their orientations, and checks all \(32\) intermediate subsets.
It replays \(144\) new elementary minor orders and verifies complete
negative certificates with six and sixteen states, respectively.
The report is \path{damp_square_attachment_data_verification.json}.

The following lemmas constrain adjacent rank-three repairs, whose
existence is established in Theorem~\ref{thm:damp-adjacent-rank-three-repair}.
They use the square attachment construction to exclude a peeling that
first removes the corner opposite the retained root.

\begin{lemma}[Two opposite corner edges of a three-cube]
\label{lem:damp-opposite-edge-endpoint}
Let \(H\) be ample and contain a three-cube of support \(\{p,q,f\}\).
Write its vertices as
\[
 c,\quad a=c^{\{f\}},\quad z=c^{\{p,q\}},\quad
 w=c^{\{p,q,f\}},\quad
 u_i=c^{\{p\}\cup F_i},\quad v_i=c^{\{q\}\cup F_i}
 \quad(i=0,1),
\]
where \(F_0=\varnothing\) and \(F_1=\{f\}\).
If \(c,a,z,w\) are all degree-three corners of \(H\), then
\[
 H\text{ peels to }a
 \quad\Longleftrightarrow\quad
 H\setminus\{c\}\text{ peels to }a.
\]
\end{lemma}

\begin{proof}
Normalize \(c=0\).  In the restriction \((p,q)=(0,0)\), the edge
\(\{c,a\}\) is an isolated component: its two vertices have no
neighbors outside the displayed cube.  Ample restrictions are connected,
so this restriction consists of that edge alone.  The restriction
\((p,q)=(1,1)\) consists of \(\{z,w\}\) for the same reason.
Let \(U,V\) be the projections of the other two restrictions, with
\((p,q)=(1,0),(0,1)\), respectively.  They are ample and contain
the same edge \(E=\{e_0,e_1\}\), whose endpoints have \(f\)-bits
zero and one and all other remaining bits zero.

We prove the following explicit criterion for the asserted endpoint:
\begin{equation}
 U\text{ and }V\text{ peel to }e_1,
 \qquad
 \text{at least one of }U,V\text{ peels relatively to }E.
 \label{eq:damp-opposite-edge-endpoint-criterion}
\end{equation}
Suppose first that \(H\) peels to \(a\).  Its root-containing
restriction \(q=0\) is the peripheral expansion of \(U\) with
duplicated site \(E\), with the root in that site.  Root-preserving
minor closure and Proposition~\ref{prop:damp-rooted-operations}
give a peeling of \(U\) to \(e_1\).  Restricting \(p=0\) gives
the corresponding statement for \(V\).

Follow the peeling until its first deletion from
\(W=\{u_0,u_1,v_0,v_1\}\).  Before this time, deletions outside
the displayed cube induce corner deletions in \(U\) or \(V\)
retaining \(E\): such vertices have no neighbors in either of the
two other restrictions.  Both residual families are therefore ample.
Among \(c,a,z,w\), the vertex \(a\) is retained.  With all of
\(W\) still present, the only possible sets of deleted vertices from
\(\{c,z,w\}\) are
\[
 \varnothing,\quad\{c\},\quad\{z\},\quad\{w\},\quad\{z,w\}.
\]
Indeed, deleting \(c\) leaves neither \(z\) nor \(w\) a corner;
deleting \(z\) permits only \(w\) next among these vertices, and
deleting \(w\) permits only \(z\).  After both are gone, \(c\)
still has three neighbors but its incident cube is incomplete.

Consequently a lower vertex of \(W\) still has a neighbor in
\(\{c,z\}\), and an upper vertex still has the neighbor \(a\).
These neighboring vertices have no outside neighbors.  Thus the first
corner deleted from \(W\) has no neighbor in its own restriction
outside \(E\): an outside edge would require a missing square through
the retained neighboring corner.

Suppose this first boundary deletion is \(u_i\); the case of \(v_i\)
is obtained by interchanging \(p,q\) and \(U,V\).
Let \(U'\) be the residual family just before this deletion.
Its \(f=i\) section is connected and \(e_i\) is isolated there,
so that section is the singleton \(\{e_i\}\).
The remaining global family still peels to \(a\).
Restrict it to \(q=0\).  If \(i=0\), this is the upper section
of \(U'\), joined to \(a\) at \(u_1\), with possibly the leaf
\(c\) also attached to \(a\).  Contracting these attachment edges
gives that upper section rooted at \(e_1\), so root-preserving minor
closure makes it peel to \(e_1\).  If \(i=1\), the lower section remains, joined
to \(a\) by the path \(u_0,c,a\).  Here \(c\) must still be
present, since otherwise this nonempty restriction would be disconnected.
Contracting that path makes the lower section peel to \(e_0\).
In either case, reattaching the deleted singleton as a leaf at the
opposite endpoint allows that section peeling while retaining both
members of \(E\).  Prepending the earlier deletions in \(U\)
gives a relative peeling of \(U\) to \(E\).  This proves
\eqref{eq:damp-opposite-edge-endpoint-criterion}.

Conversely, suppose \eqref{eq:damp-opposite-edge-endpoint-criterion}
holds, and first assume that \(U\) peels relatively to \(E\).
In \(H\setminus\{c\}\), perform those deletions in its
\((1,0)\) restriction.  Then delete
\[
 u_0,\ z,\ w,\ u_1
\]
in that order.  The first three deletions are square corners and the
last is a leaf.  The remaining graph is \(V\) with \(a\) attached
as a leaf at \(v_1\).  Peel \(V\) to \(v_1\), then delete
\(v_1\), retaining \(a\).  If instead \(V\) supplies the
relative peeling, interchange the two restrictions.  Thus the criterion
implies that \(H\setminus\{c\}\) peels to \(a\).  Prefixing
\(c\) also peels \(H\) to \(a\), proving the equivalence.
\end{proof}

\begin{theorem}[Reduction to a switch between two square attachments]
\label{thm:damp-two-square-endpoint-switch}
Failure of the adjacent-corner endpoint property in
Corollary~\ref{cor:damp-adjacent-corner-minimal-repair-test} is equivalent
to the existence of the following data.
Take an ample family \(G\), distinct coordinates \(p,q,r\), and
\(b\in G\).  Put
\[
 a=b^{\{r\}},\qquad c=a^{\{p\}},\qquad z=a^{\{q\}},
 \qquad C=Q_{\{p,q,r\}}(a).
\]
Require
\begin{equation}
 G\cap C=C\setminus\{a,c,z\},\qquad
 \{p,r\},\{q,r\}\notin\operatorname{Sh}(G).
 \label{eq:damp-two-square-switch-data}
\end{equation}
Require in addition that the two families
\[
 S=G\cup\{z,a\},\qquad T=G\cup\{c,a\}
\]
have different endpoint behavior: \(T\) peels to \(a\) and
\(S\) does not.

For all data satisfying \eqref{eq:damp-two-square-switch-data}, both
\(S\) and \(T\) are ample square extensions of \(G\), adding
the same two supports in opposite orders, and
\[
 \operatorname{Sh}(S)=\operatorname{Sh}(T)
 =\operatorname{Sh}(G)\mathbin{\dot\cup}
   \bigl\{\{p,r\},\{q,r\}\bigr\}.
\]
Their union \(H\) is ample; \(a,c,z\) are degree-three corners
of its completed cube.  In an endpoint-changing example, \(G\)
is ordinarily peelable.  One may choose such an example with \(S\)
a rooted factor whose root \(a\) is its only degree-two vertex.
The vertex \(z\) is then a noncorner degree-three neighbor of \(a\)
in \(S\).
\end{theorem}

\begin{proof}
First check the construction.  The five supplied cube vertices give
punctured squares at \(c\) and \(z\), on supports \(\{q,r\}\)
and \(\{p,r\}\), respectively.  The distinct-pair construction of
Theorem~\ref{thm:damp-square-attachment-data} makes \(G+c\),
\(G+z\), and \(G+c+z\) ample.  It also gives the two square
addition sequences \(z,a\) and \(c,a\), with respective support
sequences
\[
 \{p,r\},\{q,r\}\quad\text{and}\quad
 \{q,r\},\{p,r\}.
\]
These have exactly the prescribed punctured squares, so the construction
theorem proves the claims about \(S,T\) and their shadows.
Adding \(a\) to \(G+c+z\) completes its three-cube.  Its support
\(\{p,q,r\}\) is not yet shattered, and the punctured-cube lemma
excludes any further incident direction, as in that construction theorem.
The one-concept criterion makes \(H\) ample.  Its three displayed
vertices have exactly their three cube neighbors, so they are corners.
If \(T\) peels to \(a\), prefixing \(z\) peels \(H\) to
\(a\).  Failure in \(S=H-c\) is therefore an adjacent rank-three
repair.  Ordinary rank-two invariance also makes \(G\) peelable.

For the other direction, choose an adjacent rank-three repair \(H,S=H-c\)
of minimum cardinality, with retained root \(a\), and let \(C\)
be the common three-cube at \(a,c\).  A successful peeling cannot
start at \(a\), which is retained, or at \(c\), which would leave
the bad side \(S\).  It cannot start outside \(C\): such a first corner is also a corner of \(S\),
and deleting it leaves a smaller repair with the whole cube intact.
Nor can it start at a neighbor \(x\ne a\) of \(c\) in \(C\).
In that case \(c,x\) are adjacent degree-three corners, and their
simultaneous removal leaves an ample family to which their individual
additions are squares on the same support.  Attach the fixed bad rooted
factor at \(a\).  Propositions~\ref{prop:damp-vertex-gluing} and
\ref{prop:damp-square-pair-invariance} imply that the smaller family
peels to \(a\); adding back \(x\) would make \(S\) do so.

Write the cube directions as \(\{p',q',f\}\), where
\(a=c^{\{f\}}\), and put
\(z'=c^{\{p',q'\}}\), \(w'=z'^{\{f\}}\).
We claim that a successful peeling can start at one of
\(a^{\{p'\}},a^{\{q'\}}\).
The only other possible first vertices are \(z'\) and \(w'\).
Suppose first that it starts at \(z'\), the antipode of \(a\).
A second deletion outside \(C\) could be moved to the front: \(z'\)
has no outside neighbor, deleting a nonneighbor cannot create a corner,
and the cube of \(z'\) survives that outside deletion.  This would
give the already excluded smaller repair.  Inside \(C\), a second
corner must neighbor \(z'\), since a nonneighbor retains all three
cube directions while their cube is incomplete.

The options are \(c^{\{p'\}},c^{\{q'\}},w'\).
For either of the first two, its remaining corner square contains
\(c\), which has no outside neighbor.  Thus it has no outside
neighbor either and was already a degree-three corner of \(H\).
Interchanging the first two adjacent corner deletions gives a peeling
starting at a neighbor of \(c\), which was excluded above.
For \(w'\), its remaining corner square contains \(a\); the
same reasoning makes it a degree-three corner of \(H\).
Now \(c,a,z',w'\) satisfy
Lemma~\ref{lem:damp-opposite-edge-endpoint}, contradicting the repair.
Thus \(z'\) cannot be first.

If the first vertex is \(w'\), a second outside deletion can again
move to the front.  A second inside deletion must be one of its neighbors
\(a^{\{p'\}},a^{\{q'\}},z'\).
Each one's remaining corner square contains \(c\), so it has no
outside neighbors and was already a corner of \(H\).
Interchanging the first two deletions makes that vertex first.
The option \(z'\) was just excluded; either other option proves
the claim.

Let \(z\ne c\) be this successful first corner adjacent to \(a\).
Name the directions of \(ac,az\) as \(p,q\), and the third as
\(r\).  Put \(G=H\setminus\{a,c,z\}\).
Deleting \(a\) first leaves \(c,z\) as square corners; deleting
both therefore makes \(G\) ample.  These three deletions remove
exactly the triple support and the two distinct pairs
\(\{q,r\},\{p,r\}\).  Thus
\eqref{eq:damp-two-square-switch-data} holds with \(b=a^{\{r\}}\).
The successful first deletion of \(z\) makes \(T=H-z\) peel
to \(a\), whereas \(S=H-c\) does not.

Finally, \(S\) is ordinarily peelable by the adjacent square-pair
argument in Corollary~\ref{cor:damp-adjacent-corner-minimal-repair-test}.
Proposition~\ref{prop:damp-minimal-rank-three-endpoint-repair} would
produce a smaller repair if \(S\) were not a rooted factor, and
Proposition~\ref{prop:damp-root-protected-three-core} would do so if
it had another degree-two vertex.  Minimality excludes both possibilities.
The vertex \(z\) is not adjacent to \(c\), so it retains all three
incident directions in \(S\), but its cube now lacks \(c\).
It is therefore a noncorner of degree three adjacent to the root.
\end{proof}

The switch can change the endpoint: Theorem~\ref{thm:damp-adjacent-rank-three-repair}
supplies a $619$-word base and two $621$-word completions with opposite
outcomes, one a rooted factor. Intrinsic characterization of all
endpoint-changing switches remains open; existence is no longer the
unresolved question. The structural equivalence is not a complete
admissibility test.

The checker \path{verify_damp_two_square_endpoint_switch.py} validates
the opposite-edge criterion on two positive connectors and a connector
with two bad relative edges.  The latter has complete negative endpoint
certificates with \(25\) and \(125\) states before and after restoring
the missing cube corner.  It also checks the five possible sets of removed
cube vertices and all \(20\) first-boundary corner exclusions.
For the switch itself, the existing \(71\)-concept example gives
\(a=29,c=157,z=25\), a \(68\)-concept base, and two square
supports with bitmasks \(6,130\).  Both \(70\)-concept sides peel
to \(29\), so this is a positive control.  All \(20\) stored
endpoint orders, including the two fixed edge tests, pass.  The report is
\path{damp_two_square_endpoint_switch_verification.json}.

\begin{theorem}[A two-square switch with only small lower supports]
\label{thm:damp-small-lower-switch}
Assume the geometric switch conditions of
\eqref{eq:damp-two-square-switch-data}.  Flip coordinate values so that
\(a=0\), and identify words with their sets of nonzero coordinates.
Delete coordinate \(r\) from the two sections of \(G\), writing
\[
 G=(D\times\{0\})\cup(U\times\{1\}),\qquad
 E=D\cap U,\quad B=D\cup U,\quad d=\{p,q\}.
\]
Then \(D,U,E,B\) are ample, \(d\in E\), every word of \(D\)
has \(p=q=1\), and \(U\) contains the square
\(\{0,\{p\},\{q\},d\}\).  Moreover,
\begin{equation}
 \operatorname{Sh}(E)=\operatorname{Sh}(D)\cap\operatorname{Sh}(U),
 \qquad
 \operatorname{Sh}(B)=\operatorname{Sh}(D)\cup\operatorname{Sh}(U).
 \label{eq:damp-switch-section-shadows}
\end{equation}
If
\begin{equation}
 |I|\le2\quad\text{for every }
 I\in\operatorname{Sh}(D)\setminus\operatorname{Sh}(E),
 \label{eq:damp-switch-small-lower-gap}
\end{equation}
then the two switch sides \(S,T\) have the same endpoint behavior:
\begin{align}
 S\text{ peels to }a
 &\quad\Longleftrightarrow\quad T\text{ peels to }a\notag\\
 &\quad\Longleftrightarrow\quad
 U\text{ peels to }0\ \text{ and }\ E\text{ peels to }d\notag\\
 &\quad\Longleftrightarrow\quad
 B\text{ peels to }0\ \text{ and }\ D\text{ peels to }d.
 \label{eq:damp-switch-small-lower-endpoints}
\end{align}
In particular, a switch with \(D\subseteq U\) cannot change the
endpoint.  Every endpoint-changing switch requires a support of size
at least three in \(\operatorname{Sh}(D)\setminus\operatorname{Sh}(U)\).
Such a support avoids all three coordinates \(p,q,r\).
If the bad side \((S,a)\) is a rooted factor, there must also be
a support of size at least three in
\(\operatorname{Sh}(U)\setminus\operatorname{Sh}(D)\).
\end{theorem}

\begin{proof}
The five specified vertices of \(G\) realize the three traces
\(01,11,10\) on each of \(\{p,r\}\) and \(\{q,r\}\).
Since neither pair is shattered, the trace \(00\) is absent globally
on each pair.  Thus \(r=0\) forces \(p=q=1\).  The same five
vertices give \(d\in E\) and the stated square in \(U\).
Sections, contractions, and reductions of an ample family
are ample, so \(D,U,B,E\) are ample.

For completeness, a support avoiding \(r\) is shattered by \(G\)
exactly when it is shattered by \(B\).  A support \(I\cup\{r\}\)
is shattered by \(G\) exactly when both sections shatter \(I\).
Strong shattering in \(G\) then supplies a full
\(I\cup\{r\}\)-cube, whose projection is an \(I\)-cube in \(E\).
The reverse implication follows from \(E\subseteq D,U\).  Hence
the first equality in \eqref{eq:damp-switch-section-shadows} holds and
\[
 \operatorname{Sh}(G)=\operatorname{Sh}(B)
 \mathbin{\dot\cup}
 \{I\cup\{r\}:I\in\operatorname{Sh}(E)\}.
\]
An \(I\)-cube in \(G\) for a support avoiding \(r\) lies entirely
in one section.  Ampleness therefore also gives the second equality
in \eqref{eq:damp-switch-section-shadows}.

Assume \eqref{eq:damp-switch-small-lower-gap}.  By
Theorem~\ref{thm:damp-relative-small-supports}, delete
\(D\setminus E\) from \(D\) by corners of degree at most two.
These deletions lift to both \(S\) and \(T\), with the same degrees.
Indeed, each deleted word has no upper mate, and the two added lower
vertices form the path \(0,\{q\},d\) in \(S\), or
\(0,\{p\},d\) in \(T\), attached to \(D\) only at
\(d\in E\).  Thus no deleted word gains a neighbor from either
the upper section or the path.  Its entire incident cube remains in
the current lower section.  Proposition~\ref{prop:damp-rooted-operations}
preserves the old endpoint \(a\) at every deletion.

Write \(h=q\) for \(S\), or \(h=p\) for \(T\).
The remaining family is
\[
 ((E\cup\{0,\{h\}\})\times\{0\})
 \cup(U\times\{1\}).
\]
Its lower section is contained in its upper section.  The upper-root
case of the peripheral criterion in
Proposition~\ref{prop:damp-rooted-operations}, with the two levels
interchanged, makes its endpoint at \((0,0)\) equivalent to
\(U\) peeling to \(0\) and \(E\cup\{0,\{h\}\}\) peeling
to \(0\).  The latter is equivalent to \(E\) peeling to \(d\):
necessity follows by contracting the two edges of the attached path,
and sufficiency follows by peeling \(E\) to \(d\), then deleting
\(d,\{h\}\) and retaining \(0\).
This proves the first two equivalences in
\eqref{eq:damp-switch-small-lower-endpoints} for both choices of \(h\).

Theorem~\ref{thm:damp-relative-small-supports} also equates the
endpoint \(d\) in \(D\) and \(E\).  By
\eqref{eq:damp-switch-section-shadows},
\[
 \operatorname{Sh}(B)\setminus\operatorname{Sh}(U)
 =\operatorname{Sh}(D)\setminus\operatorname{Sh}(E).
\]
The same theorem applied to \(U\subseteq B\) equates their endpoint
\(0\), proving the last equivalence.  When \(D\subseteq U\),
the support difference is empty.  In every other endpoint-changing
case, the contrapositive supplies a support of size at least three
in that difference.  Coordinates \(p,q\) are constant on \(D\),
and coordinate \(r\) has been removed, so this support avoids all three.

Finally suppose \((S,a)\) is a rooted factor.  Its section opposite
\(r=0\) is \(U\), and the overlap is
\(E\cup\{0,\{q\}\}\).  The attached path uses only the
coordinates \(p,q\), which are constant on \(E\), so
\[
 \operatorname{Sh}(E\cup\{0,\{q\}\})
 =\operatorname{Sh}(E)\cup\bigl\{\{p\},\{q\}\bigr\}.
\]
If every support in \(\operatorname{Sh}(U)\setminus
\operatorname{Sh}(E)\) had size at most two, then
Theorem~\ref{thm:damp-relative-small-supports} would peel \(U\)
relatively to this overlap.  This contradicts
Proposition~\ref{prop:damp-rooted-section-obstructions}.
Hence there is a support of size at least three in
\(\operatorname{Sh}(U)\setminus\operatorname{Sh}(E)\), which equals
\(\operatorname{Sh}(U)\setminus\operatorname{Sh}(D)\) by
\eqref{eq:damp-switch-section-shadows}.
\end{proof}

The existing \(68\)-concept switch base has section sizes
\(|D|=|E|=11\) and \(|U|=57\), so its equal positive endpoints
are explained by the peripheral case.  The checker
\path{verify_damp_switch_section_supports.py} also treats two fixed
nonperipheral switches with a three-concept lower square attachment.
One has positive endpoints on both sides.  In the other, both sides
have a root-preserving contraction onto the bad rooted factor from
Remark~\ref{rem:damp-noncorner-rooted-factor}, and both endpoints fail.  These controls check the
two possible signs of the criterion; neither is an endpoint-changing
switch.  The remaining case requires supports of size at least three
exclusive to each section, with the lower support outside the switch
cube.  Their compatibility with rooted minimality remains open.

\begin{lemma}[Four corners along a three-cube path]
\label{lem:damp-four-corner-path}
Let \(H\) be ample and contain a three-cube \(C\).  Suppose that
\(c,a,z,v\) are four degree-three corners whose three consecutive
edges form a shortest path between opposite vertices of \(C\).
Then
\[
 H\text{ peels to }a
 \quad\Longleftrightarrow\quad
 H\setminus\{c\}\text{ peels to }a.
\]
\end{lemma}

\begin{proof}
The reverse implication follows by deleting \(c\) first.
For the forward implication, suppose there is a counterexample of
minimum cardinality.  Normalize \(a=0\) and write
\[
 c=\{p\},\quad z=\{q\},\quad v=\{q,r\},\quad
 b=\{r\},\quad u=\{p,r\},\quad d=\{p,q\},\quad
 w=\{p,q,r\}.
\]
The displayed words have zero entries outside \(P=\{p,q,r\}\).
The four assumed corners have no neighbors outside \(C\).
We first restrict the possible deletions inside this cube, and then
use a peripheral restriction to contradict minimality.

Neither \(d\) nor \(u\) can be a corner of \(H\).
Indeed, if \(x\ne a\) is a corner adjacent to \(c\), then
both \(c,x\) have degree three and
\(H\setminus\{c,x\}\) is ample.  The two individual additions
are squares on the same support.  Glue the fixed bad rooted factor
from Remark~\ref{rem:damp-noncorner-rooted-factor} at \(a\).
The vertex-gluing criterion and
Proposition~\ref{prop:damp-square-pair-invariance} make peeling
\(H\) to \(a\) equivalent to peeling
\(H\setminus\{c,x\}\) to \(a\).  Adding back \(x\) preserves
this endpoint by Proposition~\ref{prop:damp-rooted-operations},
contradicting the failure in \(H-c\).
Nor can \(w\) be a corner of \(H\): the two opposite edges
\(\{c,a\}\), \(\{w,v\}\) would then satisfy
Lemma~\ref{lem:damp-opposite-edge-endpoint}.

A successful first deletion outside \(C\) would leave a smaller
counterexample with all four marked corners intact.  Such a corner
is also a corner of \(H-c\), so root failure survives its deletion.
The first deletion is therefore \(z,v\), or \(b\).
If it is \(v\), the second cannot lie outside \(C\), since it
would commute to the front.  Here and below, consecutive nonadjacent
corner deletions commute: the second is already a corner before the
first, and it cannot lie diagonally in the first corner's cube, since
that missing first vertex would prevent it from being a corner afterward.
Inside \(C-v\), a second corner is one of \(z,b,w\).
Its remaining square contains \(c\), so it has no outside neighbor
and was already a corner of \(H\).  The option \(w\) is excluded;
interchanging the first two adjacent cube-corner deletions gives a
successful first deletion at \(z\) or \(b\).
Interchanging \(q,r\) if needed, take this first vertex to be \(z\).
The vertex \(v\) is unchanged by this interchange and remains a
degree-three corner of \(H\).

The second deletion is now \(v\).  An outside deletion would again
commute to the front.  A corner of \(C-z\) other than the retained
root must be \(d\) or \(v\).  At \(d\), the remaining square
contains \(c\), so cornerhood would make \(d\) a degree-three
corner of \(H\), already excluded.  Thus the first two deletions
are \(z,v\).

A third deletion outside \(C\) would commute past both of these
vertices, which have no outside neighbors in \(H\).
The nonroot corners of the remaining cube can only be \(b,d,w\).
Their squares contain \(a\) or \(c\), so each would have no
outside neighbor and be a corner of \(H\).
The exclusions of \(d,w\) force the third deletion to be \(b\).
In particular, \(b\) is also a degree-three corner of \(H\).

The four vertices \(a,z,v,b\) form the entire restriction \(p=0\).
They form a square with no outside neighbor in that restriction;
connectedness of an ample restriction excludes any other component.
Writing \(Q\) for their projected \(\{q,r\}\)-square, we have
\[
 H=(Q\times\{0\})\cup(V\times\{1\}),\qquad Q\subseteq V,
\]
where the displayed levels now refer to coordinate \(p\).
The image of \(c\) is \(0\), a degree-two corner of \(V\).
The peripheral endpoint criterion gives a peeling of \(V\) to \(0\).
Its first corner cannot lie in \(Q\): it is not the retained root,
and any of the other three vertices has its entire square through
\(0\), hence no outside neighbor.  It would consequently make
\(d,u\), or \(w\) a corner of \(H\), contrary to the exclusions.

Let \(x\notin Q\) be this first corner of \(V\), regarded as a
vertex in the \(p=1\) section of \(H\).  It has no mate in the
other section, so it is also a corner of \(H\) and \(H-c\).
The peripheral criterion applied to \(V-x\) makes \(H-x\) peel
to \(a\).  But \(H-c-x\) still fails at \(a\), since otherwise
prefixing \(x\) would peel \(H-c\) to \(a\).
The four marked corners and their cube survive, giving the final
contradiction to minimality.
\end{proof}

\begin{corollary}[The first two upper cubes in a smallest repair]
\label{cor:damp-first-upper-switch-cubes}
Suppose an adjacent rank-three endpoint repair exists, and choose
\(H,S=H-c\) of minimum cardinality, rooted at \(a\).
Use the switch coordinates of
Theorem~\ref{thm:damp-two-square-endpoint-switch}, normalized as
\(a=0,c=\{p\},z=\{q\}\), and let \(U\) be its upper
\(r\)-section as in Theorem~\ref{thm:damp-small-lower-switch}.
A successful peeling of \(H\) can be chosen to begin
\[
 z,\qquad v=\{q,r\},\qquad x=\{q,r,e\}.
\]
There are nonempty sets \(J,F\), disjoint from \(P=\{p,q,r\}\),
with \(e\in J\) and \(F\cap J=\varnothing\), such that the
corner supports of these second and third deletions are respectively
\begin{equation}
 R=\{p,q\}\cup J,
 \qquad R'=(\{p,q\}\cup(J\setminus\{e\}))\cup F.
 \label{eq:damp-first-upper-switch-cubes}
\end{equation}
The two corner cubes project to full cubes in the original \(U\).
Their intersection is the facet of the first cube with \(e=1\).
In particular, \(U\) has a shattered support of size at least three
containing both \(p,q\), whereas the lower section has an exclusive
support of size at least three avoiding all of \(P\).
\end{corollary}

\begin{proof}
Put \(C=Q_P(0)\) and write
\[
 b=\{r\},\qquad u=\{p,r\},\qquad
 d=\{p,q\},\qquad w=\{p,q,r\}.
\]
Choose the successful first corner \(z\) using
Theorem~\ref{thm:damp-two-square-endpoint-switch}.
The nonroot neighbors \(d,u\) of \(c\) cannot be corners of
\(H\), by the pointed square-pair argument in the preceding proof.
After deleting \(z\), an outside first corner would commute ahead
of \(z\), contradicting minimality.  Inside the cube, the only
nonroot possibilities are \(d\) and \(v\).
At \(d\), the remaining square contains \(c\); hence it would
have no outside neighbor and would already be a corner of \(H\).
This excludes \(d\), so the second vertex is \(v\).

In \(H-z\), the corner at \(v\) contains the upper
\(\{p,q\}\)-square.  It has no \(r\)-neighbor, so its support
has the form \(R=\{p,q\}\cup J\), with \(J\cap P=\varnothing\).
If \(J\) were empty, restoring \(z\) would make \(v\) a
degree-three corner of \(H\).  Then \(c,a,z,v\) would satisfy
Lemma~\ref{lem:damp-four-corner-path}, contradicting the repair.
Thus \(J\ne\varnothing\).

After deleting \(z,v\), no nonroot vertex inside \(C\) can be
the next corner.  The vertices \(c,u\) retain all three cube
directions while their cube is incomplete.  Cornerhood at \(d\)
would again force it to be a corner of \(H\).
For every \(j\in J\), the upper cube at \(v\) supplies the
outside neighbors \(b^{\{j\}},w^{\{j\}}\).
The vertex \(b\) also retains its neighbor \(a\), which has no
outside neighbor, so the corresponding square at \(b\) is missing.
The vertex \(w\) retains directions \(q,r\); a corner cube at
\(w\) containing \(j\) would require \(c^{\{j\}}\), also
absent.  Thus neither \(b\) nor \(w\) is a corner.

The third vertex \(x\) is therefore outside \(C\).  Unless it
neighbors \(v\), its deletion commutes past \(v\) and \(z\),
giving a successful first deletion outside \(C\) and a smaller
repair.  Hence \(x=v^{\{e\}}\) for some \(e\in J\).
Its corner cube contains the full \(R\setminus\{e\}\)-face
opposite \(v\).  Direction \(e\) is absent after deleting \(v\),
and direction \(r\) is absent because the lower section has
\(p=q=1\) apart from the three added vertices at outside word zero.
Its support therefore has the form \(R'\) in
\eqref{eq:damp-first-upper-switch-cubes}, allowing for now \(F=\varnothing\).

If \(F\) were empty, \(x\) would already be a corner of
\(H-z\) on the same cube as \(v\), and also a corner of \(H\),
since \(z\) lies in the opposite \(r\)-section and has no outside
neighbors.  The deletions \(z,v,x\) could be reordered as
\(x,z,v\): deleting \(x\) leaves the cube at \(z\) intact,
and after deleting \(z,x\), the remaining face at \(v\) is a
corner cube.  This again contradicts minimality.  Thus \(F\ne\varnothing\).
Both cubes are in the upper section.  The second lies entirely in
\(e=1\), so it did not use the removed vertex \(v\) and was
already present in \(U\).  Their supports and positions give exactly
the stated intersection.  Finally, \(R\) is exclusive to \(U\)
because \(p,q\) are constant on the lower section; the lower exclusive
support follows from Theorem~\ref{thm:damp-small-lower-switch}.
\end{proof}

The four-corner path lemma excludes this configuration uniformly,
without assumptions on the supports of the remaining family.
The checker \path{verify_damp_four_corner_path.py} tests it on two
nonperipheral controls with exclusive supports of size three in both
sections: the \(36\)-concept example and its corner deletion have
positive endpoints, whereas the \(317\)-concept example and its
corner deletion contract to the known bad rooted factor.
It also checks every local corner exclusion in the proof.

The new upper-cube pattern itself is realizable: a \(23\)-concept
ample family has a successful prefix with corner supports
\(\{p,q,r\},\{p,q,e\},\{p,q,f\}\), where \(e\ne f\).
Both \(22\)-concept switch sides still peel to the root.
Thus the required successive cubes do not alone imply failure.
Whether this pattern can occur in an endpoint-changing switch with
a rooted minimal bad side remains open.









\subsection{Relative peeling and quantitative reductions}
\label{sec:damp-relative-quantitative}

Proposition~\ref{prop:damp-periphery-free} lets us replace an absolute cardinality question
by a smaller relative one.  Define the \emph{relative corner threshold}
\begin{equation}
 \rho_{\Damp}=\min\bigl\{|A\setminus B|:
   \varnothing\ne B\subsetneq A,\ A,B\in\Damp,
   \operatorname{Cor}(A)\subseteq B\bigr\},
 \label{eq:damp-relative-corner-threshold}
\end{equation}
with the value $+\infty$ if there is no such pair.  Thus a pair counted by
$\rho_{\Damp}$ consists of two dismantlable ample families, but every corner
of the upper family is trapped in the lower one.

\begin{proposition}[Two-wing reduction to a relative corner threshold]
\label{prop:damp-two-wing-relative-corner-threshold}
Let $H$ have minimum cardinality among all nonempty cornerless ample
families.  Then
\begin{equation}
 |H|\ge2\rho_{\Damp}+2.
 \label{eq:damp-two-wing-relative-corner-bound}
\end{equation}
Consequently, $\rho_{\Damp}\ge22$ would exclude every such family of order
less than $46$ in one step.
\end{proposition}

\begin{proof}
Delete the constant coordinates of $H$ and fix an active coordinate $e$.
Write
\[
 A_b=\rho_e^bH\quad(b\in\{0,1\}),\qquad
 B=A_0\cap A_1,\qquad D_b=A_b\setminus B.
\]
The two sections $A_0,A_1$ are ample, and the two-layer intersection lemma
makes $B$ ample.  Isometry of $H$ gives an $e$-edge, so $B$ is nonempty.
Every one of these three families has order smaller than $|H|$; hence
Lemma~\ref{lem:damp-least-cornerless-cardinality-minimality} puts all of
them in $\Damp$.

The family $H$ is a forbidden pc-minor for $\Damp$: every proper pc-minor is
ample and has smaller order, and therefore belongs to $\Damp$ by the same
lemma.  Proposition~\ref{prop:damp-periphery-free} now shows that $e$ is
nonperipheral.  Equivalently, both $D_0$ and $D_1$ are nonempty.

If $d\in D_b$ were a corner of $A_b$, then $(d,b)$ would have no incident
$e$-edge in $H$, and all its other incident cubes would be exactly the
incident cubes of $d$ in $A_b$.  Thus $(d,b)$ would be a corner of $H$, a
contradiction.  Hence
\[
 \operatorname{Cor}(A_b)\subseteq B
 \qquad(b\in\{0,1\}).
\]
Both nested pairs $B\subsetneq A_b$ occur in the minimum defining
$\rho_{\Damp}$, so $|D_b|\ge\rho_{\Damp}$.  Finally,
\[
 |H|=|A_0|+|A_1|
     =|D_0|+|D_1|+2|B|
     \ge2\rho_{\Damp}+2,
\]
because $B$ is nonempty.
\end{proof}

\begin{proposition}[A twelve-concept trapped interval]
\label{prop:damp-twelve-concept-trapped-interval}
The relative corner threshold satisfies
\begin{equation}
 \rho_{\Damp}\le12.
 \label{eq:damp-relative-corner-threshold-upper-twelve}
\end{equation}
More precisely, the smaller section $Q\subseteq Q_{11}$ of the
$291$-concept obstruction in
Proposition~\ref{prop:damp-hall-order-291-obstruction} contains a family
$B\in\Damp$ such that
\[
 |Q|=85,\qquad |B|=73,\qquad
 \operatorname{Cor}(Q)\subseteq B,
\]
and there is no ample family strictly between $B$ and $Q$.
\end{proposition}

\begin{proof}[Computer-assisted proof]
Let $K=P\cap Q$ be the $59$-concept overlap in the construction of
Proposition~\ref{prop:damp-hall-order-291-obstruction}.  In the decimal
encoding of eleven-bit words, put
\[
\begin{split}
(w_1,\ldots,w_{14})={}&
(2038,898,910,398,386,1024,1958,1926,\\
&\hspace{35mm}1974,902,390,1942,1152,1408)
\end{split}
\]
and $B_i=K\cup\{w_1,\ldots,w_i\}$.  The exact one-concept colon test in
Lemma~\ref{lem:damp-one-concept-colon-test} verifies that every $B_i$ is
ample.  The family $K$ belongs to $\Damp$, and $w_i$ is a corner of
$B_i$, because deleting it leaves the ample family $B_{i-1}$.  Reversing
these fourteen additions and then peeling $K$ therefore proves
$B=B_{14}\in\Damp$.  The section $Q$ also belongs to $\Damp$.

The section-corner calculation gives
$\operatorname{Cor}(Q)\subseteq K\subseteq B$, and
$|Q\setminus B|=85-(59+14)=12$.  Thus $B\subsetneq Q$ occurs in the
minimum defining $\rho_{\Damp}$.  As an independent minimality check, the
twelve acquired shattered supports are
\[
 14,22,26,28,518,536,1030,1048,1538,1540,1544,1552.
\]
Testing all $2^{12}-2=4094$ proper nonempty subsets of the twelve remaining
concepts against these supports finds no ample intermediate family.  The
reconstruction, the fourteen colon certificates, the terminal blocking
vertices, and the interval scan are recorded in
\path{damp_hall_relative_leaf_core.json} and reproduced by
\path{analyze_damp_hall_relative_leaf_core.py}.  The mathematical digest is
\[
 \texttt{e019f8c6ec0679ba930d06731c11988c98793fa5e26aba330739d8b503d721a0}.
\]
\end{proof}

\begin{proposition}[The sharp interval has an octahedral support sphere]
\label{prop:damp-twelve-concept-octahedral-support-sphere}
For the twelve-concept interval \(B\subsetneq Q\) in
Proposition~\ref{prop:damp-twelve-concept-trapped-interval}, let
\[
 Y_0=\{1,2,3,4,9,10\},
 \qquad
 P_1=\{1,2\},\quad P_2=\{3,4\},\quad P_3=\{9,10\}.
\]
All twelve acquired supports have order three, and their family is
\begin{equation}
 \mathcal H_0
 =\bigl\{T\in\tbinom{Y_0}{3}:
          P_i\subseteq T\text{ for some }i\in\{1,2,3\}\bigr\}.
 \label{eq:damp-twelve-support-pair-blocks}
\end{equation}
Consequently, the pure complex whose facets are the rank-three supports in
\(\binom{Y_0}{3}\setminus\mathcal H_0\) is
\begin{equation}
 \partial 2^{P_1}*\partial 2^{P_2}*\partial 2^{P_3}
 \cong S^0*S^0*S^0\cong S^2,
 \label{eq:damp-twelve-octahedral-support-sphere}
\end{equation}
the boundary complex of the three-dimensional cross-polytope.  The three
opposite vertex pairs of this octahedral sphere are precisely the direction
pairs of the three vertex-disjoint squares in the relative graph.
\end{proposition}

\begin{proof}
Decode the twelve masks displayed in the proof of
Proposition~\ref{prop:damp-twelve-concept-trapped-interval}.  They become
\[
 \begin{gathered}
 \{1,2,3\},\{1,2,4\},\{1,3,4\},\{2,3,4\},\\
 \{1,2,9\},\{3,4,9\},\{1,2,10\},\{3,4,10\},\\
 \{1,9,10\},\{2,9,10\},\{3,9,10\},\{4,9,10\}.
 \end{gathered}
\]
These are exactly the four
triples containing each of \(P_1,P_2,P_3\), and no triple contains two of
the disjoint pairs.  This proves
\eqref{eq:damp-twelve-support-pair-blocks}.  A three-subset of \(Y_0\)
contains none of the \(P_i\) exactly when it chooses one element from each
pair.  The eight such transversals are the facets of the join of the three
zero-spheres, proving
\eqref{eq:damp-twelve-octahedral-support-sphere}.  The residual-edge
reconstruction gives square directions
\(\{1,2\},\{3,4\},\{9,10\}\); it is recorded independently in
\path{damp_terminal_artinian_owner_paths_d12.json}.
\end{proof}

\begin{remark}[The owner-cycle sphere target]
\label{rem:damp-owner-cycle-sphere-target}
Proposition~\ref{prop:damp-twelve-concept-octahedral-support-sphere}
suggests a topological route to the two-connected theorem.  It would be
enough to prove that the lower and upper gate cycles of
Corollary~\ref{cor:damp-terminal-fibre-two-root-cycles} force a rank-three
support cycle containing a join of three nonempty zero-dimensional parts.
Such a join uses at least six active coordinates.  The open-ear inequality
\eqref{eq:damp-two-connected-coordinate-cost} would then give
\(|D|\ge12\) immediately.  Proposition
\ref{prop:damp-terminal-two-connected-global-completion} obtains the same
coordinate conclusion more directly from the positioned shore concepts
and forced shielding owners, with one exact five-coordinate completion
lemma.  Thus the support sphere is not needed for the orders-ten-and-eleven
exclusion, although the sharp interval still realizes the expected equality
geometry by the octahedral sphere in
\eqref{eq:damp-twelve-octahedral-support-sphere}.
\end{remark}

\begin{lemma}[Local carrier rules through residual degree three]
\label{lem:damp-local-passive-carrier-rules}
Let \(K\subsetneq A\subseteq Q_E\) be ample, put

\[
 R=A\setminus K,
 \qquad
 \mathcal H=\operatorname{Sh}(A)\setminus\operatorname{Sh}(K),
\]

and suppose that \(c\in R\) is \emph{doubly blocked}: neither
\(K\cup\{c\}\) nor \(A\setminus\{c\}\) is ample.  Then the following
statements hold.

\begin{enumerate}
\item Every compatible lower carrier of an inclusion-minimal support
  \(S\in\mathcal H\) and upper carrier of an inclusion-maximal support
  \(T\in\mathcal H\) meet in a full cube contained in \(R\), whose free
  coordinate set is \(T\setminus S\).
\item The concept \(c\) belongs to at least two lower carriers of minimal
  supports and at least two upper carriers of maximal supports.
\item If \(c\) has residual degree two, it marks an isolated support
  \(S_c\in\mathcal H\), meaning that \(S_c\) is both minimal and maximal,
  and exactly one of the two residual edge coordinates at \(c\) belongs to
  \(S_c\).  More precisely, if
  \[
   P_c=\{e:c^{\{e\}}\in K\},
   \qquad
   W_c=\{e:c^{\{e\}}\in R\},
  \]
  then
  \begin{equation}
   S_c=P_c\mathbin{\dot\cup}\{\lambda_c\}
   \qquad\text{for one }\lambda_c\in W_c.
   \label{eq:damp-degree-two-isolated-support-form}
  \end{equation}
  If adjacent degree-two concepts \(c,d\) have edge coordinate \(e\),
  then \(e\in S_c\mathbin\triangle S_d\).
\item If \(c\) has residual degree three, there are a minimal support
  \(S^-\) and a maximal support \(S^+\) such that

  \[
   S^-\subseteq S^+,
   \qquad |S^+\setminus S^-|\le1,
   \qquad c\in C^-_{S^-}\cap C^+_{S^+}.
   \tag{A4g1c}\label{eq:damp-degree-three-isolated-or-cover}
  \]

  Thus \(c\) marks either an isolated support or a minimal--maximal
  support cover.
\end{enumerate}
\end{lemma}

\begin{proof}
Put \(B=Q_E\setminus A\), which is ample by complement duality.  For
\(S\in\mathcal H\), the two trace-defect classes

\[
 U_S=Q_S\setminus K|_S,
 \qquad
 V_S=Q_{E\setminus S}\setminus B|_{E\setminus S}
\]

are nonempty and ample, and \(U_S\mathbin\Box V_S\subseteq R\).  Sauer
equality in the two projections shows that \(U_S\) is a singleton exactly
when \(S\) is minimal in \(\mathcal H\), and \(V_S\) is a singleton exactly
when \(S\) is maximal.  If \(u_S\) and \(v_T\) denote these singleton
traces, their prescriptions agree on their common domain precisely when

\[
 u_S|_{S\setminus T}=v_T|_{S\setminus T}.
\]

When they agree, all common extensions lie in \(R\) and form the full
\((T\setminus S)\)-cube.  This proves the first assertion.

Adjoining \(c\) to \(K\) shatters exactly the minimal relative supports
whose lower missing trace is \(c|_S\).  The number of these supports is one
exactly when \(K\cup\{c\}\) is ample.  Applying the complementary statement
to \(B\cup\{c\}=Q_E\setminus(A\setminus\{c\})\) gives the upper assertion.
Double blocking therefore gives carrier depth at least two on both shores.

Translate \(c\) to the zero word and partition its coordinate directions
as

\[
 P=\{e:c^{\{e\}}\in K\},\quad
 Q=\{e:c^{\{e\}}\in B\},\quad
 W=\{e:c^{\{e\}}\in R\}.
\]

The commuting existential and universal operations for lopsided systems
give two ample classes on \(W\),

\[
 L=\exists_Q\forall_P K,
 \qquad
 U=\exists_P\forall_Q B.
\]

The full-collar form of the blocker formula shows that \(L,U\) are
nonempty, disjoint, avoid the zero word, and contain respectively every
lower and upper obstruction signature.  When \(|W|=2\), the three nonzero
words form a path, so \(L\) and \(U\) are disjoint intervals in that path.
Their order selects opposite residual directions.  The two corresponding
minimal transversals give the isolated support
\[
 S_c=P\mathbin{\dot\cup}\{\lambda_c\}
\]
for one \(\lambda_c\in W\), proving
\eqref{eq:damp-degree-two-isolated-support-form}.  Projecting the
coordinate of an edge between two such
markers, and applying the same collision argument to the complementary
pair, gives the stated symmetric-difference law.  This is the
degree-two isolated-support argument in a form that retains every
coordinate constant on \(R\).

Suppose finally that \(|W|=3\).  An exact local calculation in \(Q_3\)
shows that any two nonempty disjoint ample classes avoiding zero have
disjoint transversals of their inclusion-minimal words.  There are \(63\)
such ample classes and \(672\) ordered disjoint pairs.  Shrinking the two
transversals to minimal transversals of the actual lower and upper
obstruction clutters preserves disjointness.  They are both nonempty, so
their union leaves zero or one of the three coordinates.  The blocker
formula now gives \(S^-\subseteq S^+\), with difference equal to that
unused set, and puts \(c\) in both displayed carriers.  This proves
\eqref{eq:damp-degree-three-isolated-or-cover}.  The calculation is
reproduced by
\path{verify_damp_degree_three_two_shore_dichotomy.py}; its digest is

\[
 \texttt{51a1cbfec7834c283c3a73e23c8df9f3cf7ce22852d0e9e8102665934cf86fc0}.
\]
\end{proof}

\begin{proposition}[Component-size amplification in a terminal interval]
\label{prop:damp-terminal-relative-component-size}
Let \(K\subsetneq A\) be ample families, put \(D=A\setminus K\), and
suppose that
\[
 K\cup\{d\}\quad\hbox{is nonample for every }d\in D.
\]
Then every connected component of the graph induced by \(D\) has at least
nine vertices.  In particular, if \(|D|\le17\), then this graph is
connected.
\end{proposition}

\begin{proof}
Let \(D_1,\ldots,D_s\) be the vertex sets of the connected components of
the graph induced by \(D\), and let \(\mathcal Q(A,K)\) denote the
squarefree Alexander-dual module of the relative Yang complex for
\((A,K)\).  This module has a linear resolution because both endpoints are
ample.  The
relative linear-presentation theorem
\cite{BousquetYangComplex26} says that its degree-one relations are of two
types: an annihilator supported on one generator, coming from a neighbour
in \(K\), or an edge relation between two adjacent members of \(D\).
Every such relation is supported inside one component \(D_i\).  Since
these relations generate the full presentation, one obtains a direct-sum
decomposition
\[
 \mathcal Q(A,K)
 \cong
 \bigoplus_{i=1}^s \mathcal Q(K\cup D_i,K).
 \tag{A4g1d}\label{eq:damp-terminal-relative-component-splitting}
\]
Each summand therefore has a linear resolution, so its relative Yang module
is Cohen--Macaulay.  The exact sequence from that relative module to the
Yang modules of \(K\cup D_i\) and \(K\), together with ampleness of \(K\),
makes \(K\cup D_i\) ample.

If \(|D_i|\le8\), the arbitrary-upper relative linear-quotient theorem
through eight facets \cite{BousquetYangComplex26} supplies an ordering of
the concepts of \(D_i\) whose first concept, say \(d\), can be adjoined
amply to \(K\).  This contradicts the hypothesis.  Hence \(|D_i|\ge9\)
for every \(i\).  Two components would consequently have total order at
least eighteen, which proves the final assertion.
\end{proof}

\begin{proposition}[Small cut layers have one block shape]
\label{prop:damp-small-terminal-cut-layer-domino-wedge}
Let \(K\subsetneq A\subseteq Q_X\) be ample, put \(D=A\setminus K\), and
suppose that every \(d\in D\) is doubly blocked:
\[
 K\cup\{d\}\quad\hbox{and}\quad A\setminus\{d\}
 \quad\hbox{are nonample}.
\]
If the graph induced by \(D\) has a cut vertex and \(|D|\le11\), then
\(|D|=11\).  Moreover, its block--cut tree consists of two six-vertex
dominoes sharing one articulation vertex, and the articulation has degree
two inside each domino.
\end{proposition}

\begin{proof}
Proposition~\ref{prop:damp-terminal-relative-component-size} makes the
residual graph connected.  It has minimum degree at least two.  Indeed, an
isolated residual concept teaches \(K\), by taking the first edge of an
\(A\)-geodesic to each member of \(K\).  At a residual leaf, the relative
leaf-deletion dichotomy gives either the same teaching conclusion or a
Cohen--Macaulay one-concept deletion.  The latter makes
\(A\setminus\{d\}\) ample, since \(K\) is ample.  Both alternatives
contradict double blocking.

The block--cut tree has at least two leaf blocks.  None is a bridge, since
its nonarticulation endpoint would have residual degree one.  Choose two
leaf blocks \(L_1,L_2\).  They are two-connected and
\begin{equation}
 |D|\ge |L_1|+|L_2|-1.
 \label{eq:damp-small-cut-two-leaf-block-count}
\end{equation}
Thus one of them has at most six vertices.

We first classify a two-connected induced-cube graph on at most six
vertices.  It is bipartite and has minimum degree two.  At order four the
only possibility is \(C_4\).  There is no possibility at order five:
the bipartition would have sizes two and three and the graph would be
\(K_{2,3}\), whereas two cube vertices have at most two common neighbours.
At order six, a bipartition of sizes two and four similarly gives the
impossible graph \(K_{2,4}\).  With bipartition sizes three and three, the
missing edges form a matching.  Deleting zero or one edge from \(K_{3,3}\)
again leaves a \(K_{2,3}\).  Deleting two or three independent edges gives,
respectively, the domino and \(C_6\).  Hence the only block shapes are
\[
 C_4,\qquad C_6,\qquad\hbox{and the domino}.
\]

It remains important where the articulation lies in the domino.  Every
nonarticulation vertex of residual degree two marks an isolated support by
Lemma~\ref{lem:damp-local-passive-carrier-rules}.  Such a support contains
exactly one of the two incident block directions; adjacent markers reverse
ownership of their common direction.  Compatible lower and upper carriers
meet in a full residual cube.  Applying only these necessary rules excludes
a leaf \(C_4\), a leaf \(C_6\), and a domino rooted at a degree-three
vertex.  A domino rooted at a degree-two vertex is the sole surviving
rooted shape.

For completeness, the finite check keeps coordinates outside the block
symbolically through inclusion of their passive support parts.  A carrier
cube through a nonarticulation vertex may not use a passive direction,
whereas a cube meeting the block only at its root is deliberately allowed
to continue inward.  The four first-witness searches visit respectively
\(5,11,11,4\) nodes.  Exhausting the active masks in the last search gives
exactly the four patterns
\[
 (2,1,4),\qquad(3,6,5),\qquad(4,1,2),\qquad(5,6,3)
\]
on its three degree-two nonroot vertices.  Thus the check is a relaxation
of the full carrier system, not a conditioning that discards passive
boundary data.  It is reproduced by
\path{verify_damp_small_leaf_block_carriers.py} and
\path{damp_small_leaf_block_carriers.json}; its digest is
\[
 \texttt{3af61b6a1eeae8680d6d442d897558b5937a53543855c04b66f6bab3565083d0}.
\]

Consequently the smaller leaf block is a six-vertex domino rooted at a
degree-two vertex.  Inequality
\eqref{eq:damp-small-cut-two-leaf-block-count} now makes the other leaf
block have order at most six, so it has the same rooted shape.  The two
blocks already account for eleven vertices.  Equality throughout forces
\(|D|=11\), makes their roots the same articulation, and leaves no further
block in the block--cut tree.
\end{proof}

\begin{proposition}[Terminal cut layers through eleven are impossible]
\label{prop:damp-terminal-cut-layer-through-eleven}
Let \(K\subsetneq A\) be ample families, put \(D=A\setminus K\), and
suppose that every \(d\in D\) is doubly blocked:
\[
 K\cup\{d\}\quad\hbox{and}\quad A\setminus\{d\}
 \quad\hbox{are nonample}.
\]
If \(|D|\le11\), then the graph induced by \(D\) is two-connected.
\end{proposition}

\begin{proof}[Computer-assisted proof]
Proposition~\ref{prop:damp-terminal-relative-component-size} makes the
residual graph connected, and the proof of
Proposition~\ref{prop:damp-small-terminal-cut-layer-domino-wedge} gives it
minimum degree at least two.  It remains to exclude a cut vertex.  If one
exists, that proposition forces \(|D|=11\) and writes the graph as two
six-vertex dominoes sharing an articulation which has degree two in each
domino.

We first reduce the coordinate labels without a graph catalogue.  Translate
the common root to zero.  A rooted domino has three distinct directions
\(a,b,c\), where \(a,c\) are incident with the root, and vertex set
\begin{equation}
 \{0,a,a+b,c,a+c,a+b+c\}\subseteq\mathbb F_2^E.
 \label{eq:damp-outer-rooted-domino-coordinates}
\end{equation}
The four root directions of the two dominoes are distinct, since the
blocks meet only at their root.  The middle direction of either domino
cannot be a root direction of the other: if, for example, \(b=a'\), then
the second root neighbour \(b\) is adjacent to the first-block vertex
\(a+b\), producing an edge between the two blocks away from the root.
Inducedness gives a contradiction.  Thus the two middle directions are
either equal or distinct.  These are the only two coordinate profiles:
five active coordinates in the \emph{shared-middle} profile and six in the
\emph{disjoint} profile.

For each profile, instantiate the passive-coordinate carrier formula from
the proof of
Proposition~\ref{prop:damp-nine-concept-passive-carrier-exclusion} with all
eleven relative supports.  In addition to its support order, lower and
upper carrier anchors, carrier depths, cross-carrier cubes, degree-two
markers, and degree-three dichotomy, impose adjacent-marker reversal and
one of the four active patterns certified above on each domino.  Eleven
formal passive coordinates suffice: replacing actual passive parts
\(P_1,\ldots,P_{11}\) by
\[
 \widehat P_i=\{j:P_j\subseteq P_i\}
\]
preserves exactly every passive inclusion and equality used by the
formula.  The one-concept deletion co-pairs are deliberately omitted.
Consequently both formulas are relaxations of every genuine doubly blocked
interval, rather than conditionings which suppress passive boundary data.

Both formulas are unsatisfiable.  Their DIMACS sizes and the statistics of
the independently checked proof cores are
\begin{center}
\begin{tabular}{lrrrr}
\toprule
profile & variables & clauses & core lemmas & resolution steps\\
\midrule
shared middle & \(12{,}126\) & \(130{,}526\) & \(311{,}417\)
 & \(20{,}710{,}244\)\\
disjoint & \(13{,}476\) & \(145{,}285\) & \(749{,}165\)
 & \(51{,}150{,}761\)\\
\bottomrule
\end{tabular}
\end{center}
The shared-middle proof uses no RAT lemma; the disjoint proof uses
\(26{,}857\).  The generator, manifest, proof cores, and aggregate checker
are
\begin{center}
\begin{tabular}{c}
 \path{solve_damp_outer_domino_wedge_carrier_cnf.py},\\
 \path{damp_outer_domino_wedge_carrier_cnf_manifest.json},\\
 \path{damp_outer_domino_wedge_*_core.drat},\\
 \path{verify_damp_outer_domino_wedge_carrier_proofs.py}.
\end{tabular}
\end{center}
The aggregate mathematical digest is
\[
 \texttt{962dba99701919e0be0c25c98956e0373ff7147392124d319d272eeb3fd47942}.
\]
The two coordinate profiles exhaust the cut-vertex conclusion of
Proposition~\ref{prop:damp-small-terminal-cut-layer-domino-wedge}, and both
are impossible.  Hence the connected residual graph has no cut vertex and
is two-connected.
\end{proof}

The connected conclusion still leaves many possible residual graphs.  The
next elementary bound shows that their two-connected parts cannot use many
coordinates unless they also use many vertices.  Here a coordinate is
\emph{active} on a cube subgraph if it labels one of its edges.

\begin{lemma}[Coordinate cost of a two-connected cube graph]
\label{lem:damp-two-connected-coordinate-cost}
Let \(G\) be a two-connected subgraph of a Boolean cube, and let \(Y\) be
the set of coordinates active on \(G\).  Then
\begin{equation}
 |V(G)|\ge 2|Y|.
 \label{eq:damp-two-connected-coordinate-cost}
\end{equation}
The subgraph need not be induced or isometric.
\end{lemma}

\begin{proof}
We first recall a short construction of an open-ear decomposition.  Start
with a cycle \(G_0\) of \(G\).  Suppose that a subgraph \(G_i\) with an
open-ear decomposition has already been constructed.  If a component of
\(G-V(G_i)\) remains, it has at least two distinct neighbours in \(G_i\):
otherwise its unique attachment vertex would be a cut vertex of \(G\).
A shortest path through the component between two such attachment vertices
is an open ear.  After every vertex has been added, add the remaining edges
as ears of length one.  This constructs
\[
 G_0\subseteq G_1\subseteq\cdots\subseteq G_s=G,
\]
where \(G_0\) is a cycle and \(G_i-G_{i-1}\) is the interior of a path
whose distinct endpoints belong to \(G_{i-1}\).

Every coordinate used on the initial cycle occurs a positive even number
of times, and hence at least twice.  If \(r_0\) coordinates occur there,
then
\[
 |V(G_0)|=|E(G_0)|\ge2r_0.
\]
Now consider an ear with \(\ell\) edges, and suppose that it introduces
\(k\) coordinates not used in the preceding connected subgraph.  All old
vertices have the same value in each such coordinate.  Since the two
endpoints of the ear are old, every new coordinate must therefore occur an
even positive number of times on the ear, accounting for at least \(2k\)
edges.  At least one further edge has an old coordinate.  Indeed, otherwise
the parity condition would make the two endpoints agree in every cube
coordinate, contrary to their being distinct vertices.  Thus
\[
 \ell\ge2k+1.
\]
The ear contributes \(\ell-1\ge2k\) new vertices.  Summing this inequality
over the ears, including the length-one ears for which \(k=0\), gives
\eqref{eq:damp-two-connected-coordinate-cost}.
\end{proof}

The next localization is useful beyond the two-connected argument.  It is
the one-concept form of the minimal-nonample structure theorem from the
Yang-complex dictionary.  We state both directions so that all later uses of
antipodal blockers refer to one source-of-truth result.

\begin{proposition}[Antipodal-pair localization of failed augmentation]
\label{prop:damp-failed-augmentation-antipodal-localization}
Let \(K\subseteq Q_X\) be ample and let \(d\in Q_X\setminus K\).  Then
\(K\cup\{d\}\) is nonample if and only if there are a coordinate set
\(T\subseteq X\), with \(|T|\ge2\), and
\(h=d^{\oplus T}\in K\) such that
\begin{equation}
 (K\cup\{d\})\cap Q_T(d)=\{d,h\},
 \qquad
 K\cap Q_T(d)=\{h\}.
 \label{eq:damp-failed-augmentation-antipodal-localization}
\end{equation}
Here \(Q_T(d)\) is the face obtained by fixing every coordinate outside
\(T\) as in \(d\).  In particular, \(d\) and \(h\) are the antipodal
vertices of this face.
\end{proposition}

\begin{proof}
If the displayed face exists, it is a nonample conditioning of
\(K\cup\{d\}\): two antipodal vertices do not form an ample family in a
cube of dimension at least two.  This proves the reverse implication.

Starting with the nonample class \(K\cup\{d\}\), condition coordinates
while a nonempty nonample conditioning remains, and stop at a
conditioning-minimal nonample class \(H\).  Every nonample conditioning
encountered contains \(d\), since otherwise it would be a conditioning of
the ample class \(K\).  Consequently
\[
 H\setminus\{d\}
\]
is a nonempty conditioning of \(K\), and is ample.

We claim that a minimal nonample class \(H\) with \(H\setminus\{d\}\)
ample has VC dimension one.  Let \(n\) be its number of free coordinates
and \(k=\operatorname{VC}(H)\).  The structure theorem for minimal
nonample classes \cite{BousquetYangComplex26} gives
\[
 |H|=2\sum_{i=0}^{k-1}\binom{n-1}{i},
 \qquad
 \operatorname{SSh}(H)=
 \{P\subseteq[n]:|P|\le k-1\}.
\]
Deleting a concept cannot create a strongly shattered support.  Since
\(H\setminus\{d\}\) is ample, its shattered and strongly shattered
supports coincide, and hence its VC dimension is at most \(k-1\).
Sauer's inequality therefore gives
\[
 2\sum_{i=0}^{k-1}\binom{n-1}{i}-1
 =|H|-1
 \le \sum_{i=0}^{k-1}\binom ni.
\]
Pascal's identity reduces this inequality to
\[
 \binom{n-1}{k-1}\le1.
\]
The same structure theorem gives \(1\le k\le n-1\), so necessarily
\(k=1\).  It now gives \(|H|=2\), and its antipodality conclusion says
that the two concepts are antipodes on all \(n\) free coordinates.

Let \(T\) be those free coordinates and restore the fixed assignment of
the conditioning.  Then \(H=\{d,h\}\), where
\(h=d^{\oplus T}\in K\), and the definition of the conditioning gives
\eqref{eq:damp-failed-augmentation-antipodal-localization}.
\end{proof}

\begin{lemma}[Antipodal deletion blockers descend along their faces]
\label{lem:damp-antipodal-deletion-blocker-face-descent}
Let \(A\subseteq Q_X\), let \(v\in A\), and suppose that for some
\(J\subseteq X\), with \(|J|\ge2\), one has
\begin{equation}
 A\cap Q_J(v)=Q_J(v)\setminus\{h\},
 \qquad h=v^{\oplus J}.
 \label{eq:damp-antipodal-deletion-blocker}
\end{equation}
For every proper \(S\subset J\), put \(w=v^{\oplus S}\).  Then
\begin{equation}
 A\cap Q_{J\setminus S}(w)
 =Q_{J\setminus S}(w)\setminus\{h\}.
 \label{eq:damp-antipodal-deletion-blocker-descent}
\end{equation}
In particular, if \(|J\setminus S|\ge2\), then
\(A\setminus\{w\}\) is nonample.
\end{lemma}

\begin{proof}
Every member of \(Q_{J\setminus S}(w)\) has the form
\[
 v^{\oplus(S\cup T)}
 \qquad(T\subseteq J\setminus S).
\]
This is a member of the \(J\)-face through \(v\), and it equals \(h\)
exactly when \(T=J\setminus S\).  Since \(S\ne J\), all other displayed
vertices belong to \(A\) by
\eqref{eq:damp-antipodal-deletion-blocker}.  This proves
\eqref{eq:damp-antipodal-deletion-blocker-descent}.

If \(|J\setminus S|\ge2\), the complement of \(A\setminus\{w\}\) has,
on this face, exactly the antipodal pair \(\{w,h\}\).  This conditioning is
nonample.  Conditioning and complement duality preserve ampleness, so
\(A\setminus\{w\}\) itself is nonample.
\end{proof}

We can now replace every two-connected residual graph of order ten or
eleven by one common boundary-ownership problem.  For a family
\(L\subseteq Q_Y\) and \(d\in Q_Y\setminus L\), write
\[
 N_L(d)=\{y\in Y:d^{\oplus y}\in L\}.
\]

\begin{proposition}[Five-coordinate normalization of the two-connected
terminal frontier]
\label{prop:damp-two-connected-passive-root-cover}
Let \(K\subsetneq A\subseteq Q_X\) be ample, put \(D=A\setminus K\), and
suppose that
\begin{equation}
 K\cup\{d\}\quad\hbox{is nonample for every }d\in D.
 \label{eq:damp-terminal-no-one-extension}
\end{equation}
Suppose moreover that the graph induced by \(D\) is two-connected and that
\(|D|\le11\).  Let \(Y\) be its active coordinate set and put
\(Z=X\setminus Y\).  Then the following statements hold.

\begin{enumerate}
\item One has \(|Y|\le5\), and all members of \(D\) have a common
  restriction \(\sigma\in Q_Z\).
\item Identifying \(D\) with its projection to \(Y\), put
  \[
   L=\{k|_Y:k\in K,\ k|_Z=\sigma\},
   \qquad B=L\mathbin{\dot\cup}D.
  \]
  Both \(L\) and \(B\) are ample, with the empty family allowed for \(L\),
  and the set
  \[
   R=\{d\in D:L\cup\{d\}\text{ is ample}\}
  \]
  is nonempty.
\item For every \(d\in R\), there are a passive coordinate \(z\in Z\)
  and a concept \(c\in D\) such that
  \begin{equation}
   c^{\oplus z}\in K,
   \qquad d^{\oplus z}\notin K,
   \qquad
   N_L(d)\cap\operatorname{Diff}(d,c)=\varnothing.
   \label{eq:damp-passive-boundary-blocks-active-root}
  \end{equation}
\end{enumerate}
In particular, if
\[
 S_z=\{c\in D:c^{\oplus z}\in K\}\qquad(z\in Z),
\]
then the nonempty passive boundary types \(S_z\) collectively block every
active teaching root \(d\in R\), in the precise sense of
\eqref{eq:damp-passive-boundary-blocks-active-root}.  Thus no enumeration
of two-connected graphs of order ten or eleven is needed: every such
terminal interval has the same five-coordinate passive-owner normal form.
\end{proposition}

\begin{proof}
The graph induced by \(D\) is connected, so every coordinate not in \(Y\)
is constant on \(D\); call the common restriction \(\sigma\).  Lemma
\ref{lem:damp-two-connected-coordinate-cost} gives
\[
 2|Y|\le |D|\le11,
\]
and hence \(|Y|\le5\).  Conditioning \(A\) and \(K\) on the assignment
\(\sigma\), and then deleting the constant coordinates, gives respectively
\(B\) and \(L\).  Conditioning preserves ampleness, so both families are
ample whenever they are nonempty.

If \(L=\varnothing\), every singleton \(\{d\}\) is ample and hence
\(R=D\ne\varnothing\).  Suppose that \(L\ne\varnothing\).  Apply the
five-coordinate relative-peeling theorem,
Theorem~\ref{thm:damp-five-coordinate-relative-peeling}, to \(L\subsetneq
B\).  In the resulting peeling, all members of \(D\) are removed before
the members of \(L\).  Immediately before the last member of \(D\) is
removed, the current ample family is \(L\cup\{d\}\) for some \(d\in D\).
Thus \(d\in R\), proving the second assertion.

Fix \(d\in R\).  The one-concept colon criterion says that \(d\) teaches
\(L\): for every \(\ell\in L\),
\begin{equation}
 N_L(d)\cap\operatorname{Diff}(d,\ell)\ne\varnothing.
 \label{eq:damp-active-root-teaches-conditioned-lower}
\end{equation}
On the other hand,~\eqref{eq:damp-terminal-no-one-extension} and the same
criterion give \(h\in K\) such that
\begin{equation}
 N_K(d)\cap\operatorname{Diff}(d,h)=\varnothing,
 \qquad
 N_K(d)=\{x\in X:d^{\oplus x}\in K\}.
 \label{eq:damp-full-lower-blocks-active-root}
\end{equation}

Because \(A\) and \(K\) are ample, their relative Alexander-dual module has
a linear resolution and in particular is linearly presented.  The interval
exit lemma therefore supplies a path from \(d\) inside
\(D\cap[d,h]\), followed by an edge \(cg\) with
\[
 c\in D\cap[d,h],\qquad g\in K\cap[d,h].
\]
The coordinate of this last edge cannot belong to \(Y\).  Indeed, in that
case \(g|_Z=\sigma\), so \(g|_Y\in L\).  Formula
\eqref{eq:damp-active-root-teaches-conditioned-lower} would then give a
coordinate in
\[
 N_L(d)\cap\operatorname{Diff}(d,g)
 \subseteq N_K(d)\cap\operatorname{Diff}(d,h),
\]
contradicting~\eqref{eq:damp-full-lower-blocks-active-root}.

Let \(z\in Z\) be the coordinate of \(cg\).  Since every member of \(D\)
has restriction \(\sigma\), one has \(g=c^{\oplus z}\in K\).  Moreover,
\(c\in[d,h]\) and~\eqref{eq:damp-full-lower-blocks-active-root} give
\[
 N_L(d)\cap\operatorname{Diff}(d,c)=\varnothing.
\]
Finally \(z\in\operatorname{Diff}(d,h)\).  If \(d^{\oplus z}\in K\), then
\(z\in N_K(d)\cap\operatorname{Diff}(d,h)\), another contradiction.
Thus \(d^{\oplus z}\notin K\), and all three conclusions in
\eqref{eq:damp-passive-boundary-blocks-active-root} follow.
\end{proof}

The antipodal blocker lemma upgrades the last part of the proposition from
one boundary edge to an entire peripheral cell.

\begin{lemma}[The isolated blocker is a peripheral active cell]
\label{lem:damp-active-root-peripheral-blocker-cell}
Use the hypotheses and notation of
Proposition~\ref{prop:damp-two-connected-passive-root-cover}, and let
\(d\in R\).  Then there are a passive coordinate \(z\in Z\), a nonempty
set
\[
 W\subseteq Y\setminus N_L(d),
\]
and an ample family \(C\subseteq D\cap Q_W(d)\) such that
\begin{equation}
 d,\ d^{\oplus W}\in C,\qquad
 (d^{\oplus W})^{\oplus z}\in K,
 \label{eq:damp-peripheral-blocker-cell-endpoints}
\end{equation}
with
\[
 d\notin S_z,\qquad d^{\oplus W}\in S_z,
 \qquad C\cap S_z=\{d^{\oplus W}\},
\]
and the conditioning of \(A\) to the face
\(Q_{W\cup\{z\}}(d)\) is the peripheral expansion
\begin{equation}
 A\cap Q_{W\cup\{z\}}(d)
 =C\mathbin{\dot\cup}\{(d^{\oplus W})^{\oplus z}\},
 \label{eq:damp-isolated-blocker-peripheral-expansion}
\end{equation}
whose duplicated site is the singleton \(\{d^{\oplus W}\}\subseteq C\).
The corresponding conditioning of \(K\) is the singleton
\(\{(d^{\oplus W})^{\oplus z}\}\).  Moreover, \(C\) lies in the unique
minimal active carrier of \(d\).
\end{lemma}

\begin{proof}
Apply
Proposition~\ref{prop:damp-failed-augmentation-antipodal-localization} to
\(K\cup\{d\}\), and write its isolated face as \(Q_T(d)\), with antipodal
member \(h=d^{\oplus T}\in K\).  The relative module of the ample pair
\(K\subsetneq A\) is linearly presented.  The interval-exit lemma gives a
path from \(d\) inside \(D\cap[d,h]\), followed by an edge \(ch\) into
\(K\cap[d,h]\).  By
\eqref{eq:damp-failed-augmentation-antipodal-localization}, the endpoint in \(K\) is
necessarily \(h\).

All members of \(D\) have passive restriction \(\sigma\).  The edge
\(ch\) cannot have an active coordinate.  Otherwise \(h\) would also have
passive restriction \(\sigma\), and hence \(h|_Y\in L\).  Since \(d\)
teaches \(L\), some coordinate of
\(N_L(d)\cap\operatorname{Diff}(d,h)\) would give a member
\(d^{\oplus y}\in K\cap Q_T(d)\).  It is different from \(h\), since
equality would make \(T=\{y\}\) and the minimal nonample conditioning the
ample one-cube.  This contradicts
\eqref{eq:damp-failed-augmentation-antipodal-localization}.  Thus \(ch\) has a
passive coordinate \(z\).  Since \(c\) has passive restriction \(\sigma\)
and \(h\) is adjacent to \(c\), this is the only passive coordinate in
\(T\).  Put
\[
 W=T\setminus\{z\}\subseteq Y.
\]
The set \(W\) is nonempty: an antipodal pair on one free coordinate is the
full one-cube and is ample, whereas the conditioning which produced
\(\{d,h\}\) is nonample.  We have
\[
 c=h^{\oplus z}=d^{\oplus W}\in D.
\]
If \(y\in W\cap N_L(d)\), then \(d^{\oplus y}\in L\subseteq K\) would be
a member of \(Q_T(d)\) different from \(h\), again contradicting
\eqref{eq:damp-failed-augmentation-antipodal-localization}.  Hence
\(W\cap N_L(d)=\varnothing\).

Put
\[
 C=A\cap Q_W(d).
\]
This is a nonempty ample conditioning of \(A\), and it is contained in
\(D\), because
\eqref{eq:damp-failed-augmentation-antipodal-localization} leaves no member of \(K\)
on that section.  Now condition \(A\) only on the coordinates outside
\(T\).  Its \(\sigma_z\)-section is \(C\), while its opposite section,
after deleting \(z\), is the singleton \(\{c\}\): there are no relative
concepts on the opposite passive side, and \(h=c^{\oplus z}\) is the only
member of \(K\) in the face.  This proves
\eqref{eq:damp-isolated-blocker-peripheral-expansion}.  The interval-exit
path puts both \(d\) and \(c=d^{\oplus W}\) in \(C\).
The identity \(c^{\oplus z}=h\in K\) gives \(c\in S_z\), while
\eqref{eq:damp-failed-augmentation-antipodal-localization} gives
\(d^{\oplus z}\notin K\), and hence \(d\notin S_z\).
More generally, if \(u\in C\cap S_z\), then \(u^{\oplus z}\) belongs to
\(K\cap Q_T(d)=\{h\}\).  Thus \(u=c\), proving
\(C\cap S_z=\{c\}\).

Finally put \(P=N_L(d)\).  Every member of \(Q_W(d)\) agrees with \(d\)
on \(P\), because \(W\cap P=\varnothing\).  This common \(P\)-trace is
the missing trace defining the unique minimal carrier of \(d\).
Consequently \(C\) lies in that carrier.
\end{proof}

The passive types in the preceding proposition are not arbitrary subsets of
\(D\).  Relative Cohen--Macaulayness couples all of them through the fine
numerator.  The following elementary consequence is the first Helly-type
constraint that we shall need.

\begin{lemma}[Common owner when the active lower slice is empty]
\label{lem:damp-empty-active-slice-passive-helly}
Let \(K\subsetneq A\subseteq Q_{Y\mathbin{\dot\cup}Z}\) be ample, put
\(D=A\setminus K\), and suppose that all members of \(D\) have the same
restriction \(\sigma\in Q_Z\).  Suppose moreover that
\[
 \{k|_Y:k\in K,\ k|_Z=\sigma\}=\varnothing.
\]
For \(z\in Z\), put
\[
 S_z=\{d\in D:d^{\oplus z}\in K\},
\qquad
 Z_0=\{z\in Z:S_z\ne\varnothing\}.
\]
If \(Z_0\ne\varnothing\), then
\begin{equation}
 \bigcap_{z\in Z_0}S_z\ne\varnothing.
 \label{eq:damp-empty-slice-passive-helly}
\end{equation}
\end{lemma}

\begin{proof}
For \(Q\subseteq Y\), let \(\mathcal C_Q\) be the fibers of the projection
\(D\to D|_Q\).  The active lower slice is empty, so every such fiber is a
relative projection block.  Label a block \(E\in\mathcal C_Q\) by
\[
 \lambda(E)=\{z\in Z_0:E\cap S_z\ne\varnothing\},
\]
and, for \(R\subseteq Z_0\), put
\[
 n_Q(R)=
 \bigl|\{E\in\mathcal C_Q:\lambda(E)=R\}\bigr|,
\qquad
 h_{P,R}=\sum_{Q\subseteq P}(-1)^{|P\setminus Q|}n_Q(R).
\]
The rank-selected supporting-cube identity for relative Yang modules
\cite{BousquetYangComplex26} expresses every \(h_{P,R}\) as a sum of fine
numerator coefficients of the full relative pair \((A,K)\).  Both endpoints
are ample, so that relative module is Cohen--Macaulay and all these
coefficients are nonnegative.  Hence
\begin{equation}
 h_{P,R}\ge0
 \qquad(P\subseteq Y,\ R\subseteq Z_0).
 \label{eq:damp-passive-refined-h-nonnegative}
\end{equation}

At \(Q=\varnothing\), the unique projection block is \(D\).  It meets every
\(S_z\), by the definition of \(Z_0\), and therefore
\[
 n_\varnothing(Z_0)=h_{\varnothing,Z_0}=1.
\]
Boolean M\"obius inversion and
\eqref{eq:damp-passive-refined-h-nonnegative} now give
\[
 n_Y(Z_0)=\sum_{P\subseteq Y}h_{P,Z_0}\ge1.
\]
Since all members of \(D\) agree on \(Z\), projection to \(Y\) is
injective.  Thus the blocks in \(\mathcal C_Y\) are singletons, and a
singleton has label \(Z_0\) exactly when its concept belongs to every
\(S_z\).  This proves~\eqref{eq:damp-empty-slice-passive-helly}.
\end{proof}

We next localize the same argument to one minimal active carrier.  Continue
with the notation \(L\subseteq B=L\mathbin{\dot\cup}D\subseteq Q_Y\) of
Proposition~\ref{prop:damp-two-connected-passive-root-cover}, but only
assume that \(D\) is connected.  Put
\[
 \mathcal H_Y=\operatorname{Sh}(B)\setminus\operatorname{Sh}(L).
\]
If \(P\) is inclusion-minimal in \(\mathcal H_Y\), then \(L\) misses one
unique trace \(u_P\in Q_P\), and its \emph{minimal active carrier} is
\[
 E_P=\{d\in D:d|_P=u_P\}.
\]

\begin{proposition}[Common passive owner of a minimal active carrier]
\label{prop:damp-minimal-active-carrier-common-owner}
Let \(K\subsetneq A\) be ample, let \(D=A\setminus K\) be connected, and
use the preceding active-slice notation.  For \(z\in Z\), put
\[
 S_z=\{d\in D:d^{\oplus z}\in K\},
\qquad
 Z_P=\{z\in Z:E_P\cap S_z\ne\varnothing\}.
\]
Then, whenever \(Z_P\ne\varnothing\),
\begin{equation}
 E_P\cap\bigcap_{z\in Z_P}S_z\ne\varnothing.
 \label{eq:damp-minimal-carrier-common-passive-owner}
\end{equation}
Thus all passive boundary types incident with one minimal active carrier
have a common owner inside that carrier.
\end{proposition}

\begin{proof}
Condition both \(A\) and \(K\) on the active trace \(u_P\).  The relative
concepts of the conditioned pair are exactly \(E_P\).  They agree on the
passive assignment \(\sigma\).  Moreover, its lower slice at \(\sigma\) is
empty, because \(u_P\) is not realized by \(L\).  Conditioning preserves
ampleness, so Lemma~\ref{lem:damp-empty-active-slice-passive-helly} applies.
Its nonempty passive boundary types are precisely
\[
 E_P\cap S_z\qquad(z\in Z_P).
\]
Their common intersection is~\eqref{eq:damp-minimal-carrier-common-passive-owner}.
\end{proof}

\begin{corollary}[A blocked active root transfers to carrier depth two]
\label{cor:damp-blocked-active-root-depth-transfer}
Under the hypotheses of
Proposition~\ref{prop:damp-two-connected-passive-root-cover}, let
\(d\in R\).  Then the unique minimal active carrier \(E_P\) containing
\(d\) contains a concept \(o\ne d\) which belongs to at least two minimal
active carriers.
\end{corollary}

\begin{proof}
The one-concept colon criterion identifies \(P=N_L(d)\) as the unique
minimal relative support whose missing trace is realized by \(d\).  The
blocking conclusion
\eqref{eq:damp-passive-boundary-blocks-active-root} supplies
\(z\in Z_P\) such that \(d\notin S_z\).  Choose a common owner
\[
 o\in E_P\cap\bigcap_{w\in Z_P}S_w
\]
from Proposition~\ref{prop:damp-minimal-active-carrier-common-owner}.
Then \(o\ne d\).

Suppose that \(o\) belonged to only one minimal active carrier.  The local
carrier form of the one-concept criterion would give
\(L\cup\{o\}\) ample, with unique carrier \(E_P\).  Apply
Proposition~\ref{prop:damp-two-connected-passive-root-cover} to \(o\).
It gives a passive blocker \(S_w\) meeting \(E_P\) but not containing
\(o\).  Hence \(w\in Z_P\), contrary to the choice of the common owner.
Thus \(o\) has minimal-carrier depth at least two.
\end{proof}

The preceding owner statement records only which relative concepts have a
neighbour across a passive coordinate.  The whole opposite section is
available, however, and ampleness makes it considerably more rigid.  We
record that extra structure before continuing the cardinality argument.
For \(u\in Q_Y\), let \(\widehat u\in Q_X\) denote its lift with passive
restriction \(\sigma\).

\begin{lemma}[A passive coordinate gives two ample subintervals]
\label{lem:damp-passive-coordinate-two-ample-subintervals}
Use the notation of
Proposition~\ref{prop:damp-two-connected-passive-root-cover}.  For
\(z\in Z\), put
\begin{equation}
 M_z=\{u\in Q_Y:\widehat u^{\oplus z}\in K\}.
 \label{eq:damp-passive-opposite-collar}
\end{equation}
Then \(S_z=D\cap M_z\).  Moreover, whenever the displayed families are
nonempty, each of
\begin{equation}
 L\cap M_z,\quad B\cap M_z,\quad
 L\cup M_z,\quad B\cup M_z
 \label{eq:damp-passive-collar-four-ample-families}
\end{equation}
is ample, and
\begin{equation}
 \begin{split}
  (B\cap M_z)\setminus(L\cap M_z)&=S_z,\\
  (B\cup M_z)\setminus(L\cup M_z)&=D\setminus S_z.
 \end{split}
 \label{eq:damp-passive-collar-two-relative-shores}
\end{equation}
Thus a passive boundary type and its complement are not arbitrary subsets
of \(D\): they are the relative facets of two nested ample pairs on the
same five active coordinates.
\end{lemma}

\begin{proof}
Condition \(K\) on the assignment \(\sigma\) on \(Z\setminus\{z\}\), and
then delete those fixed coordinates.  The two \(z\)-sections of the
resulting ample class are exactly \(L\) and \(M_z\).  Apply the
two-layer shadow-intersection lemma,
Lemma~\ref{lem:damp-ample-two-layer-shadow-intersection}: its two sections,
their union, and their intersection are ample whenever nonempty.  This
gives the two families involving \(L\) in
\eqref{eq:damp-passive-collar-four-ample-families}.

The same conditioning of \(A\) has sections \(B\) and \(M_z\).  Indeed,
every member of \(A\setminus K=D\) has passive restriction \(\sigma\), so
the opposite section of \(A\) is the same as that of \(K\).  A second
application of the two-layer lemma gives the two families involving \(B\).
Finally, \(B=L\mathbin{\dot\cup}D\), and hence elementary set algebra gives
\eqref{eq:damp-passive-collar-two-relative-shores}.  The identity
\(S_z=D\cap M_z\) follows directly from the definitions.
\end{proof}

\begin{lemma}[Passive neighbours force a punctured cube]
\label{lem:damp-passive-neighbours-force-punctured-cube}
Let \(K\subsetneq A\subseteq Q_{Y\mathbin{\dot\cup}Z}\) be ample, put
\(D=A\setminus K\), and suppose that all members of \(D\) have the same
restriction \(\sigma\in Q_Z\).  Put
\[
 L=\{k|_Y:k\in K,\ k|_Z=\sigma\},
 \qquad
 S_z=\{d\in D:d^{\oplus z}\in K\}\quad(z\in Z).
\]
For \(d\in D\), define
\[
 Z(d)=\{z\in Z:d\in S_z\},
 \qquad t(d)=|Z(d)|.
\]
Then
\begin{equation}
 d^{\oplus T}\in K
 \qquad
 \text{for every }d\in D\text{ and }
 \varnothing\ne T\subseteq Z(d).
 \label{eq:damp-passive-punctured-fiber}
\end{equation}
Consequently,
\begin{equation}
 |K|\ge |L|+\sum_{d\in D}\bigl(2^{t(d)}-1\bigr).
 \label{eq:damp-passive-punctured-fiber-count}
\end{equation}
If \(L\ne\varnothing\), put
\[
 Z_0=\{z\in Z:S_z\ne\varnothing\}.
\]
For every \(z\in Z_0\), there is an \(a_z\in L\) such that the lift of
\(a_z\) and its \(z\)-flip, \(\widehat a_z^{\oplus z}\), both belong to
\(K\).  These collar vertices
sharpen the count to
\begin{equation}
 |K|\ge |L|+|Z_0|
       +\sum_{d\in D}\bigl(2^{t(d)}-1\bigr).
 \label{eq:damp-passive-punctured-fiber-collar-count}
\end{equation}
Thus \(t\) passive neighbours of one relative concept cost the whole
punctured \(t\)-cube, not merely \(t\) boundary vertices, and every
nonempty passive direction has an additional collar cost.
\end{lemma}

\begin{proof}
Fix \(d\in D\) with \(Z(d)\ne\varnothing\).  Condition \(K\) on the
active word \(d|_Y\) and on the passive assignment \(\sigma\) outside
\(Z(d)\), and delete the fixed coordinates.  The resulting family
\(F_d\subseteq Q_{Z(d)}\) is nonempty and ample.  Under the identification
of a word with the set of coordinates on which it differs from
\(\sigma|_{Z(d)}\), one has
\[
 \varnothing\notin F_d,
 \qquad
 \{z\}\in F_d\quad(z\in Z(d)).
\]
Indeed, the first assertion is \(d\notin K\), and the second is the
definition of \(Z(d)\).

By complement duality, \(Q_{Z(d)}\setminus F_d\) is ample.  It contains
the empty word.  If it contained a second word \(T\ne\varnothing\), its
one-inclusion graph would contain a geodesic from \(\varnothing\) to
\(T\), because every ample family is isometric.  The first edge of this
geodesic would end at a singleton \(\{z\}\), although every singleton
belongs to \(F_d\), a contradiction.  Hence
\[
 Q_{Z(d)}\setminus F_d=\{\varnothing\},
\]
which proves~\eqref{eq:damp-passive-punctured-fiber}.  There is nothing
to prove when \(Z(d)=\varnothing\).

The words in~\eqref{eq:damp-passive-punctured-fiber} belonging to two
distinct concepts \(d,d'\in D\) have distinct restrictions to \(Y\),
because \(d,d'\) have the same restriction to \(Z\).  They are therefore
distinct.  They are also disjoint from the lift of \(L\), since their
passive restriction differs from \(\sigma\).  For fixed \(d\), there are
\(2^{t(d)}-1\) of them.  Summing these disjoint subsets of \(K\) proves
\eqref{eq:damp-passive-punctured-fiber-count}.

Suppose that \(L\ne\varnothing\), and fix \(z\in Z_0\).  Condition \(K\)
on \(\sigma\) outside \(z\), and delete the fixed passive coordinates.
Its two \(z\)-sections are \(L\) and \(M_z\), in the notation of
Lemma~\ref{lem:damp-passive-coordinate-two-ample-subintervals}.  Both
are nonempty: the first by hypothesis, and the second because
\(S_z\ne\varnothing\).  The conditioned ample family is connected.
Since active edges stay in one \(z\)-section, it must contain a \(z\)-edge.
Equivalently, \(L\cap M_z\ne\varnothing\); choose \(a_z\) in this
intersection.

The opposite lift \(\widehat a_z^{\oplus z}\in K\) has active restriction in
\(L\), whereas every word counted in
\eqref{eq:damp-passive-punctured-fiber-count} has active restriction in
\(D\).  It is therefore a new vertex.  The chosen vertices for distinct
coordinates \(z\) have different passive restrictions, so they are
pairwise distinct.  Adding these \(|Z_0|\) vertices proves
\eqref{eq:damp-passive-punctured-fiber-collar-count}.
\end{proof}

\begin{lemma}[Every nonempty passive shore contains an active root]
\label{lem:damp-passive-shore-contains-active-root}
Under the hypotheses of
Proposition~\ref{prop:damp-two-connected-passive-root-cover}, suppose that
\eqref{eq:damp-terminal-no-one-extension} holds.  Then \(L\ne\varnothing\),
and every nonempty \(S_z\) contains a member of
\[
 R=\{d\in D:L\cup\{d\}\text{ is ample}\}.
\]
\end{lemma}

\begin{proof}
Suppose first that \(L=\varnothing\).  Then every member of \(D\) belongs
to \(R\).  Proposition~\ref{prop:damp-two-connected-passive-root-cover}
gives, for each \(d\in D\), a nonempty passive type which excludes \(d\).
In particular the set \(Z_0\) of nonempty passive types is nonempty.
Lemma~\ref{lem:damp-empty-active-slice-passive-helly} supplies
\[
 o\in\bigcap_{z\in Z_0}S_z.
\]
Applying the blocking conclusion of
Proposition~\ref{prop:damp-two-connected-passive-root-cover} to \(o\in R\)
produces \(z\in Z_0\) with \(o\notin S_z\), a contradiction.  Thus
\(L\ne\varnothing\).

Fix \(z\) with \(S_z\ne\varnothing\), and abbreviate \(M=M_z\).  The
conditioned ample class with sections \(L,M\) has both sections nonempty.
Its one-inclusion graph is connected, so it contains a \(z\)-edge.  Hence
\begin{equation}
 I:=L\cap M\ne\varnothing.
 \label{eq:damp-passive-collar-nonempty-overlap}
\end{equation}
By Lemma~\ref{lem:damp-passive-coordinate-two-ample-subintervals},
\[
 I\subsetneq (B\cap M)=I\mathbin{\dot\cup}S_z
\]
is a nested pair of ample families on at most five coordinates.  Apply
Theorem~\ref{thm:damp-five-coordinate-relative-peeling} to obtain
\(p\in S_z\) for which \(I\cup\{p\}\) is ample.  The conditioned class
with sections \(L,M\) is an ample expansion, \(p\in M\setminus L\), and
its section intersection is \(I\).  Lemma
\ref{lem:damp-expansion-separator-augmentation} therefore shows that
\(L\cup\{p\}\) is ample.  Thus \(p\in R\cap S_z\), as required.
\end{proof}

The preceding proof used only the first concept of a relative peeling of the
shore.  Keeping the whole peeling turns the shore into an actual ample
intermediate family.  This is the useful functorial form of the collar
construction.

\begin{lemma}[Passive-shore absorption and blocker-rank descent]
\label{lem:damp-passive-shore-absorption}
Under the hypotheses of
Lemma~\ref{lem:damp-passive-shore-contains-active-root}, fix \(z\in Z\)
with \(S_z\ne\varnothing\), and put
\begin{equation}
 G_z=L\cup S_z.
 \label{eq:damp-passive-shore-gate}
\end{equation}
Then \(G_z\) is ample.  More precisely, there is a saturated chain of ample
families from \(L\) to \(B\) in which every member of \(S_z\) is added
before every member of \(D\setminus S_z\).  Its first segment is
\begin{equation}
 L=G_0\subset G_1\subset\cdots\subset G_{|S_z|}=G_z,
 \qquad |G_i\setminus G_{i-1}|=1,
 \label{eq:damp-passive-shore-absorption-chain}
\end{equation}
whose added concepts are exactly the members of \(S_z\).

Suppose in addition that \(d\in R\setminus S_z\) has the peripheral
blocker cell \(C\), support \(W\), and duplicated endpoint
\(c=d^{\oplus W}\) supplied by
Lemma~\ref{lem:damp-active-root-peripheral-blocker-cell} for this same
coordinate \(z\).  Then
\begin{equation}
 C\cap G_z=\{c\},
 \qquad
 (G_z\cup\{d\})\cap Q_W(d)=\{d,c\}.
 \label{eq:damp-passive-gate-blocker-descent}
\end{equation}
Consequently, if \(|W|\ge2\), the extension \(G_z\cup\{d\}\) is
nonample, witnessed by the antipodal pair on the strictly smaller active
face \(Q_W(d)\).  If \(|W|=1\), the named transverse blocker on
\(W\cup\{z\}\) has disappeared: its trace after absorption is the full
one-cube.  This last assertion does not exclude a different blocker for
\(G_z\cup\{d\}\).
\end{lemma}

\begin{proof}
Abbreviate \(M=M_z\) and \(I=L\cap M\).  The proof of
Lemma~\ref{lem:damp-passive-shore-contains-active-root} gives
\(I\ne\varnothing\), and
Lemma~\ref{lem:damp-passive-coordinate-two-ample-subintervals} gives the
nested ample pair
\[
 I\subsetneq I\mathbin{\dot\cup}S_z=B\cap M.
\]
The five-coordinate relative-peeling theorem supplies an ordering
\(p_1,\ldots,p_{|S_z|}\) of \(S_z\) for which every
\[
 I_i=I\cup\{p_1,\ldots,p_i\}
\]
is ample.  Put \(G_i=L\cup\{p_1,\ldots,p_i\}\).  We prove inductively that
the two-layer family with sections \(G_i\) and \(M\) is ample.  The case
\(i=0\) is the conditioning of \(K\) used in the collar lemma.  Moreover,
\(G_{i-1}\cap M=I_{i-1}\), while \(I_{i-1}\cup\{p_i\}=I_i\) is ample.
Lemma~\ref{lem:damp-expansion-separator-duplication} therefore duplicates
\(p_i\) into the \(G_{i-1}\)-section and proves the induction step.  Each
\(G_i\) is a section of this ample two-layer family and hence is ample.
This proves~\eqref{eq:damp-passive-shore-absorption-chain}.
If \(G_z\ne B\), apply the five-coordinate relative-peeling theorem once
more, now to \(G_z\subsetneq B\), and append the resulting saturated
chain.  Its added concepts are exactly
\(B\setminus G_z=D\setminus S_z\).  If \(G_z=B\), the second segment is
empty.  This proves the asserted two-block ordering in both cases.

Now use the peripheral blocker cell.  Since \(C\subseteq D\), it is
disjoint from \(L\), while
\(C\cap S_z=\{c\}\).  Hence \(C\cap G_z=\{c\}\).  Also
\(C=A\cap Q_W(d)\) and \(G_z\subseteq B\subseteq A\), so
\(G_z\cap Q_W(d)=\{c\}\), which proves
\eqref{eq:damp-passive-gate-blocker-descent}.  When \(|W|\ge2\), apply
Proposition~\ref{prop:damp-failed-augmentation-antipodal-localization} to
the ample family \(G_z\).  When \(|W|=1\), the two displayed concepts are
the two vertices of the one-cube, proving exactly the stated rank-one
conclusion.
\end{proof}

Complement duality gives the simultaneous upper-shore statement.  Put
\[
 R^+=\{d\in D:B\setminus\{d\}\text{ is ample}\}.
\]
These are the active deletion roots of the conditioned upper slice.

\begin{corollary}[The two directed root cycles forced by a terminal fibre]
\label{cor:damp-terminal-fibre-two-root-cycles}
In addition to the hypotheses of
Proposition~\ref{prop:damp-two-connected-passive-root-cover}, suppose that
\begin{equation}
 A\setminus\{d\}\quad\hbox{is nonample for every }d\in D.
 \label{eq:damp-terminal-no-one-deletion}
\end{equation}
Then the following hold.

\begin{enumerate}
\item Every nonempty \(S_z\) contains a member of \(R\), and every
  nonempty \(D\setminus S_z\) contains a member of \(R^+\).  If both
  shores are nonempty, one saturated ample chain from \(L\) to \(B\)
  starts with a member of \(R\cap S_z\), adds all of \(S_z\) before its
  complement, and ends with a member of
  \(R^+\cap(D\setminus S_z)\).
\item One has
  \[
   |R|\ge2,\qquad |R^+|\ge2.
  \]
\item There is a directed graph on \(R\) with positive out-degree at
  every vertex and no loop.  An edge \(d\to e\) is witnessed by a
  passive coordinate \(z\) such that
  \begin{equation}
   d\notin S_z,\qquad e\in S_z,
   \label{eq:damp-lower-root-cycle-separation}
  \end{equation}
  by an isolated peripheral blocker cell
  \(C_d\) in the unique minimal active carrier of \(d\), and by a common
  passive owner \(o\in S_z\) which is not in \(R\).  The shore also gives
  the ample gate \(G_z=L\cup S_z\); absorbing it either lowers the named
  blocker from \(W\cup\{z\}\) to \(W\), or eliminates that blocker when
  \(|W|=1\).  Consequently this graph contains a directed cycle of length
  at least two.
\item Complementarily, there is such a directed cycle on \(R^+\), with
  every edge separated by an opposite shore \(D\setminus S_z\) and
  carrying a common owner which is not in \(R^+\).
\end{enumerate}
\end{corollary}

\begin{proof}
The assertion for \(S_z\) is
Lemma~\ref{lem:damp-passive-shore-contains-active-root}.  Apply the same
lemma to the complementary nested ample pair
\[
 Q_X\setminus A\subsetneq Q_X\setminus K.
\]
Its relative concepts are again \(D\), its active conditioned pair is
\[
 Q_Y\setminus B\subsetneq Q_Y\setminus L,
\]
and its passive boundary type in coordinate \(z\) is \(D\setminus S_z\).
Indeed, \(d^{\oplus z}\) belongs to \(A\) exactly when it belongs to
\(K\), because every member of \(A\setminus K\) has passive restriction
\(\sigma\).  Condition~\eqref{eq:damp-terminal-no-one-deletion} is
precisely the no-one-extension condition for this complementary pair.  Its
active roots are the members of \(R^+\).  This proves item~1.

For the strengthened simultaneous statement, use the two-block chain from
Lemma~\ref{lem:damp-passive-shore-absorption}.  Its first added concept is
an active lower root and belongs to \(S_z\); its last added concept is an
active upper root and belongs to \(D\setminus S_z\).

Fix \(d\in R\).  Lemma
\ref{lem:damp-active-root-peripheral-blocker-cell} supplies a nonempty
passive type \(S_z\) which excludes \(d\), together with the isolated
peripheral blocker cell in the unique minimal carrier of \(d\).  By
item~1, \(S_z\) contains some \(e\in R\), necessarily different from
\(d\); direct the edge from \(d\) to one such \(e\).  Proposition
\ref{prop:damp-minimal-active-carrier-common-owner} supplies a common owner
\(o\) of all passive types incident with the unique minimal carrier of
\(d\).  In particular \(o\in S_z\), whereas \(d\notin S_z\).  Corollary
\ref{cor:damp-blocked-active-root-depth-transfer} shows that \(o\) has
minimal-carrier depth at least two, so \(o\notin R\).  Every vertex of the
finite directed graph has an outgoing edge and there are no loops.  It
therefore contains a directed cycle of length at least two, and in
particular \(R\) has at least two members.  This proves items~2 and~3 on
the lower shore.

Apply the same argument to the complementary pair.  Minimal active
carriers there are the upper carriers of the original active interval,
and carrier depth one is exactly membership in \(R^+\).  This gives the
upper directed cycle, its non-root owners, and the remaining inequality in
item~2.
\end{proof}

The two root cycles can be coupled without choosing a common set of root
vertices.  The correct common objects are the ample gates.

\begin{corollary}[Two-sided holonomy on the passive gates]
\label{cor:damp-two-sided-passive-gate-holonomy}
Under the hypotheses of
Corollary~\ref{cor:damp-terminal-fibre-two-root-cycles}, let
\[
 \mathscr S=\{S_z:\varnothing\ne S_z\subsetneq D,\ z\in Z\},
\]
where repeated shore sets are identified.  Then \(\mathscr S\ne\varnothing\).
For every \(S\in\mathscr S\), one may choose a two-block relative peeling
of \(B\) over \(L\), with \(S\) first, whose first and last relative
concepts \(r_S,u_S\) satisfy
\[
 r_S\in R\cap S,\qquad
 u_S\in R^+\cap(D\setminus S).
\]
There are maps
\[
 \lambda,\mu:\mathscr S\longrightarrow\mathscr S
\]
such that
\begin{equation}
 r_S\in S\setminus\lambda(S),
 \qquad
 u_S\in\mu(S)\setminus S.
 \label{eq:damp-two-sided-gate-holonomy}
\end{equation}
In particular, neither map has a fixed point, both maps contain a directed
cycle, and every arrow is certified by a nonempty asymmetric difference
of its two shores.

If \(|\mathscr S|=2\), say \(\mathscr S=\{S,T\}\), then necessarily
\[
 r_S,u_T\in S\setminus T,
 \qquad
 r_T,u_S\in T\setminus S.
\]
Thus the two-shore case is one crossed lower--upper profile, rather than
two unrelated root cycles.
\end{corollary}

\begin{proof}
The active root set \(R\) is nonempty.  Applying
Lemma~\ref{lem:damp-active-root-peripheral-blocker-cell} to any of its
members produces a nonempty shore which excludes that member, and hence is
proper.  Thus \(\mathscr S\ne\varnothing\).

For \(S=S_z\in\mathscr S\), take the two-block chain from
Lemma~\ref{lem:damp-passive-shore-absorption}.  Its first relative concept
is an active lower root in \(S\), and its last relative concept is an
active upper root in \(D\setminus S\); call them \(r_S,u_S\).

Apply Lemma~\ref{lem:damp-active-root-peripheral-blocker-cell} to \(r_S\).
Its blocker shore is a nonempty \(S_w\) which excludes \(r_S\).  It is
proper because it excludes a member of \(D\), so its underlying shore set
belongs to \(\mathscr S\).  Define \(\lambda(S)=S_w\).  This gives the
first relation in~\eqref{eq:damp-two-sided-gate-holonomy}.

Apply the complementary version of the same lemma to \(u_S\in R^+\).
It supplies a nonempty proper opposite shore \(D\setminus S_v\) which
excludes \(u_S\).  Equivalently \(u_S\in S_v\), and \(S_v\) is itself
nonempty and proper.  Define \(\mu(S)=S_v\).  This gives the second
relation.  The displayed asymmetric differences exclude fixed points.
Every self-map without a fixed point on a nonempty finite set contains a
directed cycle of length at least two.

If \(\mathscr S=\{S,T\}\), both fixed-point-free maps interchange the two
members.  Substituting this into
\eqref{eq:damp-two-sided-gate-holonomy} gives the final display.
\end{proof}

The holonomy now has a quantitative cost.  The following statement also
identifies the exact zero-excess geometry, which will be the rigid branch
of the remaining comparison theorem.

\begin{corollary}[Exponential cost and equality form of gate holonomy]
\label{cor:damp-gate-holonomy-exponential-cost}
Assume the hypotheses of
Corollary~\ref{cor:damp-two-sided-passive-gate-holonomy}, and let
\[
 S_1,S_2,\ldots,S_k
\]
be a directed cycle of \(\lambda\), indexed so that
\(\lambda(S_i)=S_{i+1}\) modulo \(k\).  Thus \(k\ge2\).  For
\(d\in D\), let \(t(d)=|\{z\in Z:d\in S_z\}|\), and put
\[
 Z_0=\{z\in Z:S_z\ne\varnothing\},
\qquad
 t_{\mathcal C}(d)=|\{i:d\in S_i\}|.
\]
Define the uncounted lower-fiber remainder
\begin{equation}
 \eta=
 |K|-|L|-|Z_0|-\sum_{d\in D}\bigl(2^{t(d)}-1\bigr).
 \label{eq:damp-holonomy-uncounted-remainder}
\end{equation}
Then \(\eta\ge0\), and
\begin{align}
 |K|-|L|
 &\ge |Z_0|+\sum_{d\in D}\bigl(2^{t(d)}-1\bigr)\ge3k,
 \label{eq:damp-holonomy-exponential-cost}\\
 |K|-|L|-3k
 &=
 \eta+(|Z_0|-k)
   +\sum_{d\in D}\bigl(2^{t(d)}-1-t_{\mathcal C}(d)\bigr)
   +\sum_{i=1}^k\bigl(|S_i|-2\bigr).
 \label{eq:damp-holonomy-excess-decomposition}
\end{align}
Every summand on the right of
\eqref{eq:damp-holonomy-excess-decomposition} is nonnegative.

If the total bound in
\eqref{eq:damp-holonomy-exponential-cost} is sharp, namely
\begin{equation}
 |K|=|L|+3k,
 \label{eq:damp-holonomy-zero-excess}
\end{equation}
then the following conclusions hold.
\begin{enumerate}
\item The cycle shores are pairwise disjoint and have order two.  Every
  nonempty passive shore is one of them, and each cycle shore is realized
  by a unique passive coordinate.
\item Each \(S_i\) consists of one lower root \(r_i\in R\) and one
  common owner \(o_i\notin R\), and \(r_io_i\) is an active edge.
\item If \(z_i\) is the unique passive coordinate realizing \(S_i\), then
  \(L\cap M_{z_i}=\{a_i\}\) and
  \[
   M_{z_i}=\{a_i,r_i,o_i\}.
  \]
  This is a three-vertex active path whose terminal edge is \(r_io_i\).
  Consequently \(K\) contains the attached three-edge path
  \[
   \widehat a_i,
   \widehat a_i^{\oplus z_i},\
   r_i^{\oplus z_i},\
   o_i^{\oplus z_i},
  \]
  in one of the two possible orders after
  \(\widehat a_i^{\oplus z_i}\).
\item The upper family \(A\) contains the full square on the
  active--passive direction pair of \(r_io_i\) and \(z_i\):
  \[
   r_i,o_i\in D,
   \qquad r_i^{\oplus z_i},o_i^{\oplus z_i}\in K.
  \]
  These \(k\) squares are pairwise vertex-disjoint.
\item Apart from the lift of \(L\), the lower family \(K\) consists
  exactly of the \(3k\) nonbase vertices in the attached paths.
\end{enumerate}
Thus every departure from a system of pendant three-edge square-gate arms
has a positive, explicitly charged cost in \(K\setminus L\).
\end{corollary}

\begin{proof}
For every cycle shore \(S_i\), the shore-root lemma gives a lower root in
\(S_i\).  Moreover, \(S_i\) blocks the lower root of its predecessor
\(S_{i-1}\).  The common-owner theorem and
Corollary~\ref{cor:damp-blocked-active-root-depth-transfer} therefore give
a member of \(S_i\setminus R\).  Hence
\begin{equation}
 |S_i|\ge2,
 \qquad
 \sum_{d\in D}t_{\mathcal C}(d)=\sum_{i=1}^k|S_i|\ge2k.
 \label{eq:damp-holonomy-cycle-incidence-count}
\end{equation}
Moreover \(t(d)\ge t_{\mathcal C}(d)\), and
\(2^{t(d)}-1\ge t(d)\).  The cycle contains \(k\) distinct shore sets,
so \(Z_0\) contains at least \(k\) coordinates.  Lemma
\ref{lem:damp-passive-neighbours-force-punctured-cube} and
\eqref{eq:damp-holonomy-cycle-incidence-count} prove
\eqref{eq:damp-holonomy-exponential-cost}, including
\(\eta\ge0\).  Subtracting \(3k\) and inserting the middle incidence sum
in the definition of \(\eta\) gives
\begin{align*}
 |K|-|L|-3k
 ={}&\eta+(|Z_0|-k)+\sum_{d\in D}
      \bigl(2^{t(d)}-1-t_{\mathcal C}(d)\bigr)\\
 &+\sum_{i=1}^k(|S_i|-2),
\end{align*}
which proves~\eqref{eq:damp-holonomy-excess-decomposition}.  Its
summands are nonnegative by the two preceding inequalities.

Suppose now that~\eqref{eq:damp-holonomy-zero-excess} holds.  Every
inequality just used is then an equality.  Thus \(|S_i|=2\) for every
\(i\), and \(t(d)\le1\) for every \(d\in D\), since
\(2^t-1=t\) only for \(t\in\{0,1\}\).  Equality also gives
\(t(d)=t_{\mathcal C}(d)\).  It follows that the cycle shores are
pairwise disjoint, that no other passive coordinate has a nonempty
shore, and that no cycle shore is represented by a second passive
coordinate.  This proves item~1.

Choose a lower root \(r_i\in R\cap S_i\).  The common owner supplied
from the predecessor belongs to \(S_i\setminus R\); call it \(o_i\).
Since \(|S_i|=2\), these are its two members.  The ample gate
\(G_{z_i}=L\cup\{r_i,o_i\}\) admits the two-block chain of
Lemma~\ref{lem:damp-passive-shore-absorption}.  Its first addition is
\(r_i\), because \(o_i\notin R\).

Suppose that \(r_i\) and \(o_i\) were not adjacent.  Since
\(L\cup\{o_i\}\) is nonample, the one-concept colon criterion gives
\(y\in L\) which blocks \(o_i\).  Adding the non-neighbour \(r_i\) does
not add a neighbour of \(o_i\), so the same \(y\) blocks \(o_i\) over the
ample family \(L\cup\{r_i\}\).  The criterion would make
\(L\cup\{r_i,o_i\}=G_{z_i}\) nonample, a contradiction.  Hence
\(r_io_i\) is an active edge, proving item~2.

Equality in the collar count gives
\(|L\cap M_{z_i}|=1\); write this intersection as \(\{a_i\}\).
It also leaves no further member of \(M_{z_i}\), since such a member
would give another vertex of \(K\) in the \(z_i\)-layer.  Hence
\[
 M_{z_i}=\{a_i,r_i,o_i\}.
\]
This family is ample, and hence connected.  It contains the edge
\(r_io_i\), so its third vertex \(a_i\) is adjacent to exactly one of
\(r_i,o_i\).  Together with the \(z_i\)-edge from \(\widehat a_i\) to
\(\widehat a_i^{\oplus z_i}\), this is the attached three-edge path in
item~3.

The definition of the shore gives
\(r_i^{\oplus z_i},o_i^{\oplus z_i}\in K\), while
\(r_i,o_i\in D\subseteq A\).  These are the four vertices of the claimed
square.  Different shores have disjoint active endpoints.  Their flipped
vertices also have distinct active restrictions, and no flipped vertex has
passive restriction \(\sigma\).  The squares are therefore
vertex-disjoint.  This proves item~4.  Finally, equality in the counting
argument of Lemma~\ref{lem:damp-passive-neighbours-force-punctured-cube}
leaves no member of \(K\) beyond the lift of \(L\) and these \(3k\)
nonbase path vertices, proving item~5.
\end{proof}

\begin{lemma}[A shore outside the holonomy cycle pays its order plus one]
\label{lem:damp-off-cycle-shore-cost}
Retain the notation of
Corollary~\ref{cor:damp-gate-holonomy-exponential-cost}, and put
\[
 \varepsilon=|K|-|L|-3k.
\]
For every nonempty passive shore \(S_z\), one has
\begin{equation}
 \begin{cases}
  |S_z|\le \varepsilon+2,
    &\text{if }S_z\in\{S_1,\ldots,S_k\},\\
  |S_z|\le \varepsilon-1,
    &\text{if }S_z\notin\{S_1,\ldots,S_k\}.
 \end{cases}
 \label{eq:damp-on-off-cycle-shore-cost}
\end{equation}
Equivalently, a cycle shore costs at least \(|S_z|-2\) units of excess,
whereas a shore outside the cycle costs at least \(|S_z|+1\).
\end{lemma}

\begin{proof}
If \(S_z=S_i\) is a cycle shore, its summand
\(|S_i|-2\) in
\eqref{eq:damp-holonomy-excess-decomposition} gives the first bound.

Suppose that \(S_z\) is not a cycle shore.  The coordinate \(z\) is in
\(Z_0\) but is not one of the \(k\) coordinates needed to realize the
distinct cycle shores.  Hence its contribution gives
\(|Z_0|-k\ge1\).  For every \(d\in S_z\), the noncycle incidence in
coordinate \(z\) gives \(t(d)\ge t_{\mathcal C}(d)+1\), and therefore
\[
 2^{t(d)}-1-t_{\mathcal C}(d)\ge1.
\]
The corresponding terms for the \(|S_z|\) distinct concepts, together
with the extra coordinate term, contribute at least \(|S_z|+1\) to the
right-hand side of
\eqref{eq:damp-holonomy-excess-decomposition}.  This proves the second
bound.
\end{proof}

\begin{corollary}[A lower family of order at most nine has two-arm
holonomy]
\label{cor:damp-small-lower-family-two-arm-holonomy}
Under the hypotheses of
Corollary~\ref{cor:damp-gate-holonomy-exponential-cost}, suppose that
\(|K|\le9\).  Every directed cycle of \(\lambda\) has length two.  For
such a cycle \(S_1,S_2\), put
\[
 \varepsilon=|K|-|L|-6.
\]
Then \(0\le\varepsilon\le2\), and
\begin{equation}
 \eta+(|Z_0|-2)
 +\sum_{d\in D}\bigl(2^{t(d)}-1-t_{\mathcal C}(d)\bigr)
 +\sum_{i=1}^2(|S_i|-2)
 =\varepsilon.
 \label{eq:damp-small-lower-family-two-arm-budget}
\end{equation}
In particular, \(|Z_0|\le4\) and
\(|\mathscr S|\le4\).  If \(|K|=7\), then \(|L|=1\),
\(\varepsilon=0\), and the complete pendant two-arm normal form of
Corollary~\ref{cor:damp-gate-holonomy-exponential-cost} holds.
\end{corollary}

\begin{proof}
Lemma~\ref{lem:damp-passive-shore-contains-active-root} gives
\(|L|\ge1\).  A cycle of length \(k\) therefore gives
\[
 9\ge|K|\ge|L|+3k\ge1+3k.
\]
Since \(k\ge2\), one has \(k=2\).  The same inequalities give
\(0\le\varepsilon\le2\), and
\eqref{eq:damp-holonomy-excess-decomposition} gives
\eqref{eq:damp-small-lower-family-two-arm-budget}.  Its first summand
gives \(|Z_0|\le4\), and every member of \(\mathscr S\) is realized by
a coordinate in \(Z_0\), so \(|\mathscr S|\le4\).  If \(|K|=7\), then
the displayed chain of inequalities is sharp throughout, and the equality
part of Corollary~\ref{cor:damp-gate-holonomy-exponential-cost} applies.
\end{proof}

\begin{lemma}[A two-concept cycle shore is an active edge]
\label{lem:damp-two-concept-cycle-shore-active-edge}
Under the hypotheses of
Corollary~\ref{cor:damp-gate-holonomy-exponential-cost}, let \(S\) lie
on a directed cycle of \(\lambda\), and suppose that \(|S|=2\).
Then
\[
 S=\{r,o\},
 \qquad r\in R,\quad o\notin R,
\]
and \(ro\) is an active edge.
\end{lemma}

\begin{proof}
The shore-root lemma gives \(r\in R\cap S\).  The predecessor of \(S\)
on the \(\lambda\)-cycle is blocked by \(S\), so the common-owner theorem
and Corollary~\ref{cor:damp-blocked-active-root-depth-transfer} give
\(o\in S\setminus R\).  These are the two members of \(S\).

The gate \(L\cup S\) is ample, as is \(L\cup\{r\}\), whereas
\(L\cup\{o\}\) is nonample.  If \(r,o\) were nonadjacent, a blocking
vertex in \(L\) for the latter one-concept extension would remain a
blocking vertex after adding the non-neighbour \(r\).  The one-concept
colon criterion would make \(L\cup\{r,o\}\) nonample, a contradiction.
\end{proof}

\begin{corollary}[The one-unit holonomy anomalies form a trichotomy]
\label{cor:damp-one-unit-holonomy-anomaly-trichotomy}
Use the notation of
Corollary~\ref{cor:damp-small-lower-family-two-arm-holonomy}, and let
\(S_1,S_2\) be its directed two-cycle.  If
\(\varepsilon=1\), exactly one of the following alternatives holds:
\begin{enumerate}
\item \(\eta=1\); the two shores are disjoint two-concept sets, and the
  two cycle shores are the only nonempty passive shores, each represented
  by one coordinate.
\item \(\eta=0\); the two shores have order two and meet in one concept.
  Their common concept belongs to \(D\setminus R\): it is the common
  non-root owner of both shores.
\item \(\eta=0\); the two shores are disjoint, one has order two, and the
  other has order three.
\end{enumerate}
In particular, an extra nonempty passive coordinate cannot occur at cost
one.

If \(\varepsilon=2\) and \(|Z_0|=3\), then the third passive coordinate
has a singleton shore disjoint from \(S_1\cup S_2\), while the two cycle
shores are disjoint two-concept sets and \(\eta=0\).  Thus a new passive
direction has minimum total cost two.
\end{corollary}

\begin{proof}
For the two-cycle, the exact identity
\eqref{eq:damp-small-lower-family-two-arm-budget} is
\begin{equation}
 \varepsilon=
 \eta+(|Z_0|-2)
 +\sum_{d\in D}\bigl(2^{t(d)}-1-t_{\mathcal C}(d)\bigr)
 +(|S_1|-2)+(|S_2|-2).
 \label{eq:damp-one-unit-anomaly-identity}
\end{equation}
All summands are nonnegative integers.

Suppose first that \(\varepsilon=1\).  If \(|Z_0|>2\), choose an
additional nonempty shore coordinate \(w\), and a concept \(d\in S_w\).
If \(d\notin S_1\cup S_2\), then its summand in the middle sum is at
least one, in addition to \(|Z_0|-2\ge1\).  If \(d\) belongs to a cycle
shore, then \(t(d)\ge t_{\mathcal C}(d)+1\) and its middle summand is at
least two.  Both conclusions contradict \(\varepsilon=1\).  Hence
\(|Z_0|=2\), and each cycle shore is represented by its unique coordinate.

With these two coordinates, one has \(t(d)=t_{\mathcal C}(d)\).  The
middle summand is zero unless \(d\in S_1\cap S_2\), in which case it is
\(2^2-1-2=1\).  The final two summands count the concepts in the cycle
shores beyond their required root and non-root owner.  Equation
\eqref{eq:damp-one-unit-anomaly-identity} therefore gives exactly the
three listed alternatives.

In the overlap alternative, write \(r_i\in R\cap S_i\) for the selected
root.  Since \(\lambda(S_1)=S_2\), relation
\eqref{eq:damp-two-sided-gate-holonomy} gives \(r_1\notin S_2\);
similarly \(r_2\notin S_1\).  Each two-concept shore contains a root and
a non-root owner.  Their common concept is consequently the non-root
member of both shores.

Finally suppose that \(\varepsilon=2\) and \(|Z_0|=3\).  The extra
coordinate already contributes one through \(|Z_0|-2\).  The preceding
argument shows that every concept in its shore contributes at least one
more to the middle sum, and contributes at least two if it also lies in a
cycle shore.  Equality with total two therefore forces the extra shore to
be a singleton disjoint from \(S_1\cup S_2\), and exhausts the entire
budget.  Thus \(\eta=0\), the cycle shores are disjoint and have order
two, and no further anomaly occurs.
\end{proof}

\begin{proposition}[A shared-owner one-unit anomaly is impossible]
\label{prop:damp-shared-owner-one-unit-anomaly-impossible}
Under the hypotheses of
Corollary~\ref{cor:damp-one-unit-holonomy-anomaly-trichotomy}, the second
alternative of that corollary cannot occur.  Thus, when
\(\varepsilon=1\), the only possibilities are \(\eta=1\) or disjoint
cycle shores of orders two and three.
\end{proposition}

\begin{proof}
Suppose that
\[
 S_1=\{r_1,o\},
 \qquad
 S_2=\{r_2,o\},
\]
where \(r_1,r_2\in R\) and \(o\notin R\).  Let \(z_i\) be the unique
passive coordinate realizing \(S_i\).  By
Lemma~\ref{lem:damp-two-concept-cycle-shore-active-edge}, \(r_io\) is an
active edge; write its direction as \(f_i\).  The two directions are
distinct, since \(r_1,r_2,o\) are three distinct cube vertices.

Here \(\eta=0\), so the forced-fiber count accounts for all of \(K\).
For each \(i\), the collar intersection \(L\cap M_{z_i}\) is therefore a
singleton, say \(\{a_i\}\), and
\[
 M_{z_i}=\{a_i,r_i,o\}.
\]
This is an ample three-vertex family containing the edge \(r_io\);
hence it is a path and \(a_i\) is adjacent to exactly one of
\(r_i,o\).  The punctured passive fiber over \(o\) contributes
\[
 o^{\oplus z_1},\quad o^{\oplus z_2},\quad
 o^{\oplus\{z_1,z_2\}}\in K.
\]
Together with the two collars and the two root lifts, these are the seven
members of \(K\) outside the lift of \(L\).  Since \(|K|\le9\), one has
\(|L|\le2\).

We determine the one-inclusion graph of \(K\).  First suppose that
\(a_1=a_2=a\), which includes the case \(|L|=1\).  Translate \(o\) to
the zero active word.  If \(a\) is adjacent to \(o\) in both active
paths, the vertices
\[
 \widehat a,\ \widehat a^{\oplus z_1},\
 o^{\oplus z_1},\ o^{\oplus\{z_1,z_2\}},\
 o^{\oplus z_2},\ \widehat a^{\oplus z_2}
\]
form a six-cycle.  If \(a\) is adjacent to both \(r_1\) and \(r_2\),
then \(a=o^{\oplus\{f_1,f_2\}}\), and inserting the two root lifts in
the preceding display gives an eight-cycle.  The two mixed alternatives
are impossible: two adjacent cube vertices have no third common
neighbour.  Any remaining root lift, and the second vertex of \(L\) when
present, is attached as a tree edge.

It remains that \(|L|=2\) and \(a_1\ne a_2\).  The two vertices of the
ample family \(L\) are adjacent.  Again translate \(o\) to zero.  If
\(a_1\) and \(a_2\) were both adjacent to \(o\), they would be at Hamming
distance two; if both were adjacent to \(r_1,r_2\), respectively, their
active words would both have even weight and could not be adjacent.
After interchanging the indices, the only remaining possibility is
\[
 a_1=o^{\oplus g},
 \qquad
 a_2=o^{\oplus\{f_2,g\}}
\]
for an active direction \(g\notin\{f_1,f_2\}\).  The collar paths, the
edge \(a_1a_2\), and the mixed owner vertex then form an eight-cycle;
the unused root lift is a pendant edge.

Thus in every case the one-inclusion graph of \(K\) is a cycle
\(C_{2m}\), with \(m\in\{3,4\}\), with possibly some trees attached.
There are no further edges: two vertices in different passive layers
already differ in the corresponding passive coordinates, while within one
passive layer the active adjacencies are exactly those of the three-vertex
paths displayed above.  The only remaining possible adjacency is the edge
of the two-vertex family \(L\), which is one of the attached tree edges in
the preceding description.
It has no four-cycle.  If \(t\) is the number of attached tree edges,
then its isometric dimension is \(m+t\): the cycle has \(m\) coordinate
classes and every tree edge is a singleton coordinate class.  Since no
pair of coordinates supports a square, the only strongly shattered sets
are the empty set and these \(m+t\) singleton sets.  Hence
\[
 |\operatorname{SSh}(K)|=1+m+t
 <2m+t=|K|.
\]
This contradicts ampleness of \(K\), which requires equality in the lower
Sandwich bound.  The shared-owner alternative is impossible.
\end{proof}

\begin{corollary}[Two-arm reduction for a small absolute obstruction]
\label{cor:damp-below-forty-two-arm-reduction}
Let \(H\) have minimum cardinality among nonempty cornerless ample
families, suppose that \(|H|\le39\), and use the coordinate-section
notation
\[
 A_b=\rho_e^bH,\qquad B=A_0\cap A_1,\qquad D_b=A_b\setminus B
\]
from Proposition~\ref{prop:damp-two-wing-relative-corner-threshold}.
Then \(|B|\le9\).  If some \(D_b\) has order ten and is two-connected,
the relative pair \(B\subsetneq A_b\) is terminal under both one-concept
addition and deletion, its lower holonomy consists of a two-cycle, and
its excess satisfies
\eqref{eq:damp-small-lower-family-two-arm-budget}.  In particular,
\(|B|\ge7\).  The same conclusion holds for a two-connected
order-eleven wing whenever \(B\subsetneq A_b\) is terminal under
one-concept addition.
\end{corollary}

\begin{proof}
The two relative wings satisfy \(|D_0|,|D_1|\ge\rho_{\Damp}\ge10\).
Consequently,
\[
 2|B|=|H|-|D_0|-|D_1|\le39-20,
\]
and hence \(|B|\le9\).

Suppose that \(|D_b|=10\).  If \(B\cup\{d\}\) were ample for some
\(d\in D_b\), then \(d\) would be a corner of this one-concept extension.
Since \(B\in\Damp\), reversing a peeling of \(B\) would give
\(B\cup\{d\}\in\Damp\).  Moreover,
\(\operatorname{Cor}(A_b)\subseteq B\subseteq B\cup\{d\}\).
The pair
\[
 B\cup\{d\}\subsetneq A_b
\]
would therefore be a trapped \(\Damp\)-interval with nine relative
concepts, contrary to \(\rho_{\Damp}\ge10\).  Thus no one-concept
addition is ample.  No one-concept deletion \(A_b\setminus\{d\}\) is
ample either: the ample corner-deletion criterion would make
\(d\in D_b\) a corner of \(A_b\), contrary to
\(\operatorname{Cor}(A_b)\subseteq B\).  The pair is terminal.

If \(D_b\) is two-connected, the passive-gate normalization and
Corollary~\ref{cor:damp-small-lower-family-two-arm-holonomy}, with
\(K=B\), apply.  They give the two-cycle, the displayed excess budget,
and \(|B|\ge|L|+6\ge7\).  For an order-eleven wing, terminality is
assumed, and the same argument applies without change.
\end{proof}

\begin{proposition}[Pair-rich opposite wings have only three alternating
support graphs]
\label{prop:damp-pair-rich-opposite-wing-support-graphs}
Let
\[
 H=(A_0\times\{0\})\cup(A_1\times\{1\})
\]
be ample, and put \(B=A_0\cap A_1\).  Suppose that, for each
\(b\in\{0,1\}\), there are two disjoint direction pairs
\[
 P_{b,1},P_{b,2}
 \in\operatorname{Sh}(A_b)\setminus\operatorname{Sh}(B).
\]
Then the four pairs are distinct.  Moreover, either three of them are
pairwise disjoint, or their two-coloured support graph is an alternating
four-cycle, an alternating four-edge path, or two disjoint alternating
two-edge paths.
\end{proposition}

\begin{proof}
The two-layer shadow-intersection identity gives
\[
 \operatorname{Sh}(A_0)\cap\operatorname{Sh}(A_1)
 =\operatorname{Sh}(B).
\]
Consequently a pair acquired by one wing cannot also be acquired by the
other.  Together with disjointness inside each wing, this proves that the
four pairs are distinct.

Colour a pair by the wing which acquires it.  Each colour class is a
matching of order two, so the four-edge graph has maximum degree at most
two and every cycle is alternating.  Suppose that it has no matching of
order three.  A component cannot be an isolated edge together with a
three-edge path, a two-edge path together with two isolated edges, or any
finer partition of four edges, because each of those graphs has a
three-edge matching.  Nor can it be a longer cycle, since there are only
four edges.  The only remaining possibilities are one alternating
four-cycle, one alternating four-edge path, and two disjoint alternating
two-edge paths.
\end{proof}

\begin{lemma}[A crossed pair freezes its common carrier]
\label{lem:damp-crossed-acquired-pair-carrier-freezing}
Let
\[
 H=(A_0\times\{0\})\cup(A_1\times\{1\})
\]
be ample, put \(B=A_0\cap A_1\), and denote the layer coordinate by
\(e\).  Suppose that \(x,y,z\) are distinct directions whose singletons
are shattered by \(B\), and that
\[
 \{x,y\}\in\operatorname{Sh}(A_0)\setminus\operatorname{Sh}(B),
 \qquad
 \{x,z\}\in\operatorname{Sh}(A_1)\setminus\operatorname{Sh}(B).
 \label{eq:damp-crossed-acquired-pairs}
\]
Then there are unique traces
\[
 q_{xy}\in Q_{\{x,y\}},
 \qquad q_{xz}\in Q_{\{x,z\}}
\]
such that
\begin{align}
 B|_{\{x,y\}}=A_1|_{\{x,y\}}
   &=Q_{\{x,y\}}\setminus\{q_{xy}\},
 \label{eq:damp-crossed-pair-first-extenture}\\
 B|_{\{x,z\}}=A_0|_{\{x,z\}}
   &=Q_{\{x,z\}}\setminus\{q_{xz}\}.
 \label{eq:damp-crossed-pair-second-extenture}
\end{align}
In particular, \(\{e,x,y\}\) and \(\{e,x,z\}\) are minimal nonfaces of
\(\operatorname{Sh}(H)\).  Their unique omitted traces are respectively
\[
 (1,q_{xy})
 \qquad\hbox{and}\qquad
 (0,q_{xz}),
\]
where the first entry is the \(e\)-coordinate.

Moreover, every \(x\)-edge of \(B\) lies in the common slice
\begin{equation}
 y=1-q_{xy}(y),
 \qquad z=1-q_{xz}(z).
 \label{eq:damp-crossed-pair-carrier-slice}
\end{equation}
Consequently, if two distinguished occurrences of the direction \(x\)
are different \(x\)-edges of \(B\), an anchor path between them in the
connected \(x\)-carrier of \(B\) uses neither \(y\) nor \(z\).  Every
step of this path is therefore a square of \(B\) with support
\(\{x,t\}\) for some \(t\notin\{x,y,z\}\).
\end{lemma}

\begin{proof}
Projection is a pc-minor operation and therefore preserves ampleness.
The projection of \(B\) to \(\{x,y\}\) shatters both singletons but not
the pair.  Its shattered-set complex consequently has three members, so
the Sandwich equality makes the projection a three-vertex path in the
square.  It omits one unique trace, which we call \(q_{xy}\).

The shadow-intersection identity gives
\[
 \{x,y\}\notin\operatorname{Sh}(A_1).
\]
The projection \(A_1|_{\{x,y\}}\) contains
\(B|_{\{x,y\}}\), which already has three traces, and it cannot contain
the fourth.  This proves
\eqref{eq:damp-crossed-pair-first-extenture}.  The same argument with the
two layers interchanged proves
\eqref{eq:damp-crossed-pair-second-extenture}.

The \(e=0\) section realizes all four traces on \(\{x,y\}\), while the
\(e=1\) section realizes all except \(q_{xy}\).  Thus the projection of
\(H\) to \(\{e,x,y\}\) is the cube with the single trace
\((1,q_{xy})\) removed.  The pair \(\{x,y\}\) is shattered by \(H\),
and the pairs \(\{e,x\}\) and \(\{e,y\}\) are shattered because \(B\)
shatters \(x\) and \(y\).  Hence \(\{e,x,y\}\) is a minimal nonface.
The argument for \(\{e,x,z\}\) is symmetric; its missing trace lies in
the \(e=0\) layer.

Consider an \(x\)-edge of \(B\).  Its two endpoints have one common
\(y\)-value.  If that value were \(q_{xy}(y)\), the endpoint whose
\(x\)-value is \(q_{xy}(x)\) would realize the omitted trace \(q_{xy}\),
contrary to
\eqref{eq:damp-crossed-pair-first-extenture}.  Its common \(y\)-value is
therefore \(1-q_{xy}(y)\).  Applying the same argument to \(q_{xz}\)
fixes its \(z\)-value and proves
\eqref{eq:damp-crossed-pair-carrier-slice}.

The anchors of the \(x\)-edges of an ample family form an ample, hence
connected, carrier.  Equation
\eqref{eq:damp-crossed-pair-carrier-slice} makes both \(y\) and \(z\)
constant on this carrier.  Neither direction can label an anchor-path
edge.  Finally, an anchor edge in direction \(t\) is exactly a full
\(\{x,t\}\)-square of \(B\), which proves the last assertion.
\end{proof}

\begin{corollary}[One signed carrier object covers the three alternating
graphs]
\label{cor:damp-alternating-arm-incidence-carriers}
In the setting of
Proposition~\ref{prop:damp-pair-rich-opposite-wing-support-graphs},
suppose in addition that every direction occurring in the four acquired
pairs is shattered as a singleton by \(B\).  Form a bipartite graph
\(J\) whose four vertices are the two acquired pairs of each wing, joining
two opposite-wing vertices once for their common direction.  If the four
pairs have no matching of order three, then
\[
 J\cong K_{2,2},\qquad J\cong P_4,
 \qquad\hbox{or}\qquad J\cong2K_2.
\]
For every edge of \(J\), the common direction carries the two omitted
trace labels and the conditioned carrier slice of
Lemma~\ref{lem:damp-crossed-acquired-pair-carrier-freezing}.  If its two
distinguished arm edges do not coincide, their transport is a path of
squares in \(B\) which avoids the two partner directions.

Thus the alternating four-cycle, four-edge path, and two disjoint
two-edge paths are not three unrelated unpositioned exceptions: they are
respectively the incidence shadows \(K_{2,2}\), \(P_4\), and \(2K_2\) of
one signed carrier-transport problem.
\end{corollary}

\begin{proof}
Two acquired pairs belonging to the same wing are disjoint, so two
opposite-wing pairs meet in at most one direction and \(J\) is a simple
subgraph of \(K_{2,2}\).  The three support graphs in
Proposition~\ref{prop:damp-pair-rich-opposite-wing-support-graphs} have
pair-incidence graphs \(K_{2,2}\), \(P_4\), and \(2K_2\), respectively.
The carrier assertion is exactly
Lemma~\ref{lem:damp-crossed-acquired-pair-carrier-freezing}, applied at
each edge of \(J\).
\end{proof}

\begin{proposition}[Opposite sharp two-arm systems have three alternating
support graphs]
\label{prop:damp-opposite-sharp-two-arm-support-graphs}
Let
\[
 H=(A_0\times\{0\})\cup(A_1\times\{1\})
\]
be ample, and put \(B=A_0\cap A_1\).  Suppose that, for each
\(b\in\{0,1\}\), the terminal relative pair
\(B\subsetneq A_b\) satisfies the sharp two-arm conclusion of
Corollary~\ref{cor:damp-gate-holonomy-exponential-cost}.  For its two
arms, let
\[
 P_{b,i}=\{f_{b,i},z_{b,i}\}\qquad(i\in\{1,2\})
\]
be the direction pair of the terminal active--passive square, where
\(f_{b,i}\) is the direction of the shore edge and \(z_{b,i}\) is its
passive direction.  Then:
\begin{enumerate}
\item \(P_{b,1}\) and \(P_{b,2}\) are disjoint for each \(b\);
\item the four sets \(P_{b,i}\) are pairwise distinct;
\item either three of the four pairs are pairwise disjoint, or the
  two-coloured graph with edge set
  \(\{P_{b,i}:b\in\{0,1\},\,i\in\{1,2\}\}\) is one of the following:
  an alternating four-cycle, an alternating four-edge path, or two
  disjoint alternating two-edge paths.
\end{enumerate}
Thus, in the sharp shared-overlap branch, failure to produce three
independent direction pairs has only three coordinate-incidence forms.
\end{proposition}

\begin{proof}
Fix \(b\).  The equality description writes \(B\) as its active slice
\(L_b\) together with two pendant three-edge arms.  Every edge of either
arm is a bridge in the one-inclusion graph of \(B\).  In a partial cube,
a bridge is the unique edge in its coordinate class: if a second edge had
the same coordinate, convexity and connectivity of the two semicubes
would give a path avoiding the first edge.  The six arm edges therefore
have six distinct directions.  In particular
\(P_{b,1}\cap P_{b,2}=\varnothing\), which proves item~1.

The full \(P_{b,i}\)-square belongs to \(A_b\), so
\[
 P_{b,i}\in\operatorname{Sh}(A_b).
\]
On the other hand, the passive direction \(z_{b,i}\) labels exactly one
edge of \(B\), namely the first edge of its pendant arm.  An ample family
which shatters \(P_{b,i}\) contains a full \(P_{b,i}\)-square, and such a
square has two \(z_{b,i}\)-edges.  Hence
\[
 P_{b,i}\notin\operatorname{Sh}(B).
\]
Proposition~\ref{prop:damp-pair-rich-opposite-wing-support-graphs} now
applies to these two pairs in each wing.  It proves simultaneously that
the four pairs are distinct and that failure of a matching of order three
has exactly the three forms in item~3.
\end{proof}

\begin{lemma}[Two pendant three-edge arms have a unique small core]
\label{lem:damp-two-pendant-three-edge-arms-unique-small-core}
Let \(T\) be a tree.  Suppose that \(T\) is obtained from a path \(C\)
of order \(\ell\in\{1,2,3\}\) by attaching two internally
vertex-disjoint pendant paths of length three, each meeting \(C\) only
at its initial vertex.  Then \(C\) and the unordered pair of attached
paths are uniquely determined by \(T\) among all such decompositions
with core order \(\ell\).
\end{lemma}

\begin{proof}
For \(\ell=1\), the tree is the seven-vertex path, its core is the unique
central vertex, and its two components are the two arms.

For \(\ell=2\), write the core as the edge \(xy\).  If the two arms are
attached at different endpoints, \(T\) is the eight-vertex path and
\(xy\) is its unique central edge.  If both are attached at, say, \(x\),
then \(x\) is the unique vertex of degree three.  The three components of
\(T-x\) have orders \(1,3,3\); the one-vertex component contains \(y\),
and the two three-vertex components are the arms.

For \(\ell=3\), write the core as \(x-y-z\).  Up to reversing this path,
there are four attachment patterns.  Attaching at \(x,z\) gives the
nine-vertex path, whose central three-vertex path is unique.  Attaching
twice at \(x\) gives one degree-three vertex with branch lengths
\(2,3,3\); the length-two branch is the rest of the core.  Attaching
twice at \(y\) gives one degree-four vertex with branch lengths
\(1,1,3,3\); the two short branches complete the core.  Finally, attaching
at \(x,y\) gives one degree-three vertex, namely \(y\), with branch
lengths \(1,3,4\).  The length-one branch is the vertex \(z\), the
length-three branch is the arm at \(y\), and the first vertex on the
length-four branch is \(x\), after which its last three vertices form the
other arm.  In every case the decomposition is forced.  These patterns
exhaust the unordered pairs of attachment vertices, proving the lemma.
\end{proof}

\begin{corollary}[Opposite sharp two-arm wings do not fit below order ten]
\label{cor:damp-opposite-sharp-two-arm-small-overlap-impossible}
In the setting of
Proposition~\ref{prop:damp-opposite-sharp-two-arm-support-graphs}, one
has
\[
 |B|\ge10.
\]
Equivalently, if \(|B|\le9\), the two opposite relative wings cannot both
attain equality in the \(3k\) holonomy bound.
\end{corollary}

\begin{proof}
Suppose that \(|B|\le9\).  For either wing, the sharp two-arm description
writes the one-inclusion graph of \(B\) as the graph of its active slice
\(L_b\), with two pendant three-edge paths attached.  It also gives
\[
 |L_b|=|B|-6\in\{1,2,3\}.
\]
The ample family \(L_b\) is connected.  Its graph is therefore a vertex,
an edge, or a three-vertex path.  The whole graph of \(B\) is consequently
a tree of the form in
Lemma~\ref{lem:damp-two-pendant-three-edge-arms-unique-small-core}.

Apply that lemma to the two wing decompositions of the same graph \(B\).
Their unordered pairs of arms coincide.  The square-support pair attached
to an arm consists of the coordinate labels of its attachment edge and
its terminal edge, so the two wings have the same unordered pair of
square supports.  This contradicts item~2 of
Proposition~\ref{prop:damp-opposite-sharp-two-arm-support-graphs}, which
says that all four supports are distinct.  Thus \(|B|\ge10\).
\end{proof}

\begin{lemma}[An ample one-vertex completion of a tree is a square fan]
\label{lem:damp-ample-one-vertex-tree-square-fan}
Let \(B\) be ample, let \(v\in B\), and suppose that the graph
\[
 T=G(B)-v
\]
is a nonempty tree.  Write \(d=\deg_{G(B)}(v)\).  Define the
\emph{square link} \(\Lambda_v\) as follows.  Its vertices are the \(d\)
edge directions incident with \(v\), and two such directions are adjacent
when their square through \(v\) belongs to \(B\).  Then \(\Lambda_v\) is
a tree.  More precisely:
\begin{enumerate}
\item every cube of \(B\) of dimension at least two is a square through
  \(v\);
\item \(B\) contains exactly \(d-1\) squares;
\item a direction not incident with \(v\) labels exactly one edge of
  \(G(B)\), while a direction \(i\) incident with \(v\) labels
  \(1+\deg_{\Lambda_v}(i)\) edges.
\end{enumerate}
Thus the cycles created by the one added vertex form one connected tree
of squares; no longer even cycle or disconnected collection of square
repairs is possible.
\end{lemma}

\begin{proof}
Every square avoiding \(v\) would be a cycle in \(T\), so every square
contains \(v\).  A cube of dimension at least three has a square face
avoiding each one of its vertices.  Hence no such cube belongs to \(B\),
which proves item~1.

We shall use the elementary connected-carrier property of ample families.
For a coordinate \(i\), identify the two \(i\)-sections of \(B\) after
deleting \(i\).  Their intersection is ample by the two-layer
intersection lemma and is therefore connected.  Its vertices correspond
to the \(i\)-edges of \(B\), and its edges correspond to squares joining
two \(i\)-edges.  Consequently all edges of one direction \(i\) are joined
by a sequence of squares.

If \(i\) is not incident with \(v\), this carrier consists of one edge:
a square in a connecting sequence would contain \(v\) and would therefore
have an \(i\)-edge incident with \(v\).  If \(i\) is incident with \(v\),
every other \(i\)-edge is the edge opposite the unique \(i\)-edge at \(v\)
in a square through \(v\).  Such squares are indexed exactly by the
neighbours of \(i\) in \(\Lambda_v\).  This proves item~3.

The graph \(\Lambda_v\) is a forest.  Indeed, a cycle
\(i_1i_2\cdots i_si_1\) in \(\Lambda_v\) would give in \(T\) the
subdivision
\[
 v^{\oplus i_1},
 v^{\oplus\{i_1,i_2\}},
 v^{\oplus i_2},\ldots,
 v^{\oplus\{i_s,i_1\}},
 v^{\oplus i_1}
\]
of that cycle, contrary to the assumption that \(T\) is a tree.

Let \(r\) be the number of nonconstant coordinates of \(B\), let
\(q=|E(\Lambda_v)|\), and put \(n=|B|\).  Item~3 gives
\begin{equation}
 |E(G(B))|=r+2q.
 \label{eq:damp-tree-completion-carrier-edge-count}
\end{equation}
On the other hand, \(T\) has \(n-2\) edges and adjoining \(v\) adds
\(d\) edges, so
\[
 |E(G(B))|=n-2+d.
\]
Since \(B\) is ample, its shattered and strongly shattered sets coincide.
By item~1 they are the empty set, the \(r\) singleton directions, and the
\(q\) square supports through \(v\).  Therefore
\[
 n=1+r+q.
\]
Substitution in
\eqref{eq:damp-tree-completion-carrier-edge-count} gives
\(q=d-1\).  A forest on \(d\) vertices with \(d-1\) edges is connected,
so \(\Lambda_v\) is a tree and item~2 follows.
\end{proof}

\begin{lemma}[An ample two-vertex completion of a tree is a two-apex
square complex]
\label{lem:damp-ample-two-vertex-tree-square-complex}
Let \(B\) be ample, let \(u,v\in B\) be distinct, and suppose that

\[
 T=G(B)-\{u,v\}
\]

is a nonempty tree.  Put

\[
 d_u=\deg_{G(B)}(u),\qquad d_v=\deg_{G(B)}(v),
 \qquad
 \delta_{uv}=\begin{cases}
  1,&uv\in E(G(B)),\\
  0,&uv\notin E(G(B)).
 \end{cases}
\]

Then every cube of \(B\) of dimension at least two is a square meeting
\(\{u,v\}\), no two squares have the same direction support, and the
number \(q(B)\) of squares is

\begin{equation}
 q(B)=d_u+d_v-\delta_{uv}-2.
 \label{eq:damp-two-vertex-tree-square-count}
\end{equation}

Thus two vertices added to a tree cannot create a higher-dimensional cube,
a repeated square carrier, or an arbitrary collection of cycles: all
cycles are uniquely supported squares incident with one of the two added
vertices, and their number is determined by the two degrees.
\end{lemma}

\begin{proof}
Every square avoiding \(u,v\) would be a cycle in \(T\), so every square
meets \(\{u,v\}\).  A cube of dimension at least four has a square face
avoiding any two prescribed vertices: after choosing two free directions,
there are at least four parallel square faces and at most two of them meet
\(\{u,v\}\).  Hence no such cube belongs to \(B\).

Suppose that \(B\) contains a three-cube \(C\).  If \(C\) contains at most
one of \(u,v\), the square face opposite that vertex avoids both.  If it
contains both vertices and they are not antipodal in \(C\), they agree in
one of the three cube directions, and the opposite square face again avoids
both.  Finally, if \(u,v\) are antipodal in \(C\), then
\(C-\{u,v\}\) contains the canonical six-cycle \(Q_3^{--}\), which is a
cycle in \(T\).  All alternatives contradict the tree hypothesis.  Thus
every nontrivial cube is a square meeting \(\{u,v\}\).

Fix a direction pair \(S\) supporting a square.  Successively intersecting
the two coordinate sections in the directions of \(S\) shows that the
family of anchors of the full \(S\)-squares is ample, by the two-layer
intersection lemma.  It is therefore connected.  If two distinct
\(S\)-squares existed, the first edge on a path between their anchors would
join two opposite \(S\)-squares into a three-cube.  This has just been
excluded.  Consequently every square support occurs exactly once.

Let \(r\) be the number of nonconstant directions of \(B\), let
\(n=|B|\), and write \(q=q(B)\).  Ampleness and the preceding conclusions
give

\begin{equation}
 n=1+r+q.
 \label{eq:damp-two-vertex-tree-shadow-count}
\end{equation}

We also use the edge--shadow identity

\begin{equation}
 |E(G(B))|=r+2q.
 \label{eq:damp-two-vertex-tree-edge-shadow-count}
\end{equation}

For completeness, in one direction \(i\), the anchors of the \(i\)-edges
form an ample carrier.  Its shattered supports are exactly the sets
\(S\setminus\{i\}\) with \(i\in S\in\operatorname{Sh}(B)\).  The
Sandwich equality in that carrier therefore counts the \(i\)-edges by
those supports.  Summing over \(i\) proves the edge--shadow identity; here
only singleton and pair supports occur, which gives
\eqref{eq:damp-two-vertex-tree-edge-shadow-count}.

The tree \(T\) has \(n-2\) vertices and \(n-3\) edges.  Restoring \(u,v\)
adds \(d_u+d_v-\delta_{uv}\) edges, since the edge \(uv\), when present,
is counted twice by the two degrees.  Hence

\[
 |E(G(B))|=n-3+d_u+d_v-\delta_{uv}.
\]

Substituting
\eqref{eq:damp-two-vertex-tree-shadow-count} and
\eqref{eq:damp-two-vertex-tree-edge-shadow-count}, and cancelling \(r\),
gives

\[
 2q=q-2+d_u+d_v-\delta_{uv},
\]

which is exactly
\eqref{eq:damp-two-vertex-tree-square-count}.
\end{proof}

\begin{lemma}[A one-vertex completion of a unicyclic graph has a controlled
square link]
\label{lem:damp-ample-one-vertex-unicyclic-square-link}
Let \(B\) be ample, let \(v\in B\), and suppose that
\[
 T=G(B)-v
\]
is connected and unicyclic.  Assume that \(T\) contains at most one
square, and let \(q_0\in\{0,1\}\) be its number of squares.  If
\(d=\deg_{G(B)}(v)\), define the square link \(\Lambda_v\) as in
Lemma~\ref{lem:damp-ample-one-vertex-tree-square-fan}.  Then:
\begin{enumerate}
\item \(B\) has no cube of dimension at least three, and every square
  support occurs exactly once;
\item \(B\) has exactly \(d\) squares, of which \(d-q_0\) contain \(v\);
\item if \(q_0=0\), then \(\Lambda_v\) is connected and unicyclic;
\item if \(q_0=1\), then the unique cycle of \(T\) is that square and
  \(\Lambda_v\) is a tree.
\end{enumerate}
Thus the cycle rank missing from the tree case appears exactly once: as
the unique cycle of the square link when the base cycle is not already a
square.
\end{lemma}

\begin{proof}
A three-cube avoiding \(v\) would place six squares in \(T\).  A
three-cube containing \(v\) has three distinct square faces opposite
\(v\), again all contained in \(T\).  Both conclusions contradict the
hypothesis on \(T\).  Every cube of larger dimension contains a
three-cube, so \(B\) has no cube of dimension at least three.

As in the proof of
Lemma~\ref{lem:damp-ample-two-vertex-tree-square-complex}, the anchors of
all squares with one fixed direction support form an ample, hence
connected, carrier.  Two anchors would have an edge between consecutive
anchors and therefore create a three-cube.  Thus each square support
occurs once.

Let \(r\) be the number of nonconstant directions, let \(q\) be the
number of squares, and put \(n=|B|\).  The preceding paragraph and
ampleness give
\[
 n=1+r+q,
 \qquad
 |E(G(B))|=r+2q.
\]
The graph \(T\) has \(n-1\) vertices and, because it is unicyclic,
\(n-1\) edges.  Restoring \(v\) adds \(d\) edges, so
\[
 |E(G(B))|=n-1+d.
\]
Substitution gives \(q=d\).  The edges of \(\Lambda_v\) are precisely
the squares through \(v\), and hence
\begin{equation}
 |E(\Lambda_v)|=d-q_0.
 \label{eq:damp-unicyclic-square-link-edge-count}
\end{equation}

A cycle of length \(s\) in \(\Lambda_v\) gives its canonical
one-subdivision, a cycle of length \(2s\), in \(T\).  Consequently
\(\Lambda_v\) has cycle rank at most one.  If \(q_0=0\), then
\eqref{eq:damp-unicyclic-square-link-edge-count} gives \(d\) edges on
\(d\) vertices.  If it had \(c\) connected components, its cycle rank
would be \(c\); hence \(c=1\), and it is connected and unicyclic.

If \(q_0=1\), the unique cycle of \(T\) is its square.  A cycle in
\(\Lambda_v\) would instead give a subdivided cycle of length at least
six, so no such cycle exists.  Equation
\eqref{eq:damp-unicyclic-square-link-edge-count} now gives a forest with
\(d\) vertices and \(d-1\) edges.  It is therefore a tree.
\end{proof}

\begin{corollary}[The pure two-unit anomaly has two square apices]
\label{cor:damp-eta-two-two-apex-normal-form}
Use the notation of
Corollary~\ref{cor:damp-small-lower-family-two-arm-holonomy}.  Suppose that
\(|K|\le9\), that the directed two-cycle has \(\varepsilon=2\), and that
all of this excess lies in the uncounted remainder:

\begin{equation}
 \eta=2,
 \qquad |Z_0|=2,
 \qquad 2^{t(d)}-1=t_{\mathcal C}(d)\quad(d\in D),
 \qquad |S_1|=|S_2|=2.
 \label{eq:damp-pure-eta-two-budget}
\end{equation}

Then \(|L|=1\) and \(|K|=9\).  After choosing one collar vertex for each
passive direction, exactly two vertices \(u,v\in K\) remain uncounted, and

\[
 G(K)-\{u,v\}=P_7,
\]

the path obtained by attaching the two pendant three-edge arms to the
one-vertex active slice.  Consequently \(K\) has only uniquely supported
squares, every such square meets \(\{u,v\}\), and

\[
 q(K)=\deg_K(u)+\deg_K(v)-\mathbf1_{uv\in E(K)}-2.
\]

Let \(P_i=\{z_i,f_i\}\) be the support of the terminal square of arm
\(i\), where \(z_i\) is its passive attachment direction and \(f_i\) its
terminal direction.  At most one of \(P_1,P_2\) is shattered by \(K\).
Therefore at least one of the two terminal pairs is acquired by the upper
family.
\end{corollary}

\begin{proof}
Equality in every term of
\eqref{eq:damp-pure-eta-two-budget}, apart from \(\eta=2\), has the same
consequences as the sharp part of
Corollary~\ref{cor:damp-gate-holonomy-exponential-cost}: the two shores are
disjoint active edges, each has one passive direction and one collar, and
the counted lower vertices form two pendant three-edge arms attached to
the lift of \(L\).  Moreover

\[
 |K|=|L|+8.
\]

The shore-root lemma gives \(L\ne\varnothing\), while \(|K|\le9\), so
\(|L|=1\) and \(|K|=9\).  The seven counted vertices therefore induce
\(P_7\), and the two units of \(\eta\) are exactly the two remaining
vertices \(u,v\).  Lemma
\ref{lem:damp-ample-two-vertex-tree-square-complex} gives the asserted
square structure and count.

It remains to compare the terminal supports.  The direction of the
attachment edge of arm \(i\) is \(z_i\), and \(z_i\) occurs on no other
edge of \(P_7\).  Its two adjacent edges in the path have directions
\(z_{3-i}\) and, say, \(g_i\), neither of which is the active terminal
direction \(f_i\).  If a \(P_i\)-square of \(K\) contained exactly one of
\(u,v\), deleting that apex from the square would leave in \(P_7\) a
two-edge path with directions \(z_i,f_i\), a contradiction.  A square
cannot avoid both apices by the preceding lemma.  Hence every
\(P_i\)-square in \(K\) contains both \(u\) and \(v\).

We next show that \(f_1\ne f_2\).  Suppose instead that both terminal
edges have direction \(f\).  The two terminal leaves lie in the same
\(f\)-semicube: the path from the central vertex to either leaf crosses
\(f\) only on its terminal edge.  After deleting the two terminal
\(f\)-edges, that semicube contains no other vertex of \(P_7\).  Semicubes
of a partial cube are convex and hence connected, so its two terminal
leaves have a path through the two apices.  The leaves belong to the same
bipartition class.  With only two apices available, they must therefore
have a common apex and be at distance two.

After translating the active word of the central vertex to the empty set,
the two leaves have respective direction words
\[
 z_1\mathbin\triangle g_1\mathbin\triangle f,
 \qquad
 z_2\mathbin\triangle g_2\mathbin\triangle f.
\]
The passive directions \(z_1,z_2\) are distinct and the directions
\(g_1,g_2\) are active.  Distance two consequently forces \(g_1=g_2\).
The two active shore edges then have the same endpoint pair
\(\{g_1,g_1\mathbin\triangle f\}\), contrary to the disjointness of the
cycle shores forced by
\eqref{eq:damp-pure-eta-two-budget}.  Thus \(f_1\ne f_2\), and the
active--passive pairs \(P_1,P_2\) are disjoint.

If squares on both supports existed and \(u,v\)
were adjacent, the direction of \(uv\) would belong to both supports.  If
they were nonadjacent, they would be opposite in each square, and both
supports would equal \(\operatorname{Diff}(u,v)\).  Either conclusion is
impossible.  Thus at most one terminal pair is shattered by \(K\).  Both
pairs are shattered by the upper family, which contains the corresponding
terminal squares, so every pair not shattered by \(K\) is acquired there.
\end{proof}

\begin{corollary}[The two-unit budget has exactly seven profiles]
\label{cor:damp-two-unit-holonomy-seven-profiles}
Use the notation of
Corollary~\ref{cor:damp-gate-holonomy-exponential-cost}, let
\(S_1,S_2\) be a directed two-cycle, and assume that its excess is
\(\varepsilon=2\).
Up to interchanging \(S_1,S_2\), exactly one of the following alternatives
holds:
\begin{enumerate}
\item \(\eta=2\), and \(S_1,S_2\) are disjoint two-concept shores;
\item \(\eta=1\), and \(S_1,S_2\) are two-concept shores meeting in one
  concept;
\item \(\eta=1\), and \(S_1,S_2\) are disjoint shores of orders two and
  three;
\item \(\eta=0\), and \(S_1,S_2\) have orders two and three and meet in
  one concept;
\item \(\eta=0\), and \(S_1,S_2\) are disjoint shores of orders two and
  four;
\item \(\eta=0\), and \(S_1,S_2\) are disjoint three-concept shores;
\item \(\eta=0\), the cycle shores are disjoint two-concept sets, and a
  third passive coordinate has a singleton shore disjoint from their
  union.
\end{enumerate}
In alternatives 1--6, \(|Z_0|=2\), each cycle shore is represented by
exactly one passive coordinate, and there is no other nonempty passive
shore.  Alternative 7 is the rigid singleton-extra-shore profile.
Consequently alternatives 1--3 are precisely the two-unit profiles with
an uncounted lower vertex.
\end{corollary}

\begin{proof}
Write the exact excess identity as
\begin{equation}
 2=
 \eta+(|Z_0|-2)
 +\sum_{d\in D}\bigl(2^{t(d)}-1-t_{\mathcal C}(d)\bigr)
 +(|S_1|-2)+(|S_2|-2).
 \label{eq:damp-two-unit-profile-identity}
\end{equation}
Every summand is a nonnegative integer.  The two distinct cycle shores
require two passive coordinates.  Every further coordinate has a
nonempty shore and therefore contributes at least one additional incidence
to \(\sum_d(t(d)-t_{\mathcal C}(d))\).  Since
\(2^t-1\ge t\), the second and third terms in
\eqref{eq:damp-two-unit-profile-identity} together are at least
\(2(|Z_0|-2)\).  Hence \(|Z_0|\le3\).  If \(|Z_0|=3\), equality in both
contributions forces precisely the singleton-extra-shore profile of
Corollary~\ref{cor:damp-one-unit-holonomy-anomaly-trichotomy}, which is
alternative 7.

It remains that \(|Z_0|=2\).  Each of the two distinct cycle shores is
then represented by its unique passive coordinate, so
\(t(d)=t_{\mathcal C}(d)\) for every \(d\in D\).  The middle sum in
\eqref{eq:damp-two-unit-profile-identity} is therefore
\[
 |S_1\cap S_2|:
\]
a concept belonging to one shore contributes zero, while a concept
belonging to both contributes \(2^2-1-2=1\).  Put
\[
 a=|S_1\cap S_2|,
 \qquad x=|S_1|-2,
 \qquad y=|S_2|-2.
\]
Then
\[
 \eta+a+x+y=2.
\]
The shores have order at least two.  If \(\eta=2\), all three remaining
parameters vanish, giving alternative 1.  If \(\eta=1\), the remaining
unit is either \(a=1\) or, up to reversal, \(x=1\), giving alternatives
2 and 3.  Finally let \(\eta=0\).  The formal solution \(a=2,x=y=0\)
would make the two two-concept shores equal, contrary to their being
distinct members of the directed cycle.  The remaining solutions are
\((a,x,y)=(1,1,0),(0,2,0),(0,1,1)\), up to reversal, which are exactly
alternatives 4--6.
\end{proof}

\begin{lemma}[A disjoint passive shore cannot hide its terminal square]
\label{lem:damp-disjoint-shore-terminal-square-acquired}
Use the active--passive normalization of
Lemma~\ref{lem:damp-passive-neighbours-force-punctured-cube}, and suppose
that \(L=\{a\}\).  Let \(z\in Z\) satisfy
\[
 S_z\cap S_w=\varnothing
 \qquad\text{for every }w\in Z\setminus\{z\}.
\]
If \(xy\) is an active edge contained in \(S_z\), with direction \(f\),
then
\[
 \{z,f\}\in
 \operatorname{Sh}(A)\setminus\operatorname{Sh}(K).
\]
Thus the square over an edge of a shore separated from all the other
shores is genuinely acquired by the upper family; arbitrary additional
vertices of \(K\) cannot hide another square with the same support.
\end{lemma}

\begin{proof}
The definition of \(S_z\) gives
\[
 x,y,x^{\oplus z},y^{\oplus z}\in A,
 \qquad
 x,y\in D,\qquad x^{\oplus z},y^{\oplus z}\in K.
\]
These four vertices form a full \(\{z,f\}\)-square in \(A\), so the pair
is shattered there.

Suppose that \(K\) also shattered \(\{z,f\}\).  Since \(K\) is ample, it
would contain a full square with this support.  The anchors of all
\(\{z,f\}\)-squares of \(A\), after the two directions are deleted, form
the intersection of the four corresponding sections of \(A\).  Repeated
application of the two-layer intersection lemma makes this anchor family
ample and hence connected.  It therefore contains a path from the anchor
of the displayed split square to the anchor of a full square in \(K\).

At the first anchor every passive coordinate other than \(z\) has value
\(\sigma\).  The anchor of a full square in \(K\) cannot have this same
passive restriction: its two vertices on the \(z=\sigma(z)\) side would
be distinct members of the one-concept slice \(L\).  Hence the anchor path
has a first edge which changes some passive coordinate
\(w\ne z\).  Immediately before this edge, consider the two vertices of
the current square whose \(z\)-value agrees with \(\sigma\).  They have
the full passive restriction \(\sigma\), and therefore belong either to
\(D\) or to the lift of \(L\).  Since they are distinct and \(L\) has one
concept, at least one of them, say \(d\), belongs to \(D\).

The current square contains \(d^{\oplus z}\), and the three-cube joining
the two consecutive anchors contains \(d^{\oplus w}\).  Both vertices lie
outside the common passive layer of \(D\), so both belong to \(K\).  Thus
\[
 d\in S_z\cap S_w,
\]
contrary to the shore-separation hypothesis.  Hence \(K\) does not
shatter \(\{z,f\}\), as required.
\end{proof}

\begin{lemma}[At most one uncounted vertex separates the shore directions]
\label{lem:damp-one-uncounted-disjoint-shore-directions}
In the notation of
Corollary~\ref{cor:damp-gate-holonomy-exponential-cost}, suppose that
\(L=\{a\}\), that the nonempty shores are pairwise disjoint, and that
\(\eta\le1\).  If distinct shores \(S_z,S_w\) contain active edges of
directions \(f_z,f_w\), respectively, then
\[
 f_z\ne f_w.
\]
Consequently the two active--passive pairs
\(\{z,f_z\}\) and \(\{w,f_w\}\) are disjoint.
\end{lemma}

\begin{proof}
Suppose that both shore edges had direction \(f\).  Their \(z\)- and
\(w\)-flips are two \(f\)-edges of \(K\), lying in the passive layers
\(\sigma\oplus z\) and \(\sigma\oplus w\).  The anchor family of the
\(f\)-edges of \(K\) is ample and hence connected.  Along an anchor path
between these two edges, take the first step which changes the passive
restriction.  Before that step the passive restriction is
\(\sigma\oplus z\).  The new \(f\)-edge lies either in the base layer
\(\sigma\), if the step removes \(z\), or in a layer differing from
\(\sigma\) in at least two coordinates.

The base layer contains only the lift of the one concept \(a\), so it
contains no edge.  Because the shores are pairwise disjoint, the punctured
fibres forced by
Lemma~\ref{lem:damp-passive-neighbours-force-punctured-cube} occur only in
single-coordinate passive layers.  The collar vertices do as well.
Consequently both endpoints of an edge in a layer with at least two
passive flips would have to belong to the uncounted remainder.  This
requires at least two uncounted vertices, contrary to \(\eta\le1\).
Thus the connected \(f\)-carrier cannot join the two displayed edges, a
contradiction.
\end{proof}

\begin{corollary}[The feasible uncounted two-unit profiles are pair-rich]
\label{cor:damp-two-unit-uncounted-profiles-acquire-pair}
In the setting of
Corollary~\ref{cor:damp-two-unit-holonomy-seven-profiles}, assume in
addition that \(|K|\le9\).  Alternatives 1--3 have the following
outcomes:
\begin{enumerate}
\item alternative 1 has the two-apex \(P_7\) normal form of
  Corollary~\ref{cor:damp-eta-two-two-apex-normal-form}, and both of its
  terminal pairs are acquired;
\item alternative 2 is impossible;
\item in alternative 3, \(|L|=1\), \(|K|=9\), and the terminal pair of
  each of its two cycle shores is acquired.
\end{enumerate}
\end{corollary}

\begin{proof}
In alternative 1, the two terminal pairs are disjoint by
Corollary~\ref{cor:damp-eta-two-two-apex-normal-form}.  Here \(L\) is a
singleton and the two shores are disjoint, so
Lemma~\ref{lem:damp-disjoint-shore-terminal-square-acquired} shows that
both pairs are acquired.  Consider
alternative 2.  The exact count gives
\[
 |K|=|L|+8.
\]
The shore-root lemma and \(|K|\le9\) force \(|L|=1\) and \(|K|=9\).
Write the shores as
\[
 S_1=\{r_1,o\},
 \qquad
 S_2=\{r_2,o\},
\]
where \(r_1,r_2\) are the two lower roots and \(o\) is their common
non-root owner.  Let \(z_i\) be the passive direction of \(S_i\), let
\(f_i\) be the direction of the active edge \(r_io\), and write
\(L=\{a\}\).  Since \(r_i\in R\), the two-concept family
\(\{a,r_i\}\) is ample and hence \(ar_i\) is an active edge.  On the
other hand, \(o\notin R\), so \(a\) is not adjacent to \(o\).  Thus,
after the single uncounted vertex \(v\) is deleted, the counted graph
\(T=G(K)-v\) is the shared-owner eight-cycle.

The graph \(T\) is square-free and unicyclic, so
Lemma~\ref{lem:damp-ample-one-vertex-unicyclic-square-link} makes the
square link \(\Lambda_v\) connected and unicyclic.  Its unique link
cycle must have length four, and its subdivision uses all eight vertices
of \(T\).  If its consecutive directions are
\(i_1,i_2,i_3,i_4\), the cyclic direction sequence of the subdivided
cycle is
\[
 i_2,i_1,i_3,i_2,i_4,i_3,i_1,i_4.
\]
No direction occurs on opposite edges.  In contrast, the provenance of
the shared-owner eight-cycle gives the sequence
\[
 z_1,f_2,f_1,z_2,z_1,f_2,f_1,z_2,
\]
in which every direction occurs on opposite edges.  This contradiction
excludes alternative~2.

Finally consider alternative 3, and let \(S_1\) be its two-concept
shore, and let \(S_2\) be its three-concept shore.  Again \(|L|=1\),
\(|K|=9\), and the two shores are disjoint.  The shore \(S_1\) is an
active edge by
Lemma~\ref{lem:damp-two-concept-cycle-shore-active-edge}.  The shore
\(S_2\) also contains an active edge.  Indeed, it contains a common
non-root owner \(o\).  The two-concept family \(L\cup\{o\}\) is nonample,
so \(o\) is not adjacent to the unique member of \(L\).  A path from
\(o\) to \(L\) in the connected ample gate \(L\cup S_2\) must therefore
begin with an edge internal to \(S_2\).

Lemma~\ref{lem:damp-disjoint-shore-terminal-square-acquired} applied to
these two shore edges gives one acquired pair from each shore.  Their
active directions are distinct by
Lemma~\ref{lem:damp-one-uncounted-disjoint-shore-directions}, while their
passive directions are distinct and are disjoint from every active
direction.  The two acquired pairs are therefore disjoint, proving
item~3.
\end{proof}

\begin{corollary}[The vertex-free two-unit profiles are separated by their
passive carriers]
\label{cor:damp-two-unit-vertex-free-passive-carriers}
In the setting of
Corollary~\ref{cor:damp-two-unit-holonomy-seven-profiles}, assume in
addition that \(|K|\le9\).  The four vertex-free alternatives 4--7 have
the following outcomes:
\begin{enumerate}
\item alternative 4, whose shores have orders two and three and meet in
  one concept, is impossible;
\item in each of alternatives 5--7, the two cycle shores supply two
  disjoint acquired direction pairs.
\end{enumerate}
Consequently every feasible two-unit wing is pair-rich: it has two
disjoint terminal direction pairs acquired by the upper family.
\end{corollary}

\begin{proof}
In all four alternatives, \(\eta=0\) and the exact count gives
\(|K|=|L|+8\).  Hence \(|L|=1\), \(|K|=9\), and we write
\(L=\{a\}\).  Equality in the collar and punctured-fiber counts describes
all vertices of \(K\): for every nonempty passive direction \(z\), the
family contains the collar edge
\[
 a\,a^{\oplus z},
\]
and, for every active concept \(d\), it contains
\(d^{\oplus T}\) for every nonempty
\(T\subseteq Z(d)=\{z:d\in S_z\}\).  It contains no further vertex.

Suppose first that alternative 4 holds, and let
\(S_{z_1}\cap S_{z_2}=\{c\}\).  The \(z_1\)-edges of \(K\) are exactly
\[
 a\,a^{\oplus z_1},
 \qquad
 c^{\oplus z_2}\,c^{\oplus\{z_1,z_2\}}.
\]
After contracting \(z_1\), their two anchors are \(a\) and
\(c^{\oplus z_2}\).  They differ in the passive direction \(z_2\) and
also in at least one active direction, because \(a\in L\) and
\(c\in D\).  Thus the two anchors are nonadjacent.  On the other hand,
the anchor family of all edges in one direction is the intersection of
the two corresponding sections of \(K\); it is ample and hence connected.
This contradiction excludes alternative 4.

Now consider alternatives 5--7.  Each cycle shore is disjoint from every
other nonempty shore.  Fix one of them, write it as \(S_z\), and let
\(o\in S_z\setminus R\) be the common owner supplied by the predecessor
on the holonomy cycle.  The gate \(L\cup S_z\) is ample and hence
connected.  The two-concept family \(L\cup\{o\}\) is not ample, because
\(o\) is not a lower root.  Since \(L=\{a\}\), the vertices \(a,o\) are
not adjacent.  A path from \(o\) to \(a\) in \(G(L\cup S_z)\) therefore
starts with an active edge \(xy\) having both endpoints in \(S_z\); let
\(f\) be its direction.

Disjointness of the shores and the exact vertex description above show
that the collar \(a\,a^{\oplus z}\) is the only \(z\)-edge of \(K\).
Thus \(K\) cannot shatter \(\{z,f\}\): by ampleness it would then contain
a full \(\{z,f\}\)-square and hence two \(z\)-edges.  In the upper
family, both \(x,y\) and their \(z\)-flips are present, so this full square
does exist there.  The pair \(\{z,f\}\) is therefore acquired.  Applying
the argument to each of the two cycle shores gives two acquired pairs.
Their active directions are distinct by
Lemma~\ref{lem:damp-one-uncounted-disjoint-shore-directions}, applied with
\(\eta=0\); their passive directions are distinct and disjoint from the
active directions.  Hence the pairs are disjoint, proving item~2.
\end{proof}

\begin{corollary}[Pair-rich comparisons in the two-unit row]
\label{cor:damp-two-unit-cross-wing-reduction}
Let
\[
 H=(A_0\times\{0\})\cup(A_1\times\{1\})
\]
be ample, put \(B=A_0\cap A_1\), and suppose that \(|B|\le9\).  If the
terminal relative pair \(B\subsetneq A_0\) has two-cycle holonomy and
excess two, then it has two disjoint acquired direction pairs.  If the
opposite wing \(B\subsetneq A_1\) is sharp, has the uncounted
\(\eta=1\) one-unit form, or has any feasible two-unit profile, then
either the two wings supply three pairwise disjoint acquired pairs or
their four acquired pairs form one of the three alternating support
graphs in
Proposition~\ref{prop:damp-pair-rich-opposite-wing-support-graphs}.

At this stage the only apparent profile outside the pair-rich comparison
is the one-unit extra-owner wing with disjoint shores of orders two and
three.  Proposition~\ref{prop:damp-two-unit-extra-owner-comparison} below
shows that an extra-owner wing opposite a two-unit wing is pair-rich as
well.
\end{corollary}

\begin{proof}
Corollary~\ref{cor:damp-two-unit-uncounted-profiles-acquire-pair} excludes
alternative 2 and gives two disjoint acquired pairs in alternatives 1 and
3 of the two-unit classification.  Alternative 4 is impossible, and
Corollary~\ref{cor:damp-two-unit-vertex-free-passive-carriers} gives two
disjoint acquired pairs in alternatives 5--7.  This proves the first
assertion.

A sharp wing has two such pairs by the proof of
Proposition~\ref{prop:damp-opposite-sharp-two-arm-support-graphs}; the
uncounted one-unit wing has them by
Corollary~\ref{cor:damp-eta-one-square-fan-normal-form}; and a second
two-unit wing has them by the preceding paragraph.  Proposition
\ref{prop:damp-pair-rich-opposite-wing-support-graphs} now gives the
claimed matching or one of the three alternating graphs.  The one-unit
trichotomy and
Proposition~\ref{prop:damp-shared-owner-one-unit-anomaly-impossible}
leave only the extra-owner form outside these pair-rich cases.  Its
pair-poor subcase is localized by
Corollary~\ref{cor:damp-extra-owner-pair-poor-active-slice}.
\end{proof}

\begin{lemma}[A small three--four arm decomposition is intrinsic except
on a path]
\label{lem:damp-small-three-four-arm-decomposition}
Let \(G\) be a connected graph of order at most nine which is obtained
from a path core \(C\) of order one or two by attaching the following two
components at vertices of \(C\):
\begin{enumerate}
\item a pendant path of length three;
\item a connected four-vertex graph \(F\), attached by one edge at a
  distinguished vertex of \(F\), where \(F\) is a tree or a four-cycle.
\end{enumerate}
Call the first component the \emph{ordinary arm}.  If \(G\) admits two
such decompositions, then their ordinary arms coincide, except in the
following two cases:
\begin{enumerate}
\item \(G=P_8\), the core has order one, and the two decompositions
  interchange the length-three and length-four sides of the central edge;
\item \(G=P_9\), the core has order two, and the two decompositions
  interchange the length-three and length-four sides of the central
  three-vertex path.
\end{enumerate}
Moreover, if \(G\) also has a sharp two-arm decomposition as in
Lemma~\ref{lem:damp-two-pendant-three-edge-arms-unique-small-core}, then
the ordinary arm of every three--four decomposition is one of its two
three-edge arms.
\end{lemma}

\begin{proof}
A connected ample graph of order four is either a four-vertex tree or a
four-cycle.  If \(F\) is a four-cycle, it is the unique nontrivial block
of \(G\).  The block--cut tree then identifies \(F\), the path core, and
the ordinary arm.  Such a graph has no sharp two-arm decomposition,
because the latter is a tree.

Suppose that \(F\) is a tree.  There are four rooted possibilities for
\((F,\text{attachment vertex})\): a four-vertex path rooted at an endpoint
or at an internal vertex, and a three-leaf star rooted at its centre or
at a leaf.  In the internally rooted path case, the attachment vertex has
degree three and the two branches inside \(F\) have orders one and two.  In
the centre-rooted star case it has degree four.  In the leaf-rooted star
case, the unique degree-three centre lies one edge beyond the attachment
vertex and has two one-vertex branches away from the core.  These three
degree and branch-order profiles are different from one another and from
the path core.  They therefore identify the rooted copy of \(F\); the
remaining three-vertex component is the ordinary arm.

It remains that \(F\) is a path rooted at an endpoint.  Thus \(G\) is
formed from its core by attaching paths of lengths three and four.  For a
one-vertex core this is \(P_8\), and either side of its central edge can
be declared the ordinary arm.  For a two-vertex core, attaching at
different core endpoints gives \(P_9\), with the two choices stated in
the lemma.  Attaching both paths at the same endpoint gives a unique
degree-three vertex with branch lengths \(1,3,4\), which distinguishes
the ordinary arm.  This proves the first assertion.

Finally, compare with a sharp decomposition.  Its graph is a path core
of order at most three with two length-three arms.  The uniqueness proof
of Lemma~\ref{lem:damp-two-pendant-three-edge-arms-unique-small-core},
or the same branch-length comparison, shows that any pendant
three-vertex component in a three--four decomposition is one of those two
arms.  In the only path possibilities, this is the length-three side of
the central edge or central core.  Hence its attachment and terminal
edges, and not merely its vertex set, coincide with a sharp arm.
\end{proof}

\begin{corollary}[The \(\eta=1\) anomaly is a tree of squares]
\label{cor:damp-eta-one-square-fan-normal-form}
Use the notation of
Corollary~\ref{cor:damp-one-unit-holonomy-anomaly-trichotomy}, suppose
that \(|K|\le9\), and assume its first alternative:
\[
 \eta=1,
\]
with two disjoint two-concept cycle shores.  Then \(|L|\in\{1,2\}\).
After choosing one collar vertex for each passive direction, there is
exactly one remaining, uncounted vertex \(v\in K\), and
\(G(K)-v\) is obtained from the path \(G(L)\) by attaching two pendant
paths of length three.  The square link of \(v\) is a tree.  In
particular, if \(d=\deg_{G(K)}(v)\), then \(K\) contains exactly \(d-1\)
squares, all through \(v\).

For each arm, let \(z_i\) be the direction of its attachment edge and
let \(f_i\) be the direction of its terminal edge.  Then
\[
 P_i=\{z_i,f_i\}\notin\operatorname{Sh}(K).
\]
The two pairs \(P_1,P_2\) are disjoint.  They are precisely the supports
of the two terminal squares acquired by the upper family.
\end{corollary}

\begin{proof}
The one-unit identity gives
\[
 |K|=|L|+7,
\]
so \(|L|\le2\).  The shore-root lemma gives \(L\ne\varnothing\), and
hence \(|L|\in\{1,2\}\).

For a shore \(S_i=\{r_i,o_i\}\), the opposite active section \(M_{z_i}\)
contains the edge \(r_io_i\) and at least one collar vertex in
\(L\cap M_{z_i}\).  All but at most one vertex of \(M_{z_i}\) are
accounted for by the lower slice, the collar count, and the punctured
shore fibers.  Since \(M_{z_i}\) is connected and bipartite, one may
choose a collar vertex \(a_i\) adjacent to one endpoint of \(r_io_i\).
Thus
\[
 a_i,r_i,o_i
\]
induce a three-vertex path.  After lifting to the \(z_i\)-layer and
adjoining the collar edge from \(L\), this is a pendant path of length
three in \(K\).

The two arms have disjoint nonbase vertices: their passive restrictions
are respectively the \(z_1\)- and \(z_2\)-flips of the common passive
word, and their shore endpoints have distinct active words.  No edge joins
the two arms, since such vertices differ in both \(z_1\) and \(z_2\);
and an arm vertex is adjacent to the lower slice only along its collar
edge.  The six arm vertices together with \(L\) account for all but the
single unit \(\eta=1\).  Calling that remaining vertex \(v\), we obtain
the asserted tree \(G(K)-v\).  Lemma
\ref{lem:damp-ample-one-vertex-tree-square-fan} gives the square-link
conclusion.

Within \(G(K)-v\), the passive direction \(z_i\) occurs only on the
attachment edge of arm \(i\).  By the carrier description in
Lemma~\ref{lem:damp-ample-one-vertex-tree-square-fan}, a square of \(K\)
using \(z_i\), if one exists, has this attachment edge as the side
opposite \(v\).  Its other direction is therefore the direction of the
next edge of the arm, not the terminal direction \(f_i\).  Hence
\(\{z_i,f_i\}\) is not strongly shattered by \(K\), and ampleness gives
\(P_i\notin\operatorname{Sh}(K)\).

The directions \(z_1,z_2\) are distinct passive coordinates and do not
occur on the active edges of either arm.  If \(f_1=f_2\), connectedness
of the \(f_1\)-carrier would put both terminal edges opposite the
\(f_1\)-edge at \(v\) in two squares.  The two corresponding length-two
paths in \(G(K)-v\), together with the path between the arms through
\(L\), would form a cycle, a contradiction.  Thus \(f_1\ne f_2\), and
the two pairs are disjoint.  Finally, the four vertices over the shore
edge \(r_io_i\) form the acquired terminal square in the upper family,
as in the zero-excess case.
\end{proof}

\begin{lemma}[Comparison of one-unit two-arm decompositions]
\label{lem:damp-eta-one-two-arm-comparison}
Let \(G\) be an ample partial cube of order at most nine admitting an
\(\eta=1\) decomposition from
Corollary~\ref{cor:damp-eta-one-square-fan-normal-form}, with terminal
support pairs \(P_1,P_2\).
\begin{enumerate}
\item If \(G\) also admits a sharp two-arm decomposition, one of its
  terminal support pairs is equal to \(P_1\) or \(P_2\).
\item If \(G\) admits two \(\eta=1\) decompositions, then either their
  terminal-support systems have a common pair, three of their four pairs
  are pairwise disjoint, or \(G\) is the following exceptional graph:
  a four-cycle with a pendant path of length two attached at each vertex
  of one antipodal pair.
\item If \(G\) also admits a three--four decomposition from the
  extra-owner anomaly, then its ordinary-arm support is equal to one of
  \(P_1,P_2\), or it is disjoint from both of them.
\end{enumerate}
In the exceptional graph of item~2, write \(z_1,z_2\) for the directions
of the central square and \(f_1,f_2\) for the terminal directions of the
two pendant paths.  The two decompositions use the two opposite remaining
vertices of the square as their respective cores and extra vertices, and
their four terminal supports are
\begin{equation}
 \{z_i,f_j\}\qquad(i,j\in\{1,2\}).
 \label{eq:damp-eta-one-exceptional-four-supports}
\end{equation}
\end{lemma}

\begin{proof}
Delete the extra vertex of an \(\eta=1\) decomposition.  The remaining
tree has one of only three branch profiles:
\[
 P_7,\qquad P_8,\qquad
 \text{a three-branch tree with branch lengths }1,3,3.
\]
Lemma~\ref{lem:damp-ample-one-vertex-tree-square-fan} says that all
cycles restored by the extra vertex form a connected square fan.  We
compare decompositions using the block--cut tree of this fan.

First suppose that \(G\) is a tree.  Then the extra vertex is a leaf.
The branch lengths identify both three-edge arms unless \(G\) itself is
\(P_8\) or \(P_9\).  In the identified cases a second sharp,
\(\eta=1\), or three--four decomposition retains an entire arm and hence
its terminal support.  On \(P_8\) or \(P_9\), distinct
\(\eta=1\) decompositions move the extra leaf from one end to the other.
Reading the first and third edge of each resulting arm gives three
pairwise disjoint support pairs.  The same reading shows that a sharp
decomposition retains one arm, while an ordinary three--four arm either
coincides with an \(\eta=1\) arm or has both of its support directions
outside the two disjoint \(\eta=1\) supports.

Now suppose that \(G\) contains a cycle.  Every nontrivial block is part
of the square fan.  If the fan has at least two squares, any second vertex
whose deletion is a tree must belong to every square.  The square-link
tree is consequently a star, and that vertex is the common neighbour of
the fan apex.  Deleting either endpoint of this common fan edge leaves
the same branches in the block--cut tree.  Their lengths identify one
whole arm, unless the other decomposition uses its ordinary arm outside
both old arms; in the latter case its first and terminal directions are
disjoint from both \(P_1\) and \(P_2\).  This proves the three asserted
alternatives in the multi-square case.

If the two decompositions use the same fan apex, then they are already
decompositions of the same deleted tree.  The \(P_7,P_8\), and
\((1,3,3)\) branch profiles above identify their arms directly, so their
terminal-support systems agree.  Thus the preceding common-neighbour
argument covers every genuinely different apex.

It remains that the square fan is one four-cycle.  Removing any candidate
extra vertex turns this block into a three-edge path.  Attach to its four
vertices the components of the block--cut tree and compare their orders
with the three base profiles displayed above.  Unless one three-edge arm
is unchanged, the only way that two removals both leave a \(P_7\) core
profile is to attach two-edge paths at two opposite vertices of the
square and to use the other opposite pair as the two possible
core--extra choices.  This is the exceptional graph in item~2.
If the central square directions are \(z_1,z_2\), each of the two core
choices pairs each \(z_i\) with one of the two tail directions
\(f_1,f_2\), and the choices interchange the tails.  This gives exactly
\eqref{eq:damp-eta-one-exceptional-four-supports}.  For a sharp or
three--four comparison, the same four block--cut positions leave an old
arm unchanged or put the ordinary arm outside both old supports.  This
completes all three comparisons.
\end{proof}

\begin{corollary}[A pair-poor extra-owner wing has a rooted-claw overlap]
\label{cor:damp-extra-owner-pair-poor-active-slice}
In the one-unit extra-owner alternative of
Corollary~\ref{cor:damp-one-unit-holonomy-anomaly-trichotomy}, suppose
that \(|K|\le9\).  If \(|L|=1\), then the two disjoint shores supply two
disjoint acquired direction pairs.  Consequently an extra-owner wing
which is not pair-rich must satisfy
\[
 |L|=2,\qquad |K|=9.
\]
In that case, write \(G(L)=ab\), and choose the notation so that the
ordinary three-edge arm is
\[
 a p_1p_2p_3.
\]
Then \(G(K)\) is obtained by adjoining a vertex \(c\), three leaves
\(q_1,q_2,q_3\) adjacent to \(c\), and one edge \(rc\), where
\(r\in\{a,b\}\).  Thus there are only two rooted graphs, according as
\(r=a\) or \(r=b\).
\end{corollary}

\begin{proof}
The one-unit identity gives \(|K|=|L|+7\), so the only possible slice
orders are one and two.  Suppose that \(L=\{a\}\).  The two-concept shore
is an active edge by
Lemma~\ref{lem:damp-two-concept-cycle-shore-active-edge}.  The
three-concept shore contains a non-root owner \(o\).  Since
\(L\cup\{o\}\) is nonample, \(o\) is not adjacent to \(a\); connectivity
of the ample shore gate therefore supplies an active edge internal to
that shore.

The shores are disjoint, so
Lemma~\ref{lem:damp-disjoint-shore-terminal-square-acquired} makes both
active--passive pairs acquired.  Their active directions are distinct by
Lemma~\ref{lem:damp-one-uncounted-disjoint-shore-directions}, applied with
\(\eta=0\), and their passive directions are distinct.  Thus the two
pairs are disjoint.  The order identity now gives the final assertion.

It remains to identify the graph when \(|L|=2\) and the wing is
pair-poor.  The proof of
Lemma~\ref{lem:damp-small-three-four-arm-decomposition} writes
\(G(K)\) as the edge \(G(L)\), an ordinary pendant path of length three,
and a rooted four-vertex component \(F\).  The latter consists of its
collar vertex and the three concepts of the larger shore.  The passive
direction attaching \(F\) occurs on only the collar edge of \(K\): the
two shores are disjoint and \(\eta=0\), so the exact punctured-fibre
count contains no further edge in that passive direction.

Suppose that two concepts of the larger shore were adjacent, in active
direction \(f\).  Their passive translates in \(K\), together with the
corresponding two vertices in the upper family, give an acquired
active--passive pair.  This pair is disjoint from the ordinary-arm pair.
Indeed, the two passive directions are distinct and are disjoint from
the active directions, while the terminal edge of the ordinary pendant
arm is a bridge of \(G(K)\).  A bridge in a partial cube is the unique
edge of its direction, so its active direction cannot equal \(f\).
The wing would therefore be pair-rich, contrary to the hypothesis.

Hence the three shore concepts are pairwise nonadjacent in \(F\).
Since \(F\) is connected, each is adjacent to the collar vertex.  Thus
\(F\) is a claw rooted at its centre.  There are no edges between the
two attached components in the active--passive normal form, so the edge
list in the statement is complete.
\end{proof}

\begin{lemma}[The two rooted-claw overlaps have rigid cuts]
\label{lem:damp-extra-owner-rooted-claw-cuts}
Let \(T^\parallel\) and \(T^\times\) be the two trees with vertex set
\[
 \{a,b,p_1,p_2,p_3,c,q_1,q_2,q_3\}
\]
and common edge set
\[
 \{ab,ap_1,p_1p_2,p_2p_3,cq_1,cq_2,cq_3\},
\]
where the final edge is \(ac\) in \(T^\parallel\) and \(bc\) in
\(T^\times\).  Call \(ap_1p_2p_3\) the \emph{old arm}.  Then:
\begin{enumerate}
\item if deleting two vertices leaves a path on seven vertices, viewed
  as two pendant paths of length three attached to its centre, then the
  tree is \(T^\times\), the deleted vertices are two of
  \(q_1,q_2,q_3\), and one of the two arms is the old arm;
\item suppose that deleting one vertex leaves a connected graph obtained
  from a one-vertex core by attaching a pendant path of length three and
  a rooted four-vertex tree.  If the three nonattachment vertices of the
  four-vertex tree contain an edge, then the pendant path is the old arm;
\item the only cut producing component orders \((3,5)\) is the cut at
  \(a\) in \(T^\times\), and its three-vertex component is
  \(\{p_1,p_2,p_3\}\).  The only cut producing component orders
  \((4,4)\) is the cut at \(b\) in \(T^\times\); one of its rooted
  four-vertex components is the claw on
  \(\{c,q_1,q_2,q_3\}\), whose three nonattachment vertices are
  independent.  Neither tree has a cut with component orders
  \((2,3,3)\).
\end{enumerate}
When either tree is embedded isometrically in a cube, the occurrence of
the old arm in items~1--3 has the same attachment and terminal directions
as in the original rooted-claw decomposition.
\end{lemma}

\begin{proof}
For item~1, the claw centre \(c\) cannot be deleted: its three leaves
would be isolated, and only one further vertex may be removed.  Thus at
least two neighbours of \(c\) must be deleted in order that the remaining
graph have maximum degree two.  In \(T^\parallel\), the neighbour \(a\)
cannot be deleted, because the retained vertices in the old arm, at
\(b\), and at the claw would lie in different components.  Hence two
claw leaves must be deleted, but then \(a\) still has the three neighbours
\(b,p_1,c\).  The remainder is not a path.  In \(T^\times\), the
neighbour \(b\) likewise cannot be deleted without separating the old
arm from the claw.  Deleting two claw leaves is therefore necessary and
sufficient, and the remaining path is
\[
 q_i c b a p_1p_2p_3.
\]
Its two length-three arms meet at \(a\), and one is the old arm.

For item~2, deletion of one vertex from a tree leaves a connected graph
only when that vertex is a leaf.  Checking the leaf orbits and then the
component orders after deleting the proposed one-vertex core gives
\[
\begin{array}{c|c|c|c}
 \text{tree}&\text{deleted leaf}&\text{core}&\text{outcome}\ \\ \hline
 T^\parallel&b&a&
   \text{old arm and a centre-rooted claw}\ \\
 T^\times&p_3&b&
   \text{a three-edge arm and a centre-rooted claw}\ \\
 T^\times&q_i&a&
   \text{old arm and a rooted four-vertex tree}.
\end{array}
\]
These are all the possibilities: deleting \(p_3\) or a claw leaf in
\(T^\parallel\) gives no cut into components of orders three and four;
in the last row, the alternative core \(b\) attaches the three-vertex
component at its middle and hence does not give a pendant path of length
three.  In the first two rows the three vertices of the claw away from
its attachment are independent.  The hypothesis of item~2 therefore
leaves only the last row, where the pendant path is the old arm.

Finally, direct deletion of a vertex gives the following complete list of
component orders.  In \(T^\parallel\), deleting \(a,p_1,p_2,c\), or a
leaf gives, respectively,
\[
 (1,3,4),\qquad(2,6),\qquad(1,7),\qquad
 (1,1,1,5),\qquad(8).
\]
Thus no vertex gives one of \((3,5),(4,4),(2,3,3)\).  In
\(T^\times\), deleting \(a\) gives \((3,5)\), deleting \(b\) gives
\((4,4)\), and every other deletion gives one of
\[
 (8),\qquad(2,6),\qquad(1,7),\qquad(1,1,1,5).
\]
The components in the two relevant cuts are exactly those described in
item~3.  This proves the cut assertions.  The final direction statement
holds because the repeated arm has the same edges, and every edge of a
tree partial cube is its own direction class.
\end{proof}

\begin{proposition}[A two-unit wing excludes a pair-poor extra-owner wing]
\label{prop:damp-two-unit-extra-owner-comparison}
Let
\[
 H=(A_0\times\{0\})\cup(A_1\times\{1\})
\]
be ample, put \(B=A_0\cap A_1\), and suppose that \(|B|\le9\).  If one
terminal relative pair \(B\subsetneq A_b\) has two-cycle holonomy and
excess two, while the opposite wing has the one-unit extra-owner form,
then the extra-owner wing is pair-rich.  Consequently either the two
wings supply three pairwise disjoint acquired direction pairs, or their
four acquired pairs form one of the three alternating support graphs in
Proposition~\ref{prop:damp-pair-rich-opposite-wing-support-graphs}.
\end{proposition}

\begin{proof}
Suppose for a contradiction that the extra-owner wing is pair-poor.
Corollary~\ref{cor:damp-extra-owner-pair-poor-active-slice} gives
\(|B|=9\) and identifies \(G(B)\) with one of the two rooted-claw trees
in Lemma~\ref{lem:damp-extra-owner-rooted-claw-cuts}.  Let \(Q\) be the
acquired support determined by the attachment and terminal directions of
its old arm.  We compare this tree with the seven two-unit alternatives
of Corollary~\ref{cor:damp-two-unit-holonomy-seven-profiles}.

In alternative~1, deleting the two uncounted vertices leaves the
seven-vertex path with two length-three arms from
Corollary~\ref{cor:damp-eta-two-two-apex-normal-form}.  Item~1 of the
rooted-claw cut lemma says that one of these arms is the old arm.  Its
acquired terminal support is therefore \(Q\).

Alternative~2 is impossible because deleting its uncounted vertex leaves
the six-cycle with two pendant terminal edges described in the proof of
Corollary~\ref{cor:damp-two-unit-uncounted-profiles-acquire-pair}; an
induced subgraph of the tree \(G(B)\) cannot contain that cycle.  In
alternative~3, deleting the uncounted vertex leaves a one-vertex core, a
three-edge arm, and a rooted four-vertex component.  The three shore
vertices in the latter contain an edge, as proved in the final paragraph
of that same corollary.  Item~2 of the rooted-claw cut lemma therefore
makes the three-edge arm the old arm, so its acquired support is again
\(Q\).

Alternative~4 is impossible by
Corollary~\ref{cor:damp-two-unit-vertex-free-passive-carriers}.  In the
remaining vertex-free alternatives, deleting the one-vertex active slice
leaves components of orders \((3,5)\), \((4,4)\), and \((2,3,3)\),
respectively.  Item~3 of the rooted-claw cut lemma shows that the first
profile can occur only with its three-vertex component equal to the old
arm; its acquired support is \(Q\).  In the second profile, the only
possible cut has a rooted claw whose three shore vertices are independent,
contrary to the internal shore edge used in the proof of
Corollary~\ref{cor:damp-two-unit-vertex-free-passive-carriers}.  The third
component profile does not occur in either rooted-claw tree.

Thus every feasible comparison makes the two-unit wing acquire \(Q\).
The extra-owner wing also acquires \(Q\), whereas \(B\) does not shatter
it.  This contradicts the two-layer shadow-intersection identity
\[
 \operatorname{Sh}(A_0)\cap\operatorname{Sh}(A_1)
 =\operatorname{Sh}(B).
\]
The extra-owner wing is therefore pair-rich.  The two-unit wing is
pair-rich by Corollaries
\ref{cor:damp-two-unit-uncounted-profiles-acquire-pair} and
\ref{cor:damp-two-unit-vertex-free-passive-carriers}; applying
Proposition~\ref{prop:damp-pair-rich-opposite-wing-support-graphs} gives
the final alternative.
\end{proof}

\begin{proposition}[The extra-owner anomaly contradicts the opposite wing
or forces three pairs]
\label{prop:damp-extra-owner-anomaly-opposite-wing}
Let
\[
 H=(A_0\times\{0\})\cup(A_1\times\{1\})
\]
be ample, put \(B=A_0\cap A_1\), and suppose that \(|B|\le9\).
For each \(b\in\{0,1\}\), suppose that the terminal pair
\(B\subsetneq A_b\) has a two-cycle and is either sharp or has the
one-unit extra-owner form: \(\eta=0\), with disjoint cycle shores of
orders two and three.  Then the two wings cannot both be sharp.  If at
least one has the extra-owner form, either the two wing structures are
incompatible, or the acquired-support family contains three pairwise
disjoint direction pairs.
\end{proposition}

\begin{proof}
The sharp--sharp case is
Corollary~\ref{cor:damp-opposite-sharp-two-arm-small-overlap-impossible}.
Consider an extra-owner wing.  Its two-concept shore is an active edge by
Lemma~\ref{lem:damp-two-concept-cycle-shore-active-edge}.  Together with
its unique passive coordinate it gives a full square in the upper wing.
The passive coordinate labels exactly one edge of \(B\), because
\(\eta=0\), so this square support is not shattered by \(B\).  We call it
the ordinary-arm support.

The opposite active section \(M_z\) of the three-concept shore has order
four: it consists of the unique collar vertex and the three shore
concepts.  It is ample.  Hence its graph is a four-vertex tree or a
four-cycle.  The graph of \(B\) is obtained from the active slice, of
order
\[
 |B|-7\le2,
\]
by attaching the ordinary three-edge arm and this rooted four-vertex
component.  It therefore has the form of
Lemma~\ref{lem:damp-small-three-four-arm-decomposition}.

If the other wing is sharp, the last assertion of that lemma makes the
ordinary arm one of its two sharp arms.  If the other wing also has the
extra-owner form and the ordinary arms are intrinsic, the first assertion
makes the two ordinary arms coincide.  In either case their support pairs
coincide.  But each pair belongs to the shadow of its wing and not to
\(\operatorname{Sh}(B)\), contradicting
\[
 \operatorname{Sh}(A_0)\cap\operatorname{Sh}(A_1)
 =\operatorname{Sh}(B)
\]
from the two-layer shadow-intersection identity.

The only remaining case consists of two extra-owner decompositions of
\(P_8\) or \(P_9\).  Write the consecutive edge directions of the path
as \(x_0,x_1,\ldots,x_{q-2}\), where \(q\in\{8,9\}\).  For \(P_8\),
the two ordinary-arm supports are
\[
 \{x_0,x_2\},\qquad \{x_4,x_6\}.
\]
The four-vertex branch of the first decomposition has passive direction
\(x_3\) and a shore edge of direction \(x_5\), so it also gives the full
square support \(\{x_3,x_5\}\).  These three pairs are pairwise disjoint.
For \(P_9\), the corresponding three supports are
\[
 \{x_0,x_2\},\qquad \{x_5,x_7\},\qquad \{x_4,x_6\},
\]
again pairwise disjoint.  Each is acquired by one of the two wings,
because the relevant passive direction is a bridge of \(B\) and hence
cannot belong to a square of \(B\).  This proves the proposition.
\end{proof}

\begin{proposition}[All one-unit opposite-wing comparisons reduce to one
central square]
\label{prop:damp-one-unit-opposite-wing-central-square}
Let
\[
 H=(A_0\times\{0\})\cup(A_1\times\{1\})
\]
be ample, put \(B=A_0\cap A_1\), and suppose that \(|B|\le9\).
Assume that both terminal pairs \(B\subsetneq A_b\) have two-cycle
holonomy and excess at most one.  Then one of the following holds:
\begin{enumerate}
\item the acquired-support family of the two wings contains three
  pairwise disjoint direction pairs;
\item \(B\) is the exceptional graph of
  Lemma~\ref{lem:damp-eta-one-two-arm-comparison}: a central four-cycle
  with pendant paths of length two at two opposite vertices, and both
  wings have the \(\eta=1\) form based at the other two opposite vertices.
\end{enumerate}
In the second alternative the four terminal square supports form the
complete bipartite graph
\[
 \bigl\{\{z_i,f_j\}:i,j\in\{1,2\}\bigr\}
\]
between the two central-square directions and the two tail directions.
Thus, apart from this one graph, the entire two-connected one-unit
frontier produces three independent direction pairs in one step.
\end{proposition}

\begin{proof}
By
Proposition~\ref{prop:damp-shared-owner-one-unit-anomaly-impossible}, a
one-unit wing is either of \(\eta=1\) type or of extra-owner type.  A
zero-unit wing is sharp.

Two sharp wings are impossible by
Corollary~\ref{cor:damp-opposite-sharp-two-arm-small-overlap-impossible}.
If at least one wing has the extra-owner form and neither has the
\(\eta=1\) form, Proposition
\ref{prop:damp-extra-owner-anomaly-opposite-wing} gives either a
contradiction or three pairwise disjoint acquired supports.

It remains to compare an \(\eta=1\) wing with the opposite wing.  Its two
terminal pairs are acquired by that wing and are not shattered by \(B\),
by Corollary~\ref{cor:damp-eta-one-square-fan-normal-form}.  The ordinary
support of an extra-owner wing has the same property by the proof of
Proposition~\ref{prop:damp-extra-owner-anomaly-opposite-wing}; and the
terminal pairs of a sharp wing do so by
Proposition~\ref{prop:damp-opposite-sharp-two-arm-support-graphs}.
The two-layer shadow identity
\[
 \operatorname{Sh}(A_0)\cap\operatorname{Sh}(A_1)
 =\operatorname{Sh}(B)
\]
therefore forbids an acquired pair from occurring in both wings.

Apply Lemma~\ref{lem:damp-eta-one-two-arm-comparison}.  Against a sharp
wing it forces such a repeated pair, a contradiction.  Against an
extra-owner wing it forces either a repeated pair or an ordinary support
disjoint from both \(\eta=1\) supports; the latter gives three independent
pairs.  Against a second \(\eta=1\) wing it gives a repeated pair, three
independent pairs, or the unique exceptional central-square graph.  The
repeated case is again forbidden by the shadow identity, and
\eqref{eq:damp-eta-one-exceptional-four-supports} gives the final support
description in the exceptional case.
\end{proof}

\begin{lemma}[Passive layers can shield only nearby overlap corners]
\label{lem:damp-passive-radius-one-shielding}
Let
\[
 H=(A_0\times\{0\})\cup(A_1\times\{1\})
\]
be cornerless, put \(B=A_0\cap A_1\), and fix \(b\in\{0,1\}\).
Suppose that there is a set of directions \(Z_b\) and a word
\(\sigma_b\in Q_{Z_b}\) such that every exclusive concept in
\(D_b=A_b\setminus B\) has restriction \(\sigma_b\) to \(Z_b\).
Then every corner \(v\) of \(B\) satisfies
\begin{equation}
 d_{Q_{Z_b}}(v|_{Z_b},\sigma_b)\le1.
 \label{eq:damp-passive-radius-one-shielding}
\end{equation}
If equality holds and \(z\) is the unique coordinate of \(Z_b\) on
which \(v\) and \(\sigma_b\) differ, then the exclusive shield of \(v\)
is forced:
\begin{equation}
 v^{\oplus z}\in D_b.
 \label{eq:damp-passive-distance-one-forced-owner}
\end{equation}
Equivalently, a corner of the overlap at passive distance at least two
from the exclusive layer of a wing makes the corresponding copy a corner
of the full two-layer family, while a distance-one corner has only one
possible exclusive neighbour.
\end{lemma}

\begin{proof}
Suppose that \(v\in\operatorname{Cor}(B)\) and that the distance in
\eqref{eq:damp-passive-radius-one-shielding} is at least two.  Every
member of \(D_b\) differs from \(v\) in at least two coordinates of
\(Z_b\), so \(v\) has no neighbour in \(D_b\).  Let \(Q=Q_B(v)\) be the
unique maximal cube of \(B\) through \(v\).  Every horizontal edge of
\(H\) incident with \((v,b)\) is therefore an edge of \(B\times\{b\}\),
and its direction belongs to \(Q\).  The vertical edge through \((v,b)\)
is present because \(v\in B\), and all these incident directions belong
to the single cube
\[
 Q\times\{0,1\}.
\]
Thus \((v,b)\) belongs to a unique maximal cube of \(H\), contrary to
cornerlessness.

Now suppose that the passive distance is one, with unique differing
coordinate \(z\).  The same unique-maximal-cube argument shows that
\(v\) must have a neighbour in \(D_b\).  Every member of \(D_b\) has
passive word \(\sigma_b\), so such a neighbour must change \(z\).
Changing any further coordinate would give Hamming distance at least
two.  The only possible neighbour is therefore \(v^{\oplus z}\), which
proves \eqref{eq:damp-passive-distance-one-forced-owner}.
\end{proof}

\begin{proposition}[Connected pure two-apex incidence exposes a distant
corner]
\label{prop:damp-pure-eta-two-connected-incidence-distant-corner}
Let \(B\) be an ample family of order nine with two distinguished
decompositions of the pure \(\eta=2\) form in
Corollary~\ref{cor:damp-eta-two-two-apex-normal-form}.  For decomposition
\(j\in\{0,1\}\), write \(c_j\) for the central vertex of its \(P_7\),
\(U_j\) for its two uncounted vertices, \(Z_j=\{z_{j,1},z_{j,2}\}\)
for its passive directions, and
\[
 \mathcal P_j=
 \bigl\{\{z_{j,1},f_{j,1}\},\{z_{j,2},f_{j,2}\}\bigr\}
\]
for its terminal-support system.  Suppose that the four members of
\(\mathcal P_0\cup\mathcal P_1\) are distinct, have no matching of order
three, and form a connected support graph.  Then, for some
\(j\in\{0,1\}\), there is a vertex \(v\in U_j\) such that
\[
 v\in\operatorname{Cor}(B),
 \qquad
 d_{Q_{Z_j}}(v|_{Z_j},c_j|_{Z_j})=2.
 \label{eq:damp-pure-eta-two-distant-corner}
\]
\end{proposition}

\begin{proof}[Proof by a complete normal-form audit]
Fix one distinguished decomposition and label its seven-vertex spine
\[
 P_7=p_{-3}p_{-2}p_{-1}p_0p_1p_2p_3,
 \qquad c_0=p_0.
\]
Every family in the statement is obtained by adjoining two vertices to
this fixed path.  The first has an arbitrary neighbourhood in \(P_7\);
the second has an arbitrary neighbourhood in the resulting eight-vertex
graph.  Thus there are exactly
\[
 2^7\,2^8=32768
\]
labelled candidates, with no assumption on the positions of the two
new vertices.

The verifier
\path{verify_damp_pure_eta_two_connected_incidence.py} performs the
following direct tests.  It computes the Djokovi\'c classes and retains
exactly the isometric cube families satisfying the Sandwich equality.
For every retained family it deletes every pair of vertices, rediscovers
every induced \(P_7\) with singleton centre, and checks that each of its
two attachment directions occurs nowhere else on the path.  It then
forms the four terminal pairs for every two decompositions.  Finally, it
tests \eqref{eq:damp-pure-eta-two-distant-corner} directly.

For completeness, the corner test used in the last step is exact rather
than a degree proxy.  Lemma
\ref{lem:damp-ample-two-vertex-tree-square-complex} says that every
nontrivial cube is a uniquely supported square meeting the deleted pair.
The verifier lists all such squares through \(v\), adds every incident
edge direction contained in no such square, and declares \(v\) a corner
exactly when this list has one maximal member.

The audit retains \(317\) ample partial-cube completions, counted with
their distinguished labelled \(P_7\).  It finds respectively
\[
\begin{array}{c|ccc}
 \text{support graph}&C_4&P_5&2P_3\\ \hline
 \text{comparisons}&52&32&42\\
 \text{without the corner in
 \eqref{eq:damp-pure-eta-two-distant-corner}}&0&0&42
\end{array}
\]
with the same harmless labelling multiplicities.  Since
Proposition~\ref{prop:damp-pair-rich-opposite-wing-support-graphs} gives
only \(C_4\), \(P_5\), and \(2P_3\) when a matching of order three is
absent, the connected cases are exactly the first two columns.  The full
record is
\path{damp_pure_eta_two_connected_incidence.json}; its mathematical
digest is
\[
 \texttt{e12c192cb3e1ef387bc42655cd7486ad58b9eee3fd3085c32d3b75fcca3252e8}.
\]
\end{proof}

\begin{proposition}[Disconnected pure two-apex incidence has one path
falsifier]
\label{prop:damp-pure-eta-two-disconnected-incidence-owner-collision}
Keep the notation of
Proposition~\ref{prop:damp-pure-eta-two-connected-incidence-distant-corner},
but suppose that the four terminal supports form \(2P_3\).  Define the
forced passive-owner set
\[
 O_j=
 \left\{
   v^{\oplus z}:
   \begin{array}{l}
   v\in\operatorname{Cor}(B),\
   d_{Q_{Z_j}}(v|_{Z_j},c_j|_{Z_j})=1,\\
   z\text{ is the unique differing coordinate}
   \end{array}
 \right\}.
\]
Then one of the following alternatives holds:
\begin{enumerate}
\item
\[
 O_0\cap O_1\ne\varnothing
 \qquad\hbox{or}\qquad
 (O_0\cup O_1)\cap B\ne\varnothing;
 \label{item:damp-pure-eta-two-forced-owner-collision}
\]
\item \(G(B)=P_9\).  Up to reversing the path and interchanging the
  decompositions, write its vertices as
  \(v_0v_1\cdots v_8\), with edge directions
  \(x_i\) on \(v_{i-1}v_i\).  The two \(P_7\)'s are
  \(v_0\cdots v_6\) and \(v_2\cdots v_8\), with central vertices
  \(v_3,v_5\).  Their two inner arms have the common collar support
  \[
   \{x_4,x_5\}.
   \label{eq:damp-pure-eta-two-p9-common-collar}
  \]
\end{enumerate}
\end{proposition}

\begin{proof}[Proof by the same complete normal-form audit]
The verifier used in the preceding proposition continues the exact corner
test through the \(42\) labelled \(2P_3\) comparisons.  For every corner
at passive distance one, it computes the unique word obtained by correcting
the differing passive coordinate.  Thirty-six comparisons have either a
word forced on both sides or a forced word already in \(B\).  After
identifying repeated embedded comparisons, these are \(20\) types.

The remaining six labelled comparisons, or three embedded types, all have
the path form in item~2.  In both decompositions the shared directions of
the two components reverse attachment and terminal roles.  Reading the
directions of the two overlapping \(P_7\)'s gives
\[
\begin{array}{c|cc}
 &\text{left arm}&\text{right arm}\\ \hline
 v_0\cdots v_6
   & (x_3,x_2,x_1)&(x_4,x_5,x_6)\\
 v_2\cdots v_8
   & (x_5,x_4,x_3)&(x_6,x_7,x_8).
\end{array}
\]
The collar pair is formed by the first two directions of an arm, so the
right arm in the first row and the left arm in the second row both give
\(\{x_4,x_5\}\).  No comparison remains unresolved.  These further counts
are recorded in the same JSON artifact and covered by the digest displayed
in the preceding proof.
\end{proof}

\begin{corollary}[Two pure two-unit wings force three independent pairs]
\label{cor:damp-pure-eta-two-opposite-wings-three-independent-pairs}
Let
\[
 H=(A_0\times\{0\})\cup(A_1\times\{1\})
\]
be cornerless and ample, and put \(B=A_0\cap A_1\).  Suppose that both
relative pairs \(B\subsetneq A_b\) have the pure \(\eta=2\) two-apex
form.  Then three of their four terminal acquired pairs are pairwise
disjoint.
\end{corollary}

\begin{proof}
Both terminal pairs of either wing are acquired by
Corollary~\ref{cor:damp-two-unit-uncounted-profiles-acquire-pair}.  The
two-layer shadow identity makes the two systems disjoint.  If there is no
matching of order three, Proposition
\ref{prop:damp-pair-rich-opposite-wing-support-graphs} leaves \(C_4\),
\(P_5\), and \(2P_3\).

Suppose that one of the first two profiles occurs.  Proposition
\ref{prop:damp-pure-eta-two-connected-incidence-distant-corner} gives a
decomposition \(j\) and a corner \(v\) satisfying
\eqref{eq:damp-pure-eta-two-distant-corner}.  In the active--passive
normalization for that wing, every exclusive concept has the common
passive word \(c_j|_{Z_j}\).  Lemma
\ref{lem:damp-passive-radius-one-shielding} now contradicts
cornerlessness of \(H\).

It remains to exclude \(2P_3\).  Apply
Proposition~\ref{prop:damp-pure-eta-two-disconnected-incidence-owner-collision}.
In its first alternative, every forced word in \(O_j\) must belong to
\(D_j=A_j\setminus B\), by
\eqref{eq:damp-passive-distance-one-forced-owner}.  A word in \(B\)
cannot do so.  Nor can one word belong to both \(O_0\) and \(O_1\),
because
\[
 D_0\cap D_1=\varnothing.
\]
Thus the first alternative is impossible.

In the second alternative, the active--passive normal form makes the
first two edges of every arm one side and one lift of a full collar
square in the corresponding \(A_j\).  Hence
\eqref{eq:damp-pure-eta-two-p9-common-collar} belongs to both
\(\operatorname{Sh}(A_0)\) and \(\operatorname{Sh}(A_1)\).  On the
other hand, \(B=P_9\) is a tree and shatters no direction pair.  This
contradicts
\[
 \operatorname{Sh}(A_0)\cap\operatorname{Sh}(A_1)
 =\operatorname{Sh}(B).
\]
The no-matching assumption is therefore impossible, proving the
corollary.
\end{proof}

\begin{proposition}[One pure two-apex wing resolves the sharp, uncounted,
and two-unit rows]
\label{prop:damp-one-pure-eta-two-wing-resolves-opposite-profiles}
Let
\[
 H=(A_0\times\{0\})\cup(A_1\times\{1\})
\]
be cornerless and ample, put \(B=A_0\cap A_1\), and suppose that
\(|B|\le9\).  Assume that \(B\subsetneq A_0\) has the pure
\(\eta=2\) two-apex form.  If \(B\subsetneq A_1\) is sharp, has the
uncounted \(\eta=1\) one-unit form, or has any feasible two-unit profile,
then the acquired-support family of the two wings contains three pairwise
disjoint direction pairs.
\end{proposition}

\begin{proof}[Proof by an exact-profile audit]
We explain both the finite scope of the audit and how each machine outcome
becomes a mathematical conclusion.  Normalize the pure wing by fixing its
counted spine
\[
 P_7=p_{-3}p_{-2}p_{-1}p_0p_1p_2p_3.
\]
As in
Proposition~\ref{prop:damp-pure-eta-two-connected-incidence-distant-corner},
adjoining the two uncounted vertices with arbitrary neighbourhoods gives
all \(2^7 2^8\) labelled candidates.  After the partial-cube and Sandwich
tests, there are \(158\) distinct embedded ample overlaps.  The audit
rediscovers every pure decomposition of each overlap.

The opposite profiles are recognized without a graph catalogue.  A sharp
marker is a path core of order three with two pendant three-edge arms.  An
uncounted one-unit marker becomes a path core of order two with the same
two arms after its unique uncounted vertex is deleted.  For the two-unit
row, the verifier applies the exact identity
\eqref{eq:damp-two-unit-profile-identity}: it chooses the passive set, its
common word, and the uncounted remainder, and requires the remaining
vertices to be exactly the collar vertices and all nonempty punctured
fibres of the shore-incidence sets.  It also checks ampleness of every
active gate.  This recovers alternatives 1, 3, and 5--7 of
Corollary~\ref{cor:damp-two-unit-holonomy-seven-profiles}; alternative~4
is already impossible by
Corollary~\ref{cor:damp-two-unit-vertex-free-passive-carriers}.  The
sharp and one-unit graph markers are deliberately over-inclusive, so in
particular they contain every marker furnished by the corresponding
structural theorems.

For each genuine marker the two terminal pairs are the acquired pairs
proved in the sharp, one-unit, and two-unit results above.  The audit also
records an \emph{acquired collar}: if \(z\) is passive and an active edge
of direction \(g\) joins the active slice to a shore, the
active--passive normal form supplies the full \(\{z,g\}\)-square in the
upper family.  The pair is recorded only when \(B\) does not shatter it.
Thus every recorded terminal or collar pair really belongs to
\(\operatorname{Sh}(A_b)\setminus\operatorname{Sh}(B)\).

The following table gives all \(2268\) candidate comparisons, with the
harmless marker multiplicities retained.  We abbreviate repeated terminal,
terminal matching, distant corner, owner conflict, repeated collar, and
collar-augmented matching by
\(\mathrm{RT},\mathrm{TM},\mathrm{DC},\mathrm{OC},\mathrm{RC}\), and
\(\mathrm{CM}\), respectively:
\[
\begin{array}{c|rrrrrr}
 \text{opposite profile}
 &\mathrm{RT}&\mathrm{TM}&\mathrm{DC}&\mathrm{OC}&\mathrm{RC}&\mathrm{CM}
 \\ \hline
 \text{sharp}&16&3&0&0&0&0\\
 \eta=1\text{ one-unit}&172&42&2&0&0&0\\
 \text{feasible two-unit}&1188&592&118&95&24&16
\end{array}
\]
There is no further outcome.  Here is why the six columns suffice.  A
repeated acquired terminal or collar pair contradicts the two-layer
shadow identity
\[
 \operatorname{Sh}(A_0)\cap\operatorname{Sh}(A_1)
 =\operatorname{Sh}(B).
\]
A distant overlap corner contradicts
Lemma~\ref{lem:damp-passive-radius-one-shielding}.  At passive distance
one, equation~\eqref{eq:damp-passive-distance-one-forced-owner} forces
the correcting word into the exclusive set \(D_b=A_b\setminus B\).  An
owner conflict means either that this word already lies in \(B\), or that
the same word is forced into both \(D_0\) and \(D_1\); both alternatives
are impossible because \(D_0,D_1,B\) are pairwise disjoint.  The two
matching columns give the desired three acquired pairs directly, before
or after the acquired collars are added.

The verifier is
\path{verify_damp_pure_eta_two_mixed_profiles.py}, and its complete
output is \path{damp_pure_eta_two_mixed_profiles.json}.  The mathematical
digest is
\[
 \texttt{f6fed025e0c67f1789358d287e116bc8d46cfbdca691c78b340a446cfb312ffd}.
\]
\end{proof}

\begin{proposition}[A two-unit wing resolves every opposite profile through
two units]
\label{prop:damp-two-unit-wing-resolves-all-opposite-profiles}
Let
\[
 H=(A_0\times\{0\})\cup(A_1\times\{1\})
\]
be cornerless and ample, put \(B=A_0\cap A_1\), and suppose that
\(|B|\le9\).  Assume that \(B\subsetneq A_0\) has two-cycle holonomy
and excess two, while \(B\subsetneq A_1\) has two-cycle holonomy and
excess at most two.  Then the acquired-support family of the two wings
contains three pairwise disjoint direction pairs.
\end{proposition}

\begin{proof}[Proof by the complete exact-profile audit]
The classification and the two local impossibility results leave five
feasible two-unit profiles.  Alternative~2 of
Corollary~\ref{cor:damp-two-unit-uncounted-profiles-acquire-pair} is
impossible, as is alternative~4 by
Corollary~\ref{cor:damp-two-unit-vertex-free-passive-carriers}.  In each
of alternatives 1, 3, and 5--7, the active slice is a singleton and the
exact excess identity determines the entire counted overlap from its
rooted active gates.

This gives a short exhaustive construction of the overlap universe.  The
pure profile is obtained by adjoining two arbitrary vertices to its
\(P_7\) spine.  In alternative~3, one arbitrary vertex is adjoined to
the wedge of a rooted three-vertex path and a rooted ample gate of order
four.  Alternatives~5 and~6 are respectively the wedges of rooted ample
gates of orders three and five, and of two rooted ample gates of order
four.  Alternative~7 is the wedge of rooted paths of orders three, three,
and two.  Directly filtering all rooted graphs of orders four and five by
partial-cube isometry and the Sandwich equality gives respectively
\[
 5\qquad\hbox{and}\qquad13
\]
rooted ample gate types.  Consequently this construction is exhaustive,
without a catalogue of arbitrary nine-vertex graphs.

The accompanying exact-profile verifier constructs these normal forms and
retains \(148\) distinct embedded ample overlaps.  It
then rediscovers \(960\) two-unit markers and every possible opposite
sharp, uncounted one-unit, two-unit, or pair-rich extra-owner marker.
The root/non-root orientation is checked explicitly: every cycle shore
contains a lower root and a non-root owner, and a concept common to two
cycle shores is a non-root in both.  For an extra-owner marker, an edge
inside its three-concept shore supplies the second acquired pair.  If no
such edge exists, the marker is pair-poor; the audit finds three such
decompositions, all covered by the rooted-claw contradiction in
Proposition~\ref{prop:damp-two-unit-extra-owner-comparison}.

Opposite wings which use the same marker repeat both acquired terminal
supports and are immediately impossible by shadow intersection.  For
distinct markers, the audit makes \(8395\) comparisons.  With the
abbreviations from the preceding proof, their outcomes are
\[
\begin{array}{c|rrrrrr}
 \text{opposite row}
 &\mathrm{RT}&\mathrm{TM}&\mathrm{DC}&\mathrm{OC}&\mathrm{RC}&\mathrm{CM}
 \\ \hline
 \text{two-unit}&3131&1930&117&347&112&238\\
 \text{sharp}&121&49&0&0&0&7\\
 \eta=1\text{ one-unit}&791&260&12&66&0&0\\
 \text{pair-rich extra-owner}&718&334&0&94&0&68
\end{array}
\]
and there is no unresolved comparison.

Each column has the proof meaning established above.  Repeated acquired
terminal or collar supports contradict
\(\operatorname{Sh}(A_0)\cap\operatorname{Sh}(A_1)
=\operatorname{Sh}(B)\).  A distant corner contradicts passive-radius-one
shielding.  An owner conflict forces a word into \(B\), or into both
disjoint exclusive sets.  The two matching columns give three independent
acquired pairs directly.  Thus every possible opposite profile has the
conclusion of the proposition.

The verifier source and its complete output are, respectively,
\begin{center}
\path{verify_damp_two_unit_all_opposite_profiles.py},\\
\path{damp_two_unit_all_opposite_profiles.json}.
\end{center}
The mathematical digest is
\[
 \texttt{272d55715bd51f23c8956fb0578ccc84c804eb5a54dd4995ca0d9e9484de1b81}.
\]
\end{proof}

\begin{corollary}[The central-square profile cannot be shielded]
\label{cor:damp-one-unit-central-square-unshielded}
Under the hypotheses of
Proposition~\ref{prop:damp-one-unit-opposite-wing-central-square},
suppose in addition that \(H\) is cornerless.  Then the acquired-support
family contains three pairwise disjoint direction pairs.  In particular,
the central-square \(K_{2,2}\) alternative of that proposition cannot
occur in a cornerless family.
\end{corollary}

\begin{proof}
Suppose that the exceptional alternative occurs, and fix one of the two
wings.  In its \(\eta=1\) normal form, the active slice \(L\) is a
singleton, say \(\{c\}\), and the uncounted vertex \(v\) is the vertex
of the central square opposite \(c\).  The two directions \(z_1,z_2\)
of that square are the two passive attachment directions of the wing.
All relative concepts \(D=A_b\setminus B\) have one common passive word,
namely the passive word of \(c\).  Hence
\[
 v|_{\{z_1,z_2\}}
 =(c^{\oplus\{z_1,z_2\}})|_{\{z_1,z_2\}},
\]
so \(v\) is at passive distance two from this word.  The two pendant
paths of the exceptional graph are attached at the other
antipodal pair of vertices of the central square.  Consequently \(v\) is
a corner of \(B\), contradicting
Lemma~\ref{lem:damp-passive-radius-one-shielding}.  The exceptional
alternative is impossible.
\end{proof}

\begin{lemma}[Forced wing words certify active coordinates]
\label{lem:damp-forced-wing-words-active-coordinates}
Let
\[
 H=(A_0\times\{0\})\cup(A_1\times\{1\})
\]
be cornerless and ample, put \(B=A_0\cap A_1\), and fix one wing
\(D_b=A_b\setminus B\).  Suppose that all members of \(D_b\) have one
common word \(\sigma\) on a passive coordinate set \(Z\).  Let
\(\mathcal S\subseteq D_b\) be any collection of shore concepts whose
opposite \(Z\)-fibres occur in \(B\).  For every corner \(v\) of \(B\)
whose \(Z\)-word differs from \(\sigma\) in exactly one coordinate \(z\),
put \(o_v=v^{\oplus z}\), and set
\[
 W_b=\mathcal S\cup\{o_v:v\in\operatorname{Cor}(B),
       d_{Q_Z}(v|_Z,\sigma)=1\}.
\]
Then \(W_b\subseteq D_b\).  If the graph induced by \(D_b\) is connected,
every coordinate which is nonconstant on \(W_b\) is active on that graph.
Consequently, if the induced graph is two-connected, then
\begin{equation}
 |D_b|\ge
 2\bigl|\{x:|\{w|_x:w\in W_b\}|=2\}\bigr|.
 \label{eq:damp-forced-wing-active-coordinate-cost}
\end{equation}
\end{lemma}

\begin{proof}
The shore concepts belong to \(D_b\) by definition.  At passive distance
one, Lemma~\ref{lem:damp-passive-radius-one-shielding} says that the unique
concept which can shield the corner \(v\) is \(v^{\oplus z}\), and
cornerlessness forces that concept into \(D_b\).  Thus \(W_b\subseteq D_b\).

If two members of \(W_b\) differ in coordinate \(x\), a path between them
inside the connected graph induced by \(D_b\) must use an \(x\)-edge.
Hence \(x\) is active.  Applying
Lemma~\ref{lem:damp-two-connected-coordinate-cost} gives
\eqref{eq:damp-forced-wing-active-coordinate-cost}.
\end{proof}

\begin{proposition}[The terminal two-connected matching row closes]
\label{prop:damp-terminal-two-connected-global-completion}
Let
\[
 H=(A_0\times\{0\})\cup(A_1\times\{1\})
\]
be cornerless and ample, put \(B=A_0\cap A_1\), and suppose that
\(|B|\le9\).  For \(b\in\{0,1\}\), suppose that the relative pair
\(B\subsetneq A_b\) is terminal under both one-concept addition and
deletion, that the graph induced by \(D_b=A_b\setminus B\) is
two-connected, and that
\[
 |D_b|\in\{10,11\}.
\]
Then the two lower gate-holonomy cycles cannot both have excess at most
two.  In particular, no pair of terminal two-connected wings of orders ten
or eleven occurs in the below-forty two-wing reduction.
\end{proposition}

\begin{proof}[Computer-assisted proof]
By Corollary~\ref{cor:damp-small-lower-family-two-arm-holonomy}, both
holonomy cycles have length two.  The local comparison is already complete.
If one wing has excess two,
Proposition~\ref{prop:damp-two-unit-wing-resolves-all-opposite-profiles}
gives either an immediate contradiction or three independent acquired
pairs.  If both excesses are at most one,
Proposition~\ref{prop:damp-one-unit-opposite-wing-central-square} gives the
same alternative or the central-square profile, and
Corollary~\ref{cor:damp-one-unit-central-square-unshielded} excludes the
latter.  Thus it remains only to resolve the three-pair outcomes.

For each normal-form marker, form \(W_b\) from its actual shore concepts and
the owners forced by every overlap corner at passive distance one.  Corners
at larger passive distance are already impossible.  The complete two-unit
normal-form universe of
Proposition~\ref{prop:damp-two-unit-wing-resolves-all-opposite-profiles}
contains \(2886\) matching comparisons.  Passive distance and owner
collisions exclude \(302\) of them.  Among the remaining \(2584\), the
ordered numbers of coordinates varying on \((W_0,W_1)\) are
\[
\begin{array}{c|rrr}
 (|Y(W_0)|,|Y(W_1)|)&(5,5)&(6,5)&(6,6)\\ \hline
 \text{comparisons}&504&108&1972.
\end{array}
\]
Lemma~\ref{lem:damp-forced-wing-words-active-coordinates} excludes the last
two columns, because a two-connected wing using six active coordinates has
at least twelve concepts.

The low-excess normal forms give the same dichotomy.  The \(76\) ample
one-vertex completions contain four distinct \(\eta=1\) versus
\(\eta=1\) matchings and eight distinct \(\eta=1\) versus extra-owner
matchings.  Every positioned wing forces six active coordinates, except
when the overlap is the eight-vertex path.  The two extra-owner path
decompositions behave identically: \(P_9\) forces six coordinates, while
\(P_8\) is the sole five-coordinate residue.

It remains to check a fixed five-cube, not an ambient order interval.  Up
to cube automorphisms, the forced subsets of a wing are
\[
 \{0,1,2,3,12,28\},\qquad
 \{0,1,2,7,15,31\},\qquad
 \{0,1,2,13,29\}\subseteq Q_5.
\]
They are, respectively, a square and a disjoint edge, two disjoint
three-vertex paths, and the five-vertex \(P_8\) equality pattern.  Retaining
their positions in \(B\) gives fourteen two-unit side states and six
low-excess side states.  For each state, the finite check adjoins every
subset of the remaining \(Q_5\)-vertices which gives a relative layer of
order ten or eleven, and then tests induced two-connectivity, ampleness of
\(A_b=B\cup D_b\), and both terminality conditions.  The results are
\[
\begin{array}{c|r|r|r|r}
 &\multicolumn{2}{c|}{|D_b|=10}
 &\multicolumn{2}{c}{|D_b|=11}\\
 \text{source}&\text{two-connected}&\text{ample upper}
 &\text{two-connected}&\text{ample upper}\\ \hline
 \text{two-unit residue}&504&238&3346&1330\\
 \text{low-excess }P_8&744&198&4176&1026
\end{array}
\]
Every ample upper family in the table has both an ample one-concept
extension of \(B\) and a deletable concept in \(D_b\).  Hence none is
terminal.

All tests use the exact cube adjacency relation and the Sandwich equality
\(|\operatorname{Sh}(F)|=|F|\); no list of ambient families is used.  The
source files and compact certificates are
\begin{center}
\path{explore_damp_three_pair_global_compatibility.py}\\
\path{damp_three_pair_global_compatibility.json}\\
\path{verify_damp_low_excess_global_matching.py}\\
\path{damp_low_excess_global_matching.json}\\
\path{explore_damp_five_coordinate_matching_completions.py}\\
\path{damp_five_coordinate_matching_completions.json}.
\end{center}
Their mathematical digests are, respectively,
\[
\begin{gathered}
 \texttt{ff2d10fdf95a608da5e7cdc0dca8cc9d973c4cfc345336f3b7278e40df8e3f92},\\
 \texttt{760ff9566860cefa41b8a300dd9f99b03ecd08a79701a89803027f954e5afe08},\\
 \texttt{fc485e9c7e6550c8478248effa90bf60b84719d42ea1579e2399b3d3f3cb5b10}.
\end{gathered}
\]
This exhausts every matching outcome and proves the proposition.
\end{proof}

For a two-connected terminal layer of order ten or eleven,
Lemma~\ref{lem:damp-two-connected-coordinate-cost} supplies the five-active-
coordinate hypothesis automatically.  Corollary
\ref{cor:damp-terminal-fibre-two-root-cycles} therefore replaces all
sixty-two residual graph types by two coupled directed cycles of roots and
their carrier owners.  Corollary
\ref{cor:damp-two-sided-passive-gate-holonomy} further replaces those two
cycles by one two-sided holonomy system on ample intermediate gates.
Corollary~\ref{cor:damp-gate-holonomy-exponential-cost} adds a quantitative
dichotomy: every shore overlap or additional owner pays a positive vertex
cost, while zero excess is a system of pendant three-edge square-gate
arms.  At one unit, the shared-owner alternative is impossible, the
extra-owner comparisons give three independent pairs or a contradiction,
and the uncounted vertex is the apex of a connected tree of squares.
Comparing the two wings first leaves two native support profiles: three
independent acquired pairs, or the single central-square graph of
Proposition~\ref{prop:damp-one-unit-opposite-wing-central-square}, whose
four supports form \(K_{2,2}\).  The latter profile is now excluded by
Corollary~\ref{cor:damp-one-unit-central-square-unshielded}: for either
wing, its uncounted central-square vertex is a corner of \(B\) at passive
distance two from every exclusive concept of that wing, so it cannot be
shielded in \(H\).

The remaining two-connected step is therefore no longer an open taxonomy
of one- or two-unit anomalies.  Every one-unit comparison gives three
independent acquired pairs.  The exact two-unit identity gives seven
formal profiles; the shared-owner uncounted profile and the overlapping
vertex-free profile are impossible, while each of the other five supplies
two disjoint acquired pairs.

The three alternating support graphs \(C_4,P_5,2P_3\) were first merged
by Corollary~\ref{cor:damp-alternating-arm-incidence-carriers} into the
incidence shadows \(K_{2,2},P_4,2K_2\) of one signed carrier problem.
Passive-radius shielding, forced distance-one owners, and acquired collar
squares now resolve that problem completely.
Proposition~\ref{prop:damp-two-unit-wing-resolves-all-opposite-profiles}
excludes every alternating comparison of a two-unit wing with a sharp,
one-unit, or two-unit opposite wing.  Its audit includes the formerly
separate neither-wing-pure and pair-rich extra-owner rows; the pair-poor
extra-owner row is already impossible by the rooted-claw comparison.
Thus every local two-connected profile through two units yields three
independent acquired pairs.

Lemma~\ref{lem:damp-forced-wing-words-active-coordinates} turns those pairs
into positioned residual concepts.  Proposition
\ref{prop:damp-terminal-two-connected-global-completion} then closes the
global matching row: six varying coordinates contradict the open-ear bound,
and the unique five-coordinate residue has no terminal completion.  The
octahedral support sphere of
Remark~\ref{rem:damp-owner-cycle-sphere-target} remains the geometry of the
sharp twelve-concept example, but it is no longer a proof obligation for
the orders-ten-and-eleven frontier.  Proposition
\ref{prop:damp-terminal-cut-layer-through-eleven} also closes every
terminal cut wing in that range.  Thus all terminal wings through order
eleven are two-connected and covered by the preceding comparison.  The
only remaining below-forty compatibility issue is the marked case of
Lemma~\ref{lem:damp-small-trapped-wing-one-mark-terminalization}: one
order-eleven wing may become a terminal order-ten core only after a single
ample concept has been moved into the common overlap.

\begin{proposition}[Nine-concept passive-carrier exclusion]
\label{prop:damp-nine-concept-passive-carrier-exclusion}
There is no trapped pair

\[
 \varnothing\ne K\subsetneq A,
 \qquad A,K\in\Damp,
 \qquad \operatorname{Cor}(A)\subseteq K,
 \qquad |A\setminus K|=9.
\]

Consequently,

\begin{equation}
 10\le\rho_{\Damp}\le12.
 \label{eq:damp-relative-corner-threshold-ten-twelve}
\end{equation}
\end{proposition}

\begin{proof}[Computer-assisted proof]
The arbitrary-upper relative linear-quotient theorem through eight facets
\cite{BousquetYangComplex26} already gives \(\rho_{\Damp}\ge9\).  Suppose
that the displayed pair exists.  Choose \(C\in\Damp\) inclusion-maximal
subject to \(K\subseteq C\subsetneq A\).  The pair \(C\subsetneq A\) is
still trapped, so the lower bound nine forces \(C=K\).  Hence
\(K\cup\{c\}\) is not ample for every \(c\in R:=A\setminus K\): if it
were, one could peel \(c\) first and contradict maximality.  Also
\(A\setminus\{c\}\) is not ample, because \(c\) is not a corner of \(A\).
Thus every member of \(R\) is doubly blocked.  Proposition
\ref{prop:damp-terminal-block-curvature-tax} gives minimum residual degree
two.

The graph induced by \(R\) is connected by
Proposition~\ref{prop:damp-terminal-relative-component-size}.

Delete the coordinates constant on \(R\).  An exact cut-system
classification tests all \(160\) connected bipartite graph types of order
nine and minimum degree two.  Ten admit an induced-cube embedding, and
quotienting their coordinate cuts by graph automorphisms gives thirteen
coordinate-coloured profiles.  Applying the degree-two marker and
cross-carrier rules of Lemma~\ref{lem:damp-local-passive-carrier-rules}
excludes ten of them, leaving only the following active residual words:

\begin{center}
\begin{tabular}{c@{\qquad}l}
\toprule
profile & active words in \(Q_4\)\\
\midrule
5.1 & \(0,2,4,6,7,8,10,11,15\)\\
8.1 & \(0,2,4,6,8,10,11,14,15\)\\
9.1 & \(0,2,3,4,6,7,8,10,11\)\\
\bottomrule
\end{tabular}
\end{center}

\begin{sloppypar}
The classification and first reduction are reproduced by
\path{enumerate_damp_nine_relative_graphs.py},
\path{enumerate_damp_nine_colored_embeddings.py}, and
\path{verify_damp_nine_degree_two_carrier_frontier.py}.  The last verifier
keeps passive support parts symbolically through their inclusion preorder;
its digest is

\[
 \texttt{e5ea4fb30cb6611290a6555182a2af36dbd1339598d921ba1b04f37356b114ea}.
\]
\end{sloppypar}

It remains to exclude the three displayed profiles without conditioning
away passive coordinates.  Since \(A,K\) are ample, the relative support
set \(\mathcal H\) has exactly nine members.  For each member, the final
formula records its four-bit active part, its passive part, its position in
the support order, and residual anchors for its lower and upper missing
traces.  It imposes support distinctness, carrier depth at least two at
every residual vertex, the exact full-cube conclusion for every compatible
minimal--maximal carrier pair, the degree-two isolated markers, and the
degree-three dichotomy
\eqref{eq:damp-degree-three-isolated-or-cover}.  These are only necessary
conditions: in particular, the formula omits the adjacent-marker reversal
law, so it is a relaxation of every genuine terminal interval.

Nine formal passive coordinates suffice.  If \(P_1,\ldots,P_9\) are the
actual passive parts, replace \(P_i\) by the down-set

\[
 \widehat P_i=\{j:P_j\subseteq P_i\}.
\]

Then \(\widehat P_i\subseteq\widehat P_j\) if and only if
\(P_i\subseteq P_j\).  Thus this replacement faithfully preserves every
passive inclusion and equality used in the formula.

The three structural formulas are unsatisfiable.  CaDiCaL generated core
proofs that were checked independently by \texttt{drat-trim}; the proof
statistics are

\begin{center}
\begin{tabular}{crrr}
\toprule
profile & variables & clauses & resolution steps\\
\midrule
5.1 & \(6{,}366\) & \(71{,}937\) & \(9{,}796{,}884\)\\
8.1 & \(6{,}615\) & \(76{,}490\) & \(6{,}854{,}743\)\\
9.1 & \(6{,}627\) & \(74{,}647\) & \(7{,}636{,}141\)\\
\bottomrule
\end{tabular}
\end{center}

All three checks use zero RAT lemmas.  The generator, DIMACS manifest,
proofs, and aggregate checker are

\[
\begin{gathered}
 \texttt{solve\_damp\_nine\_passive\_carrier\_cnf.py},\\
 \texttt{damp\_nine\_passive\_carrier\_cnf\_manifest.json},\\
 \texttt{damp\_nine\_passive\_carrier\_*\_core.drat},\\
 \texttt{verify\_damp\_nine\_passive\_carrier\_proofs.py}.
\end{gathered}
\]

The aggregate mathematical digest is

\[
 \texttt{54b27131c083db9074d7c79ffcb303de8a8b36fe6d0c5de2ab6edd3c4f937d8b}.
\]

This excludes the last three profiles, proving that no trapped difference
of order nine exists.  The upper bound in
\eqref{eq:damp-relative-corner-threshold-ten-twelve} is Proposition
\ref{prop:damp-twelve-concept-trapped-interval}.
\end{proof}

\begin{lemma}[Every sub-ten dismantlable-ample interval has a first concept]
\label{lem:damp-subten-interval-first-concept}
Let \(C\subsetneq H\) belong to \(\Damp\), and suppose that
\[
 |H\setminus C|\le9.
\]
Then some \(h\in H\setminus C\) satisfies \(C\cup\{h\}\in\Damp\).
\end{lemma}

\begin{proof}
Suppose not, and choose \(J\in\Damp\) inclusion-minimal subject to
\(C\subsetneq J\subseteq H\).  No corner of \(J\) lies outside \(C\).
Indeed, deleting such a corner \(j\) leaves a member of \(\Damp\).  If
\(J\setminus\{j\}=C\), then \(C\cup\{j\}=J\), contrary to the
assumption; otherwise it contradicts the minimal choice of \(J\).
Consequently
\[
 \operatorname{Cor}(J)\subseteq C,
 \qquad
 1\le |J\setminus C|\le9.
\]
This is a trapped pair smaller than the lower bound
\(\rho_{\Damp}\ge10\) from
Proposition~\ref{prop:damp-nine-concept-passive-carrier-exclusion}, a
contradiction.
\end{proof}

\begin{lemma}[Small trapped wings have at most one marked synchronization]
\label{lem:damp-small-trapped-wing-one-mark-terminalization}
Let \(B\subsetneq A\) be ample, suppose that
\[
 \operatorname{Cor}(A)\subseteq B,
 \qquad |A\setminus B|\le11,
\]
and assume that every ample family \(C\) with
\(B\subseteq C\subseteq A\) belongs to \(\Damp\).  Then there is an ample
family \(C\) with
\[
 B\subseteq C\subsetneq A,
 \qquad |C\setminus B|\le1,
\]
such that, on putting \(R=A\setminus C\),
\begin{equation}
 |R|\in\{10,11\},
 \qquad
 C\cup\{r\}\text{ and }A\setminus\{r\}
 \text{ are nonample for every }r\in R.
 \label{eq:damp-one-mark-terminal-residual}
\end{equation}
Moreover, the graph induced by \(R\) is two-connected.  In particular, a
trapped wing of order ten is already terminal, whereas a trapped wing of
order eleven is either terminal or consists of a terminal ten-concept core
over one marked ample addition to \(B\).
\end{lemma}

\begin{proof}
Choose \(C\) inclusion-maximal among the proper ample intermediate families
between \(B\) and \(A\).  Such a family exists because \(B\subsetneq A\).
For every \(r\in R=A\setminus C\), maximality makes
\(C\cup\{r\}\) nonample.  If \(A\setminus\{r\}\) were ample, the ample
corner-deletion criterion would make \(r\) a corner of \(A\).  This is
impossible because \(r\notin C\supseteq B\supseteq\operatorname{Cor}(A)\).
Thus every member of \(R\) is doubly blocked.

The hypothesis puts both \(C\) and \(A\) in \(\Damp\), and
\(\operatorname{Cor}(A)\subseteq B\subseteq C\).  Hence \(C\subsetneq A\)
is a trapped pair in the definition of \(\rho_{\Damp}\).  Proposition
\ref{prop:damp-nine-concept-passive-carrier-exclusion} gives
\[
 10\le|R|\le|A\setminus B|\le11.
\]
It follows that
\(|C\setminus B|=|A\setminus B|-|R|\le1\).  Finally,
Proposition~\ref{prop:damp-terminal-cut-layer-through-eleven} makes the
residual graph two-connected.  The last assertion is the specialization to
the two possible initial orders.
\end{proof}

\begin{proposition}[A holonomy shore has a contracted split-square marker]
\label{prop:damp-holonomy-shore-split-square-marker}
Let \(K\subsetneq A\) belong to \(\Damp\), put \(D=A\setminus K\), and
suppose that
\[
 K\cup\{d\}\text{ is nonample for every }d\in D,
 \qquad |D|\le11.
\]
Let \(z\) be constant on \(D\), put
\[
 S_z=\{d\in D:d^{\oplus z}\in K\},
\]
and suppose that \(2\le|S_z|<|D|\).  After contracting \(z\), the relative
concepts of the nested pair \(\pi_zK\subsetneq\pi_zA\) are exactly the
images of \(D\setminus S_z\), and there is a concept
\(u\in D\setminus S_z\) whose image can be adjoined amply to \(\pi_zK\).
This root cannot lift to \(K\).  Consequently there is an active edge
\begin{equation}
 uu^{\oplus f}\in E(D),
 \qquad
 u\notin S_z,\quad u^{\oplus f}\in S_z,
 \label{eq:damp-holonomy-shore-crossing-edge}
\end{equation}
whose \(z,f\)-square is split as
\begin{equation}
 (u^{\oplus f})^{\oplus z}\in K,
 \qquad
 u^{\oplus z}\notin A.
 \label{eq:damp-holonomy-shore-split-square}
\end{equation}
\end{proposition}

\begin{proof}
All members of \(D\) have the same \(z\)-value.  Their projections remain
distinct.  A projected member already belongs to \(\pi_zK\) exactly when
its original concept lies in \(S_z\).  Hence
\[
 \pi_zA\setminus\pi_zK=(D\setminus S_z)|_{X\setminus\{z\}}.
\]
Both contracted endpoints belong to \(\Damp\) by pc-minor closure, and
\[
 1\le|D\setminus S_z|\le11-2=9.
\]
Lemma~\ref{lem:damp-subten-interval-first-concept} supplies the asserted
contracted root; denote its original lift by \(u\).

The extension \(K\cup\{u\}\) is nonample by hypothesis.  Apply
Proposition~\ref{prop:damp-contracted-teaching-root-wrong-side} with
contracted coordinate \(z\), lower family \(K\), and upper family \(A\).
The other member \(u^{\oplus z}\) of the \(z\)-fibre does not belong to
\(A\): it cannot belong to \(D\), whose \(z\)-value is constant, and if it
belonged to \(A\) it would belong to \(K\), contrary to
\(u\notin S_z\).  Thus \(u\) is the unique relative member of its fibre,
as required by that proposition.  Its wrong-side conclusion supplies an
\(f\) for which
\[
 u^{\oplus f}\in D,\qquad
 u^{\oplus\{z,f\}}\in K.
\]
The second membership says exactly that \(u^{\oplus f}\in S_z\), and
also proves~\eqref{eq:damp-holonomy-shore-split-square}; this gives
\eqref{eq:damp-holonomy-shore-crossing-edge}.
\end{proof}

\begin{corollary}[Lower gate holonomy is carried by split squares]
\label{cor:damp-lower-gate-holonomy-split-squares}
Assume in addition the hypotheses of
Corollary~\ref{cor:damp-two-sided-passive-gate-holonomy}, with
\(K,A\in\Damp\) and \(|D|\le11\).  Every shore on a directed cycle of
the map \(\lambda\) has at least two concepts and carries a split-square
marker of the form
\eqref{eq:damp-holonomy-shore-crossing-edge}--\eqref{eq:damp-holonomy-shore-split-square}.
Thus the remaining two-connected obstruction is a cyclic system of
oriented split squares, not a cyclic system of arbitrary boundary subsets.
\end{corollary}

\begin{proof}
Let \(S\) lie on a \(\lambda\)-cycle and let \(T\) be its predecessor on
that cycle.  By construction, \(S=\lambda(T)\) blocks the lower root
\(r_T\).  The shore-root lemma gives a member of \(R\cap S\), while the
common-owner theorem gives a member of \(S\setminus R\).  These two
concepts are distinct, so \(|S|\ge2\).  Also \(S\subsetneq D\), by the
definition of \(\mathscr S\).  Proposition
\ref{prop:damp-holonomy-shore-split-square-marker} now supplies the
claimed marker.
\end{proof}

\begin{remark}[Why the relative threshold does not give the next absolute bound]
\label{rem:damp-relative-threshold-route-limit}
The present calibration is
\[
 10\le\rho_{\Damp}\le12.
\]
In particular, the proposed estimate $\rho_{\Damp}\ge22$ is false.
Proposition~\ref{prop:damp-two-wing-relative-corner-threshold} remains a
valid reduction, but by itself gives only order at least \(22\).  Its
coordinatewise form is stronger: it now excludes order \(34\) without an
ambient-family census.  Reaching the upper calibration
\(\rho_{\Damp}=12\) would make the same first-moment argument exclude every
order below \(40\).  The alternative bulk route is the uniform path-core
repair theorem of
Corollary~\ref{cor:damp-path-only-reduction-below-forty-eight}.
\end{remark}

The next proposition retains more of this compatibility.  Rather than
recording only the number of concepts in a trapped interval, it extracts a
terminal block and charges separately for its cycles and for the cube corners
that must be attached to the lower family.

\begin{proposition}[Terminal block-curvature tax]
\label{prop:damp-terminal-block-curvature-tax}
Let
\[
 \varnothing\ne B\subsetneq A,
 \qquad A,B\in\Damp,
 \qquad \operatorname{Cor}(A)\subseteq B.
\]
Choose \(C\in\Damp\) inclusion-maximal subject to
\(B\subseteq C\subsetneq A\), and put \(D=A\setminus C\).  Write
\(e(D)\) for the number of edges induced by \(D\), and define
\begin{equation}
 \tau(D)=e(D)-|D|+
   \max\{1,|\operatorname{Cor}(D)|\}.
 \label{eq:damp-terminal-block-curvature}
\end{equation}
Then the graph induced by \(D\) has minimum degree at least two, every
corner of the cube family \(D\) has a neighbor in \(C\), and the relative
Yang module for \((A,C)\) is deletion-minimal Cohen--Macaulay: it is
Cohen--Macaulay, whereas the module for
\((A\setminus\{d\},C)\) is not Cohen--Macaulay for any \(d\in D\).
Moreover,
\begin{equation}
 \sum_{S\in\operatorname{Sh}(A)\setminus\operatorname{Sh}(B)}|S|
 \ge |A\setminus B|+\tau(D).
 \label{eq:damp-terminal-block-rank-tax}
\end{equation}
In particular, \(\tau(D)\ge1\).
\end{proposition}

\begin{proof}
No \(d\in D\) satisfies the \(C\)-teaching condition.  Indeed, that
condition says that the one-concept extension \(C\cup\{d\}\) is ample.
It then belongs to \(\Damp\): first peel \(d\), and then peel \(C\).  This
contradicts the maximal choice of \(C\).

Because \(A\) and \(C\) are ample, their relative Yang module is
Cohen--Macaulay.  On the other hand, \(A\setminus\{d\}\) is not ample for
any \(d\in D\): otherwise \(d\) would be a corner of the ample family
\(A\), contrary to
\(\operatorname{Cor}(A)\subseteq C\).  The relative Cohen--Macaulay
criterion over the ample lower family \(C\) proves the deletion-minimal
assertion.

Suppose first that \(d\) is isolated in the graph induced by \(D\).  The
ample family \(A\) is isometric.  For every \(c\in C\), the first edge of
an isometric path from \(d\) to \(c\) therefore joins \(d\) to \(C\) in a
coordinate on which \(d\) and \(c\) differ.  Thus the \(C\)-neighbors of
\(d\) meet \(\operatorname{Diff}(d,c)\) for every \(c\in C\), which is
exactly the \(C\)-teaching condition, a contradiction.

Suppose next that \(d\) has degree one in the graph induced by \(D\).
The relative leaf-deletion dichotomy
\cite{BousquetYangComplex26} says that either \(d\) satisfies the
\(C\)-teaching condition, already excluded, or the relative module for
\((A\setminus\{d\},C)\) is Cohen--Macaulay.  Since \(C\) is ample, the
relative Cohen--Macaulay criterion makes \(A\setminus\{d\}\) ample.
Hence \(d\) is a corner of \(A\), contrary to
\(d\notin C\supseteq\operatorname{Cor}(A)\).  This proves the minimum-degree
assertion and, consequently, \(e(D)\ge |D|\).

Let \(d\) be a corner of \(D\), and let \(Q\) be the unique maximal cube of
\(D\) containing it.  If \(d\) had no neighbor in \(C\), then every edge of
\(A\) at \(d\) would already be an edge of \(D\), and hence would lie in
\(Q\).  Every cube of \(A\) through \(d\) would consequently lie in \(Q\),
making \(d\) a corner of \(A\).  This is again impossible.  Distinct
corners of \(D\) give distinct boundary edges because their endpoints in
\(D\) are distinct.  Since the one-inclusion graph of \(A\) is connected,
there is also at least one edge between the two nonempty sets \(C\) and
\(D\).  If \(e_\partial(C,D)\) denotes the number of these boundary edges,
then
\begin{equation}
 e_\partial(C,D)\ge
 \max\{1,|\operatorname{Cor}(D)|\}.
 \label{eq:damp-terminal-block-boundary-corners}
\end{equation}

For every ample cube family \(X\), the edge--shadow identity gives
\[
 |E(X)|=\sum_{S\in\operatorname{Sh}(X)}|S|.
\]
Therefore
\begin{align*}
 \sum_{S\in\operatorname{Sh}(A)\setminus\operatorname{Sh}(C)}|S|
 &=|E(A)|-|E(C)|\\
 &=e(D)+e_\partial(C,D)\\
 &\ge |D|+\tau(D).
\end{align*}
Finally, every support in
\(\operatorname{Sh}(C)\setminus\operatorname{Sh}(B)\) is nonempty,
because \(B\ne\varnothing\), and ampleness gives
\[
 |\operatorname{Sh}(C)\setminus\operatorname{Sh}(B)|=|C|-|B|.
\]
Adding these earlier acquired supports proves
\eqref{eq:damp-terminal-block-rank-tax}.
\end{proof}

\begin{proposition}[Wrong-side witnesses localize on relative edges]
\label{prop:damp-contracted-teaching-root-wrong-side}
Let \(C\subsetneq A\subseteq Q_E\), suppose that \(A\) is isometrically
embedded, and fix a coordinate \(e\in E\).  Put
\[
 D=A\setminus C,
 \qquad \overline A=\pi_eA,
 \qquad \overline C=\pi_eC.
\]
For a cube family \(K\) and a vertex \(x\), write
\(N_K(x)=\{f:x^{\{f\}}\in K\}\).  Let \(d\in D\), write
\(y=d|_{E\setminus\{e\}}\), and suppose that
\(y\in\overline A\setminus\overline C\) and that \(d\) is the unique member
of \(D\) in its \(e\)-fibre.  If \(y\) satisfies the
\(\overline C\)-teaching condition but \(d\) does not satisfy the
\(C\)-teaching condition, then there is \(c\in C\) such that, for every
\[
 f\in N_{\overline C}(y)
       \cap\operatorname{Diff}(y,c|_{E\setminus\{e\}}),
\]
one has
\begin{equation}
 d^{\{f\}}\in D,
 \qquad
 d^{\{e,f\}}\in C.
 \label{eq:damp-contracted-root-wrong-side-square}
\end{equation}
In particular, every witness supplied by the contracted teaching root but
lost on lifting is the direction of an edge induced by \(D\); the opposite
vertex of the corresponding \(e,f\)-square lies in \(C\).
\end{proposition}

\begin{proof}
The hypotheses on the \(e\)-fibre imply that \(d^{\{e\}}\notin A\).
Indeed, the other lift cannot belong to \(C\), since
\(y\notin\overline C\), and it cannot belong to \(D\) by uniqueness.

Because \(d\) fails the \(C\)-teaching condition, choose \(c\in C\) such that
\[
 N_C(d)\cap\operatorname{Diff}(d,c)=\varnothing.
\]
The contracted root \(y\) teaches \(\overline C\), so the set displayed in
the statement is nonempty.  Fix one of its members \(f\).  Since
\(y^{\{f\}}\in\overline C\), one of the two lifts
\(d^{\{f\}}\) and \(d^{\{e,f\}}\) belongs to \(C\).  Moreover,
\(f\in\operatorname{Diff}(d,c)\).  The choice of \(c\) therefore excludes
\(d^{\{f\}}\in C\), and hence
\(d^{\{e,f\}}\in C\).

The two vertices \(d\) and \(d^{\{e,f\}}\) are at Hamming distance two.
Isometry of \(A\) supplies a two-edge path between them in \(A\).  Its
intermediate vertex is either \(d^{\{e\}}\) or \(d^{\{f\}}\).  The first is
not in \(A\), as observed above, so \(d^{\{f\}}\in A\).  It is not in \(C\),
and consequently it belongs to \(D\).  This proves
\eqref{eq:damp-contracted-root-wrong-side-square} for every such \(f\).
\end{proof}

\begin{remark}[The first two cyclic block shapes]
\label{rem:damp-first-terminal-block-curvatures}
The sole nine-concept relative graph not excluded by the current
linear-quotient arguments is, up to signed coordinate permutation,
\[
 D_9=\{0,1,2,3,32,48,52,56,60\}\subseteq Q_6.
\]
It consists of two squares joined by a path of two edges.  Thus
\(e(D_9)=10\), and its six nonattachment square vertices are precisely its
cube-family corners.  If \(D_9\) occurs as a terminal block, then
\[
 \tau(D_9)=10-9+6=7,
 \qquad
 \sum_{S\in\operatorname{Sh}(A)\setminus\operatorname{Sh}(B)}|S|
 \ge16.
\]
For comparison, the twelve-concept terminal block in
Proposition~\ref{prop:damp-twelve-concept-trapped-interval} consists of
three squares coupled cyclically.  It has fifteen induced edges and six
corners, so its curvature tax is \(9\); its twelve acquired supports have
total rank \(36\), well above the lower bound \(12+9=21\).

The gap in the last comparison is informative.  Cycle rank and the number
of boundary attachments are necessary data, but they still forget which
square owns each boundary direction.  Excluding \(D_9\), or realizing it,
therefore requires a compatibility theorem for these side-coloured
boundary owners rather than a further uncoloured moment inequality.
Proposition~\ref{prop:damp-contracted-teaching-root-wrong-side} makes this
requirement local.  At the level of the relative graph, contracting either
edge of the two-edge bridge identifies its endpoints and gives the
eight-vertex barbell.  The contracted relative pair may have fewer facets
when another projected vertex already belongs to the contracted lower
class, but it is nonempty and has at most eight facets.  The relative
linear-quotient theorem therefore supplies a contracted teaching root.
Away from the contracted bridge fibre, such a root can fail to lift only
by transferring every available witness along an internal edge of \(D_9\)
and completing the corresponding square into the lower class.  Thus the
remaining obstruction is a signed transfer along the block tree, not an
unrestricted choice in the ambient cube.
\end{remark}

\section{Exact forbidden pc-minors}
\label{sec:damp}

This section proves the exact basis criterion.  The algebraic dictionary
translates corner peelings into linear quotients and pc-minors into
color-respecting specializations.  The Yang-complex dictionary supplies
the equivalent topological form.  Pc-minor closure then reduces all
proper pc-minors to the three elementary maps in each coordinate.

\subsection{Alexander-dual formulation}

The same problem has a concise algebraic form.  This formulation is useful
because it preserves both structures that matter here: a corner peeling is
an ordering property of monomial generators, and a pc-minor is a
color-respecting specialization rather than an arbitrary ideal minor.

Fix a field \(k\).  For an irredundant cube family
\(X\subseteq\{0,1\}^E\), put
\[
 S_E=k[v_{e,0},v_{e,1}:e\in E].
\]
The \emph{Yang complex} of \(X\) is the pure simplicial complex
\[
 \Delta_X
 =\left\langle
    F_x:=\{v_{e,x_e}:e\in E\}:x\in X
  \right\rangle.
\]
Its Alexander-dual ideal is the transversal squarefree ideal
\begin{equation}
J_X
 :=\left(m_x:x\in X\right),
 \qquad
 m_x:=\prod_{e\in E}v_{e,1-x_e}.
 \label{eq:damp-yang-dual-ideal}
\end{equation}
Thus every minimal generator chooses exactly one variable from each color
pair \(\{v_{e,0},v_{e,1}\}\).

Recall that an equigenerated monomial ideal
\(J=(u_1,\ldots,u_t)\) has \emph{linear quotients} in the displayed order
if
\[
 (u_1,\ldots,u_{j-1}):u_j
\]
is generated by variables for every \(j>1\).  It has linear quotients if
some ordering of its minimal generators has this property.

\begin{proposition}[Algebraic pc-minor dictionary]
\label{prop:damp-algebraic-minor-dictionary}
Let \(X\subseteq\{0,1\}^E\) be isometric in its given cube embedding.
Then the following statements hold.
\begin{enumerate}
 \item \(X\in\Damp\) if and only if \(J_X\) has linear quotients.
 \item Let \(F,D\subseteq E\) be disjoint, let
       \(f\in\{0,1\}^F\) be realized by \(X\), and set
       \[
        K=\pi_D(X_f),
        \qquad R=E\setminus(F\cup D).
       \]
       Define the \(k\)-algebra specialization
       \(\phi_{f,D}:S_E\to S_R\) by
       \begin{equation}
\begin{array}{c|cc}
        &v_{e,f(e)}&v_{e,1-f(e)}\\ \hline
        e\in F&0&1\\
        e\in D&1&1
       \end{array}
       \qquad\text{and}\qquad
       \phi_{f,D}(v_{e,b})=v_{e,b}\quad(e\in R).
       \label{eq:damp-colored-specialization}
\end{equation}
       Then
       \begin{equation}
J_K=\phi_{f,D}(J_X).
        \label{eq:damp-minor-dual-ideal}
\end{equation}
 \item Consequently, \(X\in\Obs(\Damp)\) if and only if \(J_X\) has no
       linear-quotient order while the image under every proper
       specialization of the form
       \eqref{eq:damp-colored-specialization} has linear quotients.
\end{enumerate}
\end{proposition}

\begin{proof}
For an order \(x_1,\ldots,x_t\) of \(X\), the monomial-colon
formula gives
\begin{equation}
(m_{x_1},\ldots,m_{x_{j-1}}):m_{x_j}
 =\left(
   \frac{m_{x_i}}{\gcd(m_{x_i},m_{x_j})}:i<j
  \right).
 \label{eq:damp-linear-quotient-colon}
\end{equation}
The quotient belonging to \(x_i\) is the product of the literals indexed
by the coordinates on which \(x_i\) and \(x_j\) differ.  The right-hand
side of \eqref{eq:damp-linear-quotient-colon} is generated by variables
exactly when, for every \(i<j\), there is \(q<j\) such that \(x_q\) is a
cube neighbor of \(x_j\) lying in the Boolean interval between \(x_i\)
and \(x_j\).  Induction on Hamming distance shows that this is equivalent
to every prefix \(\{x_1,\ldots,x_j\}\) being isometric.
Proposition~4.2 of Chalopin--Chepoi--Moran--Warmuth makes this a
prefix-building corner order; reversing it gives a corner peeling in
our deletion convention
\cite[Proposition~4.2]{ChalopiChepoiMoran22}.  This proves item~1.
Equivalently, it is the usual shellability--linear-quotients
correspondence applied to \(\Delta_X\).

For item~2, consider a generator \(m_x\).  If \(x|_F\ne f\), one of its
factors is sent to zero.  If \(x|_F=f\), all its \(F\)-factors are sent to
one; the two factors of every contracted color are also sent to one.
The nonzero image is therefore
\[
 \phi_{f,D}(m_x)
 =\prod_{e\in R}v_{e,1-x_e}
 =m_{x|_R}.
\]
Repeated residual concepts give repeated monomials and hence do not
change the generated ideal.  This proves
\eqref{eq:damp-minor-dual-ideal}.  The normal form for pc-minors as a
conditioning followed by coordinate contractions now gives item~3 from
item~1.
\end{proof}

\begin{proposition}[Ball form of ampleness]
\label{prop:damp-yang-homology-ball}
Let \(\varnothing\ne X\subsetneq Q_E\), and fix a field \(k\).  Then
the following are equivalent:
\begin{enumerate}
\item \(X\) is ample;
\item \(\Delta_X\) is a constructible simplicial ball of dimension
      \(|E|-1\);
\item \(\Delta_X\) is a \(k\)-homology ball of dimension \(|E|-1\).
\end{enumerate}
When these conditions hold,
\begin{equation}
\partial\Delta_X
 =\Delta_X\cap\Delta_{Q_E\setminus X}.
 \label{eq:damp-yang-ball-boundary}
\end{equation}
Moreover, \(X\) is corner-peelable if and only if this ball is shellable.
Thus the corner-peelable ample families are exactly the cube families whose
Yang complexes are shellable simplicial balls.
\end{proposition}

\begin{proof}
The ample--Cohen--Macaulay Yang dictionary says that \(X\) is ample if
and only if \(\Delta_X\) is Cohen--Macaulay over \(k\)
\cite{BousquetYangComplex26}.  Because
\(\varnothing\ne X\subsetneq Q_E\), the Yang complex is a nonempty
proper full-dimensional subcomplex of the boundary of the
cross-polytope.  Theorem~4.3 of Fl{\o}ystad--Gabrielsen--Mafi applies
to precisely such a subcomplex and makes it a constructible
triangulated ball \cite[Theorem~4.3]{FloystadGabrielsenMafi26}.  This
proves that item~1 implies item~2.  Conversely, a constructible
simplicial ball is Cohen--Macaulay over every field, so the same Yang
dictionary gives item~1.  A simplicial ball is a \(k\)-homology ball,
and a \(k\)-homology ball is Cohen--Macaulay over \(k\); the Yang
dictionary therefore also proves the equivalence with item~3.

It remains to identify the boundary.  Write
\(\Sigma_E=\Delta_{Q_E}\), the boundary of the
\(|E|\)-dimensional cross-polytope.  The complementary-ball theorem for
polarizations identifies \(\Delta_{Q_E\setminus X}\) as the
complementary full-dimensional ball in \(\Sigma_E\), and the two balls
meet in their common boundary
\cite[Corollaries~8.5--8.6]{FloystadGabrielsenMafi26}.  This proves
\eqref{eq:damp-yang-ball-boundary}.  Finally, the
shellability--corner-peeling correspondence
\cite{BousquetYangComplex26,ChalopiChepoiMoran22} gives the last assertion.
\end{proof}

The specialization criterion in
Proposition~\ref{prop:damp-algebraic-minor-dictionary} appears to quantify
over the whole pc-minor order.  In fact, pc-minor closure reduces it to
three substitutions per color.  We record the resulting one-color
parametrization explicitly.

For \(e\in E\) and \(b\in\{0,1\}\), let
\(r_{e,b}:S_E\to S_{E\setminus\{e\}}\) and
\(c_e:S_E\to S_{E\setminus\{e\}}\) be the homomorphisms that fix all
other variables and satisfy
\begin{equation}
\begin{aligned}
  r_{e,b}(v_{e,b})&=0,
  &r_{e,b}(v_{e,1-b})&=1,\\
  c_e(v_{e,0})&=1,
  &c_e(v_{e,1})&=1.
 \end{aligned}
 \label{eq:damp-one-color-maps}
\end{equation}
Thus \(r_{e,b}\) retains the concepts on the side \(x_e=b\), whereas
\(c_e\) contracts color \(e\).  Repeated or divisible monomials are, as
usual, removed when the image ideal is written with its minimal
generators.  We use the vacuous convention that an ideal with at
most one minimal generator, including the zero ideal, has linear
quotients; the zero colon is generated by the empty set of variables.

Let \(\mathfrak B_d^{\mathrm{col}}\) be the set of squarefree monomial
ideals \(J\subseteq k[v_{e,0},v_{e,1}:e\in[d]]\) satisfying the following
four conditions on \(J\) and its one-color images.
\begin{enumerate}
 \item The minimal generators are distinct monomials of the form
       \(\prod_{e\in[d]}v_{e,1-x_e}\), and both values of \(x_e\) occur
       among them for every \(e\).
 \item \(J\) has a \(d\)-linear resolution.
 \item \(J\) has no linear-quotient order.
 \item Each of the \(3d\) ideals
       \begin{equation}
r_{e,0}(J),\qquad r_{e,1}(J),\qquad c_e(J)
        \quad(e\in[d])
       \label{eq:damp-one-color-tests}
\end{equation}
       has linear quotients.
\end{enumerate}
The first condition lets us recover an irredundant cube family \(X(J)\)
from the exponent choices of the generators.  The hyperoctahedral group
\(C_2\wr S_d\) acts by permuting colors and interchanging the two variables
within a color.

\begin{theorem}[Colored-ideal parametrization]
\label{thm:damp-descendant-free-parametrization}
Write \(\operatorname{Iso}(\mathcal C)\) for the set of partial-cube
isomorphism classes in a class \(\mathcal C\).  Then
\begin{equation}
\operatorname{Iso}\bigl(\Obs(\Damp)\bigr)
 =
 \{[Q_n^{--}]:n\ge3\}
 \mathbin{\dot\cup}
 \coprod_{d\ge1}
 \mathfrak B_d^{\mathrm{col}}/(C_2\wr S_d).
 \label{eq:damp-descendant-free-basis}
\end{equation}
On the second summand the bijection is \([J]\mapsto[X(J)]\).
Equivalently, for an irredundantly embedded ample family
\(X\subseteq Q_d\),
\begin{equation}
X\in\Obs(\Damp)
 \quad\Longleftrightarrow\quad
 \begin{cases}
  X\notin\Damp,\\
  \rho_e^0X,\rho_e^1X,\pi_eX\in\Damp
     &\text{for every }e\in[d].
 \end{cases}
 \label{eq:damp-descendant-free-cubical}
\end{equation}
In particular, membership in the ample part of the forbidden basis is
decided from \(J\) and its \(3d\) one-color images in
\eqref{eq:damp-one-color-tests}; no additional proper pc-minor has to
be inspected.

Within this parametrization, the cornerless basis members are exactly
those \(J\in\mathfrak B_d^{\mathrm{col}}\) for which deletion of every
minimal generator destroys \(d\)-linearity.  The remaining ideals
parametrize the cornered basis members.
\end{theorem}

\begin{proof}
Let \(J\in\mathfrak B_d^{\mathrm{col}}\).  The first condition defines an
irredundant family \(X=X(J)\subseteq Q_d\).  By the Yang-complex
correspondence, the second condition is equivalent to ampleness of \(X\);
in particular, \(X\) is isometric.  Proposition
\ref{prop:damp-algebraic-minor-dictionary} and the third condition give
\(X\notin\Damp\).

The three maps belonging to \(e\) are precisely the Alexander-dual maps
of the elementary pc-minors
\[
 \rho_e^0X,\qquad \rho_e^1X,\qquad \pi_eX.
\]
The fourth condition and the same dictionary therefore put every
elementary pc-minor of \(X\) in \(\Damp\).  Every proper pc-minor of \(X\)
is a pc-minor of at least one elementary pc-minor.  Since \(\Damp\) is
pc-minor closed, every proper pc-minor of \(X\) belongs to \(\Damp\), and
hence \(X\in\Obs(\Damp)\cap\Ample\).

Conversely, let \(X\in\Obs(\Damp)\cap\Ample\) have irredundant dimension
\(d\).  Its ideal \(J_X\) satisfies the first condition.  Ampleness gives
the \(d\)-linear resolution, and \(X\notin\Damp\) rules out linear
quotients.  Every elementary pc-minor is proper and belongs to \(\Damp\),
so all the ideals in \eqref{eq:damp-one-color-tests} have linear
quotients.  Thus \(J_X\in\mathfrak B_d^{\mathrm{col}}\).  The uniqueness of
an irredundant partial-cube embedding up to relabelling and reorientation
of its \(\Theta\)-classes gives the quotient by \(C_2\wr S_d\).
Proposition~\ref{prop:nonample-layer} supplies the first family in
\eqref{eq:damp-descendant-free-basis} and shows that the union is disjoint.

Finally, the ample corner-deletion criterion identifies a corner \(x\)
with preservation of ampleness after deleting \(x\).  Under Alexander
duality this is preservation of \(d\)-linearity after deleting \(m_x\),
which proves the last assertion.
\end{proof}

\begin{corollary}[Yang-ball parametrization]
\label{cor:damp-descendant-free-yang-ball-parametrization}
Let \(X\subseteq Q_d\) be irredundant and ample, put
\(\Delta=\Delta_X\), and write
\(V_e=\{v_{e,0},v_{e,1}\}\) for the two vertices of color \(e\).
Then \(X\) is a forbidden pc-minor for \(\Damp\) if and only if
\begin{enumerate}
 \item \(\Delta\) is a nonshellable constructible simplicial
       \((d-1)\)-ball;
 \item for every \(e\in[d]\) and \(b\in\{0,1\}\), the link
       \(\operatorname{lk}_{\Delta}(v_{e,b})\) is shellable; and
 \item for every \(e\in[d]\), the color-pair deletion
       \(\Delta[V(\Delta)\setminus V_e]\) is shellable.
\end{enumerate}
Consequently the full forbidden basis consists of the antipodally
twice-punctured cubes \(Q_n^{--}\), \(n\ge3\), and these
color-minimal nonshellable Yang balls, modulo permutations and
reorientations of the colors.  No iterated link or additional pc-minor test is required.
\end{corollary}

\begin{proof}
The elementary Yang-complex identities are
\begin{equation}
\Delta_{\rho_e^bX}
   =\operatorname{lk}_{\Delta_X}(v_{e,b}),
 \qquad
 \Delta_{\pi_eX}
   =\Delta_X[V(\Delta_X)\setminus V_e].
 \label{eq:damp-elementary-yang-minors}
\end{equation}
Indeed, taking the link selects precisely the facets whose concepts
have \(e\)-value \(b\) and then removes that literal.  Deleting both
vertices of color \(e\) projects every facet to the remaining colors;
repeated projected facets are immaterial, exactly as in contraction.

By Proposition~\ref{prop:damp-yang-homology-ball}, ampleness of \(X\)
makes \(\Delta\) a constructible ball, and \(X\notin\Damp\) is
equivalent to nonshellability of that ball.  The same proposition and
\eqref{eq:damp-elementary-yang-minors} identify items~2 and~3 with
corner peelability of all \(3d\) elementary pc-minors.  Theorem
\ref{thm:damp-descendant-free-parametrization} now proves the
equivalence.  Proposition~\ref{prop:nonample-layer} supplies the
nonample part of the final statement.
\end{proof}

\section{Core normal forms}
\label{sec:damp-cores}

\begin{definition}[Critical cores and linear-extension words]
\label{def:damp-critical-core-word}
Let \(\mathfrak C_d^{\mathrm{col}}\) be the set of squarefree monomial
ideals \(K\subseteq k[v_{e,0},v_{e,1}:e\in[d]]\) with the following
properties:
\begin{enumerate}
 \item every minimal generator chooses exactly one variable from each
       color pair;
 \item \(K\) has a \(d\)-linear resolution but has no linear-quotient
       order; and
 \item \(K\setminus u\) does not have a \(d\)-linear resolution for any
       minimal generator \(u\) of \(K\).
\end{enumerate}
Unlike the ideals in \(\mathfrak B_d^{\mathrm{col}}\), a member of
\(\mathfrak C_d^{\mathrm{col}}\) need not use both variables of every
color.

For any full-transversal ideal \(K\), a
\emph{linear-extension word} over \(K\) is a sequence of distinct
full-transversal monomials \(\mathbf u=(u_1,\ldots,u_t)\), none a
generator of \(K\), such that, on putting
\begin{equation}
K_0=K,\qquad K_i=K_{i-1}+(u_i)\quad(1\le i\le t),
 \label{eq:damp-critical-core-chain}
\end{equation}
the colon \(K_{i-1}:u_i\) is generated by variables for every \(i\).
The empty word is allowed.
\end{definition}

\begin{theorem}[Critical-core extension normal form]
\label{thm:damp-critical-core-normal-form}
The ideal family in Theorem~\ref{thm:damp-descendant-free-parametrization}
is exactly
\begin{equation}
\mathfrak B_d^{\mathrm{col}}
 =
 \left\{
 \begin{array}{@{}c|l@{}}
 K_t&
 \begin{array}{l}
 K\in\mathfrak C_d^{\mathrm{col}},\
 (u_1,\ldots,u_t)\text{ is a linear-extension word over }K,\\
 K_t\text{ is irredundant and has no linear quotients},\\
 r_{e,0}(K_t),r_{e,1}(K_t),c_e(K_t)
   \text{ have linear quotients for every }e\in[d]
 \end{array}
 \end{array}
 \right\}.
 \label{eq:damp-critical-core-normal-form}
\end{equation}
Every \(K_i\) in \eqref{eq:damp-critical-core-chain} has a
\(d\)-linear resolution.  Moreover, the word has length zero exactly
for the cornerless members of \(\mathfrak B_d^{\mathrm{col}}\); the
cornered members have \(t\ge1\).

In simplicial language, every ample forbidden pc-minor is obtained by
adjoining a sequence of shelling facets to a minimal Cohen--Macaulay,
nonshellable Yang subcomplex.  Pc-minimality is imposed only at the final
complex by the \(3d\) one-color tests.
\end{theorem}

\begin{proof}
First let \(K\in\mathfrak C_d^{\mathrm{col}}\) and let
\((u_1,\ldots,u_t)\) be a linear-extension word.  Suppose inductively
that \(K_{i-1}\) has a \(d\)-linear resolution.  The exact
sequence
\begin{equation}
0\longrightarrow
 S/(K_{i-1}:u_i)(-d)
 \xrightarrow{\ \cdot u_i\ }
 S/K_{i-1}
 \longrightarrow S/K_i\longrightarrow0
 \label{eq:damp-linear-extension-exact-sequence}
\end{equation}
and the fact that \(K_{i-1}:u_i\) is generated by variables show that
\(\operatorname{reg}(S/K_i)\le d-1\).  Since \(K_i\) is generated in
degree \(d\), it has a \(d\)-linear resolution.  Thus every \(K_i\) is
linear.  If the endpoint has the remaining properties displayed in
\eqref{eq:damp-critical-core-normal-form}, then it satisfies the four
defining conditions of \(\mathfrak B_d^{\mathrm{col}}\).

Conversely, take \(J\in\mathfrak B_d^{\mathrm{col}}\).  Starting from
\(J\), delete a minimal generator whenever the deletion still has a
\(d\)-linear resolution.  The process terminates at an ideal \(K\) for
which no further deletion preserves \(d\)-linearity.  Reverse the
deleted generators and call them \(u_1,\ldots,u_t\).  Every ideal
\(K_i\) in the resulting chain has a \(d\)-linear resolution.
Theorem~2.3 of Lu says that, after deleting any one minimal generator
from a monomial ideal with a linear resolution, the colon of the
deleted-generator ideal by that generator is generated by variables
\cite[Theorem~2.3]{Lu23quasilinear}.  Applied to \(K_i\) and its
generator \(u_i\), it gives that each \(K_{i-1}:u_i\) is generated by
variables.  Hence the reversed sequence is a linear-extension word.

It remains to check that \(K\) has no linear quotients.  Otherwise take
a linear-quotient order of \(K\) and append
\(u_1,\ldots,u_t\).  The defining colon condition makes the resulting
order a linear-quotient order of \(J\), contrary to
\(J\in\mathfrak B_d^{\mathrm{col}}\).  Thus
\(K\in\mathfrak C_d^{\mathrm{col}}\), proving the reverse inclusion in
\eqref{eq:damp-critical-core-normal-form}.

The endpoint \(J\) is critical linear exactly when the deletion process
stops immediately.  By the ample corner-deletion criterion, this is
equivalent to the corresponding concept class being cornerless.  This
proves the assertion about the word length.  The final simplicial
statement is the shelling--linear-extension translation; it is also the
Yang-complex specialization of the theorem that every Cohen--Macaulay
complex is shelled over a minimal Cohen--Macaulay subcomplex
\cite[Theorem~2.1]{DaoDoolittleLyle20minimalCM}.
\end{proof}


\section{Signed minimal nonfaces and one-concept extensions}
\label{sec:damp-extensions}

Using the local extension test from
Section~\ref{sec:damp-local-extension-foundations}, this section
gives its algebraic and topological forms.  It then encodes every ample
family with a fixed shattered-set complex by the traces omitted on
its minimal nonfaces and expresses corner peelings in the resulting
trace table.

\subsection{Algebraic and topological forms of one-concept additions}

\begin{proposition}[Algebraic and topological form of one-facet attachment]
\label{prop:damp-one-facet-attachment}
Retain the hypotheses of
Lemma~\ref{lem:damp-one-concept-extension-test}, and suppose that
\(A=S\cup\{a\}\) is ample.  Put \(I=I_S(a)\), let
\[
 F_a=\{v_{e,a_e}:e\in F\}
\]
be the new Yang facet, and write \(2^W\) for the simplex on a vertex set
\(W\).  Then
\begin{align}
 J_S:m_a&=(v_{e,a_e}:e\in I),
 \label{eq:damp-one-facet-colon}\\
 2^{F_a}\cap\Delta_S
 &=\partial 2^{\{v_{e,a_e}:e\in I\}}
   *2^{\{v_{e,a_e}:e\in F\setminus I\}}.
 \label{eq:damp-one-facet-intersection}
\end{align}
Equivalently, the intersection is the union of the ridges
\[
 F_a\setminus\{v_{e,a_e}\},\qquad e\in I.
\]
If \(S\in\Obs(\Damp)\cap\Ample\), then
\(\varnothing\ne I\subsetneq F\), so the complex in
\eqref{eq:damp-one-facet-intersection} is a shellable
\((|F|-2)\)-ball.
\end{proposition}

\begin{proof}
The family \(A\) is isometric.  For every \(x\in S\), take an
\(a\)--\(x\) geodesic in \(A\).  Its first edge has some direction
\(e\in I\).  A hypercube geodesic uses each direction at most once, so
\(x_e\ne a_e\).  Hence the monomial quotient
\[
 \frac{m_x}{\gcd(m_x,m_a)}
\]
is divisible by \(v_{e,a_e}\).  Conversely, the neighbor \(a^{\{e\}}\)
belongs to \(S\) for every \(e\in I\), and its quotient is precisely
\(v_{e,a_e}\).  The monomial-colon formula proves
\eqref{eq:damp-one-facet-colon}.

A face \(G\subseteq F_a\) belongs to \(\Delta_S\) precisely when
\(G\subseteq F_x\) for some \(x\in S\).  By the preceding geodesic
argument, \(x\) differs from \(a\) in some \(e\in I\), and therefore
\(G\subseteq F_a\setminus\{v_{e,a_e}\}\).  Conversely, this ridge is
contained in the facet \(F_{a^{\{e\}}}\) of \(\Delta_S\).  This proves
the ridge-union description, which is exactly the join in
\eqref{eq:damp-one-facet-intersection}.

If \(S\) is an ample obstruction, it is nonempty, and connectedness of
the isometric extension \(A\) gives \(I\ne\varnothing\).  If \(I=F\),
condition~\eqref{eq:damp-one-extension-punctured-cube} gives
\(S=Q_F\setminus\{a\}\), which is corner-peelable.  Thus \(I\subsetneq
F\).  The displayed join is the join of a sphere with a nonempty
simplex, hence a shellable ball of dimension \(|F|-2\).
\end{proof}

\subsection{Signed minimal nonfaces and first-loss orders}

\begin{definition}[Signed minimal-nonface presentation]
\label{def:damp-signed-minimal-nonface-presentation}
Here a simplicial complex is a nonvoid downward-closed family, so it
contains the empty face.  Let \(\Sigma\subseteq 2^F\) be such a complex,
and let
\[
 \mathcal N(\Sigma)
 =\{I\subseteq F:I\notin\Sigma\text{ and every proper subset of }I
                         \text{ belongs to }\Sigma\}
\]
be its clutter of minimal nonfaces.  A \emph{signing} of
\(\mathcal N(\Sigma)\) is a choice
\(\boldsymbol\sigma=(\sigma_I)_{I\in\mathcal N(\Sigma)}\), where
\(\sigma_I\in Q_I\).  Its satisfying family is
\begin{equation}
 H(\Sigma,\boldsymbol\sigma)
 =\{x\in Q_F:x|_I\ne\sigma_I
                 \text{ for every }I\in\mathcal N(\Sigma)\}.
 \label{eq:damp-signed-clutter-solutions}
\end{equation}
We call the signing \emph{admissible} when
\begin{equation}
 |H(\Sigma,\boldsymbol\sigma)|=|\Sigma|.
 \label{eq:damp-signing-admissibility}
\end{equation}
Thus \((I,\sigma_I)\) forbids one Boolean trace on each minimal
nonface.  Admissibility is a single cardinality condition on this signed
clutter; it contains no recursive minor or peeling test.
\end{definition}

\begin{theorem}[Parametrization by signed minimal nonfaces]
\label{thm:damp-fixed-shadow-signed-clutter}
For every simplicial complex \(\Sigma\subseteq2^F\), write
\(H{\upharpoonright}I=\{x|_I:x\in H\}\).  The map
\begin{equation}
 H\longmapsto
 \boldsymbol\sigma_H=(\sigma_I(H))_{I\in\mathcal N(\Sigma)},
 \qquad
 Q_I\setminus(H{\upharpoonright}I)=\{\sigma_I(H)\},
 \label{eq:damp-fixed-shadow-signing-map}
\end{equation}
is a bijection from the ample families \(H\subseteq Q_F\) with
\(\operatorname{Sh}(H)=\Sigma\) to the admissible signings of
\(\mathcal N(\Sigma)\).  Its inverse is
\(\boldsymbol\sigma\mapsto H(\Sigma,\boldsymbol\sigma)\).
In particular, the signed minimal nonfaces determine the entire ample
family, not merely its shattered-set complex.
\end{theorem}

\begin{proof}
Let \(H\) be ample with shattered-set complex \(\Sigma\), and take
\(I\in\mathcal N(\Sigma)\).  Coordinate projection preserves
ampleness \cite[Theorem~2]{Lawrence83}, and
\[
 \operatorname{Sh}(H{\upharpoonright}I)=2^I\setminus\{I\}.
\]
The Sandwich equality for the ample family \(H{\upharpoonright}I\)
therefore gives \(|H{\upharpoonright}I|=2^{|I|}-1\).  Hence this
projection omits a unique trace
\(\sigma_I(H)\).

Plainly \(H\subseteq H(\Sigma,\boldsymbol\sigma_H)\).  The latter family
does not shatter any member of \(\mathcal N(\Sigma)\), and hence does not
shatter any nonface of \(\Sigma\).  The upper Sandwich inequality gives
\[
 |H(\Sigma,\boldsymbol\sigma_H)|
 \le |\operatorname{Sh}(H(\Sigma,\boldsymbol\sigma_H))|
 \le |\Sigma|=|H|.
\]
Thus equality holds throughout and
\(H=H(\Sigma,\boldsymbol\sigma_H)\); in particular the signing is
admissible and is uniquely determined.

Conversely, let \(\boldsymbol\sigma\) be admissible and put
\(H=H(\Sigma,\boldsymbol\sigma)\).  Every nonface of \(\Sigma\)
contains a member of \(\mathcal N(\Sigma)\), whose prescribed trace is
absent from \(H\).  Consequently
\(\operatorname{Sh}(H)\subseteq\Sigma\).  Now
\[
 |\Sigma|=|H|
 \le |\operatorname{Sh}(H)|
 \le |\Sigma|.
\]
Equality in the Sandwich bound characterizes ample families, so \(H\)
is ample and \(\operatorname{Sh}(H)=\Sigma\).  Its omitted traces are
the prescribed \(\sigma_I\), proving that the two maps are inverse.
\end{proof}

\begin{proposition}[Lawrence's signed elimination criterion and minor rules]
\label{prop:damp-signed-route-calculus}
A \emph{literal} is a pair $(e,b)\in F\times\{0,1\}$; its opposite is
$(e,1-b)$. A literal set is consistent if it contains at most one literal
at each coordinate. Let $\mathcal R$ be a clutter of consistent literal
sets, and let $H(\mathcal R)$ consist of the words containing no member of
$\mathcal R$. Then $H(\mathcal R)$ is ample and $\mathcal R$ is its full
clutter of minimally missing patterns if and only if
\begin{equation}
 \text{for every }P,Q\in\mathcal R\text{ and }(e,b)\in P,
 \ (e,1-b)\in Q,\quad
 \text{some }U\in\mathcal R\text{ satisfies }
 U\subseteq(P\cup Q)\setminus\{(e,0),(e,1)\}.
 \label{eq:damp-lawrence-route-elimination}
\end{equation}
This is Lawrence's route-system characterization
\cite[Theorems~5--6]{Lawrence83}. It applies also when $P,Q$ disagree
at other coordinates. For the signed minimal-nonface data of
Definition~\ref{def:damp-signed-minimal-nonface-presentation}, it is
equivalent to the cardinality admissibility condition
\eqref{eq:damp-signing-admissibility}.

Write $\min\mathcal S$ for inclusion-minimal members of a set family,
and put $e^b=(e,b)$. For an ample $H$ with route clutter $\mathcal R$,
the complete route clutters of its elementary images are
\begin{align}
 \mathcal R(\pi_eH)
 &=\{P\in\mathcal R:e^0,e^1\notin P\},
 \label{eq:damp-route-contraction}\\
 \mathcal R(\rho_e^bH)
 &=\min\{P\setminus\{e^b\}:P\in\mathcal R,\ e^{1-b}\notin P\},
 \label{eq:damp-route-restriction}\\
 \mathcal R(H^{\{e\}})
 &=\min\{P\setminus\{e^0,e^1\}:P\in\mathcal R\}.
 \label{eq:damp-route-reduction}
\end{align}
The empty clutter represents the full cube; the clutter consisting of the
empty pattern represents the empty family. Thus these formulas retain
constant coordinates and cover empty images as well.
\end{proposition}
\begin{proof}
The equivalence for an arbitrary consistent clutter is exactly the cited
route-system theorem. In the fixed-shadow setting, the literal sets
associated with $(I,\sigma_I)$ form a clutter. If they satisfy signed
elimination, the theorem makes their avoidance class ample and makes them
its full minimally missing patterns. A support is unshattered precisely
when it contains the support of a missing pattern. Its shattered complex
is therefore $\Sigma$, so its cardinality is $|\Sigma|$.
Conversely, the cardinality criterion gives an ample class with the
prescribed missing patterns by
Theorem~\ref{thm:damp-fixed-shadow-signed-clutter}. To see that these
are all its minimally missing patterns, let $p$ be minimally missing on
$I$. Every one-bit neighbour of $p$ occurs in the projection to $I$,
by minimality. The complement of that projection is ample and hence
connected when nonempty, by the projection and complement closure
properties of lopsided sets \cite[Theorem~2 and Lemma~1]{Lawrence83}. Since $p$ is
isolated in this complement, the complement is $\{p\}$. Thus $I$ is
a minimal nonface of the shattered complex, and $p$ is the prescribed
pattern there. Lawrence's necessity now gives
\eqref{eq:damp-lawrence-route-elimination}.

For contraction, a pattern not using $e$ is missing from $\pi_eH$
exactly when it is missing from $H$. Taking its minimal subpatterns proves
\eqref{eq:damp-route-contraction}. For restriction, a visible pattern $T$
is missing exactly when $T\cup\{e^b\}$ contains a route. That route
cannot contain $e^{1-b}$, and deleting $e^b$ gives a missing subpattern
of $T$. Conversely each surviving trimmed route is missing in the
restriction. Minimalization proves~\eqref{eq:damp-route-restriction}.
These two arguments work for the full route clutter of any family.

For reduction, a visible word belongs to both sections precisely when it
avoids every route after erasing its $e$-literal. It remains to show that
the displayed minimal trimmed routes are already the complete route
clutter; arbitrary incomplete Boolean formulas would not justify this.
Let $P',Q'$ be two minimal trimmed routes with opposite literals at a
visible coordinate $f$. Choose parent routes $P,Q$ trimming to them.
Eliminate at $f$ in the parent clutter, obtaining a route $U$ contained
in their union without the two $f$-literals. Its trim is contained in
$(P'\cup Q')\setminus\{f^0,f^1\}$ and contains a minimal trimmed route.
Thus the trimmed clutter again satisfies signed elimination. Lawrence's
sufficiency says that it is the full route clutter of its avoidance
family, which we identified as $H^{\{e\}}$.
\end{proof}

For example, the ample path $\{000,100,110,111\}$ on coordinates
$x,y,z$ has routes $\{x^0,y^1\}$, $\{x^0,z^1\}$ and
$\{y^0,z^1\}$. Eliminating the opposite $y$-literals in the first
and third routes gives the second. Fixing $y=0$ leaves only the minimal
route $\{z^1\}$; projecting away $y$ leaves $\{x^0,z^1\}$; reducing
in $y$ leaves the two singleton routes $\{x^0\}$ and $\{z^1\}$.
These respectively describe $\{00,10\}$, $\{00,10,11\}$ and $\{10\}$
on the visible coordinates $x,z$.

The quantifier over \emph{every} opposite-literal pair is essential.
For $\{00,11\}$ the two minimal missing patterns are $01$ and $10$.
Ordinary Boolean resolution discards their tautological resolvents, whereas
\eqref{eq:damp-lawrence-route-elimination} fails: eliminating either
coordinate leaves no route. Thus ordinary resolution closure is not a
replacement for the signed elimination axiom.

This supplies structural admissibility for the ample support and exact
signed data for every elementary image, with no enumeration of all allowed
words. It is existing ample-set theory, not a criterion for $\Damp$:
signed elimination holds for both peelable classes and ample obstructions.
For the forbidden-minor generator, the additional unresolved condition must
exclude an acyclic USO on the original support while guaranteeing one on
every proper image described above. Requiring peeling orders of those images
would remain the existing recognition test.

\begin{corollary}[Signed sufficient conditions for every prescribed sink]
\label{cor:damp-large-pattern-positive-filter}
Let $H\ne\varnothing$ be ample and let $\mathcal R_{\ge3}(H)$ denote
its minimally missing patterns with support size at least three.
Each of the following conditions implies that every vertex of $H$
is the sink of an acyclic USO:
\begin{enumerate}
 \item $|\mathcal R_{\ge3}(H)|\le2$;
 \item every pair of members of $\mathcal R_{\ge3}(H)$ agrees on
 the intersection of its supports.
\end{enumerate}
In particular $H\in\Damp$. Consequently every ample forbidden pc-minor,
and every rooted factor of Theorem~\ref{thm:damp-cut-vertex-obstructions},
has at least three minimally missing patterns of support size at least
three, and two of these patterns disagree on a common coordinate.
Binary missing patterns are unrestricted in both sufficient conditions.
\end{corollary}
\begin{proof}
We apply two established geometric criteria from the companion manuscript
\emph{Maximum parents of ample classes: constructions and signed criteria},
\path{../CompressionSchemes/maximum_parent_methods.tex}, specifically
\path{thm:two-large-patterns} and \path{thm:agreeing-large-patterns}.
Their precise input is a nonempty ample class with, respectively, at most
two larger minimally missing patterns or agreement of all larger patterns
on their overlaps. Constant coordinates are allowed. Their conclusion is
that $H$ is the strict coordinate-sign class of the interior of a bounded
full-dimensional convex polytope. Thus the coordinate hyperplanes on that
open polyhedron give geometric realizability in the sense of
Proposition~\ref{prop:damp-realizable-prescribed-endpoint}.
That proposition supplies a peeling to each prescribed vertex, and the
CCMW correspondence turns it into an acyclic USO with that sink.
An ample forbidden class is not in $\Damp$, while a rooted factor fails
to admit its specified root as a sink. Neither can satisfy either
sufficient condition, proving both necessary conclusions.
\end{proof}

The geometric input has more content than feasibility of an unspecified
system of inequalities. In the first case the companion proof uses the
fixed forms
\[
 \ell_p(z)=\sum_{i\in\operatorname{supp}(p)}(2p_i-1)z_i.
\]
For every allowed concept it constructs a point in its strict orthant
where all these forms are negative, using the partial order imposed by
binary missing patterns and signed elimination. The agreement theorem
allows arbitrary prescribed positive magnitudes for the coefficients.
Finitely many witness points can then be scaled and bounded to obtain
the stated polytope; every forbidden orthant is excluded by its missing
pattern. These are the proved geometric criteria being imported, not a
claim that every ample class is realizable.

Together with Proposition~\ref{prop:damp-signed-route-calculus}, the
corollary supplies direct sufficient positivity tests on the signed data
of each elementary image. It does not say that positivity forces either
condition, or that every proper minor of a forbidden input satisfies one.
The remaining generator problem cannot discard an image merely because
both tests fail: such a failure gives no negative certificate.

\begin{theorem}[Realization when each missing pattern conflicts with at most one other]
\label{thm:damp-matching-route-realization}
Let $H\subseteq Q_F$ be nonempty and ample. Form a graph on its full
clutter $\mathcal R(H)$ of minimally missing patterns, joining two patterns
when they assign opposite values to a common coordinate. Suppose this graph
has maximum degree at most one. For arbitrary positive coefficients
$\lambda_{p,i}>0$, the forms
\[
 \ell_p(z)=\sum_{i\in\operatorname{supp}(p)}
       \lambda_{p,i}(2p_i-1)z_i
 \qquad(p\in\mathcal R(H))
\]
realize $H$ as the strict coordinate-sign class of a bounded open convex
polyhedron defined by a box and inequalities $\ell_p(z)<-\varepsilon$,
for some $\varepsilon>0$. Consequently every vertex of $H$ is the sink
of an acyclic USO. In particular, the full missing-pattern conflict graph
of an ample forbidden pc-minor or a rooted factor has a vertex of degree
at least two.
\end{theorem}
\begin{proof}
The zero-coordinate singleton is immediate, so assume $|F|>0$.
Fix $c\in H$ and reflect coordinates so that $c=0$. The conflict graph
and coefficient magnitudes are unchanged. Every missing pattern contains
a literal with value one, since the zero word is allowed.
Call a coordinate \emph{pure} for this reflected system if some pattern
specifies value one there and no pattern specifies value zero there.
Call a pattern \emph{uncovered} if its support contains no pure coordinate.
These terms are used only to construct a witness for the chosen $c$.

Two conflicting patterns $p,q$ cannot both be uncovered. To prove this,
choose an opposite-literal coordinate $e$ between them. Every value-one
literal in $p$ must have its opposite in some pattern because $p$ is
uncovered; the degree bound forces that pattern to be $q$. The symmetric
statement holds for $q$. Signed elimination at $e$, from
Proposition~\ref{prop:damp-signed-route-calculus}, gives a pattern $u$
contained in the union of their literals with both $e$-literals removed.
This pattern differs from $p$ and $q$, which both use $e$.
If $u$ contains a value-one literal from $p$, it conflicts with $q$,
giving $q$ a second neighbour. If that literal comes from $q$, it gives
$p$ a second neighbour. Thus $u$ contains only value-zero literals.
The zero word would contain this missing pattern, a contradiction.

For each uncovered $p$, choose a coordinate $i(p)$ where $p$ specifies
one. Its opposite occurs in a unique conflicting pattern $q$, which is
not uncovered by the preceding argument. Moreover, no third pattern
uses coordinate $i(p)$: value one would give $q$ a second neighbour,
and value zero would give $p$ a second neighbour. Therefore changing
the magnitude at $i(p)$ affects no other uncovered pattern.

Start with positive weights $w_i=1$ and put
\[
 L_p(w)=\sum_{i\in\operatorname{supp}(p)}
             \lambda_{p,i}(1-2p_i)w_i=\ell_p(-w).
\]
For each uncovered $p$, increase $w_{i(p)}$ until $L_p(w)<0$.
Its coefficient is negative, so this is possible; the preceding paragraph
shows that it preserves all earlier uncovered inequalities. Now increase
all pure-coordinate weights by the same nonnegative amount $t$.
An uncovered pattern uses none of these coordinates, so its form does
not change. Every other pattern uses at least one pure coordinate, all
with value one, so its form decreases with a strictly negative slope.
Taking $t$ sufficiently large makes every remaining form negative.
All weights stay positive.

Undo the reflection, obtaining a witness
$z_i^c=(2c_i-1)w_i$ in the strict orthant of $c$ with
$\ell_p(z^c)<0$ for every original missing pattern $p$. Although the
witness weights depend on $c$, the forms and their prescribed coefficients
do not. Scale the finitely many witnesses so that all their form values
are below $-1$, and choose $R$ larger than every witness coordinate
magnitude. Then
\[
 P=\{z\in\mathbb R^F: |z_i|<R\ (i\in F),\quad
                       \ell_p(z)<-\tfrac12\ (p\in\mathcal R(H))\}
\]
is a bounded nonempty open convex polyhedron meeting every allowed
strict orthant. Every forbidden word contains a missing pattern $p$;
on its closed orthant all terms of $\ell_p$ are nonnegative, so that
orthant misses $P$. Thus $P$ has exactly the strict sign class $H$.
If there are no missing patterns, the box alone suffices.
Proposition~\ref{prop:damp-realizable-prescribed-endpoint} supplies a
peeling to every prescribed concept, hence the claimed acyclic USOs.
The obstruction and rooted-factor conclusions follow by contraposition.
\end{proof}

\begin{corollary}[Conflict graphs under elementary images]
\label{cor:damp-route-conflict-image}
For a nonempty ample $H$, the missing-pattern conflict graph of any
contraction, coordinate restriction, or reduction is isomorphic to a
subgraph of the missing-pattern conflict graph of $H$. Consequently its
maximum degree cannot increase. In particular, the hypothesis of
Theorem~\ref{thm:damp-matching-route-realization} is preserved by these
operations and their compositions.
\end{corollary}
\begin{proof}
Use the exact route formulas in
Proposition~\ref{prop:damp-signed-route-calculus}. For every route of the
image choose one parent route that survives or trims to it. Distinct image
routes have distinct parents, since each parent has a unique trim. If two
image routes conflict at a remaining coordinate, their chosen parents
conflict there as well. This injective map preserves every image edge,
which gives the asserted subgraph embedding. An empty image has only the
empty missing pattern and hence no edges; the same argument still applies.
Iterate for compositions.
\end{proof}

The degree condition concerns \emph{all} minimally missing patterns,
including binary ones. It is therefore different from agreement of only
the larger patterns in Corollary~\ref{cor:damp-large-pattern-positive-filter}.
It permits arbitrarily many mutually disjoint conflicting pairs.
For a concrete case outside both conditions of that corollary, take the
four-word path on three coordinates and replace each coordinate by a block
of size two. Its three routes are
\[
 \{x_1^0,x_2^0,y_1^1,y_2^1\},\quad
 \{x_1^0,x_2^0,z_1^1,z_2^1\},\quad
 \{y_1^0,y_2^0,z_1^1,z_2^1\}.
\]
There are three larger patterns and a conflict between the first and third,
but the conflict graph is a single edge and an isolated vertex. Their
avoidance family has $54$ words in six coordinates. Its ordinary positivity
also follows from existing block invariance; the example demonstrates the
additional coverage of the signed criterion, not a new forbidden family.
The replay \path{research/matching_route_realization.py/json} verifies
witness inequalities on this example and all nonempty ample classes in
$Q_3$, with unit and one nonuniform positive coefficient assignment.
The arbitrary-coefficient and arbitrary-dimension statements are proved
above. No converse to the conflict-degree condition is asserted.

\begin{theorem}[Bipartite route conflicts give a peripheral construction]
\label{thm:damp-bipartite-route-peripheral}
Let $H\subseteq Q_F$ be nonempty and ample. If its full minimally
missing-pattern conflict graph is bipartite, then $H$ is obtained from
singletons by peripheral expansions
\[
 (B\times\{0\})\cup(S\times\{1\}),\qquad
 \varnothing\ne S\subseteq B,
\]
with both children recursively obtained in the same way; constant
coordinates may be inserted freely. Thus $H$ belongs to
$\mathcal R_{\mathrm{per}}$ and every vertex of $H$ is an attainable
sink of an acyclic USO. The construction is recovered directly from its
signed patterns by repeatedly choosing a coordinate at which at most one
literal occurs.

Consequently the full missing-pattern conflict graph of every ample
forbidden pc-minor for either $\mathcal R_{\mathrm{per}}$ or $\Damp$,
and of every rooted factor of
Theorem~\ref{thm:damp-cut-vertex-obstructions}, contains an odd cycle.
\end{theorem}
\begin{proof}
We first prove that, if $F\ne\varnothing$, some coordinate occurs with
at most one polarity in the route clutter $\mathcal R(H)$. Suppose
otherwise, and fix a bipartition $\mathcal R_0,\mathcal R_1$ of its
conflict graph. At any coordinate $e$, every route using value zero is
adjacent to every route using value one. As both sets are nonempty,
all occurrences in $\mathcal R_0$ have one value $\sigma_e$, and all
occurrences in $\mathcal R_1$ have value $1-\sigma_e$. In particular,
each route in the first part is a restriction of the same word $\sigma$,
and each route in the second is a restriction of its opposite.

Choose two routes $P,Q$ with opposite literals at a coordinate $e$.
Signed elimination gives a route
\[
 U\subseteq(P\cup Q)\setminus\{e^0,e^1\}.
\]
If $U\in\mathcal R_0$, all its literals agree with $\sigma$, so none
can come from $Q$. Hence $U$ is properly contained in $P$.
If $U\in\mathcal R_1$, it is properly contained in $Q$.
Both conclusions contradict the clutter property. This proves the claim.
It also covers unused coordinates, which already have zero polarities.

Choose such a coordinate $e$. If its only occurring literal is $e^b$,
flipping a word of $H$ from value $b$ to $1-b$ cannot introduce a
missing pattern: routes not using $e$ are unchanged, and those using it
cease to match. Thus
\[
 \rho_e^bH\subseteq\rho_e^{1-b}H.
\]
If neither literal occurs, the two sections are equal. After possibly
reversing $e$, write the larger section as $B$ and the smaller as $S$.
Both are ample and their conflict graphs are bipartite by
Corollary~\ref{cor:damp-route-conflict-image}. The original family is
exactly the displayed peripheral expansion when $S\ne\varnothing$;
when $S=\varnothing$, coordinate $e$ is constant and can be suppressed.

Induct on $|F|$. The zero-coordinate case is a singleton. Apply the
same construction to the nonempty children $B,S$, each on fewer
coordinates. This gives the stated peripheral construction tree and
hence membership in $\mathcal R_{\mathrm{per}}$, whose defining rule
is precisely this recursion. It is the convention used in the
\href{../cover-recursive/research/component_label_recursion.pdf}
{peripheral-recursion companion}, Section~1.
Inductively both children admit every prescribed sink; the peripheral
endpoint criterion in Proposition~\ref{prop:damp-rooted-operations}
gives every sink of the parent. A constant coordinate does not change
the one-inclusion graph. This proves the acyclic-USO assertion.

An ample forbidden input for either class cannot have the displayed
positive construction, and a rooted factor cannot admit its specified
root as a sink. Their conflict graphs are therefore nonbipartite,
equivalently contain an odd cycle.
\end{proof}

This closes the \emph{Damp positivity} question for forest conflict
graphs, indeed for all bipartite full conflict graphs, without proving
geometric realizability for those classes. The condition includes binary
patterns. It does not resolve the companion question in which only larger
patterns have forest conflicts and binary patterns are unrestricted.
In particular it bypasses the failed operation of deleting a leaf from
an incomplete clause list: every recursion step here is an actual
coordinate operation with its complete route clutter retained.

Bipartiteness is sufficient, not necessary. The six-word path
$\{00000,10000,11000,11100,11110,11111\}$ has routes
$P_{ij}=\{i^0,j^1\}$ for $1\le i<j\le5$. Its conflict graph contains
\[
 P_{12},P_{23},P_{34},P_{45},P_{24},P_{12}
\]
as a five-cycle, while the one-inclusion graph is a path and is built
by peripheral leaf attachments. Thus an odd cycle is not a negative
certificate. The replay
\path{research/bipartite_route_recursion.py/json} constructs peripheral
trees and every prescribed endpoint order for this control, a five-word
path, and the eleven-word maximum class whose conflict graph is $P_4$.
The last control is the companion's unit-coefficient failure: the new
proof establishes its peripheral recursion without needing to repair those
coefficients. The general generator still requires a selection rule for
the nonbipartite inputs, together with soundness and exhaustive recovery.

The distinction already matters for proper minors of the terminal
$267$-word obstruction. Each of its twelve coordinate contractions has
$208$ vertices and no peripheral coordinate, although each admits every
prescribed acyclic-USO sink by
Corollary~\ref{cor:damp-267-all-proper-minor-sinks}. For every contraction
and every remaining active coordinate, the replay
\path{research/terminal_minor_peripheral_gap.py/json} records a word in
each of the two exclusive sections. These pairs certify that neither
section contains the other, so no peripheral construction can have a
first reverse step. Thus all twelve contractions lie in
$\Damp\setminus\mathcal R_{\mathrm{per}}$; in particular all their route
conflict graphs are nonbipartite. The same audit finds no bipartite route
graph among any of the seed's $36$ elementary pc-minors. It leaves the
$\mathcal R_{\mathrm{per}}$ status of the $24$ coordinate restrictions
undecided. Consequently peripheral construction trees, including those
supplied by the bipartite criterion, cannot replace the required Damp
positivity witnesses on all proper minors.

Nor does replacing those peripheral certificates by unit-coefficient
missing-pattern inequalities handle all these positive contractions.
For the coordinate-zero contraction, take the allowed word $c=204$ in
the certificate's integer encoding (coordinate zero is suppressed but
remaining indices are retained). Four of its missing patterns, written
as $(\text{support mask},\text{value mask})$, are
\[
 (102,36),\quad(2188,136),\quad(2272,128),\quad(3140,2116).
\]
For a point in the strict orthant of $c$, write
$z_i=(2c_i-1)w_i$ with $w_i>0$. The unit-coefficient signed forms for
these patterns become
\[
\begin{split}
 L_1&=w_1+w_2-w_5-w_6,\\
 L_2&=-w_2+w_3+w_7+w_{11},\\
 L_3&=w_5-w_6+w_7+w_{11},\\
 L_4&=w_2+w_6+w_{10}-w_{11}.
\end{split}
\]
Their exact combination is
\[
 L_1+3L_2+L_3+2L_4
   =w_1+3w_3+4w_7+2w_{10}+2w_{11}>0.
\]
They therefore cannot all be negative. The certificate
\path{research/contraction_unit_realization.py/json} checks membership
of the word and uses the complete projected route list; the displayed
identity is an exact obstruction to this fixed coefficient assignment.
It does not rule out a free-coefficient realization and does not
contradict the established acyclic USOs of the contraction.

\paragraph{Coordinate expansions and their signed clauses.}
The cardinality test in
Theorem~\ref{thm:damp-fixed-shadow-signed-clutter} parametrizes all ample
realizations at once.  Fixing a contraction and a reduction gives
a complementary constructive description: the remaining choices are
independent labels on connected components.  The following theorem identifies
those labels with the new bits of the signed clauses.  Its conclusion concerns ampleness. In the USO formulation, membership in
$\Damp$ asks whether the constructed support admits any acyclic USO;
forbiddenness additionally asks for independent acyclic USOs on all proper
pc-minors. A labeling alone supplies neither orientation nor those witnesses.
Fixing a restriction and reduction representation map is a more constrained
extension problem than choosing the support labels. The completion theory
in the \href{research/ample_filtrations.pdf}{representation-map writeup}
treats that additional prescribed data. Its fixed-map criteria do not
replace the existential conditions in the present support classification.

\begin{proposition}[Prescribed lifts need a union of equality partitions]
\label{prop:damp-joint-switch-conflict}
For a fixed pair of acyclic USOs on a contraction and its reduction,
the feasible upper-layer assignments need not be the assignments constant
on the blocks of a single partition. Even when every crossing edge is
individually compatible, their combined constraints can be cyclic.
This already occurs when both input families are $Q_4$.
\end{proposition}
\begin{proof}
We give explicit data and an exact description of all feasible assignments.
Encode vertices by integers $0,\ldots,15$, with adjacency given by changing
one binary digit. Orient edges forward in the two orders
\[
\begin{split}
 \alpha&=(9,1,11,10,14,15,13,3,12,7,2,6,8,5,0,4),\\
 \gamma&=(11,3,15,2,13,10,7,5,14,9,8,12,6,1,0,4),
\end{split}
\]
obtaining $D_A,D_C$. Both orders are corner peelings: at each deletion,
the remaining incident directions span a contained cube. Thus both
orientations are acyclic CCMW USOs.

For precision, the upper bit $s(x)$ specifies which copy of $x$ is deleted
last in a peeling of $Q_4\times\{0,1\}$. The orders of last and first
copies must induce $D_A$ and $D_C$, respectively. The exact orientation
gluing theorem in the representation-map writeup says that such a peeling
exists precisely when the following graph $L_s$ is acyclic: start with
$D_A$, and for every edge $x\to y$ with $s(x)\ne s(y)$ add $x\to z$
for every $D_C$-descendant $z$ of $y$, including $y$ itself.
Call an edge individually compatible if adding just its descendant
comparisons leaves $D_A$ acyclic.

The individually incompatible edges force $s$ to be constant on
\[
 \{0\},\quad B=\{1,3,5,7\},\quad\{4\},\quad\{6\},\quad
 E=\{2,10\},\quad F=\{9,11,13,14,15\},\quad\{8,12\}.
\]
Subject to these equalities, the exact additional condition is
\begin{equation}
 s(E)=s(B)\quad\text{or}\quad s(F)=s(B).
 \label{eq:damp-joint-switch-disjunction}
\end{equation}
Here $s(B)$ denotes the common value on the indicated block.
To prove necessity, if both equalities fail, the crossing edges
$3\to2$ and $15\to7$ add $3\to14$ and $15\to1$:
indeed $2\to10\to14$ and $7\to5\to1$ are paths in $D_C$.
Together with edges of $D_A$ this gives
\[
 1\longrightarrow3\longrightarrow14\longrightarrow15\longrightarrow1.
\]
For sufficiency, merge $B$ with $E$, or merge $B$ with $F$, respectively.
Add all descendant comparisons for every edge joining distinct blocks
of that partition. The resulting graphs have topological orders
\[
\begin{split}
 &(9,11,10,14,15,13,1,3,2,7,5,6,12,8,0,4),\\
 &(9,1,11,3,10,2,14,15,7,6,13,5,12,8,0,4),
\end{split}
\]
respectively. Each possible $L_s$ in that branch is a subgraph, proving
sufficiency for every assignment of the remaining bits. All finite checks
in this paragraph, the input corner checks, and the individually compatible
edge orders are replayed in \path{research/joint_switch_conflict.py};
the explicit certificates are in its accompanying JSON file.

In particular, the indicator assignments of $E$ and of $F$ are feasible,
but their sum modulo two, the indicator of $E\cup F$, is infeasible.
Every crossing edge for the latter assignment is individually compatible.
Assignments constant on the blocks of any one partition are closed under
addition modulo two, so no one partition describes the feasible set.
In this example exactly two partitions suffice, with $96$ feasible
assignments among the $128$ satisfying the forced equalities.
\end{proof}

This is a limitation of gluing prescribed maps, not a forbidden support:
the support is always $Q_5$ and is dismantlable. It strengthens the earlier
cyclic-lift example by allowing every individual crossing edge and showing
that a single equality partition cannot replace the union of partitions
in the exact lifting criterion. It neither excludes alternative parent
maps nor settles the intrinsic reduction-terminal inputs.

\begin{theorem}[Component labels and expansion signs]
\label{thm:damp-component-expansion-signs}
Let $\varnothing\ne C\subseteq A\subseteq Q_F$ be ample, let $e\notin F$,
and let $K_1,\ldots,K_k$ be the connected components of the graph induced
by $A\setminus C$.  For $b=(b_1,\ldots,b_k)\in\{0,1\}^k$, put
\[
 X_b=(C\times\{0,1\})\ \cup\
       \bigcup_{i=1}^k(K_i\times\{b_i\}).
\]
Then $X_b$ is ample, and
\begin{equation}
 \operatorname{Sh}(X_b)=\operatorname{Sh}(A)
       \cup\{T\cup\{e\}:T\in\operatorname{Sh}(C)\}.
 \label{eq:damp-component-expansion-shadow}
\end{equation}
Conversely, these are exactly the ample families on $F\cup\{e\}$ whose
contraction in $e$ is $A$ and whose reduction in $e$ is $C$.
The labels are recovered uniquely.

Write $\mathcal I=\mathcal N(\operatorname{Sh}(C))\cap
\operatorname{Sh}(A)$.  For $T\in\mathcal I$, let $\tau_T\in Q_T$
be the unique trace missing from $C|_T$.  The nonempty fiber
\[
 A(T,\tau_T)=\{x\in A:x|_T=\tau_T\}
\]
lies in a unique component $K_{i(T)}$.  The signed minimal nonfaces of
$X_b$ are exactly
\begin{enumerate}
 \item the signed minimal nonfaces of $A$, on supports not containing $e$;
 \item for each $T\in\mathcal I$, the support $T\cup\{e\}$ with missing
       trace $(\tau_T,1-b_{i(T)})$.
\end{enumerate}
Every component occurs as some $i(T)$.  Thus, with these supports, old
clauses, and traces $\tau_T$ fixed, a choice of new-coordinate bits is
admissible for the shadow in \eqref{eq:damp-component-expansion-shadow}
if and only if its bits agree whenever $i(T)=i(U)$.  There are exactly
$2^k$ such choices.  If $A=C$, the component list and $\mathcal I$ are
empty and the construction is $A\times\{0,1\}$.
\end{theorem}

\begin{proof}
We first prove ampleness directly by comparing shattered and strongly
shattered supports.  We use that complements and face sections of ample
families are ample, and that nonempty ample families are connected
\cite[Theorem~2 and Lemma~1]{Lawrence83}.

Let $Q$ be any cube contained in $A$.  The family $C\cap Q$ is ample
inside $Q$, so its complement $Q\setminus C$, when nonempty, is connected.
Consequently it belongs to one $K_i$.  The entire cube $Q$ therefore
occurs in at least one layer of $X_b$.  A support $T\subseteq F$ is
shattered by $X_b$ exactly when it is shattered by its contraction $A$;
since $A$ is ample, the preceding cube argument shows that it is also
strongly shattered by $X_b$.

A support $T\cup\{e\}$ is strongly shattered by $X_b$ exactly when $T$
is strongly shattered by $C$.  If $T$ is shattered by $C$, both layers
already realize every $T$-trace.  Conversely, suppose $T$ is not shattered
by $C$.  If it is not shattered by $A$, it cannot give a shattered support
of $X_b$.  Otherwise choose a trace $q$ missing from $C|_T$.  The nonempty
face section $A(T,q)$ is connected and avoids $C$, so it lies in one
component $K_i$.  Only layer $b_i$ realizes $q$; hence $T\cup\{e\}$ is
not shattered.  This proves \eqref{eq:damp-component-expansion-shadow}
and equality with the strongly shattered supports, proving ampleness.

For the converse, let $X$ have the stated contraction and reduction.  Every member of $C$ has both lifts, and every member of
$A\setminus C$ has a unique lift.  Two adjacent vertices of
$A\setminus C$ cannot have different lift values: their two lifts would
form an antipodal pair in a square face, with the other two vertices
absent.  Such a face section is not ample.  Thus lift values are constant
on each $K_i$, which gives $X=X_b$ and uniqueness.

The minimal nonfaces of the shadow in
\eqref{eq:damp-component-expansion-shadow} which avoid $e$ are exactly
those of $\operatorname{Sh}(A)$.  A minimal nonface containing $e$ has
the form $T\cup\{e\}$, where $T$ is shattered by $A$, is not shattered
by $C$, and every proper subset of $T$ is shattered by $C$.  These are
precisely the members of $\mathcal I$.  The unique missing trace
$\tau_T$ exists by Theorem~\ref{thm:damp-fixed-shadow-signed-clutter}.
Its fiber in $A$ is a nonempty connected face section outside $C$.
Only its component's layer realizes that trace, giving the stated sign.

To see that every component occurs, take $x\in A\setminus C$.
The signed-clause presentation of $C$ supplies a minimal nonface $T$ of
$\operatorname{Sh}(C)$ with $x|_T=\tau_T$.  All proper subsets of $T$
are shattered by both $C$ and $A$.  If $T$ were not shattered by $A$,
the two projections $C|_T\subseteq A|_T$ would both have
$2^{|T|}-1$ members.  They would be equal, contradicting the presence
of $x|_T$ in $A|_T$.  Hence $T\in\mathcal I$ and $x\in A(T,\tau_T)$.

Finally, bits constant on each component's clause group produce $X_b$
and are admissible by the construction.  Conversely, any admissible
signing with the displayed fixed data defines an ample family $Y$ with
that shadow.  The old clauses imply that its contraction $A'$ is
contained in $A$.  Every point in its reduction avoids each
$(T,\tau_T)$ for $T\in\mathcal I$, because otherwise one of its two
lifts is forbidden.  Every minimal nonface of $\operatorname{Sh}(C)$
outside $\mathcal I$ is a minimal nonface of $\operatorname{Sh}(A)$;
the two projections have the same unique missing trace, by the preceding
cardinality argument.  Thus the old clauses and these additional
conditions imply all clauses defining $C$, so $C'\subseteq C$.  Since
\[
 |Y|=|A'|+|C'|=|A|+|C|,
\]
both inclusions are equalities.  The converse already proved identifies
$Y=X_b$ and forces equality of the bits within each clause group.
\end{proof}

The theorem gives an intrinsic recursion for ample families: choose nested
ample families on fewer coordinates and label their complementary components.
Fixed coordinates may be added separately.  Every active coordinate of a
nonempty ample family has a nonempty reduction, by connectedness,
so the converse supplies exhaustive recovery.  This construction of the
ambient ample class does not identify the forbidden labelings.

\begin{corollary}[Constructing realizations in a fixed coordinate order]
\label{cor:damp-component-reflection-generation}
Fix a simplicial complex $\Sigma\subseteq2^F$ and an ordering
$e_1,\ldots,e_n$ of $F$.  Start with the down-set
\[
 D_\Sigma=\{x\in Q_F:\{f:x_f=1\}\in\Sigma\}.
\]
At stage $j$, form the contraction $A$ and reduction $C$ in
$e_j$ of the current family.  Its unique lifts outside $C$ are all in
layer zero.  Independently choose a label for every component of
$A\setminus C$ and apply the construction of
Theorem~\ref{thm:damp-component-expansion-signs}.  If $C$ is empty,
$e_j$ is constant and the choice is simply to retain or reflect that
constant coordinate.  After these $n$ stages, the possible outputs are
exactly all ample families with shadow $\Sigma$.

For the prescribed coordinate order, every output determines its entire
sequence of intermediate families and component choices uniquely.  In
particular, the graph joining two realizations when one component is
reflected in one coordinate is connected; reflection of a constant
coordinate is included as a move.
\end{corollary}

\begin{proof}
Define $D_e(H)$ by moving every unique lift in coordinate $e$ to layer
zero and retaining every double lift.  The theorem shows that $D_e(H)$
is ample with the same shadow and that its signed-clause presentation
is obtained by setting the $e$-bit to one in every clause containing $e$.
Other bits are unchanged.  For a constant coordinate the same assertion
holds for its singleton clause, and the operation is a coordinate
reflection or the identity.

Consequently the $D_e$ commute.  Applying every one of them makes every
clause the all-one trace, which gives exactly $D_\Sigma$.  This is the
signed-clause form of simultaneous shift normalization.  A family for
which all clause bits in coordinate $e$ are one is downward closed in
$e$: changing that coordinate from one to zero cannot create a forbidden
trace.  Thus all unique lifts at a stage involving an as-yet unprocessed
coordinate lie in layer zero, as asserted.

Reversing the normalization steps, in the prescribed order, reconstructs
any target realization by the component choices in the theorem.  Each
individual component may be reflected separately; all intervening families
are ample with the same shadow.  Conversely, every choice allowed by the
recipe preserves those two properties.  If $H$ is the final output,
the intermediate family after stage $j$ is
\[
 D_{e_{j+1}}\cdots D_{e_n}(H).
\]
Its signed clauses retain exactly the final bits in processed coordinates
and have all-one bits in the remaining coordinates.  These intermediate
families are therefore uniquely determined, as are their recovered
component labels.  Reversing a path to $D_\Sigma$ proves connectivity.
\end{proof}

\begin{proposition}[Alternative repairs in component labelings]
\label{prop:damp-component-label-repair-disjunction}
Let $C\subseteq S\subseteq Q_F$ be nonempty ample families, with $C$
corner-peelable and $S$ not corner-peelable.  Choose distinct points
$a_1,\ldots,a_r$ outside $S$, where $r\ge1$, and put
$R=\{a_1,\ldots,a_r\}$.  Assume:
\begin{enumerate}
 \item $A=S\cup R$ is ample;
 \item each $\{a_i\}$ is a connected component of $A\setminus C$;
 \item every $S\cup\{a_i\}$ is corner-peelable.
\end{enumerate}
In the expansion of Theorem~\ref{thm:damp-component-expansion-signs}, give
all components contained in $S\setminus C$ label zero and give
$\{a_i\}$ label $\varepsilon_i$.  Then
\[
 X_{\varepsilon}\in\Damp
 \quad\Longleftrightarrow\quad
 \varepsilon_i=0\text{ for at least one }i.
\]
If all $\varepsilon_i=1$, the nonpeelable family $S$ is a proper section
of $X_{\varepsilon}$, so this negative expansion is not a forbidden minor.
\end{proposition}

\begin{proof}
Every family $S\cup U$, with $U\subseteq R$, is a component interpolant
between $C$ and $A$.  It is ample: assign its selected components one
label and the others the opposite label in
Theorem~\ref{thm:damp-component-expansion-signs}, and take the appropriate
section.  Thus deleting any available $a_i$ preserves ampleness and is a
corner deletion.  If $U\ne\varnothing$, delete all but one of its members
and finish with the given peeling of $S\cup\{a_i\}$.  In particular,
$A$ and every such $S\cup U$ are corner-peelable.

Suppose some $\varepsilon_i=0$, and let $U$ consist of these tips.
Delete the unique lifts of the tips with label one.  At each step the
remaining family is an ample expansion over a component interpolant
$S\cup U'$ and the same $C$, by the theorem, so each deletion is a
corner deletion.  The remainder is the peripheral expansion with base
$S\cup U$ and site $C$.  Both peel, and
Proposition~\ref{prop:damp-exact-expansion-product-tests} completes its
peeling.  If instead all labels are one, the zero section is exactly
$S$; minor closure of corner peelability proves the negative assertion
and also excludes pc-minimality.
\end{proof}

\begin{remark}[Ample sign choices and peelable sign choices differ]
\label{ex:damp-component-signs-nonlinear-peeling}
The fixed 291-concept cornerless family $S$ has two independent repairing
tips $91$ and $157$, with incident supports $102$ and $390$ in binary
encoding.  Let $C$ be the union of their punctured incident four-cubes.
The fixed replay verifies that $|C|=26$, that $C$ peels, and that
$A=S\cup\{91,157\}$ is ample and peels.  Its complement of $C$ has
components of orders $265,1,1$, the last two being the tips.

There are twelve new-coordinate clauses: ten belong to the first component
and one to each tip.  All eight component labelings give ample expansions
of order $319$.  Precisely six peel.  With the first component labelled
zero, the labelings $(0,1,0)$ and $(0,0,1)$ peel, whereas their
coordinatewise maximum $(0,1,1)$ does not.  Thus the equality constraints
describing admissible expansion signs do not describe the peelable subset.
The negative examples have the proper nonpeelable section $S$, as the
proposition requires; they do not test a criterion restricted to families
whose proper minors all peel.

The independent checker \path{verify_component_expansion_signs.py} also
tests all $2244$ nested nonempty ample pairs through $Q_3$ and all $9656$
new-clause bit assignments.  Exactly $5384$ assignments satisfy the
component equalities and the asserted shadow formula.  It also normalizes
these small families in every coordinate order and replays the inverse
component reflections.  The fixed larger
example is checked by full positive orders and a cornerless proper-section
certificate, without a peeling search for negative answers.  Inputs,
orders, signs and scope are in
\path{component_expansion_signs_verification.json}.  These finite checks
support the examples and the implementation; the universal claims use
the preceding proofs.
\end{remark}



\begin{lemma}[Corners as unique facet traces]
\label{lem:damp-corner-unique-facet-trace}
Let $H\subseteq Q_F$ be ample, let
$\Sigma=\operatorname{Sh}(H)$.  Then $x\in H$ is a corner if and only if there is a facet
$T$ of $\Sigma$ such that $x|_T$ is realized by no other member of $H$.
\end{lemma}

\begin{proof}
Deleting $x$ loses a shattered face $D\in\Sigma$ exactly when
$x|_D$ is realized by no other member of $H$.  The ample
corner-deletion criterion says that $x$ is a corner exactly when
$H\setminus\{x\}$ is ample.  By the Sandwich inequalities, this happens
exactly when deletion loses one shattered face: if none is lost, then
$H\setminus\{x\}$ still has $|H|$ shattered faces and hence does not
satisfy the Sandwich equality for its $|H|-1$ members, while losing two
or more faces would violate the upper Sandwich inequality
$|H\setminus\{x\}|\le
 |\operatorname{Sh}(H\setminus\{x\})|$.

It remains to identify the possible unique face.  If $x|_D$ is unique
and $D\subseteq T\in\Sigma$, then $x|_T$ is also unique.  If $D$ is not
a facet, choose a facet $T$ with $D\subsetneq T$; there are then at
least two uniquely realized faces, namely $D$ and $T$.  Therefore exactly
one face is lost if and only if the lost face is a facet, proving the
claim.
\end{proof}

\begin{proposition}[A simplicial maximal cube exposes a corner]
\label{prop:damp-simplicial-maximal-cube-corner}
Let $X\subseteq Q_E$ be a nonempty cube family, and let
$\Gamma(X)$ be the graph whose vertices are the inclusion-maximal cubes
of $X$, with two such cubes adjacent when they intersect.  If a maximal
cube $C$ is a simplicial vertex of $\Gamma(X)$, then $C$ contains a
vertex that belongs to no other maximal cube of $X$.  In particular,
$X$ has a corner.  Consequently, if $X$ is cornerless, then
$\Gamma(X)$ has no simplicial vertex; in particular, it is nonchordal
and contains an induced cycle of length at least four.
\end{proposition}

\begin{proof}
Recall that a graph vertex is simplicial when its neighbors form a
clique.  First suppose that $C$ has no neighbor in $\Gamma(X)$.  Every
vertex of $C$ then belongs to no other maximal cube and is a corner.

Suppose next that $C$ has neighbors $C_1,\ldots,C_k$.  The cubes in
$\{C,C_1,\ldots,C_k\}$ intersect pairwise: each $C_i$ meets $C$, and
the $C_i$ meet one another because $C$ is simplicial.  Subcubes of a
hypercube have the Helly property.  Indeed, each subcube is described by
fixing some coordinates, and pairwise intersection says precisely that
none of these fixed-coordinate prescriptions conflict.  Hence there is
a vertex
\[
 z\in C\cap C_1\cap\cdots\cap C_k.
\]
Let $y$ be the antipode of $z$ in $C$, obtained by reversing every
coordinate that is free in $C$.  For every $i$, the intersection
$C\cap C_i$ is a proper subcube of $C$: otherwise $C\subseteq C_i$,
contrary to the distinctness and maximality of $C$ and $C_i$.  Since
$C\cap C_i$ is a proper subcube containing $z$, it fixes at least one
free coordinate of $C$ to its value at $z$ and therefore does not
contain $y$.  Thus no neighbor $C_i$ contains $y$, while a nonneighbor
is disjoint from $C$ and cannot contain $y$ either.  Hence $C$ is the
unique maximal cube containing $y$, so $y$ is a corner.

It follows that the maximal-cube intersection graph of a cornerless
family has no simplicial vertex.  Every finite chordal graph has a
simplicial vertex, whereas a nonchordal graph contains an induced cycle
of length at least four.  This proves the final assertions.
\end{proof}

\begin{lemma}[Square-free edges are bridges]
\label{lem:damp-square-free-edge-bridge}
Let \(X\) be an ample family.  Every cycle in the one-inclusion graph of
\(X\), regarded as an edge chain over \(\mathbb F_2\), is a sum of
boundaries of squares contained in \(X\).  Consequently, every edge of
\(X\) that belongs to no square is a bridge.
\end{lemma}

\begin{proof}
We prove the first assertion by induction on the number of active
coordinates.  It is immediate when there is at most one such coordinate.
Fix an active coordinate \(e\), and write \(X_0,X_1\) for the two
\(e\)-sections, with the \(e\)-coordinate deleted.  Both sections are
ample.  Their intersection

\[
 I=X_0\cap X_1
\]

is also ample.  Indeed, the contraction \(X_0\cup X_1\) is ample, and the
two-layer shattering formula gives
\[
 |X_0|+|X_1|
 =|\operatorname{Sh}(X_0\cup X_1)|
  +|\operatorname{Sh}(X_0)\cap\operatorname{Sh}(X_1)|.
\]
The first term is \(|X_0\cup X_1|\), so the second is \(|I|\).
Since
\[
 \operatorname{Sh}(I)
 \subseteq\operatorname{Sh}(X_0)\cap\operatorname{Sh}(X_1)
\]
and Sauer's inequality gives \(|\operatorname{Sh}(I)|\ge|I|\), equality
holds and \(I\) is ample.  It is therefore connected whenever it is
nonempty.

Let \(Z\) be a cycle of \(X\).  Its \(e\)-edges occur in an even number.
Pair them arbitrarily.  For a paired pair with projected endpoints
\(u,v\in I\), choose a path from \(u\) to \(v\) in \(I\).  The two lifts
of this path, together with the two paired \(e\)-edges, form a ladder.
Its boundary is the sum of the boundaries of the squares corresponding
to the edges of the chosen path.  Adding all these ladder boundaries to
\(Z\) cancels every \(e\)-edge.  The remaining even edge chain is the
disjoint sum of an even edge chain in the lower section and one in the
upper section.  Decompose each into cycles and apply induction inside
\(X_0\) and \(X_1\).  This expresses \(Z\) as a sum of square boundaries.

Now suppose that an edge \(a\) belongs to no square but is not a bridge.
Then \(a\) lies on a graph cycle \(Z\).  In the square-boundary expression
for \(Z\), the coefficient of \(a\) is zero because no summand contains
\(a\), whereas its coefficient in \(Z\) is one.  This contradiction proves
that \(a\) is a bridge.
\end{proof}

\begin{lemma}[Local trace covers]
\label{lem:damp-local-maximal-cube-trace-cover}
Let $X\subseteq Q_E$ be a nonempty cube family, let $C$ be a maximal cube
of $X$, and let $C_1,\ldots,C_k$ be the other maximal cubes that meet
$C$.  Then
\begin{equation}
\operatorname{Cor}(X)\cap C
 =C\setminus\bigcup_{i=1}^k(C\cap C_i).
 \label{eq:damp-local-corner-complement}
\end{equation}
Moreover, for distinct $i,j$, one has
\begin{equation}
C_i\cap C_j\ne\varnothing
 \quad\Longleftrightarrow\quad
 (C\cap C_i)\cap(C\cap C_j)\ne\varnothing.
 \label{eq:damp-local-neighbor-intersection}
\end{equation}
Consequently, if $X$ is cornerless, the nonempty proper subcubes
$C\cap C_i$ cover $C$, and the graph induced by the neighbors of $C$ in
$\Gamma(X)$ is their intersection graph.  After deleting repeated and
redundant traces, one obtains an irredundant cover of $C$ by proper
subcubes.
\end{lemma}

\begin{proof}
A vertex $x\in C$ is not a corner precisely when it lies in some maximal
cube other than $C$.  Any such cube meets $C$ and is one of the $C_i$,
which proves \eqref{eq:damp-local-corner-complement}.

If the two traces on the right-hand side of
\eqref{eq:damp-local-neighbor-intersection} meet, then $C_i$ and $C_j$
plainly meet.  Conversely, if $C_i$ and $C_j$ meet, then the three
subcubes $C,C_i,C_j$ intersect pairwise.  The hypercube Helly property
used in Proposition~\ref{prop:damp-simplicial-maximal-cube-corner}
gives $C\cap C_i\cap C_j\ne\varnothing$.  This proves the equivalence.
Each trace is proper because $C$ and $C_i$ are distinct maximal cubes.
When $X$ is cornerless, \eqref{eq:damp-local-corner-complement} says that
the traces cover $C$.  Finally choose an inclusion-minimal subfamily of
the distinct traces that still covers $C$; every member of the resulting
cover has a private vertex, so the cover is irredundant.
\end{proof}

\begin{proposition}[The rank-three local trace catalogue]
\label{prop:damp-rank-three-local-trace-catalogue}
Up to coordinate permutations and reversals, the numbers of irredundant
covers of $Q_r$ by nonempty proper subcubes are as follows:
\begin{equation}
\begin{array}{c|c|c|c|c}
r&\text{proper subcubes}&\text{labelled covers}&\text{orbits}
  &\text{cover-size range}\\ \hline
1&2 &1   &1 &2\\
2&8 &11  &4 &2\text{--}4\\
3&26&869 &43&2\text{--}8
\end{array}
\label{eq:damp-rank-three-local-trace-catalogue}
\end{equation}
Consequently, Lemma~\ref{lem:damp-local-maximal-cube-trace-cover}
reduces the local cornerlessness condition for every family of VC
dimension at most three to forty-eight fixed cover types.
\end{proposition}

\begin{proof}[Computer-assisted proof]
A proper subcube of $Q_r$ is encoded by a nonempty set of fixed
coordinates and their prescribed bits, giving $3^r-1$ possibilities.
The verifier branches on the first uncovered cube vertex, retains a cover
only when every selected subcube has a private vertex, and then takes the
lexicographically least image under all $2^r r!$ cube automorphisms.  The
three rows of \eqref{eq:damp-rank-three-local-trace-catalogue} are the
resulting exhaustive counts.

The generator and output are
\[
 \begin{gathered}
 \texttt{explore\_damp\_local\_maximal\_cube\_covers.py},\\
 \texttt{damp\_local\_maximal\_cube\_covers.json},
 \end{gathered}
\]
and the output digest is
\[
 \texttt{c570301788f52c46a536d0905c627cc05dfcf530625afc3124f6971d7613e931}.
\]
\end{proof}

\begin{proposition}[Triangle-carrier trees in VC dimension three]
\label{prop:damp-maximum-vc-three-carrier-trees}
Let $H\subseteq Q_E$ be ample of VC dimension at most three, put
$\Sigma=\operatorname{Sh}(H)$, and fix a two-set $S\in\Sigma$.  Define
\[
 L_S=\{e\in E\setminus S:S\cup\{e\}\in\Sigma\}
\]
and the $S$-cube carrier family by
\begin{equation}
\mathcal K_S(H)
 =\{y\in Q_{E\setminus S}:Q_S\times\{y\}\subseteq H\}.
 \label{eq:damp-maximum-carrier-family}
\end{equation}
After its constant coordinates are deleted, $\mathcal K_S(H)$ is a
maximum VC-dimension-one class on $L_S$.  Its one-inclusion graph is a
tree with $|L_S|+1$ vertices and exactly one edge in each direction
$e\in L_S$.

The edge in direction $e$ corresponds to the unique three-cube
$C_{S\cup\{e\}}$ of $H$ with support $S\cup\{e\}$, and its endpoints
correspond to the two opposite $S$-facets of that cube.  An endpoint is
incident with another edge of the carrier tree if and only if the
corresponding $S$-facet lies in another maximal three-cube of $H$.
Consequently, both opposite $S$-facets of
$C_{S\cup\{e\}}$ are covered by other maximal cubes if and only if the
$e$-edge has no leaf endpoint in the carrier tree.
In particular, if $H$ is maximum of VC dimension three, then
$L_S=E\setminus S$, so the carrier tree has $|E|-1$ vertices and
$|E|-2$ edges.
\end{proposition}

\begin{proof}
For $A\subseteq E\setminus S$, the carrier family strongly shatters $A$
if and only if $H$ strongly shatters $S\cup A$.  Likewise, if
$\mathcal K_S(H)$ shatters $A$, then $H$ shatters $S\cup A$: for every
$A$-trace supplied by the carrier, the accompanying full $S$-cube
supplies all traces on $S\cup A$.  Since $H$ is ample, these observations
give
\[
 \operatorname{SSh}(\mathcal K_S(H))
 =\operatorname{Sh}(\mathcal K_S(H))
 =\{\varnothing\}\cup\{\{e\}:e\in L_S\}.
\]
Thus the carrier is ample and has $1+|L_S|$ members.  A coordinate
outside $L_S$ is not shattered and is therefore constant on the carrier;
delete all such coordinates.

Its one-inclusion graph is connected because every ample family is an
isometric cube family.  Choose a spanning tree.  Every coordinate of
$L_S$ is shattered and hence occurs on an edge; moreover every
path between the two coordinate semicubes must use an edge in that
direction.  The spanning tree therefore has at least one edge in each of
the $|L_S|$ directions.  It has exactly $|L_S|$ edges, so it uses every
direction exactly once.  There can be no further edge of the
one-inclusion graph: if an additional $e$-edge joined $u$ and $v$, the
tree path from $u$ to $v$ would have to change only coordinate $e$.
Because every tree-edge direction occurs at most once, that path would be
the existing $e$-edge itself.  Hence the one-inclusion graph equals the
spanning tree and has exactly one edge of each direction.

By \eqref{eq:damp-maximum-carrier-family}, an $e$-edge of the carrier is
precisely a full cube of $H$ with support $S\cup\{e\}$.  Its endpoints
are its two $S$-facets.  Since the carrier has one $e$-edge, this
three-cube is unique.  A second carrier edge is incident with one of its
endpoints precisely when the corresponding $S$-cube extends in a second
direction, or equivalently lies in another maximal three-cube.  The last
assertion is now exactly the statement that neither endpoint of the
$e$-edge is a leaf.
\end{proof}

\begin{proposition}[Reconstruction from anchored facet cubes]
\label{prop:damp-facet-cube-reconstruction}
Let $\mathcal F\subseteq2^E$ be a nonempty antichain, let
\[
 \Sigma=\{S\subseteq E:S\subseteq T\text{ for some }T\in\mathcal F\},
\]
and, for each $T\in\mathcal F$, choose an affine cube $C_T\subseteq Q_E$
with support $T$.  Put $H=\bigcup_{T\in\mathcal F}C_T$.  Then
\begin{equation}
H\text{ is ample with }\operatorname{Sh}(H)=\Sigma
 \quad\Longleftrightarrow\quad
 |H|=|\Sigma|.
 \label{eq:damp-facet-cube-sandwich-test}
\end{equation}
When these conditions hold, the cubes $C_T$, for $T\in\mathcal F$, are
exactly the maximal cubes of $H$.  Consequently,
\begin{equation}
H\text{ is cornerless}
 \quad\Longleftrightarrow\quad
 |\{T\in\mathcal F:x\in C_T\}|\ne1
 \quad\text{for every }x\in Q_E.
 \label{eq:damp-facet-cube-cornerless-test}
\end{equation}
Thus a fixed-rank minimum-order search is a static anchored-cube union
problem: it requires one cardinality equality and forbids incidence
multiplicity one.
\end{proposition}

\begin{proof}
Every face of $\Sigma$ is strongly shattered by one of the selected
cubes, so
\[
 |\Sigma|\le|\operatorname{SSh}(H)|\le|H|.
\]
If $|H|=|\Sigma|$, equality holds throughout.  Equality in the lower
Sandwich inequality makes $H$ ample, and
$\operatorname{SSh}(H)=\Sigma$ then gives
$\operatorname{Sh}(H)=\Sigma$.  The converse is the Sandwich equality.

It remains to identify the maximal cubes.  For any $S\in\Sigma$, let
$\mathcal K_S(H)$ be the carrier family defined as in
\eqref{eq:damp-maximum-carrier-family}.  The argument in the proof of
Proposition~\ref{prop:damp-maximum-vc-three-carrier-trees}, without the
rank-three specialization, gives
\[
 \operatorname{Sh}(\mathcal K_S(H))
 =\operatorname{SSh}(\mathcal K_S(H))
 =\operatorname{link}_{\Sigma}(S).
\]
If $S=T$ is a facet, this link consists only of the empty face, so the
carrier has one member: there is a unique $T$-cube, namely $C_T$.  If
$S$ is not a facet, its link contains a singleton.  The carrier is then
a nontrivial connected ample family, so every one of its vertices is
incident with an edge along some carrier path.  Equivalently, every
$S$-cube of $H$ extends to a larger cube and is not maximal.  Hence the
selected facet cubes are exactly the maximal cubes.  A vertex is a corner
precisely when it belongs to exactly one maximal cube, which proves
\eqref{eq:damp-facet-cube-cornerless-test}.
\end{proof}

\begin{corollary}[Triangle cap at the current lower endpoint]
\label{cor:damp-order-twenty-seven-triangle-cap}
Let $H$ be an irredundant ample family of VC dimension three and order
$27$, and put $d=|E|$.  If $f_2$ is the number of triangular faces of
$\operatorname{Sh}(H)$, then
\begin{equation}
f_2\le
 \begin{cases}
 8,&d=8,\\
 7,&d=9,\\
 7,&d=10.
 \end{cases}
 \label{eq:damp-order-twenty-seven-triangle-cap}
\end{equation}
In particular, each of the three candidate profiles in
Theorem~\ref{prop:damp-minimum-order-profile} is an anchored-cube union
with at most eight maximal three-cubes.
\end{corollary}

\begin{proof}
Write $f_1$ for the number of two-element faces.  Irredundancy and the
Sandwich equality give
\[
 27=1+d+f_1+f_2.
\]
The Kruskal--Katona theorem \cite{Kruskal63,Katona68} says that the pair
shadow of $f_2$ triples has at least as many pairs as the first $f_2$
triples in colexicographic order.  For $f_2=8$ this minimum is $10$, while
for $f_2=9$ it is still $10$; for $f_2=7$ it is $9$, while for $f_2=8$
it is $10$.  Thus
\[
\begin{array}{c|c|c}
d&f_1+f_2=26-d&\text{first impossible }f_2\\ \hline
8&18&9\quad(10+9>18),\\
9&17&8\quad(10+8>17),\\
10&16&8\quad(10+8>16).
\end{array}
\]
This proves the displayed bounds.  The final assertion follows from
Proposition~\ref{prop:damp-facet-cube-reconstruction}.
\end{proof}

\begin{lemma}[Four-coordinate compatibility]
\label{lem:damp-four-coordinate-forbidden-tree}
Let $\Sigma$ be a two-dimensional simplicial complex with complete
one-skeleton such that every four-set contains at most two triangular
faces, let $H\subseteq Q_F$ be ample with
$\operatorname{Sh}(H)=\Sigma$, and let
$\boldsymbol\sigma=(\sigma_I)_{I\in\mathcal N(\Sigma)}$ be its signing.
Fix a four-set $U\subseteq F$.  For each missing triple $I\subset U$, put
\[
 E_U(I)=\{z\in Q_U:z|_I=\sigma_I\}.
\]
Then the edges $E_U(I)$, indexed by the missing triples in $U$, have
connected union.  In particular:
\begin{enumerate}
 \item if $U$ contains exactly two missing triples $I,J$, then
       $\sigma_I|_{I\cap J}=\sigma_J|_{I\cap J}$;
 \item if $U$ contains exactly three missing triples and
       $e$ belongs to all three, then their three omitted traces have the
       same value at $e$.
\end{enumerate}
\end{lemma}

\begin{proof}
Let $t$ be the number of triangular faces of $\Sigma$ contained in $U$.
There are $m=4-t$ missing triples in $U$.  Projection preserves ampleness,
and the shattered sets of $H{\upharpoonright}U$ are exactly
$\Sigma[U]=\{D\in\Sigma:D\subseteq U\}$.  Hence
\[
 |H{\upharpoonright}U|=|\Sigma[U]|=1+4+6+t=11+t.
\]
The minimal nonfaces of $\Sigma[U]$ are precisely its $m$ missing triples,
with the same omitted traces as in $H$: projection cannot create a trace
that was absent from $H|_I$, and both projections omit exactly one
$I$-trace by ampleness.  Applying
Theorem~\ref{thm:damp-fixed-shadow-signed-clutter} to the ample projection
therefore identifies the complement of $H{\upharpoonright}U$ with
$\bigcup_{I\in\mathcal N(\Sigma),\ I\subset U}E_U(I)$.  Thus this union has
\[
 16-(11+t)=5-t=m+1
\]
vertices.

Every $E_U(I)$ is an edge in direction $U\setminus I$, so the $m$ edges
have distinct directions.  A cycle in a hypercube uses every direction an
even number of times; consequently these edges form a forest.  A forest
with $m$ edges and $m+1$ vertices is connected, proving the first claim.

Two cube edges in distinct directions meet exactly when their prescribed
traces agree on the common coordinates.  This gives item~1.  For item~2,
the three edges form a tree.  Along each pair of incident edges the
prescribed traces agree at their common coordinate $e$, so connectedness
propagates one value of $\sigma_I(e)$ through all three edges.
\end{proof}

\begin{definition}[Coordinate-propagation graph]
\label{def:damp-coordinate-propagation-graph}
Under the hypotheses of
Lemma~\ref{lem:damp-four-coordinate-forbidden-tree}, fix $e\in F$.
The graph $P_e(\Sigma)$ has as vertices the minimal nonfaces
$I\in\mathcal N(\Sigma)$ containing $e$.  Join $I$ and $J$ if either
\begin{enumerate}
 \item $|I\cup J|=4$ and $I,J$ are the only missing triples in
       $I\cup J$, or
 \item $|I\cup J|=4$ and $I\cup J$ contains exactly three missing
       triples, all containing $e$.
\end{enumerate}
\end{definition}

\begin{proposition}[Rigidity from coordinate propagation]
\label{prop:damp-coordinate-propagation-rigidity}
Let $\Sigma$ be a two-dimensional simplicial complex with complete
one-skeleton such that every four-set contains at most two triangular
faces.  If $P_e(\Sigma)$ is connected whenever it is nonempty, then every
ample family $H\subseteq Q_F$ with
$\operatorname{Sh}(H)=\Sigma$ is a cube translate of the canonical face
family
\[
 \{\mathbf1_D:D\in\Sigma\}.
\]
\end{proposition}

\begin{proof}
Let $\boldsymbol\sigma$ be the signing of $H$.  Each edge of
$P_e(\Sigma)$ equates the two values of the corresponding omitted traces at
$e$ by Lemma~\ref{lem:damp-four-coordinate-forbidden-tree}.  Connectivity
therefore gives a bit $h_e$ such that $\sigma_I(e)=h_e$ for every minimal
nonface $I$ containing $e$.  If no minimal nonface contains $e$, choose
$h_e$ arbitrarily.  With $h=(h_e)_{e\in F}$, we have
$\sigma_I=h|_I$ for every $I\in\mathcal N(\Sigma)$.  Let
$g=h+\mathbf1_F$, with addition in $Q_F$.  Translating $Q_F$ by $g$
changes every omitted trace to the all-one trace.  The resulting satisfying
family is
\[
 \{x\in Q_F:\operatorname{supp}(x)\text{ contains no member of }
                  \mathcal N(\Sigma)\}
 =\{\mathbf1_D:D\in\Sigma\},
\]
which proves the claim.
\end{proof}

\begin{remark}[A three-coordinate example]
Let \(F=\{1,2,3\}\) and
\[
 \Sigma=\{D\subseteq F:\{1,2\}\nsubseteq D\}.
\]
Then \(\mathcal N(\Sigma)=\{\{1,2\}\}\).  Choosing the omitted trace
\(\sigma_{\{1,2\}}=(1,1)\) gives
\[
 H(\Sigma,\boldsymbol\sigma)
 =Q_F\setminus\{(1,1,0),(1,1,1)\}.
\]
Both this family and \(\Sigma\) have six members, so the signing is
admissible and the theorem makes \(H\) ample with shattered complex
\(\Sigma\).  The same count works for any of the four choices of the
omitted \(\{1,2\}\)-trace.  Reorienting the single clause can therefore
change four concepts, which is why clause reorientation alone is
coarser than adjacency in the mutation graph below.
\end{remark}

\begin{definition}[Oriented one-hole cover]
\label{def:damp-oriented-one-hole-cover}
Let \(E\) be finite.  An \emph{oriented one-hole datum} is a pair
\((y,B)\) with \(y\in Q_E\) and \(B\subseteq E\), \(|B|\ge2\).  Its
punctured affine cube and distinguished antipode are
\begin{equation}
 \mathsf P(y,B)=\{y^U:\varnothing\ne U\subseteq B\},
 \qquad a(y,B)=y^B.
 \label{eq:damp-oriented-one-hole-cube}
\end{equation}
For a nonempty finite family \(\mathcal W\) of such data, put
\[
 H(\mathcal W)=\bigcup_{(y,B)\in\mathcal W}\mathsf P(y,B).
\]
We call \(\mathcal W\) an \emph{admissible oriented one-hole cover} if
\begin{enumerate}
 \item \(y\notin H(\mathcal W)\) for every \((y,B)\in\mathcal W\);
 \item the distinguished antipodes cover the constructed family:
       \begin{equation}
        H(\mathcal W)=\{a(y,B):(y,B)\in\mathcal W\};
        \label{eq:damp-one-hole-antipode-cover}
       \end{equation}
 \item the Sandwich equality
       \begin{equation}
        |\operatorname{Sh}(H(\mathcal W))|=|H(\mathcal W)|
        \label{eq:damp-one-hole-sandwich}
       \end{equation}
       holds.
\end{enumerate}
All three conditions are finite conditions on the displayed punctured
cubes; none asks for a corner or a peeling order.
\end{definition}

\begin{theorem}[Intrinsic one-hole construction of cornerless ample families]
\label{thm:damp-one-hole-cornerless-construction}
Every admissible oriented one-hole cover \(\mathcal W\) constructs a
nonempty cornerless ample family \(H(\mathcal W)\).  Conversely, every
nonempty cornerless ample family \(H\subseteq Q_E\) is obtained in this
way.

More precisely, for \(x\in H\) define its local cube-support complex by
\begin{equation}
 \Lambda_H(x)=
 \{A\subseteq E:\{x^C:C\subseteq A\}\subseteq H\}.
 \label{eq:damp-local-cube-support-complex}
\end{equation}
Then the canonical recovered cover is
\begin{equation}
 \mathcal W(H)=
 \{(x^B,B):x\in H,\ B\in\mathcal N(\Lambda_H(x)),\ |B|\ge2\}.
 \label{eq:damp-canonical-one-hole-cover}
\end{equation}
Thus oriented one-hole covers give an explicit, order-free, and
canonically recoverable parametrization of all cornerless ample families.
They do not by themselves characterize which such families are
pc-minor-minimal outside \(\Damp\).
\end{theorem}

\begin{proof}
Let \(\mathcal W\) be admissible and put \(H=H(\mathcal W)\).
Condition~\eqref{eq:damp-one-hole-sandwich} and the Sandwich inequalities
make \(H\) ample.  Fix \(x\in H\).  By
\eqref{eq:damp-one-hole-antipode-cover}, there is \((y,B)\in\mathcal W\)
with \(x=y^B\).  Every proper subset \(C\subsetneq B\) satisfies
\[
 x^C=y^{B\setminus C}\in\mathsf P(y,B)\subseteq H,
 \]
whereas \(x^B=y\notin H\).  Hence \(B\) is a nontrivial minimal nonface
of \(\Lambda_H(x)\).  Equivalently, the cubes of \(H\) through \(x\)
cannot all lie in one maximal cube: such a cube would contain every
direction of \(B\) and therefore also \(x^B=y\).  Thus every concept lies
in at least two maximal cubes, so \(H\) is cornerless.

Conversely, let \(H\) be nonempty, ample, and cornerless.  The complex
\(\Lambda_H(x)\) is not a simplex for any \(x\in H\), because its maximal
faces are precisely the supports of the maximal cubes through \(x\).
It therefore has a minimal nonface \(B\) of order at least two.  For every
proper \(C\subsetneq B\), one has \(x^C\in H\), while \(x^B\notin H\).
Putting \(y=x^B\), we obtain
\[
 \mathsf P(y,B)
 =\{x^{B\setminus U}:\varnothing\ne U\subseteq B\}
 \subseteq H,
 \qquad a(y,B)=x.
\]
Taking all such minimal nonfaces for all \(x\) gives
\(\mathcal W(H)\).  Every punctured cube in this family lies in \(H\),
and every \(x\in H\) occurs as its distinguished antipode.  Hence
\(H(\mathcal W(H))=H\); the missing-tip and antipode-cover conditions
hold by construction, and the Sandwich equality follows from ampleness.
This proves admissibility and the recovery formula.
\end{proof}

The one-hole construction does not become a forbidden-minor construction
by selecting its pc-minimal cornerless outputs.  A cornerless ample class
can have a corner in every nonempty proper pc-minor while still having a
proper minor that is not corner-peelable.  The following gluing statement
isolates the distinction.

\begin{proposition}[Cornerlessness can be minimal without forbiddenness]
\label{prop:damp-cornerless-minor-minimal-gluing}
Let \(S\subseteq Q_F\) be an ample forbidden pc-minor whose only
corner is \(0\), and let \(P\subseteq Q_E\) be a nontrivial
corner-peelable ample family whose only corner is \(0\).
Assume that every proper root-preserving pc-minor of each family
peels to the image of \(0\).  The coordinate sets are disjoint and
the embeddings irredundant.  Set
\[
 X=S\vee P
   =(S\times\{0\}^E)\cup(\{0\}^F\times P).
\]
Then \(X\) is ample and cornerless, every nonempty proper pc-minor
of \(X\) has a corner, and every non-corner-peelable pc-minor of
\(X\) contains \(S\) as a pc-minor.  In particular, \(S\) is the
unique forbidden pc-minor of \(X\), up to isomorphism, and \(X\)
has no cornerless forbidden pc-minor.
\end{proposition}

\begin{proof}
The vertex-gluing formula in
Proposition~\ref{prop:damp-vertex-gluing} proves ampleness.
A nonzero vertex has exactly the same incident cubes as in its own
piece, so it is not a corner.  Both nontrivial pieces are connected
and have a neighbor of \(0\).  Two such neighbors, one from each
piece, have no completing mixed square.  Thus \(0\) is not a
corner either.  Contracting all coordinates of \(P\) recovers
\(S\), so \(X\) is not forbidden.

We describe all pc-minors, rather than only the elementary ones.
A coordinate contraction acts within one of the two blocks and retains
the common root.  A convex restriction meeting both blocks away from
that root must contain the root, since every path between the blocks
passes through it.  Its intersection with each block is a convex
restriction of that block.  A nonempty restriction omitting the root
lies wholly in one block: it is then a proper pc-minor of \(S\)
or a pc-minor of \(P\).  Both are corner-peelable, and every
subsequent minor remains corner-peelable.  Such a nonempty minor
therefore has a corner.

All other minors retain the common root and have the form
\(S'\vee P'\), where \(S'\) and \(P'\) are root-preserving
pc-minors of the respective pieces.  Suppose first that \(S'\)
is proper.  It peels to \(0\) by hypothesis.  The family \(P'\)
is corner-peelable, since \(P\) is and this property is pc-minor
closed.  Peel \(S'\) down to \(0\), then peel \(P'\)
ordinarily.  This gives a complete peeling of \(S'\vee P'\).
Therefore every nonpeelable minor must retain \(S'=S\), and
contains \(S\) by contracting the other block.

It remains to check corners when \(S'=S\) and the minor is
proper.  Irredundancy forces \(P'\) to be proper.  If
\(P'=\{0\}\), the minor is \(S\), which has its corner
\(0\).  Otherwise \(P'\) peels to \(0\); the first vertex
of such a peeling is a nonzero corner of \(P'\), and remains a
corner of the union.  This proves the assertion about every
nonempty proper minor.  Finally, a forbidden minor containing
\(S\) must be isomorphic to \(S\), by its proper-minor
minimality.  Since \(S\) has a corner, no cornerless forbidden
minor occurs.
\end{proof}

\begin{corollary}[A cornerless class whose forbidden minor has one corner]
\label{cor:damp-cornerless-minor-minimal-not-forbidden}
There is an ample class \(X\subseteq Q_{24}\) with \(580\)
concepts, VC dimension three, and minimum graph degree four, such
that \(X\) is cornerless and every nonempty proper pc-minor has a
corner.  Its unique forbidden pc-minor has \(292\) concepts and
one corner.  Thus pc-minimality among cornerless ample classes is
strictly weaker than being a cornerless forbidden pc-minor, even
among graph three-cores of VC dimension three.
\end{corollary}

\begin{proof}
Let \(H\) be the \(292\)-concept class of
Theorem~\ref{thm:damp-empty-core-counterexample}, and let \(G\)
be the fixed \(291\)-concept family used there.  In its twelve-bit
integer notation, let \(\oplus\) denote coordinatewise addition
modulo two and put
\[
 S=\{x\oplus3959:x\in H\},\qquad
 P=\{x\oplus29:x\in(G\cup\{91\})\setminus\{61,125,93\}\}.
\]
The cited theorem makes \(S\) ample and forbidden, with \(0\)
as its only corner.  Remark~\ref{rem:damp-noncorner-rooted-factor}
shows that \(P\) is ample and ordinarily peelable, has \(0\)
as its only corner, and has all proper root-preserving minors
peelable to their roots.  Both families have VC dimension three.

The additional condition on \(S\) is verified directly: all
\(24\) elementary root-preserving minors have supplied corner
peelings ending at \(0\).  Root-preserving minor closure of
peeling to the root gives the condition for every proper rooted
minor.  The checker
\path{verify_damp_cornerless_minor_minimal_counterexample.py}
reconstructs both families, recomputes their shadows and corners,
and checks these \(24\) orders, the \(24\) rooted orders for
\(P\), its ordinary order, and the \(36\) ordinary proper-minor
orders for \(S\).  In particular, the new rooted assertion does
not rely on the ordinary orders ending at arbitrary vertices.

Use disjoint coordinate blocks and write the union as
\[
 X=S\cup\{2^{12}x:x\in P\}.
\]
The preceding proposition applies.  It gives
\(|X|=292+289-1=580\) and all the assertions about its proper
minors and forbidden minor.  The union has dimension \(24\) and
VC dimension three because its shattered supports lie in the two
blocks.  The checker verifies minimum degree four and cornerlessness
directly.  It also inspects all \(72\) elementary images: \(48\)
have checked complete peelings and \(24\) project onto \(S\),
so do not peel, although each has a corner.  These finite checks are
recorded in
\path{damp_cornerless_minor_minimal_counterexample_verification.json}.
The assertion for all proper minors follows from the proposition's
block decomposition, not from extrapolating the elementary corner
checks.
\end{proof}

This counterexample rules out selecting forbidden one-hole covers by
testing only that their nonempty proper pc-minors have corners.
It does not affect the minimum-cardinality argument: a globally smallest
nonempty cornerless ample family is still forbidden.  Every smaller ample
family peels, since otherwise maximal corner deletion would leave a
smaller nonempty cornerless ample family.  Proper pc-minors are smaller.
Minimality under pc-minors and global minimum cardinality are therefore
different requirements.

\paragraph{Restrictions on small covers.}
The antipode convention can be compressed by allowing one punctured cube
to protect several concepts.  For \(P_i=\mathsf P(y_i,L_i)\), write
\[
 I_i=\{y_i^T:T\subseteq L_i,\ |T|\ge2\},
 \qquad r_i=|L_i|.
\]
Consider a nonempty family of \(m\) distinct such pieces satisfying
\begin{equation}
 H=\bigcup_{i=1}^m P_i=\bigcup_{i=1}^m I_i,
 \qquad y_i\notin H\quad(1\le i\le m).
 \label{eq:damp-protected-cover}
\end{equation}
If \(|\operatorname{Sh}(H)|=|H|\), then \(H\) is ample and cornerless:
replace each \((y_i,L_i)\) by all \((y_i,T)\) with
\(T\subseteq L_i\), \(|T|\ge2\), and apply
Theorem~\ref{thm:damp-one-hole-cornerless-construction}.
The following restrictions use only~\eqref{eq:damp-protected-cover}.

For each \(q\in L_i\), choose a piece \(j\ne i\) with
\(y_i^{\{q\}}\in I_j\), and write \(i\to j\), with label
\(q_{ij}=q\).  Call \(j\) the chosen protector of that neighbour.
If \(D_{ij}\) is the set of coordinates on which \(y_i,y_j\) differ,
then
\begin{equation}
 q_{ij}\in L_i\setminus L_j,\qquad
 D_{ij}\setminus L_j=\{q_{ij}\},\qquad
 |D_{ij}\cap L_j|\ge2.
 \label{eq:damp-protector-equations}
\end{equation}
Indeed, if \(q\in L_j\), membership of \(y_i^{\{q\}}\) in \(I_j\)
would imply that \(y_i\) lies in the full carrier of \(P_j\).
Its absence from \(P_j\) would then force \(y_i=y_j\), contradicting
the distance-at-least-two condition defining \(I_j\).
Thus \(q\notin L_j\), and the other assertions follow by comparing
\(y_i\) and \(y_i^{\{q\}}\) outside \(L_j\).
In particular, one piece cannot protect two neighbours of the same tip,
so all chosen targets are distinct and \(r_i\le m-1\).

\begin{theorem}[Two supports of largest possible size]
\label{thm:damp-two-largest-supports}
Suppose~\eqref{eq:damp-protected-cover} holds and
\(r_1=r_2=m-1\).  Put \(a=q_{12}\), \(b=q_{21}\).
There is a set \(K\) of size \(m-2\) such that
\[
 L_1=K\mathbin{\dot\cup}\{a\},\qquad
 L_2=K\mathbin{\dot\cup}\{b\}.
\]
For \(k\notin\{1,2\}\), put \(u_k=q_{1k}\), \(v_k=q_{2k}\).
Both lists enumerate \(K\), so \(\pi(u_k)=v_k\) defines a permutation
of \(K\).  If \(C=\{u\in K:\pi(u)\ne u\}\) and \(c=|C|\), then
\(c\ge2\) and
\begin{equation}
 r_k\ge
 \begin{cases}
 c,&u_k\in C,\\
 c+2,&u_k\notin C.
 \end{cases}
 \qquad\text{Moreover, }r_k\ge3\text{ for every }k\notin\{1,2\}.
 \label{eq:damp-two-largest-rank-bounds}
\end{equation}
\end{theorem}

\begin{proof}
Both pieces choose every other piece as a protector.  Applying
\eqref{eq:damp-protector-equations} in both directions gives
\[
 D_{12}=\{a,b\}\mathbin{\dot\cup} C_0,
 \qquad\varnothing\ne C_0\subseteq L_1\cap L_2.
\]
The common part is nonempty because \(D_{12}\cap L_2\) contains
\(b\) and has size at least two.
For every third index \(k\), the identity
\(D_{12}=D_{1k}\mathbin\triangle D_{2k}\) gives
\begin{equation}
 D_{12}\setminus L_k=\{u_k\}\mathbin\triangle\{v_k\}.
 \label{eq:damp-two-largest-difference}
\end{equation}
The coordinate \(a\) differs from \(u_k\), since piece 1 uses distinct
labels for distinct targets; it also differs from \(v_k\in L_2\),
since \(a\notin L_2\).  Thus \(a\in L_k\).  The same argument with
1 and 2 interchanged gives \(b\in L_k\).

We next show that \(u_k\in L_2\).  Otherwise \(u_k\ne a\) and
\(D_{12}\setminus L_2=\{a\}\) imply \(u_k\notin D_{12}\).
But \(u_k\notin L_k\) and \(u_k\ne v_k\in L_2\), contradicting
\eqref{eq:damp-two-largest-difference} at coordinate \(u_k\).
Similarly \(v_k\in L_1\).
The \(m-2\) distinct labels \(u_k\) exhaust \(L_1\setminus\{a\}\)
and belong to \(L_2\); since \(a\notin L_2\), this set is exactly
\(K=L_1\cap L_2\).  Interchanging the indices proves the assertion
about the \(v_k\)'s and the two supports.

Since \(a,b\in L_k\), equation
\eqref{eq:damp-two-largest-difference} becomes
\[
 C_0\setminus L_k=\{u_k\}\mathbin\triangle\{v_k\}.
\]
Both \(u_k,v_k\) are absent from \(L_k\).  If \(u_k=v_k\), then
\(u_k\notin C_0\) and \(C_0\subseteq L_k\).  If \(u_k\ne v_k\),
then \(u_k,v_k\in C_0\) and \(C_0\setminus\{u_k,v_k\}\subseteq L_k\).
As the \(u_k\)'s enumerate \(K\), this proves that \(C_0=C\) is
exactly the set of moved points of \(\pi\).  A permutation cannot move
exactly one point, so \(c\ge2\).  Counting these required coordinates
together with \(a,b\) proves the two displayed rank bounds.

It remains to exclude \(r_k=2\).  The bounds already proved force
\(c=2\), \(u_k\ne v_k\), \(C=\{u_k,v_k\}\), and \(L_k=\{a,b\}\).
Translate all tips by \(y_1\), so that \(y_1\) is the zero word, and
identify words with their supports.  Equation
\eqref{eq:damp-protector-equations} for \(1\to k\) then forces
\(y_k=\{a,b,u_k\}\): its part outside \(L_k\) is \(\{u_k\}\),
and its part inside has at least two elements.
But \(y_2=\{a,b,u_k,v_k\}\), so \(y_k=y_2^{\{v_k\}}\in P_2\),
contradicting the missing-tip condition in~\eqref{eq:damp-protected-cover}.
\end{proof}

\begin{corollary}[Capacity with two largest supports]
\label{cor:damp-six-piece-two-largest-capacity}
For a cover satisfying~\eqref{eq:damp-protected-cover} with six pieces,
if two supports have size five, then every other support has size at
least three and
\[
 \sum_{i=1}^6|I_i|\ge82.
\]
In particular, the rank multiset \((5,5,4,2,2,2)\) is impossible.
\end{corollary}

\begin{proof}
A rank-\(r\) piece has \(f(r)=2^r-r-1\) protected concepts, and
\(f\) is increasing for \(r\ge2\).
In Theorem~\ref{thm:damp-two-largest-supports}, \(|K|=4\), so
\(c\in\{2,3,4\}\).  If \(c=2\), the two moved indices have rank
at least three and the two fixed indices have rank at least four;
the capacity is at least \(2f(5)+2f(3)+2f(4)=82\).
If \(c=3\), three ranks are at least three and the remaining rank is
at least five, giving \(2f(5)+3f(3)+f(5)=90\).
If \(c=4\), all four ranks are at least four, giving
\(2f(5)+4f(4)=96\).  Each case proves the stated bound.
\end{proof}

This is a bound on the sum of the protected-interior sizes, not on
\(|H|\): different interiors may overlap.  Nor does it give a
pc-minimality criterion for the constructed family.  Its role is to
exclude cover parameters using their signed neighbour incidences before
examining any elementary image.

\begin{proposition}[Elementary images of an oriented one-hole cover]
\label{prop:damp-one-hole-elementary-images}
For \(z\in Q_E\) and \(C\subseteq E\), write
\[
 \mathsf Q(z,C)=\{z^U:U\subseteq C\}.
\]
Let \((y,B)\) be an oriented one-hole datum, fix \(e\in E\), and let
\(\bar y\) be the word obtained by suppressing coordinate \(e\).  If
\(e\notin B\), then
\begin{align}
 \rho_e^i\mathsf P(y,B)
 &=
 \begin{cases}
  \mathsf P(\bar y,B),&i=y_e,\\
  \varnothing,&i\ne y_e,
 \end{cases}
 &
 \pi_e\mathsf P(y,B)&=\mathsf P(\bar y,B).
 \label{eq:damp-one-hole-image-outside-support}
\end{align}
If \(e\in B\) and \(C=B\setminus\{e\}\), then
\begin{align}
 \rho_e^{y_e}\mathsf P(y,B)&=\mathsf P(\bar y,C),
 &
 \rho_e^{1-y_e}\mathsf P(y,B)&=\mathsf Q(\bar y,C),
 &
 \pi_e\mathsf P(y,B)&=\mathsf Q(\bar y,C).
 \label{eq:damp-one-hole-image-inside-support}
\end{align}
Here \(\mathsf P(\bar y,C)\) has the same formula as in
\eqref{eq:damp-oriented-one-hole-cube}; it may have \(|C|=1\) and hence
consist of one vertex.

Consequently, for every finite cover \(\mathcal W\), each elementary
restriction or contraction of \(H(\mathcal W)\) is the union of the
corresponding images in
\eqref{eq:damp-one-hole-image-outside-support}--
\eqref{eq:damp-one-hole-image-inside-support}.  Thus it has a directly
computed mixed cover by full affine cubes and oriented one-hole cubes, with
the original missing-tip and support provenance retained.
\end{proposition}

\begin{proof}
If \(e\notin B\), every vertex of \(\mathsf P(y,B)\) has the fixed
\(e\)-value \(y_e\), which gives
\eqref{eq:damp-one-hole-image-outside-support}.  Suppose \(e\in B\).
On the side with value \(y_e\), the flipped set \(U\) is a nonempty
subset of \(B\setminus\{e\}\), giving the first formula in
\eqref{eq:damp-one-hole-image-inside-support}.  On the opposite side,
\(U\) contains \(e\), while \(U\setminus\{e\}\) is arbitrary; after
coordinate \(e\) is suppressed, this gives the full cube
\(\mathsf Q(\bar y,C)\).  Contraction is the union of the two restrictions
and hence gives the same full cube.  Finally, restrictions, contractions,
and unions commute, proving the cover formula for \(H(\mathcal W)\).
\end{proof}

\begin{definition}[Punctured expansion profile]
\label{def:damp-punctured-expansion-profile}
Let \(H\subseteq Q_E\), fix \(e\in E\), and put
\(\bar E=E\setminus\{e\}\).  For \(z\in\pi_eH\), \(A\subseteq\bar E\),
and \(C\subseteq A\), let
\begin{equation}
 \varepsilon_{H,e}(z,A;C)
 =\{i\in\{0,1\}:\iota_e^i(z^C)\in H\},
 \label{eq:damp-punctured-side-profile}
\end{equation}
where \(\iota_e^i\) inserts the value \(i\) in coordinate \(e\).
We call this a \emph{punctured expansion profile} if
\begin{equation}
 \varepsilon_{H,e}(z,A;A)=\varnothing,
 \qquad
 \varepsilon_{H,e}(z,A;C)\ne\varnothing
 \quad(C\subsetneq A).
 \label{eq:damp-profile-punctured-condition}
\end{equation}
It is \emph{direct} if one side contains the entire punctured cube, that is,
\begin{equation}
 \bigcap_{C\subsetneq A}\varepsilon_{H,e}(z,A;C)\ne\varnothing,
 \label{eq:damp-profile-direct}
\end{equation}
and \emph{switched} otherwise.  For a switched profile write
\begin{equation}
 \begin{split}
  P_A&=Q_A\setminus\{A\},\\
  D&=\{C\in P_A:\varepsilon_{H,e}(z,A;C)=\{0,1\}\},\\
  Z_i&=\{C\in P_A:\varepsilon_{H,e}(z,A;C)=\{i\}\}
       \qquad(i\in\{0,1\}).
 \end{split}
 \label{eq:damp-profile-separator-parts}
\end{equation}
Thus \(P_A=D\mathbin{\dot\cup}Z_0\mathbin{\dot\cup}Z_1\).
\end{definition}

The distinction is necessary because contraction takes a union of two
sections.  A punctured cube in the union need not occur in either section:
the two sections may cover different parts of it and change sides through
their common fibres.

\begin{theorem}[Canonical transport of local holes]
\label{thm:damp-canonical-hole-transport}
Let \(H\subseteq Q_E\) be ample, let \(e\in E\), and put
\(\bar E=E\setminus\{e\}\).
\begin{enumerate}
 \item If \(z\in\rho_e^iH\), then
       \begin{equation}
        \mathcal N(\Lambda_{\rho_e^iH}(z))
        =
        \{A\in\mathcal N(\Lambda_H(\iota_e^i(z))):A\subseteq\bar E\}.
        \label{eq:damp-restriction-local-hole-transport}
       \end{equation}
       Hence restrictions transport the canonical one-hole data simply by
       discarding the supports that contain \(e\).
 \item If \(z\in\pi_eH\), then
       \begin{equation}
        A\in\mathcal N(\Lambda_{\pi_eH}(z))
        \quad\Longleftrightarrow\quad
        \varepsilon_{H,e}(z,A;\cdot)
        \text{ is a punctured expansion profile}.
        \label{eq:damp-contraction-local-hole-transport}
       \end{equation}
 \item For every switched profile, \(D,Z_0,Z_1\) from
       \eqref{eq:damp-profile-separator-parts} satisfy:
       \(D,Z_0,Z_1\ne\varnothing\); no edge of \(P_A\) joins \(Z_0\) to
       \(Z_1\); every component of \(P_A\setminus D\) is contained in one
       \(Z_i\); and
       \begin{equation}
        D\text{ is ample},
        \qquad
        \operatorname{Sh}(D\cup Z_0)
        \cap\operatorname{Sh}(D\cup Z_1)
        =\operatorname{Sh}(D).
        \label{eq:damp-profile-separator-shadow-meet}
       \end{equation}
\end{enumerate}

Conversely, suppose that
\(P_A=D\mathbin{\dot\cup}Z_0\mathbin{\dot\cup}Z_1\), that \(D\) is ample,
and that the shadow equality in
\eqref{eq:damp-profile-separator-shadow-meet} holds.  Then the two-layer
family
\begin{equation}
 L(D,Z_0,Z_1)
 =((D\cup Z_0)\times\{0\})
  \cup((D\cup Z_1)\times\{1\})
 \label{eq:damp-profile-two-layer-construction}
\end{equation}
is ample, its contraction is \(P_A\), and \(A\) is a canonical local hole
at the vertex \(\varnothing\) of that contraction.  If both \(Z_i\) are
nonempty, this hole is switched.  Thus
\eqref{eq:damp-profile-separator-parts}--
\eqref{eq:damp-profile-separator-shadow-meet} give an intrinsic local
construction and recovery theorem for the contraction holes not visible in
the independent images of the original one-hole pieces.
\end{theorem}

\begin{proof}
For \(A\subseteq\bar E\), the \(A\)-cube through \(z\) lies in
\(\rho_e^iH\) exactly when the same cube through \(\iota_e^i(z)\) lies in
\(H\).  The local cube-support complex of the restriction is therefore the
induced subcomplex of \(\Lambda_H(\iota_e^i(z))\) on \(\bar E\).  Its
minimal nonfaces are precisely the minimal nonfaces of the original
complex that avoid \(e\), proving
\eqref{eq:damp-restriction-local-hole-transport}.

An \(A\)-cube through \(z\) lies in \(\pi_eH\) exactly when
\(\varepsilon_{H,e}(z,A;C)\ne\varnothing\) for every \(C\subseteq A\).
It is a minimal missing cube exactly when this holds for every proper
\(C\), but fails for \(C=A\).  This is
\eqref{eq:damp-contraction-local-hole-transport}.

Now fix a switched profile and abbreviate
\(D,Z_0,Z_1\) as in
\eqref{eq:damp-profile-separator-parts}.  Switchedness says that neither
side occurs at every proper vertex, so both \(Z_0\) and \(Z_1\) are
nonempty.  Suppose that adjacent vertices \(C\in Z_0\) and \(C'\in Z_1\)
existed.  Their chosen lifts in \(H\) would differ in coordinate \(e\) and
in the one coordinate of \(C\triangle C'\).  The only two intermediate
vertices on a length-two cube path are the opposite lift of \(C\) and the
opposite lift of \(C'\), but both are absent by the definitions of
\(Z_0,Z_1\).  This contradicts the fact that the ample family \(H\) is an
isometric cube subgraph.  Hence no such edge exists, and every component
of \(P_A\setminus D\) is monochromatic.  Since the punctured cube \(P_A\)
is connected and both colours occur, \(D\) is nonempty.

Restrict \(H\) outside \(A\cup\{e\}\) to the values prescribed by \(z\).
On the remaining coordinates the resulting ample family is exactly
\[
 L=((D\cup Z_0)\times\{0\})
   \cup((D\cup Z_1)\times\{1\}),
\]
and its contraction is \(P_A\).  For any two-layer family with sections
\(L_0,L_1\), union \(P=L_0\cup L_1\), and intersection
\(D=L_0\cap L_1\), one has the exact formula
\[
 \operatorname{Sh}(L)
 =\operatorname{Sh}(P)
  \mathbin{\dot\cup}
  \{J\cup\{e\}:J\in
    \operatorname{Sh}(L_0)\cap\operatorname{Sh}(L_1)\}.
\]
Moreover, \(|L|=|P|+|D|\).  Here \(L\) and
\(P=P_A\) are ample.  Consequently
\[
 |D|
 =|\operatorname{Sh}(L_0)\cap\operatorname{Sh}(L_1)|.
\]
Since
\(\operatorname{Sh}(D)\subseteq
  \operatorname{Sh}(L_0)\cap\operatorname{Sh}(L_1)\),
the Sandwich inequalities give
\[
 |D|\le|\operatorname{Sh}(D)|
 \le|\operatorname{Sh}(L_0)\cap\operatorname{Sh}(L_1)|
 =|D|.
\]
This proves that \(D\) is ample and establishes
\eqref{eq:damp-profile-separator-shadow-meet}.

Conversely, construct \(L(D,Z_0,Z_1)\) from the displayed partition.
Its contraction is \(P_A\), its order is \(|P_A|+|D|\), and the same
two-layer shadow formula, the ampleness of \(P_A,D\), and
\eqref{eq:damp-profile-separator-shadow-meet} give
\[
 |\operatorname{Sh}(L(D,Z_0,Z_1))|
 =|P_A|+|D|
 =|L(D,Z_0,Z_1)|.
\]
The Sandwich equality makes \(L(D,Z_0,Z_1)\) ample.  The tip \(A\) is
absent from its contraction while every proper \(C\subsetneq A\) is
present, so \(A\) is a minimal nonface of the local cube-support complex
at \(\varnothing\).  Both exclusive parts being nonempty is exactly the
switched condition.
\end{proof}

\begin{remark}[The first switched hole]
\label{ex:damp-first-switched-hole}
Let \(A=\{p,q\}\), put
\[
 D=\{\varnothing\},\qquad
 Z_0=\{\{p\}\},\qquad
 Z_1=\{\{q\}\}.
\]
The family \(L(D,Z_0,Z_1)\) is a four-vertex path.  Its contraction is
the punctured square
\(Q_{\{p,q\}}\setminus\{\{p,q\}\}\), but neither section contains that
punctured square: the two arms occur on opposite sides and meet at the
duplicated vertex \(\varnothing\).  Thus the missing square is a canonical
one-hole certificate of the contraction and is not the image of a
one-hole certificate supported away from the contracted coordinate.
\end{remark}

\begin{remark}[The switch phenomenon in Hall's obstruction]
\label{rem:damp-hall-switched-hole-audit}
The distinction is not confined to the preceding path.  A worker audit of
all twelve contractions of Hall's \(299\)-concept obstruction found
\(8{,}560\) canonical local holes.  Of these, \(7{,}691\) are direct and
\(869\) are switched; every contraction contains switched holes.  The
same audit verified
\eqref{eq:damp-restriction-local-hole-transport} for all \(9{,}159\)
canonical holes in the twenty-four restrictions.  The durable diagnostic is
\texttt{damp\_hall\_contraction\_hole\_transport.json}, generated by
\texttt{analyze\_damp\_hall\_contraction\_hole\_transport.py}.  These counts
are not used in the proof.
\end{remark}

\begin{remark}[The remaining seed condition]
\label{rem:damp-one-hole-seed-scope}
Theorem~\ref{thm:damp-one-hole-cornerless-construction} supplies the
nonpeelable part of a seed structurally: the one-hole antipodes cover the
family, so no first peeling step exists.  To obtain the cornerless seeds of
the forbidden basis, one must still characterize, in the same parameters,
when all elementary restrictions and contractions are corner-peelable.
Proposition~\ref{prop:damp-one-hole-elementary-images} gives a useful
piecewise image cover, but Theorem~\ref{thm:damp-canonical-hole-transport}
shows that it is not the canonical local-hole data after contraction:
switched holes arise by jointly assembling the two sections.  The exact
minor parameter must therefore retain the recovered punctured expansion
profiles, including their ample separators and component colours.  The
missing theorem must make peelability follow from the global incidence and
compatibility of those profiles.
Encoding the image conditions by orders, acyclic trace-carrier schemes, or
SAT would return to recognition and is not the desired next theorem.
\end{remark}

\begin{proposition}[Two-cone decomposition and the gluing dichotomy]
\label{prop:damp-yang-two-cone-gluing-dichotomy}
Let \(H\subseteq Q_E\) be ample, let \(e\in E\) be active, and write

\begin{equation}
 B_0=\rho_e^0H,\qquad B_1=\rho_e^1H,\qquad
 A=B_0\cap B_1,\qquad C=B_0\cup B_1.
 \label{eq:damp-yang-two-cone-sections}
\end{equation}

Then \(B_0,B_1,A,C\) are ample and

\begin{equation}
 \Delta_H
 =\bigl(v_{e,0}*\Delta_{B_0}\bigr)
  \underset{\Delta_A}{\cup}
  \bigl(v_{e,1}*\Delta_{B_1}\bigr).
 \label{eq:damp-yang-two-cone-pushout}
\end{equation}

For nested ample families \(A\subseteq B\), call an ordering
\(p_1,\ldots,p_k\) of \(B\setminus A\) an
\emph{\(A\)-relative ample chain} if

\begin{equation}
 A\cup\{p_1,\ldots,p_j\}
 \quad\hbox{is ample for every }0\le j\le k.
 \label{eq:damp-relative-ample-chain}
\end{equation}

Equivalently, this is a relative shelling of
\((\Delta_B,\Delta_A)\).  If \(A,B_{1-b}\in\Damp\) and
\(A\subseteq B_b\) has an \(A\)-relative ample chain for one
\(b\in\{0,1\}\), then \(H\in\Damp\).

Consequently, if \(H\in\Obs(\Damp)\cap\Ample\), then every active
coordinate has exactly one of the following two signatures.
\begin{enumerate}
 \item \emph{Inherited-site signature:} \(A\notin\Damp\).
 \item \emph{Gluing-critical signature:} \(A\in\Damp\), but neither
       relative pair
       \((\Delta_{B_0},\Delta_A)\) nor
       \((\Delta_{B_1},\Delta_A)\) is relatively shellable.
\end{enumerate}
In the second case the three absolute Yang complexes
\(\Delta_{B_0},\Delta_A,\Delta_{B_1}\) are shellable; the
obstruction lies in their relative gluing rather than in the topology of
an individual piece.
\end{proposition}

\begin{proof}
The two restrictions, their union, and their intersection are ample by
Lemma~\ref{lem:damp-ample-two-layer-shadow-intersection}.  Every facet of
\(\Delta_H\) contains exactly one of \(v_{e,0},v_{e,1}\), so

\[
 \Delta_H
 =\bigl(v_{e,0}*\Delta_{B_0}\bigr)
  \cup
  \bigl(v_{e,1}*\Delta_{B_1}\bigr).
\]

The two cones have distinct apex vertices.  Their intersection is therefore
\(\Delta_{B_0}\cap\Delta_{B_1}\), which equals \(\Delta_A\) by the
signed trace identity
\eqref{eq:damp-two-layer-yang-intersection}.  This proves
\eqref{eq:damp-yang-two-cone-pushout}.

Suppose, after interchanging the signs if necessary, that
\(p_1,\ldots,p_k\) is an \(A\)-relative ample chain in \(B_1\).
Starting from \(H\), delete the upper concepts
\((p_k,1),\ldots,(p_1,1)\) in this order.  Immediately before
\((p_j,1)\) is deleted, the upper section is
\(A\cup\{p_1,\ldots,p_j\}\), and deleting \(p_j\) preserves its
ampleness.  The ample corner-deletion criterion therefore makes \(p_j\)
a corner of that section.  Since \(p_j\notin A\), its lower mate is
absent, so the same concept is a corner of the current two-layer family.
After these deletions the remaining family is

\[
 (B_0\times\{0\})\cup(A\times\{1\}),
\]

the peripheral expansion with base \(B_0\) and site \(A\).  The exact
peripheral-expansion test and \(A,B_0\in\Damp\) give a peeling of this
remainder.  Reversing signs proves the symmetric assertion.

Now assume that \(H\) is a forbidden pc-minor.  The sections \(B_0,B_1\)
and the contraction \(C\) are proper pc-minors and hence belong to
\(\Damp\).  If \(A\notin\Damp\), the first signature holds.  Otherwise,
the preceding implication shows that a relative shelling on either side
would put \(H\) in \(\Damp\), a contradiction.  This gives the second
signature and proves the dichotomy.
\end{proof}

\begin{theorem}[Exact descent through a nonpeelable overlap]
\label{thm:damp-inherited-overlap-descent}
Let \(H\in\Obs(\Damp)\cap\Ample\), let \(e\) be active, and retain the
notation
\[
 B_b=\rho_e^bH,\qquad
 A=B_0\cap B_1=H^{\{e\}},\qquad
 C=B_0\cup B_1.
\]
If \(A\notin\Damp\), then
\[
 A\in\Obs(\Damp)\cap\Ample.
\]
Thus the inherited-site signature always exhibits a smaller member of the
forbidden basis, not merely an unspecified nonpeelable subfamily.  Notice
that \(A\) is a reduction of \(H\), not necessarily a pc-minor of
\(H\).

Consequently, repeated descent through nonpeelable overlaps terminates at
an obstruction \(S\) for which
\[
 S^{\{q\}}\in\Damp
 \qquad\text{for every active coordinate }q.
\]
Every coordinate of \(S\) is gluing-critical.  At each inverse step the
larger obstruction is recovered in the explicit form
\begin{equation}
 H=(A\times\{0,1\})
   \cup(R_0\times\{0\})
   \cup(R_1\times\{1\}),
 \qquad R_b=B_b\setminus A,
 \label{eq:damp-inherited-overlap-repair-expansion}
\end{equation}
where \(R_0,R_1\ne\varnothing\) and
\[
 A\cup R_0,\quad A\cup R_1,\quad
 A\cup R_0\cup R_1
 \quad\text{belong to }\Damp.
\]
The core \(A\) and the two signed repair wings \(R_0,R_1\) are recovered
canonically once \(e\) is chosen.
\end{theorem}

\begin{proof}
The ampleness of \(A\) is part of
Lemma~\ref{lem:damp-ample-two-layer-shadow-intersection}.  We prove that
all its elementary pc-minors belong to \(\Damp\).  Fix a coordinate
\(q\ne e\) active in \(A\).  Restrictions commute with reduction:
\[
 (\rho_q^iH)^{\{e\}}=\rho_q^iA
 \qquad(i\in\{0,1\}).
\]
The family \(\rho_q^iH\) is a proper pc-minor of \(H\), hence belongs to
\(\Damp\); Proposition~\ref{prop:damp-strong-projection-closure} puts
\(\rho_q^iA\) in \(\Damp\).

Contraction commutes as well, but here ampleness is essential.  The two
\(e\)-sections of \(\pi_qH\) are \(\pi_qB_0,\pi_qB_1\).  Applying the
signed trace identity
\eqref{eq:damp-two-layer-signed-trace-intersection} with the coordinates
\(e,q\) suppressed gives
\[
 (\pi_qH)^{\{e\}}
 =\pi_qB_0\cap\pi_qB_1
 =\pi_q(B_0\cap B_1)
 =\pi_qA.
\]
Again \(\pi_qH\) is a proper pc-minor in \(\Damp\), so its reduction \(\pi_qA\) belongs to \(\Damp\).  Therefore every elementary
pc-minor of \(A\) is in \(\Damp\), and pc-minor closure handles every
proper pc-minor.  The assumed \(A\notin\Damp\) proves
\(A\in\Obs(\Damp)\).

Each descent removes one coordinate, so iteration terminates.  At
termination every overlap is peelable, and
Proposition~\ref{prop:damp-yang-two-cone-gluing-dichotomy} makes every
coordinate gluing-critical.  Formula
\eqref{eq:damp-inherited-overlap-repair-expansion} is the section
decomposition itself.  The two repair wings are nonempty because an ample
forbidden pc-minor is periphery-free by
Proposition~\ref{prop:damp-periphery-free}.  Finally, the two sections and
their union are proper pc-minors of \(H\), proving the three displayed
\(\Damp\)-memberships.
\end{proof}

\begin{proposition}[Marked reductions of proper obstructions]
\label{prop:damp-inverse-proper-obstruction-reduction-test}
Let $X\subseteq Q_{F\cup\{e\}}$ be ample, using all $n+1$ displayed
coordinates, and suppose that its two $e$-sections and its $e$-contraction
belong to $\Damp$. Put $O=X^{\{e\}}$ and $n=|F|$.
Every ample forbidden proper pc-minor $H$ of $X$ has a presentation
$H=\nu X$ in which $\nu$ is a nonempty sequence of elementary operations
using only coordinates of $F$. The coordinate $e$ remains active, and
\[
 H^{\{e\}}=\nu O
\]
after identifying surviving coordinates. If $O$ is itself forbidden
and uses every coordinate of $F$, this marked reduction belongs to
$\Damp$.

If $H$ has $n$ active coordinates, $\nu$ consists of one elementary
operation. In that case the two sets of cardinalities
\[
 \{|H^{\{f\}}|:f\text{ active in }H\},\qquad
 \{|\rho_q^bO|,|\pi_qO|:q\in F,\ b\in\{0,1\}\}
\]
must intersect. These are necessary conditions, not an inverse
construction criterion.
\end{proposition}
\begin{proof}
Coordinate restrictions and contractions on distinct coordinates commute.
If a minor sequence producing $H$ removes $e$, move that operation first.
The resulting class is a positive section or contraction of $X$, and
pc-minor closure would give $H\in\Damp$, a contradiction. Thus the
sequence uses only old coordinates. The coordinate $e$ cannot become
constant either: the resulting family would equal an old-coordinate
minor of one positive $e$-section. Omitting operations in already constant
coordinates, we therefore obtain the stated nonempty sequence $\nu$.

Restriction in an old coordinate $q$ commutes immediately with reduction
in $e$. Contraction commutes for ample classes, as in the proof of
Theorem~\ref{thm:damp-inherited-overlap-descent}: if projected $e$-mates
are represented on opposite $q$-sides, their representatives have Hamming
distance two. Isometry supplies one of the intermediate vertices, hence
an actual $e$-edge on one $q$-side. Consequently
$(\pi_qX)^{\{e\}}=\pi_q(X^{\{e\}})$.
All intermediate pc-minors remain ample, so iteration gives
$H^{\{e\}}=\nu O$. If $O$ uses every old coordinate, the first
operation is proper on $O$; forbidden-minor minimality then makes
$\nu O$ positive. Finally each nontrivial elementary operation removes
at least one active coordinate. If $H$ has $n$ active coordinates,
there is exactly one such operation, giving the cardinality condition.
\end{proof}

\begin{corollary}[A cube-dimension exclusion for inverse counterexamples]
\label{cor:damp-inverse-cube-dimension-exclusion}
Under the hypotheses of
Proposition~\ref{prop:damp-inverse-proper-obstruction-reduction-test},
every forbidden proper pc-minor $H$ has an active coordinate $f$ with
\[
 \operatorname{VC}(H^{\{f\}})\le\operatorname{VC}(O).
\]
In particular, if $\operatorname{VC}(O)\le d$, an ample forbidden class
in which every active coordinate occurs in a cube of dimension at least
$d+2$ cannot be a proper pc-minor of $X$.
\end{corollary}
\begin{proof}
Take $f$ to be the surviving marked coordinate. Its reduction is the
pc-minor $\nu O$, so its VC dimension is at most that of $O$.
A cube of dimension at least $d+2$ containing $f$ reduces to a cube
of dimension at least $d+1$ in $H^{\{f\}}$, contradicting this bound.
\end{proof}

Consequently no member of the $1364$-vertex maximum-pair rectangular
family can be a proper pc-minor of an expansion satisfying the proposition
with overlap of VC dimension at most three, regardless of the dimension
or order of that overlap. Indeed, in
$(P_1\times O_2)\cup(O_1\times P_2)$ each coordinate of the first
block lies in a three-cube of $P_1$, whose product with a two-cube of
$O_2$ is a five-cube. The other block is symmetric. This applies to the
whole family described after
Corollary~\ref{cor:damp-rectangular-boundary-1364}; no assumption about
its particular signed realization is used.

\begin{corollary}[Known obstructions excluded from expansions of the recorded $267$-class]
\label{cor:damp-order267-inverse-adversary-exclusion}
In the preceding proposition, take $O=Y_{267}$ from
Proposition~\ref{prop:damp-order-267-obstruction}. Then no proper
pc-minor of $X$ is isomorphic to $Y_{267}$ or to any of the
$1364$-vertex maximum-pair rectangular obstructions described after
Corollary~\ref{cor:damp-rectangular-boundary-1364}.
\end{corollary}
\begin{proof}
The recorded $Y_{267}$ uses twelve coordinates. Every contraction has
$208$ vertices and every nonempty coordinate restriction has between
$77$ and $190$, as verified in
Proposition~\ref{prop:damp-order-267-obstruction}. Every coordinate
reduction has $267-208=59$ vertices, using the section union--intersection
identity. Thus $59$ is absent from the elementary-minor cardinalities
of $O$, excluding a proper minor isomorphic to $Y_{267}$.

For the rectangular family, write
$H=(P_1\times O_2)\cup(O_1\times P_2)$, with maximum rank-three
$P_i\subseteq Q_6$ and maximum rank-two $O_i\subset P_i$.
An $e$-reduction of $P_1$ has $16$ vertices and an $e$-reduction of
$O_1$ has $6$: their shattered supports are respectively all subsets
of at most two and at most one of the other five coordinates.
Over a vertex of $O_2$ the first factor is $P_1$, and over a vertex of
$P_2\setminus O_2$ it is $O_1$. Hence
\[
 |H^{\{e\}}|=16\cdot22+6\cdot(42-22)=472.
\]
The same count holds in the other block. These classes also use twelve
coordinates, and $472$ is absent from the elementary-minor cardinalities
of $O$, all of which are at most $208$. The proposition applies.
\end{proof}

This excludes familiar witnesses for a failure of the unrestricted
inverse construction; it does not prove that construction sound.
If an expansion in the corollary is not forbidden-minimal, a forbidden
proper minor on twelve coordinates must have a reduction matching an
elementary minor of $Y_{267}$. A forbidden witness on fewer coordinates
is not excluded by the one-operation argument.

\begin{theorem}[An inverse construction with one relatively peelable wing]
\label{thm:damp-one-relative-wing-inverse}
Let $O\in\Obs(\Damp)\cap\Ample$ use all coordinates of $Q_F$.
Choose $B_0,B_1\in\Damp$ with $B_0\cap B_1=O$, such that
$A=B_0\cup B_1$ is ample and no edge joins $B_0\setminus O$ to
$B_1\setminus O$. If at least one $B_i$ peels to $O$, then
\[
 X=(B_0\times\{0\})\cup(B_1\times\{1\})
 \quad\text{belongs to }\Obs(\Damp).
\]
No peeling tests on the old-coordinate elementary minors of $X$ are
required. The marked expansion and its inputs are recovered by its
two sections, contraction, and reduction. Consequently this rule recovers
every ample forbidden expansion with negative reduction and at least
one relatively peelable section.
\end{theorem}
\begin{proof}
The component-expansion theorem gives ampleness. Its reduction is $O$,
so reduction closure of $\Damp$ makes $X$ nonpeelable. Relabel the
layers so that $B_0$ peels to $O$. Lift this relative order to the
exclusive vertices of layer zero. Each deleted vertex has no vertical
mate, and its horizontal corner cube is exactly its corner cube in
the current section. This is a corner-deletion sequence from $X$ to
\[
 Y=(O\times\{0\})\cup(B_1\times\{1\}).
\]
Projecting the same sequence in the new coordinate gives a relative
peeling from $A$ to $B_1$: alternatively, each exclusive vertex has
no neighbours in the other exclusive wing. Hence $A$ is positive.
The three new-coordinate elementary minors $B_0,B_1,A$ are positive.

Now apply an elementary operation $\mu$ in an old coordinate to this
entire deletion sequence. Consecutive images are equal or differ by
one concept. They remain ample, so omitting repeated images gives a
corner-deletion sequence from $\mu X$ to
\[
 \mu Y=(\mu O\times\{0\})\cup(\mu B_1\times\{1\}).
\]
This argument transports a relative ample chain; it does not presume
an acyclic USO on the nonpeelable retained family $O$.
The family $\mu O$ is positive by forbidden-minor minimality of $O$,
and $\mu B_1$ is positive by pc-minor closure. The remainder is a
peripheral expansion of these positive inputs and is therefore positive.
Thus every elementary minor of $X$ is positive, proving forbiddenness.
For recovery, the negative reduction of a forbidden expansion is itself
forbidden by Theorem~\ref{thm:damp-inherited-overlap-descent}; its sections
are positive and recover the stipulated data.
\end{proof}

The relative-peeling hypothesis has a local sufficient replacement:
for at least one index $i$, require that every
$O\cup\{a\}$, $a\in B_i\setminus O$, is ample. Since $B_i$ is
positive, Corollary~\ref{cor:damp-initial-tip-minor-positivity} forces
$B_i$ to peel to $O$. Thus the inverse theorem applies with no supplied
relative order. Each one-concept condition is the punctured-cube test of
Lemma~\ref{lem:damp-one-concept-extension-test}. This applies to wings
of unrestricted size and dimension, but does not claim that every
inherited wing has individually admissible vertices.

\begin{corollary}[Two relative wings preserve all proper-minor sink choices]
\label{cor:damp-two-wing-all-minor-sinks}
Under the hypotheses of Theorem~\ref{thm:damp-one-relative-wing-inverse},
suppose both $B_0$ and $B_1$ peel to $O$, and every vertex of either
section is an attainable sink. Then every vertex of every nonempty proper
pc-minor of the forbidden expansion $X$ is an attainable sink. In
particular, $X$ has no rooted-factor pc-minor.
\end{corollary}
\begin{proof}
The theorem already proves forbiddenness. Each section is all-sinks by
hypothesis. The contraction $B_0\cup B_1$ peels to either section by
clearing the opposite wing, since no edge joins the exclusive wings.
It therefore also attains every prescribed sink.

For an elementary operation $\mu$ in an old coordinate, reductions
commute with $\mu$ on ample classes, as in
Theorem~\ref{thm:damp-inherited-overlap-descent}. Thus the overlap is
$\mu O$, nonempty and positive by minimality of $O$. The sections
$\mu B_i$ attain every sink by rooted minor transport, and both relative
wing orders transport to $\mu O$. To end at a prescribed vertex in one
section of $\mu X$, clear the opposite wing, delete that layer's copy
of $\mu O$ using an ordinary order with all mates retained, then use
the prescribed endpoint order in the retained section. This gives every
endpoint of $\mu X$. Rooted transport extends the assertion from all
elementary minors to every proper pc-minor. A rooted factor has an
unattainable root, whereas $X$ itself is not ordinarily positive, so no
rooted factor occurs as any pc-minor.
\end{proof}

\begin{corollary}[Forest wings remove the minor tests]
\label{cor:damp-forest-wing-inverse}
In the preceding theorem, replace the relative-peeling hypothesis by
the condition that at least one induced graph $G(B_i\setminus O)$ is
a forest. The same conclusion holds. This gives sound construction
and marked recovery for all ample forbidden expansions with negative
reduction and a forest exclusive wing.
\end{corollary}
\begin{proof}
For nested nonempty ample families $O\subseteq B_i$, the relative Yang
module is Cohen--Macaulay by
Lemma~\ref{lem:damp-relative-yang-cm-flag-h}. The relative forest theorem
of \cite{BousquetYangComplex26}, Theorem
\texttt{relative-cm-forest-linear-quotients}, supplies an $O$-relative
ample chain in $B_i$. Reversing it peels $B_i$ to $O$. Apply the theorem.
\end{proof}

\begin{corollary}[Small graph cores remove the minor tests]
\label{cor:damp-small-component-wing-inverse}
Under the structural and positive-section hypotheses of
Theorem~\ref{thm:damp-one-relative-wing-inverse}, suppose that, for
one $i\in\{0,1\}$, every connected component of the graph
$2$-core of $G(B_i\setminus O)$ has at most eight vertices.
Here the $2$-core is obtained by repeatedly deleting vertices of degree
at most one; an empty $2$-core satisfies the condition.
Then $X$ is an ample forbidden pc-minor. There is no bound on the
number of components or on the trees attached to their $2$-cores.
Consequently an expansion with neither section relatively peelable
to $O$ must have, in each exclusive wing, a connected component of its
$2$-core with at least nine vertices.
\end{corollary}
\begin{proof}
We first establish the corresponding statement for any relative Yang
pair $\varnothing\ne H\subseteq A$ whose relative face module is
Cohen--Macaulay: if every component of the $2$-core of $G(A\setminus H)$
has at most eight vertices, its relative facets admit linear quotients.
The intermediate pairs need not be ample. We use the leaf deletion
dichotomy, the deletion-append corollary, the supporting-cube teaching
lift, and the eight-facet theorem of \cite{BousquetYangComplex26}
(respectively \texttt{relative-cm-leaf-deletion-dichotomy},
\texttt{relative-cm-deletion-append},
\texttt{supporting-cube-teaching-lift}, and
\texttt{general-eight-facet-linear-quotients}). A concept teaches
relative to $H$ when its generator has a colon generated by variables
modulo the generators of $H$; it may then be placed first.

Induct on $|D|$, where $D=A\setminus H$, allowing $D=\varnothing$.
Deleting any vertex of $G(D)$ preserves the hypothesis: every
minimum-degree-two subgraph of the remaining graph lies in the old
$2$-core, so each new core component lies in an old one.
Suppose first that $d\in D$ has degree one. If $d$ teaches relative
to $H$, the teaching lift makes the pair $(A,H\cup\{d\})$
Cohen--Macaulay; prepend $d$ to the ordering supplied by induction.
Otherwise the leaf dichotomy makes $(A\setminus\{d\},H)$
Cohen--Macaulay. Induction orders that pair, and the deletion-append
corollary appends $d$ to its order.

It remains to find a teaching concept when there is no degree-one
vertex. The Alexander-dual relative module has a linear resolution and
its linear presentation splits as a direct sum over the components of
$G(D)$: its linear relations involve either one generator or two cube
neighbours. Each component summand therefore has a linear resolution.
An isolated vertex gives a cyclic summand, whose annihilator is
generated by variables, and hence teaches relative to $H$.
If there is no isolated vertex either, every component has minimum
degree at least two and is a component of the $2$-core, of order at
most eight. Apply the eight-facet theorem to one such summand to
obtain a concept teaching relative to $H$. In both cases the teaching
lift and induction allow that concept to be prepended, as above.
This completes the induction.

Apply this statement to $(B_i,O)$, whose relative face module is
Cohen--Macaulay by Lemma~\ref{lem:damp-relative-yang-cm-flag-h}.
That lemma also translates the relative linear-quotient order into an
ample addition chain from $O$ to $B_i$. Its reversal peels $B_i$ to
$O$, so Theorem~\ref{thm:damp-one-relative-wing-inverse} applies.
The final assertion is the contrapositive, applied to each section.
\end{proof}

\begin{lemma}[Removing a uniform inactive coordinate]
\label{lem:damp-uniform-inactive-coordinate}
Let $\varnothing\ne O\subsetneq B\subseteq Q_F$ be ample. Suppose
every concept of $B\setminus O$ has the same value $b$ at coordinate
$j$, and that the section $O_j^b$ is nonempty. If $j$ belongs to
every member of $\mathcal T=\operatorname{Sh}(B)\setminus
\operatorname{Sh}(O)$, or to none of them, then restriction to the
$j=b$ section preserves exactly the relative corner deletion orders.
The relative shattered supports of this section pair are
\[
 \{S\setminus\{j\}:S\in\mathcal T\},
\]
with no repetitions. Thus removing such a coordinate preserves both
relative peeling and the residual graph, with its remaining edge labels.
\end{lemma}
\begin{proof}
Suppress $j$ in section words and write $A=O_j^b$, $D=B_j^b$ and
$C=O_j^{1-b}=B_j^{1-b}$. For any ample family $H$ with sections
$H_0,H_1$, the supports containing $j$ are exactly
$\{S\cup\{j\}:S\in\operatorname{Sh}(H_0)\cap
\operatorname{Sh}(H_1)\}$. The other supports are
$\operatorname{Sh}(H_0\cup H_1)$. Section ampleness and
$|\operatorname{Sh}(H)|=|H_0|+|H_1|$ therefore give
\[
 \operatorname{Sh}(H_0\cup H_1)
 =\operatorname{Sh}(H_0)\cup\operatorname{Sh}(H_1):
\]
the left side contains the right side and the cardinalities agree.
Apply this identity to $O$ and $B$. Each
$S\in\operatorname{Sh}(D)\setminus\operatorname{Sh}(A)$
contributes exactly one relative support: $S\cup\{j\}$ if
$S\in\operatorname{Sh}(C)$, and $S$ otherwise. This proves the
asserted bijection of supports.

Coordinate projections of ample families are ample. Consequently
\[
 |(D\setminus A)\setminus C|
 =|D\cup C|-|A\cup C|
 =|\{S\in\mathcal T:j\notin S\}|.
\]
If no relative support contains $j$, this equals
$|\mathcal T|=|D\setminus A|$. Thus no residual concept has a
neighbour across coordinate $j$ in $B$. During any relative deletion
order in $D$, every cube through the next residual concept in the
whole family is confined to the $b$ section. A section corner is
therefore a corner of the whole family, so the order lifts.

If every relative support contains $j$, the displayed count is zero,
so $D\setminus A\subseteq C$ and $D\cup C=A\cup C$. Let
$A\subseteq D'\subseteq D$ be any ample intermediate family in a
section deletion order. Then $D'\cup C=A\cup C$ and, by the support
bijection,
\[
 \operatorname{Sh}(D')\cup\operatorname{Sh}(C)
 =\operatorname{Sh}(A)\cup\operatorname{Sh}(C)
 =\operatorname{Sh}(A\cup C).
\]
The family with sections $D',C$ thus has exactly
\[
 |\operatorname{Sh}(D')\cup\operatorname{Sh}(C)|
 +|\operatorname{Sh}(D')\cap\operatorname{Sh}(C)|
 =|D'|+|C|
\]
shattered supports, and is ample. Each singleton deletion in the lifted
chain is a corner deletion by the ample corner-deletion criterion.
Conversely, restricting any relative deletion order of $B$ to its
$b$ section gives the same sequence of singleton deletions through
ample families. It is therefore a relative corner order of $D$ to $A$.
\end{proof}

\begin{corollary}[A unique corner cannot have codimension two in seven coordinates]
\label{cor:damp-unique-corner-codimension-two}
Let $X$ be an ample family on $q+2$ coordinates, where $1\le q\le5$.
If a corner of $X$ has degree $q$, it is not the unique corner.
\end{corollary}
\begin{proof}
Suppose the corner is unique, and translate it to zero. Use its incident
directions as a set $F$ of size $q$, and call the other two directions
$u,v$. The sections of $X$ have the form
\[
 X=(Q_F\times\{00\})\cup(P\times\{10\})
   \cup(Q\times\{01\})\cup(C\times\{11\}),
\]
where $P,Q,C$ are ample, possibly empty. Neither $P$ nor $Q$ contains
zero, because $u,v$ are not incident with the root. Neither does $C$:
otherwise the root and $(0,11)$ would have distance two with both
intermediate vertices absent, contrary to isometry. Every nonzero
vertex of the full $00$ section must have an outside neighbour, or it
would also be a corner. Hence
\[
 P\cup Q=Q_F\setminus\{0\}.
\]

We claim $P,Q\subseteq C$. If $P\setminus C\ne\varnothing$, then
$P\cap C$ is an ample reduction of the $u=1$ face. If this intersection
is nonempty, the five-coordinate relative theorem gives a corner of
$P$ outside it. If it is empty, ordinary peeling in at most five
coordinates gives a corner of $P$. In either case the chosen vertex
has no mate in $C$ and has all mates of its old corner cube in the
full $00$ section. Its copy in $10$ is consequently another corner
of $X$, a contradiction. The same argument applies to $Q$. It follows
that
\[
 C=P\cup Q=Q_F\setminus\{0\}.
\]
If $P$ or $Q$ is empty, the other equals $C$. A corner of $C$ in
the $11$ section then lifts with its entire corner cube into that
other section, and is a corner of $X$. Thus both $P,Q$ are nonempty.

Consider the ample $u=1$ face and the ample $u$-reduction,
respectively:
\[
 U=(P\times\{0\})\cup(C\times\{1\}),\qquad
 K=(P\times\{0\})\cup(Q\times\{1\}),
\]
where the displayed last coordinate is $v$. If $U=K$, the required
relative peeling is empty. Otherwise their projections along
$v$ both equal $C$, so every relative shattered support of
$K\subseteq U$ contains $v$: the supports avoiding $v$ are exactly
the shattered supports of the projection. Lemma
\ref{lem:damp-uniform-inactive-coordinate} applies to this interval
and identifies its relative orders with those of $Q\subseteq C$.
The latter interval peels relatively by the five-coordinate theorem.
Therefore $U$ peels to $K$.

Lift these deletions into $X$: each deleted concept is outside the
$u$-overlap, so has no $u$-mate. The remaining upper layer is $K$,
contained in the lower layer
$L=(Q_F\times\{0\})\cup(Q\times\{1\})$.
The ample family $K$ uses at most six coordinates and peels ordinarily;
delete its upper copy using that order, with all mates retained in $L$.
Next delete the $Q$ layer of $L$, with all mates in $Q_F$, and peel
the final cube to zero. This is a peeling of $X$ ending at its unique
corner. Since $X$ is nontrivial, its first deleted corner would be
different from zero, the final contradiction.
\end{proof}

\begin{corollary}[Only degree three remains in seven coordinates]
\label{cor:damp-seven-coordinate-unique-corner-degree}
If an ample class using seven coordinates has a unique corner, its
corner has degree three. The exclusion of degree four is computer-assisted.
\end{corollary}
\begin{proof}
Degrees one and two are excluded by the six-coordinate prescribed-endpoint
result (Corollary~\ref{cor:damp-six-coordinate-peeling}) and Corollary~\ref{cor:damp-small-rooted-overlaps}, respectively.
Degrees six and seven are excluded by the full-cube argument: in the
codimension-one case the opposite section must contain every unit vertex
but not zero, so its ample connected complement is the singleton zero;
the resulting cube minus one vertex has other corners. Degree five is
excluded by Corollary~\ref{cor:damp-unique-corner-codimension-two}.

For degree four, write the root cube as $Q_4\times\{0\}$ and its
three-coordinate fibers as $A_x$. Every $A_x$ with $x\ne0$ is nontrivial:
otherwise $(x,0)$ would be another corner. At least one unit fiber is
proper. Indeed, if all four unit fibers were full, every nonzero
opposite section would contain the four unit vertices and exclude zero.
Its ample complement would be connected with zero isolated, forcing
that section to be $Q_4\setminus\{0\}$. The entire class would be
$(Q_4\times\{0\})\cup((Q_4\setminus\{0\})\times Q_3)$, and every
$(e_i,y)$ with $y\ne0$ would be another corner.

Normalize a proper unit fiber by coordinate permutations. Its $62$
possible labelled ample patterns form $17$ orbits under permutations
of the three fiber coordinates. The companion paper
\path{../ample-corners/main.tex}, Theorem
\texttt{thm:unique-corner-degree-four}, gives the exact Lawrence
total-asymmetry encoding and exhaustive orbit list. All $17$ cases have
independently checked DRAT proofs, retained with their formulas and
fingerprints in
\path{../ample-corners/evidence/unique_corner_q7_degree4_certificates/}.
Thus no proper unit fiber is possible either, excluding degree four.
\end{proof}

\begin{corollary}[Prescribed endpoints in seven coordinates, computer-assisted]
\label{cor:damp-seven-coordinate-peeling}
Every nonempty ample class of isometric dimension at most seven can
be corner-peeled to any prescribed concept. Consequently every rooted
factor has isometric dimension at least eight, and every ample forbidden
pc-minor of $\Damp$ with a cut vertex has isometric dimension at least
sixteen.
\end{corollary}
\begin{proof}
Corollary~\ref{cor:damp-six-coordinate-peeling} covers dimensions at
most six. In dimension seven, Corollary
\ref{cor:damp-seven-coordinate-unique-corner-degree} leaves degree
three as the only possible unique-corner degree. The companion
\path{../ample-corners/main.tex}, Theorem
\texttt{thm:unique-corner-degree-three}, excludes it by $185$
independently DRAT-checked formulas. Its proper-unit-fiber lemma reduces
the target to the $2763$ proper nontrivial ample four-coordinate fibers
containing zero, partitioned into $185$ coordinate-permutation orbits.
Each case is encoded exactly by Lawrence's total-asymmetry criterion,
the unique-corner clauses, and its fixed fiber. Formulas, proofs, logs,
source snapshots and fingerprints are retained under
\path{../ample-corners/evidence/unique_corner_q7_degree3_certificates/}.

Proposition~\ref{prop:damp-seven-coordinate-ample-corners} supplies a
corner in every nonempty ample class in this dimension. Together with
the unique-corner exclusion, this gives at least two corners whenever
the class is nontrivial. Delete a corner other than the prescribed
endpoint and repeat; ampleness and the dimension bound persist.
A rooted factor fails prescribed-root peeling, so it cannot have
dimension at most seven. Finally,
Theorem~\ref{thm:damp-cut-vertex-obstructions} expresses every ample
forbidden class with a cut vertex as two rooted factors on disjoint
coordinate sets. Their dimensions add, giving at least $8+8=16$.
\end{proof}

\begin{corollary}[Uniform-support necessity for a nonpeelable cyclic interval]
\label{cor:damp-marked-cycle-wing-inverse}
Let $\varnothing\ne O\subsetneq B\subseteq Q_F$ be ample, and suppose
that $G(B\setminus O)$ is a cycle of length greater than four. Put
$\mathcal T=\operatorname{Sh}(B)\setminus\operatorname{Sh}(O)$, ordered
by inclusion. Call $S\in\mathcal T$ isolated if it is incomparable
with every other member of $\mathcal T$.

At least as many residual vertices admit an ample
addition to $O$ or deletion from $B$ as there are nonisolated members
of $\mathcal T$.
Let $Y\subseteq F$ be the coordinates varying on $R=B\setminus O$.
If $B$ does not peel to $O$, there is a cyclic listing
$c_0,\ldots,c_{m-1}$ of $R$ and a bijective listing
$S_0,\ldots,S_{m-1}$ of $\mathcal T$ such that, with indices modulo $m$,
\[
 S_i\setminus S_{i+1}=\{e_i\},\qquad
 |S_{i+1}\setminus S_i|=1,
 \qquad c_{i+1}=c_i^{\{e_i\}}.
\]
In particular all $S_i$ have the same cardinality and agree on
$F\setminus Y$. Writing $k=|S_i\cap Y|$ gives
$m\le\binom{|Y|}{k}$. More precisely, restriction to the smallest
coordinate face containing $R$ preserves relative peeling and these
cyclic exchanges, and
\[
 3\le k\le |Y|-3,\qquad
 2|Y|\le m\le\binom{|Y|}{k}.
\]
In particular a nonpeelable residual cycle has at least six active
coordinates and at least twelve vertices.

Thus $B$ peels to $O$ whenever its relative supports have different
cardinalities, or do not all agree outside $Y$. Either alternative
is a sufficient condition on a section
of Theorem~\ref{thm:damp-one-relative-wing-inverse}, with no bound on
the cycle length or ambient dimension and no relative order supplied.
Every residual cycle of length at most ten also peels relatively,
in arbitrary ambient dimension.
\end{corollary}
\begin{proof}
Write $R=B\setminus O$ and let $I$ be the number of isolated members
of $\mathcal T$. Every vertex of $R$ has degree two in its induced
graph. The lemma ``Degree-two marker direction'' of
\cite{BousquetResidualExtension26} says that a vertex $c$ for which
neither $O\cup\{c\}$ nor $B\setminus\{c\}$ is ample determines an
isolated relative support $S_c$, and distinct such vertices give
distinct supports. Thus at most $I$ vertices are blocked at both
ends. Ampleness and nesting give
\[
 |\mathcal T|=|\operatorname{Sh}(B)|-|\operatorname{Sh}(O)|
 =|B|-|O|=|R|.
\]
At least $|\mathcal T|-I$ vertices therefore admit one of the two
operations, proving the quantitative assertion. If there is an available
vertex $c$, the following argument supplies a full relative peeling.

In the first case the residual graph of $O\cup\{c\}\subseteq B$
is a path. The relative forest theorem gives a peeling of $B$ to
$O\cup\{c\}$, after which deleting $c$ leaves $O$.
In the second case delete the corner $c$ of $B$ first, and apply the
relative forest theorem to $O\subseteq B\setminus\{c\}$.
Both constructions peel $B$ to $O$.

Now suppose no such peeling exists. Every vertex is blocked, so the
injection above is a bijection onto all of $\mathcal T$, and every
relative support is isolated. Denote by $S(c)$ the support assigned
to $c$. The degree-two lemma further gives exactly one of the two
incident cycle directions in $S(c)$. The antichain carrier-duality
theorem and the exact-edge-orientation lemma of
\cite{BousquetResidualExtension26} say that, for an edge of direction
$e$ with endpoints $c,d$, exactly one of $S(c),S(d)$ contains $e$.
Orient the edge away from that endpoint. Every vertex has one outgoing
and one incoming edge, so the cycle is directed. List its vertices
in that order and put $S_i=S(c_i)$.

For an isolated support $S$, let $p_S$ be its assigned concept and set
\[
 C^-_S=\{x\in R:x|_S=p_S|_S\},\qquad
 C^+_S=\{x\in R:x|_{F\setminus S}=p_S|_{F\setminus S}\}.
\]
The cross-carrier cube proposition of the same source says that a
nonempty $C^-_S\cap C^+_T$ is a full cube in $R$ with free coordinates
$T\setminus S$. Because $e_i\in S_i\setminus S_{i+1}$, both endpoints
of the edge $c_ic_{i+1}$ belong to
$C^-_{S_{i+1}}\cap C^+_{S_i}$. A cycle longer than four contains no
square, so this cube is that edge, and
$S_i\setminus S_{i+1}=\{e_i\}$.
Incomparability gives $|S_{i+1}\setminus S_i|\ge1$, whence
$|S_{i+1}|\ge|S_i|$. Going around the cycle forces equality everywhere,
and every reverse difference is a singleton as asserted.

All removed coordinates $e_i$ belong to $Y$. A coordinate outside $Y$
therefore cannot leave a support along this cyclic listing; cyclicity
also prevents it from entering. The supports consequently agree
outside $Y$. Their distinct intersections with $Y$ all have size $k$,
which proves the binomial bound. Remove the coordinates outside $Y$
one at a time using Lemma~\ref{lem:damp-uniform-inactive-coordinate}.
Its nonempty-section hypothesis holds at every step: otherwise the
upper section would be the residual cycle itself, whereas an ample
square-free graph is a tree. Indeed it has VC dimension at most one,
each coordinate labels at most one edge by connectedness of reductions,
and Sandwich equality gives one more vertex than edges.
The resulting ample interval $O'\subsetneq B'\subseteq Q_Y$ has
the same relative orders and relative supports
$\{S\cap Y:S\in\mathcal T\}$, all of size $k$.

If $k\le2$, Theorem~\ref{thm:damp-relative-small-supports} peels
$B'$ to $O'$, a contradiction. For the other inequality consider the
complementary ample interval
$Q_Y\setminus B'\subsetneq Q_Y\setminus O'$. Its lower family is
nonempty: otherwise $B'=Q_Y$, and the uniform relative support family
would contain $Y$, hence consist only of $Y$, contrary to $m>4$.
Its relative supports are $\{Y\setminus T:T\in
\operatorname{Sh}(B')\setminus\operatorname{Sh}(O')\}$.
This is the usual shattering duality: a support is not shattered by
a family exactly when its complementary support is strongly shattered
by the cube complement; for ample families strong shattering equals
shattering. If $|Y|-k\le2$, the same small-support theorem supplies
a relative peeling of the complementary interval. Taking complements
and reversing the chain supplies one of the original interval,
again a contradiction. Hence $3\le k\le |Y|-3$.
Each active coordinate occurs a positive even number of times along
a closed cube walk, so $m\ge2|Y|\ge12$.
The sufficient conditions follow by
contraposition, and the inverse conclusion follows from
Theorem~\ref{thm:damp-one-relative-wing-inverse}.
\end{proof}

The singleton-absorption conclusion is essential here. The
general representation-map extension theorem for residual cycles in
\cite{BousquetResidualExtension26} does not assert relative peeling
or acyclicity. It therefore does not settle the remaining case of
equal-sized supports agreeing outside the active residual coordinates
and admitting the displayed cyclic exchanges. These conditions are
necessary for relative failure, not sufficient. In particular every
non-antichain relative-support family is covered by the sufficient
criterion, since a strict comparison gives different cardinalities.

\begin{remark}[The leaf dichotomy does not extend vertexwise to cycles]
\label{rem:damp-cycle-endpoint-dichotomy}
The degree-one hypothesis in the preceding induction is essential.
Encode subsets of four coordinates by binary integers, and put
\[
 A=\{0,1,2,3,4,5,6,8,10,11\},\qquad H=\{0,8,10,11\}.
\]
Both families are ample. Indeed, $L=\{0\}$, $S=\{0,2,3\}$, and
$U=\{0,1,2,3,4,5,6\}$ are ample with $L\subseteq S\subseteq U$;
$H$ has sections $L,S$ in the fourth coordinate and $A$ has sections
$U,S$, so each is a peripheral expansion of ample families.
The exclusive graph is the induced cycle
$1,3,2,6,4,5,1$. Nevertheless $d=1$ cannot be placed either first or
last in an ample addition chain from $H$ to $A$.
In $H\cup\{1\}$ the vertices $1,11$ differ in two coordinates but
both possible intermediate vertices $3,9$ are absent. In
$A\setminus\{1\}$ the vertices $3,5$ have the same failure, with
missing intermediates $1,7$. Neither family is isometric, hence neither
is ample.

This does not obstruct relative peeling of the pair: deleting
$5,6,4,1,3,2$ from $A$ retains $H$, and every deletion is a corner.
The identities and this order are replayed in
\path{research/relative_cycle_endpoint_dichotomy.py/json}.
Thus a general cycle argument would need a suitable choice of vertex;
it cannot apply the leaf dichotomy to an arbitrary cycle vertex.
No failure of the small-core corollary or of a general cycle-peeling
statement is asserted.
\end{remark}

\begin{proposition}[Simultaneous extensions retaining a negative reduction]
\label{prop:damp-fixed-negative-reduction-extension}
Let $X\subseteq Y\subseteq Q_F$ be ample, with $X$ an irredundant
forbidden pc-minor. Suppose an active coordinate $e$ satisfies
\[
 X^{\{e\}}=Y^{\{e\}}\notin\Damp.
\]
If $Y$ peels to $X$, then $Y$ is a forbidden pc-minor.
In particular, the same conclusion holds under the purely graph-theoretic
condition that every connected component of $G(Y\setminus X)$ is a tree
or has at most eight vertices. The number of components and the sizes of
the tree components are unrestricted. No addition order or output
proper-minor peeling certificate is required in this latter construction.
\end{proposition}
\begin{proof}
Fix an elementary restriction or contraction $\mu$ in $F$. Apply it to
the relative deletion sequence from $Y$ to $X$. Its images are ample,
and consecutive images are equal or differ by one concept. After omitting
repetitions, the ample corner-deletion criterion gives a relative peeling
from $\mu Y$ to $\mu X$. Since all coordinates are active in $X$,
$\mu X$ is a proper pc-minor of $X$ and belongs to $\Damp$.
Appending its peeling proves $\mu Y\in\Damp$. Hence every proper
pc-minor of $Y$ is positive. The displayed negative reduction prevents
$Y$ itself from being positive, proving forbiddenness.

For the graph condition, apply the componentwise relative-chain argument
in Corollary~\ref{cor:damp-small-component-wing-inverse} to the ample
inclusion $X\subseteq Y$. Each component $K$ gives an ample family
$X\cup K$, by the component-expansion theorem. The relative forest or
eight-facet theorem supplies its deletion order down to $X$; orders for
different components concatenate because their vertices are nonadjacent.
Thus $Y$ peels to $X$, and the first part applies.
\end{proof}

The first part is the relative-chain extension of the one-concept atlas:
transporting the entire chain proves all elementary minor tests at once.
The second part gives a simultaneous, nonrecursive choice of added words
for this subclass. Ampleness is tested by the established local or shadow
criteria, and equality of the two reductions is an explicit fibre condition.
The forbidden input still has to be supplied. No converse claiming that
all extensions preserving the reduction have such component graphs, or
that all reduced obstructions contain a suitable smaller forbidden input,
is asserted.

For an explicit application, use the fixed rooted factor $C$ of order
$289$ and its repairs $C\cup\{128,160\}$ and $C\cup\{290\}$, both
peeling to root zero. Put $\widehat C=\{2^{12}x:x\in C\}$ and
\[
 O=C\cup\widehat C,\quad
 B_0=O\cup\{128,160\},\quad
 B_1=O\cup\{2^{12}290\}.
\]
The cut-vertex theorem makes $O$ a $577$-vertex forbidden minor.
Vertex gluing makes both sections positive. Their exclusive wings
are respectively an edge and a singleton, with no edge between them;
their union is ample by the same vertex-gluing construction.
The resulting expansion has $1157$ vertices in $25$ coordinates and
is forbidden by the corollary. The verifier
\path{research/one_relative_wing.py/json} replays the two-step relative
order $(160,128)$, both positive section orders, and the $24$ archived
root-preserving elementary orders of $C$. The old-coordinate minor
conclusion follows uniformly from the theorem, not from a new search.
This extends the one-concept repair rule to a nontrivial wing; it does
not classify positive repairs of arbitrary seeds or expansions whose
two wings both fail relative peeling.

\begin{proposition}[Opposite peripheral completions of rooted factors]
\label{prop:damp-opposite-peripheral-completions}
Let $(A_i,0)$, $i\in\{0,1\}$, be rooted factors as in
Theorem~\ref{thm:damp-cut-vertex-obstructions}, irredundantly embedded
on disjoint coordinate sets $F_i$. Choose ample supersets
$A_i\subseteq T_i\subseteq Q_{F_i}$. With all unused coordinates
set to zero and a fresh coordinate $e$, put
\[
 B_0=T_0\vee A_1,\qquad B_1=A_0\vee T_1,\qquad
 X=(B_0\times\{0\})\cup(B_1\times\{1\}),
\]
where $\vee$ identifies the two zero vertices. Then $X$ is ample and
\[
 X\in\Obs(\Damp)
 \quad\Longleftrightarrow\quad
 T_0\text{ and }T_1\text{ each admit an acyclic USO with sink }0.
\]
Equivalently, the right-hand condition says $B_0,B_1\in\Damp$;
it also implies $B_0\cup B_1\in\Damp$. No relative peeling from
$T_i$ to $A_i$ is required. The overlap is the forbidden class
$O=A_0\vee A_1$, and every proper-minor test on $X$ follows from
the rooted-factor contracts and the two supplied endpoint orders.
For the marked coordinate $e$ and coordinate blocks $F_0,F_1$,
all four inputs are recovered uniquely from the two sections.
The order is $|A_0|+|A_1|+|T_0|+|T_1|-2$.
\end{proposition}
\begin{proof}
The useful decomposition consists of two peripheral expansions meeting
at the edge $D=\{(0,0),(0,1)\}$:
\[
 P_0=(T_0\times\{0\})\cup(A_0\times\{1\}),\qquad
 P_1=(A_1\times\{0\})\cup(T_1\times\{1\}).
\]
Each piece is ample by the peripheral-expansion theorem. Their private
coordinate sets are disjoint. A shattered support cannot use a private
coordinate from both pieces, since its trace assigning one to both is
absent. The shadows of the pieces intersect in $\{\varnothing,\{e\}\}$.
Thus the shadow and vertex cardinalities both add with subtraction of two,
proving ampleness of $X=P_0\cup P_1$. Its $e$-reduction is $O$, which
is forbidden by the cut-vertex theorem. Reduction closure of $\Damp$
therefore makes $X$ negative for every choice of the two supersets.

Suppose first that $X$ is forbidden. Its two $e$-sections are positive.
The vertex-gluing test applied to $T_0\vee A_1$ forces $T_0$ to
peel to zero, since $A_1$ cannot peel to zero. Applying the same test
to $A_0\vee T_1$ gives the other required endpoint order.
These arguments also prove necessity in the stated equivalence with
positivity of the two sections.

Conversely, suppose both endpoint orders are supplied. Vertex gluing
makes the two sections and the contraction $T_0\vee T_1$ positive.
Each $P_i$ is positive by peripheral expansion, since both of its
sections are positive. It remains to check old-coordinate minors.

Fix $f\in F_0$. Restriction to $f=1$ removes $P_1$ entirely and
leaves a pc-minor of positive $P_0$, so that image is positive.
A contraction in $f$, or restriction to $f=0$, replaces $A_0,T_0$
by $A'_0\subseteq T'_0$. Rooted minor minimality makes $A'_0$
peel to zero. Root-preserving minor transport of the supplied order
does the same for $T'_0$. In the changed piece $P'_0$, delete the
upper nonzero vertices using an order of $A'_0$ ending at zero.
Each deletion has its whole required mate cube in the untouched lower
$T'_0$. Then delete the lower nonzero vertices using an order of
$T'_0$ ending at zero. The single retained upper vertex affects only
the retained lower root. This is a relative peeling of $P'_0$ to $D$.
Every deleted vertex has no neighbour in $P_1\setminus D$, because
the private coordinate blocks are disjoint. The same deletions are
therefore valid in the whole minor, leaving the positive $P_1$.

For $f\in F_1$, restriction to $f=1$ instead leaves a minor of
$P_1$. Contraction or zero restriction gives $A'_1\subseteq T'_1$
with both endpoint orders. Delete first the nonzero lower $A'_1$
and then the nonzero upper $T'_1$, retaining $D$ and leaving positive
$P_0$. This treats all three elementary operations in every old
coordinate. Pc-minor closure now proves positivity of every proper
pc-minor of $X$, hence forbiddenness.

Finally the marked sections recover
$T_0=B_0\cap Q_{F_0}$, $A_1=B_0\cap Q_{F_1}$,
$A_0=B_1\cap Q_{F_0}$ and $T_1=B_1\cap Q_{F_1}$.
Each section has a single shared zero, giving the order formula.
\end{proof}

This is an exact admissibility rule within a prescribed two-block
presentation. Unlike the one-relative-wing proof, its minor orders clear
one changed peripheral piece to the common edge; they do not clear a wing
to the negative overlap. The completion sizes are unrestricted: one may,
for example, take $T_i=Q_{F_i}$ with its standard order ending at zero.
The preceding $1157$-vertex construction is a common special case, not a
new obstruction. We do not assert that this rule adds supports outside
every earlier construction. Rooted-factor selection remains an input,
and neither general inherited admissibility nor exhaustive recovery follows.

\begin{corollary}[Recovery when all separating-edge pieces are peripheral]
\label{cor:damp-peripheral-separator-recovery}
The construction of Proposition~\ref{prop:damp-opposite-peripheral-completions}
generates exactly the ample forbidden pc-minors having a separating edge
$D$ such that, in every canonical piece at $D$, the two sections in the
direction of $D$ are nested. Every such class has exactly two canonical
pieces. Their larger sections lie on opposite sides of $D$, their smaller
sections are rooted factors, and their larger sections attain the common
projected root as a sink. The separating edge and all four inputs are
recovered uniquely, up to exchanging the pieces and reversing its direction.
In particular, a three-piece forbidden amalgam has a piece whose two
sections are not nested.
\end{corollary}
\begin{proof}
Let $H$ be an ample forbidden class with the stated separating edge,
normalized to $D=\{(0,0),(0,1)\}$ in direction $e$.
Write its canonical pieces as $P_1,\ldots,P_m$. Their private
coordinate blocks are disjoint by
Theorem~\ref{thm:damp-edge-separator-obstructions}.
In each piece choose a larger section $T_i$ and a smaller section
$S_i\subseteq T_i$. The two sections are positive, since each is
a pc-minor of the positive proper piece $P_i$. Say a section attains
its root when it peels to the projected zero vertex.

If both $S_i$ and $T_i$ attain their roots, their endpoint orders
peel $P_i$ to $D$: first clear the nonzero vertices of its smaller
layer and then those of its larger layer. The separator theorem
forbids this, so at least one of these two sections fails to attain
its root. On the other hand, each of the two $e$-sections of $H$
and its $e$-contraction is positive. The vertex-gluing test therefore
allows at most one root failure in each of these three images.

Suppose some $T_k$ fails to attain its root, and let $b$ be its
layer. Positivity of the contraction forces every $T_i$, $i\ne k$,
to attain its root; hence every corresponding $S_i$ fails to do so.
Since $m\ge2$, there is at least one such index. The section $S_k$
must attain its root: otherwise piece $k$ contributes a failure in
both layers, and any other $S_i$ gives a second failure in one layer.
Every $S_i$, $i\ne k$, must lie in layer $1-b$, since layer $b$
already contains the failure of $T_k$. That layer can contain only
one failure, so $m=2$. Both larger sections then lie in layer $b$.
The entire $H$ is a peripheral expansion whose base and site are its
positive $e$-sections, making $H$ positive, a contradiction.

Consequently all $T_i$ attain their roots, and all $S_i$ fail to do
so. Each layer can contain at most one smaller section; hence $m=2$
and the two smaller sections occupy opposite layers. Their vertex
gluing is $H^{\{e\}}$ and is negative by the vertex-gluing test.
The inherited-overlap descent theorem makes it forbidden. The
cut-vertex characterization therefore identifies $S_1,S_2$ as its
two rooted factors. The full-coordinate assertion in
Theorem~\ref{thm:damp-separator-strong-descent} makes each $S_i$
use every private coordinate of its piece. We have recovered exactly the construction of
Proposition~\ref{prop:damp-opposite-peripheral-completions}.

Conversely, take an output of that proposition with root-attaining
completions. Its two peripheral pieces meet only in $D$. Each rooted
factor $A_i$ is two-connected, so $A_i\setminus\{0\}$ is connected.
The larger section $T_i\setminus\{0\}$ is connected as well.
Indeed, for a nonzero $x\in T_i$ choose a coordinate $f$ with $x_f=1$.
Full coordinate support of $A_i$ gives $y\in A_i$ with $y_f=1$.
An isometric path from $x$ to $y$ in $T_i$ stays in the $f=1$
halfspace and avoids zero. Thus every vertex joins
$A_i\setminus\{0\}$ without passing through zero. It follows that
each peripheral piece minus $D$ is connected, and these are precisely
the two components of $H\setminus D$. They are its canonical pieces.

Finally the negative $e$-reduction and
Theorem~\ref{thm:damp-separator-strong-descent} make $D$ the unique
separating edge. Its components recover the two coordinate blocks,
and the nested sections recover the smaller and larger inputs.
Equality of those sections is impossible: the smaller input is
root-negative and the larger one is root-positive. This proves the
claimed uniqueness and the three-piece exclusion.
\end{proof}

The conclusion is exhaustive for the nested-section separator branch.
The root-attaining completions have explicit positive certificates, but
the rooted factors still need an intrinsic generator. Nonnested pieces
and classes without separating edges remain outside this recovery theorem.

\begin{theorem}[Reduction descent in the separator branch]
\label{thm:damp-separator-strong-descent}
Let \(H\subseteq Q_F\) be an irredundantly embedded ample forbidden
pc-minor, and let \(\Gamma=\operatorname{Cr}(H)\).
For every \(D\subseteq F\), a nonpeelable reduction \(H^D\)
is an ample forbidden pc-minor using all coordinates of \(F\setminus D\).
In particular, if \(H^{\{e\}}\) is nonpeelable, then \(e\) is a
\emph{universal vertex} of \(\Gamma\), meaning that it is adjacent
to every other vertex.

Suppose that \(\Gamma\) is not two-connected, or equivalently that
\(H\) has a cut vertex or a separating edge.  Then exactly one of
the following holds:
\begin{enumerate}
 \item Every reduction \(H^D\) with \(D\ne\varnothing\)
 belongs to \(\Damp\).
 \item There is a unique coordinate \(e\) such that
 \begin{equation}
  H^D\notin\Damp
  \quad\Longleftrightarrow\quad
  D\in\{\varnothing,\{e\}\}.
  \label{eq:damp-separator-negative-strong-supports}
 \end{equation}
 In this case \(e\) is the unique universal vertex and the unique
 cut vertex of \(\Gamma\); the graph \(\Gamma-e\) has exactly
 two components.  The family \(H\) has a unique separating
 \(e\)-edge, and \(H^{\{e\}}\) is an ample forbidden pc-minor
 with a cut vertex.  Its two rooted blocks are the strong
 \(e\)-projections of the two canonical pieces at that edge.
\end{enumerate}
Consequently every forbidden class with a cut vertex has only peelable
nonempty-coordinate reductions.  Throughout the separator branch,
descent through nonpeelable overlaps has length at most one, and its
terminal obstruction is uniquely determined.
\end{theorem}

\begin{proof}
Write \(K=H^D\notin\Damp\).  Restrictions and contractions in a
coordinate outside \(D\) commute with reduction along \(D\).
For contractions, iterate the ample signed-trace identity used in
Theorem~\ref{thm:damp-inherited-overlap-descent}; all intermediate
reductions are ample.  Suppose a coordinate \(q\in F\setminus D\)
were inactive in \(K\).  The family \(\pi_qH\) is a proper pc-minor
of \(H\), so reduction closure gives
\[
 \pi_qK=(\pi_qH)^D\in\Damp.
\]
The nonempty ample family \(K\) is connected.  An inactive coordinate
is therefore constant, making \(K\) isomorphic to \(\pi_qK\), a
contradiction.  This proves the full-coordinate assertion.  Every
intermediate projection along a subset of \(D\) is nonpeelable,
since its further reduction is \(K\).  Repeated application
of Theorem~\ref{thm:damp-inherited-overlap-descent} now makes \(K\)
a forbidden pc-minor.

If \(D=\{e\}\), an edge of direction \(q\ne e\) in \(K\)
lifts to an \(\{e,q\}\)-square in \(H\).  All such \(q\)
are active by the preceding paragraph, so \(e\) is universal in
\(\Gamma\).  More generally, if \(H^D\) is nonpeelable, then
\(H^{\{e\}}\) is nonpeelable for every \(e\in D\), by
reduction closure.

An ample forbidden class has VC dimension at least three, so
\(\Gamma\) has at least three vertices.  A graph of this order
with two universal vertices is two-connected: after any single vertex
is removed, a remaining universal vertex connects the rest.  Thus
\(\Gamma\) has at most one universal vertex under the hypothesis.
It follows that either all nonempty-coordinate reductions peel,
or their only nonpeelable member is \(H^{\{e\}}\) for a unique
\(e\), proving \eqref{eq:damp-separator-negative-strong-supports}.
If \(\Gamma\) is disconnected it has no universal vertex, which
also proves the assertion for a cut vertex of \(H\).

In the remaining case, \(e\) is universal and \(\Gamma\) is
connected.  Removing any vertex other than \(e\) leaves a connected
graph.  Since \(\Gamma\) is not two-connected, \(e\) is its
unique cut vertex.  Put \(K=H^{\{e\}}\).  Its crossing graph uses
all vertices of \(\Gamma-e\), and every edge of
\(\operatorname{Cr}(K)\) is an edge of \(\Gamma-e\): a shattered
pair in \(K\) lifts to a shattered triple in \(H\).
Thus \(\operatorname{Cr}(K)\) is disconnected, so \(K\) has a
cut vertex by Proposition~\ref{prop:damp-crossing-separator-recovery}.
Theorem~\ref{thm:damp-cut-vertex-obstructions} gives exactly two
blocks in \(K\).  Their coordinate sets are precisely the two
components of \(\operatorname{Cr}(K)\), as proved in the separator
recovery proposition.  The disconnected spanning supergraph
\(\Gamma-e\) consequently has exactly two components as well.

The graph \(\Gamma\) is connected, so \(H\) has no cut vertex.
Proposition~\ref{prop:damp-crossing-separator-recovery} now recovers a
unique separating \(e\)-edge and its two canonical convex pieces
\(H_1,H_2\).  Normalize that edge to \(\{0,\mathbf1_{\{e\}}\}\).
Taking its reduction gives
\[
 K=H_1^{\{e\}}\vee H_2^{\{e\}},
\]
with common root \(0\) and disjoint exclusive coordinate sets.
Both factors are nontrivial: the full-coordinate assertion for \(K\)
retains every coordinate assigned to each piece.  The cut-vertex theorem
therefore identifies these factors with the two rooted blocks of \(K\).
Equation~\eqref{eq:damp-separator-negative-strong-supports} says that
all further nonempty-coordinate reductions peel, proving the
termination and uniqueness claims.
\end{proof}

The universal-vertex condition locates the only possible descent coordinate;
it does not decide whether that carrier peels.  Both outcomes occur in the
existing constructions.  The \(1158\)-concept terminal edge amalgam of
Remark~\ref{ex:damp-terminal-edge-amalgam-1158} has one universal crossing
coordinate, but its \(261\)-concept carrier peels.  The one-layer tower
over the \(1158\)-concept terminal cut-vertex obstruction in
Remark~\ref{ex:damp-cornered-maximal-seed-2318} has \(2318\) concepts; its layer coordinate is the unique universal crossing vertex,
and its reduction is that nonpeelable cut-vertex obstruction.
The fixed checks, including explicit orders transported from proper
contractions when a carrier loses a coordinate, are recorded by
\path{verify_damp_separator_strong_descent.py} in
\path{damp_separator_strong_descent_verification.json}.
This bounds and makes canonical an existing descent within the separator
branch.  It does not classify its rooted factors intrinsically or settle
descent choices for cores with two-connected crossing graphs.

\begin{proposition}[Square attachments leave nonpeelable carriers unchanged]
\label{prop:damp-square-attachments-fixed-carriers}
Let \(H\subseteq Q_F\) be an irredundantly embedded ample forbidden
pc-minor, and let \(K\) be its canonical graph three-core.
For every nonempty \(D\subseteq F\),
\[
 H^D\notin\Damp\quad\Longleftrightarrow\quad K^D\notin\Damp,
\]
and in this case \(H^D=K^D\) as embedded families.
Moreover, no acquired coordinate pair in the square-attachment data
meets a coordinate \(e\) for which \(H^{\{e\}}\notin\Damp\).
For such an $e$, put $O=H^{\{e\}}=K^{\{e\}}$. For each
$b\in\{0,1\}$, the section $\rho_e^bH$ peels to $O$ if and only if
$\rho_e^bK$ peels to $O$. Thus the failure of relative peeling on
both sides of a negative overlap occurs in $H$ exactly when it occurs
at the same marked coordinate in its three-core $K$.

Consequently, if \(\operatorname{Cr}(K)\) is not two-connected,
descent from \(H\) through nonpeelable reductions has length
at most one and a unique terminal obstruction.  When a descent exists,
its result is the same cut-vertex obstruction obtained from \(K\).
This conclusion does not require a separator in \(H\) itself.
\end{proposition}

\begin{proof}
Theorem~\ref{thm:damp-three-core} and
Theorem~\ref{thm:damp-square-attachment-data} reconstruct \(H\)
from \(K\) by ample one-concept additions of support size two.
Every intermediate family is forbidden and uses all of \(F\).
Consider one such addition \(V=U\cup\{x\}\), with acquired pair
\(I\).  Proposition~\ref{prop:damp-repair-strong-projections} gives
\[
 V^D=
 \begin{cases}
  U^D\cup\{x|_{F\setminus D}\},&D\subseteq I,\\
  U^D,&D\not\subseteq I,
 \end{cases}
\]
and the addition in the first case is ample, with acquired support
\(I\setminus D\).  Therefore a peelable \(U^D\) always gives
a peelable \(V^D\): when the families differ, delete the new corner
first.  This includes the empty-family and singleton cases.

Suppose \(U^D\) is nonpeelable.  If \(D\subseteq I\), then
\(|I\setminus D|\le1\).  But
Theorem~\ref{thm:damp-separator-strong-descent} makes \(U^D\)
nonempty and active on every coordinate of \(F\setminus D\).
Its shattered complex already contains the empty set and every singleton,
so none can be an acquired support.  This contradiction proves
\(D\not\subseteq I\) and hence \(V^D=U^D\).
Thus the nonpeelable projections and their exact embedded families are
unchanged at every attachment, proving the first assertion.
For \(D=\{e\}\), the same argument gives \(e\notin I\)
at every step, proving the assertion about acquired pairs.

Fix such an $e$ and one square addition $V=U\cup\{x\}$.
Since $e\notin I$, the new vertex has no $e$-neighbour in $V$.
It therefore changes only the section indexed by $x_e$; the other
section and the reduction $O$ are unchanged. In the changed section
it is an ample square addition with the same acquired pair $I$.
Proposition~\ref{prop:damp-rank-two-extension-invariance}, with retained
target $O$, preserves relative peelability in both directions. This
proposition permits a nonpeelable target, so no orientation on $O$ is
assumed. Iterating over all attachments proves the two section
invariances and their simultaneous-failure consequence.

The three-core has the same active coordinates as \(H\).  Apply
Theorem~\ref{thm:damp-separator-strong-descent} to \(K\) and use
the just-proved equality of all nonpeelable projections.  If every
nonempty-coordinate projection of \(K\) peels, the same is true of
\(H\).  Otherwise their only such nonpeelable projection is the
same cut-vertex obstruction \(K^{\{e\}}=H^{\{e\}}\), all of
whose further nonempty-coordinate reductions peel.
\end{proof}

For a concrete example where the extension matters, take the
\(2318\)-concept one-layer tower just described, in the integer
coordinates of Remark~\ref{ex:damp-cornered-maximal-seed-2318}.
The word \(32769=1+2^{15}\) completes the square at \(0\) in
directions \(0,15\).  Its addition has degree two, so the tower is
the three-core of the resulting \(2319\)-concept obstruction.
The new shattered pair joins the two components left after deleting
the layer coordinate \(26\) from the crossing graph.  That coordinate
is still universal, so the new crossing graph is two-connected.
The full obstruction consequently has neither a cut vertex nor a
separating edge, yet its strong \(26\)-projection is still the
\(1158\)-concept cut-vertex obstruction.  The same verifier checks
this square, the three-core, the crossing graph, the unchanged carrier,
and all proper elementary-minor peelings of the enlarged family.

\begin{theorem}[Replacing a coordinate by a block]
\label{thm:damp-coordinate-block-replacement}
Let \(H\subseteq Q_{G\cup\{z\}}\) be ample, and write
\[
 B_i=\rho_z^iH,\qquad C=B_0\cup B_1,\qquad A=B_0\cap B_1.
\]
Let \(E\) be a new nonempty coordinate set, disjoint from \(G\),
and put \(m=|E|\).  Define
\begin{equation}
 W_E(H)=
 (B_0\times\{0_E\})\cup(B_1\times\{1_E\})
 \cup\bigl(C\times(Q_E\setminus\{0_E,1_E\})\bigr).
 \label{eq:damp-coordinate-block-construction}
\end{equation}
For \(m=1\) this is \(H\) with its coordinate \(z\) renamed.
For every \(m\ge1\), the following assertions hold.
\begin{enumerate}
 \item The family \(W_E(H)\) is ample, and
 \begin{align}
 \operatorname{Sh}(W_E(H))
 &=\{D\cup T:D\in\operatorname{Sh}(C),\ T\subsetneq E\}
   \mathbin{\dot\cup}
   \{D\cup E:D\in\operatorname{Sh}(A)\},
   \label{eq:damp-coordinate-block-shadow}\\
 |W_E(H)|&=(2^m-1)|C|+|A|.
 \label{eq:damp-coordinate-block-order}
 \end{align}
 If \(H\ne\varnothing\), its VC dimension increases by \(m-1\).
 In its signed minimal-nonface presentation, leave every clause avoiding
 \(z\) unchanged.  In every other clause, replace \(z\) by all of
 \(E\), repeating the forbidden sign at \(z\) on every new coordinate.
 These are exactly the signed minimal nonfaces of \(W_E(H)\).
 \item \(W_E(H)\in\Damp\) if and only if \(H\in\Damp\).
 If \(H\) is irredundantly embedded and \(z\) is active, then
 \[
 W_E(H)\in\Obs(\Damp)
 \quad\Longleftrightarrow\quad H\in\Obs(\Damp).
 \]
 \item If \(H\) is cornerless, then \(W_E(H)\) is cornerless.
 If \(H\) is an irredundant ample forbidden pc-minor and \(m\ge2\),
 then \(W_E(H)\) has minimum degree at least three and is its own
 graph three-core.  Every coordinate in \(E\) is universal in its
 crossing graph, which is two-connected.
\end{enumerate}
The marked construction has unique recovery: its two constant sections
recover \(B_0,B_1\), and hence recover \(H\).  Equivalently, the
marked signed clauses recover the original presentation by replacing
\(E\) by \(z\).  No resolving tips or peeling orders are construction
parameters.
\end{theorem}

\begin{proof}
We first determine the shadow, then lift corner deletions and check the
proper minors.  Empty families and empty sections are allowed in the
shadow and peeling assertions; their shadows are empty.

Suppose first that \(m\ge2\).  If \(T\subsetneq E\), every word on
\(T\) extends to a nonconstant word on \(E\).  Consequently a support
\(D\cup T\), with \(D\subseteq G\), is shattered exactly when
\(D\) is shattered by \(C\).  A support containing all of \(E\)
is shattered exactly when its old part is shattered by both \(B_0\)
and \(B_1\).  Identity
\eqref{eq:damp-residual-section-shadow-intersection}, with its residual
family \(K=C\), identifies their common shadow with
\(\operatorname{Sh}(A)\).  If a section is empty, this identity
is immediate.
This proves \eqref{eq:damp-coordinate-block-shadow}; for \(m=1\)
it is the usual two-layer formula.  Its support count is
\((2^m-1)|C|+|A|\), also the vertex count in
\eqref{eq:damp-coordinate-block-construction}.  The Sandwich equality
therefore proves ampleness.  The largest support has size
\[
 (m-1)+\max\{\operatorname{VC}(C),1+\operatorname{VC}(A)\}
 =(m-1)+\operatorname{VC}(H),
\]
with the second term omitted when \(A\) is empty.

The shadow formula also identifies its minimal nonfaces.  Those avoiding
\(z\) in the original shadow remain unchanged.  Each old minimal
nonface containing \(z\) becomes \((I\setminus\{z\})\cup E\).
Indeed, a nonface whose new-coordinate part is proper must already
contain an old nonface of \(\operatorname{Sh}(C)\); a minimal remaining
nonface must contain \(E\), and deleting any element of its old part
must put that part in \(\operatorname{Sh}(A)\).  These are precisely
the old minimality conditions for a nonface containing \(z\).
All nonconstant new-coordinate traces occur over \(C\); the two
constant traces have exactly the original sections.  Thus the unique
missing trace repeats the old sign, proving the clause assertion using
Theorem~\ref{thm:damp-fixed-shadow-signed-clutter}.

For the peeling assertion, first take \(E=\{e,f\}\).  Let
\((x,i)\) be a corner of \(H\).  If its mate \((x,1-i)\) belongs to
\(H\), the cube at that corner has all its old directions in the
intersection \(A\).  Hence \((x,ii)\) is a corner of \(W_E(H)\);
its deletion gives exactly \(W_E(H\setminus\{(x,i)\})\).
If its mate is absent, delete
\[
 (x,01),\quad (x,10),\quad (x,ii)
\]
in that order.  To verify these are corners, no old neighbor of
\(x\in B_i\setminus B_{1-i}\) can lie in
\(B_{1-i}\setminus B_i\).  Otherwise the corresponding vertices
of \(H\), at Hamming distance two, would have neither possible
intermediate vertex, contradicting isometry.  Thus the old neighbor
set of \(x\) in \(C\) equals its neighbor set in \(B_i\).
The corner cube of \((x,i)\) supplies that entire old cube in
\(B_i\).  Before each of the first two deletions its product with
the remaining incident new edge is present.  After those deletions,
the old cube alone certifies the last corner.  Again the remainder is
exactly the replacement of the source after its corner deletion.
Lifting every step of a peeling proves the forward implication.
The reverse implication follows from
\[
 W_{\{e,f\}}(H)^{\{e\}}=H
\]
with \(z\) renamed \(f\), and reduction closure.
Repeatedly replace one member of the new block by two coordinates.
The constant sections remain \(B_0,B_1\), and every nonconstant
section is \(C\), so this gives exactly \(W_E(H)\) for every
\(m\).  Induction proves the peeling equivalence, including empty
and constant-coordinate intermediate families.

Now let \(H\) be an irredundant ample forbidden pc-minor.  For a
binary replacement, every restriction or contraction \(\mu\) in an
old coordinate \(q\in G\) commutes with the construction: the formula
uses only the two sections and their union.  Thus
\(\mu W_{\{e,f\}}(H)=W_{\{e,f\}}(\mu H)\) peels by the preceding
paragraph.  A contraction in \(e\) or \(f\) is \(C\square Q_1\).
A restriction in either coordinate is a peripheral expansion of \(C\)
along \(B_0\) or \(B_1\).  All three input families are proper
pc-minors of \(H\), so Proposition
\ref{prop:damp-exact-expansion-product-tests} makes these images peel.
The replacement itself does not peel.  This proves forbiddenness;
iteration proves it for every block size.  All coordinates remain active.

Conversely, suppose \(W_E(H)\) is forbidden, with \(m\ge2\).
The peeling equivalence makes \(H\) nonpeelable.  Every proper
old-coordinate image of \(H\) is a reduction of the corresponding
proper image of \(W_E(H)\), so it peels.  The two constant sections
\(B_i\) are proper face restrictions of \(W_E(H)\), and \(C\)
is a face restriction after contracting one new coordinate.  They peel
as well.  These are all elementary pc-minors of \(H\), proving the
reverse implication.  For \(m=1\) it is just coordinate renaming.

For cornerlessness, each vertex \((x,t)\) of the replacement belongs
to a full fibre on \(E\setminus\{e\}\) for some \(e\in E\).
If \(x\in A\), any \(e\) works.  If
\(x\in B_i\setminus B_{1-i}\), choose \(e\) with \(t_e=i\);
such an \(e\) exists because the excluded constant word is absent.
The reduction along that fibre is \(H\), with \(z\) renamed
\(e\).  A corner in a full fibre projects to a corner: every neighbor
of the projected vertex lifts beside the given vertex, and the corner
cube contains its product with the fibre.  A corner of the replacement
would therefore give a corner of \(H\).

Finally, the degree of every new vertex is at least
\(\delta(H)+m-1\).  For an exclusive old word, the old neighbor set
is unchanged by the isometry argument above, and its punctured cube row
has degree at least \(m-1\).  For an overlap word, a constant section
retains its old neighbors and replaces its one \(z\)-edge by \(m\)
edges; a nonconstant section has at least as many old neighbors.
An ample forbidden family has minimum degree at least two, as shown in
the proof of Theorem~\ref{thm:damp-three-core}.  Thus the replacement
has minimum degree at least three.  Each old coordinate is active in
\(C\), and \(A\ne\varnothing\) since \(H\) is connected and
\(z\) is active.  The shadow formula now makes every new coordinate
adjacent to all others in the crossing graph.  Two universal vertices
make a graph of this order two-connected.  The recovery assertions
follow directly from the constant sections and the signed clause rule.
\end{proof}

\begin{proposition}[Relative peeling under coordinate-block replacement]
\label{prop:damp-block-relative-transfer}
Let $T\subseteq H$ be nonempty ample families in the same cube, with
marked coordinate $z$, and use the block construction
$W_E$ of Theorem~\ref{thm:damp-coordinate-block-replacement} on both.
Then
\[
 W_E(H)\text{ peels to }W_E(T)
 \quad\Longleftrightarrow\quad H\text{ peels to }T.
\]
For an ample forbidden $H$, retain the section notation $B_0,B_1$,
$A=B_0\cap B_1$, $C=B_0\cup B_1$ from that theorem.
If $|E|\ge2$ and $e\in E$, then
\begin{equation}
 \rho_e^bW_E(H)\text{ peels to }W_E(H)^{\{e\}}
 \quad\Longleftrightarrow\quad B_b\text{ peels to }A
 \qquad(b=0,1).
 \label{eq:damp-block-relative-new-coordinate}
\end{equation}
For a coordinate $q$ outside the replaced block, relative peelability
of either $q$-section to its $q$-reduction is also unchanged by the
replacement.
\end{proposition}
\begin{proof}
The corner-lifting argument in the block-replacement theorem works for
a deletion prefix as well as for a full peeling. Stop it when $T$
remains to obtain the forward construction of the relative order.
Conversely reduce $W_E(H)$ along all but one coordinate of $E$.
This recovers $H$ and sends $W_E(T)$ to $T$. Applied to a supplied
relative deletion sequence, reduction gives ample families with
consecutive images equal or differing by one concept. Omitting repeated
images gives a relative corner-deletion sequence from $H$ to $T$.
No peelability assumption on the retained target is needed.

For the second assertion, first let $E=\{e,f\}$ and consider $e=0$.
Its section and retained overlap, written in coordinate $f$, are
\[
 L_0=(B_0\times\{0\})\cup(C\times\{1\}),\qquad
 O=(B_0\times\{0\})\cup(B_1\times\{1\}).
\]
Only vertices $(x,1)$ with $x\in B_0\setminus A$ can be deleted.
After deleting those indexed by a subset $R$ of this difference,
the old neighbours of any remaining such $x$ in $C\setminus R$
are exactly its neighbours in $B_0\setminus R$: there are no edges
between the two exclusive parts of an ample expansion. Its lower mate
is present. Consequently $(x,1)$ is a corner of the current $L_0$
exactly when $x$ is a corner of $B_0\setminus R$. In one direction,
the latter corner cube has its full lower copy in $B_0$, supplying
the product corner cube. In the other direction, the product corner
cube projects to the required cube in $B_0\setminus R$.
Thus the two relative orders correspond term by term. Interchanging
zero and one proves the other case.

For larger $E$, replace $f$ in this binary pair $(L_b,O)$ by the
block $E\setminus\{e\}$. The resulting families are precisely the
$e$-section and the $e$-reduction of $W_E(H)$, so the first assertion
proves \eqref{eq:damp-block-relative-new-coordinate}.
For an old coordinate $q$, sections and reductions commute with block
replacement by \eqref{eq:damp-coordinate-block-strong-projections}
and its immediate restriction counterpart. Apply the first assertion
to that section and its retained reduction.
\end{proof}

We next record a limitation of certificates that stop at protected missing
cubes. A \emph{protected set} in a family $L$ is a set satisfying the
hypotheses of Proposition~\ref{prop:damp-protected-uso-cycle}; equivalently,
no first deletion from that set can be a corner deletion. The union of
all protected sets is its \emph{protected core}. A corner-choice proof
tree contains one child for every available corner deletion and stops
at a node only when its remaining family has a protected set. For a
nonpeelable ample family $L$, let $d_{\rm prot}(L)$ be the least possible
maximum number of deletions on a branch of such a tree. This number is
finite: a maximal corner-deletion sequence ends at a nonempty cornerless
family, which is itself protected. The following result concerns this
specific certificate measure, not the length of arbitrary proofs.

\begin{theorem}[Unbounded depth before a protected certificate]
\label{thm:damp-unbounded-protected-depth}
For every integer $m\ge1$ there is an ample forbidden pc-minor $K_m$ with
\[
 |K_m|=560\,2^m-493,\qquad
 \operatorname{idim}(K_m)=m+23,\qquad
 \operatorname{VC}(K_m)=m+5,
 \qquad d_{\rm prot}(K_m)\ge2^m.
\]
All these families have empty protected core. They are coordinate-block
replacements of one fixed $627$-vertex forbidden class. For $m>1$ they
have negative coordinate reductions, so the assertion does not establish
unbounded certificate depth among reduction-terminal inputs or among
bases fixed by column compression.
\end{theorem}
\begin{proof}
We first give two elementary certificate observations, then show that
the fixed construction preserves them under arbitrarily large blocks.
If $L$ peels relatively to $T$, every protected set of $L$ lies in $T$:
its first vertex could not be a corner when deleted, since all its
specified neighbours would still remain. Its missing witnesses are also
absent from $T$, so it is protected there. Thus relative extensions of
a family with empty protected core have empty protected core.

An ordering of all vertices of $L$ also certifies empty protected core
if, at each vertex $x$, the directions to its later neighbours span a
cube contained in the \emph{original} $L$. A protected set would violate
this condition at its first vertex in the ordering. This condition does
not require that the cube remain present at that step, and hence does
not assert an actual corner peeling. The restriction of a positive
family's corner order to any subset satisfies this condition, tested
against that positive family.

Here are the fixed inputs, in the binary encoding of
\path{research/root_cube_completion.json}. Let $A\subseteq Q_{12}$ be
the certified positive rooted factor of order $289$, whose only corner
is $0$, with incident directions $D_A=\{1,2,8\}$. Let $B$ be its copy
on coordinates $12,\ldots,23$, with incident directions
$D_B=\{13,14,20\}$. Extend both by zero on the other block, put
$D=D_A\cup D_B$, and write $Q_D$ for the cube supported on $D$ at $0$.
Set
\[
 S=A\cup B,\qquad H=S\cup Q_D,\qquad K=H\cup\{32\}.
\]
The established root-face filling construction proves that $H$ is an
ample forbidden class of order $626$ and VC dimension six. Its full
boundary case and minor proof are in
\path{research/root_cube_completion.tex}, Proposition~3. The compatible
addition at $32$, of acquired support $38$ (directions $1,2,5$), gives
the forbidden $K$ of order $627$, still of VC dimension six; the full
proof and fixed missing-cube witnesses are in
\path{research/protected_core_delay.tex}, Theorem~2. These inputs use
all $24$ coordinates. The new verifier
\path{research/protected_block_delay.py} replays the factor order,
checks the root directions, the positioned cube, the acquired support,
and the projection counts used below. The existing forbiddenness
certificates are reused.

Replace coordinate $0$ by a block $E$ of size $m$, and put
$H_m=W_E(H)$ and $K_m=W_E(K)$. Coordinate $0$ is outside $D$ and is
not incident at the root. Put $P=Q_E\setminus\{1_E\}$. The block
formula gives the exact decomposition
\[
 H_m=W_E(A)\ \cup\ (B\times P)\ \cup\ (Q_D\times P),
\]
where the old coordinate $0$ is erased from the factors on the right.
Both $W_E(A)$ and $B\times P$ are positive, by block and product
invariance. The only old word they share is the root, over the fibre
$P$. Let $M=Q_D\setminus S$ be the mixed root-cube words.

We construct a fixed-original-family ordering for $H_m$. First remove
all fibres $\{x\}\times P$ for $x\in M\cup\{0\}$, using a corner
order of $P$ within each fibre, for example decreasing Hamming weight.
Every old-coordinate neighbour of such a vertex has direction in $D$.
For mixed words this follows directly from the displayed decomposition.
For the root, projection in coordinate $0$ creates no new incident old
direction: if $\{0,f\}\in A$ while $\{0\}\notin A$, isometry of
$A$ forces $\{f\}\in A$, so $f\in D_A$.
The analogous assertion for $B$ is immediate. The new-block directions
to later neighbours span a cube contained in $P$, by the chosen fibre
order. Their union with any old incident directions therefore spans a
cube in $Q_D\times P\subseteq H_m$. Earlier deletions of other fibres
can only remove incident directions, so the same verification applies
throughout this first part of the ordering.

After it, the remaining vertices lie in $W_E(A)$ or $B\times P$,
with no edges between the two remaining parts: their old nonzero
supports lie in disjoint coordinate blocks. Finish with corner orders
of these two positive families, restricted to their remaining vertices.
Each restricted order has its later-neighbour cubes in its original
positive family, hence in $H_m$. This proves that $H_m$ has empty
protected core for every $m$.

The added word $32$ has no neighbour in coordinate $0$ in $K$; its
mate $33$ is absent. Hence
\[
 K_m\setminus H_m=\{32\}\times P,
 \qquad |K_m\setminus H_m|=2^m-1.
\]
The deletion of $32$ from $K$ lifts, by
Proposition~\ref{prop:damp-block-relative-transfer}, to a relative
corner peeling of this entire fibre, ending at $H_m$. Every
intermediate family still peels relatively to $H_m$ and therefore has
empty protected core, by the first observation. Any corner-choice
proof tree must include this branch of $2^m-1$ deletions and cannot
stop even at its endpoint $H_m$. Consequently its depth is at least
$2^m$. In particular $K_m$ itself has empty protected core.

Finally, Theorem~\ref{thm:damp-coordinate-block-replacement} makes
every $K_m$ forbidden. The fixed input has
$|\pi_0H|=559$ and $|H^{\{0\}}|=67$. The tip increases the first
count by one and leaves the second unchanged. Thus the block order
formula gives $|K_m|=560(2^m-1)+67$, and the block dimension formulas
give the displayed dimensions. For $m>1$, reduction in a member of
$E$ gives $K_{m-1}$, which is negative. This proves the stated scope.
\end{proof}

The finite replay checks $m=1,2$ using both the explicit
fixed-original-family order and independent protected-core pruning,
as well as the actual tip-fibre corner deletions. The proof for all
$m$ is the symbolic decomposition above; no large-block enumeration
or exact-depth assertion is needed. The examples rule out any uniform
bound for protected-leaf depth on the full forbidden class, even with
isometric dimension minus VC dimension fixed at $18$. They do not
rule out a method that first performs canonical column compression.

\begin{lemma}[Rooted factors and frozen witnesses under block replacement]
\label{lem:damp-rooted-block-witness-transfer}
Let $(A,0)$ satisfy the three rooted-factor conditions of
Theorem~\ref{thm:damp-cut-vertex-obstructions}, and replace an active
coordinate $z$ by a block $E$. The constant lift of $0$ again makes
$W_E(A)$ a rooted factor.

Suppose also that $0^{\{z\}}\in A$, and that $C\supseteq A$ has,
for every $x\in A\setminus\{0\}$, a set $J_x$ with
\[
 x^{\{f\}}\in A\quad(f\in J_x),\qquad x^{J_x}\notin C.
\]
Then every nonroot vertex of $W_E(A)$ has such a witness with neighbours
in $W_E(A)$ and missing word outside $W_E(C)$. In particular, if $0$
is the only corner of $A$, it is the only corner of $W_E(A)$.
\end{lemma}
\begin{proof}
For a constant lift of a prescribed endpoint, endpoint attainability is
preserved by block replacement. In one direction, lift a relative peeling
to the singleton by Proposition~\ref{prop:damp-block-relative-transfer},
then peel its punctured-cube fibre to the prescribed lift, as in
Corollary~\ref{cor:damp-block-all-sinks}. In the other direction, reduce
along all but one block coordinate. Images of a corner chain remain
ample and change by at most one vertex, giving an order to the old root.
Thus ordinary positivity and failure at the root both persist.

For rooted minimality it suffices to split $z$ into two coordinates and
iterate. Old-coordinate root-preserving minors commute with the block
construction and inherit their prescribed-root orders. A new-coordinate
contraction is $(\pi_z A)\times Q_1$. A new root-side restriction is a
peripheral expansion with base $\pi_z A$ and site $\rho_z^0 A$;
both inputs peel to the root by rooted minimality of $A$. The product
and prescribed-endpoint peripheral criteria of
Proposition~\ref{prop:damp-rooted-operations} therefore give the required
root orders. These are all elementary root-preserving minors.

For the witness assertion, write a nonroot lifted vertex as $v=(y,t)$.
Choose $x=(y,s)\in A\setminus\{0\}$ whose expanded singleton contains
$v$. For constant $t$ use its corresponding preimage. For nonconstant
$t$ either existing preimage works; when $y=0$, use the nonroot
preimage $0^{\{z\}}$, which exists by assumption. Set
$r=s\mathbin\oplus\mathbf1_{z\in J_x}$, and put
\[
 L_v=(J_x\setminus\{z\})\cup\{e\in E:t_e\ne r\}.
\]
Each old-direction neighbour in $L_v$ is a lift of $x^{\{f\}}\in A$.
If $z\in J_x$, both old preimages are present, so the full block fibre
supplies all required new-direction neighbours. Otherwise these flips
move toward the available constant $s$-pole and cannot reach the missing
opposite pole in one step. The opposite vertex $v^{L_v}$ has constant
block trace $r$ and projects to $x^{J_x}\notin C$, so is absent from
$W_E(C)$. This proves the witness assertion. For the last conclusion
apply it with $C=A$; the root itself is a corner because its incident
cube is the old root cube with $z$ replaced by the full block.
\end{proof}

\begin{theorem}[Unbounded protected depth among reduced terminal inputs]
\label{thm:damp-terminal-unbounded-protected-depth}
For every $m\ge1$ there is an ample forbidden pc-minor $L_m$ with
\[
 |L_m|=253\,2^m+121,\qquad
 \operatorname{idim}(L_m)=m+23,\qquad
 \operatorname{VC}(L_m)=m+5,
 \qquad d_{\rm prot}(L_m)\ge2^m.
\]
Every one-coordinate reduction of $L_m$ belongs to $\Damp$.
Moreover $L_m$ has minimum degree at least three and pairwise distinct
signed minimal-nonface columns up to sign. Thus it is fixed by both
canonical reductions of Theorem~\ref{thm:damp-alternating-reductions}.
Its protected core is empty.
\end{theorem}
\begin{proof}
Use the fixed factors $A,B$, root directions $D_A,D_B$, and tip $32$
from Theorem~\ref{thm:damp-unbounded-protected-depth}. This time replace
coordinate $z=8$ \emph{inside $A$ only}, by a block $E$ of size $m$.
Put
\[
 A_m=W_E(A),\quad D_m=(D_A\setminus\{8\})\cup E,\quad
 S_m=A_m\cup B,\quad H_m=S_m\cup Q_{D_m\cup D_B},
 \quad P=Q_E\setminus\{1_E\}.
\]
The rooted-block lemma makes $A_m$ a rooted factor with unique corner
$0$. The established full root-face filling theorem therefore makes
$H_m$ forbidden with positive one-coordinate reductions, just as for
$H$ in Theorem~\ref{thm:damp-unbounded-protected-depth}.
This uses Proposition~3 of \path{research/root_cube_completion.tex}
with $A_m,B$ in place of $A,B$.
Define the tip set and output by
\[
 T_m=\{(32,t):t\in P\},\qquad L_m=H_m\cup T_m,
\]
where coordinate $8$ is erased from the displayed old word $32$.
The tip's acquired support $I=\{1,2,5\}$ avoids coordinate $8$;
its mate in that coordinate is absent. Thus
$W_E(A\cup\{32\})=A_m\cup T_m$.

The fixed input includes witnesses for all $x\in A\setminus\{0\}$,
with neighbours in $A$ and missing words outside $A\cup\{32\}$.
They are recorded in \path{research/protected_core_delay.json} and
replayed by \path{research/terminal_protected_depth.py}. The rooted-block
lemma lifts them to every nonroot vertex of $A_m$, still missing from
$A_m\cup T_m$. Filling the mixed root cube adds no new pure-factor
word, so the witnesses remain missing from $L_m$. The same fixed
witnesses, shifted to $B$, protect all its nonroot vertices. Consequently
no nonroot vertex of $S_m$ is a corner in any $H_m\cup U$, $U\subseteq T_m$.

We check all remaining first-move possibilities; this proves
nonpeelability for every subset $U$, without any row-ampleness assumption.
Write $M_m=Q_{D_m\cup D_B}\setminus S_m$ for the mixed root words.
They have no neighbours in $T_m$, because the tip uses coordinate $5$
outside $D_m$, whereas each mixed word also uses a coordinate of $B$.
They remain corners with their full root cubes. Deleting $w\in M_m$
protects $S_m$: use the fixed nonroot witnesses and, at $0$, the
support of $w$, whose unit neighbours are root neighbours in $S_m$
and whose opposite vertex $w$ has just disappeared.

The root is a corner exactly when $(32,0_E)\notin U$. If that tip
is present, it and any root neighbour from $B$ span a missing mixed
square at $0$. If the root can be deleted, the entire old set
$H_m\setminus\{0\}$ is protected in the child. For a mixed word
use its support to reach the missing root. An old word not adjacent
to the root retains its fixed nonroot witness. For a root neighbour
$x$ in either factor, choose a direction $g$ outside that factor's
root directions with $x^{\{g\}}$ still in the factor. Such $g$ exists:
otherwise its incident cube would lie in the full root cube and $x$
would be another corner of that factor. Choose a root direction $f$
in the other factor. Both $x^{\{g\}}$ and $x^{\{f\}}$ remain in
$H_m\setminus\{0\}$, whereas $x^{\{g,f\}}$ is mixed and outside
the completed root cube, hence absent even from $L_m$.
These witnesses cover every old word after root deletion.

Any other first corner lies in $U$ and its deletion gives a smaller
subset extension. Induction on $|U|$, starting at the negative $H_m$,
therefore proves that every $H_m\cup U$ is nonpeelable.

Each tip is individually an ample addition to $H_m$ with acquired
support $I$. Its punctured $I$-cube lies in $A_m$, and $I$ remains
unshattered in $H_m$: the pure-factor missing trace has coordinate
$5$ equal to one, while the completed root cube and $B$ have it zero.
There are no additional old neighbours. This last fact also follows
from isometry of $A\cup\{32\}$: since the tip has no $8$-mate,
a projected neighbour in a new direction would require an old neighbour
in that direction. The only equal-support row is $P$, which is positive.
Theorem~\ref{thm:damp-equal-support-extensions} now makes $L_m$ ample
and supplies a relative peeling to $H_m$, removing all $2^m-1$ tips.
Every intermediate family is ample and negative. Reversing this chain
and using the one-concept atlas gives forbiddenness at every stage.
Positive coordinate reductions also persist: at each ample one-word
addition a reduction stays unchanged or gains a removable corner.
Thus $L_m$ is reduction-terminal.

The base $H_m$ has empty protected core. To verify this, first remove
its root and mixed words in a fixed-original-family pruning order:
all their incident cubes lie in the original completed root cube.
Then use the positive factor orders restricted to their nonroot
vertices, which have no edges between factors. This is exactly the
empty-core ordering criterion proved in
Theorem~\ref{thm:damp-unbounded-protected-depth}. Every stage of the
tip-fibre relative peeling also has empty protected core, since it
peels relatively to $H_m$. A corner-choice tree must follow that
branch and cannot stop at its endpoint, proving $d_{\rm prot}(L_m)\ge2^m$.

It remains to verify the size and the claimed reduced status. The fixed
factor has $|\pi_8A|=224$ and $|A^{\{8\}}|=65$, so
$|A_m|=224(2^m-1)+65$. Its root has $m+2$ incident directions, while
$B$ has three. The mixed filling contributes
$7(2^{m+2}-1)$ words. Adding the $2^m-1$ tips gives
$|L_m|=253\,2^m+121$. All $m+23$ coordinates are active. The completed
root cube has dimension $m+5$; the two factor shadows have dimensions
$m+2$ and three, and the tip row adds supports $I\cup J$ of size at
most $m+2$. Hence the VC dimension is exactly $m+5$.

Each rooted factor is two-connected and has no degree-two nonroot
vertex: in an ample family such a vertex is a corner or a cut vertex.
Their roots have degrees $m+2$ and three. Old degrees only increase
on assembly; mixed vertices have degree $m+5$, and every tip has at
least its three incident directions in $I$. Thus the minimum degree
is at least three.

Finally there are nine coordinates of each factor outside its root
cube. For every left coordinate $f$ and every right coordinate $q$
outside $D_B$, the pair $\{f,q\}$ is a minimal nonface, missing trace
$11$. Likewise $\{p,g\}$ is such a minimal nonface for every left
$p$ outside $D_m$ and every right $g$. These pairs distinguish any
two columns: for two left columns use a right coordinate outside $D_B$;
for two right columns use a left coordinate outside $D_m$; for a left
column and a right column $g$, choose a right coordinate outside
$D_B$ different from $g$. In each case a row is nonzero in one column
and zero in the other. No two signed columns agree up to sign.
The canonical reduction theorem now proves the stated fixed-point property.
\end{proof}

The replay checks the fixed data and the cases $m=1,2$, including lifted
frozen witnesses, all old first-move certificates, the tip-fibre order,
empty-core pruning, minimum degree, and the column-separating missing
pairs. The unbounded statement follows from the uniform proof above.
Unlike the preceding block-only family, these examples cannot be
removed by canonical graph or column reduction. They rule out a
uniform protected-leaf depth bound even for reduced, reduction-terminal
inputs. They remain one explicit construction family, not an exhaustive
generator of those inputs.

\begin{corollary}[Prescribed sinks under coordinate-block replacement]
\label{cor:damp-block-all-sinks}
If every vertex of a nonempty ample $H$ is an attainable sink, then the
same holds for $W_E(H)$. If $H$ is an irredundant ample forbidden class
whose every proper pc-minor admits every prescribed sink, then $W_E(H)$
is also forbidden and has the same proper-minor sink property.
\end{corollary}
\begin{proof}
A cube with one vertex removed admits every remaining vertex as a sink.
Indeed, choose a coordinate face containing the desired endpoint and
omitting the missing vertex. Clear the opposite section using an ordinary
order of its punctured cube, with all mates in the retained full face,
and then peel that face to the endpoint. Induction on the dimension
justifies the ordinary punctured-cube order; the dimension-one case is
a singleton.

For $v=(x,t)\in W_E(H)$ choose $r\in H$ projecting to $x$ such that
$v\in W_E(\{r\})$. For a constant $t$ choose its indicated section;
for a nonconstant $t$ choose either available preimage. The expanded
singleton $W_E(\{r\})$ is a punctured block cube, or a singleton when
$|E|=1$. An order of $H$ ending at $r$, together with
Proposition~\ref{prop:damp-block-relative-transfer}, gives a relative
order from $W_E(H)$ to this fibre. Finish at $v$ by the preceding
punctured-cube argument. This proves the first assertion.

For the second, it suffices to replace one coordinate by two and iterate.
Forbiddenness is already Theorem~\ref{thm:damp-coordinate-block-replacement}.
Every old-coordinate elementary minor is the same block replacement of
a proper minor of $H$, so attains every sink by the first assertion.
A new-coordinate contraction is $C\times Q_1$, and a new-coordinate
restriction is a peripheral expansion with base $C$ and site $B_0$ or
$B_1$. These three inputs are proper pc-minors of $H$ and attain every
sink. Products and the prescribed-endpoint peripheral criterion in
Proposition~\ref{prop:damp-rooted-operations} make each new image
all-sinks as well. Rooted transport covers every longer proper minor,
and iteration proves the assertion for every block size.
\end{proof}

In particular, the inherited branch with neither section relatively
peelable is nonempty. Take an ample forbidden $H$ and a coordinate
$z$ whose reduction $A$ is positive. Neither $B_b$ peels to $A$, by
Proposition~\ref{prop:damp-yang-two-cone-gluing-dichotomy}. Yet in
$W_{\{e,f\}}(H)$ the $e$-reduction is the negative family $H$, and
neither $e$-section peels to it by
\eqref{eq:damp-block-relative-new-coordinate}. The block theorem
already supplies forbiddenness. More simply, for cornerless $H$ the
replacement is cornerless, so an exclusive corner in either section
would give an impossible first corner of the whole replacement.

For the certified $267$-concept cornerless seed, coordinate zero has
section orders $149,118$, overlap order $59$, and contraction order
$208$. Its binary block replacement has $683$ concepts in $13$
coordinates, negative overlap of order $267$, and exclusive wings of
orders $90$ and $59$. Neither section has a corner outside the overlap.
The checker \path{research/block_relative_transfer.py/json} replays all
$36$ seed minor orders, checks ampleness and cornerlessness of the
replacement, and supplies positive orders of its two sections.
More generally block size $m\ge2$ gives such examples of order
$208\cdot2^m-149$ in dimension $11+m$.
These are applications of the existing block generator, not a new
family claimed outside it. They refute exhaustiveness of the
one-relative-wing rule alone. Duplicated-coordinate examples alone do not decide whether the same
failure occurs after both canonical reductions. Proposition~\ref{prop:damp-block-row-completion}
below answers that question affirmatively by a single nonconstant-row
completion. Intrinsic generation of all reduced forbidden bases remains open.

\begin{theorem}[Arbitrary enlargement of one nonconstant row]
\label{thm:damp-one-row-enlargement}
Let $H\subseteq Q_{F\cup\{z\}}$ be nonempty and ample, with
$z$-sections $B_0,B_1$ and contraction $C=B_0\cup B_1$.
Let $D\subseteq Q_F$ be ample with $C\subseteq D$. On two fresh
coordinates $(e,f)$ define
\begin{equation}
 V(H,D)=(B_0\times\{00\})\cup(D\times\{01\})
       \cup(C\times\{10\})\cup(B_1\times\{11\}).
 \label{eq:damp-one-row-enlargement}
\end{equation}
Then $V(H,D)$ is ample, its $e$-reduction is $H$ with $z$ renamed $f$,
and
\begin{align}
 V(H,D)\searrow D\times\{01\}
   &\quad\Longleftrightarrow\quad H\in\Damp,
   \label{eq:damp-one-row-relative}\\
 V(H,D)\in\Damp
   &\quad\Longleftrightarrow\quad H,D\in\Damp.
   \label{eq:damp-one-row-positive}
\end{align}
If $H$ uses all displayed coordinates, then
\begin{equation}
 V(H,D)\in\Obs(\Damp)
 \quad\Longleftrightarrow\quad
 H\in\Obs(\Damp)\ \hbox{and}\ D\in\Damp.
 \label{eq:damp-one-row-forbidden}
\end{equation}
Thus any positive ample superset $D$ of the contraction gives a sound
inverse step from a forbidden $H$. No relative peeling from $D$ to $C$
is required. The inputs are recovered uniquely for the marked, oriented
pair $(e,f)$: read the four rows, require that row $10$ equal the union
of rows $00,11$, and that row $01$ contain that union.
\end{theorem}

\begin{proof}
We first prove ampleness by counting shattered supports. We then lift
an ordinary peeling of $H$ while keeping the whole enlarged row fixed.
Applying this lift to proper minors supplies forbidden-minor minimality.

Put $A=B_0\cap B_1$. Since $H$ is ample,
$\operatorname{Sh}(B_0)\cap\operatorname{Sh}(B_1)=\operatorname{Sh}(A)$ and all sections, the contraction,
and the reduction are ample. The supports of $V(H,D)$ containing
neither $e$ nor $f$ are $\operatorname{Sh}(D)$. Those containing exactly one of
them correspond to $\operatorname{Sh}(C)$, in each of the two cases. Those containing
both correspond to $\operatorname{Sh}(A)$. Therefore
\[
 |\operatorname{Sh}(V(H,D))|=|D|+2|C|+|A|
              =|D|+|C|+|H|=|V(H,D)|,
\]
proving ampleness. Intersecting the two $e$-sections gives rows
$B_0,B_1$, so its reduction is the stated copy of $H$.

For the forward construction in \eqref{eq:damp-one-row-relative},
maintain the displayed row form for the remainder $H'$ of a peeling
of $H$, always keeping $D$ fixed. Its contraction $C'$ is contained
in $D$. Let $(x,b)$ be the next corner, and let $I$ be its incident
old directions, excluding $z$.

If $(x,1-b)\in H'$, the corner cube is the full
$I\cup\{z\}$ cube. Delete only $(x,bb)$ from $V(H',D)$.
Its incident cube has support $I\cup\{e,f\}$: the two constant rows
contain the old corner cube, and both nonconstant rows contain it
because $C'\subseteq D$. This is a valid corner deletion, and the
remaining rows are exactly $V(H'\setminus\{(x,b)\},D)$.

If $(x,1-b)\notin H'$, the neighbors of $x$ in $C'$ are exactly its
neighbors in the $b$-section. Indeed, if $x^i$ occurs on the opposite
side, isometry applied to $(x,b)$ and $(x^i,1-b)$ forces $(x^i,b)$,
since the other intermediate vertex $(x,1-b)$ is absent.
The full $I$-cube through $x$ lies in the $b$-section. First delete
$(x,10)$. Its only block neighbor is $(x,bb)$, so its incident cube
is this $I$-cube times the corresponding block edge. Next delete
$(x,bb)$, whose only remaining block neighbor is $(x,01)$ in $D$;
its incident cube is the $I$-cube times this other edge. Both are
corners. The resulting family again has the row form for
$H'\setminus\{(x,b)\}$, with $x$ removed from $C'$ but retained in $D$.
Iteration leaves exactly $D\times\{01\}$.

Conversely, reduce any such relative peeling in coordinate $e$.
Consecutive reductions are ample and either equal or differ by one
concept. Omitting repeated images gives a corner peeling of $H$ to
the empty family, since the retained row has no $e$-edge.
This proves \eqref{eq:damp-one-row-relative}.
Following the constructed relative order by a peeling of $D$ proves
the sufficient direction of \eqref{eq:damp-one-row-positive}.
Its necessity follows from pc-minor closure for the face $D$ and
reduction closure for $H$.

Suppose now $H$ is forbidden and $D$ is positive. The reduction $H$
makes $V(H,D)$ negative. For an elementary operation $\mu$ in an old
coordinate, $\mu H$ and $\mu D$ are positive, and
$\mu V(H,D)=V(\mu H,\mu D)$. Here sections and unions commute with
$\mu$, and the union of the sections of $\mu H$ is $\mu C$.
The positive equivalence proves this old-coordinate image positive.
The $e$-restrictions and contraction are peripheral expansions with
base/site pairs $(D,B_0)$, $(C,B_1)$, and $(D,C)$, respectively.
The corresponding $f$-images have pairs $(C,B_0)$, $(D,B_1)$, and
$(D,C)$. All bases and sites are positive: $B_0,B_1,C$ are proper
pc-minors of $H$. Thus all elementary images are positive, proving
forbiddenness.

Conversely, if $V(H,D)$ is forbidden, its proper face $D$ is positive.
Equation~\eqref{eq:damp-one-row-positive} makes $H$ negative.
Theorem~\ref{thm:damp-inherited-overlap-descent}, applied to its
$e$-reduction, makes $H$ forbidden. This proves
\eqref{eq:damp-one-row-forbidden}. Reading the rows recovers all
marked data and their stated equalities, proving the recovery assertion.
\end{proof}

\begin{corollary}[A closed formula for simultaneous row enlargements]
\label{cor:damp-simultaneous-row-enlargement}
Let $\varnothing\ne H\subseteq D\subseteq Q_F$ be ample, and choose
$\varnothing\ne K\subseteq F$. Replace every $i\in K$ by an ordered
pair $(e_i,f_i)$. For $t\in\{00,01,10,11\}^{K}$ and $A\subseteq Q_F$,
define the row
\[
 A[t]=\{x\in Q_{F\setminus K}:\text{ some }a\in A\text{ has }
 a|_{F\setminus K}=x\text{ and }a_i=b\text{ whenever }t_i=bb\}.
\]
Thus nonconstant pairs leave the corresponding old coordinate free.
Define $\mathcal V_K(H,D)$ by the explicit rows
\begin{equation}
 \mathcal V_K(H,D)_t=
 \begin{cases}
 H[t],&t_i\ne01\text{ for every }i\in K,\\
 D[t],&t_i=01\text{ for at least one }i\in K.
 \end{cases}
 \label{eq:damp-simultaneous-row-formula}
\end{equation}
Write $\pi_S$ for deletion of all coordinates in $S$. Then
$\mathcal V_K(H,D)$ is ample and
\begin{equation}
 \mathcal V_K(H,D)\in\Damp
 \quad\Longleftrightarrow\quad
 H\in\Damp\ \text{and}\ \pi_iD\in\Damp\text{ for every }i\in K.
 \label{eq:damp-simultaneous-row-positive}
\end{equation}
If $H$ uses every coordinate of $F$, then also
\begin{equation}
 \mathcal V_K(H,D)\in\Obs(\Damp)
 \quad\Longleftrightarrow\quad
 H\in\Obs(\Damp)\ \text{and}\ \pi_iD\in\Damp\text{ for every }i\in K.
 \label{eq:damp-simultaneous-row-forbidden}
\end{equation}
The construction can be performed by the one-row rule in any order of
$K$, with the common completion $D$ determining every enlarged row.
Its order is
\begin{equation}
 |\mathcal V_K(H,D)|=
 \sum_{S\subseteq K}\bigl(|\pi_SH|+(2^{|S|}-1)|\pi_SD|\bigr).
 \label{eq:damp-simultaneous-row-order}
\end{equation}
The marked output recovers $H$ and each $\pi_iD$. These recovered
families determine the output uniquely, although they need not determine
$D$ itself.
\end{corollary}

\begin{proof}
Put $\mathcal W_J(A)=\mathcal V_J(A,A)$, with
$\mathcal V_\varnothing(H,D)=H$. Its rows are the ordinary simultaneous
two-coordinate block replacement of $A$: each constant pair prescribes
the old bit, and each nonconstant pair allows either bit. Replacing
distinct coordinates in any order gives these same rows. Consequently
Theorem~\ref{thm:damp-coordinate-block-replacement} gives ampleness and
the equivalence $\mathcal W_J(A)\in\Damp\iff A\in\Damp$.

Suppose $J\subsetneq K$ has already been replaced and $i\in K\setminus J$.
The row formula gives
\[
 \mathcal V_J(H,D)\subseteq\mathcal W_J(D),\qquad
 \pi_i\mathcal V_J(H,D)\subseteq
 \pi_i\mathcal W_J(D)=\mathcal W_J(\pi_iD).
\]
Apply Theorem~\ref{thm:damp-one-row-enlargement} to the current family
$\mathcal V_J(H,D)$, enlarging its new $01$ row to
$\mathcal W_J(\pi_iD)$. The two constant rows retain the old bit of $i$,
the $10$ row forgets it in the current family, and the $01$ row uses $D$
regardless of every other pair. These are exactly the rows in
\eqref{eq:damp-simultaneous-row-formula} for $J\cup\{i\}$.
This proves the iteration identity, its independence of order, and
ampleness by induction. The positive equivalence in the one-row theorem
adds precisely the condition $\mathcal W_J(\pi_iD)\in\Damp$, which is
equivalent to $\pi_iD\in\Damp$. Induction proves
\eqref{eq:damp-simultaneous-row-positive}. When $H$ has full coordinate
support, every intermediate family does too on its displayed coordinates.
Applying the forbidden equivalence in the same theorem at every step
proves \eqref{eq:damp-simultaneous-row-forbidden} in both directions.
In particular, $D$ itself need not be dismantlable.

For the cardinality formula fix the set $S$ of nonconstant pairs. There
is one choice with all these pairs equal to $10$, which uses $H$;
the other $2^{|S|}-1$ choices use $D$. Summing the row sizes over all
assignments to the constant pairs gives respectively $|\pi_SH|$ and
$|\pi_SD|$. Summing over $S$ proves
\eqref{eq:damp-simultaneous-row-order}.

Reduction in all coordinates $e_i$, $i\in K$, recovers $H$, with each
$f_i$ renamed $i$. Indeed, the all-constant choice in such a full fiber
requires membership in $H$; conversely membership in $H\subseteq D$
supplies every member of that fiber. Restricting one pair to $01$ gives
$\mathcal W_{K\setminus\{i\}}(\pi_iD)$, and reducing its remaining
$e_j$ coordinates recovers $\pi_iD$. Finally, any row using $D$ has a
pair $t_i=01$, so its membership can be read from $\pi_iD$ after
forgetting any other nonconstant coordinates. Thus the recovered
$H$ and selected projections determine every row of the output.
\end{proof}

This is a nonrecursive description of a common-completion subclass of
iterated one-row enlargements. It reorganizes that proved operation;
it does not enlarge its closure or generate the unknown primitive seeds.
For $|K|=1$, any enlarged row $D_0$ in the one-row theorem is recovered
by taking the common completion $D=D_0\times Q_1$. Taking $D=H$
recovers ordinary coordinate-block replacement. No assertion is made
that arbitrary independently chosen enlarged rows admit a common ample
completion, or that every forbidden obstruction has this presentation.

The limitation on recovery is essential even for forbidden inputs.
If $a\notin H$, the ample completions $Q_F\setminus\{a\}$ and $Q_F$
have the same one-coordinate projections and give identical marked
outputs. Nor must the largest family having the recovered projections
be ample. For example, encode $Q_4$ by integers with coordinate zero
least significant, and take
\[
 D=Q_4\setminus\{0,1,3,7,15\},\qquad K=\{0,3\}.
\]
Here $D$ shatters exactly the eleven supports of size at most two, so
it is ample. The family
$\bigcap_{i\in K}\pi_i^{-1}(\pi_iD)=D\cup\{3\}$ has twelve
vertices but thirteen shattered supports: the two additional supports
are $\{0,1,3\}$ and $\{0,2,3\}$. Hence it is not ample, despite
the given ample completion $D$. Thus testing that maximal intersection
for ampleness is not a valid test for the existence of an ample
completion: it would discard these admissible projection data.
The check \path{research/simultaneous_row_enlargement.py/json} verifies
the formula, all iteration orders, cardinality and marked recovery for
all nested nonempty ample pairs on at most three coordinates, as well
as this four-coordinate example. The general soundness proof above
uses the established block and one-row theorems.

\begin{remark}[A row enlargement with no relative first move]
\label{ex:damp-one-row-nonrelative}
Take the certified $42/22$ pair $O\subset P\subset Q_6$ from
Proposition~\ref{prop:damp-six-coordinate-relative-failure}, and put
$H=(P\times O)\cup(O\times P)$, the $1364$-vertex rectangular
obstruction. Replace a coordinate $z$ of the first block. Writing
$P'=\pi_zP$ and $O'=\pi_zO$, choose
\[
 C=(P'\times O)\cup(O'\times P),\qquad D=P'\times P.
\]
The maximum-class contraction counts give $|P'|=26$, $|O'|=16$,
so $|C|=892$, $|D|=1092$, and $|V(H,D)|=3348$ in thirteen
coordinates. The factors of $D$ are positive proper minors of the
certified inputs, or the positive input $P$ itself, so $D$ is positive.
Every corner of $P$ lies in $O$. Product corners therefore give
$\operatorname{Cor}(D)\subseteq P'\times O\subseteq C$.
Since $C\ne D$, no relative peeling from $D$ to $C$ can start.
Nevertheless Theorem~\ref{thm:damp-one-row-enlargement} makes
$V(H,D)$ forbidden.

This output is cornerless. The original block replacement $W(H)$ is
cornerless, and every one of its vertices has a block-direction neighbor.
If an old vertex became a corner after enlarging row $01$, its corner
cube would contain a new row vertex and that vertex's mate in such a
block direction. The mate is absent, since the new old word lies outside
$C$. A new vertex has only old-coordinate neighbors, and could be a
corner only if it belonged to $\operatorname{Cor}(D)\setminus C$,
which is empty. Thus this enlargement cannot be obtained from $W(H)$
by any nonempty sequence of ample one-concept additions: the last added
vertex would be a corner of the final family.

The script \path{research/one_row_enlargement.py} verifies the lifting
rule on all $357$ small ample input pairs with $H\subseteq Q_3$ and
$D\subseteq Q_2$, and constructs orders for all $39$ elementary minors
of the $3348$-vertex output. The independent replay
\path{research/verify_one_row_enlargement.py} recomputes its shattering
count, cornerlessness, and all minor USO certificates.
This strictly extends the single-row corner-addition mechanism; no claim
is made that the example lies outside every other known construction.
\end{remark}

\begin{proposition}[Two enlarged rows: ampleness and a transfer obstruction]
\label{prop:damp-two-row-transfer-obstruction}
Let $H\subseteq Q_{F\cup\{z\}}$ be ample, with sections $B_0,B_1$
and contraction $C=B_0\cup B_1$. For arbitrary supersets
$D_0,D_1\subseteq Q_F$ of $C$, put
\[
 V=(B_0\times\{00\})\cup(D_0\times\{01\})
   \cup(D_1\times\{10\})\cup(B_1\times\{11\}),\qquad
 K=(D_1\times\{0\})\cup(D_0\times\{1\}).
\]
Then $V$ is ample if and only if $K$ is ample and $D_0\cap D_1=C$.
In this case $V$ has reduction $H$ and contraction $K$ in the first
new coordinate. Nevertheless, $H,K\in\Damp$ does not imply
$V\in\Damp$: the certified $267$-vertex obstruction has such a
presentation with $|H|=59$, $|K|=208$, and $|C|=48$.
\end{proposition}
\begin{proof}
The contraction in the first new coordinate is $K$. If $V$ is ample,
then $K$ is ample. An $x\in(D_0\cap D_1)\setminus C$ would give,
by fixing the old coordinates to $x$, a two-vertex diagonal of a square
as a face of $V$, contrary to ampleness. Hence $D_0\cap D_1=C$.

Conversely, suppose $K$ is ample and its section intersection is $C$.
Put $A=B_0\cap B_1$. The ample section identities for $H$ and $K$
give
\[
 \operatorname{Sh}(B_0)\cap\operatorname{Sh}(B_1)
 =\operatorname{Sh}(A),\qquad
 \operatorname{Sh}(D_0)\cap\operatorname{Sh}(D_1)
 =\operatorname{Sh}(C).
\]
The shattered supports of $V$ avoiding the new coordinates are those
of $D_0\cup D_1$; those using exactly one new coordinate correspond
to $\operatorname{Sh}(C)$, twice; those using both correspond to
$\operatorname{Sh}(A)$. Therefore
\[
 |\operatorname{Sh}(V)|=|\operatorname{Sh}(K)|+|C|+|A|
 =|K|+|H|=|V|.
\]
This proves ampleness. Intersecting the two sections in the first new
coordinate gives the claimed reduction $H$.

For the counterexample use the labelled $Y_{267}$ of
Proposition~\ref{prop:damp-order-267-obstruction}. On its coordinates
$0,1$, flip coordinate $1$ and read the four rows. Their sizes in order
$00,01,10,11$ are $13,136,72,46$. Their sets satisfy
$B_0\cup B_1=D_0\cap D_1$, of order $48$, and
$|B_0\cap B_1|=11$. Thus the displayed reduction and contraction
have orders $59$ and $208$. Explicit corner orders for both are
replayed as acyclic USO certificates in
\path{research/verify_two_row.py}, with the data in
\path{research/two_row_verification.json}. The same verifier checks
that $Y_{267}$ is ample and cornerless. This proves the stated failure.
\end{proof}

The failure is already present inside the two-row presentation, rather
than only in an unrestricted expansion. It prevents extending the
one-row proof by assuming automatic positivity transfer from $H$ and $K$.
It does not refute the inverse conjecture with a common \emph{forbidden}
reduction: here $H$ is positive. The ampleness equivalence is checked on
all $1739$ small input triples with $H\subseteq Q_3$ and
$D_0,D_1\subseteq Q_2$; the proof above has no dimension bound.

There are six such unordered coordinate pairs in this labelled
$Y_{267}$. They also give a bounded shrinking test: keep $B_0,B_1$
fixed and delete from either $D_i$ one nonempty coordinate-trace fiber
disjoint from $C$. The scripts
\path{research/two_row_fiber_shrink.py} and
\path{research/verify_two_row.py} enumerate $7148$ distinct outputs,
allowing deletions of up to $56$ words. Every output is nonample.
For each one the independent verifier supplies a support shattered but
not strongly shattered; the separate shadow-cardinality check in
\path{research/two_row_fiber_positivity.py} agrees, including the
$5886$ outputs with a corner. This excludes these batch deletions, not
arbitrary deletions, simultaneous changes of both rows, or repairs by
insertions. The minimum-order interval remains $64$--$267$.

\begin{lemma}[Corners in a two-row assembly]
\label{lem:damp-two-row-corner-compatibility}
Use the ample two-row assembly of
Proposition~\ref{prop:damp-two-row-transfer-obstruction}. For $x\in C$ put
\[
 T_x=\bigcap_{b\in\{0,1\}:\,x\in B_b} B_b.
\]
For a corner $(x,b)$ of $K=(D_1,D_0)$, let $M_K(x,b)$ be the projection
of its corner cube onto the old coordinates $F$. The corners of $V$
are exactly the following vertices:
\begin{enumerate}
 \item $(x,b,b)$, where $(x,b)$ is a corner of $H$ incident with $z$;
 \item $(x,1-b,b)$, where $(x,b)$ is a corner of $K$ and either
 $x\notin C$ or $M_K(x,b)\subseteq T_x$.
\end{enumerate}
Thus $V$ is cornerless exactly when the first set is empty and every
corner of $K$ lies over $C$ with its projected cube not contained in $T_x$.
\end{lemma}
\begin{proof}
A vertex in constant row $bb$ has both new-coordinate neighbours,
since $B_b\subseteq D_0\cap D_1=C$. Its full incident cube exists
precisely when its old-coordinate incident cube in $B_b$ exists and
lies in the other section $B_{1-b}$. This is exactly the corner
condition at $(x,b)$ in $H$ together with incidence with $z$.

Consider a vertex in row $(1-b,b)$. If $x\notin C$, it has no
new-coordinate neighbour in either $V$ or $K$, so the two corner
conditions coincide. If $x\in C$, it has a neighbour in each constant
row indexed by a section containing $x$. Its full incident cube requires
the old-coordinate cube in $D_{1-b}$ to lie in every such section.
This is precisely containment in $T_x$. Such containment implies the
cube lies in $C$, so it also gives the mate cube needed for $(x,b)$
to be a corner of $K$. Conversely, a corner of $K$ with projected cube
inside $T_x$ supplies every required vertex of the incident cube in $V$.
These four rows exhaust its vertices.
\end{proof}

This is a first-corner condition, not an acyclic-USO gluing theorem.
In CCMW's language, a source of an acyclic USO must be a corner because
its full incident direction set is its outgoing cube. The lemma detects
when no source can exist; it does not assert that a permitted corner
extends to an acyclic USO. For the marked $267$ example, $H$ has three
corners, none incident with $z$. The contraction $K$ has two corners;
both lie over $C$, and each projected corner cube fails containment in
$T_x$. These are the five excluded possibilities for a first move.
The script \path{research/two_row_corner_compatibility.py} records the
two missing-cube witnesses and checks the formula on all $610$ ample
assemblies in the preceding small test domain.

\begin{corollary}[Two enlarged rows with one relative peeling]
\label{cor:damp-two-row-relative-enlargement}
Use the ample two-row presentation of
Proposition~\ref{prop:damp-two-row-transfer-obstruction}. Suppose that
$H$ is a full-support forbidden pc-minor and $K\in\Damp$.
If $D_0\searrow C$ or $D_1\searrow C$, then $V$ is forbidden.
In particular, it suffices that every connected component of the graph
$2$-core of $G(D_i\setminus C)$ have at most eight vertices for one
$i\in\{0,1\}$.
Consequently, any counterexample to the two-row inverse under these
hypotheses must have a component of its $2$-core of order at least nine
in each of $G(D_0\setminus C)$ and $G(D_1\setminus C)$.
\end{corollary}
\begin{proof}
First suppose $D_1\searrow C$. Delete the corresponding concepts
only in row $10$. Each such concept lies outside $C$, so has no
neighbour in either constant row; the opposite nonconstant row is
at distance two in the new coordinates. Its incident cube in $V$
is therefore precisely its incident cube in the current $D_1$.
The supplied relative order consists of valid corner deletions and leaves
\[
 V_0=(B_0\times\{00\})\cup(D_0\times\{01\})
      \cup(C\times\{10\})\cup(B_1\times\{11\}).
\]
Since $D_0$ is a section of positive $K$, it is positive.
Theorem~\ref{thm:damp-one-row-enlargement} makes $V_0$ forbidden.
Both $V$ and $V_0$ have the same negative reduction $H$ in coordinate
$e$. Proposition~\ref{prop:damp-fixed-negative-reduction-extension}
therefore makes $V$ forbidden. Explicitly, any elementary minor of the
relative chain from $V$ to $V_0$ is an ample chain with successive
images equal or differing by one concept. Removing repetitions and
appending a peeling of the proper minor of $V_0$ proves positivity.
Full support makes each such minor of $V_0$ proper.

If instead $D_0\searrow C$, delete its exclusive concepts in row
$01$, leaving rows $(B_0,C,D_1,B_1)$. Exchanging $e$ and $f$ gives
the preceding one-row construction with enlarged row $D_1$.
This exchange fixes both constant rows and swaps the two nonconstant
rows. The same extension argument proves forbiddenness.

The relative small-core argument in the proof of
Corollary~\ref{cor:damp-small-component-wing-inverse} applies to any
nested nonempty ample pair, hence to $(C,D_i)$. It supplies the required
relative peeling under the displayed graph condition. Taking the
contrapositive on both sides gives the necessary condition for a
counterexample.
\end{proof}

This is a sufficient condition for the two-row inverse, obtained by
combining the proved one-row rule with relative-extension transport.
It does not enlarge the closure of those operations. In particular,
positivity of $K$ alone still does not supply either relative peeling.
The graph condition is on the enlarged rows over the positive contraction
$C$, not on the sections of $V$ over its negative reduction $H$.

\begin{remark}[A product completion cannot supply two failing rows]
\label{rem:damp-product-two-row-eligibility}
The relative-attachment theorem already rules out a natural attempt to
escape the preceding sufficient condition. Let $0\in P\subseteq T$
be ample families on one coordinate block, both peeling to $0$.
On a disjoint block let $R$ be ample and peelable, with unique corner
$0$, and suppose $R$ does not peel to $0$. Write
\[
 P\vee R=(P\times\{0\})\cup(\{0\}\times R),\qquad
 C=P\vee R,\quad D_0=P\times R,\quad D_1=T\vee R.
\]
Assume $|P|,|R|>1$. Both enlarged rows and $C$ are positive and ample,
and $D_0\cap D_1=C$. The product row cannot peel to $C$, since
\[
 \operatorname{Cor}(D_0)
 =\operatorname{Cor}(P)\times\{0\}\subseteq C
 \quad\hbox{and}\quad D_0\ne C.
\]
Nevertheless this does not yield eligible rows with both relative
peelings failing: their contraction satisfies exactly
\[
 K=(D_1\times\{0\})\cup(D_0\times\{1\})\in\Damp
 \quad\Longleftrightarrow\quad T\searrow P
 \quad\Longleftrightarrow\quad D_1\searrow C.
\]

To see this, on the first block and a fresh coordinate put
$B=(T\times\{0\})\cup(P\times\{1\})$ and
$L=\{(0,0)\}\cup(P\times\{1\})$. These are ample peripheral
families, and $L$ is positive. After permuting blocks,
$K=P_R(B,L)$ in the notation of
Theorem~\ref{thm:damp-relative-attachment-test}. Thus $K$ is positive
exactly when $B\searrow L$. Projecting any such relative order along
the fresh coordinate gives $T\searrow P$: omit repeated images,
and each remaining step is a corner deletion between ample families.
Conversely, first apply $T\searrow P$ in layer zero of $B$;
then apply $P\searrow\{0\}$ in that layer, keeping layer one fixed.
The incident cubes are respectively the old corner cubes and their
products with the fresh edge, so all deletions are valid and leave $L$.
Finally, $D_1\searrow C$ implies $T\searrow P$ by restricting to
the zero face of the $R$ block. The converse deletes $T\setminus P$
in its relative order; these vertices have no neighbours with a nonzero
$R$ coordinate. This proves all equivalences.

This is an application of the existing relative-attachment test, not
a proof of the unrestricted two-row inverse. It excludes this entire
product-completion approach to producing two failing relative rows with
a positive contraction, independently of the selected reduced input.
\end{remark}

For cornerless reduced inputs the same observation describes all possible
corner deletions, rather than just a supplied relative order.
\begin{lemma}[Corner deletions with a fixed cornerless reduced input]
\label{lem:damp-two-row-corner-deletion-factorization}
In the ample two-row presentation, suppose $H$ is cornerless. Every
corner of $V$ lies in a nonconstant row outside $C$, and
\[
 \operatorname{Cor}(V)=
 ((\operatorname{Cor}(D_0)\setminus C)\times\{01\})\cup
 ((\operatorname{Cor}(D_1)\setminus C)\times\{10\}).
\]
Every corner-deletion sequence of $V$ is an interleaving of relative
corner-deletion sequences in $D_0,D_1$ retaining $C$. Conversely every
such interleaving is a corner-deletion sequence in $V$. Thus a maximal
sequence ends exactly when neither remaining row has a corner outside
$C$. This does not assert uniqueness of the remaining rows.
\end{lemma}
\begin{proof}
Constant-row corners would give corners of $H$ by
Lemma~\ref{lem:damp-two-row-corner-compatibility}, and hence are absent.
Suppose a nonconstant-row vertex over $x\in C$ were a corner.
Its old-coordinate incident cube $Q$ lies in every section $B_b$
containing $x$, by that lemma. Every old neighbour of $x$ in any such
$B_b$ is also a neighbour in the enlarged row; conversely all neighbours
of the enlarged row lie in $Q\subseteq B_b$. These neighbour sets
therefore coincide. If $x$ belongs to only one $B_b$, $Q$ is its full
incident cube in $H$, giving a corner. If it belongs to both sections,
$Q\times Q_1$ is the full incident cube in $H$, again a corner.
Both alternatives contradict the hypothesis. Outside $C$ there are no
block-coordinate neighbours, so the corner condition is exactly the
corner condition in the relevant enlarged row. This proves the formula.

Deleting such a corner retains the constant rows and the intersection
$C$, and leaves an ample assembly with the same cornerless $H$.
Induction applies the formula after every deletion. Deletions in one
nonconstant row do not affect old-coordinate neighbours in the other,
so the relative orders may be interleaved arbitrarily. The criterion
for maximality follows from the same formula.
\end{proof}

For a concrete two-row application, use the $42/22$ pair $O\subset P$
and the $1364$-vertex rectangular input of
Remark~\ref{ex:damp-one-row-nonrelative}. Put
$P'=\pi_0P$, $O'=\pi_0O$ and
\[
 C=(P'\times O)\cup(O'\times P),\quad
 D_0=P'\times P,\quad
 D_1=(Q_5\times O)\cup(O'\times P).
\]
The enlarged rows have orders $1092,1024$ and intersection $C$ of
order $892$. Both strictly contain $C$. A relative order from $Q_5$
to $P'$, combined with a peeling of $O$, removes the $132$ words of
$D_1\setminus C$. The resulting assembly is the certified $3348$-word
one-row example. Hence the $3480$-word two-row output is forbidden by
the preceding relative-extension proof. Its $18$ corners are exactly
those supplied by the lemma, and the displayed order leaves the
cornerless $3348$-word family. The check
\path{research/two_row_relative_extension.py/json} verifies ampleness,
all corners, the relative order, and all $39$ transported elementary-minor
orders. This illustrates the criterion with two strict enlarged rows;
it is not a new primitive seed or a size improvement.

\begin{proposition}[Nested ample triples give two nonrelative enlarged rows]
\label{prop:damp-nested-triple-two-row}
Let $H=(B_0,B_1)\subseteq Q_{F\cup\{z\}}$ be nonempty ample with
full coordinate support, and put $C=B_0\cup B_1$. Choose ample families
$C\subseteq D\subseteq E\subseteq Q_F$. For pairs of families on
$F$, use $(A,B)=(A\times\{0\})\cup(B\times\{1\})$ on a fresh
coordinate $a$. Set
\[
 U_0=(B_0,C),\qquad U_1=(D,B_1),\qquad
 R_0=(D,D),\qquad R_1=(E,C),\qquad J=(D,C).
\]
On two further coordinates $(b,f)$ define the four rows of $Z$ by
\[
 (Z_{00},Z_{01},Z_{10},Z_{11})=(U_0,R_0,R_1,U_1).
\]
Then $Z$ is ample, $U_0\cup U_1=R_0\cap R_1=J$, and
\[
 Z\in\Obs(\Damp)
 \quad\Longleftrightarrow\quad
 H\in\Obs(\Damp),\ D\in\Damp,\ E\in\Damp.
\]
Its $b$-reduction is $H_*=(U_0,U_1)$, a copy of $V(H,D)$, and
\[
 |Z|=|H|+2|C|+3|D|+|E|.
\]
The relative conditions in this two-row presentation satisfy exactly
\[
 R_0\searrow J\iff D\searrow C,\qquad
 R_1\searrow J\iff E\searrow D.
\]
Thus two successive relative failures in a positive nested triple
produce a forbidden two-row assembly with neither enlarged row
relatively peelable to their intersection.
\end{proposition}
\begin{proof}
First construct $X=V(H,D)$ with coordinates $(e,f)$ replacing $z$;
its rows are $(B_0,D,C,B_1)$. Its $e$-contraction has $f$-sections
$(C,D)$. The ample peripheral family $L=(E,D)$, now indexed by $f$,
contains that contraction. Apply the one-row theorem again, to $X$
in coordinate $e$, using enlarged row $L$, and replace $e$ by $(a,b)$.
The eight rows, indexed in order by $(a,b,f)$, are
\[
\begin{array}{c|cccccccc}
(a,b,f)&000&001&010&011&100&101&110&111\\\hline
\text{row}&B_0&D&E&D&C&D&C&B_1.
\end{array}
\]
Reading this table instead as four $(b,f)$-rows, with $a$ retained,
gives exactly $Z$. This proves ampleness and the construction identity.
Direct unions and intersections give $J$ and the claimed reduction.
The latter is the original $X$ with its coordinate $e$ renamed $a$.
Both applications retain full coordinate support. The one-row forbidden
criterion says that the second output is forbidden precisely when $X$
is forbidden and $L$ is positive. The first application says that $X$
is forbidden precisely when $H$ is forbidden and $D$ is positive.
Finally, the peripheral expansion $L=(E,D)$ is positive precisely when
$E,D$ are positive, by the peripheral criterion and face closure.
This proves the equivalence. Summing the eight row sizes gives the
cardinality formula.

A relative order from $D$ to $C$ can be deleted in the $a=1$ layer
of $R_0=D\times Q_1$, with its entire mate layer retained. It leaves
$J=(D,C)$. Conversely, reducing such a relative chain in $a$ yields
an ample deletion chain from $D$ to $C$, after equal consecutive images
are omitted. This proves the first equivalence.
For the second, delete a relative order from $E$ to $D$ only in the
$a=0$ layer of $R_1=(E,C)$. The deleted vertices lie outside $D$
and hence outside $C$, so have no $a$-mates and have exactly their
old corner cubes. Conversely, restricting any relative order from
$R_1$ to $J$ to the $a=0$ face yields a relative order from $E$ to $D$.
This proves the second equivalence.
\end{proof}

The nested-triple rule is an alternative marked presentation of two
successive one-row steps. It does not enlarge their closure. Its
additional point is that a different two-row presentation can have
\emph{both} relative obstructions, so the sufficient relative-row
condition cannot be treated as necessary. The ordered three new
coordinates recover all inputs directly from the eight displayed rows;
no assertion is made about arbitrary two-row presentations.

For an explicit instance, use the certified pair $O\subset P\subset Q_6$
and the $1364$-vertex rectangular $H$, replacing its first coordinate.
Put $P'=\pi_0P$, $O'=\pi_0O$, and take
\[
\begin{split}
 C&=(P'\times O)\cup(O'\times P),\\
 D&=(P'\times P)\cup(Q_5\times O),\\
 E&=Q_5\times P.
\end{split}
\]
Their orders are $892,1224,1344$. A relative order from $Q_5$ to
$P'$, with a peeling of $O$ inside each removed slice, peels $D$
to the positive product $P'\times P$; thus $D$ is positive. The
same ample slice chain also establishes ampleness of $D$: its reverse
runs from $P'\times P$ through unions of the two displayed product
slabs, which are ample by the rectangular-boundary identity.
The product $E$ is positive and ample.

Choose $u\in P'\setminus O'$ and $v\in Q_5\setminus P'$.
The faces at $u$ of $(D,C)$ are $(P,O)$, and those at $v$ of
$(E,D)$ are also $(P,O)$. Since $P$ cannot peel to $O$, relative
minor transport proves both failures. In the recorded integer
coordinates one can take $u=4$ and $v=9$.
The resulting $Z$ has $8164$ vertices in $14$ coordinates, forbidden
reduced input $H_*$ of order $3480$, enlarged-row orders $2448,2236$,
and intersection order $2116$. Its contraction is positive by
forbiddenness; all input eligibility conditions hold. The input $H_*$
has $18$ corners, so this example does not settle the analogous question
restricted to cornerless reduced inputs.

The check \path{research/two_row_both_nonrelative.py/json} verifies the
two construction identities, direct ampleness, both face witnesses,
the supplied positive input orders, and all $42$ elementary-minor
orders of $Z$. The failures use the certified relative obstruction
$(P,O)$, not a failed greedy search. This supplies an explicit instance
of the both-nonrelative regime; it neither improves the minimum order
nor provides an exhaustive inverse rule or a new primitive input.

The fixed reduced input $H_*$ also restricts potential counterexamples
beyond this one construction. Let $Y$ be any ample class on fourteen
active coordinates whose marked reduction is this $H_*$, and whose two
marked sections and marked contraction are positive. Neither the recorded
$Y_{267}$ nor any of the $1364$-vertex rectangular obstructions can be
a proper pc-minor of $Y$. Indeed, each target has twelve active
coordinates, so Proposition~\ref{prop:damp-inverse-proper-obstruction-reduction-test}
would identify a coordinate reduction of the target with an image of
$H_*$ under one or two old-coordinate elementary operations. The target
reduction orders are respectively $59$ and $472$. The exact check
\path{research/two_row_parent_minor_filter.py/json} enumerates all
$39+\binom{13}{2}3^2=741$ labelled choices; neither order occurs.
A separate direct restriction/projection expression agrees with each
successive-operation image. This applies to every such $Y$ with the
fixed reduced input, not just the displayed $8164$-vertex output.
It excludes these known witnesses only, and does not prove the inverse
sound or exclude unknown forbidden proper minors.

A stronger finite exclusion permits arbitrary changes of both diagonal
rows while fixing $D_0,D_1$. For each of the six presentations of the
labelled $Y_{267}$, the only cornerless ample assemblies with
$B_0\cup B_1=D_0\cap D_1$ are the original assembly and the exchange
of its diagonal rows. In particular no smaller cornerless assembly has
that fixed contraction. The companion ample-corners manuscript proves
this computer-assisted rigidity statement, including the exact Boolean
encoding and coverage argument; the source label is
\texttt{prop:record-fixed-contraction-rigidity}. The complete data and
six independently replayed refutations are in
\path{research/two_row_fixed_contraction/}. This statement does not
classify assemblies with changed large rows or cornered negative outputs.

There is also an obstruction to using this particular positive-input
counterexample as a source of forbidden reduced inputs. For its marked
coordinate pair $0,1$ (with coordinate $1$ flipped), put $C=D_0\cap D_1$.
Every ample subfamily of the $96$-vertex family $C\times Q_1$ is
in $\Damp$, including subfamilies with corners and those whose projection
is smaller than $C$. The companion ample-corners manuscript proves this
by an exact all-subfamily encoding and independently replayed DRAT
refutation, at \texttt{prop:record-two-row-overlap-envelope}; evidence
is in \path{../ample-corners/evidence/two_row_overlap_envelope/}.
Therefore changing the diagonal input while retaining this common
intersection cannot produce a forbidden $H$. This exclusion even permits
changes to the large rows that preserve their intersection. It does not
settle the two-row inverse for other intersections or improve the global
$64$--$267$ minimum interval.

The following one-concept case retains additional information about
corners and signed columns that the arbitrary enlargement does not assert.

\begin{proposition}[Completing one nonconstant block row]
\label{prop:damp-block-row-completion}
Let $H$ be an irredundant cornerless ample forbidden pc-minor, and
write $B_0,B_1,C=B_0\cup B_1$ for its sections and contraction at $z$.
Suppose $C\cup\{x\}$ is an ample one-concept extension with acquired
support $P$ of size at least three. In the binary block expansion
$W=W_{\{e,f\}}(H)$, adjoin $a=(x,0,1)$, where the final two entries
are the $e,f$ coordinates, and put $V=W\cup\{a\}$. Then:
\begin{enumerate}
 \item $V$ is an ample forbidden pc-minor with unique corner $a$,
 and minimum degree at least three;
 \item $V^{\{e\}}=H$ with $z$ renamed $f$, and neither $e$-section
 peels to this negative overlap;
 \item if the signed-clause columns of $H$ are distinct up to sign,
 the same is true of $V$. In this case $V$ is unchanged by either
 canonical graph or signed-column reduction.
\end{enumerate}
The construction uses only the local ample-addition conditions on $x$;
no output peeling or proper-minor tests are required.

Conversely, let $V$ be an irredundant ample forbidden pc-minor with a
unique corner $a$ of degree at least three. Suppose that in
$W=V\setminus\{a\}$ the signed-clause columns of two coordinates
$e,f$ agree up to sign. Normalize their signs so that the columns agree,
and suppose that $a_e\ne a_f$. Then $V$ has precisely the preceding
row-completion presentation for this marked pair, with cornerless forbidden
input $H=W^{\{e\}}=V^{\{e\}}$ (and $z$ renamed $f$).
The input, added old word, and acquired support are recovered uniquely
for the marked pair. Different eligible pairs are not asserted to give
the same unmarked presentation.
\end{proposition}
\begin{proof}
Since $x\notin C$, the new word $a$ has no neighbour in either block
direction. Its old incident cube is the completed $P$-cube in the
nonconstant row $C$. The support $P$ is not shattered by $W$, by
\eqref{eq:damp-coordinate-block-shadow}. The local extension criterion
therefore makes $V$ ample with acquired support $P$. Its $e$-reduction
is unchanged, since $a$ has no $e$-mate. This reduction is the negative
family $H$, so reduction closure makes $V$ nonpeelable. The one-concept
extension theorem, Proposition~\ref{prop:damp-one-concept-atlas}, now
makes $V$ forbidden.

The block theorem makes $W$ cornerless. Every old vertex has a neighbour
in at least one of the block directions. If an old vertex became a
corner in $V$, its corner cube would have to contain $a$, since otherwise
it was already a corner in $W$. That cube would also contain the neighbour
of $a$ in this incident block direction, which is absent. This is
impossible. Thus $a$ is the only corner. The old minimum degree is at
least three by the block theorem, and $a$ has degree $|P|\ge3$.

Before the addition, neither $e$-section has a corner outside its overlap:
such an exclusive corner would also be a corner of $W$. The one section
is unchanged. In the zero section, every old exclusive vertex has its
$f$-mate. If it became a corner after adjoining $a$, its corner cube
would contain $a$ and its absent $f$-mate, by the same argument.
Therefore a relative deletion in that section must first delete $a$,
after which it cannot continue. This proves item 2.

For item 3, let $\tau=x|_P$. The old signed minimal nonface on $P$
is removed. Every other old signed minimal nonface persists: its support
cannot contain $P$, by minimality. Among the new minimal nonfaces are
$P\cup\{e\}$ and $P\cup\{f\}$, with missing traces $(\tau,1)$
and $(\tau,0)$ respectively. Indeed all proper supports were shattered
already, except $P$, which the addition now shatters; and only $a$
realizes the formerly missing $P$-trace. These two rows distinguish the
columns $e,f$.

No two old-coordinate columns become equal up to sign, since replacing
the old $P$-row by either new row preserves all of their entries.
To compare an old coordinate with $e$, use the new $P\cup\{f\}$ row:
on these columns it reproduces the removed $P$-row, including its zero
entry at $e$. The other old rows persist, so any former distinction is
preserved. Use $P\cup\{e\}$ for the analogous comparison with $f$.
Before the addition the only equal columns up to sign were $e,f$, by
the block clause rule and the hypothesis on $H$. All columns of $V$
are therefore distinct up to sign. Together with its minimum degree,
this proves the last assertion.

For recovery, deleting the corner leaves $W$ ample. It is nonpeelable,
since a peeling of $W$ prefixed by $a$ would peel $V$. The equal-column
argument in Theorem~\ref{thm:damp-canonical-signed-column-reduction}
uses only ampleness to recover the block presentation: all clauses either
avoid $e,f$ or contain both with matching signs, every nonconstant row
is their common contraction $C$, and the constant rows are $B_0,B_1$
with union $C$. Thus $W=W_{\{e,f\}}(H)$.
Because the omitted word $a$ has a nonconstant block trace, its old word
$x$ lies outside $C$. It has neither block mate. Hence $V^{\{e\}}$
is exactly $W^{\{e\}}=H$. The block peeling equivalence makes $H$
nonpeelable, and negative-overlap descent applied to the forbidden $V$
makes $H$ an irredundant ample forbidden minor.

We must verify cornerlessness of this recovered input. If a corner of
$H$ has its $z$-mate, its corresponding constant-row lift is a corner
of $W$, by the corner-lifting proof of the block theorem. It is not
adjacent to $a$: its old word lies in $C$, and its block trace differs
from that of $a$. If the corner of $H$ has no mate, both nonconstant
lifts are corners of $W$. Choose the lift in the nonconstant row
opposite to that of $a$; it also is not adjacent to $a$.
In either case that corner survives in $V$, since adjoining a
nonadjacent word cannot change its incident directions or remove its
corner cube. This contradicts uniqueness of $a$. Therefore $H$ is
cornerless. Finally, contracting the two block coordinates in $V$
gives the ample family $C\cup\{x\}$. Its acquired support is exactly
the old incident support of $a$, of size at least three. The sections
and the deleted word determine all marked data uniquely, completing
recovery.
\end{proof}

For a fixed instance take the $267$-concept seed and the coordinate-zero
block expansion above. In its eleven old coordinates, choose $x=119$.
The contraction admits this ample addition with acquired support
$P=265=2^0+2^3+2^8$. With $f=11$ and $e=12$, the added word is
$a=2167$. The resulting $684$-concept forbidden class has minimum degree
three, unique corner $2167$, and thirteen pairwise distinct signed-clause
columns up to sign. Its $e$-overlap still has $267$ concepts. The zero
section must delete $2167$ first and then stalls with no allowable
relative deletion; the one section has no allowable first deletion.
The fixed checker \path{research/block_row_completion.py/json} verifies
the local addition, all signed clauses, both positive section orders,
and these forced relative failures. Output minor minimality follows
from the uniform proposition and the certified block input.

This example settles the existence question even on bases reduced under
both canonical operations. It does not refute generation by all existing
operations: it is a block expansion followed by one explicitly certified
nonresolving addition. The new row-completion rule gives a structural
sufficient condition for that addition and needs no minorwise tests.
Intrinsic generation of all such reduced bases remains open.

\begin{proposition}[Two complementary components force antipodal cube rows]
\label{prop:damp-two-component-cube-code}
Let $A\subseteq C$ belong to $\Damp$, and suppose that $C\setminus A$
has exactly two connected components $K_0,K_1$. Assume that
\[
 H=(A\times\{0,1\})\cup(K_0\times\{0\})\cup(K_1\times\{1\})
\]
is an ample forbidden pc-minor. For a new nonempty coordinate set $E$
and words $q_0,q_1\in Q_E$, set
\[
 Z=(A\times Q_E)\cup
 \bigcup_{i=0}^1\bigl(K_i\times(Q_E\setminus\{q_i\})\bigr).
\]
Then $Z$ is forbidden if and only if $q_0,q_1$ are antipodal.
In that case it is a coordinate-block replacement of $H$, up to
reorienting and renaming coordinates. If $q_0=q_1$, then $Z\in\Damp$.
If the two words are distinct and not antipodal, $Z$ has a nonpeelable
proper coordinate-face restriction and is not forbidden.
\end{proposition}
\begin{proof}
If the words are antipodal, reorient them to $1_E,0_E$, respectively.
The displayed family is then exactly $W_E(H)$, so the block theorem
proves forbiddenness. If both words equal $q$, reorient $q$ to $0_E$.
The family becomes $W_E(P)$ for the peripheral family
$P=(A\times\{0\})\cup(C\times\{1\})$. Since $A,C\in\Damp$,
$P$ and hence its block replacement are positive.

Finally suppose the words are distinct but not antipodal. Choose
$e,f\in E$ such that $q_0(e)=q_1(e)$ and $q_0(f)\ne q_1(f)$.
Restrict to the common $e$-value, and then reduce along all surviving
coordinates of $E$ except $f$. Over $A$ the $f$-fibre stays full.
Over $K_i$ its only surviving value is $1-q_i(f)$, since the other
value has a missing word in its reduced fibre. The result is $H$ up
to reorientation of $f$. Thus the proper face restriction has a
negative reduction and cannot belong to $\Damp$. This rules out
forbiddenness. The restriction uses an active coordinate, because
$A\ne\varnothing$ and its $E$-fibre is full.
\end{proof}

For the $267$-concept seed, its coordinate-zero contraction and overlap
have precisely these two complementary components, of orders $90$ and
$59$. Hence all sixteen choices of two omitted square words are settled:
four equal pairs are positive, eight non-antipodal distinct pairs have a
negative proper face, and four antipodal pairs are the forbidden block
replacements. This is uniform for every block dimension, not evidence
for a general statement about arbitrary cornerless forbidden classes.
A separate bounded test in \path{research/explore_block_sign_break.py/json}
changes one block-coordinate sign in each of the $93$ mixed clauses of
the fixed $683$-word presentation. Exactly four resulting families are
ample, of orders $649,639,645,634$, with respectively $2,2,1,2$ corners.
None meets the proposed cornerless-input condition. No claim about their
forbiddenness or about all sign changes is made.

\begin{proposition}[One relative section suffices for a mixed puncture]
\label{prop:damp-one-wing-puncture-retention}
Let $H$ be ample, containing the marked edge
$J=\{a_0,a_1\}=\{0,\mathbf1_{\{z\}}\}$. Write $A_0,A_1$ for its
$z$-sections, $D=A_0\cap A_1$, and $C=A_0\cup A_1$. Suppose at least
one of $A_0,A_1$ peels to $D$. Replace $z$ by $E=\{e,f\}$, and let
$T$ be the root square in $W_E(H)$. For either mixed word
$q\in\{01,10\}$,
\[
 W_E(H)\searrow T\setminus\{q\}
 \quad\Longleftrightarrow\quad
 H\searrow\{a_0\}\ \hbox{and}\ H\searrow\{a_1\}.
\]
The forward implication does not require the relative-section hypothesis.
Consequently, a counterexample to the converse without that hypothesis
must have both sections failing to peel to $D$. Each exclusive wing must
then contain a component with a cycle and at least nine vertices.
\end{proposition}
\begin{proof}
Reduce a relative peeling in each of the two block coordinates.
Consecutive reductions are ample and either equal or differ by one word;
discarding repetitions gives relative corner-deletion orders. The two
reductions of $W_E(H)$ are $H$, while those of the mixed-punctured root
square are its two different marked endpoints. This proves necessity.

For sufficiency, reversing the old marked bit and both block bits, and
interchanging $e,f$ if needed, reduces to $A_0\searrow D$ and
$q=(0,1)$ in the ordered coordinates $(e,f)$. Initially the four rows
$00,01,10,11$ of the block are $A_0,C,C,A_1$.
Use a relative order from $A_0$ to $D$ to delete the words of
$A_0\setminus D$ \emph{only in row $01$}. There are no old-coordinate
edges between $A_0\setminus D$ and $A_1\setminus D$: such an edge
would produce two vertices of $H$ at distance two with neither intermediate
vertex present. Thus the old neighbours of each deleted word are those
in the current $A_0$-relative order. Its only block-coordinate neighbour
is in row $00$, which still contains its whole old corner cube. Every
deletion is therefore a corner deletion.

The rows are now $A_0,A_1,C,A_1$. The $e=0$ section is a copy of $H$;
the $e=1$ section is the peripheral family
\[
 L=(C\times\{0\})\cup(A_1\times\{1\}),
\]
on the old coordinates and $f$, and it contains that copy of $H$.
Use a peeling of $H$ ending at $a_0$ to delete every nonroot vertex
of the $e=0$ section. All its corner cubes have full mates in $L$,
which is retained throughout, so these are valid deletions. This step
deletes the omitted root word $01$ and retains the root word $00$.

The remaining family is $L$ in the $e=1$ section, together with $00$.
Restriction of the endpoint-$a_1$ order gives $A_1\searrow\{0\}$;
contraction of either endpoint order gives $C\searrow\{0\}$.
Delete the nonroot words of the smaller row $A_1$ along its rooted
order, using their full mates in $C$. Then delete the nonroot words
of $C$ along its rooted order. The retained vertex $00$ has no neighbour
among these nonroot words, so it does not change their corner tests.
Exactly $00,10,11$ remain, proving the required relative peeling.

For the last assertion, the preceding implication forces failure of both
relative sections in any counterexample. For either section, every union
of $D$ with one exclusive component is ample, by the component-expansion
theorem. If all those components were trees or had at most eight vertices,
the relative forest/eight-facet argument using
Lemma~\ref{lem:damp-relative-yang-cm-flag-h}, as in the proof of
Corollary~\ref{cor:damp-small-component-wing-inverse}, would peel them
successively to $D$. Hence each wing has a cyclic component of order at
least nine.
\end{proof}

This strictly extends the two-relative-wing sufficient argument for the
mixed punctures of $Y_2$ in
Theorem~\ref{thm:damp-cube-boundary-retention}. The archived positive
$294$-word class has certified endpoint orders at the adjacent words
$1639,1655$. In their marked coordinate $4$, its sections have sizes
$136,158$, intersection $65$, and union $229$. The $71$-word exclusive
wing peels to the overlap, whereas none of the $93$ words of the other
wing is initially deletable with the overlap retained. The new schedule
retains a mixed-punctured square in its $752$-word block. The replay
\path{research/one_wing_puncture_retention.py/json} checks all four
stages, of lengths $71,293,157,228$, and the completed acyclic nonclashing
map. These input endpoint and one-wing orders are reused from the
coordinate-stripping audit; the mixed-puncture schedule is the new conclusion.

The following application of the earlier product-attachment theorems
identifies the exact relative target in a square assembly. Retaining the
whole root square is sufficient, but the exact target omits its hole.

\begin{proposition}[Exact relative criterion for an adjacent square hole]
\label{prop:damp-adjacent-square-relative-criterion}
Let $A\subseteq Q_F$ be ample and corner-peelable, containing
$J=\{0,\mathbf1_{\{z\}}\}$, and let $(B,0)$ be a rooted factor on
disjoint coordinates. All coordinates of both factors are active.
Put $P=W_E(A)$, where $E=\{e,f\}$ replaces $z$, and let $T$ be its
root square. For $q\in\{01,10\}$ put $K=T\setminus\{q\}$ and
\[
 X=P\cup(K\times B),
\]
embedding $P$ at the root of $B$. Thus this is precisely the rooted
square assembly used below: the missing root word already belongs to $P$.
Then $X$ is ample and
\[
 X\in\Damp\quad\Longleftrightarrow\quad P\searrow K.
\]
Writing $A_0,A_1,C$ for the $z$-sections and contraction, the assembly
is forbidden if and only if all three conditions hold:
\begin{enumerate}
 \item $P\not\searrow K$;
 \item $A_0,A_1,C$ all peel to the projected root $0$;
 \item for every $j\in F\setminus\{z\}$ and
 $\mu\in\{\rho_j^0,\pi_j\}$, the block $W_E(\mu A)$ peels to $K$.
\end{enumerate}
If $A$ has unique corner $0$, the first condition is automatic.
The criterion is independent of which rooted factor $B$ is used.
It is an exact relative-peeling criterion, not an intrinsic selection of
all admissible $A$.
\end{proposition}
\begin{proof}
The identity $X=P_B(P,K)$ makes ampleness and the positivity equivalence
an application of Theorem~\ref{thm:damp-relative-attachment-test}; the
three-vertex path $K$ is ample and corner-peelable. By
Theorem~\ref{thm:damp-relative-attachment-minimality}, forbiddenness is
equivalent to failure of $P\searrow K$ and restoration after every
elementary operation on $P$ whose image of $K$ is nonempty. We spell
out those operations to obtain the displayed conditions.

For $j\in F\setminus\{z\}$, a one restriction removes $K$, while
a zero restriction or contraction leaves $K$ unchanged and sends $P$
to $W_E(\mu A)$. These are exactly item~3.
A contraction in $e$ or $f$ sends the pair $(P,K)$ to
$(C\times Q_1,\{0\}\times Q_1)$. This pair peels relatively exactly
when $C$ peels to $0$: one direction follows by restriction to a row,
and the other by deleting nonroot fibres along a rooted order of $C$.

A square-coordinate restriction fixing its bit to $b$ leaves a
peripheral family whose two rows are $A_b\subseteq C$. Its target is
either the root edge, or the root in the smaller row $A_b$. Indeed,
if the fixed bit differs from the corresponding bit of $q$, no word
is removed from the root edge; otherwise the surviving root has the
other bit equal to $b$, since the bits of $q$ differ. For either target,
relative peeling implies a rooted order on $A_b$ by row restriction
and a rooted order on $C$ by contraction. Conversely, given both
orders, delete the nonroot words in the smaller row first, using their
full mates in $C$, and then the nonroot words of $C$. This retains
the root edge. For a singleton target, delete its other endpoint last.
Each deletion is a corner: in the smaller row its corner cube has all
mates, and outside the root the larger row has no remaining vertical
neighbour. Both values of $b$ occur, proving that all square-coordinate
operations are equivalent to item~2.

Finally, when $A$ has unique corner $0$, its block $P$ has unique
corner $00$ by the block corner calculation below. That vertex lies
in $K$, so no relative peeling of $P$ to $K$ can start.
\end{proof}

\begin{corollary}[A square generator from rooted endpoints and small minor wings]
\label{cor:damp-rooted-endpoint-square-generator}
Let $(A,a_0)$ be a rooted factor with unique corner $a_0=0$, and mark
an incident edge $a_0a_1$ of direction $z$. Suppose either $a_1$ is
an attainable endpoint of $A$, or $(A,a_1)$ is also a rooted factor.
For each $j\ne z$ and $\mu\in\{\rho_j^0,\pi_j\}$ put $M=\mu A$.
Assume that in at least one of the two $z$-sections of $M$, every
component outside their intersection is a tree or has at most eight
vertices. Then for every rooted factor $(B,0)$ on disjoint coordinates,
and either adjacent square hole $q$, the assembly in
Proposition~\ref{prop:damp-adjacent-square-relative-criterion} is forbidden.
No relative orders or output-minor tests are required in this rule;
the rooted-factor and endpoint conditions remain input assumptions.
\end{corollary}
\begin{proof}
Rooted minimality at $a_0$ makes the $z=0$ section and the contraction
peel to the projected root. The assumption at $a_1$ does the same for
the $z=1$ section, by rooted minor closure or rooted minimality there.
For every private operation $\mu$, both images of the marked endpoints
are attainable in $M$, by the same two properties. The componentwise
relative forest/eight-facet argument gives a section of $M$ peeling to
its overlap. Proposition~\ref{prop:damp-one-wing-puncture-retention}
therefore makes $W_E(M)$ retain the mixed-punctured root square.
These are all conditions of
Proposition~\ref{prop:damp-adjacent-square-relative-criterion}; its negative
condition is automatic from the unique corner of $A$.
\end{proof}

This graph-only rule is a sufficient subclass, not a replacement for all
known input certificates. For the fixed $A289$ input with marked coordinate
$1$, only one of its $22$ private elementary images meets the component
condition. The audit is included in
\path{research/one_wing_puncture_retention.py/json}. Failure of this graph
condition does not imply failure of relative peeling.

The relative test in item~3 of the proposition is genuinely weaker than retention of the
whole square on unrestricted input pieces. The $1165$-word family
$Y_2$ of Theorem~\ref{thm:damp-cube-boundary-retention} is the block
replacement of the $581$-word edge piece of
Proposition~\ref{prop:damp-edge-boundary-transfer}. It retains every
punctured root square but not the full root square. Attaching any rooted
factor along one of these punctured squares therefore gives a positive
family by the attachment theorem. With a $289$-word factor this has
$2029$ words; \path{research/square_attachment_criterion.py/json}
replays its complete positive order from the archived relative prefix.
This shows strictness of the local relative test. It does not show that
such a piece occurs as a proper minor of an otherwise admissible
unique-corner first factor; necessity of the stronger edge conditions
under all those additional hypotheses remains open.

\begin{theorem}[A rooted square assembly with automatic minor minimality]
\label{thm:damp-rooted-square-assembly}
Let $A\subseteq Q_F$ be ample and corner-peelable, with unique corner
$0$, and let $z\in F$ be incident with $0$. Write
$J=\{0,\mathbf1_{\{z\}}\}$ and denote the two $z$-sections and
contraction by $A_0,A_1,C$. Assume:
\begin{enumerate}
 \item every elementary contraction or zero restriction in a coordinate
 of $F\setminus\{z\}$ sends $A$ to a family that peels to the image of $J$;
 \item $A_0,A_1,C$, with $z$ suppressed, all peel to $0$.
\end{enumerate}
Let $(B,0)$ be any rooted factor from
Theorem~\ref{thm:damp-cut-vertex-obstructions}, on coordinates disjoint
from $F$. Replace $z$ in $A$ by a
square coordinate set $E=\{e,f\}$, giving $P=W_E(A)$.
For any $q\in Q_E\setminus\{00\}$ put
\[
 Q=(B\times(Q_E\setminus\{q\}))\cup(\{0\}\times Q_E),
 \qquad X=P\cup Q.
\]
Here each piece is embedded with the other's private coordinates zero,
so their intersection is the common root square $T=\{0\}\times Q_E$.
Then $X$ is an ample forbidden pc-minor. Its corners are precisely
$(b,t)$ with $b$ a nonroot corner of $B$ and $t$ a square neighbour of
$q$. In particular it is cornerless exactly when $0$ is the unique
corner of $B$. Its order is
\[
 |X|=3|C|+|A_0\cap A_1|+3|B|-3.
\]
The marked pieces and inputs are recovered by their coordinate faces,
block reduction, and square rows. No peeling tests on output minors
are required. The two displayed input conditions are assumptions;
rooted minimality alone is not asserted to imply edge retention.
\end{theorem}
\begin{proof}
We first check positivity and corners of the pieces and nonpeelability
of their union, then treat the
three types of elementary minor: private coordinates of $A$, private
coordinates of $B$, and the two square coordinates.

The block theorem makes $P$ ample and positive. The family $Q$ is
ample: it is the corner extension at $(0,q)$ of the ample product
$B\times(Q_E\setminus\{q\})$. Its incident support is $E$, which
that product does not shatter. Deleting this corner leaves a product
of two positive classes, so $Q$ is positive. The union over their
common coordinate square is ample. Indeed no shattered support can
use private coordinates from both pieces (the trace assigning one
on a coordinate of each piece is missing), and their shadows intersect
exactly in the square shadow. The vertex and shadow cardinalities
therefore both add with the same subtraction $|T|=4$.

The only corner of $P$ is $(0,00)$. A corner with a constant block
trace projects to a corner of $A$ at that trace. At a nonconstant
trace over $x\in A_0\cap A_1$, its corner cube projects to corner
cubes at both endpoints of the old $z$-edge, since both new-coordinate
neighbours are present. If instead $x$ belongs to only one section,
the edge to that section's constant trace forces its old neighbour
cube into that section, making $x$ a corner of $A$ there. This last
case is impossible because the unique corner $0$ of $A$ has its
$z$-mate; the full-fibre case would require two corners of $A$.
Thus only $(0,00)$ can occur, and its corner cube is present.
The corners of $Q$ are $(0,q)$ and the pairs $(b,t)$
with $b$ a nonroot corner of $B$ and $t$ an endpoint of the path
$Q_E\setminus\{q\}$. At a nonroot old word the incident cube test
is exactly the product corner test for $B$ and that path; the extra
root word at $q$ is not adjacent to this word. At a root word in another
row, its old neighbour and its two square neighbours would force a
missing copy over $q$ in any corner cube. A corner of the union must
be a corner in each coordinate-face piece containing it. The two
possible boundary corners are distinct because $q\ne00$. The nonroot
corners of $Q$ have no neighbours in $P\setminus T$, so they remain
corners and give exactly the stated list.

Put $K=T\setminus\{q\}$. The assembly is $P_B(P,K)$, so
Theorem~\ref{thm:damp-relative-attachment-test} says it peels exactly
when $P$ peels to $K$. Since $P$ has unique corner $00\in K$, no
such relative peeling can start. This proves nonpeelability even when
$B$ has nonroot corners. For adjacent holes the preceding proposition
also gives the exact minor criterion; the following schedules verify
the sufficient hypotheses uniformly for all $q\ne00$.

Consider an elementary operation in a private coordinate of $A$.
A zero restriction or contraction sends $A$ to a family peeling to $J$,
by assumption. Relative block transport makes the resulting $P$ peel
to $T$. These deletions remain corners in the union: outside $T$,
vertices of the two pieces have no adjacent private vertices in the
other piece. The whole minor therefore peels to the unchanged positive
$Q$. A one restriction removes $Q$ entirely, leaving a proper minor
of the positive $P$.

For a private coordinate of $B$, zero restriction or contraction makes
$B$ peel to $0$ by rooted minimality. The corresponding $Q$ then peels
to $T$. To see this directly, follow the rooted order of $B$; for each
deleted nonroot word, delete its three square-row copies in a corner
order of $Q_E\setminus\{q\}$. Their corner cubes are products of
the current old corner cube and the current row corner cube. The root
row stays intact. This reduces the whole minor to positive $P$.
A one restriction leaves only a product of a positive minor of $B$
with the punctured square, which is positive.

A restriction in $e$ or $f$ gives a peripheral section of $P$, with
nested sections $A_b\subseteq C$. Both peel to the root by assumption.
First delete the nonroot words of the smaller section in its rooted
order, using mates in the larger section; then delete the nonroot words
of the larger section. This retains exactly the shared root edge.
In the union, the same deletions leave the corresponding proper minor
of positive $Q$, proving positivity. A contraction in $e$ or $f$ gives
\[
 (C\vee B)\times Q_1.
\]
Since $C$ peels to the gluing root and $B$ is positive, this product
is positive. These cases cover every elementary minor. Pc-minor closure
and the negative attachment test prove forbiddenness. The order formula follows
from the two piece sizes and their four-vertex intersection.

Finally the marked coordinate faces recover $P,Q$. Reduction of $P$
in either square coordinate recovers $A$ with its marked direction
renamed. Any nonmissing row of $Q$ recovers $B$, and the unique missing
row over a nonroot word recovers $q$. This is marked recovery; uniqueness
of an unmarked decomposition is not claimed.
\end{proof}

The fixed $289$-concept factor used in the cornerless double satisfies
these stronger marked-edge conditions for $z=1$. The replay
\path{research/rooted_square_assembly.py/json} checks its $22$
edge-relative elementary orders and the three projected endpoint orders.
Here $|C|=224$ and $|A_0\cap A_1|=65$. Thus every admissible second
rooted factor $B$ gives a forbidden class of order
$3|B|+734$ in dimension $\dim(B)+13$, cornerless exactly when $B$ has
unique corner at its root. The certified $D289$ factor, whose root is
not a corner, instead gives an output with exactly two corners. The
replay \path{research/square_attachment_criterion.py/json} verifies
its forced-prefix rooted failure, all $24$ rooted minor orders, the
attachment identity, and the predicted corner set; output minimality
is supplied by the theorem. Taking the same $289$-concept
factor for $B$ and a square hole adjacent to $00$ gives the $1601$-concept
example below, up to reversing both square bits. This replaces its
output-specific minor checks by a uniform boundary argument. Selection
of all admissible rooted and marked-edge inputs remains open.

\begin{proposition}[The two reductions of a rooted square assembly]
\label{prop:damp-rooted-square-recovery}
Use the square assembly $X=P\cup Q$ above, but assume only that $A,B$
are nontrivial corner-peelable ample families, that $J\subseteq A$, and
that $(B,0)$ is a rooted factor. Write $a_0=0$,
$a_1=\mathbf1_{\{z\}}$, and $q=(q_e,q_f)\ne00$.
If $A\vee_{a_b}B$ denotes gluing the root of $B$ to $a_b$ in $A$, then,
after renaming the surviving square coordinate to $z$,
\[
 X^{\{e\}}=A\vee_{a_{1-q_f}}B,
 \qquad X^{\{f\}}=A\vee_{a_{1-q_e}}B.
\]
Each displayed reduction is corner-peelable exactly when $A$ peels to
its indicated attachment vertex. If $X$ is forbidden and an indicated
vertex is not attainable, then $A$ rooted at that vertex is a rooted
factor of Theorem~\ref{thm:damp-cut-vertex-obstructions}.

For the opposite hole $q=11$, one has the stronger identity
\[
 X=W_E(A\vee_{a_0}B).
\]
Consequently this opposite-hole assembly is forbidden if and only if
$(A,a_0)$ is a rooted factor. No marked-edge retention or projected
endpoint conditions are needed in this case. For an adjacent hole the
two reductions instead attach $B$ to the two different endpoints of $J$.
\end{proposition}
\begin{proof}
Reduction of $P$ in either square coordinate recovers $A$. In the
$e$-reduction of $Q$, every nonroot word of $B$ survives precisely in
the row $f=1-q_f$: in the other row its pair includes the missing word
$q$. The root survives in both rows, giving $B$ at the indicated
endpoint together with the edge $J$, already contained in $A$.
No additional paired word can be formed from one private word of each
piece, since their private coordinate sets are disjoint. This proves
the first identity; interchanging $e,f$ proves the second.
Proposition~\ref{prop:damp-vertex-gluing} gives the positivity criterion,
since $B$ does not peel to its root. If $X$ is forbidden, every negative
displayed reduction is forbidden by
Theorem~\ref{thm:damp-inherited-overlap-descent}. Its two pieces meet at
a cut vertex, so Theorem~\ref{thm:damp-cut-vertex-obstructions} gives
the necessary rooted-factor condition on $A$.

When $q=11$, put $H=A\vee_{a_0}B$. Its $z$-sections are $A_0\vee B$
and $A_1$, with contraction $C\vee B$. Thus the four rows of $W_E(H)$
are exactly those of $P\cup Q$: the nonroot words of $B$ occur in
every row except $11$. The block theorem
(Theorem~\ref{thm:damp-coordinate-block-replacement}) transfers
forbiddenness in both directions between $X$ and $H$, and the
cut-vertex theorem characterizes forbiddenness of $H$ as claimed.
\end{proof}

This separates reuse of the block construction from the adjacent-hole
case of Theorem~\ref{thm:damp-rooted-square-assembly}. In particular,
if $A$ has unique corner $0$, then any forbidden assembly with
$q\ne00$ necessarily has $(A,0)$ as a rooted factor: at least one of
the two reductions attaches at $0$, and no peeling of a nontrivial
family can retain its unique corner. It is not proved that every such
rooted factor satisfies the adjacent-hole theorem's stronger assumptions.
The bounded audit \path{research/explore_rooted_square_inputs.py/json}
checks all incident root edges of the certified factors $A288$ and
$B289$. Their eligible directions are respectively $1,2,7$ and
$1,2,8$ in that certificate's coordinates. All six pass the $22$
edge-relative and three projected endpoint tests, with replayed orders.
These are finite input certificates, not a universal implication from
rooted minimality.

\begin{proposition}[Square attachments with only positive coordinate reductions]
\label{prop:damp-square-positive-reductions}
Let $A$ be ample and corner-peelable with marked edge
$J=\{a_0,a_1\}$, normalized to $a_0=0$ and
$a_1=\mathbf1_{\{z\}}$. Let $(B,0)$ be a rooted factor on disjoint coordinates,
and take an adjacent hole $q\in\{01,10\}$. Put
$P=W_{\{e,f\}}(A)$, let $T$ be its root square, and set
$K=T\setminus\{q\}$ and $X=P_B(P,K)$.
Every reduction in a private coordinate of either factor is positive.
Moreover,
\[
 X^{\{d\}}\in\Damp\text{ for every active }d
 \quad\Longleftrightarrow\quad
 A\searrow\{a_0\}\text{ and }A\searrow\{a_1\}.
\]
Here empty reductions may be omitted. In particular, the right-hand
condition does not by itself assert positivity of $X$.
If $A$ has a private active coordinate, the two square coordinates are exactly the
universal vertices of $\operatorname{Cr}(X)$, and deleting both leaves
the disjoint crossing graphs of $\pi_zA$ and $B$. Thus they form a
two-vertex cut of the crossing graph, which has no cut vertex.
\end{proposition}
\begin{proof}
If $j$ is a private coordinate of $A$, all words in the attached part
have $j=0$. A pair contributing to the $j$-reduction must therefore
have its $j=1$ word in $P$, with all private $B$ coordinates zero.
At those coordinates every attachment word already belongs to $P$.
Consequently $X^{\{j\}}=P^{\{j\}}$, which is positive by the block
theorem and reduction closure
(Proposition~\ref{prop:damp-strong-projection-closure}).

If $j$ is private to $B$, the $j=1$ word of any contributing pair lies
over $K$. Over $K$ the complete family is exactly $K\times B$, including
the root copy already in $P$. Hence
\[
 X^{\{j\}}=K\times B^{\{j\}},
\]
which is positive by reduction and product closure. The two remaining
reductions attach $B$ at the two different endpoints of $J$, by
Proposition~\ref{prop:damp-rooted-square-recovery}. Since $B$ cannot peel
to its root, vertex gluing makes each such reduction positive exactly
when $A$ can peel to its corresponding endpoint. This proves the equivalence.
All further reductions are positive or empty by reduction closure.

For the crossing graph, the block shadow formula makes each square
coordinate cross every private coordinate of $A$, and their root square
makes them cross one another. The product $K\times B$ makes each cross
every private coordinate of $B$. No support using private coordinates from
both factors is shattered, since their simultaneous nonzero trace is absent.
On the private coordinates the shadows are those of $\pi_zA$ and $B$.
This proves the stated graph description. No private coordinate is universal,
since it misses every coordinate of the other nonempty private block.
Both private blocks are nonempty;
deleting the square coordinates disconnects them, whereas after any single
vertex deletion at least one universal square coordinate remains.
\end{proof}

\begin{proposition}[Recovering mixed-puncture square attachments without markings]
\label{prop:damp-square-unmarked-recovery}
Let $X$ be an ample forbidden pc-minor with all coordinates active.
All presentations
\[
 X=P_B(W_{\{e,f\}}(A),T\setminus\{q\})
\]
with $A$ positive, a nonempty private coordinate set for $A$, a rooted
factor $B$, and a mixed hole $q$, can be recovered by the following
finite procedure. It uses only crossing graphs and coordinate fibres;
no marked root, coordinate partition, or peeling order is supplied.
\begin{enumerate}
 \item Require that $\operatorname{Cr}(X)$ have exactly two universal
 vertices, and let $E$ be this pair. For each connected component $L$
 of $\operatorname{Cr}(X)-E$, put $I=F\setminus(E\cup L)$, where $F$
 is the coordinate set of $X$, and require $I\ne\varnothing$.
 \item Project $X$ onto $L\cup E$. For each projected $L$-word $b$,
 let $S_b\subseteq Q_E$ be its fibre. Require a unique word $r_B$ with
 $S_{r_B}=Q_E$, and require all other fibres to be the same
 three-element set $Q_E\setminus\{q\}$. Require at least one such
 other fibre, and let $B$ be the set of fibre indices.
 \item Require that every word of $X$ whose $L$-word differs from
 $r_B$ have the same $I$-word $r_A$. In the face with $L$-word $r_B$,
 write $R_t\subseteq Q_I$ for the row indexed by $t\in Q_E$.
 Require $r_A\in R_t$ for all four rows.
 \item Test each opposite pair $\{u,\bar u\}\subset Q_E$. Require
 $q\notin\{u,\bar u\}$ and
 \[
 R_t=R_u\cup R_{\bar u}\qquad(t\notin\{u,\bar u\}).
 \]
 After independently reversing the square coordinates to make
 $u=00$, recover
 \[
 A=(R_u\times\{0\})\cup(R_{\bar u}\times\{1\}),
 \qquad J=\{(r_A,0),(r_A,1)\}.
 \]
\end{enumerate}
Every passing choice gives such a presentation. The square coordinate
pair is unique; different components $L$ may give different presentations,
and uniqueness of the factors is not asserted. Reversing the marked
edge or interchanging the square coordinates merely changes its notation.
\end{proposition}
\begin{proof}
First consider a presentation of the stated kind. The preceding crossing
calculation shows that $E$ is exactly the set of universal vertices and
that its deletion leaves the disjoint crossing graphs of $\pi_zA$ and
$B$. A rooted factor has no cut vertex by
Theorem~\ref{thm:damp-cut-vertex-obstructions}; hence its crossing graph
is connected by Proposition~\ref{prop:damp-crossing-separator-recovery}.
There is no low-dimensional exception here: an ample family on at most
two active coordinates is a tree or a square, and every vertex is an
attainable sink, so it cannot be a rooted factor.
Thus the private coordinate set of $B$ is one of the components $L$
examined in step~1. Over its root the four square rows are present.
Over every other word of $B$, precisely the punctured square is present,
at the fixed first-factor root. This gives steps~2 and~3.
The definition of the coordinate-block replacement gives step~4,
including the condition that the omitted word is mixed. Consequently
no presentation of the stated kind is missed.

Conversely, take a passing choice. Steps~2 and~3 determine all words
outside the face with $L$-word $r_B$: they are exactly
\[
 \{r_A\}\times(B\setminus\{r_B\})\times
                (Q_E\setminus\{q\}),
\]
with coordinates ordered as $I,L,E$. Step~4 says that the remaining
face $P$ is $W_E(A)$. The four occurrences of $r_A$ are its root square
$T$. These identities recover the asserted product attachment exactly.

It remains to verify the input properties, which the fibre tests did
not assume. The face $P$ is proper because $L$ is nonempty and active;
it is therefore ample and positive by forbiddenness of $X$.
The recovered $A$ is a coordinate reduction of $P$, so it too is ample
and positive by Proposition~\ref{prop:damp-strong-projection-closure}.
For any nonmissing square row, fixing also the $I$-word to $r_A$
recovers $B$ as a proper face of $X$. Thus $B$ is ample and positive.
If $B$ peeled to $r_B$, its product over $T\setminus\{q\}$ could be
peeled to the root copy, leaving positive $P$. This would make $X$
positive, a contradiction. Hence its root is unattainable.

For rooted minimality, let $\mu$ be any elementary restriction or
contraction of $B$ retaining its root. The image $\mu X$ is a proper
pc-minor, hence positive, and is the same attachment with second
factor $\mu B$. If $\mu B$ still failed to peel to its root,
Theorem~\ref{thm:damp-relative-attachment-test} would imply
$P\searrow T\setminus\{q\}$. Applying that theorem to the original
positive root-negative $B$ would then make $X$ positive, again a
contradiction. Every such elementary image therefore attains its root;
transport of rooted orders under further pc-minors proves the full
rooted-minimality condition. This completes the converse.
\end{proof}

This is exhaustive recovery within the mixed-puncture attachment branch,
not a test that arbitrary ample inputs are forbidden. In particular,
positivity of the recovered first factor and rooted minimality of the
second factor were deduced from the hypothesis that $X$ is forbidden.
The independent fibre replay in
\path{research/square_reduction_terminal.py/json} recovers exactly one
component and opposite pair on each of the $1601$-word forbidden and
$2029$-word positive controls. The latter is a geometric presentation
only; its positivity illustrates why presentation recovery alone is
not an obstruction test. The same replay excludes the $267$-word seed
(all twelve crossing-graph vertices are universal) and the terminal
$627$- and $630$-word examples (neither has a universal vertex). Thus
none of these witnesses lies in the recovered branch.

\begin{corollary}[The endpoint converse and a terminal forbidden branch]
\label{cor:damp-puncture-converse-terminal-branch}
Fix any rooted factor $(B,0)$. The following statements are equivalent:
\begin{enumerate}
 \item For every ample $A$ with a marked edge $a_0a_1$, attainability of
 both endpoints implies that $W_2(A)$ peels to its root square with a
 mixed word removed.
 \item No forbidden square attachment $P_B(W_2(A),T\setminus\{q\})$,
 with $q$ mixed and $A$ corner-peelable, has all its coordinate reductions
 in $\Damp$.
\end{enumerate}
If item~1 fails, choose a counterexample minimal under pc-minors retaining
the whole marked edge. It generates, with this fixed $B$, a forbidden
attachment with all coordinate reductions positive. The same first factor
works for every other rooted factor on disjoint coordinates.
Every such counterexample $A$ uses at least seven coordinates, fails to
retain its full marked edge, and has a cyclic component of at least nine
vertices in each exclusive wing. If it is minimal under pc-minors retaining
the marked edge, it has no cut vertex and its marked edge is not separating.
\end{corollary}
\begin{proof}
Assume item~1. If an attachment has all coordinate reductions positive,
Proposition~\ref{prop:damp-square-positive-reductions} makes both
endpoints attainable in $A$.
Item~1 gives retention of the punctured square, and the product-attachment
test makes the whole attachment positive. It is therefore not forbidden.

Conversely suppose item~1 fails. Both endpoint orders, and hence ordinary
positivity, survive every pc-minor retaining the whole edge. Choose a
counterexample $A$ with as few active coordinates as possible among these
minors, discarding constant coordinates. For every private coordinate
$j\ne z$, its zero restriction and contraction retain both endpoints and
are proper. By minimality their blocks retain the mixed-punctured square.
The two $z$-sections and the contraction of $A$ peel to the projected root,
by restricting or contracting the two endpoint orders. Thus all conditions
of Proposition~\ref{prop:damp-adjacent-square-relative-criterion} hold:
the full block fails retention, every private image restores it, and the
three projected endpoint conditions hold. That proposition proves the
attachment forbidden for every rooted factor $B$. All its coordinate
reductions are positive by Proposition~\ref{prop:damp-square-positive-reductions}.
This contradicts item~2 and proves the equivalence and construction claim.

If $A$ retained the whole edge, relative block transport followed by deleting
the missing square word would retain the punctured square, a contradiction.
Proposition~\ref{prop:damp-one-wing-puncture-retention} supplies the two
cyclic-wing conclusions. Finally if $A$ used at most six coordinates,
each marked section would use at most five. Their nonempty ample overlap
could then be retained by
Theorem~\ref{thm:damp-five-coordinate-relative-peeling}. The one-section
proposition would again give puncture retention. Thus at least seven
coordinates are necessary.

Suppose now that this minimal $A$ has a cut vertex $c$. Let $U$ be
$c$ together with the component of $A-c$ containing the marked edge
(or its other endpoint if $c$ lies on the edge), and let $V$ be $c$
together with all remaining components. These are convex coordinate
restrictions on disjoint private coordinate sets, as in the proof of
Theorem~\ref{thm:damp-cut-vertex-obstructions}. Restrict an acyclic
USO with sink $a_0$ to $V$. A sink in $V\setminus\{c\}$ would be
a sink of the whole family, since it has no neighbours outside $V$.
Thus the restricted sink is $c$, and $V$ peels to $c$.
Its nonroot deletions remain valid in $A$: outside $c$ there are no
neighbours in $U\setminus\{c\}$. Hence $A\searrow U$.
Both marked endpoints remain attainable in $U$ by restriction, so
minimality gives $W_2(U)\searrow K$. Relative block transport
(Proposition~\ref{prop:damp-block-relative-transfer}) gives
$W_2(A)\searrow W_2(U)\searrow K$, a contradiction.

Finally suppose that the marked edge $J$ is separating. Adjoin $J$
to each component of $A-J$, obtaining the canonical convex pieces
$A_1,\ldots,A_k$ of
Proposition~\ref{prop:damp-crossing-separator-recovery}, with $k\ge2$.
Use an acyclic USO of $A$ with sink $a_0$; it orients $J$ from $a_1$
to $a_0$. Its restriction to each piece has sink $a_0$: an exclusive
sink would be a global sink, and $a_1$ has its outgoing edge to $a_0$.
If a piece fails to retain $J$, its restricted out-map at $a_1$ must
have a private outgoing direction. Otherwise both endpoints have
inward boundary at $J$, which implies relative peeling to $J$ by
Proposition~\ref{prop:damp-relative-uso-extension}.
Two failing pieces would therefore give private outgoing directions
from different pieces at $a_1$. Their mixed square is absent, violating
the outgoing-cube condition. At most one piece fails to retain $J$.
Choose it as $A_i$ if it exists, and choose any piece otherwise.
Peel every other piece to $J$. These deletions remain corners in the
whole union, because exclusive vertices have no neighbours in other
pieces. Thus $A\searrow A_i$. This is a proper coordinate restriction
retaining both attainable endpoints, and minimality and relative block
transport again imply $W_2(A)\searrow W_2(A_i)\searrow K$, a
contradiction. The marked edge is consequently not separating.
\end{proof}

This identifies the endpoint question with existence of a specific branch
of forbidden inputs at which negative-reduction descent stops. It does not
assert existence of that branch, nor that every forbidden class with positive
coordinate reductions has this presentation. In particular, an existing
terminal example does not refute the endpoint converse unless this marked
block-and-attachment decomposition is recovered. The unique-root-corner
first inputs of Theorem~\ref{thm:damp-rooted-square-assembly} cannot lie
in the terminal branch: their root endpoint is unattainable.
The crossing-graph condition also excludes the already known terminal
cut-vertex and separating-edge examples from this presentation: their
crossing graphs are not two-connected, by
Proposition~\ref{prop:damp-crossing-separator-recovery}.
The replay \path{research/square_reduction_terminal.py/json} checks all
$51$ reduction identities on the two archived controls: the $1601$-word
forbidden assembly has exactly one negative coordinate reduction, whereas
the positive $2029$-word attachment has all $26$ positive. The latter's
ordinary order is also replayed. Neither control witnesses the hypothetical
terminal forbidden branch.

\begin{theorem}[A cornerless reduced obstruction with a negative reduction]
\label{thm:damp-three-component-cornerless-obstruction}
Let $S$ be the certified $577$-concept cornerless double from
Remark~\ref{ex:damp-square-completed-578}, before the square addition.
Contract its coordinate $1$ and denote the contraction by $U$ and
reduction by $D$, numbering the remaining coordinates $0,\ldots,22$.
The graph $G(U\setminus D)$ has three components $K_0,K_1,K_2$,
ordered by increasing least binary word. For $q=(q_0,q_1,q_2)\in Q_2^3$
put
\[
 Z_q=(D\times Q_2)\cup\bigcup_{i=0}^2
       \bigl(K_i\times(Q_2\setminus\{q_i\})\bigr).
\]
Encode square words by $0,1,2,3$, with their least significant bit
in coordinate $23$. All these families are ample and have $1601$
concepts. Their exact classification is
\[
 Z_q\in\Obs(\Damp)
 \quad\Longleftrightarrow\quad
 q_0\mathbin\oplus q_1=3\quad\hbox{and}\quad q_2\ne q_1.
\]
There are $28$ positive assignments, $24$ assignments with a negative
proper face, and $12$ forbidden assignments.

For $q=(0,3,1)$, the family $Z_q$ is cornerless, has minimum degree five,
and has $25$ distinct signed-clause columns up to sign. Its reduction
in coordinate $23$ is $S$, up to a coordinate isometry. It is therefore
a forbidden class fixed by both canonical reductions, with a negative
coordinate reduction and neither relatively peelable section.
\end{theorem}
\begin{proof}[Computer-assisted proof with explicit certificates]
The input $S$ is the previously certified rooted double, not merely
an arbitrary cornerless family. The source comparison and coordinate
maps are checked in \path{research/verify_three_component_cornerless.py}.
One has $|U|=512$ and $|D|=65$. The component cube construction in
\href{research/component_label_partitions.pdf}{the component-partition note}
gives ampleness for all the displayed codes and order
$4|D|+3|U\setminus D|=1601$. Alternatively its shadow formula follows
by considering old trace fibres: a connected fibre of $U$ avoiding $D$
lies in one $K_i$ and therefore misses exactly its designated square word.
Every proper square support is shattered over the shadow of $U$, and
the full square support over the shadow of $D$.

For the representative $(0,3,1)$, the certificate
\path{research/three_component_code_minors.json} lists the family and
one full corner order for each of its $75$ elementary pc-minors.
The verifier reconstructs every minor and replays every deletion using
the incident-cube definition of a corner. It also computes actual cubes,
checks cornerlessness, and checks minimum degree and every signed clause.
Cornerlessness proves nonpeelability; the $75$ orders and pc-minor closure
independently prove forbiddenness. Theorem~\ref{thm:damp-rooted-square-assembly}
also proves it uniformly: its fixed input orders are replayed and the
assembled family is compared with this representative by a square isometry. The reduction in coordinate $23$ is compared
explicitly with $S$, after relocating its old coordinate $1$ to $24$
and, if necessary, reversing that bit. The output record is
\path{research/verify_three_component_cornerless.json}.
Since an exclusive section corner would also be a corner of the whole
family, neither relative peeling can begin.

Here are complete certificates for the finite classification. If
$q_0=q_1$ or $q_1=q_2$, the two established binary traces for labels
$001$ and $011$ lift to full-word traces, respectively. The verifier
constructs and replays all $28$ resulting positive orders.
If $q_0\oplus q_1=3$ and $q_2=q_0$, the four assignments are coordinate
block replacements of $S$. If $q_0\oplus q_1=3$ and $q_2$ differs
from both, the eight assignments are the images of $(0,3,1)$ under
square-coordinate permutations and sign changes; the verifier checks
the whole-family isometries. These are the twelve forbidden assignments.
In the remaining $24$ cases it records a proper square-coordinate face
and checks directly that the face is cornerless. They are not forbidden.
These three disjoint cases exhaust all $64$ assignments.
\end{proof}

Thus neither cornerlessness nor distinct signed columns forces all
coordinate reductions of a forbidden class to be positive. The example
is outside the one-relative-wing inverse branch even before any corner
completion. It is generated by the existing component cube construction
with a newly verified admissible assignment; the result does not supply
a structural selector for arbitrary inputs or prove exhaustive generation.
The distinction from the earlier $(0,2,3)$ square control is minimality:
that control has a cornerless proper face, whereas $(0,3,1)$ has explicit
positive orders for every elementary minor.

\begin{corollary}[Arbitrarily many maximal descent supports]
\label{cor:damp-coordinate-block-descent}
With the notation of Theorem~\ref{thm:damp-coordinate-block-replacement},
for every \(D\subseteq G\) and \(T\subseteq E\),
\begin{equation}
 W_E(H)^{D\cup T}=
 \begin{cases}
 W_{E\setminus T}(H^D),&T\subsetneq E,\\
 H^{D\cup\{z\}},&T=E.
 \end{cases}
 \label{eq:damp-coordinate-block-strong-projections}
\end{equation}
Consequently its nonpeelable reductions are determined exactly
by those of \(H\): in the first case the projection is nonpeelable
if and only if \(H^D\) is nonpeelable, and in the second case if and
only if \(H^{D\cup\{z\}}\) is nonpeelable.

In particular, suppose \(H\) is an ample forbidden pc-minor whose
nonempty-coordinate reductions all peel.  Then
\[
 \{U:W_E(H)^U\notin\Damp\}=\{T:T\subsetneq E\}.
\]
There are exactly \(m\) maximal descent supports, namely
\(E\setminus\{e\}\) for \(e\in E\).  Every maximal sequence of
one-coordinate nonpeelable descents has length \(m-1\), and every
endpoint is \(H\) with a coordinate renamed.  Thus neither uniqueness
of the maximal support nor a uniform bound on descent length extends
to all forbidden three-cores with two-connected crossing graphs.
This does not refute uniqueness of the terminal isomorphism type.
\end{corollary}

\begin{proof}
Reduction along \(D\subseteq G\) commutes with the construction.
The only point requiring justification is
\[
 (B_0\cup B_1)^D=B_0^D\cup B_1^D.
\]
If a full \(D\)-cube \(Q\) occurs in the union, apply
\eqref{eq:damp-residual-section-shadow-union} with residual family
\(K=Q\).  One of \(B_0\cap Q,B_1\cap Q\) shatters \(D\),
and hence contains the whole cube \(Q\).  This proves the nontrivial inclusion;
the reverse inclusion is immediate.  The argument also covers an empty
section directly.

It remains to take the reduction in \(T\).  If \(T\subsetneq E\),
a constant surviving word has fibre \(B_i\), and a nonconstant surviving
word has fibre \(C\).  For \(T=E\), a full fibre requires the old
word to lie in \(A\).  Applying this description after the projection
in \(D\) proves \eqref{eq:damp-coordinate-block-strong-projections}.
The peeling equivalence in the theorem gives the exact status rule.
Under the additional hypothesis, the first case is nonpeelable precisely
when \(D=\varnothing\), and the second case always peels.  Maximal
proper subsets of \(E\) have size \(m-1\); projection along any one
of them leaves the original family on its remaining marked coordinate.
This proves all the descent assertions.
\end{proof}

\begin{theorem}[Canonical reduction of repeated signed-clause columns]
\label{thm:damp-canonical-signed-column-reduction}
Let \(H\subseteq Q_F\) be an irredundantly embedded ample forbidden
pc-minor.  Index its signed minimal nonfaces by \(\mathcal N\), and
associate to each coordinate \(e\in F\) the column
\[
 v_e(I)=
 \begin{cases}
  0,&e\notin I,\\
  2\sigma_I(e)-1,&e\in I,
 \end{cases}
 \qquad I\in\mathcal N.
\]
Declare \(e\sim f\) if \(v_e=v_f\) or \(v_e=-v_f\).
This partitions \(F\) into classes \(E_1,\ldots,E_r\).
Reorient coordinates so that the columns within each class agree,
choose representatives \(e_j\in E_j\), and take the reduction
\begin{equation}
 R=H^{F\setminus\{e_1,\ldots,e_r\}}
 \label{eq:damp-canonical-signed-column-reduction}
\end{equation}
of the reoriented family.  Then \(R\) is an irredundant ample forbidden
pc-minor whose signed-clause columns are pairwise distinct up to sign.

The family \(H\) is recovered from \(R\) by replacing \(e_j\)
by a block of size \(m_j=|E_j|\), for every \(j\).  The reduced
family and these multiplicities are unique up to coordinate renaming
and reorientation.  Every sequence of compressions of equal columns
up to sign gives this same result when no such pair remains.

Conversely, every irredundant ample forbidden \(R\) with pairwise
distinct columns up to sign, together with arbitrary positive integers
\(m_1,\ldots,m_r\), gives an ample forbidden class by these replacements.
Thus the forbidden inputs with distinct columns, and the positive integer
multiplicities on their coordinates, form an exact parametrization.
The intrinsic description of those reduced forbidden inputs remains open.
\end{theorem}

\begin{proof}
Equality of columns up to sign is an equivalence relation.  Reversing
one coordinate negates its column, so the indicated normalization is
possible.  Consider a class \(E\) of size at least two after normalization.
Every clause either avoids \(E\) or contains all of \(E\), with a
constant forbidden sign on that block.  Let \(C\) be the set of old
words satisfying the clauses that avoid \(E\), and let \(B_0,B_1\)
be the two constant sections of \(H\).  Its nonconstant sections
are all \(C\).  Also \(B_0\cup B_1\subseteq C\).

Here ampleness forces the reverse inclusion.  If
\(x\in C\setminus(B_0\cup B_1)\), its face restriction in \(H\)
would be exactly \(Q_E\setminus\{0_E,1_E\}\).  This family has
\(2^{|E|}-2\) concepts and shatters every proper subset of \(E\),
so its shadow has \(2^{|E|}-1\) members.  It is not ample,
contradicting closure of ampleness under face restriction.  Hence
\(C=B_0\cup B_1\).  The section pattern is exactly
\eqref{eq:damp-coordinate-block-construction}.

Keeping one representative of \(E\), the reduction along the
other coordinates has sections \(B_0,B_1\), and is ample.  Call it
\(U\).  We have \(H=W_E(U)\), so the peeling equivalence of
Theorem~\ref{thm:damp-coordinate-block-replacement} makes \(U\)
nonpeelable.  The full-coordinate and descent assertion in
Theorem~\ref{thm:damp-separator-strong-descent}, which requires no
separator for that assertion, makes \(U\) an irredundant ample forbidden
pc-minor.  Its signed minimal nonfaces are obtained by replacing
\(E\) by its representative in every clause.

This last operation changes neither the set of clause rows nor the
entries in the remaining columns.  The rows correspond bijectively:
a support was a union of whole coordinate classes, so compressing a
class cannot identify two different supports or change their inclusion
relations.  Consequently no two previously inequivalent columns become
equivalent.  The remaining equivalence classes are exactly the original
classes with their compressed members removed.  Induction gives
\eqref{eq:damp-canonical-signed-column-reduction}, forbiddenness of
\(R\), and distinctness of its columns up to sign.  Reductions
commute, and choosing a different representative in a class merely renames
its coordinate.  Changing the normalization reverses that coordinate.
This proves independence of all compression choices and orders.

Conversely, the block-replacement theorem preserves forbiddenness at
every step.  Its signed clause rule shows that replacements on different
coordinates commute: each original column is simply repeated the chosen
number of times.  Since the input columns are distinct up to sign, the
resulting equivalence classes are exactly these repeated blocks.
Compression therefore recovers precisely \(R\) and the prescribed
multiplicities, proving the converse and uniqueness assertions.
\end{proof}

This reduction is canonical in the signed clause presentation and removes
all freedom due to repetition of signed-clause columns.  Its interaction
with graph three-core reduction is settled by
Theorem~\ref{thm:damp-alternating-reductions} below: the reductions
need not commute, but their alternation is forced and canonical.
Column reduction also differs from terminality under reduction.  A repeated
block supplies at least two nonpeelable one-coordinate projections, by
\eqref{eq:damp-coordinate-block-strong-projections}.  The verified
\(2318\)-concept tower has exactly one such projection, so its columns
are already distinct up to sign although it still admits a nonpeelable
reduction descent.

For an explicit control, use the fixed \(291\)-concept cornerless
obstruction and replace coordinate \(0\).  Its contraction has \(224\)
concepts and its reduction has \(67\).  All its nonempty-coordinate
reductions peel, because its VC dimension is three and those
projections have VC dimension at most two
\cite[Proposition~4.11]{ChalopiChepoiMoran22}.  Blocks of sizes two and three
therefore give cornerless forbidden three-cores of orders
\(3\cdot224+67=739\) and \(7\cdot224+67=1635\), with respectively
two and three maximal descent supports.  In the binary case both
singleton projections are nonpeelable but their joint projection peels;
the nonpeelable supports need not be closed under union.
The verifier \path{verify_damp_coordinate_block_replacement.py} checks
all \(832\) replacements of ample inputs through \(Q_3\) at block
sizes two and three, including \(1108\) corner lifts and the signed
clause and reduction formulas.  It independently checks the
fixed inputs and replays \(117\) elementary-minor orders, including
all proper elementary minors of the two displayed obstructions.
A further \(36\) positive carrier orders, together with cornerlessness
of all proper-block projections, check their complete descent-support
descriptions using reduction closure.
The same checks recover the canonical signed-column reduction of every
ample input through \(Q_3\), test every representative choice, and recover
the \(291\)-concept reduced input from both displayed replacements.
The record is \path{damp_coordinate_block_replacement_verification.json}.
The construction provides an intrinsic operation over every forbidden
input, together with exact recovery of its marked outputs.  It does
not characterize the forbidden inputs that admit no such reduction.

\begin{theorem}[Alternating graph and signed-column reductions]
\label{thm:damp-alternating-reductions}
Let \(H\subseteq Q_F\) be an irredundantly embedded ample forbidden
pc-minor, and put \(d=\operatorname{VC}(H)\).  Allow either of the
following reductions whenever it is available:
\begin{enumerate}
 \item delete a vertex of degree two;
 \item compress two signed minimal-nonface columns that agree up to sign,
 keeping one representative coordinate after reorientation.
\end{enumerate}
Every maximal sequence of these reductions ends at the same ample
forbidden family \(H_\ast\), up to coordinate renaming and reorientation.
It has minimum graph degree at least three and pairwise distinct signed
minimal-nonface columns up to sign.

More precisely, group each maximal consecutive run of reductions of the
same type into a \emph{phase}.  The phase sequence and its intermediate
families are canonical: a graph phase takes the full graph three-core,
and a column phase compresses all repeated columns.  The types alternate.
If there are \(q\) column phases and \(s\) graph phases, then
\begin{equation}
 q\le d-\operatorname{VC}(H_\ast)\le d-3,
 \qquad s\le q+1,
 \qquad q+s\le2d-5.
 \label{eq:damp-alternating-phase-bound}
\end{equation}
Every phase preserves the difference
\(\operatorname{idim}(H)-\operatorname{VC}(H)\).

Conversely, start from any irredundant ample forbidden family \(B\)
of minimum degree at least three whose signed-clause columns are distinct
up to sign.  Apply a finite alternating sequence of the following two
constructions, starting with either type, or apply neither:
\begin{enumerate}
 \item nonempty admissible square attachment data of
 Definition~\ref{def:damp-square-attachment-data};
 \item coordinate-block replacements with positive integer multiplicities,
 at least one of which is greater than one.
\end{enumerate}
All outputs are ample forbidden families.  Their canonical phase reductions
recover exactly the construction sequence, its square attachment data,
its coordinate multiplicities, and its base \(B\), up to the renamings
allowed in the two constituent recovery theorems.  Every ample forbidden
family is obtained this way.
In particular, a minimum-cardinality ample forbidden family is already
fixed by both reductions.
\end{theorem}

\begin{proof}
The decisive fact is that the two types of reduction cannot be available
simultaneously.  If a forbidden family has repeated signed-clause columns,
Theorem~\ref{thm:damp-canonical-signed-column-reduction} expresses it as
a nontrivial coordinate-block replacement of another ample forbidden
family.  Theorem~\ref{thm:damp-coordinate-block-replacement} makes the
output have minimum degree at least three.  Thus any forbidden family
with a degree-two vertex has no repeated columns.

Every allowed reduction preserves ample forbidden-minor status, by
Theorems~\ref{thm:damp-three-core} and
\ref{thm:damp-canonical-signed-column-reduction}.  It strictly reduces
cardinality.  This is immediate for a vertex deletion.  For compression
of a pair, write the family before compression as \(W_{\{e,f\}}(U)\).
Its excess over the compressed input \(U\) is
\[
 |W_{\{e,f\}}(U)|-|U|=2|\pi_z U|>0
\]
by \eqref{eq:damp-coordinate-block-order}.  Here \(z\) is the
representative coordinate of \(U\), which is nonempty and forbidden.
Consequently every reduction sequence terminates.

Whenever a graph reduction is possible, no column reduction is possible
until all degree-two vertices have been removed.  The graph three-core
theorem makes the resulting family independent of every deletion choice.
Whenever a column reduction is possible, no graph reduction is possible
as long as any repeated pair remains.  The canonical column-reduction
theorem makes the resulting family independent of the choices and order
of all pair compressions.  At the end of a phase, that type is exhausted;
either the other type is available or both are exhausted.  Induction over
the strictly decreasing cardinalities proves that the phase sequence,
its intermediate families, and its final family are canonical.  A forbidden
family has no vertex of degree zero or one, by the graph three-core theorem.
Exhaustion of both types therefore gives the two asserted properties
of \(H_\ast\).

A graph phase changes neither isometric dimension nor VC dimension.
A column phase with multiplicities \(m_1,\ldots,m_r\) decreases each
of these dimensions by \(\sum_j(m_j-1)\), by the coordinate-block
shadow formula.  This sum is at least one.  Hence their difference is
invariant and \(q\le d-\operatorname{VC}(H_\ast)\).  Every ample
forbidden class has VC dimension at least three
\cite[Proposition~4.11]{ChalopiChepoiMoran22}.  Alternation permits
at most one more graph phase than column phases, proving
\eqref{eq:damp-alternating-phase-bound}.

We finally verify that the inverse description needs no additional
compatibility conditions.  A nonempty square attachment over a forbidden
graph three-core has that core as its unique graph reduction, by
Theorem~\ref{thm:damp-square-attachment-data}.  Its larger output has a
degree-two vertex, so the mutual-exclusion argument gives distinct
signed-clause columns.  It can therefore serve as the input to the next
coordinate-block phase, and canonical column compression will recover
it exactly.  A nontrivial block replacement of an input with distinct
columns has that input as its canonical column reduction; its output
has minimum degree at least three and can serve as the next square
attachment base.  Both statements hold at the initial base \(B\),
which has both required properties.  Induction proves soundness and
exact phase-by-phase recovery for every alternating construction sequence.
Conversely, reversing the canonical reductions of \(H\) supplies precisely
such a sequence and the already proved intrinsic attachment data and
multiplicities.  The endpoint satisfies both required base conditions.
If a minimum-cardinality obstruction admitted either reduction, the
strictly smaller forbidden output would contradict minimality.
\end{proof}

The individual full reductions need not commute.  Let \(A\) be the
\(579\)-concept square extension of the \(577\)-concept three-core
described after Theorem~\ref{thm:damp-square-completed-rooted-blocks},
in the coordinates of \path{damp_three_core_verification.json}.
Replace coordinate \(0\)
by \(\{0,24\}\), obtaining a \(1603\)-concept family \(X\).
Then add the word \(24578\), which completes a square on
\(\{1,14\}\) in the constant \(00\)-section.  The resulting
\(1604\)-concept obstruction \(Y\) has the forced phase sequence
\begin{equation}
 Y_{1604}\xrightarrow{\text{graph}}X_{1603}
 \xrightarrow{\text{columns}}A_{579}\xrightarrow{\text{graph}}B_{577},
 \label{eq:damp-alternating-fixed-example}
\end{equation}
with graph, column, and graph phases, respectively.  The source has
isometric dimension \(25\) and VC dimension \(4\); the final base
has dimensions \(24\) and \(3\).  Thus the dimension difference
is \(21\), and three phases attain \eqref{eq:damp-alternating-phase-bound}
at VC dimension four.  On \(X\), graph reduction is the identity,
whereas column reduction gives \(A\), whose graph reduction has
\(577\) concepts.  Applying the two full reductions once in opposite
orders therefore gives different outputs, of orders \(579\) and \(577\).

The standalone verifier \path{verify_damp_alternating_reductions.py}
checks the four ample families from actual cubes, replays \(294\)
proper elementary-minor orders, and verifies nonpeelability by the
original negative-state certificate or its unchanged reduction.
Both column representatives and both graph-deletion priorities give the
same final base.  The recorded abstract square data reconstruct the two
graph phases in either available dependency order.  The report is
\path{damp_alternating_reductions_verification.json}.

The theorem accounts for all square attachment and coordinate repetition
layers, including their interaction.  The remaining intrinsic construction
problem lies in the final forbidden bases with minimum degree at least
three and distinct signed-clause columns.  These two conditions alone do
not imply forbiddenness.  For example, take the \(42\)-concept down-set
of words of Hamming weight at most three in \(Q_6\).  Its shattered
complex is that same down-set, so it is ample; deleting maximal words
preserves this property and gives a peeling.  Vertices of weight three
have degree three, and all other vertices have degree six.  Its minimal
nonfaces are all four-element supports, each with missing trace \(1111\).
For any two coordinates there is such a support containing one and
omitting the other, so their columns cannot agree up to sign.  The same
verifier checks this positive control.  Nor does being fixed by these two reductions imply
that every nonempty-coordinate reduction peels: the preceding
\(2318\)-concept example remains a counterexample to that implication.


The earlier opposite-status clause pairs give stronger tests of this
remaining problem.  Take the intermediate families \(X'_j,Y'_j\) in
Corollary~\ref{cor:damp-remote-forbidden-clause-pivot} for \(j=0,1\),
of orders \(577,1158\).  All four are periphery-free and already
fixed by both reductions.  Their respective minimum-degree pairs are
\((4,3)\) and \((4,4)\).  Within each pair the full shadow is identical,
only one four-coordinate forbidden trace changes, every proper pc-minor
peels, and the two original families have opposite peeling status.
Thus the reduced-base problem still retains the signed distinction in
these existing examples.  The independent replay
\path{verify_damp_reduced_base_clause_pairs.py} checks actual cubes,
all signed clauses, both fixed-point conditions, the two positive orders,
the complete negative certificates, and all \(300\) proper-minor orders.
Its report is \path{damp_reduced_base_clause_pairs_verification.json}.
This is a verification of the remaining difficulty, not a new
construction or an intrinsic admissibility criterion.

\begin{remark}[What overlap descent does and does not classify]
\label{rem:damp-overlap-descent-scope}
Theorem~\ref{thm:damp-inherited-overlap-descent} separates two genuinely
different parts of the requested parametrization.  Its terminal objects
are the \emph{overlap-atomic obstructions}, namely those forbidden
pc-minors whose one-coordinate reductions all belong to
\(\Damp\).  Every other obstruction is obtained from a smaller obstruction
by one or more inverse steps of
\eqref{eq:damp-inherited-overlap-repair-expansion}.  The antipodal
two-repair towers are explicit examples of such inverse steps.

This is a completeness decomposition, not yet a generative
classification.  In particular, the three displayed \(\Damp\)-memberships
in the theorem are consequences for a recovered obstruction; using them
as admissibility tests for arbitrary wings would be circular.  Two
intrinsic tasks remain: classify the overlap-atomic seeds by their signed
Yang gluing data, and replace those three memberships by local attachment
axioms that guarantee a valid inverse repair step.
\end{remark}

\begin{remark}[The candidate global Yang invariant]
\label{rem:damp-relative-yang-gluing-invariant}
Proposition~\ref{prop:damp-yang-two-cone-gluing-dichotomy} identifies a
smaller global object than a peeling order.  By
Theorem~\ref{thm:damp-inherited-overlap-descent}, it is needed only for
the terminal overlap-atomic seeds.  For every coordinate of such a seed,
retain the paired relative Yang complexes

\begin{equation}
 \mathfrak R_e(H)=
 \bigl((\Delta_{B_0},\Delta_A),
       (\Delta_{B_1},\Delta_A)\bigr).
 \label{eq:damp-paired-relative-yang-object}
\end{equation}

Each relative face module is fully Cohen--Macaulay and has \(0/1\) flag
\(h\)-support
\(\operatorname{Sh}(B_b)\setminus\operatorname{Sh}(A)\), by
Lemma~\ref{lem:damp-relative-yang-cm-flag-h}.  Thus ordinary homology and
the flag \(h\)-vector cannot detect the gluing-critical signature.  What is
missing is a finite signed incidence invariant describing how the two
relative modules attach to their common Yang ball.  Such an invariant
must be extracted directly from the oriented one-hole and punctured
expansion profiles and must force the absence or presence of a relative
shelling without including a shelling order in its data.

This is a research reduction, not the requested intrinsic
parametrization.  In particular, replacing
\eqref{eq:damp-paired-relative-yang-object} by a search over relative
shellings would merely rename the original problem.  The expression
\emph{extension/holonomy class} is therefore only a specification for the
missing invariant, not the name of an invariant already constructed.
\end{remark}

The relative Yang theory nevertheless extracts a first, genuinely static
part of the missing invariant.  For \(A\subseteq B\subseteq Q_F\), write

\[
 G(B,A):=Q_F[B\setminus A]
\]

for the graph induced by the relative facets.  Its edges are naturally
labelled by their cube coordinates.

\begin{proposition}[Cycle carriers at an overlap-atomic coordinate]
\label{prop:damp-atomic-wing-cycle-carriers}
Let \(H\in\Obs(\Damp)\cap\Ample\) be overlap-atomic, let \(e\) be active,
and write

\[
 B_b=\rho_e^bH,\qquad A=B_0\cap B_1,\qquad R_b=B_b\setminus A.
\]

Then both graphs \(G(B_0,A)=Q[R_0]\) and
\(G(B_1,A)=Q[R_1]\) contain a cycle.  In particular,

\begin{equation}
 |R_0|,|R_1|\ge4,
 \qquad
 |H|=2|A|+|R_0|+|R_1|\ge2|A|+8.
 \label{eq:damp-atomic-wing-cycle-order}
\end{equation}
\end{proposition}

\begin{proof}
The overlap \(A\) and the two sections \(B_0,B_1\) belong to \(\Damp\).
In particular, the relative Yang module of each nested pair
\(A\subseteq B_b\) is Cohen--Macaulay.  The forest theorem for relative
Yang modules says that if \(G(B_b,A)\) is a forest, then
\((\Delta_{B_b},\Delta_A)\) has a relative shelling
\cite{BousquetYangComplex26}.  Proposition
\ref{prop:damp-yang-two-cone-gluing-dichotomy} would then make \(H\)
corner-peelable, a contradiction.  Thus both relative-facet graphs contain
cycles.  A cube graph is bipartite, so each cycle has at least four
vertices.  The section decomposition gives

\[
 |H|=|B_0|+|B_1|=2|A|+|R_0|+|R_1|,
\]

which proves \eqref{eq:damp-atomic-wing-cycle-order}.
\end{proof}

\begin{remark}[The rooted signed cycle system]
\label{rem:damp-rooted-signed-cycle-system}
Proposition~\ref{prop:damp-atomic-wing-cycle-carriers} identifies the first
nonzero layer of a possible global invariant.  It is useful to retain more
than the two unlabelled cycle spaces.  For a nested pair \(A\subseteq B\),
adjoin a root vertex \(*\) to \(G(B,A)\), and for every
\(x\in B\setminus A\) and every coordinate \(q\) with
\(x^{\{q\}}\in A\), add a root edge \(x*\) labelled by the signed
boundary incidence \((q,x^{\{q\}})\).  Together with the coordinate
labels on the internal
edges, this gives a finite rooted signed graph

\[
 \widehat G(B,A).
\]

It is recovered directly from the facet--ridge incidence of the relative
Yang pair and involves no ordering.  Thus an overlap-atomic coordinate
canonically supplies the paired rooted cycle system

\[
 \bigl(\widehat G(B_0,A),\widehat G(B_1,A)\bigr).
\]

The proposition proves that its two internal cycle spaces are nonzero.
Cycle rank alone is not sufficient: shellable relative pairs can also have
cyclic relative-facet graphs.  The full signed rooted graph carries more
information.  Proposition~\ref{prop:damp-rooted-relative-presentation}
below shows that it is the complete first-syzygy presentation of the
relative Alexander-dual module and hence determines relative shellability.
Thus higher-dimensional Yang incidences are not needed to decide the
individual relative problem.  The remaining difficulty is to describe the
nonshellable rooted presentations structurally, and then to describe the
compatibility of the two rooted presentations at a gluing, without hiding a
linear-quotient order in admissibility.

Retaining every concept label and boundary endpoint would still be too
close to copying the two sections to serve as the final parameters.  The
presentation theorem instead directs attention to smaller labelled block,
cut, or circuit data extracted from the graph.  The root-cut criterion in
the same proposition is the first such compression.
\end{remark}

For the next statement, fix a field \(k\), put

\[
 S_F=k[v_{q,0},v_{q,1}:q\in F],
 \qquad
 m_x=\prod_{q\in F}v_{q,1-x_q},
 \qquad
 J_X=(m_x:x\in X),
\]

and, for \(A\subseteq B\subseteq Q_F\), write

\[
 \mathcal Q(B,A)=J_B/J_A.
\]

The vertices of \(\widehat G(B,A)\) are signed by their cube words.  Thus
an edge in coordinate \(q\) incident with \(x\) carries both its coordinate
and the literal \(v_{q,x_q}\).

\begin{proposition}[Rooted presentation and root-cut criterion]
\label{prop:damp-rooted-relative-presentation}
Let \(A\subsetneq B\subseteq Q_F\) be ample and put
\(D=B\setminus A\), with \(J_\varnothing=0\).  Let

\[
 \varphi:\bigoplus_{x\in D}S_F\varepsilon_x\longrightarrow
             \mathcal Q(B,A),
 \qquad
 \varepsilon_x\longmapsto\overline{m_x}.
\]

Then

\begin{align}
 \ker\varphi={}&
 \sum_{\substack{x\in D,\ q\in F\\x^{\{q\}}\in A}}
 S_Fv_{q,x_q}\varepsilon_x
 \notag\\
 &+\sum_{\substack{\{x,y\}\in E(Q[D])\\y=x^{\{q\}}}}
 S_F\bigl(v_{q,x_q}\varepsilon_x
          -v_{q,y_q}\varepsilon_y\bigr).
 \label{eq:damp-rooted-relative-presentation}
\end{align}

Consequently, the signed rooted graph \(\widehat G(B,A)\) determines the
relative module with its distinguished generators and therefore determines
all relative shelling orders of \((\Delta_B,\Delta_A)\).

More explicitly, for \(x\in D\), let

\[
 L_A(x)=\{q\in F:x^{\{q\}}\in A\}
\]

be the coordinate labels of the root edges incident with \(x\).  Delete
from \(\widehat G(B,A)\) every internal or root edge whose coordinate label
belongs to \(L_A(x)\).  Then

\begin{equation}
 A\cup\{x\}\text{ is ample}
 \quad\Longleftrightarrow\quad
 x\text{ and the root lie in different components of the resulting graph}.
 \label{eq:damp-root-cut-first-relative-facet}
\end{equation}
\end{proposition}

\begin{proof}
Since \(A\) and \(B\) are ample, the relative Yang face module is
Cohen--Macaulay.  Its squarefree Alexander dual
\(\mathcal Q(B,A)\) therefore has an \(|F|\)-linear resolution; in
particular, it is linearly presented.  The general monomial presentation
has one annihilator relation for every pair \(x\in D,a\in A\), and one
pair relation for every two members of \(D\).  A relation of the first kind
is linear exactly when \(x\) and \(a\) are cube neighbours, and a relation
of the second kind is linear exactly when its two concepts are cube
neighbours.  Since the entire kernel is generated by its linear part, these
are precisely the root and internal relations displayed in
\eqref{eq:damp-rooted-relative-presentation}.

The presentation matrix is read verbatim from the signed rooted graph.
An isomorphism of such graphs, including the coordinate and sign labels,
therefore identifies the relative modules and their distinguished minimal
generators.  Linear quotients of those generators are the
Alexander-dual form of a relative shelling, proving the comparison
assertion.

It remains to prove the cut criterion.  The monomial colon formula gives

\begin{equation}
 J_A:m_x
 =\left(
   \prod_{q\in\operatorname{Diff}(x,a)}v_{q,x_q}:a\in A
  \right).
 \label{eq:damp-root-cut-colon}
\end{equation}

Its linear monomials are exactly
\(v_{q,x_q}\) for \(q\in L_A(x)\).  By the one-concept colon criterion,
\(A\cup\{x\}\) is ample precisely when every generator in
\eqref{eq:damp-root-cut-colon} is divisible by one of these variables.

Suppose first that the edge-deleted graph contains an \(x\)-to-root path.
Its internal edges and final root edge all have labels outside \(L_A(x)\).
The endpoint \(a\in A\) of the final root edge consequently agrees with
\(x\) on every coordinate in \(L_A(x)\).  The corresponding generator in
\eqref{eq:damp-root-cut-colon} is divisible by none of the displayed
linear variables, so \(A\cup\{x\}\) is not ample.

Conversely, suppose that the colon is not generated by those variables.
Then some \(a\in A\) agrees with \(x\) on \(L_A(x)\).  Apply
\eqref{eq:damp-rooted-relative-presentation} to the annihilator relation
indexed by \(x,a\), in its signed multidegree \(\mu\).  A relative generator
divides \(\mu\) exactly when its cube word lies in the Boolean interval
\([x,a]\).  After dividing each surviving basis vector by the unique
monomial needed to reach \(\mu\), an internal edge relation becomes
\(e_z-e_w\), a root relation becomes \(e_z\), and the annihilator relation
indexed by \(x,a\) becomes \(e_x\).  Consequently \(e_x\) lies in the span
of the edge differences and rooted unit vectors of this finite graph.  In
any component containing no rooted unit vector, the coordinate-sum
functional vanishes on every edge difference but not on \(e_x\).  The
component of \(x\) must therefore contain a root edge.  This gives an
\(x\)-to-root path whose labels all belong to
\(\operatorname{Diff}(x,a)\), hence avoid \(L_A(x)\).  Thus \(x\) and the
root remain connected after the prescribed edge deletion, proving
\eqref{eq:damp-root-cut-first-relative-facet}.
\end{proof}

\begin{definition}[Root-locked and biroot-locked relative pairs]
\label{def:damp-root-locked-relative-pair}
The nested ample pair \(A\subsetneq B\) is \emph{lower-root-locked} if,
for every \(x\in B\setminus A\), the root remains connected to \(x\) after
deleting all edges labelled in \(L_A(x)\).  Equivalently, no one-concept
ample extension of \(A\) lies in \(B\).

Apply the same construction to the complementary nested pair

\[
 Q_F\setminus B\subsetneq Q_F\setminus A.
\]

We call \(A\subsetneq B\) \emph{upper-root-locked} if that complementary
pair is lower-root-locked, and \emph{biroot-locked} if it is both
lower- and upper-root-locked.
\end{definition}

\begin{corollary}[Static terminal-pair certificate]
\label{cor:damp-biroot-locked-terminal-pair}
Lower-root-locking is a finite, ordering-free certificate that
\((\Delta_B,\Delta_A)\) is not relatively shellable.  Upper-root-locking
is equivalent to

\[
 B\setminus\{x\}\text{ being nonample for every }x\in B\setminus A.
\]

Thus a biroot-locked pair admits neither a first nor a last relative
shelling facet.  Both assertions are decided entirely by labelled cut
tests in the two rooted graphs.
\end{corollary}

\begin{proof}
The first assertion follows immediately from
Proposition~\ref{prop:damp-rooted-relative-presentation}: a relative
shelling must begin with an ample one-concept extension of \(A\).  For the
second, observe that

\[
 B\setminus\{x\}\text{ is ample}
 \quad\Longleftrightarrow\quad
 (Q_F\setminus B)\cup\{x\}\text{ is ample},
\]

and apply the same root-cut criterion to the complementary pair.
\end{proof}

\begin{remark}[What the rooted presentation accomplishes]
\label{rem:damp-rooted-presentation-scope}
The presentation theorem changes the proof contract in two ways.  First,
the signed rooted graph is not merely a cycle-count proxy: it is a complete
invariant of each individual relative shelling problem.  Second,
lower- or upper-root-locking gives a direct nonshellability certificate
without a shelling or carrier order.  The remaining classification problem
is now sharper.  One must extract a smaller family of labelled cut blocks
whose intrinsic assembly axioms force the required locking, describe how
two such presentations attach to the same overlap, and recover those blocks
from every overlap-atomic obstruction.  The full rooted graph nearly records
the original pair, while merely asking whether its presentation has linear
quotients would still be circular.  Neither is an admissible final
parametrization.
\end{remark}

\begin{definition}[Relative one-hole block cover]
\label{def:damp-relative-one-hole-cover}
Let \(A\subseteq Q_F\), and let \(\mathcal W\) be a finite family of
oriented one-hole data \((y,S)\) on \(F\).  Put
\[
 \mathsf I(y,S)=\{y^T:T\subseteq S,\ |T|\ge2\},
 \qquad
 B(A,\mathcal W)
 =A\cup\bigcup_{(y,S)\in\mathcal W}\mathsf P(y,S).
\]
We call \((A,\mathcal W)\) an \emph{admissible relative one-hole cover} if
\begin{enumerate}
 \item \(A\) and \(B(A,\mathcal W)\) satisfy their Sandwich equalities;
 \item every displayed tip is absent:
       \(y\notin B(A,\mathcal W)\) for \((y,S)\in\mathcal W\); and
 \item \(B(A,\mathcal W)\setminus A\ne\varnothing\); and
 \item the interiors of the displayed punctured cubes protect every
       relative concept:
       \begin{equation}
        B(A,\mathcal W)\setminus A
        \subseteq\bigcup_{(y,S)\in\mathcal W}\mathsf I(y,S).
        \label{eq:damp-relative-hole-protection-cover}
       \end{equation}
\end{enumerate}
These are finite conditions on \(A\) and the displayed punctured cubes.
In particular, admissibility contains no deletion, corner, shelling,
linear-quotient, or rooted-connectivity predicate.
\end{definition}

\begin{theorem}[One-hole normal form for upper-root-locked pairs]
\label{thm:damp-relative-one-hole-upper-lock}
Every admissible relative one-hole cover constructs a nested ample pair
\[
 A\subseteq B(A,\mathcal W)
\]
that is upper-root-locked.  Conversely, every upper-root-locked nested
ample pair \(A\subsetneq B\subseteq Q_F\) is obtained in this way.

More precisely, put \(D=B\setminus A\).  Its canonical recovered block
family is
\begin{equation}
 \mathcal W_{\max}(B,A)=
 \left\{(y,S):
 \begin{array}{l}
   y\notin B,\ \mathsf P(y,S)\subseteq B,\ |S|\ge2,\\
   \mathsf I(y,S)\cap D\ne\varnothing,\text{ and }S
   \text{ is inclusion-maximal with }\mathsf P(y,S)\subseteq B
 \end{array}\right\}.
 \label{eq:damp-relative-one-hole-recovery}
\end{equation}
Consequently, upper-root-locked relative Yang presentations admit an
explicit, descendant-free construction and recovery by punctured-cube
blocks.
\end{theorem}

\begin{proof}
Let \((A,\mathcal W)\) be admissible and put
\(B=B(A,\mathcal W)\).  The two Sandwich equalities make \(A\) and \(B\)
ample.  Fix \(x\in B\setminus A\).  By
\eqref{eq:damp-relative-hole-protection-cover}, some
\((y,S)\in\mathcal W\) and \(T\subseteq S\), \(|T|\ge2\), satisfy
\(x=y^T\).  For every proper \(C\subsetneq T\),
\[
 x^C=y^{T\mathbin{\triangle}C}\in\mathsf P(y,S)\subseteq B,
 \qquad
 x^T=y\notin B.
\]
Thus \(T\) is a nontrivial minimal nonface of \(\Lambda_B(x)\).
The concept \(x\) is therefore contained in more than one maximal cube
of \(B\), so it is not a corner of the ample family \(B\).  Equivalently,
\(B\setminus\{x\}\) is nonample.  This holds for every
\(x\in B\setminus A\), which is upper-root-locking.

Conversely, let \(A\subsetneq B\) be upper-root-locked and ample.  For
each \(x\in B\setminus A\), the family \(B\setminus\{x\}\) is nonample.
Hence \(x\) is not a corner of \(B\), so its local cube-support complex
\(\Lambda_B(x)\) is not a simplex and has a minimal nonface \(S\) of
order at least two.  For every such \(S\), put \(y=x^S\).  Minimality
gives
\[
 \mathsf P(y,S)\subseteq B,
 \qquad y\notin B,
 \qquad x\in\mathsf I(y,S).
\]
Extend each such \(S\), with its tip \(y\) fixed, to an inclusion-maximal
support \(S'\) satisfying \(\mathsf P(y,S')\subseteq B\).  Then
\(x\in\mathsf I(y,S')\), so the finite family
\(\mathcal W_{\max}(B,A)\) protects all of \(B\setminus A\).  Every one
of its punctured cubes lies in \(B\).  It follows that
\[
 B(A,\mathcal W_{\max}(B,A))=A\cup(B\setminus A)=B.
\]
The missing-tip conditions hold by construction, while both Sandwich
equalities follow from the ampleness of \(A\) and \(B\).  Hence the
recovered cover is admissible and reconstructs the original pair.
\end{proof}

\begin{corollary}[Root cuts and one-hole blocks]
\label{cor:damp-root-cut-one-hole-comparison}
For nested ample \(A\subsetneq B\), the following are equivalent:
\begin{enumerate}
 \item \(A\subsetneq B\) is upper-root-locked;
 \item every \(x\in B\setminus A\) belongs to
       \(\mathsf I(y,S)\) for some punctured cube
       \(\mathsf P(y,S)\subseteq B\) with \(y\notin B\) and \(|S|\ge2\);
 \item in the complementary rooted presentation, every relative vertex
       remains joined to the root after deleting the labels incident with
       it at the root.
\end{enumerate}
If \(x=y^T\in\mathsf I(y,S)\) is protected as in (2), then
\(T\cap\{q:x^{\{q\}}\notin B\}=\varnothing\).  The block therefore gives,
through Proposition~\ref{prop:damp-rooted-relative-presentation}, a rooted
path avoiding all deleted labels.
\end{corollary}

\begin{proof}
The equivalence of (1) and (2) is
Theorem~\ref{thm:damp-relative-one-hole-upper-lock}; the equivalence of
(1) and (3) is the complementary root-cut criterion.  For the final
assertion, \(x^{\{q\}}\in B\) for every \(q\in T\), because
\(T\setminus\{q\}\) is a nonempty subset of \(S\).  Hence no coordinate
of \(T\) is deleted in the complementary cut test.  The root-cut
criterion converts the corresponding nonlinear missing-tip relation into
the asserted path.
\end{proof}

\begin{remark}[What has and has not been parametrized]
\label{rem:damp-relative-one-hole-scope}
Theorem~\ref{thm:damp-relative-one-hole-upper-lock} closes the individual
upper-locking row: one need not classify arbitrary rooted graphs or use
locking as an admissibility oracle.  The primitive blocks may be taken to
be the inclusion-maximal punctured cubes, one of which can protect many
relative concepts, and they are canonically recovered.  This is not yet the
basis classification.  The Hall audit shows that both wings at every
tested coordinate have this normal form, but the missing theorem must
still characterize which paired covers over a common overlap have all
elementary pc-minors peelable.  In particular, appending arbitrary
shelling certificates for those images would violate the non-circularity
gate.
\end{remark}

\begin{definition}[Internal and cross-layer protection]
\label{def:damp-internal-cross-protection}
For \(B\subseteq Q_F\), let
\begin{equation}
 \operatorname{Int}(B)=
 \{x\in B:\text{there are }y\notin B,\ S\subseteq F
   \text{ with }\mathsf P(y,S)\subseteq B
   \text{ and }x\in\mathsf I(y,S)\}.
 \label{eq:damp-internal-protection}
\end{equation}
For \(B_0,B_1\subseteq Q_F\), put \(A=B_0\cap B_1\).  For
\(b\in\{0,1\}\), let \(\operatorname{Cross}_b(B_0,B_1)\) consist of the
words \(a\in A\) for which there is a nonempty \(T\subseteq F\) such that
\begin{equation}
 \{a^U:U\subseteq T\}\subseteq B_b,
 \qquad
 \{a^U:U\subsetneq T\}\subseteq B_{1-b},
 \qquad
 a^T\notin B_{1-b}.
 \label{eq:damp-cross-protection}
\end{equation}
Thus internal protection is supplied by an absent-tip punctured cube in
one section.  Cross-layer protection is supplied by a full \(T\)-cube on
one side and the same cube punctured at \(a^T\) on the other.
\end{definition}

\begin{theorem}[Paired one-hole criterion for a cornerless gluing]
\label{thm:damp-paired-one-hole-cornerless-gluing}
Let \(B_0,B_1\subseteq Q_F\), and form the two-layer family
\[
 H=(B_0\times\{0\})\cup(B_1\times\{1\})
 \subseteq Q_{F\cup\{e\}}.
\]
Suppose that \(H\) is nonempty and satisfies the Sandwich equality.  Then \(H\) is
cornerless if and only if
\begin{equation}
 B_b=\operatorname{Int}(B_b)
      \cup\operatorname{Cross}_b(B_0,B_1)
 \qquad(b\in\{0,1\}).
 \label{eq:damp-paired-protection-cover}
\end{equation}
All the sets in \eqref{eq:damp-paired-protection-cover} are computed
directly from full and punctured affine cubes in the two sections.  Hence
this is an explicit, ordering-free construction and recovery theorem for
the nonpeelable root of a two-layer gluing.
\end{theorem}

\begin{proof}
The Sandwich equality makes \(H\) ample.  Fix \(x\in B_b\) and write
\(\widetilde x=(x,b)\).  By the local-hole characterization,
\(\widetilde x\) is not a corner of \(H\) exactly when there are a missing
tip \(z\notin H\) and a support \(R\), \(|R|\ge2\), such that
\[
 \widetilde x=z^R,
 \qquad
 \mathsf P(z,R)\subseteq H.
 \tag{*}\label{eq:damp-paired-protection-hole}
\]

Suppose first that \(e\notin R\).  The entire punctured cube in
\eqref{eq:damp-paired-protection-hole} lies in layer \(b\).  Suppressing
\(e\) gives a punctured cube in \(B_b\) whose protected interior contains
\(x\).  Thus \(x\in\operatorname{Int}(B_b)\).  Conversely, every block
in \eqref{eq:damp-internal-protection} lifts within layer \(b\) and gives
\eqref{eq:damp-paired-protection-hole}.

Now suppose that \(e\in R\), and put \(T=R\setminus\{e\}\).  Since
\(|R|\ge2\), the set \(T\) is nonempty.  The tip is
\[
 z=(x^T,1-b).
\]
The layer-\(b\) section of \(\mathsf P(z,R)\) is the full \(T\)-cube
through \(x\), whereas its layer-\((1-b)\) section is that cube punctured
at \(x^T\).  The assertion \(z\notin H\) says precisely that
\(x^T\notin B_{1-b}\).  These are exactly the three conditions in
\eqref{eq:damp-cross-protection}.  The converse follows by reversing the
same construction: the full/punctured pair in
\eqref{eq:damp-cross-protection} is
\(\mathsf P((x^T,1-b),T\cup\{e\})\).

We have proved, for every \(x\in B_b\),
\[
 (x,b)\text{ is not a corner of }H
 \quad\Longleftrightarrow\quad
 x\in\operatorname{Int}(B_b)
       \cup\operatorname{Cross}_b(B_0,B_1).
\]
The family \(H\) is cornerless exactly when this holds for both layers,
which is \eqref{eq:damp-paired-protection-cover}.
\end{proof}

\begin{remark}[The remaining paired problem]
\label{rem:damp-paired-one-hole-scope}
Equation~\eqref{eq:damp-paired-protection-cover} is the missing
nonpeelability compatibility law for the two sides: internal holes protect
a section concept on their own, while full/punctured cross blocks protect
the concepts exposed by the section split.  It does not assert that the
\(3|F\cup\{e\}|\) elementary images are peelable.  The intrinsic
generative classification now requires attachment axioms on these same
blocks whose transported direct and switched profiles construct peelable
elementary images without storing or searching for their shellings.
\end{remark}


\begin{definition}[Latent punctured-face system]
\label{def:damp-latent-punctured-face-system}
Let \(H\subseteq Q_F\).  Its \emph{latent punctured-face system} is
\begin{equation}
 \mathcal L(H)
 =
 \{(y,S):y\in Q_F,\ S\subseteq F,\ |S|\ge2,
             \mathsf P(y,S)\subseteq H\}.
 \label{eq:damp-latent-punctured-face-system}
\end{equation}
The tip \(y\) is allowed to belong to \(H\).  If it does, then
\(\{y\}\cup\mathsf P(y,S)\) is a full \(S\)-cube; deleting \(y\)
activates the punctured face.  Thus \(\mathcal L(H)\) records both the
holes already present in \(H\) and the holes that can appear later during
a peeling.

For an ordering \(\mathbf x=(x_1,\ldots,x_m)\) of \(H\), put
\[
 \operatorname{pos}_{\mathbf x}(y)=
 \begin{cases}
  0,&y\notin H,\\
  j,&y=x_j,
 \end{cases}
\]
and let \(b_{\mathbf x}(y,S)\) be the first member of
\(\mathsf P(y,S)\) in this ordering.
\end{definition}

\begin{theorem}[Latent punctured-face word criterion]
\label{thm:damp-latent-punctured-face-word-criterion}
Let \(H\subseteq Q_F\) be ample.  An ordering
\(\mathbf x=(x_1,\ldots,x_m)\) is a corner peeling if and only if, for
every \((y,S)\in\mathcal L(H)\),
\begin{equation}
 \operatorname{pos}_{\mathbf x}(y)
 <
 \operatorname{pos}_{\mathbf x}\bigl(b_{\mathbf x}(y,S)\bigr)
 \quad\Longrightarrow\quad
 b_{\mathbf x}(y,S)=y^{\{q\}}
 \text{ for some }q\in S.
 \label{eq:damp-latent-punctured-face-word-rule}
\end{equation}
In words: if the tip has disappeared before any displayed vertex of a
punctured cube, then the first displayed vertex must be a neighbor of the
tip.
\end{theorem}

\begin{proof}
Suppose first that \(\mathbf x\) is a corner peeling.  Fix
\((y,S)\in\mathcal L(H)\), let
\(x=b_{\mathbf x}(y,S)=y^T\), and suppose that \(y\) occurs before \(x\),
with a missing tip understood to occur at time zero.  Immediately before
\(x\) is deleted, the tip is absent and every member of
\(\mathsf P(y,S)\) is still present.  If \(|T|\ge2\), then
\(\mathsf P(y,T)\) is still present.  Equivalently, every proper
\(T\)-flip of \(x\) is present while \(x^T=y\) is absent.  The local
cube-support complex at \(x\) is therefore not a simplex, so \(x\) is
not a corner, a contradiction.  Hence \(|T|=1\), which is
\eqref{eq:damp-latent-punctured-face-word-rule}.

Conversely, suppose that the displayed rule holds.  We prove inductively
that \(x_j\) is a corner of
\(H_j=\{x_j,\ldots,x_m\}\).  The preceding deletions preserve ampleness
by the induction hypothesis, so \(H_j\) is ample.  If \(x_j\) were not a
corner, its local cube-support complex would have a nonface
\(T\subseteq F\) of order at least two.  Replacing \(T\) by a minimal
nonface if necessary and putting \(y=x_j^T\), we obtain
\[
 y\notin H_j,
 \qquad
 \mathsf P(y,T)\subseteq H_j.
\]
Thus \((y,T)\in\mathcal L(H)\), the tip \(y\) is either absent from
\(H\) or occurs before \(x_j\), and \(x_j\) is the first member of
\(\mathsf P(y,T)\).  Since \(x_j=y^T\) and \(|T|\ge2\), this contradicts
\eqref{eq:damp-latent-punctured-face-word-rule}.  Therefore every
\(x_j\) is a corner, and \(\mathbf x\) is a corner peeling.
\end{proof}

The word criterion admits a particularly small finite state space.  Write
\(a(y,S)=y^S\) for the antipode of a latent block.  Once the tip has
disappeared, a corner deletion of the antipode forces one of its
\(S\)-neighbors to have disappeared already.

\begin{corollary}[Finite facet-protector criterion]
\label{cor:damp-latent-facet-protector-criterion}
Let \(H\subseteq Q_F\) be ample.  For each
\((y,S)\in\mathcal L(H)\), choose one of the following states:
\begin{enumerate}
 \item if \(y\in H\), the \emph{escape state}, contributing the arc
       \[
        y^S\longrightarrow y;
       \]
 \item a \emph{facet-protector state} marked by a coordinate \(q\in S\),
       contributing the arc
       \[
        y^{S\setminus\{q\}}\longrightarrow y^S.
       \]
\end{enumerate}
Then \(H\) is corner-peelable if and only if the states can be chosen so
that the union of their arcs is acyclic.
\end{corollary}

\begin{proof}
Let \(\mathbf x\) be a corner peeling and fix a latent block
\((y,S)\).  If \(y^S\) precedes the present tip \(y\), choose the escape
state.  Otherwise the tip is absent or precedes the antipode \(a=y^S\).
Immediately before \(a\) is deleted, suppose that every neighbor
\(y^{S\setminus\{q\}}\), \(q\in S\), is still present.  Then every
coordinate of \(S\) is incident with \(a\) in the remaining family.
Because \(a\) is a corner, the unique maximal cube through \(a\) contains
the cube generated by all its incident coordinates, including the
\(S\)-cube.  This would put the tip \(y\) in the remaining family, a
contradiction.  Hence some \(S\)-neighbor precedes \(a\); choose its
coordinate as the facet-protector state.  Every selected arc points
forward in \(\mathbf x\), so their union is acyclic.

Conversely, topologically order an acyclic selected-arc digraph and
suppose that the word rule fails for a latent block \((y,S)\).  Write its
first displayed vertex as \(x=y^T\), where \(|T|\ge2\).  The subblock
\((y,T)\) is also latent, its antipode is \(x\), and its tip is absent or
precedes \(x\).  It cannot carry the escape state, whose arc
\(x\to y\) would reverse these positions.  Its selected state is
therefore a facet protector
\[
 y^{T\setminus\{q\}}\longrightarrow x
 \qquad(q\in T).
\]
The protector is another displayed vertex of the original block and
precedes \(x\), contradicting the choice of \(x\).  Thus the topological
order satisfies \eqref{eq:damp-latent-punctured-face-word-rule}, and
Theorem~\ref{thm:damp-latent-punctured-face-word-criterion} makes it a
corner peeling.
\end{proof}


\begin{proposition}[Contraction transport of facet protectors]
\label{prop:damp-latent-facet-protector-transport}
Let \(H\subseteq Q_{F\cup\{e\}}\) be ample and put \(K=\pi_eH\).
For \(z\in K\), define its lift set
\[
 L_e(z)=\{b\in\{0,1\}:(z,b)\in H\}.
\]
Then every edge \(zz^{\{q\}}\) of \(K\) has a nonempty \emph{owner set}
\begin{equation}
 \Omega_e(z,q)
 =L_e(z)\cap L_e(z^{\{q\}})
 \ne\varnothing.
 \label{eq:damp-latent-facet-protector-owner}
\end{equation}
Consequently, if \((y,S)\in\mathcal L(K)\) has antipode \(z=y^S\),
then every facet-protector state \(q\in S\) lifts in each layer belonging
to \(\Omega_e(z,q)\).  Contraction preserves the coordinate \(q\); its
only additional state is the one- or two-element owner set
\(\Omega_e(z,q)\).

For a switched punctured expansion profile
\(P_S=D\mathbin{\dot\cup}Z_0\mathbin{\dot\cup}Z_1\), written in the
coordinates based at its antipode \(z\), one has
\begin{equation}
 \Omega_e(z^C,q)=
 \begin{cases}
  \{0,1\},
    &C,C\triangle\{q\}\in D,\\
  \{b\},
    &\{C,C\triangle\{q\}\}\not\subseteq D
      \text{ and its non-\(D\) endpoints lie in }Z_b.
 \end{cases}
 \label{eq:damp-switched-profile-protector-owner}
\end{equation}
Thus the separator and component colours recover the transport of every
facet-protector state without a peeling order.
\end{proposition}

\begin{proof}
Suppose that \(L_e(z)\cap L_e(z^{\{q\}})=\varnothing\).  Both lift sets
are nonempty because their projected vertices belong to \(K\); hence they
are opposite singletons, say
\[
 (z,0),(z^{\{q\}},1)\in H,
 \qquad
 (z,1),(z^{\{q\}},0)\notin H.
\]
The two present vertices have Hamming distance two.  Their only two
possible intermediate cube vertices are precisely the two absent
vertices displayed above, contradicting isometry of the ample family
\(H\).  This proves \eqref{eq:damp-latent-facet-protector-owner}.  The
protector of a latent block is the edge
\(z^{\{q\}}z\), so the first consequence follows immediately.

In a punctured expansion profile, the vertices of \(D\) have both lifts,
whereas the vertices of \(Z_b\) have only their layer-\(b\) lift.  The
canonical profile theorem shows that no edge joins \(Z_0\) to \(Z_1\).
Intersecting the two endpoint lift sets now gives exactly
\eqref{eq:damp-switched-profile-protector-owner}.
\end{proof}

\begin{remark}[What remains after protector transport]
\label{rem:damp-latent-facet-protector-transport-scope}
Facet-protector transport is canonical and coordinate preserving.  The
remaining compatibility is smaller but still global: owner choices must
agree when blocks share protector edges, and the resulting protector arcs
must synchronize with escape arcs \(y^S\to y\) of blocks whose tips are
present.  Selecting owners and testing the resulting digraph for
acyclicity would again be recognition.  The desired construction theorem
must instead derive these choices from the incidence of the ample
separators \(D\), their monochromatic components \(Z_b\), and the paired
section attachments.  This is the precise state-transport problem on which
a signed holonomy invariant can now be tested.
\end{remark}

\begin{definition}[Colored Yang protector atlas]
\label{def:damp-yang-protector-atlas}
Let \(H\subseteq Q_F\) be ample and write \(F_x\) for the Yang facet
belonging to \(x\in H\).  For \((y,S)\in\mathcal L(H)\), define the
\emph{punctured cross-polytopal chart}
\[
 \mathcal C_H(y,S)
 =\left\langle F_{y^T}:\varnothing\ne T\subseteq S\right\rangle
 \subseteq\Delta_H.
\]
It is the join of the simplex on the literals fixed outside \(S\) with
the boundary of the \(S\)-cross-polytope after the relative interior of
the facet indexed by \(y|_S\) has been removed.  Its distinguished
\emph{antipodal facet} is \(A(y,S)=F_{y^S}\).  For \(q\in S\), its
\emph{\(q\)-protector facet} is
\[
 P_q(y,S)=F_{y^{S\setminus\{q\}}};
\]
this is the facet of the chart sharing the \(q\)-ridge with \(A(y,S)\).
If \(y\in H\), the additional facet \(E(y,S)=F_y\) is called the
\emph{escape facet}.  The family of these labelled charts and facets is
the \emph{colored Yang protector atlas} \(\mathscr A(H)\).
\end{definition}

\begin{corollary}[Yang protector-atlas criterion]
\label{cor:damp-yang-protector-atlas}
Let \(H\subseteq Q_F\) be ample.  In each chart of
\(\mathscr A(H)\), choose either
\[
 A(y,S)\prec E(y,S)
 \quad\text{when the escape facet is present},
\]
or
\[
 P_q(y,S)\prec A(y,S)
 \quad\text{for some }q\in S.
\]
Then the Yang complex \(\Delta_H\) is shellable if and only if these
local precedence relations can be chosen so that their union is
acyclic.  In that case every linear extension is the facet order
corresponding, under the shellability--corner-peeling dictionary, to a
corner peeling of \(H\).
\end{corollary}

\begin{proof}
Under the facet--concept bijection, the displayed relations are exactly
the escape and facet-protector arcs of
Corollary~\ref{cor:damp-latent-facet-protector-criterion}.  Apply that
corollary and the shellability--corner-peeling correspondence.
\end{proof}

An elimination order is \emph{distance-preserving} if every nonempty
suffix induces an isometric subgraph of the original graph.  A vertex
\(v\) is \emph{weakly \(4\)-simplicial} if every two of its neighbors
have distance at most two after \(v\) is deleted.  A weakly
\(4\)-simplicial elimination order tests this condition in the graph
remaining at each step.

\begin{proposition}[Yang shellability is metric elimination]
\label{prop:damp-yang-shellability-metric-elimination}
Let \(\varnothing\ne H\subseteq Q_F\) be isometric, and let \(G_H=Q_F[H]\) be its
concept graph.  Then the following are equivalent:
\begin{enumerate}
 \item the Yang complex \(\Delta_H\) is shellable;
 \item \(\Delta_H\) is strongly shellable;
 \item \(G_H\) has a distance-preserving elimination order; and
 \item \(G_H\) has a weakly \(4\)-simplicial elimination order.
\end{enumerate}
Every isometric family satisfying these conditions is itself ample
and corner-peelable.  More precisely, its distance-preserving elimination
orders are exactly its corner-peeling orders, whose reversals give the
corresponding Yang shellings.
\end{proposition}

\begin{proof}
All Yang facets have the same cardinality.  Use the standard pairwise
criterion for a pure shelling order.  For earlier and later facets
\(F_x,F_y\), respectively, it supplies an earlier facet \(F_z\) such that

\[
 |F_y\setminus F_z|=1,
 \qquad
 F_x\cap F_y\subseteq F_z.
\]

Suppose that \(F_z\) and \(F_y\) differ in color \(e\).  If \(x_e=y_e\),
then the \(e\)-literal belongs to \(F_x\cap F_y\) but not to \(F_z\), a
contradiction.  Hence \(x_e\ne y_e\), and the new \(e\)-literal of \(F_z\)
belongs to \(F_x\).  Therefore \(F_z\subseteq F_x\cup F_y\), which is the
additional condition for a strong shelling.  Thus ordinary and strong
shellability coincide for Yang complexes.

The codimension-one graph of \(\Delta_H\) is \(G_H\): two Yang facets meet
in a ridge exactly when their concepts differ in one coordinate.  Moreover,
the facet distance between \(F_x\) and \(F_y\) is their Hamming distance.
Since \(H\) is isometric, this equals the distance between \(x\) and \(y\)
in \(G_H\).  Thus the codimension-one graph is harmonious.  The graph
characterization of pure strong shellability now identifies item~2 with a
distance-preserving elimination order
\cite[Theorem~4.7]{GuoShenWu19strong}.  Deleting a vertex preserves all
distances exactly when that vertex is weakly \(4\)-simplicial, so items~3
and~4 are equivalent \cite[Lemma~3.3 and Proposition~3.4]{SmithZahedi25distance}.
To obtain the last assertion without assuming ampleness, reverse a
distance-preserving elimination order, starting from its last singleton.
Every next family is isometric in the cube.  By
Lemma~\ref{lem:damp-one-concept-extension-test}, each addition preserves
ampleness and adds a corner.  Thus the original order is a corner
peeling.  Conversely, every suffix of a corner peeling is ample and
therefore isometric.  This proves the order-level assertion as well.
\end{proof}

\begin{proposition}[Metric criterion for a relative peeling]
\label{prop:damp-relative-metric-peeling}
Let \(\varnothing\ne A\subseteq H\subseteq Q_F\) be ample,
and list \(H\setminus A\) as \(x_1,\ldots,x_k\).
For a distance-two pair \(u,v\in H\), put
\(C_H(u,v)=N_H(u)\cap N_H(v)\).
The following conditions are equivalent:
\begin{enumerate}
 \item Deleting \(x_1,\ldots,x_k\) is a relative corner peeling
 of \(H\) to \(A\), with all of \(A\) retained.
 \item Each \(H\setminus\{x_1,\ldots,x_j\}\), \(0\le j\le k\),
 is isometric in the original concept graph of \(H\).
 \item For every distance-two pair \(u,v\in H\), either
 \(C_H(u,v)\cap A\ne\varnothing\), or there are
 \(z\in\{u,v\}\setminus A\) and \(w\in C_H(u,v)\)
 such that \(z\) precedes \(w\) in the displayed deletion list.
\end{enumerate}
The target \(A\) need not be corner-peelable.  A complete peeling
of \(H\) with \(A\) as its final block exists exactly when such
a relative peeling exists and \(A\) itself is corner-peelable.
\end{proposition}

\begin{proof}
A relative corner peeling leaves ample families, hence isometric
families, proving \(1\Rightarrow2\).
Conversely, reverse a sequence in item~2, starting at the ample
family \(A\).  At each addition the enlarged family is isometric
in the cube.  The metric extension criterion in
Lemma~\ref{lem:damp-one-concept-extension-test} makes it ample
and makes the added concept a corner.  Reversing these additions
proves \(2\Rightarrow1\).  This argument uses ampleness of
\(A\), but no peeling order for \(A\).

We prove the relative version of the original-graph criterion
\cite[Corollary~2]{CoudertDucoffeNisseSoto18}.
Assume item~2 and take \(u,v\) with \(C_H(u,v)\cap A=\varnothing\).
The pair cannot lie wholly in \(A\), since \(A\) is isometric.
Immediately before the first of its endpoints is deleted, an
isometric residual graph still contains a common neighbor \(w\).
This neighbor is outside \(A\) and is deleted later than that
endpoint, giving item~3.

Conversely, suppose item~3 and consider a first deletion \(x_j\)
that fails to preserve distances.  Its preceding graph is isometric
in \(H\).  There must be two nonadjacent neighbors \(u,v\)
of \(x_j\) with no common neighbor after the deletion: otherwise
each two-edge segment through \(x_j\) on a shortest path could
be replaced by a surviving two-edge segment, preserving every
distance.  The pair has distance two in \(H\), all of
\(C_H(u,v)\) has been deleted by step \(j\), and both endpoints
remain.  Neither alternative in item~3 can hold, a contradiction.
This proves \(3\Rightarrow2\).

Finally, concatenate a relative peeling with a peeling of \(A\).
Conversely, the two portions of a complete peeling with final block
\(A\) give exactly those two orders.
\end{proof}

Both target qualifications matter.  In \(Q_3\), the deletion list
\(000,111\) leaves isometric graphs and satisfies item~3 for
\(A=Q_3\setminus\{000,111\}\), but its second deletion is not
a corner deletion.  The target is an isometric six-cycle, which is
not ample.  In the opposite direction, take the cornerless ample
\(291\)-concept seed \(G\) and its resolving extension
\(G\cup\{91\}\) from Remark~\ref{rem:damp-noncorner-rooted-factor}.
Deleting \(91\) is a relative peeling to \(G\), although no
complete peeling can have \(G\) as its final block.

The checker \path{verify_damp_relative_metric_peeling.py} verifies
all \(504\) one-concept extension candidates over the \(127\)
nonempty ample families in \(Q_3\), all \(720\) orders from
\(Q_3\) to a fixed edge, and both target qualifications above.
Exactly \(344\) of the extensions and \(176\) of the relative
orders pass their equivalent tests.  Its report is
\path{damp_relative_metric_peeling_verification.json}.
This gives an exact way to retain an ample core while expressing
the exterior constraints in the original graph.  Existence of the
deletion list is still an ordering predicate, so it is not intrinsic
admissibility for the generative classification.


The canonical overlap is a particularly restrictive choice of retained
target.  Its two wings cannot repair one another while the overlap is
kept fixed.  The next proposition makes this statement at the level of the
original metric constraints, including partial deletion lists.

\begin{proposition}[Metric separation across an ample overlap]
\label{prop:damp-overlap-metric-separation}
Let \(H\subseteq Q_{F\cup\{e\}}\) be ample, with \(e\) active, and put
\[
 B_b=\rho_e^bH,\qquad A=B_0\cap B_1,\qquad
 C=B_0\cup B_1,\qquad D_b=B_b\setminus A.
\]
These four families \(B_0,B_1,A,C\) are nonempty and ample.  For
\(X\in\{B_0,B_1,C\}\), let \(C_X(u,v)\) be the common-neighbor set
of a distance-two pair in \(X\), and define
\[
 \mathcal T_A(X)=
 \{(\{u,v\},C_X(u,v)):
       d_{G_X}(u,v)=2,\ C_X(u,v)\cap A=\varnothing\}.
\]
Then
\begin{equation}
 \mathcal T_A(C)
 =\mathcal T_A(B_0)\mathbin{\dot\cup}\mathcal T_A(B_1).
 \label{eq:damp-overlap-relative-constraint-splitting}
\end{equation}
All vertices outside \(A\) in a member of \(\mathcal T_A(B_b)\)
belong to \(D_b\).  Thus these are independent constraint systems on
the two wings, with their vertices in \(A\) retained permanently.

More precisely, let \(\sigma\) be any list of distinct vertices of
\(C\setminus A\), and let \(\sigma_b\) be its subsequence in \(D_b\).
Every successive deletion in \(\sigma\) preserves distances in \(C\)
if and only if every successive deletion in each \(\sigma_b\) preserves
distances in \(B_b\).  In particular, when \(\sigma\) lists all of
\(C\setminus A\), it is a relative corner peeling to \(A\) if and
only if both \(\sigma_b\) are relative corner peelings of \(B_b\) to
\(A\).  Every interleaving of two such wing orders is valid.

There is also an exact local description of new midpoints under
contraction.  Suppose \(u,v\in B_b\) have distance two and
\(C_C(u,v)\setminus C_{B_b}(u,v)\ne\varnothing\).  Then there are
vertices \(w,z\) such that
\begin{equation}
 C_{B_b}(u,v)=\{w\},\qquad C_C(u,v)=\{w,z\},\qquad
 \{u,v,w\}\subseteq A,\quad z\in D_{1-b}.
 \label{eq:damp-contraction-new-midpoint-overlap}
\end{equation}
The four vertices \(Q=\{u,v,w,z\}\) form a square, and the corresponding
three-dimensional face of \(H\) is exactly
\begin{equation}
 H\cap(Q\times\{0,1\})
 =(Q\times\{0,1\})\setminus\{(z,b)\}.
 \label{eq:damp-contraction-new-midpoint-punctured-cube}
\end{equation}
Conversely, each such punctured three-cube supplies the displayed new
midpoint.  Its old distance-two constraint already has a common neighbor
in the retained overlap.
\end{proposition}

\begin{proof}
The ampleness assertions follow from
Proposition~\ref{prop:damp-yang-two-cone-gluing-dichotomy}; activity of
\(e\) gives an \(e\)-edge and hence \(A\ne\varnothing\).
No edge of \(C\) joins \(D_0\) to \(D_1\).  Otherwise the exclusive
lifts of its endpoints would have Hamming distance two in \(H\), with
both possible intermediate vertices absent, contradicting isometry.

If \(u\in D_0\) and \(v\in D_1\) have distance two in \(C\), every
common neighbor belongs to \(A\): a neighbor outside \(A\) would give
an edge between the wings.  If both endpoints belong to \(A\), its
isometry supplies a common neighbor in \(A\) as well.  Every remaining
constraint therefore has an endpoint in exactly one wing, say \(D_b\),
and its other endpoint and all common neighbors lie in \(B_b\).
Its common-neighbor set is unchanged on passing to \(B_b\).
Conversely, a constraint of \(\mathcal T_A(B_b)\) has an endpoint in
\(D_b\), since \(A\) is isometric, and cannot acquire a common
neighbor in the other wing.  This proves
\eqref{eq:damp-overlap-relative-constraint-splitting}, including
disjointness.

A sequence of deletions from an isometric graph preserves distances at
every step exactly when no original distance-two pair loses all its
common neighbors while both endpoints remain.  Necessity follows by
looking at that pair; sufficiency follows at the first failed deletion
by rerouting every two-edge passage through its deleted vertex, as in
Proposition~\ref{prop:damp-relative-metric-peeling}.  While \(A\) is
retained, every pair with a common neighbor in \(A\) satisfies this
condition permanently.  The other pairs split by
\eqref{eq:damp-overlap-relative-constraint-splitting}.  Deletions in
one wing do not affect any vertex of a constraint in the other wing.
This proves the assertion for partial lists.  For full lists,
Proposition~\ref{prop:damp-relative-metric-peeling} identifies the
three metric conditions with the respective relative corner peelings.

Finally, a cube distance-two pair has just two possible common neighbors.
The side \(B_b\) has at least one, so a new common neighbor means
\(C_{B_b}(u,v)=\{w\}\) and
\(C_C(u,v)=\{w,z\}\), with \(z\in D_{1-b}\).
Neither \(u\) nor \(v\) can lie in \(D_b\), because each is adjacent
to \(z\); thus both lie in \(A\).  Isometry of \(A\) then forces
\(w\in A\).  The vertices \(u,v,w\) have both lifts, and \(z\) has
only its \((1-b)\)-lift.  These are exactly the seven vertices in
\eqref{eq:damp-contraction-new-midpoint-punctured-cube}.
Reading that face in the opposite direction proves the converse.
\end{proof}

\begin{corollary}[Contraction peelings must enter the overlap]
\label{cor:damp-contraction-must-enter-overlap}
In the notation of
Proposition~\ref{prop:damp-overlap-metric-separation}, suppose that
\(H\in\Obs(\Damp)\cap\Ample\).  The contraction \(C\) is peelable,
but no complete corner peeling of \(C\) has \(A\) as its final block.
If \(A\in\Damp\), the first deleted vertex of \(A\) occurs while
both wings \(D_0,D_1\) still have vertices remaining.
If \(H\) is cornerless, then
\begin{equation}
 \operatorname{Cor}(C)\subseteq A\setminus\operatorname{Cor}(A).
 \label{eq:damp-contraction-corners-destroy-overlap}
\end{equation}
Consequently every peeling of \(C\) enters \(A\) at its first
deletion, and that deletion makes the remaining overlap nonample.
\end{corollary}

\begin{proof}
The contraction is a proper pc-minor and hence belongs to \(\Damp\).
If \(A\notin\Damp\), it cannot be the final block of a complete
peeling.  If \(A\in\Damp\), the gluing-critical alternative of
Proposition~\ref{prop:damp-yang-two-cone-gluing-dichotomy} says that
neither \(B_b\) peels relatively to \(A\).  Suppose one wing were
exhausted before the first deletion in \(A\).  Projecting that prefix
to the exhausted wing, the partial-list assertion of
Proposition~\ref{prop:damp-overlap-metric-separation} gives a
distance-preserving deletion list from \(B_b\) to \(A\).
Proposition~\ref{prop:damp-relative-metric-peeling} makes it a relative
corner peeling, a contradiction.  This proves both assertions in this
case.

For the last assertion, every cube \(Q\subseteq C\) lies wholly in
one side \(B_b\).  Indeed, \(A\cap Q\) is ample, so its complement
in \(Q\) is either empty or connected.  In the nonempty case the
absence of edges between the wings puts that complement in one \(D_b\),
and then \(Q\subseteq B_b\); in the empty case both sides contain
\(Q\).  If a corner \(x\) of \(C\) lay in \(D_b\), its unique
maximal cube would therefore lie in \(B_b\).  Its lift has no
\(e\)-neighbor, and all its incident cubes in \(H\) lie in that
lifted cube.  It would be a corner of \(H\), a contradiction.
Now let \(x\in A\) be a corner of \(C\), and let its unique maximal
cube \(Q\) lie in \(B_b\).  Every cube of \(B_{1-b}\) through \(x\)
lies in \(Q\) and hence in \(A\).  Suppose \(x\) were a corner of
\(A\), with maximal cube \(R\).  Then all those cubes lie in \(R\),
and \(R\) has both lifts.  The cube \(R\times\{0,1\}\) would contain
every cube of \(H\) incident with \((x,1-b)\), making that lift a
corner.  This again contradicts cornerlessness of \(H\), proving
\eqref{eq:damp-contraction-corners-destroy-overlap}.  The ample
corner-deletion criterion makes \(A\setminus\{x\}\) nonample.
\end{proof}

The checker \path{verify_damp_overlap_metric_separation.py} verifies the
constraint splitting and midpoint faces on all active splits in \(Q_3\),
including partial and complete deletion lists.  It also replays the
\(72\) elementary-minor orders of the fixed \(291\)- and
\(292\)-concept families from
Theorem~\ref{thm:damp-empty-core-counterexample}, and verifies the forced
overlap entry at all \(24\) coordinate splits.  For the cornerless source
it also checks that every contraction corner is a noncorner of its overlap.
Its report is
\path{damp_overlap_metric_separation_verification.json}.
The structural consequence is a restriction on the next construction:
retaining the canonical overlap cannot produce the required contraction
peelings.  Their attachment rules must instead justify deletions inside
that overlap, together with their effect on the two remaining wings.
For a cornerless source the overlap already ceases to be ample at the
first such step, so keeping an ample overlap throughout is also excluded.

The loss of overlap ampleness has an exact description in terms of
shattered sets.  On every ample contraction residual, the two section
shadows intersect in precisely the overlap shadow.  Consequently, all
failure of ampleness in the lifted family is measured by the overlap.

\begin{proposition}[Overlap and section deficits]
\label{prop:damp-contraction-section-deficits}
Retain the notation \(H,B_0,B_1,A,C\) of
Proposition~\ref{prop:damp-overlap-metric-separation}.  For any cube family
\(X\), including the empty family, write
\[
 \delta(X)=|\operatorname{Sh}(X)|-|X|.
\]
For an ample subfamily \(K\subseteq C\), put
\[
 B_b^K=B_b\cap K,\qquad A^K=A\cap K,\qquad
 H[K]=\{(v,b)\in H:v\in K\},
\]
where \(K\) may be empty.  The section residuals and their overlap
need not be ample.  For every \(J\subseteq F\), their signed traces satisfy
\begin{equation}
 (B_0^K)|_J\cap(B_1^K)|_J=(A^K)|_J.
 \label{eq:damp-residual-section-trace-intersection}
\end{equation}
Moreover,
\begin{align}
 \operatorname{Sh}(K)
   &=\operatorname{Sh}(B_0^K)\cup\operatorname{Sh}(B_1^K),
       \label{eq:damp-residual-section-shadow-union}\\
 \operatorname{Sh}(A^K)
   &=\operatorname{Sh}(B_0^K)\cap\operatorname{Sh}(B_1^K),
       \label{eq:damp-residual-section-shadow-intersection}\\
 \delta(H[K])
   &=\delta(B_0^K)+\delta(B_1^K)
     =\delta(A^K).
       \label{eq:damp-residual-section-deficit-balance}
\end{align}
In particular, \(H[K]\) is ample if and only if \(A^K\) is ample.

Let \(x\) be a corner of \(K\), put \(K'=K\setminus\{x\}\), and
let \(m(x)=|\{b:x\in B_b\}|\).  Define
\begin{align*}
 L_b(x;K)&=
 \operatorname{Sh}(B_b^K)\setminus\operatorname{Sh}(B_b^{K'}),\\
 L_A(x;K)&=
 \operatorname{Sh}(A^K)\setminus\operatorname{Sh}(A^{K'}).
\end{align*}
Then
\begin{align}
 \delta(H[K'])-\delta(H[K])
 &=m(x)-|L_0(x;K)|-|L_1(x;K)|\notag\\
 &=
 \begin{cases}
  0,&x\notin A,\\
  1-|L_A(x;K)|,&x\in A.
 \end{cases}
 \label{eq:damp-residual-section-deficit-transition}
\end{align}
An exterior deletion preserves both side deficits individually.  An
overlap deletion increases the deficit by at most one, and decreases it
exactly when it loses at least two shattered supports of the overlap.
\end{proposition}

\begin{proof}
Suppose a trace on \(J\) is realized by \(u\in B_0^K\) and
\(v\in B_1^K\).  If either concept lies in \(A\), the trace is
already realized in \(A^K\).  Otherwise, \(u,v\) lie in opposite
wings.  Since \(K\) is ample, it is isometric; a shortest path in
\(K\) from \(u\) to \(v\) flips no coordinate on which they agree,
so every concept of the path has the same \(J\)-trace.  There are no
edges between the wings, by
Proposition~\ref{prop:damp-overlap-metric-separation}.  The path must
therefore pass through \(A^K\), proving the nontrivial inclusion in
\eqref{eq:damp-residual-section-trace-intersection}.  The reverse
inclusion is immediate, and the assertion is also valid when \(K\)
is empty.  A support \(J\) shattered in both sections has all its
traces in their intersection, so it is shattered in \(A^K\).  This
proves \eqref{eq:damp-residual-section-shadow-intersection}.

Every cube of \(C\) lies wholly in one side.  To see this, intersect
such a cube \(Q\) with the ample family \(A\).  The complement of
\(A\cap Q\) in \(Q\) is ample and hence, if nonempty, connected.
There are no edges between the two wings, so this complement lies in
one wing and \(Q\) lies in that side.  If the complement is empty,
both sides contain \(Q\).

Since \(K\) is ample, every support shattered by \(K\) is realized
by a cube in \(K\).  That cube lies in one \(B_b^K\), proving one
inclusion in \eqref{eq:damp-residual-section-shadow-union}; the reverse
inclusion follows from \(B_b^K\subseteq K\).  This includes \(K=\varnothing\).
For arbitrary two-layer families the shadow formula is
\[
 \operatorname{Sh}(H[K])
 =\operatorname{Sh}(K)\mathbin{\dot\cup}
   \{J\cup\{e\}:J\in
     \operatorname{Sh}(B_0^K)\cap\operatorname{Sh}(B_1^K)\}.
\]
Use \(|H[K]|=|K|+|A^K|\) and
\(|\operatorname{Sh}(K)|=|K|\), together with
\eqref{eq:damp-residual-section-shadow-intersection}, to obtain
\(\delta(H[K])=\delta(A^K)\).
Inclusion--exclusion in
\eqref{eq:damp-residual-section-shadow-union} also gives
\(\delta(B_0^K)+\delta(B_1^K)=\delta(A^K)\), proving
\eqref{eq:damp-residual-section-deficit-balance}.
The ampleness equivalence follows because a family is ample exactly
when its deficit is zero.

Deleting \(x\) removes \(m(x)\) concepts from the two sections and
loses exactly \(|L_0(x;K)|+|L_1(x;K)|\) section supports.  Taking
the difference of the first equality in
\eqref{eq:damp-residual-section-deficit-balance} gives the equality
in \eqref{eq:damp-residual-section-deficit-transition}.
Both \(K\) and \(K'\) are ample, so the same balance identity
identifies this difference with \(\delta(A^{K'})-\delta(A^K)\).
If \(x\notin A\), the overlap is unchanged, and so is one of the
two sections.  The sum of the side deficits is fixed; hence each side
deficit is fixed.  If \(x\in A\), its deletion removes one overlap
concept and exactly \(|L_A(x;K)|\) overlap supports.  This proves the
second equality and all the final assertions.
\end{proof}

\begin{proposition}[The first nonample overlap has one-sided deficit]
\label{prop:damp-first-contraction-deficit}
In the same notation, let
\(x\in\operatorname{Cor}(C)\cap
 (A\setminus\operatorname{Cor}(A))\), and let \(Q\) be the unique
maximal cube of \(C\) through \(x\), with support \(I\).
There is a unique \(b\in\{0,1\}\) for which \(Q\subseteq B_b\).
The deletion of \(x\) then has the following effects:
\begin{align}
 \operatorname{Sh}(B_b\setminus\{x\})
     &=\operatorname{Sh}(B_b)\setminus\{I\},
       \label{eq:damp-first-deficit-owner-loss}\\
 \operatorname{Sh}(B_{1-b}\setminus\{x\})
     &=\operatorname{Sh}(B_{1-b}),\qquad
 \operatorname{Sh}(A\setminus\{x\})=\operatorname{Sh}(A),
       \label{eq:damp-first-deficit-unchanged-shadows}\\
 \operatorname{Sh}(H[C\setminus\{x\}])
     &=\operatorname{Sh}(H)\setminus\{I\}.
       \label{eq:damp-first-deficit-lift-shadow}
\end{align}
In particular, \(B_b\setminus\{x\}\) is ample, whereas
\[
 \delta(B_{1-b}\setminus\{x\})
 =\delta(A\setminus\{x\})
 =\delta(H[C\setminus\{x\}])=1.
\]

The nonample section and overlap have the same conditioning-minimal
nonample faces.  Define
\[
 \mathcal T_x=
 \{T\subseteq F:|T|\ge2,\ x^T\notin A,
                       \ x^U\in A\text{ for every }U\subsetneq T\}.
\]
This family is nonempty.  For every \(T\in\mathcal T_x\), let
\(Q_T(x)=\{x^U:U\subseteq T\}\).  Then
\begin{equation}
 (A\setminus\{x\})\cap Q_T(x)
 =(B_{1-b}\setminus\{x\})\cap Q_T(x)
 =Q_T(x)\setminus\{x,x^T\},
 \qquad x^T\in B_b\setminus A.
 \label{eq:damp-first-deficit-positioned-copairs}
\end{equation}
These are exactly the conditioning-minimal nonample faces of each
displayed residual family.  Thus their two missing vertices are the
deleted overlap concept and a surviving concept of the opposite wing.
\end{proposition}

\begin{proof}
The cube containment proved above puts \(Q\) in at least one side.
It cannot lie in both, because then it would be the unique maximal cube
of \(A\) through \(x\), making \(x\) an \(A\)-corner.
Let \(b\) be its unique side.  All cubes of \(B_b\) through \(x\)
lie in \(Q\), so \(x\) is a corner of \(B_b\), with support \(I\).
Every cube of \(B_{1-b}\) through \(x\) also lies in \(Q\) and hence
in \(A\).  Therefore the local cube-support complexes of \(A\) and
\(B_{1-b}\) at \(x\) coincide, and \(x\) is a noncorner of both.

A corner deletion from an ample family loses its unique cube support,
by Lemma~\ref{lem:damp-one-concept-extension-test} applied in reverse.
A noncorner deletion loses no shattered support: losing any support would,
by the upper Sandwich inequality, make the remaining family ample and
hence make the deleted concept a corner.  This proves
\eqref{eq:damp-first-deficit-owner-loss} and
\eqref{eq:damp-first-deficit-unchanged-shadows}.
Since \(C\setminus\{x\}\) loses \(I\) while the other side loses
nothing, \eqref{eq:damp-residual-section-shadow-union} implies
\(I\notin\operatorname{Sh}(B_{1-b})\).
The common section shadow consequently remains
\(\operatorname{Sh}(A)\).  The two-layer shadow formula now proves
\eqref{eq:damp-first-deficit-lift-shadow}, and counting the deleted
concepts gives all the deficit identities.

The supports in \(\mathcal T_x\) are precisely the minimal nonfaces
of order at least two in the common local cube-support complex at \(x\).
There is at least one such support because \(x\) is a noncorner.
Every such support is
contained in \(I\); hence its opposite vertex \(x^T\) lies in
\(Q\subseteq B_b\).  It is absent from both \(A\) and \(B_{1-b}\),
which proves \eqref{eq:damp-first-deficit-positioned-copairs}.

For completeness, let \(M\) be either original ample family \(A\)
or \(B_{1-b}\).  A nonample conditioning of \(M\setminus\{x\}\)
must come from a face containing \(x\), since other faces are unchanged
conditionings of \(M\).  In that face, deleting \(x\) from an ample
family fails exactly when its local cube-support complex at \(x\)
has a minimal nonface of order at least two.  Such a support is some
\(T\in\mathcal T_x\) and supplies the smaller face in
\eqref{eq:damp-first-deficit-positioned-copairs}.  Minimality of the
nonample conditioning forces equality with that face.  Conversely,
removing two antipodes from a cube of dimension at least two is nonample,
and every proper nonempty conditioning is a punctured or full cube,
hence ample.  This proves the exact minimal-face assertion.
\end{proof}

\begin{remark}[Completing the corner cube need not repair the overlap]
\label{rem:damp-corner-cube-completion-failure}
The preceding local description does not imply that adjoining all of
\(Q\) to a lower ample family preserves ampleness.  Here is an explicit
failure even with a marked noncorner of the lower family.  Encode the
vertices of \(Q_4\) by \(v=\sum_{i=0}^3v_i2^i\), and put
\[
 C_0=\{0,1,2,3,6,7,8,9\},\quad
 A_0=\{0,2,6,7,8\},\quad Q_0=\{0,1,8,9\}.
\]
The graph of \(C_0\) is a four-vertex path times an edge, and \(A_0\)
is the geodesic path \(8,0,2,6,7\), so both are ample.
The union \(A_0\cup Q_0=C_0\setminus\{3\}\) is nonample:
fixing coordinates \(0\) and \(3\) to \(1\) and \(0\), respectively,
leaves the antipodal pair \(\{1,7\}\) on coordinates \(1,2\).
Now put
\[
 C=C_0\times Q_2,\qquad
 A=A_0\times\{00,10,01\},\qquad x=(8,00).
\]
The families \(C,A\) are ample, of orders \(32,15\).  The vertex
\(x\) is a corner of \(C\) and a noncorner of \(A\), and its corner
cube in \(C\) is \(Q=Q_0\times Q_2\).  Nevertheless, the \(00\)
section of \(A\cup Q\) is the nonample family \(A_0\cup Q_0\), so
\(A\cup Q\) is nonample.  Its order is \(25\), and its shattered
shadow has order \(28\).  This refutes the unconditional local
cube-completion step; no forbidden-minor assertion is needed for this
control.
\end{remark}

The checker \path{verify_damp_contraction_section_deficits.py} verifies
the signed trace and shadow identities on all ample residuals of the
active \(Q_3\) splits, and the deficit laws on their corner transitions.
It checks all \(24\) first contraction
corners of the fixed cornerless \(291\)-concept family, including their
positioned antipodal pairs, and a full saved \(224\)-step contraction
peeling.  In that particular order the lifted deficit reaches \(14\)
before returning to zero; this is not a claim about the minimum possible
peak over all orders.  The report
\path{damp_contraction_section_deficits_verification.json} also verifies
the cube-completion counterexample.  The deficit identities record losses
exactly, but their scalar values do not determine which losses are
compatible.  An exterior deletion can change which overlap concepts
are corners of the contraction, but it cannot lower either side deficit.
Intrinsic construction still requires compatible overlap losses to be
derived from the positioned supports and the two section shadows.
\begin{theorem}[Metric Yang obstruction class]
\label{thm:damp-yang-metric-obstruction-class}
Let \(H\subseteq Q_F\) be ample, put \(m=|H|\ge2\), and let

\[
 \mathscr O(H)
 =\operatorname{Ord}\bigl(2^H\setminus\{\varnothing,H\}\bigr)
\]

be the order sphere whose chambers are the orders of \(H\).  Let

\[
 \mathcal D_2(H)=
 \bigl\{\{u,v\}\subseteq H:d_{G_H}(u,v)=2\bigr\},
 \qquad
 C_H(u,v)=N_{G_H}(u)\cap N_{G_H}(v).
\]

For \(D=\{u,v\}\in\mathcal D_2(H)\), let \(P_D\) be the poset on \(H\)
with generating relations

\begin{equation}
 w\mathrel{\triangleleft_D}u,
 \qquad
 w\mathrel{\triangleleft_D}v
 \qquad(w\in C_H(u,v)),
 \label{eq:damp-yang-metric-failure-poset}
\end{equation}

and define

\[
 \mathscr R_D^{\rm met}
 =
 \operatorname{Ord}\bigl(
   \mathcal J(P_D)\setminus\{\varnothing,H\}
 \bigr),
 \qquad
 \mathscr U_{\rm met}(H)
 =
 \bigcup_{D\in\mathcal D_2(H)}\mathscr R_D^{\rm met}.
\]

Then

\begin{equation}
 H\text{ is nonpeelable}
 \quad\Longleftrightarrow\quad
 \mathscr U_{\rm met}(H)=\mathscr O(H)
 \quad\Longleftrightarrow\quad
 \widetilde H_{m-2}
 \bigl(\mathscr U_{\rm met}(H);\mathbb F_2\bigr)
 \cong\mathbb F_2.
 \label{eq:damp-yang-metric-obstruction-class}
\end{equation}

Consequently

\[
 \omega_{\rm met}(H)
 =
 \dim_{\mathbb F_2}
 \widetilde H_{m-2}
 \bigl(\mathscr U_{\rm met}(H);\mathbb F_2\bigr)
 \in\{0,1\}
\]

is an intrinsic global invariant of the metric graph \(G_H\), equal to one
exactly when the Yang complex is nonshellable.
\end{theorem}

\begin{proof}
A static characterization of a distance-preserving elimination order
\(\prec\) says that, for every pair \(u,v\) at distance two, there is a
common neighbor \(w\in C_H(u,v)\) such that

\[
 u\prec w
 \qquad\text{or}\qquad
 v\prec w;
\]

see \cite[Corollary~2]{CoudertDucoffeNisseSoto18}.  An order fails this
condition at \(D=\{u,v\}\) exactly when every common neighbor precedes both
\(u\) and \(v\).  These are precisely the linear extensions of the poset
\(P_D\), hence precisely the chambers of \(\mathscr R_D^{\rm met}\).
Proposition~\ref{prop:damp-yang-shellability-metric-elimination} therefore
gives the first equivalence in
\eqref{eq:damp-yang-metric-obstruction-class}.

The chamber-adjacency graph of the order sphere is connected.  As in the
proof of Theorem~\ref{thm:damp-yang-order-sphere-obstruction}, a proper
subcomplex of this sphere has zero top-dimensional homology over
\(\mathbb F_2\), whereas the full sphere has a one-dimensional fundamental
class.  This proves the second equivalence and the assertion about
\(\omega_{\rm met}\).
\end{proof}

\begin{remark}[What the metric class does and does not solve]
\label{rem:damp-yang-metric-obstruction-scope}
The invariant \(\omega_{\rm met}\) is substantially more intrinsic than the
latent-chart recognition theorem: it uses only distance-two pairs and their
common-neighbor sets in the concept graph.  It contains no corners,
punctured faces, shelling orders, linear quotients, or descendant tests in
its definition.  It is nevertheless still an exact recognition invariant,
not the requested generative parametrization.  Computing whether its
fundamental class survives is equivalent to deciding shellability.

The useful new target is now precise.  One must describe the minimal
supports of this metric fundamental class and prove how those supports are
destroyed by restriction and contraction.  A successful construction would
recover an obstruction from finite distance-two interval data with local
boundary and transport axioms.  Merely requiring
\(\omega_{\rm met}(H)=1\) and its vanishing on proper pc-minors would again
be circular and is not an admissibility condition.
\end{remark}

\begin{theorem}[Boolean-prefix cutset and exact pc-minor transport]
\label{thm:damp-boolean-prefix-cutset}
Let \(H\subseteq Q_F\) be ample.  For each
\(D=\{u,v\}\in\mathcal D_2(H)\), define the Boolean interval
\[
 \mathcal I_D(H)
 =
 \bigl[C_H(u,v),\,H\setminus\{u,v\}\bigr]
 =
 \bigl\{P\subseteq H:
 C_H(u,v)\subseteq P,\ P\cap\{u,v\}=\varnothing\bigr\},
\]
and put
\[
 \mathcal U_{\rm pre}(H)
 =\bigcup_{D\in\mathcal D_2(H)}\mathcal I_D(H)
 \subseteq 2^H.
\]
Then \(H\) is nonpeelable if and only if
\(\mathcal U_{\rm pre}(H)\) meets every maximal chain of the Boolean
lattice \(2^H\).

Moreover, this interval system has the following exact pc-minor transport
laws.
\begin{enumerate}
 \item If \(A=H_{e=b}\) is a coordinate restriction and
       \(D\in\mathcal D_2(H)\) has \(D\subseteq A\), then
       \[
        \mathcal I_D(A)=\mathcal I_D(H)\cap 2^A.
       \]
       Thus the bad-prefix system of \(A\) is obtained by retaining the
       intervals whose owners lie in \(A\) and restricting them to the
       Boolean sublattice \(2^A\subseteq2^H\).
 \item Let \(q\colon H\to\overline H=\pi_eH\) be a coordinate
       contraction.  Pullback identifies \(2^{\overline H}\) with the
       sublattice
       \[
        \operatorname{Sat}_q(H)
        =\{q^{-1}(P):P\subseteq\overline H\}\subseteq2^H
       \]
       of unions of \(q\)-fibers.  Under this identification, the interval
       owned by
       \(\overline D=\{\overline u,\overline v\}
         \in\mathcal D_2(\overline H)\)
       becomes
       \[
        \bigl[
          q^{-1}C_{\overline H}(\overline u,\overline v),
          H\setminus q^{-1}\{\overline u,\overline v\}
        \bigr]_{\operatorname{Sat}_q(H)}.
       \]
       Hence contraction is represented without choosing an order: pass
       to the fiber-saturated Boolean sublattice and rebuild the finitely
       many owner intervals from the quotient metric.
\end{enumerate}
\end{theorem}

\begin{proof}
Write a total order of \(H\) as \(x_1,\ldots,x_m\), and let
\(P_i=\{x_1,\ldots,x_i\}\).  The order fails the distance-two condition at
\(D=\{u,v\}\) precisely when every member of \(C_H(u,v)\) precedes both
\(u\) and \(v\).  If this happens, the prefix ending with the last member
of \(C_H(u,v)\) belongs to \(\mathcal I_D(H)\).  Conversely, if some
prefix belongs to \(\mathcal I_D(H)\), then every common neighbor precedes
both endpoints, so the order fails at \(D\).  Orders of \(H\) are in
bijection with maximal chains of \(2^H\).  The static metric criterion used
in Theorem~\ref{thm:damp-yang-metric-obstruction-class} now proves the
cutset assertion.

For a restriction, the common-neighbor identity in
Proposition~\ref{prop:damp-metric-support-coordinate-span} gives
\(C_A(u,v)=C_H(u,v)\); intersecting the two endpoints of the Boolean
interval with \(2^A\) gives the displayed formula.  For a contraction,
pullback is a lattice isomorphism from \(2^{\overline H}\) onto
\(\operatorname{Sat}_q(H)\).  Applying this isomorphism to the lower and
upper endpoints of each quotient interval gives the second formula.
\end{proof}

\begin{proposition}[Unique-midpoint core and the square remainder]
\label{prop:damp-unique-midpoint-core}
Let \(V\) be finite and let \(\mathcal T\) be a family of rooted-pair
constraints \((w;\{u,v\})\) on three distinct members of \(V\).  Call a
nonempty set \(W\subseteq V\) a \emph{\(\mathcal T\)-core} if, for every
\(w\in W\), there is a constraint
\((w;\{u,v\})\in\mathcal T\) with \(u,v\in W\).  There is an order
\(\prec\) of \(V\) satisfying
\[
 u\prec w\quad\hbox{or}\quad v\prec w
 \qquad\text{for every }(w;\{u,v\})\in\mathcal T
 \tag{\(\dagger\)}
\]
if and only if no \(\mathcal T\)-core exists.

This criterion is constructive and order-free.  Starting with \(W=V\),
repeatedly delete any \(w\in W\) for which no constraint rooted at \(w\)
has both heads in \(W\).  The final set is the unique largest
\(\mathcal T\)-core; when it is empty, the deletion order satisfies
\((\dagger)\).

Apply this to
\[
 \mathcal T_1(H)
 =\bigl\{(w;\{u,v\}):
        \{u,v\}\in\mathcal D_2(H),\ C_H(u,v)=\{w\}\bigr\}.
\]
If the \(\mathcal T_1(H)\)-core is nonempty, then \(H\) is nonpeelable.
If it is empty, the unique-midpoint constraints admit an order, and the
only remaining metric conditions come from full squares.  More precisely,
a total order of \(H\) is a distance-preserving elimination order if and
only if it satisfies \((\dagger)\) for \(\mathcal T_1(H)\) and, in every
square of \(G_H\),
the first two vertices are adjacent.

The square conditions alone never obstruct elimination.  For every cube
vertex \(z\in Q_F\), any total order of \(H\) that is nondecreasing in
\(d_{Q_F}(z,\cdot)\) has adjacent first two vertices in every square.
Consequently, if \(\mathcal T_1(H)=\varnothing\), then \(H\) is peelable.
\end{proposition}

\begin{proof}
Suppose first that \(W\) is a \(\mathcal T\)-core, and take the first
vertex \(w\) of \(W\) in an arbitrary total order.  Some constraint
\((w;\{u,v\})\) has both heads in \(W\).  Both heads follow \(w\), so this
constraint violates \((\dagger)\).

Conversely, perform the stated deletion procedure.  A union of
\(\mathcal T\)-cores is again a \(\mathcal T\)-core, so the set left at
termination is the unique largest one.  If the final set is empty, list
the vertices in the order in which they were deleted.  When \(w\) was
deleted, every constraint rooted at \(w\) had a head already deleted.
That head precedes \(w\) in the deletion order, proving \((\dagger)\).

For a unique common neighbor \(w\), the metric condition for
\(D=\{u,v\}\) is exactly \((\dagger)\).  On a square, the two diagonals
give the two reciprocal two-root constraints.  One of them fails exactly
when the first two vertices of the square are the opposite diagonal that
forms its common-neighbor set.  Thus both diagonal constraints hold
exactly when the first two square vertices are adjacent.  Item~1 of
Proposition~\ref{prop:damp-metric-support-coordinate-span} says that these
are all possible distance-two constraints.

Finally, fix \(z\in Q_F\).  On any square there is a unique vertex nearest
to \(z\), the two vertices at the next distance are its two square
neighbors, and the remaining vertex is farthest from \(z\).  Hence the
first two square vertices in every distance-nondecreasing order are
adjacent.  If \(\mathcal T_1(H)\) is empty, such an order satisfies all
metric constraints.  The static metric criterion completes the proof.
\end{proof}

\begin{proposition}[A one-sided criterion for empty unique-midpoint core]
\label{prop:damp-one-sided-empty-core}
For nested ample families \(A\subseteq B\), let \(\kappa_A(B)\)
be the largest set \(W\subseteq B\setminus A\) such that every
\(w\in W\) roots a constraint
\[
 (w;\{u,v\})\in\mathcal T_1(B)
 \quad\text{with}\quad u,v\in A\cup W.
\]
The empty set is allowed.  Equivalently, start with \(W=B\setminus A\)
and repeatedly delete a root having no such witness, keeping all of
\(A\) present and keeping \(\mathcal T_1(B)\) fixed.

Let \(H\subseteq Q_F\) be ample, fix an active coordinate \(e\),
and write \(B_i=\rho_e^iH\), \(A=B_0\cap B_1\), and
\(D_i=B_i\setminus A\).  Suppose that
\(\kappa(B_0)=\kappa(B_1)=\varnothing\), where \(\kappa\) denotes
the ordinary unique-midpoint core.  If
\(\kappa_A(B_i)=\varnothing\) for either sign \(i\), then
\(\kappa(H)=\varnothing\).

In particular, the conclusion holds if either \(B_i\) can be reduced
to \(A\) by deleting corners outside \(A\).  Equivalently, it holds
if either inclusion \(A\subseteq B_i\) admits an \(A\)-relative
ample chain as in \eqref{eq:damp-relative-ample-chain}.
\end{proposition}

\begin{proof}
Unions of sets satisfying the displayed witness condition again satisfy
it.  Every such set survives the stated pruning, and the set left at
termination satisfies the condition.  This proves existence, uniqueness,
and the pruning description of \(\kappa_A(B)\).

Suppose, after exchanging the signs if necessary, that
\(\kappa_A(B_0)=\varnothing\), and let \(K\) be any
unique-midpoint core of \(H\).  A word \((w,0)\) with \(w\in D_0\)
has no neighbor in direction \(e\).  Its witnessing constraint in
\(K\) must therefore be horizontal.  Its heads project into
\(A\cup\{v\in D_0:(v,0)\in K\}\).  Thus
\[
 \{w\in D_0:(w,0)\in K\}\subseteq\kappa_A(B_0)=\varnothing.
\]
Every remaining member of \(K\) in the zero section projects into
\(A\).  A unique-midpoint constraint rooted at \((w,0)\) and using
\(e\) would have heads \((w,1)\) and \((v,0)\), where \(v\) is an
old-coordinate neighbor of \(w\).  Uniqueness requires
\((v,1)\notin H\), so \(v\in D_0\), whose copy cannot belong to
\(K\).  Hence every constraint witnessing a member of the zero section
of \(K\) is horizontal.  That section would be a core of \(B_0\),
which is impossible.  The zero section of \(K\) is therefore empty.
Every witness for its one section is now horizontal as well, making that
section a core of \(B_1\).  Thus \(K\) is empty, as required.

Finally, suppose a corner-deletion sequence reduces \(B\) to \(A\),
and that \(W\) satisfies the witness condition defining
\(\kappa_A(B)\).  At the first deletion of a member of a nonempty
\(W\), both heads of its witness are still present.  Their other
possible midpoint was absent from \(B\), and is still absent, so the
root cannot be a corner.  This contradiction proves
\(\kappa_A(B)=\varnothing\).  Reversing the deletion sequence gives
the stated equivalence with a relative ample chain.
\end{proof}

\begin{corollary}[Nonempty cores force two large cyclic wings]
\label{cor:damp-nonempty-core-wing-obstructions}
Let \(H\in\Obs(\Damp)\cap\Ample\) be irredundantly embedded, and
suppose \(\kappa(H)\ne\varnothing\).  For every active coordinate
\(e\), with sections, overlap, and wings as in
Proposition~\ref{prop:damp-one-sided-empty-core}, both
\(\kappa_A(B_0)\) and \(\kappa_A(B_1)\) are nonempty.
Neither inclusion \(A\subseteq B_i\) admits an \(A\)-relative ample
chain.  Both induced wing graphs \(Q[D_i]\) contain cycles, and
\begin{equation}
 |D_0|,|D_1|\ge9,
 \qquad |H|\ge2|A|+18.
 \label{eq:damp-nonempty-core-wing-size}
\end{equation}
These conclusions hold whether or not \(A\) is corner-peelable.
If \(A\notin\Damp\), then additionally
\[
 |H|\ge2m_{\mathrm{amp}}+18.
\]
\end{corollary}

\begin{proof}
The two sections are proper pc-minors and hence corner-peelable, so their
unique-midpoint cores are empty.  Proposition
\ref{prop:damp-one-sided-empty-core} proves the first two assertions.
An edge of \(H\) in direction \(e\) projects into the overlap, so
\(A\) is nonempty.  Proposition~\ref{prop:damp-yang-two-cone-gluing-dichotomy}
makes it ample.  Each relative Yang module of
\(A\subseteq B_i\) is Cohen--Macaulay by
Lemma~\ref{lem:damp-relative-yang-cm-flag-h}.
The forest and eight-facet linear-quotient theorems for relative Yang
modules \cite{BousquetYangComplex26} would give an \(A\)-relative
ample chain if \(Q[D_i]\) were a forest or if \(|D_i|\le8\).
Both alternatives are excluded.  The identity
\(|H|=2|A|+|D_0|+|D_1|\) gives
\eqref{eq:damp-nonempty-core-wing-size}.
If \(A\) is nonpeelable,
Theorem~\ref{thm:damp-inherited-overlap-descent} makes it an ample
forbidden pc-minor, so \(|A|\ge m_{\mathrm{amp}}\), proving
the last bound.
\end{proof}

One side suffices.  For example, let \(B_0,B_1\)
be Hall's two coordinate-zero sections, and put \(A=B_0\cap B_1\).
The peripheral family
\((B_0\times\{0\})\cup(A\times\{1\})\) has \(224\) concepts
and empty core.  Its zero section cannot even start a peeling down to
\(A\): all nine corners of \(B_0\) lie in \(A\), and
\(|\kappa_A(B_0)|=52\).  Its one section already equals \(A\).
Conversely, two nonempty sets \(\kappa_A(B_i)\) at one coordinate do
not force a nonempty global core: at coordinate zero of the
\(292\)-concept example of Theorem~\ref{thm:damp-empty-core-counterexample},
their orders are \(54\) and \(17\),
while the global core is empty.  The fixed checker
\path{verify_damp_relative_midpoint_wings.py} verifies these controls
and both wing conditions at every coordinate of Hall's family and its
\(300\)-concept extension at \(373\).  These are necessary conditions
on a possible terminal counterexample to layer existence, not a
classification of such counterexamples.

\begin{theorem}[An ample forbidden pc-minor with empty unique-midpoint core]
\label{thm:damp-empty-core-counterexample}
Let \(G\subseteq Q_{12}\) be the certified \(291\)-concept family stored
as \texttt{base} in \path{damp_fixed_291_repair_creation.json}, with words
encoded as \(x=\sum_{i=0}^{11}x_i2^i\).  Put
\begin{align*}
 D&=\{96,126,394,1922,1934\},\\
 A&=\{553,763,2968,4008,4026\},\\
 \widehat G&=(G\setminus D)\cup A,\qquad
 H=\widehat G\cup\{3959\}.
\end{align*}
Both \(\widehat G\) and \(H\) are ample forbidden pc-minors.
The family \(\widehat G\) is cornerless and has \(291\) concepts;
its largest unique-midpoint core has \(76\) concepts.
The family \(H\) has \(292\) concepts, its only corner is \(3959\),
and its largest unique-midpoint core is empty.  The corner cube at
\(3959\) has support \(\{4,7,11\}\).

In particular, an ample forbidden pc-minor with empty unique-midpoint
core need not peel, and it need not have a corner whose deletion leaves
an empty unique-midpoint core.
\end{theorem}

\begin{proof}
This is a finite certificate proof.  The complete two vertex sets, their
source hashes, and the elementary-minor peelings are stored in
\path{damp_empty_core_counterexample.json}.  The checker
\path{verify_damp_empty_core_counterexample.py} recomputes all shattered
sets and all corners directly.  It verifies
\[
\begin{array}{c|c|c|c|c}
\text{family}&\text{order}&|\operatorname{Sh}|&
 \text{corners}&\text{largest core order}\\ \hline
\widehat G&291&291&\varnothing&76\\
H&292&292&\{3959\}&0
\end{array}
\]
The two Sandwich equalities prove ampleness.  Cornerlessness makes
\(\widehat G\) nonpeelable.  Every peeling of \(H\) would have to
start at \(3959\), leaving \(\widehat G\), so \(H\) is nonpeelable
as well.

For completeness of the core calculation, the checker enumerates every
distance-two endpoint pair and its common neighbors.  It obtains exactly
\(766\) unique-midpoint constraints in \(H\), also checked by independently
enumerating rooted neighbor pairs.  The verification report stores a
\(292\)-word pruning order for this \emph{fixed} constraint family:
when a root is removed, each constraint rooted there already has a
removed head.  Proposition~\ref{prop:damp-unique-midpoint-core} proves
that its core is empty.  Repeating the deterministic calculation on
\(\widehat G\) leaves \(76\) vertices; one internal rooted-pair
witness for each of them is recorded in
\path{damp_empty_core_counterexample_verification.json}.
Thus deleting the only corner activates a nonempty core.

Finally, every coordinate is active in both families.  For each family,
coordinate, and elementary restriction or contraction, the witness gives
a complete corner peeling of the resulting image.  The checker validates
all \(72\) orders, testing the corner cube at every deletion and checking
that the image is exhausted.  Hence every elementary pc-minor of each
family peels.  Pc-minor closure of \(\Damp\), together with their
nonpeelability, proves that both families are forbidden pc-minors.
\end{proof}

\begin{remark}[A compatible completion that remains nonpeelable]
\label{rem:damp-fixed-compatible-completion}
For the same \(G\) as in Theorem~\ref{thm:damp-empty-core-counterexample},
put
\[
\begin{split}
T=\{&186,553,680,763,898,910,1958,\\
    &2184,2592,2968,3137,3959,4008,4026\}.
\end{split}
\]
Every family \(G\cup U\), \(U\subseteq T\), is an ample forbidden
pc-minor.  Thus these fourteen additions do not repair \(G\), even
collectively.  This asserts \(2^{14}\) distinct embedded families,
without an isomorphism count.

Here the finite negative certificate concerns only the largest family
\(C=G\cup T\), of order \(305\).  It is ample and has empty
unique-midpoint core.  Assign a nine-bit unsigned rank \(r_x\) to each
\(x\in C\).  For every distance-two pair \(u,v\), impose
\[
 \bigvee_{\substack{x\in\{u,v\}\\w\in C_C(u,v)}} r_x<r_w.
\]
By Proposition~\ref{prop:damp-unique-midpoint-core}, these conditions
admit a solution exactly when \(C\) peels: a peeling supplies distinct
ranks, and arbitrarily refining tied ranks preserves every satisfied
strict inequality.  The explicit comparator encoding has \(17721\)
Boolean variables and \(87217\) clauses.  Its stored DRAT refutation
is independently checked by \texttt{drat-trim}.  The checker
\path{verify_damp_fixed_305_completion.py} regenerates the formula,
checks the proof hashes, and validates complete peelings of all
thirty-six elementary minors of both \(G\) and \(C\).
Its report is \path{damp_fixed_305_completion_verification.json}.

It remains to justify the simultaneous assertion about all subsets.
The acquired supports of the fourteen individually ample additions
are distinct.  Each added tip is therefore a corner after any subset
has been retained, by the union-shadow formula and the
ample corner-deletion criterion.  If \(G\cup U\) peeled, first deleting
\(T\setminus U\) from \(C\) would give a peeling of \(C\), a contradiction.
All intermediate families are consequently nonpeelable and ample;
iterating Proposition~\ref{prop:damp-one-concept-atlas} from the
forbidden family \(G\) proves their pc-minimality.  No subset
enumeration is needed.  This fixed conclusion does not prove a
general prohibition on cooperating additions over a forbidden seed.
\end{remark}

\begin{remark}[Failure of the core-or-nonample criterion]
\label{conj:damp-core-or-nonample}
\label{rem:damp-square-repair-conjecture}
The proposed equivalence between nonpeelability of an ample family and
nonemptiness of its unique-midpoint core is false by
Theorem~\ref{thm:damp-empty-core-counterexample}.
The counterexample also refutes the equivalent core-safe-corner assertion.
It is ample, so all its pc-minors are ample; failure to repair its square
constraints cannot produce a nonample antipodally punctured pc-minor.
The sufficient implication from a nonempty core to nonpeelability in
Proposition~\ref{prop:damp-unique-midpoint-core} remains valid.

The earlier bounded checks explain why that equivalence was plausible:
it has no mismatch among the \(5529\) nonempty ample subfamilies of
\(Q_4\), and Hall's \(299\)-concept family and the original
\(291\)-concept example have cores of orders \(158\) and \(76\).
These checks do not extend to the new one-corner example.  Its only
corner has degree three, so it also satisfies the exclusions of
Proposition~\ref{prop:damp-rank-two-extension-invariance} and
Corollary~\ref{cor:damp-low-degree-corner-descent}.
The example is not corner-descent terminal: deleting its corner leaves
\(\widehat G\), another ample forbidden pc-minor.  It therefore does
not settle the layer-existence or cornerless-seed questions for terminal
obstructions.
\end{remark}

\begin{remark}[What the global invariant isolates]
\label{rem:damp-boolean-cutset-scope}
Theorem~\ref{thm:damp-boolean-prefix-cutset} is the Boolean-lattice shadow
of nonshellability of the Yang complex: the metric intervals form a cutset
for all chambers of its order sphere.  It is finite, intrinsic, and its
behavior under both pc-minor operations is explicit.  However, ``meets
every maximal chain'' is still equivalent to the forbidden ordering
search, so it is not an admissibility axiom for the desired construction.

Proposition~\ref{prop:damp-unique-midpoint-core} removes that circularity
for the entire incomplete-square part: its obstruction is the largest
rooted-pair core, obtained by deterministic pruning.  Consequently the
remaining structural problem is narrower.  One must parameterize how full
squares couple otherwise prunable rooted-pair systems, and prove a
minor-cancellation theorem for that coupling.  This square-coupling datum,
rather than another certificate that a shelling order does not exist, is
the candidate global invariant that a complete generative classification
must expose.  Theorem~\ref{thm:damp-empty-core-counterexample} shows
that this coupling creates an additional obstruction even in an ample
forbidden pc-minor: the initial rooted-pair system is prunable, but
deleting its sole corner activates a nonempty core.
\end{remark}

\begin{proposition}[Local alphabet and pc-minor span laws for metric supports]
\label{prop:damp-metric-support-coordinate-span}
Let \(H\subseteq Q_F\) be isometric.  For
\(\mathcal S\subseteq\mathcal D_2(H)\), put
\[
 V(\mathcal S)
 =\bigcup_{\{u,v\}\in\mathcal S}
   \bigl(\{u,v\}\cup C_H(u,v)\bigr).
\]
Call \(\mathcal S\) \emph{order-infeasible} if there is no total order
\(\prec\) of \(V(\mathcal S)\) such that, for every
\(\{u,v\}\in\mathcal S\), some \(w\in C_H(u,v)\) satisfies
\[
 u\prec w\qquad\text{or}\qquad v\prec w.
\]
An inclusion-minimal order-infeasible family is called a \emph{metric
support}.  Then the following statements hold.
\begin{enumerate}
 \item Every \(D=\{u,v\}\in\mathcal D_2(H)\) has either one or two
       common neighbors.  Thus the metric failure cover has only two local
       types: a three-vertex constraint with one common neighbor and a
       four-vertex constraint supported on a square.
 \item If \(e\in F\), \(b\in\{0,1\}\), and
       \(H_{e=b}=\{x\in H:x_e=b\}\), then, after suppressing the constant
       coordinate \(e\),
       \[
        \mathcal D_2(H_{e=b})
        =\{\{u,v\}\in\mathcal D_2(H):u_e=v_e=b\},
        \qquad
        C_{H_{e=b}}(u,v)=C_H(u,v).
        \tag{\(\ast\)}
       \]
 \item If \(H\) is ample, then \(H\) is nonpeelable if and only if its
       distance-two interval system contains a metric support.
 \item Suppose that \(H\) is an ample forbidden pc-minor for corner
       peelability.  Then every metric support is coordinate-spanning:
       for every effective coordinate \(e\),
       \[
        V(\mathcal S)\cap H_{e=0}\ne\varnothing
        \qquad\text{and}\qquad
        V(\mathcal S)\cap H_{e=1}\ne\varnothing.
       \]
 \item Under the hypotheses of item~4, put \(\overline H=\pi_eH\) and
       \(\overline x=\pi_e(x)\).  For every effective coordinate \(e\)
       and every metric support \(\mathcal S\), contraction causes at least
       one of the following three changes:
       \begin{enumerate}
        \item two vertices of \(V(\mathcal S)\) have the same projection;
        \item some \(\{u,v\}\in\mathcal S\) satisfies \(u_e\ne v_e\),
              so that \(\overline u\) and \(\overline v\) are adjacent;
        \item for some \(\{u,v\}\in\mathcal S\) with \(u_e=v_e\), the
              contraction creates a new common neighbor:
              \[
               C_{\overline H}(\overline u,\overline v)
               \not\subseteq
               \pi_e\bigl(C_H(u,v)\bigr).
              \]
       \end{enumerate}
\end{enumerate}
\end{proposition}

\begin{proof}
If \(d_{G_H}(u,v)=2\), isometry says that \(u\) and \(v\) differ in
exactly two cube coordinates, say \(p\) and \(q\).  Their only possible
common neighbors in the ambient cube are \(u^{\{p\}}\) and
\(u^{\{q\}}\).  At least one belongs to \(H\), because a length-two
geodesic exists in \(G_H\).  This proves item~1.

The coordinate halfspace \(H_{e=b}\) is convex in \(G_H\).  Indeed, a
shortest path between two of its vertices has Hamming length and therefore
cannot flip the coordinate \(e\).  Hence distances inside the halfspace
are unchanged.  Moreover, a common neighbor of two distinct vertices with
\(e\)-coordinate \(b\) also has \(e\)-coordinate \(b\): otherwise both
vertices would be obtained from that common neighbor by flipping only
\(e\), and would coincide.  This proves \((\ast)\).

By the static distance-two criterion used in
Theorem~\ref{thm:damp-yang-metric-obstruction-class}, \(H\) is
nonpeelable exactly when the full family \(\mathcal D_2(H)\) is
order-infeasible.  Since this family is finite, it then contains an
inclusion-minimal order-infeasible subfamily.  Conversely, every order of
\(H\) restricts to an order of \(V(\mathcal S)\), so an order-infeasible
subfamily already prevents a peeling.  This proves item~3.

Finally, let \(e\) be effective and suppose that
\(V(\mathcal S)\subseteq H_{e=b}\).  By \((\ast)\), the same constraints
form an order-infeasible subfamily of the distance-two interval system of
\(H_{e=b}\).  Item~3 makes this restriction nonpeelable.  It is a proper
pc-minor of \(H\), contradicting pc-minimality.  Therefore every metric
support meets both coordinate halfspaces, proving item~4.

Suppose, finally, that none of the three changes in item~5 occurs.
Projection is then injective on \(V(\mathcal S)\), every constrained pair
remains at distance two, and
\[
 C_{\overline H}(\overline u,\overline v)
 =\pi_e\bigl(C_H(u,v)\bigr)
 \qquad(\{u,v\}\in\mathcal S).
\]
Indeed, projection preserves every old same-layer common neighbor, while
the absence of item~5(c) excludes all new ones.  Thus projection identifies
the constraint system on \(\mathcal S\) with an order-infeasible subsystem
of \(\mathcal D_2(\overline H)\).  Item~3 makes the proper contraction
\(\overline H\) nonpeelable, again contradicting pc-minimality.  This proves
item~5.
\end{proof}

\begin{remark}[Structural content and remaining contraction problem]
\label{rem:damp-metric-support-contraction-wall}
The proposition replaces an arbitrary failure region by a finite family of
ternary and quaternary interval constraints and proves two global
minimality laws: no support can hide in a coordinate halfspace, and every
contraction must hit every support by a collision, a distance drop, or a
new common neighbor.  These are necessary structural conditions, not yet
admissibility for a generative construction, because order-infeasibility
remains equivalent to the original shellability failure.  The next missing
functorial statement is a cancellation theorem: one must give intrinsic
attachment axioms under which the three local changes jointly destroy all
minimal supports.  The direct and switched punctured-expansion profiles
retain precisely the owner data that the uncolored quotient interval system
forgets.
\end{remark}

\begin{theorem}[Canonical Yang order-sphere obstruction]
\label{thm:damp-yang-order-sphere-obstruction}
Let \(H\subseteq Q_F\) be ample, put \(m=|H|\ge2\), and let
\[
 \mathscr O(H)
 =\operatorname{Ord}\bigl(2^H\setminus\{\varnothing,H\}\bigr)
\]
be the order complex of the proper nonempty subsets of \(H\).  Thus
\(\mathscr O(H)\) is the barycentric subdivision of the boundary of an
\((m-1)\)-simplex, and its chambers are the orders of \(H\), represented
by their chains of proper prefixes.

For a latent chart \(B=(y,S)\), put
\[
 a_B=y^S,
 \qquad
 p_{B,q}=y^{S\setminus\{q\}}
 \quad(q\in S),
\]
and define a poset \(P_B\) on \(H\) by the following generating
relations:
\begin{equation}
 \begin{array}{ll}
  a_B\mathrel{\triangleleft_B}p_{B,q}
       &(q\in S),\qquad y\notin H,\\[2mm]
  y\mathrel{\triangleleft_B}a_B
       \mathrel{\triangleleft_B}p_{B,q}
       &(q\in S),\qquad y\in H.
 \end{array}
 \label{eq:damp-yang-failure-poset}
\end{equation}
Let \(\mathcal J(P_B)\) be its lattice of order ideals, and define the
\emph{failure region}
\[
 \mathscr R_B
 =
 \operatorname{Ord}\bigl(
   \mathcal J(P_B)\setminus\{\varnothing,H\}
 \bigr)
 \subseteq\mathscr O(H).
\]
Finally put
\[
 \mathscr U(H)=
 \bigcup_{B\in\mathcal L(H)}\mathscr R_B.
\]
Then
\begin{equation}
 \Delta_H\text{ is nonshellable}
 \quad\Longleftrightarrow\quad
 \mathscr U(H)=\mathscr O(H)
 \quad\Longleftrightarrow\quad
 \widetilde H_{m-2}\bigl(\mathscr U(H);\mathbb F_2\bigr)
       \cong\mathbb F_2.
 \label{eq:damp-yang-order-sphere-obstruction}
\end{equation}
If \(\Delta_H\) is shellable, the last homology group is zero.
Consequently
\[
 \omega_{\mathrm{ord}}(H)
 =
 \dim_{\mathbb F_2}
 \widetilde H_{m-2}\bigl(\mathscr U(H);\mathbb F_2\bigr)
 \in\{0,1\}
\]
is a canonical global invariant of the colored Yang protector atlas,
equal to one exactly for nonpeelable \(H\).

More precisely, let \(\varnothing\ne\mathcal I\subseteq\mathcal L(H)\).
Write \(P_{\mathcal I}\) for the preorder generated by all relations in
\(P_B\), \(B\in\mathcal I\), and let \(\overline P_{\mathcal I}\) be
the poset obtained by contracting the strongly connected classes of this
preorder.  If \(k_{\mathcal I}\) is the number of these classes, then
\begin{equation}
 \bigcap_{B\in\mathcal I}\mathscr R_B
 =
 \operatorname{Ord}\bigl(
   \mathcal J(\overline P_{\mathcal I})
       \setminus\{\varnothing,\overline P_{\mathcal I}\}
 \bigr).
 \label{eq:damp-yang-failure-region-intersection}
\end{equation}
This intersection is homotopy equivalent to
\[
 \begin{cases}
  S^{k_{\mathcal I}-2},
     &\overline P_{\mathcal I}\text{ is an antichain},\\
  \{\mathrm{pt}\},
     &\overline P_{\mathcal I}\text{ is not an antichain},
 \end{cases}
\]
with \(S^{-1}=\varnothing\).  Hence the augmented
Mayer--Vietoris spectral sequence of the canonical cover
\((\mathscr R_B)_{B\in\mathcal L(H)}\) has
\begin{equation}
 E^1_{p,q}
 =
 \bigoplus_{\substack{
      \mathcal I\subseteq\mathcal L(H)\\
      |\mathcal I|=p+1}}
 \widetilde H_q\left(
      \bigcap_{B\in\mathcal I}\mathscr R_B;\mathbb F_2
 \right)
 \ \Longrightarrow\
 \widetilde H_{p+q}\bigl(\mathscr U(H);\mathbb F_2\bigr),
 \label{eq:damp-yang-failure-cover-spectral-sequence}
\end{equation}
and its only nonzero \(E^1\)-summands are the one-dimensional summands
with \(\overline P_{\mathcal I}\) an antichain and
\(q=k_{\mathcal I}-2\).
\end{theorem}

\begin{proof}
An order \(\mathbf x=(x_1,\ldots,x_m)\) is a chamber of
\(\mathscr R_B\) precisely when every prefix of \(\mathbf x\) is an
order ideal of \(P_B\), or equivalently when \(\mathbf x\) is a linear
extension of \(P_B\).  If \(y\notin H\), this says that the antipode
\(a_B\) precedes every protector \(p_{B,q}\).  If \(y\in H\), it says
that
\[
 y\prec a_B\prec p_{B,q}
 \qquad\text{for every }q\in S.
\]
These are exactly the two ways in which all allowed gate relations in
the chart \(B\) point backwards.  Corollary
~\ref{cor:damp-yang-protector-atlas} therefore says that an order is a
shelling order if and only if its chamber belongs to none of the
\(\mathscr R_B\).  Thus \(\Delta_H\) is nonshellable exactly when the
failure regions contain every chamber.  Since they are subcomplexes,
they then contain every face, proving the first equivalence in
\eqref{eq:damp-yang-order-sphere-obstruction}.

The chamber-adjacency graph of \(\mathscr O(H)\) is connected: adjacent
chambers differ by interchanging two consecutive members of an order.
Over \(\mathbb F_2\), the boundary equation at their common ridge forces
the two chamber coefficients of every top-dimensional cycle to agree.
If a subcomplex of \(\mathscr O(H)\) omits one chamber, give that chamber
coefficient zero; connectivity then forces every coefficient to be zero.
Hence a proper subcomplex has zero homology in dimension \(m-2\), whereas
\(\mathscr O(H)\cong S^{m-2}\) has one-dimensional top homology.  This
proves the second equivalence and the assertion about
\(\omega_{\mathrm{ord}}\).

A subset of \(H\) is an order ideal of every \(P_B\),
\(B\in\mathcal I\), exactly when it is an order ideal of the generated
preorder \(P_{\mathcal I}\).  Such an ideal is a union of strongly
connected classes, and these unions are precisely the order ideals of
\(\overline P_{\mathcal I}\).  This proves
\eqref{eq:damp-yang-failure-region-intersection}.

It remains to identify its homotopy type.  If
\(\overline P_{\mathcal I}\) is an antichain on \(k_{\mathcal I}\)
members, its ideal lattice is Boolean, so its proper-part order complex
is \(S^{k_{\mathcal I}-2}\).  Otherwise the atoms of its ideal lattice
are the singleton minimal elements.  Even the join of all these atoms is
then a proper ideal.  The atomic crosscut complex is consequently a full
simplex, and the crosscut argument contracts the proper-part order
complex.  Finally, the standard augmented intersection double complex
of a finite cover by subcomplexes has the \(E^1\)-page displayed in
\eqref{eq:damp-yang-failure-cover-spectral-sequence}: for each simplex
of the union, the horizontal summands containing it form a simplex and
hence an exact augmented chain complex.  The intersection calculation
gives the stated support of the \(E^1\)-page.
\end{proof}

\begin{corollary}[Balanced circuit carrier]
\label{cor:damp-yang-balanced-circuit-carrier}
If \(H\) is nonpeelable, there is a nonempty collection
\(\mathcal I\subseteq\mathcal L(H)\) such that
\begin{enumerate}
 \item every generating relation of \(P_{\mathcal I}\) lies on a directed
       cycle; equivalently,
       \(\overline P_{\mathcal I}\) is an antichain; and
 \item
       \begin{equation}
        |\mathcal I|+k_{\mathcal I}=|H|+1.
        \label{eq:damp-yang-balanced-circuit}
       \end{equation}
\end{enumerate}
If \(r_{\mathcal I}\) concepts are incident with a generating relation
and these concepts form \(c_{\mathcal I}\) nontrivial strongly connected
components, then the balance equation is equivalently
\[
 |\mathcal I|=r_{\mathcal I}-c_{\mathcal I}+1.
\]
Such balanced circuit collections are necessary carriers of the global
obstruction class; their existence alone is not asserted to be
sufficient.
\end{corollary}

\begin{proof}
When \(H\) is nonpeelable, the top class in
\eqref{eq:damp-yang-order-sphere-obstruction} survives to the abutment of
\eqref{eq:damp-yang-failure-cover-spectral-sequence}.  Some nonzero
\(E^1_{p,q}\)-summand of total degree \(p+q=m-2\) must therefore support
it.  Theorem~\ref{thm:damp-yang-order-sphere-obstruction} makes
\(\overline P_{\mathcal I}\) an antichain,
\(p=|\mathcal I|-1\), and \(q=k_{\mathcal I}-2\).  The total-degree
identity is exactly
\eqref{eq:damp-yang-balanced-circuit}.  An antichain condensation has no
edge between distinct strongly connected classes, so every generating
relation lies within one class and hence on a directed cycle.  Finally,
the \(m-r_{\mathcal I}\) isolated concepts and the
\(c_{\mathcal I}\) nontrivial components give
\(k_{\mathcal I}=m-r_{\mathcal I}+c_{\mathcal I}\), yielding the last
formula.
\end{proof}

\begin{proposition}[The first cancellation page is a cut-nerve homology]
\label{prop:damp-yang-first-cancellation-cut-nerve}
Let \(\Pi(H)\) be the set of partitions of \(H\).  For
\(\pi\in\Pi(H)\), let \(\mathcal A_\pi\) be the latent charts whose every
generating failure relation has both endpoints in one part of \(\pi\).
Define \(\mathcal N_\pi\subseteq 2^{\mathcal A_\pi}\) by
\[
 \mathcal N_\pi
 =
 \left\{
  \mathcal I\subseteq\mathcal A_\pi:
  \text{the relations of \(\mathcal I\) fail to make some part of
        \(\pi\) strongly connected}
 \right\}.
\]
This is a simplicial complex.  If \(|\pi|=k\), then the row
\(q=k-2\) of the first differential in
\eqref{eq:damp-yang-failure-cover-spectral-sequence} decomposes as
\begin{equation}
 \bigl(E^1_{\bullet,k-2},d^1\bigr)
 \cong
 \bigoplus_{\substack{\pi\in\Pi(H)\\|\pi|=k}}
 C_\bullet\bigl(
    2^{\mathcal A_\pi},\mathcal N_\pi;\mathbb F_2
 \bigr).
 \label{eq:damp-yang-first-page-partition-decomposition}
\end{equation}
Consequently,
\begin{equation}
 E^2_{p,k-2}
 \cong
 \bigoplus_{\substack{\pi\in\Pi(H)\\|\pi|=k}}
 \widetilde H_{p-1}(\mathcal N_\pi;\mathbb F_2).
 \label{eq:damp-yang-second-page-connectivity-complex}
\end{equation}

The complexes on the right have a canonical cut-nerve model.  For a
part \(C\in\pi\) and a nonempty proper subset \(U\subsetneq C\), put
\[
 \mathcal S_\pi(C,U)
 =
 \left\{
  B\in\mathcal A_\pi:
  P_B\text{ has no generating arc from \(U\) to \(C\setminus U\)}
 \right\}.
\]
After discarding empty \(\mathcal S_\pi(C,U)\), let
\(\mathcal K_\pi\) be the nerve of the family of simplices
\[
 \left(
 2^{\mathcal S_\pi(C,U)}
 \right)_{\substack{C\in\pi\\
                    \varnothing\ne U\subsetneq C}}.
\]
Then
\begin{equation}
 \mathcal N_\pi\simeq\mathcal K_\pi.
 \label{eq:damp-yang-connectivity-cut-nerve}
\end{equation}
In particular, if \(H\) is nonpeelable, then for some partition
\(\pi\) with \(k=|\pi|\),
\begin{equation}
 \widetilde H_{|H|-k-1}(\mathcal K_\pi;\mathbb F_2)\ne0.
 \label{eq:damp-yang-cut-nerve-necessary-class}
\end{equation}
\end{proposition}

\begin{proof}
A nonzero summand of \(E^1_{p,k-2}\) is indexed by a chart collection
\(\mathcal I\) whose strongly connected quotient is an antichain of
order \(k\).  Let \(\pi\) be its strongly connected partition.  Every
failure arc of \(\mathcal I\) lies inside a part of \(\pi\), and the
arcs strongly connect every part.  Conversely, these two conditions
make the strongly connected partition exactly \(\pi\).  Thus the
summands belonging to \(\pi\) are indexed by the faces of the simplex
\(2^{\mathcal A_\pi}\) that do not belong to \(\mathcal N_\pi\).

Remove one chart from \(\mathcal I\).  If every part remains strongly
connected, the strongly connected partition is still \(\pi\).
Both intersections in
\eqref{eq:damp-yang-failure-region-intersection} are then the same
Boolean sphere on the parts of \(\pi\), so the inclusion induces the
identity on its one-dimensional homology over \(\mathbb F_2\).  If a
part splits, the new sphere has a different dimension and the component
of \(d^1\) in row \(q=k-2\) is zero.  This is exactly the relative
simplicial boundary of
\((2^{\mathcal A_\pi},\mathcal N_\pi)\), proving
\eqref{eq:damp-yang-first-page-partition-decomposition}.  Since the
first member of the pair is a simplex, its relative homology is the
shifted reduced homology of the second, which proves
\eqref{eq:damp-yang-second-page-connectivity-complex}.

A directed graph on \(C\) is strongly connected exactly when every
nonempty proper \(U\subsetneq C\) has an arc leaving \(U\).  Therefore a
chart collection belongs to \(\mathcal N_\pi\) exactly when it is
contained in \(2^{\mathcal S_\pi(C,U)}\) for some displayed cut:
\[
 \mathcal N_\pi
 =
 \bigcup_{\substack{C\in\pi\\
                    \varnothing\ne U\subsetneq C}}
 2^{\mathcal S_\pi(C,U)}.
\]
Every nonempty intersection in this cover is again a simplex, on the
intersection of the corresponding chart sets.  The nerve lemma gives
\eqref{eq:damp-yang-connectivity-cut-nerve}.

Finally, a surviving top class has total degree \(m-2\).  It must have
a nonzero representative on the second page, so
\eqref{eq:damp-yang-second-page-connectivity-complex} is nonzero for
some \(p+k-2=m-2\), namely \(p=m-k\).  Apply
\eqref{eq:damp-yang-connectivity-cut-nerve} to obtain
\eqref{eq:damp-yang-cut-nerve-necessary-class}.
\end{proof}

\begin{remark}[Higher differentials are balanced partition splittings]
\label{rem:damp-yang-higher-differential-rank-law}
The bidegree of the homological differential gives a sharp remaining
contract.  A possible nonzero map out of a spherical row has the form
\[
 d^r:E^r_{p,k-2}\longrightarrow E^r_{p-r,k+r-3}.
\]
By Theorem~\ref{thm:damp-yang-order-sphere-obstruction}, the target row
can be supported only on strongly connected partitions with
\[
 k'=k+r-1
\]
parts.  Thus an \(r\)-th differential can only record a net splitting
into exactly \(r-1\) additional strongly connected classes.  Dually, an
incoming \(r\)-th differential can only originate from a partition with
\(k-r+1\) parts.  Hence a second-page class with no compatible balanced
refinement or coarsening of these ranks is spectrally isolated and
survives to the global obstruction class.  The remaining generative
problem is to read these balanced splitting maps, or their absence, from
the direct and switched chart attachments.
\end{remark}

\begin{proposition}[A binary partition split has a canonical filling torsor]
\label{prop:damp-yang-binary-split-filling-torsor}
Let \(\pi\) be a partition and let \(\sigma\) be obtained by replacing
one part \(C\in\pi\) by two nonempty parts
\(C_0,C_1\).  Put
\[
 \mathbb S_\pi
 =\operatorname{Ord}\bigl(2^\pi\setminus\{\varnothing,\pi\}\bigr),
 \qquad
 \mathbb S_\sigma
 =\operatorname{Ord}\bigl(2^\sigma\setminus
      \{\varnothing,\sigma\}\bigr).
\]
Identifying a union of parts of \(\pi\) with the corresponding union of
parts of \(\sigma\) embeds
\[
 \iota_{\pi,\sigma}:\mathbb S_\pi\hookrightarrow\mathbb S_\sigma.
\]
For \(b\in\{0,1\}\), let \(\mathbb B_b(\pi,\sigma)\) be the order
complex of the nonempty proper subsets
\(\varnothing\ne A\subsetneq\sigma\) satisfying
\begin{equation}
 C_{1-b}\in A\quad\Longrightarrow\quad C_b\in A.
 \label{eq:damp-yang-binary-split-hemisphere-rule}
\end{equation}
Then \(\mathbb B_0(\pi,\sigma)\) and
\(\mathbb B_1(\pi,\sigma)\) are PL balls with
\begin{equation}
 \mathbb B_0(\pi,\sigma)\cap\mathbb B_1(\pi,\sigma)
 =\iota_{\pi,\sigma}(\mathbb S_\pi),
 \qquad
 \mathbb B_0(\pi,\sigma)\cup\mathbb B_1(\pi,\sigma)
 =\mathbb S_\sigma.
 \label{eq:damp-yang-binary-split-two-hemispheres}
\end{equation}
Thus the old sphere is the common equator of the refined sphere.

Over \(\mathbb F_2\), let \(h_b(\pi,\sigma)\) be the sum of the
top-dimensional simplices of \(\mathbb B_b(\pi,\sigma)\).  The two
chains are precisely the two top-dimensional fillings of the equatorial
fundamental cycle:
\begin{equation}
 \partial h_0(\pi,\sigma)
 =\partial h_1(\pi,\sigma)
 =[\mathbb S_\pi],
 \qquad
 h_0(\pi,\sigma)+h_1(\pi,\sigma)
 =[\mathbb S_\sigma].
 \label{eq:damp-yang-binary-split-filling-torsor}
\end{equation}
Consequently the fillings form a canonical two-element
\(\mathbb F_2\)-torsor indexed by the two new parts.
\end{proposition}

\begin{proof}
The subsets satisfying
\eqref{eq:damp-yang-binary-split-hemisphere-rule} are the order ideals
of the one-relation poset
\(C_b<C_{1-b}\), together with the unaffected parts of \(\pi\).
Their proper-part order complex is the closed half of the Coxeter sphere
\(\mathbb S_\sigma\) in which \(C_b\) precedes \(C_{1-b}\); it is a PL
ball.  A subset satisfies both implications exactly when it contains
either both \(C_0,C_1\) or neither, which is precisely the embedded
Boolean lattice on the parts of \(\pi\).  Moreover, a chain of subsets
cannot contain both a set containing \(C_0\) but not \(C_1\) and a set
containing \(C_1\) but not \(C_0\), because those two sets are
incomparable at the first change of the split coordinates.  Hence every
simplex belongs to one of the two balls.  This proves
\eqref{eq:damp-yang-binary-split-two-hemispheres}.

The boundary of the fundamental chain of either PL ball is the
fundamental cycle of its common equator, giving the first identity in
\eqref{eq:damp-yang-binary-split-filling-torsor}.  Their sum is the
fundamental cycle of the sphere obtained by gluing the two balls, giving
the second identity.  Finally, two top-dimensional chains with the
displayed boundary differ by a top cycle of
\(\mathbb S_\sigma\).  Its top cycle space over \(\mathbb F_2\) is
one-dimensional, so the two displayed chains are the only two fillings.
\end{proof}

\begin{proposition}[Iterated binary fillings factor through a permutation]
\label{prop:damp-yang-iterated-filling-permutation}
Let \(\tau\) refine a partition \(\pi\) by replacing one part
\(C\in\pi\) with nonempty parts indexed by a finite set \(L\).
An \emph{oriented binary refinement tree} for \(C\) is a rooted full
binary tree whose leaves are \(L\), together with a choice of first child
at every internal node.  Reading the first child before the second
recursively induces a total order
\begin{equation}
 \lambda_1\prec_T\lambda_2\prec_T\cdots\prec_T\lambda_{|L|}
 \label{eq:damp-yang-refinement-tree-leaf-order}
\end{equation}
of \(L\).

Starting from the parts of \(\pi\), perform the binary splits prescribed
by the tree and at each split take the hemisphere filling indexed by its
first child.  The resulting iterated filling region in
\(\mathbb S_\tau\) is
\begin{equation}
 \mathbb B(T)
 =
 \operatorname{Ord}\bigl(
   \mathcal J(P_T)\setminus\{\varnothing,\tau\}
 \bigr),
 \label{eq:damp-yang-iterated-filling-region}
\end{equation}
where \(P_T\) makes the members of \(L\) the chain in
\eqref{eq:damp-yang-refinement-tree-leaf-order} and leaves all parts of
\(\pi\setminus\{C\}\) incomparable with that chain and with one another.
In particular:
\begin{enumerate}
 \item \(\mathbb B(T)\) is a PL ball;
 \item it depends only on the induced permutation of \(L\), not on the
       binary refinement tree; and
 \item every permutation of \(L\) is obtained.
\end{enumerate}
Consequently all coherence changes between binary refinement trees factor
through the natural symmetric-group action of
\(\mathfrak S_L\).  Independent refinements of distinct parts commute,
whereas interacting refinements of one part carry permutation-valued
holonomy.
\end{proposition}

\begin{proof}
Proceed by induction on the refinement tree.  A leaf contributes the
one-element chain.  At an internal node, the chosen hemisphere orders
every leaf of the first child before every leaf of the second child.
Thus it takes the ordinal sum of the two child posets.  By induction both
child posets are chains in their recursively read leaf orders, and their
ordinal sum is the concatenated chain
\eqref{eq:damp-yang-refinement-tree-leaf-order}.  Iterating the
one-relation order-ideal description from
Proposition~\ref{prop:damp-yang-binary-split-filling-torsor} gives
\eqref{eq:damp-yang-iterated-filling-region}.

The linear-extension chambers of a poset consisting of one chain and
otherwise incomparable elements form a closed convex region of the
Coxeter complex, hence a PL ball.  The poset \(P_T\), and therefore its
order-ideal complex, depends only on the final leaf order.  Conversely,
given any order of \(L\), a binary comb whose successive first children
are read in that order realizes it.  Finally, trees on disjoint parts
impose relations on disjoint leaf sets and their ordinal-sum operations
commute.  Trees refining the same part change only the induced order of
its leaves, giving the stated symmetric-group action.
\end{proof}

\begin{remark}[The exact finite local state group]
\label{rem:damp-yang-local-gate-symmetric-group}
For a latent chart \(B=(y,S)\), define its gate-label set by
\[
 G_B=
 \begin{cases}
  S,&y\notin H,\\
  S\mathbin{\dot\cup}\{\ast\},&y\in H,
 \end{cases}
\]
where \(q\in S\) labels the \(q\)-protector and \(\ast\) labels the
escape.  The local gate states are the points of \(G_B\).  Proposition
~\ref{prop:damp-yang-iterated-filling-permutation} shows that iterated
binary fillings on this label set factor through the candidate coherence
group \(\mathfrak S(G_B)\), with gate alphabet the homogeneous space
\[
 G_B\cong
 \mathfrak S(G_B)/
 \mathfrak S(G_B\setminus\{g\})
\]
for any fixed \(g\in G_B\).  Thus an absent rank-\(r\) chart has the
natural \(\mathfrak S_r/\mathfrak S_{r-1}\) alphabet, while a present
rank-\(r\) chart has the
\(\mathfrak S_{r+1}/\mathfrak S_r\) alphabet.  The binary owner switch
is a transposition.  This explains simultaneously why missing squares
are binary and why ranks three and four cannot be compressed to a scalar
parity before their permutation holonomy has been controlled.
This identifies the smallest natural finite group containing the local
binary moves; proving that the actual chart-cancellation maps factor
through this action remains part of the coherence theorem.
\end{remark}

\begin{remark}[Owner transport is the local filling torsor]
\label{rem:damp-yang-owner-filling-torsor}
Proposition~\ref{prop:damp-yang-binary-split-filling-torsor} supplies the
missing topological meaning of the owner bit in
Proposition~\ref{prop:damp-latent-facet-protector-transport}.  Once the
protector coordinate \(q\) is fixed, splitting a projected class into
its two lift layers creates exactly the two equatorial fillings
\(h_0,h_1\).  A singly owned protector edge permits only its owner's
filling; a doubly owned separator edge permits both.  Thus the
\(\mathbb F_2\)-torsor is canonical within each coordinate-labelled
protector sector, while
Proposition~\ref{prop:damp-yang-iterated-filling-permutation} supplies
the ambient symmetric-group coherence in which transitions among sectors
must land.  What remains is
to prove that the cancellation maps
of the Yang failure cover factor through these local filling choices and
that their compositions around partition-refinement diamonds give all,
and only, the higher differentials.  That coherence theorem would turn
the present local torsors into the desired global holonomy invariant.
\end{remark}

\begin{remark}[Circuit existence is not the invariant]
\label{rem:damp-yang-balanced-circuit-not-sufficient}
The balance condition in
Corollary~\ref{cor:damp-yang-balanced-circuit-carrier} is only the support
condition on the first page of the spectral sequence.  It is not
sufficient even for the full square \(Q_2\).  Take all four latent square
charts, one for each present tip.  Their failure relations form a strongly
connected digraph on the four concepts, so
\[
 |\mathcal I|=4=|Q_2|,
 \qquad
 k_{\mathcal I}=1,
\]
and the balance equation holds.  Nevertheless \(Q_2\) peels.  The
corresponding \(E^1\)-generator is killed by a later differential.

The same warning appears at scale in the known obstructions.  A worker
diagnostic finds strongly connected generating subfamilies of \(202\)
latent charts for Hall's order-\(299\) seed and \(201\) charts for the
certified order-\(291\) seed.  Padding them to \(299\) and \(291\) charts
gives balanced carriers, but their mere existence does not distinguish the
two seeds from shellable systems such as \(Q_2\).  The durable artifact is
\path{damp_yang_balanced_circuits.json}, generated by
\path{analyze_damp_yang_balanced_circuits.py}.  Thus the smaller invariant
must encode survival of the canonical fundamental class---equivalently,
the cancellation maps between circuit carriers---rather than cycle rank or
strong connectivity alone.
\end{remark}

\begin{remark}[The candidate global invariant]
\label{rem:damp-yang-protector-holonomy-target}
Ordinary topology of \(\Delta_H\) cannot detect the obstruction: every
ample Yang complex in the present problem is already a constructible
Cohen--Macaulay ball.  Theorem
~\ref{thm:damp-yang-order-sphere-obstruction} nevertheless gives an exact
global invariant of its colored protector atlas.  It is not yet the final
generative answer.  The order sphere has \(m!\) chambers, and merely
computing its top class is an order-free recognition oracle equivalent to
shellability.  The real gain is the spectral compression: every possible
carrier of the fundamental class is a balanced collection of latent charts
whose local failure relations close into directed circuits.

Associate local gate states to the charts and transport them through the
color links and color-pair deletions.  Proposition
~\ref{prop:damp-latent-facet-protector-transport} shows that contraction
preserves the protector coordinate \(q\); only the owner layer may change.
Thus the natural smaller invariant is the monodromy of this finite-state
transport system on the balanced circuit carriers.  If the transition
relations can be proved to be canonical bijections, they define a
holonomy representation into a finite permutation group.  If, more
strongly, the remaining owner states form two-element torsors, the
representation reduces to a class
\[
 [\omega_H]\in H^1(\Gamma_H;\mathbb F_2),
\]
whose value is obtained by multiplying transition signs around cycles.

Neither conclusion is currently automatic.  Rank-three and rank-four
charts have three and four coordinate-labelled protector states, present
tips add escape states, and shared protector edges impose compatibility.
Merely defining a global section to be an acyclic selection would restate
the shelling search.  The required structural theorem is therefore a
\emph{torsor-reduction theorem}: from the signed chart attachments alone,
it must canonically quotient the gate states, prove functorial transport
under all elementary pc-minors, and show that vanishing of the resulting
holonomy constructs the minor shellings.  Only after that theorem would
\([\omega_H]\), or its finite-group analogue, be a genuine order-free
global invariant rather than new notation for shellability.
\end{remark}

\begin{remark}[Why initially missing blocks do not suffice]
\label{rem:damp-latent-square}
The full square \(Q_2\) has no missing-tip one-hole block.  Nevertheless,
after deleting \(00\), the punctured square
\(\mathsf P(00,\{1,2\})\) becomes active and prevents the antipode \(11\)
from being the next corner.  The latent word rule says exactly that if
\(00\) is first, then the next vertex must be \(01\) or \(10\).
Equivalently, the first two vertices of a peeling of \(Q_2\) must be
adjacent.  Thus an invariant built only from the holes present in the
initial family omits essential transition data even in dimension two.
As a bounded regression, the static word rule and direct corner deletion
agree on all \(93{,}248\) orders of all \(127\) nonempty ample families in
\(Q_3\); this computation is not used in the proof.  The durable artifact is
\path{damp_latent_punctured_word_q3.json}, generated by
\path{verify_damp_latent_punctured_word.py}.
The same implementation was checked on all \(36\) elementary images of
Hall's obstruction.  It examined \(104{,}839\) latent blocks and found a
facet protector for every nonescape state, as the corollary requires.  The
artifact is \path{damp_hall_latent_protectors.json}, generated on the worker
by \path{analyze_damp_hall_latent_protectors.py}.  Neither computation is
used in the proof.
\end{remark}

\begin{remark}[Consequences for the holonomy program]
\label{rem:damp-latent-facet-protector-scope}
Theorem~\ref{thm:damp-latent-punctured-face-word-criterion} identifies
the canonical local states that a global invariant must synchronize.
They are coordinate-labelled rather than binary in general.  For a present
tip, an \(S\)-block has one escape state and \(|S|\) facet-protector
states; for an absent tip it has the \(|S|\) protector states.  Thus an
absent square is binary, while absent blocks of ranks three and four have
three and four states.  Moreover, all latent
blocks are needed: inclusion-maximal holes present at the root do not
record the blocks activated by later deletions.

Corollary~\ref{cor:damp-latent-facet-protector-criterion} is still a
recognition theorem, not the desired generative classification.  Choosing
states and asking whether their precedence digraph is acyclic is another
finite form of the shelling search.  The remaining structural problem is
to prove that, on the recovered paired one-hole presentations of
overlap-atomic obstructions and their elementary images, these states
factor through canonical block transitions.  A binary
\(\mathbb F_2\)-holonomy can be complete only after proving a quotient
theorem that collapses these coordinate states without losing
compatibility; it cannot be assumed from the direct/switched dichotomy
alone.
\end{remark}

\begin{remark}[Known Hall seeds are gluing-critical in every coordinate]
\label{rem:damp-hall-overlap-shellability-audit}
A worker audit reconstructed the fixed \(299\)-concept Hall obstruction,
the certified \(291\)-concept obstruction, and the \(27\) cornerless
families in the radius-two Hall exchange atlas.  Every one of their \(348\)
coordinate overlaps \(A=B_0\cap B_1\) is an ample corner-peelable family
of VC dimension two.  Their orders are \(67\) in \(339\) cases, \(65\) in
eight cases, and \(59\) in one case.  Hence every tested coordinate has
the gluing-critical signature.  The same audit applied the two root-cut
tests to both relative wings at every coordinate.  All \(696\) wings are
upper-root-locked, none is lower-root-locked, and the direct ampleness test
agrees with the root-cut prediction in every case.  More precisely, every
wing has between two and fourteen possible first relative facets, but no
possible last relative facet.  Thus the known Hall mechanism is one-sided
terminality, not biroot-locking.

The canonical relative one-hole recovery contains \(164{,}762\) minimal
blocks over the \(57{,}372\) relative concepts.  Their support orders are
\(2,3,4\).  At \(49{,}404\) concepts the least block has order two, at
\(3{,}288\) it has order three, and at \(4{,}680\) it has order four; in
particular, no tested wing is covered by punctured squares alone.  Passing
to inclusion-maximal blocks gives \(86{,}226\) canonical blocks.  A
deterministic greedy subcover uses \(18{,}352\) of them, between \(26\%\)
and \(40\%\) as many blocks as relative concepts in each wing, and one
order-four block protects as many as eleven relative concepts.  The greedy
counts are compression evidence only; the construction and recovery theorem
does not depend on their optimality.  The corresponding artifact is
\path{damp_relative_one_hole_blocks.json}.

By contrast, if \(K\notin\Damp\), then the product
\(K\square Q_1\) has \(B_0=B_1=A=C=K\) in the product coordinate
and therefore has the inherited-site signature.  The overlap signature
thus separates the product adversary from all presently known Hall seeds,
although this finite audit is not a proof that every Cartesian-prime seed
is gluing-critical.  The finite audit does not itself supply the desired
construction: upper locking is equivalent to the absence of an
ample-preserving deletion and therefore cannot simply be installed as an
admissibility oracle.  Its role is to identify the structural fact that the
future cut-block axioms must force.  The reproducible artifact is
\path{damp_seed_overlap_shellability.json}, generated by
\path{analyze_damp_seed_overlap_shellability.py}.
\end{remark}

\begin{proposition}[First-loss criterion for a corner peeling]
\label{prop:damp-trace-death-peeling}
Let \(\varnothing\ne H\subseteq Q_F\) be ample, put
\(\Sigma=\operatorname{Sh}(H)\), and write \(m=|H|=|\Sigma|\).
For an ordering \(\mathbf x=(x_1,\ldots,x_m)\) of \(H\), define
\begin{equation}
 \ell_{\mathbf x}(D,q)
 =\max\{i:x_i|_D=q\},
 \qquad
 \delta_{\mathbf x}(D)
 =\min_{q\in Q_D}\ell_{\mathbf x}(D,q)
 \quad(D\in\Sigma).
 \label{eq:damp-trace-death-time}
\end{equation}
The maxima exist because \(H\) shatters \(D\).  Then \(\mathbf x\) is
a corner-peeling order if and only if
\begin{equation}
 \delta_{\mathbf x}:\Sigma\longrightarrow[m]
 \quad\text{is a bijection}.
 \label{eq:damp-trace-death-bijection}
\end{equation}
Thus corner peelability is an ordering property of the original trace
table: no descendant family has to be constructed in order to state or
check the criterion.
\end{proposition}

\begin{proof}
After deleting \(x_1,\ldots,x_k\), a support \(D\in\Sigma\) is still
shattered precisely when every \(D\)-trace has a representative whose
index is greater than \(k\).  Equivalently,
\begin{equation}
 D\in\operatorname{Sh}(\{x_{k+1},\ldots,x_m\})
 \quad\Longleftrightarrow\quad
 \delta_{\mathbf x}(D)>k.
 \label{eq:damp-trace-death-suffix-shadow}
\end{equation}
If \(\delta_{\mathbf x}\) is bijective, the right-hand side holds for
exactly \(m-k\) supports.  The suffix has \(m-k\) concepts, so equality
in the Sandwich bound makes every suffix ample.  Successive deletion
from an ample family preserves ampleness exactly when the deleted
concept is a corner; hence \(\mathbf x\) is a corner peeling.

Conversely, if \(\mathbf x\) is a corner peeling, every suffix is ample.
Taking cardinalities in \eqref{eq:damp-trace-death-suffix-shadow} gives
\[
 |\{D\in\Sigma:\delta_{\mathbf x}(D)>k\}|=m-k
 \qquad(0\le k\le m).
\]
It follows successively that exactly one support has each death time
\(1,\ldots,m\), which is \eqref{eq:damp-trace-death-bijection}.
\end{proof}

\begin{definition}[Acyclic trace-carrier scheme]
\label{def:damp-acyclic-trace-carrier-scheme}
Let \(H\subseteq Q_F\) be ample and put
\(\Sigma=\operatorname{Sh}(H)\).  For \(D\in\Sigma\) and
\(q\in Q_D\), write
\[
 H(D,q)=\{x\in H:x|_D=q\}
\]
for the corresponding trace fiber.  Every such fiber is nonempty because
\(D\) is shattered.

A \emph{trace-carrier scheme} on \(H\) consists of
\begin{enumerate}
 \item a bijection \(\kappa:\Sigma\to H\); and
 \item for every \(D\in\Sigma\) and every
       \(q\in Q_D\setminus\{\kappa(D)|_D\}\), a protector
       \[
        p(D,q)\in H(D,q).
       \]
\end{enumerate}
The \emph{precedence digraph} \(G(\kappa,p)\) has vertex set \(H\) and
the following arcs, for every \(D\in\Sigma\):
\begin{align}
 y&\longrightarrow\kappa(D)
 &&\bigl(y\in H(D,\kappa(D)|_D)\setminus\{\kappa(D)\}\bigr),
 \label{eq:damp-carrier-precedence}\\
 \kappa(D)&\longrightarrow p(D,q)
 &&\bigl(q\in Q_D\setminus\{\kappa(D)|_D\}\bigr).
 \label{eq:damp-protector-precedence}
\end{align}
The scheme is \emph{acyclic} when \(G(\kappa,p)\) has no directed cycle.
Thus its data are a support--concept bijection and marked concepts in
specified Boolean trace fibers; no total order of \(H\) is part of the
definition.
\end{definition}

\begin{theorem}[Acyclic trace-carrier criterion]
\label{thm:damp-acyclic-trace-carrier-criterion}
Let \(H\subseteq Q_F\) be ample.  Then \(H\) has an ordering with
bijective first-loss map if and only if \(H\) admits an acyclic
trace-carrier scheme.
\end{theorem}

\begin{proof}
Put \(\Sigma=\operatorname{Sh}(H)\) and \(m=|H|=|\Sigma|\).
Suppose first that
\(\mathbf x=(x_1,\ldots,x_m)\) has bijective first-loss map
\(\delta_{\mathbf x}\).  For \(D\in\Sigma\), let
\[
 i_D=\delta_{\mathbf x}(D),
 \qquad
 \kappa(D)=x_{i_D}.
\]
The values \(i_D\) are \(1,\ldots,m\) without repetition, so \(\kappa\)
is a bijection from \(\Sigma\) to \(H\).

Let \(q_D\) be a trace whose last occurrence has index \(i_D\).  Its last
representative is \(x_{i_D}\), and hence
\(q_D=\kappa(D)|_D\).  Every other member of \(H(D,q_D)\) occurs before
\(\kappa(D)\).  For each \(q\ne q_D\), choose as \(p(D,q)\) the last
member of \(H(D,q)\) in the order \(\mathbf x\).  Its index is strictly
larger than \(i_D\): it cannot be smaller by the definition of
\(i_D\), and equality would make the same concept carry two different
\(D\)-traces.  Consequently every arc in
\eqref{eq:damp-carrier-precedence}--\eqref{eq:damp-protector-precedence}
points forward in \(\mathbf x\).  The resulting trace-carrier scheme is
acyclic.

Conversely, let \((\kappa,p)\) be an acyclic trace-carrier scheme and
choose any topological ordering \(\mathbf x\) of its precedence digraph.
Fix \(D\in\Sigma\), put \(z=\kappa(D)\), and write
\(q_D=z|_D\).  The arcs in \eqref{eq:damp-carrier-precedence} place every
other member of \(H(D,q_D)\) before \(z\).  Thus the last occurrence of
the trace \(q_D\) is the position of \(z\).  For every \(q\ne q_D\), the
arc \(z\to p(D,q)\) places a member of \(H(D,q)\) strictly after \(z\).
Therefore
\[
 \delta_{\mathbf x}(D)=\operatorname{pos}_{\mathbf x}(\kappa(D)).
\]
Both \(\kappa:\Sigma\to H\) and the position map
\(H\to[m]\) are bijections.  Hence \(\delta_{\mathbf x}\) is bijective,
as required.
\end{proof}

\begin{remark}[What the trace-carrier criterion removes]
\label{rem:damp-trace-carrier-criterion-scope}
For a specified scheme, all conditions are read directly from the trace
table: fiber membership, bijectivity of \(\kappa\), and acyclicity of a
finite digraph.  The criterion therefore removes the total-order variable.
It does not yet give the desired generative classification of the
obstructions, because one must still characterize those signed
minimal-nonface presentations for which every trace-carrier scheme is
cyclic.  The remaining target is an intrinsic circuit or assembly
description of that failure, rather than an enumeration of schemes.
\end{remark}

\begin{remark}[Protector choices cannot be saturated]
\label{rem:damp-trace-carrier-saturation-fails}
One cannot make the precedence digraph canonical by replacing each chosen
arc
\(\kappa(D)\to p(D,q)\) with arcs from \(\kappa(D)\) to every member of
\(H(D,q)\).  This stronger rule already fails for \(H=Q_2\).  Indeed, in
any topological order of the saturated digraph, the carrier of a singleton
support must occur after the other concept with the same trace and before
both concepts with the opposite trace.  It is therefore forced into the
second position.  The two singleton supports have distinct carriers under
the bijection \(\kappa\), so they cannot both occupy that position.  Thus no
saturated scheme exists, although \(Q_2\) admits a trace-carrier scheme.
The selected protector, rather than the whole competing fiber, is therefore
essential incidence data.
\end{remark}

\begin{corollary}[Order-free Boolean recognition criterion for the ample forbidden basis]
\label{cor:damp-static-signed-clutter-basis}
Let \(\boldsymbol\sigma\) be a signing of
\(\mathcal N(\Sigma)\), and put
\(H=H(\Sigma,\boldsymbol\sigma)\).  Suppose that \(H\subseteq Q_F\) is
irredundantly embedded.  Then \(H\) is an ample forbidden pc-minor for
\(\Damp\) if and only if all three conditions below hold.
\begin{enumerate}
 \item The signing is admissible.
 \item The family \(H\) admits no acyclic trace-carrier scheme.
 \item Each of the \(3|F|\) families
       \(\rho_e^0H,\rho_e^1H,\pi_eH\), for \(e\in F\), admits an
       acyclic trace-carrier scheme for its own shattered-set complex.
\end{enumerate}
Together with the nonample family \(\{Q_n^{--}:n\ge3\}\), these signed
presentations, modulo coordinate permutations and reorientations, give
the full forbidden basis.  The conditions are static: they quantify over
finite trace-fiber data of the original family and its explicit
one-coordinate images, not over total orders, arbitrary proper pc-minors,
or recursive peeling descendants.
\end{corollary}

\begin{proof}
Theorem~\ref{thm:damp-fixed-shadow-signed-clutter} identifies
admissibility with ampleness and fixes the shattered-set complex.
Theorem~\ref{thm:damp-acyclic-trace-carrier-criterion}, together with
Proposition~\ref{prop:damp-trace-death-peeling}, identifies the second
condition with failure of corner peelability and the third with corner
peelability of every elementary pc-minor.  Theorem
\ref{thm:damp-descendant-free-parametrization} reduces pc-minimality to
those elementary images and supplies the quotient by cube symmetries.
Proposition~\ref{prop:nonample-layer} supplies the nonample part.
\end{proof}

\begin{remark}[Recognition is not a generative construction]
\label{rem:damp-recognition-not-construction}
Corollary~\ref{cor:damp-static-signed-clutter-basis} is an exact
recognition criterion, but it does not solve the structural construction
problem for the ample forbidden basis.  Conditions~2 and~3 still ask
whether specified families admit acyclic trace-carrier schemes.  By
Theorem~\ref{thm:damp-acyclic-trace-carrier-criterion}, these are precisely
the original corner-peelability questions with the total-order variable
removed.

Encoding scheme existence as SAT, applying a fixed SAT or resolution
procedure, or adjoining satisfying assignments and refutations as
certificates does not change this point.  Each is a finite decision or
verification format for the same property, not structural data that explain
why the root fails and all its elementary images succeed.  We therefore do
not regard such encodings or certificates as a parametrized construction of
the obstructions.
\end{remark}

\begin{problem}[Intrinsic generative classification]
\label{prob:damp-intrinsic-generative-classification}
Write
\[
 \mathfrak B_{\rm amp}
 =\{H:H\in\Ample\setminus\Damp\text{ and every proper pc-minor of }H
       \text{ belongs to }\Damp\}/\!\cong .
\]
Find a finite list of explicit combinatorial data types
\((\mathcal P_\tau)_{\tau\in T}\), intrinsic admissibility predicates
\(\mathcal A_\tau\), and direct constructions
\(P\mapsto H_\tau(P)\) such that
\[
 \mathfrak B_{\rm amp}
 =\{[H_\tau(P)]:P\in\mathcal P_\tau,
                    \ \mathcal A_\tau(P)\}.
\]
Here each individual parameter is finite, although its size and ambient
coordinate set may vary.  The displayed equality is required to hold for
the following structural reasons.
\begin{enumerate}
 \item \emph{Explicit construction.}
       The concepts and coordinates of \(H_\tau(P)\) are obtained directly
       from the finite combinatorial data \(P\).  Constructing the output
       does not call a dismantlability, corner-peeling, or pc-minimality
       oracle.
 \item \emph{Uniform soundness.}
       Every \(P\) satisfying \(\mathcal A_\tau(P)\) produces an ample
       forbidden pc-minor for \(\Damp\).  Nonpeelability of the root and
       peelability of its proper pc-minors must follow uniformly from the
       structural form of \(P\), rather than from running a recognition
       test on the output.
 \item \emph{Completeness and recovery.}
       Every ample forbidden pc-minor is isomorphic to some
       \(H_\tau(P)\), and its parameter \(P\) is recovered from the
       obstruction, canonically or with explicitly controlled
       multiplicity, by a stated structural decomposition.  In particular,
       completeness is not merely a proof that a suitable parameter exists.
 \item \emph{Non-circular admissibility.}
       The conditions \(\mathcal A_\tau(P)\) are stated in native finite
       data such as signed incidences, supports, carriers, attachment maps,
       or prescribed assembly pieces.  They do not invoke membership in
       \(\Damp\), the existence or nonexistence of a corner peeling,
       first-loss order, shelling order, linear-quotient order, or acyclic
       trace-carrier scheme; nor may they replace one of these tests by
       SAT/UNSAT of an equivalent formula or by an arbitrary proof
       certificate for the same assertion.  In particular, one may not take
       \(P\) to be merely a lossless copy of \(H_\tau(P)\) and hide the
       obstruction predicate in the definition of admissibility.
 \item \emph{Descendant-free verification.}
       Admissibility does not quantify over all proper pc-minors.  The
       construction itself supplies a uniform structural reason why the
       required restrictions and contractions are peelable, for example
       through explicit parameter transformations or peelings prescribed
       directly from \(P\).
\end{enumerate}
These restrictions concern the description of the parameters and the
definition of admissibility.  They do not prevent a soundness proof from
\emph{deriving}, uniformly from \(P\), a canonical corner peeling of an
elementary image.  What is excluded is supplying such a peeling as auxiliary
witness data, searching for one, or making its existence an admissibility
condition.
The nonample family \(\{Q_n^{--}:n\ge3\}\) is already explicit; the problem
concerns the ample part of the basis.

We regard the classification as solved only when the construction,
admissibility theorem, soundness theorem, and completeness-and-recovery
theorem are all stated and proved.  A parametrization of all ample or all
cornerless ample families, a recognition criterion for the forbidden basis,
or a construction proved only relative to an already classified seed is a
component of such a solution, not the solution itself.
\end{problem}

\subsection{One-concept exchanges with fixed shattered sets}

\begin{proposition}[Exchange law for signed minimal nonfaces]
\label{prop:damp-signed-clause-exchange}
Let \(H=H(\Sigma,\boldsymbol\sigma)\) be the ample family of an
admissible signing.  For \(a\in Q_F\setminus H\) and \(c\in H\), put
\begin{align}
 \mathcal V_H(a)
 &=\{I\in\mathcal N(\Sigma):a|_I=\sigma_I\},
 \label{eq:damp-violated-signed-clauses}\\
 \mathcal R_H(a,c)
 &=\left\{I\in\mathcal N(\Sigma):
   \begin{array}{l}
    a|_I\ne c|_I,\\[-2pt]
    \{x\in H:x|_I=c|_I\}=\{c\}
   \end{array}\right\}.
 \label{eq:damp-replaced-private-traces}
\end{align}
Then
\begin{equation}
 (H\setminus\{c\})\cup\{a\}
 \text{ is ample with shattered complex }\Sigma
 \quad\Longleftrightarrow\quad
 \mathcal V_H(a)=\mathcal R_H(a,c).
 \label{eq:damp-exact-signed-exchange-law}
\end{equation}
When these conditions hold, the new signing is
\begin{equation}
 \sigma'_I=
 \begin{cases}
  c|_I,&I\in\mathcal V_H(a),\\
  \sigma_I,&I\notin\mathcal V_H(a).
 \end{cases}
 \label{eq:damp-signed-clause-pivot}
\end{equation}
Moreover,
\begin{equation}
 H\cup\{a\}\text{ is ample}
 \quad\Longleftrightarrow\quad
 |\mathcal V_H(a)|=1.
 \label{eq:damp-one-extension-one-clause}
\end{equation}
In the latter case, if \(\mathcal V_H(a)=\{I\}\), then
\(\operatorname{Sh}(H\cup\{a\})=\Sigma\mathbin{\dot\cup}\{I\}\).
\end{proposition}

\begin{proof}
For every \(I\in\mathcal N(\Sigma)\), the trace set
\(H{\upharpoonright}I\) is
\(Q_I\setminus\{\sigma_I\}\).  On replacing \(c\) by \(a\), this
projection still omits exactly one trace in precisely the following two
cases: either \(a|_I=\sigma_I\) and \(c\) is the unique old
representative of the distinct trace \(c|_I\), or
\(a|_I\ne\sigma_I\) and removal of \(c\) does not create a second
missing trace.  Simultaneously for all minimal nonfaces, these cases are
exactly the equality
\(\mathcal V_H(a)=\mathcal R_H(a,c)\).  They also give
\eqref{eq:damp-signed-clause-pivot}.

If the equality holds, the exchanged family omits one trace on every
minimal nonface, so its shattered complex is contained in \(\Sigma\).
It has \(|\Sigma|\) concepts.  The Sandwich inequality therefore forces
equality and ampleness.  Conversely, an ample exchanged family with
shadow \(\Sigma\) omits exactly one trace on every minimal nonface, so
the preceding two-case analysis forces the displayed equality.  This
proves \eqref{eq:damp-exact-signed-exchange-law}.

The set \(\mathcal V_H(a)\) is nonempty because
Theorem~\ref{thm:damp-fixed-shadow-signed-clutter} says that \(H\) is
exactly the common satisfying family.  Adjoining \(a\) shatters every
member of \(\mathcal V_H(a)\).  Hence, if \(H\cup\{a\}\) is ample, its
shadow has exactly one more member than \(\Sigma\), and
\(|\mathcal V_H(a)|=1\).

Conversely suppose \(\mathcal V_H(a)=\{I_0\}\), and let a support
\(J\notin\Sigma\) become shattered after adjoining \(a\).  Since a
single trace was added,
\(H{\upharpoonright}J=Q_J\setminus\{a|_J\}\).  Choose a minimal nonface
\(I\subseteq J\).  If \(I\subsetneq J\), there is a trace on \(J\)
that extends \(\sigma_I\) but differs from \(a|_J\); it cannot be
realized by \(H\), a contradiction.  Thus \(J=I\), and necessarily
\(a|_I=\sigma_I\).  Therefore \(J=I_0\), and
\[
 \operatorname{Sh}(H\cup\{a\})=\Sigma\mathbin{\dot\cup}\{I_0\}.
\]
Its cardinality equals \(|H|+1\), so the Sandwich equality proves
ampleness and completes the proof.
\end{proof}

\begin{proposition}[Fiber decomposition of the punctured-face parametrization]
\label{prop:damp-punctured-face-fibers}
Let \(\mathcal M(S)\) be the minimal nonfaces of
\(\operatorname{Sh}(S)\).  By
Theorem~\ref{thm:damp-fixed-shadow-signed-clutter}, every
\(I\in\mathcal M(S)\) has a unique omitted trace; denote it by
\(\sigma_I\).  Define
\begin{equation}
\mathcal Y_I(S)
 =\left\{y\in Q_{F\setminus I}:
 \begin{array}{ll}
  (\sigma_I^U,y)\in S
    &\text{for every }\varnothing\ne U\subseteq I,\\
  (\sigma_I,y^{\{e\}})\notin S
    &\text{for every }e\in F\setminus I
 \end{array}\right\}.
 \label{eq:damp-punctured-face-outside-fibers}
\end{equation}
Then
\begin{equation}
\mathcal P(S)
 =\mathop{\dot\bigcup}_{I\in\mathcal M(S)}
   \{((\sigma_I,y),I):y\in\mathcal Y_I(S)\}.
 \label{eq:damp-punctured-face-fiber-decomposition}
\end{equation}
In particular, all punctured-face data with support \(I\) have the same
missing trace on \(I\).  Their remaining datum is the outside assignment
\(y\): it locates a punctured \(I\)-cube and requires the missing tip to
have no neighbor outside \(I\).
\end{proposition}

\begin{proof}
Take \((a,I)\in\mathcal P(S)\).  The punctured-cube condition realizes
every \(I\)-trace except possibly \(a|_I\), while
\(I\notin\operatorname{Sh}(S)\) says that at least one trace is missing.
Thus \(a|_I\) is the unique missing trace \(\sigma_I\).  Put
\(y=a|_{F\setminus I}\).  The punctured-cube condition is the first line
of \eqref{eq:damp-punctured-face-outside-fibers}.  Since
\(I=I_S(a)\), no coordinate outside \(I\) is incident with \(a\), which
is its second line.

Conversely, take \(I,\sigma_I,y\) on the right-hand side of
\eqref{eq:damp-punctured-face-fiber-decomposition} and put
\(a=(\sigma_I,y)\).  The missing trace ensures \(a\notin S\).
The two lines in
\eqref{eq:damp-punctured-face-outside-fibers} say respectively that the
punctured \(I\)-cube around \(a\) lies in \(S\) and that
\(I_S(a)=I\).  Since \(I\) is a nonface, all conditions in
\eqref{eq:damp-punctured-face-atlas} hold.  The two constructions are
inverse.
\end{proof}

\begin{proposition}[Adjacent square completions]
\label{prop:damp-square-pair-invariance}
Let \(G\subseteq Q_F\) be nonempty and ample.  Suppose that
\(c,u\notin G\) are adjacent and that both individual additions
are ample with the same acquired support \(I\), where \(|I|=2\).
Then \(H=G\cup\{c,u\}\) is ample and
\[
 H\in\Damp\quad\Longleftrightarrow\quad G\in\Damp.
\]
\end{proposition}

\begin{proof}
Write \(I=\{p,q\}\) and \(u=c^{\{f\}}\), so \(f\notin I\).
The two punctured squares in \(G\), together with \(c\), form the
punctured \(I\cup\{f\}\)-cube at \(u\).  This support is absent
from \(\operatorname{Sh}(G\cup\{c\})
=\operatorname{Sh}(G)\cup\{I\}\), so
Lemma~\ref{lem:damp-one-concept-extension-test} makes \(H\) ample.  If \(G\) peels, delete \(u,c\) in that order
and then peel \(G\).

For the converse, suppose that \(H\) peels and \(G\) does not,
and choose such a counterexample with \(|G|\) minimal.  Put
\[
 C=Q_{I\cup\{f\}}(c),\qquad
 x=c^{\{p,q\}},\quad y=u^{\{p,q\}},\qquad
 W=\{c^{\{p\}},c^{\{q\}},u^{\{p\}},u^{\{q\}}\}.
\]
Thus \(G\cap C=C\setminus\{c,u\}\).
Neither added tip has a neighbor outside \(C\).  Moreover, a corner
of \(H\) lying in \(C\) has no such neighbor: its incident cube
already contains all directions of \(C\), and an additional direction
\(e\) would force the absent concept \(c^{\{e\}}\).

Let \(z\) be the first deletion in a peeling of \(H\).
It cannot be \(c\) or \(u\), because either deletion leaves a
nonpeelable square completion of \(G\), by
Proposition~\ref{prop:damp-rank-two-extension-invariance}.
If \(z\notin C\), its corner cube avoids both added tips, whose
unique maximal cube is \(C\).  Thus \(z\) is also a corner of
\(G\).  The ample family \(G'=G\setminus\{z\}\) is nonpeelable,
since otherwise deleting \(z\) first would peel \(G\).
Both punctured squares remain intact in \(G'\), and both tips still
acquire \(I\).  The peeling of \(H\setminus\{z\}\) therefore
contradicts the minimal choice of \(|G|\).

If \(z\in W\), interchange the two layers and the names \(p,q\)
as necessary to write \(z=c^{\{p\}}\).  Since \(z\) has no
outside neighbor, its incident \(\{q,f\}\)-square lies in \(G\),
so it is a corner of \(G\).  Again \(G'=G\setminus\{z\}\)
is ample and nonpeelable.  Adding \(u\) to \(G'\) completes its
unchanged \(I\)-square, whose support is still unshattered because
\(G'\subseteq G\).  The subsequent addition of \(c\) gives
\(H\setminus\{z\}\) and has incident support \(\{q,f\}\).
Both steps are ample and acquire supports of size two.  Applying
Proposition~\ref{prop:damp-rank-two-extension-invariance} twice
contradicts the peelability of \(H\setminus\{z\}\).

Only \(z=x\) or \(z=y\) remains.  Exchanging the layer names
handles the latter case, so assume \(z=x\), and let \(w\) be
the second deletion.  The tips \(c,u\) retain every direction of
\(C\) but their cube lacks \(x\), so neither is a corner.
If \(w\notin C\), then \(w\) was already a corner of \(H\):
it is not adjacent to \(x\), and deleting a nonneighbor cannot
create a corner.  Deleting \(w\) first leaves the entire corner
cube of \(x\) intact.  Swapping the first two deletions therefore
gives a peeling starting outside \(C\), which was just excluded.

Suppose instead that \(w\in W\).  It retains its neighbor in
\(\{c,u\}\), so any outside neighbor would still give a missing
square through that tip.  Hence \(w\) has no outside neighbor.
It must also be adjacent to \(x\), since otherwise all its cube
directions remain incident while \(x\) is missing.  Thus \(x,w\)
are adjacent corners of the full cube, both without outside neighbors.
Deleting either leaves the other as a corner of the opposite face.
Swapping these two deletions gives a peeling starting in \(W\),
which was also excluded.

The only remaining choice is \(w=y\).  Its full \(I\)-square
contains \(u\), so an outside neighbor would require an outside
neighbor of \(u\).  Therefore \(y\), like \(x\), has no
neighbor outside \(C\).  The square automorphism interchanging
\(c,x\) and fixing \(c^{\{p\}},c^{\{q\}}\), applied in both
\(f\)-layers, interchanges the edges \(\{c,u\}\) and \(\{x,y\}\)
and fixes \(W\).  It extends by the identity outside \(C\), since
neither exchanged edge has outside neighbors.  Consequently
\(H\setminus\{x,y\}\cong G\), contradicting the remaining
peeling.  All first and second deletions have been accounted for.
\end{proof}

The checker \path{verify_damp_square_pair_invariance.py} exercises
these reductions on three ample bases with corners.  Each has a proper
restriction equal to the certified cornerless \(291\)-concept family.
It checks the two support-two additions after a common corner, both
commutations, and the opposite-edge graph isomorphism.  These fixed
controls are nonminimal and require no peeling search.


\begin{proposition}[Resolving additions and reductions]
\label{prop:damp-repair-strong-projections}
Let \(S\subseteq Q_F\) be nonempty and ample, and suppose that
\(H=S\cup\{a\}\) is an ample one-concept extension with acquired
support \(I\).  For \(D\subseteq F\), put
\(a_D=a|_{F\setminus D}\).  Then
\begin{equation}
 H^D=
 \begin{cases}
  S^D\cup\{a_D\},&D\subseteq I,\\
  S^D,&D\not\subseteq I.
 \end{cases}
 \label{eq:damp-repair-strong-projection}
\end{equation}
In the first case the added concept is new and has acquired support
\(I\setminus D\); if \(S^D\) is empty, the extension is a singleton
and this support is empty.

Consequently, if \(H\in\Damp\) and \(S^D\notin\Damp\), then
\begin{equation}
 D\subseteq I,\qquad |I\setminus D|\ge3.
 \label{eq:damp-repair-nonpeelable-projection-support}
\end{equation}
In particular, a resolving addition with \(|I|=3\) forces every
reduction of \(S\) on a nonempty coordinate set to belong
to \(\Damp\).
\end{proposition}

\begin{proof}
The only full \(D\)-fibre that the addition can create is the fibre
through \(a\).  If \(D\subseteq I\), all its other concepts belong
to the punctured \(I\)-cube in \(S\), so that fibre becomes full.
Its projection \(a_D\) was absent from \(S^D\), because its old
fibre lacked \(a\).  If instead \(e\in D\setminus I\), the same
fibre still lacks \(a^{\{e\}}\), which is absent by the definition
of \(I=I_S(a)\).  This proves
\eqref{eq:damp-repair-strong-projection}, including \(D=\varnothing\).

Assume \(D\subseteq I\).  For each \(e\in I\setminus D\), the
entire \(D\)-fibre through \(a^{\{e\}}\) belongs to \(S\), so
\(e\) is incident with \(a_D\) in the extension.  No
\(e\in F\setminus I\) can be incident: its fibre would require
\(a^{\{e\}}\in S\).  The incident support is therefore exactly
\(I\setminus D\).  Reductions preserve ampleness, so
Lemma~\ref{lem:damp-one-concept-extension-test} identifies this as
the acquired support whenever \(S^D\ne\varnothing\).  If \(S^D\)
is empty, there are no incident directions, and the singleton acquires
only the empty support.

Now suppose that \(H\) peels and \(S^D\) does not.  To see directly
that \(H^D\) peels, reverse a peeling of \(H\) to an ample prefix
order, and order the full \(D\)-fibres by their completion times.
The resulting prefixes are reductions of ample prefixes and
hence are ample.  Reversing this order peels \(H^D\).
Thus the two reductions cannot be equal,
which forces \(D\subseteq I\).  Their difference is a resolving
one-concept addition of support \(I\setminus D\), and
Proposition~\ref{prop:damp-rank-two-extension-invariance} gives
\(|I\setminus D|\ge3\).  For \(|I|=3\), this is impossible when
\(D\ne\varnothing\), proving the last assertion.
\end{proof}

The prescribed-endpoint analogue of this projection argument is false,
even when the root has a full fibre; see
Proposition~\ref{prop:damp-rooted-strong-projection-failure}.



\begin{corollary}[Deletion of adjacent degree-three corners]
\label{cor:damp-degree-three-corner-pair-descent}
Let \(H\subseteq Q_F\) be an irredundantly embedded ample forbidden
pc-minor.  If \(c,u\) are adjacent corners of degree three, then
\[
 H\setminus\{c\},\qquad H\setminus\{u\},\qquad
 H\setminus\{c,u\}
 \quad\text{belong to }\Obs(\Damp)\cap\Ample.
\]
In particular, a degree-three corner of a corner-descent-terminal ample
forbidden pc-minor has no corner as a neighbor.
\end{corollary}

\begin{proof}
The adjacent corners have the same unique maximal cube \(C\), because
their edge belongs to a maximal cube containing both.  Write \(P\)
for its three-coordinate support, \(u=c^{\{f\}}\), and
\(I=P\setminus\{f\}\).  Neither corner has an outside neighbor.
Deleting \(c\) leaves \(u\) as a corner of its \(I\)-face, so
\(G=H\setminus\{c,u\}\) is ample.  Each individual addition to
\(G\) acquires \(I\).  Proposition~\ref{prop:damp-square-pair-invariance}
makes \(G\) nonpeelable.  Its embedding remains irredundant: the
punctured cube \(C\setminus\{c,u\}\) retains every direction of
\(P\), and a surviving cube vertex has the same outside coordinates
as the two deleted concepts.

We check every elementary pc-minor \(\mu G\), using
\(\mu H\in\Damp\).  Under a restriction in coordinate \(e\),
either both tips disappear, only one is retained, or both are retained.
The single-tip case occurs when \(e=f\), and its acquired support is
\(I\).  If both survive, their individual acquired supports are
\(I\setminus\{e\}\), and they remain adjacent in direction \(f\).
All these intermediate images are ample, being images of
\(G,G\cup\{c\},G\cup\{u\},H\).

For a contraction in \(e\in I\), each tip is identified with a
neighbor in \(G\), so \(\mu H=\mu G\).  Contraction in \(f\)
identifies the two tips with each other and adds one concept to
\(\pi_fG\), with acquired support \(I\).  A contraction outside
\(P\) retains two distinct new tips: neither has a neighbor in
that direction.  Their individual acquired supports remain \(I\),
since contraction restricts shattered supports to those avoiding \(e\)
and hence preserves the one-support gain
\eqref{eq:damp-one-extension-shadow}.

Thus \(\mu H\) either equals \(\mu G\), adds one concept with
support at most two, or adds an adjacent pair with common individual
support of size one or two.  In the single-concept case,
Proposition~\ref{prop:damp-rank-two-extension-invariance} makes
\(\mu G\) peelable.  For a pair with common support of size two,
Proposition~\ref{prop:damp-square-pair-invariance} does the same.
For a pair with common support of size one, the two successive
additions have supports of sizes one and two, so the one-concept
invariance applies twice.  The equal-image case is immediate.
Every elementary pc-minor of \(G\) therefore peels, proving
\(G\in\Obs(\Damp)\).

The two single deletions from \(H\) are nonpeelable: a peeling of
either could be prefixed by the deleted corner to peel \(H\).
They are ample one-concept extensions of the irredundant forbidden
family \(G\), so Proposition~\ref{prop:damp-one-concept-atlas}
makes both forbidden.  Finally, adjacent corners necessarily share
their maximal cube and hence their degree.  A corner neighbor of a
degree-three corner would therefore give an obstruction-preserving
deletion, contradicting corner-descent terminality.
\end{proof}


\begin{proposition}[Same-shadow mutations]
\label{prop:damp-same-shadow-corner-exchange}
Let \(S\in\Obs(\Damp)\cap\Ample\) be cornerless, let
\((a,I)\in\mathcal P(S)\), and put \(A=S\cup\{a\}\).  Define the set of
\emph{exposed tips} of this punctured face by
\begin{equation}
 \mathcal E_S(a,I)
 =\{a^U:\varnothing\ne U\subseteq I,\quad
           a^{U\triangle\{e\}}\notin S
           \text{ for every }e\in F\setminus I\}.
 \label{eq:damp-exposed-punctured-face-tips}
\end{equation}
Then
\begin{equation}
 \operatorname{Cor}(A)
 =\{a\}\mathbin{\dot\cup}\mathcal E_S(a,I).
 \label{eq:damp-one-extension-corners}
\end{equation}
For every \(c\in\mathcal E_S(a,I)\), the exchanged family
\[
 T_c=(S\setminus\{c\})\cup\{a\}=A\setminus\{c\}
\]
is ample and has the same shattered-set complex as \(S\).  Moreover,
\begin{equation}
 (a,I)\in\mathcal P_{\mathrm{peel}}(S)
 \quad\Longleftrightarrow\quad
 T_c\in\Damp\text{ for some }c\in\mathcal E_S(a,I).
 \label{eq:damp-same-shadow-exchange-criterion}
\end{equation}

No exposed tip is at Hamming distance one from \(a\).  If \(|I|=2\),
then every exposed tip \(c\) satisfies \(T_c\cong S\), and hence
\((a,I)\in\mathcal P_{\mathrm{obs}}(S)\).  Consequently every peelable
punctured-face datum over a cornerless forbidden family has
\begin{equation}
 |I|\ge3.
 \label{eq:damp-cornerless-repair-rank-three}
\end{equation}
\end{proposition}

\begin{proof}
Lemma~\ref{lem:damp-one-concept-extension-test} makes \(a\) a corner of
\(A\), with unique maximal cube the completed \(I\)-cube.  Let
\(c\in S\) be any other corner of \(A\).  Its unique maximal cube in
\(A\) must contain \(a\).  Indeed, if it avoided \(a\), it would remain
a cube of \(S\) containing every cube of \(S\) through \(c\), making
\(c\) a corner of the cornerless family \(S\).  The maximal cube through
\(c\) is therefore the completed \(I\)-cube.  Thus
\(c=a^U\) for a nonempty \(U\subseteq I\), and \(c\) has no neighbor in
a direction outside \(I\).  This proves that the left-hand side of
\eqref{eq:damp-one-extension-corners} is contained in the right-hand
side.

Conversely, take \(c\in\mathcal E_S(a,I)\).  The completed \(I\)-cube
contains every neighbor of \(c\) in a direction in \(I\), and the
definition of an exposed tip excludes every other direction.  Hence this
cube is the unique maximal cube of \(A\) through \(c\), so \(c\) is a
corner.  This proves \eqref{eq:damp-one-extension-corners}.

Deleting the corner \(c\) from the ample family \(A\) gives the ample
family \(T_c\).  The unique maximal cube of \(c\) has support \(I\), so
the shattered-support form of corner deletion gives
\[
 \operatorname{Sh}(T_c)
 =\operatorname{Sh}(A)\setminus\{I\}
 =\operatorname{Sh}(S),
\]
where the second equality is
\eqref{eq:damp-one-extension-shadow}.  This proves the same-shadow
assertion.

A nonempty ample family is corner-peelable precisely when some corner can
be removed leaving a corner-peelable family.  Removing \(a\) from \(A\)
leaves the nonpeelable family \(S\); by
\eqref{eq:damp-one-extension-corners}, all other possible first vertices
are exactly the exposed tips.  This proves
\eqref{eq:damp-same-shadow-exchange-criterion}.

Suppose that \(a^{\{e\}}\) were an exposed tip.  In \(S\), its incident
directions would be exactly \(I\setminus\{e\}\), and the punctured
\(I\)-cube contains the entire cube on those directions.  It would
therefore be a corner of \(S\), a contradiction.  Hence exposed tips
have Hamming distance at least two from \(a\).

Finally let \(I=\{p,q\}\).  The only possible exposed tip is
\(c=a^{\{p,q\}}\).  In \(A\), the vertices \(a\) and \(c\) have the
same two neighbors \(a^{\{p\}},a^{\{q\}}\) and no others.  Interchanging
\(a\) and \(c\) is consequently a graph automorphism of \(A\), and it
maps \(T_c=A\setminus\{c\}\) to \(S=A\setminus\{a\}\).  Thus
\(T_c\cong S\notin\Damp\).  The exchange criterion and the parametrization
dichotomy in Proposition~\ref{prop:damp-one-concept-atlas} now give the
last two assertions.
\end{proof}

\begin{proposition}[Corner elimination by a maximum-class mutation]
\label{prop:damp-maximum-mutation-corner-elimination}
Let $C,C'\subseteq Q_n$ be maximum classes of VC dimension $d$, with
$C'=(C\setminus\{r\})\cup\{a\}$, where $r\in C$ and $a\notin C$.
For a cube vertex $v$, let $N(v)$ be its Hamming-distance-one neighbors,
and put $\deg_C(v)=|N(v)\cap C|$ for $v\in C$. Define
\[
 U=(C\setminus\{r\})\cap(N(r)\setminus N(a)),\qquad
 V=(C\setminus\{r\})\cap(N(a)\setminus N(r)).
\]
The incident coordinate directions at $r$ in $C$ and at $a$ in $C'$ agree.
Writing their common degree as $q$, one has
\[
 |\operatorname{Cor}(C')|-|\operatorname{Cor}(C)|
 =|\{v\in U:\deg_C(v)=d+1\}|
  -|\{v\in V:\deg_C(v)=d\}|.
\]
More precisely, $C'$ is cornerless if and only if
\[
 q>d,\qquad \operatorname{Cor}(C)\subseteq V,\qquad
 \{v\in U:\deg_C(v)=d+1\}=\varnothing.
\]
In particular, if $C$ has exactly two corners, a mutation to a cornerless
maximum class is possible only when those corners have Hamming distance two.
These conditions test a legal mutation; they do not assert that the
exchanged class is maximum.
\end{proposition}
\begin{proof}
We first recall two properties of maximum classes and justify their use.
Every reduction on a support $S$ of size at most $d$ is maximum of VC
dimension $d-|S|$ on the remaining coordinates. For one coordinate this
follows from $|C|=|\pi_e C|+|C^{\{e\}}|$: the two Sauer bounds, of ranks
$d$ and $d-1$, sum to the maximum cardinality of $C$, so both are attained.
Iteration gives the assertion for $S$. Consequently each $d$-support has
exactly one full cube. Also every maximal cube has dimension $d$. Otherwise
a maximal $S$-cube with $|S|<d$ gives an isolated vertex in $C^S$, although
that reduction is connected and has more than one vertex. Reductions are
ample and hence connected by \cite[Section~3.1]{ChalopiChepoiMoran22}.
Thus every vertex and edge lies in a $d$-cube, and a vertex is a corner
exactly when its degree is $d$.

For each $d$-support, its unique cube changes precisely when it contains
$r$; its replacement must contain $a$. Conversely every cube through $a$
replaces a cube through $r$, since a cube avoiding the changed vertex
survives the exchange and the parallel cube is unique. The sets of supports
of incident $d$-cubes at the two exchanged vertices are therefore equal.
Taking their unions proves equality of the incident direction sets and of
$q$. In particular, $r$ and $a$ contribute equally to the corner counts.

Among surviving vertices, exactly those in $U$ lose one neighbor and
exactly those in $V$ gain one; every other degree is unchanged. All degrees
in either maximum class are at least $d$. Hence new corners among surviving
vertices are exactly the degree-$(d+1)$ vertices of $U$, and corners lost
are exactly the degree-$d$ vertices of $V$. This proves the formula and,
by also excluding a corner at $a$, the stated equivalence. If there are two
old corners and $C'$ is cornerless, both lie in $V\subseteq N(a)$. Two
distinct neighbors of a cube vertex differ in exactly two coordinates.
\end{proof}

The two corners of the recorded $232$-vertex Hall projection are $1600$
and $1894$ in its binary integer encoding, at Hamming distance four.
The verified three-mutation paths change their pair to either
$\{1537,1894\}$ or $\{1600,2023\}$; performing both paths gives
$\{1537,2023\}$. Each of these pairs has distance six. The proposition
therefore excludes a cornerless one-mutation continuation from each of
these endpoints, without a further search of its exchanges. It does not
exclude a longer path that first brings the corners to distance two.
The paths and independent trace/cube checks are recorded in
\texttt{papers/ample-corners/evidence/maximum\_two\_corner\_component/}.

The distance-two requirement is attained by actual corner-eliminating
mutations. In the recorded binary encoding of Hall's $299$-vertex
cornerless maximum class $C_H$, put
$C=(C_H\setminus\{2\})\cup\{2049\}$. This is maximum and its two
corners are $1026$ and $2050$, both neighbors of $2$. The reverse exchange
removes the degree-four vertex $2049$ and adds $2$, returning $C_H$
without creating any new corner. The finite verification is recorded in
\texttt{papers/ample-corners/evidence/hall\_reverse\_mutation\_criterion/}.
This example reuses Hall's class; it does not improve the obstruction
order bound.

\begin{proposition}[Two external section types for the marked four-cube]
\label{prop:damp-final-mutation-two-types}
Let $X\subseteq Q_{\{1,2,3,4\}\sqcup J}$ be a cornerless maximum
class of VC dimension three. Encode the first four coordinates by integers
$0,\ldots,15$, and suppose that the face with all $J$-coordinates zero is
$P=\{0,\ldots,15\}\setminus\{7\}$. Suppose also that the three vertices
$(0,0),(1,0),(8,0)$ have no neighbors in directions in $J$.
For $j\in J$, let $R_j$ be the projection onto the first four coordinates
of the section $X\cap\{x_j=1\}$. Then exactly one of the following holds:
\[
 R_j=Q_4\setminus\{0,1,3,7,8\},\qquad
 R_j=Q_4\setminus\{0,1,5,7,8\}.
\]
Let $A$ and $B$ partition $J$ according to these two alternatives, in the
displayed order. Then $|A|,|B|\ge2$. The fiber over pattern $3$ uses only
coordinates in $B$, and the fiber over pattern $5$ uses only coordinates
in $A$. In particular, for $|J|=7$ the type sizes are $2+5$ or $3+4$,
up to exchanging $A$ and $B$ and permuting coordinates within each type.
If $|B|=2$, the fiber over $3$ is the full $B$-square, with all
$A$-coordinates zero, and both fibers over $2$ and $11$ contain that same
$B$-square. Symmetrically, if $|A|=2$, the fiber over $5$ is the full
$A$-square and the fibers over $4,13$ contain it.
These are necessary conditions on the marked configuration, not an
existence or sufficiency assertion.
\end{proposition}
\begin{proof}
The projection onto the first four coordinates is $P$: the base face
already realizes its fifteen patterns, and realizing $7$ anywhere would
shatter four coordinates. Each of the fibers over $0,1,8$ is an ample
coordinate face, hence connected. Its zero vertex has no neighbor in that
fiber, so it is its only vertex.

Project $X$ onto the first four coordinates and $j$. Projections of a
maximum rank-three class are maximum of the same rank when at least three
coordinates remain, by the Sauer-bound equality argument in the preceding
proposition. This five-coordinate projection has $26$ vertices. Its zero
$j$-section is $P$, of size $15$, and its one section $R_j$ is a subset
of $P$. Thus $R_j$ has $11$ vertices and is also the $j$-reduction, a
maximum rank-two class on four coordinates. The singleton fibers exclude
$0,1,8$ from $R_j$, and $7$ is excluded everywhere. Its complement is
therefore $\{0,1,7,8,t\}$ for one further pattern $t$. This complement
is ample and connected. The vertices $0,1,8$ form a connected component
before $t$ is added, whereas $7$ has no neighbor among them. Hence $t$
must neighbor both $7$ and one of $0,1,8$. Their common-neighbor lists
leave precisely $t=3$ and $t=5$.

The fiber-support assertions follow immediately from the definitions of
$A,B$. At the base vertex $(3,0)$, the three horizontal neighbors have
patterns $1,2,11$. Since a cornerless maximum rank-three class has minimum
degree at least four, this vertex has an external neighbor, necessarily
in $B$. If $B$ had only one coordinate $j$, the vertex $(3,e_j)$ would
have at most one external neighbor. Its only possible horizontal
neighbors would have patterns $2,11$: the $1$-fiber is a singleton and
the $7$-fiber is empty. Its total degree would be at most three, a
contradiction. Thus $|B|\ge2$. Applying the same argument at pattern $5$,
whose nonzero-fiber horizontal neighbors can only have patterns $4,13$,
gives $|A|\ge2$. Finally, interchanging the first-block directions with
weights $2$ and $4$ interchanges the two types while preserving all the
hypotheses. This gives the asserted two size cases when $|J|=7$.

For the additional assertion, let $|B|=2$. The fiber over $3$ contains
its zero vertex and at least one other vertex, as proved above. At each
nonzero vertex it has at most two horizontal neighbors, so cornerlessness
forces both possible $B$-neighbors and both horizontal neighbors to be
present. A subset of a square containing zero and a nonzero vertex whose
nonzero vertices all have degree two must be the full square. The two
horizontal neighbors have patterns $2$ and $11$, proving their asserted
containments; at external word zero these follow from the base face.
The symmetric argument gives the other assertion.
\end{proof}

\begin{remark}[What a leaf of the double reduction does and does not force]
\label{rem:damp-final-mutation-tree-leaves}
Reducing a maximum rank-three class $X$ in a two-coordinate set $D$
gives a maximum rank-one class $T=X^D$, hence a tree. Let $y$ be a leaf
of $T$ and let $k$ be its unique incident direction. Its fiber is a full
square $Q_D$, and every vertex of that square has the two $D$-neighbors
and its $k$-neighbor in $X$. For each remaining direction
$j\notin D\cup\{k\}$, put
\[
 E_j=\{u\in Q_D:(y^{\{j\}},u)\in X\}.
\]
Each $E_j$ is an ample coordinate face, possibly empty, and is a proper
subset of $Q_D$: equality would give a second neighbor of $y$ in $T$.
Nevertheless the square over $y$ contains no corner of $X$ precisely when
\[
 \bigcup_{j\notin D\cup\{k\}} E_j=Q_D.
\]
Indeed, a vertex of a maximum rank-three class is a corner exactly when
its degree is three. The displayed union records exactly which square
vertices have a neighbor in addition to the three already specified.

The recorded Hall controls realize such coverage even when $T$ is a
path: all twelve two-corner reverse-mutation configurations have an
$11$-vertex path as their double reduction, and at each of its two leaves
two proper sets $E_j$ already cover the square. The exact sets and trees
are recorded in
\texttt{papers/ample-corners/evidence/final\_mutation\_tree/}.
Thus a tree, or even a path, does not by itself force a lifted corner.
A dimension-sensitive exclusion of the $2+5$ case must constrain these
sets simultaneously using the rest of $X$; the Hall examples have type
sizes $2+6$ and do not settle that exclusion.
\end{remark}

\begin{remark}[Relation with oriented-matroid mutations]
\label{rem:damp-same-shadow-mutation-terminology}
For a uniform oriented matroid, replacing a simplicial tope and its
antipode by the two missing opposite tips is called a mutation, and the
resulting mutation graph is a well-studied object
\cite[Section~3]{KnauerMarc23corners}.  The operation in Proposition
\ref{prop:damp-same-shadow-corner-exchange} is the one-sided ample
analogue: it replaces one exposed tip of a completed cube by the opposite
missing tip and preserves the shattered-set complex.  We therefore call
it a \emph{same-shadow mutation}.  Our exchange graph is more general
than the oriented-matroid mutation graph: neither antipodality nor
uniformity is assumed, and its relevant local ranks need not equal the
VC dimension.  The known connectivity theorems for rank-three oriented
matroids consequently do not settle the shellability question here.
\end{remark}


\section{Layered extensions and irreducible families}
\label{sec:damp-two-repair}

We first give an exact trace formula for arbitrary marked layer systems.
For the natural independent construction---one new concept in each
selected layer over a common seed---we then determine ampleness, the
signed presentation, and all inverse parameters.  Allowing repair tips to
persist through arbitrary layer families yields a labelled occurrence-cover
normal form in which ampleness and the signed clauses separate exactly.
Pc-minimality makes the singleton-occurrence specialization collapse to
the two-extension lift over one new coordinate.  We finish by reducing
every ample nonpeelable family to a cornerless, periphery-free,
Cartesian-prime family.

\begin{proposition}[Trace transport for an arbitrary layer system]
\label{prop:damp-layer-system-trace-transport}
Let \(E\) and \(F\) be disjoint coordinate sets.  For each
\(t\in Q_E\), choose a section \(H_t\subseteq Q_F\), and put
\[
 X=\mathop{\dot\bigcup}_{t\in Q_E}(H_t\times\{t\})
 \subseteq Q_{F\cup E}.
\]
For any cube family \(Y\), support \(D\), and trace \(q\in Q_D\), write
\[
 Y[D,q]=\{y\in Y:y|_D=q\}.
\]
For \(L\subseteq E\) and \(r\in Q_L\), define the projected union of
the layers with \(L\)-trace \(r\) by
\begin{equation}
 H_{L,r}=\bigcup_{\substack{t\in Q_E\\t|_L=r}}H_t.
 \label{eq:damp-layer-fiber-union}
\end{equation}
Then, for \(D\subseteq F\), \(q\in Q_D\), \(L\subseteq E\), and
\(r\in Q_L\),
\begin{equation}
 X[D\cup L,(q,r)]
 =\mathop{\dot\bigcup}_{\substack{t\in Q_E\\t|_L=r}}
   \bigl(H_t[D,q]\times\{t\}\bigr).
 \label{eq:damp-layer-system-trace-fiber}
\end{equation}
Consequently,
\begin{equation}
 D\cup L\in\operatorname{Sh}(X)
 \quad\Longleftrightarrow\quad
 D\in\bigcap_{r\in Q_L}\operatorname{Sh}(H_{L,r}).
 \label{eq:damp-layer-system-shattering}
\end{equation}
Conversely, the marked layer set \(E\) recovers every parameter \(H_t\)
as the projected \(t\)-section of \(X\).  Thus
\eqref{eq:damp-layer-system-trace-fiber} gives an exact, canonically
recoverable trace table for any finite number of binary assembly layers.
\end{proposition}

\begin{proof}
A vertex \((x,t)\in X\) has \((D\cup L)\)-trace \((q,r)\) precisely
when \(t|_L=r\) and \(x|_D=q\).  Partitioning the corresponding trace
fiber by the full layer word \(t\) proves
\eqref{eq:damp-layer-system-trace-fiber}.  The support \(D\cup L\) is
shattered precisely when this fiber is nonempty for every
\((q,r)\in Q_D\times Q_L\).  For fixed \(r\), that condition says
exactly that \(H_{L,r}\) shatters \(D\), proving
\eqref{eq:damp-layer-system-shattering}.  The recovery statement is
the definition of a section.
\end{proof}

\begin{proposition}[Independent repair layers]
\label{prop:damp-independent-repair-layers}
Let \(E\ne\varnothing\) and \(F\) be disjoint coordinate sets, let
\(S\subseteq Q_F\) be ample, and let \(T\subseteq Q_E\).  For every
\(t\in T\), choose a vertex \(a_t\in Q_F\setminus S\), with the vertices
\((a_t)_{t\in T}\) pairwise distinct, such that

\[
 A_t=S\cup\{a_t\}
\]

is ample.  Write

\begin{equation}
 \operatorname{Sh}(A_t)\setminus\operatorname{Sh}(S)=\{P_t\},
 \qquad
 P_t=\{f\in F:a_t^{\{f\}}\in S\}.
 \label{eq:damp-independent-repair-supports}
\end{equation}

Place one repair tip in each selected layer and no repair tip in the
other layers:

\begin{equation}
 X_T=(S\times Q_E)\cup\{(a_t,t):t\in T\}.
 \label{eq:damp-independent-repair-assembly}
\end{equation}

Then

\begin{equation}
 X_T\text{ is ample}
 \quad\Longleftrightarrow\quad
 t\longmapsto P_t\text{ is injective on }T.
 \label{eq:damp-independent-repair-ample-iff}
\end{equation}

Assume henceforth that the supports are pairwise distinct, and put

\[
 \Sigma=\operatorname{Sh}(S),
 \qquad
 \Sigma^+=\Sigma\mathbin{\dot\cup}\{P_t:t\in T\},
 \qquad
 C_T=S\cup\{a_t:t\in T\}.
\]

The shadow, the signed minimal nonfaces, and the inverse parameters of
the assembly are all explicit.  Namely,

\begin{equation}
 \operatorname{Sh}(X_T)
 =\{D\cup L:D\in\Sigma,\ L\subseteq E\}
  \mathbin{\dot\cup}\{P_t:t\in T\},
 \label{eq:damp-independent-repair-shadow}
\end{equation}

\begin{equation}
 \mathcal N(\operatorname{Sh}(X_T))
 =\mathcal N(\Sigma^+)\mathbin{\dot\cup}
   \{P_t\cup\{e\}:t\in T,\ e\in E\}.
 \label{eq:damp-independent-repair-minimal-nonfaces}
\end{equation}

For the corresponding unique missing traces,

\begin{align}
 \sigma_{X_T}(I)&=\sigma_{C_T}(I)
 &&(I\in\mathcal N(\Sigma^+)),
 \label{eq:damp-independent-repair-old-signs}\\
 \sigma_{X_T}(P_t\cup\{e\})
   &=\bigl(a_t|_{P_t},1-t_e\bigr)
 &&(t\in T,\ e\in E).
 \label{eq:damp-independent-repair-layer-signs}
\end{align}

In particular, these signed clauses reconstruct \(X_T\).  Conversely,
the marked layer set \(E\) canonically recovers all the displayed
parameters.  If \(H_v\subseteq Q_F\) is the projected \(v\)-section of
\(X_T\), then

\begin{equation}
 S=\bigcap_{v\in Q_E}H_v,
 \qquad
 T=\{v\in Q_E:H_v\ne S\},
 \qquad
 \{a_t\}=H_t\setminus S,
 \label{eq:damp-independent-repair-recovery}
\end{equation}

and each \(P_t\) is recovered by the local formula in
\eqref{eq:damp-independent-repair-supports}.
\end{proposition}

\begin{proof}
For \(U\subseteq T\), put \(C_U=S\cup\{a_t:t\in U\}\).  Since \(S\)
is ample, it strongly shatters every member of \(\Sigma\).  The
one-concept extension criterion shows that \(A_t\) strongly shatters
\(P_t\).  If the supports \(P_t\), \(t\in U\), are distinct, then
\(C_U\) therefore strongly shatters at least

\[
 |\Sigma|+|U|=|S|+|U|=|C_U|
\]

distinct supports.  The Sandwich inequalities force equality, so
\(C_U\) is ample and

\begin{equation}
 \operatorname{Sh}(C_U)
 =\Sigma\mathbin{\dot\cup}\{P_t:t\in U\}.
 \label{eq:damp-independent-repair-union-shadow}
\end{equation}

Apply Proposition~\ref{prop:damp-layer-system-trace-transport} to the
sections \(H_t=A_t\) for \(t\in T\) and \(H_t=S\) otherwise.  If
\(L=\varnothing\), their union is \(C_T\), whose shadow is given by
\eqref{eq:damp-independent-repair-union-shadow}.  If
\(L\ne\varnothing\), every union \(H_{L,r}\) contains \(S\), while a
given new support \(P_t\) occurs only in the unique union with
\(r=t|_L\).  Since \(Q_L\) has at least two words, the intersection of
the shadows of the \(H_{L,r}\) is exactly \(\Sigma\).  The shattering
criterion \eqref{eq:damp-layer-system-shattering} now proves
\eqref{eq:damp-independent-repair-shadow}.  Its right-hand side has

\[
 2^{|E|}|\Sigma|+|T|=2^{|E|}|S|+|T|=|X_T|
\]

members, giving ampleness.

Conversely, suppose \(P_t=P_u=P\) for distinct \(t,u\in T\), and choose
\(e\in E\) with \(t_e\ne u_e\).  Each of the two unions of sections on
the two sides of \(e\) contains an extension that shatters \(P\).
Hence \(P\cup\{e\}\in\operatorname{Sh}(X_T)\).  The product
\(S\times Q_E\) already shatters all \(D\cup L\) with \(D\in\Sigma\)
and \(L\subseteq E\).  Moreover, the old-coordinate projection \(C_T\)
has at least \(|C_T|=|S|+|T|\) shattered supports by Sauer's inequality,
and hence at least \(|T|\) old supports outside \(\Sigma\).  The support
\(P\cup\{e\}\) is an additional one.  Therefore

\[
 |\operatorname{Sh}(X_T)|
 \ge 2^{|E|}|S|+|T|+1=|X_T|+1,
\]

so \(X_T\) is not ample.  This proves
\eqref{eq:damp-independent-repair-ample-iff}.

It remains to read the signed presentation from
\eqref{eq:damp-independent-repair-shadow}.  Each \(P_t\) is a minimal
nonface of \(\Sigma\).  A nonface using a layer coordinate is minimal
precisely when it has the form \(P_t\cup\{e\}\); nonfaces contained in
\(F\) are exactly the minimal nonfaces of \(\Sigma^+\).  This proves
\eqref{eq:damp-independent-repair-minimal-nonfaces}.  Projecting to the
old coordinates gives \(C_T\), which proves
\eqref{eq:damp-independent-repair-old-signs}.  The trace
\(a_t|_{P_t}\) is missing from \(S\).  No other tip \(a_u\) supplies it:
otherwise \(A_u\) would shatter \(P_t\), contradicting

\[
 \operatorname{Sh}(A_u)=\Sigma\mathbin{\dot\cup}\{P_u\}
 \quad\text{and}\quad P_u\ne P_t.
\]

Thus this trace occurs only in the full layer \(t\), and the missing
trace on \(P_t\cup\{e\}\) is
\((a_t|_{P_t},1-t_e)\).  This proves
\eqref{eq:damp-independent-repair-layer-signs}; the fixed-shadow
signed-clutter theorem reconstructs \(X_T\).

Finally, \(E\ne\varnothing\), every tip occurs in exactly one section,
and the tips are pairwise distinct.  The intersection of all sections is
therefore \(S\), after which the remaining formulas in
\eqref{eq:damp-independent-repair-recovery} and the local support formula
are immediate.
\end{proof}

\begin{proposition}[Replicated repair layers]
\label{prop:damp-replicated-repair-layers}
Let \(E\ne\varnothing\) and \(F\) be disjoint, let
\(S\subseteq Q_F\) be ample, and choose a finite set
\(\mathcal A\subseteq Q_F\setminus S\).  Suppose that every
\(S\cup\{a\}\), \(a\in\mathcal A\), is ample, and that their new
supports
\[
 \operatorname{Sh}(S\cup\{a\})
       \setminus\operatorname{Sh}(S)=\{P_a\}
\]
are pairwise distinct.  For each \(a\in\mathcal A\), choose a
nonempty occurrence family \(R_a\subseteq Q_E\), and put
\begin{equation}
 X(\mathbf R)
 =(S\times Q_E)\cup
   \bigcup_{a\in\mathcal A}(\{a\}\times R_a).
 \label{eq:damp-replicated-repair-assembly}
\end{equation}
Then its shattered-set complex is
\begin{equation}
 \operatorname{Sh}(X(\mathbf R))
 =\{D\cup L:D\in\operatorname{Sh}(S),\ L\subseteq E\}
  \mathbin{\dot\cup}
  \mathop{\dot\bigcup}_{a\in\mathcal A}
  \{P_a\cup L:L\in\operatorname{Sh}(R_a)\}.
 \label{eq:damp-replicated-repair-shadow}
\end{equation}
Consequently,
\begin{equation}
 X(\mathbf R)\text{ is ample}
 \quad\Longleftrightarrow\quad
 R_a\text{ is ample for every }a\in\mathcal A.
 \label{eq:damp-replicated-repair-ample-iff}
\end{equation}

Assume these equivalent conditions.  Set
\[
 \Sigma=\operatorname{Sh}(S),\qquad
 \Sigma^+=\Sigma\mathbin{\dot\cup}\{P_a:a\in\mathcal A\},\qquad
 \Gamma_a=\operatorname{Sh}(R_a),\qquad
 C=S\cup\mathcal A.
\]
Then the signed minimal nonfaces are
\begin{equation}
 \mathcal N(\operatorname{Sh}(X(\mathbf R)))
 =\mathcal N(\Sigma^+)\mathbin{\dot\cup}
  \mathop{\dot\bigcup}_{a\in\mathcal A}
  \{P_a\cup J:J\in\mathcal N(\Gamma_a)\},
 \label{eq:damp-replicated-repair-minimal-nonfaces}
\end{equation}
with signs
\begin{align}
 \sigma_{X(\mathbf R)}(I)&=\sigma_C(I)
 &&(I\in\mathcal N(\Sigma^+)),
 \label{eq:damp-replicated-repair-old-signs}\\
 \sigma_{X(\mathbf R)}(P_a\cup J)
  &=\bigl(a|_{P_a},\sigma_{R_a}(J)\bigr)
 &&(a\in\mathcal A,\ J\in\mathcal N(\Gamma_a)).
 \label{eq:damp-replicated-repair-layer-signs}
\end{align}
Thus the old repair clause and the layer clause compose without any
additional compatibility condition.

If every \(R_a\) is proper in \(Q_E\), the marked layer set recovers the
parameters canonically.  For the projected sections \(H_t\subseteq Q_F\),
\begin{equation}
 S=\bigcap_{t\in Q_E}H_t,\qquad
 \mathcal A=\left(\bigcup_{t\in Q_E}H_t\right)\setminus S,\qquad
 R_a=\{t\in Q_E:a\in H_t\},
 \label{eq:damp-replicated-repair-recovery}
\end{equation}
and \(P_a=\{f\in F:a^{\{f\}}\in S\}\).  Without the properness
assumption, the same formulas give the unique normalized presentation:
every tip with \(R_a=Q_E\) is absorbed into the common core.
\end{proposition}

\begin{proof}
For \(U\subseteq\mathcal A\), the proof of
Proposition~\ref{prop:damp-independent-repair-layers} gives
\begin{equation}
 \operatorname{Sh}(S\cup U)
 =\operatorname{Sh}(S)\mathbin{\dot\cup}\{P_a:a\in U\}.
 \label{eq:damp-replicated-repair-union-shadow}
\end{equation}
Indeed, these are \(|S|+|U|\) distinct strongly shattered supports, so
the Sandwich inequalities force equality.

For \(L\subseteq E\) and \(r\in Q_L\), the union of the sections with
layer trace \(r\) is
\[
 H_{L,r}
 =S\cup\{a\in\mathcal A:
          R_a\cap\{t\in Q_E:t|_L=r\}\ne\varnothing\}.
\]
Equation~\eqref{eq:damp-replicated-repair-union-shadow} therefore shows
that \(H_{L,r}\) shatters \(P_a\) exactly when \(R_a\) contains a word
with \(L\)-trace \(r\).  Hence every \(H_{L,r}\) shatters \(P_a\)
exactly when \(R_a\) shatters \(L\).  Every one of these unions shatters
all members of \(\Sigma\), and their total union shatters no old support
outside \(\Sigma^+\).  Proposition
\ref{prop:damp-layer-system-trace-transport} now gives
\eqref{eq:damp-replicated-repair-shadow}.

The two sides of the Sandwich inequality differ by
\[
 |\operatorname{Sh}(X(\mathbf R))|-|X(\mathbf R)|
 =\sum_{a\in\mathcal A}
   \bigl(|\operatorname{Sh}(R_a)|-|R_a|\bigr).
\]
Every summand is nonnegative by Sauer's inequality.  It vanishes exactly
when \(R_a\) is ample, proving
\eqref{eq:damp-replicated-repair-ample-iff}.

Now suppose all occurrence families are ample.  A nonface contained in
\(F\) is minimal exactly when it lies in \(\mathcal N(\Sigma^+)\).
If a minimal nonface uses a layer coordinate, its old part must be one
of the minimal nonfaces \(P_a\) added to \(\Sigma\), and its layer part
must be a minimal nonface of \(\Gamma_a\).  Conversely, every
\(P_a\cup J\) with \(J\in\mathcal N(\Gamma_a)\) is minimal: deleting an
old coordinate puts its old part in \(\Sigma\), while deleting a layer
coordinate puts its layer part in \(\Gamma_a\).  This proves
\eqref{eq:damp-replicated-repair-minimal-nonfaces}.

The old-coordinate projection of \(X(\mathbf R)\) is \(C\), which gives
\eqref{eq:damp-replicated-repair-old-signs}.  The trace
\(a|_{P_a}\) is missing from \(S\) and from every other repair tip, by
the distinctness of the supports.  It occurs in \(X(\mathbf R)\)
precisely over the layer words in \(R_a\).  Thus a trace on
\(P_a\cup J\) is missing precisely when its old part is
\(a|_{P_a}\) and its layer part is the unique trace
\(\sigma_{R_a}(J)\) missing from \(R_a\).  This proves
\eqref{eq:damp-replicated-repair-layer-signs}.

Finally, if all occurrence families are proper, every tip is absent
from some section, whereas every member of \(S\) occurs in every
section.  This proves \eqref{eq:damp-replicated-repair-recovery}, and
the local support formula follows from the one-concept extension
criterion.  In general,
\[
 \bigcap_{t\in Q_E}H_t
 =S\cup\{a\in\mathcal A:R_a=Q_E\}.
\]
Absorbing these full rows into the core preserves ampleness and the
distinct-support union formula
\eqref{eq:damp-replicated-repair-union-shadow}; all remaining occurrence
families are proper.  The displayed intersection makes this normalized
presentation unique.
\end{proof}

\begin{remark}[The first genuinely coupled parameters]
\label{rem:damp-replicated-repair-coupling}
The independent construction in
Proposition~\ref{prop:damp-independent-repair-layers} is the special
case in which every \(R_a\) is a singleton.  Proposition
\ref{prop:damp-replicated-repair-layers} replaces a search through
layered cube families by two finite pieces of data: a repair tip labelled
by \(P_a\), and an ample occurrence family \(R_a\) on the layer cube.
Overlap among the \(R_a\) is exactly the first coupling mechanism:
several tips may occur in one section and one tip may persist through
several sections.  The remaining obstruction problem is no longer to
guess which such systems are ample, but to characterize which ample
labelled covers \((R_a)_{a\in\mathcal A}\) are pc-minimal and preserve
the unavoidable trace circuits of the seed.
\end{remark}

\begin{corollary}[Layer pc-minors act on the occurrence cover]
\label{cor:damp-replicated-repair-layer-minors}
Use the notation of
Proposition~\ref{prop:damp-replicated-repair-layers}.  For \(e\in E\)
and \(i\in\{0,1\}\), let
\[
 \rho_e^iR_a
 =\{t|_{E\setminus\{e\}}:t\in R_a,\ t_e=i\},
 \qquad
 \pi_eR_a
 =\{t|_{E\setminus\{e\}}:t\in R_a\}.
\]
Then
\begin{align}
 \rho_e^iX(\mathbf R)
  &=(S\times Q_{E\setminus\{e\}})
    \cup\bigcup_{a\in\mathcal A}
      (\{a\}\times\rho_e^iR_a),
 \label{eq:damp-replicated-repair-restriction}\\
 \pi_eX(\mathbf R)
  &=(S\times Q_{E\setminus\{e\}})
    \cup\bigcup_{a\in\mathcal A}
      (\{a\}\times\pi_eR_a).
 \label{eq:damp-replicated-repair-contraction}
\end{align}
Empty rows are discarded, and rows equal to
\(Q_{E\setminus\{e\}}\) are absorbed into the common core.  After this
normalization, both images are again replicated repair assemblies with
the surviving labels \(P_a\).

Thus restrictions and contractions in marked layer coordinates are
computed entirely inside the finite labelled occurrence cover
\((R_a,P_a)_{a\in\mathcal A}\).  In particular, if \(X(\mathbf R)\) is
a forbidden pc-minor, all \(3|E|\) normalized covers in
\eqref{eq:damp-replicated-repair-restriction}--
\eqref{eq:damp-replicated-repair-contraction} are corner-peelable.
\end{corollary}

\begin{proof}
A word \((x,t)\) survives the \(e=i\) restriction precisely when
\(t_e=i\), and suppressing \(e\) replaces its layer word by
\(t|_{E\setminus\{e\}}\).  This proves
\eqref{eq:damp-replicated-repair-restriction}.  Contraction takes the
union of the two restricted sections, so it projects each \(R_a\), proving
\eqref{eq:damp-replicated-repair-contraction}.  The normalization and
label statements follow from the recovery part of Proposition
\ref{prop:damp-replicated-repair-layers}.  The final assertion is the
elementary pc-minor condition for a forbidden pc-minor.
\end{proof}

\begin{proposition}[Corners and layer-minimality of a repair cover]
\label{prop:damp-replicated-repair-corners}
Use the notation and hypotheses of
Proposition~\ref{prop:damp-replicated-repair-layers}, assume that all
occurrence families are ample and proper, and let
\[
 C=S\cup\mathcal A,
 \qquad
 Q_a=\{a^U:U\subseteq P_a\}\quad(a\in\mathcal A).
\]
Then \(Q_a\) is the unique maximal cube of \(C\) through \(a\).  The
maximal cubes of \(X(\mathbf R)\) through \((a,t)\) are exactly
\begin{equation}
 Q_a\times K,
 \qquad
 K\text{ a maximal cube of }R_a\text{ through }t.
 \label{eq:damp-replicated-repair-tip-cubes}
\end{equation}
If \(S\) is cornerless, this gives the exact corner set
\begin{equation}
 \operatorname{Cor}(X(\mathbf R))
 =\mathop{\dot\bigcup}_{a\in\mathcal A}
   \bigl(\{a\}\times\operatorname{Cor}(R_a)\bigr).
 \label{eq:damp-replicated-repair-corners}
\end{equation}

Suppose in addition that \(S\) is nonpeelable, every one-concept repair
\(S\cup\{a\}\) belongs to \(\Damp\), and every \(R_a\) belongs to
\(\Damp\).  For a layer restriction or contraction
\[
 \nu\in\{\rho_e^0,\rho_e^1,\pi_e\},
 \qquad e\in E,
\]
one then has
\begin{equation}
 \nu X(\mathbf R)\in\Damp
 \quad\Longleftrightarrow\quad
 \nu R_a=Q_{E\setminus\{e\}}
 \text{ for some }a\in\mathcal A.
 \label{eq:damp-replicated-repair-layer-peeling-iff-full-row}
\end{equation}
Consequently, all elementary pc-minors in marked layer coordinates are
corner-peelable if and only if the occurrence rows own every signed facet
of the layer cube:
\begin{equation}
 \text{for every }e\in E\text{ and }i\in\{0,1\},
 \quad
 \{t\in Q_E:t_e=i\}\subseteq R_a
 \text{ for some }a\in\mathcal A.
 \label{eq:damp-replicated-repair-signed-facet-cover}
\end{equation}
Thus the entire layer-coordinate part of pc-minimality is a finite
containment test on the labelled cover; it involves neither a corner order
nor a trace-carrier scheme.
\end{proposition}

\begin{proof}
For every \(U\subseteq\mathcal A\), equation
\eqref{eq:damp-replicated-repair-union-shadow} says that \(S\cup U\) is
ample.  In particular, both \(C\) and \(C\setminus\{a\}\) are ample, so
\(a\) is a corner of \(C\).  The one-concept extension criterion identifies
its unique maximal cube as \(Q_a\).

Let \(L\) be the support of a cube of \(X(\mathbf R)\) through
\((a,t)\) in the layer coordinates.  Its old-coordinate projection is a
cube of \(C\) through \(a\), hence is contained in \(Q_a\).  Moreover,
all vertices with old word \(a\) show that the layer face through \(t\)
on \(L\) lies in \(R_a\).  Conversely, whenever a layer cube
\(K\subseteq R_a\) contains \(t\), the product \(Q_a\times K\) lies in
the assembly: its vertices with old word different from \(a\) belong to
\(S\times Q_E\), while its vertices with old word \(a\) are supplied by
\(R_a\).  Maximality on both sides proves
\eqref{eq:damp-replicated-repair-tip-cubes}.

Now let \(s\in S\), and let \(M\) be a maximal cube of \(S\) through
\(s\).  The cube \(M\times Q_E\) is maximal in the assembly.  Indeed, a
proper extension of it can use only old coordinates, and any new old word
in that extension would have to occur in every layer.  Such a word would
be a repair tip with full occurrence row, contrary to properness.  If
\(S\) is cornerless, \(s\) lies in at least two maximal cubes of \(S\),
so \((s,t)\) lies in at least two maximal cubes of the assembly.  Together
with \eqref{eq:damp-replicated-repair-tip-cubes}, this proves
\eqref{eq:damp-replicated-repair-corners}.

Fix \(\nu\) as in the statement and normalize the rows of \(\nu X\) as
in Corollary~\ref{cor:damp-replicated-repair-layer-minors}.  If no image
row is full, strongly projecting all remaining layer coordinates gives
exactly \(S\).  Since \(S\) is nonpeelable,
Proposition~\ref{prop:damp-strong-projection-closure} gives
\(\nu X\notin\Damp\).  This argument also covers a zero-dimensional
remaining layer cube and requires no assumption on the corners of \(S\).

Suppose instead that the set
\[
 \mathcal A_0=\{a:\nu R_a=Q_{E\setminus\{e\}}\}
\]
is nonempty.  Absorb these rows into the new core
\(C_0=S\cup\mathcal A_0\).  Choose \(a_0\in\mathcal A_0\).  Every other
member of \(\mathcal A_0\) is a corner of the current old-coordinate
union, so these tips may be peeled until \(S\cup\{a_0\}\) remains.
The latter belongs to \(\Damp\), and hence \(C_0\in\Damp\).  Every
remaining image row is a pc-minor of some \(R_a\), so it belongs to
\(\Damp\).  Formula \eqref{eq:damp-replicated-repair-tip-cubes}, applied
successively, lifts a corner peeling of each proper row to the assembly.
After all proper rows have been removed, the family is
\(C_0\times Q_{E\setminus\{e\}}\), which belongs to \(\Damp\) by the
product test.  This proves
\eqref{eq:damp-replicated-repair-layer-peeling-iff-full-row}.

For a restriction, the image row \(\rho_e^iR_a\) is full exactly when
the signed facet \(\{t:t_e=i\}\) is contained in \(R_a\).  If this holds
for either sign, then \(\pi_eR_a\) is also full.  Thus the
\(2|E|\) restriction conditions imply all \(|E|\) contraction conditions,
and \eqref{eq:damp-replicated-repair-signed-facet-cover} follows.
\end{proof}

\begin{lemma}[Compatible repairs do not collide]
\label{lem:damp-compatible-repairs-do-not-collide}
Let \(S\subseteq Q_F\) be ample, and let
\(\mathcal A\subseteq Q_F\setminus S\) be a set of one-concept ample
extensions whose new shattered supports \(P_a\), \(a\in\mathcal A\), are
pairwise distinct.  Then no two members of \(\mathcal A\) are adjacent in
\(Q_F\).

More generally, let \(\mu\) be a restriction or contraction in one old
coordinate, put \(S_\mu=\mu S\), and discard the tips which do not survive
the restriction or whose image belongs to \(S_\mu\).  The map
\(a\mapsto\mu a\) is injective on the remaining tips.  Moreover, after any
subset of these projected tips has been retained, every retained tip is a
corner of the resulting old-coordinate family.
\end{lemma}

\begin{proof}
Suppose that \(b=a^{\{f\}}\) for distinct \(a,b\in\mathcal A\).  Since
neither tip belongs to \(S\), the repair-support formula gives
\(f\notin P_a\cup P_b\).  The family
\(C=S\cup\{a,b\}\) is ample by
\eqref{eq:damp-replicated-repair-union-shadow}.  Under contraction of \(f\),
the two tips merge into one word \(c\), and \(c\notin\pi_fS\); hence
\[
 |\pi_fC|=|\pi_fS|+1.
\]
On the other hand, both distinct supports \(P_a,P_b\subseteq F\setminus
\{f\}\) are shattered by \(\pi_fC\) and not by \(\pi_fS\).  Indeed, the
missing trace on \(P_a\) is \(a|_{P_a}\), contraction does not change this
trace, and the same argument applies to \(P_b\).  Both contracted families
are ample, so the equality between cardinality and number of shattered
supports makes a gain of at least two shattered supports incompatible with
the displayed gain of one concept.  This proves nonadjacency.

A restriction is injective on the tips which survive it.  Two tips can have
the same image under a contraction only when they are adjacent in the
contracted coordinate, which has just been excluded.  This proves the
injectivity assertion.  Finally, let \(\mathcal U\) be any subset of the
remaining tips.  The families \(S\cup\mathcal U\) and
\(S\cup(\mathcal U\setminus\{a\})\) are ample by
\eqref{eq:damp-replicated-repair-union-shadow}; their images under \(\mu\)
are ample as well.  Injectivity and the exclusion of images in \(S_\mu\)
show that the second image is obtained from the first by deleting precisely
\(\mu a\).  The ample corner-deletion criterion therefore makes \(\mu a\)
a corner, independently of the chosen subset.
\end{proof}

\begin{theorem}[Simultaneous extensions grouped by acquired support]
\label{thm:damp-equal-support-extensions}
Let $S\subseteq Q_F$ be nonempty and ample, and let
$U\subseteq Q_F\setminus S$ be a finite set such that each
$S\cup\{a\}$, $a\in U$, is ample. Write
\[
 \operatorname{Sh}(S\cup\{a\})\setminus\operatorname{Sh}(S)=\{I_a\}.
\]
For each occurring support $I$, put
\[
 U_I=\{a\in U:I_a=I\},\qquad
 R_I=\{a|_{F\setminus I}:a\in U_I\},\qquad X=S\cup U.
\]
Then
\begin{enumerate}
 \item $X$ is ample if and only if every $R_I$ is ample.
 \item Under these equivalent conditions,
 \[
  X^I=R_I,\qquad
  \operatorname{Sh}(X)=\operatorname{Sh}(S)
   \mathbin{\dot\cup}\mathop{\dot\bigcup}_{I}
    \{I\cup J:J\in\operatorname{Sh}(R_I)\}.
 \]
 \item Under the same conditions, $X$ peels to $S$ if and only if
 every $R_I$ belongs to $\Damp$.
\end{enumerate}
The empty set of tips is allowed, with no row conditions.
\end{theorem}
\begin{proof}
The one-concept extension criterion makes $I_a$ a minimal nonface of
$\operatorname{Sh}(S)$ and identifies it with the directions from $a$
to its neighbours in $S$. Let $\tau_I$ be the unique missing trace of
$S$ on $I$. Every member of $U_I$ has trace $\tau_I$, and its punctured
$I$-cube lies in $S$. A tip of another acquired support cannot have
trace $\tau_I$: otherwise its one-concept extension would shatter $I$
as well as its distinct acquired support. Projection is therefore
injective on $U_I$, and the vertices of $X$ with $I$-trace $\tau_I$
are precisely $U_I$.

It follows, without assuming ampleness of $X$, that a full $I$-fibre
occurs exactly above a member of $R_I$. This proves $X^I=R_I$.
If $X$ is ample, reduction closure of ampleness gives the necessity
in item~1.

Conversely, suppose all $R_I$ are ample. Each contained $J$-cube in
$R_I$ lifts to a full $(I\cup J)$-cube in $X$: its vertices with
$I$-trace $\tau_I$ are the corresponding tips, and every other vertex
belongs to one of their punctured $I$-cubes in $S$. Thus $X$ strongly
shatters every support in the displayed formula. None of these new
supports belongs to $\operatorname{Sh}(S)$, since it contains $I$.
No support contributed by $I$ can contain another occurring support
$K\ne I$. Indeed, its lifted cube lies already in $S\cup U_I$,
whereas this family still misses $\tau_K$. In particular the
contributions of different groups are disjoint. Their total number is
\[
 |S|+\sum_I|\operatorname{Sh}(R_I)|
 =|S|+\sum_I|R_I|=|X|.
\]
The Sandwich inequalities force ampleness and the exact shadow formula.

For item~3, tips from different groups are nonadjacent by
Lemma~\ref{lem:damp-compatible-repairs-do-not-collide}, applied to each
pair of distinct supports. Fix a tip $a\in U_I$ and write
$z=a|_{F\setminus I}$. Its neighbours in $S$ have precisely the
directions $I$, and its remaining neighbours correspond precisely to
the neighbours of $z$ in $R_I$. If those latter directions form $J$,
the whole incident $(I\cup J)$-cube at $a$ is present exactly when
the incident $J$-cube at $z$ is present in $R_I$. This follows from
the punctured-fibre description above. Hence $a$ is a corner of $X$
if and only if $z$ is a corner of $R_I$.

This equivalence remains true after any set of tips has been deleted:
the unchanged base $S$ still supplies every punctured fibre. Peelings
of the row families therefore lift and concatenate to a relative
peeling of $X$ to $S$. Conversely, a relative peeling, restricted to
each group and projected outside $I$, is a complete corner peeling
of $R_I$. This proves item~3.
\end{proof}

The theorem extends the distinct-support union formula to repeated
acquired supports in the original coordinates. The replicated-layer
construction of Proposition~\ref{prop:damp-replicated-repair-layers}
uses additional independent layer coordinates and assumes distinct old
supports; neither assumption is needed here. It is an exact simultaneous
extension rule, not a claim that arbitrary ample supersets can be reached
by tips individually admissible over their original base.

\begin{corollary}[Proper-minor positivity for simultaneous initial tips]
\label{cor:damp-initial-tip-minor-positivity}
Use the notation of Theorem~\ref{thm:damp-equal-support-extensions}
and assume $X$ is ample. Each $R_I$ is a proper coordinate face of $X$.
Consequently, if $X\in\Damp$, or more generally if all proper coordinate
faces of $X$ belong to $\Damp$, then $X$ peels to $S$.

If in addition $S$ is an ample forbidden pc-minor using every coordinate
of $F$, then the following are equivalent:
\begin{enumerate}
 \item every $R_I$ belongs to $\Damp$;
 \item $X$ peels to $S$;
 \item every nonempty proper pc-minor of $X$ belongs to $\Damp$;
 \item $X\in\Damp\cup\Obs(\Damp)$.
\end{enumerate}
If $X$ is forbidden, every intermediate family in any relative peeling
from $X$ to $S$ is forbidden. In particular, a corner-descent-terminal
forbidden class cannot properly contain a forbidden $S$ for which all
its additional vertices are individually ample additions to $S$.
\end{corollary}
\begin{proof}
As shown in the theorem's proof, the vertices of $X$ with $I$-trace
$\tau_I$ are exactly $U_I$. Fixing these coordinates therefore gives
$R_I$. The face is proper because $S$ is nonempty and contains no vertex
with that trace. It is a pc-minor of $X$, so positivity of the proper
faces implies positivity of every $R_I$. Item~3 of the theorem now
supplies the relative peeling. This also applies when $X$ itself is
positive, by pc-minor closure.

For the equivalences, the theorem gives $1\Longleftrightarrow2$, and
the proper-face observation gives $3\Longrightarrow1$. To prove
$2\Longrightarrow3$, apply any elementary coordinate restriction or
contraction $\mu$ to a relative peeling from $X$ to $S$. Consecutive
images are ample and either identical or differ by one vertex.
Omitting repetitions gives a relative corner peeling to $\mu S$.
Every coordinate is active in $S$, so this last image is a proper
pc-minor of the forbidden class $S$ and is positive. Thus $\mu X$
is positive. Every proper pc-minor is obtained through such elementary
operations, and pc-minor closure proves item~3. Equivalence of items~3
and~4 is the definition of forbiddenness together with pc-minor closure.

Finally, let $X$ be forbidden and let $Y$ occur on a relative peeling
to $S$. If $Y$ were positive, the preceding corner deletions followed
by a peeling of $Y$ would peel $X$, a contradiction. The remainder
of the relative chain transports to every elementary minor of $Y$
just as above, proving that all its proper minors are positive.
Hence $Y$ is forbidden. If $X\ne S$, its first corner deletion
already contradicts corner-descent terminality.
\end{proof}

The corollary removes the need to supply a relative peeling when a
positive completion consists entirely of individually admissible tips.
It also separates the two selection questions over a forbidden base:
positive row faces settle all proper-minor conditions, while deciding
whether the full assembly is nonpeelable remains a separate obligation.
It gives no augmentation theorem for supersets containing vertices that
cannot individually be added amply to the base.

\begin{theorem}[Simple complete common-seed cover criterion]
\label{thm:damp-common-seed-cover-classification}
Let \(S\in\Obs(\Damp)\cap\Ample\) be irredundantly
embedded in \(Q_F\).  Let \(\mathcal A\subseteq Q_F\setminus S\) be a
finite set such that every \(S\cup\{a\}\) belongs to \(\Damp\) and the
new supports \(P_a\) are pairwise distinct.  Let \(E\ne\varnothing\),
and assign to every \(a\in\mathcal A\) a nonempty proper ample family
\(R_a\subsetneq Q_E\).  Form \(X(\mathbf R)\) as in
\eqref{eq:damp-replicated-repair-assembly}.  Then
\begin{equation}
 X(\mathbf R)\in\Obs(\Damp)
 \quad\Longleftrightarrow\quad
 \begin{gathered}
 R_a\in\Damp\quad\text{for every }a\in\mathcal A,\\
 \{t:t_e=i\}\subseteq R_a\quad\text{for some }a\in\mathcal A
       \quad(e\in E,\ i\in\{0,1\}).
 \end{gathered}
 \label{eq:damp-common-seed-cover-classification}
\end{equation}
Thus no condition in the old coordinates is needed: compatible repairs
cannot collide under an elementary old-coordinate map, and their projected
tips remain successively peelable.  The two displayed conditions are finite and
descendant-free, but the row condition \(R_a\in\Damp\) is still a
recognition condition and not intrinsic structural data.

The normalized marked sections canonically recover
\(S,(a,P_a,R_a)_{a\in\mathcal A}\).  Hence
\eqref{eq:damp-common-seed-cover-classification} is exact within the
common-seed class over any fixed ample forbidden seed.  It is not yet an
intrinsic generative classification: the seed is assumed to be an
obstruction, and peelability of the occurrence rows remains to be explained
structurally.  The terminal specialization below removes the latter defect.
\end{theorem}

\begin{proof}
Proposition~\ref{prop:damp-replicated-repair-layers} makes the assembly
ample.  Since every row is proper, its reduction along \(E\)
is exactly the nonpeelable seed \(S\).  Proposition
\ref{prop:damp-strong-projection-closure} therefore gives
\(X(\mathbf R)\notin\Damp\), whether or not \(S\) has corners.

Suppose first that \(X(\mathbf R)\) is a forbidden pc-minor.  Fixing all
old coordinates to the word \(a\) realizes \(R_a\) as a proper pc-minor,
so \(R_a\in\Damp\).  Proposition
\ref{prop:damp-replicated-repair-corners} applied to every layer restriction
gives the signed-facet condition in
\eqref{eq:damp-common-seed-cover-classification}.  This proves necessity.

Conversely, assume the two displayed conditions.  The signed-facet
condition and
\eqref{eq:damp-replicated-repair-layer-peeling-iff-full-row} put every
elementary layer restriction and contraction in \(\Damp\).  It remains to
consider an old-coordinate map \(\mu\).  Put \(S_\mu=\mu S\), discard the
tips which do not survive a restriction or whose image belongs to
\(S_\mu\), and denote the remaining index set by \(\mathcal A_\mu\).
Lemma~\ref{lem:damp-compatible-repairs-do-not-collide} gives
\begin{equation}
 \mu X(\mathbf R)
 =(S_\mu\times Q_E)
  \cup
  \bigcup_{a\in\mathcal A_\mu}
    (\{\mu a\}\times R_a),
 \label{eq:damp-common-seed-old-image}
\end{equation}
and says that every projected tip is a corner after any subset of the other
projected tips has been removed.

Choose \(a\in\mathcal A_\mu\).  A corner peeling of \(R_a\), lifted by the
product argument in the proof of
\eqref{eq:damp-replicated-repair-tip-cubes}, removes its entire occurrence
row and then its projected tip.  The lemma permits this operation to be
repeated in any order.  Eventually only \(S_\mu\times Q_E\) remains.
Because the embedding of \(S\) is irredundant, \(S_\mu\) is a proper
pc-minor of the obstruction \(S\), so it belongs to \(\Damp\); the product
test finishes the peeling.  Hence every old-coordinate elementary image
lies in \(\Damp\).

All elementary pc-minors of \(X(\mathbf R)\) are therefore
corner-peelable.  Pc-minor closure proves sufficiency in
\eqref{eq:damp-common-seed-cover-classification}.  The trace-carrier
reformulation follows from Theorem
\ref{thm:damp-acyclic-trace-carrier-criterion}, and parameter recovery is
the recovery statement of Proposition
\ref{prop:damp-replicated-repair-layers}.
\end{proof}

\begin{proposition}[Unique-midpoint cores separate by layer]
\label{prop:damp-layer-midpoint-cores}
Use the hypotheses and notation of
Proposition~\ref{prop:damp-replicated-repair-layers}, with all occurrence
rows ample and proper.  Write \(\kappa(Y)\) for the largest
unique-midpoint core of a family \(Y\), allowing it to be empty, and put
\[
 H_t=S\cup\{a\in\mathcal A:t\in R_a\}\qquad(t\in Q_E).
\]
If \(\kappa(R_a)=\varnothing\) for every \(a\in\mathcal A\), then
\begin{equation}
 \kappa(X(\mathbf R))
 =\mathop{\dot\bigcup}_{t\in Q_E}
      \bigl(\kappa(H_t)\times\{t\}\bigr).
 \label{eq:damp-layer-midpoint-cores}
\end{equation}
In particular, if the rows and all the sections \(H_t\) are
corner-peelable, then \(X(\mathbf R)\) has empty unique-midpoint core.
\end{proposition}

\begin{proof}
We first exclude occurrence words from a core, and then show that all
remaining witnesses stay in one layer.  The compatible tips are pairwise
nonadjacent by Lemma~\ref{lem:damp-compatible-repairs-do-not-collide}.
The cube \(Q_a\) of
Proposition~\ref{prop:damp-replicated-repair-corners} is the unique
maximal cube through \(a\) in \(S\cup\mathcal A\), and
\(Q_a\setminus\{a\}\subseteq S\).

Consider a unique-midpoint constraint rooted at an occurrence word
\((a,t)\).  If its two incident directions are old coordinates, the
two edges lie in \(Q_a\times\{t\}\), so their opposite midpoint is
present.  If one direction is old and the other is a layer coordinate,
the product description \eqref{eq:damp-replicated-repair-tip-cubes}
again supplies that midpoint: the old neighbor belongs to \(S\),
which occurs in every layer.  Thus every unique-midpoint constraint
rooted at \((a,t)\) uses two layer directions and is exactly a
constraint of \(R_a\), lifted at the word \(a\).

Let \(K\) be any unique-midpoint core of \(X(\mathbf R)\).
For fixed \(a\), a nonempty set
\(K_a=\{t:(a,t)\in K\}\) would consequently be a unique-midpoint
core of \(R_a\).  The hypothesis excludes this, so
\(K\subseteq S\times Q_E\).  A witnessing constraint for
\((s,t)\in K\) now has both heads in \(S\times Q_E\).
Two layer directions span a square in \(\{s\}\times Q_E\).
An old and a layer direction also span a full square in
\(S\times Q_E\).  Hence a unique-midpoint witness with both heads
in \(K\) must use two old directions and lie entirely in \(H_t\).
It follows that
\(\{s:(s,t)\in K\}\subseteq\kappa(H_t)\), proving one inclusion
in \eqref{eq:damp-layer-midpoint-cores}.

Conversely, a horizontal distance-two pair has both possible midpoints
in its own layer.  Its unique-midpoint constraint in \(H_t\) therefore
remains exactly the same constraint in \(X(\mathbf R)\).
Each nonempty \(\kappa(H_t)\times\{t\}\) is consequently a core of
\(X(\mathbf R)\), and their union is a core as well.  This proves the
reverse inclusion.  Finally, a corner-peelable family has empty
unique-midpoint core by Proposition~\ref{prop:damp-unique-midpoint-core},
which gives the last assertion.
\end{proof}

\begin{corollary}[Layer presentations force empty unique-midpoint core]
\label{cor:damp-layered-obstructions-empty-core}
Every forbidden pc-minor in
Theorem~\ref{thm:damp-common-seed-cover-classification} has empty
unique-midpoint core.  More generally, if an ample forbidden pc-minor
has a successful layer set, its unique-midpoint core is empty.
Thus a nonempty unique-midpoint core excludes every successful layer set.
The converse fails: empty core alone does not guarantee a successful
layer set, even for an ample forbidden pc-minor.
\end{corollary}

\begin{proof}
The common-seed criterion makes each occurrence row corner-peelable.
For every \(t\in Q_E\), the section \(H_t\) is obtained by fixing
all layer coordinates.  Since \(E\ne\varnothing\) and the nonempty
seed occurs in every layer, this is a proper pc-minor of the
obstruction, and hence is corner-peelable.  Proposition
\ref{prop:damp-layer-midpoint-cores} applies.  Theorem
\ref{thm:damp-terminal-clause-recovery} places every successful layer
presentation of an ample forbidden pc-minor in this common-seed class,
proving the necessary condition.

For the failure of the converse, use the one-corner obstruction of
Theorem~\ref{thm:damp-empty-core-counterexample}.  A successful layer
presentation has at least two nonempty proper occurrence rows: one
proper row cannot own both opposite facets of a layer coordinate.
Each row peels, so it has a corner, which lifts to a distinct occurrence
corner by \eqref{eq:damp-replicated-repair-tip-cubes}.  Such a
presentation would therefore give at least two corners, contrary to
the example's unique corner.
\end{proof}

\begin{corollary}[Terminal empty-core obstructions in unbounded dimension]
\label{cor:damp-terminal-empty-core-towers}
Let \(G\subseteq Q_{12}\) be the certified \(291\)-concept family
stored as \texttt{base} in \path{damp_fixed_291_repair_creation.json}.
For every integer \(r\ge1\), put
\[
 X_r=(G\times Q_r)
       \cup\bigl(\{91\}\times(Q_r\setminus\{0^r\})\bigr)
       \cup\bigl(\{157\}\times(Q_r\setminus\{1^r\})\bigr).
\]
Then \(X_r\) is a corner-descent-terminal ample forbidden pc-minor
with empty unique-midpoint core, isometric dimension \(12+r\),
\(293\cdot2^r-2\) concepts, and exactly \(2r\) corners.
For \(r=1\), both corner deletions have unique-midpoint core of
order \(76\).  In particular, terminality does not restore either the
core-or-nonample criterion or its core-safe-corner form.
\end{corollary}

\begin{proof}
The stored family \(G\) is cornerless and forbidden.  The two extensions
at \(91\) and \(157\) have recorded complete corner peelings and
acquire the distinct supports
\[
 P_{91}=\{1,2,5,6\},\qquad
 P_{157}=\{1,2,7,8\}.
\]
Theorem~\ref{thm:damp-antipodal-repair-towers} proves ampleness,
forbidden status, and the order and dimension formulas for every \(r\).
The rows partition the signed layer facets between the two owners, so
Theorem~\ref{thm:damp-terminal-common-seed-facet-partitions} proves
terminality.  Each punctured row has exactly \(r\) corners; since
\(G\) is cornerless, the corner formula
\eqref{eq:damp-replicated-repair-corners} gives exactly \(2r\) corners
of \(X_r\).  Corollary
\ref{cor:damp-layered-obstructions-empty-core} proves that their
unique-midpoint cores are empty.

When \(r=1\), deleting either occurrence corner empties its row.
The two remaining horizontal sections are \(G\), whose core has
order \(76\), and a one-concept repair of \(G\), whose core is empty.
Discarding the empty row and applying
\eqref{eq:damp-layer-midpoint-cores} gives core order \(76\) for both
deletions.  Thus neither of the two corners of the terminal
\(584\)-concept example is core-safe.

The fixed checker \path{verify_damp_layer_midpoint_cores.py} validates
the source hashes and both stored repair peelings.  It checks the
layer-core identity on the unrepaired product, on \(X_1,X_2\), and on
all their corner deletions.  Direct shattered-set counts verify these
finite controls, and the report
\path{damp_layer_midpoint_cores_verification.json} records complete
static pruning orders and surviving-core witnesses.  The uniform
conclusion uses the preceding proofs, not extrapolation from the controls.
\end{proof}

\begin{theorem}[Universality of peelable occurrence rows]
\label{thm:damp-universal-peelable-occurrence-row}
Suppose that the cornerless obstruction \(S\in\Obs(\Damp)\cap\Ample\)
has three one-concept extensions
\[
 S\cup\{a_0\},\qquad S\cup\{a_+\},\qquad S\cup\{a_-\}
\]
in \(\Damp\), and that their three new shattered supports are pairwise
distinct.  Then every nonempty family \(K\in\Damp\) is isomorphic to a
proper pc-minor of an explicitly constructed member
\[
 X_K\in\Obs(\Damp)\cap\Ample.
\]

More precisely, if \(K\subseteq Q_D\), adjoin one layer coordinate \(z\),
put \(E=D\mathbin{\dot\cup}\{z\}\), choose any word \(u\in Q_E\), and use
the three occurrence rows
\begin{equation}
 R_0=K\times\{0\},\qquad
 R_+=Q_E\setminus\{u\},\qquad
 R_-=Q_E\setminus\{\bar u\}
 \label{eq:damp-universal-occurrence-rows}
\end{equation}
at \(a_0,a_+,a_-\), respectively.  The common-seed assembly determined by
\eqref{eq:damp-universal-occurrence-rows} is such an \(X_K\), and fixing
the old coordinates to \(a_0\) and the coordinate \(z\) to zero recovers
\(K\).
\end{theorem}

\begin{proof}
The row \(R_0\) is isomorphic to \(K\), and the two other rows are singly
punctured cubes.  Hence all three rows are nonempty, proper, ample, and
corner-peelable.  It remains only to check the signed-facet condition in
Theorem~\ref{thm:damp-common-seed-cover-classification}.

Fix \(e\in E\).  The facet
\(\{t\in Q_E:t_e=1-u_e\}\) avoids \(u\), so it is contained in \(R_+\).
The opposite facet \(\{t:t_e=u_e\}\) avoids \(\bar u\), so it is contained
in \(R_-\).  Thus the two punctured rows own all \(2|E|\) signed facets,
independently of \(R_0\).  The common-seed criterion gives
\(X_K\in\Obs(\Damp)\cap\Ample\).

Finally, no member of \(S\), \(a_+\), or \(a_-\) has old-coordinate word
\(a_0\).  The section of \(X_K\) at that word is therefore exactly \(R_0\).
Restricting further to \(z=0\) gives \(K\).  Since the construction has the
nonempty old-coordinate factor \(S\times Q_E\), this is a proper pc-minor.
\end{proof}

\begin{remark}[What shelling universality rules out]
\label{rem:damp-shelling-universality-scope}
The verified Hall extension atlas supplies the required triple: in zero-based
coordinate numbering one may take
\[
 a_0=183,\qquad a_+=315,\qquad a_-=1836,
\]
whose new support labels are, respectively,
\[
 \{2,3,6,7\},\qquad \{2,3,8,9\},\qquad \{4,5,10,11\}.
\]
Thus the conclusion of the theorem is unconditional.  The underlying
extension record is
\begin{center}
\path{computations/hall_corner_exchange/source_H.json}.
\end{center}

Under the Yang dictionary, Theorem
\ref{thm:damp-universal-peelable-occurrence-row} says that an arbitrary
shellable Yang complex of an ample family can occur as an iterated
color-link inside the nonshellable Yang ball of a forbidden pc-minor.
Consequently, the elementary-image part of a generative classification
cannot be reduced to a special finite catalogue of easy shellable rows.
Either the parameters must carry a structural construction of arbitrary
corner-peelable ample families, or a global attachment theorem must make
their shellability automatic without naming the rows individually.  The
theorem is a universality barrier, not the missing construction: using
the condition \(K\in\Damp\) as admissibility would remain circular.
\end{remark}

\begin{theorem}[Universality within the reduced bases]
\label{thm:damp-reduced-base-universality}
For every nonempty ample family \(K\) of isometric dimension \(n\),
there is an explicit ample family \(X_K\) with the following properties.
\begin{enumerate}
 \item \(X_K\notin\Damp\), and \(K\) is isomorphic to a proper
 coordinate face of \(X_K\).
 \item The minimum graph degree of \(X_K\) is at least four, its
 signed minimal-nonface columns are pairwise distinct up to sign, and
 its crossing graph is complete.  In particular, it has no cut vertex
 or separating edge.
 \item Its parameters are
 \begin{equation}
 \begin{split}
 |X_K|&=301\,2^{n+4}+2|K|+1,\\
 \operatorname{idim}(X_K)&=n+16,\qquad
 \operatorname{VC}(X_K)=n+7.
 \end{split}
 \label{eq:damp-reduced-base-universality-size}
 \end{equation}
 \item \(X_K\in\Obs(\Damp)\) if and only if \(K\in\Damp\).
\end{enumerate}
Consequently, every nonempty corner-peelable ample family is a proper
coordinate face of a forbidden base already fixed by both reductions
of Theorem~\ref{thm:damp-alternating-reductions}.  These bases can all
be chosen with isometric dimension minus VC dimension equal to nine.
There are infinitely many such bases, even with that difference fixed.
\end{theorem}

\begin{proof}
We first enlarge the input so that each of its coordinates has a
distinct incidence pattern.  We then insert it into the three-repair
construction, whose two other rows ensure all required signed-facet
coverage.  The input is allowed to be nonpeelable throughout the
construction; its peeling status will decide precisely the minimality
of the output.

Embed \(K\subseteq Q_D\) irredundantly, with \(|D|=n\), and
reverse coordinates so that \(0\in K\).  Adjoin four distinct
coordinates \(w,p,q,z\), and put
\[
 A=D\mathbin{\dot\cup}\{w\},\qquad P=\{p,q\},\qquad
 E=A\mathbin{\dot\cup}P\mathbin{\dot\cup}\{z\},\qquad m=|E|=n+4.
\]
Let \(U=K\times Q_{\{w\}}\subseteq Q_A\) and form the vertex gluing
\begin{equation}
 L=(U\times\{0_P\})\cup(\{0_A\}\times Q_P)
 \subseteq Q_{A\cup P}.
 \label{eq:damp-reduced-base-separating-row}
\end{equation}
Products and vertex gluing preserve ampleness, so \(L\) is ample.
Moreover, \(|L|=2|K|+3\).  If \(K\in\Damp\), first delete the
three nonzero vertices of the attached square, in decreasing Hamming
weight, and then peel the product \(U\).  This proves \(L\in\Damp\)
using Proposition~\ref{prop:damp-vertex-gluing} and the product test
in Proposition~\ref{prop:damp-exact-expansion-product-tests}.

Fix the \(299\)-concept Hall obstruction \(S\subseteq Q_F\),
\(|F|=12\), and the three repairs \(a_0,a_+,a_-\) of
Remark~\ref{rem:damp-shelling-universality-scope}, taking \(F\)
disjoint from \(E\).  Its shadow consists of all subsets of \(F\)
of size at most three, and the three acquired supports
\(P_0,P_+,P_-\) have size four.  Define
\begin{equation}
 \begin{split}
 R_0&=L\times\{0_z\},\qquad
 R_+=Q_E\setminus\{0_E\},\qquad
 R_-=Q_E\setminus\{1_E\},\\
 X_K&=(S\times Q_E)
       \cup(\{a_0\}\times R_0)
       \cup(\{a_+\}\times R_+)
       \cup(\{a_-\}\times R_-).
 \end{split}
 \label{eq:damp-reduced-base-universal-construction}
\end{equation}
The three rows are nonempty, proper, and ample.  Proposition
\ref{prop:damp-replicated-repair-layers} makes \(X_K\) ample.
Since all rows are proper, its full layer reduction is
\(X_K^E=S\).  If \(X_K\) had a corner peeling, reverse it to an
ample prefix order and order the complete \(E\)-fibres by the
appearance of their last vertex.  Every resulting prefix is a reduction of an ample prefix, hence is ample
\cite[Section~3.1]{ChalopiChepoiMoran22}.  This would give an ample
prefix order of \(S\), and therefore a peeling of \(S\), a
contradiction.  Thus \(X_K\notin\Damp\), irrespective of the status
of \(K\).
Fixing the old coordinates to \(a_0\) and \(w,p,q,z\) to zero
recovers exactly \(K\).  This is a proper coordinate face, since
the entire nonempty product \(S\times Q_E\) is also present.

If \(K\in\Damp\), then all three occurrence rows peel.  The two
punctured rows contain opposite facets in every layer direction,
exactly as in the proof of
Theorem~\ref{thm:damp-universal-peelable-occurrence-row}.
Theorem~\ref{thm:damp-common-seed-cover-classification} now proves
\(X_K\in\Obs(\Damp)\).  Conversely, if \(X_K\) is forbidden,
its proper coordinate face \(K\) must belong to \(\Damp\).
This proves the equivalence in item~4.

We next verify both reduction stopping conditions directly.  Every
vertex of \(S\times Q_E\) has at least \(m\ge4\) neighbors in
its layer cube.  A vertex at repair tip \(a_b\) has its four
old-coordinate neighbors in \(S\times Q_E\), one for each
coordinate in \(P_b\).  Hence the minimum degree is at least four.

For column distinctness it suffices to distinguish the unsigned
support-incidence columns.  The old clauses of \(X_K\) have supports
\[
 \binom F4\setminus\{P_0,P_+,P_-\}.
\]
Indeed, these are the minimal nonfaces after adding the three acquired
supports to the seed shadow: a five-element set has five four-element
subsets, so cannot become a new minimal nonface after just three
additions.  For distinct \(f,g\in F\), there are
\(\binom{10}{3}=120\) four-element supports containing \(f\)
and omitting \(g\); removing three leaves at least one.  These old
rows distinguish all old coordinates from one another.  Each old
coordinate occurs in an old row, whereas no layer coordinate does,
so they also distinguish old coordinates from layer coordinates.

It remains to distinguish the layer coordinates.  The shadow of
the vertex gluing \(L\) is the union of the shadows of its two
factors.  Every coordinate of \(A\) is active in \(U\), including
\(w\) even when \(n=0\).  Thus \(\{t,p\}\) and \(\{t,q\}\)
are minimal nonfaces of \(L\) for each \(t\in A\); their missing
traces are \(11\).  The constant coordinate \(z\) supplies the
additional minimal nonface \(\{z\}\) of \(R_0\).  By the exact
clause formula
\eqref{eq:damp-replicated-repair-minimal-nonfaces}, \(X_K\) therefore
has clauses on
\[
 P_0\cup\{t,p\},\quad P_0\cup\{t,q\}\quad(t\in A),
 \qquad P_0\cup\{z\}.
\]
For distinct \(t,t'\in A\), the first row for \(t\) distinguishes
them.  The row using \(w,p\) distinguishes \(p\) from \(q\).
The row using \(t,q\) distinguishes \(t\) from \(p\), and the
row using \(t,p\) distinguishes \(t\) from \(q\).  Finally the
last row distinguishes \(z\) from every other layer coordinate.
All support columns are consequently distinct, which excludes equality
of signed columns up to sign.

The product \(S\times Q_E\) already shatters every coordinate
pair, so the crossing graph is complete.  Proposition
\ref{prop:damp-crossing-separator-recovery} excludes cut vertices
and separating edges.  All \(12+m\) coordinates are active.
The disjoint vertex count in
\eqref{eq:damp-reduced-base-universal-construction} gives
\[
 |X_K|=299\,2^m+(2|K|+3)+2(2^m-1)
       =301\,2^m+2|K|+1.
\]
In the shadow formula \eqref{eq:damp-replicated-repair-shadow}, the
seed product contributes maximum rank \(3+m\).  Every proper ample
row has VC dimension at most \(m-1\), so each repair contribution
has rank at most \(4+(m-1)=m+3\).  This proves
\eqref{eq:damp-reduced-base-universality-size}.  When \(K\) peels,
items~2 and~4 put \(X_K\) among the reduced forbidden bases.
Taking \(K=Q_D\) for arbitrarily large \(|D|\) gives pairwise
nonisomorphic such bases with dimension difference nine, proving the
last assertion.
\end{proof}

The verifier \path{verify_damp_reduced_base_universality.py} replays
all \(36\) elementary-minor orders of the fixed seed and the three
repair orders.  It checks the input and enlarged-row peelings for a
singleton, an edge, and the ordinarily peelable \(289\)-concept
rooted factor.  A cornerless \(299\)-concept input checks the opposite
direction: its output is nonpeelable but cannot be forbidden because
that input is a proper coordinate face.  The \(4819\)- and
\(9637\)-concept outputs are materialized and checked from actual
cubes, complete signed presentations, degrees, and recovered faces.
The two larger outputs are checked through the exact formulas and
compact inputs; no output-minor orders or output peeling search are
claimed.  The record is
\path{damp_reduced_base_universality_verification.json}.

This theorem shows that the reduced-base problem retains the full
peelability question even with complete crossing graph and fixed
isometric dimension minus VC dimension.  In particular, reduction does
not leave only finitely many bases at each value of that difference,
nor does it restrict proper coordinate faces to a smaller class of
peelable ample families.  It does not rule out the requested
description by unrestricted finite structural parameters: the missing
step is still intrinsic admissibility of those parameters.

\begin{proposition}[Realizing corner-deletion pairs inside reduced forbidden classes]
\label{prop:damp-corner-pair-universality}
Let $K$ be a nonempty ample family, let $c$ be a corner, and suppose
$K$ and $A=K\setminus\{c\}$ are both nonempty members of $\Damp$.
There are explicit ample forbidden classes $X$ and $X'=X\setminus\{v\}$,
where $v$ is a corner of $X$, and a common coordinate-face restriction
whose images are $K$ and $A$, respectively. Both parents have minimum
degree at least four, distinct signed minimal-nonface columns up to sign,
and complete crossing graph. If $K$ has $n$ active coordinates, then
\[
 |X|=301\,2^{n+4}+2|K|+1,\quad |X'|=|X|-1,
 \quad\operatorname{idim}(X)=\operatorname{idim}(X')=n+16.
\]
\end{proposition}
\begin{proof}
Normalize a vertex of $A$ to zero and use the construction of
Theorem~\ref{thm:damp-reduced-base-universality}. Retain its coordinate
names $w,p,q,z$ and occurrence row
\[
 L=(K\times Q_{\{w\}})\vee_0 Q_{\{p,q\}},
 \qquad R_0=L\times\{0_z\}.
\]
Let $u=(c,1_w,0_p,0_q)$ and replace $L$ by $L'=L\setminus\{u\}$.
The square attached at zero does not meet $u$, since $u_w=1$.
Consequently $u$ is a corner of $L$, with corner cube the product
of the corner cube of $c$ in $K$ and the $w$-edge.
Moreover
\[
 L'=\big((K\times\{0_w\})\cup(A\times\{1_w\})\big)
       \vee_0 Q_{\{p,q\}}.
\]
This row is ample and peelable: delete the three nonzero square vertices,
then the upper copy of $A$ using an ordinary peeling with its lower mates
retained, and finally the lower copy of $K$.

Keep the two other occurrence rows $R_+,R_-$ unchanged. Their opposite
facets still cover every layer direction. The common-seed classification
therefore makes both resulting parents forbidden, since all three rows
are proper, ample, nonempty, and peelable. They differ only at
$v=(a_0,u,0_z)$. At $v$ the corner cube is the product of the old
repair cube of support $P_0$ with the corner cube at $u$ in $L$:
all old-coordinate vertices except the repair tip lie in $S\times Q_E$.
Thus $v$ is a corner, and the modified parent is exactly $X-v$.
Fix the old coordinates to $a_0$, set $w=1$, and set $p=q=z=0$.
The two resulting faces are $K$ and $A$.

The product $S\times Q_E$ remains intact, so its degree bound and
complete crossing graph are unchanged. Every repair-row vertex still
has its four old-coordinate neighbours, giving minimum degree four
in both parents. To check distinct columns, reuse the explicit separating
clauses in the preceding theorem. In $L'$, all input coordinates and $w$
remain active: its lower layer is $K$, and its upper layer is nonempty.
The cross pairs $\{t,p\},\{t,q\}$ for $t\in D\cup\{w\}$ remain
minimal nonfaces with missing trace $11$, and the constant $z$ remains.
Hence the clauses with supports $P_0\cup\{t,p\}$,
$P_0\cup\{t,q\}$ and $P_0\cup\{z\}$ still distinguish every layer
column. The old-coordinate separating clauses are unchanged. This proves
both reduction stopping conditions. The size and dimension formulas
follow from the preceding theorem and the deletion of one vertex.
\end{proof}

Apply this proposition to the all-sinks $289$-vertex class and its
$288$-vertex corner deletion from
Proposition~\ref{prop:damp-all-sinks-corner-failure}. Endpoint loss then
occurs in corresponding proper faces during a corner deletion between
two reduced forbidden parents. The parents have $19\,726\,915$ and
$19\,726\,914$ vertices and $28$ coordinates; these large families are
specified by the formula, not enumerated. The verifier
\path{research/corner_pair_universality.py/json} checks the two positive
occurrence rows of orders $581,580$, the marked corner, and recovery of
the $289,288$ faces. Parent forbiddenness follows symbolically from the
common-seed theorem. Thus merely requiring a forbidden parent, even a
reduced one with complete crossing graph, does not prevent this local
endpoint loss. This does not assert that every proper minor of either
parent has every sink attainable, and by itself does not settle descent under that stronger
hypothesis; Theorem~\ref{thm:damp-all-sinks-corner-terminal} now refutes it. The proposition extends universality to matched deletion
pairs; it does not provide intrinsic generation of the input $K$.

\begin{proposition}[Clause changes inside nonpeelable families]
\label{prop:damp-clause-transport-minimality}
Fix a finite coordinate set \(D\) and a simplicial complex
\(\Sigma\subseteq2^D\) containing every singleton of \(D\).  For each ample
\(K\subseteq Q_D\) with \(0\in K\) and
\(\operatorname{Sh}(K)=\Sigma\), use the same fixed coordinates and
repairs in the construction of
Theorem~\ref{thm:damp-reduced-base-universality}.
All resulting \(X_K\) have the same shadow.  Their inherited input
clauses are exactly
\begin{equation}
 \sigma_{X_K}(P_0\cup I)
   =\bigl(a_0|_{P_0},\sigma_K(I)\bigr)
 \qquad(I\in\mathcal N(\Sigma));
 \label{eq:damp-universal-clause-transport}
\end{equation}
every other output clause is independent of \(K\).
The map \(K\mapsto X_K\) is injective and preserves both the
number of changing clauses and the total number of changing sign bits
between any two inputs.  All outputs are nonpeelable; precisely those
with peelable inputs are forbidden pc-minors.

In particular, there are two nonpeelable ample families with the same
shadow whose presentations differ in one bit of one clause, such that
exactly one is a forbidden pc-minor.  Both have minimum graph degree
at least four, distinct signed columns up to sign, and complete crossing
graph.  Thus an adjacent clause change can destroy pc-minimality while
preserving ampleness, nonpeelability, and both reduction stopping
conditions.
\end{proposition}

\begin{proof}
Retain \(A=D\mathbin{\dot\cup}\{w\}\), \(P=\{p,q\}\),
\(E=A\mathbin{\dot\cup}P\mathbin{\dot\cup}\{z\}\), and
\(L,R_0,R_+,R_-\) from the preceding theorem.
The product \(U=K\times Q_{\{w\}}\) has shadow
\(\{I\cup T:I\in\Sigma,\ T\subseteq\{w\}\}\).
The two factors in the vertex gluing defining \(L\) use disjoint
coordinates and meet at zero.  A trace containing a nonzero coordinate
from each factor is absent; traces supported in either factor are
exactly its traces.  Hence
\[
 \operatorname{Sh}(L)=\operatorname{Sh}(U)\cup2^P.
\]
Every coordinate of \(A\) is active in \(U\).  It follows that
the minimal nonfaces of \(R_0=L\times\{0_z\}\) are precisely
\[
 \mathcal N(\Sigma)\mathbin{\dot\cup}
 \{\{t,p\},\{t,q\}:t\in A\}
 \mathbin{\dot\cup}\{\{z\}\}.
\]
On a support in \(\mathcal N(\Sigma)\), the trace family is that
of \(K\): the extra square contributes only the already present
zero trace.  Its missing trace is therefore \(\sigma_K(I)\).
The missing trace on each displayed pair is \(11\), and that on
\(\{z\}\) is \(1\).  These are all the clauses of \(R_0\).

The other two rows are fixed punctured cubes.  Applying
\eqref{eq:damp-replicated-repair-shadow} shows that the output shadow
depends only on \(\Sigma\).  Equations
\eqref{eq:damp-replicated-repair-minimal-nonfaces}--
\eqref{eq:damp-replicated-repair-layer-signs} now give
\eqref{eq:damp-universal-clause-transport}; all remaining signs are
fixed.  Distinct input supports give distinct output supports, and
adjoining the same old trace changes no sign-bit distance.  This proves
both preservation assertions.  The proper coordinate face recovering
\(K\) proves injectivity.  Nonpeelability of every output and the
exact forbiddenness equivalence are the preceding theorem.

For the last assertion, take the adjacent pair of order \(1188\)
from Corollary~\ref{cor:damp-forbidden-adjacent-sign-pivots}.
Call its peelable member \(K\) and its forbidden member \(G\).
Both contain zero on the same twenty-six active coordinates, and
their only changing clause has support of bitmask \(6531\), with
missing traces \(2306\) and \(2307\), respectively.
Thus \(X_K\) and \(X_G\) differ in one bit of one clause.
Both are nonpeelable and have all the stated degree, column, and
crossing-graph properties.  But \(X_K\) is forbidden, whereas
\(X_G\) has the nonpeelable \(G\) as a proper coordinate face
and is not pc-minimal.  This proves the claimed failure of preservation.
\end{proof}

The checker \path{verify_damp_clause_changes_minimality.py} replays
the positive \(1188\)-concept input order, all \(78\) elementary
orders of the negative input, and the two ordinary factor orders.
It checks the negative endpoint certificates from the unique corner
and the two fixed overlaps with a missing root neighbor.  Actual cube
shadows and full signed presentations verify the input exchange.
The enlarged rows have \(2379\) concepts; their clauses give the
complete \(1589\)-clause presentations of both outputs.  Exactly one
bit of one ten-coordinate clause changes.  The outputs have dimension
\(42\), VC dimension \(33\), and order \(323\,196\,291\,401\);
they are not materialized.  Their nonpeelability and minimality statuses
use the uniform theorem, not an output-minor enumeration.
The record is \path{damp_clause_changes_minimality_verification.json}.

Thus pc-minimality can change between adjacent nonpeelable ample families.
The full signed input variation survives in the outputs.
Intrinsic admissibility of the reduced forbidden bases remains open,
and the minimum-cardinality interval is unchanged.

\begin{lemma}[Descent from the ambient cube]
\label{lem:damp-ambient-cube-descent}
Let \(A\subsetneq Q_E\) be nonempty and ample.  There is a sequence of
corner deletions from \(Q_E\) to \(A\) if and only if
\(Q_E\setminus A\in\Damp\).
\end{lemma}

\begin{proof}
Put \(B=Q_E\setminus A\), which is ample by complement duality.
Suppose that \(b_1,\ldots,b_m\) are the successive deleted concepts in
a descent from \(Q_E\) to \(A\).  Each residual
\(Q_E\setminus\{b_1,\ldots,b_j\}\) is ample, so its complement
\(\{b_1,\ldots,b_j\}\) is ample.  Reversing these ample prefixes gives
a corner peeling \(b_m,\ldots,b_1\) of \(B\).
Conversely, reverse a peeling of \(B\) to obtain ample prefixes
\(B_j=\{b_1,\ldots,b_j\}\).  All their complements are ample.
The ample corner-deletion criterion therefore makes \(b_j\) a corner
of \(Q_E\setminus B_{j-1}\), giving the desired descent.  Here
\(B_0=\varnothing\), and the initial residual is the full cube.
\end{proof}

\begin{theorem}[Prescribed sections in complements of obstructions]
\label{thm:damp-complement-section-construction}
Let \(S\in\Obs(\Damp)\cap\Ample\) be irredundantly embedded in
\(Q_F\), with three resolving one-concept additions
\(a_0,a_+,a_-\) having pairwise distinct acquired supports.
Let \(\varnothing\ne K\subsetneq Q_D\) belong to \(\Damp\), where
\(F\cap D=\varnothing\).  Choose \(u\in Q_D\), and define
\begin{equation}
 \begin{split}
 X_K={}&(S\times Q_D)\cup(\{a_0\}\times K)\\
       &\cup(\{a_+\}\times(Q_D\setminus\{u\}))
        \cup(\{a_-\}\times(Q_D\setminus\{\bar u\})).
 \end{split}
 \label{eq:damp-complement-section-construction}
\end{equation}
Then \(X_K\in\Obs(\Damp)\cap\Ample\), and the section of
\(X_K^c=Q_{F\mathbin{\dot\cup}D}\setminus X_K\) obtained by fixing
the old word to \(a_0\) is exactly \(Q_D\setminus K\).
In particular, if \(Q_D\setminus K\notin\Damp\), then neither
\(X_K\) nor \(X_K^c\) is corner-peelable, and \(X_K\) cannot be
obtained from its ambient cube by corner deletions.  Its order is
\begin{equation}
 |X_K|=(|S|+2)2^{|D|}+|K|-2.
 \label{eq:damp-complement-section-order}
\end{equation}
\end{theorem}

\begin{proof}
All three occurrence rows in
\eqref{eq:damp-complement-section-construction} are nonempty, proper,
ample, and corner-peelable.  For every \(e\in D\), the facet
\(t_e=1-u_e\) lies in the row at \(a_+\), and the facet
\(t_e=u_e\) lies in the row at \(a_-\).  Thus these two rows alone
supply the signed-facet condition of
Theorem~\ref{thm:damp-common-seed-cover-classification}, which proves
that \(X_K\) is an ample forbidden pc-minor.

At the old word \(a_0\), only the row \(K\) occurs in \(X_K\),
so the complementary section is exactly \(Q_D\setminus K\).
This section is a proper pc-minor of the ample family \(X_K^c\):
the latter also contains \((a_+,u)\), outside that section.
If \(X_K^c\) had a peeling, intersecting its reversed ample prefixes
with this section would give an ample prefix order of
\(Q_D\setminus K\).  Thus \(Q_D\setminus K\notin\Damp\) implies
\(X_K^c\notin\Damp\).  Lemma~\ref{lem:damp-ambient-cube-descent}
then rules out the asserted cube descent.  Finally, the four displayed
parts of \(X_K\) have different old words, or old words in \(S\),
and are disjoint; counting them gives
\eqref{eq:damp-complement-section-order}.
\end{proof}

\begin{corollary}[Forbidden minors with nonpeelable cube complements]
\label{cor:damp-nonpeelable-complement-obstructions}
For every integer \(r\ge0\), there is an ample forbidden pc-minor
\(X_r\) of isometric dimension \(24+r\) and order
\[
 |X_r|=1{,}236{,}693\,2^r-2
\]
whose cube complement is nonpeelable.  In particular, ample forbidden
pc-minors are not all obtainable by corner deletions from their ambient
cubes.
\end{corollary}

\begin{proof}
Let \(H\subseteq Q_{12}\) be the fixed \(299\)-concept Hall obstruction
used in Remark~\ref{rem:damp-shelling-universality-scope}.  Its
complement \(K_0=Q_{12}\setminus H\) is corner-peelable: the shelling
of the complementary Hall facets described in
\cite[Section~4.1]{ChalopiChepoiMoran22} gives such a peeling.
An independent complete order is supplied by the fixed check below.
Use \(S=H\), the three resolving additions in that remark, and
\(K=K_0\times Q_r\) on a disjoint set of \(12+r\) coordinates in
Theorem~\ref{thm:damp-complement-section-construction}.
The product test puts \(K\) in \(\Damp\), whereas its complement is
\(H\times Q_r\), which has the nonpeelable pc-minor \(H\).
Since \(|K_0|=4096-299=3797\), the order formula follows from
\eqref{eq:damp-complement-section-order}.  The full product
\(S\times Q_{12+r}\) inside \(X_r\) supplies edges in all
\(24+r\) coordinates, proving the dimension claim.
\end{proof}

\begin{remark}[Verification and scope of the cube-descent obstruction]
\label{rem:damp-cube-descent-verification}
The fixed verifier \path{verify_damp_cube_descent_obstruction.py}
checks the Hall source's exact shattered-set complex and absence of
corners, all \(36\) elementary-minor peeling orders, the three resolving
additions and their distinct supports, and a \(3797\)-concept peeling
of its complement.  It also checks both punctured-cube row orders,
all \(24\) signed facets, and the complementary section identity.
The report \path{damp_cube_descent_obstruction_verification.json}
records hashes of these inputs and the checker's source.  For \(r=0\),
the output has \(1{,}236{,}691\) concepts and its complement has
\(15{,}540{,}525\); neither large family is enumerated.  Their asserted
properties follow from the uniform theorem and the verified inputs.

This closes the proposed exhaustive generation by descent from cubes.
The examples need not be terminal under obstruction-preserving corner
deletion: the row at \(a_0\) can be peeled away while the other two
rows continue to cover all signed facets.  Intrinsic generation with
other primitive data remains open, as does the complete classification.
The minimum-cardinality bounds are unchanged.
\end{remark}

\begin{corollary}[Deletion test for a common-seed cover]
\label{cor:damp-common-seed-obstruction-deletions}
Let \(X(\mathbf R)\in\Obs(\Damp)\) satisfy the hypotheses of
Theorem~\ref{thm:damp-common-seed-cover-classification}.  Fix
\(a\in\mathcal A\) and \(t\in R_a\), and put
\[
 R_a^-=R_a\setminus\{t\}.
\]
In the modified cover, replace \(R_a\) by \(R_a^-\), discarding the label
\(a\) if \(R_a^-=\varnothing\).  Then
\begin{equation}
 X(\mathbf R)\setminus\{(a,t)\}\in\Obs(\Damp)
 \quad\Longleftrightarrow\quad
 \begin{gathered}
 R_a^-=\varnothing\text{ or }R_a^-\in\Damp,\\
 \text{the modified rows own every signed facet of }Q_E.
 \end{gathered}
 \label{eq:damp-common-seed-obstruction-deletion-test}
\end{equation}
In particular, the deletion on the left is necessarily a corner deletion;
this need not be imposed as a separate condition.

If \(S\) is cornerless, a common-seed cover admits no obstruction-preserving corner
deletion if and only if, for every \(a\in\mathcal A\) and \(t\in R_a\),
either \(R_a^-\ne\varnothing\) and \(R_a^-\notin\Damp\), or deleting \(t\)
destroys signed-facet ownership.  This is an exact recognition test on the
row data; the next result replaces it by a structural normal form.
\end{corollary}

\begin{proof}
After deleting \((a,t)\), the family remains a replicated repair assembly
over the same seed, with the single row \(R_a\) replaced by \(R_a^-\).
Proposition~\ref{prop:damp-replicated-repair-layers} says that this family
is ample exactly when the nonempty modified row is ample.  Theorem
\ref{thm:damp-common-seed-cover-classification}, applied after discarding an
empty row, now gives exactly the two conditions on the right of
\eqref{eq:damp-common-seed-obstruction-deletion-test}.  When they hold, both
the original and the deleted family are ample, so the ample
corner-deletion criterion makes \((a,t)\) a corner.  The final
characterization follows by negating the displayed equivalence, since
for a cornerless seed all assembly corners lie in occurrence rows.
\end{proof}

\begin{lemma}[Signed-facet shell and residual core of a row]
\label{lem:damp-row-facet-shell-core}
Let \(R\subsetneq Q_E\) be a proper cube family, and for
\(e\in E\), \(i\in\{0,1\}\), put
\[
 F_{e,i}=\{t\in Q_E:t_e=i\}.
\]
Define its set of full signed facets and the face opposite all of them by
\begin{equation}
 \mathcal O(R)=\{(e,i):F_{e,i}\subseteq R\},
 \qquad
 F_R=\bigcap_{(e,i)\in\mathcal O(R)}F_{e,1-i},
 \label{eq:damp-row-full-facets-opposite-face}
\end{equation}
where an empty intersection is \(Q_E\), and put \(B_R=R\cap F_R\).
Then \(\mathcal O(R)\) contains at most one of \((e,0),(e,1)\) for each
\(e\), and
\begin{equation}
 R=\left(\bigcup_{(e,i)\in\mathcal O(R)}F_{e,i}\right)
   \mathbin{\dot\cup} B_R
  =(Q_E\setminus F_R)\mathbin{\dot\cup}B_R.
 \label{eq:damp-row-facet-shell-decomposition}
\end{equation}
After the coordinates fixed on \(F_R\) are suppressed, one has
\begin{equation}
 R\in\Ample\quad\Longleftrightarrow\quad B_R\in\Ample,
 \qquad
 R\in\Damp\quad\Longleftrightarrow\quad B_R\in\Damp,
 \label{eq:damp-row-facet-shell-tests}
\end{equation}
with the empty family allowed as a vacuous terminal row.  Moreover,
\(B_R\) contains no full signed facet of its residual cube.
\end{lemma}

\begin{proof}
If both \(F_{e,0}\) and \(F_{e,1}\) were contained in \(R\), then
\(R=Q_E\), contrary to properness.  The union of the facets in
\(\mathcal O(R)\) is the complement of their common opposite face
\(F_R\).  This union lies in \(R\), and every remaining member of \(R\)
lies in \(F_R\), proving
\eqref{eq:damp-row-facet-shell-decomposition}.

After reorienting coordinates, write
\(K=\{e:(e,0)\in\mathcal O(R)\}\) and \(L=E\setminus K\).  Then
\[
 R=\bigl((Q_K\setminus\{\mathbf1_K\})\times Q_L\bigr)
   \cup\bigl(\{\mathbf1_K\}\times B_R\bigr).
\]
Choosing one coordinate of \(K\) at a time expresses this family as an
iterated peripheral expansion: one section is a full cube and the other is
the same construction with one fewer shell coordinate.  The exact
peripheral-expansion test therefore gives both equivalences in
\eqref{eq:damp-row-facet-shell-tests}, starting from \(B_R\).  The reverse
implications also follow by restricting all coordinates of \(K\) to~1.

Finally, if \(B_R\) contained a full signed facet in a residual coordinate,
then adjoining the already present shell \(Q_E\setminus F_R\) would make
the corresponding full signed facet of \(Q_E\) a subset of \(R\).  That
facet would belong to \(\mathcal O(R)\), contradicting that its coordinate
is residual.
\end{proof}

\begin{theorem}[Terminal covers as signed-facet partitions]
\label{thm:damp-terminal-common-seed-facet-partitions}
Let \(X(\mathbf R)\in\Obs(\Damp)\) satisfy the hypotheses of
Theorem~\ref{thm:damp-common-seed-cover-classification}.  For each
\(a\in\mathcal A\), put
\[
 \mathcal O_a=\{(e,i)\in E\times\{0,1\}:F_{e,i}\subseteq R_a\}.
\]
If \(X(\mathbf R)\) is corner-descent terminal, then
\begin{enumerate}
 \item the sets \((\mathcal O_a)_{a\in\mathcal A}\) form a partition of
       \(E\times\{0,1\}\); and
 \item every occurrence row is exactly its signed-facet shell:
       \begin{equation}
        R_a=\bigcup_{(e,i)\in\mathcal O_a}F_{e,i}.
        \label{eq:damp-terminal-row-facet-shell}
       \end{equation}
\end{enumerate}
If \(S\) is cornerless, these two conditions are also sufficient.
In particular, each block \(\mathcal O_a\) is nonempty and contains at
most one sign of any coordinate.

Conversely, retain cornerlessness and the seed and repair-tip hypotheses of
Theorem~\ref{thm:damp-common-seed-cover-classification}.  Let
\((\mathcal O_a)_{a\in\mathcal A}\) be any partition of
\(E\times\{0,1\}\) into nonempty blocks, no block containing both
\((e,0)\) and \((e,1)\), and define the rows by
\eqref{eq:damp-terminal-row-facet-shell}.  Then the resulting
\(X(\mathbf R)\) is a corner-descent-terminal forbidden pc-minor.  The
marked sections recover every block canonically as the set of full signed
facets of its row.
\end{theorem}

\begin{proof}
Assume first that \(X(\mathbf R)\) is corner-descent terminal.  Apply
Lemma~\ref{lem:damp-row-facet-shell-core} to a row \(R_a\), and suppose
that its residual core \(B_{R_a}\) is nonempty.  Because \(R_a\in\Damp\),
the lemma gives \(B_{R_a}\in\Damp\).  Choose the first concept \(t\) in a
corner peeling of this residual core.  The same lemma, applied after
deleting \(t\), shows that \(R_a\setminus\{t\}\in\Damp\).  Moreover,
\(t\) lies in the opposite face \(F_{R_a}\), hence in none of the full
signed facets of \(R_a\).  Deleting it therefore leaves every signed-facet
containment in the cover unchanged.  Corollary
\ref{cor:damp-common-seed-obstruction-deletions} would then make
\(X(\mathbf R)\setminus\{(a,t)\}\) another forbidden pc-minor, contrary
to terminality.  Thus every residual core is empty, proving
\eqref{eq:damp-terminal-row-facet-shell}.

The signed-facet ownership condition in
Theorem~\ref{thm:damp-common-seed-cover-classification} says that the
sets \(\mathcal O_a\) cover \(E\times\{0,1\}\).  Suppose that one signed
facet \(F_{e,i}\) belongs to both \(\mathcal O_a\) and
\(\mathcal O_b\), with \(a\ne b\).  Choose a word \(t\in R_a\) which has
value \(i\) at \(e\) and, for every other \((f,j)\in\mathcal O_a\), has
value \(1-j\) at \(f\).  Thus \(t\) lies in exactly one maximal facet of
the union~\eqref{eq:damp-terminal-row-facet-shell}, namely \(F_{e,i}\),
and is a corner of \(R_a\).  Deleting \(t\) leaves a corner-peelable row:
after the other shell facets are removed by
Lemma~\ref{lem:damp-row-facet-shell-core}, its residual row is a punctured
cube in the coordinates not fixed by \(\mathcal O_a\).  All other facets
of \(\mathcal O_a\) remain full, while \(F_{e,i}\) is still owned by
\(R_b\).  The deletion corollary again produces a forbidden pc-minor,
contradicting terminality.  Hence the cover is disjoint and the
\(\mathcal O_a\) form a partition.

Conversely, assume that \(S\) is cornerless and start with such a partition.  A row in
\eqref{eq:damp-terminal-row-facet-shell} is the complement of the face
opposite its block.  Lemma~\ref{lem:damp-row-facet-shell-core}, with empty
residual core, makes it ample and corner-peelable.  The partition owns
every signed facet, so Theorem
\ref{thm:damp-common-seed-cover-classification} gives
\(X(\mathbf R)\in\Obs(\Damp)\).

The maximal cubes of a row are exactly its full signed facets.  Hence a
row concept is a corner precisely when it lies in exactly one of them.
Deleting such a concept destroys that signed facet in its row; the
partition says that no other row owns it.  The deletion criterion therefore
shows that no corner deletion preserves membership in \(\Obs(\Damp)\).
Thus the assembly is corner-descent terminal.  Finally,
\(\mathcal O_a\) is read directly from \(R_a\), proving canonical
recovery.
\end{proof}

\begin{remark}[What the facet-partition theorem closes]
\label{rem:damp-terminal-facet-partition-scope}
For cornerless seeds,
Theorem~\ref{thm:damp-terminal-common-seed-facet-partitions} removes both
row peelability and the row-deletion test from the terminal common-seed
parameters.  The occurrence data are only a partition of the signed layer
facets into consistent nonempty blocks; soundness and terminality follow
uniformly from the facet shells.  For two blocks, consistency forces the
two blocks to choose opposite signs in every coordinate, recovering the
antipodally punctured two-row towers.

This is a relative generative theorem, not the full classification.  The
seed is still assumed to be a cornerless forbidden pc-minor, and it is not
yet intrinsic which one-concept extensions of the seed are admissible repair
tips.  Universal common-seed completeness is false by
Theorem~\ref{thm:damp-terminal-cut-vertex}.  The remaining intrinsic
debts are seed classification, repair-tip admissibility, and a complete
description of the additional terminal types.  The structure of terminal
occurrence rows within the common-seed class is settled.
\end{remark}

\begin{proposition}[Thickness of terminal facet partitions]
\label{prop:damp-terminal-facet-thickness}
Use the construction in
Theorem~\ref{thm:damp-terminal-common-seed-facet-partitions}, put
\(r=|E|\), and write \(K_a\) for the coordinates represented in the
owner block \(\mathcal O_a\).  Denote by \(\lambda(e,i)\) the unique
owner of \((e,i)\).  For the two projected sections \(H_i=\rho_e^iX\)
of a layer coordinate \(e\), one has
\begin{equation}
 |H_i\setminus H_{1-i}|=2^{r-|K_{\lambda(e,i)}|}
 \qquad(i=0,1).
 \label{eq:damp-terminal-facet-thickness}
\end{equation}
Thus a layer coordinate has \(|H_0\triangle H_1|=2\) if and only if
there are exactly two repair tips.  In that case every layer coordinate
has this property.  If there are at least three repair tips, then every
coordinate of \(X\), including every old coordinate, has two nonempty
exclusive sections whose total order is at least three.
\end{proposition}

\begin{proof}
For \(a=\lambda(e,i)\), the \(i\)-section of its row is the full
\((r-1)\)-cube.  The opposite section misses the face fixing each
coordinate in \(K_a\setminus\{e\}\) opposite to its owned sign.
That face has \(2^{r-|K_a|}\) vertices.  Apart from the two owners of
\(e\), all rows have equal \(e\)-sections.  This proves
\eqref{eq:damp-terminal-facet-thickness}.
Both exclusive orders are one exactly when the two owners each represent
every coordinate.  They then exhaust the partition of the \(2r\) signed
facets, leaving no third owner.  Conversely, a partition into two
consistent blocks assigns one sign of every coordinate to each block.

For an old coordinate \(f\), the seed is periphery-free by
Proposition~\ref{prop:damp-periphery-free}.  Choose
\(s_i\in\rho_f^iS\setminus\rho_f^{1-i}S\), for \(i=0,1\).
The product core supplies \(2^r\) occurrences of \(s_i\) in the
\(i\)-section.  These remain exclusive unless the opposite old word is
a repair tip \(b_i\).  If it is a tip, its row fills all but
\(2^{r-|K_{b_i}|}\) of those occurrences.  Thus each exclusive section
is nonempty.  If their total order were two, both selected contributions
would be one.  Since \(r\ge1\), both opposite words would have to be
repair tips, and both owner supports would equal \(E\).  The two tips
are distinct because their \(f\)-values differ.  Their blocks would
exhaust all signed facets, again excluding a third owner.
\end{proof}

\begin{remark}[A terminal obstruction with every coordinate thick]
\label{ex:damp-thick-terminal-three-repair}
Use Hall's seed \(C_H\) and the three repair tips
\(a_0=183\), \(a_+=315\), \(a_-=1836\) from
Remark~\ref{rem:damp-shelling-universality-scope}.  On
\(E=\{p,q\}\), assign
\[
 \lambda(p,0)=\lambda(q,0)=a_0,\qquad
 \lambda(p,1)=a_+,\qquad \lambda(q,1)=a_-.
\]
The rows are \(R_0=Q_E\setminus\{(1,1)\}\),
\(R_+=\{t:t_p=1\}\), and \(R_-=\{t:t_q=1\}\).
The facet-partition theorem gives a corner-descent-terminal ample
forbidden pc-minor with
\[
 |X|=4\cdot299+3+2+2=1203,\qquad \operatorname{idim}(X)=14.
\]
Proposition~\ref{prop:damp-terminal-facet-thickness} proves that all
fourteen coordinates are thick.  Consequently thickness does not separate
terminal obstructions inside this construction from those outside it.
\end{remark}

\paragraph{Recovering layer coordinates from forbidden traces.}
The signed minimal nonfaces distinguish layer coordinates more directly
than section sizes do.  Write \(\eta_a(e)=i\) when
\((e,i)\in\mathcal O_a\), and put \(C=S\cup\mathcal A\).
The row \(R_a\) misses exactly the face with trace \(1-\eta_a\) on
\(K_a\).  Its shattered complex therefore has the single minimal
nonface \(K_a\), with that missing trace.  The replicated-repair formulas
\eqref{eq:damp-replicated-repair-minimal-nonfaces}--
\eqref{eq:damp-replicated-repair-layer-signs} give
\begin{align}
 \mathcal N(\operatorname{Sh}(X))
 &=\mathcal N(\operatorname{Sh}(C))
   \mathbin{\dot\cup}\{P_a\cup K_a:a\in\mathcal A\},
 \label{eq:damp-terminal-facet-clauses}\\
 \sigma_X(P_a\cup K_a)&=(a|_{P_a},1-\eta_a),
 \label{eq:damp-terminal-facet-clause-signs}\\
 \sigma_X(J)&=\sigma_C(J)
       \quad(J\in\mathcal N(\operatorname{Sh}(C))).\notag
\end{align}
These identities require only an ample seed and compatible one-concept
extensions; the seed need not be an obstruction.
Every layer coordinate occurs in exactly two members of this signed
clause family, with opposite missing-trace values.

\begin{definition}[Coordinates occurring twice with opposite signs]
\label{def:damp-balanced-two-clause-coordinates}
For a nonempty ample family \(X\subseteq Q_V\), put
\(\mathcal J=\mathcal N(\operatorname{Sh}(X))\), with missing traces
\((\sigma_J)_{J\in\mathcal J}\), and define
\[
 \mathcal B(X)=\left\{e\in V:
   \begin{array}{l}
    \exists J_0,J_1\in\mathcal J:\\
    \{J\in\mathcal J:e\in J\}=\{J_0,J_1\},\\
    J_0\ne J_1,\quad \sigma_{J_0}(e)\ne\sigma_{J_1}(e)
   \end{array}\right\}.
\]
For a nonempty \(E\subseteq\mathcal B(X)\), put \(F=V\setminus E\),
\(\mathcal J_E=\{J\in\mathcal J:J\cap E\ne\varnothing\}\), and
\begin{align}
 C_E&=\{x\in Q_F:x|_J\ne\sigma_J
                    \text{ for every }J\in\mathcal J\text{ with }J\subseteq F\},
 \label{eq:damp-clause-recovery-old-family}\\
 P_J&=J\cap F,\qquad K_J=J\cap E,\qquad
 T_J=\{x\in C_E:x|_{P_J}=\sigma_J|_{P_J}\}
        \quad(J\in\mathcal J_E).
 \label{eq:damp-clause-recovery-fibres}
\end{align}
The family \(C_E\) is defined by the old clauses alone.  The theorem below
identifies it with the projection of \(X\) onto \(F\).
\end{definition}

\begin{theorem}[Recovery from signed clauses and singleton trace fibres]
\label{thm:damp-terminal-clause-recovery}
A marked nonempty ample family \((X,E)\), with \(E\ne\varnothing\),
has a common-seed presentation
with a nonempty ample seed, distinct one-concept repair supports, and rows
given by a partition of the signed \(E\)-facets if and only if
\(E\subseteq\mathcal B(X)\) and every \(T_J\), \(J\in\mathcal J_E\),
is a singleton \(\{a_J\}\), with these singletons pairwise distinct.
In that case \(C_E\) is the projection of \(X\) onto \(F\), and
\[
 S_E=C_E\setminus\{a_J:J\in\mathcal J_E\}=X^E
\]
is automatically nonempty and ample.  Every \(S_E\cup\{a_J\}\) is ample,
its unique new shattered support is \(P_J\), and these supports are
pairwise distinct.  Thus the recovery criterion uses only signed clauses
and their trace fibres; seed and repair ampleness require no additional
tests.  The presentation is unique for the
marked set \(E\): its tips are \(a_J\), its owner supports are \(K_J\),
and its owner signs are \(1-\sigma_J|_{K_J}\).

If \(X\) is an irredundantly embedded ample forbidden pc-minor, then
\(S_E\) is itself a forbidden pc-minor and every
\(S_E\cup\{a_J\}\) is corner-peelable.  If the recovered seed is also
cornerless, this presentation belongs to the terminal common-seed
obstruction class of
Theorem~\ref{thm:damp-terminal-common-seed-facet-partitions}.
\end{theorem}

\begin{proof}
First suppose such a presentation is given.  Equations
\eqref{eq:damp-terminal-facet-clauses}--
\eqref{eq:damp-terminal-facet-clause-signs} show that
\(E\subseteq\mathcal B(X)\), that \(C_E=C=S\cup\mathcal A\), and
that its mixed clauses correspond bijectively to the repairs.
For the clause associated with \(a\), the missing trace on \(P_a\)
is supplied in \(C\) only by \(a\).  Indeed, the union-shadow formula
\eqref{eq:damp-replicated-repair-union-shadow} shows that
\(C\setminus\{a\}\) is ample and
does not shatter \(P_a\); it contains \(S\), which already supplies
every other trace on that minimal nonface.  Hence \(T_J=\{a\}\).

Conversely, assume the distinct-singleton condition and
\(E\subseteq\mathcal B(X)\).  Write \(C=C_E\),
\(S=C\setminus\{a_J:J\in\mathcal J_E\}\).
We first identify the seed and its one-tip extensions as ample families,
then show that their acquired supports are exactly the old clause parts.

Projection onto \(F\) preserves precisely the shattered sets contained
in \(F\).  Its minimal nonfaces are consequently the members of
\(\mathcal J\) contained in \(F\), with the same missing traces.
This projection is ample, so
Theorem~\ref{thm:damp-fixed-shadow-signed-clutter} identifies it with
\(C\).  In particular, \(C\) is nonempty and ample.

A word \((x,t)\) satisfying
all signed clauses of \(X\) must have \(x\in C_E\).
If \(x\in S_E\), it matches none of the old parts of the mixed
forbidden traces, so every \(t\in Q_E\) is allowed.
If \(x=a_J\), the distinct-singleton condition says that the only
mixed clause it can violate is \(J\).  Its allowed row is therefore
\[
 R_J=\{t\in Q_E:t|_{K_J}\ne\sigma_J|_{K_J}\}
     =\bigcup_{e\in K_J}\{t:t_e=1-\sigma_J(e)\}.
\]
The same reconstruction theorem now gives
\[
 X=(S_E\times Q_E)\cup
       \bigcup_{J\in\mathcal J_E}(\{a_J\}\times R_J).
\]
Every \(K_J\) is nonempty.  Each \(e\in E\) occurs in exactly two
mixed clauses with opposite signs, so these displayed owner blocks
partition all signed \(E\)-facets.  Every row is proper, so \(S=X^E\).
Reductions of ample families are ample: if \(X^E\) shatters
\(D\subseteq F\), then \(X\) shatters \(D\cup E\), hence contains
a cube on that support; its reduction is a \(D\)-cube in
\(X^E\).  The same argument proves
\begin{equation}
 D\in\operatorname{Sh}(S)
 \quad\Longleftrightarrow\quad
 D\cup E\in\operatorname{Sh}(X)
 \qquad(D\subseteq F).
 \label{eq:damp-clause-recovery-strong-shadow}
\end{equation}

The seed cannot be empty.  Otherwise \(X\) would not strongly shatter
\(E\), hence would not shatter it, so some \(J\in\mathcal J\) would
satisfy \(J\subseteq E\).  Then \(P_J=\varnothing\) and \(T_J=C\).
But any \(e\in E\) belongs to two distinct mixed clauses, whose
distinct singleton fibres already give \(|C|\ge2\), a contradiction.

For each \(J\), restrict \(X\) to a signed facet owned by \(J\) and
strongly project the remaining layer coordinates.  The row \(R_J\)
becomes full.  For any other row, its missing face still meets the
chosen facet: its owner block does not contain that signed facet.
Thus the resulting family is exactly \(S\cup\{a_J\}\).  Restrictions
and reductions preserve ampleness, proving that each one-tip
extension is ample.

By Lemma~\ref{lem:damp-one-concept-extension-test}, write \(Q_J\) for
its unique acquired support, a minimal nonface of
\(\operatorname{Sh}(S)\).  Since \(T_J=\{a_J\}\), adjoining \(a_J\)
adds exactly one word to the projection of \(S\) onto \(P_J\).
Both projected families are ample, so their shadows differ by one set.
The one-tip shadow formula \eqref{eq:damp-one-extension-shadow} therefore
forces \(Q_J\subseteq P_J\).

It remains to exclude a proper inclusion and coincident acquired supports.
By \eqref{eq:damp-clause-recovery-strong-shadow}, \(Q_J\cup E\) is not
shattered in \(X\), so it contains some minimal nonface \(I\in\mathcal J\).
This cannot be an old clause: \(Q_J\) is shattered in
\(S\cup\{a_J\}\subseteq C\), and hence in the projection of \(X\).
Consequently \(I\in\mathcal J_E\), and
\[
 Q_I\subseteq P_I\subseteq Q_J.
\]
Minimality of \(Q_I,Q_J\) as nonfaces of \(\operatorname{Sh}(S)\)
gives \(Q_I=P_I=Q_J\).  Every tip \(a_L\) with \(Q_L=Q_J\)
supplies the unique missing trace of \(S\) on this support.  That trace
is \(a_I|_{P_I}=\sigma_I|_{P_I}\), so every such tip lies in \(T_I\).
Since \(T_I=\{a_I\}\) and all tips are distinct, necessarily \(L=I\).
In particular \(J=I\), which proves \(Q_J=P_J\), and no other tip
has that acquired support.  This completes the presentation proof.
All its data were forced by the clauses, proving uniqueness.

Finally, suppose \(X\) is a forbidden pc-minor.
Each row \(R_J\) is a product of a singly punctured cube on \(K_J\)
and a cube on \(E\setminus K_J\), hence is corner-peelable.
Equation~\eqref{eq:damp-replicated-repair-tip-cubes} lifts these row
peelings to successive corner deletions in \(X\), leaving
\(S_E\times Q_E\).  If \(S_E\) were peelable, the product test
\eqref{eq:damp-exact-product-test} would finish a peeling of \(X\),
a contradiction.  Thus \(S_E\) is nonpeelable.
Every intermediate reduction
is nonpeelable: otherwise its further reduction \(S_E\) would
be peelable by Proposition~\ref{prop:damp-strong-projection-closure}.
Repeated application of Theorem~\ref{thm:damp-inherited-overlap-descent}
therefore makes \(S_E\) a forbidden pc-minor, after suppressing any
inactive coordinates.

For a repair \(a_J\), choose a signed facet owned by its block.
Restrict \(X\) to that facet and strongly project all the other layer
coordinates.  Its row becomes full.  No other row becomes full, since
no other block owns that signed facet.  Thus the resulting family is
exactly \(S_E\cup\{a_J\}\).  The initial restriction is a proper
pc-minor of \(X\), hence is peelable, and reduction closure
makes \(S_E\cup\{a_J\}\) peelable as well.
There can be no inactive old coordinate in \(S_E\): irredundancy of
\(X\) would require a tip on its opposite side, whose one-tip extension
would have the nonpeelable \(S_E\) as a proper restriction.
If \(S_E\) is cornerless, the hypotheses of the terminal facet-partition
theorem now all hold.
\end{proof}

\begin{theorem}[One-concept repairs of a recovered facet partition]
\label{thm:damp-recovered-partition-resolving-tips}
Let \(X\subseteq Q_V\) be an irredundantly embedded ample forbidden
pc-minor, and let \(E\) satisfy the singleton recovery criterion of
Theorem~\ref{thm:damp-terminal-clause-recovery}.  Retain its recovered
seed \(S\), tips \(a_J\), supports \(P_J\), and owner supports \(K_J\).
Then
\begin{equation}
 \{z\in Q_V\setminus X:X\cup\{z\}\in\Damp\}
 =\{(a_J,\sigma_J|_E):J\in\mathcal J_E,\ K_J=E\}.
 \label{eq:damp-recovered-partition-resolving-tips}
\end{equation}
The one-concept extension at the displayed word acquires the support
\(P_J\cup E=J\).  Thus \(X\) has at most two one-concept extensions
in \(\Damp\).  It has two exactly when its recovered presentation has
two repair tips and their rows are oppositely punctured layer cubes.

If the recovered seed \(S=X^E\) itself has a nonempty layer set \(D\)
satisfying the singleton recovery criterion, then both presentations have
exactly two repair tips and flatten into one antipodally punctured
two-repair presentation with layer set \(E\cup D\) and seed \(S^D\).
In particular, \(E\cup D\) satisfies the singleton recovery criterion
for \(X\).
\end{theorem}

\begin{proof}
Suppose \(X\cup\{z\}\in\Damp\).  Its reduction along \(E\)
must be peelable by Proposition~\ref{prop:damp-strong-projection-closure},
whereas \(S=X^E\) is not.  Adding \(z\) must therefore complete a full
\(E\)-fibre.  Since \(E\ne\varnothing\), an old word outside the
projection of \(X\) cannot acquire a full fibre from this single addition.
An old word in \(S\) already has a full fibre.  The only possibility is
an old repair word \(a_J\) whose row is missing exactly one layer word.
The complement of that row is a face of order
\(2^{|E|-|K_J|}\), so necessarily \(K_J=E\) and
\(z=(a_J,\sigma_J|_E)\).

Conversely, add this word for a clause with \(K_J=E\).
The completed row joins the product seed, giving
\((S\cup\{a_J\})\times Q_E\).  Peel each other occurrence row using
\eqref{eq:damp-replicated-repair-tip-cubes}; all these rows are products
of punctured cubes with cubes.  The recovered one-tip extension
\(S\cup\{a_J\}\) is peelable by
Theorem~\ref{thm:damp-terminal-clause-recovery}, so the product test
\eqref{eq:damp-exact-product-test} finishes a peeling.
The old-coordinate neighbours of the added word have support \(P_J\),
and every layer neighbour is present.  The local extension test
therefore identifies its acquired support as \(P_J\cup E\).

For any fixed \(e\in E\), only two owner blocks contain \(e\).
Hence at most two blocks can have support \(E\).  If two do, they
exhaust all signed facets, so there are no other blocks; their signs
are opposite in every coordinate.  This proves the count and the
two-row characterization.

For the last assertion, put \(T=S^D\) and apply the result just proved
to the recovered presentation of the obstruction \(S\).
The outer presentation of \(X\) supplies at least two distinct
one-concept repairs of \(S\), since every layer coordinate has two
different owners.  Therefore these are exactly the two possible repairs
of \(S\).  Its inner presentation has the form
\[
 S=(T\times Q_D)
   \cup(\{a\}\times(Q_D\setminus\{u\}))
   \cup(\{b\}\times(Q_D\setminus\{\bar u\})),
\]
and its two resolving concepts are \((a,u)\) and \((b,\bar u)\).
These must be all the tips of the outer presentation.  Its two rows
are consequently \(Q_E\setminus\{v\}\) and
\(Q_E\setminus\{\bar v\}\) for some \(v\in Q_E\).
Substituting the inner presentation into the outer one gives
\begin{align*}
 X={}&(T\times Q_{D\cup E})\\
 &\cup\bigl(\{a\}\times
       (Q_{D\cup E}\setminus\{(u,v)\})\bigr)\\
 &\cup\bigl(\{b\}\times
       (Q_{D\cup E}\setminus\{(\bar u,\bar v)\})\bigr).
\end{align*}
The inner repair supports are distinct, so this is a presentation of
the kind recognized by Theorem~\ref{thm:damp-terminal-clause-recovery}.
That theorem proves the asserted success of \(D\cup E\).
\end{proof}

\begin{remark}[A terminal obstruction with no one-concept repair]
\label{rem:damp-terminal-no-resolving-tip}
Use the Hall seed and the three tips \(a_0,a_+,a_-\) from
Remark~\ref{ex:damp-thick-terminal-three-repair}.  On
\(E=\{p,q,r\}\), assign their owner blocks as
\[
 \mathcal O_0=\{(p,0),(q,0)\},\qquad
 \mathcal O_+=\{(q,1),(r,0)\},\qquad
 \mathcal O_-=\{(p,1),(r,1)\}.
\]
The facet-partition theorem gives a terminal ample forbidden pc-minor
of order \(8\cdot299+3\cdot6=2410\) and dimension \(15\).
Every owner support is a proper subset of \(E\), so
Theorem~\ref{thm:damp-recovered-partition-resolving-tips} shows that no
one-concept extension of this obstruction is corner-peelable.
For comparison, the order-\(1203\) example has exactly one such extension,
obtained by completing the row of \(a_0\), and an antipodal two-repair
tower has exactly two.  The fixed verifier
\path{verify_damp_facet_partition_repairs.py} checks all three cases and
records explicit peelings of every positive extension.  Hence a
completeness argument cannot assume that a terminal obstruction itself
has two resolving tips; that property belongs to the seed of a recovered
presentation.
\end{remark}

\begin{remark}[Scope of common-seed factorization]
\label{rem:damp-clause-recovery-completeness}
For an ample forbidden pc-minor \(X\), let \(\mathcal E(X)\) consist
of the nonempty subsets of \(\mathcal B(X)\) satisfying the distinct-singleton
trace-fibre condition in
Theorem~\ref{thm:damp-terminal-clause-recovery} and yielding a cornerless
seed.  The universal claim that \(\mathcal E(X)\ne\varnothing\)
for every cornered, corner-descent-terminal \(X\) is refuted by
Theorem~\ref{thm:damp-terminal-cut-vertex}: that example has no
nontrivial successful layer set at all.  Even the restriction to
two-connected terminal obstructions fails, by the edge amalgam in
Proposition~\ref{prop:damp-terminal-edge-amalgam}.  The recovery theorem still
gives an explicit map for every successful choice.  Even when a successful choice exists, the largest recovered seed
need not be cornerless: Proposition~\ref{prop:damp-terminal-cut-seed-towers}
gives two-connected terminal counterexamples.  Theorem~\ref{thm:damp-obstruction-layer-union}
and Corollary~\ref{cor:damp-cornerless-recovery-maximal} below show that
\(\mathcal E(X)\) has at most one member: only the largest successful
layer set can yield a cornerless seed.  In particular, a cornered terminal obstruction with
\(\mathcal B(X)=\varnothing\) cannot have such a presentation.
The mere existence of a coordinate in \(\mathcal B(X)\) is only
necessary; distinct singleton fibres and cornerlessness of the recovered
seed must still be proved.  Nonemptiness and ampleness of that seed,
ampleness of each one-tip extension, and distinctness of the acquired
supports follow from the singleton condition.
Theorem~\ref{thm:damp-recovered-partition-resolving-tips} gives another
direct falsifier: a cornered terminal obstruction with three distinct
one-concept extensions in \(\Damp\) cannot have such a presentation.
Theorem~\ref{thm:damp-layer-recovery-dichotomy} sharpens this test:
two resolving extensions already suffice if every coordinate is thick.
This does not yet provide intrinsic admissibility for arbitrary forbidden
seeds or their resolving repairs.
\end{remark}


\begin{theorem}[Forced enlargement of candidate layer sets]
\label{thm:damp-forced-layer-closure}
Let \(X\subseteq Q_V\) be nonempty and ample, and use the signed clauses
and balanced coordinate set \(\mathcal B(X)\) of
Definition~\ref{def:damp-balanced-two-clause-coordinates}.
Write \(T_J^E\) for the fibre in
\eqref{eq:damp-clause-recovery-fibres} when the marked set is \(E\).
For a nonempty family of words \(T\subseteq Q_{V\setminus E}\), put
\[
 \operatorname{var}(T)=
 \{f\in V\setminus E:\exists x,y\in T\text{ with }x_f\ne y_f\},
 \qquad
 W(E)=\bigcup_{J\in\mathcal J_E}\operatorname{var}(T_J^E).
\]
Starting from any nonempty \(E_0\subseteq\mathcal B(X)\), perform these
steps with \(E=E_0\):
\begin{enumerate}
 \item If \(W(E)\not\subseteq\mathcal B(X)\), stop with failure.
 \item If two distinct mixed clauses have intersecting fibres
       \(T_I^E\cap T_J^E\ne\varnothing\), stop with failure.
 \item If \(W(E)=\varnothing\), return \(E\).
       Otherwise replace \(E\) by \(E\cup W(E)\) and repeat.
\end{enumerate}
The procedure terminates after at most
\(|\mathcal B(X)|-|E_0|+1\) passes.
On success it returns the unique smallest superset of \(E_0\) satisfying
the singleton recovery criterion of
Theorem~\ref{thm:damp-terminal-clause-recovery}.
On failure, no superset of \(E_0\) satisfies that criterion.
The successful sets are closed under nonempty intersections.

The first failure has a certificate consisting of a mixed clause,
two words in its fibre, and a coordinate outside \(\mathcal B(X)\)
on which they differ.  The second has a certificate consisting of two
mixed clauses and a word in both fibres.  Thus existence of any nonempty
layer set satisfying the singleton criterion is decided by starting the
procedure at each of the \(|\mathcal B(X)|\) one-coordinate sets.
This assertion does not impose cornerlessness on the recovered seed.
\end{theorem}

\begin{proof}
For every \(E\subseteq\mathcal B(X)\), the family \(C_E\) is the
projection of \(X\) onto \(V\setminus E\): its minimal forbidden
supports are exactly the old clauses, and
Theorem~\ref{thm:damp-fixed-shadow-signed-clutter} identifies its words.
If \(J\in\mathcal J_E\), then \(J\setminus E\) is a proper subset
of a minimal nonface, hence is shattered by \(X\).  In particular,
its prescribed trace occurs in \(C_E\), so \(T_J^E\ne\varnothing\).

Suppose \(E\subseteq D\subseteq\mathcal B(X)\).  Projecting a word
of \(T_J^E\) onto \(V\setminus D\) gives a word in \(C_D\) that
still matches \(\sigma_J\) on \(J\setminus D\).  Therefore
\begin{equation}
 \{x|_{V\setminus D}:x\in T_J^E\}\subseteq T_J^D
 \qquad(J\in\mathcal J_E).
 \label{eq:damp-layer-fibre-projection-inclusion}
\end{equation}
If \(D\) satisfies the singleton criterion, all coordinates on which
two words in \(T_J^E\) differ must consequently belong to \(D\).
Thus \(W(E)\subseteq D\).  Furthermore, any word in two distinct
fibres at \(E\) projects to a word in the corresponding two fibres at
\(D\), contrary to their being distinct singletons.

It follows inductively that every successful superset of \(E_0\)
contains every set produced by the procedure.  Either failure therefore
excludes every such superset, with the stated certificate.  Otherwise
each nonterminal pass adds at least one coordinate of \(\mathcal B(X)\).
At termination all fibres are nonempty with no varying coordinate, hence
are singletons; the second test makes them distinct.  The returned set
is therefore successful and is contained in every successful superset
of \(E_0\), proving existence, minimality, and uniqueness.
Finally, any nonempty successful set contains some coordinate \(e\),
so the procedure started at \(\{e\}\) must succeed.  The converse is
immediate from the success statement.
For a nonempty intersection \(E_0\) of successful sets, the procedure
must return a superset of \(E_0\) contained in each of those sets.
It therefore returns \(E_0\) itself, proving the intersection assertion.
\end{proof}

\begin{theorem}[A largest recovered layer set for an ample obstruction]
\label{thm:damp-obstruction-layer-union}
Let \(X\subseteq Q_V\) be an irredundantly embedded ample forbidden
pc-minor.  Call a nonempty \(E\subseteq\mathcal B(X)\) successful if
it satisfies the singleton criterion of
Theorem~\ref{thm:damp-terminal-clause-recovery}, without requiring its
recovered seed to be cornerless.  The successful sets are closed under
union.  Consequently, if any successful set exists, there is a unique
largest one, given by
\begin{equation}
 E_{\max}(X)=
 \{e\in\mathcal B(X):\text{the forced-enlargement procedure
 started at }\{e\}\text{ succeeds}\}.
 \label{eq:damp-largest-recovered-layer-set}
\end{equation}
Thus the one-coordinate starts of
Theorem~\ref{thm:damp-forced-layer-closure} recover this largest set
directly.  Its seed \(X^{E_{\max}(X)}\) is determined by the marked
family \(X\), independently of a choice of presentation.
\end{theorem}

\begin{proof}
We first control how signed clauses change when a layer presentation is
recovered.  We then treat disjoint successful sets; reduction of
their intersection will reduce the general case to that one.

Every coordinate of an irredundantly embedded ample forbidden pc-minor
occurs with both missing signs among its minimal forbidden traces.
Indeed, if a coordinate occurs with at most one sign, the signed-clause
description in Theorem~\ref{thm:damp-fixed-shadow-signed-clutter} makes
one section contain the other.  Both nonempty sections are proper
pc-minors and hence peelable.  The peripheral-expansion test
\eqref{eq:damp-exact-peripheral-expansion-test} would make the whole
family peelable, a contradiction.

Fix a successful \(E\), put \(F=V\setminus E\), and write its
presentation with seed \(S=X^E\), repair tips \(a\), acquired supports
\(P_a\), and owner supports \(K_a\subseteq E\).
The seed is an irredundantly embedded ample forbidden pc-minor by
Theorem~\ref{thm:damp-terminal-clause-recovery}.  Each minimal nonface
\(I\) of \(\operatorname{Sh}(S)\) supplies a distinct clause of
\(X\): it supplies \(I\) itself if it is not an acquired support,
and supplies \(I\cup K_a\) if \(I=P_a\).  In the first case \(I\)
remains a minimal nonface after the selected supports are adjoined to the
shadow.  The clause formulas
\eqref{eq:damp-replicated-repair-minimal-nonfaces}--
\eqref{eq:damp-replicated-repair-layer-signs} prove the assertion and
preserve all signs on \(I\).  Since every seed coordinate occurs with
both signs, an old coordinate occurring in exactly two clauses of \(X\)
must already occur in exactly two oppositely signed clauses of \(S\).
In particular,
\begin{equation}
 \mathcal B(X)\setminus E\subseteq\mathcal B(S).
 \label{eq:damp-balanced-coordinate-seed-inheritance}
\end{equation}

Suppose first that \(D\) is another successful set and \(D\cap E
=\varnothing\).  Then \(D\subseteq\mathcal B(S)\).  Let \(I\)
be a seed clause meeting \(D\).  Because \(I\setminus D\) is a
proper subset of a minimal nonface, some word \(x\) in the projection
of \(S\) onto \(F\setminus D\) has its prescribed missing trace.
The projection of \(X\) along \(D\) contains
\(\{x\}\times Q_E\), supplied by the seed product.  If the clause
supplied by \(I\) is the unchanged old clause \(I\), its recovery
fibre for \(D\) therefore contains all these words, contradicting
the singleton criterion.  Hence \(I=P_a\) for an outer repair.  Its
clause \(P_a\cup K_a\) has at least
\(2^{|E\setminus K_a|}\) words in that fibre, obtained by fixing the
missing trace on \(K_a\).  Thus \(K_a=E\).

Choose any \(f\in D\).  Its two incident seed clauses consequently
give two distinct owner blocks whose supports are all of \(E\).
They exhaust the signed \(E\)-facets.  The outer presentation therefore
has exactly two tips, say \(a,b\), and antipodally punctured rows.
Moreover, the only seed clauses meeting \(D\) are \(P_a,P_b\).
Every coordinate of \(D\) belongs to both, with opposite missing signs.

Consider the recovery fibres for these two clauses in the projection of
\(S\) along \(D\).  Each is nonempty by the proper-subset shattering
argument above.  Appending the missing \(E\)-word of its respective
outer row injects it into the corresponding singleton recovery fibre of
\(X\) for \(D\).  Hence these seed fibres are singletons, say
\(\{x_a\}\) and \(\{x_b\}\).  They are distinct.  Otherwise the
restriction of \(S\) obtained by fixing \(F\setminus D\) to their
common word would be
\[
 Q_D\setminus\{u,\bar u\}.
\]
Indeed, the old clauses hold at this projected word, and the only two
clauses meeting \(D\) exclude precisely the two opposite words.  For
\(|D|=1\) this contradicts membership in the projection.  For
\(|D|\ge2\) this restriction is not ample: it has \(2^{|D|}-2\)
words but shatters all \(2^{|D|}-1\) proper subsets of \(D\).
Restrictions of ample families are ample, giving the contradiction.
Thus \(D\) satisfies the distinct-singleton criterion for \(S\).
The nested-recovery conclusion of
Theorem~\ref{thm:damp-recovered-partition-resolving-tips} now shows
that \(E\cup D\) is successful for \(X\).

For arbitrary successful \(E,D\), containment makes the conclusion
immediate.  Otherwise put \(I=E\cap D\), \(E'=E\setminus I\),
\(D'=D\setminus I\), and \(Y=X^I\).  In the presentation marked
by \(E\), strongly projecting a row along \(I\) leaves the union
of its owned facets whose coordinates lie in \(E'\); the row vanishes
if it owns none of these facets.  Discard those empty rows.  The surviving
tips remain distinct compatible repairs, and the surviving owner blocks
partition all signed \(E'\)-facets.  This gives a presentation of
\(Y\) with layer set \(E'\) and seed \(X^E\), so \(E'\) is
successful.  The presentation marked by \(D\) likewise proves success
of \(D'\) in \(Y\).

The family \(Y\) is irredundant: its seed \(X^E\) uses all its old
coordinates, and the nonempty product with \(Q_{E'}\) uses the others.
It is nonpeelable because its further reduction \(X^E\) is
nonpeelable.  Reduction closure of \(\Damp\) and repeated
application of Theorem~\ref{thm:damp-inherited-overlap-descent} make
\(Y\) an ample forbidden pc-minor; the same seed product keeps the
remaining coordinates active at each intermediate projection.
Apply the disjoint argument to
\(E',D'\) in \(Y\).  It proves that \(D'\) is successful in
\(Y^{E'}=X^E\).  Nested recovery in the original \(X\) then makes
\(E\cup D'=E\cup D\) successful.

Finally, the union of all successful sets is successful because \(V\)
is finite.  A coordinate belongs to that union exactly when it belongs
to some successful set, which the forced-enlargement theorem decides
by starting at its singleton.  This proves
\eqref{eq:damp-largest-recovered-layer-set} and uniqueness.
\end{proof}

\begin{theorem}[All successful layer choices and the thin coordinates]
\label{thm:damp-layer-recovery-dichotomy}
Let \(X\subseteq Q_V\) be an irredundantly embedded ample forbidden
pc-minor.  Write \(\mathcal L(X)\) for its family of successful layer
sets, without requiring a cornerless seed, and define its thin
coordinates by
\[
 \mathcal T(X)=\{e\in V:
   |\rho_e^0X\setminus\rho_e^1X|
   =|\rho_e^1X\setminus\rho_e^0X|=1\}.
\]
Then exactly one of the following alternatives holds.
\begin{enumerate}
 \item If \(\mathcal T(X)\ne\varnothing\), then
 \begin{equation}
   E_{\max}(X)=\mathcal T(X),\qquad
   \mathcal L(X)=2^{\mathcal T(X)}\setminus\{\varnothing\}.
   \label{eq:damp-thin-coordinate-largest-recovery}
 \end{equation}
 The largest presentation has two repair tips and antipodally punctured
 rows.  The obstruction \(X\) has exactly two one-concept extensions
 in \(\Damp\).
 \item If \(\mathcal T(X)=\varnothing\), then either
 \(\mathcal L(X)=\varnothing\), or
 \(\mathcal L(X)=\{E_{\max}(X)\}\).  In the latter case
 \(|E_{\max}(X)|\ge2\), its presentation has at least three repair
 tips, and \(X\) has at most one one-concept extension in \(\Damp\).
\end{enumerate}
In particular, a cornered terminal obstruction with every coordinate
thick and two distinct one-concept extensions in \(\Damp\) has
no common-seed presentation.
\end{theorem}

\begin{proof}
First consider any recovered presentation marked by \(E\), with seed
\(S\), tips \(\mathcal A\), and clauses
\(J_a=P_a\cup K_a\).  For a nonempty \(D\subsetneq E\), put
\(L=E\setminus D\).  The recovery fibre of a clause meeting \(D\)
has the exact form
\begin{equation}
 T_{J_a}^{D}
 =\{a\}\times
   \{t\in Q_L:t|_{K_a\cap L}=\sigma_{J_a}|_{K_a\cap L}\}
 \qquad(K_a\cap D\ne\varnothing).
 \label{eq:damp-sublayer-recovery-fibre}
\end{equation}
Indeed, in the full old projection \(C=S\cup\mathcal A\), only
\(a\) has the missing trace on \(P_a\), as proved in
Theorem~\ref{thm:damp-terminal-clause-recovery}.  Since
\(K_a\cap D\ne\varnothing\), projecting its row along \(D\)
fills every \(L\)-word: choose a coordinate of \(K_a\cap D\)
opposite to its missing sign.  The remaining clause condition fixes
exactly the displayed coordinates of \(L\), proving the formula.
The fibres for distinct tips are disjoint.  Thus \(D\) is successful
if and only if
\[
 L\subseteq K_a\qquad
 \text{for every }a\text{ with }K_a\cap D\ne\varnothing.
\]

If this condition holds, choose a coordinate of \(D\) and let \(a,b\)
be its two owners.  Both owner supports contain all of \(L\), so these
two blocks exhaust the signed \(L\)-facets.  Any other owner meeting
\(D\) would also have to own a facet in each coordinate of \(L\),
which is impossible.  No other owner can be supported inside \(L\)
either.  Hence there are exactly two owners.  Conversely, with exactly
two owners both supports equal \(E\), and the displayed fibre formula
makes every nonempty \(D\subseteq E\) successful.  We have proved:
a recovered presentation admits a proper nonempty successful subset
exactly when it has two owners, whenever \(|E|\ge2\).

A singleton \(\{e\}\) is successful exactly when \(e\) is thin.
One direction follows immediately from its two one-point rows; the other
is the thin-coordinate recovery in
Proposition~\ref{prop:damp-thin-coordinate-fragile-kernel}.
If a thin coordinate exists, the largest successful set exists by
Theorem~\ref{thm:damp-obstruction-layer-union}.  Its singleton subset
forces two owners by the preceding argument, unless the largest set is
already a singleton, in which case its two signed facets give two owners
directly.  Every nonempty subset of that largest set is therefore
successful, and its singleton members identify it with
\(\mathcal T(X)\).

If no thin coordinate exists but recovery is possible, the largest
presentation cannot have two owners.  It has at least three, so the
fibre argument rules out every proper nonempty successful subset.
Every successful set lies in the largest one, proving uniqueness and
\(|E_{\max}(X)|\ge2\).  Finally,
Theorem~\ref{thm:damp-recovered-partition-resolving-tips} identifies
the resolving additions with owner blocks supported on the entire layer
set.  There are exactly two such blocks in the first case and at most
one when there are at least three owners.  This proves the extension
counts and the stated falsifier.
\end{proof}

\begin{corollary}[Cornerless recovery requires a maximal layer set]
\label{cor:damp-cornerless-recovery-maximal}
If \(E\) satisfies the singleton recovery criterion and \(X^E\) is
cornerless, then no proper superset of \(E\) satisfies that criterion.
If \(X\) is an irredundantly embedded ample forbidden pc-minor and
\(E\) is maximal among successful layer sets, its recovered seed
\(X^E\) has no nonempty layer set satisfying the singleton criterion.
In particular, it has no thin coordinate.  That seed is an obstruction
with at least two one-concept repairs in
\(\Damp\), and can be cornered as in
Remark~\ref{rem:damp-cornered-maximal-hall-seed}.
For such an obstruction \(X\), a successful recovery with cornerless
seed exists if and only if \(E_{\max}(X)\) exists and
\(X^{E_{\max}(X)}\) is cornerless.
\end{corollary}

\begin{proof}
Suppose a proper superset \(D\) also satisfies the criterion, and use its
recovered presentation with seed \(S_D\), tips \(a_J\), owner supports
\(K_J=J\cap D\), and rows \(R_J\).  Put \(L=D\setminus E\).
Strongly projecting the rows along \(E\) gives
\[
 X^E=(S_D\times Q_L)\cup
 \bigcup_{J:K_J\cap L\ne\varnothing}
 \bigl(\{a_J\}\times
       \{t\in Q_L:t|_{K_J\cap L}\ne\sigma_J|_{K_J\cap L}\}\bigr).
\]
Indeed, a row survives the reduction exactly when one of its
owned facets uses a coordinate outside \(E\); the displayed remaining
facets then describe it.  Each surviving row is nonempty, proper, and a
product of a singly punctured cube with a cube.  At least one survives
because \(L\ne\varnothing\) and every signed \(L\)-facet has an owner.
Such a row has a corner, which lifts to a corner of \(X^E\) by
\eqref{eq:damp-replicated-repair-tip-cubes}, a contradiction.
For the second assertion, a successful layer set in the recovered seed
would enlarge \(E\) by
Theorem~\ref{thm:damp-recovered-partition-resolving-tips}, contradicting
maximality.  Obstruction descent and the two distinct facet owners at any
\(e\in E\) give the remaining assertions.
The seed has no thin coordinate because any such coordinate would give
a successful singleton by Theorem~\ref{thm:damp-layer-recovery-dichotomy}.
The last equivalence follows because
Theorem~\ref{thm:damp-obstruction-layer-union} makes the maximal
successful set unique.
\end{proof}

\begin{theorem}[Seed-corner deletions in a signed-facet assembly]
\label{thm:damp-seed-corner-facet-deletion}
Let \(S\in\Obs(\Damp)\cap\Ample\) and the compatible resolving tips
\(\mathcal A\) satisfy
Theorem~\ref{thm:damp-common-seed-cover-classification}; the seed may have
corners.  Let \(E\ne\varnothing\), let
\[
 \lambda:E\times\{0,1\}\longrightarrow\mathcal A
\]
be a surjection with \(\lambda(e,0)\ne\lambda(e,1)\), and put
\[
 R_a=\bigcup_{\lambda(e,i)=a}\{t\in Q_E:t_e=i\},
 \qquad
 X=(S\times Q_E)\cup
       \bigcup_{a\in\mathcal A}(\{a\}\times R_a).
\]
Then \(X\) is an ample forbidden pc-minor.  For each
\(s\in\operatorname{Cor}(S)\), define
\[
 G_s=S\setminus\{s\},\qquad
 \mathcal A_s=\{a\in\mathcal A:
       a\text{ is not adjacent to }s\text{ in }Q_F,
       \quad G_s\cup\{a\}\in\Damp\}.
\]
Every obstruction-preserving corner deletion of \(X\) removes a seed
concept \((s,t)\), and its exact criterion is
\begin{equation}
 \begin{aligned}
 &(s,t)\in\operatorname{Cor}(X),\quad
       X\setminus\{(s,t)\}\in\Obs(\Damp)\\
 &\qquad\Longleftrightarrow\quad
 s\in\operatorname{Cor}(S),\quad G_s\in\Obs(\Damp),\quad
 \lambda(e,t_e)\in\mathcal A_s\quad\text{for every }e\in E.
 \end{aligned}
 \label{eq:damp-seed-corner-facet-deletion}
\end{equation}
Here \(s\in S\) and \(t\in Q_E\).  For each fixed \(s\) with
\(G_s\in\Obs(\Damp)\), its allowable layer words form the possibly
empty face
\begin{equation}
 D_s=\prod_{e\in E}\{i\in\{0,1\}:\lambda(e,i)\in\mathcal A_s\}.
 \label{eq:damp-seed-deletion-layer-face}
\end{equation}
Consequently \(X\) is corner-descent terminal if and only if
\begin{equation}
 \begin{gathered}
 \text{for every }s\in\operatorname{Cor}(S)
       \text{ with }G_s\in\Obs(\Damp),\\
 \text{some }e\in E\text{ satisfies }
       \lambda(e,0),\lambda(e,1)\notin\mathcal A_s.
 \end{gathered}
 \label{eq:damp-seed-corner-opposite-owner-test}
\end{equation}
In particular, every corner-descent-terminal seed gives terminal
signed-facet assemblies, including when that seed has corners.
\end{theorem}

\begin{proof}
The common-seed criterion proves that \(X\) is forbidden: each facet
shell is peelable and the rows own every signed facet.  We first locate
the seed corners and show that occurrence-row peelings survive their
deletion.  Reductions will then prove necessity, while projected
row peelings prove sufficiency for all elementary pc-minors.

Every maximal seed cube through \(s\), multiplied by \(Q_E\), remains
maximal in \(X\), because all repair rows are proper.  Thus a corner
\((s,t)\) requires \(s\in\operatorname{Cor}(S)\).  All layer directions
are incident at \((s,t)\).  If a repair tip adjacent to \(s\) occurs at
\(t\), its old-coordinate direction would also have to belong to the
corner cube.  That cube would contain the tip in every layer, contradicting
properness of its row.  Conversely, if no adjacent repair occurs at
\(t\), the incident directions are precisely those of the corner cube
of \(s\) in \(S\), together with \(E\), and their full product is present.
We have proved
\begin{equation}
 (s,t)\in\operatorname{Cor}(X)
 \quad\Longleftrightarrow\quad
 s\in\operatorname{Cor}(S),\quad
 t\notin R_a\text{ for every repair }a\text{ adjacent to }s.
 \label{eq:damp-seed-corner-occurrences}
\end{equation}

Fix such a corner and write \(H=X\setminus\{(s,t)\}\).
The family \(G_s\) is ample and nonpeelable: a peeling of \(G_s\)
could be preceded by \(s\) to peel \(S\).  It still uses every old
coordinate.  Otherwise it would be a proper signed restriction of
\(S\) and hence peelable.  In particular, \(G_s\times Q_E\subseteq H\)
makes the embedding of \(H\) irredundant as well.

We will use the following observation about deleting occurrence rows
from \(H\).  If the repair cube
\(Q_a=\{a^U:U\subseteq P_a\}\) contains \(s\), then \(a\) is adjacent
to \(s\).  Indeed, at distance at least two from \(a\), all
\(P_a\)-neighbors of \(s\) lie in the punctured cube contained in
\(S\).  Since \(s\) is a corner of \(S\), its corner cube would then
contain the entire \(P_a\)-cube through \(s\), including \(a\), which
is impossible.  Thus whenever \(s\in Q_a\),
\eqref{eq:damp-seed-corner-occurrences} gives \(t\notin R_a\).
Every cube \(Q_a\times K\), with \(K\subseteq R_a\), therefore avoids
the deleted concept.  The row peelings supplied by
\eqref{eq:damp-replicated-repair-tip-cubes} remain valid in \(H\), in
any order and with any chosen rows retained.  Applying a restriction or
contraction to this sequence again gives corner deletions whenever its
image changes: consecutive images are ample and differ by one concept.

Suppose now that \(H\) is forbidden.  Its strong \(E\)-projection is
exactly \(G_s\), which is nonpeelable.  Repeated inherited-overlap descent
in Theorem~\ref{thm:damp-inherited-overlap-descent} makes \(G_s\) a
forbidden pc-minor.  For each \(e\in E\), put \(a=\lambda(e,t_e)\).
Its row contains \(t\), so \(a\) is not adjacent to \(s\) by
\eqref{eq:damp-seed-corner-occurrences}.  Restricting \(H\) to the
facet \(t_e\) and strongly projecting the remaining layer coordinates
gives exactly \(G_s\cup\{a\}\): only the owner row becomes full, and
the \(s\)-fibre is missing \(t\).  The restriction is a proper
pc-minor and hence peelable; reduction closure makes
\(G_s\cup\{a\}\) peelable.  This proves necessity in
\eqref{eq:damp-seed-corner-facet-deletion}.

Conversely, assume its right-hand side.  Each adjacent tip is excluded
from \(\mathcal A_s\), so none of its owned facets contains \(t\).
Equation~\eqref{eq:damp-seed-corner-occurrences} makes \((s,t)\) a
corner; thus \(H\) is ample.  Its nonpeelable reduction \(G_s\)
makes \(H\) nonpeelable by
Proposition~\ref{prop:damp-strong-projection-closure}.

Consider an elementary old-coordinate map \(\mu\).  Peel every
occurrence row in \(H\), and project this sequence by \(\mu\), using
the observation above.  The remaining family is \(\mu Z\), where
\[
 Z=(S\times Q_E)\setminus\{(s,t)\}
   =(G_s\times Q_E)\cup
      (\{s\}\times(Q_E\setminus\{t\})).
\]
The base \(\mu G_s\) is peelable, since \(G_s\) is an irredundant
forbidden pc-minor.  If \(\mu\) erases \(s\) or identifies it with a
word of \(\mu G_s\), then \(\mu Z=(\mu G_s)\times Q_E\) peels.
Otherwise \(\mu S=(\mu G_s)\cup\{\mu s\}\) is an ample one-concept
extension, so \(\mu s\) is a corner.  Peel its punctured-cube occurrence
row and then \((\mu G_s)\times Q_E\).  This proves that \(\mu H\)
peels in every old coordinate.

Every layer contraction fills the single deleted concept from its
\(s\)-mate, and every layer restriction to the side opposite \(t\)
is unchanged from \(X\).  These images therefore peel.  For the
remaining restriction \(\nu=\rho_e^{t_e}\), put
\(a=\lambda(e,t_e)\in\mathcal A_s\).  Peel all occurrence rows except
\(a\) in \(H\), and project the sequence by \(\nu\).  Because the
\(a\)-row becomes full, its final family is
\[
 ((S\cup\{a\})\times Q_{E\setminus\{e\}})
       \setminus\{(s,t|_{E\setminus\{e\}})\}.
\]
Nonadjacency of \(a,s\) leaves the corner cube of \(s\) unchanged
in \(S\cup\{a\}\).  Its deletion is the peelable family
\(G_s\cup\{a\}\).  If no layer coordinate remains, this is already
the displayed final family.  Otherwise peel the punctured layer row of
\(s\), then the product over \(G_s\cup\{a\}\).  All elementary
pc-minors of \(H\) now peel, completing the sufficiency proof.

Finally, an occurrence-row corner in a facet shell lies in exactly one
of its maximal facets.  Deleting it destroys that facet's unique owner.
The row-deletion criterion
\eqref{eq:damp-common-seed-obstruction-deletion-test} therefore rules out
every obstruction-preserving occurrence-row deletion.  Formula
\eqref{eq:damp-seed-corner-facet-deletion} accounts for all remaining
corners.  Its choices of \(t_e\) are independent, giving the face
\eqref{eq:damp-seed-deletion-layer-face}.  This face is empty exactly
when some coordinate offers no owner in \(\mathcal A_s\), proving
\eqref{eq:damp-seed-corner-opposite-owner-test}.  A terminal seed has
no corner \(s\) with \(G_s\) forbidden, so the last assertion follows.
\end{proof}

For a nonadjacent pair \(s,a\) in this theorem, deleting \(s\) leaves
the punctured repair cube at \(a\) intact.  Thus \(G_s\cup\{a\}\)
is automatically ample and acquires the same support \(P_a\).
The set \(\mathcal A_s\) asks exactly which of these inherited ample
repairs remain resolving.  Formula
\eqref{eq:damp-seed-corner-opposite-owner-test} is therefore a complete
terminality test on the signed-facet data once these seed-deletion and
repair data are known.  The necessary direction of
Theorem~\ref{thm:damp-terminal-common-seed-facet-partitions} requires
these facet partitions in every terminal common-seed cover, so the two
results characterize terminality throughout that class, including
cornered seeds.  This is not intrinsic seed admissibility: it still
uses \(G_s\in\Obs(\Damp)\) and \(G_s\cup\{a\}\in\Damp\).
The fixed verifier \path{verify_damp_seed_corner_facet_deletion.py}
checks the corner formula on the \(602\)- and \(1202\)-concept Hall
assemblies.  For each of the two deletions of a copy of \(373\) from
the first family, it records peelings of all \(39\) elementary images.
For the cornered terminal \(600\)-concept seed of the nested second
family, it checks that each seed-corner deletion has Hall's obstruction
as a proper layer restriction.  These controls check both surviving
deletions and terminal-seed inheritance.  The next construction settles
the global cornerless-recovery question negatively.

\begin{proposition}[Largest layer recovery can retain a cut vertex]
\label{prop:damp-terminal-cut-seed-towers}
Let \(S\subseteq Q_F\) be an irredundantly embedded,
corner-descent-terminal ample forbidden pc-minor with a cut vertex.  Suppose that \(a,b\notin S\) are
distinct resolving tips: both \(S\cup\{a\}\) and
\(S\cup\{b\}\) belong to \(\Damp\), and their acquired shattered
supports are distinct.  For a nonempty coordinate set \(E\) disjoint
from \(F\), put
\begin{equation}
 X_E=(S\times Q_E)
     \cup\bigl(\{a\}\times(Q_E\setminus\{0^E\})\bigr)
     \cup\bigl(\{b\}\times(Q_E\setminus\{1^E\})\bigr).
 \label{eq:damp-terminal-cut-seed-tower}
\end{equation}
Then \(X_E\) is a two-connected, corner-descent-terminal ample forbidden
pc-minor.  Its unique largest successful layer set is exactly \(E\),
and its recovered seed is \(X_E^E=S\).  In particular, if \(S\)
has corners, no successful layer choice for \(X_E\) gives a
cornerless seed.  Every \(X_E\) is already its graph three-core.
\end{proposition}

\begin{proof}
The two punctured rows are a signed-facet partition, with opposite
facets assigned to different tips.  The common-seed construction
theorem and Theorem~\ref{thm:damp-seed-corner-facet-deletion} show that
\(X_E\) is ample, forbidden, and terminal.  It is two-connected by
Proposition~\ref{prop:damp-replicated-two-connected}.  The displayed
presentation also makes \(E\) a successful layer set and gives
\(X_E^E=S\).

By Theorem~\ref{thm:damp-obstruction-layer-union}, there is a largest
successful set \(D\), and \(E\subseteq D\).  Suppose the inclusion
is strict, and consider the facet-partition presentation marked by
\(D\), with seed \(T=X_E^D\).  Strongly projecting this presentation
along \(E\) retains its product seed \(T\times Q_{D\setminus E}\).
An occurrence row becomes the union of its owned facets whose
coordinates lie in \(D\setminus E\); it vanishes if there are no
such facets.  Discard the vanishing rows.  The remaining owner blocks
still partition all signed facets of \(Q_{D\setminus E}\), and
each remaining row is nonempty and proper.  There are at least two
rows, since opposite facets have distinct owners.  Their tips are
distinct resolving tips of \(T\), by the recovery theorem.
Thus \(S=X_E^E\) has a nontrivial replicated common-seed presentation
with resolving tips.  Proposition~\ref{prop:damp-replicated-two-connected}
would make \(S\) two-connected, contrary to its cut vertex.  Hence
\(D=E\).

Any successful choice giving a cornerless seed must be the largest
one, by Corollary~\ref{cor:damp-cornerless-recovery-maximal}.  This
proves the assertion when \(S\) is cornered.  Finally, a degree-two
vertex of an ample forbidden minor is a corner whose deletion preserves
forbiddenness, by Theorem~\ref{thm:damp-three-core}.  Terminality excludes
such a vertex; the same theorem excludes smaller degrees.  Hence
\(X_E\) has minimum degree at least three and equals its three-core.
\end{proof}

\begin{remark}[Two-connected terminal towers]
\label{ex:damp-cornered-maximal-seed-2318}
Use \(S=Y\vee C\) from
Remark~\ref{ex:damp-terminal-cut-vertex-1158}.  Its two corners are in
\(Y\), and its copy of \(C\) uses coordinates \(14,\ldots,25\).
In that copy, take
\[
 a=128\cdot2^{14},\qquad b=290\cdot2^{14}.
\]
These are resolving tips for \(S\).  Indeed,
\(S\cup\{a\}=Y\vee(C\cup\{128\})\) is ample.  First peel
\(C\cup\{128\}\) to the common root, using its prescribed endpoint
peeling, then peel \(Y\) ordinarily.  The same argument applies to
\(b\).  The acquired supports have bitmasks \(390\cdot2^{14}\)
and \(480\cdot2^{14}\), respectively, so they are distinct.
All hypotheses of Proposition~\ref{prop:damp-terminal-cut-seed-towers}
therefore hold.

For \(|E|=r\ge1\), the resulting family has
\[
 |X_E|=1160\cdot2^r-2,\qquad
 \dim(X_E)=26+r,\qquad \operatorname{VCdim}(X_E)=4+r.
\]
Its minimum degree is \(r+3\), and it has exactly
\(2^{r+1}+2r\) corners.  To see the corner count, neither new tip
is adjacent to either seed corner: the tips and seed corners lie in
the two disjoint coordinate blocks.  Both seed corners therefore lift
to all \(2^r\) layers, by
\eqref{eq:damp-seed-corner-occurrences}.  Each punctured row has
exactly \(r\) corners, namely the neighbors of its missing word.
These row corners have degree \(4+(r-1)=r+3\), while every seed
concept has degree at least \(4+r\).  The shadow formula
\eqref{eq:damp-replicated-repair-shadow} gives the stated VC dimension.

Thus the first examples have \(2318\) and \(4638\) concepts, with
six and twelve corners.  Their largest recovered seed is the same
\(1158\)-concept terminal family, with two corners and a cut vertex.
This refutes cornerlessness of the largest recovered seed even for
two-connected terminal obstructions of minimum degree at least three.

The checker \path{verify_damp_cornered_maximal_seed.py} constructs both
examples and replays all \(165\) elementary-minor peelings, as well as
four resolving-extension peelings.  It checks every initial corner
deletion by exhibiting a proper minor whose reduction is either
the nonpeelable \(1158\)-concept seed or the cornerless
\(577\)-concept double.  It also checks two-connectivity, exact
shadows, and that the only thin coordinates are the new layer
coordinates, which independently identifies the largest layer set by
Theorem~\ref{thm:damp-layer-recovery-dichotomy}.  The report is
\path{damp_cornered_maximal_seed_verification.json}.
\end{remark}

The preceding family still admits layer recovery.  The construction below
shows that even the existence of such a recovery fails for two-connected,
cornered terminal obstructions.  Its separator is an edge shared by two
coordinate blocks.  The input conditions concern peeling to this entire
edge, not just to one of its endpoints.

\begin{proposition}[Terminal obstructions from an edge amalgam]
\label{prop:damp-terminal-edge-amalgam}
Let \(C\subseteq Q_F\) be an irredundantly embedded, two-connected
ample family which is corner-peelable and has \(0\) as its only corner.
Choose \(e\in F\) with \(D=\{0,0^{\{e\}}\}\subseteq C\).
A peeling \emph{to \(D\)} deletes corners outside \(D\) until precisely
\(D\) remains.  Suppose that
\begin{enumerate}
 \item for \(f\ne e\), the restriction \(\rho_f^0 C\) and the
 contraction \(\pi_f C\) peel to their copies of \(D\);
 \item each of \(\rho_e^0 C,\rho_e^1 C,\pi_e C\) peels to the
 singleton image of \(D\); and
 \item there are distinct \(a,b\notin C\) such that \(C\cup\{a\}\)
 and \(C\cup\{b\}\) are ample, peel to \(D\), and acquire distinct
 shattered supports \(P_a,P_b\).
\end{enumerate}
Take copies \(C_1,C_2\) whose coordinate sets \(F_1,F_2\) intersect
only in \(e\), identifying their edges \(D\).  Put zeros in the
coordinates outside each copy.  With two new coordinates \(p,q\), form
\begin{equation}
 X=C_2\ \cup\
   \bigl(C_1\times(Q_{\{p,q\}}\setminus\{00\})\bigr)
   \ \cup\ \{(a_1,10),(b_1,01)\},
 \label{eq:damp-terminal-edge-amalgam}
\end{equation}
where \(C_2\) occupies the \(00\)-section.  Then \(X\) is a
two-connected, corner-descent-terminal ample forbidden pc-minor, of order
\(4|C|+2\) and dimension \(2|F|+1\).  Its only corners are
\((a_1,10)\) and \((b_1,01)\).

If \(|F|\ge4\) and \(e\) belongs to at least three minimal nonfaces
of \(\operatorname{Sh}(C)\), no coordinate of \(X\) is balanced in
the sense of Definition~\ref{def:damp-balanced-two-clause-coordinates}.
Consequently \(X\) has no successful common-seed layer set, even with
cornered recovered seeds allowed.
\end{proposition}

\begin{proof}
We first establish ampleness and ordinary peelings of the pieces.  Then
we check the three kinds of elementary minors, locate all corners, and
exclude balanced coordinates.  Write \(L=Q_{\{p,q\}}\setminus\{00\}\)
and, temporarily working with just the first old coordinate block, put
\[
 Y=(D\times\{00\})\cup(C\times L)
       \cup\{(a,10),(b,01)\}.
\]
Compatibility of the two acquired supports gives
\(\operatorname{Sh}(C\cup\{a,b\})
 =\operatorname{Sh}(C)\cup\{P_a,P_b\}\).
The shadows of the two sections in either new coordinate intersect
in \(\operatorname{Sh}(C)\).  A support using both new coordinates
must already be shattered in the \(00\)-section \(D\), and that
condition is sufficient since \(D\times Q_{\{p,q\}}\subseteq Y\).
It follows that \(\operatorname{Sh}(Y)\) is the disjoint union of
\[
 \operatorname{Sh}(C)\cup\{P_a,P_b\},\quad
 \{I\cup\{p\}:I\in\operatorname{Sh}(C)\},\quad
 \{I\cup\{q\}:I\in\operatorname{Sh}(C)\},\quad
 \{I\cup\{p,q\}:I\in\operatorname{Sh}(D)\}.
\]
Its cardinality is \(3|C|+4=|Y|\), proving ampleness.
In \(X=Y\cup C_2\), a support using a coordinate of
\((F_1\setminus\{e\})\cup\{p,q\}\) and a coordinate of
\(F_2\setminus\{e\}\) cannot be shattered: their joint trace
\(11\) is absent.  Thus its shadow is precisely the union of the
two block shadows, with intersection \(\operatorname{Sh}(D)\).
Again shadow cardinality equals family cardinality, proving ampleness
of \(X\) and its asserted order.  Every indicated coordinate is active.

There is an ordinary peeling of \(Y\): delete the two concepts in
\(D\times\{00\}\), then its two added tips, and then peel
\(C\times L\) using any ordinary peeling of \(C\) and the path
peeling \(01,11,10\) of \(L\).  The first two concepts lie on a
full \(D\times Q_{\{p,q\}}\)-cube and have no additional incident
directions in their section.  The tip corner cubes lie in their own
sections.  These observations justify the initial deletions as well
as the subsequent product peeling.

We now construct peelings of all elementary minors of \(X\).
A corner outside the shared edge in either block remains a corner of
the union: it has no additional neighbors in the other block.  Hence
a relative peeling of one block to the intersection can be followed
by an ordinary peeling of the other.

For a coordinate \(f\in F_2\setminus\{e\}\), restriction to
\(f=1\) discards \(Y\) and leaves a peelable restriction of \(C_2\).
Restriction to zero or contraction first peels the changed \(C_2\)
to \(D\), by item~1, and then peels \(Y\).
For \(f\in F_1\setminus\{e\}\), restriction to one similarly
leaves a peelable restriction of \(Y\).  Under restriction to zero
or contraction, write \(C'=\mu C\).  Delete any surviving images
of the tips that lie outside \(C'\times L\); these are corner
deletions, being the changing images of the original tip deletions.
The remaining family is
\((D\times\{00\})\cup(C'\times L)\).
Use the peeling of \(C'\) to \(D\) from item~1, deleting each
outside concept in the three layers in path-peeling order.  Finally
peel \(D\times Q_{\{p,q\}}\) to \(D\times\{00\}\), and follow
with an ordinary peeling of \(C_2\).

For a new coordinate, restriction to one again discards \(C_2\)
and is a restriction of the peelable family \(Y\).  Restriction to
zero leaves \(D\) in one layer and either \(C\cup\{a\}\) or
\(C\cup\{b\}\) in the other.  Item~3 peels the latter to its
copy of \(D\); delete that copy and then peel \(C_2\).
Under contraction, the two layers are \(C\cup\{a\}\) and
\(C\cup\{b\}\), with the intersection edge in one of them.
Delete the added tip in the layer not containing that edge, peel
its remaining copy of \(C\) ordinarily, and peel the other layer
to \(D\) using item~3.  All corner cubes in the first copy of
\(C\) lift across the two layers, since the other layer still
contains all of \(C\).  Finish with \(C_2\).

For any of the three elementary operations in the shared coordinate
\(e\), the intersection becomes a singleton.  Item~2 peels the
image of \(C_2\) to this singleton; the image of \(Y\) is peelable
by pc-minor closure.  This treats every elementary minor.

It remains to prove nonpeelability and terminality.  Every vertex of
\(C_2\) outside \(D\) retains its old incident directions and is
not a corner.  At either vertex of \(D\), two-connectivity of
\(C_2\) supplies an incident direction outside \(e\), while
\(p\) and \(q\) are incident from \(Y\); the required mixed
squares are absent, so neither vertex is a corner.
Consider \((v,10)\) or \((v,01)\) with \(v\in C\).
Its neighbor in the \(11\)-section forces every old direction
of a hypothetical corner cube to be supported in \(C\).  Thus
\(v\) would be a corner of \(C\), hence \(v=0\).  But at zero
the other new direction is also incident, and its full corner cube
would require the corner cube of \(C\) at zero to lie inside
the edge \(D\).  This is impossible since \(C\) is two-connected
and zero has degree at least two.  The same argument in the
\(11\)-section first forces \(v=0\) and then fails in the
\(00\)-section.  The two tips are corners, each with its acquired
support as its incident directions.  This argument remains valid
after either or both tips are deleted.  Consequently the complete
corner-deletion process is confined to these two tips and stops at
a nonempty cornerless family.  Therefore \(X\) is nonpeelable,
and the elementary-minor checks make it forbidden.

After deleting \((a_1,10)\), restrict to \(q=0\) and contract
\(p\).  After deleting \((b_1,01)\), interchange \(p,q\).
Both operations give \(C_1\cup C_2\), with intersection \(D\).
This family is cornerless: each copy has only zero as an old corner,
and at zero there are incident exclusive directions in both blocks
but no mixed square.  It has \(2|C|-2\) concepts, so it is a
proper nonpeelable minor of each corner deletion.  Neither deletion
is forbidden.  Hence \(X\) is corner-descent terminal.

For two-connectivity, the Cartesian product \(C\times L\) has no
cut vertex.  Indeed, after deleting one vertex, every remaining
vertex can reach a layer \(C\times\{t\}\) avoiding that vertex,
using another vertex of \(C\) if necessary to pass the deleted
layer position; that intact layer is connected.  Each tip has at
least two old neighbors: its acquired support cannot have size zero
or one in an irredundant nonempty ample base.  Each vertex of
\(D\times\{00\}\) has two neighbors in \(C\times L\).
Adding these vertices therefore preserves two-connectivity of
\(Y\).  Its union with the two-connected graph \(C_2\), along
the two vertices of \(D\), is two-connected as well.

Finally suppose the extra hypotheses hold, and put \(m=|F|\).
Each coordinate in \((F_1\setminus\{e\})\cup\{p,q\}\)
belongs to \(m-1\ge3\) missing cross-pair supports, one for
each coordinate in \(F_2\setminus\{e\}\).  Each coordinate
in the latter set belongs to \(m+1\) such supports.  These
are minimal nonfaces since all singleton supports are shattered.
The three assumed minimal nonfaces through \(e\) in the second
copy of \(C\) remain minimal nonfaces of \(X\), because its
projection to \(F_2\) is exactly \(C_2\).  Thus every coordinate
belongs to at least three minimal nonfaces.  No coordinate can be
balanced, and Theorem~\ref{thm:damp-terminal-clause-recovery}
rules out every successful layer set.
\end{proof}

\begin{remark}[A fixed two-connected terminal obstruction without layers]
\label{ex:damp-terminal-edge-amalgam-1158}
Take the translated \(289\)-concept family \(C\) used in
Remark~\ref{ex:damp-terminal-cut-vertex-1158}, choose the edge
\(D=\{0,2\}\), and use \(a=128,b=290\).  Their acquired supports
have bitmasks \(390,480\).  The input minor conditions require
exactly \(22+3=25\) relative peelings; the two repairs also peel
to \(D\).  These \(27\) orders are explicitly recorded in
\path{computations/round806_edge_amalgam_orders.json} and replayed
by \path{verify_damp_terminal_edge_amalgam.py}.
The coordinate of the edge belongs, for example, to the three
minimal nonfaces with bitmasks \(38,70,15\); their omitted traces
have bitmasks \(32,64,8\), respectively.

For an explicit embedding, shift the first copy's words by two bits,
put the two tips at \(514\) and \(1161\), and embed the second
copy by the coordinate map
\[
 1\longmapsto3,\qquad
 (0,2,3,4,5,6,7,8,9,10,11)
       \longmapsto(14,15,16,17,18,19,20,21,22,23,24).
\]
The resulting obstruction has \(1158\) concepts, dimension \(25\),
VC dimension \(4\), minimum degree \(4\), and exactly the two
displayed corners.  Each corner deletion has the cornerless
\(576\)-concept edge amalgam as a proper pc-minor.
The verifier checks all \(75\) elementary-minor peelings, the
complete four-state negative certificate, two-connectivity, both
terminality witnesses, and the exact shadow.  It recovers all
\(1004\) minimal signed clauses; every coordinate occurs in at
least \(24\) of them.  Hence there is no balanced coordinate
and no common-seed layer presentation.  The full report is
\path{damp_terminal_edge_amalgam_verification.json}.

This refutes layer-existence completeness even for two-connected,
cornered terminal graph three-cores.  The edge-amalgam construction
still assumes relative peelings of its input minors and repairs;
it does not solve their intrinsic admissibility.  Theorem
\ref{thm:damp-edge-separator-obstructions} now gives exact decomposition
and recovery throughout the separating-edge branch, using the five
boundary peeling properties of its pieces.  Their intrinsic
admissibility, and the full generative classification beyond the known
branches, remain open.
\end{remark}

\begin{proposition}[One-corner extensions and a two-neighbor falsifier]
\label{prop:damp-one-corner-two-neighbor}
Let \(G\subseteq Q_F\) be an irredundantly embedded cornerless ample
forbidden pc-minor.  Suppose that \(c,h\notin G\) are adjacent, both
one-concept extensions are ample, and both acquire the same support
\(I\).  Then \(S=G\cup\{c\}\) is an ample forbidden pc-minor
whose only corner is \(c\).  It has no successful layer set and no
common-seed cover with nonempty proper rows satisfying
Theorem~\ref{thm:damp-common-seed-cover-classification}.

Suppose, more generally, that an irredundantly embedded ample forbidden
pc-minor \(S\subseteq Q_F\) has only one corner \(c\), and that two
distinct words \(u,v\in Q_F\setminus S\)
are adjacent to \(c\) and satisfy
\(S\cup\{u\},S\cup\{v\}\in\Damp\).  For a new coordinate \(e\),
the family
\begin{equation}
 X=((S\cup\{u\})\times\{0\})
     \cup((S\cup\{v\})\times\{1\})
 \label{eq:damp-one-corner-two-neighbor-lift}
\end{equation}
is a corner-descent-terminal ample forbidden pc-minor with exactly two
corners.  Its only successful layer set is \(\{e\}\), and its recovered
seed is the cornered family \(S\).  Thus such \(S,u,v\) would give a counterexample with a one-corner
seed.  Proposition~\ref{prop:damp-terminal-cut-seed-towers} already
refutes general cornerless recovery using a two-corner seed.
\end{proposition}

\begin{proof}
For the first assertion, write \(h=c^{\{f\}}\).  The acquired-support
formula gives \(f\notin I\), and the punctured \(I\)-cubes at both
\(c\) and \(h\) lie in \(G\).  The added word \(c\) is a corner
of \(S\), with completed cube \(Q_c\) on support \(I\).
Any other corner of \(S\) would have to have \(Q_c\) as its unique
maximal cube: otherwise its maximal cube avoids \(c\) and makes it
a corner of \(G\).  But every \(x\in Q_c\setminus\{c\}\) has
the neighbor \(x^{\{f\}}\in G\), supplied by the punctured cube at
\(h\).  This direction lies outside \(I\), so \(Q_c\) cannot be
its unique maximal cube.  Hence \(c\) is the only corner of \(S\).
Deleting it leaves the nonpeelable \(G\), so \(S\notin\Damp\);
Proposition~\ref{prop:damp-one-concept-atlas} makes it forbidden.

Every proper-row common-seed obstruction in
Theorem~\ref{thm:damp-common-seed-cover-classification} has at least
two corners.  There are at least two rows, since a proper row cannot
contain both opposite facets of a layer coordinate.  Each row is
nonempty and peelable, and a corner of each row lifts to a distinct
occurrence corner by \eqref{eq:damp-replicated-repair-tip-cubes}.
Successful layer presentations of an obstruction are such covers by
Theorem~\ref{thm:damp-terminal-clause-recovery}.  This proves both
nonrepresentation assertions for \(S\).

For the second assertion, write \(u=c^{\{i\}}\) and
\(v=c^{\{j\}}\), with \(i\ne j\).  The word
\(c^{\{i,j\}}\) does not belong to \(S\): otherwise isometry of
\(S\) would require one of the missing words \(u,v\) on a length-two
path to \(c\).  The support acquired by \(u\) therefore contains
\(i\) and omits \(j\); the support acquired by \(v\) contains
\(j\) and omits \(i\).  These supports are distinct.
The common-seed theorem now makes \(X\) an ample forbidden pc-minor.
Its two occurrence words \((u,0),(v,1)\) are corners.  No seed word
is a corner: only copies of \(c\) could be, and at each layer an
adjacent repair occurs, which excludes them by
\eqref{eq:damp-seed-corner-occurrences}.  Deleting either occurrence
corner leaves \(S\) as a proper layer restriction, so neither deletion
is a forbidden pc-minor.  Thus \(X\) is terminal.

The coordinate \(e\) is thin and hence successful.  By
Theorem~\ref{thm:damp-layer-recovery-dichotomy}, the largest presentation
has two antipodally punctured rows on \(E_{\max}(X)\).  Each row has
at least \(|E_{\max}(X)|\) corners, namely the neighbors of its missing
word.  These lift to distinct corners of \(X\).  Since \(X\) has
exactly two corners, \(|E_{\max}(X)|=1\), and therefore
\(E_{\max}(X)=\{e\}\), proving uniqueness and the cornered-seed
assertion.  Existence of the two resolving neighbors is an additional
hypothesis here, not a consequence of the first assertion.
\end{proof}

\begin{corollary}[No resolving neighbor after a square completion]
\label{cor:damp-square-completion-neighbors}
Let \(G\subseteq Q_F\) be nonempty, ample, and nonpeelable, and let
\(S=G\cup\{c\}\) be ample with acquired support \(I\), where
\(|I|\le2\).  No word adjacent to \(c\) is a resolving one-concept
addition to \(S\).
Consequently, the one-corner construction in
Proposition~\ref{prop:damp-one-corner-two-neighbor} can supply its
two resolving neighbors only when \(|I|\ge3\).
\end{corollary}

\begin{proof}
Suppose that \(u=c^{\{f\}}\notin S\) and \(H=S\cup\{u\}\)
is ample.  Then \(f\notin I\), and the punctured-cube test gives
its acquired support as \(P=J\cup\{f\}\), where \(J\subseteq I\):
each \(j\in J\) supplies \(u^{\{f,j\}}=c^{\{j\}}\in G\).
The family \(S\) is nonpeelable by
Proposition~\ref{prop:damp-rank-two-extension-invariance} applied
to \(G,S\).  The same proposition excludes a resolving addition
when \(|J|\le1\).

Otherwise \(I=J\) has size two.  Removing \(c\) leaves the
punctured \(I\)-square at \(u\) intact and removes its only
neighbor in direction \(f\).  Since \(I\notin\operatorname{Sh}(G)\),
the one-concept test makes \(G\cup\{u\}\) ample with acquired
support \(I\).  Proposition~\ref{prop:damp-square-pair-invariance}
now makes \(H\) nonpeelable as well.
\end{proof}


The fixed verifier \path{verify_damp_rank_two_additions.py} checks two
nonpeelable controls over the certified \(291\)-concept seed.  Completing
a square attached along a path changes the order from \(293\) to \(294\)
and exposes the vertex opposite the addition; their complete deletion
certificates have three and seven states.  In the product of the seed
with a punctured square, one and two adjacent completions give orders
\(874\) and \(875\), with two- and four-state certificates.  These
controls contain the original seed as a proper pc-minor; they check the
local mechanisms used above.

\begin{proposition}[Two forced deletions at a rank-three neighbor]
\label{prop:damp-rank-three-neighbor-first-deletions}
Let \(S\subseteq Q_F\) be ample and nonpeelable, with exactly one
corner \(c\), and let
\[
 I=\{k\in F:c^{\{k\}}\in S\},\qquad |I|\ge3.
\]
Suppose that \(u=c^{\{i\}}\notin S\) and that \(H=S\cup\{u\}\)
is ample with acquired support \(P\) of size three.  Then
\(P=J\cup\{i\}\) for a two-element proper subset \(J\subsetneq I\).
Put
\[
 x=u^J,\qquad y=c^J.
\]
If \(H\) peels, every peeling starts by deleting \(x\) and then \(y\).
In particular, \(x\) has no neighbor outside the completed \(P\)-cube,
and
\begin{equation}
 \{y^{\{k\}}:k\in F\setminus I,\ y^{\{k\}}\in S\}=\{x\}.
 \label{eq:damp-rank-three-neighbor-private-edge}
\end{equation}

Consequently, \(u\) is not resolving if there is another adjacent
ample addition \(v=c^{\{j\}}\ne u\) whose acquired support is
\(K\cup\{j\}\), with \(J\subseteq K\subseteq I\).
\end{proposition}

\begin{proof}
The punctured-cube test for the ample addition \(u\) gives \(i\in P\).
For \(k\in P\setminus\{i\}\), it also gives
\(u^{\{i,k\}}=c^{\{k\}}\in S\), so \(k\in I\).
This proves the asserted form of \(P\).  Since \(u\notin S\),
one has \(i\notin I\).  The corner cube at \(c\) is the full
\(I\)-cube \(Q_I(c)\).

We first locate every possible first deletion.  The vertex \(c\)
is not a corner of \(H\): for \(k\in I\setminus J\), its neighbors
\(c^{\{k\}}\) and \(u\) have the missing opposite vertex \(u^{\{k\}}\).
Any other old corner \(z\) of \(H\) must have its unique maximal cube
contain \(u\); otherwise that cube would already make \(z\) a corner
of \(S\).  Since \(u\) is a corner with cube \(Q_P(u)\), this forces
\(z\in Q_P(u)\) and excludes every neighbor of \(z\) outside that cube.
No vertex \(c^U\), \(U\subseteq J\), can satisfy this: its neighbor
in any direction of \(I\setminus J\) belongs to \(Q_I(c)\), outside
\(Q_P(u)\).

Write \(J=\{p,q\}\).  On the other \(i\)-side of \(Q_P(u)\), the
remaining candidates are \(u^{\{p\}},u^{\{q\}},x\), in addition
to the new corner \(u\).  Each of \(u^{\{p\}},u^{\{q\}}\) has a
neighbor outside \(Q_P(u)\).  Indeed, without such a neighbor its
incident directions in \(S\) would span the full face of \(Q_P(u)\)
opposite \(u\) in the corresponding coordinate, making it a corner
of \(S\), different from \(c\).
Thus the only possible corners of \(H\) are \(u\) and \(x\).
Deleting \(u\) returns the nonpeelable \(S\), so a peeling must
start with \(x\).  In particular, \(x\) has no outside neighbor.

Removing a vertex can create a corner only at one of its neighbors:
elsewhere the incident directions do not change, and deleting a point
cannot fill a missing cube.  The three neighbors of \(x\) are
\(u^{\{p\}},u^{\{q\}},y\).
The first two retain their outside neighbors and their neighbor \(u\).
If an outside direction is \(k\notin P\), the square joining these
two directions would require \(u^{\{k\}}\notin H\).  Neither vertex
can therefore become a corner after deleting \(x\).
The vertex \(u\) also ceases to be a corner: its incident directions
remain \(P\), but their cube now misses \(x\).
Hence \(y\) is the only possible second deletion.

At \(y\), the full \(I\)-cube remains intact.  Its \(i\)-neighbor
was \(x\), which has been removed.  If another direction
\(k\notin I\) remained incident, the corner cube would have to
contain \(c^{\{k\}}\).  Here \(k\ne i\), so this vertex is absent
even after adding \(u\), by the definition of \(I\).
Thus \(y\) is a corner after deleting \(x\) exactly when
\eqref{eq:damp-rank-three-neighbor-private-edge} holds.
This proves the forced prefix and the stated necessary condition.

Finally, a second addition \(v\) as in the proposition supplies
\(v^J\in S\), since \(\varnothing\ne J\subseteq K\) and its
punctured repair cube lies in \(S\).  This vertex is \(y^{\{j\}}\).
It is an outside-\(I\) neighbor of \(y\) distinct from \(x=y^{\{i\}}\),
contradicting \eqref{eq:damp-rank-three-neighbor-private-edge}.
\end{proof}

\begin{corollary}[A companion tip excludes rank-three resolving neighbors]
\label{cor:damp-companion-neighbor-rank-four}
In the first construction of
Proposition~\ref{prop:damp-one-corner-two-neighbor}, every resolving
addition adjacent to the unique corner \(c\) has acquired support
of size at least four.
If the initial support \(I\) has size three, each such resolving
tip \(u=c^{\{i\}}\) completes the entire \(I\cup\{i\}\)-cube.
Moreover, \(G\cup\{u\}\) is itself an ample one-corner forbidden
pc-minor acquiring \(I\).

Thus, at initial support three, the proposed two-neighbor falsifier
requires three tips \(c,u,v\) with the same acquired support over \(G\):
their individual extensions are nonpeelable, while both
\(G\cup\{c,u\}\) and \(G\cup\{c,v\}\) must peel.
\end{corollary}

\begin{proof}
The earlier square-completion result
(Corollary~\ref{cor:damp-square-completion-neighbors}) excludes
\(|I|\le2\), so assume \(|I|\ge3\).
Write the companion tip as \(h=c^{\{f\}}\).
The two punctured \(I\)-cubes in \(G\), together with \(c\), form
the punctured \(I\cup\{f\}\)-cube at \(h\).
The one-concept test makes \(S\cup\{h\}\) ample, with acquired
support \(I\cup\{f\}\): this support is not shattered by \(S\),
whose shadow is \(\operatorname{Sh}(G)\cup\{I\}\).

A resolving neighbor cannot have support of size at most two by
Proposition~\ref{prop:damp-rank-two-extension-invariance}.
If it had size three, it would differ from \(h\), whose support
has size at least four.  The last assertion of
Proposition~\ref{prop:damp-rank-three-neighbor-first-deletions}
would then exclude it, using \(h\) with \(K=I\).
This proves the lower bound four.

For any adjacent ample addition, the punctured-cube argument above gives
\(P_u\subseteq I\cup\{i\}\), regardless of its support size.
If \(|I|=3\), the lower bound four therefore forces equality.
Removing \(c\) leaves the entire punctured
\(I\)-cube at \(u\) and removes precisely its neighbor in direction
\(i\).  Since \(I\notin\operatorname{Sh}(G)\), the one-concept test
makes \(G\cup\{u\}\) ample with acquired support \(I\).
The adjacent same-support tips \(u,c\) satisfy the first part of
Proposition~\ref{prop:damp-one-corner-two-neighbor}, so this
extension is one-corner and forbidden.  Applying the same argument
to a second resolving neighbor \(v\) proves the final assertion.
\end{proof}

The fixed checker \path{verify_damp_rank_three_corner_neighbors.py}
tests both stages of the forced-prefix argument.  It attaches copies of
the certified cornerless \(291\)-concept family at selected vertices of
two small ample families, leaving each with exactly one corner.
In the first resulting family,
the rank-three addition permits the forced deletions \(x,y\); in the
second, a companion tip blocks \(y\) after \(x\) is deleted.
Complete negative deletion certificates are checked, together with
fifteen rank-three neighbors in the previously recorded one-corner
examples.  The two constructed controls have proper cornerless minors
and are not forbidden witnesses.  The result excludes support three
in the companion construction; it does not exclude rank-three
resolving neighbors of arbitrary one-corner families or settle the
remaining support-four cooperation problem.

\begin{corollary}[Reductions under adjacent pair completions]
\label{cor:damp-three-tip-strong-projections}
Let \(G\subseteq Q_F\) be a nonempty ample family.
Suppose that \(c,u,v\notin G\) are distinct tips whose individual
extensions are ample with the same acquired support \(I\), where
\(|I|=3\).  Suppose also that \(u=c^{\{i\}}\),
\(v=c^{\{j\}}\), and
\[
 G\cup\{c,u\},\quad G\cup\{c,v\}\in\Damp.
\]
Then \(G^D\in\Damp\) for every nonempty coordinate set
\(D\subseteq F\).
\end{corollary}

\begin{proof}
Put \(S=G\cup\{c\}\).  The same-support hypothesis gives
\(i,j\notin I\) and \(i\ne j\).  The punctured cubes at \(c,u\)
make the addition of \(u\) to \(S\) acquire \(I\cup\{i\}\);
the addition of \(v\) similarly acquires \(I\cup\{j\}\).
Both resulting families are peelable by assumption.

Let \(D\ne\varnothing\) and \(G^D\notin\Damp\).  Apply
Proposition~\ref{prop:damp-repair-strong-projections} first to
\(G,S\).  Either \(S^D=G^D\), or \(S^D\) adds one concept
with support \(I\setminus D\) of size at most two.  In the second
case Proposition~\ref{prop:damp-rank-two-extension-invariance}
preserves nonpeelability.  Hence \(S^D\notin\Damp\) in both cases.
Applying \eqref{eq:damp-repair-nonpeelable-projection-support} to
both additions now gives
\[
 D\subseteq(I\cup\{i\})\cap(I\cup\{j\})=I,
 \qquad 4-|D|\ge3.
\]
Thus \(D=\{k\}\) with \(k\in I\).  Put \(A=G^{\{k\}}\).
By \eqref{eq:damp-repair-strong-projection}, the projected tips
\(c_k,u_k\) give individual ample additions to \(A\) with the
same acquired support \(I\setminus\{k\}\), of size two.  They
remain distinct and adjacent because \(i\ne k\).  Both have their
\(k\)-mates in \(G\), so
\[
 (G\cup\{c,u\})^{\{k\}}=A\cup\{c_k,u_k\}\in\Damp.
\]
Proposition~\ref{prop:damp-square-pair-invariance} forces \(A\)
to peel, contradicting the choice of \(D\).  Hence no nonempty
\(D\) has a nonpeelable reduction.
\end{proof}

The checker \path{verify_damp_repair_strong_projections.py} verifies
\eqref{eq:damp-repair-strong-projection} on every coordinate subset
in six fixed controls, including the recorded rank-three tip, its two
adjacent pair completions, and repaired terminal towers with one and
two layer coordinates.  It also extracts twelve overlap peelings of
the certified \(291\)-concept seed from three stored repair orders:
the supports at tips \(91,157,319\) have empty intersection, so for
each coordinate one repair leaves that reduction unchanged.
Thus every reduction of this seed on a nonempty coordinate set
is peelable.  This seed meets the necessary projection condition.
The next proposition nevertheless excludes its three-tip falsifier
profile by using VC dimension and the positioned punctured cubes.


\begin{proposition}[A third companion blocks a pair in VC dimension three]
\label{prop:damp-vc-three-companion}
Let \(G\subseteq Q_F\) be a nonempty cornerless ample family with
\(\operatorname{VCdim}(G)\le3\).  Suppose that the individual
extensions at three distinct missing tips \(c,u,v\) are ample and
acquire the same support \(I\), where \(|I|=3\).  If
\(u=c^{\{i\}}\) and \(v=c^{\{j\}}\), then
\[
 G\cup\{c,u\}\notin\Damp.
\]
The same conclusion holds for \(G\cup\{c,v\}\).
\end{proposition}

\begin{proof}
The two punctured cubes at \(c,u\) complete a four-dimensional
cube when both tips are added.  We show that the third punctured cube
protects its entire \(c\)-side.  After the first old vertex is
deleted, the VC-dimension bound confines all further deletions to four
vertices on the other side; they cannot reach the attached part of \(G\).

Put \(P=I\cup\{i\}\),
\(C=\{c^U:U\subseteq P\}\), and \(H=G\cup\{c,u\}\).
The acquired-support hypothesis implies \(i,j\notin I\),
\(i\ne j\), and
\[
 C\setminus\{c,u\}\subseteq G,\qquad
 \operatorname{Sh}(H)=\operatorname{Sh}(G)\cup\{I,P\}.
\]
Indeed, the punctured cube at each tip has support \(I\), and after
adding \(c\) the punctured cube at \(u\) has support \(P\).
The one-concept test (Lemma~\ref{lem:damp-one-concept-extension-test})
applies successively, since \(I\) is not shattered by \(G\) and
no support containing \(I\) is shattered by \(G\).
Thus \(H\) is ample, its only shattered four-set is \(P\), and
neither added tip has a neighbor outside \(C\).

Partition the vertices of \(C\) into
\[
 C_0=\{c^U:U\subseteq I\},\qquad \{u\},\qquad
 W=\{u^{\{p\}}:p\in I\},\qquad B\cup\{z\},
\]
where \(B=\{u^{I\setminus\{p\}}:p\in I\}\) and \(z=u^I\).
Every \(c^U\in C_0\setminus\{c\}\) has the neighbor
\(v^U\in G\setminus C\), supplied by the punctured cube at \(v\).
Every \(w=u^{\{p\}}\in W\) also has a neighbor outside \(C\).
Otherwise its incident directions in \(G\) would be exactly
\(P\setminus\{p\}\), whose entire cube through \(w\) lies in
\(C\setminus\{c,u\}\); this would make \(w\) a corner of \(G\).

Any corner of \(H\) belonging to \(G\) has its corner cube
containing \(c\) or \(u\), since otherwise it was already a
corner of \(G\).  It therefore lies in \(C\); all four directions
of \(P\) are incident there, and the added tips have no further
incident directions.  Its corner cube is exactly \(C\).
Consequently every such corner has no neighbor outside \(C\), and
the preceding paragraph confines it to \(B\cup\{z\}\).
Deleting \(c\) first leaves \(G\cup\{u\}\), whose only corner
is \(u\): each old vertex of the completed \(I\)-cube at \(u\)
has its \(i\)-neighbor in \(G\).  Its only next deletion leaves
the cornerless \(G\).  The same argument applies with \(c,u\)
interchanged.  Hence a successful peeling would have to start in
\(B\cup\{z\}\).

Such a first deletion has cube support \(P\), so it removes precisely
\(P\) from the shadow.  The remaining ample family has VC dimension
at most three, as does every family obtained by further corner deletions.
We prove that any further deletion must again lie in \(B\cup\{z\}\)
and have no neighbor outside \(C\) in \(H\).  This maintains all
vertices outside \(C\), as well as the twelve vertices of
\(C_0\cup\{u\}\cup W\).

For the induction, assume that all deletions so far have the stated
form.  No vertex outside \(C\) can have become a corner: none was
initially a corner, and deleting a nonneighbor cannot create a corner,
because its incident directions remain unchanged.
Each vertex of \(C_0\setminus\{c\}\) retains its three
\(I\)-neighbors and its outside \(j\)-neighbor.  The vertex
\(c\) retains its three \(I\)-neighbors and \(u\); the vertex
\(u\) retains \(W\) and \(c\).  Each has degree at least four,
whereas a corner in an ample family of VC dimension at most three
has degree at most three.
Every \(w=u^{\{p\}}\in W\) retains \(u\) and an outside
neighbor \(w^{\{k\}}\), with \(k\notin P\).
The square on these two edges would require \(u^{\{k\}}\), which
is absent from \(H\).  Thus \(w\) cannot be a corner either.

Each remaining vertex \(b\in B\) retains its two neighbors in
\(W\) and its mate in \(C_0\).  If it has an outside neighbor,
these four surviving neighbors prevent its deletion.  Moreover, while
\(z\) remains, every such \(b\) has four surviving neighbors
even without an outside neighbor.  Hence if the first deletion was
a vertex of \(B\), no second vertex of \(B\) can be deleted
before \(z\).  At any subsequent deletion of \(z\), it therefore
still has its mate in \(C_0\) and two neighbors in \(B\).
An outside neighbor would again give degree at least four.  Thus every
possible deletion of \(z\), including a first deletion, has no
outside neighbor.  After \(z\) is deleted, the same conclusion
for the remaining vertices of \(B\) follows from their three
protected neighbors.  This completes the induction.

The protected twelve vertices survive every deletion sequence that
starts at an old vertex.  Neither possible type of first deletion can
therefore begin a complete peeling.  Interchanging \(i,u\) with
\(j,v\) proves the second assertion.
\end{proof}

\begin{corollary}[At most one resolving neighbor in VC dimension three]
\label{cor:damp-vc-three-one-corner-neighbors}
\label{cor:damp-vc-three-companion-neighbors}
Let \(S\subseteq Q_F\) be an ample nonpeelable family with exactly
one corner \(c\), and put
\(I=\{k\in F:c^{\{k\}}\in S\}\).
If \(\operatorname{VCdim}(S)\le3\), at most one missing neighbor
\(u\) of \(c\) satisfies \(S\cup\{u\}\in\Damp\).
Consequently the one-corner two-neighbor falsifier in
Proposition~\ref{prop:damp-one-corner-two-neighbor} requires
\(\operatorname{VCdim}(S)\ge4\), whether or not its source comes
from the companion construction.

Without a VC-dimension assumption, if two resolving neighbors
\(u=c^{\{i\}}\), \(v=c^{\{j\}}\) have acquired supports
\(J\cup\{i\}\), \(K\cup\{j\}\) of size three, where
\(J,K\subseteq I\), then
\begin{equation}
 J\cup K\subsetneq I.
 \label{eq:damp-two-rank-three-unused-direction}
\end{equation}

In that companion construction, if
\(\operatorname{VCdim}(G)\le3\), the specified companion \(h\)
is the only possible resolving neighbor of \(c\).
If \(c\) has two distinct same-support companions over \(G\),
it has no resolving neighbor in \(S=G\cup\{c\}\).
\end{corollary}

\begin{proof}
The predecessor \(G=S\setminus\{c\}\) is ample and nonpeelable:
otherwise deleting \(c\) first would peel \(S\).
Adding \(c\) to \(G\) acquires \(I\).
If \(|I|\le2\), Corollary~\ref{cor:damp-square-completion-neighbors}
excludes every resolving neighbor, so assume \(|I|\ge3\).

We first prove \eqref{eq:damp-two-rank-three-unused-direction}.
The punctured-cube test gives \(J,K\subseteq I\), each of size
two, and distinct \(i,j\notin I\).
Proposition~\ref{prop:damp-rank-three-neighbor-first-deletions}
excludes \(J=K\) and makes every peeling of
\(H_u=S\cup\{u\}\) start at
\(x_u=u^J\) and \(y_u=c^J\).
Write \(J=\{p,q\}\).  After deleting \(x_u\), the only
corner is \(y_u\).  Its remaining neighbors, by
\eqref{eq:damp-rank-three-neighbor-private-edge}, are
\[
 c^{\{p\}},\quad c^{\{q\}},\quad
 t_r=c^{J\cup\{r\}}\quad(r\in I\setminus J).
\]
Deleting a nonneighbor cannot create a corner, so these are the only
third-deletion candidates.  For any \(r\in I\setminus J\),
the vertex \(c^{\{p\}}\) retains the neighbors
\(c,c^{\{p,r\}},u^{\{p\}}\).  A cube spanning those three
directions would contain \(u^{\{r\}}\), which is absent because
\(r\notin J\cup\{i\}\).  The argument for
\(c^{\{q\}}\) uses
\(c,c^{\{q,r\}},u^{\{q\}}\) and the same missing word.
Thus the third deletion must be some \(t_r\).

Every \(t_r\) has an outside-\(I\) neighbor in \(S\), since
otherwise the full \(I\)-cube would make it a second corner of
\(S\).  After deleting \(x_u,y_u\), it retains all these
outside neighbors and every direction in \(I\setminus\{r\}\):
\(x_u\) is not adjacent to \(t_r\), and \(y_u\) was its
neighbor along \(r\).
If it is the third corner deleted, choose any retained outside
direction \(k\notin I\).  Its corner cube contains the face
\begin{equation}
 \{t_r^U:U\subseteq (I\setminus\{r\})\cup\{k\}\}
 \subseteq S.
 \label{eq:damp-rank-three-forced-third-cube}
\end{equation}
The face avoids \(u\), since its coordinate \(r\) is flipped
from the value at \(c\), whereas \(u\) agrees with \(c\) there.

If \(r\in K\), then \(y_v=c^K\) lies in this face, as does
\(y_v^{\{k\}}\).  Applying
\eqref{eq:damp-rank-three-neighbor-private-edge} to \(v\)
forces \(k=j\).  Choose \(s\in I\setminus K\), which exists
because \(|I|\ge3>|K|\).  The same face then contains
\(x_v^{\{s\}}\), where \(x_v=v^K=y_v^{\{j\}}\).
Indeed, \(s\ne r\), so direction \(s\) belongs to its support.
This gives a neighbor of \(x_v\) outside its completed
\(K\cup\{j\}\)-cube, contrary to its required first deletion.
Hence every possible third deletion has
\(r\in I\setminus(J\cup K)\), proving
\eqref{eq:damp-two-rank-three-unused-direction}.

Now impose \(\operatorname{VCdim}(S)\le3\).
The full corner cube has support \(I\), so \(|I|=3\).
For any two putative resolving neighbors, write their supports as
\(J\cup\{i\}\), \(K\cup\{j\}\), with \(J,K\subseteq I\).
Proposition~\ref{prop:damp-rank-two-extension-invariance} excludes
supports of size at most two.  Thus each has size three or four.
If one has size three and the other size four, the last assertion of
Proposition~\ref{prop:damp-rank-three-neighbor-first-deletions}
excludes the former.  If both have size three, that assertion gives
\(J\ne K\), whence \(J\cup K=I\), contradicting
\eqref{eq:damp-two-rank-three-unused-direction}.

The remaining case has both acquired supports of size four, so
\(J=K=I\).  We show that the predecessor \(G\) is cornerless.
Any corner created by deleting \(c\) must be a neighbor
\(c^{\{p\}}\), with \(p\in I\), because \(c\) was the
only corner of \(S\).  Such a vertex retains its two neighbors
along \(I\setminus\{p\}\), as well as \(u^{\{p\}}\)
and \(v^{\{p\}}\), supplied by the two punctured repair cubes.
These are four distinct neighbors in the ample \(G\), whose VC
dimension is at most three, so it cannot be a corner.
Moreover, the individual extensions of \(G\) at \(c,u,v\)
all acquire \(I\).  For \(u\), removing \(c\) from its
punctured repair cube removes its neighbor along \(i\) and leaves
its entire punctured \(I\)-cube; the support \(I\) is not
shattered by \(G\), since it was the support removed with the
corner \(c\).  The one-concept test therefore applies, and the
same argument handles \(v\).
Proposition~\ref{prop:damp-vc-three-companion} now contradicts
\(G\cup\{c,u\}=H_u\in\Damp\).
This finishes the general assertion.

For the more specific companion conclusion, write
\(h=c^{\{f\}}\).  Its common acquired support \(I\) over the
nonempty ample \(G\) is nonempty.  For every \(p\in I\), the
two punctured \(I\)-cubes contain the full face on
\((I\setminus\{p\})\cup\{f\}\) with coordinate \(p\)
flipped from its value at \(c\).  Hence
\(|I|\le\operatorname{VCdim}(G)\le3\).
Supports of size at most two are already excluded.  At size three,
Corollary~\ref{cor:damp-companion-neighbor-rank-four} makes every
resolving neighbor an individual addition to \(G\) acquiring
\(I\).  If it differs from \(h\), the three tips violate
Proposition~\ref{prop:damp-vc-three-companion}.  Thus only \(h\)
can resolve.  Applying this conclusion to two companions leaves no
possible resolving neighbor.
\end{proof}

The fixed checker \path{verify_damp_vc_three_companion.py} attaches
copies of the certified cornerless \(291\)-concept family at the
corners of two small ample families.  The VC-three control has
\(1761\) concepts in \(77\) coordinates; its complete negative
deletion certificate has fifteen states and realizes the four possible
old first deletions.  The VC-four control permits the upper antipode
to become a corner while retaining a neighbor outside the completed
four-cube.  It isolates the step that uses the VC bound.  Both controls
have proper nonpeelable minors, and the latter still does not peel:
it is not a counterexample to a dimension-free nonpeelability claim.
The checker \path{verify_damp_vc_three_one_corner.py} tests the
additional argument for incomparable three-supports.  In one fixed
VC-three control, both additions permit their first two old deletions,
but the third is blocked.  In another, the third corner cube for one
addition supplies the outside neighbor that blocks the other's second
deletion.  A VC-four control permits both three-deletion prefixes along
a direction outside \(J\cup K\).  Each extension still has a proper
cornerless nonpeelable minor; these prefixes are not resolving examples.
The general corollary excludes all one-corner sources of VC dimension
at most three.  The two-neighbor existence question remains open in
higher VC dimension.


\begin{remark}[A peelable class with a unique corner]
\label{rem:damp-one-corner-peelable}
The fixed family in \path{damp_one_corner_peelable_control.json}
is an ample corner-peelable class with \(294\) concepts in twelve
coordinates and VC dimension three.  Its only corner is \(1894\),
with support \(\{0,8,11\}\), in the binary-word convention of the
certificate.  Thus having one corner does not imply nonpeelability.

The checker \path{verify_damp_one_corner_peelable_control.py}
reconstructs this family from the stored owner colouring of the
\(232\)-concept Hall contraction and its \(62\)-concept overlap.
The two sections have orders \(89\) and \(205\).  It verifies
ampleness of the sections, overlap, and lift, checks the unique corner,
and replays all \(294\) steps of the stored peeling, which begins
\(1894,1638,1634,1632\).

Every peeling of this nontrivial family must start at \(1894\), so none
can end there.  A corner-peelable ample family therefore cannot in general
be peeled to an arbitrarily prescribed remaining vertex.  VC dimension
three is the smallest dimension where this can fail: the proof of
\cite[Proposition~4.11]{ChalopiChepoiMoran22} constructs an ample prefix
order starting at any chosen concept in VC dimension at most two.
The checker also verifies that none of the nine missing neighbors of
\(1894\) is an ample addition.  This positive control does not
supply the two resolving neighbors in
Proposition~\ref{prop:damp-one-corner-two-neighbor}.
\end{remark}

\begin{remark}[A fixed one-corner source and a disconnected placement fibre]
\label{rem:damp-one-corner-fixed-neighbors}
For the certified \(291\)-concept cornerless seed \(G\), the support
with integer mask \(I=2312\) has precisely five ample one-concept
extensions, at
\[
 2184,\quad2696,\quad3720,\quad3818,\quad3822.
\]
Use the convention \(x=\sum_{f=0}^{11}x_f2^f\).  The adjacency graph
on these tips consists of the path \(2184\!-\!2696\!-\!3720\) and
the edge \(3818\!-\!3822\).  They have the same trace on \(I\), so
the outside-placement fibre \(\mathcal Y_I(G)\) from
\eqref{eq:damp-punctured-face-outside-fibers} is disconnected; it has
five words but seven shattered supports, so it is not ample.
Thus even for a cornerless forbidden seed one cannot assume
that the possible placements of a single acquired support form an ample
family.

Proposition~\ref{prop:damp-one-corner-two-neighbor} makes
\(S=G\cup\{2696\}\) a \(292\)-concept forbidden pc-minor with
exactly one corner, outside the proper-row common-seed cover class.
Its only adjacent ample additions are \(2568,2184,3720\), and all
three are nonpeelable.  Their complete corner-deletion certificates
have respectively three, five, and four states.
Corollary~\ref{cor:damp-vc-three-companion-neighbors} now gives a
uniform reason that this source has no resolving neighbor: the seed
has VC dimension three and \(2696\) has two companions.
The fixed verifier \path{verify_damp_one_corner_neighbors.py} checks
the complete placement fibre, the one-corner conclusion, and all adjacent
ample additions at this source.  The strengthened corollary also
excludes two resolving neighbors for arbitrary one-corner sources of
VC dimension three.  Higher VC dimension and nonadjacent cooperative
repairs remain outside that conclusion.
\end{remark}

\begin{proposition}[Separated cooperative additions]
\label{prop:damp-separated-cooperative-additions}
For every integer \(\ell\ge0\), there is a cornerless ample family
\(G_\ell\) of order \(597+\ell\) and two outside words
\(a_\ell,b_\ell\) such that all four families
\[
 G_\ell,\qquad G_\ell\cup\{a_\ell\},\qquad
 G_\ell\cup\{b_\ell\},\qquad G_\ell\cup\{a_\ell,b_\ell\}
\]
are ample.  The first three are nonpeelable, while the last is peelable.
The two additions acquire disjoint four-coordinate supports, and their
completed cubes have Hamming distance \(8+\ell\).
Each of the first three families has Hall's cornerless obstruction as
a proper pc-minor; in particular none is a forbidden pc-minor.
\end{proposition}

\begin{proof}
Use Hall's \(299\)-concept family \(C_H\), its resolving addition
\(183\), and the certified peeling of \(C_H\cup\{183\}\) ending
at \(4095\), from the data in
Remark~\ref{ex:damp-thick-terminal-three-repair}.
Translate every word by \(4095\), so that the last word becomes \(0\).
Write \(C\) for the translated cornerless family and
\(v=183\mathbin{\oplus}4095\) for the translated repair tip, where
\(\oplus\) denotes bitwise addition modulo two.  Its acquired support
is \(P=\{2,3,6,7\}\).  The word \(v\) has four nonzero coordinates
outside \(P\), so the completed \(P\)-cube has distance four from
\(0\).  Denote the translated peeling by
\(w_1,\ldots,w_{300}=0\).

Take two disjoint copies \(F_1,F_2\) of the twelve-coordinate set,
and a further set \(Z\) of \(\ell\) coordinates.  Put
\[
 T_\ell=\{1^k0^{\ell-k}:0\le k\le\ell\}\subseteq Q_Z.
\]
The desired family consists of two copies of \(C\) joined at their
zero words by this path:
\begin{align}
 G_\ell={}&(C\times\{0\}\times\{0_Z\})
     \cup(\{0\}\times C\times\{1_Z\})\notag\\
 &\quad\cup(\{(0,0)\}\times T_\ell),
 \label{eq:damp-separated-cooperative-family}\\
 a_\ell={}&(v,0,0_Z),\qquad b_\ell=(0,v,1_Z).
 \label{eq:damp-separated-cooperative-tips}
\end{align}
For \(\ell=0\), the two endpoint words coincide and the copies meet
at one vertex.  For positive \(\ell\), the path has \(\ell-1\)
internal vertices.  In both cases \(|G_\ell|=597+\ell\).

No pair of coordinates from different blocks \(F_1,F_2,Z\) is
shattered.  Between \(F_1\) and \(F_2\) the trace \(11\) is missing;
between \(F_1\) and \(Z\) it is again \(11\); between \(F_2\)
and \(Z\) the trace \(10\), with the \(F_2\)-coordinate first,
is missing.  The \(Z\)-projection is a path and shatters only the
empty set and its singletons.  Consequently the full shattered complex
is the union of the two block shadows and these \(\ell\) singletons,
with the empty set counted once.  Its size is \(597+\ell\), proving
ampleness by the Sandwich equality.  The same calculation after replacing
either or both copies of \(C\) by \(C\cup\{v\}\) proves ampleness
of the three extensions.

Every word in a copy of \(C\) retains its failure to be a corner:
no cube on that copy's coordinates is added by the other copy or the
path.  An internal path vertex has two incident directions which do
not span a square.  Thus \(G_\ell\) is cornerless.  Restricting to
\(F_1=0,Z=1_Z\) leaves the second copy of \(C\), even after adding
\(a_\ell\).  Restricting to \(F_2=0,Z=0_Z\) leaves the first copy,
even after adding \(b_\ell\).  These are proper pc-minors and are
nonpeelable, proving nonpeelability and failure of pc-minimality for
the first three families.

After both additions, peel the first completed copy using
\(w_1,\ldots,w_{299}\), leaving its attachment vertex.  Every deleted
word has exactly its old incident cubes, since it is not the attachment
vertex.  Do the same in the second completed copy.  The remaining path
\(T_\ell\) peels from one endpoint; for \(\ell=0\) only the common
vertex remains.  This is an explicit peeling of the final family.
The two added supports are the copies of \(P\) in \(F_1,F_2\).
Their cubes are each at distance four from the respective attachment
vertex, and all \(\ell\) bridge coordinates differ.  Their distance
is therefore \(4+\ell+4\), as asserted.
\end{proof}

The verifier \path{verify_damp_separated_cooperative_repairs.py} checks
the cases \(\ell=0,3\), including the original peeling certificate,
all four Sandwich equalities, cornerlessness, the proper Hall restrictions,
and the combined peelings of lengths \(599,602\).  It computes the
shattered sets block by block using the missing mixed pairs, without
an ambient-cube enumeration.
This construction rules out geometric separation of completed cubes as
a substitute for pc-minimality in a repair-inheritance argument.
Whether such cooperation can occur over a cornerless forbidden pc-minor
remains open.

\begin{theorem}[Terminality of iterated common-seed covers]
\label{thm:damp-iterated-common-seed-terminal}
Let \(G\subseteq Q_F\) be an irredundantly embedded ample forbidden
pc-minor, and let
\[
 S=(G\times Q_D)\cup
       \bigcup_{a\in\mathcal A}(\{a\}\times R_a),
 \qquad D\ne\varnothing,
\]
be a common-seed cover satisfying
Theorem~\ref{thm:damp-common-seed-cover-classification}, with \(S\)
itself forbidden.  Thus the inner tips resolve \(G\), their acquired
supports are distinct, and their occurrence rows are nonempty and proper.
The inner rows need not be facet shells or have disjoint facet ownership.
Let \(E\ne\varnothing\) be disjoint from \(F\cup D\), and let
\(X\subseteq Q_{F\cup D\cup E}\) be a second common-seed cover over
\(S\), satisfying the same theorem with nonempty proper outer rows.

If \(X\) is a corner-descent-terminal forbidden pc-minor, then both
stages have exactly two repair tips and antipodally punctured rows.
More precisely, for some distinct \(a,b\in\mathcal A\),
\(u\in Q_D\), and \(v\in Q_E\),
\begin{align}
 S={}&(G\times Q_D)
   \cup(\{a\}\times(Q_D\setminus\{u\}))
   \cup(\{b\}\times(Q_D\setminus\{\bar u\})),\notag\\
 X={}&(G\times Q_{D\cup E})
   \cup(\{a\}\times(Q_{D\cup E}\setminus\{(u,v)\}))\notag\\
 &\hspace{25mm}\cup
   (\{b\}\times(Q_{D\cup E}\setminus\{(\bar u,\bar v)\})).
 \label{eq:damp-iterated-terminal-flattening}
\end{align}
In particular, \(D\cup E\) is successful for \(X\).
Consequently, the seed recovered from the largest layer set of a terminal
obstruction admits no further common-seed cover satisfying the displayed
hypotheses, even with arbitrary proper occurrence rows.
More generally, a terminal output of two or more successive covers of
this kind is a single antipodal two-repair tower over the initial seed.
\end{theorem}

\begin{proof}
We first determine all resolving additions to the inner cover, without
assuming that its rows are facet shells.  Put
\[
 \mathcal A_*=
 \{a\in\mathcal A:R_a=Q_D\setminus\{u_a\}
                       \text{ for some }u_a\in Q_D\}.
\]
Then
\begin{equation}
 \{z\in Q_{F\cup D}\setminus S:S\cup\{z\}\in\Damp\}
   =\{(a,u_a):a\in\mathcal A_*\}.
 \label{eq:damp-proper-cover-resolving-additions}
\end{equation}
Indeed, \(S^D=G\notin\Damp\), so a resolving addition must change
this reduction and hence complete a full \(D\)-fibre.
An old word outside the projection cannot acquire a full fibre from one
addition, since \(D\ne\varnothing\).  The only possible additions
are therefore the displayed missing words.  Conversely, completing such
a row gives the product \((G\cup\{a\})\times Q_D\), together with
the other occurrence rows.  These other rows peel by the common-seed
criterion, and their peelings lift by
\eqref{eq:damp-replicated-repair-tip-cubes}.  The remaining product
peels because \(G\cup\{a\}\in\Damp\).  The added word's incident
directions are exactly \(P_a\cup D\), where \(P_a\) is the support
acquired over \(G\): every layer neighbor is present, and compatible
old tips are nonadjacent.  The local extension test gives this new
support, and these supports remain pairwise distinct.

Index the chosen outer tips by a subset \(\mathcal B\subseteq
\mathcal A_*\), and write \(T_a\subsetneq Q_E\) for the outer row
of \((a,u_a)\).  Substitution expresses \(X\) over the original seed
\(G\), with one row for every \(a\in\mathcal A\):
\begin{equation}
 X=(G\times Q_{D\cup E})\cup
       \bigcup_{a\in\mathcal A}(\{a\}\times U_a),\qquad
 U_a=
 \begin{cases}
 (R_a\times Q_E)\cup(\{u_a\}\times T_a),&a\in\mathcal B,\\
 R_a\times Q_E,&a\notin\mathcal B.
 \end{cases}
 \label{eq:damp-arbitrary-cover-flattening}
\end{equation}
Every \(U_a\) is nonempty and proper: an unselected row retains a
missing \(D\)-word, and a selected row misses
\(\{u_a\}\times(Q_E\setminus T_a)\).  Each is ample, being the
restriction of the ample \(X\) obtained by fixing its \(F\)-word
to \(a\).  Hence this is another common-seed cover meeting the
hypotheses of Theorem~\ref{thm:damp-common-seed-cover-classification}.

Terminality of \(X\) and
Theorem~\ref{thm:damp-terminal-common-seed-facet-partitions} force the
rows \(U_a\) to be a signed-facet partition on \(D\cup E\).
Their reductions along \(E\) are exactly \(R_a\), because
every outer row is proper.  Reduction of a facet shell retains
precisely its owned facets whose coordinates lie in \(D\).
Every \(R_a\) is nonempty, so each owner block retains at least one
such facet.  The retained blocks partition all signed \(D\)-facets.
Thus the original inner rows \(R_a\) themselves form a facet partition,
and \(D\) is successful for \(S\).

Terminality also makes the original outer rows a facet partition, so
\(E\) is successful for \(X\).  The nested-recovery conclusion of
Theorem~\ref{thm:damp-recovered-partition-resolving-tips} now applies.
It forces two owners at each stage and gives
\eqref{eq:damp-iterated-terminal-flattening} and success of \(D\cup E\).
If \(E\) was the largest successful set of \(X\), this strict
enlargement is impossible, proving the canonical-seed assertion.
For a longer chain, apply the two-stage result to the last two covers,
then regard their flattened output as the outer cover over the preceding
stage.  Repeating this operation proves the assertion by induction on
the number of stages.
\end{proof}

\begin{remark}[A redundant path row and its outer lift]
\label{rem:damp-iterated-path-row}
Take Hall's seed \(G=C_H\), its three certified repairs
\(a=183,b=315,d=1836\), and \(D=\{p,q,r\}\).  In Boolean-word
notation, set
\[
 R_a=Q_D\setminus\{000\},\qquad
 R_b=Q_D\setminus\{111\},\qquad
 R_d=\{000,100,110,111\}.
\]
These rows give a \(2410\)-concept ample forbidden pc-minor \(S\).
The path row owns no signed facet, while the other two rows partition
all signed facets.  The two resolving additions to \(S\) are
\((a,000)\) and \((b,111)\).  Every layer coordinate belongs to
four minimal forbidden supports: the two punctured-row clauses and two
path-row clauses.  Old coordinates belong to \(165\) or \(167\)
clauses.  Thus \(\mathcal B(S)=\varnothing\).
The only obstruction-preserving corner deletions remove the two
endpoints of the path row.  Each resulting three-word path has one
constant coordinate and one missing trace on the other two coordinates;
every layer coordinate then belongs to three clauses.  Their balanced
coordinate sets are again empty, so neither predecessor has layer
recovery.

Use the two resolving additions as opposite owners on a new coordinate
\(e\).  The resulting \(4822\)-concept outer lift has
\(\mathcal B(X)=\{e\}\) and unique layer recovery to \(S\).
Formula~\eqref{eq:damp-arbitrary-cover-flattening} identifies its
extra row over \(G\) as \(R_d\times Q_{\{e\}}\).  Peel these eight
occurrences.  Every intermediate family is a forbidden pc-minor by the
common-seed cover criterion, and the endpoint is the terminal
\(4814\)-concept antipodal tower over \(G\) on four layers.
The fixed verifier \path{verify_damp_iterated_common_seed_covers.py}
checks the complete inner repair list with two explicit peelings, all
eight inner corners, both preserving predecessors, the flattening
identity, and all eight outer corner deletions.  Here the iterated-cover
exclusion applies even though neither preserving corner predecessor
admits layer recovery.
\end{remark}

\begin{theorem}[Repair growth after extending a recovered obstruction]
\label{thm:damp-layered-predecessor-repair-growth}
Let \(H\subseteq Q_V\) be an irredundantly embedded ample forbidden
pc-minor with a successful nonempty layer set \(E\), and put
\(F=V\setminus E\).  Suppose that \(c\notin H\) and
\(S=H\cup\{c\}\) is ample and nonpeelable.  For an ample family
\(Y\subseteq Q_V\), write
\[
 \mathcal R(Y)=\{a\in Q_V\setminus Y:Y\cup\{a\}\in\Damp\}.
\]
Then
\begin{equation}
 \mathcal R(H)\subseteq\mathcal R(S),\qquad
 |\mathcal R(S)\setminus\mathcal R(H)|\le1,
 \qquad |\mathcal R(S)|\le3.
 \label{eq:damp-layered-predecessor-repair-growth}
\end{equation}
Every member of \(\mathcal R(H)\) is nonadjacent to \(c\) and retains
its acquired support over \(S\).  If a new resolving tip \(a\) exists,
then \(a|_F=c|_F\), and the missing part of that fibre in \(H\) is
the edge with endpoints \(a|_E,c|_E\).  In particular, \(a\) is adjacent
to \(c\).

No signed-facet assembly over \(S\), using compatible resolving tips,
is corner-descent terminal.  Consequently, if a terminal facet assembly
has a cornered seed \(S\), every obstruction-preserving corner deletion
\(S\setminus\{s\}\) has no successful layer presentation.
\end{theorem}

\begin{proof}
Proposition~\ref{prop:damp-one-concept-atlas} makes \(S\) a forbidden
pc-minor.  Recover the facet-partition presentation of \(H\), and write
\(K=H^E\) for its nonpeelable seed.  By
Theorem~\ref{thm:damp-recovered-partition-resolving-tips}, the
one-concept additions that complete an entire \(E\)-fibre are exactly
\(\mathcal R(H)\); each has an acquired support containing \(E\).
Since \(S\) is nonpeelable, adding \(c\) completes no such fibre,
and hence \(S^E=K\).  Its acquired support \(Q_c\) cannot contain
\(E\): the new corner cube on \(Q_c\) would otherwise complete the
\(E\)-fibre through \(c\).

For \(a\in\mathcal R(H)\), its acquired support \(Q_a\) therefore
differs from \(Q_c\).  The compatible-repair union formula
\eqref{eq:damp-replicated-repair-union-shadow} makes
\(H\cup\{a,c\}\) ample and preserves both acquired supports.
Deleting \(c\) leaves the peelable family \(H\cup\{a\}\), so the
union peels as well.  This proves
\(\mathcal R(H)\subseteq\mathcal R(S)\).
Lemma~\ref{lem:damp-compatible-repairs-do-not-collide} makes \(a,c\)
nonadjacent.

Now take \(a\in\mathcal R(S)\setminus\mathcal R(H)\).
The strong \(E\)-projection of \(S\) is still the nonpeelable family
\(K\), whereas that of \(S\cup\{a\}\) is peelable by
Proposition~\ref{prop:damp-strong-projection-closure}.  Thus adding
\(a\) completes an \(E\)-fibre.  If its old word differed from
\(c|_F\), the same fibre would already be completed in
\(H\cup\{a\}\), implying \(a\in\mathcal R(H)\), a contradiction.
The two additions therefore share their old word, and the original
fibre has exactly the two missing words \(c|_E,a|_E\).

For an old word in the projection of \(H\), this missing set is the
face opposite a recovered owner block; a two-word face is an edge.
For an old word outside that projection, the entire \(Q_E\) is missing,
so two missing words force \(|E|=1\) and again form an edge.  There is
at most one other missing word in the fibre through \(c\), proving the
second inequality and adjacency.  The previous repair theorem gives
\(|\mathcal R(H)|\le2\), proving the last inequality.

Finally, \(c\) is a corner of \(S\), and deleting it returns the
forbidden family \(H\).  Among any compatible chosen resolving tips
of \(S\), exactly the inherited ones belong to \(\mathcal A_c\) in
Theorem~\ref{thm:damp-seed-corner-facet-deletion}: all inherited tips
are nonadjacent to \(c\) and resolve \(H\), while any new tip is
adjacent to \(c\).  At most one owner is therefore outside
\(\mathcal A_c\).  Opposite facets have distinct owners, so they
cannot both be outside that set.  The exact deletion criterion supplies
an obstruction-preserving copy of \(c\) in every facet assembly over
\(S\), proving nonterminality.  A nonpeelable corner deletion of \(S\)
still uses every coordinate, as proved in
Theorem~\ref{thm:damp-seed-corner-facet-deletion}.  Thus applying this
conclusion with \(H=S\setminus\{s\}\) proves the final assertion.
\end{proof}

\begin{remark}[A sharp three-repair example without layer recovery]
\label{rem:damp-sharp-layered-repair-growth}
Let \(K=C_H\) be Hall's \(299\)-concept obstruction, and use its
three compatible repairs \(a=183,b=315,d=1836\) from
Remark~\ref{ex:damp-thick-terminal-three-repair}.  On one new coordinate
\(p\), put
\[
 H=((K\cup\{b\})\times\{0\})
       \cup((K\cup\{a\})\times\{1\}),\qquad
 c=(d,0),\qquad S=H\cup\{c\}.
\]
The family \(H\) is the terminal \(600\)-concept two-repair lift.
The family \(S\) has \(601\) concepts and is an ample forbidden
pc-minor: its three singleton occurrence rows are peelable and own both
signed \(p\)-facets.  Completing any one of these rows gives a
peelable extension, while any resolving addition must change the
nonpeelable strong \(p\)-projection \(K\).  Hence
\[
 \mathcal R(H)=\{(a,0),(b,1)\},\qquad
 \mathcal R(S)=\{(a,0),(b,1),(d,1)\}.
\]
The new repair \((d,1)\) is adjacent to \(c\), attaining the bound
in \eqref{eq:damp-layered-predecessor-repair-growth}.
Three resolving tips preclude any successful layer set for \(S\),
by Theorem~\ref{thm:damp-recovered-partition-resolving-tips}.
Equivalently, every old coordinate occurs in \(165\) minimal forbidden
traces and \(p\) occurs in three, so \(\mathcal B(S)=\varnothing\).
Nevertheless every facet assembly over \(S\) is nonterminal by the
new theorem.

For example, put \(u=(d,1)\), \(v=(b,1)\), and form the
\(1204\)-concept opposite-layer lift
\[
 X=((S\cup\{u\})\times\{0\})
       \cup((S\cup\{v\})\times\{1\})
\]
on a second coordinate \(q\).  Relative to deleting \(c\), the owner
\(u\) is ineligible and \(v\) is inherited.  Thus deleting \((c,1)\)
preserves obstruction status.  In the remaining family, the Hall repair
rows own the disjoint facet blocks
\[
 a:\{(p,1)\},\qquad
 b:\{(p,0),(q,1)\},\qquad
 d:\{(q,0)\}.
\]
The result is a terminal \(1203\)-concept Hall assembly.
Before deletion, \(p\) occurs in three minimal forbidden traces and
\(q\) in two, while every old coordinate occurs in \(165\).
Thus \(\{q\}\) is the unique successful layer set of \(X\), with
recovered seed \(S\).  The deletion makes \(p\) balanced as well,
allowing recovery of the cornerless Hall seed on layers \(\{p,q\}\).
The fixed verifier \path{verify_damp_layered_predecessor_repair_growth.py}
records all five resolving-extension peelings for \(H,S\), checks the
three-repair list and clause degrees, and verifies this deletion and its
terminal facet partition.
\end{remark}

\begin{proposition}[Removing a seed extension throughout the layer cube]
\label{prop:damp-inherited-repair-fibre-descent}
Let \(G\subseteq Q_F\) be an irredundantly embedded cornerless ample
forbidden pc-minor.  Suppose that \(c\notin G\), that
\(S=G\cup\{c\}\) is ample and nonpeelable, and that
\(\mathcal A\subseteq Q_F\setminus S\) is a set of tips with
\(G\cup\{a\}\in\Damp\) for every \(a\in\mathcal A\).
Assume that the one-concept acquired supports
\(P_c\) and \((P_a)_{a\in\mathcal A}\) over \(G\) are pairwise distinct.
Let \((R_a)_{a\in\mathcal A}\) be rows given by a partition of the
signed facets of a nonempty layer cube \(Q_E\), as in
Theorem~\ref{thm:damp-terminal-common-seed-facet-partitions}, and put
\[
 Y=(G\times Q_E)\cup
      \bigcup_{a\in\mathcal A}(\{a\}\times R_a),
 \qquad
 X=Y\cup(\{c\}\times Q_E).
\]
Then \(X\) is an ample forbidden pc-minor, and there is a sequence of
\(2^{|E|}\) obstruction-preserving corner deletions removing precisely
\(\{c\}\times Q_E\) and ending at \(Y\).
In particular, \(X\) is not corner-descent terminal, although its marked
layer set \(E\) recovers the cornered seed \(S\).
\end{proposition}

\begin{proof}
The facet-partition theorem makes \(Y\) a terminal ample forbidden
pc-minor.  Choose a corner peeling of the cube \(Q_E\), and denote its
successive suffix families by
\(T_0=Q_E,T_1,\ldots,T_m=\varnothing\), where \(m=2^{|E|}\).
Every \(T_i\) is ample.  The compatible-repair union formula
\eqref{eq:damp-replicated-repair-union-shadow} and
Proposition~\ref{prop:damp-replicated-repair-layers} make
\[
 Y_i=Y\cup(\{c\}\times T_i)
 \qquad(0\le i\le m)
\]
ample.  Their reductions along \(E\) are \(G\) for \(i>0\)
and \(S\) for \(i=0\).  All are therefore nonpeelable by
Proposition~\ref{prop:damp-strong-projection-closure}.

Starting with \(Y_m=Y\), traverse these families in reverse order.
Each step is an ample one-concept extension in the same ambient cube,
whose coordinates are already active in \(Y\).  The extension remains
nonpeelable, so Proposition~\ref{prop:damp-one-concept-atlas} makes it
another forbidden pc-minor.  Thus every \(Y_i\) is a forbidden pc-minor.
Successive families differ by one concept and are ample, so traversing
them from \(Y_0=X\) to \(Y_m=Y\) gives the asserted corner deletions.
Finally, all original rows \(R_a\) are proper, whence \(X^E=S\).
The compatible-repair union formula says that the acquired supports over
\(S\) are still the distinct \(P_a\), so this is the marked
presentation recognized by
Theorem~\ref{thm:damp-terminal-clause-recovery}.
\end{proof}

\begin{remark}[Maximal recovery need not give a cornerless seed]
\label{rem:damp-cornered-maximal-hall-seed}
Let \(H=C_H\) be Hall's \(299\)-concept obstruction, and put
\(S=H\cup\{373\}\).  Its added support is
\(P_c=\{1,2,7,9\}\).  The two corners of \(S\) are \(373\) and \(1011\);
deleting either leaves a nonempty cornerless family.  Hence \(S\) is
nonpeelable and, by Proposition~\ref{prop:damp-one-concept-atlas}, is
itself a forbidden pc-minor.  It is not terminal because deleting \(373\)
recovers \(H\).

Take the Hall repairs \(a=183\), \(b=315\) from
Remark~\ref{ex:damp-thick-terminal-three-repair}, whose acquired supports
are distinct from \(P_c\), and form
\[
 X=((S\cup\{b\})\times\{0\})
   \cup((S\cup\{a\})\times\{1\}).
\]
Proposition~\ref{prop:damp-inherited-repair-fibre-descent} makes \(X\)
a \(602\)-concept ample forbidden pc-minor.  Its signed minimal nonfaces
are all old four-sets except the three acquired supports, together with
the two mixed repair clauses.  Thus an old coordinate occurs in \(164\)
clauses if it belongs to the added support of \(373\), and in \(165\)
otherwise; the new layer coordinate occurs twice with opposite signs.
Consequently that coordinate is the only member of \(\mathcal B(X)\).
Its singleton layer set is successful and maximal, and its recovered
seed is the cornered obstruction \(S\).  Deleting the two layer copies
of \(373\), in the order supplied by the proposition, reaches the
terminal \(600\)-concept Hall lift.

The fixed diagnostic
\path{explore_damp_cornered_hall_repair_signatures.py} checks the four
nonpeelable Hall extensions at words \(373,1105,1360,1525\).
Each has nineteen ample one-concept extensions: twelve peelable and
seven nonpeelable.  The twelve resolving tips, with their acquired
supports, are exactly those of \(H\).  The record contains explicit
positive peelings and complete corner-deletion certificates for the
negative decisions.  Thus every facet-partition assembly using one of
these four cornered seeds satisfies the proposition and is nonterminal.
This rules out that source of counterexamples; it does not show that
every cornered recovered seed permits the same descent.  Proposition~\ref{prop:damp-terminal-cut-seed-towers} supplies terminal
counterexamples using a different, two-corner seed.
\end{remark}

\begin{remark}[Union closure requires the obstruction hypothesis]
\label{rem:damp-layer-union-collision}
Union closure in Theorem~\ref{thm:damp-obstruction-layer-union} does
not hold for arbitrary ample families.
Using the convention \(x=\sum_{i=0}^3x_i2^i\), consider
\[
 X=\{0,1,2,3,7,8\}\subseteq Q_{\{0,1,2,3\}}.
\]
It consists of a square and two pendant edges.  Its shattered sets are
the empty set, the four singletons, and \(\{0,1\}\), so it is ample.
Its minimal forbidden traces, written as support and missing word in the
same integer convention, are
\[
 (12,12),\quad(10,10),\quad(9,9),\quad(6,4),\quad(5,4).
\]
Thus \(\mathcal B(X)=\{0,1\}\).  Marking coordinate \(0\) gives
fibres \(T_9=\{8\},T_5=\{6\}\); marking coordinate \(1\) gives
\(T_{10}=\{8\},T_6=\{5\}\).  Both choices are successful.
For their union, however,
\[
 T_9=T_{10}=\{8\},\qquad T_5=T_6=\{4\},
\]
so recovery fails by a collision.  This ample example is peelable.
For obstruction inputs, in contrast, the union theorem permits all
successful one-coordinate closures to be united and recovers the unique
largest layer set.  Neither universal layer existence nor cornerlessness of the largest
recovered seed holds: Theorem~\ref{thm:damp-terminal-cut-vertex} and
Proposition~\ref{prop:damp-terminal-cut-seed-towers} give the respective
counterexamples.

The fixed verifier \path{verify_damp_obstruction_layer_unions.py} checks
the full recovery families of six Hall constructions.  Antipodal towers
on two and three layers have respectively three and seven successful
sets; the three-layer check includes overlapping incomparable choices.
It also checks the order-\(1203\) three-owner example and the
\(1204\)-to-\(1203\) descent in
Remark~\ref{rem:damp-sharp-layered-repair-growth}, including the
cornered seed of the larger input.  The cyclic order-\(2410\) example
in Remark~\ref{rem:damp-terminal-no-resolving-tip} supplies the case
with no resolving addition.  The verifier checks the sublayer fibre
formula, all thin coordinates, and the zero-, one-, and two-extension
cases of Theorem~\ref{thm:damp-layer-recovery-dichotomy}.
These finite regressions accompany
the proof and retain the six-word example as a control for its hypothesis.
\end{remark}

\begin{remark}[Unique layer recovery for the three-repair example]
\label{rem:damp-thick-example-unique-layers}
For the order-\(1203\) family in
Remark~\ref{ex:damp-thick-terminal-three-repair}, one has
\(\mathcal B(X)=\{p,q\}\), and the only successful layer set in
Theorem~\ref{thm:damp-terminal-clause-recovery} is \(\{p,q\}\).
Indeed, the Hall shadow consists of all subsets of \(F\) of order at
most three.  Adjoining the three repairs adds exactly their three
four-element supports.  No five-element set has all its four-subsets
among these three additions, so the old minimal nonfaces are exactly
\(\binom{F}{4}\setminus\{P_0,P_+,P_-\}\).
The mixed clauses replace the three removed clauses while retaining
their old supports.  Each old coordinate therefore occurs in exactly
\(\binom{11}{3}=165\) clauses, whereas \(p,q\) each occur twice with
opposite signs.  If only \(p\) is marked, the old part of the clause
\(P_+\cup\{p\}\) has the two solutions \((a_+,0)\) and \((a_+,1)\)
in \(C_{\{p\}}\), so its trace fibre is not a singleton.
Marking only \(q\) fails in the same way at \(a_-\).
Marking both coordinates recovers the Hall seed and the displayed three
repairs.  In Theorem~\ref{thm:damp-forced-layer-closure}, starting at
\(\{p\}\) therefore forces \(q\), and starting at \(\{q\}\) forces
\(p\); both procedures stop at \(\{p,q\}\).
The diagnostic \path{explore_damp_forced_layer_closure.py} records these
trajectories and the two failure certificates on fixed small examples.
The fixed witness verifier
\path{verify_damp_terminal_clause_recovery.py} checks these traces and
all fourteen section differences directly; its output also contains
explicit corner peelings of the three seed extensions.
\end{remark}

\begin{remark}[The reduced completeness target]
\label{rem:damp-corner-descent-terminal-target}
Recall that an ample forbidden pc-minor is \emph{corner-descent terminal}
if deleting none of its corners gives another forbidden pc-minor.  Starting from any
ample obstruction and repeatedly making an obstruction-preserving corner
deletion must terminate.  Reversing the deletions expresses the original
family as a finite sequence of the punctured-face extensions from
Proposition~\ref{prop:damp-one-concept-atlas}, starting at a
corner-descent-terminal obstruction.

The terminal objects are not all cornerless:
Corollary~\ref{cor:damp-two-repair-refutes-direct-descent} gives
two-extension lifts with exactly two corners and no preserving deletion.
They do not all admit common-seed presentations either:
Theorem~\ref{thm:damp-terminal-cut-vertex} constructs a terminal
cut-vertex obstruction, while every nontrivial common-seed assembly
with resolving tips is two-connected.  Proposition
\ref{prop:damp-terminal-edge-amalgam} further constructs two-connected
terminal obstructions with no successful layer set.  Thus splitting off
cut vertices does not make the common-seed construction exhaustive;
shared-edge assemblies must also be accounted for.

Inside the common-seed class,
Corollary~\ref{cor:damp-common-seed-obstruction-deletions} gives the
exact deletion criterion.  The thin/thick dichotomy in
Corollary~\ref{cor:damp-terminal-thin-thick-dichotomy} remains useful:
a thin coordinate canonically exposes a two-sided fragile kernel;
otherwise every coordinate has symmetric section difference at least
three.  The new terminal cut-vertex example lies outside all successful
layer recoveries, and hence has no thin coordinate.  Its two resolving extensions also yield the terminal towers of
Proposition~\ref{prop:damp-terminal-cut-seed-towers}, whose largest
recovered seed is cornered.  Thus the seed class cannot be restricted
to cornerless obstructions even after terminality and largest-layer
recovery.  Intrinsic admissibility of the recovered seeds, rooted
factors, and their repairing extensions remains necessary for the
full generative classification.
\end{remark}

\begin{theorem}[Antipodally punctured two-repair towers]
\label{thm:damp-antipodal-repair-towers}
Let \(S\in\Obs(\Damp)\cap\Ample\) be cornerless and irredundantly
embedded in \(Q_F\).  Let \(a,b\in Q_F\setminus S\) be distinct, suppose
that \(S\cup\{a\},S\cup\{b\}\in\Damp\), and suppose that their new
supports \(P_a,P_b\) are distinct.  Let \(E\ne\varnothing\) be disjoint
from \(F\), and let \(R_a,R_b\subsetneq Q_E\) be nonempty ample families.
Then
\begin{equation}
 X=(S\times Q_E)
   \cup(\{a\}\times R_a)
   \cup(\{b\}\times R_b)
 \quad\text{belongs to }\Obs(\Damp)
 \label{eq:damp-two-row-cover-obstruction}
\end{equation}
if and only if there is a word \(u\in Q_E\) such that, after possibly
interchanging \(a\) and \(b\),
\begin{equation}
 R_a=Q_E\setminus\{u\},
 \qquad
 R_b=Q_E\setminus\{\bar u\}.
 \label{eq:damp-antipodal-punctured-rows}
\end{equation}
In that case
\begin{equation}
 |X|=2^{|E|}(|S|+2)-2,
 \qquad
 \operatorname{idim}(X)=|F|+|E|.
 \label{eq:damp-antipodal-repair-tower-size}
\end{equation}
The normalized marked presentation recovers \(S,a,b\) and the antipodal
pair \(\{u,\bar u\}\).  Modulo layer-coordinate permutations and
reorientations, the only new parameter is \(|E|\).  For \(|E|=1\), the
construction is the opposite-layer two-extension lift.
\end{theorem}

\begin{proof}
Suppose first that \(X\) is a forbidden pc-minor.  Restricting all old
coordinates to the word \(a\), respectively \(b\), realizes \(R_a\) and
\(R_b\) as proper pc-minors of \(X\).  Hence both rows belong to
\(\Damp\).  Proposition~\ref{prop:damp-replicated-repair-corners} and
pc-minimality imply that the two rows own every signed facet of \(Q_E\).
For each \(e\in E\), a proper row cannot contain both opposite
\(e\)-facets.  Consequently there is a word \(v\in Q_E\) such that
\[
 \{t:t_e=v_e\}\subseteq R_a,
 \qquad
 \{t:t_e=1-v_e\}\subseteq R_b
 \quad(e\in E).
\]
The union of the facets on the left is
\(Q_E\setminus\{\bar v\}\).  Since \(R_a\) is proper, it equals this
punctured cube.  Similarly \(R_b=Q_E\setminus\{v\}\), which is
\eqref{eq:damp-antipodal-punctured-rows} with \(u=\bar v\).

Conversely, suppose \eqref{eq:damp-antipodal-punctured-rows} holds.
A punctured cube belongs to \(\Damp\), as follows inductively from the
peripheral-expansion test.  Proposition
\ref{prop:damp-replicated-repair-layers} therefore makes \(X\) ample,
and Proposition~\ref{prop:damp-replicated-repair-corners} makes it
non-corner-peelable: all its corners lie in the two proper rows, and after
any rowwise corner deletions one either gets stuck in a row or reaches the
cornerless product \(S\times Q_E\).

Every elementary layer restriction makes exactly one of the two rows
full and leaves the other row empty when \(|E|=1\), or punctured when
\(|E|>1\).  Every layer contraction makes both rows full.  Equation
\eqref{eq:damp-replicated-repair-layer-peeling-iff-full-row} therefore
puts all these layer images in \(\Damp\).

It remains to consider an elementary restriction or contraction \(\mu\)
in an old coordinate.  Put \(S'=\mu(S)\).  This is a proper pc-minor of
the obstruction \(S\), so \(S'\in\Damp\).  By Lemma
\ref{lem:damp-compatible-repairs-do-not-collide}, the two repair tips
cannot merge outside \(S'\), and every surviving outside tip is a corner
after the other has been removed.  Each such tip retains its punctured
occurrence row.  The product description
\eqref{eq:damp-replicated-repair-tip-cubes} lifts a peeling of its
punctured row.  Peeling all surviving rows leaves \(S'\times Q_E\), which
belongs to \(\Damp\).  Thus all old-coordinate elementary images are
corner-peelable as well.  Pc-minor closure now proves
\(X\in\Obs(\Damp)\).

The two rows each have \(2^{|E|}-1\) members, giving the order formula.
The product core realizes every old and layer direction, giving the
isometric dimension.  Recovery follows from
\eqref{eq:damp-replicated-repair-recovery}; the complements of the two
recovered rows are the two antipodal words.  Finally, the automorphism
group of \(Q_E\) is transitive on its words, so all choices of \(u\) are
equivalent in the marked layer cube.
\end{proof}

\begin{remark}[The first Hall towers]
\label{ex:damp-hall-antipodal-repair-towers}
For Hall's \(299\)-concept obstruction, the repair tips
\[
 a=183,\qquad b=315
\]
are corner-peelable one-concept extensions with respective support labels
in the verifier's zero-based coordinate numbering
\[
 P_a=\{2,3,6,7\},
 \qquad
 P_b=\{2,3,8,9\}.
\]
They therefore satisfy Theorem~\ref{thm:damp-antipodal-repair-towers}.
Layer dimensions one, two, and three give explicit forbidden pc-minors of
orders
\[
 600,\qquad 1202,\qquad 2406,
\]
and isometric dimensions \(13,14,15\), respectively.  More generally,
the order in layer dimension \(r\ge1\) is \(301\cdot2^r-2\).  These are
the first members of one tower; every pair of Hall repair tips with
distinct support labels gives another.
\end{remark}

\begin{lemma}[Exact ampleness test for a pair of one-concept extensions]
\label{lem:damp-two-repair-ampleness}
Let \(S\subseteq Q_F\) be ample, let \(a,b\in Q_F\setminus S\) be
distinct, and suppose that
\[
 A=S\cup\{a\},\qquad B=S\cup\{b\}
\]
are ample.  Each extension has a unique new shattered coordinate set;
write
\begin{equation}
\operatorname{Sh}(A)\setminus\operatorname{Sh}(S)=\{P_a\},
 \qquad
 \operatorname{Sh}(B)\setminus\operatorname{Sh}(S)=\{P_b\}.
 \label{eq:damp-two-repair-new-supports}
\end{equation}
These supports are local: explicitly,
\begin{equation}
P_a=\{f\in F:a^{\{f\}}\in S\},
 \qquad
 P_b=\{f\in F:b^{\{f\}}\in S\}.
 \label{eq:damp-two-repair-local-supports}
\end{equation}
Put \(C=S\cup\{a,b\}\), and, for a new coordinate \(e\), put
\[
 X=(B\times\{0\})\cup(A\times\{1\}).
\]
Then
\begin{equation}
X\text{ is ample}\quad\Longleftrightarrow\quad P_a\ne P_b.
 \label{eq:damp-two-repair-ample-iff-distinct-supports}
\end{equation}
When these equivalent conditions hold, \(C\) is ample and
\begin{equation}
 \operatorname{Sh}(C)
 =\operatorname{Sh}(S)\mathbin{\dot\cup}\{P_a,P_b\},
 \label{eq:damp-two-repair-union-shadow}
\end{equation}
while
\begin{equation}
\operatorname{Sh}(A)\cap\operatorname{Sh}(B)
 =\operatorname{Sh}(S).
 \label{eq:damp-two-repair-independent-shadows}
\end{equation}
\end{lemma}

\begin{proof}
Lemma~\ref{lem:damp-one-concept-extension-test}, applied to \(a\) and
\(b\), proves the uniqueness and local formulas in
\eqref{eq:damp-two-repair-new-supports}--
\eqref{eq:damp-two-repair-local-supports}.  In particular, \(P_a\) and
\(P_b\) are strongly shattered in \(A\) and \(B\), respectively.

Suppose first that \(P_a\ne P_b\).  The family \(C\) strongly shatters
every member of \(\operatorname{Sh}(S)\), as well as \(P_a\) and
\(P_b\).  Hence
\[
 |\operatorname{SSh}(C)|\ge |S|+2=|C|.
\]
The lower Sandwich inequality gives the reverse inequality.  Equality
in that inequality characterizes ample families, so \(C\) is ample.
Since \(|\operatorname{Sh}(C)|=|C|=|S|+2\), the shattered supports
already exhibited are all of them, proving
\eqref{eq:damp-two-repair-union-shadow}.  The distinctness of the two
new supports also gives
\eqref{eq:damp-two-repair-independent-shadows}.

A set of old coordinates is shattered by \(X\) exactly when it is
shattered by \(C\), while \(D\cup\{e\}\) is shattered exactly when
\(D\) is shattered by both \(A\) and \(B\).  Therefore
\begin{equation}
\operatorname{Sh}(X)
 =\operatorname{Sh}(C)\mathbin{\dot\cup}
  \{D\cup\{e\}:D\in
       \operatorname{Sh}(A)\cap\operatorname{Sh}(B)\}.
 \label{eq:damp-two-repair-shattering}
\end{equation}
Using the ampleness of \(C,S\) and
\eqref{eq:damp-two-repair-independent-shadows}, we obtain
\[
 |\operatorname{Sh}(X)|
 =|C|+|S|=2|S|+2=|X|,
\]
so \(X\) is ample.

If \(P_a=P_b\), then the intersection in
\eqref{eq:damp-two-repair-shattering} contains
\(\operatorname{Sh}(S)\cup\{P_a\}\).  Sauer's inequality gives
\(|\operatorname{Sh}(C)|\ge|C|\), and consequently
\[
 |\operatorname{Sh}(X)|
 \ge |C|+|S|+1=|X|+1.
\]
Thus \(X\) is not ample, proving
\eqref{eq:damp-two-repair-ample-iff-distinct-supports}.
\end{proof}

\begin{proposition}[Signed trace transport and recovery for a two-repair lift]
\label{prop:damp-two-repair-trace-transport}
Use the notation and hypotheses of
Lemma~\ref{lem:damp-two-repair-ampleness}, assume that
\(P_a\ne P_b\), and put
\[
 \Sigma=\operatorname{Sh}(S),
 \qquad
 \Sigma^+=\Sigma\mathbin{\dot\cup}\{P_a,P_b\}.
\]
Then the shattered-set complex of the lift is
\begin{equation}
 \operatorname{Sh}(X)
 =\Sigma^+\mathbin{\dot\cup}
   \{D\cup\{e\}:D\in\Sigma\}.
 \label{eq:damp-two-repair-transport-shadow}
\end{equation}
Its entire trace table is obtained from the two sections as follows:
\begin{align}
 X[D,q]
  &=(B[D,q]\times\{0\})\mathbin{\dot\cup}
    (A[D,q]\times\{1\})
 &&(D\in\Sigma^+,\ q\in Q_D),
 \label{eq:damp-two-repair-old-trace-transport}\\
 X[D\cup\{e\},(q,0)]&=B[D,q]\times\{0\},
 \notag\\
 X[D\cup\{e\},(q,1)]&=A[D,q]\times\{1\}
 &&(D\in\Sigma,\ q\in Q_D).
 \label{eq:damp-two-repair-vertical-trace-transport}
\end{align}

Let \(\mathcal N(\Gamma)\) denote the minimal nonfaces of a simplicial
complex \(\Gamma\), and let \(\sigma_Y\) denote the unique missing trace
on each minimal nonface of an ample family \(Y\).  Then
\begin{equation}
 \mathcal N(\operatorname{Sh}(X))
 =\mathcal N(\Sigma^+)\mathbin{\dot\cup}
   \{P_a\cup\{e\},P_b\cup\{e\}\},
 \label{eq:damp-two-repair-transport-minimal-nonfaces}
\end{equation}
and its signing is
\begin{align}
 \sigma_X(I)&=\sigma_C(I)
 &&(I\in\mathcal N(\Sigma^+)),
 \label{eq:damp-two-repair-transport-old-signs}\\
 \sigma_X(P_a\cup\{e\})&=(a|_{P_a},0),
 &
 \sigma_X(P_b\cup\{e\})&=(b|_{P_b},1).
 \label{eq:damp-two-repair-transport-vertical-signs}
\end{align}
In particular, the signed presentation in
\eqref{eq:damp-two-repair-transport-minimal-nonfaces}--
\eqref{eq:damp-two-repair-transport-vertical-signs} reconstructs \(X\).

Conversely, the marked pair \((X,e)\) recovers all the assembly data,
up to reversing \(e\).  Indeed, after projecting the two \(e\)-sections
to \(Q_F\),
\begin{equation}
 B=\rho_e^0X,
 \qquad
 A=\rho_e^1X,
 \qquad
 S=A\cap B,
 \qquad
 \{a\}=A\setminus B,
 \qquad
 \{b\}=B\setminus A,
 \label{eq:damp-two-repair-parameter-recovery}
\end{equation}
and
\begin{equation}
 P_a=\{f\in F:a^{\{f\}}\in S\},
 \qquad
 P_b=\{f\in F:b^{\{f\}}\in S\}.
 \label{eq:damp-two-repair-support-recovery}
\end{equation}
\end{proposition}

\begin{proof}
Equation~\eqref{eq:damp-two-repair-union-shadow} gives
\(\operatorname{Sh}(C)=\Sigma^+\), and
\eqref{eq:damp-two-repair-independent-shadows} gives
\(\operatorname{Sh}(A)\cap\operatorname{Sh}(B)=\Sigma\).
Substitution in \eqref{eq:damp-two-repair-shattering} proves
\eqref{eq:damp-two-repair-transport-shadow}.  The two trace formulas
follow directly from whether the traced support contains the layer
coordinate \(e\).

A set not containing \(e\) is a minimal nonface of
\(\operatorname{Sh}(X)\) precisely when it is a minimal nonface of
\(\Sigma^+\).  Now consider \(D\cup\{e\}\).  It is a nonface precisely
when \(D\notin\Sigma\).  It is minimal precisely when
\(D\in\mathcal N(\Sigma)\) and deletion of \(e\) leaves a face of
\(\Sigma^+\).  Since \(P_a\) and \(P_b\) are the two distinct minimal
nonfaces added to \(\Sigma\), this happens exactly for
\(D\in\{P_a,P_b\}\).  This proves
\eqref{eq:damp-two-repair-transport-minimal-nonfaces}.

On an old minimal nonface, the union of the projected traces of the two
sections is exactly \(C\), which proves
\eqref{eq:damp-two-repair-transport-old-signs}.  The trace
\(a|_{P_a}\) is the unique trace on \(P_a\) missing from \(S\).  It is
also missing from \(B\): otherwise adding \(b\) would shatter \(P_a\),
contrary to
\(\operatorname{Sh}(A)\cap\operatorname{Sh}(B)=\Sigma\).
The upper section contains that trace through \(a\), so the unique trace
of \(P_a\cup\{e\}\) missing from \(X\) is \((a|_{P_a},0)\).  The
argument for \(P_b\) in the lower section is symmetric.  This proves
\eqref{eq:damp-two-repair-transport-vertical-signs}.  The fixed-shadow
signed-clutter theorem now reconstructs \(X\) from these clauses.

Finally, the two sections of the displayed lift are \(B\) and \(A\).
Their common concepts are exactly \(S\), and their two differences are
the singleton repair tips.  This proves
\eqref{eq:damp-two-repair-parameter-recovery}; the local formula from
Lemma~\ref{lem:damp-two-repair-ampleness} gives
\eqref{eq:damp-two-repair-support-recovery}.  Reversing \(e\) merely
interchanges \((a,A,P_a)\) with \((b,B,P_b)\).
\end{proof}

\begin{proposition}[A thin coordinate recovers a two-sided fragile kernel]
\label{prop:damp-thin-coordinate-fragile-kernel}
Let \(X\subseteq Q_{F\mathbin{\dot\cup}\{e\}}\) be ample and not
corner-peelable.  Project its two \(e\)-sections to \(Q_F\), and write
\[
 H_i=\rho_e^iX\quad(i\in\{0,1\}),
 \qquad K=H_0\cap H_1.
\]
Suppose that the two section differences are singletons, say
\begin{equation}
 H_0\setminus K=\{b\},
 \qquad
 H_1\setminus K=\{a\}.
 \label{eq:damp-thin-coordinate}
\end{equation}
Then \(K\), \(H_0=K\cup\{b\}\), \(H_1=K\cup\{a\}\), and
\(C=K\cup\{a,b\}\) are ample, whereas \(K\notin\Damp\).  Moreover, the
two one-concept extensions acquire distinct shattered supports:
\begin{equation}
 \operatorname{Sh}(K\cup\{a\})\setminus\operatorname{Sh}(K)=\{P_a\},
 \qquad
 \operatorname{Sh}(K\cup\{b\})\setminus\operatorname{Sh}(K)=\{P_b\},
 \qquad P_a\ne P_b.
 \label{eq:damp-thin-coordinate-distinct-supports}
\end{equation}
The family is exactly the two-extension lift
\begin{equation}
 X=((K\cup\{b\})\times\{0\})
   \cup((K\cup\{a\})\times\{1\}).
 \label{eq:damp-thin-coordinate-recovery}
\end{equation}
Thus the marked family \((X,e)\) recovers \(K,a,b,P_a,P_b\), up to
reversing \(e\), by intersections, differences, and the local formulas
in \eqref{eq:damp-two-repair-support-recovery}.

If, in addition, \(X\in\Obs(\Damp)\), then
\begin{equation}
 K\cup\{a\},\quad K\cup\{b\},\quad K\cup\{a,b\}
 \quad\text{all belong to }\Damp.
 \label{eq:damp-thin-coordinate-fragile-diamond}
\end{equation}
In particular, the thin coordinate exposes a finite, intrinsic
two-sided fragile kernel.  In this obstruction case,
\(K\in\Obs(\Damp)\cap\Ample\) as well.
\end{proposition}

\begin{proof}
A set of old coordinates is shattered by \(X\) exactly when it is
shattered by \(C=H_0\cup H_1\).  A set \(D\cup\{e\}\) is shattered
exactly when \(D\) is shattered by both sections.  Hence
\begin{equation}
 \operatorname{Sh}(X)
 =\operatorname{Sh}(C)\mathbin{\dot\cup}
   \{D\cup\{e\}:D\in
        \operatorname{Sh}(H_0)\cap\operatorname{Sh}(H_1)\}.
 \label{eq:damp-thin-coordinate-shattering}
\end{equation}
Since every support shattered by \(K\) is shattered by both sections,
Sauer's inequality and the ampleness of \(X\) give
\begin{align*}
 2|K|+2
 &=|X|=|\operatorname{Sh}(X)|\\
 &\ge |C|+|\operatorname{Sh}(K)|
 \ge (|K|+2)+|K|=2|K|+2.
\end{align*}
Equality holds throughout.  It follows that \(C\) and \(K\) are ample
and that
\begin{equation}
 \operatorname{Sh}(H_0)\cap\operatorname{Sh}(H_1)
 =\operatorname{Sh}(K).
 \label{eq:damp-thin-coordinate-shadow-intersection}
\end{equation}
The sections \(H_0,H_1\) are signed restrictions of the ample family
\(X\), so they are ample as well.  Each one-concept extension of \(K\)
therefore acquires exactly one shattered support.  These supports are
distinct by \eqref{eq:damp-thin-coordinate-shadow-intersection}, proving
\eqref{eq:damp-thin-coordinate-distinct-supports}.  The section
decomposition gives \eqref{eq:damp-thin-coordinate-recovery} and the
stated recovery formulas.

It remains to prove that \(K\notin\Damp\).  Deleting \((a,1)\) from \(X\)
leaves the peripheral expansion
\[
 ((K\cup\{b\})\times\{0\})\cup(K\times\{1\}),
\]
which is ample because both its base and its duplicated site are ample.
Thus \((a,1)\) is a corner.  After deleting it, \((b,0)\) is a corner,
and deleting both exclusive tips leaves \(K\square Q_1\).  If
\(K\in\Damp\), the product equivalence would make this last family
corner-peelable, and the two preceding corner deletions would give a
corner peeling of \(X\), a contradiction.  Hence \(K\notin\Damp\).

Finally, if \(X\in\Obs(\Damp)\), then \(H_0,H_1\), and \(C=X/e\) are
proper pc-minors of \(X\).  Pc-minimality gives
\eqref{eq:damp-thin-coordinate-fragile-diamond}.
Since \(K=X^{\{e\}}\notin\Damp\),
Theorem~\ref{thm:damp-inherited-overlap-descent} gives
\(K\in\Obs(\Damp)\cap\Ample\).
\end{proof}

\begin{corollary}[Thin--thick terminal dichotomy]
\label{cor:damp-terminal-thin-thick-dichotomy}
Let \(O\in\Obs(\Damp)\cap\Ample\) be corner-descent terminal and suppose
that \(O\) has a corner.  For every coordinate \(e\), put
\[
 H_i(e)=\rho_e^iO,
 \qquad
 D_i(e)=H_i(e)\setminus H_{1-i}(e)
 \quad(i\in\{0,1\}).
\]
Exactly one of the following alternatives holds.
\begin{enumerate}
 \item \emph{Thin fragile-kernel branch.}
       Some coordinate \(e\) satisfies
       \(|D_0(e)|=|D_1(e)|=1\).  Proposition
       \ref{prop:damp-thin-coordinate-fragile-kernel} canonically recovers
       an ample forbidden pc-minor \(K\) and the peelable diamond
       \eqref{eq:damp-thin-coordinate-fragile-diamond}.  The two exclusive
       concepts are corners of \(O\).  Deleting either one leaves a family
       outside \(\Damp\), but terminality says that neither deletion is a
       forbidden pc-minor.
 \item \emph{Thick-section branch.}
       For every coordinate \(e\),
       \begin{equation}
        |D_0(e)|\ge1,
        \qquad |D_1(e)|\ge1,
        \qquad |D_0(e)|+|D_1(e)|\ge3.
        \label{eq:damp-terminal-thick-section}
       \end{equation}
\end{enumerate}
\end{corollary}

\begin{proof}
Proposition~\ref{prop:damp-periphery-free} implies that neither section
of \(O\) is contained in the other.  Thus \(D_0(e)\) and \(D_1(e)\) are
nonempty for every \(e\).  If both have order one for some coordinate,
Proposition~\ref{prop:damp-thin-coordinate-fragile-kernel} applies.  Its
proof shows that the two exclusive concepts are corners and that deleting
either one gives a peripheral expansion with duplicated site
\(K\notin\Damp\); the exact peripheral-expansion test therefore puts both
deletions outside \(\Damp\).  Corner-descent terminality says that they
are not forbidden pc-minors.  If there is no such coordinate, two positive
integers \(|D_0(e)|,|D_1(e)|\) cannot have sum two, which gives
\eqref{eq:damp-terminal-thick-section}.
\end{proof}

\begin{corollary}[Protector-incidence transport alphabet]
\label{cor:damp-two-repair-protector-transport}
For a trace-carrier scheme on \(X\), every precedence arc is read from
\eqref{eq:damp-two-repair-old-trace-transport}--
\eqref{eq:damp-two-repair-vertical-trace-transport}.  More explicitly:
\begin{enumerate}
 \item an old-support carrier has same-trace predecessors in the union of
       the two corresponding section fibers, and each of its protector
       choices lies in such a two-section union;
 \item if the carrier of \(D\cup\{e\}\) lies in layer \(i\), all of its
       same-trace predecessors lie in layer \(i\), while for every
       \(q\in Q_D\) it has a protector demand in the opposite-layer fiber
       \[
        X[D\cup\{e\},(q,1-i)].
       \]
\end{enumerate}
Consequently every directed trace circuit of the lift is a finite word in
these explicitly listed old-support, same-layer, and cross-layer incidence
types.  Its number of cross-layer arcs is even.
\end{corollary}

\begin{proof}
The carrier and protector arcs are defined solely by membership in the
relevant trace fibers.  The two assertions therefore follow from
\eqref{eq:damp-two-repair-old-trace-transport} and
\eqref{eq:damp-two-repair-vertical-trace-transport}.  Traversing a closed
directed walk changes layers an even number of times.
\end{proof}

\begin{theorem}[Two-extension lift]
\label{thm:damp-two-repair-lift}
Let \(S\in\Obs(\Damp)\cap\Ample\) be irredundantly embedded in
\(Q_F\). Let \(a,b\in Q_F\setminus S\) be distinct, and put
\[
 A=S\cup\{a\},\qquad B=S\cup\{b\},\qquad C=S\cup\{a,b\}.
\]
Assume that \(A,B\in\Damp\) and that their new shattered supports
\(P_a,P_b\) from \eqref{eq:damp-two-repair-new-supports} are distinct.
Then, for a new coordinate \(e\), the family
\begin{equation}
 X=(B\times\{0\})\cup(A\times\{1\})
 \label{eq:damp-two-repair-lift}
\end{equation}
is an ample forbidden pc-minor for \(\Damp\).

For a corner \(z\) of an ample family \(Z\), denote its unique maximal
cube by \(Q_Z(z)\). If in addition
\begin{equation}
\begin{aligned}
 Q_B(z)&\nsubseteq A&&\text{for every }z\in S
   \text{ that is a corner of }B,\\
 Q_A(z)&\nsubseteq B&&\text{for every }z\in S
   \text{ that is a corner of }A,
\end{aligned}
\label{eq:damp-two-repair-no-reflected-corner}
\end{equation}
then the only corners of \(X\) are \((b,0)\) and \((a,1)\).
\end{theorem}
\begin{proof}
Lemma~\ref{lem:damp-two-repair-ampleness} makes \(C\) and \(X\) ample.
An edge between the exclusive tips would give a diagonal-only square
face of \(X\), so no such edge exists. Each section peels to \(S\)
by deleting its added corner. Theorem~\ref{thm:damp-one-relative-wing-inverse}
therefore proves forbiddenness, without the additional corner condition.

The two exclusive tips are corners of \(X\), since they have no
vertical mates and are corners of their sections. For \(z\in S\),
the lower lift has all the old incident directions of \(z\) in \(B\)
and the vertical direction. It is a corner exactly when \(z\) is a
corner of \(B\) and its corner cube \(Q_B(z)\) lies in the opposite
section \(A\). The upper lift satisfies the symmetric criterion.
Condition~\eqref{eq:damp-two-repair-no-reflected-corner} excludes all
these additional corners and proves the final assertion.
\end{proof}

\begin{corollary}[Complete multipartite extension graph]
\label{cor:damp-cornerless-seed-repair-graph}
Let \(S\in\Obs(\Damp)\cap\Ample\) be cornerless and irredundantly
embedded in \(Q_F\).  Define its set of peelable extension vertices by
\[
 \mathcal R(S)
 =\{a\in Q_F\setminus S:S\cup\{a\}\in\Damp\}.
\]
For \(a\in\mathcal R(S)\), let \(P(a)\) be the unique member of
\[
 \operatorname{Sh}(S\cup\{a\})\setminus\operatorname{Sh}(S).
\]
By Lemma~\ref{lem:damp-two-repair-ampleness}, this label is the local
neighbour-direction signature
\begin{equation}
P(a)=\{f\in F:a^{\{f\}}\in S\}.
 \label{eq:damp-cornerless-repair-local-signature}
\end{equation}
and a minimal nonface of \(\operatorname{Sh}(S)\).
For distinct \(a,b\in\mathcal R(S)\), form the opposite-layer lift
\begin{equation}
X_{a,b}
 =((S\cup\{b\})\times\{0\})
  \cup((S\cup\{a\})\times\{1\}).
 \label{eq:damp-cornerless-seed-repair-lift}
\end{equation}
Then
\begin{equation}
X_{a,b}\text{ is ample}
 \quad\Longleftrightarrow\quad P(a)\ne P(b).
 \label{eq:damp-cornerless-repair-graph-adjacency}
\end{equation}
Whenever these conditions hold,
\[
 X_{a,b}\in\Obs(\Damp)\cap\Ample
\]
and \(X_{a,b}\) has exactly the two corners \((b,0)\) and \((a,1)\).
Consequently, the graph on \(\mathcal R(S)\) in which \(a,b\) are
adjacent precisely when their lift is ample is the complete multipartite
graph whose parts are the nonempty fibers of \(P\).  Every edge of this graph gives a forbidden pc-minor without any
pc-minor test.  In particular, the compatibility graph is computed from the
Hamming neighbours of the extension vertices, without computing shattered sets
or proper pc-minors.
\end{corollary}

\begin{proof}
Lemma~\ref{lem:damp-two-repair-ampleness} proves
\eqref{eq:damp-cornerless-repair-graph-adjacency}.  Suppose that the
supports are distinct.  We verify that the reflection hypothesis in
Theorem~\ref{thm:damp-two-repair-lift} is automatic.

Let \(z\in S\) be a corner of \(S\cup\{b\}\).  If its unique maximal
cube did not contain \(b\), that cube would lie in \(S\).  Every cube of
\(S\) through \(z\) is also a cube of \(S\cup\{b\}\) through \(z\), and
hence is contained in this unique maximal cube.  Thus \(z\) would be a
corner of \(S\), contrary to the hypothesis.  The corner cube therefore
contains \(b\), which does not belong to \(S\cup\{a\}\).  It is not
contained in the opposite one-concept extension.  The symmetric argument applies to
old corners of \(S\cup\{a\}\).

Theorem~\ref{thm:damp-two-repair-lift} now gives the asserted
obstruction and corner statements.  The adjacency condition
depends only on inequality of the support labels, which is exactly the
complete multipartite description.
\end{proof}

\begin{corollary}[Independent repair assemblies stop at two layers]
\label{cor:damp-independent-repair-terminality}
Let \(S\in\Obs(\Damp)\cap\Ample\) be cornerless and irredundantly
embedded in \(Q_F\).  Use the construction of
Proposition~\ref{prop:damp-independent-repair-layers}, with
\(E\ne\varnothing\), \(T\subseteq Q_E\), and
\(a_t\in\mathcal R(S)\) for every \(t\in T\).  Assume that the tips and
their supports are pairwise distinct, so that \(X_T\) is ample.  Then

\begin{equation}
 X_T\in\Obs(\Damp)
 \quad\Longleftrightarrow\quad
 |E|=1\ \text{ and }\ T=Q_E.
 \label{eq:damp-independent-repair-terminality}
\end{equation}

Thus the two-extension lift is the only forbidden pc-minor in the entire
class of independent one-tip-per-layer assemblies over a common seed.
In particular, placing three tips in three layers of a square cannot
produce a forbidden pc-minor, and filling all four layers cannot produce
a new one.
\end{corollary}

\begin{proof}
If \(|E|=1\) and \(T=Q_E\), the family is the two-extension lift associated
with its two tips.  Their support labels are distinct, so
Corollary~\ref{cor:damp-cornerless-seed-repair-graph} places \(X_T\) in
\(\Obs(\Damp)\).

Suppose next that \(T\ne Q_E\), and choose \(v\in Q_E\setminus T\).
Restricting all layer coordinates to \(v\) gives exactly \(S\).  This is
a proper pc-minor of \(X_T\) outside \(\Damp\), so \(X_T\) is not
pc-minimal outside \(\Damp\).

It remains to consider \(T=Q_E\) and \(|E|\ge2\).  Choose \(e\in E\)
and fix all coordinates in \(E\setminus\{e\}\).  The resulting proper
pc-minor has two sections \(S\cup\{a_t\}\) and
\(S\cup\{a_u\}\), where \(t,u\) are the endpoints of the selected
layer edge.  Their support labels are distinct, so this restriction is
itself a forbidden two-extension lift by
Corollary~\ref{cor:damp-cornerless-seed-repair-graph}.  Again \(X_T\)
is not pc-minimal.  These three cases prove
\eqref{eq:damp-independent-repair-terminality}.
\end{proof}

\begin{corollary}[Two-extension lifts refute direct corner descent]
\label{cor:damp-two-repair-refutes-direct-descent}
Every family \(X\) produced by
Theorem~\ref{thm:damp-two-repair-lift} is a counterexample to the
corner-deletion assertion: neither of its corners \(c\) satisfies
\begin{equation}
X\setminus\{c\}\in\Obs(\Damp).
 \label{eq:damp-two-repair-no-direct-descent}
\end{equation}
Thus direct corner descent fails even for forbidden endpoints with
exactly two corners.
\end{corollary}

\begin{proof}
Theorem~\ref{thm:damp-two-repair-lift} says that the only corners of
\(X\) are \((b,0)\) and \((a,1)\).  Deleting the upper corner gives
\[
 X\setminus\{(a,1)\}
  =(B\times\{0\})\cup(S\times\{1\}).
\]
This family is not corner-peelable by the exact peripheral-expansion
test, because \(S\notin\Damp\).  Its restriction to the upper
semicube is \(S\), a proper pc-minor outside \(\Damp\); hence the
deletion is not pc-minimal outside \(\Damp\), and therefore does not
belong to \(\Obs(\Damp)\).  The same argument after reversing the two
layers applies to the deletion of \((b,0)\).  This proves
\eqref{eq:damp-two-repair-no-direct-descent}.
\end{proof}

\begin{theorem}[Reduction to an irreducible cornerless family]
\label{thm:damp-prime-cornerless-seed-reduction}
Every finite ample non-corner-peelable cube family \(X\) can be reduced
to an ample, cornerless, periphery-free, Cartesian-prime,
non-corner-peelable family \(H\) by a finite sequence of the following
operations:
\begin{enumerate}
 \item delete a corner;
 \item if the current family is a peripheral expansion
       \[
        (B\times\{0\})\cup(S\times\{1\}),\qquad S\subseteq B,
       \]
       replace it by a nonempty non-corner-peelable member of
       \(\{B,S\}\);
 \item if the current family is \(A\square B\) with both factors
       nontrivial, replace it by a non-corner-peelable factor.
\end{enumerate}
Every intermediate family selected by this procedure is ample and
non-corner-peelable.  In ideal language, the terminal family is a
periphery-free Cartesian-prime member of
\(\mathfrak C_r^{\mathrm{col}}\)
for some \(r\le d\).  Thus cornerless prime cores are the starting objects that cannot be
reduced further by these three operations.
\end{theorem}

\begin{proof}
Deleting a corner preserves ampleness.  It also preserves failure of a
complete corner peeling: otherwise the deleted corner followed by a
peeling of the remainder would peel the original family.  Therefore we
may delete corners until a cornerless ample nonpeelable family is
reached.

If that family has a peripheral coordinate, write it in the form in
item~2.  Its nonempty projected sections are ample.  The same
expansion test in Proposition
\ref{prop:damp-exact-expansion-product-tests} shows, when \(S\ne
\varnothing\), that at least one of \(B,S\) is non-corner-peelable.  If
\(S=\varnothing\), deleting the redundant coordinate identifies the
family with \(B\), which is therefore non-corner-peelable.  Select a
nonempty non-corner-peelable section.  If the
family has a nontrivial Cartesian factorization, the factors are ample,
and the same proposition shows that at least one factor is
non-corner-peelable; select one.  Either selection strictly decreases the
isometric dimension.  Starting from the selected family, delete corners
again and repeat.

The process terminates because the pair consisting of isometric dimension
and order, ordered lexicographically, strictly decreases: a corner
deletion decreases the order without increasing dimension, whereas a
section or factor selection decreases the dimension.  At termination the
family is cornerless and has neither a
peripheral coordinate nor a nontrivial Cartesian factor.  It is therefore
periphery-free and Cartesian-prime.  Cornerlessness says that deletion of
every generator destroys linearity, so its ideal belongs to \(\mathfrak C_r^{\mathrm{col}}\).
\end{proof}


\section{Examples and the minimum-order problem}
\label{sec:damp-examples}

The preceding recognition criteria characterize the full forbidden basis,
while the construction theorems describe specified subclasses.  We now
record representative members and study the independent extremal question
of the least ample member.

\subsection{The Hall obstruction and its first exchanges}

The theorem of Chalopin, Chepoi, Moran, and Warmuth produces a maximum
VC-dimension-three family with no corner.  Cornerlessness proves that the
family is not in \(\Damp\), but does not by itself prove pc-minor
minimality.  The following proposition records the additional finite check
needed here and the first neighboring obstructions it reveals.  Integers in
the statement encode their twelve-bit binary expansions in the coordinate
convention of \cite[Theorem~4.5]{ChalopiChepoiMoran22}.

\begin{proposition}[The Hall obstruction and four first exchanges]
\label{prop:damp-hall-first-exchanges}
Let \(C_H\subseteq Q_{12}\) be the \(299\)-concept maximum
VC-dimension-three class of
\cite[Theorem~4.5]{ChalopiChepoiMoran22}.  Then
\[
 C_H\in\Obs(\Damp)\cap\Ample.
\]
It has exactly sixteen ample one-concept extensions.  Twelve are
corner-peelable and have four corners.  The other four are nonpeelable,
have exactly two corners, and are described by
\begin{equation}
\begin{array}{c|c}
\text{adjoined concept }a&\text{other corner }c\\ \hline
373&1011\\
1105&192\\
1360&3120\\
1525&252
\end{array}
\label{eq:damp-hall-first-exchanges}
\end{equation}
For every row, both
\[
 C_H\cup\{a\}
 \quad\text{and}\quad
 (C_H\setminus\{c\})\cup\{a\}
\]
belong to \(\Obs(\Damp)\cap\Ample\); the second family is cornerless,
has order \(299\), and has the same shattered-set complex as \(C_H\).
The four second families are pairwise nonisomorphic and are not isomorphic
to \(C_H\).
\end{proposition}

\begin{proof}[Computer-assisted proof]
The cited theorem proves that \(C_H\) is maximum of VC dimension three and
has no corner, hence is ample and not corner-peelable.  Exact evaluation of
its twenty-four signed restrictions and twelve contractions gives a complete
corner peeling of every elementary pc-minor.  The criterion in
Theorem~\ref{thm:damp-descendant-free-parametrization} therefore proves
\(C_H\in\Obs(\Damp)\).

For each of the \(2^{12}-299=3797\) outside concepts, the finite check first
tests the one-concept isometry condition and then the Sandwich equality,
the complete corner set, and corner peelability.  This gives the split
\(16=12+4\) and the four pairs in
\eqref{eq:damp-hall-first-exchanges}.  Proposition
\ref{prop:damp-one-concept-atlas} makes each nonpeelable union
\(C_H\cup\{a\}\) a forbidden pc-minor.  Its only possible first peeling
steps are its two corners.  Deleting \(a\) returns the nonpeelable family
\(C_H\), so deleting the other corner \(c\) also leaves a nonpeelable ample
family.  The same-shadow criterion in Proposition
\ref{prop:damp-same-shadow-corner-exchange} proves that this deletion
preserves the shattered-set complex.  Direct corner evaluation makes the
resulting family cornerless.  Finally, the thirty-six elementary minors of
each exchanged family are corner-peelable, and exact canonicalization under
coordinate permutations and independent coordinate reversals distinguishes
the five cornerless types.
\end{proof}

\begin{proposition}[The carrier paths in Hall's class]
\label{prop:damp-hall-carrier-paths}
For each of the sixty-six coordinate pairs $S\subseteq[12]$, the carrier
tree $\mathcal K_S(C_H)$ from
Proposition~\ref{prop:damp-maximum-vc-three-carrier-trees} is a path on
eleven vertices.  It has ten edges, of which eight are internal.  Among
the $220$ maximal three-cubes of $C_H$, exactly $36$ are internal edges
in one of their three carrier paths, $60$ are internal edges in two, and
$124$ are internal edges in all three.  In particular, every maximal
cube of $C_H$ is covered by two opposite square facets belonging to other
maximal cubes.
\end{proposition}

\begin{proof}[Computer-assisted proof]
The exact verifier reconstructs the unique three-cube on every triple
support from the fixed list of the $299$ concepts.  For each pair $S$, it
forms the carrier graph whose vertices are the $S$-cube anchors and whose
$e$-edge is the cube with support $S\cup\{e\}$.  All sixty-six graphs have
eleven vertices, ten edges, and degree sequence
\[
 1,1,2,2,2,2,2,2,2,2,2,
\]
so each is a path.  Counting, for every triple, in how many of its three
carrier paths its labeled edge has no leaf endpoint gives the asserted
histogram.  Proposition~\ref{prop:damp-maximum-vc-three-carrier-trees}
then identifies such an internal edge with an opposite-facet cover.

The durable verifier and output are
\[
 \begin{gathered}
 \texttt{verify\_damp\_hall\_carrier\_trees.py},\\
 \texttt{damp\_hall\_carrier\_trees.json},
 \end{gathered}
\]
and the mathematical-profile digest is
\[
 \texttt{65c031beb26e738618488674026fd5308367c99c513773edf2793e09c2ba7dcf}.
\]
\end{proof}

\begin{remark}[What the Hall neighborhood shows]
\label{rem:damp-hall-neighborhood}
The obstruction \(C_H\) is not isolated.  Repeating the same-shadow exchange
through distance two gives twenty-seven nonisomorphic cornerless
obstructions of order \(299\) and thirty-two nonisomorphic two-corner
obstructions of order \(300\).  Continuing one further one-concept layer
yields \(454\) nonisomorphic forbidden pc-minors in dimension twelve, of
orders \(299\), \(300\), and \(301\).  These exact finite atlases are useful
evidence about the geometry of the basis, but the paper uses them neither
as a classification nor as a substitute for the symbolic parametrizations
above.
\end{remark}

\begin{proposition}[A maximum fibre inside every ample family]
\label{prop:damp-antipodal-double-maximum-fibre}
Let $H\subsetneq Q_E$ be a nonempty ample family, and let $z\notin E$.
Write $\overline H=\{\overline h:h\in H\}$ and form the antipodal double
\[
 \widehat H=(\{0\}\times H)\cup(\{1\}\times\overline H)
 \subseteq Q_{\{z\}\cup E}.
\]
There are a set $R\subseteq E$ and a word $f\in Q_R$ such that, after
deleting the fixed coordinates $R$,
\begin{equation}
 H_f=H_{\overline f}=:A,
 \label{eq:damp-opposite-equal-maximum-fibres}
\end{equation}
the common family $A\subseteq Q_{E\setminus R}$ is maximum, and
\[
 (\{0\}\times A)\cup(\{1\}\times\overline A)
\]
is minimal nonample.  In particular, $A$ is the tope set on one side of a
distinguished element of a simple uniform oriented matroid, and hence is a
pseudogeometric maximum class.

The set $R$ can be taken empty precisely when $\widehat H$ is itself
minimal nonample.  Put $r=|R|$.  If $R$ is nonempty, then
\begin{equation}
 (r+1)|A|\leq |H|.
 \label{eq:damp-antipodal-double-geodesic-order}
\end{equation}
Equality holds only if, for every $a\in A$, the section
\[
 P_a=\{g\in Q_R:(g,a)\in H\}
\]
is a geodesic path from $f$ to $\overline f$, and every section over a
word outside $A$ is empty.  In fact all these paths are the same: equality
holds if and only if
\begin{equation}
 H=A\times P
 \label{eq:damp-antipodal-double-equality-product}
\end{equation}
for one antipodal geodesic path $P\subseteq Q_R$.
\end{proposition}

\begin{proof}
Put $\Sigma=\operatorname{Sh}(H)$.  Since $H$ is proper, $\Sigma$ has a
minimal nonface $I$.  Every proper subset of $I$ is shattered by $H$.
The projection $H|_I$ is ample, so it has exactly $2^{|I|}-1$ traces and
therefore omits a unique word $q\in Q_I$.  Its antipodal image omits
$\overline q$, and $q\neq\overline q$.  Consequently
$(H\cup\overline H)|_I=Q_I$, so $\widehat H$ shatters $I$.

For every $S\in\Sigma$, both $S$ and $S\cup\{z\}$ are shattered by
$\widehat H$: the two $z$-layers each realize every trace on $S$.  These
$2|\Sigma|$ sets do not include $I$.  Since $H$ is ample,
\[
 |\operatorname{Sh}(\widehat H)|
 \geq 2|\Sigma|+1
 =2|H|+1
 =|\widehat H|+1.
\]
Thus $\widehat H$ is nonample.

Choose a conditioning of $\widehat H$ that is nonample and has as few
unfixed coordinates as possible.  It is minimal nonample.  The coordinate
$z$ remains unfixed: fixing it produces a conditioning of $H$ or of
$\overline H$, which is ample because ampleness is closed under
conditionings.  Hence the chosen conditioning fixes a word $f$ on some
$R\subseteq E$ and has the two $z$-layers
\[
 H_f
 \qquad\text{and}\qquad
 \overline{H_{\overline f}}.
\]
The classification of minimal nonample families
\cite{BousquetYangComplex26} says that such a family is antipodally closed
and that each nonempty one-coordinate conditioning is maximum.  Antipodal
closure identifies the displayed layers, giving
$H_f=H_{\overline f}=A$, and the $z=0$ conditioning shows that $A$ is
maximum.  The same classification identifies the minimal nonample double
as a uniform-oriented-matroid tope set, so its side $A$ is
pseudogeometric.  The assertion about $R=\varnothing$ follows from the
choice of the conditioning.

For each $a\in A$, the section $P_a$ is a nonempty conditioning of the
ample family $H$, and hence is itself ample and isometric.  It contains the
antipodal words $f$ and $\overline f$.  Their Hamming distance is $r$, so
$P_a$ contains a shortest path of length $r$ and therefore at least $r+1$
vertices.  The sections over different words of $Q_{E\setminus R}$ are
disjoint, and summing over $a\in A$ proves
\eqref{eq:damp-antipodal-double-geodesic-order}.  Equality leaves no
nonempty section outside $A$ and forces every $P_a$ to consist of exactly
the vertices of one shortest antipodal path.  Such a path flips every
coordinate of $R$ exactly once.

It remains to use ampleness.  For every $S\in\operatorname{Sh}(A)$, the
family $H$ shatters $S$ and $S\cup\{e\}$ for every $e\in R$: the endpoint
layers $f$ and $\overline f$ both contain all of $A$.  These are
$(r+1)|\operatorname{Sh}(A)|=(r+1)|A|=|H|$ distinct shattered sets, so in
the equality case they exhaust $\operatorname{Sh}(H)$.  If two paths
$P_a$ and $P_b$ were different, their coordinate orders would reverse
some pair $\{e,g\}\subseteq R$.  One path realizes three traces on that
pair, including the trace obtained by flipping $e$ but not $g$, while the
other realizes the complementary intermediate trace obtained by flipping
$g$ but not $e$.  Together with their common antipodal endpoints, their
union therefore shatters $\{e,g\}$.  This is an additional shattered set,
a contradiction.  Hence all $P_a$ equal one path $P$, which proves
\eqref{eq:damp-antipodal-double-equality-product}.  Conversely, the product
$A\times P$ is ample and attains equality.
\end{proof}

\begin{corollary}[Maximum-core dichotomy below Hall's order]
\label{cor:damp-subhall-maximum-core-dichotomy}
Let $H$ be cornerless and ample with $|H|<299$.  Then either $H$ is itself
a pseudogeometric maximum class, or $H$ contains two opposite equal
conditionings whose common value is a pseudogeometric maximum class $A$
with
\[
 |A|\leq\left\lfloor\frac{298}{|R|+1}\right\rfloor\leq149.
\]
\end{corollary}

\begin{proof}
Apply Proposition~\ref{prop:damp-antipodal-double-maximum-fibre}.  If its
set $R$ is empty, then $A=H$.  Otherwise
\eqref{eq:damp-antipodal-double-geodesic-order} gives the displayed bound.
\end{proof}

\begin{corollary}[Maximum core in a least cornerless family]
\label{cor:damp-least-cornerless-maximum-core}
Let $H$ have minimum cardinality among all nonempty cornerless ample
families, and suppose $|H|<299$.  Then either $H$ is a pseudogeometric
maximum class, or Proposition
\ref{prop:damp-antipodal-double-maximum-fibre} gives a set $R$ with
$r=|R|\geq2$ and a pseudogeometric maximum class $A$ such that
\begin{equation}
 (r+1)|A|+1\leq |H|,
 \qquad
 |A|\leq\left\lfloor\frac{297}{r+1}\right\rfloor\leq99.
 \label{eq:damp-least-core-order-bound}
\end{equation}
\end{corollary}

\begin{proof}
Assume $R\neq\varnothing$.  Equality in
\eqref{eq:damp-antipodal-double-geodesic-order} would give
$H=A\times P$.  If $A$ were cornerless, it would be a smaller nonempty
cornerless ample family.  Otherwise $A$ has a corner, and its product with
either endpoint of the path $P$ is a corner of $A\times P$.  Both
alternatives contradict the choice of $H$.  The inequality is therefore
strict.  If $r=1$, there are only the two opposite layers, so
$H=A\times Q_1$ and equality is automatic; hence $r\geq2$.  Finally,
$|H|\leq298$ gives $(r+1)|A|\leq297$, and therefore the displayed
bound.
\end{proof}

\begin{proposition}[Geodesic-disagreement tax]
\label{prop:damp-geodesic-disagreement-tax}
Use the maximum core $A$, the coordinate set $R$, and the sections $P_a$
from Proposition~\ref{prop:damp-antipodal-double-maximum-fibre}, and put
\[
 \Delta(H;A,R)=|H|-(|R|+1)|A|.
\]
Let $G_R(H)$ be the graph on $R$ in which $eg$ is an edge exactly when
the two-coordinate projection $H|_{\{e,g\}}$ is the full square.  Then
\begin{equation}
 |E(G_R(H))|\leq\Delta(H;A,R).
 \label{eq:damp-geodesic-disagreement-tax}
\end{equation}
Moreover,
\begin{equation}
 \Delta(H;A,R)
 =\bigl|H\setminus(Q_R\times A)\bigr|
  +\sum_{a\in A}\bigl(|P_a|-(|R|+1)\bigr).
 \label{eq:damp-geodesic-excess-decomposition}
\end{equation}
In particular, if shortest antipodal paths chosen inside $P_a$ and
$P_b$ flip $e$ and $g$ in opposite orders, then $eg\in E(G_R(H))$ and
this disagreement spends one unit of the excess in
\eqref{eq:damp-geodesic-disagreement-tax}.
\end{proposition}

\begin{proof}
As in the proof of Proposition
\ref{prop:damp-antipodal-double-maximum-fibre}, $H$ shatters
\[
 S\quad\text{and}\quad S\cup\{e\}
 \qquad(S\in\operatorname{Sh}(A),\ e\in R).
\]
These are $(|R|+1)|A|$ distinct supports and each uses at most one
coordinate of $R$.  For every edge $eg$ of $G_R(H)$, the additional
support $\{e,g\}$ is shattered.  Ampleness now gives
\[
 |H|=|\operatorname{Sh}(H)|
 \geq (|R|+1)|A|+|E(G_R(H))|,
\]
which is \eqref{eq:damp-geodesic-disagreement-tax}.  The sections over
the words in $A$ are disjoint, each has at least $|R|+1$ vertices, and
all remaining vertices of $H$ lie over words outside $A$.  Partitioning
$H$ by these sections proves
\eqref{eq:damp-geodesic-excess-decomposition}.

Finally, orient the two coordinates from the common endpoint $f$.  One
of the two paths contains the intermediate trace $10$ on $\{e,g\}$ and
the other contains $01$; their common antipodal endpoints supply $00$
and $11$.  Hence the projection is the full square, as claimed.
\end{proof}

\begin{lemma}[Everything below a least cornerless family is dismantlable]
\label{lem:damp-least-cornerless-cardinality-minimality}
Let $H$ have minimum cardinality among all nonempty cornerless ample
families.  If $K$ is ample and $|K|<|H|$, then $K\in\Damp$.
\end{lemma}

\begin{proof}
Suppose that $K\notin\Damp$, and delete corners successively for as long as
possible.  Every remaining family is ample.  The process cannot delete all
of $K$, since the reverse order would then be a corner peeling of $K$.
It therefore stops at a nonempty cornerless ample family $K'$.  But
$|K'|\le|K|<|H|$, contrary to the choice of $H$.
\end{proof}

\begin{proposition}[Every coordinate shift resolves a least obstruction]
\label{prop:damp-least-cornerless-every-shift-dismantlable}
Let $H$ have minimum cardinality among all nonempty cornerless ample
families, and let $e$ be an active coordinate of $H$.  Write its two
$e$-sections as $B_0,B_1$, and put
\[
 A=B_0\cap B_1,
 \qquad
 C=B_0\cup B_1.
\]
For either orientation of $e$, replace the two sections by $A$ on the
lower side and $C$ on the upper side.  The resulting coordinate shift
\begin{equation}
 U_e(H)=(A\times\{0\})\cup(C\times\{1\})
 \label{eq:damp-least-cornerless-coordinate-shift}
\end{equation}
is ample, has the same order and the same shattered-set complex as $H$,
and belongs to $\Damp$.

Consequently a least cornerless ample family is shift-critical in every
active coordinate: every one-coordinate move to the standard
polarization is already corner-peelable.  In the Yang-complex language,
every paired coordinate shift of its strongly nonshellable Yang ball is
shellable.
\end{proposition}

\begin{proof}
The two-layer shadow-intersection lemma makes $A$ and $C$ ample.  Since
$e$ is active and $H$ is connected, both sections are nonempty and
$A\ne\varnothing$; hence
\[
 |A|<|H|,
 \qquad
 |C|=|H|-|A|<|H|.
\]
Lemma~\ref{lem:damp-least-cornerless-cardinality-minimality} therefore
gives $A,C\in\Damp$.

The family in \eqref{eq:damp-least-cornerless-coordinate-shift} is the
peripheral expansion with lower section $A\subseteq C$.  It is ample by
the two-layer shadow formula: its supports avoiding $e$ are indexed by
$\operatorname{Sh}(C)$, while those containing $e$ are indexed by
$\operatorname{Sh}(A)$.  Their total number is
$|C|+|A|=|U_e(H)|$.

To peel the shifted family, first take any corner peeling of $A$ in its
lower copy.  If $a$ is the next corner of the current lower section and
$Q$ is its unique maximal cube there, then $Q\times\{0,1\}$ is present,
because the upper section contains the lower one.  It contains every edge
of the shifted family incident with $(a,0)$, so $(a,0)$ is a corner.
After the lower copy has been deleted, peel the remaining upper copy of
$C$.  Thus $U_e(H)\in\Damp$.

The equality of the shattered complexes is the paired coordinate-shift
theorem; it also follows directly from the two-layer identities
\[
 \operatorname{Sh}(C)=
 \operatorname{Sh}(B_0)\cup\operatorname{Sh}(B_1),
 \qquad
 \operatorname{Sh}(A)=
 \operatorname{Sh}(B_0)\cap\operatorname{Sh}(B_1).
\]
Interchanging the two signs of $e$ gives the other orientation.  The
Yang-complex assertion is the shellability--corner-peeling dictionary.
\end{proof}

\begin{corollary}[Every overlap corner has two exclusive shields]
\label{cor:damp-least-cornerless-overlap-two-sided-shields}
Keep the notation of Proposition
\ref{prop:damp-least-cornerless-every-shift-dismantlable}.  If $a$ is a
corner of $A$, then for each $b\in\{0,1\}$ there is a coordinate
$f_b$ such that
\begin{equation}
 a^{\{f_b\}}\in B_b\setminus A.
 \label{eq:damp-overlap-corner-two-sided-shields}
\end{equation}
The two shield directions are distinct.  Thus every corner of the
contraction overlap is owned on both sides by two different exclusive
wing concepts.
\end{corollary}

\begin{proof}
Let $Q$ be the unique maximal cube of $A$ through $a$.  Suppose that
$a$ has no neighbour in $B_b\setminus A$.  Then its incident directions
in $B_b$ are exactly the directions of $Q$.  The other copy of every
vertex of $Q$ belongs to $H$, because $Q\subseteq A$.  Hence all
directions incident with $(a,b)$ generate the single cube
$Q\times\{0,1\}$, making $(a,b)$ a corner of $H$, a contradiction.
This proves \eqref{eq:damp-overlap-corner-two-sided-shields} on both
sides.  If $f_0=f_1$, the same neighbour of $a$ would belong to both
sections and hence to $A$, contrary to the displayed inclusions.
\end{proof}

\begin{lemma}[Repeated ownership forces a punctured cube]
\label{lem:damp-exclusive-shield-owner-capacity}
Let $A$ be an isometric cube family, let $z\notin A$, and suppose that
the distinct neighbours
\[
 a_i=z^{\{f_i\}}\in A
 \qquad(i\in I)
\]
are corners of $A$.  Then the face through $z$ on
$F=\{f_i:i\in I\}$ satisfies
\begin{equation}
 Q_F(z)\setminus\{z\}\subseteq A.
 \label{eq:damp-exclusive-owner-punctured-cube}
\end{equation}
In particular,
\begin{equation}
 2^{|I|}-1\le |A|.
 \label{eq:damp-exclusive-owner-capacity}
\end{equation}
Consequently, if $|H|<48$ in Corollary
\ref{cor:damp-least-cornerless-overlap-two-sided-shields}, no one
exclusive wing concept can shield four corners of $A$.
\end{lemma}

\begin{proof}
For distinct $i,j\in I$, the two vertices $a_i,a_j$ have Hamming
distance two.  Isometry of $A$ supplies a two-edge path between them.
Their only common cube neighbours are $z$ and
$z^{\{f_i,f_j\}}$; the first is absent from $A$, so the second belongs
to $A$.

Fix $i\in I$.  The corner $a_i$ is therefore incident in every direction
$f_j$, $j\ne i$.  All these directions lie in its unique maximal cube,
which contains
\[
 z^{\{f_i\}\mathbin\triangle J}
 \qquad(J\subseteq F\setminus\{f_i\}).
\]
As $i$ varies, these are all nonempty vertices of $Q_F(z)$, proving
\eqref{eq:damp-exclusive-owner-punctured-cube} and the cardinality bound.

For the last assertion, Proposition
\ref{prop:damp-nine-concept-passive-carrier-exclusion} shows that the two
exclusive wings in any active coordinate have at least ten concepts each.
Hence
\[
 |A|\le\frac{|H|-20}{2}<14.
\]
Four corners with one owner would force $|A|\ge2^4-1=15$, a
contradiction.
\end{proof}

\begin{proposition}[A fibre corner forces a transverse puncture]
\label{prop:damp-fibre-corner-transverse-puncture}
Let $H\subseteq Q_{R\mathbin{\dot\cup}X}$ be ample and cornerless, let
$g\in H|_R$, and put
\[
 H_g=\{a\in Q_X:(g,a)\in H\}.
\]
If $a$ is a corner of $H_g$, then there is a set
$J\subseteq I_H((g,a))$ such that
\begin{equation}
 |J|\ge2,
 \qquad J\cap R\ne\varnothing,
 \qquad
 (g,a)^U\in H\ (U\subsetneq J),
 \qquad
 (g,a)^J\notin H.
 \label{eq:damp-fibre-corner-transverse-puncture}
\end{equation}
Thus the punctured cube rooted at $(g,a)$ is not contained in the fibre:
at least one of its directions is a projection direction.
\end{proposition}

\begin{proof}
Put $v=(g,a)$.  Since $H$ is cornerless, the directions incident with
$v$ do not generate a cube of $H$.  Choose an inclusion-minimal set
$J\subseteq I_H(v)$ for which $v^J\notin H$.  Every singleton incident
with $v$ is present, so $|J|\ge2$, and minimality gives
$v^U\in H$ for every proper subset $U\subsetneq J$.

It remains to prove that $J$ uses a coordinate of $R$.  If $J\subseteq X$,
then $J\subseteq I_{H_g}(a)$.  Because $a$ is a corner of $H_g$, all
directions incident with $a$ belong to its unique maximal cube.  That cube
would contain $a^J$, and hence $v^J=(g,a^J)$ would belong to $H$, a
contradiction.  Therefore $J\cap R\ne\varnothing$.
\end{proof}

\begin{proposition}[A column corner forces a core-transverse puncture]
\label{prop:damp-column-corner-core-transverse-puncture}
Let $H\subseteq Q_{R\mathbin{\dot\cup}X}$ be ample and cornerless, let
$a\in H|_X$, and put
\[
 P_a=\{g\in Q_R:(g,a)\in H\}.
\]
If $g$ is a corner of $P_a$, then there is a set
$J\subseteq I_H((g,a))$ such that
\begin{equation}
 |J|\ge2,
 \qquad J\cap X\ne\varnothing,
 \qquad
 (g,a)^U\in H\ (U\subsetneq J),
 \qquad
 (g,a)^J\notin H.
 \label{eq:damp-column-corner-core-transverse-puncture}
\end{equation}
Thus a punctured cube rooted at a column corner must use a core direction.
\end{proposition}

\begin{proof}
Put $v=(g,a)$.  Since $H$ is cornerless, the directions incident with
$v$ do not generate a cube of $H$.  Choose an inclusion-minimal
$J\subseteq I_H(v)$ for which $v^J\notin H$.  As in Proposition
\ref{prop:damp-fibre-corner-transverse-puncture}, one has $|J|\ge2$ and
$v^U\in H$ for every proper subset $U\subsetneq J$.

If $J\subseteq R$, then $J\subseteq I_{P_a}(g)$.  Because $g$ is a
corner of $P_a$, all those directions lie in the unique maximal cube of
$P_a$ through $g$.  That cube contains $g^J$, so
$v^J=(g^J,a)\in H$, a contradiction.  Hence $J\cap X\ne\varnothing$.
\end{proof}

\begin{corollary}[Fibrewise transverse-puncture cover]
\label{cor:damp-fibrewise-transverse-puncture-cover}
Let $H$ have minimum cardinality among all nonempty cornerless ample
families, and use the nonempty repair set $R$ supplied by
Proposition~\ref{prop:damp-antipodal-double-maximum-fibre}.  Put
$W=H|_R$.  Every nonempty fibre $H_g$ belongs to $\Damp$.  In particular,
for every $g\in W$ there are a corner $a_g$ of $H_g$ and a
set $J_g$ satisfying
\eqref{eq:damp-fibre-corner-transverse-puncture}.  The resulting rooted
punctures have pairwise distinct repair roots $g$.  If $|H|<48$, then
\[
 2\le |J_g|\le5\qquad(g\in W).
\]
\end{corollary}

\begin{proof}
Every nonempty fibre $H_g$ is ample, because it is a conditioning of $H$.
It is strictly smaller than $H$: the two opposite repair fibres from
Proposition~\ref{prop:damp-antipodal-double-maximum-fibre} are nonempty and
$R\ne\varnothing$.
Lemma~\ref{lem:damp-least-cornerless-cardinality-minimality} therefore gives
$H_g\in\Damp$.

Choose the first member $a_g$ of a corner peeling of $H_g$ and apply
Proposition~\ref{prop:damp-fibre-corner-transverse-puncture}.  Distinct
words $g$ give distinct roots.  Finally, the punctured $J_g$-face contains
$2^{|J_g|}-1$ vertices of $H$.  If $|H|<48$, this excludes
$|J_g|\ge6$.
\end{proof}

\begin{corollary}[Bidirectional transverse-puncture cover]
\label{cor:damp-bidirectional-transverse-puncture-cover}
Let $H\subseteq Q_{R\mathbin{\dot\cup}X}$ have minimum cardinality among
all nonempty cornerless ample families, and suppose that both projections
\[
 W=H|_R,
 \qquad Z=H|_X
\]
have at least two members.  Then every nonempty row $H_g$ and every
nonempty column $P_a$ belongs to $\Damp$.  One may therefore choose
\begin{itemize}
 \item for every $g\in W$, a row corner $a_g$ and a rooted puncture
       $J_g$ with $J_g\cap R\ne\varnothing$;
 \item for every $a\in Z$, a column corner $g_a$ and a rooted puncture
       $L_a$ with $L_a\cap X\ne\varnothing$.
\end{itemize}
The roots within the first family have distinct repair words, and the
roots within the second have distinct core words.  If $|H|<48$, every
puncture in both families has rank between two and five.
\end{corollary}

\begin{proof}
Rows and columns are conditionings of $H$, hence are ample.  The hypothesis
$|W|,|Z|\ge2$ makes every nonempty one strictly smaller than $H$.
Lemma~\ref{lem:damp-least-cornerless-cardinality-minimality} therefore puts
all of them in $\Damp$.

Choose a corner in each row and column.  Proposition
\ref{prop:damp-fibre-corner-transverse-puncture} gives the row-indexed
punctures, while Proposition
\ref{prop:damp-column-corner-core-transverse-puncture} gives the
column-indexed punctures.  Their stated root distinctions follow from their
indices.  Finally, each punctured $J$-face contains $2^{|J|}-1$ vertices of
$H$, so $|H|<48$ excludes $|J|\ge6$ in either family.
\end{proof}

\begin{proposition}[Section corners force a mixed defect or an incidence cycle]
\label{prop:damp-section-corner-incidence-cycle-dichotomy}
Let $H\subseteq Q_{R\mathbin{\dot\cup}X}$ be ample and cornerless, put
$W=H|_R$ and $Z=H|_X$, and suppose that every nonempty row $H_g$ and
column $P_a$ has a corner.  Choose one row corner $a_g\in H_g$ for every
$g\in W$ and one column corner $g_a\in P_a$ for every $a\in Z$.
Then at least one of the following holds.
\begin{enumerate}
 \item Some cell $(g,a)\in H$ is both selected: $a=a_g$ and $g=g_a$.
       At this cell there is a rooted puncture $J$ satisfying
       \begin{equation}
        J\cap R\ne\varnothing,
        \qquad J\cap X\ne\varnothing.
        \label{eq:damp-doubly-selected-mixed-puncture}
       \end{equation}
 \item The incidence graph
       \begin{equation}
        \mathcal I(H)=\bigl(W\mathbin{\dot\cup}Z,
          \{ga:(g,a)\in H\}\bigr)
        \label{eq:damp-row-column-incidence-graph}
       \end{equation}
       contains an even cycle of length at least four whose edges alternate
       between selected row corners and selected column corners.
\end{enumerate}
In particular, a least cornerless ample family in the path-spine regime
always contains either a mixed transverse puncture at a common section
corner or an alternating cycle in its row--column incidence graph.
\end{proposition}

\begin{proof}
Direct the selected incidence edge from $g\in W$ to $a_g\in Z$, and direct
the selected incidence edge from $a\in Z$ to $g_a\in W$.  Every vertex of
the resulting finite directed bipartite graph has out-degree one.  Following
outgoing edges therefore eventually enters a directed cycle.  Its length is
even.  A directed two-cycle is exactly an incidence edge $(g,a)$ selected
from both endpoints, giving item~1.  Every longer directed cycle has distinct
vertices and, after orientations are forgotten, gives the alternating
incidence cycle in item~2.

It remains to verify the mixed puncture assertion in item~1.  Put
$v=(g,a)$.  Cornerlessness of $H$ gives an inclusion-minimal
$J\subseteq I_H(v)$ such that $v^J\notin H$; as before, $|J|\ge2$ and
$v^U\in H$ for every proper $U\subsetneq J$.  If $J\subseteq X$, then the
unique maximal cube through the row corner $a$ in $H_g$ contains $a^J$, a
contradiction.  If $J\subseteq R$, the unique maximal cube through the
column corner $g$ in $P_a$ gives the same contradiction.  Hence $J$ meets
both coordinate blocks, proving
\eqref{eq:damp-doubly-selected-mixed-puncture}.
\end{proof}

\begin{conjecture}[Mutual section-corner principle]
\label{conj:damp-mutual-section-corner}
Let $H\subseteq Q_{R\mathbin{\dot\cup}X}$ be ample, with $R$ and $X$
nonempty.  Suppose that every nonempty row $H_g$ and every nonempty column
$P_a$ belongs to $\Damp$.  Then some concept $(g,a)\in H$ is a corner in
both of its sections: $a$ is a corner of $H_g$ and $g$ is a corner of
$P_a$.
\end{conjecture}

For a least cornerless ample family, the section hypothesis holds by
Lemma~\ref{lem:damp-least-cornerless-cardinality-minimality}.  The
conjecture would therefore force item~1 of Proposition
\ref{prop:damp-section-corner-incidence-cycle-dichotomy}, and hence a
genuinely mixed rooted puncture, without analyzing longer alternating
cycles.  The dismantlability hypothesis cannot simply be omitted: if $C$
is cornerless and ample, then $C\times Q_1$ is ample but has no mutual
section corner for the split $Q_1\mid C$, because its two $C$-rows are
cornerless.  Exact tests verify the stated conjecture for every ample family
in $Q_4$, for the induced-six-cycle witness described in the campaign
ledger, and for every coordinate split of the order-$291$ Hall obstruction.
These tests are not used in any proof.

\begin{lemma}[Fixed-support punctures have disjoint carriers]
\label{lem:damp-fixed-support-puncture-disjointness}
Let $J$ be a coordinate set and let $v_1,\ldots,v_s\in H$ be distinct
roots such that
\[
 v_i^U\in H\quad(U\subsetneq J),
 \qquad v_i^J\notin H
 \qquad(1\le i\le s).
\]
Then the parallel $J$-faces $Q_J(v_1),\ldots,Q_J(v_s)$ are pairwise
disjoint.  Consequently
\begin{equation}
 s(2^{|J|}-1)\le |H|.
 \label{eq:damp-fixed-support-puncture-capacity}
\end{equation}
\end{lemma}

\begin{proof}
Two parallel $J$-faces are either equal or disjoint.  Suppose that
$Q_J(v_i)=Q_J(v_j)$.  The puncture rooted at $v_i$ contains every vertex
of this face except $v_i^J$, whereas the puncture rooted at $v_j$ omits
$v_j^J$.  If the two omitted vertices were different, each would be
required to be both present and absent.  Hence $v_i^J=v_j^J$, and flipping
$J$ gives $v_i=v_j$.  Distinct roots therefore have disjoint faces.  Each
face contributes $2^{|J|}-1$ distinct vertices of $H$, which proves
\eqref{eq:damp-fixed-support-puncture-capacity}.
\end{proof}

\begin{lemma}[Projected cube carriers conserve support]
\label{lem:damp-projected-cube-carrier-support-conservation}
Let $H\subseteq Q_{R\mathbin{\dot\cup}X}$ be ample, let
$\varnothing\ne K\subseteq X$ be strongly shattered by $H$, and define
the $K$-cube carrier and its repair projection by
\begin{equation}
 \mathcal K_K(H)
 =\{y\in Q_{(R\mathbin{\dot\cup}X)\setminus K}:
       Q_K\times\{y\}\subseteq H\},
 \qquad
 W_K=\mathcal K_K(H)|_R.
 \label{eq:damp-projected-core-carrier}
\end{equation}
Then $\mathcal K_K(H)$ and $W_K$ are nonempty ample families.  Moreover,
\begin{equation}
 \operatorname{Sh}(W_K)
 =\{T\subseteq R:K\mathbin{\dot\cup}T
                    \in\operatorname{Sh}(H)\}.
 \label{eq:damp-projected-carrier-support-lift}
\end{equation}
Consequently, if $m$ pairwise distinct repair words occur as projections
of $K$-cube anchors, then $H$ shatters at least $m$ supports whose
intersection with $X$ is exactly $K$.
\end{lemma}

\begin{proof}
For $A\subseteq(R\mathbin{\dot\cup}X)\setminus K$, the carrier
$\mathcal K_K(H)$ strongly shatters $A$ if and only if $H$ strongly
shatters $K\mathbin{\dot\cup}A$.  Indeed, a full $A$-cube of carrier
anchors and the full $K$-cube over each of its vertices together form a
full $(K\mathbin{\dot\cup}A)$-cube of $H$, and the converse follows by
contracting the $K$-directions of such a cube.

If the carrier merely shatters $A$, then for every trace on $A$ it
contains an anchor of a full $K$-cube.  Those cubes supply all traces on
$K\mathbin{\dot\cup}A$, so $H$ shatters that union.  Since $H$ is ample,
the union is strongly shattered, and the preceding equivalence shows
that $A$ is strongly shattered by the carrier.  Thus the shattered and
strongly shattered sets of $\mathcal K_K(H)$ coincide, so the carrier is
ample.  It is nonempty because $K$ is strongly shattered.  Projections
are pc-minors, hence preserve ampleness, and therefore $W_K$ is ample as
well.

The trace of the carrier on a set $T\subseteq R$ is the same as the
trace of its projection $W_K$ on $T$.  Thus
$T\in\operatorname{Sh}(W_K)$ implies that the carrier shatters $T$, and
the preceding argument gives
$K\mathbin{\dot\cup}T\in\operatorname{Sh}(H)$.  Conversely, if $H$
shatters $K\mathbin{\dot\cup}T$, ampleness makes that support strongly
shattered.  Contracting the $K$-directions of the resulting cube gives a
full $T$-cube in $\mathcal K_K(H)$, and hence $T$ is shattered by $W_K$.
This proves \eqref{eq:damp-projected-carrier-support-lift}.

Finally, if $m$ distinct
repair words occur among the projected anchors, then $|W_K|\ge m$.
Ampleness gives $|\operatorname{Sh}(W_K)|=|W_K|$, and
\eqref{eq:damp-projected-carrier-support-lift} supplies the asserted
$m$ supports.  Their core part is exactly $K$ because every
$T\in\operatorname{Sh}(W_K)$ is contained in $R$.
\end{proof}

\begin{corollary}[Every core carrier below a least obstruction is dismantlable]
\label{cor:damp-least-obstruction-core-carriers-dismantlable}
Let $H\subseteq Q_{R\mathbin{\dot\cup}X}$ have minimum cardinality among
all nonempty cornerless ample families, and suppose that $H$ uses at least
two core words.  For a core word $a\in Q_X$, put
\[
 P_a=\{g\in Q_R:(g,a)\in H\}.
\]
Every nonempty column $P_a$ belongs to $\Damp$.  For every nonempty
strongly shattered $K\subseteq X$, the carrier $\mathcal K_K(H)$ and its
repair projection $W_K$ in
\eqref{eq:damp-projected-core-carrier} also belong to $\Damp$.
\end{corollary}

\begin{proof}
The column $P_a$ is a conditioning of $H$, so it is ample.  It is strictly
smaller than $H$ because $H$ contains a vertex with a different core word.
Lemma~\ref{lem:damp-least-cornerless-cardinality-minimality} gives
$P_a\in\Damp$.

Lemma~\ref{lem:damp-projected-cube-carrier-support-conservation} makes the
carrier and its projection ample.  The full $K$-cubes over distinct carrier
anchors are disjoint, and hence
\[
 2^{|K|}|\mathcal K_K(H)|\le |H|.
\]
Thus both $\mathcal K_K(H)$ and $W_K$ have order strictly below $|H|$.
Another application of
Lemma~\ref{lem:damp-least-cornerless-cardinality-minimality} proves the
claim.
\end{proof}

\begin{proposition}[Anchored coface ascent]
\label{prop:damp-anchored-coface-ascent}
Let \(H\subseteq Q_E\) have minimum cardinality among all nonempty
cornerless ample families, and put
\(\Sigma=\operatorname{Sh}(H)\).  For a nonempty \(K\in\Sigma\), a
\emph{positioned \(K\)-cube} is a full \(K\)-cube \(Q_K(y)\subseteq H\),
where \(y\in Q_{E\setminus K}\) is its anchor.  The following statements
hold.

First, the positioned \(K\)-cubes can be matched bijectively with the
cofaces of \(K\) in \(\Sigma\): there is a bijection
\begin{equation}
 \psi_K:\{Q_K(y)\subseteq H\}
 \longrightarrow
 \{L\in\Sigma:K\subseteq L\}
 \label{eq:damp-anchored-carrier-coface-bijection}
\end{equation}
such that \(Q_K(y)\) is contained in a positioned
\(\psi_K(Q_K(y))\)-cube.  Exactly one value of \(\psi_K\) is \(K\);
all other values are strict cofaces of \(K\).

Second, let \(\mathcal A\) be any finite multiset of anchored faces
\((K,Q_K(y))\), with \(K\ne\varnothing\).  Each member may be enlarged,
by repeatedly replacing it with a containing positioned coface, so that
the following normal form holds:
\begin{equation}
 \text{for every final support \(L\), all members labelled by \(L\)
 occupy one positioned \(L\)-cube.}
 \label{eq:damp-anchored-coface-ascent-normal-form}
\end{equation}
In particular, either the members of \(\mathcal A\) acquire pairwise
distinct support labels, or at least two of them end in one common
positioned cube.  Thus collisions between distinct cube positions never
form a terminal obstruction: they ascend to strict cofaces.  The only
terminal multiplicity is a same-positioned-cube packet.
\end{proposition}

\begin{proof}
Fix \(K\ne\varnothing\).  Regard \(K\) as one coordinate block and
\(E\setminus K\) as the other.  Lemma
\ref{lem:damp-projected-cube-carrier-support-conservation} makes the full
carrier
\[
 \mathcal K_K(H)
 =\{y\in Q_{E\setminus K}:Q_K(y)\subseteq H\}
\]
ample and gives the full link identity
\begin{equation}
 \operatorname{Sh}(\mathcal K_K(H))
 =\{S\subseteq E\setminus K:K\mathbin{\dot\cup}S\in\Sigma\}.
 \label{eq:damp-full-carrier-link-identity}
\end{equation}
Indeed, both implications are proved before projection in that lemma: a
full \(S\)-cube of carrier anchors is the same thing as a full
\(K\mathbin{\dot\cup}S\)-cube of \(H\).

The disjoint positioned \(K\)-cubes show that
\(2^{|K|}|\mathcal K_K(H)|\le |H|\).  Hence the carrier is smaller than
\(H\), and minimum-cardinality makes it a member of \(\Damp\), as in
Corollary~\ref{cor:damp-least-obstruction-core-carriers-dismantlable}.
Choose a corner peeling
\[
 F_1=\mathcal K_K(H)\supset F_2\supset\cdots\supset F_s=\{y_s\},
 \qquad F_{i+1}=F_i\setminus\{y_i\}.
\]
For \(i<s\), let \(S_i\) be the support of the unique maximal cube of
\(F_i\) through \(y_i\), and put \(S_s=\varnothing\).  Deleting the
corner \(y_i\) removes exactly the shattered support \(S_i\).  Therefore
the \(S_i\) are distinct and, by ampleness, exhaust
\(\operatorname{Sh}(\mathcal K_K(H))\).  Moreover the positioned
\(S_i\)-cube of carrier anchors through \(y_i\) gives a positioned
\(K\mathbin{\dot\cup}S_i\)-cube of \(H\) containing \(Q_K(y_i)\).
Together with~\eqref{eq:damp-full-carrier-link-identity}, the assignment
\[
 Q_K(y_i)\longmapsto K\mathbin{\dot\cup}S_i
\]
is the required bijection.  Only \(S_s\) is empty, which proves the
strict-coface assertion.

We now prove the normal form.  Start with the given multiset.  Whenever
some support \(K\) labels members occupying at least two distinct
positioned \(K\)-cubes, choose one member at each occupied position and
apply the preceding bijection to those distinct positions.  Replace each
chosen member by the containing positioned coface supplied by the
bijection.  The new support labels are pairwise distinct; at most one
chosen member keeps the label \(K\), while every other chosen member moves
to a strict coface.

Use
\[
 \Phi=\sum_{(L,Q_L(z))\in\mathcal A}|L|
\]
as a potential, counted with multiplicity.  Every such operation strictly
increases \(\Phi\), and \(\Phi\le |E||\mathcal A|\).  Hence the procedure
terminates.  At termination no support is represented at two different
positions, which is precisely
\eqref{eq:damp-anchored-coface-ascent-normal-form}.  If no final position
contains two members, then every support labels at most one member and the
labels are pairwise distinct.  Otherwise the repeated members form the
asserted same-positioned-cube packet.
\end{proof}

\begin{lemma}[Compatible signed traces in an ample two-layer family]
\label{lem:damp-ample-two-layer-shadow-intersection}
Let $x$ be a coordinate and let
\[
 Y=(P\times\{0\})\cup(Q\times\{1\})
\]
be ample, with both sections nonempty.  Then $P$, $Q$, $P\cup Q$, and
$P\cap Q$ are ample.  More precisely, for every set $T$ of coordinates
different from $x$,
\begin{equation}
 \pi_T(P)\cap\pi_T(Q)=\pi_T(P\cap Q).
 \label{eq:damp-two-layer-signed-trace-intersection}
\end{equation}
Equivalently, their Yang complexes satisfy the signed identity
\begin{equation}
 \Delta_P\cap\Delta_Q=\Delta_{P\cap Q}.
 \label{eq:damp-two-layer-yang-intersection}
\end{equation}
Consequently,
\begin{equation}
 \operatorname{Sh}(P)\cap\operatorname{Sh}(Q)
 =\operatorname{Sh}(P\cap Q).
 \label{eq:damp-two-layer-shadow-intersection}
\end{equation}
Moreover,
\begin{equation}
 \operatorname{Sh}(Y)
 =\operatorname{Sh}(P\cup Q)
  \mathbin{\dot\cup}
  \{S\cup\{x\}:S\in\operatorname{Sh}(P\cap Q)\}.
 \label{eq:damp-two-layer-shadow-decomposition}
\end{equation}
Furthermore,
\begin{equation}
 \operatorname{Sh}(P\cup Q)
 =\operatorname{Sh}(P)\cup\operatorname{Sh}(Q).
 \label{eq:damp-two-layer-shadow-union}
\end{equation}
\end{lemma}

\begin{proof}
Fix a trace
\(\alpha\in\pi_T(P)\cap\pi_T(Q)\), and choose
\(p\in P\) and \(q\in Q\) which both realize \(\alpha\).  The ample
family $Y$ is isometric.  A hypercube geodesic in $Y$ from $(p,0)$ to
$(q,1)$ crosses the $x$-coordinate in an edge
\[
 (r,0)(r,1),\qquad r\in P\cap Q.
\]
The endpoints $p$ and $q$ agree on every coordinate of $T$, so a
hypercube geodesic between them changes no coordinate of $T$.  Hence
$r|_T=\alpha$.  This proves the nontrivial inclusion in
\eqref{eq:damp-two-layer-signed-trace-intersection}; the reverse inclusion
is immediate.  Faces of a Yang complex are precisely signed traces, so
\eqref{eq:damp-two-layer-yang-intersection} is equivalent to the trace
identity.  Taking all traces on a coordinate support gives
\eqref{eq:damp-two-layer-shadow-intersection}.

The two sections and their union are restrictions and a contraction of
$Y$, so they are ample.  Direct inspection of the two sections gives
\[
 \operatorname{Sh}(Y)
 =\operatorname{Sh}(P\cup Q)
  \mathbin{\dot\cup}
  \{S\cup\{x\}:S\in
       \operatorname{Sh}(P)\cap\operatorname{Sh}(Q)\}.
\]
Using the Sandwich equality for $Y$ and $P\cup Q$ now gives
\[
 |\operatorname{Sh}(P)\cap\operatorname{Sh}(Q)|
 =|P|+|Q|-|P\cup Q|=|P\cap Q|.
\]
On the other hand, the Sandwich inequality gives
$|\operatorname{Sh}(P\cap Q)|\ge|P\cap Q|$.  Equality holds throughout.
Thus $P\cap Q$ is ample,
and substitution gives \eqref{eq:damp-two-layer-shadow-decomposition}.
Finally,
$\operatorname{Sh}(P)\cup\operatorname{Sh}(Q)$ is contained in
$\operatorname{Sh}(P\cup Q)$, while
\[
 \left|\operatorname{Sh}(P)\cup\operatorname{Sh}(Q)\right|
 =|P|+|Q|-|P\cap Q|
 =|P\cup Q|
 =|\operatorname{Sh}(P\cup Q)|.
\]
This proves \eqref{eq:damp-two-layer-shadow-union}.
\end{proof}

\begin{corollary}[Monochromatic upper stars outside the overlap]
\label{cor:damp-two-layer-monochromatic-upper-star}
In the setting of
Lemma~\ref{lem:damp-ample-two-layer-shadow-intersection}, put
$U=P\cup Q$ and $K=P\cap Q$.  If
\[
 S\in\operatorname{Sh}(U)\setminus\operatorname{Sh}(K),
\]
then exactly one of $P,Q$ shatters $S$.  The same section shatters every
$T\in\operatorname{Sh}(U)$ with $S\subseteq T$.

In particular, if $U$ is a maximum class of VC dimension three on $d$
coordinates and a pair $S$ is not shattered by $K$, then all $d-2$
triples containing $S$ are shattered by one common section.
\end{corollary}

\begin{proof}
Equation~\eqref{eq:damp-two-layer-shadow-union} puts $S$ in at least one
section shadow, whereas
\eqref{eq:damp-two-layer-shadow-intersection} prevents it from lying in
both.  Suppose, say, that $P$ shatters $S$.  If $Q$ shattered a set
$T\supseteq S$, downward closure of the shattering complex would put
$S$ in $\operatorname{Sh}(Q)$, a contradiction.  Thus every such
$T$ is shattered by $P$, again by
\eqref{eq:damp-two-layer-shadow-union}.  A maximum VC-dimension-three
class shatters every triple, which gives the final assertion.
\end{proof}

\begin{proposition}[The bichromatic core must be nonample]
\label{prop:damp-two-layer-bichromatic-core}
Let $U\subseteq Q_E$ be a maximum class of VC dimension three whose
cubical complex is pure, and let
\[
 Y=(P\times\{0\})\cup(Q\times\{1\})
 \subseteq Q_{E\cup\{e\}}
\]
be ample.  Put $K=P\cap Q$, and suppose that
$\operatorname{VCdim}(K)\le2$.  For every triple $T\subseteq E$, let
$C_T$ be the unique $T$-cube of $U=P\cup Q$.  The cube $C_T$
has one of the two \emph{owners} $P,Q$, according as
\[
 T\in\operatorname{Sh}(P)\setminus\operatorname{Sh}(Q)
 \quad\hbox{or}\quad
 T\in\operatorname{Sh}(Q)\setminus\operatorname{Sh}(P).
\]
Define the \emph{bichromatic core}
\[
 M=\{v\in U:\text{$v$ belongs to triple cubes of both owners}\}.
\]
If $M$ is nonempty and ample and $M\subsetneq K$, then $Y$ has a
corner.

Consequently, if $U$ has a corner and $Y$ is cornerless, then $M$
is nonempty and nonample.
\end{proposition}

\begin{proof}
The two-layer shadow identities imply that every triple belongs to
exactly one section shadow, since $K$ has VC dimension at most two.
Maximum classes have one cube on each top-dimensional support, which
defines the asserted owner.  A vertex incident with cubes of both owners
belongs to both sections, and hence
\[
 M\subseteq K.
\]

Suppose first that $M$ is nonempty and ample and that $M\subsetneq K$.
The relative VC-two peeling lemma,
Lemma~\ref{lem:damp-relative-peeling-vc-two}, supplies a corner
$v\in K\setminus M$.  All maximal three-cubes of $U$ through $v$
have one owner; interchange the layers so that this owner is $P$.
Let $D$ be the unique maximal cube of $K$ through $v$.

We claim that the upper copy $(v,1)$ lies in the unique maximal cube
$D\times\{0,1\}$ of $Y$.  Indeed, any cube of $Q$ through $v$
is a cube of $U$.  Purity extends it to a maximal three-cube of $U$
through $v$, whose owner is $P$.  Every vertex of the original
$Q$-cube therefore lies in both $P$ and $Q$, so that cube is
contained in $K$, and then in $D$.  The same is automatic for every
vertical cube through $(v,1)$.  Thus $(v,1)$ is a corner of $Y$.

Now assume that $U$ has a corner $u$ and that $Y$ is cornerless.
The unique maximal cube of $U$ through $u$ has one owner, so
$u\notin M$.  If $u\notin K$, its sole copy in the owner layer lies
in a unique maximal cube of $Y$, contrary to cornerlessness.  Hence
$u\in K\setminus M$, and in particular $M\subsetneq K$.
If $M$ were empty, any corner of the nonempty ample VC-two family
$K$ would be monochromatic, and the preceding unique-maximal-cube
argument would again give a corner of $Y$.  Thus $M\ne\varnothing$.
The first part now shows that $M$ cannot be ample.
\end{proof}

\begin{proposition}[Exact relative-corner criterion]
\label{prop:damp-two-layer-relative-corner-criterion}
In the setting of
Proposition~\ref{prop:damp-two-layer-bichromatic-core}, write
$v^P=(v,0)$ and $v^Q=(v,1)$.  If $v\notin M$, let $o(v)\in\{P,Q\}$ be
the common owner of all maximal three-cubes of $U$ through $v$, and let
$\overline{o}(v)$ be the other owner.  Then
\[
 \operatorname{Cor}(Y)
 =
 \{u^{o(u)}:u\in\operatorname{Cor}(U)\setminus K\}
 \,\dot\cup\,
 \{v^{\overline{o}(v)}:
       v\in\operatorname{Cor}(K)\setminus M\}.
 \label{eq:damp-two-layer-relative-corner-identity}
\]
Consequently,
\[
 Y\text{ is cornerless}
 \quad\Longleftrightarrow\quad
 \operatorname{Cor}(U)\subseteq K
 \ \text{ and }\
 \operatorname{Cor}(K)\subseteq M.
 \label{eq:damp-two-layer-relative-corner-criterion}
\]
\end{proposition}

\begin{proof}
First let $v\notin K$.  Then $v\notin M$, so all maximal three-cubes of
$U$ through $v$ have owner $o(v)$, and only the copy $v^{o(v)}$ belongs
to $Y$.  Purity of $U$ identifies the maximal cubes of $Y$ through this
copy with the maximal three-cubes of $U$ through $v$.  Hence
$v^{o(v)}$ is a corner of $Y$ exactly when $v$ is a corner of $U$.

Now let $v\in K$.  Every maximal cube $D$ of $K$ through $v$ gives a
vertical maximal cube $D\times\{0,1\}$ of $Y$.  If $v$ belongs to at
least two maximal cubes of $K$, these vertical cubes already prevent
both copies of $v$ from being corners.

It remains to suppose that $v$ is a corner of $K$, with unique maximal
cube $D$.  If $v\in M$, then a $P$-owned and a $Q$-owned maximal
three-cube of $U$ pass through $v$.  Each copy of $v$ therefore lies in
both a horizontal maximal cube of its own layer and the vertical cube
$D\times\{0,1\}$, so neither copy is a corner.

Finally suppose that $v\notin M$, and interchange the layers if
necessary so that $o(v)=P$.  Every cube of $Q$ through $v$ extends in
$U$ to a $P$-owned maximal three-cube.  Its vertices consequently
belong to both sections, so the cube lies in $K$, and then in $D$.
Thus $v^Q$ lies only in the vertical maximal cube
$D\times\{0,1\}$ and is a corner.  The copy $v^P$ also lies in a
horizontal $P$-owned maximal three-cube and is not a corner.  These
three cases prove the identity and its stated consequence.
\end{proof}

\begin{proposition}[Relative leaf-charge identity]
\label{prop:damp-relative-leaf-charge}
Let \(K\subseteq Q_E\) be ample of VC dimension at most two, and suppose that
every maximal cube of \(K\) is a square.  For \(i\in E\), define the
\emph{singleton carrier} \(T_i(K)\) as follows: its vertices are the
\(i\)-edges of \(K\), and two such edges are adjacent when their union is a
square of \(K\).  Every nonempty \(T_i(K)\) is a tree.

For an \(i\)-edge \(e\), an endpoint \(x\in e\), and a set \(M\subseteq K\),
let \(h_i(x)\) be the number of maximal squares of \(K\) that contain \(x\)
but whose support does not contain \(i\).  Then
\begin{equation}
 2|\operatorname{Cor}(K)\setminus M|
 =
 \sum_{i\in E}
 \sum_{\substack{e\text{ leaf of }T_i(K)}}
 \bigl|\{x\in e\setminus M:h_i(x)=0\}\bigr|.
 \label{eq:damp-relative-leaf-charge}
\end{equation}
In particular,
\(\operatorname{Cor}(K)\subseteq M\) exactly when every endpoint outside
\(M\) of every carrier leaf has a maximal-square shield not using the
carrier coordinate.
\end{proposition}

\begin{proof}
The carrier of a strongly shattered singleton in an ample family is ample.
After contracting its \(i\)-edges, its one-inclusion graph has VC dimension
at most one and is connected.  An ample family of VC dimension at most one
has a tree as its one-inclusion graph.  This is precisely \(T_i(K)\).

Fix an \(i\)-edge \(e\) and \(x\in e\).  Every maximal square through \(x\)
whose support contains \(i\) contains the whole edge \(e\), and corresponds
to one edge of \(T_i(K)\) incident with \(e\).  Hence the number of maximal
squares of \(K\) through \(x\) is
\[
 \deg_{T_i(K)}(e)+h_i(x).
\]
If \(e\) is a leaf, the endpoint \(x\) is therefore a corner of \(K\)
exactly when \(h_i(x)=0\).  Conversely, if \(x\) is a corner and its unique
maximal square has support \(\{i,j\}\), then the \(i\)-edge and the
\(j\)-edge of that square through \(x\) are carrier leaves and have no
off-carrier shield at \(x\).  Thus every corner outside \(M\) is counted
twice on the right of \eqref{eq:damp-relative-leaf-charge}, once for each
direction of its unique square, and no other endpoint is counted.
\end{proof}

\begin{proposition}[Owner-deficit identity]
\label{prop:damp-two-layer-owner-deficit}
In the setting of
Proposition~\ref{prop:damp-two-layer-bichromatic-core}, let $A_P$ be the
set of vertices of $U$ incident only with $P$-owned maximal cubes, and
define $A_Q$ symmetrically.  Let $t_P$ and $t_Q$ be the numbers of
triples owned by $P$ and $Q$, respectively, and put
\[
 \ell_P=
 \bigl|\{S\in\operatorname{Sh}(P)\setminus\operatorname{Sh}(K):
              |S|\le2\}\bigr|,
\]
with the analogous definition of $\ell_Q$.  Then
\begin{equation}
 |A_P|-t_P=|A_P\cap K|+\ell_P,
 \qquad
 |A_Q|-t_Q=|A_Q\cap K|+\ell_Q.
 \label{eq:damp-two-layer-owner-deficit}
\end{equation}
In particular, both owner deficits are nonnegative.  If $Y$ is
cornerless, then
\begin{equation}
 |A_P|-t_P\ge|A_P\cap\operatorname{Cor}(U)|,
 \qquad
 |A_Q|-t_Q\ge|A_Q\cap\operatorname{Cor}(U)|.
 \label{eq:damp-two-layer-owner-corner-surcharge}
\end{equation}
Consequently, an owner with zero deficit cannot own a corner of $U$.
\end{proposition}

\begin{proof}
Purity implies that every vertex of $U$ is incident with a maximal
three-cube.  Every vertex in a $P$-owned cube belongs to $P$, every
vertex in a $Q$-owned cube belongs to $Q$, and every vertex incident
with both owners belongs to $K$.  It follows that
\[
 P\setminus K=A_P\setminus K,
 \qquad
 Q\setminus K=A_Q\setminus K.
\]
The nested families $K\subseteq P$ are both ample, so their Sandwich
equalities give
\[
 |A_P\setminus K|=|P\setminus K|
 =|\operatorname{Sh}(P)\setminus\operatorname{Sh}(K)|.
\]
Because $K$ has VC dimension at most two, the three-element supports on
the right are precisely the $t_P$ triples owned by $P$; the remaining
supports are the $\ell_P$ lower-dimensional ones.  Therefore
\[
 |A_P|-|A_P\cap K|=t_P+\ell_P,
\]
which is the first identity in
\eqref{eq:damp-two-layer-owner-deficit}.  The second is symmetric.

If $Y$ is cornerless, the proof of
Proposition~\ref{prop:damp-two-layer-bichromatic-core} shows that every
corner of $U$ belongs to $K$.  Such a corner is incident with one
maximal cube and hence belongs to the corresponding one of $A_P,A_Q$.
Dropping the nonnegative terms $\ell_P,\ell_Q$ from
\eqref{eq:damp-two-layer-owner-deficit} proves
\eqref{eq:damp-two-layer-owner-corner-surcharge}.
\end{proof}

\begin{corollary}[Forced-core shadow bound]
\label{cor:damp-two-layer-forced-core-shadow}
In the setting of
Proposition~\ref{prop:damp-two-layer-bichromatic-core}, suppose that $Y$
is cornerless and put
\[
 F=M\cup\operatorname{Cor}(U).
\]
Then
\begin{equation}
 F\subseteq K,
 \qquad
 |K|\ge|\operatorname{Sh}(F)|,
 \qquad
 |Y|\ge|U|+|\operatorname{Sh}(F)|.
 \label{eq:damp-two-layer-forced-core-shadow}
\end{equation}
Thus a target order for $Y$ imposes one owner-only shadow bound before
the overlap is chosen.
\end{corollary}

\begin{proof}
The bichromatic core lies in $K$, and the proof of
Proposition~\ref{prop:damp-two-layer-bichromatic-core} puts every corner
of $U$ in $K$.  Hence $F\subseteq K$.  Since $K$ is ample, the Sandwich
equality and monotonicity of shattering give
\[
 |K|=|\operatorname{Sh}(K)|
 \ge|\operatorname{Sh}(F)|.
\]
Finally, the two-layer decomposition gives $|Y|=|U|+|K|$.
\end{proof}

\begin{corollary}[Cubical-completion bound]
\label{cor:damp-two-layer-cubical-completion}
For $F\subseteq U$, define
\[
 \kappa_U(F)=
 \min\bigl\{|Z|:F\subseteq Z\subseteq U
 \text{ and every }S\in\operatorname{Sh}(F)
 \text{ is strongly shattered by }Z\bigr\}.
\]
In the setting of
Corollary~\ref{cor:damp-two-layer-forced-core-shadow}, every cornerless lift
satisfies
\[
 |K|\ge\kappa_U(F)
 \qquad\hbox{and}\qquad
 |Y|\ge|U|+\kappa_U(F).
\]
Equivalently, $Z$ must contain a positioned $S$-cube for every support
$S$ shattered by the forced core.  This can be strictly stronger than merely
counting the supports in $\operatorname{Sh}(F)$.
\end{corollary}

\begin{proof}
Because $F\subseteq K$, monotonicity gives
$\operatorname{Sh}(F)\subseteq\operatorname{Sh}(K)$.  Ampleness identifies
shattering and strong shattering, so $K$ strongly shatters every support in
$\operatorname{Sh}(F)$.  Thus $K$ is feasible in the minimization defining
$\kappa_U(F)$, and the first inequality follows.  The second again uses
$|Y|=|U|+|K|$.
\end{proof}

\begin{corollary}[Relative-corner completion bound]
\label{cor:damp-two-layer-relative-corner-completion}
For $M\subseteq F\subseteq U$, define
\[
 \lambda_U(F,M)=
 \min\bigl\{|Z|:F\subseteq Z\subseteq U,\,
 Z\text{ is ample},\,
 \operatorname{VCdim}(Z)\le2,\,
 \operatorname{Cor}(Z)\subseteq M\bigr\},
\]
with value $\infty$ if no such $Z$ exists.  In the setting of
Corollary~\ref{cor:damp-two-layer-forced-core-shadow}, every cornerless lift
satisfies
\[
 |K|\ge\lambda_U(F,M)
 \qquad\hbox{and}\qquad
 |Y|\ge|U|+\lambda_U(F,M).
\]
\end{corollary}

\begin{proof}
The overlap $K$ contains $F$, is ample of VC dimension at most two, and
Proposition~\ref{prop:damp-two-layer-relative-corner-criterion} gives
$\operatorname{Cor}(K)\subseteq M$.  Thus $K$ is feasible in the
minimization defining $\lambda_U(F,M)$.
\end{proof}

\begin{corollary}[Corner-cube cover obstruction]
\label{cor:damp-two-layer-corner-cube-cover}
In the setting of
Proposition~\ref{prop:damp-two-layer-bichromatic-core}, suppose that $Y$
is cornerless.  If $u$ is a corner of $U$ and $C$ is the unique maximal
cube of $U$ containing $u$, then
\[
 u\in K
 \qquad\hbox{and}\qquad
 C\setminus\bigl(M\cup\{u\}\bigr)\ne\varnothing.
\]
Thus every contraction-corner cube has a vertex distinct from its corner
that is incident only with maximal cubes of its own owner.
\end{corollary}

\begin{proof}
The proof of
Proposition~\ref{prop:damp-two-layer-bichromatic-core} shows that
$u\in K$.  If every other vertex of $C$ belonged to $M$, then
$M\subseteq K$ would give $C\subseteq K$.  The support of $C$ would then
be shattered by $K$, contrary to $\operatorname{VCdim}(K)\le2$.
\end{proof}

\begin{remark}
Proposition~\ref{prop:damp-two-layer-bichromatic-core} and
Corollary~\ref{cor:damp-two-layer-corner-cube-cover} can be checked before
the overlap $K$ is chosen.  They therefore eliminate complete owner
colourings, simultaneously for every possible overlap order, rather than
merely eliminating one cardinality layer of the reverse-expansion search.
\end{remark}

\begin{corollary}[Monochromatic components of the overlap-defect graph]
\label{cor:damp-two-layer-defect-components}
In the setting of
Corollary~\ref{cor:damp-two-layer-monochromatic-upper-star}, suppose that
$U$ is maximum of VC dimension three on $d$ coordinates.  Form a graph
$M_K$ on the coordinates by declaring $ij$ to be an edge when
$\{i,j\}\notin\operatorname{Sh}(K)$.  Every nontrivial connected component
of $M_K$ has one owner: all its edges, and all triples containing any of
those edges, belong to the same one of the two section shadows.

If $K$ shatters every singleton, let $m$ be the number of edges of $M_K$
and let $c$ be the number of triples shattered by $K$.  Then
\begin{equation}
 |K|=1+d+\binom d2-m+c.
 \label{eq:damp-overlap-defect-budget}
\end{equation}
In particular, an overlap smaller than the maximum VC-two order satisfies
$m\ge c+1$.  Every one of its common triples is supported on a triangle of
the complement graph $\overline{M_K}$.
\end{corollary}

\begin{proof}
Two incident missing pairs have a three-element union $T$.  The
monochromatic-upper-star property assigns $T$ the owner of either pair, so
the two owners must agree.  Propagation along a path in $M_K$ proves the
component statement, and the same upper-star property assigns every triple
containing a component edge to that component's owner.

The shadow of $K$ is a down-set contained in the rank-three shadow of $U$.
When all singletons are present, it therefore consists of the empty set,
the $d$ singletons, the $\binom d2-m$ common pairs, and the $c$ common
triples.  Since $K$ is ample, the Sandwich equality gives
\eqref{eq:damp-overlap-defect-budget}.  Comparing this identity with
$1+d+\binom d2$ gives $m\ge c+1$ in the stated range.  Finally, downward
closure forbids a common triple from containing an edge missing from the
common shadow.
\end{proof}

\begin{proposition}[Carrier-subtree and one-cut principles]
\label{prop:damp-two-layer-carrier-one-cut}
Let $U\subseteq Q_E$ be an ample family of VC dimension at most three and let
$P\subseteq U$ be ample.  For a pair $R\in\operatorname{Sh}(P)$, the
edges of the $R$-cube carrier of $U$ whose triple supports belong to
$\operatorname{Sh}(P)$ form a connected subtree of the carrier of $U$.
If there is no such triple, the carrier of $P$ consists of one vertex.

Now suppose that $U$ is maximum of VC dimension three, every pair
carrier of $U$ is a path, and
\[
 Y=(P\times\{0\})\cup(Q\times\{1\})
\]
is ample with $P\cup Q=U$.  If, for some pair $S$,
\begin{equation}
 \operatorname{Sh}(P\cap Q)
 =\{R\subseteq E:|R|\leq2\}\setminus\{S\},
 \label{eq:damp-one-pair-deficit-shadow}
\end{equation}
then the following assertions hold.
\begin{enumerate}
 \item All triples containing $S$ belong to one section shadow.
 \item Every other pair-carrier path is cut into at most two monochromatic
       edge intervals, one belonging to $P$ and one to $Q$.
 \item If both colours occur, their unique common pair cube is the carrier
       vertex at the cut.  If the path is monochromatic, the section with
       no edge chooses one carrier vertex, and this is again the unique
       common pair cube.
 \item If $Y$ is cornerless and $v$ is a corner of $U$ with unique maximal
       triple cube $C_T$, then at least two of the three pair facets of
       $C_T$ containing $v$ belong to $P\cap Q$.  Equivalently, at least two
       of the three corresponding pair carriers choose the vertex $v$.
 \item If $Y$ is cornerless and $v,w$ are two corners of $U$ whose unique
       triple supports $T_v,T_w$ meet in one coordinate, then $T_v$ and
       $T_w$ have the same owner.
 \item More generally, if a missing common pair $S'$ meets the unique
       triple support $T_v$ of a corner $v$ in one coordinate, then $S'$
       and $T_v$ have the same owner.
\end{enumerate}
Thus a lift satisfying \eqref{eq:damp-one-pair-deficit-shadow} is specified
at rank three by one colour per triple, at most one cut per pair-carrier,
and one vertex choice only on a monochromatic carrier.  No independent
choice of the vertices of $P$ and $Q$ remains.
\end{proposition}

\begin{proof}
The $R$-cube carrier $\mathcal K_R(P)$ is an ample subfamily of
$\mathcal K_R(U)$.  Its one-inclusion graph is connected.  An edge labelled
$e$ belongs to this subfamily exactly when $P$ contains the unique
$(R\cup\{e\})$-cube of $U$, or equivalently when
$R\cup\{e\}\in\operatorname{Sh}(P)$.  Since the carrier of $U$ is a tree,
these edges and their incident vertices form a connected subtree.  With no
edge, ampleness and the Sandwich equality leave one carrier vertex.

Apply Lemma~\ref{lem:damp-ample-two-layer-shadow-intersection}.  Under
\eqref{eq:damp-one-pair-deficit-shadow}, every triple of $U$ belongs to
exactly one of the two section shadows.  The first assertion is
Corollary~\ref{cor:damp-two-layer-monochromatic-upper-star}.  Fix
$R\ne S$.  Both sections shatter $R$, and the preceding paragraph says
that the edges of either colour form a connected subpath.  The two edge
sets partition the carrier path, so they are two intervals with at most one
cut.  If both are nonempty, the corresponding vertex subpaths meet exactly
at the cut vertex.  A full $R$-cube lies in $P\cap Q$ precisely when its
carrier vertex belongs to both subpaths, so that cut vertex is the unique
common $R$-cube.  If, say, every edge belongs to $P$, then the carrier of
$P$ is the full path, whereas the edge-free carrier of $Q$ is one vertex;
their intersection is that vertex.  This proves the first three assertions
and the rank-three parametrization.

It remains to prove the corner assertions.  The unique triple cube $C_T$ is
owned by one section, say $P$.  The copy of $v$ in that layer belongs to
one horizontal maximal cube, so cornerlessness requires a vertical maximal
cube through it.  Such a cube is the product with the new coordinate of a
common pair facet of $C_T$ through $v$.  This vertical cube also puts the
opposite-layer copy of $v$ in $Y$.  That copy belongs to no horizontal
triple cube, because $C_T$ is the only maximal cube of $U$ containing $v$
and is not owned by $Q$.  A second vertical pair cube is therefore required
to keep the opposite-layer copy from being a corner.  Hence at least two of
the three pair facets through $v$ are common.

For the two-corner assertion, a common pair cube through the corner $v$ is a
leaf of its carrier path.  It cannot be the cut vertex of two nonempty
colour intervals, because such a cut is internal.  Consequently the path
is monochromatic, with the same owner as $T_v$.  Among any two pair facets
of $T_v$, one contains the coordinate in $T_v\cap T_w$; the same holds at
$w$.  If the two corner triples had opposite owners, choose such a common
pair at each corner.  Their union is a triple and would belong to both
section shadows by the monochromatic-upper-star property, a contradiction.
Thus the owners agree.

The same argument proves the more general statement.  Of the two common
pair facets forced through $v$, at least one contains the coordinate in
$S'\cap T_v$.  Its carrier is monochromatic with the owner of $T_v$, while
Corollary~\ref{cor:damp-two-layer-monochromatic-upper-star} makes the entire
upper star of $S'$ monochromatic with the owner of $S'$.  The union of the
two pairs is a triple, so the two owners must coincide.
\end{proof}

\begin{corollary}[Two-corner trace obstruction]
\label{cor:damp-two-layer-two-corner-trace-obstruction}
Assume the one-pair-deficit hypotheses of
Proposition~\ref{prop:damp-two-layer-carrier-one-cut}, and write the missing
pair as $S=\{a,b\}$.  Suppose that $U$ has corners $v_0,v_1$ with unique
triple supports $T_0,T_1$ such that
\[
 T_0\cap S=T_1\cap S=\{a\}
 \qquad\text{and}\qquad
 v_0(b)\ne v_1(b).
\]
Then the lift $Y$ has a corner.
\end{corollary}

\begin{proof}
Suppose that $Y$ is cornerless and put $K=P\cap Q$.  At least two of the
three pair facets through $v_i$ are common, by
Proposition~\ref{prop:damp-two-layer-carrier-one-cut}.  Since $a\in T_i$,
at least one of these two pair facets contains $a$.  Its common pair cube
contains both $v_i$ and $v_i\mathbin{\triangle}\{a\}$.  Thus the trace of
$K$ on $S$ contains both labels having $b$-coordinate $v_i(b)$.  The two
values $v_0(b)$ and $v_1(b)$ are different, so together the two corners
force all four labels on $S$.  This says that $K$ shatters $S$, contrary to
\eqref{eq:damp-one-pair-deficit-shadow}.
\end{proof}

\begin{remark}[The first Hall shield profile closes without search]
\label{ex:damp-hall-first-trace-shield-profile}
In the coordinate-zero projection of Hall's class, the two corners are the
projected vertices $1600$ and $1894$.  Their unique triple supports are
$\{0,5,6\}$ and $\{0,7,8\}$, while their values on coordinate $1$ are
$0$ and $1$, respectively.  Taking $S=\{0,1\}$ in
Corollary~\ref{cor:damp-two-layer-two-corner-trace-obstruction} excludes
every one-pair-deficit lift with this missing pair.  In particular, the
apparent $3\times3$ choices of two shields at the two corners are not nine
separate cases: the same two-line trace argument eliminates all of them.
\end{remark}

\begin{proposition}[Exterior-component normal form for a one-pair deficit]
\label{prop:damp-two-layer-exterior-component-normal-form}
Let $U\subseteq Q_E$ be maximum of VC dimension three, and write $C_T$ for
its unique triple cube with support $T\in\binom E3$.  Let $K\subseteq U$ be
ample with
\[
 \operatorname{Sh}(K)
 =\{R\subseteq E:|R|\leq2\}\setminus\{S\},
 \qquad S\in\binom E2.
\]
Form a graph on the triple supports by joining $T,T'$ whenever
$C_T\cap C_{T'}$ contains a vertex outside $K$, and additionally join all
triples containing $S$.  Let $\mathcal B$ be its set of connected components,
let $B_S$ be the component containing the triples above $S$, and put
\[
 V_B=\left(\bigcup_{T\in B}C_T\right)\setminus K,
 \qquad
 \omega(B)=|V_B|-|B|-\mathbf1_{B=B_S}.
 \label{eq:damp-exterior-component-weight}
\]
Then the sets $V_B$ partition $U\setminus K$ and
\begin{equation}
 \sum_{B\in\mathcal B}\omega(B)=0.
 \label{eq:damp-exterior-component-total-balance}
\end{equation}

For $\mathcal R\subseteq\mathcal B$ with $B_S\in\mathcal R$, define
\[
 P_{\mathcal R}=K\cup\bigcup_{B\in\mathcal R}\bigcup_{T\in B}C_T,
 \qquad
 Q_{\mathcal R}=K\cup\bigcup_{B\notin\mathcal R}\bigcup_{T\in B}C_T.
\]
The following are equivalent:
\begin{enumerate}
 \item $P_{\mathcal R}$ and $Q_{\mathcal R}$ are ample;
 \item $\sum_{B\in\mathcal R}\omega(B)=0$.
\end{enumerate}
Under these conditions, $P_{\mathcal R}\cup Q_{\mathcal R}=U$,
$P_{\mathcal R}\cap Q_{\mathcal R}=K$, and the corresponding two-layer
lift is ample.  Conversely, every ample two-layer lift with contraction
$U$ and overlap $K$ is obtained in this way, up to interchanging its two
layers.

For $v\in U$, let $\nu_K(v)$ be the number of maximal pair cubes of $K$
containing $v$, and let $\rho_{\mathcal R}(v)$ and
$\rho_{\overline{\mathcal R}}(v)$ count the triple cubes through $v$ whose
supports lie in $\mathcal R$ and its complement.  The lift is cornerless if
and only if
\begin{equation}
 \nu_K(v)+\rho_{\mathcal R}(v)\ne1,
 \qquad
 \nu_K(v)+\rho_{\overline{\mathcal R}}(v)\ne1
 \quad(v\in U).
 \label{eq:damp-exterior-component-no-isolated-incidence}
\end{equation}
Equivalently, the lift is cornerless if and only if
\begin{enumerate}
 \item every corner of $U$ belongs to at least two maximal pair cubes of
       $K$; and
 \item every corner of $K$ is incident with triple cubes from both sides of
       the component partition.
\end{enumerate}
If $\mathcal B$ has one component, no cornerless lift exists.
\end{proposition}

\begin{proof}
If $v\notin K$, all triple cubes containing $v$ belong to the same graph
component.  The triple cubes cover $U$, so the sets $V_B$ are nonempty,
pairwise disjoint, and cover $U\setminus K$.  Since
\[
 |U\setminus K|
 =\left(1+|E|+\binom{|E|}{2}+\binom{|E|}{3}\right)
  -\left(1+|E|+\binom{|E|}{2}-1\right)
 =1+\binom{|E|}{3},
\]
their orders give \eqref{eq:damp-exterior-component-total-balance}.

The definitions and the exterior-incidence edges give
$P_{\mathcal R}\cup Q_{\mathcal R}=U$ and
$P_{\mathcal R}\cap Q_{\mathcal R}=K$.  The selected facet cubes strongly
shatter exactly
\[
 \operatorname{Sh}(K)\cup\{S\}\cup
   \{T:T\in B,\ B\in\mathcal R\}
\]
on the $P_{\mathcal R}$ side, and exactly
\[
 \operatorname{Sh}(K)\cup
   \{T:T\in B,\ B\notin\mathcal R\}
\]
on the other side.  Here all triples above $S$ lie in $B_S\subseteq
\mathcal R$, so the full $S$-carrier belongs to $P_{\mathcal R}$, whereas
$Q_{\mathcal R}$ does not shatter $S$.  Therefore
\[
 |P_{\mathcal R}|-|\operatorname{SSh}(P_{\mathcal R})|
 =\sum_{B\in\mathcal R}\omega(B).
\]
The corresponding difference for $Q_{\mathcal R}$ is the negative of this
sum by \eqref{eq:damp-exterior-component-total-balance}.  Each difference is
nonnegative by the lower Sandwich inequality.  Thus either one vanishes if
and only if both vanish, and Proposition~\ref{prop:damp-facet-cube-reconstruction}
gives the stated ampleness equivalence.  The horizontal and vertical strongly
shattered supports of the lift have total order $|U|+|K|$.  This is also the
order of the lift, so equality in the lower Sandwich inequality proves that
the lift is ample.

Conversely, consider any ample two-layer lift with contraction $U$ and
overlap $K$.  Lemma~\ref{lem:damp-ample-two-layer-shadow-intersection}
shows that its sections are ample and that exactly one of them shatters
$S$.  Every triple support is shattered by at least one section, so its
unique triple cube in $U$ lies in that section; it cannot lie in both,
because then its support would be shattered by $K$.  Triple cubes meeting
at a vertex outside $K$ must lie in the same section, and the section that
shatters $S$ contains every triple cube above $S$.  After interchanging
the two layers if necessary, its triple supports are therefore the union
of a set $\mathcal R$ of graph components with $B_S\in\mathcal R$.  Hence
its sections are exactly $P_{\mathcal R}$ and $Q_{\mathcal R}$.

The maximal cubes of the lift are precisely the horizontal triple cubes and
the products of the new-coordinate edge with the maximal pair cubes of $K$.
Thus the two copies of $v$ have maximal-cube incidence multiplicities
$\nu_K(v)+\rho_{\mathcal R}(v)$ and
$\nu_K(v)+\rho_{\overline{\mathcal R}}(v)$; multiplicity zero means that the
copy is absent.  Proposition~\ref{prop:damp-facet-cube-reconstruction} now
proves \eqref{eq:damp-exterior-component-no-isolated-incidence}.

At a corner of $U$, exactly one of the two $\rho$-terms is one and the other
is zero.  The two inequalities therefore force $\nu_K(v)\ge2$.  At a corner
of $K$, one has $\nu_K(v)=1$, so both $\rho$-terms must be positive.  Finally,
these two conditions are also sufficient.  Indeed, a vertex outside $K$
that is not a corner of $U$ belongs to at least two triple cubes, all on
the same side, while a vertex in $K$ that is not a corner of $K$ belongs
to at least two maximal pair cubes.  Thus every vertex not covered by the
two conditions already satisfies
\eqref{eq:damp-exterior-component-no-isolated-incidence}.  If
$\mathcal B$ has one component, one section is $U$ and the other is $K$.
The latter has a corner by the VC-two corner-peeling theorem, and its copy
in the lift belongs only to the corresponding vertical maximal cube.
\end{proof}

\begin{proposition}[Exterior-component normal form for an arbitrary overlap]
\label{prop:damp-two-layer-general-exterior-component-normal-form}
Let $U\subseteq Q_E$ be maximum of VC dimension three, let $C_T$ be its
unique triple cube of support $T\in\binom E3$, and let $K\subseteq U$ be a
nonempty ample family.  Put $\Sigma=\operatorname{Sh}(K)$ and
\[
 \mathcal N=\binom E3\setminus\Sigma,
 \qquad
 \mathcal D=
 \{A\subseteq E:|A|\le2\text{ and }A\notin\Sigma\}.
\]
Form a graph on $\mathcal N$ by joining two triple supports when their cubes
meet at a vertex outside $K$, and, for every $A\in\mathcal D$, joining all
triple supports that contain $A$.  Let $\mathcal B$ be the set of connected
components.  For $B\in\mathcal B$, define
\[
 V_B=\left(\bigcup_{T\in B}C_T\right)\setminus K,
 \qquad
 \mathcal D_B=\{A\in\mathcal D:A\subseteq T\text{ for some }T\in B\},
\]
and
\begin{equation}
 \omega_K(B)=|V_B|-|B|-|\mathcal D_B|.
 \label{eq:damp-general-exterior-component-weight}
\end{equation}
Then the $V_B$ partition $U\setminus K$, the $\mathcal D_B$ partition
$\mathcal D$, and every component is balanced individually:
\begin{equation}
 \omega_K(B)=0\quad(B\in\mathcal B).
 \label{eq:damp-general-exterior-component-balance}
\end{equation}

For $\mathcal R\subseteq\mathcal B$, put
\[
 P_{\mathcal R}=K\cup\bigcup_{B\in\mathcal R}\bigcup_{T\in B}C_T,
 \qquad
 Q_{\mathcal R}=K\cup\bigcup_{B\notin\mathcal R}\bigcup_{T\in B}C_T.
\]
Then $P_{\mathcal R}$ and $Q_{\mathcal R}$ are both ample for every
$\mathcal R\subseteq\mathcal B$, and the corresponding two-layer family is
ample.  Every ample two-layer lift with contraction $U$ and overlap $K$
arises from such a component cut, up to interchanging the layers.

For $v\in U$, let $\nu_K(v)$ be the number of maximal cubes of $K$ through
$v$, and let $\rho_{\mathcal R}(v)$ and
$\rho_{\overline{\mathcal R}}(v)$ count the noncommon triple cubes through
$v$ on the two sides of the component cut.  The corresponding lift is
cornerless if and only if
\begin{equation}
 \nu_K(v)+\rho_{\mathcal R}(v)\ne1,
 \qquad
 \nu_K(v)+\rho_{\overline{\mathcal R}}(v)\ne1
 \quad(v\in U).
 \label{eq:damp-general-exterior-no-isolated-incidence}
\end{equation}
Equivalently, every corner of $U$ belongs to at least two maximal cubes of
$K$, and, for every corner $v$ of $K$, the component set
\[
 \Gamma_K(v)=
 \{B\in\mathcal B:v\in C_T\text{ for some }T\in B\}
\]
meets both sides of the cut.
\end{proposition}

\begin{proof}
Every support in $\mathcal D$ lies below at least one triple support.  All
triples above it belong to $\mathcal N$, by downward closure of $\Sigma$,
and the second kind of graph edge puts them in one component.  Thus the
$\mathcal D_B$ partition $\mathcal D$.  Every vertex of the maximum class
$U$ lies in a triple cube.  If a vertex outside $K$ belongs to cubes with
two noncommon supports, the first kind of graph edge joins those supports.
It follows that the sets $V_B$ partition $U\setminus K$.

Both $U$ and $K$ are ample.  Since the supports missing from $\Sigma$ are
exactly $\mathcal N\mathbin{\dot\cup}\mathcal D$, the two partitions first
give
\[
 \sum_{B\in\mathcal B}|V_B|
 =|U|-|K|
 =|\operatorname{Sh}(U)|-|\Sigma|
 =|\mathcal N|+|\mathcal D|
 =\sum_{B\in\mathcal B}(|B|+|\mathcal D_B|),
\]
and hence $\sum_B\omega_K(B)=0$.

Fix one component $B$.  The families $P_{\{B\}}$ and $Q_{\{B\}}$ strongly
shatter, respectively, at least the supports
\[
 \Sigma\mathbin{\dot\cup}B\mathbin{\dot\cup}\mathcal D_B
 \quad\hbox{and}\quad
 \Sigma\mathbin{\dot\cup}(\mathcal N\setminus B)
 \mathbin{\dot\cup}(\mathcal D\setminus\mathcal D_B).
\]
Their vertex sets are $K\mathbin{\dot\cup}V_B$ and
$K\mathbin{\dot\cup}\bigcup_{B'\ne B}V_{B'}$.  The lower Sandwich
inequality therefore gives
\[
 \omega_K(B)\ge0
 \quad\hbox{and}\quad
 -\omega_K(B)\ge0.
\]
Thus \eqref{eq:damp-general-exterior-component-balance} holds component by
component.  This two-sided argument is the point at which ampleness of the
overlap eliminates every possible cancellation between exterior
components.

The exterior-incidence edges give
$P_{\mathcal R}\cup Q_{\mathcal R}=U$ and
$P_{\mathcal R}\cap Q_{\mathcal R}=K$.  The positioned cubes contained in
$P_{\mathcal R}$ strongly shatter at least
\[
 \Sigma
 \cup\bigcup_{B\in\mathcal R}B
 \cup\bigcup_{B\in\mathcal R}\mathcal D_B,
\]
and the complementary family of supports is strongly shattered by
$Q_{\mathcal R}$.  Componentwise balance gives
\[
 |P_{\mathcal R}|
 =|\Sigma|+\sum_{B\in\mathcal R}(|B|+|\mathcal D_B|),
\]
and the complementary equality holds for $Q_{\mathcal R}$.  Equality in
the lower Sandwich inequality now proves that both families are ample.
The corresponding two-layer family has order $|U|+|K|$ and strongly
shatters the disjoint support families
\[
 \operatorname{Sh}(U)
 \quad\hbox{and}\quad
 \{A\cup\{e_0\}:A\in\Sigma\},
\]
where $e_0$ is the new coordinate.  These have total order $|U|+|K|$,
so the lower Sandwich equality proves that the lift is ample as well.

Conversely, let two ample sections have union $U$ and intersection $K$.
The exact shadow intersection and union identities put every triple in
$\Sigma$ in both sections and every triple in $\mathcal N$ in exactly one.
Every missing lower support $A\in\mathcal D$ and its entire upper star have
one owner.  Triple cubes meeting outside $K$ also have one owner.  Thus the
owner classes are unions of the components in $\mathcal B$, and the two
sections are the displayed $P_{\mathcal R},Q_{\mathcal R}$.

Finally, the maximal cubes of the lift are the horizontal noncommon triple
cubes and the products of the new-coordinate edge with the maximal cubes of
$K$.  Hence the two copies of $v$ have maximal-cube incidence multiplicities
shown in \eqref{eq:damp-general-exterior-no-isolated-incidence}; multiplicity
zero means that the copy is absent.  At a corner of $U$ the condition reduces
to $\nu_K(v)\ge2$, and at a corner of $K$ it requires a noncommon triple from
each side.  Every remaining present vertex already has incidence at least
two, proving the last equivalence.
\end{proof}

\begin{corollary}[Property-B criterion for arbitrary overlaps]
\label{cor:damp-two-layer-general-property-b}
In the setting of
Proposition~\ref{prop:damp-two-layer-general-exterior-component-normal-form},
form the hypergraph $\mathcal H_K$ on $\mathcal B$ whose hyperedges are the
sets $\Gamma_K(v)$ over the corners $v$ of $K$.  There exists a cornerless
ample two-layer lift with contraction $U$ and overlap $K$ if and only if
\begin{enumerate}
 \item every corner of $U$ belongs to at least two maximal cubes of $K$;
       and
 \item $\mathcal H_K$ has Property B, that is, its vertices admit a
       two-colouring with no monochromatic hyperedge.
\end{enumerate}
No weight or section-ampleness condition remains to be checked.  If the
lift is a minimum-cardinality ample non-corner-peelable family, then $K$ is
corner-peelable and has a corner; consequently the colouring uses at least
two exterior components.
\end{corollary}

\begin{proof}
Every component cut gives ample sections by the proposition.  Its exact
relative-corner criterion says that the two copies of a corner of $K$ are
shielded precisely when its hyperedge meets both colours, while a corner of
$U$ is shielded precisely when it lies in at least two maximal cubes of
$K$.  This proves the equivalence.  In the minimum-cardinality case, $K$ is
an ample family of smaller order.  It must therefore be corner-peelable,
since otherwise it would itself be a smaller ample non-corner-peelable
family.  Its first vertex in a corner peeling is a corner, so
$\mathcal H_K$ has a nonempty hyperedge; Property B then forces both colours.
\end{proof}

\begin{proposition}[All-rank corner localization in Hall projections]
\label{prop:damp-hall-common-triple-corner-localization}
Let $U$ be an eleven-coordinate contraction of Hall's maximum class, let
$K\subseteq U$ be a non-singleton ample family, and let $v$ be a corner of
$K$.  Write $C$ for its unique maximal cube and $A$ for the support of
$C$; thus $1\le |A|\le3$.  Form a graph $L_U(A,v)$ on the triple supports
$T$ for which $v\in C_T$, omitting $T=A$ when $|A|=3$.  Join $T,T'$ when,
for some $e\notin A$, both cubes contain the neighbour
$v\mathbin\triangle\{e\}$.  Then
\begin{equation}
 |\Gamma_K(v)|\le \operatorname{cc}(L_U(A,v))
 \le
 \begin{cases}
  1,&|A|=1,\\
  2,&|A|\in\{2,3\}.
 \end{cases}
 \label{eq:damp-common-triple-corner-localization}
\end{equation}
Consequently a cornerless two-layer lift has no corner of $K$ supported by
a maximal edge.  Every corner supported by a maximal square or a common
triple cube is one of the listed two-component incidences, and its two local
components have opposite owners.  In particular, every hyperedge of
$\mathcal H_K$ has order two, and the owner constraints form a bipartite
graph on the exterior components.

Across all twelve contractions, the audit covers $14{,}784$ positioned
edge--vertex incidences, $26{,}400$ square--vertex incidences, and $15{,}840$
triple-cube--vertex incidences.  The respective numbers with two local
components are $0$, $480$, and $184$; no local graph has more than the bound
in \eqref{eq:damp-common-triple-corner-localization}.  In the coordinate-zero
contraction there are forty binary square incidences and fifteen binary
triple incidences.
\end{proposition}

\begin{proof}[Proof with a finite local audit]
Because $C$ is the unique maximal cube of $K$ through $v$, every neighbour
$v\mathbin\triangle\{e\}\in U$ with $e\notin A$ lies outside $K$.
Therefore all triple cubes containing such a neighbour belong to one
exterior component.  Each component of $L_U(A,v)$ is consequently contained
in one member of $\Gamma_K(v)$, which proves the first inequality in
\eqref{eq:damp-common-triple-corner-localization}.  The arbitrary-overlap
normal form, Proposition~\ref{prop:damp-two-layer-general-exterior-component-normal-form},
then shows that cornerlessness requires at least two members of
$\Gamma_K(v)$.

The remaining statement is a fixed local computation, independent of the
order of $K$.  The verifier constructs every positioned edge, square, and
triple cube in every Hall projection, forms $L_U(A,v)$ for every incidence,
and computes its connected components directly.  It records the complete
distributions and exceptional support lists in
\path{damp_hall_all_corner_localization.json}.  The mathematical
digest is
\[
 \texttt{c32cc3bbdc908fe80c5886762d64f3c58b2e69d9dea58a21a730fdeeb9760ddf}.
\]
The rank-one bound rules out maximal-edge corners.  In ranks two and three,
Property B rules out local component count at most one and forces the two
components at every remaining corner to have opposite colours.  Hence the
corner hypergraph is a graph, and its Property-B colourings are exactly its
bipartite two-colourings, unique up to reversal on each connected component.
\end{proof}

\begin{proposition}[Incidence separators and signed holonomy]
\label{prop:damp-hall-incidence-separator-holonomy}
Let $U$ be a Hall projection on coordinate set $E$.  Form the bipartite
incidence graph $\mathfrak I_U$ with one class $\binom E3$ and the other
class
\[
 U\mathbin{\dot\cup}\binom E1\mathbin{\dot\cup}\binom E2.
\]
A triple support $T$ is adjacent to $x\in U$ when $x\in C_T$, and it is
adjacent to a singleton or pair $A$ when $A\subseteq T$.  For an ample
overlap $K$ with shadow $\Sigma$, delete from $\mathfrak I_U$
\begin{equation}
 \{T\in\tbinom E3:T\in\Sigma\}\ \mathbin{\dot\cup}\ K\
 \mathbin{\dot\cup}\
 \{A\in\tbinom E1\cup\tbinom E2:A\in\Sigma\}.
 \label{eq:damp-hall-incidence-deletion-set}
\end{equation}
The connected components of the surviving triple nodes are exactly the
exterior components $\mathcal B$ of
Proposition~\ref{prop:damp-two-layer-general-exterior-component-normal-form}.
Moreover, the deletion set in
\eqref{eq:damp-hall-incidence-deletion-set} has order
\begin{equation}
 |K|+|\Sigma\setminus\{\varnothing\}|=2|K|-1.
 \label{eq:damp-hall-incidence-deletion-order}
\end{equation}

Suppose now that every corner of $K$ is one of the two-component incidences
in Proposition~\ref{prop:damp-hall-common-triple-corner-localization}.  Give
all paths in the residual incidence graph label zero, and, for each corner,
add one label-one edge between representatives of its two local classes.
Then the corner-component graph is bipartite if and only if this signed graph
is balanced, that is, every cycle contains an even number of label-one
edges.  Equivalently, the affine equations
\[
 z_T+z_{T'}=0\quad\hbox{on residual incidence paths},
 \qquad
 z_T+z_{T'}=1\quad\hbox{at overlap corners}
 \]
have a solution over $\mathbb F_2$.

For the coordinate-zero Hall projection, the minimum unit vertex-separator
widths in $\mathfrak I_U$ between the two local classes of the fifty-five
possible binary corners have distribution
\begin{equation}
 13^{(10)},\qquad 14^{(43)},\qquad 21^{(2)}.
 \label{eq:damp-hall-binary-separator-widths}
\end{equation}
Here terminal triple nodes are undeletable, while every other triple,
projection-vertex, singleton, and pair node has unit capacity.
If one additionally imposes the cube-closure condition that the eight
vertices of a terminal triple cube cannot all be deleted while its terminal
triple node survives, the distribution becomes
\begin{equation}
 13^{(10)},\qquad 14^{(42)},\qquad 16^{(1)},\qquad 22^{(2)}.
 \label{eq:damp-hall-binary-closed-separator-widths}
\end{equation}
\end{proposition}

\begin{proof}
A surviving projection-vertex node joins precisely the noncommon triple
cubes that meet at a vertex outside $K$.  A surviving singleton or pair node
joins precisely the noncommon upper star of a support missing from $\Sigma$.
Taking the two-section of the residual bipartite graph therefore gives the
two kinds of edges in the definition of $\mathcal B$, proving the first
claim.  The three parts of the deletion set are disjoint and together contain
one node for every member of $K$ and one node for every nonempty member of
$\Sigma$.  Since $K$ is ample, $|K|=|\Sigma|$, which proves
\eqref{eq:damp-hall-incidence-deletion-order}.

A zero-labelled connected component must have one owner.  A corner asks its
two incident components to have opposite owners, and hence contributes the
displayed affine equation of label one.  Such an affine system on a graph is
solvable exactly when the sum of its edge labels on every cycle is zero:
necessity follows by summing the vertex equations around the cycle, and
sufficiency follows by propagating a value from one root in each connected
component.  This proves the signed-holonomy criterion.

The last statement is a fixed max-flow/min-cut computation in the ambient
incidence graph and is independent of $K$ and of its order.  The verifier
splits every nonterminal node into an in-node and an out-node of capacity one,
computes a maximum flow for each of the fifty-five terminal pairs, and checks
the corresponding cut directly.  It imposes terminal-cube closure by
branching on a vertex retained from every terminal cube completed by a
candidate cut.  The complete cuts are recorded in
\path{damp_hall_binary_separator_cuts.json}; their mathematical digest is
\[
 \texttt{d3daab1128456dc68c5f356f93f696c1618275373cc7a963f74dada147195827}.
\]
\end{proof}

\begin{remark}[The remaining Hall separator inequality]
Fix the coordinate-zero Hall contraction.  Componentwise balance makes
section ampleness automatic, and the proposition reduces every overlap
corner to one of fifty-five binary local incidences.  Thus the common-triple
branch is the following single problem: can the deletion set
\eqref{eq:damp-hall-incidence-deletion-set} of an ample overlap $K$ of order
at most $58$ double-shield both contraction corners and leave a balanced
signed incidence graph?  By Menger's theorem, each binary corner requires
one of the separator packages in
\eqref{eq:damp-hall-binary-separator-widths}.  The missing step is an
isoperimetric transfer from these ambient blocker packages to the coupled
identity $|K|=|\Sigma|$: prove that double shielding and balanced holonomy
force $|\Sigma|\ge59$.  This formulation ranges over all overlap orders at
once and contains neither section variables nor a cardinality-by-cardinality
census.
\end{remark}

\begin{proposition}[A $298$-concept ample obstruction]
\label{prop:damp-hall-order-298-obstruction}
There is an ample forbidden pc-minor for corner peelability with \(298\)
concepts.  Consequently,
\[
 m_{\mathrm{amp}}\leq298.
\]
\end{proposition}

\begin{proof}
Use the coordinate-zero contraction \(U\) of Hall's class.  There is an
ample overlap \(K\) of order \(66\) whose shadow is the full rank-two
shadow except for the pair \(\{6,9\}\).  Its exterior graph has two
components.  The distinguished component contains \(95\) triple cubes and
\(96\) exterior vertices, while the other contains \(70\) triple cubes and
\(70\) exterior vertices.  Thus both component weights are zero.  Assigning
the two components to opposite sections gives section orders
\[
 |P|=66+96=162,\qquad |Q|=66+70=136.
\]
The two contraction corners \(1600\) and \(1894\) each lie in three maximal
pair cubes of \(K\).  The three corners of \(K\), namely
\[
 412,\qquad1025,\qquad2039,
\]
are each incident with both exterior components.  The arbitrary-overlap
Property-B criterion therefore shows that
\[
 Y=(P\times\{0\})\cup(Q\times\{1\})
\]
is ample and cornerless, and \(|Y|=162+136=298\).

The exact section lists are recorded in
\texttt{damp\_hall\_overlap\_component\_c0\_cnf.json}; their mathematical
digest is
\[
 \texttt{c9a32cff523db600553d42ef6a288d986be49bd2e6aa7df78d50625d03b575d3}.
\]
The independent verifier
\texttt{verify\_damp\_hall\_order298\_obstruction.py} reconstructs \(Y\)
from those lists, recomputes its shattered sets and corners, and constructs
a complete corner peeling for each of its \(24\) restrictions and \(12\)
contractions.  Its mathematical digest is
\[
 \texttt{bfa75ccdb16a493fa3fc805fb5516d80f3a060edaa14051b57935f927d472772}.
\]
Every proper pc-minor factors through one of these elementary minors, and
corner peelability is pc-minor closed.  Hence every proper pc-minor of
\(Y\) is corner-peelable, whereas \(Y\) itself has no first corner.  Thus
\(Y\) is an ample forbidden pc-minor.
\end{proof}

\begin{corollary}[Owner-optimality of the $298$-concept lift]
\label{cor:damp-hall-order-298-owner-optimal}
Fix the coordinate-zero Hall contraction and the ownership of its maximal
triple cubes used in Proposition~\ref{prop:damp-hall-order-298-obstruction}.
Every cornerless ample two-layer lift realizing this ownership has at least
$298$ concepts.
\end{corollary}

\begin{proof}
For this ownership, the bichromatic core has $59$ vertices.  Adjoining the
two contraction corners gives a forced core $F$ of order $61$, and a direct
shadow calculation gives
\[
 |\operatorname{Sh}(F)|=66.
\]
The contraction has order $232$, so
Corollary~\ref{cor:damp-two-layer-forced-core-shadow} gives
\[
 |Y|\ge232+66=298.
\]
The verifier \path{verify_damp_hall_bichromatic_core.py} performs the shadow
calculation directly from the owner list; its mathematical digest is
\[
 \texttt{81fa19a62fb7e653593621f72048b59e961336f7be360c39371e6a19d8f2bf29}.
\]
Proposition~\ref{prop:damp-hall-order-298-obstruction} attains equality.
\end{proof}

\begin{proposition}[A $291$-concept ample obstruction]
\label{prop:damp-hall-order-291-obstruction}
There is an ample forbidden pc-minor for corner peelability with \(291\)
concepts.  Consequently,
\[
 m_{\mathrm{amp}}\leq291.
\]
\end{proposition}

\begin{proof}
Again let \(U\) be the coordinate-zero contraction of Hall's class.  There
is an ample overlap \(K\subseteq U\) of order \(59\) whose shadow contains
every singleton and all pairs except
\[
\begin{split}
 &\{3,7\},\ \{4,7\},\ \{3,8\},\ \{4,8\},\\
 &\{5,9\},\ \{6,9\},\ \{5,10\},\ \{6,10\}.
\end{split}
\]
Thus the missing-pair graph is the disjoint union of the two four-cycles
\(K_{\{3,4\},\{7,8\}}\) and \(K_{\{5,6\},\{9,10\}}\).  Assigning \(139\)
of the \(165\) maximal three-cubes of \(U\) to the first section gives
sections \(P,Q\) of orders
\[
 |P|=206,\qquad |Q|=85.
\]
Their intersection is \(K\), both sections are ample, and their union is
\(U\).  The two corners of \(U\) belong to \(K\), while the three corners
\(412,1025,2039\) of \(K\) all belong to the bichromatic core.  The exact
relative-corner criterion therefore shows that the lift
\[
 Y=(P\times\{0\})\cup(Q\times\{1\})
\]
is ample and cornerless, with \(|Y|=206+85=291\).

The overlap, owner set, and reconstructed sections are recorded in
\path{damp_hall_carrier_defect_min8.json}.  The generating formula uses one
variable for the missing status of each pair and one choice among the
eleven vertices of its Hall carrier path.  Hence its single threshold
constraint covers all full-singleton, square-pure overlaps in this carrier
model with at least eight missing pairs; it is not an order-by-order
enumeration.  The stored witness is checked directly by
\path{solve_damp_hall_carrier_defect_cnf.py}; its mathematical digest is
\[
 \texttt{f2ea815743a1d96b8ac3cdc00efb0ab6dd80777483dd38ebe91c972df4184ffa}.
\]

The two four-cycles are coupled rather than independent.  The complete
\(3^{11}\)-conditioning audit compares the mixed cores of the certified
orders \(298,297,295,291\), whose orders are respectively
\(59,56,52,42\).  All four have exactly the same unique
conditioning-minimal nonample trace.  After fixing
\[
 x_0=1,\qquad x_3=x_4=0,\qquad x_9=x_{10}=1,
\]
and leaving \(\{1,2,5,6,7,8\}\) free, this trace is the antipodal pair
\(\{0^6,1^6\}\).  Thus one coherent rank-six antipodal bridge persists while
the missing-pair graph grows from one edge to two edges, then four, then
eight.  In the last stage it spans the two defect cycles and carries zero
relative-corner charge.  The audit is recorded in
\path{damp_hall_mixed_core_minimal_nonample_traces.json}, with mathematical
digest
\[
 \texttt{b584060c83d61722762f0c9f785f7ff5a6d8ee0a62c99c60ca1407694c38e9e1}.
\]
The companion comparison
\path{analyze_damp_hall_bridge_defect_chain.py} verifies that the missing
pair sets are strictly nested, with orders \(1,2,4,8\), and have successive
shapes
\[
 K_2,\qquad P_3,\qquad P_3\mathbin{\dot\cup}P_3,\qquad
 C_4\mathbin{\dot\cup}C_4.
\]
The corresponding mixed-core orders are \(59,56,52,42\).  Its artifact
\path{damp_hall_bridge_preserving_defect_chain.json} has mathematical digest
\[
 \texttt{acbfb63e3dc1c42196847a47315a7b98e90c3a753836cfe3d6c9d67d35d1faa7}.
\]

The independent verifier
\path{verify_damp_hall_carrier_defect_witness.py} reconstructs \(Y\) from
the overlap and owner list, recomputes ampleness and cornerlessness, and
constructs a complete corner peeling for each of its \(24\) restrictions
and \(12\) contractions.  The verification artifact
\path{damp_hall_carrier_defect_min8_verification.json} has mathematical
digest
\[
 \texttt{681a0bc8ba1584ad7732783da3526ae5e6eb11271f7216caa1fe6ee4254b0f56}.
\]
Thus every elementary pc-minor of \(Y\) is corner-peelable.  Every proper
pc-minor factors through an elementary one, whereas \(Y\) itself has no
corner, so \(Y\) is an ample forbidden pc-minor.
\end{proof}

\begin{proposition}[Exact capacity of the rank-six Hall bridge]
\label{prop:damp-hall-rank-six-bridge-capacity}
Let \(U\) be the coordinate-zero Hall contraction, and let
\[
 Y=(P\times\{0\})\cup(Q\times\{1\})
\]
be a cornerless ample lift with overlap \(K=P\cap Q\).  Suppose that \(K\)
shatters all singletons, has VC dimension at most two, and is square-pure.
Let \(M\) be the bichromatic core, and condition on
\[
 x_0=1,\qquad x_3=x_4=0,\qquad x_9=x_{10}=1.
\]
If the resulting trace of \(M\) on the free coordinates
\(\{1,2,5,6,7,8\}\) is exactly \(\{0^6,1^6\}\), then \(K\) misses at most
eight pairs.  Consequently
\[
 |Y|\ge291.
\]
The bound is sharp.
\end{proposition}

\begin{proof}
Let \(m\) be the number of pairs not shattered by \(K\).  Ampleness and the
full singleton shadow give
\[
 |K|=1+11+\binom{11}{2}-m=67-m,
 \qquad |Y|=|U|+|K|=299-m.
\]
For completeness, square-purity also implies that every shattered pair has
exactly one maximal square.  Indeed, summing
\(|E(T_i(K))|=|V(T_i(K))|-1\) over the eleven singleton carrier trees and
using the cubical Euler identity gives that the number of maximal squares is
the number of shattered pairs.

The formula in \path{solve_damp_hall_carrier_defect_cnf.py} is therefore
exact for this class.  It chooses the unique positioned square for every
shattered pair and the owner of every maximal three-cube of \(U\).
Proposition~\ref{prop:damp-two-layer-carrier-one-cut} gives the one-transition
condition on every Hall pair carrier.  The three Sandwich equalities impose
ampleness of \(K,P,Q\), and
Proposition~\ref{prop:damp-two-layer-relative-corner-criterion} imposes
cornerlessness.  Finally, unit clauses prescribe the displayed trace of
\(M\).

With \(m\ge9\), the resulting CNF has \(15{,}032\) variables and \(80{,}469\)
clauses and is unsatisfiable.  CaDiCaL generated a \(440{,}994\)-line DRAT
proof.  The independent checker \texttt{drat-trim} returned
\texttt{VERIFIED}; its backward core uses \(9{,}696\) original clauses,
\(32{,}496\) lemmas, and \(1{,}792{,}976\) resolution steps.  The formula,
proof, checker hashes, and verification statistics are recorded in
\path{damp_hall_rank6_bridge_capacity_proof_verification.json}, whose
mathematical digest is
\[
 \texttt{db9ba53d639a8e1b87fa80e5d5f2133c51ebb08f99cedd224fcba96049cbd799}.
\]
Thus \(m\le8\), and the displayed order identity gives \(|Y|\ge291\).
The construction in
Proposition~\ref{prop:damp-hall-order-291-obstruction} has the prescribed
rank-six trace and \(m=8\), so equality is attained.
\end{proof}

\begin{corollary}[Capacity of the observed Hall bridge types]
\label{cor:damp-hall-observed-bridge-capacity}
In the full-singleton square-pure Hall carrier model, suppose that the mixed
core has one of the five exact conditioning-minimal antipodal traces appearing
in the forced-shadow-\(61\), forced-shadow-\(55\), forced-shadow-\(46\), or
certified order-\(298\) owner cores.  Then \(m\le8\), and hence
\[
 |Y|\ge291.
\]
\end{corollary}

\begin{proof}
Besides the common rank-six trace of
Proposition~\ref{prop:damp-hall-rank-six-bridge-capacity}, the complete
conditioning audit gives two adjacent rank-four traces for the
forced-shadow-\(61\) core, one rank-five trace for the forced-shadow-\(55\)
core, and one rank-four trace for the forced-shadow-\(46\) core.  Prescribing
each trace separately and imposing \(m\ge9\) gives four further
unsatisfiable carrier formulas.  Independent \texttt{drat-trim} checks
return \texttt{VERIFIED} for all four proof logs.  Together with the
rank-six check, the aggregate verification covers all five observed bridge
signatures.  It is recorded in
\path{damp_hall_observed_bridge_capacity_proof_verification.json}, with
mathematical digest
\[
 \texttt{a2174b2e8107e2912eae634dbc916381202fba62ea3561c89b66edf15d20a4f0}.
\]
\end{proof}

\begin{remark}
Corollary~\ref{cor:damp-hall-observed-bridge-capacity} is a finite
classification of the bridge signatures that occur in the present Hall
owner-core artifacts.  It does not assert that every possible Hall owner
colouring has one of these five traces.  The next theorem handles the
complementary owner colourings directly, without making such an assertion.
\end{remark}

\begin{theorem}[Exact optimum in the Hall carrier model]
\label{thm:damp-hall-full-carrier-optimum}
Let \(U\) be the coordinate-zero Hall contraction, and let
\[
 Y=(P\times\{0\})\cup(Q\times\{1\})
\]
be a cornerless ample lift.  If \(K=P\cap Q\) shatters all singletons, has
VC dimension at most two, and is square-pure, then
\[
 |Y|\ge291.
\]
This bound is sharp.
\end{theorem}

\begin{proof}
As in Proposition~\ref{prop:damp-hall-rank-six-bridge-capacity}, write \(m\)
for the number of missing pairs.  It is enough to exclude \(m\ge9\).
Proposition~\ref{prop:damp-two-layer-bichromatic-core} shows that the mixed
core \(M\) is nonample.

Partition the owner colourings according to whether \(M\) has one of the
five exact antipodal traces in
Corollary~\ref{cor:damp-hall-observed-bridge-capacity}.  The five independently
checked formulas cited there exclude \(m\ge9\) in the first part of the
partition.  In the residual part, compile the nonampleness of \(M\) by the
strict Sandwich inequality
\[
 |\operatorname{Sh}(M)|\ge |M|+1
\]
and forbid all five exact traces.  This residual formula has \(19{,}229\)
variables and \(150{,}552\) clauses and is unsatisfiable.  Lingeling
generated the DRAT proof, and an independent \texttt{drat-trim} check
returned \texttt{VERIFIED}.  Its backward core contains \(21{,}951\)
original clauses and \(131{,}531\) lemmas, using \(25{,}223{,}694\)
resolution steps.

The five profile certificates and the residual certificate form an
exhaustive logical partition of the exact carrier formula.  Their hashes,
checker statistics, and the sharp \(m=8\) witness are collected in
\path{damp_hall_full_carrier_capacity_proof_verification.json}, with
mathematical digest
\[
 \texttt{5758b4f9df00f6e627cdbc172e32ee15e0e812b735e2d616f1b67bb6d53b7993}.
\]
Thus \(m\le8\), and \(|Y|=299-m\ge291\).  Sharpness follows from
Proposition~\ref{prop:damp-hall-order-291-obstruction}.
\end{proof}

\begin{lemma}[Cubical normal form in VC dimension two]
\label{lem:damp-vc-two-cubical-normal-form}
Let \(K\) be a nonempty ample family of VC dimension at most two.  Every
shattered pair supports exactly one square of \(K\).  Moreover, \(K\) has a
maximal edge of support \(\{i\}\) if and only if \(\{i\}\) is shattered and
\(i\) is isolated in the graph of shattered pairs; in that case the maximal
\(i\)-edge is unique.  Finally, \(K\) has a maximal vertex if and only if
\(|K|=1\).
\end{lemma}

\begin{proof}
Fix a shattered pair \(S\).  Apply
Proposition~\ref{prop:damp-maximum-vc-three-carrier-trees} to \(K\) and
\(S\).  Since \(K\) shatters no triple, the set \(L_S\) in that proposition
is empty.  The \(S\)-carrier therefore consists of one vertex, which is
exactly one square of support \(S\).

Now fix a shattered singleton \(\{i\}\), and form the \(i\)-carrier by
contracting each full \(i\)-edge to a vertex.  The argument in the proof of
Proposition~\ref{prop:damp-maximum-vc-three-carrier-trees}, with \(S=\{i\}\),
shows that this carrier is an ample VC-dimension-one family: its shattered
coordinates are precisely the \(j\) for which \(\{i,j\}\) is shattered by
\(K\).  Hence its one-inclusion graph is a tree, with one vertex for each
\(i\)-edge of \(K\) and one edge for each square whose support contains
\(i\).

An \(i\)-edge of \(K\) is maximal precisely when the corresponding carrier
vertex is isolated.  A tree has an isolated vertex precisely when it consists
of that single vertex.  Thus a maximal \(i\)-edge exists precisely when no
shattered pair contains \(i\), and in that case it is unique.  Finally, the
one-inclusion graph of an ample family is connected.  A maximal zero-cube is
an isolated vertex of that graph, so it can occur precisely when \(K\) is a
singleton.
\end{proof}

\begin{theorem}[Exact Hall optimum for every VC-two overlap]
\label{thm:damp-hall-vc-two-optimum}
Let \(U\) be the coordinate-zero Hall contraction and let

\[
 Y=(P\times\{0\})\cup(Q\times\{1\})
\]

be a cornerless ample lift.  If \(K=P\cap Q\) has VC dimension at most two,
then

\[
 |Y|\ge291.
\]

This bound is sharp.
\end{theorem}

\begin{proof}[Computer-assisted proof]
Theorem~\ref{thm:damp-hall-full-carrier-optimum} handles the case in which
\(K\) shatters every singleton and every maximal cube of \(K\) is a square.
By Lemma~\ref{lem:damp-vc-two-cubical-normal-form}, every other VC-two overlap
either misses a singleton or has an isolated shattered coordinate with one
maximal edge.
These alternatives form an exact structural partition, not a partition by
order.

The all-profile formula in
\path{solve_damp_hall_overlap65_cnf.py} chooses \(K\), every triple owner,
and both sections simultaneously.  It imposes the three Sandwich equalities,
the exact relative-corner criterion, the nonampleness of the bichromatic
core, and the cubical normal form of
Lemma~\ref{lem:damp-vc-two-cubical-normal-form}.  Requiring one of the two
exceptional alternatives and \(|K|\le58\), equivalently \(|Y|\le290\), gives
a formula with \(26{,}406\) variables and \(250{,}636\) clauses.  It is
unsatisfiable.  Lingeling generated a DRAT proof, and the independent checker
\texttt{drat-trim} returned \texttt{VERIFIED}.  Its backward core uses
\(13{,}448\) original clauses and \(3{,}102\) lemmas, with \(262{,}558\)
resolution steps.

The formula, proof, hashes, and checker statistics are recorded in
\path{damp_hall_vc2_normal_form_capacity_proof_verification.json}, whose
mathematical digest is

\[
 \texttt{6d0ef7268aa38ca55a2a19ce16f9f551f95420cb050fe3a39e3547649ff6f0b0}.
\]

Thus no exceptional VC-two overlap has order at most \(58\).  Together with
Theorem~\ref{thm:damp-hall-full-carrier-optimum}, this gives
\(|K|\ge59\), and hence \(|Y|=|U|+|K|\ge232+59=291\).  The construction in
Proposition~\ref{prop:damp-hall-order-291-obstruction} attains equality.
\end{proof}

\begin{remark}
Theorem~\ref{thm:damp-hall-vc-two-optimum} closes every Hall lift whose
overlap has VC dimension at most two, including missing-singleton and
non-pure overlaps.  A smaller lift over this Hall contraction must therefore
have a common triple cube, equivalently an overlap of VC dimension three.
The theorem does not normalize an arbitrary least obstruction to a Hall
two-layer lift.
\end{remark}

The remaining VC-three overlaps are excluded by a direct search over all
two-layer lifts, using the facet-cube description of ample families.  The
same search applies to many further bases near Hall's class.

\begin{theorem}[Hall lifts with arbitrary overlaps; computer-assisted]
\label{thm:damp-hall-all-overlaps-optimum}
Let \(B\) be one of the following ample families:
\begin{enumerate}
\item the contraction of Hall's class \(C_H\) at any one of its twelve
      coordinates;
\item the contraction of \(C_H\) at any two of its coordinates;
\item the contraction, at any one coordinate, of any of the twenty-seven
      cornerless order-\(299\) classes of Remark~\ref{rem:damp-hall-neighborhood}.
\end{enumerate}
If \(Y=(P\times\{0\})\cup(Q\times\{1\})\) is a cornerless ample family with
\(P\cup Q=B\), then \(|Y|\geq291\).  For the coordinate-zero contraction of
\(C_H\), the bound is attained by
Proposition~\ref{prop:damp-hall-order-291-obstruction}.
\end{theorem}

\begin{proof}[Computer-assisted proof]
Every listed base uses all of its coordinates and has a corner; these finite
properties are checked directly by
\texttt{scripts/verify\_hall\_lift\_audit\_inputs.py}.
Since the projection of \(Y\) is \(B\), all base coordinates remain active
in \(Y\).  The added coordinate is active as well: otherwise \(Y\) is a
copy of \(B\) in a single layer and inherits a corner, contrary to the
hypothesis.  Thus forcing every singleton into the encoded shattering
complex loses no possible cornerless lift.

Every vertex of such a lift \(Y\) has the form \((x,0)\) or \((x,1)\) with
\(x\in B\).  By Lemma~\ref{lem:damp-ample-facet-cubes}, \(Y\) is the union
of one cube per facet of \(\operatorname{Sh}(Y)\), every such cube lies
inside \(B\times\{0,1\}\), and cornerlessness says that every vertex of
\(Y\) lies in at least two of these cubes.  The formula has a variable for
each coordinate set, saying that it belongs to a down-closed complex
\(\Sigma\).  It has one cube-choice variable for each facet of \(\Sigma\)
and each cube inside \(B\times\{0,1\}\), and a covering variable for each of
the \(2|B|\) candidate vertices.  It requires that the union shatters no set
outside \(\Sigma\), that every covered vertex is covered twice, that for
every \(x\in B\) at least one of \((x,0),(x,1)\) is covered, and that at
most \(290\) vertices are covered.  By the converse half of
Lemma~\ref{lem:damp-ample-facet-cubes}, its models are exactly the cornerless
ample lifts of \(B\) of order at most \(290\).

The \(12+66+324=402\) labeled cases, comprising \(272\) distinct formulas,
are generated by
\texttt{explore\_damp\_lift\_cornerless\_sat.py} with
\texttt{--max-order=290}.  For every one of them, CaDiCaL~2.1.2 reports
\textsc{unsat}, and its DRAT proof is accepted by \texttt{drat-trim}.  The same
formula for the coordinate-zero contraction with bound \(291\) is
satisfiable.  Its decoded model is a cornerless ample family of order \(291\),
which recovers the equality case independently.  The fresh replay retains all \(272\) formulas and checked proofs, together
with input maps, stage logs and tool identities, in the split archive
\path{evidence/hall_lift_optimum/replay-20260924/}.
Its manifest records the part and archive hashes; the atlas reconciliation
records every payload hash and the coverage of all \(402\) labeled bases.
\end{proof}

\begin{remark}[Scope of the lift bound]
\label{rem:damp-hall-lift-scope}
Theorem~\ref{thm:damp-hall-all-overlaps-optimum} removes the overlap hypothesis of
Theorem~\ref{thm:damp-hall-vc-two-optimum} and extends the bound from one
base to \(402\) bases.  It is still a statement about two-layer lifts of
these bases.  Cornerless ample families below order \(291\) exist
(Proposition~\ref{prop:damp-order-267-obstruction} below), but none of them is a
one-coordinate lift of any of these bases.  Each single contraction of
\(C_H\) has exactly two corners.  Its double contractions have two, three,
or four corners, and exactly six of them have two: those over the
coordinate pairs \(\{0,1\},\{2,3\},\ldots,\{10,11\}\), which form a perfect
matching of the coordinates.  In the other
direction, Proposition~\ref{prop:damp-seven-coordinate-ample-corners}
requires at least eight coordinates for any cornerless ample family.
\end{remark}

The same step can be iterated: contract a cornerless family at one
coordinate and search for a cornerless ample lift of least order of the
contraction.  Starting from the \(291\)-concept lift of the
coordinate-zero contraction, three further steps, each contracting at
coordinate zero, give cornerless ample families of orders \(283\), \(275\) and
\(267\).  None of these lifts contains the previous family.

\begin{proposition}[A \(267\)-concept ample obstruction; computer-assisted]
\label{prop:damp-order-267-obstruction}
There is an ample forbidden pc-minor for corner peelability with \(267\)
concepts on twelve coordinates, of VC dimension three.  Consequently
\[
 m_{\mathrm{amp}}\leq267.
\]
\end{proposition}

\begin{proof}[Computer-assisted proof]
The family \(Y_{267}\subseteq Q_{12}\) is listed in
\path{evidence/order267_obstruction/record.json}, with SHA-256 digest
\[
 \texttt{45eebc6ef09bf97ffa814295b621bb09d9788f87258031ffc25c5181f879573c}.
\]
It was found by the search just described, but its properties are checked
directly, without a SAT solver, by
\path{../ample-corners/scripts/verify_ample_corners_data.py}.  First,
\(|\operatorname{Sh}(Y_{267})|=|\operatorname{SSh}(Y_{267})|=|Y_{267}|=267\),
so \(Y_{267}\) is ample by the Sandwich equality, and all twelve coordinates
are active.  Second, no vertex lies in a unique maximal cube, so
\(Y_{267}\notin\Damp\).  Third, each of the \(36\) elementary minors
\(\rho_e^0Y_{267}\), \(\rho_e^1Y_{267}\) and \(\pi_eY_{267}\), with inactive
coordinates removed, is ample and is emptied by repeatedly deleting a
corner, so it lies in \(\Damp\).  The contractions all have \(208\)
concepts, and the restrictions have between \(77\) and \(190\).  By the
recognition criterion \eqref{eq:damp-introduction-recognition},
\(Y_{267}\) is an ample forbidden pc-minor.
\end{proof}

\begin{corollary}[Minimality of the recorded seed under ample inclusion]
\label{cor:damp-267-ample-inclusion-minimality}
Every proper ample subfamily of $Y_{267}$ belongs to $\Damp$.
Moreover every maximal sequence of corner deletions from such a subfamily
empties it. Thus no smaller ample obstruction can be obtained from this
labelled seed by deleting vertices alone.
\end{corollary}
\begin{proof}
The companion ample-corners manuscript proves
\texttt{thm:record267-ample-inclusion-minimality} by an exact membership
encoding over all subsets of $Y_{267}$ and an independently checked DRAT
refutation. It excludes every nonempty proper ample cornerless subfamily;
the formula does not fix rows or require the original shattered supports
to survive. Any nonempty proper ample subfamily therefore has a corner.
Deleting it preserves ampleness and proper containment, so induction
empties the family under any choices. The formula, compressed proof,
coverage argument, replay commands and hashes are retained in
\path{../ample-corners/evidence/seed267_all_subfamilies/}.
\end{proof}

This inclusion statement is separate from positivity of the coordinate
minors. It also rules out cornered negative subfamilies, because a maximal
unsuccessful corner peeling would leave a smaller ample cornerless one.
It is a property of this seed, not a general property of forbidden
pc-minors: the $268$-word example below is forbidden and contains
$Y_{267}$ as a proper ample negative subfamily.

\begin{remark}[Deleting a corner of the contraction before lifting]
\label{rem:damp-267-corner-deleted-relifts}
There is also a certified exclusion allowing new words in the lifted
class. For every coordinate $e$ and every corner $c$ of $\pi_eY_{267}$,
put $B=(\pi_eY_{267})\setminus\{c\}$. No cornerless ample class
$Y\subseteq B\times\{0,1\}$ has last-coordinate contraction exactly $B$.
This has no order or VC-dimension bound. Each contraction has two corners;
the twenty-four choices give twelve bases up to explicit translations.
The companion ample-corners manuscript proves
\texttt{prop:record-corner-deleted-relifts} by exact membership formulas
with twelve independently checked DRAT refutations. All formulas, proofs
and coverage maps are retained in
\path{../ample-corners/evidence/corner_deleted_relifts_unbounded/}.
This closes the operation of deleting one corner from a contraction and
then replacing the entire lift. It does not classify further deletions,
other contraction bases, or cornered negative lifts, and does not change
the interval $64\le m_{\mathrm{amp}}\le267$.
Local corner-cube compatibility alone does not explain these exclusions:
all twelve bases admit compatible ample local lifts at their corners.
For one base an isometric $272$-word cornerless lift satisfies every
ampleness equality on supports of size at most three, but has three
shattered four-supports that are not strongly shattered. The companion paper
records the explicit nonample diagnostic and its direct checks; the
higher-support conditions cannot be discarded in that lift formulation.
The companion proposition \texttt{prop:native272-cycle-fiber} further
localizes this diagnostic to a six-cycle minor with five singleton trace
fibers and one $17$-word fiber. Every ample subfamily must empty a trace
fiber. Every ample subfamily of the entire $272$-word family is positive.
The fiber-empty case has a checked refutation on its $255$-word envelope.
For a cornerless ample subfamily meeting the fiber, split into five disjoint
cases according to the first missing singleton word. All five exact
membership formulas have independently replayed DRAT refutations, archived
in \path{../ample-corners/evidence/native272_remote_split/}.
Every nonempty ample subfamily therefore has a corner, and induction
supplies a full peeling. This closes arbitrary deletion searches in this
fixed envelope; it is not a global exclusion below $267$.
\end{remark}

\begin{corollary}[Every sink is attainable in every proper minor of the
267-concept obstruction; computer-assisted]
\label{cor:damp-267-all-proper-minor-sinks}
For every nonempty proper pc-minor $M$ of $Y_{267}$ and every $v\in M$,
there is an acyclic USO of $M$ with sink $v$. Consequently no rooted
factor occurs as a pc-minor of $Y_{267}$, with any choice of root.
\end{corollary}
\begin{proof}[Computer-assisted proof]
The certificate \path{research/explore_seed267_endpoints.json} lists,
for each of the $36$ elementary minors $N$ of the preceding proposition
and each $u\in N$, a complete corner order ending at $u$.
There are exactly
\[
 12\cdot208+12\cdot267=5700
\]
such orders: every contraction has $208$ words, while the two sections
in each coordinate partition the $267$ words of the original family.
The verifier \path{research/explore_seed267_endpoints.py} checks every
order as a permutation of its prescribed minor. At each word, it computes
the directions to later unit neighbours and checks that their entire
cube occurs later in the order. Thus every deletion is a corner deletion.
It also verifies that the resulting out-sets are distinct and satisfy
nonclashing for every pair of words. The fixed input is identified by
the SHA-256 digest in Proposition~\ref{prop:damp-order-267-obstruction};
the report records that digest and all orders. No failed greedy search
is used as evidence: all $5700$ tests have positive witnesses.

Now write a nonempty proper minor $M$ as a sequence beginning with an
elementary minor $N$. For a prescribed $v\in M$, choose a preimage
$u\in N$ surviving every restriction in the sequence. The certificate
provides a corner order of $N$ ending at $u$. Rooted minor transport
restricts and contracts that order to one of $M$ ending at $v$.
CCMW's corner-peeling/acyclic-USO correspondence gives the claimed
orientation. A rooted factor has an unattainable prescribed sink,
so it cannot be a proper minor. The original $Y_{267}$ cannot be a
rooted factor either, since it is not ordinarily corner-peelable.
\end{proof}

This is stronger than ordinary positivity of all proper minors or exclusion
from one root-face assembly. It rules out any proposed exhaustive construction
in which every output retains a rooted factor as a pc-minor. It does not rule
out using rooted data in a construction that subsequently loses that minor,
and it does not classify the primitive inputs. In this example the failure
of an acyclic USO on the whole support persists despite complete sink
freedom on every proper minor; these independent minor maps have not been
asserted to be simultaneously compatible.

\begin{corollary}[A cornered obstruction with all proper-minor sinks attainable;
computer-assisted]
\label{cor:damp-268-all-proper-minor-sinks}
In the integer encoding of $Y_{267}$, put $X=Y_{267}\cup\{238\}$.
Then $X$ is an ample forbidden pc-minor on twelve coordinates, of VC
dimension three and minimum degree three. Its only corner is $238$.
Every vertex of every nonempty proper pc-minor is an attainable sink of
an acyclic USO. Every coordinate reduction belongs to $\Damp$.
In particular, $X$ has no rooted-factor pc-minor despite having a corner.
\end{corollary}
\begin{proof}[Computer-assisted proof]
The verifier \path{research/seed268_all_sinks.py}, with output in the
corresponding JSON file, uses the fixed family and endpoint certificates
of the preceding corollary. Coordinates are numbered from zero. The
incident directions of $a=238$ in $X$ are $I=\{1,4,9\}$, whose bit mask
is $530$. All seven words $a^J$ with $\varnothing\ne J\subseteq I$
belong to $Y_{267}$, and $I\notin\operatorname{Sh}(Y_{267})$. The
one-concept extension test therefore gives ampleness and
\[
 \operatorname{Sh}(X)=\operatorname{Sh}(Y_{267})\cup\{I\}.
\]
The verifier independently checks this equality, the minimum degree,
and that $a$ is the unique corner. Deleting $a$ leaves the cornerless
$Y_{267}$, so $X$ cannot peel. Proposition~\ref{prop:damp-one-concept-atlas}
now implies forbiddenness.

For each elementary minor operation $\mu$, set $B=\mu(Y_{267})$ and
$M=\mu(X)$. Either $M=B$, or $M=B\cup\{a_\mu\}$ with $a_\mu$ a
corner. Every old endpoint remains attainable: delete $a_\mu$ first,
when present, then use its order in $B$. This imports the $5700$
endpoint orders of Corollary~\ref{cor:damp-267-all-proper-minor-sinks}.
There are exactly $21$ new endpoints: one in the section containing $a$
for each coordinate, and one in each contraction outside $I$.
For every new endpoint the certificate supplies a complete order ending
there. The verifier checks its suffix cubes, distinct out-sets, and
pairwise nonclashing. Thus all $5721$ elementary endpoint choices have
positive witnesses. Rooted minor transport, as in the preceding proof,
gives the assertion for every longer proper-minor sequence and excludes
rooted-factor minors.

Finally a coordinate reduction of an ample VC-three family is ample
of VC dimension at most two: any cube in the reduction lifts to a cube
with one additional direction in $X$. Such families are corner-peelable
by CCMW \cite[Proposition~4.11]{ChalopiChepoiMoran22}. The certificate
also supplies independently checked orders for all twelve reductions.
\end{proof}

The conclusion refutes the implication that an ample forbidden class
without a rooted-factor pc-minor must be cornerless. It does not refute
generation from cornerless inputs by ample additions: this very example
is a one-concept extension of $Y_{267}$. Nor does it provide a general
corner deletion preserving forbiddenness or all proper-minor sink choices.

\begin{proposition}[Reducing all-sinks corner deletion to forbidden inputs]
\label{prop:damp-all-sinks-deletion-reduction}
The following assertions are equivalent.
\begin{enumerate}
\item If every vertex of a nonempty ample $H$ is an attainable sink,
then deleting any corner $c$ with $H\setminus\{c\}\ne\varnothing$
leaves a member of $\Damp$.
\item No ample forbidden pc-minor has an ample one-concept extension
in which every vertex is an attainable sink.
\end{enumerate}
These equivalent assertions are false by
Theorem~\ref{thm:damp-all-sinks-forbidden-repair}. They assert ordinary peelability
after deletion, not attainability of every remaining endpoint.
\end{proposition}
\begin{proof}
An ample one-concept extension adds a corner, so the first assertion
immediately implies the second.

Conversely suppose the first assertion fails, and put $S=H\setminus\{c\}$.
The family $S$ is ample and not corner-peelable. Choose a pc-minor sequence
$\mu$ taking $S$ to an ample forbidden minor $O$. Apply the same
restrictions and contractions to $H$, writing $M=\mu(H)$.
The all-sinks property is pc-minor closed: for any desired endpoint in
a nonempty minor, choose a surviving preimage in $H$ and transport an
order ending at that preimage. Thus every vertex of $M$ is attainable.
Each restriction or contraction preserves containment and can send the
one extra word to at most one extra image. Consequently
\[
 O\subseteq M,\qquad |M\setminus O|\le1.
\]
Equality $M=O$ is impossible, since $M$ is corner-peelable and $O$ is not.
Thus $M$ is an ample one-concept extension of $O$ with every sink attainable,
contradicting the second assertion. Constant coordinates cause no problem:
if a coordinate is constant on $O$ but not on $M$, its constant-side face
is $O$, contradicting pc-minor closure of $\Damp$. Hence both may be
viewed on the same irredundant coordinate set.
\end{proof}

\begin{proposition}[Corner deletion can destroy attainability of a sink;
computer-assisted]
\label{prop:damp-all-sinks-corner-failure}
There is an ample class $X$ on twelve coordinates with $289$ vertices
such that every vertex is an attainable sink, but deletion of a corner
leaves a $288$-vertex rooted factor with an unattainable sink.
In particular, the property that every vertex is an attainable sink is
not closed under corner deletion, even though it is pc-minor closed.
\end{proposition}
\begin{proof}[Computer-assisted proof]
Use $C,b,H$ from Proposition~\ref{prop:damp-forced-prefix-rooted-inputs},
and translate $H$ by $b=256$. Write $A=H\mathbin\triangle b$, where
translation means bitwise exclusive-or on every word. Its root is now
zero, its unique corner, and its other $287$ vertices are attainable.
Put $a=34$ and $X=A\cup\{a\}$. The verifier
\path{research/all_sinks_corner_failure.py}, with certificate in the
corresponding JSON file, checks
\[
 \operatorname{Sh}(X)=\operatorname{Sh}(A)\cup\{\{5,6,7,8\}\},
 \qquad I_X(a)=\{5,6,7,8\}.
\]
Thus $X$ is ample and $a$ is a corner of degree four.
For every $r\in A\setminus\{0\}$, delete $a$ first and follow the
archived order of $A$ ending at $r$. This supplies $287$ endpoint orders.
The certificate supplies two additional complete orders ending at $0$
and $a$. All $289$ orders are checked through their suffix cubes,
distinct out-sets, and pairwise nonclashing, so every vertex of $X$
is an attainable sink of an acyclic USO.

Deleting $a$ leaves $A$, whose only corner is zero. Every corner
order of $A$ must therefore start at zero and cannot end there, since
$|A|>1$. The other $287$ endpoints remain attainable. The input proposition
already proves the stronger rooted-factor minimality of $(A,0)$;
no new minor census is needed.
\end{proof}

This counterexample concerns preservation of all prescribed endpoints.
This example alone does not refute the ordinary-positivity assertion in
Proposition~\ref{prop:damp-all-sinks-deletion-reduction}: its deleted
class $A$ still peels. The next theorem supplies that stronger refutation.
Nor does it refute corner descent for forbidden parents with all proper-minor
sinks attainable. It rules out proving that descent by an unrestricted
all-sinks corner-deletion closure theorem.

\begin{theorem}[An all-sinks one-concept repair of a forbidden class]
\label{thm:damp-all-sinks-forbidden-repair}
Let $(A,0)$ be an irredundant rooted factor on $n$ coordinates. Suppose
$a,b\notin A$ are distinct and
\[
 Z=A\cup\{a\},\qquad W=A\cup\{b\}
\]
are ample, every vertex of each is an attainable sink, $Z\cup W$ is
ample, and $\operatorname{Sh}(Z)\cap\operatorname{Sh}(W)
=\operatorname{Sh}(A)$. On two disjoint copies of the old coordinates
and one common coordinate $e$, put
\[
 C=Z\times Q_{\{e\}},\qquad
 L=(Z\times\{0_e\})\cup(W\times\{1_e\}),\qquad
 X=C\cup_D L,
\]
where $D=\{(0,0_e),(0,1_e)\}$ and $C,L$ use the two different old
coordinate blocks. Let $c=(a,1_e)$ in the $C$ block and put $S=X-c$.
Then every vertex of $X$ is an attainable sink, while $S$ is an ample
forbidden pc-minor. Their orders are $4|A|+2$ and $4|A|+1$, and their
isometric dimension is $2n+1$. The vertex $c$ is a corner of $X$.
The reduction $S^{\{e\}}$ is the forbidden vertex gluing of two copies
of $(A,0)$.

The hypotheses hold for the normalized $288$-vertex factor of
Proposition~\ref{prop:damp-all-sinks-corner-failure}, with $a=34$ and
$b=860$. Consequently an all-sinks ample class of order $1154$ has a
corner deletion that is forbidden of order $1153$, on $25$ coordinates.
\end{theorem}
\begin{proof}
The extension criterion makes the added vertices corners of $Z,W$.
In particular, both sections peel to $A$. The shadow assumptions give
ampleness of $L$ by Proposition~\ref{prop:damp-edge-boundary-transfer}.
Products and edge gluing give ampleness of $X$. The product corner
$(a,1_e)$ of $C$ lies outside $D$, so it remains a corner in $X$.
Thus $S$ is ample as well. Write
\[
 T=C-c=(Z\times\{0_e\})\cup(A\times\{1_e\}),
 \qquad S=T\cup_D L.
\]
Both $T$ and $L$ peel: $T$ is a peripheral expansion with positive
base $Z$ and positive site $A$, while $L$ has positive endpoint orders
from the edge-boundary transfer proposition.

We first prove the stronger all-sinks assertion for $X$. Every endpoint
of $L$ is attainable. To retain a desired vertex in one section, delete
the other section's added tip, delete its whole copy of $A$ using an
ordinary order with all opposite mates present, and finish with the
prescribed order in the retained section $Z$ or $W$.
Also $C$ peels to $D$: use a root-ending order of $Z$, deleting both
copies of each nonroot word before proceeding to the next word.
Therefore any endpoint in $L$ is attained in $X$ by clearing $C-D$
and then using its order in $L$.

For a desired endpoint in layer $j$ of $C-D$, first clear the opposite
layer of $C$ down to its root using a root-ending order of $Z$. Next
follow an order of $L$ ending at $(0,j)$, omitting its last vertex.
At the other boundary vertex the only remaining $C$-edge is the common
edge, so these deletions are still valid in the union. Finally peel the
retained layer $Z$ to the desired endpoint. This handles every vertex
of $X$ and proves the all-sinks assertion.

Nonpeelability has a shorter certificate using the common reduction:
\[
 S^{\{e\}}=A\vee_0 A.
\]
The right-hand side is forbidden by
Theorem~\ref{thm:damp-cut-vertex-obstructions}. Reduction closure of
$\Damp$ (Proposition~\ref{prop:damp-strong-projection-closure}) therefore
implies $S\notin\Damp$.

We check every elementary pc-minor of $S$. For an exclusive coordinate
in either block, restriction to the side omitting zero discards the
other piece and is a minor of the positive $T$ or $L$. For restriction
to the root side or contraction, the image $A'$ of $A$ peels to its
root by rooted minimality. Each changed section is $A'$ with at most
one extra word, which can be deleted first if present. These relative
orders are also transported from the original one-word wings. Clearing
the wings and then $A'\times Q_{\{e\}}$ retains $D$. Append any ordinary
order of the unchanged piece. Thus all exclusive-coordinate minors peel.

Restriction in $e$ gives either $Z\vee_0 Z$ or $A\vee_0 W$ on disjoint
blocks. Both peel by the vertex-gluing criterion: $Z,W$ attain zero,
while $A$ is ordinarily positive. Contraction in $e$ gives
$Z\vee_0(Z\cup W)$. The second factor peels to zero by deleting $b$
and using a root-ending order of $Z$; here $b$ is a corner because
both $Z$ and its one-word extension $Z\cup W$ are ample. This contraction
also peels. Every elementary minor is positive, so $S$ is forbidden.
The counts follow by subtracting the common edge, and its $e$-reduction
is $A\vee_0 A$, forbidden by the rooted vertex-gluing theorem.

For the finite input, \path{research/all_sinks_forbidden_repair.py/json}
reuses the established $24$ rooted minor orders of $A$ and all $289$
endpoint orders of $Z$. It verifies the shadow hypotheses for $W$ and
$Z\cup W$, and all $289$ endpoint orders of $W$: $287$ are inherited
from $A$, while the orders ending at $0$ and $860$ are supplied separately.
Each order is checked by its outgoing cubes, distinct out-sets, and
nonclashing. Representative assembled orders on $X$ are also checked.
The argument above covers all $1154$ endpoint choices and every proper
minor of $S$ symbolically; no exhaustive output-order search is claimed.
\end{proof}

This refutes ordinary corner-deletion closure even when every prescribed
sink is attainable before deletion. Both equivalent assertions in
Proposition~\ref{prop:damp-all-sinks-deletion-reduction} are therefore
false. It does not refute descent for a forbidden parent whose proper
minors all admit every sink: the parent $X$ here is positive, and $S$
has the negative reduction $A\vee_0 A$. The construction uses established
edge and vertex gluing; its additional conclusion is an all-sinks repair
of an ordinary forbidden class, not an exhaustive input generator.

\begin{theorem}[A corner-descent-terminal obstruction with all proper-minor sinks]
\label{thm:damp-all-sinks-corner-terminal}
There is an ample forbidden pc-minor $Y$ with $2308$ vertices on $26$
coordinates such that every vertex of every nonempty proper pc-minor is
an attainable sink, but no corner deletion preserves forbiddenness.
It has exactly eight corners. Consequently, even an ample forbidden
class with no rooted-factor pc-minor need not admit a sequence of
forbiddenness-preserving corner deletions to a cornerless class.
\end{theorem}
\begin{proof}
Take the rooted factor $(A,0)$ and its two all-sinks repairs $Z=A+a$,
$W=A+b$ from Theorem~\ref{thm:damp-all-sinks-forbidden-repair}.
Here $|A|=288$, $a=34$, $b=860$, and zero is the only corner of $A$.
Keep disjoint copies of its twelve old coordinates and introduce
coordinates $e,f$. First work without $f$, writing
\[
 L=(Z\times\{0_e\})\cup(W\times\{1_e\})
\]
on the second old block. On the first old block set
\[
 H_R=(R\times Q_{\{e\}})\cup_D L,
 \qquad D=\{(0,0_e),(0,1_e)\},
\]
for $R\in\{A,Z,W,U\}$, where $U=Z\cup W$.
Put $O=H_A$. Then
\[
 |O|=1152,\quad |H_Z|=|H_W|=1154,\quad
 H_Z\cap H_W=O,\quad H_Z\cup H_W=H_U.
\]
The order construction in the preceding theorem works with either
all-sinks product factor $Z$ or $W$, and shows that $H_Z,H_W$ attain
every sink. The union $U$ also attains every sink: delete the other
tip and use the endpoint order in $Z$ or $W$. Thus $H_U$ is ample
and all-sinks by the same construction.

We verify that $O$ is forbidden. Its reduction in $e$ is $A\vee_0 A$,
so it is nonpeelable. Each opposite-side restriction in an exclusive
coordinate is a minor of one of its positive pieces $A\times Q_1,L$.
A root-side restriction or contraction makes that piece retain $D$,
because the corresponding proper rooted minor of $A$ attains zero;
clearing its wings and then its core prism supplies the order. Append
an ordinary order of the other piece. The two $e$-sections are
$A\vee_0 Z$ and $A\vee_0 W$, and the contraction is $A\vee_0 U$.
All three peel by vertex gluing, since $Z,W,U$ attain zero and $A$
is ordinarily positive. This covers every elementary pc-minor of $O$.

Now set
\[
 Y=(H_Z\times\{0_f\})\cup(H_W\times\{1_f\}).
\]
Each wing consists of the two copies of its added tip along $e$.
Delete them successively to peel $H_Z,H_W$ to $O$. No edge joins
the two wings, since $a,b$ differ in more than one old coordinate.
The union $H_U$ is ample as observed above. Therefore
Corollary~\ref{cor:damp-two-wing-all-minor-sinks} proves that $Y$ is
forbidden and every proper-minor endpoint is attainable. Its order is
$2\cdot1154=2308$, and all $26$ coordinates are active.

For the corner and deletion assertions, write the same family symmetrically
as the union of two prisms over the two-section class $Z/W$, intersecting
in the full $(e,f)$-square. Its corners are exactly the four first-block
repair-tip vertices
\[
 (a,0_f,j_e),\ (b,1_f,j_e)\qquad(j=0,1)
\]
and the four second-block repair-tip vertices
\[
 (a,0_e,j_f),\ (b,1_e,j_f)\qquad(j=0,1).
\]
Every such tip is a product corner. To exclude other corners, a nonzero
old core word in either block has its mate in the coordinate separating
$Z$ and $W$. If it were a corner, its complete corner cube and that mate
would force its incident old-coordinate cube to lie in $A$, making it
a corner of $A$. The only such word is zero. At a zero word on the common
square, incident directions from both old blocks have a missing mixed
square, so that word is not a corner either.

Delete a first-block tip and fix $f$ to its layer. The resulting proper
face is $H_Z$ or $H_W$ with one copy of its product tip deleted.
Its $e$-reduction is exactly $A\vee_0 A$, so it is nonpeelable.
For a second-block tip, interchange the two blocks and $e,f$.
Thus every corner deletion leaves a negative proper face, and cannot
be forbidden. This proves all assertions.

The replay \path{research/all_sinks_corner_terminal.py/json} constructs
all $2308$ vertices, verifies exactly the eight listed corners, and checks
all eight negative-face reduction identities. It also checks representative
orders in $H_Z,H_W,H_U$ using the archived input endpoint orders.
All proper-minor sink choices follow from the uniform corollary, not from
an exhaustive endpoint search on the output.
\end{proof}

This refutes the remaining corner-descent proposal even with complete
sink freedom on every proper pc-minor. It also supplies a construction
from rooted inputs whose output retains no rooted-factor pc-minor; the
exclusion following Corollary~\ref{cor:damp-267-all-proper-minor-sinks}
therefore does not exclude this operation. The construction uses the
existing negative-reduction inverse theorem, strengthened to preserve
all proper-minor sink choices. It does not generate every admissible
rooted input or every terminal base. Here ``corner-descent terminal''
does not mean terminal under coordinate reductions: $Y^{\{f\}}=O$
and $Y^{\{e,f\}}=A\vee_0 A$ are both negative.

\begin{corollary}[An explicit two-parameter family without rooted-factor minors]
\label{cor:damp-corner-terminal-block-family}
For every $p,q\ge1$ there is an explicit ample forbidden class $Y_{p,q}$
with every proper-minor sink attainable and no forbiddenness-preserving
corner deletion. Its parameters are
\[
 |Y_{p,q}|=579\,2^{p+q}-2^{p+1}-2^{q+1},\qquad
 \operatorname{idim}(Y_{p,q})=24+p+q,
\]
and it has exactly $2p\,2^q+2q\,2^p$ corners. The case $p=q=1$ is
Theorem~\ref{thm:damp-all-sinks-corner-terminal}.
\end{corollary}
\begin{proof}
This is a specialization of the existing signed-facet partition generator,
with its sink property supplied by Corollary~\ref{cor:damp-block-all-sinks}.
To specify it without recursion, take the cornerless forbidden seed
$S=A\vee_0 A$ of order $575$ and its four tips $a_1,b_1,a_2,b_2$,
where $a=34,b=860$ and the subscript indicates the old coordinate block.
Each single tip repairs $S$, by vertex gluing, and their acquired supports
are pairwise distinct: within a block this follows from the verified
shadow intersection of $A+a,A+b$, and between blocks the nonempty supports
use disjoint coordinates.

Let $F,E$ be disjoint layer blocks of sizes $p,q$. Define four rows in
$Q_{F\cup E}$ by
\[
\begin{array}{ll}
 R_{a_1}=(Q_F\setminus\{1_F\})\times Q_E,&
 R_{b_1}=(Q_F\setminus\{0_F\})\times Q_E,\\
 R_{a_2}=Q_F\times(Q_E\setminus\{1_E\}),&
 R_{b_2}=Q_F\times(Q_E\setminus\{0_E\}),
\end{array}
\]
and set
\[
 Y_{p,q}=(S\times Q_{F\cup E})\cup
 \bigcup_{t\in\{a_1,b_1,a_2,b_2\}}(\{t\}\times R_t).
\]
These rows are unions of signed facets, with each signed facet owned
by exactly one row. Theorem~\ref{thm:damp-terminal-common-seed-facet-partitions}
therefore proves forbiddenness and corner-descent terminality. This
identification also shows that these two conclusions for $Y_{1,1}$
already follow from the established cover theorem; its additional feature
is complete proper-minor sink freedom.

Equivalently, $Y_{p,q}$ is obtained from $Y_{1,1}$ by replacing $f$
with $F$ and $e$ with $E$. The block formula replaces a singleton signed
facet row by precisely the punctured-block row displayed above.
Corollary~\ref{cor:damp-block-all-sinks} proves the proper-minor sink
assertion, with no new input or minor tests. Summing the four disjoint
tip rows and the core gives the stated size. The core has no corners;
a punctured $p$-cube has exactly $p$ corners, and every vertex of its
full $q$-cube factor is a corner. The row-corner criterion therefore
gives $2p2^q+2q2^p$ corners. All coordinates are active.
\end{proof}

The replay \path{research/all_sinks_block_transfer.py/json} checks
punctured-fibre controls through dimension four and the direct row formula
for $(p,q)=(2,1)$: $4620$ vertices, $16$ corners, and recovery of the
$2308$ support on constant blocks. The general claims are symbolic.
This is a nonrecursive family with fixed certified inputs, not a complete
generator of the forbidden basis. For larger blocks its repeated columns
compress to the existing example, and its full layer reduction remains
$S$. Thus it does not supply new reduction-terminal inputs.

\begin{remark}[Prescribed sinks and corner deletion]
\label{rem:damp-267-corner-deletion-test}
An attainable sink at a corner does not ensure that deleting that corner
preserves ordinary peelability. In the fixed encoding above,
$Z=Y_{267}\cup\{220\}$ is ample and has a corner order ending at $220$,
yet $Z\setminus\{220\}=Y_{267}$ does not peel. This does not refute
the stronger hypothesis requiring every vertex of $Z$ to be an attainable
sink: vertex $236$ is unattainable. The certificate
\path{research/seed267_corner_descent.json} supplies the positive order and
an exhaustive five-state deletion graph for this failed endpoint. Each
state retains $236$, lists every other corner as an allowed deletion,
and has all its successors among the listed states; no singleton occurs.
Induction on the state size proves the failure, independently of any
greedy search outcome.

The same verifier enumerates every word of $Q_{12}\setminus Y_{267}$
and applies the punctured-cube extension criterion. There are exactly
$35$ ample one-concept extensions, and every one has an unattainable sink,
certified by an exhaustive deletion graph with at most nine states.
For each of the twelve extensions with a stored positive order, the
vertex opposite the added corner in its corner cube is also certified
unattainable, with a five-state deletion graph.
The ordinary status of all $35$ extensions is determined in
Proposition~\ref{prop:damp-267-one-word-atlas} below.
Thus this fixed neighbourhood supplies no all-sinks one-word repair of
$Y_{267}$. It does not support a general closure conclusion: both endpoint
closure and ordinary positivity after deletion are refuted above using
different inputs. The stronger forbidden-parent descent is refuted by
Theorem~\ref{thm:damp-all-sinks-corner-terminal}.
\end{remark}

\begin{proposition}[Complete one-concept atlas of the fixed $267$-vertex seed;
computer-assisted]
\label{prop:damp-267-one-word-atlas}
In the encoding of Proposition~\ref{prop:damp-order-267-obstruction},
$Y_{267}$ has exactly $35$ ample one-concept extensions. Precisely the
following twelve added words give dismantlable classes:
\[
 220,370,551,749,1022,1456,1547,2353,2579,3268,3656,3968.
\]
The other $23$ extensions are ample forbidden pc-minors, and every
one-coordinate reduction of each admits every vertex as an acyclic-USO sink.
Within this fixed
atlas, the positive extensions have four corners; the negative extensions
have one or two.
\end{proposition}
\begin{proof}[Computer-assisted proof]
The punctured-face criterion of
Proposition~\ref{prop:damp-one-concept-atlas} enumerates the $35$ admissible
words among $Q_{12}\setminus Y_{267}$. The twelve positive orders are
already stored in \path{research/seed267_corner_descent.json}; the verifier
\path{research/seed267_one_word_atlas.py} replays every deletion.
It also checks the complete corner sets, not merely corners encountered
by those orders.

For $21$ of the other additions, the added word is the only corner.
Deleting it leaves the original cornerless seed, so no peeling exists.
For addition $597$ the corners are precisely $597,707$, and for addition
$1872$ they are precisely $1872,256$. In either case, deleting the new
word returns $Y_{267}$, while deleting the indicated old word leaves
another cornerless $267$-vertex family. The verifier checks all these
cornerless children directly. These exhaust the possible first moves
and prove nonpeelability without a heuristic search. The one-concept
atlas theorem then supplies forbidden-minor minimality.

Finally, a coordinate reduction of an ample one-concept extension is
unchanged or gains one concept. Every old reduction admits every prescribed
sink, since $Y_{267}$ has VC-dimension three and its nonempty reductions
have VC-dimension at most two; use the VC-two endpoint result recalled
in Remark~\ref{rem:damp-one-corner-peelable}. If a new concept survives,
it is a corner of the ample reduced extension. Deleting it first therefore
preserves attainability of every old endpoint.

For the $21$ forbidden additions of support size three, the parent still
has VC-dimension three, so the same VC-two theorem settles every endpoint
of every reduction. The remaining two additions, $597$ and $1872$, each
have acquired support of size four. Only the four reductions in that
support change. In each of these eight cases the changed reduction has
$60$ concepts, and the sole new endpoint has an explicit corner order in
\path{research/inherited_rooted_reduction_probe.json}. The companion
verifier checks all eight orders through their final prescribed vertices.
Reductions outside the support are unchanged. This covers every endpoint
of every one-coordinate reduction of the $23$ negative extensions. The
complete input and corner certificates remain in
\path{research/seed267_one_word_atlas.json}.
\end{proof}

Each of the twelve positive extensions is an inclusion-minimal positive
completion of this seed. This does not assert that every such completion
has one added word. In contrast to the canonical completion theorem for
$\mathcal R_{\rm per}$, none of the $24$ one-sided coordinate saturations
of this seed is inclusion-minimal as a Damp completion: each contains
exactly two of the listed one-word repairs and adds at least $18$ words.
This containment statement is checked in the same certificate. The
classification above concerns the fixed embedded one-word atlas; it makes
no claim of new isomorphism types or exhaustive primitive-input generation.

\begin{proposition}[All simultaneous batches from the initial $267$-vertex atlas]
\label{prop:damp-267-batch-atlas}
Let $H=Y_{267}$, let $U$ be its $35$ ample one-concept tips from
Proposition~\ref{prop:damp-267-one-word-atlas}, and let $P\subset U$
be the twelve positive tips listed there. For any $T\subseteq U$,
group its tips by their acquired support over $H$, and project each
group outside that support, as in
Theorem~\ref{thm:damp-equal-support-extensions}. Then:
\begin{enumerate}
 \item $H\cup T$ is ample exactly when every projected group is ample.
 Each group has at most seven words.
 \item For every such ample batch,
 \[
 H\cup T\in\Damp\quad\Longleftrightarrow\quad T\cap P\ne\varnothing.
 \]
 \item If $T\cap P=\varnothing$ and $H\cup T$ is ample, then
 $H\cup T$ is a forbidden pc-minor and all its one-coordinate
 reductions belong to $\Damp$.
\end{enumerate}
Consequently, among positive completions contained in $H\cup U$,
the inclusion-minimal ones are exactly the twelve singleton repairs.
\end{proposition}
\begin{proof}[Proof with a finite certificate]
The general theorem proves item~1. The fixed atlas has $22$ acquired
supports: $18$ singleton groups and four further groups of sizes
$2,4,4,7$. Every positive tip lies in a singleton group. These facts
are checked from the one-word atlas by
\path{research/seed267_batch_atlas.py}.

Every nonempty ample row of size at most seven has VC-dimension at
most two: shattering three coordinates requires at least eight words.
It therefore belongs to $\Damp$ by CCMW
\cite[Proposition~4.11]{ChalopiChepoiMoran22}. The general theorem
supplies a relative peeling of every ample batch to $H$. If $T$
contains $p\in P$, keep its singleton group and peel every other
group. This gives a relative peeling to the positive class
$H\cup\{p\}$, proving the positive implication in item~2.

For the converse, the same verifier constructs a set $K\subseteq H$
of $257$ words. For every $x\in K$ it records a coordinate set $J_x$
with
\[
 x^{\{e\}}\in K\quad(e\in J_x),\qquad
 x^{J_x}\notin H\cup(U\setminus P).
\]
All these finite membership checks are replayed directly. Thus in any
family between $H$ and $H\cup(U\setminus P)$, the first vertex of
$K$ to be deleted would still have all its indicated neighbours and
would lack the opposite vertex of their cube. It could not be a
corner. No such family admits a complete peeling, irrespective of
whether the full envelope is ample. This is the protected-missing-cube
argument of Proposition~\ref{prop:damp-protected-uso-cycle}.

For item~3, apply an elementary pc-minor operation $\mu$ in any of
the twelve coordinates to the relative ample chain from $H\cup T$
to $H$. Consecutive images are ample and either identical or differ
by one vertex. Omitting repetitions gives a relative corner peeling
to $\mu H$, which is positive because $H$ is forbidden. Hence every
elementary proper pc-minor is positive, and pc-minor closure handles
all proper pc-minors. The protected set proves that the parent is
negative. The same chain argument applies to each one-coordinate
reduction, since each reduction is ample and changes by at most one
vertex at a deletion. Its final reduction of $H$ has VC-dimension at
most two and is positive. Thus every reduction of the batch is positive.

Finally, every positive completion in this envelope contains one of
the twelve positive singleton completions by item~2. Each singleton
completion is minimal because $H$ itself is negative.
\end{proof}

This is a factorized, nonrecursive selection rule for a fixed envelope.
Its admissibility tests concern only the projected groups, and its
forbidden outputs are selected simply by excluding the twelve resolving
tips. The certificate lists every allowed subset of each group and its
local peeling, as well as the protected set and all its missing-cube
witnesses. Multiplying the independent group counts gives $356\,352$
forbidden embedded outputs, including $H$; this is not an isomorphism
count. No enumeration of all $2^{35}$ batches is used. All these outputs
peel relatively to the already known seed. Completions using vertices
outside the initial atlas, and exhaustive primitive-input generation,
remain unclassified.

We next separate incompatibility between proper face orientations from
cycles that appear only after compatible orientations have been glued.
The coordinate-face restriction theorem applies here; arbitrary ample
subpieces cannot be substituted, as
Remark~\ref{rem:damp-nonconvex-peripheral-cover} shows.

\begin{proposition}[Gluing USOs from all coordinate faces]
\label{prop:damp-all-face-uso-gluing}
Let $\varnothing\ne X\subseteq Q_F$ be ample, with $n=|F|\ge2$ and
$|X|\le2^n-2$. Give each nonempty elementary coordinate face
$X\cap\{x_i=b\}$ a USO. Suppose these orientations agree on every
edge belonging to two of the faces. Then their union is a USO of $X$.
If every supplied face orientation is acyclic, every directed cycle
in the union uses every coordinate of $F$, and hence has length at
least $2n$.

Let $\mathcal A(X)$ denote the finite set of edge-consistent choices
of acyclic USOs on these faces. Then $X\in\Damp$ if and only if
some member of $\mathcal A(X)$ has acyclic union. In particular, when
all faces are positive, the two distinct possibilities for nonpeelability are
\begin{enumerate}
 \item $\mathcal A(X)=\varnothing$;
 \item $\mathcal A(X)\ne\varnothing$, but every member has a directed
 cycle using every coordinate.
\end{enumerate}
\end{proposition}
\begin{proof}
Every edge belongs to an elementary coordinate face because $n\ge2$.
Agreement therefore defines one orientation on all edges of $X$.
Every contained cube is proper in $Q_F$, since $X\ne Q_F$, and hence
lies in an elementary face. It has a unique sink by the supplied USO
there. This proves the contained-cube axiom (C2).

For the outgoing-cube axiom (C1), let $I$ be the outgoing directions
at $x\in X$. If $I\subsetneq F$, choose $j\in F\setminus I$.
The face fixing $x_j$ has exactly the same outgoing directions at $x$,
so its USO supplies the required $I$-cube in $X$.
If instead $I=F$, then for every $j\in F$ the face fixing $x_j$
supplies the full cube at $x$ on $F\setminus\{j\}$. These cubes
together contain every vertex of $Q_F$ except possibly $x^F$.
This would give $|X|\ge2^n-1$, a contradiction. Thus (C1) holds,
and CCMW's USO characterization proves the first assertion.

A directed cycle omitting a coordinate lies entirely in a face fixing
that coordinate. This contradicts the assumed acyclicity of the supplied
face orientation. Each coordinate occurs a positive even number of times
on a closed cube walk, proving the length bound. Finally, an acyclic USO
of $X$ restricts to an acyclic USO on every coordinate face by CCMW
Lemma~6.9, and its restrictions agree on common edges. Conversely,
an edge-consistent choice with acyclic union gives an acyclic USO of $X$
by the first assertion. CCMW's peeling correspondence proves the
positivity equivalence and the two alternatives.
\end{proof}

The cardinality hypothesis is essential for the unconditional USO-gluing
assertion. On $Q_3\setminus\{111\}$ orient the three edges at $000$
outward, and orient the remaining six edges as the cycle
\[
 001\longrightarrow011\longrightarrow010\longrightarrow110
 \longrightarrow100\longrightarrow101\longrightarrow001.
\]
Here binary words are written with coordinate zero on the right.
Every elementary face has an acyclic USO, and their edge orientations
agree. The union fails (C1) at $000$, whose outgoing three-cube is
missing $111$. The check \path{research/proper_face_uso_gluing.py/json}
verifies this control and exhaustively checks the proposition on all
$118$ nonempty ample subfamilies of $Q_3$ of order at most six:
$2784$ edge assignments include $1192$ choices with acyclic face USOs.
The general statement uses the preceding proof, not this enumeration.

Every ample forbidden pc-minor meets the cardinality hypothesis: a full
cube is positive, and a singly punctured cube is positive by induction
using its nested full and singly punctured sections. Consequently the
two alternatives distinguish all ample forbidden supports whose proper
face data are being considered. They do not themselves generate these
supports. If all elementary contractions are positive as well as all
faces, either alternative certifies forbiddenness by pc-minor closure;
selecting all such inputs intrinsically remains open.

Both alternatives occur in the established examples below.
Proposition~\ref{prop:damp-267-minor-acyclic-uso} supplies a nonempty
$\mathcal A(Y_{267})$; cornerlessness makes every compatible union
cyclic. In contrast, the $598$-vertex class of
Proposition~\ref{prop:damp-proper-face-source-conflict} has
$\mathcal A(X)=\varnothing$: a compatible choice would glue to a USO,
while its restrictions to the stated proper factor faces would be acyclic,
contradicting that proposition. Each factor face is contained in an
elementary face, so this use of restriction is legitimate.
Thus compatible proper-face maps cannot be required of every forbidden
input, nor can their existence be treated as a positivity certificate.
The all-coordinate cycle observation for the recorded $267$ map is reused;
the added statement is the general gluing criterion and this explicit
distinction between the two failures.

\begin{proposition}[A cyclic representation map with acyclic proper induced
minors; computer-assisted]
\label{prop:damp-267-minor-acyclic-uso}
The ample forbidden class $Y_{267}$ admits a CCMW representation map
whose orientation is cyclic but whose induced orientation under every
nonempty proper sequence of coordinate faces, contractions, and reductions
is acyclic. One such map has a unique source strongly connected component
of order $40$; deleting it leaves an acyclic representation map on $227$
vertices. Every directed cycle of the parent map uses all twelve coordinate
directions. The certificate includes a directed cycle of length $28$.
\end{proposition}
\begin{proof}[Computer-assisted proof]
The map is listed in
\path{research/explore_seed267_minor_acyclic_uso.json}. The solver-free
verifier \path{research/verify_seed267_uso.py} first checks that its domain
is exactly the fixed family of
Proposition~\ref{prop:damp-order-267-obstruction}. It enumerates contained
cubes, checks that the out-map is a bijection onto their $267$ distinct
supports, and verifies
\[
 (x\mathbin\triangle y)\cap(r(x)\mathbin\triangle r(y))
 \ne\varnothing\qquad(x\ne y),
\]
where words are identified with their sets of one-coordinates. It also
checks containment of every outgoing cube. These are CCMW's
representation-map and USO conditions; the cited dictionary supplies
the orientation.

For each of the twelve coordinates the verifier constructs four induced
maps: the two coordinate-face maps, the contraction map using the sink
of each fibre, and the reduction map using its source. It checks each
map by the same global nonclashing and cube conditions, then produces
a topological order and independently checks every resulting corner
deletion. All $48$ maps are acyclic. Any proper sequence begins with
one of these operations, and acyclic minor transport proves acyclicity
under all subsequent operations. This last step is symbolic and does
not rely on enumerating longer sequences.

The parent map has the stored directed $28$-cycle. Direct strong-component
computation identifies its unique source component of size $40$; the
restricted map on the other $227$ words is checked to be a representation
map and acyclic. This also agrees with CCMW's source-component deletion
result. Finally a directed cycle omitting a coordinate direction would
lie in a coordinate face, contradicting the verified acyclicity of both
face maps in that direction. Thus every directed cycle uses all twelve
directions, and in particular has length at least $24$.
\end{proof}

This supplies compatible positive oriented minors on an actual forbidden
support, strengthening the earlier positive-cube control and the Hall-map
example with positive reductions. It does not assert that an arbitrary
tuple of positive maps on this support extends: the earlier incompatible
tuple remains incompatible. Nor is the existence of this particular cyclic
map the proof of forbiddenness; cornerlessness excludes every acyclic USO,
independently of the displayed map. In particular, having a representation
map whose proper induced minors are acyclic is insufficient to characterize
positive ample families. A general structural generator of the negative
supports with this behavior remains open.

\begin{remark}[The acyclic-proper-map property excludes small separators]
\label{rem:damp-minor-acyclic-separator-scope}
The map property in Proposition~\ref{prop:damp-267-minor-acyclic-uso}
cannot be imposed on all forbidden inputs, even those with only positive
coordinate reductions. More precisely, if an ample forbidden class $X$
has a cut vertex or a separating edge, every USO of $X$, if one exists,
is cyclic on some proper coordinate face. In particular some elementary
coordinate-face map is cyclic.

Indeed, the canonical pieces at a cut vertex or separating edge are
proper convex coordinate restrictions; see
Theorem~\ref{thm:damp-cut-vertex-obstructions} and
Proposition~\ref{prop:damp-crossing-separator-recovery}. A parent USO
restricts to mutually agreeing USOs on these pieces. If all pieces were
acyclic, their union would be acyclic by the singleton and edge cases
of Proposition~\ref{prop:damp-uso-cube-gluing}, contradicting forbiddenness.
Thus one piece contains a directed cycle. That proper coordinate face
fixes at least one active coordinate, so the cycle also occurs in an
elementary face map. This argument uses the existing acyclic gluing law;
it is not a new obstruction to the existence of representation maps.

The $1237$-word example of Theorem~\ref{thm:damp-three-piece-witness}
has a separating edge and all coordinate reductions positive. It therefore
refutes a proposed extension of the acyclic-proper-map property to every
input terminal for negative-reduction descent. A necessary condition for
a forbidden support to admit this property is consequently that its
crossing graph be two-connected. The square gluing example of
Remark~\ref{rem:damp-square-gluing-cycle} explains why the acyclic-union
argument does not extend automatically to larger separating cubes.

The cornered $627$-word terminal example in
\path{research/protected_core_delay.tex} passes this necessary test:
it has $50$ corners and its crossing graph is two-connected. The bounded
search recorded in \path{research/explore_terminal627_minor_acyclic_uso.json}
examined $60$ candidate maps and added $4004$ cycle exclusions, without
finding a map having all proper induced maps acyclic. It reached its
refinement limit; this is neither an unsatisfiability certificate nor a
negative answer. In particular, two-connectivity is only a necessary
condition here, and no sufficiency claim is made.
\end{remark}

\begin{proposition}[Forced sources in proper faces]
\label{prop:damp-proper-face-source-conflict}
Let $X$ be ample and let $A_1,\ldots,A_k$ be proper coordinate faces
containing the same vertex $v$. Suppose $v$ is the unique corner of each
$A_i$. Let $D_i$ be the incident directions of $v$ in $A_i$.
If the cube at $v$ on $\bigcup_iD_i$ is not contained in $X$, no USO
of $X$ has acyclic restrictions to all these faces.

There is a $598$-vertex ample forbidden pc-minor with this property,
all coordinate reductions positive, and a two-connected crossing graph.
Thus the absence of cut vertices and separating edges does not suffice
for the acyclic-proper-map property of
Proposition~\ref{prop:damp-267-minor-acyclic-uso}.
\end{proposition}
\begin{proof}
Every finite acyclic orientation has a source. In a CCMW USO a source
is a corner, since its full incident cube is its outgoing cube.
Consequently an acyclic USO on $A_i$ must have source $v$ and direct
every edge in $D_i$ away from $v$. If all face maps were acyclic,
the parent would have every direction of $\bigcup_iD_i$ outgoing at
$v$, contrary to the outgoing-cube condition. This proves the general
statement directly from CCMW's axioms.

For the example, use the certified rooted factor $A$ of order $289$
with unique corner $0$ and incident directions $D_1=\{1,2,8\}$, from
\path{research/root_cube_completion.json}. Put a copy $B$ in coordinates
$12,\ldots,23$, so its root directions are $D_2=\{13,14,20\}$.
All unspecified coordinates are zero. Set
\[
 S=A\vee B,\qquad P=\{1,2,13,14,20\},\qquad X=S\cup Q_P(0).
\]
This is a particular instance of the established root-face filling
construction in \href{research/root_cube_completion.pdf}{the root-face note}.
We recall its soundness argument to make the application explicit.
The cut-vertex theorem makes $S$ forbidden. Add the mixed words of
$Q_P(0)$ in increasing support size. Each added word fills a punctured
cube, has no neighbour outside that cube at its addition, and acquires
its previously unshattered support. Lemma~\ref{lem:damp-one-concept-extension-test}
therefore makes every intermediate family ample.
No old nonroot word of $S$ can become a corner: its original missing
corner-cube witness is confined to its own factor and is not supplied
by mixed additions. The root cannot become a corner either, since its
incident directions still include $D_1\cup D_2$ and the full cube on
those directions is not filled. Any peeling would need a first deletion
from $S$, which is impossible. All intermediate families are therefore
negative; Proposition~\ref{prop:damp-one-concept-atlas} makes each forbidden.

Every coordinate reduction of $S$ is a reduction of one positive factor,
so is positive. At each ample one-word addition, a reduction is unchanged
or gains one word. In the latter case the added word is a corner of the
ample reduction and can be deleted first. Empty reductions can only gain
a singleton. Thus every reduction of $X$ remains positive.
The two pure factor faces remain exactly $A$ and $B$, with their unique
corner at zero, whereas $Q_{D_1\cup D_2}(0)$ is missing its top word.
The first part proves the claimed impossibility of acyclic proper face maps.

Finally $|S|=577$ and the new mixed words number
$(2^2-1)(2^3-1)=21$, giving $|X|=598$.
The verifier \path{research/verify_partial_root_cube_map_obstruction.py}
checks the factor, its unique root corner, the two factor faces, and
that the crossing graph consists of two copies of $K_{12}$, with all
cross edges between $\{1,2\}$ and $\{13,14,20\}$.
After deleting any one vertex each clique stays connected and at least
one cross edge survives. This graph is two-connected, so $X$ has neither
a cut vertex nor a separating edge by
Proposition~\ref{prop:damp-crossing-separator-recovery}.
\end{proof}

\begin{proposition}[Representation of the root-face obstruction]
\label{prop:damp-root-face-representation}
The $598$-vertex obstruction in
Proposition~\ref{prop:damp-proper-face-source-conflict} admits a
representation map with sink at its shared root. Neither the existence
of a representation map nor the absence of small separators therefore
implies the existence of a map whose proper face restrictions are all
acyclic.
\end{proposition}
\begin{proof}
For the fixed $289$-vertex factor $A$ used above, the certificate
\path{research/root_sink_representation.json} gives a representation
out-map $r$ with $r(0)=\varnothing$. The solver-free replay in
\path{research/root_sink_representation.py}, invoked with
\texttt{--replay}, verifies the outgoing cube at every vertex, injectivity,
and the global condition
\[
 (x\mathbin{\triangle}y)\cap
 (r(x)\mathbin{\triangle}r(y))\ne\varnothing
 \qquad(x\ne y).
\]
Here vertices are identified with their coordinate supports. Since $A$
is ample, outgoing cubes and injectivity identify the image with
$\operatorname{Sh}(A)$; the displayed condition implies nonclashing.
The recorded orientation is cyclic, as required by the previously proved
failure of peeling to the root.

Place the same rooted map on each of the two disjoint-coordinate copies.
Both maps have empty out-map at the common root, so they agree there and
supply no outgoing exclusive directions. By
Proposition~\ref{prop:damp-uso-cube-gluing}, they give a USO on the
$577$-vertex union. Add the $21$ mixed words of the specified root face
in increasing support size, breaking ties by their integer encoding.
As in the preceding proof, each is a corner when added. Extend the
orientation by making the added corner a source, using
Proposition~\ref{prop:damp-relative-uso-extension}. This keeps the
out-map of the root empty and gives the desired USO on all $598$ vertices.
The certificate also records the complete parent out-map, and the replay
checks it directly by the same global conditions, independently of SAT.
\end{proof}

This is an application of the existing gluing and corner-extension
results, with a newly verified prescribed-root input. More generally,
two disjoint-coordinate copies of a rooted ample family admit a USO
on their one-vertex union if and only if the factor admits a USO with
that root as sink. Indeed, the gluing criterion forces at least one
of the two factor maps to have no outgoing edge at the shared root;
conversely, use a root-sink map on both copies. Every ample extension
that peels back to this union then inherits representability by relative
extension. This does not assert that arbitrary rooted ample factors
admit root-sink maps, or that every ample obstruction is representable.

The preceding source argument excludes maps with the specified acyclic
face restrictions, not representation maps altogether. Forbiddenness
still comes from the filling construction, independently of the map.
More generally, unique-corner deletion prefixes in a proper face force
all edges from each deleted corner to its later neighbours. Opposite
forced directions on one edge, or a missing cube spanned by forced
outgoing directions, certify incompatibility of positive face maps.
The replay \path{research/face_forced_prefix.py/json} checks these
constraints on the $48$ elementary faces and the two factor faces.
For $598$ it finds $12$ forced edges and the missing cube at zero.
For $627$ it finds $15$ forced edges without either contradiction, so
that case remains unknown. This explains both the successful obstruction
and the limit of the same test on the earlier search input.

A single contract-and-lift step does not go below \(267\) from the families
found so far.  For every single and double contraction of the families of
order at most \(275\) found in this search, \(573\) distinct bases in all,
CaDiCaL reports that no cornerless ample lift of order at most \(266\) exists.
These runs kept no proofs, and we make no claim based on them.

\begin{remark}[A relative-corner obstruction below the shadow threshold]
\label{rem:damp-hall-relative-corner-computation}
The owner-only shadow bound is not uniformly \(66\).  An exact all-owner
formula finds a colouring with
\[
 |M|=44,\qquad |F|=46,\qquad |\operatorname{Sh}(F)|=61.
\]
For this fixed colouring, there is an ample overlap \(K\) of order \(62\)
whose three corners consist of two vertices of \(M\) and one vertex outside
\(M\).  The induced sections are ample, and their lift has order \(294\)
and exactly one corner.  Thus neither raw shadow size nor ample completion
alone controls cornerlessness.

Remark~\ref{rem:damp-one-corner-peelable} verifies that this lift is
corner-peelable, although no peeling can leave its unique corner last.

The relative-corner criterion gives the exact obstruction.  Since an ample
VC-dimension-two family in \(Q_{11}\) has order at most \(67\), one formula
with upper bound \(67\) covers every possible overlap.  For the colouring
above, this formula is unsatisfiable when
\(\operatorname{Cor}(K)\subseteq M\).  Its DRAT certificate was independently
checked by \texttt{drat-trim}; the reduced certificate contains \(10{,}685\)
original clauses and \(550{,}770\) proof lines.  Together with the
order-\(62\) witness, this proves that the minimum number of corners of \(K\)
outside \(M\) is exactly one.

The corner defect is visible along a nested sequence
\[
 K_{62}\subset K_{65}\subset K_{66}\subset K_{67}.
\]
At every stage the overlap and both induced sections are ample, while the
unique corner outside \(M\) moves along
\[
 1894\xrightarrow{\,7\,}2022
 \xrightarrow{\,4\,}2038
 \xrightarrow{\,3\,}2046.
\]
Its unique maximal-square support is, respectively,
\[
 \{0,8\},\qquad\{0,7\},\qquad\{0,4\},\qquad\{0,3\}.
\]
Thus repairing one endpoint exposes the next one along a singleton carrier.
The last overlap is maximum of VC dimension two.  This nested chain is
verified by
\path{analyze_damp_hall_relative_corner_transport.py}; the corresponding
artifact has mathematical digest
\[
 \texttt{5f25fec458bd3750d63cd5c3c8678cc7387b3980e635614d31cb479d96c37a0b}.
\]

The witness, proof-verification record, and generating program are,
respectively,
\begin{itemize}
\item {\small\path{damp_hall_forced_core_relative_defect1_f61_o67.json}},
\item
  {\small\path{damp_hall_forced_core_relative_completion_f61_o67_proof_verification.json}},
  and
\item {\small\path{analyze_damp_hall_forced_core_completion.py}}.
\end{itemize}
Their mathematical and verification digests are recorded in the accompanying
campaign ledger.  This computation excludes one complete owner-colouring
class at every overlap order; it does not establish the corresponding
all-owner bound.
\end{remark}

\begin{proposition}[Two-ended core-corner surcharge]
\label{prop:damp-maximum-core-corner-surcharge}
Use the maximum core $A$, the repair set $R$, and the sections $P_a$ of
Proposition~\ref{prop:damp-antipodal-double-maximum-fibre}, and put
$r=|R|$.  If $H$ is cornerless, then
\begin{equation}
 |P_a|\ge r+2
 \qquad\text{for every }a\in\operatorname{corn}(A).
 \label{eq:damp-core-corner-section-surcharge}
\end{equation}
Consequently
\begin{equation}
 \Delta(H;A,R)
 \ge |\operatorname{corn}(A)|.
 \label{eq:damp-core-corner-total-surcharge}
\end{equation}
If $r\ge3$, then the stronger bounds
\begin{equation}
 |P_a|\ge r+3
 \qquad\text{for every }a\in\operatorname{corn}(A),
 \qquad
 \Delta(H;A,R)\ge2|\operatorname{corn}(A)|
 \label{eq:damp-two-ended-core-corner-surcharge}
\end{equation}
hold.  Both charges are independent of the number of repair coordinates.
\end{proposition}

\begin{proof}
Fix a corner $a$ of $A$ and let $Q_A(a)$ be its unique maximal cube, with
support $S$.  The copies of $Q_A(a)$ in the opposite fibres $H_f$ and
$H_{\overline f}$ give two vertices of the $S$-cube carrier of $H$.  That
carrier is ample and isometric, by the carrier argument in the proof of
Proposition~\ref{prop:damp-maximum-vc-three-carrier-trees}.  The two carrier
vertices differ precisely on the coordinates of $R$.  A shortest path
between them therefore has length $r$ and changes no old coordinate.  In
other words, it supplies an antipodal geodesic $P\subseteq Q_R$ such that
\begin{equation}
 Q_A(a)\subseteq H_g
 \qquad(g\in P).
 \label{eq:damp-core-corner-carrier-geodesic}
\end{equation}
In particular, $P\subseteq P_a$ and $|P|=r+1$.

Suppose that $|P_a|=r+1$.  Then $P_a=P$ is itself this geodesic.  At the
copy $(f,a)$, the repair section has exactly one incident direction, namely
the first direction of $P$.  The old incident directions are exactly the
directions of $Q_A(a)$, because $a$ is a corner of $A$.  Equation
\eqref{eq:damp-core-corner-carrier-geodesic} says that all these directions
together generate the cube $Q_A(a)$ times the first edge of $P$.  The local
corner criterion therefore makes $(f,a)$ a corner of $H$, a contradiction.
This proves \eqref{eq:damp-core-corner-section-surcharge}.

Now suppose that $r\ge3$ and $|P_a|=r+2$.  Since the carrier geodesic
$P$ already has $r+1$ vertices, write $P_a=P\cup\{w\}$.  If $w$ is not
adjacent to $f$, then the first edge of $P$ is again the only repair edge
at $(f,a)$, and the preceding corner argument applies unchanged.  Thus
cornerlessness forces $d_{Q_R}(w,f)=1$.  Applying the same argument at
the other endpoint forces $d_{Q_R}(w,\overline f)=1$.  The triangle
inequality would then give
\[
 r=d_{Q_R}(f,\overline f)\le2,
\]
contrary to $r\ge3$.  Hence $|P_a|\ge r+3$.

Finally sum the excess $|P_a|-(r+1)$ over the core corners and use the exact
decomposition \eqref{eq:damp-geodesic-excess-decomposition}.  All other
terms in that decomposition are nonnegative.  This proves
\eqref{eq:damp-core-corner-total-surcharge}, and, when $r\ge3$, summing
the two units supplied by each core corner proves
\eqref{eq:damp-two-ended-core-corner-surcharge}.
\end{proof}

\begin{proposition}[VC-one path-core propagation surcharge]
\label{prop:damp-vc-one-path-core-propagation-surcharge}
Use the maximum core $A$, repair set $R$, and sections $P_a$ from
Proposition~\ref{prop:damp-antipodal-double-maximum-fibre}.  Suppose that
$A$ has VC dimension one and uses $n\ge2$ active coordinates.  Thus its
one-inclusion graph is a path
\[
 a_0a_1\cdots a_n.
\]
If $H$ is cornerless and $r=|R|\ge3$, then
\begin{equation}
 \Delta(H;A,R)\ge
 \begin{cases}
  6,&n=2,\\
  7,&n\ge3.
 \end{cases}
 \label{eq:damp-vc-one-path-core-propagation-surcharge}
\end{equation}
\end{proposition}

\begin{proof}
For an endpoint $a\in\{a_0,a_n\}$, let $b$ be its unique neighbour in
$A$ and put
\[
 C=P_a\cap P_b.
\]
Fixing all core coordinates except the coordinate of the edge $ab$
produces an ample two-layer conditioning of $H$ with sections $P_a$ and
$P_b$.  Lemma~\ref{lem:damp-ample-two-layer-shadow-intersection} therefore
makes $C$ ample.

By Proposition~\ref{prop:damp-maximum-core-corner-surcharge}, the section
excess
\[
 e_a=|P_a|-(r+1)
\]
is at least two.  Suppose first that equality holds.  Since $P_a$ is ample
and contains two antipodes, its shadow already contains the empty set and
all $r$ singletons.  The equality $|P_a|=r+3$ leaves exactly two further
shattered supports.  No support of size at least three can occur, because
it and its three two-element subsets would already give four further
supports.  Hence $\operatorname{VCdim}(P_a)\le2$.

If $C\subsetneq P_a$, relative peeling in VC dimension two
(Lemma~\ref{lem:damp-relative-peeling-vc-two}) supplies a corner
$w\in P_a\setminus C$.  If $C=P_a$, take any corner $w$ of $P_a$.
Let $K$ be the unique maximal cube of $P_a$ through $w$.  In the first
case the core edge $ab$ is absent over $w$; in the second case
$K\times ab$ is present.  Because $a$ is an endpoint of the core path,
these are all incident core directions belonging to $A$.  Consequently,
unless $H_w$ contains a core neighbour of $a$ outside $A$, all directions
incident with $(w,a)$ generate respectively $K$ or $K\times ab$, making
$(w,a)$ a corner of $H$.  Cornerlessness therefore forces at least one
vertex of $H$ over $w$ whose core word lies outside $A$.

The same conclusion holds when $e_a=3$.  Indeed, in that case the shadow
of $P_a$ has only three members beyond the empty set and the $r$
singletons.  A shattered support of size at least three would bring with
it that support and at least three two-element subsets, so
$\operatorname{VCdim}(P_a)\le2$ again, and the preceding relative-peeling
argument applies unchanged.  We have therefore proved the following
alternative independently at the two path endpoints:
\[
 e_a\ge4
 \quad\text{or}\quad
 \text{an external core vertex adjacent to }a\text{ is present}.
\]
The two endpoint sections already contribute four units to
$\Delta(H;A,R)$.  The preceding argument in particular recovers the
preliminary bounds $\Delta\ge5$ for $n=2$ and $\Delta\ge6$ for
$n\ge3$: when $n\ge3$, an external core word cannot be adjacent to both
$a_0$ and $a_n$, whereas for $n=2$ the unique missing word of the core
square is adjacent to both endpoints.  We now exclude equality in each
preliminary bound.

Suppose first that $n=2$ and $\Delta=5$.  Put
\[
 L=P_{a_0},\qquad M=P_{a_1},\qquad N=P_{a_2}.
\]
The excess decomposition and the endpoint alternatives force
\[
 |L|=|N|=r+3,\qquad |M|=r+1,
\]
and exactly one vertex of $H$ has core word outside $A$.  The two endpoint
corners exposed by relative peeling must both be hidden by that vertex.
Thus the outside core word is the fourth word of the core square and, for
some repair word $w$,
\begin{equation}
 w\in (L\cap N)\setminus M.
 \label{eq:damp-two-edge-equality-shield-word}
\end{equation}

Let
\[
 s=|\operatorname{Sh}(L)\cap\operatorname{Sh}(N)|.
\]
Both endpoint sections contain the two repair antipodes, so their shadows
contain the empty set and all $r$ singletons; hence $s\ge r+1$.  Partition
the shattered supports of $H$ according to their intersection with the
two core coordinates.  Supports using neither core coordinate contain
\[
 \operatorname{Sh}(L)\cup\operatorname{Sh}(N),
\]
of size $2(r+3)-s$.  Supports using only the first core coordinate contain
a copy of $\operatorname{Sh}(L)\cap\operatorname{Sh}(N)$: on its two
sides the projected repair families are $L$ and $M\cup N$, because
$w\in L$.  The second core coordinate gives another such copy, using
$L\cup M$ and $N$, because $w\in N$.  Finally the two-coordinate support
itself is shattered, since all four core sections are nonempty.  Therefore
\[
 |\operatorname{Sh}(H)|
 \ge 2(r+3)-s+2s+1
 =2r+7+s
 \ge3r+8=|H|.
\]
Ampleness forces equality throughout, and in particular $s=r+1$.

On the other hand, $L\cap M$ is ample by
Lemma~\ref{lem:damp-ample-two-layer-shadow-intersection}; it contains the
two repair antipodes and is contained in the $(r+1)$-vertex family $M$.
Isometry therefore gives $L\cap M=M$.  Similarly $M\cap N=M$, so
$M\subseteq L\cap N$.  Together with
\eqref{eq:damp-two-edge-equality-shield-word}, this gives
$|L\cap N|\ge r+2$.  The Sandwich inequality now yields
\[
 r+2\le|\operatorname{Sh}(L\cap N)|
 \le|\operatorname{Sh}(L)\cap\operatorname{Sh}(N)|
 =r+1,
\]
a contradiction.  Hence $\Delta\ge6$ when $n=2$.

Now suppose that $n\ge3$ and $\Delta=6$.  The two endpoint alternatives
show that both endpoint excesses equal two, every interior section has
excess zero, and there are exactly two external-core vertices, one
adjacent to each endpoint.  Indeed, an endpoint excess of three still
requires an external vertex, an excess of at least four already spends
two additional units, and external neighbours of the two endpoints are
distinct.

Write $P_i=P_{a_i}$.  Every adjacent intersection
$C_i=P_{i-1}\cap P_i$ is ample and contains the two repair antipodes.
Since each interior section has order $r+1$, isometry gives
\[
 P_1=P_2=\cdots=P_{n-1}=:P,\qquad
 P\subseteq P_0\cap P_n.
 \label{eq:damp-long-path-equality-common-geodesic}
\]
The relative endpoint corners lie at words
\[
 w_0\in P_0\setminus P,\qquad
 w_n\in P_n\setminus P.
\]
Write their external core neighbours as
\[
 b_0=a_0^{\{x_j\}}\quad(2\le j\le n),\qquad
 b_n=a_n^{\{x_k\}}\quad(1\le k\le n-1).
 \label{eq:damp-long-path-equality-external-neighbours}
\]

Let
\[
 t=|(P_0\cap P_n)\setminus P|.
\]
Every word counted by $t$ occurs at both endpoints and at no internal
path vertex.  Isometry of its core fibre forces at least $n-1$ external
core vertices between $a_0$ and $a_n$.  These vertices are disjoint for
different repair words, so
\[
 2\ge(n-1)t.
 \label{eq:damp-long-path-equality-gap-count}
\]

For the residual shadow count needed below, put
\[
 W=H|_R=P_0\cup P_n
\]
and remove from $\operatorname{Sh}(H)$ the disjoint family
\[
 \mathcal B=
 \operatorname{Sh}(W)
 \mathbin{\dot\cup}
 \mathop{\dot\bigcup}_{i=1}^n
 \{T\cup\{x_i\}:T\in\operatorname{Sh}(P)\}.
\]
The projection $W$ and the common carrier $P$ are ample, and every
displayed support is shattered by $H$.  Since
$|W|=r+5-t$ and
$|H|=(n+1)(r+1)+6$, the residual family
$\Omega=\operatorname{Sh}(H)\setminus\mathcal B$ has
\[
 |\Omega|
 =|H|-|W|-n|P|=t+2.
 \label{eq:damp-long-path-equality-residual-size}
\]

Each external endpoint neighbour forces two kinds of residual supports.
For the left endpoint, the four core traces supplied by
$a_0,a_{j-1},a_j,b_0$ shatter the pure core pair
$\{x_1,x_j\}$.  Project $H$ onto $R\cup\{x_j\}$ and let $D_j$ be the
intersection of its two $x_j$-sections.  The two-layer shadow lemma makes
$D_j$ ample, and $P\subseteq D_j$.  Moreover
$w_0\in D_j\setminus P$.  Hence we may choose
\[
 T_0\in\operatorname{Sh}(D_j)\setminus\operatorname{Sh}(P).
\]
Then $T_0\cup\{x_j\}\in\Omega$.  Notice that $T_0$ contains at least two
repair coordinates, because $P$ contains the repair antipodes and hence
shatters the empty set and every repair singleton.  Symmetrically, the
right endpoint forces the pure core pair $\{x_k,x_n\}$ and a support
$T_n\cup\{x_k\}\in\Omega$ with a nonempty repair part.

If $t=0$, these four named supports have at least three distinct values.
Indeed, if the two pure core pairs are distinct, either mixed support is
a third value.  If they coincide, then necessarily $j=n$ and $k=1$, so
the two mixed supports use the distinct core coordinates $x_n$ and $x_1$.
This contradicts $|\Omega|=2$.

It remains to consider $t>0$.  Equation
\eqref{eq:damp-long-path-equality-gap-count} forces $n=3$ and $t=1$.
The unique common word outside $P$ must be both $w_0$ and $w_n$, and its
core fibre must use the two external vertices as the internal vertices of
a shortest $a_0$--$a_3$ path.  Hence $b_0$ and $b_n$ are adjacent and
$j\ne k$.  The original core path and this second antipodal geodesic have
coordinate orders $(1,2,3)$ and, respectively,
$(2,3,1)$, $(3,2,1)$, or $(3,1,2)$.  The latter order has at least two
inversions relative to the former, and every inverted pair is shattered
by the union of the two geodesics.  Thus $\Omega$ contains at least two
distinct pure core pairs.  The two mixed supports above are also distinct,
because their unique core coordinates are $x_j\ne x_k$, and neither is a
pure core pair.  Hence $|\Omega|\ge4$, contradicting
$|\Omega|=t+2=3$.

Thus equality is impossible for $n\ge3$, and $\Delta\ge7$.
\end{proof}

\begin{proposition}[Path-spine Euler identity]
\label{prop:damp-path-spine-euler-identity}
Keep the VC-one path core
$a_0a_1\cdots a_n$ and its edge coordinates
$x_1,\ldots,x_n$ from Proposition
\ref{prop:damp-vc-one-path-core-propagation-surcharge}.  Put
\[
 P_i=P_{a_i},\qquad
 C_i=P_{i-1}\cap P_i\quad(1\le i\le n),
\]
let
\[
 W=H|_R,\qquad U=\bigcup_{i=0}^nP_i,
\]
and let
\[
 V=\bigl|H\setminus(Q_R\times A)\bigr|.
\]
For $g\in U$, let $c(g)$ be the number of connected components of
\[
 \{i\in\{0,\ldots,n\}:g\in P_i\}
\]
in the path on $\{0,\ldots,n\}$, and put
\[
 \Phi=\sum_{g\in U}\bigl(c(g)-1\bigr).
\]
Finally define the disjoint spine family
\begin{equation}
 \mathcal B=
 \operatorname{Sh}(W)
 \mathbin{\dot\cup}
 \mathop{\dot\bigcup}_{i=1}^n
 \{T\cup\{x_i\}:T\in\operatorname{Sh}(C_i)\}
 \subseteq\operatorname{Sh}(H)
 \label{eq:damp-path-spine-shadow-family}
\end{equation}
and its residual shadow
\[
 \Omega=\operatorname{Sh}(H)\setminus\mathcal B.
\]
Then
\begin{equation}
 |\Omega|
 =\Phi+V-|W\setminus U|.
 \label{eq:damp-path-spine-euler-identity}
\end{equation}
In particular, both terms on the right are nonnegative: $\Phi$ counts
disconnected reappearances of repair words along the core path, while
$V-|W\setminus U|$ counts external-core multiplicity after subtracting
the one vertex needed to introduce each repair word outside $U$.
\end{proposition}

\begin{proof}
The projection $W$ is ample.  Each $C_i$ is ample by applying
Lemma~\ref{lem:damp-ample-two-layer-shadow-intersection} to the
conditioning that fixes all core coordinates except $x_i$.  That
conditioned family also shows that every
$T\cup\{x_i\}$ with $T\in\operatorname{Sh}(C_i)$ is shattered by $H$.
The supports in \eqref{eq:damp-path-spine-shadow-family} are pairwise
distinct, so
\[
 |\mathcal B|=|W|+\sum_{i=1}^n|C_i|.
\]
Since $H$ is ample,
\begin{equation}
 |\Omega|
 =|H|-|W|-\sum_{i=1}^n|C_i|
 =V+\sum_{i=0}^n|P_i|-|W|-\sum_{i=1}^n|C_i|.
 \label{eq:damp-path-spine-euler-count}
\end{equation}

Fix $g\in U$.  If it belongs to $m(g)$ of the sections and to both
sections at $b(g)$ adjacent path edges, then its contribution to
$\sum_i|P_i|-\sum_i|C_i|$ is $m(g)-b(g)$.  The induced subgraph on its
section indices is a disjoint union of $c(g)$ paths, and hence
$m(g)-b(g)=c(g)$.  Summing over $g$ gives
\[
 \sum_{i=0}^n|P_i|-\sum_{i=1}^n|C_i|
 =\sum_{g\in U}c(g)=|U|+\Phi.
\]
Substitution in \eqref{eq:damp-path-spine-euler-count}, together with
$U\subseteq W$, proves \eqref{eq:damp-path-spine-euler-identity}.

Finally, every word of $W\setminus U$ supports at least one vertex of $H$
whose core word lies outside $A$.  These vertices are distinct for
distinct repair words, so $V\ge|W\setminus U|$.  This proves the final
assertion.
\end{proof}

\begin{proposition}[The support-interval exact sequence]
\label{prop:damp-path-spine-support-interval-exact-sequence}
Keep the notation of Proposition
\ref{prop:damp-path-spine-euler-identity}, and put
\[
 \mathcal L_R=\{\varnothing\}\cup\binom R1,
 \qquad
 \mathcal S=\bigcup_{j=0}^n\operatorname{Sh}(P_j),
 \qquad
 \mathcal T=\mathcal S\setminus\mathcal L_R.
\]
For \(T\in\mathcal T\), let
\[
 I_T=\{j:T\in\operatorname{Sh}(P_j)\},
 \qquad m_T=|I_T|,
\]
and let \(c_T\) be the number of connected components of \(I_T\) in the
section path.  Finally, put
\[
 q_{\rm ext}
 =\bigl|\operatorname{Sh}(W)\setminus\mathcal S\bigr|,
 \qquad
 \Sigma_0=\operatorname{Sh}(A)*\mathcal L_R,
 \qquad
 \Xi=\mathcal B\setminus\Sigma_0.
\]
Then
\begin{align}
 \sum_{j=0}^n\bigl(|P_j|-(r+1)\bigr)
   &=\sum_{T\in\mathcal T}m_T,
 \label{eq:damp-support-interval-section-excess}\\
 |\Xi|
   &=\sum_{T\in\mathcal T}(m_T-c_T+1)+q_{\rm ext},
 \label{eq:damp-support-interval-spine-excess}\\
 \Delta(H;A,R)
   &=|\Xi|+|\Omega|,
 \label{eq:damp-support-interval-total-excess}\\
 |\Omega|
   &=V+\sum_{T\in\mathcal T}(c_T-1)-q_{\rm ext}.
 \label{eq:damp-support-interval-residual-balance}
\end{align}

There is also an exact linear form of these identities.  Fix a field
\(k\).  For each \(T\in\mathcal T\), choose one component of \(I_T\)
and one index \(j_T\) in that component.  Let \(C_0(T)\) have basis
\(e_{j,T}\), \(j\in I_T\), and let \(C_1(T)\) have one basis vector
\(u_T\), together with one basis vector \(u_{i,T}\) for every path edge
\(a_{i-1}a_i\) whose two ends belong to \(I_T\).  The map
\begin{equation}
 \partial_Tu_T=e_{j_T,T},
 \qquad
 \partial_Tu_{i,T}=e_{i,T}-e_{i-1,T}
 \label{eq:damp-support-interval-boundary-map}
\end{equation}
is injective and gives an exact sequence
\begin{equation}
 0\longrightarrow C_1(T)
 \xrightarrow{\ \partial_T\ }C_0(T)
 \longrightarrow k^{c_T-1}\longrightarrow0.
 \label{eq:damp-support-interval-exact-sequence}
\end{equation}
The basis of \(C_0(T)\) records the occurrences of \(T\) in the section
shadows.  The basis of \(C_1(T)\) records exactly the spine supports
\[
 T,
 \qquad T\cup\{x_i\}
 \quad\bigl(T\in\operatorname{Sh}(C_i)\bigr).
\]
Consequently a successful adjacent duplication, which transports one
support \(T\) across one section edge, changes the occurrence vector by a
member of \(\operatorname{im}\partial_T\).  Successful acquisitions are
therefore relations in the support-interval complex, not an additional
cardinality charge.  The only supportwise cokernel is the one-dimensional
defect of every component of \(I_T\) after the distinguished component.
\end{proposition}

\begin{proof}
Every \(P_j\), every \(C_i\), and \(W\) contains the two repair
antipodes.  Their shadows therefore contain the common repair baseline
\(\mathcal L_R\), of order \(r+1\).  Since each \(P_j\) is ample,
counting every nonbaseline repair support in every section gives
\eqref{eq:damp-support-interval-section-excess}.

The two-layer shadow-intersection identity gives
\[
 \operatorname{Sh}(C_i)
 =\operatorname{Sh}(P_{i-1})
   \cap\operatorname{Sh}(P_i).
\]
For a fixed \(T\in\mathcal T\), the number of path edges whose two ends
belong to \(I_T\) is \(m_T-c_T\).  Thus the carrier blocks of
\(\mathcal B\setminus\Sigma_0\) contain exactly
\(m_T-c_T\) supports of the form \(T\cup\{x_i\}\).  The pure repair
block contains \(T\) itself for every \(T\in\mathcal T\), together with
the \(q_{\rm ext}\) supports that occur in
\(\operatorname{Sh}(W)\) but in no section shadow.  The blocks are
disjoint, proving \eqref{eq:damp-support-interval-spine-excess}.

The baseline \(\Sigma_0\) has order \((n+1)(r+1)\).  Since \(H\) is
ample and
\[
 \operatorname{Sh}(H)
 =\Sigma_0\mathbin{\dot\cup}\Xi
  \mathbin{\dot\cup}\Omega,
\]
subtracting the baseline proves
\eqref{eq:damp-support-interval-total-excess}.  On the other hand, the
geodesic excess decomposition
\eqref{eq:damp-geodesic-excess-decomposition} and
\eqref{eq:damp-support-interval-section-excess} give
\[
 \Delta(H;A,R)=V+\sum_{T\in\mathcal T}m_T.
\]
Subtracting \eqref{eq:damp-support-interval-spine-excess} proves
\eqref{eq:damp-support-interval-residual-balance}.

For the linear statement, adjoin a new root vertex to the component of
the path induced by \(I_T\) that contains \(j_T\), using one edge from
the root to \(j_T\).  Keep all internal edges of the induced path and no
other edges.  The resulting graph is a forest with \(m_T+1\) vertices,
\(m_T-c_T+1\) edges, and \(c_T\) components.  After deleting the row of
the new root from its oriented incidence matrix, one obtains exactly
\eqref{eq:damp-support-interval-boundary-map}.  The columns are linearly
independent, and their cokernel has dimension
\[
 m_T-(m_T-c_T+1)=c_T-1.
\]
This proves \eqref{eq:damp-support-interval-exact-sequence}.  Finally,
Lemma~\ref{lem:damp-adjacent-duplication-shadow-transport} says that a
successful duplication across \(a_{i-1}a_i\) transports one existing
support \(T\) between the two corresponding section shadows.  Its change
in \(C_0(T)\) is therefore, up to orientation,
\(e_{i,T}-e_{i-1,T}=\partial_Tu_{i,T}\), proving the last assertion.
\end{proof}

\begin{proposition}[Projected carriers double support fragmentation]
\label{prop:damp-projected-carrier-doubled-support-fragmentation}
Keep the notation of Proposition
\ref{prop:damp-path-spine-support-interval-exact-sequence}.  For
\(T\in\mathcal T\), define the carrier of all positioned \(T\)-cubes and
its core projection by
\[
 \mathcal K_T(H)
 =\{y\in Q_{(R\mathbin{\dot\cup}X)\setminus T}:
       Q_T\times\{y\}\subseteq H\},
 \qquad
 Z_T=\mathcal K_T(H)|_X.
\]
Write the components of \(I_T\) as
\[
 [\ell_1,r_1],\ldots,[\ell_{c_T},r_{c_T}],
 \qquad r_h+1<\ell_{h+1},
\]
and put
\[
 \lambda_T
 =\sum_{h=1}^{c_T-1}(\ell_{h+1}-r_h-1).
\]
Finally let
\[
 \Omega[T]=\{S\in\Omega:S\cap R=T\}.
\]
Then \(Z_T\) is ample,
\begin{equation}
 Z_T\cap A=\{a_j:j\in I_T\},
 \qquad
 |Z_T\setminus A|\ge\lambda_T\ge c_T-1,
 \label{eq:damp-projected-carrier-path-gap}
\end{equation}
and
\begin{equation}
 |\Omega[T]|
 =|Z_T\setminus A|+c_T-1
 \ge\lambda_T+c_T-1
 \ge2(c_T-1).
 \label{eq:damp-projected-carrier-doubled-support-fragmentation}
\end{equation}
Consequently
\begin{align}
 |\Omega|
 &\ge2\sum_{T\in\mathcal T}(c_T-1),
 \label{eq:damp-total-doubled-support-fragmentation}\\
 V-q_{\rm ext}
 &\ge\sum_{T\in\mathcal T}(c_T-1).
 \label{eq:damp-external-support-fragmentation-balance}
\end{align}
Thus changing the positioned cube that realizes \(T\) between two section
intervals creates no additional multiplicity problem: all positions are
already assembled by the ample projected carrier \(Z_T\), and each extra
interval component has two residual units before any ranks are used.
\end{proposition}

\begin{proof}
The support \(T\) is shattered by some \(P_j\), hence by \(H\), and
ampleness makes it strongly shattered.  Apply
Lemma~\ref{lem:damp-projected-cube-carrier-support-conservation} with the
roles of \(R\) and \(X\) interchanged.  It makes
\(\mathcal K_T(H)\) and \(Z_T\) ample and gives
\begin{equation}
 \operatorname{Sh}(Z_T)
 =\{K\subseteq X:T\mathbin{\dot\cup}K
                    \in\operatorname{Sh}(H)\}.
 \label{eq:damp-repair-carrier-core-support-lift}
\end{equation}
A path word \(a_j\) belongs to \(Z_T\) exactly when some positioned full
\(T\)-cube occurs in the repair section \(P_j\), which is equivalent to
\(T\in\operatorname{Sh}(P_j)\).  This proves the first equality in
\eqref{eq:damp-projected-carrier-path-gap}.

The ample family \(Z_T\) is isometric.  Consider two consecutive
components of \(I_T\), ending at \(p=r_h\) and beginning at
\(q=\ell_{h+1}\).  A shortest \(Z_T\)-path from \(a_p\) to \(a_q\) has
length \(q-p\) and lies in their interval subcube.  None of its
\(q-p-1\) internal vertices belongs to \(A\), because the internal path
vertices \(a_{p+1},\ldots,a_{q-1}\) are absent from \(Z_T\).  The interval
subcubes belonging to different gaps have disjoint interiors.  Summing
their internal vertices gives
\[
 |Z_T\setminus A|\ge
 \sum_{h=1}^{c_T-1}(\ell_{h+1}-r_h-1)
 =\lambda_T.
\]
Every gap is nonempty, so \(\lambda_T\ge c_T-1\), completing
\eqref{eq:damp-projected-carrier-path-gap}.

It remains to count supports.  Equation
\eqref{eq:damp-repair-carrier-core-support-lift} and ampleness of \(Z_T\)
show that \(H\) has exactly \(|Z_T|\) shattered supports whose exact
repair part is \(T\).  Among them, the path-spine family contains \(T\)
itself and precisely the \(m_T-c_T\) supports
\[
 T\cup\{x_i\}
 \qquad
 \bigl(T\in\operatorname{Sh}(C_i)\bigr)
\]
coming from the internal edges of \(I_T\).  Every other support with
exact repair part \(T\) belongs to \(\Omega[T]\).  Since
\(|Z_T\cap A|=m_T\), subtraction gives
\[
 |\Omega[T]|
 =|Z_T|-1-(m_T-c_T)
 =|Z_T\setminus A|+c_T-1.
\]
Together with \eqref{eq:damp-projected-carrier-path-gap}, this proves
\eqref{eq:damp-projected-carrier-doubled-support-fragmentation}.

The families \(\Omega[T]\), \(T\in\mathcal T\), are pairwise disjoint
subsets of \(\Omega\).  Summing their lower bounds proves
\eqref{eq:damp-total-doubled-support-fragmentation}.  Finally compare this
inequality with the exact residual balance
\eqref{eq:damp-support-interval-residual-balance} and cancel
\(\sum_T(c_T-1)\).  The result is
\eqref{eq:damp-external-support-fragmentation-balance}.
\end{proof}

\begin{proposition}[Punctured antipodal-section surcharge]
\label{prop:damp-punctured-antipodal-section-surcharge}
Let $P\subseteq Q_R$ be ample, let $f,\overline f\in P$ be antipodes,
and let $T\subsetneq R$ have $|T|=k\ge2$.  Suppose that
\begin{equation}
 f^U\in P\quad(U\subsetneq T),
 \qquad f^T\notin P.
 \label{eq:damp-punctured-antipodal-section}
\end{equation}
Then
\begin{equation}
 |P|\ge |R|+2^k.
 \label{eq:damp-punctured-antipodal-section-surcharge}
\end{equation}
More precisely, for every $t\in T$ there is a coordinate
$c_t\in R\setminus T$ such that $P$ shatters $\{t,c_t\}$, and these
$k$ pairs are distinct.
\end{proposition}

\begin{proof}
Translate the repair cube so that $f=0_R$.  The family $P$ is isometric.
Fix $t\in T$ and consider a shortest path in $P$ from
$T\setminus\{t\}$ to $R$, where words are identified with their supports.
The coordinates changed by this path are exactly
$(R\setminus T)\cup\{t\}$.  Its first edge cannot have direction $t$,
because the resulting vertex would be the absent word $T$.  Hence that
edge has a direction $c_t\in R\setminus T$, and
\[
 (T\setminus\{t\})\cup\{c_t\}\in P.
\]
The four words
\[
 \varnothing,\qquad \{t\},\qquad
 (T\setminus\{t\})\cup\{c_t\},\qquad R
\]
realize the four traces on $\{t,c_t\}$.  Thus this pair is shattered.
The pairs obtained for distinct $t$ are distinct, since their coordinate
in $T$ is different and their other coordinate lies outside $T$.

It remains to count shattered supports.  The empty set and the $|R|$
singletons are shattered because $P$ contains two antipodes.  Every proper
subset of $T$ is strongly shattered by the punctured face in
\eqref{eq:damp-punctured-antipodal-section}, while $T$ itself is shattered
because $\overline f$ supplies its one missing trace.  Hence $P$ shatters
all $2^k-k-1$ subsets of $T$ of order at least two.  The $k$ pairs
$\{t,c_t\}$ are distinct from these supports.  Therefore
\[
 |\operatorname{Sh}(P)|
 \ge (|R|+1)+(2^k-k-1)+k
 =|R|+2^k.
\]
Ampleness gives $|P|=|\operatorname{Sh}(P)|$, proving the result.
\end{proof}

\begin{proposition}[Opposite-puncture projection signature]
\label{prop:damp-opposite-puncture-projection-signature}
Let $P\subseteq Q_R$ be ample, let $f,\overline f\in P$ be antipodes,
and let $T\subsetneq R$ have $|T|\ge2$.  Suppose that
\begin{equation}
 \begin{aligned}
  f^U&\in P&&(U\subsetneq T),& f^T&\notin P,\\
  \overline f^U&\in P&&(U\subsetneq T),&
  \overline f^T&\notin P.
 \end{aligned}
 \label{eq:damp-opposite-puncture-pair}
\end{equation}
For every $c\in R\setminus T$, the projection
$P|_{T\cup\{c\}}$ is exactly one of
\begin{equation}
 Q_{T\cup\{c\}},\qquad
 Q_{T\cup\{c\}}\setminus\{(f|_T,\overline f_c)\},\qquad
 Q_{T\cup\{c\}}\setminus\{(\overline f|_T,f_c)\}.
 \label{eq:damp-opposite-puncture-projection-trichotomy}
\end{equation}
In the first case $T\cup\{c\}$ is shattered by $P$.  In particular,
$|R\setminus T|\ge2$.  Moreover, writing $k=|T|$, one has
\begin{equation}
 |P|\ge |R|+2^{k+1}-k-2.
 \label{eq:damp-opposite-puncture-collision-surcharge}
\end{equation}
One also has the independent shadow bound
\begin{equation}
 |P|\ge 2^k+(2^k-1)(|R|-k).
 \label{eq:damp-opposite-puncture-cross-section-surcharge}
\end{equation}
\end{proposition}

\begin{proof}
Fix $c\in R\setminus T$.  The two punctured faces in
\eqref{eq:damp-opposite-puncture-pair} realize every trace on
$T\cup\{c\}$ except possibly the antipodal pair
\[
 (\overline f|_T,f_c),\qquad(f|_T,\overline f_c).
\]
Thus the projection contains the twice-punctured cube obtained by deleting
these two traces.  Projections of ample families are ample.  The twice-
punctured cube is not ample, whereas adjoining either missing trace gives a
singly punctured cube and adjoining both gives the full cube.  These are all
possibilities, proving
\eqref{eq:damp-opposite-puncture-projection-trichotomy}.  The full-cube case
shatters its entire coordinate set.

If $R\setminus T=\{c\}$, the projection is $P$ itself.  Both exceptional
traces are absent by
\eqref{eq:damp-opposite-puncture-pair}, so $P$ would be the nonample twice-
punctured cube, a contradiction.  Hence $|R\setminus T|\ge2$.

For the cardinality bound, put $C=R\setminus T$ and translate so that
$f|_C=0_C$.  Let
\[
 F_0=\{g\in Q_C:(f|_T,g)\in P\},\qquad
 F_1=\{g\in Q_C:(\overline f|_T,g)\in P\}.
\]
These are nonempty ample conditionings of $P$.  The two punctured faces in
\eqref{eq:damp-opposite-puncture-pair} contribute
$2(2^k-1)$ vertices, and among them the only member counted by $F_0$ is
$0_C$, while the only member counted by $F_1$ is $1_C$.

For every $c\in C$, the projection trichotomy says that $c$ is active in
at least one of $F_0,F_1$.  Assign it to one such fiber and write
$C=C_0\mathbin{\dot\cup}C_1$ for the resulting partition.  Since an ample
family is isometric and hence connected, a family active on $s$ coordinates
has at least $s+1$ vertices.  Therefore
\[
 |F_0|\ge |C_0|+1,
 \qquad |F_1|\ge |C_1|+1.
\]
Apart from their displayed base vertices, these two fibers contribute at
least $|C_0|+|C_1|=|C|$ further vertices, disjoint from the two punctured
faces.  Consequently
\[
 |P|\ge2(2^k-1)+|C|
 =|R|+2^{k+1}-k-2,
\]
as required.

For the shadow bound, every subset of $T$ is shattered: the two punctured
faces supply all proper traces, and either antipodal face supplies the trace
missing from the other.  More strongly, fix $c\in C$ and a proper subset
$S\subsetneq T$.  The lower punctured face, on which $c=f_c$, realizes every
trace on $S$: any prescribed $S$-trace has an extension to $T$ different
from the one missing lower tip, because $T\setminus S\ne\varnothing$.
The upper punctured face gives the same conclusion with
$c=\overline f_c$.  Together the two faces realize every trace on
$S\cup\{c\}$, so this support is shattered.  The supports wholly contained
in $T$ give $2^k$ distinct members of the shadow, and, independently for
each $c\in C$, the proper subsets $S\subsetneq T$ give another $2^k-1$.
Consequently
\[
 |\operatorname{Sh}(P)|
 \ge 2^k+(2^k-1)|C|
 =2^k+(2^k-1)(|R|-k).
\]
Ampleness gives the claimed order bound.
\end{proof}

\begin{corollary}[Bichromatic collisions lie in an ordered two-cube collar]
\label{cor:damp-bichromatic-opposite-puncture-collar}
Keep the hypotheses of Proposition
\ref{prop:damp-opposite-puncture-projection-signature}, and suppose that no
$T\cup\{c\}$ is shattered.  Translate so that $f=0_R$.  Partition
$C=R\setminus T$ as $L\mathbin{\dot\cup}U$, where $c\in L$ when
$P|_{T\cup\{c\}}$ contains $(1_T,0_c)$ and omits $(0_T,1_c)$, and
$c\in U$ in the opposite case.  For every $c\in L$ and $d\in U$, either
$P$ shatters $\{c,d\}$ or every $g\in P$ satisfies $g_c\le g_d$.
Consequently, if no such pair is shattered, then
\begin{equation}
 P|_C\subseteq
 (\{0_L\}\times Q_U)\cup(Q_L\times\{1_U\}).
 \label{eq:damp-bichromatic-opposite-puncture-collar}
\end{equation}
\end{corollary}

\begin{proof}
Fix $c\in L$ and $d\in U$.  The two antipodal carrier words give the
traces $00$ and $11$ on $\{c,d\}$.  Since $c\in L$, there is a carrier
word with $T$-trace $1_T$ and $c$-trace zero.  The defining omission for
$d\in U$ forces its $d$-trace to be one, so the trace $01$ is also present.
If $10$ occurs, the pair is shattered.  Otherwise every carrier word avoids
$10$, which is exactly the implication $g_c\le g_d$.

Suppose that no bichromatic pair is shattered.  If a word in $P|_C$ has a
one in some coordinate of $L$, all its coordinates in $U$ are one by the
preceding implication, and it belongs to $Q_L\times\{1_U\}$.  If it has no
one in $L$, it belongs to $\{0_L\}\times Q_U$.  This proves
\eqref{eq:damp-bichromatic-opposite-puncture-collar}.
\end{proof}

\begin{proposition}[Shadow-tight opposite collisions are serial peripheral]
\label{prop:damp-shadow-tight-opposite-collision-normal-form}
Keep the hypotheses of Proposition
\ref{prop:damp-opposite-puncture-projection-signature}, translate so that
$f=0_R$, and put $C=R\setminus T$.  Suppose that
\begin{equation}
 T\cup\{c\}\notin\operatorname{Sh}(P)
 \quad(c\in C),
 \qquad
 \{c,d\}\notin\operatorname{Sh}(P)
 \quad(c,d\in C,\ c\ne d).
 \label{eq:damp-shadow-tight-opposite-collision}
\end{equation}
Then there are an ordering $c_1,\ldots,c_m$ of $C$, where $m=|C|$, and
an index $h$ with $1\le h\le m-1$ such that, on writing
\[
 v_j=1_{\{c_1,\ldots,c_j\}}\in Q_C
 \qquad(0\le j\le m),
\]
the family $P$ is exactly
\begin{equation}
 \begin{split}
 P={}&\bigcup_{j<h}
     \bigl((Q_T\setminus\{1_T\})\times\{v_j\}\bigr)\\
 &\mathbin{\cup}\bigl(Q_T\times\{v_h\}\bigr)
 \mathbin{\cup}\bigcup_{j>h}
     \bigl((Q_T\setminus\{0_T\})\times\{v_j\}\bigr).
 \end{split}
 \label{eq:damp-shadow-tight-opposite-collision-normal-form}
\end{equation}
In the notation of Corollary
\ref{cor:damp-bichromatic-opposite-puncture-collar}, one has
\[
 U=\{c_1,\ldots,c_h\},
 \qquad
 L=\{c_{h+1},\ldots,c_m\}.
\]
Contracting every coordinate of $C\setminus\{c_h,c_{h+1}\}$ sends $P$ to
the canonical collar
\begin{equation}
 \mathcal K_T
 =((Q_T\setminus\{1_T\})\times\{00\})
  \cup(Q_T\times\{10\})
  \cup((Q_T\setminus\{0_T\})\times\{11\}).
 \label{eq:damp-canonical-opposite-puncture-collar}
\end{equation}
In particular, both orientation classes are nonempty, the outside projection
is a coordinate-simple path, and $P\in\Damp$.  Thus every opposite collision
either shatters one of the supports excluded in
\eqref{eq:damp-shadow-tight-opposite-collision}, or its entire carrier has
the explicit corner-peelable normal form
\eqref{eq:damp-shadow-tight-opposite-collision-normal-form}.
\end{proposition}

\begin{proof}
The supports already shown to be shattered in the proof of Proposition
\ref{prop:damp-opposite-puncture-projection-signature} are precisely
\begin{equation}
 \Sigma_0
 =2^T\mathbin{\cup}
   \{S\cup\{c\}:S\subsetneq T,\ c\in C\}.
 \label{eq:damp-shadow-tight-mandatory-complex}
\end{equation}
Every support outside $\Sigma_0$ either contains two coordinates of $C$, or
contains $T\cup\{c\}$ for some $c\in C$.  By downward closure and
\eqref{eq:damp-shadow-tight-opposite-collision}, none is shattered.
Therefore $\operatorname{Sh}(P)=\Sigma_0$.  Its minimal nonfaces are exactly
the pairs $\{c,d\}\subseteq C$ and the sets $T\cup\{c\}$, $c\in C$.

The antipodes realize the traces $00$ and $11$ on every outside pair.
Theorem~\ref{thm:damp-fixed-shadow-signed-clutter} therefore says that the
unique omitted trace on each $\{c,d\}$ is either $01$ or $10$.  Regard the
corresponding forbidden trace as one of the implications
$x_c\le x_d$ and $x_d\le x_c$, and let $F\subseteq Q_C$ be their common
satisfying family.  For a trace
$\alpha\in Q_T\setminus\{0_T,1_T\}$, none of the signed clauses on
$T\cup\{c\}$ is active.  The same signed-minimal-nonface theorem gives
\begin{equation}
 \{z\in Q_C:(\alpha,z)\in P\}=F.
 \label{eq:damp-shadow-tight-middle-fiber}
\end{equation}
Both punctured faces show that $0_C,1_C\in F$, and this conditioning of the
ample family $P$ is ample.

The implications orient every pair of coordinates of $C$.  Collapse the
strong components of the resulting tournament.  Coordinates in one strong
component have equal values in every member of $F$, while the condensation
is an acyclic tournament and hence a total order.  If it has $q$ components,
its satisfying family has exactly $q+1$ threshold words.  On the other hand,
$F$ contains two antipodes, so it shatters every singleton of $C$.  Ampleness
therefore gives
\[
 q+1=|F|=|\operatorname{Sh}(F)|\ge m+1.
\]
Since $q\le m$, equality holds and every strong component is a singleton.
After relabelling the coordinates, $F$ is the coordinate-simple path
\begin{equation}
 F=\{v_0,v_1,\ldots,v_m\}.
 \label{eq:damp-shadow-tight-outside-path}
\end{equation}

For $\beta\in Q_T$, write
$P_\beta=\{z\in Q_C:(\beta,z)\in P\}$.  Use the partition
$C=L\mathbin{\dot\cup}U$ from Corollary
\ref{cor:damp-bichromatic-opposite-puncture-collar}.  The signed clause on
$T\cup\{c\}$ forces $c=0$ in the $0_T$-fiber when $c\in L$, and forces
$c=1$ in the $1_T$-fiber when $c\in U$.  It imposes no other restriction.
Consequently
\begin{equation}
 P_{0_T}=\{v_j:c=0\text{ on }v_j\text{ for every }c\in L\},
 \quad
 P_{1_T}=\{v_j:c=1\text{ on }v_j\text{ for every }c\in U\},
 \label{eq:damp-shadow-tight-extreme-fibers}
\end{equation}
while every other $T$-fiber is the full path $F$.  The missing tips
$(1_T,0_C)$ and $(0_T,1_C)$ show respectively that $U$ and $L$ are
nonempty.

Let $p$ be the first position occupied by a member of $L$, and let $q$ be
the last position occupied by a member of $U$ in the order
$c_1,\ldots,c_m$.  Equations
\eqref{eq:damp-shadow-tight-middle-fiber} and
\eqref{eq:damp-shadow-tight-extreme-fibers} give
\[
 |P|=(2^{|T|}-2)(m+1)+p+(m-q+1).
\]
On the other hand,
\eqref{eq:damp-shadow-tight-mandatory-complex} and ampleness give
\[
 |P|=|\Sigma_0|=2^{|T|}+(2^{|T|}-1)m.
\]
Thus $p-q=1$.  Since $L$ and $U$ partition $C$, every member of $U$ occurs
before every member of $L$.  Put $h=q=p-1$.  The fiber description now
becomes exactly
\eqref{eq:damp-shadow-tight-opposite-collision-normal-form}.
Its projection to $T\cup\{c_h,c_{h+1}\}$ has outside traces $00,10,11$;
the three corresponding $T$-sections are respectively
$Q_T\setminus\{1_T\}$, $Q_T$, and $Q_T\setminus\{0_T\}$.  This proves
\eqref{eq:damp-canonical-opposite-puncture-collar}.

It remains to verify the final dismantlability assertion.  Delete the
layers $v_0,v_1,\ldots,v_{h-1}$ successively from the left.  At each step the
coordinate changed by the first path edge is peripheral: its smaller section
is a fixed copy of $Q_T\setminus\{1_T\}$, contained in the next section.
Delete the layers $v_m,v_{m-1},\ldots,v_{h+1}$ symmetrically from the right;
the duplicated site is now a fixed copy of
$Q_T\setminus\{0_T\}$.  The remaining layer is $Q_T$.  The cube and the two
singly punctured cubes are corner-peelable, so repeated application of the
exact peripheral-expansion test
\eqref{eq:damp-exact-peripheral-expansion-test} proves $P\in\Damp$.
\end{proof}

\begin{corollary}[Relative peeling above the canonical low-rank collar]
\label{cor:damp-canonical-collar-relative-peeling}
Let $2\le |T|\le3$, let $\mathcal K_T$ be the family in
\eqref{eq:damp-canonical-opposite-puncture-collar}, and let
\[
 \mathcal K_T\subseteq B\subseteq Q_{T\cup\{u,\ell\}}
\]
be ample.  For every $s\in\mathcal K_T$, the family $B$ has a corner
peeling in which every member of $B\setminus\mathcal K_T$ precedes every
member of $\mathcal K_T$, and $s$ is last.
\end{corollary}

\begin{proof}
The ambient cube has $|T|+2\le5$ coordinates.  Apply the five-coordinate
relative-peeling theorem,
Theorem~\ref{thm:damp-five-coordinate-relative-peeling}, to the nested ample
pair $\mathcal K_T\subseteq B$; ampleness of $\mathcal K_T$ also follows
from Proposition~\ref{prop:damp-shadow-tight-opposite-collision-normal-form}
or directly from its peripheral construction.
\end{proof}

\begin{proposition}[The canonical collar has exactly two bad extensions]
\label{prop:damp-canonical-collar-two-bad-extensions}
Let $|T|\ge2$, write the last two coordinates of
$\mathcal K_T$ in the order used in
\eqref{eq:damp-canonical-opposite-puncture-collar}, and put
\[
 b_0=(0_T,01),
 \qquad
 b_1=(1_T,01).
\]
For every $x\in Q_{T\cup\{u,\ell\}}\setminus\mathcal K_T$,
\begin{equation}
 \mathcal K_T\cup\{x\}\text{ is nonample}
 \quad\Longleftrightarrow\quad
 x\in\{b_0,b_1\}.
 \label{eq:damp-canonical-collar-two-bad-extensions}
\end{equation}
More precisely, for each $t\in T$, the two bad tips have the antipodal
blockers
\begin{equation}
 (1_T^{\{t\}},00)\longleftrightarrow b_1
 \quad\text{on }\{t,\ell\},
 \qquad
 (0_T^{\{t\}},11)\longleftrightarrow b_0
 \quad\text{on }\{t,u\}.
 \label{eq:damp-canonical-collar-explicit-blockers}
\end{equation}
Thus every other missing word is an allowable first step in a relative
peeling above $\mathcal K_T$.
\end{proposition}

\begin{proof}
Each pair in
\eqref{eq:damp-canonical-collar-explicit-blockers} is antipodal in the
indicated square, and that square meets the corresponding one-concept
extension in exactly those two vertices.  Proposition
\ref{prop:damp-failed-augmentation-antipodal-localization} shows that the two
extensions are nonample.

Conversely, suppose that adjoining $x$ is nonample.  The same proposition
gives a blocker $y\in\mathcal K_T$ and a set $J$ of at least two coordinates
such that the $J$-face through $y$ contains no other collar vertex and
$x=y^J$.  If a coordinate of $J$ were incident with $y$ in
$\mathcal K_T$, the corresponding neighbor would be a second collar vertex
in that face.  Hence
\begin{equation}
 J\subseteq (T\cup\{u,\ell\})\setminus I_{\mathcal K_T}(y).
 \label{eq:damp-canonical-collar-blocker-nonincident}
\end{equation}

The nonincident directions can be read directly from the three layers in
\eqref{eq:damp-canonical-opposite-puncture-collar}.  In the $00$-layer, a
word has at least two nonincident directions only when it is
$1_T^{\{t\}}$ for some $t\in T$; its two nonincident directions are then
$t$ and $\ell$.  Every other word in that layer has only $\ell$
nonincident.  Symmetrically, in the $11$-layer only the words
$0_T^{\{t\}}$ have two nonincident directions, namely $t$ and $u$.
Every word in the middle $10$-layer has at most one nonincident direction.
Since $|J|\ge2$, inclusion
\eqref{eq:damp-canonical-collar-blocker-nonincident} is therefore an equality
in one of the two displayed cases.  Flipping those two directions gives
respectively $b_1$ and $b_0$, completing the proof.
\end{proof}

\begin{corollary}[No ample interval gap immediately above the canonical collar]
\label{cor:damp-canonical-collar-no-first-gap}
Let $|T|\ge2$ and
$\mathcal K_T\subsetneq B\subseteq Q_{T\cup\{u,\ell\}}$ be ample.  Then
there is an $x\in B\setminus\mathcal K_T$ such that
$\mathcal K_T\cup\{x\}$ is ample.
\end{corollary}

\begin{proof}
Otherwise Proposition~\ref{prop:damp-canonical-collar-two-bad-extensions}
gives
$B\setminus\mathcal K_T\subseteq\{b_0,b_1\}$.  A one-tip extension is
nonample by the same proposition.  If both tips occur, conditioning
$u=0$ and $\ell=1$ leaves the antipodal pair $\{0_T,1_T\}$, which is
nonample because $|T|\ge2$.  Both alternatives contradict ampleness of
$B$.
\end{proof}

\begin{corollary}[High-rank opposite collisions exceed the bulk threshold]
\label{cor:damp-high-rank-opposite-collision-threshold}
Keep the path-spine notation of Proposition
\ref{prop:damp-path-spine-euler-identity}, with $n\ge2$.  If one of
$W,C_1,\ldots,C_n$ contains opposite punctured faces with the same repair
support $T$ and $|T|\ge4$, then $|H|\ge46$.
\end{corollary}

\begin{proof}
Put $k=|T|$.  Proposition
\ref{prop:damp-opposite-puncture-projection-signature} gives
$r-k=|R\setminus T|\ge2$ and says that the exceptional carrier has order at
least
\[
 r+2^{k+1}-k-2=(r+1)+(2^{k+1}-k-3).
\]
Every other member of the list $W,C_1,\ldots,C_n$ is ample and contains the
two repair antipodes, so it has at least $r+1$ vertices.  The disjoint
path-spine shadow decomposition therefore gives
\[
 |H|\ge(n+1)(r+1)+2^{k+1}-k-3.
\]
The stronger shadow estimate
\eqref{eq:damp-opposite-puncture-cross-section-surcharge} says that the
exceptional carrier has tax at least
\[
 2^k-k-1+(2^k-2)(r-k)
\]
above its $r+1$ antipodal-geodesic baseline.  Hence
\[
 |H|\ge(n+1)(r+1)+2^k-k-1+(2^k-2)(r-k).
\]
Using $n\ge2$ and $r\ge k+2$, the right-hand side is at least
\[
 3(k+3)+2^k-k-1+2(2^k-2)=3\cdot2^k+2k+4.
\]
This equals $60$ at $k=4$ and is increasing for $k\ge4$.
\end{proof}

\begin{corollary}[An adjacent-carrier collision pays twice]
\label{cor:damp-adjacent-carrier-collision-double-tax}
Keep the path-spine notation of Proposition
\ref{prop:damp-path-spine-euler-identity}, with $n\ge2$, and suppose that
some $C_i$ contains opposite punctured faces with common support $T$, where
$k=|T|\ge2$.  Put
\begin{equation}
 \tau_k(r)=2^k-k-1+(2^k-2)(r-k).
 \label{eq:damp-opposite-collision-carrier-tax}
\end{equation}
Then
\begin{equation}
 |H|\ge(n+1)(r+1)+2\tau_k(r).
 \label{eq:damp-adjacent-carrier-collision-double-tax}
\end{equation}
Consequently, if $|H|<46$, no $C_i$ contains a collision of rank at least
three.  A rank-two collision in a $C_i$ has $4\le r\le6$ and satisfies
\begin{equation}
 \begin{array}{c|ccc}
 r&4&5&6\\ \hline
 n&\le6&\le4&=2.
 \end{array}
 \label{eq:damp-adjacent-rank-two-collision-profiles}
\end{equation}
Thus every rank-three collision below order $46$ lies in the pure repair
projection $W$, not in an adjacent path carrier.
\end{corollary}

\begin{proof}
The inclusion $C_i\subseteq W$ puts the same two punctured faces in both
ample carriers.  By
\eqref{eq:damp-opposite-puncture-cross-section-surcharge}, each carrier has
order at least $(r+1)+\tau_k(r)$.  Every other member of the disjoint
path-spine list $W,C_1,\ldots,C_n$ has order at least $r+1$.  Summing their
shadows in
\eqref{eq:damp-path-spine-shadow-family} proves
\eqref{eq:damp-adjacent-carrier-collision-double-tax}.

For $k=3$, one has $r\ge5$ and $\tau_3(r)=6r-14$.  Hence
\[
 |H|\ge3(r+1)+2(6r-14)=15r-25\ge50.
\]
The lower bound increases with $k$, so every $k\ge3$ is excluded below
$46$.  For $k=2$, one has $r\ge4$ and $\tau_2(r)=2r-3$, giving
\[
 |H|\ge(n+1)(r+1)+4r-6.
\]
Substitution of $r=4,5,6$ gives the displayed bounds on $n$, while $r=7$
already gives $3\cdot8+22=46$.
\end{proof}

\begin{corollary}[A projection collision forces a global puncture cover]
\label{cor:damp-projection-collision-global-puncture-cover}
Keep the path-spine notation of Proposition
\ref{prop:damp-path-spine-euler-identity}, suppose that $H$ has minimum
cardinality among all nonempty cornerless ample families, and suppose that
$W$ contains opposite punctured faces with common repair support $T$ of
order $k\ge2$.  Then $W\in\Damp$, every nonempty fibre $H_g$ belongs to
$\Damp$, and $H$ contains at least
\begin{equation}
 2^k+(2^k-1)(r-k)
 \label{eq:damp-projection-collision-puncture-count}
\end{equation}
rooted transverse punctures whose repair roots are pairwise distinct.
If $|H|<48$, every one of these punctures has rank between two and five.
Moreover, a fixed support of rank $j$ can own at most
\begin{equation}
 \left\lfloor\frac{47}{2^j-1}\right\rfloor
 =\begin{cases}
 15,&j=2,\\
 6,&j=3,\\
 3,&j=4,\\
 1,&j=5
 \end{cases}
 \label{eq:damp-below-forty-eight-puncture-owner-capacity}
\end{equation}
of these roots.
\end{corollary}

\begin{proof}
Proposition~\ref{prop:damp-opposite-puncture-projection-signature} gives
\[
 |W|\ge2^k+(2^k-1)(r-k).
\]
The projection $W$ is ample.  It is strictly smaller than $H$, because the
endpoint repair fibre contains the path core with at least three concepts.
Lemma~\ref{lem:damp-least-cornerless-cardinality-minimality} gives
$W\in\Damp$, while Corollary
\ref{cor:damp-fibrewise-transverse-puncture-cover} gives the same conclusion
for every nonempty $H_g$.
That corollary also supplies one
rooted transverse puncture over every word of $W$, with distinct repair
roots, and bounds their ranks by five when $|H|<48$.  Finally apply
Lemma~\ref{lem:damp-fixed-support-puncture-disjointness} to all roots that
choose the same support $J$ and use $|H|\le47$.
\end{proof}

\begin{proposition}[Cross-support conservation for the fibrewise punctures]
\label{prop:damp-cross-support-puncture-conservation}
Keep the path-spine notation of Proposition
\ref{prop:damp-path-spine-euler-identity}, and let the rooted transverse
punctures $(g,a_g),J_g$ be chosen as in Corollary
\ref{cor:damp-fibrewise-transverse-puncture-cover}.  Write
\[
 \Omega_{\ge2}=\{S\in\Omega:|S\cap X|\ge2\}.
\]
For $K\subseteq X$, put
\begin{equation}
 m_K=|\{g\in W:J_g\cap X=K\}|.
 \label{eq:damp-puncture-core-owner-count}
\end{equation}
For $1\le i\le n$, split the singleton-core owners into the
spine-aligned and off-carrier counts
\begin{align}
 m_i^{\parallel}
  &=|\{g\in C_i:J_g\cap X=\{x_i\}\}|,\notag\\
 m_i^{\perp}
  &=|\{g\in W\setminus C_i:J_g\cap X=\{x_i\}\}|.
 \label{eq:damp-singleton-core-owner-split}
\end{align}
Then
\begin{equation}
 |\Omega_{\ge2}|
 \ge \sum_{\substack{K\subseteq X\\|K|\ge2}}m_K
 \label{eq:damp-high-core-owner-conservation}
\end{equation}
and, more precisely,
\begin{equation}
 |\Omega|
 \ge
 \sum_{\substack{K\subseteq X\\|K|\ge2}}m_K
 +\sum_{i=1}^n m_i^{\perp}.
 \label{eq:damp-exact-cross-support-owner-conservation}
\end{equation}
In particular,
\begin{equation}
 |\Omega|
 \ge
 \sum_{\substack{K\subseteq X\\|K|\ge2}}m_K
 +\sum_{i=1}^n
   \max\{0,m_{\{x_i\}}-|C_i|\}.
 \label{eq:damp-cross-support-owner-conservation}
\end{equation}
Equivalently, the only puncture anchors that can be absorbed without a
distinct residual support have no core direction, or are spine-aligned
singleton-core owners:
\begin{equation}
 m_\varnothing+
 \sum_{i=1}^n m_i^{\parallel}
 \ge |W|-|\Omega|.
 \label{eq:damp-exact-shadow-neutral-owner-capacity}
\end{equation}
Consequently,
\begin{equation}
 m_\varnothing+
 \sum_{i=1}^n\min\{m_{\{x_i\}},|C_i|\}
 \ge |W|-|\Omega|.
 \label{eq:damp-shadow-neutral-owner-capacity}
\end{equation}
Thus punctures with two or more core directions admit no cross-support
reuse at all: each such repair root pays a different member of the
residual shadow, even when many roots use the same core support.
\end{proposition}

\begin{proof}
Fix a nonempty $K\subseteq X$ and a repair word counted by $m_K$.
Write $J_g=S_g\mathbin{\dot\cup}K$, where
$S_g=J_g\cap R$ is nonempty.  Every subset of $K$ is then a proper
subset of $J_g$.  The puncture relations
\eqref{eq:damp-fibre-corner-transverse-puncture} show that the full
$K$-cube through $(g,a_g)$ lies in $H$.  Hence $g$ occurs in the repair
projection $W_K$ of the $K$-cube carrier.  The repair roots are pairwise
distinct, so $|W_K|\ge m_K$.

Lemma~\ref{lem:damp-projected-cube-carrier-support-conservation} now gives
at least $m_K$ shattered supports with exact core part $K$.  For distinct
sets $K$, these collections are disjoint.  If $|K|\ge2$, every one of
them belongs to $\Omega_{\ge2}$, which proves
\eqref{eq:damp-high-core-owner-conservation}.

Suppose next that $K=\{x_i\}$.  Among the supports supplied by the
projected carrier, the only ones that can belong to the spine family
\eqref{eq:damp-path-spine-shadow-family} are
\[
 T\cup\{x_i\},\qquad T\in\operatorname{Sh}(C_i).
\]
There are $|\operatorname{Sh}(C_i)|=|C_i|$ of these.  Consequently at
least
$\max\{0,m_{\{x_i\}}-|C_i|\}$ of the supplied supports belong to
$\Omega$.  Their exact core parts distinguish different indices $i$ and
also distinguish them from the supports counted in
\eqref{eq:damp-high-core-owner-conservation}.  Summing proves
\eqref{eq:damp-cross-support-owner-conservation}.

The same argument retains the positions of the roots.  Every member of
$C_i$ anchors the path edge $a_{i-1}a_i$, and hence belongs to the repair
projection $W_{\{x_i\}}$ of the $x_i$-cube carrier.  All
$m_i^{\perp}$ off-carrier roots are further, distinct members of that
projection.  Therefore
\[
 |W_{\{x_i\}}|-|C_i|\ge m_i^{\perp}.
\]
Because both nested families are ample, the same difference counts the
supports in
$\operatorname{Sh}(W_{\{x_i\}})\setminus\operatorname{Sh}(C_i)$.
After adjoining $x_i$, all those supports belong to $\Omega$.  These
families are disjoint for different $i$, and disjoint from the high-core
families.  This proves
\eqref{eq:damp-exact-cross-support-owner-conservation}.

Finally the owner classes partition $W$, so
\[
 |W|=m_\varnothing+
      \sum_{i=1}^n(m_i^{\parallel}+m_i^{\perp})+
      \sum_{|K|\ge2}m_K.
\]
Subtracting the right-hand side of
\eqref{eq:damp-exact-cross-support-owner-conservation} from this identity
gives \eqref{eq:damp-exact-shadow-neutral-owner-capacity}.  The bound
$m_i^{\parallel}\le\min\{m_{\{x_i\}},|C_i|\}$ then gives
\eqref{eq:damp-shadow-neutral-owner-capacity}.
\end{proof}

\begin{proposition}[The column cover gives the dual support-capacity law]
\label{prop:damp-column-cover-dual-support-capacity}
Keep the path-spine notation of Proposition
\ref{prop:damp-path-spine-euler-identity}, and choose the column-indexed
punctures $(g_a,a),L_a$ from Corollary
\ref{cor:damp-bidirectional-transverse-puncture-cover}.  For
$S\subseteq R$, put
\begin{equation}
 n_S=|\{a\in Z:L_a\cap R=S\}|,
 \qquad
 b_S=\mathbf 1_{S\in\operatorname{Sh}(W)}
     +\sum_{i=1}^n\mathbf 1_{S\in\operatorname{Sh}(C_i)}.
 \label{eq:damp-column-owner-base-capacity}
\end{equation}
Then
\begin{equation}
 |\Omega|
 \ge
 \sum_{\varnothing\ne S\subseteq R}
       \max\{0,n_S-b_S\}.
 \label{eq:damp-column-owner-residual-capacity}
\end{equation}
Equivalently, the column owners that can be absorbed by the path-spine
shadow satisfy
\begin{equation}
 n_\varnothing+
 \sum_{\varnothing\ne S\subseteq R}\min\{n_S,b_S\}
 \ge |Z|-|\Omega|.
 \label{eq:damp-column-owner-neutral-capacity}
\end{equation}
Together with Proposition
\ref{prop:damp-cross-support-puncture-conservation}, this gives the
bidirectional bound
\begin{equation}
 |\Omega|\ge
 \max\left\{
  \sum_{\substack{K\subseteq X\\|K|\ge2}}m_K+
       \sum_{i=1}^n m_i^\perp,
  \quad
  \sum_{\varnothing\ne S\subseteq R}\max\{0,n_S-b_S\}
 \right\}.
 \label{eq:damp-bidirectional-owner-capacity}
\end{equation}
Thus a residual support may pay one row demand and one column demand at
the same Boolean-matrix entry.  Absorption by the spine is recorded by
$b_S$, while every remaining row--column reuse is localized to the exact
pair $(S,K)$ of that residual support.
\end{proposition}

\begin{proof}
Fix a nonempty $S\subseteq R$ with $n_S>0$, and define the repair-cube
carrier and its core projection by
\[
 \mathcal K_S(H)
 =\{y\in Q_{(R\mathbin{\dot\cup}X)\setminus S}:
       Q_S\times\{y\}\subseteq H\},
 \qquad
 Z_S=\mathcal K_S(H)|_X.
\]
If $a$ is counted by $n_S$, write
$L_a=S\mathbin{\dot\cup}K_a$.  Proposition
\ref{prop:damp-column-corner-core-transverse-puncture} gives
$K_a\ne\varnothing$.  Hence every subset of $S$, including $S$ itself,
is a proper subset of $L_a$.  The puncture relations show that the full
$S$-cube through $(g_a,a)$ lies in $H$, so $a\in Z_S$.  Distinct column
indices give distinct words $a$, and therefore
\begin{equation}
 |Z_S|\ge n_S.
 \label{eq:damp-column-owner-carrier-order}
\end{equation}

Apply Lemma
\ref{lem:damp-projected-cube-carrier-support-conservation} after
interchanging $R$ and $X$.  It gives the exact row--column identity
\begin{equation}
 \operatorname{Sh}(Z_S)
 =\{K\subseteq X:S\mathbin{\dot\cup}K
                    \in\operatorname{Sh}(H)\}.
 \label{eq:damp-repair-carrier-core-shadow-row}
\end{equation}
Thus the number of shattered supports of $H$ whose exact repair part is
$S$ equals $|\operatorname{Sh}(Z_S)|=|Z_S|$.

The members of the spine family
\eqref{eq:damp-path-spine-shadow-family} with exact repair part $S$ are
precisely $S$ itself when $S\in\operatorname{Sh}(W)$ and the supports
$S\cup\{x_i\}$ for which $S\in\operatorname{Sh}(C_i)$.  Their number is
exactly $b_S$.  Consequently the residual shadow has exactly
$|Z_S|-b_S$ members with repair part $S$, and
\eqref{eq:damp-column-owner-carrier-order} makes this at least
$\max\{0,n_S-b_S\}$.  Supports with distinct exact repair parts are
different, so summing over nonempty $S$ proves
\eqref{eq:damp-column-owner-residual-capacity}.

The owner classes partition $Z$, giving
\[
 |Z|=n_\varnothing+
      \sum_{\varnothing\ne S\subseteq R}n_S.
\]
Subtracting the right-hand side of
\eqref{eq:damp-column-owner-residual-capacity} proves
\eqref{eq:damp-column-owner-neutral-capacity}.  Finally combine
\eqref{eq:damp-column-owner-residual-capacity} with
\eqref{eq:damp-exact-cross-support-owner-conservation} to obtain
\eqref{eq:damp-bidirectional-owner-capacity}.
\end{proof}

\begin{corollary}[Forced row--column collisions in the residual shadow]
\label{cor:damp-forced-row-column-residual-collisions}
In the setting of Proposition
\ref{prop:damp-column-cover-dual-support-capacity}, put
\begin{align*}
 \rho&=
  \sum_{\substack{K\subseteq X\\|K|\ge2}}m_K+
  \sum_{i=1}^n m_i^\perp,\\
 \kappa&=
  \sum_{\varnothing\ne S\subseteq R}\max\{0,n_S-b_S\}.
\end{align*}
There are subsets $\mathcal R,\mathcal C\subseteq\Omega$ with
\[
 |\mathcal R|=\rho,
 \qquad
 |\mathcal C|=\kappa,
\]
such that every member of $\mathcal R$ pays one row-owner demand with its
exact core part and every member of $\mathcal C$ pays one column-owner
demand with its exact repair part.  Consequently
\begin{equation}
 |\mathcal R\cap\mathcal C|
 \ge \max\{0,\rho+\kappa-|\Omega|\}.
 \label{eq:damp-forced-row-column-residual-collisions}
\end{equation}
Every member of the intersection is a shattered support
$S\mathbin{\dot\cup}K$ that is simultaneously selected by the projected
$K$-carrier and the projected $S$-carrier.  Hence all unavoidable reuse of
the two transverse puncture covers is localized in common full
$(S\mathbin{\dot\cup}K)$-cubes of $H$.
\end{corollary}

\begin{proof}
The proof of Proposition
\ref{prop:damp-cross-support-puncture-conservation} selects, for each
$|K|\ge2$, at least $m_K$ residual supports with exact core part $K$, and
for $K=\{x_i\}$ at least $m_i^\perp$ members of the relative carrier
shadow
\[
 \{T\cup\{x_i\}:T\in
   \operatorname{Sh}(W_{\{x_i\}})
       \setminus\operatorname{Sh}(C_i)\}.
\]
These groups have different exact core parts, so their selected members
form a set $\mathcal R$ of order $\rho$.

Likewise, the proof of Proposition
\ref{prop:damp-column-cover-dual-support-capacity} supplies, for each
nonempty $S$, at least $\max\{0,n_S-b_S\}$ residual supports with exact
repair part $S$.  The groups for different $S$ are disjoint, and selected
members form a set $\mathcal C$ of order $\kappa$.
Both sets lie in $\Omega$, so inclusion--exclusion gives
\eqref{eq:damp-forced-row-column-residual-collisions}.

If $S\mathbin{\dot\cup}K\in\mathcal R\cap\mathcal C$, ampleness of $H$
makes this support strongly shattered.  A full cube on that support is
simultaneously an $S$-cube in the $K$-carrier and a $K$-cube in the
$S$-carrier, which proves the final assertion.
\end{proof}

\begin{lemma}[The shadow-neutral defects reduce to square tips]
\label{lem:damp-shadow-neutral-square-tip-reduction}
Let $v,J$ be a rooted puncture in $H$:
\[
 v^U\in H\quad(U\subsetneq J),
 \qquad v^J\notin H,
 \qquad |J|\ge2.
\]
For every two-set $L\subseteq J$, the vertex
$u=v^{J\setminus L}$ is the root of a punctured $L$-square:
\begin{equation}
 u,u^{\{p\}},u^{\{q\}}\in H,
 \qquad u^L\notin H
 \qquad(L=\{p,q\}).
 \label{eq:damp-terminal-punctured-square}
\end{equation}

In particular, keep the path-spine setting of Proposition
\ref{prop:damp-cross-support-puncture-conservation} and write $v=(g,a)$.
\begin{enumerate}
 \item A repair-only puncture contains a repair-only punctured square in
       the column $P_a$.
 \item Suppose that $J=S\mathbin{\dot\cup}\{x_i\}$ with
       $\varnothing\ne S\subseteq R$, and choose $s\in S$.  The
       punctured square on $\{s,x_i\}$ has repair root
       $h=g^{S\setminus\{s\}}$, and $h$ belongs to the repair projection
       $W_{\{x_i\}}$ of the $x_i$-cube carrier.  In the full carrier, its
       anchor $y$ is present and the $s$-neighbour $y^{\{s\}}$ is absent.
       Either $h\in C_i$, so
       this is a spine-aligned rank-two defect, or $H$ has a residual
       shattered support with exact core part $\{x_i\}$.
\end{enumerate}
Consequently the local shadow-neutral transport problem has only two
terminal defect types: a punctured repair square, or a punctured
$\{s,x_i\}$-square rooted over $C_i$.  Both are boundary data inside
smaller dismantlable carriers, by Corollary
\ref{cor:damp-least-obstruction-core-carriers-dismantlable}.  Each square
is retained together with the parent puncture $(v,J)$ from which it was
obtained.
\end{lemma}

\begin{proof}
If $A\subsetneq L$, then
$(J\setminus L)\cup A$ is a proper subset of $J$.  Hence
\[
 u^A=v^{(J\setminus L)\cup A}\in H.
\]
On the other hand, $u^L=v^J\notin H$.  This proves
\eqref{eq:damp-terminal-punctured-square} and immediately gives the first
special case.

For the second, take $L=\{s,x_i\}$.  The root $u$ has repair word
$h=g^{S\setminus\{s\}}$ and the same core word as $v$.  The two vertices
$u$ and $u^{\{x_i\}}$ form an $x_i$-edge of $H$, so
$h\in W_{\{x_i\}}$.  If $y$ denotes the word obtained from $u$ by
deleting its $x_i$-coordinate, then
$y\in\mathcal K_{\{x_i\}}(H)$, whereas
$y^{\{s\}}\notin\mathcal K_{\{x_i\}}(H)$ because the latter anchor would
require the absent vertex $u^{\{s,x_i\}}$.  If $h\notin C_i$, then
$C_i\subsetneq W_{\{x_i\}}$.  The family $C_i$ is ample by Proposition
\ref{prop:damp-path-spine-euler-identity}, while
$W_{\{x_i\}}$ is ample by Lemma
\ref{lem:damp-projected-cube-carrier-support-conservation}; therefore
\[
 \operatorname{Sh}(C_i)
 \subsetneq\operatorname{Sh}(W_{\{x_i\}}).
\]
Choose $T$ in the difference.  Lemma
\ref{lem:damp-projected-cube-carrier-support-conservation} makes
$T\cup\{x_i\}$ a shattered support of $H$, while its absence from
$\operatorname{Sh}(C_i)$ keeps it out of the spine family
\eqref{eq:damp-path-spine-shadow-family}.  Thus it belongs to $\Omega$ and
has exact core part $\{x_i\}$.  The alternative $h\in C_i$ is precisely
the asserted spine-aligned square.
\end{proof}

\begin{remark}[The terminal squares retain their provenance]
\label{rem:damp-shadow-neutral-square-tip-provenance}
The map $(v,J,L)\mapsto v^{J\setminus L}$ need not be injective when the
parent supports $J$ differ.  Thus the square-tip reduction does not replace
the distinct-root statement in Corollary
\ref{cor:damp-fibrewise-transverse-puncture-cover} by a new distinct-root
statement.  Its output is a family of provenance-tagged square tips.  The
remaining transport theorem must either separate coincident tips by their
carrier provenance, charging a new residual support, or show that the
coincidence survives a contraction.  This is precisely the synchronization
issue; forgetting the parent punctures would discard the global
multiplicity information.
\end{remark}

\begin{proposition}[Isometric path-fibre gap tax]
\label{prop:damp-isometric-path-fibre-gap-tax}
Keep the notation of Proposition
\ref{prop:damp-path-spine-euler-identity}.  For $g\in U$, write the
connected components of
\[
 I(g)=\{i:g\in P_i\}
\]
as index intervals
\[
 [\ell_1,r_1],\ldots,[\ell_c,r_c],
 \qquad r_s+1<\ell_{s+1},
\]
and put
\[
 \lambda(g)=\sum_{s=1}^{c-1}
  \bigl(\ell_{s+1}-r_s-1\bigr),
 \qquad
 \Lambda=\sum_{g\in U}\lambda(g).
\]
Then
\begin{equation}
 V\ge |W\setminus U|+\Lambda,
 \qquad
 \Lambda\ge\Phi.
 \label{eq:damp-isometric-path-fibre-gap-tax}
\end{equation}
Consequently the residual shadow satisfies the doubled fragmentation
bound
\begin{equation}
 |\Omega|\ge\Phi+\Lambda\ge2\Phi.
 \label{eq:damp-doubled-path-fragmentation-tax}
\end{equation}
More precisely, every gap
$r_s+1,\ldots,\ell_{s+1}-1$ forces that many distinct external core
concepts over $g$, and the concepts forced by different gaps are
disjoint.
\end{proposition}

\begin{proof}
Fix $g\in U$ and two consecutive components of $I(g)$.  Put
$p=r_s$ and $q=\ell_{s+1}$.  The core fibre
\[
 H_g=\{z\in Q_X:(g,z)\in H\}
\]
is a nonempty conditioning of the ample family $H$, and hence is
isometric.  It contains $a_p$ and $a_q$, whose Hamming distance is
$q-p$.  A shortest $H_g$-path between them therefore has $q-p-1$
internal vertices and changes exactly the coordinates
$x_{p+1},\ldots,x_q$.

No internal vertex of this path belongs to $A$.  Indeed, the only members
of the core path $A$ in the interval subcube between $a_p$ and $a_q$ are
$a_p,a_{p+1},\ldots,a_q$, while none of the intermediate concepts
$a_{p+1},\ldots,a_{q-1}$ belongs to $H_g$ by the choice of consecutive
components.  Thus all $q-p-1$ internal vertices are external core
concepts over $g$.

The open coordinate intervals
\[
 \{x_{r_s+1},\ldots,x_{\ell_{s+1}}\}
\]
for different gaps are disjoint.  Their interval subcubes can meet only
at a common endpoint belonging to $A$, so their internal vertices are
disjoint.  Summing the forced internal vertices over the gaps proves that
at least $\lambda(g)$ external core concepts lie over $g$.

For a repair word in $W\setminus U$, every vertex over it has core word
outside $A$, so at least one external concept is present.  Summing over
all repair words gives the first inequality in
\eqref{eq:damp-isometric-path-fibre-gap-tax}.  Every one of the
$c(g)-1$ gaps has positive length, so
$\lambda(g)\ge c(g)-1$; summing gives $\Lambda\ge\Phi$.  Finally,
Proposition~\ref{prop:damp-path-spine-euler-identity} gives
\[
 |\Omega|=\Phi+V-|W\setminus U|
 \ge\Phi+\Lambda\ge2\Phi,
\]
which proves \eqref{eq:damp-doubled-path-fragmentation-tax}.
\end{proof}

\begin{proposition}[Transverse path-spine decomposition]
\label{prop:damp-transverse-path-spine-decomposition}
Keep the notation of Proposition
\ref{prop:damp-path-spine-euler-identity}.  Orient the path coordinates so
that
\[
 a_i=a_0^{\{x_1,\ldots,x_i\}}\qquad(0\le i\le n).
\]
For $1\le i\le n$ and $b\in\{0,1\}$, put
\[
 B_i^b=
 \{g\in Q_R:\text{some }(g,z)\in H\text{ has }z_{x_i}=b\},
 \qquad
 D_i=B_i^0\cap B_i^1.
\]
Then $D_i$ is ample, $C_i\subseteq D_i$, and the residual shadow in
\eqref{eq:damp-path-spine-euler-identity} decomposes as
\begin{equation}
 \Omega=
 \mathop{\dot\bigcup}_{i=1}^n
 \{T\cup\{x_i\}:
     T\in\operatorname{Sh}(D_i)\setminus\operatorname{Sh}(C_i)\}
 \mathbin{\dot\cup}\Omega_{\ge2},
 \label{eq:damp-transverse-path-spine-decomposition}
\end{equation}
where every member of $\Omega_{\ge2}$ contains at least two coordinates
of the core path.  Consequently
\begin{equation}
 \Phi+V-|W\setminus U|
 =\sum_{i=1}^n|D_i\setminus C_i|+|\Omega_{\ge2}|.
 \label{eq:damp-transverse-path-spine-count}
\end{equation}

Moreover, let $g\in P_0$ and suppose that $H_g$ contains a core neighbour
$a_0^{\{x_j\}}$ outside $A$.  Then $j\ge2$, and either
\[
 g\in D_j\setminus C_j
\]
or $\{x_1,x_j\}\in\Omega_{\ge2}$.  Symmetrically, an external core
neighbour $a_n^{\{x_j\}}$ of $a_n$ either belongs over a word in
$D_j\setminus C_j$ or forces
$\{x_j,x_n\}\in\Omega_{\ge2}$.
\end{proposition}

\begin{proof}
Project $H$ onto $R\cup\{x_i\}$.  This projection is ample and its two
$x_i$-sections are $B_i^0$ and $B_i^1$.  Lemma
\ref{lem:damp-ample-two-layer-shadow-intersection} makes
$D_i=B_i^0\cap B_i^1$ ample and shows that the shattered supports of the
projection containing $x_i$ are exactly
\[
 \{T\cup\{x_i\}:T\in\operatorname{Sh}(D_i)\}.
\]
Every repair word supporting both adjacent path concepts
$a_{i-1},a_i$ lies in $D_i$, so $C_i\subseteq D_i$.  The family $C_i$ is
ample by the same two-layer lemma applied after fixing the other core
coordinates.  Therefore the one-core-coordinate members of $\Omega$ are
exactly the first family on the right of
\eqref{eq:damp-transverse-path-spine-decomposition}; all remaining
members use at least two core coordinates.  The union is disjoint.
Taking cardinalities, using ampleness of $C_i,D_i$, and applying
\eqref{eq:damp-path-spine-euler-identity} gives
\eqref{eq:damp-transverse-path-spine-count}.

For the final assertion, the only neighbour of $a_0$ in the core path is
$a_1=a_0^{\{x_1\}}$, so an external core neighbour has $j\ge2$.  The two
words $a_0$ and $a_0^{\{x_j\}}$ over $g$ put $g$ in $D_j$.  If
$g\notin C_j$, the first alternative holds.  If $g\in C_j$, then $H_g$
also contains $a_{j-1}$ and $a_j$.  Their traces on
$\{x_1,x_j\}$, together with the traces of
$a_0,a_0^{\{x_j\}}$, are respectively all four binary words.  Hence
$H$ shatters $\{x_1,x_j\}$.  This pair is not in the spine family
\eqref{eq:damp-path-spine-shadow-family}, so it belongs to
$\Omega_{\ge2}$.  The proof at $a_n$ is the same after reversing the
path: the four words
$a_{j-1},a_j,a_n^{\{x_j\}},a_n$ realize all traces on
$\{x_j,x_n\}$.
\end{proof}

\begin{proposition}[Endpoint trapped-wing-or-escape dichotomy]
\label{prop:damp-path-endpoint-trapped-wing-or-escape}
Keep the path-spine notation of Proposition
\ref{prop:damp-path-spine-euler-identity}, and suppose that \(H\) has
minimum cardinality among all nonempty cornerless ample families.  At the
left endpoint, at least one of the following two alternatives holds:
\begin{enumerate}
\item \(C_1\subsetneq P_0\), every corner of \(P_0\) belongs to \(C_1\),
  and
  \begin{equation}
   |P_0\setminus C_1|\ge\rho_{\Damp}\ge10;
   \label{eq:damp-left-endpoint-trapped-wing-tax}
  \end{equation}
\item there are \(g\in P_0\) and \(j\in\{2,\ldots,n\}\) such that
  \((g,a_0^{\{x_j\}})\in H\).
\end{enumerate}
The symmetric alternatives hold at the right endpoint, with \(C_n\),
\(P_n\), and an external neighbour \(a_n^{\{x_j\}}\),
\(1\le j\le n-1\).

Consequently, if
\(\Delta(H;A,R)<13\), then both endpoints have external core neighbours.
If both endpoints pay the trapped-wing tax, then
\(\Delta(H;A,R)\ge20\).
\end{proposition}

\begin{proof}
The families \(P_0\) and \(C_1=P_0\cap P_1\) are ample.  They are proper
pc-minors of \(H\), and the minimum-cardinality hypothesis therefore puts
both of them in \(\Damp\).  The two opposite repair fibres of \(H\) belong
to every \(P_i\), so \(C_1\) contains two repair antipodes.  Isometry gives
\begin{equation}
 |C_1|\ge r+1.
 \label{eq:damp-endpoint-intersection-antipodal-order}
\end{equation}

Suppose first that \(C_1\subsetneq P_0\) and
\(\operatorname{Cor}(P_0)\subseteq C_1\).  The pair
\(C_1\subsetneq P_0\) occurs in the minimum defining
\(\rho_{\Damp}\), and hence
\[
 |P_0\setminus C_1|\ge\rho_{\Damp}\ge10.
\]
This is the first alternative.

Otherwise choose a corner \(g\) of \(P_0\) outside \(C_1\) when the
inclusion is proper, and choose any corner \(g\) of \(P_0\) when
\(C_1=P_0\).  Let \(K\) be the unique maximal repair cube of \(P_0\)
through \(g\).  In the first case the core edge \(a_0a_1\) is absent over
\(g\), while in the second case the incident cube in that direction is
\(K\times a_0a_1\).  Because \(a_0\) is an endpoint of the core path,
these account for every incident core direction whose other endpoint lies
in \(A\).  If no core neighbour of \(a_0\) outside \(A\) occurred over
\(g\), the local corner criterion would therefore make \((g,a_0)\) a
corner of \(H\), a contradiction.  Such a neighbour must have the form
\(a_0^{\{x_j\}}\) with \(j\ge2\), proving the second alternative.  Reversing
the path proves the assertion at \(a_n\).

It remains to account for the excess.  In a trapped alternative,
\eqref{eq:damp-endpoint-intersection-antipodal-order} gives
\[
 |P_0|-(r+1)\ge|P_0\setminus C_1|\ge10.
\]
In an escape alternative, Proposition
\ref{prop:damp-maximum-core-corner-surcharge} gives endpoint-section
excess at least two, and the escaping concept contributes at least one
vertex to \(H\setminus(Q_R\times A)\).  Thus two trapped endpoints cost at
least \(20\), while one trapped and one escaping endpoint cost at least
\(10+2+1=13\) in the exact decomposition
\eqref{eq:damp-geodesic-excess-decomposition}.  Therefore excess below
\(13\) forces an escape at each endpoint.
\end{proof}

\begin{proposition}[Exact two-coordinate repair calculus]
\label{prop:damp-two-coordinate-maximum-core-repair}
Let $u,v$ be two coordinates, let $A$ be ample, and let $B,C\subseteq A$
be nonempty ample families with $B\cup C=A$.  Define
\begin{equation}
 H_{00}=H_{11}=A,\qquad H_{01}=B,\qquad H_{10}=C.
 \label{eq:damp-two-coordinate-core-fibres}
\end{equation}
Then
\begin{equation}
 \operatorname{Sh}(H)
 =\operatorname{Sh}(A)
  \mathbin{\dot\cup}
  \{S\cup\{u\}:S\in\operatorname{Sh}(A)\}
  \mathbin{\dot\cup}
  \{S\cup\{v\}:S\in\operatorname{Sh}(A)\}
  \mathbin{\dot\cup}
  \{S\cup\{u,v\}:S\in
       \operatorname{Sh}(B)\cap\operatorname{Sh}(C)\}.
 \label{eq:damp-two-coordinate-core-shadow}
\end{equation}
Consequently, $H$ is ample if and only if
\begin{equation}
 B\cap C\text{ is ample}\qquad\text{and}\qquad
 \operatorname{Sh}(B)\cap\operatorname{Sh}(C)
 =\operatorname{Sh}(B\cap C).
 \label{eq:damp-two-coordinate-compatible-shadow-cover}
\end{equation}
Equivalently,
\begin{equation}
 \operatorname{Sh}(B)\cup\operatorname{Sh}(C)
 =\operatorname{Sh}(A).
 \label{eq:damp-two-coordinate-shadow-union-cover}
\end{equation}
In that case
\begin{equation}
 |H|=3|A|+|B\cap C|.
 \label{eq:damp-two-coordinate-core-order}
\end{equation}

Suppose now that $H$ is ample.  If $a$ is a corner of $A$, denote its
unique maximal cube by $Q_A(a)$ and put
\[
 K_a=
 \begin{cases}
  B,&a\in B\setminus C,\\
  C,&a\in C\setminus B,\\
  B\cap C,&a\in B\cap C.
 \end{cases}
\]
Then both endpoint copies $(00,a)$ and $(11,a)$ are corners of $H$ if
and only if
\begin{equation}
 Q_A(a)\subseteq K_a.
 \label{eq:damp-two-coordinate-surviving-corner}
\end{equation}
\end{proposition}

\begin{proof}
Let $S$ use only coordinates of $A$.  Since the two endpoint fibres in
\eqref{eq:damp-two-coordinate-core-fibres} are both $A$, the sets $S$,
$S\cup\{u\}$, and $S\cup\{v\}$ are shattered by $H$ exactly when $S$
is shattered by $A$.  The set $S\cup\{u,v\}$ is shattered exactly when
all four fibres realize every trace on $S$; because $B,C\subseteq A$,
this is equivalent to
$S\in\operatorname{Sh}(B)\cap\operatorname{Sh}(C)$.  This proves
\eqref{eq:damp-two-coordinate-core-shadow}.

Since $A,B,C$ are ample and $B\cup C=A$,
\[
 |H|=2|A|+|B|+|C|=3|A|+|B\cap C|.
\]
Comparing this identity with
\eqref{eq:damp-two-coordinate-core-shadow} shows that $H$ is ample
exactly when
\[
 |\operatorname{Sh}(B)\cap\operatorname{Sh}(C)|=|B\cap C|.
\]
Now
\[
 \operatorname{Sh}(B\cap C)
 \subseteq\operatorname{Sh}(B)\cap\operatorname{Sh}(C),
\]
whereas the Sandwich inequality gives
$|\operatorname{Sh}(B\cap C)|\geq|B\cap C|$.  Equality therefore holds
throughout, which proves
\eqref{eq:damp-two-coordinate-compatible-shadow-cover}.
Inclusion--exclusion then gives
\[
 |\operatorname{Sh}(B)\cup\operatorname{Sh}(C)|
 =|B|+|C|-|B\cap C|=|A|.
\]
The union on the left is contained in $\operatorname{Sh}(A)$, which also
has size $|A|$, proving
\eqref{eq:damp-two-coordinate-shadow-union-cover}.  Conversely, that
shadow cover and inclusion--exclusion give
$|\operatorname{Sh}(B)\cap\operatorname{Sh}(C)|=|B\cap C|$, so the
Sandwich inequality gives
\eqref{eq:damp-two-coordinate-compatible-shadow-cover}.  Finally,
\eqref{eq:damp-two-coordinate-core-order} has already been established.

It remains to inspect a core corner.  At either endpoint copy of $a$, the
old-coordinate edge directions generate $Q_A(a)$.  If
$a\in B\setminus C$, there is exactly one incident repair direction and
the cube generated by all incident directions is present precisely when
the corresponding second copy of $Q_A(a)$ lies in $B$, that is, precisely
when $Q_A(a)\subseteq B$.  The case $a\in C\setminus B$ is symmetric.
If $a\in B\cap C$, both repair directions are incident and they generate
the full repair square; its product with $Q_A(a)$ is present precisely
when $Q_A(a)\subseteq B\cap C$.  The local corner criterion now proves
\eqref{eq:damp-two-coordinate-surviving-corner} at both endpoints.
\end{proof}

\begin{proposition}[Unrestricted two-coordinate shadow identity]
\label{prop:damp-unrestricted-two-coordinate-repair}
Let $u,v$ be two coordinates and suppose that an ample family $H$ has
fibres
\begin{equation}
 H_{00}=H_{11}=A,\qquad H_{01}=B,\qquad H_{10}=C,
 \label{eq:damp-unrestricted-two-coordinate-fibres}
\end{equation}
where $A$ is nonempty and $A\subseteq B\cup C$.  Put $K=B\cup C$.  Then
\begin{equation}
 D:=B\cap C\subseteq A,
 \qquad D\text{ is ample},
 \label{eq:damp-unrestricted-two-coordinate-overlap}
\end{equation}
and
\begin{align}
 \operatorname{Sh}(A\cup B)\cap\operatorname{Sh}(A\cup C)
   &=\operatorname{Sh}(A),
 \label{eq:damp-unrestricted-two-coordinate-single-link}\\
 \operatorname{Sh}(A)\cap\operatorname{Sh}(B)
       \cap\operatorname{Sh}(C)
   &=\operatorname{Sh}(D).
 \label{eq:damp-unrestricted-two-coordinate-double-link}
\end{align}
Consequently
\begin{equation}
\begin{split}
 \operatorname{Sh}(H)
 ={}&\operatorname{Sh}(K)\\
 &\mathbin{\dot\cup}
   \{S\cup\{u\},S\cup\{v\}:S\in\operatorname{Sh}(A)\}\\
 &\mathbin{\dot\cup}
   \{S\cup\{u,v\}:S\in\operatorname{Sh}(D)\},
\end{split}
 \label{eq:damp-unrestricted-two-coordinate-shadow}
\end{equation}
and
\begin{equation}
 |H|=|K|+2|A|+|D|.
 \label{eq:damp-unrestricted-two-coordinate-order}
\end{equation}
Thus ampleness forces the external parts of the two middle fibres to be
disjoint; it does not merely constrain their cardinalities.
\end{proposition}

\begin{proof}
The families $A,B,C,K,A\cup B$, and $A\cup C$ are conditionings or
contractions of $H$, and hence are ample.  For a set $S$ of old
coordinates, direct inspection of the four fibres gives the disjoint
decomposition
\[
 \operatorname{Sh}(H)
 =\operatorname{Sh}(K)
 \mathbin{\dot\cup}
 \{S\cup\{u\},S\cup\{v\}:S\in I\}
 \mathbin{\dot\cup}
 \{S\cup\{u,v\}:S\in J\},
\]
where
\[
 I=\operatorname{Sh}(A\cup B)\cap\operatorname{Sh}(A\cup C),
 \qquad
 J=\operatorname{Sh}(A)\cap\operatorname{Sh}(B)
                 \cap\operatorname{Sh}(C).
\]
Let $W=B\cap C$ and $D_0=A\cap W$.  Because
$(A\cup B)\cup(A\cup C)=K$ and all three unions are ample,
inclusion--exclusion gives
\begin{align*}
 |I|
 &=|A\cup B|+|A\cup C|
   -|\operatorname{Sh}(A\cup B)\cup
       \operatorname{Sh}(A\cup C)|\\
 &\ge |A\cup B|+|A\cup C|-|K|
  =|A\cup W|=|A|+|W|-|D_0|.
\end{align*}
Moreover, $\operatorname{Sh}(D_0)\subseteq J$, so the Sandwich inequality
gives $|J|\ge|D_0|$.

On the other hand, ampleness of $H$ and $K$ and the ordinary
inclusion--exclusion identity $|B|+|C|=|K|+|W|$ give
\[
 2|I|+|J|
 =|H|-|K|
 =2|A|+|W|.
\]
Combining the last three displays yields
\[
 2|A|+|W|
 \ge2(|A|+|W|-|D_0|)+|D_0|
 =2|A|+2|W|-|D_0|.
\]
Thus $|W|\le|D_0|$.  Since $D_0\subseteq W$, equality holds and
$W=D_0\subseteq A$.  We may therefore write $D=W$.  Now
$\operatorname{Sh}(A)\subseteq I$ and
$\operatorname{Sh}(D)\subseteq J$, while
\[
 2|I|+|J|=2|A|+|D|.
\]
The Sandwich inequality for $D$ forces equality in both inclusions:
$I=\operatorname{Sh}(A)$, $J=\operatorname{Sh}(D)$, and
$|\operatorname{Sh}(D)|=|D|$.  Hence $D$ is ample.  Substitution gives
\eqref{eq:damp-unrestricted-two-coordinate-shadow}; the order identity
follows either from that formula or directly from the four fibres.
\end{proof}

\begin{theorem}[Two-coordinate repair over a relatively peelable core]
\label{thm:damp-two-coordinate-low-vc-core-corner}
In Proposition~\ref{prop:damp-unrestricted-two-coordinate-repair}, suppose
that $A$ has a corner and that every nonempty proper ample subfamily
$D'\subsetneq A$ omits some corner of $A$.  Then $H$ has a corner.

The hypothesis holds, in particular, if either
\begin{enumerate}
 \item $\operatorname{VCdim}(A)\le2$, or
 \item after deleting constant coordinates, $A$ uses at most four
       coordinates.
\end{enumerate}
Thus no cornerless ample family can have such a core separated by two
repair coordinates, regardless of whether the repair is internal or
external.
\end{theorem}

\begin{proof}
By Proposition~\ref{prop:damp-unrestricted-two-coordinate-repair}, the
overlap $D=B\cap C$ is an ample subfamily of $A$.  We first choose a corner
$a$ of $A$ such that either $D=A$ or $a\notin D$.  If $D$ is a nonempty
proper subfamily, this is exactly the relative-corner hypothesis.  If $D$
is empty or $D=A$, use any corner of $A$.  Denote the unique maximal cube
of $A$ through $a$ by $Q_A(a)$.

Suppose first that $D=A$.  Every vertex of $Q_A(a)$ belongs to all four
fibres in \eqref{eq:damp-unrestricted-two-coordinate-fibres}.  Hence
$Q_{\{u,v\}}\times Q_A(a)$ lies in $H$.  At the endpoint copy $(00,a)$,
the old incident directions are exactly those of $Q_A(a)$, because $a$ is
a corner of $A$, and both repair directions are incident.  All incident
directions therefore generate the displayed cube, and the local corner
criterion makes $(00,a)$ a corner of $H$.

It remains to consider $a\notin D$.  The carrier family of the copies of
$Q_A(a)$ in $H$ is ample and isometric, by the carrier argument in the
proof of Proposition~\ref{prop:damp-maximum-vc-three-carrier-trees}.  It
contains the anchors of $Q_A(a)$ in the $00$- and $11$-fibres.  Those two
anchors differ precisely in $u$ and $v$, so a shortest carrier path between
them starts with a $u$- or a $v$-edge.  Equivalently, the copy of
$Q_A(a)$ in the $00$-fibre extends to a cube in one repair direction.

Since $A\subseteq B\cup C$ and $a\notin B\cap C$, the concept $a$ belongs
to exactly one middle fibre.  Thus $(00,a)$ has exactly one repair neighbor,
and the carrier path says that this repair direction together with all old
incident directions spans the extension of $Q_A(a)$.  There are no other
old directions because $a$ is a corner of $A$.  Hence all incident
directions at $(00,a)$ generate a single cube, so $(00,a)$ is a corner of
$H$.

It remains to verify the two sufficient conditions.  For VC dimension at
most two, translate a member of a nonempty $D'\subseteq A$ to the empty
set and use the relative extension theorem of M\'esz\'aros and R\'onyai
\cite[Theorem~4.1 and Lemma~4.1]{meszaros2014vc2}.  The last concept in a
valid one-concept construction from $D'$ to $A$ is a corner of $A$ outside
$D'$.  Taking $D'$ to be a singleton also shows that $A$ has a corner.
For at most four active coordinates, the same conclusions follow from
Theorem~\ref{thm:damp-nested-ample-relative-peeling-dimension-four}.
\end{proof}

\begin{proposition}[Maximum-core corner shields]
\label{prop:damp-maximum-core-corner-shield}
In Proposition~\ref{prop:damp-two-coordinate-maximum-core-repair}, suppose
that $A$ is maximum of VC dimension $k$, put $D=B\cap C$, and assume that
$H$ is ample.  Then every maximal cube of $A$ has dimension $k$, and for
each $k$-set $S$ there is exactly one $S$-cube in $A$.  Consequently,
every corner $a\in A\setminus D$ satisfies
\eqref{eq:damp-two-coordinate-surviving-corner}.

If $H$ is cornerless, it follows that
\begin{equation}
 \operatorname{corn}(A)\subseteq D,
 \qquad
 Q_A(a)\setminus D\neq\varnothing
 \quad(a\in\operatorname{corn}(A)).
 \label{eq:damp-maximum-core-corner-shield}
\end{equation}
Moreover, every corner cube $Q_A(a)$ is contained wholly in exactly one
of $B,C$.  Thus $A\setminus D$ meets every corner cube although it
contains no core corner.

If, in addition, $H$ has minimum cardinality among nonempty cornerless
ample families and $|H|<299$, then $D$, $B$, and $C$ are nonempty proper
subfamilies of $A$.  In particular, the internal two-coordinate branch is
a genuine two-sided cover by smaller ample families, not a one-sided
repair with $B=A$ or $C=A$.
\end{proposition}

\begin{proof}
The complete-cube and reduction theorems for maximum classes say that $A$
is the union of a complete collection of $k$-cubes and that reduction in
the directions of a $t$-cube is maximum of VC dimension $k-t$
\cite[Section~2]{RubinsteinRubinsteinBartlett14embeddings}.  The vertex of
the reduction representing that $t$-cube therefore lies in a
$(k-t)$-cube; lifting it shows that the original cube lies in a $k$-cube.
Thus every maximal cube has dimension $k$.  For a $k$-set $S$, the same
simultaneous reduction has VC dimension zero, so it has one concept.
Equivalently, $A$ has exactly one $S$-cube.

Let $a\in\operatorname{corn}(A)\setminus D$, and let $S$ support its
unique maximal cube.  By
\eqref{eq:damp-two-coordinate-shadow-union-cover}, one of $B,C$ shatters
$S$.  That family is ample, so it contains an $S$-cube; uniqueness forces
this cube to be $Q_A(a)$.  Since $a$ belongs to exactly one of $B,C$, the
family containing the cube is precisely $K_a$.  Equation
\eqref{eq:damp-two-coordinate-surviving-corner} now gives two corners of
$H$.  Hence cornerlessness forces
$\operatorname{corn}(A)\subseteq D$.  For a corner in $D$, the same
shadow-cover argument puts its unique cube in at least one of $B,C$.
Equation~\eqref{eq:damp-two-coordinate-surviving-corner} says that it
cannot lie in both, proving
\eqref{eq:damp-maximum-core-corner-shield} and the asserted
monochromaticity.

Now assume that $H$ is a least cornerless family.  The family $D$ is
nonempty by
\eqref{eq:damp-two-coordinate-compatible-shadow-cover}.  If $D=A$, then
$B=C=A$ and $H=A\times Q_2$, the equality product already excluded in
Corollary~\ref{cor:damp-least-cornerless-maximum-core}.  Thus $D\subsetneq
A$.  Suppose, say, that $B=A$.  Then $C=D$.  The smaller nonempty ample
family $D$ cannot be cornerless, so choose a corner $a$ of $D$.  At its
$C$-layer copy, both repair directions and exactly the old directions of
$Q_D(a)$ are incident, and the cube generated by them is
$Q_2\times Q_D(a)\subseteq H$.  This copy of $a$ is a corner of $H$, a
contradiction.  Therefore $B\subsetneq A$, and symmetry gives
$C\subsetneq A$.
\end{proof}

\begin{corollary}[Compatible corner for a VC-two core]
\label{cor:damp-two-coordinate-vc-two-compatible-corner}
Under the hypotheses of Proposition
\ref{prop:damp-two-coordinate-maximum-core-repair}, suppose in addition
that $A$ is maximum of VC dimension at most two and that $H$ is ample.
Then $H$ has a corner.  Consequently, the internal two-coordinate branch
of a least cornerless ample family cannot have a maximum core of VC
dimension at most two.
\end{corollary}

\begin{proof}
Put $D=B\cap C$.  This family is ample by
\eqref{eq:damp-two-coordinate-compatible-shadow-cover}, and it is
nonempty because otherwise the empty support would belong to both
$\operatorname{Sh}(B)$ and $\operatorname{Sh}(C)$ but not to
$\operatorname{Sh}(D)$.  Translate a concept of $D$ to the empty set.
This preserves inclusion, ampleness, maximumness, VC dimension, cubes,
and corners, and ensures that $\varnothing\in D\subseteq A$.

M\'esz\'aros and R\'onyai proved that every ample family of VC dimension
at most two containing the empty set belongs to their building class,
and, more strongly, that whenever two such families are nested, the
smaller one can be extended to the larger one by valid building
steps~\cite[Theorem~4.1 and Lemma~4.1]{meszaros2014vc2}.  Each step adds
one concept: a step of type A attaches it by one new coordinate edge,
whereas a step of type B completes a square and attaches it by the two
edges at the new vertex.

If $D=A$, take the concept added last in a valid construction of $A$
from the singleton $\{\varnothing\}$; if $A$ itself is a singleton, take
its only concept.  In either case this gives a corner $a$ of $A$ and
$Q_A(a)\subseteq D=K_a$, so
\eqref{eq:damp-two-coordinate-surviving-corner} applies.

Suppose now that $D\subsetneq A$, and take the concept $a$ added by the
last step of a valid relative construction from $D$ to $A$.  Thus
$a\notin D$, and $a$ is a corner.  Let $S$ be the support of its unique
maximal cube $Q_A(a)$.  If $A$ has VC dimension one, the last step is of
type A and its new coordinate occurs on no other edge of $A$; hence
$Q_A(a)$ is the unique $S$-cube.  If $A$ has VC dimension two, the last
step cannot be of type A: after that last step its new coordinate would
belong to no square, whereas a maximum VC-two family shatters every pair
of active coordinates.  The step is therefore of type B.  It completes
the square $Q_A(a)$, and this is the unique $S$-square in $A$.  The last
uniqueness also follows from Proposition
\ref{prop:damp-maximum-vc-three-carrier-trees}, because the carrier of an
$S$-cube in a VC-two family has no active coordinate.  The case of VC
dimension zero cannot occur when $D\subsetneq A$.

By \eqref{eq:damp-two-coordinate-shadow-union-cover}, at least one of
$B$ and $C$ shatters $S$.  Since that family is ample, it contains an
$S$-cube, necessarily the unique cube $Q_A(a)$.  On the other hand,
$a\notin D$ and $B\cup C=A$, so $a$ belongs to exactly one of $B,C$;
only that family can contain $Q_A(a)$.  Consequently
$Q_A(a)\subseteq K_a$, and
\eqref{eq:damp-two-coordinate-surviving-corner} again gives two corners
of $H$.  A least cornerless family cannot contain such an internal
repair, which proves the final assertion.
\end{proof}

\begin{corollary}[Finite-rank remainder of the two-coordinate branch]
\label{cor:damp-internal-two-coordinate-dimension-bound}
Let $H$ be as in Corollary
\ref{cor:damp-least-cornerless-maximum-core} and let $r=2$.  After the
constant coordinates of $A$ are deleted,
\[
 \operatorname{VCdim}(A)\geq3
 \qquad\text{and}\qquad
 5\leq\operatorname{idim}(A)\leq8.
\]
More precisely, if $k=\operatorname{VCdim}(A)$ and
$n=\operatorname{idim}(A)$, then
\[
 k=3\Longrightarrow n\leq8,
 \qquad
 k=4\Longrightarrow n\leq7,
 \qquad
 k\geq5\Longrightarrow n\leq6.
\]
\end{corollary}

\begin{proof}
The VC-dimension bound follows from Theorem
\ref{thm:damp-two-coordinate-low-vc-core-corner}.  The four-coordinate
alternative of the same theorem gives the lower bound
\(\operatorname{idim}(A)\ge5\).  If the active dimension $n$ were at
least nine, maximumness and
$\operatorname{VCdim}(A)\geq3$ would give
\[
 |A|=\sum_{i=0}^{\operatorname{VCdim}(A)}\binom ni
 \geq1+9+\binom92+\binom93=130,
\]
contrary to $|A|\leq99$ in
\eqref{eq:damp-least-core-order-bound}.
For $k=4$, the first excluded value is
$\sum_{i=0}^4\binom8i=163$, and for $k\geq5$ it is already
$\sum_{i=0}^5\binom7i=120$.  These give the two sharper bounds.
\end{proof}

\begin{proposition}[Core--repair shadow budget at order twenty-seven]
\label{prop:damp-order-twenty-seven-core-shadow-budget}
Let $H$ have minimum cardinality among nonempty cornerless ample families,
and suppose $|H|=27$.  Write $d=\operatorname{idim}(H)$, and use the core
$A$ and repair set $R$ from Corollary
\ref{cor:damp-least-cornerless-maximum-core}.  Then $H$ is not maximum,
$A$ is a pseudogeometric maximum VC-one path, and, with
\[
 r=|R|,\qquad n=d-r,
\]
the path has exactly $n+1$ vertices.

Let $X$ be the $n$ active coordinates of $A$.  The shattered-set complex
of $H$ contains every pair $xe$ with $x\in X$ and $e\in R$.  If $q$ is
the number of its other two-element faces and $t$ is the number of its
three-element faces, then
\begin{equation}
 27=(n+1)(r+1)+q+t,
 \qquad q\geq1,\quad t\geq1.
 \label{eq:damp-order-twenty-seven-bipartite-shadow-budget}
\end{equation}
Consequently the only numerical possibilities are
\begin{equation}
\begin{array}{c|c|c|c}
d&n&r&q+t\\ \hline
8&5,4,3,2&3,4,5,6&3,2,3,6\\
9&2&7&3.
\end{array}
\label{eq:damp-order-twenty-seven-core-parameter-table}
\end{equation}
In particular, the isometric-dimension-ten profile and the maximum
VC-two core are impossible.
\end{proposition}

\begin{proof}
A maximum class of order $27$ would have
$27=\sum_{i=0}^k\binom ni$.  The only solution is $k=1,n=26$, and a
maximum VC-one class is a tree and has a corner.  Thus $H$ is not maximum.
Corollary~\ref{cor:damp-least-cornerless-maximum-core} gives
\[
 |A|\leq\left\lfloor\frac{26}{r+1}\right\rfloor.
\]
The class $A$ is not a full cube: otherwise
$\operatorname{Susp}(A)$ would also be a full cube and hence ample,
contrary to Proposition
\ref{prop:damp-antipodal-double-maximum-fibre}.  If $r=2$, then $|A|\leq
8$.  A proper maximum class of VC dimension at least three has at least
$\sum_{i=0}^3\binom4i=15$ concepts, so
$\operatorname{VCdim}(A)\leq2$.  Theorem
\ref{thm:damp-two-coordinate-low-vc-core-corner} then gives a corner of
$H$, a contradiction.  Hence $r\ge3$.

Now $|A|\leq6$.  A proper maximum VC-two class already has
$1+3+\binom32=7$ concepts, while VC dimension zero gives a singleton full
cube and is excluded by the suspension condition.  Hence $A$ has VC
dimension one.

The family $A$ is maximum on all $n=d-r$ coordinates left after the
conditioning; these are not auxiliary constant coordinates.  A maximum
VC-one family on $n$ coordinates has $n+1$ concepts and, in the
pseudogeometric case, its one-inclusion graph is a path.  Its shattered
sets are the empty set and the $n$ coordinate singletons.  The proof of
Proposition~\ref{prop:damp-geodesic-disagreement-tax} shows that $H$
shatters $S$ and $S\cup\{e\}$ for every
$S\in\operatorname{Sh}(A)$ and $e\in R$.  Thus it contains all $nr$
cross-pairs between $X$ and $R$.

The preliminary facet-excess bound gives order at least $34$ in VC
dimension at least four, so the endpoint has VC dimension three.  Relative
peeling and Proposition
\ref{prop:damp-complete-indexed-seventh-slack-q4} give $d\geq8$, whereas
Corollary~\ref{cor:damp-fourth-strict-facet-excess} gives $d\leq10$.
Its Sandwich equality now
gives
\[
 27=1+d+nr+q+t=(n+1)(r+1)+q+t.
\]
There is at least one triangle because the VC dimension is three.  The
complete bipartite graph on $X\mathbin{\dot\cup}R$ is triangle-free, so
the pair shadow of that triangle contains at least one noncross pair;
hence $q\geq1$.  Therefore $(r+1)(d-r+1)\leq25$.  For $d=10$ the left
side is at least $27$ for $2\leq r\leq8$.  For $d=9$ only $r=7$
remains, with value $24$.  For $d=8$, the values at
$r=3,4,5,6$ are $24,25,24,21$.  Subtracting from $27$ proves
\eqref{eq:damp-order-twenty-seven-core-parameter-table}.
\end{proof}

\begin{proposition}[Small-excess bipartite shadows]
\label{prop:damp-small-excess-bipartite-shadows}
Let $H$ be an irredundant ample family of VC dimension three and order
$27$.  Suppose that, for a partition $E=X\mathbin{\dot\cup}R$, every
cross-pair $xe$ with $x\in X$ and $e\in R$ belongs to
$\operatorname{Sh}(H)$.  No cornerless such family exists in any of the
following cases, or after exchanging the two blocks:
\[
 (|X|,|R|,q+t)=(4,4,2),\qquad(5,3,3),\qquad(7,2,3),
\]
where $q$ and $t$ have the meaning of Proposition
\ref{prop:damp-order-twenty-seven-core-shadow-budget}.
\end{proposition}

\begin{proof}[Computer-assisted proof]
For fixed block sizes and excess, the pair shadow of every chosen triangle
forces its noncross pairs.  After these pairs are inserted, the verifier
chooses the remaining noncross pairs and quotients the resulting facet
antichains by permutations within $X$ and within $R$.  It obtains two
block-coloured types for $(4,4,2)$, nine for $(5,3,3)$, and seven for
$(7,2,3)$.  The two balanced types become isomorphic if the blocks are
exchanged.

For each of these eighteen support systems, Proposition
\ref{prop:damp-facet-cube-reconstruction} reduces realization to the
anchors of one affine cube on every facet support.  The exact CNF chooses
one anchor per facet, requires the union to have order $27$, and forbids
incidence one at every covered vertex.  It also includes the carrier
cardinality equalities, which follow from ampleness and improve
propagation.  The formulas have between $29{,}155$ and $60{,}248$
variables and between $67{,}857$ and $138{,}340$ clauses.  CaDiCaL returns
\textsc{unsat} on all eighteen formulas.

The deterministic verifier and its full list of facet representatives,
formula sizes, and CNF digests are
\[
\begin{gathered}
\texttt{verify\_damp\_small\_excess\_bipartite\_shadows.py},\\
\texttt{damp\_small\_excess\_bipartite\_shadows.json}.
\end{gathered}
\]
The mathematical-profile digest is
\[
\texttt{e3383a65257458ea5d27bfdc27bef2e9f61beaeac9e295945de59e29a8a2b764}.
\]
Thus none of the listed support systems has a cornerless ample anchored
realization.
\end{proof}

\begin{corollary}[Single oriented residual shadow at order twenty-seven]
\label{cor:damp-order-twenty-seven-maximum-core-localization}
Under the hypotheses of Proposition
\ref{prop:damp-order-twenty-seven-core-shadow-budget}, one has $d=8$ and
\[
 (n,r,q+t)=(2,6,6).
\]
Thus the maximum core is a three-vertex path on the two-coordinate block,
the repair set is the six-coordinate block, and the shattered complex
contains $K_{2,6}$ with exactly six further pair or triangle faces.
\end{corollary}

\begin{proof}
Proposition~\ref{prop:damp-small-excess-bipartite-shadows} excludes the
dimension-nine row and the three middle dimension-eight rows of
\eqref{eq:damp-order-twenty-seven-core-parameter-table}.  The displayed
row is the only remaining possibility.
\end{proof}

\begin{proposition}[The oriented $K_{2,6}$ endpoint is impossible]
\label{prop:damp-oriented-k26-endpoint-exclusion}
Let $H$ be an irredundant ample family of VC dimension three and order
$27$.  Suppose that its coordinates split as $X\mathbin{\dot\cup}R$ with
$|X|=2$ and $|R|=6$, that its shattered complex contains every cross-pair
of $K_{X,R}$ and exactly six further pair or triangle faces, and that two
antipodal $R$-fibres of $H$ are the same maximum VC-one family on $X$.
Then $H$ has a corner.
\end{proposition}

\begin{proof}[Computer-assisted proof]
Coordinate reversals normalize the two antipodal repair words to $0^6$
and $1^6$ and their common three-vertex path to
$Q_X\setminus\{10\}$.  The bounded support generator first chooses the
triangles, inserts their forced noncross pair shadows, and then chooses the
remaining noncross pairs.  A coloured incidence-graph certificate quotients
by permutations within $X$ and within $R$.  This gives $321$ support
systems; it is a census of the single oriented shadow layer in the
hypotheses, not of cube families or of another cardinality.

For each support system, Proposition
\ref{prop:damp-facet-cube-reconstruction} gives an exact anchored-cube CNF.
It chooses one affine cube on every facet support, requires their union to
have order $27$, forbids incidence one, and imposes the carrier-cardinality
equalities forced by ampleness.  Eight unit clauses fix the two normalized
endpoint fibres to $Q_X\setminus\{10\}$.  The formulas have between
$29{,}123$ and $30{,}046$ variables and between $64{,}909$ and $68{,}784$
clauses.  CaDiCaL returns \textsc{unsat} on all $321$ formulas.

The support generator, verifier, and merged certificate are
\[
\begin{gathered}
 \texttt{explore\_damp\_k26\_residual\_structure.py},\\
 \texttt{verify\_damp\_k26\_residual\_anchors.py},\\
 \texttt{damp\_k26\_oriented\_maximum\_core.json}.
\end{gathered}
\]
The mathematical-profile digest is
\[
 \texttt{21129e2baf496b277be92b60b369acd132dd05a50f3211bca5314ee9549ffb16}.
\]
The CNF is equivalent to a cornerless ample anchored realization with the
prescribed maximum core, so unsatisfiability proves the claim.
\end{proof}

\begin{corollary}[Exclusion of order twenty-seven]
\label{cor:damp-order-twenty-seven-excluded}
The least order of a nonempty cornerless ample family is not $27$.
\end{corollary}

\begin{proof}
If a least family had order $27$, Corollary
\ref{cor:damp-order-twenty-seven-maximum-core-localization} would put it
under Proposition~\ref{prop:damp-oriented-k26-endpoint-exclusion}, which
supplies a corner.
\end{proof}

In the two-coordinate case of Corollary
\ref{cor:damp-least-cornerless-maximum-core}, isometry first gives
$A\subseteq H_{01}\cup H_{10}$: a shortest path between the two copies
of $(a,00)$ and $(a,11)$ must use one of the intermediate copies of $a$.
Proposition~\ref{prop:damp-two-coordinate-maximum-core-repair} handles
the internal-repair branch in which both intermediate fibres lie in $A$.
It reduces that entire branch to one explicit statement: under the
compatible shadow cover
\eqref{eq:damp-two-coordinate-compatible-shadow-cover}, find a corner
$a$ of $A$ for which
\eqref{eq:damp-two-coordinate-surviving-corner} holds.  Corollary
\ref{cor:damp-two-coordinate-vc-two-compatible-corner} proves this in
all dimensions through VC dimension two.  Thus the remaining internal
branch has VC dimension at least three and active dimension at most
eight.  Proposition~\ref{prop:damp-maximum-core-corner-shield} sharpens
its obstruction: the proper ample overlap $D=B\cap C$ must contain every
core corner, while the complement $A\setminus D$ meets every corner cube
and each such cube is assigned to exactly one cover.  It is therefore a
bounded-rank relative-corner problem rather than a
cardinality-by-cardinality search.  The remaining external-repair branch is the
case in which an intermediate fibre contains a word outside $A$; those
words already contribute strict geodesic-bundle excess and should be
charged to the core cubes whose corners they repair.

The compatible-corner statement has a direct bounded check on the
canonical maximum VC-dimension-two core
\[
 A=\{x\in Q_4:|x|\leq2\}.
\]
Among its $426$ ample subfamilies there are $2{,}169$ ordered pairs
$(B,C)$ satisfying the hypotheses of Proposition
\ref{prop:damp-two-coordinate-maximum-core-repair}.  Every pair has at
least three of the six core corners satisfying
\eqref{eq:damp-two-coordinate-surviving-corner}; the resulting repaired
family has at least eight corners.  Moreover, repeatedly deleting the
least available corner completely peels all $2{,}169$ families, whose
orders range from $33$ to $44$.  Thus this test produces neither a
cornerless repair nor a cornerless descendant.  The exhaustive profile is
generated by
\path{forbidden_pc_minor_campaign/explore_damp_two_coordinate_maximum_core_repairs.py}
and has mathematical digest
\begin{center}
\small
\path{b1d4d540fa0d0852f6910b8baad08c468c15d37dd992025e2c639d82f03f865c}
\end{center}
This finite result is evidence for the compatible-corner lemma, not its
proof in arbitrary dimension.

A second audit removes the choice of core.  There are $832$ proper
labelled maximum families in $Q_4$, forming eight orbits under cube
automorphisms.  Across one representative of each orbit, the script
\path{forbidden_pc_minor_campaign/verify_damp_two_cover_compatible_corner_q4.py}
checks all $10{,}820$ compatible ordered ample covers.  Every cover has a
surviving core corner, and the same deterministic rule completely peels
every repaired family.  There are no cornerless repairs or cornerless
descendants.  The mathematical digest is
\begin{center}
\small
\path{ff73f017733c75767269a7f81f276a26c0b9ff3a2daa0f4cfd65eef6a353b41e}
\end{center}
This broader check also includes the punctured-cube maximum
VC-dimension-three core; it remains a finite-dimensional falsifier rather
than a proof of the proposed lemma.

The corollary is a dimension-free reduction, but it does not by itself
preserve cornerlessness: a corner of the maximum fibre $A$ may be covered
in $H$ by a cube using one of the conditioned coordinates in $R$.  The
remaining nonmaximum problem is therefore a labelled repair problem.  One
must show that these repairs either leave a proper cornerless pc-minor or
require at least $299$ concepts.  This is the point at which the
minimal-nonample classification supplies a genuine maximum core without
silently assuming that the original family was maximum.  Moreover, the
near-extremal objects are now concrete: their sections over $A$ are short
antipodal paths, equality is a literal product with one common path, and
the first excess measures deviations from that common path.  Compatibility
of these deviations across the cubes of $A$ is the next dimension-free
object to control.

\subsection{Codimension-two tree wedges}

The preceding example lives in twelve coordinates, but a maximum class of
codimension two already has a rigid local form.  Its complement is a tree,
and its corners are read directly from the wedges of that tree.  Besides
controlling every five-coordinate projection of a maximum
VC-dimension-three family, the formula will identify the only
pseudogeometric maximum VC-two core on four active coordinates.

\begin{proposition}[Codimension-two tree-wedge formula]
\label{prop:damp-five-coordinate-tree-wedges}
Let \(Z\) be a coordinate set of order \(d\ge3\), let
\(A\subseteq Q_Z\) be maximum of VC dimension \(d-2\), and put
\(D=Q_Z\setminus A\).  The one-inclusion graph \(T_D\) is a tree on
\(d+1\) vertices with exactly one edge in each coordinate of \(Z\).
Moreover,
there is a bijection
\[
 \operatorname{corn}(A)
 \longleftrightarrow
 \bigl\{(v,\{e,f\}):
       \text{the \(e\)- and \(f\)-edges of \(T_D\) meet at \(v\)}\bigr\}.
\]
Under this bijection, the corner corresponding to
\((v,\{e,f\})\) is \(v^{\{e,f\}}\), and its unique maximal cube has
support \(Z\setminus\{e,f\}\).  Consequently,
\begin{equation}
 |\operatorname{corn}(A)|
   =\sum_{v\in D}\binom{\deg_{T_D}(v)}2
   \geq d-1.
 \label{eq:damp-five-coordinate-tree-wedge-count}
\end{equation}
Equality holds if and only if \(T_D\) is a path; otherwise \(A\) has at
least \(d\) corners.

For \(X\subseteq Q_Z\), let
\(\overline X=\{\overline x:x\in X\}\), where \(\overline x\) is the
antipode of \(x\) in \(Q_Z\).  Then
\begin{equation}
 |A\cap\overline A|=2^d-2d-2+|D\cap\overline D|
 =\begin{cases}
   2^d-2d,&T_D\text{ is a path},\\
   2^d-2d-2,&T_D\text{ is not a path}.
  \end{cases}
 \label{eq:damp-five-coordinate-antipodal-core}
\end{equation}
In particular, by the affine-oriented-matroid criterion of
G\"artner and Welzl~\cite[Theorem~51]{GartnerWelzl1994}, \(A\) is the
tope set of a simple affine uniform oriented matroid if and only if
\(T_D\) is a path.
\end{proposition}

\begin{proof}
Complementation takes a maximum class of VC dimension three on five
coordinates---and, in general, a maximum class of VC dimension \(d-2\)
on \(d\) coordinates---to a maximum class of VC dimension one
\cite[Theorem~3.5]{ChalopiChepoiMoran22}.  Thus \(D\) is ample, has six
vertices when \(d=5\), and has \(d+1\) vertices in general.  It shatters
every singleton, and therefore contains an edge in each coordinate
direction.  No direction occurs twice:
in a one-dimensional ample family every edge is maximal, while distinct
maximal cubes have distinct supports
\cite[Corollary~3.3]{ChalopiChepoiMoran22}.  Since an ample family is
isometric and hence connected, these \(d\) edges form the asserted tree.

For \(e\in Z\) and \(x\in Q_Z\), remove the unique \(e\)-edge from \(T_D\)
and let \(H_e(x)\) be the component on which coordinate \(e\) has value
\(x_e\).  For distinct \(e,f\in Z\), the trace \(x|_{\{e,f\}}\) is
realized by \(D\) exactly when
\begin{equation}
 H_e(x)\cap H_f(x)\neq\varnothing.
 \label{eq:damp-tree-halfspace-trace}
\end{equation}
The restriction of \(D\) to \(\{e,f\}\) realizes three of the four binary
patterns.  Hence the missing pattern fixes the unique cube of \(A\) with
support \(Z\setminus\{e,f\}\).  It follows from
\eqref{eq:damp-tree-halfspace-trace} that the \((d-2)\)-cubes of \(A\)
through \(x\) are indexed by the pairs \(\{e,f\}\) for which
\(H_e(x)\) and \(H_f(x)\) are disjoint.

If \(x\in A\), the \(d\) subtrees \(H_e(x)\), \(e\in Z\), have empty
total intersection, since a vertex in their intersection would be a
member of \(D\) equal to \(x\).  Subtrees of a tree have the Helly
property, so at least one pair is disjoint.  Thus every vertex of \(A\)
belongs to one of the displayed \((d-2)\)-cubes.

Suppose first that the \(e\)- and \(f\)-edges meet at \(v\), and set
\(x=v^{\{e,f\}}\).  The two halfspaces \(H_e(x)\) and \(H_f(x)\) are the
two branches directed away from \(v\), so they are disjoint.  Every other
pair intersects: if neither coordinate is \(e\) or \(f\), both halfspaces
contain \(v\); and if the pair is \(\{e,g\}\) with \(g\neq f\), both
halfspaces contain the neighbor of \(v\) across the \(e\)-edge.  The case
\(\{f,g\}\) is symmetric.  Moreover, the \(e\)-neighbor of \(x\) is the
neighbor of \(v\) across the \(f\)-edge of \(T_D\), and the \(f\)-neighbor
of \(x\) is the neighbor of \(v\) across the \(e\)-edge.  Both belong to
\(D\).  All remaining neighbors of \(x\) lie in the displayed
\((Z\setminus\{e,f\})\)-cube.  Every cube of \(A\) through \(x\) is
therefore contained in that cube, so \(x\) is a corner with the stated
unique maximal cube.

Conversely, suppose that \(H_e(x)\) and \(H_f(x)\) are the unique disjoint
pair.  If the \(e\)- and \(f\)-edges were not incident, choose an edge
\(g\) strictly between them in \(T_D\).  Disjointness forces
\(H_e(x)\) and \(H_f(x)\) to be the two outward branches.  Neither
component of \(T_D-g\) meets both branches, contradicting the required
intersections of \(H_g(x)\) with each of them.  Hence the two edges meet at
a vertex \(v\).  For every \(g\notin\{e,f\}\), the component \(H_g(x)\)
must meet both outward branches and therefore contains \(v\).  Thus
\(x_g=v_g\), while \(x_e\neq v_e\) and \(x_f\neq v_f\), proving
\(x=v^{\{e,f\}}\).  If two pairs of halfspaces are disjoint, the
corresponding distinct \((d-2)\)-cubes both contain \(x\), so \(x\) is not a
corner.  Together with the preceding paragraph and the Helly argument,
this proves the claimed bijection.

Counting incident pairs gives the equality in
\eqref{eq:damp-five-coordinate-tree-wedge-count}.  Since every vertex of
the tree has positive degree,
\[
 \sum_{v\in D}\binom{\deg_{T_D}(v)}2
 \geq\sum_{v\in D}(\deg_{T_D}(v)-1)
 =2d-(d+1)=d-1.
\]
Equality holds exactly when every degree is at most two, that is, when the
tree is a path.  Every nonpath therefore has at least \(d\) wedges.

Finally,
\[
 |A\cap\overline A|
 =2^d-|D\cup\overline D|
 =2^d-2d-2+|D\cap\overline D|.
\]
If \(D\) contains \(u\) and its antipode, isometry gives a path of length
\(d\) between them.  This path uses all \(d\) edges of \(T_D\), so
\(T_D\) is a path and \(u,\overline u\) are its endpoints.  Conversely,
the endpoints of a path \(T_D\) differ in every coordinate, because its
edges have distinct coordinate labels.  These are the only antipodal
vertices in the path, so
\(|D\cap\overline D|=2\) in the path case and is zero otherwise.  This proves
\eqref{eq:damp-five-coordinate-antipodal-core}.  The final equivalence is
exactly the cited criterion, whose required antipodal-core cardinality is
\[
 2\sum_{i=0}^{d-3}\binom{d-1}{i}=2^d-2d.
\]
\end{proof}

\begin{corollary}[Five-coordinate repair obligation]
\label{cor:damp-five-coordinate-repair-obligation}
Let \(H\subseteq Q_E\) be a cornerless maximum family of VC dimension
three.  For every five-set \(Z\subseteq E\), the projection
\(A=H|_Z\) has at least four corners.  If \(A\) is not the tope set of a
simple affine uniform oriented matroid, it has at least five.

More precisely, let \(p\) be one of these projected corners, let \(S\) be
the support of its unique maximal cube in \(A\), and let \(x\in H\) be the
unique lift of \(p\) in the \(S\)-cube of \(H\).  Then \(x\) lies in a
second maximal cube of \(H\), with support \(R\), such that
\begin{equation}
 R\cap Z\subseteq S
 \qquad\text{and}\qquad
 R\setminus Z\neq\varnothing.
 \label{eq:damp-five-coordinate-crossing-repair}
\end{equation}

Suppose in addition that \(d=|E|\geq6\), and let \(b_5(H)\) be the number
of five-sets \(Z\subseteq E\) for which \(H|_Z\) is not an affine uniform
oriented-matroid tope set.  For \(x\in H\), let \(t_x\) be the number of
maximal cubes of \(H\) containing \(x\).  Then
\begin{equation}
 \binom{d-4}{2}\sum_{x\in H}t_x(t_x-1)
 \ \geq\ 4\binom d5+b_5(H).
 \label{eq:damp-five-coordinate-repair-capacity}
\end{equation}
More precisely, if \(\mathcal I(H)\) is the set of triples \((S,x,R)\)
such that the distinct \(S\)- and \(R\)-cubes of \(H\) both contain \(x\),
then the left-hand side may be replaced by
\begin{equation}
 \sum_{(S,x,R)\in\mathcal I(H)}
       \binom{d-|S\cup R|}{2}.
 \label{eq:damp-five-coordinate-sharp-repair-capacity}
\end{equation}
\end{corollary}

\begin{proof}
Restrictions of maximum families are maximum, so the numerical assertions
follow from Proposition~\ref{prop:damp-five-coordinate-tree-wedges}.
There is a unique \(S\)-cube of \(H\), and its projection is the unique
\(S\)-cube of \(A\); hence it contains a unique lift \(x\) of \(p\).
Cornerlessness supplies another maximal cube \(B\) through \(x\); write
\(R\) for its support.  The projection of \(B\) is a cube of \(A\) through
\(p\), so it is contained in the unique maximal cube through \(p\).  This
gives \(R\cap Z\subseteq S\).  If \(R\subseteq Z\), then \(R\subseteq S\).
Since both cubes contain \(x\), the cube \(B\) is contained in the global
\(S\)-cube, contradicting its maximality.
Thus \(R\setminus Z\neq\varnothing\).

Choose one such support \(R\) for every projected corner.  This assigns
each local corner to a triple \((S,x,R)\in\mathcal I(H)\).  For a fixed
triple, an eligible five-set has the form \(Z=S\cup W\), where
\[
 W\in\binom{E\setminus(S\cup R)}2.
\]
Thus \((S,x,R)\) receives at most
\(\binom{d-|S\cup R|}{2}\) assignments.  Summing these capacities and
using the lower bounds of four for every \(Z\), with one additional unit
for each of the \(b_5(H)\) non-pseudogeometric projections, proves
\eqref{eq:damp-five-coordinate-sharp-repair-capacity}.  Maximality gives
\(R\not\subseteq S\), so \(|S\cup R|\geq4\).  Moreover,
\(|\mathcal I(H)|\leq\sum_x t_x(t_x-1)\).  Replacing every summand by
\(\binom{d-4}{2}\) and then using this inequality gives
\eqref{eq:damp-five-coordinate-repair-capacity}.
\end{proof}

The preceding capacity bound deliberately forgets which projected corners
are repaired together.  The next description retains that information and
turns all five-coordinate corner obligations into collisions in a family
of eight-colourings.

\begin{proposition}[Tip-colouring identity]
\label{prop:damp-five-coordinate-tip-colouring}
Let \(H\subseteq Q_E\) be maximum of VC dimension three, with
\(d=|E|\geq5\).  For each \(S\in\binom E3\), let \(C_S\) be the unique
three-cube of \(H\) with support \(S\).  If \(e\in E\setminus S\), then
\(H|_{S\cup\{e\}}\) is a four-cube with one trace \(q_{S,e}\) deleted.
Define the \emph{\(e\)-tip of \(C_S\)} to be the unique vertex
\(\tau_S(e)\in C_S\) satisfying
\[
 \tau_S(e)|_S=q_{S,e}|_S.
\]
For \(x\in C_S\), put
\[
 A_S(x)=\{e\in E\setminus S:\tau_S(e)=x\},
 \qquad a_S(x)=|A_S(x)|.
\]
Then the eight sets \(A_S(x)\), \(x\in C_S\), partition
\(E\setminus S\).  Moreover, for distinct \(e,f\in E\setminus S\), the
projection \(H|_{S\cup\{e,f\}}\) has a corner whose unique maximal cube
has support \(S\) if and only if
\[
 \tau_S(e)=\tau_S(f).
\]
Consequently,
\begin{equation}
 \sum_{Z\in\binom E5}|\operatorname{corn}(H|_Z)|
 =\sum_{S\in\binom E3}\ \sum_{x\in C_S}\binom{a_S(x)}2,
 \qquad
 \sum_{x\in C_S}a_S(x)=d-3.
 \label{eq:damp-tip-collision-identity}
\end{equation}

Let
\[
 \epsilon_5(H)=
 \sum_{Z\in\binom E5}
 \left(
   \sum_{v\in T_Z}\binom{\deg_{T_Z}(v)}2-4
 \right),
\]
where \(T_Z\) is the complement tree of \(H|_Z\).  If \(H\) is
cornerless, then
\begin{equation}
 4\binom d5+\epsilon_5(H)
 =\sum_{S\in\binom E3}\ \sum_{x\in C_S}\binom{a_S(x)}2
 \leq\binom d3\binom{d-4}2.
 \label{eq:damp-cornerless-tip-collision-bound}
\end{equation}
In particular, no cornerless maximum VC-dimension-three family exists on
five or six coordinates.
\end{proposition}

\begin{proof}
The projection to \(S\cup\{e\}\) is maximum of VC dimension three and
therefore has fifteen vertices.  Its unique missing trace lies on the side
opposite the projected \(S\)-cube, so its restriction to \(S\) selects a
unique vertex of the global cube \(C_S\).  This proves that \(\tau_S(e)\)
is well-defined.  Every outside coordinate has one tip, which gives the
partition and the second identity in
\eqref{eq:damp-tip-collision-identity}.

Write \(Z=S\cup\{e,f\}\).  In the punctured four-cube
\(H|_{S\cup\{e\}}\), the vertex of the projected \(S\)-cube whose
\(S\)-trace is \(q_{S,e}|_S\) is the unique corner supported by \(S\).
Thus, if a vertex \(p\) of \(H|_Z\) is a corner supported by \(S\), its
projections after deleting \(f\) and after deleting \(e\) show that its
unique lift in \(C_S\) equals both \(\tau_S(e)\) and \(\tau_S(f)\).

Conversely, suppose \(\tau_S(e)=\tau_S(f)=x\), and put \(p=x|_Z\).
After deleting either outside coordinate, \(p\) is a corner supported by
\(S\).  If a cube of \(H|_Z\) through \(p\) used direction \(e\) or
direction \(f\), projecting away the other outside coordinate would give
a cube through one of these four-coordinate corners that is not contained
in its unique \(S\)-cube.  This is impossible.  Every cube through \(p\)
therefore lies in the projected \(S\)-cube, so \(p\) is a corner supported
by \(S\).  We have proved the claimed equivalence.  Counting the pairs in
each fibre gives the first identity in
\eqref{eq:damp-tip-collision-identity}.

Proposition~\ref{prop:damp-five-coordinate-tree-wedges} identifies the
left-hand side of that identity with
\(4\binom d5+\epsilon_5(H)\).  If some fibre has size \(d-3\), then its
vertex \(x\) is the tip for every direction outside \(S\).  The
four-coordinate corner characterization just used shows that no cube
through \(x\) has a direction outside \(S\); hence the global \(S\)-cube
is the unique maximal cube through \(x\), and \(x\) is a corner.  Thus, in
a cornerless family every part of the partition of \(E\setminus S\) has
size at most \(d-4\).  Among nontrivial partitions of a \((d-3)\)-element
set, the convex sum \(\sum_i\binom{a_i}2\) is maximized by the partition
\((d-4)+1\).  Summing over the \(\binom d3\) supports proves
\eqref{eq:damp-cornerless-tip-collision-bound}.  For \(d=5\) its right-hand
side is zero while its left-hand side is at least four; for \(d=6\), the
two sides are at least $24$ and at most $20$, respectively.
\end{proof}

The collision identity is the second member of a whole hierarchy.  This
hierarchy is useful because a lower-dimensional corner theorem can be fed
directly into the next ambient dimension, without enumerating the classes
in that dimension.

\begin{proposition}[Projection-corner moment identity]
\label{prop:damp-projection-corner-moments}
Keep the notation of
Proposition~\ref{prop:damp-five-coordinate-tip-colouring}.  For every
integer $r$ with $4\leq r\leq d$,
\begin{equation}
 \sum_{Z\in\binom Er}|\operatorname{corn}(H|_Z)|
 =\sum_{S\in\binom E3}\ \sum_{x\in C_S}
     \binom{a_S(x)}{r-3}.
 \label{eq:damp-projection-corner-moments}
\end{equation}
In particular, the full tip fibres $a_S(x)=d-3$ are in bijection with
the corners of $H$.

For $r\geq4$, let $c_r$ denote the minimum number of corners among
maximum VC-dimension-three families on $r$ coordinates.  If $H$ is
cornerless and $5\leq r<d$, then
\begin{equation}
 c_r\binom dr
 \leq \binom d3\binom{d-4}{r-3},
 \qquad\text{or equivalently}\qquad
 c_r\leq \binom r3\frac{d-r}{d-3}.
 \label{eq:damp-projection-corner-transfer}
\end{equation}
If $H$ is cornerless, then for $r=d-1$ the left-hand side of
\eqref{eq:damp-projection-corner-moments} counts exactly the
near-full fibres $a_S(x)=d-4$.  More precisely, a corner of
$H|_{E\setminus\{e\}}$ supported by $S$ corresponds to the unique
fibre
\begin{equation}
 A_S(x)=E\setminus(S\cup\{e\}).
 \label{eq:damp-near-full-tip-fibre}
\end{equation}
Thus any lower bound on corners in dimension $d-1$ becomes, in one
step, a labelled near-full-fibre obligation in dimension $d$.
\end{proposition}

\begin{proof}
Fix $Z\in\binom Er$ and $S\in\binom Z3$.  The projection $H|_Z$
is again maximum of VC dimension three, and its unique $S$-cube is the
projection of $C_S$.  A vertex $p$ of that cube has a unique lift
$x\in C_S$.

We claim that $p$ is a corner of $H|_Z$ supported by $S$ if and
only if
\begin{equation}
 \tau_S(e)=x\qquad\text{for every }e\in Z\setminus S.
 \label{eq:damp-projected-corner-all-tips}
\end{equation}
Suppose first that $p$ is such a corner.  If
$\tau_S(e)\neq x$ for some $e\in Z\setminus S$, then the trace that
agrees with $p$ on $S$ and has the opposite $e$-value is present in
$H|_{S\cup\{e\}}$.  Projections commute, so it has a lift
$y\in H|_Z$.  The vertices $p$ and $y$ agree on $S$.  Since the
maximum class $H|_Z$ is ample, its one-inclusion graph is isometric; a
shortest $p,y$-path therefore begins with an edge at $p$ in a
direction outside $S$.  A maximal cube containing this edge is not
contained in the $S$-cube, contradicting that $p$ is a corner
supported by $S$.  Hence $\tau_S(e)=x$ for every such $e$.

Conversely, assume
\eqref{eq:damp-projected-corner-all-tips}.  A cube of $H|_Z$ through
$p$ that is not contained in the $S$-cube uses some direction
$e\in Z\setminus S$.  Its projection to $S\cup\{e\}$ would then give
a cube through the $e$-tip $x$ that is not contained in the
$S$-cube, contradicting the four-coordinate tip characterization.
This proves the claim.

For fixed $S$ and $x$, condition
\eqref{eq:damp-projected-corner-all-tips} holds for precisely
$\binom{a_S(x)}{r-3}$ choices of $Z$.  Every corner of a maximum
VC-dimension-three projection lies in a unique maximal three-cube, so
summing over $S$ and $x$ proves
\eqref{eq:damp-projection-corner-moments}.  Taking $r=d$ gives the
assertion about full fibres.

Now suppose that $H$ is cornerless.  Then every fibre has size at most
$d-4$.  For $r\geq5$, superadditivity of
$a\mapsto\binom a{r-3}$ on the nonnegative integers shows that, among
nontrivial partitions of $d-3$, the sum of these binomial moments is
maximized by the partition $(d-4)+1$.  Hence
\[
 \sum_{S,x}\binom{a_S(x)}{r-3}
 \leq\binom d3\binom{d-4}{r-3}.
\]
Every $r$-coordinate projection has at least $c_r$ corners.  Combining
this observation with
\eqref{eq:damp-projection-corner-moments} gives the first inequality in
\eqref{eq:damp-projection-corner-transfer}; cancelling $\binom dr$
and simplifying the binomial coefficients gives the second.

Finally take $r=d-1$.  Since all fibres have size at most $d-4$, the
binomial coefficient $\binom{a_S(x)}{d-4}$ is one exactly for a
near-full fibre and zero otherwise.  Such a fibre omits a unique coordinate
$e\in E\setminus S$, and the already proved corner criterion identifies
its vertex with a corner of $H|_{E\setminus\{e\}}$.  This is precisely
\eqref{eq:damp-near-full-tip-fibre}.
\end{proof}

For $r=5$, Proposition~\ref{prop:damp-five-coordinate-tree-wedges}
gives $c_5=4$.  Substituting this value into
\eqref{eq:damp-projection-corner-transfer} forces $d\geq7$, recovering
the simultaneous exclusion of dimensions five and six.  Higher values of
$r$ retain strictly more of the projection structure: in particular, the
case $r=d-1$ records not merely how many projected corners exist, but
which deleted coordinate owns each one.

\begin{corollary}[Exact six-coordinate corner number]
\label{cor:damp-six-coordinate-corner-bound}
Every maximum VC-dimension-three family on six coordinates has at least
four corners, and this bound is attained.  Thus $c_6=4$.
\end{corollary}

\begin{proof}[Computer-assisted proof]
The six five-coordinate projections contribute at least $6c_5=24$
corners.  There are $\binom63=20$ tip partitions, each partitioning three
labels.  A partition with no full part contributes at most one to the
second binomial moment, whereas a full part contributes three and
corresponds to a global corner.  If $c$ is the number of global corners,
then

\[
 24\leq
 \sum_{S,x}\binom{a_S(x)}2
 \leq20+2c,
\]

and hence $c\geq2$.  This illustrates the intended use of
\eqref{eq:damp-projection-corner-moments}: a corner theorem in one
dimension becomes a simultaneous constraint on all larger dimensions.

For the sharp bound, Proposition
\ref{prop:damp-maximum-vc-three-missing-traces} below is compiled with one
corner marker for every source-cube vertex and a cardinality constraint
allowing at most three markers.  A vertex belonging to no second cube
forces its marker, so any family with at most three corners satisfies the
formula.  Cube translation, the first outside tip, and its remaining
coordinate orbit are normalized.  The deterministic formula has $3{,}935$
variables and $16{,}617$ clauses.  CaDiCaL versions 1.9.5 and 2.1.2 both
return \textsc{unsat}.  The latter proof is accepted by
\texttt{drat-trim}; its backward core contains $6{,}528$ original clauses
and $24{,}836$ lemmas and uses $2{,}317{,}233$ resolution steps.  The
formula and proof digests are
\[
 \begin{gathered}
 \texttt{e5c70eb2330bb0d221451aa67a087164c6ed0ec6a17fbb704ed93a96ba9ecef4},\\
 \texttt{dec117586483eb7eb3c293566fa44cb3ff1c103fc21f39b17457cc7eb7515503}.
 \end{gathered}
\]
Finally, every six-coordinate projection of Hall's class has exactly four
corners, as independently checked by
\texttt{analyze\_damp\_projection\_corner\_moments.py}.  Hence the lower
bound is attained.
\end{proof}

The tip maps admit an exact description using only the traces missing on
four coordinates.  Besides shortening the computational formulation, this
separates the maximum-class compatibility conditions from cornerlessness.

\begin{proposition}[Local missing-trace characterization]
\label{prop:damp-maximum-vc-three-missing-traces}
Let $E$ have $d\geq4$ elements.  For each
$A\in\binom E4$, choose a word $q_A\in Q_A$, and define
\[
 H(q)=\{x\in Q_E:x|_A\neq q_A
                 \text{ for every }A\in\tbinom E4\}.
\]
For $S\in\binom E3$, let $C_S(q)$ be the three-cube with support
$S$ whose fixed coordinates are
\begin{equation}
 x_e=1-q_{S\cup\{e\}}(e),\qquad e\in E\setminus S.
 \label{eq:damp-missing-trace-cube-anchor}
\end{equation}
Then $H(q)$ is maximum of VC dimension three, and the $q_A$ are its
missing four-traces, if and only if
\begin{equation}
 \text{for all }S\in\tbinom E3, A\in\tbinom E4
 \text{ with }|A\setminus S|\geq2,
 \quad
 q_A(e)=q_{S\cup\{e\}}(e)
 \text{ for some }e\in A\setminus S.
 \label{eq:damp-missing-trace-local-compatibility}
\end{equation}
When these conditions hold,
\begin{equation}
 H(q)=\bigcup_{S\in\binom E3}C_S(q).
 \label{eq:damp-missing-trace-cube-union}
\end{equation}

More explicitly, for $u\in Q_S$, let $x_{S,u}\in C_S(q)$ be the
vertex with $x_{S,u}|_S=u$.  The maximum class $H(q)$ is cornerless if
and only if, for every pair $(S,u)$, there is a triple
$T\neq S$ such that $x_{S,u}\in C_T(q)$.  In missing-trace
coordinates this last condition says that, for every $h\in E\setminus T$,
\begin{equation}
 \begin{cases}
 u(h)=1-q_{T\cup\{h\}}(h),&h\in S,\\
 q_{S\cup\{h\}}(h)=q_{T\cup\{h\}}(h),&h\notin S.
 \end{cases}
 \label{eq:damp-missing-trace-cornerless-compatibility}
\end{equation}
Thus the existence of a cornerless maximum VC-dimension-three class is an
exact Boolean compatibility problem on the missing four-traces; no
cardinality constraint or enumeration of concept classes is required.
\end{proposition}

\begin{proof}
Assume first that
\eqref{eq:damp-missing-trace-local-compatibility} holds.  We claim that
$C_S(q)\subseteq H(q)$ for every $S$.  Fix $A\in\binom E4$.  If
$A=S\cup\{e\}$, then every vertex of $C_S(q)$ differs from $q_A$
at $e$ by \eqref{eq:damp-missing-trace-cube-anchor}.  Otherwise
$|A\setminus S|\geq2$, and
\eqref{eq:damp-missing-trace-local-compatibility} supplies
$e\in A\setminus S$ with
$q_A(e)=q_{S\cup\{e\}}(e)$.  The fixed value of $C_S(q)$ at $e$
is its complement, so again no vertex of the cube realizes $q_A$.
This proves the claim.

The family $H(q)$ omits a trace on every four-set, so its VC dimension is
at most three.  On the other hand, the cubes $C_S(q)\subseteq H(q)$
strongly shatter every set of size at most three.  The Sandwich inequality
and Sauer's inequality therefore give
\[
 \sum_{i=0}^3\binom di
 \leq |H(q)|
 \leq \sum_{i=0}^3\binom di.
\]
Thus $H(q)$ is maximum.  The same argument applied to the union of the
cubes proves \eqref{eq:damp-missing-trace-cube-union}.

Conversely, let $H$ be maximum of VC dimension three and let $q_A$ be
the unique trace missing on $A$.  A word outside $H$ cannot avoid all
the $q_A$, because adjoining it would preserve VC dimension at most
three and violate Sauer's bound.  Hence $H=H(q)$.  Its unique
$S$-cube has the anchor in
\eqref{eq:damp-missing-trace-cube-anchor}.  Since this cube avoids $q_A$,
one of its fixed coordinates in $A\setminus S$ differs from $q_A$
whenever $|A\setminus S|\geq2$; rewriting that difference gives
\eqref{eq:damp-missing-trace-local-compatibility}.

Finally, \eqref{eq:damp-missing-trace-cube-union} lists all maximal cubes:
each has dimension three and distinct maximal cubes have distinct
supports.  Therefore $x_{S,u}$ is a corner precisely when it belongs to
no $C_T(q)$ with $T\neq S$.  Comparing the fixed coordinates of the
two anchored cubes gives
\eqref{eq:damp-missing-trace-cornerless-compatibility}, completing the
proof.
\end{proof}

\begin{corollary}[Seven-coordinate maximum classes]
\label{cor:damp-seven-coordinate-maximum-corner}
Every maximum VC-dimension-three family on seven coordinates has a corner.
\end{corollary}

\begin{proof}[Computer-assisted proof]
Proposition~\ref{prop:damp-maximum-vc-three-missing-traces} gives an exact
Boolean formulation with $4\binom74=140$ mathematical missing-trace
bits.  We normalize a cube translation, the first outside tip, and the
order of the remaining outside coordinates.  Tseitin variables record the
choice of a second cube through every source-cube vertex and the redundant
five-coordinate tip-collision inequalities from
Proposition~\ref{prop:damp-five-coordinate-tip-colouring}.  The resulting
deterministic formula has $10{,}374$ variables and $68{,}677$ clauses.

Lingeling and CaDiCaL independently return \textsc{unsat}.  Lingeling's
DRAT proof is independently accepted by \texttt{drat-trim}; the backward
core contains $5{,}281$ original clauses and $17{,}677$ lemmas, and
uses $474{,}870$ resolution steps.  The durable files are
\[
 \begin{gathered}
 \texttt{solve\_damp\_cornerless\_maximum\_vc3\_missing\_traces\_cnf.py},\\
 \texttt{damp\_cornerless\_maximum\_vc3\_missing\_traces\_d7.cnf},\\
 \texttt{damp\_cornerless\_maximum\_vc3\_missing\_traces\_d7.drat},\\
 \texttt{damp\_cornerless\_maximum\_vc3\_missing\_traces\_d7\_cnf.json},\\
 \texttt{damp\_cornerless\_maximum\_vc3\_missing\_traces\_d7\_proof\_verification.json}.
 \end{gathered}
\]
The SHA-256 digests of the DIMACS formula and DRAT proof are, respectively,
\[
 \begin{gathered}
 \texttt{900eb029ad9947f89982d45d474dd7cfa62bfcde56c5e3c79c34c07210b54ef0},\\
 \texttt{f3effdb6192b724325e07b6a1b99471d8294a999dab63fb8ab1a0732a0bedd5b}.
 \end{gathered}
\]
Since the clauses are equivalent to maximum-class compatibility and
cornerlessness, unsatisfiability proves the claim.
\end{proof}

\begin{corollary}[Eight-coordinate maximum classes]
\label{cor:damp-eight-coordinate-maximum-corner}
Every maximum VC-dimension-three family on eight coordinates has a corner.
\end{corollary}

\begin{proof}[Computer-assisted proof]
We use the same exact missing-trace formulation as in
Corollary~\ref{cor:damp-seven-coordinate-maximum-corner}.  In dimension
eight it has $4\binom84=280$ mathematical missing-trace bits.  After the
repair and tip-collision witnesses are introduced, the deterministic
formula has $26{,}824$ variables and $230{,}952$ clauses.  CaDiCaL and
Kissat~4.0.4 independently return \textsc{unsat}.

Kissat's binary DRAT trace was written directly in gzip-compressed form and
independently checked by streaming its decompression to \texttt{drat-trim}.
The checker retains $28{,}741$ original clauses and $710{,}850$ lemmas in
the backward core, and uses $27{,}205{,}283$ resolution steps.  The durable
files are
\[
 \begin{gathered}
 \texttt{damp\_cornerless\_maximum\_vc3\_missing\_traces\_d8.cnf},\\
 \texttt{damp\_cornerless\_maximum\_vc3\_missing\_traces\_d8\_cnf.json},\\
 \texttt{damp\_cornerless\_maximum\_vc3\_missing\_traces\_d8\_kissat.proof.gz},\\
 \texttt{damp\_cornerless\_maximum\_vc3\_missing\_traces\_d8\_proof\_verification.json}.
 \end{gathered}
\]
The SHA-256 digests of the DIMACS formula and compressed proof are,
respectively,
\[
 \begin{gathered}
 \texttt{c254620be1f44e24332b1e75f3e30882d52c9472d2c0250f50de952f1ebddced},\\
 \texttt{d5d4d2e66f240a777c3b28d8bd9b1e67d924a000a764c215cfdefe661616c816}.
 \end{gathered}
\]
The independent verification record has mathematical digest
\[
 \texttt{d9c94c7387253dff8a2be493bce0d7b9d805dcb3c1815c8c2739bc47838f7d40}.
\]
Proposition~\ref{prop:damp-maximum-vc-three-missing-traces} identifies the
formula's unsatisfiability with the asserted corner conclusion.
\end{proof}

The same question makes sense in every VC dimension.  By
Corollary~\ref{cor:damp-least-cornerless-maximum-core}, a least nonempty
cornerless ample family of order below \(299\) is either a pseudogeometric
maximum class or has a maximum core of order at most \(99\).  For the first
alternative, only finitely many pairs \((d,n)\) are relevant.  An ample
family of VC dimension at most two is corner-peelable, and a maximum
family of VC dimension \(n-1\) on \(n\) coordinates is a cube with one
vertex deleted, which also peels.  Hence \(d\geq3\), \(n\geq d+2\), and
\[
 \Phi_d(n)=\sum_{i=0}^{d}\binom ni<299 .
\]
For every maximum family \(H\) of VC dimension \(d\), each maximal cube
has dimension \(d\).  Suppose that a \(k\)-cube with \(k<d\) and support
\(A\) is maximal at \(x\).  The reduction \(H^A\) is maximum of VC
dimension \(d-k\geq1\), hence a connected family with at least two
vertices, so \(x|_{E\setminus A}\) has a neighbour in \(H^A\).  This
neighbour extends the cube, contradicting maximality.  Consequently, by
\cite[Lemma~4.1]{ChalopiChepoiMoran22}, a vertex of \(H\) is a corner
exactly when it lies in one \(d\)-cube of \(H\).

\begin{proposition}[Further maximum classes with corners; computer-assisted]
\label{prop:damp-small-maximum-corners}
Every maximum family of VC dimension \(d\) on \(n\) coordinates has a corner
when
\[
 (d,n)\in\{(4,6),(4,7),(4,8),(5,7),(5,8),(6,8)\}.
\]
\end{proposition}

\begin{proof}[Computer-assisted proof]
Choose one \(d\)-cube for every \(d\)-set of coordinates.  Such a choice
determines the union \(U\) of the chosen cubes.  A maximum family of VC
dimension \(d\) contains exactly one \(d\)-cube with each support, and
these cubes cover it: their union strongly shatters every set of size at
most \(d\), hence has at least \(\Phi_d(n)\) vertices.  Conversely, if \(U\)
shatters no \((d+1)\)-set, then its shattered and strongly shattered sets
coincide, so \(U\) is ample, and by the Sandwich equality it is maximum.  By
the preceding paragraph, a cornerless maximum family is therefore exactly
such a choice in which \(U\) shatters no \((d+1)\)-set and every vertex of
\(U\) lies in at least two chosen cubes.  These conditions are encoded in
propositional logic with one variable per cube choice, one covering
variable per vertex, a sequential counter for double coverage, and one
auxiliary variable for each \((d+1)\)-set and missing trace.

The formula is invariant under the hyperoctahedral group acting on the
coordinates.  We add lexicographic-leader constraints \(v\leq_{\rm lex}g(v)\)
on the cube-choice variables for the generators \(g\): the adjacent
transpositions and the single-coordinate reflections.  The
lexicographically least member of an orbit of solutions satisfies these
constraints, so they preserve satisfiability.

For each listed pair, CaDiCaL~2.1.2 reports \textsc{unsat}, and its DRAT
proof is accepted by \texttt{drat-trim}.  Without the double-coverage
constraint, each formula is satisfiable, and the solver's model is a
maximum family.  The formulas are generated by
\texttt{explore\_damp\_group\_cornerless\_sat.py} with options
\texttt{none --vc=\(d\) --vc-encoding --lex}.  The compressed formulas and
proofs, the checker logs, and SHA-256 digests are in
\path{evidence/cornerless_maximum_small/}.  For \((d,n)=(4,8)\) the
\(1.16\)\,GB proof is not archived.  Its digest and checker log are
recorded there; \texttt{drat-trim} retains \(3{,}010{,}081\) lemmas in its
backward core.  To recheck it, one regenerates the formula, compares its
digest, reruns the solver and checks the fresh proof; a fresh proof need not
coincide byte for byte with the original.  Their sizes are as follows.
\begin{center}
\begin{tabular}{@{}cccrr@{}}
 \(d\) & \(n\) & \(\Phi_d(n)\) & variables & clauses\\ \hline
 4&6&57&2\,571&5\,299\\
 4&7&99&12\,205&27\,252\\
 4&8&163&49\,858&118\,878\\
 5&7&120&6\,519&13\,232\\
 5&8&219&34\,752&75\,548\\
 6&8&247&16\,386&32\,958
\end{tabular}
\end{center}
For \(d=3\) the same encoding returns \textsc{unsat} for \(n=8\), in agreement
with Corollary~\ref{cor:damp-eight-coordinate-maximum-corner}, and for \(n=7\) its DRAT
proof is again accepted by \texttt{drat-trim}.
\end{proof}

The same cube-choice description holds for every ample family, with the
facets of its shattered complex in place of the \(d\)-sets.

\begin{lemma}[Facet cubes of an ample family]
\label{lem:damp-ample-facet-cubes}
Let \(C\subseteq Q_E\) be a nonempty ample family, and call the inclusion-maximal
members of \(\operatorname{Sh}(C)\) its \emph{facets}.  For every facet \(S\), the
family \(C\) contains exactly one cube with support \(S\), the
\emph{facet cube} of \(S\).  The maximal cubes of \(C\) are exactly its
facet cubes.  Consequently \(C\) is the union of its facet cubes, and a
vertex of \(C\) is a corner if and only if it lies in exactly one facet
cube.

Conversely, let \(\Sigma\) be a down-closed family of subsets of \(E\), and
choose one cube with support \(S\) for every inclusion-maximal
\(S\in\Sigma\).  If the union \(U\) of the chosen cubes shatters no set
outside \(\Sigma\), then \(U\) is ample and \(\operatorname{Sh}(U)=\Sigma\).
\end{lemma}

\begin{proof}
For \(T\subseteq E\), the placements of the \(T\)-cubes of \(C\) form the
reduction \(C^T\).  It is ample and shatters exactly the sets \(U\)
with \(T\cup U\in\operatorname{Sh}(C)\).  If \(S\) is a facet, then \(C^S\) shatters
only the empty set, so \(|C^S|=|\operatorname{Sh}(C^S)|=1\).  Now let a cube with
support \(T\) be maximal in \(C\).  Its placement has no neighbour in
\(C^T\), since a neighbour would extend the cube.  Every nonempty ample
family is connected, so \(C^T\) is this one point.  Hence \(C^T\)
shatters only the empty set, and \(T\) is a facet.  Every vertex lies in
a maximal cube, so the union statement follows.  The corner statement
is then \cite[Lemma~4.1]{ChalopiChepoiMoran22}.

For the converse, every member of \(\Sigma\) is a subset of a chosen
cube's support, so it is strongly shattered by \(U\), and
\(\Sigma\subseteq\operatorname{SSh}(U)\).  By hypothesis \(\operatorname{Sh}(U)\subseteq\Sigma\).
The Sandwich inequalities give
\[
 |\Sigma|\leq|\operatorname{SSh}(U)|\leq|U|\leq|\operatorname{Sh}(U)|\leq|\Sigma|,
\]
so equality holds throughout, \(U\) is ample, and \(\operatorname{Sh}(U)=\Sigma\).
\end{proof}

\begin{proposition}[Ample families on seven coordinates have corners; computer-assisted]
\label{prop:damp-seven-coordinate-ample-corners}
Every nonempty ample family on at most seven coordinates has a corner.
Consequently every nonempty cornerless ample family uses at least eight
coordinates.
\end{proposition}

\begin{proof}[Computer-assisted proof]
A family on at most six coordinates has at most \(64\) vertices, and the
only one with \(64\) is a full cube, which has corners.  By
Theorem~\ref{prop:damp-minimum-order-profile} and its proof, every
nonempty cornerless ample family has at least \(64\) vertices, so the claim
holds on at most six coordinates.  A family on seven coordinates with a
constant coordinate lives on six coordinates, so we may assume all seven
coordinates vary.

Let \(C\) be a nonempty cornerless ample family on seven coordinates and let
\(d\) be its VC dimension.  Ample families of VC dimension at most two
admit corner peelings \cite[Proposition~4.11]{ChalopiChepoiMoran22}, so
\(3\leq d\leq7\).  By Lemma~\ref{lem:damp-ample-facet-cubes}, \(C\) is the
union of its facet cubes, and every vertex lies in at least two of them.
The formula for fixed \(d\) has one variable for each subset of
coordinates, saying that the subset belongs to a down-closed complex
\(\Sigma\) containing all singletons, together with facet indicators.  It
has one cube-choice variable per facet, a covering variable per vertex, and a
sequential counter demanding double coverage.  It also requires that some
set of size \(d\) lies in \(\Sigma\), that none of size \(d+1\) does, and
that the union shatters no set outside \(\Sigma\).  By the converse part of
Lemma~\ref{lem:damp-ample-facet-cubes}, the satisfying assignments are
exactly the cornerless ample families of VC dimension \(d\), with
\(\Sigma\) their shattered complex.  Lexicographic-leader constraints for
the adjacent transpositions and the coordinate reflections are added, as in
Proposition~\ref{prop:damp-small-maximum-corners}.

For each \(d=3,\ldots,7\), CaDiCaL~2.1.2 reports \textsc{unsat}, and
\texttt{drat-trim} accepts its DRAT proof.  The backward cores contain
\(1{,}299{,}213\), \(7{,}972{,}306\), \(4{,}607{,}445\), \(72{,}680\), and \(1\)
lemmas, respectively.  The formulas have \(58{,}531\) variables and at most
\(148{,}463\) clauses and are generated by
\texttt{explore\_damp\_cornerless\_ample\_sat.py} with options
\texttt{7 --lex --vc-exact=\(d\)}.  Formula and proof digests, the checker
logs, and the proofs for \(d=6,7\) are in
\path{evidence/cornerless_ample_n7/}.  The proofs for \(d=3,4,5\) are
between \(0.5\) and \(3.9\)\,GB and are not archived.  They are rechecked
by regenerating each formula, comparing its digest, rerunning the solver and
checking the fresh proof.  The same formula without the double-coverage constraint is
satisfiable on small instances, and its models were checked independently
to be ample.  With \(\Sigma\) fixed to the sets of size at most three, the
encoding reproduces the maximum-class results above.
\end{proof}

\begin{lemma}[Four-coordinate pseudogeometric VC-two core]
\label{lem:damp-qfour-maximum-vc-two-three-corners}
Let $A\subseteq Q_E$ be maximum of VC dimension two and use all four
coordinates of $E$.  Then $A$ has at least three corners.  If $A$ is
pseudogeometric, its complement is a five-vertex path and it has exactly
three corners.  Consequently, up to a cube automorphism, the only
pseudogeometric possibility is
\begin{equation}
 A^{\mathrm{pg}}_4
 =Q_4\setminus\{0000,1000,1100,1110,1111\}.
 \label{eq:damp-qfour-pseudogeometric-vctwo-core}
\end{equation}
\end{lemma}

\begin{proof}
Apply Proposition~\ref{prop:damp-five-coordinate-tree-wedges} with
$d=4$.  It identifies the corners with the two-edge paths of the
complement tree and gives at least $d-1=3$ of them.  The same proposition
and the G\"artner--Welzl criterion say that pseudogeometricity is
equivalent to the complement tree being a path.  A path whose four edges
have distinct coordinate labels is carried by a cube translation and a
coordinate permutation to
\[
 0000,1000,1100,1110,1111.
\]
Taking its complement proves the displayed normal form, and the path has
exactly three two-edge subpaths.
\end{proof}

\begin{proposition}[Bulk reduction below order forty-eight]
\label{prop:damp-bulk-reduction-below-forty-eight}
Let $H$ have minimum cardinality among all nonempty cornerless ample
families, and suppose that $|H|<48$.  Then $H$ is not maximum.  For the
maximum core $A$ and repair set $R$ supplied by
Corollary~\ref{cor:damp-least-cornerless-maximum-core}, after deleting the
constant coordinates of $A$, one has
\[
 r:=|R|\ge3,
 \qquad
 1\le \operatorname{VCdim}(A)\le2.
\]
More precisely, exactly the following two structural regimes remain.
\begin{enumerate}
 \item If $\operatorname{VCdim}(A)=1$, then $A$ is a path on $n+1$
       vertices for some $2\le n\le9$.  If $n=2$, then $3\le r\le12$;
       if $n\ge3$, then
       \begin{equation}
        (n+1)(r+1)\le40.
        \label{eq:damp-bulk-forty-eight-path-profiles}
       \end{equation}
 \item If $\operatorname{VCdim}(A)=2$, then
       $A\cong Q_3\setminus\{v\}$ and $r\in\{3,4\}$.
\end{enumerate}
Consequently, excluding these two multi-coordinate repair regimes would
raise the lower bound for the least ample forbidden pc-minor directly to
$48$ without examining the intermediate cardinalities separately.
\end{proposition}

\begin{proof}
Suppose first that $H$ is maximum, and put
$k=\operatorname{VCdim}(H)$ and $d=\operatorname{idim}(H)$.  Maximum
families of VC dimension at most two have a corner by
Lemma~\ref{lem:damp-relative-peeling-vc-two}.  For $k=3$, the inequality
\[
 \sum_{i=0}^3\binom di<48
\]
forces $d\le6$.  The cases $d\le4$ are full or once-punctured cubes,
the five-coordinate case has at least four corners by
Proposition~\ref{prop:damp-five-coordinate-tree-wedges}, and the
six-coordinate case has at least four corners by
Corollary~\ref{cor:damp-six-coordinate-corner-bound}.  For $k=4$, order
less than $48$ forces $d\le5$, again leaving only a full or once-punctured
cube.  Finally, if $k\ge5$, then either $d=k$ and $H$ is a full cube, or
\[
 |H|\ge\sum_{i=0}^k\binom{k+1}{i}=2^{k+1}-1\ge63.
\]
Every case contradicts cornerlessness.  Thus $H$ is not maximum.

Corollary~\ref{cor:damp-least-cornerless-maximum-core} now gives a proper
pseudogeometric maximum core $A$ and $r\ge2$ with
\begin{equation}
 (r+1)|A|+1\le |H|\le47.
 \label{eq:damp-bulk-forty-eight-core-budget}
\end{equation}
If $\operatorname{VCdim}(A)\ge3$, then, after its constant coordinates
are deleted, properness and maximumness give
\[
 |A|\ge1+4+\binom42+\binom43=15.
\]
Together with \eqref{eq:damp-bulk-forty-eight-core-budget}, this forces
$r=2$, $|A|=15$, and
$A\cong Q_4\setminus\{z\}$.  Write $B,C$ for the two middle repair
fibres and put $K=B\cup C$, $D=B\cap C$.  Proposition
\ref{prop:damp-unrestricted-two-coordinate-repair} gives

\[
 D\subseteq A,\qquad D\text{ ample},\qquad
 |H|=|K|+2|A|+|D|.
\]

Both middle fibres are nonempty.  Otherwise the other one contains all
of $A$; any core corner then has an endpoint copy whose unique maximal
core cube extends in the single incident repair direction, making that
copy a corner of $H$.

Since $A\subseteq K\subseteq Q_4$, either $K=A$ or $K=Q_4$.  In
the first case the repair is internal.  Proposition
\ref{prop:damp-maximum-core-corner-shield} makes $D$ a proper ample
subfamily of $A$ containing every corner of $A$, contrary to Theorem
\ref{thm:damp-nested-ample-relative-peeling-dimension-four}.

Suppose that $K=Q_4$.  The order identity gives $|D|\le1$.  The empty
support belongs to the three nonempty shadows of $A,B,C$, so
\eqref{eq:damp-unrestricted-two-coordinate-double-link} makes $D$
nonempty; hence $D$ is a singleton.  Moreover

\[
 \operatorname{Sh}(B)\cap\operatorname{Sh}(C)=\{\varnothing\}.
\]

Indeed, the double-link identity excludes every proper nonempty support,
because $A$ shatters all proper subsets of the four-coordinate set; the
full support cannot be shattered by both $B$ and $C$, since that would
make both families equal to $Q_4$.  As $B,C$ are ample, inclusion--
exclusion now gives

\[
 |\operatorname{Sh}(B)\cup\operatorname{Sh}(C)|
 =|B|+|C|-1=|K|+|D|-1=16.
\]

Thus one of $B,C$ shatters the full coordinate set and is $Q_4$; the
other has only the empty shattered support and is the singleton $D$.
Assume $B=Q_4$.  Choose a corner $a$ of $A$ outside $D$, which is
possible because $A=Q_4^-$ has four corners.  At either endpoint repair
copy of $a$, the only incident repair direction leads into $B$, and
the unique maximal core cube through $a$ is contained in $B$.  The
local corner criterion exposes both endpoint copies of $a$, a
contradiction.  The case $C=Q_4$ is symmetric.

Thus the core has VC dimension at most two.  A maximum core of VC
dimension zero is a singleton and has no active coordinate; it is the full
zero-dimensional cube, again contradicting properness.  Theorem
\ref{thm:damp-two-coordinate-low-vc-core-corner} excludes $r=2$, so
$r\ge3$.

Let first $\operatorname{VCdim}(A)=1$, and let $n$ be its active
dimension.  Pseudogeometric maximumness makes its one-inclusion graph a
path on $n+1$ vertices, and properness gives $n\ge2$.  Proposition
\ref{prop:damp-vc-one-path-core-propagation-surcharge} gives
\[
 |H|\ge
 \begin{cases}
  3(r+1)+6,&n=2,\\
  (n+1)(r+1)+7,&n\ge3.
 \end{cases}
\]
Combining this with $|H|\le47$ gives $r\le12$ in the first case and
\eqref{eq:damp-bulk-forty-eight-path-profiles} in the second.  Since
$r\ge3$, the latter inequality also gives $n\le9$.

Finally let $\operatorname{VCdim}(A)=2$ and let $n$ again be its active
dimension.  The two-ended surcharge
\eqref{eq:damp-two-ended-core-corner-surcharge} gives
$\Delta(H;A,R)\ge2|\operatorname{corn}(A)|\ge2$.  If $n\ge5$, then
$|A|\ge1+5+\binom52=16$, already contradicting
\eqref{eq:damp-bulk-forty-eight-core-budget}.  If $n=4$, then
$|A|=11$, and Lemma
\ref{lem:damp-qfour-maximum-vc-two-three-corners} together with the
two-ended surcharge gives
\[
 |H|\ge4\cdot11+2\cdot3=50,
\]
again a contradiction.  Thus $n=3$, and the only maximum VC-two family is a three-cube with one
vertex deleted.  It has three corners, so the same surcharge gives
\[
 |H|\ge7(r+1)+6.
\]
The inequality $|H|\le47$ now forces $r\le4$, completing the reduction.
\end{proof}

\begin{proposition}[Bulk reduction through order sixty-seven]
\label{prop:damp-bulk-reduction-below-sixty-four}
Let $H$ have minimum cardinality among all nonempty cornerless ample
families, and suppose that $|H|\le67$.  Then $H$ is not maximum.  For the
maximum core $A$ and repair set $R$ supplied by
Corollary~\ref{cor:damp-least-cornerless-maximum-core}, after deleting the
constant coordinates of $A$, one has
\[
 r:=|R|\ge3,
 \qquad
 1\le \operatorname{VCdim}(A)\le2.
\]
More precisely, exactly the following profiles remain:
\begin{enumerate}
 \item If $\operatorname{VCdim}(A)=1$, then $A$ is a path on $n+1$
       vertices for some $2\le n\le14$.  If $n=2$, then
       $3\le r\le19$; if $n\ge3$, then
       \begin{equation}
        (n+1)(r+1)\le60.
        \label{eq:damp-bulk-sixty-four-path-profiles}
       \end{equation}
 \item If $\operatorname{VCdim}(A)=2$ and $A$ uses three coordinates,
       then $A\cong Q_3\setminus\{v\}$ and $3\le r\le7$.
 \item If $\operatorname{VCdim}(A)=2$ and $A$ uses four coordinates,
       then $A$ is cube-isomorphic to the pseudogeometric core
       $A^{\mathrm{pg}}_4$ in
       \eqref{eq:damp-qfour-pseudogeometric-vctwo-core}, and
       $r\in\{3,4\}$.
\end{enumerate}
There are no further core isomorphism types through order \(67\).
\end{proposition}

\begin{proof}
Suppose first that $H$ is maximum, and put
$k=\operatorname{VCdim}(H)$ and $d=\operatorname{idim}(H)$.  For
$k\le2$, relative VC-two peeling gives a corner.  If $k=3$, then
\(\sum_{i=0}^3\binom di\le67\) forces $d\le7$; the cases through four
coordinates are full or once-punctured cubes, the cases $d=5,6$ have
corners by Proposition~\ref{prop:damp-five-coordinate-tree-wedges} and
Corollary~\ref{cor:damp-six-coordinate-corner-bound}, and the case $d=7$
has a corner by Corollary~\ref{cor:damp-seven-coordinate-maximum-corner}.
If $k=4$, then
$d\le6$.  The cases $d\le5$ are again full or once-punctured cubes, while
the case $d=6$ is codimension two and has at least five corners by
Proposition~\ref{prop:damp-five-coordinate-tree-wedges}.  If $k=5$, only
the full five-cube and the once-punctured six-cube have order at most $67$,
and both have corners.  Finally, if $k\ge6$, either $d=k$ and $H$ is a
full cube, or
\[
 |H|\ge\sum_{i=0}^k\binom{k+1}{i}=2^{k+1}-1\ge127.
\]
Thus $H$ cannot be maximum.

Corollary~\ref{cor:damp-least-cornerless-maximum-core} now gives a proper
pseudogeometric maximum core $A$ and a repair set of order $r\ge2$ with
\begin{equation}
 (r+1)|A|+1\le |H|\le67.
 \label{eq:damp-bulk-sixty-four-core-budget}
\end{equation}
The case $r=2$ is impossible in one step.  If
$\operatorname{VCdim}(A)\le2$, apply Theorem
\ref{thm:damp-two-coordinate-low-vc-core-corner}.  If its VC dimension is
at least three and its active dimension is at most four, apply the
four-coordinate alternative of the same theorem.  If instead its active
dimension is at least five, then
\[
 |A|\ge1+5+\binom52+\binom53=26,
\]
so $3|A|+1\ge79$, contrary to
\eqref{eq:damp-bulk-sixty-four-core-budget}.  Hence $r\ge3$.

Suppose that $\operatorname{VCdim}(A)\ge3$.  If the active dimension is
at least five, then $4|A|+1\ge105$, again impossible.  The only remaining
proper maximum core is $Q_4\setminus\{v\}$, of order fifteen.  It has four
corners, so the two-ended surcharge
\eqref{eq:damp-two-ended-core-corner-surcharge} gives
\[
 |H|\ge4\cdot15+2\cdot4=68,
\]
another contradiction.  Thus the core has VC dimension one or two.

In VC dimension one, pseudogeometric maximumness makes $A$ a path on
$n+1$ vertices, with $n\ge2$.  Proposition
\ref{prop:damp-vc-one-path-core-propagation-surcharge} gives
\[
 |H|\ge
 \begin{cases}
  3(r+1)+6,&n=2,\\
  (n+1)(r+1)+7,&n\ge3.
 \end{cases}
\]
Together with $|H|\le67$ this gives $r\le19$ in the first case and
\eqref{eq:damp-bulk-sixty-four-path-profiles} in the second.  Since
$r\ge3$, the latter inequality gives $n\le14$.

Finally suppose that the core has VC dimension two, and let $n$ be its
active dimension.  If $n\ge5$, then $|A|\ge16$.  Relative VC-two peeling
gives at least two core corners.  Indeed, apply
Lemma~\ref{lem:damp-relative-peeling-vc-two} with lower family \(\{v\}\)
to obtain a corner \(q\ne v\), and apply it again with lower family
\(\{q\}\) to obtain a second corner distinct from \(q\).  The two-ended
surcharge therefore strengthens the core
budget to \(|H|\ge4|A|+4\ge68\), impossible.  For $n=3$ the core is
$Q_3\setminus\{v\}$;
its three corners and the two-ended surcharge give
\[
 |H|\ge7(r+1)+6,
\]
and hence $r\le7$.  For $n=4$, Lemma
\ref{lem:damp-qfour-maximum-vc-two-three-corners} gives both the displayed
pseudogeometric normal form and three corners.  Thus
\[
 |H|\ge11(r+1)+6,
\]
which leaves exactly $r=3,4$.
\end{proof}

\begin{proposition}[Path-core exclusion through order sixty-seven]
\label{prop:damp-path-core-exclusion-below-sixty-four}
Let $H$ have minimum cardinality among all nonempty cornerless ample
families, suppose that $|H|\le67$, and let the maximum core supplied by
Corollary~\ref{cor:damp-least-cornerless-maximum-core} have VC dimension
one.  Then $H$ has a corner.
\end{proposition}

\begin{proof}
Use the path-section system and baseline shadow
\(\Sigma_0\) from Corollary
\ref{cor:damp-joint-path-budget-anchored-normal-form}.  The
support-interval exact sequence, the two-sweep reduction, and
Corollary~\ref{cor:damp-terminal-remainders-no-residual-debt} cancel every
successful acquisition and place every remaining neutral or failure
demand in the original excess
\(\operatorname{Sh}(H)\setminus\Sigma_0\).  Corollaries
\ref{cor:damp-joint-path-budget-anchored-normal-form} and
\ref{cor:damp-neutral-cycle-failure-provenance-normal-form} reduce a
saturated cornerless state to a same-position anchored packet or to
packet-free equality.

If a packet ascends to support $L$, some facet $T$ of
\(\operatorname{Sh}(H)\) contains $L$.  The uniform facet-rank bound gives
\[
 |H|\ge2^{|T|+1}+4\ge68
\]
when $|L|\ge5$; here we use
Corollary~\ref{cor:damp-uniform-facet-rank-bound}.  Hence every packet has
rank three or four.  The
mixed rank-three profile separates by Corollary
\ref{cor:damp-mixed-rank-three-packet-separates}; the pure profile either
ascends or has a maximal repair three-cube owner by Corollary
\ref{cor:damp-pure-rank-three-packet-maximal-owner}, and the latter has
strict excess by Propositions
\ref{prop:damp-maximal-repair-three-cube-exterior-tax} and
\ref{prop:damp-three-two-cycles-seven-section-units}.  Rank four is closed
by Corollary~\ref{cor:damp-rank-four-packets-closed}.  In the packet-free
case, Corollary~\ref{cor:damp-packet-free-saturation-exposes-corner} makes
the whole shattered complex two-dimensional and exposes a corner.  All
alternatives contradict cornerlessness.
\end{proof}

\begin{corollary}[VC-two core frontier through order sixty-seven]
\label{cor:damp-vctwo-core-frontier-below-sixty-four}
Let $H$ have minimum cardinality among all nonempty cornerless ample
families and suppose that $|H|\le67$.  Then its maximum core has VC dimension
two and is one of exactly the following normalized types:
\[
 \begin{array}{c|c|c}
 \text{core}&|A|&\text{repair rank}\\ \hline
 Q_3\setminus\{v\}&7&3\le r\le7,\\
 A^{\mathrm{pg}}_4&11&r\in\{3,4\}.
 \end{array}
\]
\end{corollary}

\begin{proof}
Apply Proposition~\ref{prop:damp-bulk-reduction-below-sixty-four} and
exclude its path row by Proposition
\ref{prop:damp-path-core-exclusion-below-sixty-four}.
\end{proof}

\begin{proposition}[Two-ended facet-orbit quotient]
\label{prop:damp-vctwo-two-ended-facet-orbit-quotient}
Keep the hypotheses and notation of
Corollary~\ref{cor:damp-vctwo-core-frontier-below-sixty-four}, and write
the two equal antipodal core fibres as
\[
 H_{0_R}=H_{1_R}=A.
\]
Choose a corner $a$ of the normalized core whose unique maximal core cube
is a square, and denote its support by $S$.  At each of the two endpoint
copies of this square, choose a maximal cube of $H$ that contains it.  The
two supports have the form
\begin{equation}
 S\mathbin{\dot\cup}U,
 \qquad
 S\mathbin{\dot\cup}V,
 \label{eq:damp-vctwo-endpoint-extension-supports}
\end{equation}
where
\begin{equation}
 1\le |U|,|V|\le2,
 \qquad
 U\not\subseteq V,
 \qquad
 V\not\subseteq U.
 \label{eq:damp-vctwo-endpoint-extension-incomparability}
\end{equation}
After permuting the repair coordinates, the ordered pair $(U,V)$ is
therefore exactly one of
\begin{equation}
\begin{array}{c|c|c}
\text{type}&U&V\\ \hline
1/1&\{1\}&\{2\},\\
1/2&\{1\}&\{2,3\},\\
2/1&\{1,2\}&\{3\},\\
2/2\text{, intersecting}&\{1,2\}&\{1,3\},\\
2/2\text{, disjoint}&\{1,2\}&\{3,4\}.
\end{array}
\label{eq:damp-vctwo-endpoint-extension-orbits}
\end{equation}
The last row occurs only when $r\ge4$.  Thus the entire VC-two frontier
through order $67$ reduces to four ordered extension types in repair rank
three and five in every larger repair rank.  If the two antipodal repair
endpoints are left unoriented, the $1/2$ and $2/1$ rows are equivalent.
Hence there are only three genuine types in repair rank three and four in
every larger repair rank; the reduction does not branch on the cardinality
of $H$.
\end{proposition}

\begin{proof}
The chosen square in the endpoint fibre is contained in a maximal cube of
$H$.  Its support contains $S$ and contains no further core coordinate,
because restriction to that endpoint fibre would otherwise give a core
cube properly containing the unique maximal core cube through $a$.
Consequently its support is $S\mathbin{\dot\cup}U$ for some $U\subseteq R$;
the same argument at the opposite endpoint gives
$S\mathbin{\dot\cup}V$.

The least cornerless family $H$ is periphery-free.  If either displayed
maximal cube had rank at least five, the uniform facet-rank bound,
Corollary~\ref{cor:damp-uniform-facet-rank-bound}, would give
\[
 |H|\ge 2^6+4=68,
\]
contrary to $|H|\le67$.  Since $|S|=2$, this proves
$|U|,|V|\le2$.

The two repair supports are nonempty.  Indeed, if $U=\varnothing$, the
core square at the first endpoint would itself be a maximal cube.  The
copy of that square at the opposite endpoint lies in some maximal cube.
Its support either properly contains $S$, contrary to the antichain
property of maximal-cube supports, or equals $S$, contrary to the fact
that distinct maximal cubes of an ample family have distinct supports.
The same argument excludes $V=\varnothing$.

The supports in \eqref{eq:damp-vctwo-endpoint-extension-supports} belong
to distinct maximal cubes: their repair projections contain the two
endpoints of $Q_R$, whose distance is $r\ge3$, whereas either cube uses at
most two repair coordinates.  Maximal-cube supports are an antichain, and
distinct maximal cubes of an ample family cannot have the same support.
Hence the two supports, and therefore $U$ and $V$, are incomparable.

It remains only to take orbits under permutations of $R$.  Normalize $U$
to $\{1\}$ or $\{1,2\}$.  If $|U|=1$, incomparability makes $V$ disjoint
from $U$, leaving the first two rows.  If $|U|=2$ and $|V|=1$, the unique
element of $V$ lies outside $U$.  If both sets have order two, their
intersection has order one or zero, giving the final two rows.  These are
all the stabilizer orbits listed in
\eqref{eq:damp-vctwo-endpoint-extension-orbits}.

Finally, complement every repair coordinate and leave every core
coordinate fixed.  This cube automorphism exchanges the two antipodal
repair fibres, preserves their common core $A$, and replaces $(U,V)$ by
$(V,U)$.  It therefore identifies the $1/2$ and $2/1$ rows.  The other
three rows are fixed as types by endpoint exchange, proving the unoriented
count.
\end{proof}

\begin{proposition}[Core-coordinate wing squeeze]
\label{prop:damp-vctwo-core-coordinate-wing-squeeze}
Keep the hypotheses and notation of
Corollary~\ref{cor:damp-vctwo-core-frontier-below-sixty-four}.  Let \(X\)
be the active coordinate set of the maximum VC-dimension-two core \(A\),
put \(n=|X|\), and write \(m=|H|\).  For every \(x\in X\),
\begin{equation}
 \deg_{\operatorname{Sh}(H)}(x)
 \le \left\lfloor\frac{m-20}{2}\right\rfloor.
 \label{eq:damp-vctwo-core-coordinate-degree-cap}
\end{equation}
In particular,
\begin{equation}
 m\ge 20+2n(r+1).
 \label{eq:damp-vctwo-core-coordinate-baseline-bound}
\end{equation}

Choose a square core corner and the two endpoint extensions
\(S\mathbin{\dot\cup}U\) and \(S\mathbin{\dot\cup}V\) from
Proposition~\ref{prop:damp-vctwo-two-ended-facet-orbit-quotient}.  Let
\[
 t=\bigl|\{W\in\{U,V\}:|W|=2\}\bigr|.
\]
Then
\begin{equation}
 m\ge20+2n(r+1)+4t.
 \label{eq:damp-vctwo-endpoint-rank-charge}
\end{equation}
Consequently, below order \(64\), the initial \(26\) unoriented endpoint
orbits reduce analytically to the following \(14\):
\begin{equation}
\begin{array}{c|c|c}
\text{core}&r&\text{possible endpoint types}\\ \hline
Q_3\setminus\{v\}&3&1/1,\ 1/2,\ 2/2\text{ intersecting}\\
Q_3\setminus\{v\}&4&1/1,\ 1/2,\ 2/2\text{ intersecting or disjoint}\\
Q_3\setminus\{v\}&5&1/1,\ 1/2\\
Q_3\setminus\{v\}&6&1/1\\
A^{\mathrm{pg}}_4&3&1/1,\ 1/2,\ 2/2\text{ intersecting}\\
A^{\mathrm{pg}}_4&4&1/1.
\end{array}
\label{eq:damp-vctwo-wing-squeezed-orbits}
\end{equation}
There is no repair-rank-seven row.  At this stage the unresolved frontier
consists of one parametric \(1/1\) branch at repair ranks four, five, and
six, rather than unrelated cardinality cases.  The next proposition removes
the rank-six endpoint of this branch without a finite search.
\end{proposition}

\begin{proof}
Fix an active coordinate \(e\) of \(H\).  Write its two sections as
\(A_e^0,A_e^1\) and put \(B_e=A_e^0\cap A_e^1\).  The two-layer shadow
identity makes \(B_e\) ample and gives
\[
 |B_e|=\deg_{\operatorname{Sh}(H)}(e).
\]
All three families have smaller order than \(H\), so
Lemma~\ref{lem:damp-least-cornerless-cardinality-minimality} puts them in
\(\Damp\).  Since \(H\) is periphery-free, both inclusions
\(B_e\subsetneq A_e^b\) are proper.  No corner of either upper section lies
outside \(B_e\), for such a corner would lift to a corner of \(H\).
Proposition~\ref{prop:damp-nine-concept-passive-carrier-exclusion} therefore
gives
\[
 |A_e^b\setminus B_e|\ge10\qquad(b\in\{0,1\}).
\]
Hence
\[
 m=|A_e^0\setminus B_e|+|A_e^1\setminus B_e|+2|B_e|
 \ge20+2\deg_{\operatorname{Sh}(H)}(e),
\]
which proves~\eqref{eq:damp-vctwo-core-coordinate-degree-cap}.

Both normalized cores are maximum of VC dimension two, so
\(\operatorname{Sh}(A)=\binom X{\le2}\).  The equal antipodal core fibres
force
\[
 \binom X{\le2}*
 \left(\{\varnothing\}\cup\binom R1\right)
 \subseteq\operatorname{Sh}(H).
\]
For a fixed \(x\in X\), exactly \(n\) members of
\(\binom X{\le2}\) contain \(x\).  The forced join therefore contributes
\(n(r+1)\) faces through \(x\).  Substitution in the degree cap proves
\eqref{eq:damp-vctwo-core-coordinate-baseline-bound}.

Now fix \(x\in S\).  If \(W\in\{U,V\}\) has order two, the facet support
\(S\mathbin{\dot\cup}W\) is shattered.  Downward closure supplies the two
faces
\[
 \{x\}\mathbin{\dot\cup}W,
 \qquad
 S\mathbin{\dot\cup}W.
\]
Neither belongs to the forced join, because each contains two repair
coordinates.  Distinct two-element members of \(\{U,V\}\) give distinct
faces; if both have order two, their incomparability in
\eqref{eq:damp-vctwo-endpoint-extension-incomparability} makes them
distinct.  Thus
\[
 \deg_{\operatorname{Sh}(H)}(x)\ge n(r+1)+2t.
\]
The coordinate-degree cap now gives
\eqref{eq:damp-vctwo-endpoint-rank-charge}.

For \(Q_3\setminus\{v\}\) one has \(n=3\), so the lower bounds for
\(t=0,1,2\) are \(6r+26,6r+30,6r+34\).  For \(A^{\mathrm{pg}}_4\) one has
\(n=4\), and they are \(8r+28,8r+32,8r+36\).  Requiring these quantities
to be at most \(63\), and remembering that the disjoint \(2/2\) type needs
at least four repair coordinates, gives exactly the table in
\eqref{eq:damp-vctwo-wing-squeezed-orbits}.  Its row counts are
\(3+4+2+1+3+1=14\).  The only rows in that table not already covered by
the stored proof certificates are the \(1/1\) rows with
\(Q_3\setminus\{v\}\) and \(r\in\{4,5,6\}\), proving the final assertion.
\end{proof}

\begin{corollary}[The order-sixty-four--sixty-five endpoint slab]
\label{cor:damp-vctwo-endpoint-slab-sixty-four-sixty-five}
Let \(H\) have minimum cardinality among all nonempty cornerless ample
families and suppose that \(64\le |H|\le65\).  After exchanging the two
repair endpoints, its core, repair rank, and endpoint type belong to one
of the following seventeen orbits, grouped into six rows:
\[
\begin{array}{c|c|c|c}
\text{core}&r&\text{available degree slots}&
\text{possible endpoint types}\\ \hline
Q_3\setminus\{v\}&3&10&
 1/1,\ 1/2,\ 2/2\text{ intersecting}\\
Q_3\setminus\{v\}&4&7&
 1/1,\ 1/2,\ 2/2\text{ intersecting or disjoint}\\
Q_3\setminus\{v\}&5&4&
 1/1,\ 1/2,\ 2/2\text{ intersecting or disjoint}\\
Q_3\setminus\{v\}&6&1&1/1\\
A^{\mathrm{pg}}_4&3&6&
 1/1,\ 1/2,\ 2/2\text{ intersecting}\\
A^{\mathrm{pg}}_4&4&2&1/1,\ 1/2.
\end{array}
\]
In particular, orders \(64\) and \(65\) form one structural slab: the
core-coordinate degree allowance and the endpoint-orbit list do not
depend on which of the two cardinalities occurs.
\end{corollary}

\begin{proof}
The strengthened core frontier and the two-ended facet-orbit quotient
apply through order \(65\).  For both possible orders,
\[
 \left\lfloor\frac{|H|-20}{2}\right\rfloor=22.
\]
After the forced join is removed, the number of available supports through
one core coordinate is therefore
\[
 22-n(r+1),
\]
where \(n=3\) for \(Q_3\setminus\{v\}\) and \(n=4\) for
\(A^{\mathrm{pg}}_4\).  This gives the third column.

The lower bound
\eqref{eq:damp-vctwo-endpoint-rank-charge} is
\(6r+26+4t\) for the three-coordinate core and \(8r+28+4t\) for the
four-coordinate core, where \(t\) counts endpoint repair supports of order
two.  Retain the types whose bound is at most \(65\), and remember that a
disjoint \(2/2\) type needs at least four repair coordinates.  The six row
counts are \(3,4,4,1,3,2\), which sum to seventeen.
\end{proof}

\begin{corollary}[The order-sixty-six--sixty-seven endpoint slab]
\label{cor:damp-vctwo-endpoint-slab-sixty-six-sixty-seven}
Let \(H\) have minimum cardinality among all nonempty cornerless ample
families and suppose that \(66\le |H|\le67\).  After exchanging the two
repair endpoints, its core, repair rank, and endpoint type belong to one
of the following fifteen orbits, grouped into five rows:
\[
\begin{array}{c|c|c|c}
\text{core}&r&\text{available degree slots}&
\text{possible endpoint types}\\ \hline
Q_3\setminus\{v\}&4&8&
 1/1,\ 1/2,\ 2/2\text{ intersecting or disjoint}\\
Q_3\setminus\{v\}&5&5&
 1/1,\ 1/2,\ 2/2\text{ intersecting or disjoint}\\
Q_3\setminus\{v\}&6&2&1/1,\ 1/2\\
A^{\mathrm{pg}}_4&3&7&
 1/1,\ 1/2,\ 2/2\text{ intersecting}\\
A^{\mathrm{pg}}_4&4&3&1/1,\ 1/2.
\end{array}
\]
Thus orders \(66\) and \(67\) form one structural slab.  Relative to the
preceding slab, the three repair-rank-three punctured-core orbits disappear,
and the only new endpoint orbit is the \(1/2\) type for
\(Q_3\setminus\{v\}\) at repair rank six.
\end{corollary}

\begin{proof}
Corollary~\ref{cor:damp-vctwo-core-frontier-below-sixty-four} and
Proposition~\ref{prop:damp-vctwo-two-ended-facet-orbit-quotient} apply
through order \(67\).  For either possible order,
\[
 \left\lfloor\frac{|H|-20}{2}\right\rfloor=23.
\]
The number of degree slots beyond the forced join is therefore
\(23-n(r+1)\), with \(n=3\) for \(Q_3\setminus\{v\}\) and \(n=4\) for
\(A^{\mathrm{pg}}_4\).  This gives the third column.

The endpoint-rank charge
\eqref{eq:damp-vctwo-endpoint-rank-charge} retains the listed types.  It
removes the repair-rank-three punctured-core row because that profile uses
only \(3+3=6\) active coordinates and hence has order at most \(2^6=64\).
It excludes repair rank seven because its three-coordinate-core baseline is
\(20+2\cdot3\cdot8=68\), and it excludes both \(2/2\) types in the
repair-rank-six row because either would cost at least
\(20+2\cdot3\cdot7+8=70\).  The five row counts are
\(4,4,2,3,2\), which sum to fifteen.
\end{proof}

\begin{corollary}[The unified order-sixty-four--sixty-seven endpoint slab]
\label{cor:damp-vctwo-endpoint-slab-sixty-four-sixty-seven}
Let \(H\) have minimum cardinality among all nonempty cornerless ample
families and suppose that \(64\le |H|\le67\).  After exchanging the two
repair endpoints, its core, repair rank, and endpoint type belong to one
of the following eighteen orbits:
\[
\begin{array}{c|c|c}
\text{core}&r&\text{possible endpoint types}\\ \hline
Q_3\setminus\{v\}&3&1/1,\ 1/2,\ 2/2\text{ intersecting}\\
Q_3\setminus\{v\}&4&1/1,\ 1/2,\ 2/2\text{ intersecting or disjoint}\\
Q_3\setminus\{v\}&5&1/1,\ 1/2,\ 2/2\text{ intersecting or disjoint}\\
Q_3\setminus\{v\}&6&1/1,\ 1/2\\
A^{\mathrm{pg}}_4&3&1/1,\ 1/2,\ 2/2\text{ intersecting}\\
A^{\mathrm{pg}}_4&4&1/1,\ 1/2.
\end{array}
\]
Thus all four orders form one structural slab with six core--rank rows and
eighteen endpoint orbits.  In particular, a certificate for this common
list can move the lower bound from \(64\) directly to \(68\).
\end{corollary}

\begin{proof}
Apply Corollary
\ref{cor:damp-vctwo-endpoint-slab-sixty-four-sixty-five} when
\(|H|\le65\), and Corollary
\ref{cor:damp-vctwo-endpoint-slab-sixty-six-sixty-seven} otherwise.
The displayed table is the union of their orbit lists.  It has row counts
\(3,4,4,2,3,2\), whose sum is eighteen.  Among the five shared rows, the
only type present in just one subrange is the rank-six \(1/2\) type.  The
entire repair-rank-three punctured-core row can occur only at order \(64\),
since it uses six active coordinates.  Hence the union is exhaustive
throughout the stated range.
\end{proof}

\begin{lemma}[Core-carrier defect ledger]
\label{lem:damp-qthree-minus-core-carrier-defect-ledger}
Let \(H\) have minimum cardinality among all nonempty cornerless ample
families, suppose that \(|H|<64\), and suppose that its maximum core is

\[
 A=Q_3\setminus\{v\}
\]

on the active coordinate set \(X\), with a repair set \(R\) of order
\(r\in\{4,5,6\}\).  Put \(\Sigma=\operatorname{Sh}(H)\).  For each
\(S\in\binom X2\), let \(C_S\subseteq Q_R\) be the repair projection of
the full positioned \(S\)-cube carrier, and put

\[
 q_S=
 \left|
  \operatorname{Sh}(C_S)
  \setminus
  \left(\{\varnothing\}\cup\binom R1\right)
 \right|.
\]

For \(x\in X\), also put

\[
 p_x=
 \left|
  \left\{T\subseteq R:
   |T|\ge2,
   \{x\}\mathbin{\dot\cup}T\in\Sigma
  \right\}
 \right|,
 \qquad
 h=\boldsymbol 1_{\{X\in\Sigma\}}.
\]

Define the defect allowance

\[
 \delta(H,r)=
 \left\lfloor\frac{|H|-20}{2}\right\rfloor-3(r+1).
\]

Then \(0\le\delta(H,r)\le3(6-r)\).  If
\(S,S'\in\binom X2\) are the two core squares containing \(x\), then

\begin{equation}
 h+q_S+q_{S'}+p_x\le\delta(H,r).
 \label{eq:damp-qthree-minus-coordinate-defect-ledger}
\end{equation}

Consequently,

\begin{equation}
 3h+2\sum_{S\in\binom X2}q_S+
 \sum_{x\in X}p_x
 \le3\delta(H,r)\le9(6-r).
 \label{eq:damp-qthree-minus-total-defect-ledger}
\end{equation}

For later use, the allowance at the four odd order thresholds is

\[
\begin{array}{c|cccc}
 \text{upper bound on }|H|&57&59&61&63\\ \hline
 r=4&3&4&5&6\\
 r=5&0&1&2&3.
\end{array}
\]

For every \(S\in\binom X2\) and \(x\in S\), one also has

\begin{equation}
 q_S\le p_x,
 \qquad
 h+2q_S+q_{S'}\le\delta(H,r),
 \label{eq:damp-qthree-minus-carrier-to-singleton-defect}
\end{equation}

where \(S'\) is the other core square containing \(x\).

If an endpoint extension of the core square \(S\) uses a two-element
repair support \(T\), then \(T\) contributes one unit to \(q_S\) and one
unit to \(p_x\) for each \(x\in S\).  If \(C_S,C_{S'}\) are antipodal
geodesics and their path orders disagree on a repair pair \(T\), then \(T\)
contributes one unit to \(p_x\), where \(\{x\}=S\cap S'\).  Thus at
repair rank five all carrier defects,
missing-core concepts, endpoint two-supports, and carrier-order disagreements
must share a total defect budget of nine.  At repair rank six every quantity
in the two displays vanishes.
\end{lemma}

\begin{proof}
The forced join in the proof of Proposition
\ref{prop:damp-vctwo-core-coordinate-wing-squeeze} contributes exactly
\(3(r+1)\) shattered supports through each \(x\in X\).  The coordinate
degree cap therefore leaves at most

\[
 \left\lfloor\frac{|H|-20}{2}\right\rfloor-3(r+1)
 =\delta(H,r)
\]

further supports through \(x\).  This quantity is nonnegative because the
forced join is present.  Since \(|H|\le63\), it is at most \(3(6-r)\).

The support-conservation identity for a projected cube carrier,
Lemma~\ref{lem:damp-projected-cube-carrier-support-conservation}, says that
every

\[
 T\in
 \operatorname{Sh}(C_S)
 \setminus
 \left(\{\varnothing\}\cup\binom R1\right)
\]

supplies the shattered support \(S\mathbin{\dot\cup}T\) of \(H\).
For the two squares \(S,S'\) containing \(x\), these give
\(q_S+q_{S'}\) distinct nonbaseline supports through \(x\).  The
\(p_x\) supports have core part exactly \(\{x\}\), while \(X\), when it
is shattered, has core part \(X\).  These three groups are pairwise
disjoint and none belongs to the forced join.  Counting them inside the
available \(\delta(H,r)\) slots proves
\eqref{eq:damp-qthree-minus-coordinate-defect-ledger}.

Sum that inequality over \(x\in X\).  Each square \(S\) contains two
core coordinates, so its term \(q_S\) occurs twice; \(h\) occurs three
times.  This proves
\eqref{eq:damp-qthree-minus-total-defect-ledger}.

For each support \(T\) counted by \(q_S\), support conservation puts
\(S\mathbin{\dot\cup}T\) in \(\Sigma\).  Downward closure then puts
\(\{x\}\mathbin{\dot\cup}T\) in \(\Sigma\) for either \(x\in S\).
Distinct supports \(T\) remain distinct, so \(q_S\le p_x\).  Substitution
in~\eqref{eq:damp-qthree-minus-coordinate-defect-ledger} gives the second
inequality in~\eqref{eq:damp-qthree-minus-carrier-to-singleton-defect}.

Finally, if a maximal endpoint cube has support
\(S\mathbin{\dot\cup}T\) with \(|T|=2\), support conservation puts
\(T\) in the acquired shadow of \(C_S\).  Downward closure also puts
\(\{x\}\mathbin{\dot\cup}T\) in \(\Sigma\) for both \(x\in S\).
This proves the endpoint charge.

Now suppose that \(C_S,C_{S'}\) are antipodal geodesics and their orders
disagree on \(T=\{e,f\}\).  Their union realizes both orders of \(e,f\)
and therefore shatters \(T\).  Both full square carriers contain the
common core direction \(x\in S\cap S'\).  Whichever carrier supplies a
given trace on \(T\) also supplies that trace on both \(x\)-sides.  Hence
\(\{x\}\mathbin{\dot\cup}T\in\Sigma\), so \(T\) is counted by \(p_x\).
When \(r=6\), the right side of
\eqref{eq:damp-qthree-minus-total-defect-ledger} is zero, giving the final
assertion.
\end{proof}

\begin{lemma}[Zero-carrier endpoint charge]
\label{lem:damp-qthree-minus-zero-carrier-endpoint-charge}
In the setting of
Lemma~\ref{lem:damp-qthree-minus-core-carrier-defect-ledger}, suppose that
\(H\) is cornerless and that

\[
 q_S=0\qquad\text{for every }S\in\binom X2.
\]

Then

\begin{equation}
 \sum_{x\in X}p_x\ge
 \begin{cases}
  3,&r=4,\\
  2,&r=5.
 \end{cases}
 \label{eq:damp-qthree-minus-zero-carrier-endpoint-charge}
\end{equation}
In particular, the zero-carrier branch is controlled by the bounded defect
intervals

\[
\begin{array}{c|cc}
 &h=0&h=1\\ \hline
 r=4&3\le\sum_xp_x\le18&3\le\sum_xp_x\le15\\
 r=5&2\le\sum_xp_x\le 9&2\le\sum_xp_x\le 6.
\end{array}
\]
\end{lemma}

\begin{proof}
Translate the core so that its missing word is \(111\).  Label its three
maximal squares \(S_0,S_1,S_2\), and let \(a_i\) be the corner of the core
belonging only to \(S_i\).  Thus \(a_i\) is the characteristic word of
\(S_i\).  For \(i\ne j\), the squares \(S_i,S_j\) share one core
coordinate, say \(x_{ij}\), whereas

\begin{equation}
 a_i\mathbin\triangle a_j=X\setminus\{x_{ij}\}.
 \label{eq:damp-zero-carrier-corner-difference}
\end{equation}

Since \(q_{S_i}=0\), support conservation shows that the repair projection
of the full \(S_i\)-carrier has shadow

\[
 \{\varnothing\}\cup\binom R1.
\]

It is ample and contains \(0_R,1_R\), so it is an antipodal geodesic
\(P_i\).  Write \(\pi_i\) for the order in which its repair coordinates
are crossed.  The repair section

\[
 F_i=\{g\in Q_R:(g,a_i)\in H\}
\]

contains \(P_i\).  At either endpoint of \(P_i\), the two core directions
in \(S_i\) together with the incident path direction already span the
product of the unique core square through \(a_i\) and the endpoint path
edge.  Since \(H\) is cornerless and there is no further core direction
at \(a_i\), the section \(F_i\) contains both

\[
 \{j_i\},\qquad R\setminus\{k_i\},
\]

where \(j_i\) is not the first and \(k_i\) is not the last coordinate in
the order \(\pi_i\).  If \(E_i\subseteq\binom R2\) is the pair shadow of
\(F_i\), it therefore contains the two endpoint stars

\begin{equation}
 \begin{aligned}
 L_i&=\{\{e,j_i\}:e<_{\pi_i}j_i\},\\
 R_i&=\{\{k_i,e\}:k_i<_{\pi_i}e\}.
 \end{aligned}
 \label{eq:damp-zero-carrier-endpoint-stars}
\end{equation}

We record how these pair shadows are charged to the defect ledger.  Let
\(I_{ij}\) be the inversion set of the two orders \(\pi_i,\pi_j\).
Lemma~\ref{lem:damp-qthree-minus-core-carrier-defect-ledger} puts every
member of \(I_{ij}\) among the pairs counted by \(p_{x_{ij}}\).  On the
other hand, if \(T\in E_i\cap E_j\), then each of the two repair sections
realizes every trace on \(T\).  Equation
\eqref{eq:damp-zero-carrier-corner-difference} consequently puts
\(\{x\}\mathbin{\dot\cup}T\) in the shadow of \(H\) for both
\(x\in X\setminus\{x_{ij}\}\).  Hence, on defining

\begin{equation}
 D_x=I_{ij}\cup
 \bigcup_{\substack{k<\ell\\x\ne x_{k\ell}}}(E_k\cap E_\ell),
 \qquad x=x_{ij},
 \label{eq:damp-zero-carrier-charged-pairs}
\end{equation}

we have

\begin{equation}
 \sum_{x\in X}p_x\ge\sum_{x\in X}|D_x|.
 \label{eq:damp-zero-carrier-charge-lower-bound}
\end{equation}

Suppose first that \(r=5\).  If the three orders are not all equal, their
three pairwise Kendall distances have sum at least two.  The inversion
sets in~\eqref{eq:damp-zero-carrier-charged-pairs} then already give the
required lower bound.  If the three orders are equal, assume that the
sets \(E_0,E_1,E_2\) are pairwise disjoint.  Number positions in the
common order from \(0\) to \(4\), and write \(j_i,k_i\) also for the
positions of these coordinates.  Pairwise disjointness of the stars in
\eqref{eq:damp-zero-carrier-endpoint-stars} forces the \(j_i\) to be
distinct, the \(k_i\) to be distinct, and

\[
 j_i\le k_\ell\qquad(i\ne\ell).
\]

After relabelling, let \(j_0<j_1<j_2\).  Then \(j_2\ge3\), so the two
distinct values \(k_0,k_1\) must both belong to
\(\{3\}\), a contradiction.  Some \(E_i,E_j\) therefore share a pair.
That pair belongs to the two sets \(D_x\) with
\(x\ne x_{ij}\), and again the total charge is at least two.

It remains to prove the rank-four bound.  If the sum of the three Kendall
distances is at least three, there is nothing to prove.  We treat the only
two smaller possibilities explicitly.  For one fixed order \(0123\),
the nine possible unions \(L_i\cup R_i\) are displayed below; we abbreviate
\(\{a,b\}\) to \(ab\).

\begin{equation}
\begin{array}{c|ccc}
 &k=0&k=1&k=2\\ \hline
 j=1&01,02,03&01,12,13&01,23\\
 j=2&01,02,03,12&02,12,13&02,12,23\\
 j=3&01,02,03,13,23&03,12,13,23&03,13,23.
\end{array}
\label{eq:damp-rank-four-endpoint-star-table}
\end{equation}

If all three orders agree, three entries of this table contribute at
least three to the right side of
\eqref{eq:damp-zero-carrier-charge-lower-bound}.  Here is a short check
that does not conceal a case.  A total charge at most two would allow at
most one pair of entries to meet, and their intersection would have order
one; the third entry would have to be disjoint from both.  The unique
two-edge entry is \(\{01,23\}\).  Any two two-edge entries therefore meet
in two edges.  If the disjoint third entry is \(\{01,23\}\), the table has
only one entry disjoint from it, namely \(\{02,12,13\}\), so the other two
entries coincide in at least three edges.  If the disjoint third entry
has at least three edges, its complement has at most three edges.  The
other two entries would then both have to be the unique two-edge entry
in order to meet in only one edge and fit inside that complement; they
instead meet in two edges.  Thus the charge is at least three.

Finally, if the sum of the Kendall distances is two, two orders agree and
the third differs from them by one adjacent transposition.  Reversal of
the common order identifies the first and last transpositions.  Relabel
the coordinates in
\eqref{eq:damp-rank-four-endpoint-star-table} by the transposition and
compare the nine entries.  The resulting minima are

\[
\begin{array}{c|ccc}
 \text{third order}&1023&0213&0132\\ \hline
 \min\sum_x|D_x|&3&5&3.
\end{array}
\]

This covers every order triple with total inversion charge below three
and proves the rank-four assertion.  The upper endpoints in the final
table of the statement are exactly
\eqref{eq:damp-qthree-minus-total-defect-ledger} with \(q_S=0\).
\end{proof}

\begin{lemma}[Rank-five allowance-one normal form]
\label{lem:damp-qthree-minus-rank-five-allowance-one-normal-form}
In the setting of
Lemma~\ref{lem:damp-qthree-minus-core-carrier-defect-ledger}, suppose that
the repair rank is five and that \(\delta(H,5)\le1\).  Then
\[
 \delta(H,5)=1,\qquad
 h=0,\qquad q_S=0\quad\left(S\in\binom X2\right),
 \qquad
 2\le\sum_{x\in X}p_x\le3,
 \qquad p_x\le1\quad(x\in X).
\]
Consequently, up to permuting the three core coordinates, the vector
\((p_x)_{x\in X}\) is either \((1,1,0)\) or \((1,1,1)\).
\end{lemma}

\begin{proof}
Suppose that \(q_S>0\) for some core square \(S\), and choose \(x\in S\).
Equation~\eqref{eq:damp-qthree-minus-carrier-to-singleton-defect} gives
\(p_x\ge q_S\ge1\).  The coordinate defect inequality at \(x\) would give
\[
 2\le q_S+p_x
 \le h+q_S+q_{S'}+p_x
 \le\delta(H,5)\le1,
\]
where \(S'\) is the other core square containing \(x\), a contradiction.
Thus every \(q_S\) vanishes.

Lemma~\ref{lem:damp-qthree-minus-zero-carrier-endpoint-charge} now gives
\(\sum_xp_x\ge2\).  Since \(\delta(H,5)\) is a nonnegative integer, the
total defect inequality first forces \(\delta(H,5)=1\) and then gives
\[
 3h+\sum_xp_x\le3\delta(H,5)\le3,
\]
so \(h=0\) and \(\sum_xp_x\le3\).  Finally, the coordinate defect
inequality reduces to \(p_x\le1\) for every \(x\).  The two listed vectors
are the only possibilities.
\end{proof}

\begin{lemma}[Rank-five allowance-one carrier orders]
\label{lem:damp-qthree-minus-rank-five-allowance-one-carrier-orders}
Under the hypotheses of
Lemma~\ref{lem:damp-qthree-minus-rank-five-allowance-one-normal-form},
label the three maximal core squares \(S_0,S_1,S_2\), and let
\(\pi_i\) be the repair-coordinate order of the antipodal geodesic
carrier of \(S_i\).  If \(x_{ij}\) is the core coordinate common to
\(S_i\) and \(S_j\), then
\[
 d_{\mathrm K}(\pi_i,\pi_j)\le p_{x_{ij}}\le1,
\]
where \(d_{\mathrm K}\) denotes Kendall distance.  Consequently the three
orders assume at most two values; if two values occur, they differ by one
adjacent transposition.

In the normalized \(1/1\) endpoint branch, let the prescribed near and far
repair coordinates of \(S_0\) be \(0\) and \(1\).  After permuting the
other three repair coordinates, one may take
\[
 \pi_0=(0,2,3,4,1).
\]
There are then only thirteen labelled possibilities for
\((\pi_1,\pi_2)\): the pair \((\pi_0,\pi_0)\), and, for each of the four
orders \(\tau\) obtained from \(\pi_0\) by one adjacent transposition, the
three pairs
\[
 (\pi_0,\tau),\qquad(\tau,\pi_0),\qquad(\tau,\tau).
\]
\end{lemma}

\begin{proof}
The allowance-one normal form gives \(q_{S_i}=0\) for every \(i\), so the
proof of
Lemma~\ref{lem:damp-qthree-minus-zero-carrier-endpoint-charge} identifies
each carrier with an antipodal geodesic.  The inversion set of
\((\pi_i,\pi_j)\) injects into the repair supports counted by
\(p_{x_{ij}}\), by
Lemma~\ref{lem:damp-qthree-minus-core-carrier-defect-ledger}.  Its order is
the Kendall distance, proving the displayed inequalities.

The graph of permutations joined by adjacent transpositions is bipartite,
and in particular triangle-free.  Three permutations at pairwise distance
at most one therefore assume at most two values, and two distinct values
are adjacent.  In the fixed endpoint representative, the selected maximal
cube at \(0_R\) makes coordinate \(0\) the first edge of the \(S_0\)
carrier, while the selected maximal cube at \(1_R\) makes coordinate \(1\)
its last edge.  The remaining coordinates are symmetric, giving the stated
normalization of \(\pi_0\).  Relative to this order, each of
\(\pi_1,\pi_2\) is either \(\pi_0\) or an adjacent order; triangle-freeness
forces the same adjacent order whenever both differ from \(\pi_0\).
This gives \(1+4\cdot3=13\) labelled pairs.
\end{proof}

\begin{lemma}[Rank-five allowance-two normal form]
\label{lem:damp-qthree-minus-rank-five-allowance-two-normal-form}
In the setting of
Lemma~\ref{lem:damp-qthree-minus-core-carrier-defect-ledger}, suppose that
the repair rank is five and that \(\delta(H,5)=2\).  Then exactly one of
the following holds.
\begin{enumerate}
 \item \(h=1\), every \(q_S\) vanishes, and, up to permuting \(X\),
       \((p_x)_{x\in X}\) is \((1,1,0)\) or \((1,1,1)\).
 \item \(h=0\), every \(q_S\) vanishes, and, up to permuting \(X\),
       \((p_x)_{x\in X}\) belongs to
       \[
        \{(0,0,2),(0,1,1),(0,1,2),(0,2,2),
          (1,1,1),(1,1,2),(1,2,2),(2,2,2)\}.
       \]
 \item \(h=0\), there is a unique square \(S\) with \(q_S=1\), all
       other \(q_{S'}\) vanish, and
       \[
        p_x=1\quad(x\in S),
        \qquad p_z\in\{0,1,2\}\quad(X\setminus S=\{z\}).
       \]
       Moreover, the unique repair support counted by \(q_S\) is the
       unique such support counted by each \(p_x\), \(x\in S\).
\end{enumerate}
Thus allowance two has thirteen numerical defect shapes modulo core
symmetry.
\end{lemma}

\begin{proof}
Equation~\eqref{eq:damp-qthree-minus-carrier-to-singleton-defect} implies
\(q_S\le1\) for every \(S\): if \(q_S\ge2\), then for either \(x\in S\)
the coordinate ledger spends at least \(q_S+p_x\ge4\).  At most one square
can have \(q_S=1\).  Indeed, any two core squares meet in a coordinate
\(x\), and two nonzero carrier defects together with \(p_x\ge1\) would
spend at least three units at \(x\).

If \(h=1\), even one carrier defect would, at either coordinate of its
square, spend at least \(h+q_S+p_x=3\) units.  Hence all \(q_S\) vanish.
The zero-carrier endpoint charge gives \(\sum_xp_x\ge2\), while the total
and coordinate ledgers give \(\sum_xp_x\le3\) and \(p_x\le1\).  This is
the first case.

Suppose next that \(h=0\).  If every \(q_S\) vanishes, the endpoint charge
and the two ledgers give
\[
 2\le\sum_xp_x\le6,
 \qquad 0\le p_x\le2.
\]
The eight displayed triples are exactly the integer solutions modulo
permutation.  Finally, suppose that \(q_S=1\).  For \(x\in S\), one has
\(p_x\ge q_S=1\), while the coordinate ledger gives
\(q_S+p_x\le2\); hence \(p_x=1\).  The third coordinate satisfies only
\(0\le p_z\le2\).  Each support counted by \(q_S\) is counted by both
\(p_x\), and all three counts are one, proving the final assertion.
\end{proof}

\begin{lemma}[Rank-five allowance-two carrier orders]
\label{lem:damp-qthree-minus-rank-five-allowance-two-carrier-orders}
In either of the first two cases of
Lemma~\ref{lem:damp-qthree-minus-rank-five-allowance-two-normal-form},
label the three maximal core squares \(S_0,S_1,S_2\), and let
\(\pi_i\) be the repair-coordinate order of the carrier of \(S_i\).
If \(x_{ij}\) is the core coordinate common to \(S_i\) and \(S_j\), then
\begin{equation}
 d_{\mathrm K}(\pi_i,\pi_j)\le p_{x_{ij}}\le2-h.
 \label{eq:damp-rank-five-allowance-two-order-distance}
\end{equation}
After normalizing the designated carrier order to
\[
 \pi_0=(0,2,3,4,1),
\]
there are thirteen labelled possibilities for \((\pi_1,\pi_2)\) when
\(h=1\), and 106 when \(h=0\).  In the latter case, the sorted triple of
pairwise Kendall distances has one of the following five profiles, with
the indicated multiplicity:
\begin{equation}
\begin{array}{c|ccccc}
 \text{distance profile}&000&011&022&112&222\\ \hline
 \text{number of labelled pairs}&1&12&27&36&30.
\end{array}
\label{eq:damp-rank-five-allowance-two-distance-profiles}
\end{equation}
\end{lemma}

\begin{proof}
Since \(q_{S_i}=0\), the repair projection of the full \(S_i\)-cube
carrier is ample, contains the two antipodal endpoint words, and has
shadow \(\{\varnothing\}\cup\binom R1\).  It is therefore an antipodal
geodesic and has a well-defined order \(\pi_i\).  The inversion injection
from Lemma~\ref{lem:damp-qthree-minus-core-carrier-defect-ledger} gives
the first inequality in
\eqref{eq:damp-rank-five-allowance-two-order-distance}.  The coordinate
defect ledger gives \(h+p_{x_{ij}}\le2\), proving the second.

The endpoint normalization fixes the first and last directions of
\(\pi_0\), and a permutation of the remaining three directions gives the
displayed representative.  When \(h=1\), all three pairwise distances
are at most one.  As in
Lemma~\ref{lem:damp-qthree-minus-rank-five-allowance-one-carrier-orders},
the triangle-free permutation graph leaves the base pair and three
labelled pairs for each of the four adjacent transpositions, hence
thirteen pairs.

It remains to count the radius-two case.  Write \(s_1,\ldots,s_4\) for
the adjacent transpositions relative to \(\pi_0\).  The ball of Kendall
radius two consists of the identity, the four \(s_i\), and the following
nine elements of length two:
\[
 s_1s_3, s_1s_4, s_2s_4,
 \qquad
 s_is_{i+1}, s_{i+1}s_i\quad(1\le i\le3).
\]
Each of the three pairwise distances is at most two, and their sum is
even by permutation parity.  Hence the five profiles in
\eqref{eq:damp-rank-five-allowance-two-distance-profiles} are exhaustive.
The profiles \(011\) and \(022\) arise by making two of the three orders
equal, and therefore occur \(3\cdot4=12\) and \(3\cdot9=27\) times.
Among the length-two elements, the three products of commuting generators
have two neighbours of length one, whereas each of the other six has one;
there are consequently twelve incidences between the first and second
spheres.  The profile \(112\) comprises the twelve ordered pairs of
distinct length-one elements and the \(2\cdot12\) orientations of these
incidences, for a total of 36.  Finally, within the second sphere the three
commuting products have degree four at Kendall distance two, and the other
six have degree three.  Thus profile \(222\) occurs
\(3\cdot4+6\cdot3=30\) times.  Adding the identity pair gives the table,
whose entries sum to 106.
\end{proof}

\begin{lemma}[One-defect antipodal carrier]
\label{lem:damp-rank-five-one-defect-antipodal-carrier}
Let \(R\) be a five-element set, and let \(C\subseteq Q_R\) be ample with
\(0_R,1_R\in C\).  Suppose that, for some \(e,f\in R\),
\[
 \operatorname{Sh}(C)=
 \{\varnothing\}\cup\binom R1\cup\{\{e,f\}\}.
\]
Then \(C\) is the union of exactly two antipodal geodesics.  The two
geodesics agree outside one square, traverse that square in the two
opposite orders \(ef\) and \(fe\), and the directions \(e,f\) are
consecutive in either path order.  Equivalently, \(C\) is one antipodal
geodesic together with the missing fourth vertex of a square on two
consecutive path directions.  The square may occur in any of the four
positions along the path.
\end{lemma}

\begin{proof}
Ampleness gives \(|C|=|\operatorname{Sh}(C)|=7\).  The family \(C\) is
isometric and contains two vertices at distance five, so it contains an
antipodal geodesic \(P\) of order six.  Write \(z\) for the unique member
of \(C\setminus P\).

In an ample family, every shattered support is the support of a positioned
cube.  Hence \(C\) contains an \(\{e,f\}\)-square.  The geodesic \(P\)
realizes three of the four traces on \(\{e,f\}\), so this square contains
\(z\) and three vertices of \(P\).  A monotone antipodal path contains
three vertices of one positioned \(\{e,f\}\)-square only when its
\(e\)- and \(f\)-edges are consecutive.  The vertex \(z\) is then the
fourth vertex of that square.  Replacing the two-edge \(ef\)-arc of
\(P\) by the \(fe\)-arc through \(z\) gives a second antipodal geodesic.
Every antipodal geodesic in \(C\) must use one of these two arcs, so these
are the only two.  A five-edge path has four consecutive pairs, giving the
last assertion.
\end{proof}

\begin{corollary}[One-defect three-carrier synchronization]
\label{cor:damp-rank-five-one-defect-three-carrier-synchronization}
In case~3 of
Lemma~\ref{lem:damp-qthree-minus-rank-five-allowance-two-normal-form}, let
\(S\) be the unique square with \(q_S=1\), and let \(T\) be its unique
extra repair support.  Write \(\pi,\pi'\) for the two geodesic orders in
the carrier of \(S\).  Then \(T\) consists of two consecutive directions
of \(\pi\), and \(\pi'\) is obtained by transposing them.  Each of the
other two square carriers has order \(\pi\) or \(\pi'\).  If the two clean
carriers have common core coordinate \(z\) and \(p_z=0\), then their
orders are equal.  After normalizing the designated endpoint carrier
(choosing one of its two orders when it is defective), there are 32
labelled three-carrier patterns: sixteen when the designated carrier is
defective, and eight when either one of the other two carriers is
defective.
\end{corollary}

\begin{proof}
The first assertion is
Lemma~\ref{lem:damp-rank-five-one-defect-antipodal-carrier}.  Let \(S'\)
be a clean square and \(x\in S\cap S'\).  The allowance-two normal form
says that \(p_x=1\) and that its unique support is \(T\).  Compare the
geodesic carrier of \(S'\) first with the \(\pi\)-geodesic contained in
the carrier of \(S\), and then with the \(\pi'\)-geodesic.  The inversion
argument in
Lemma~\ref{lem:damp-qthree-minus-core-carrier-defect-ledger} applies to
either contained geodesic: every inversion belongs to the singleton-core
defect family at \(x\).  Thus both inversion sets are contained in
\(\{T\}\).  Since \(\pi,\pi'\) differ by the single adjacent
transposition \(T\), the only order having this property is \(\pi\) or
\(\pi'\).  This proves the assertion for both clean carriers.  Finally,
an inversion between the two clean orders would be counted by \(p_z\), so
\(p_z=0\) forces equality.  The square ear has four possible adjacent
positions.  If the designated carrier is defective, each clean carrier
independently chooses one of its two geodesics, giving
\(4\cdot2\cdot2=16\) patterns.  If one of the other carriers is defective,
the remaining clean carrier has two choices, giving \(4\cdot2=8\) patterns
for either defective-carrier position.
\end{proof}

\begin{corollary}[Rank-five punctured-core exclusion through order fifty-seven]
\label{cor:damp-qthree-minus-rank-five-through-fifty-seven}
Let \(H\) have minimum cardinality among all nonempty cornerless ample
families, suppose that \(|H|\le57\), and suppose that its maximum core is
\(Q_3\setminus\{v\}\) with a repair set of order five.  Then \(H\) has a
corner.
\end{corollary}

\begin{proof}
The threshold table in
Lemma~\ref{lem:damp-qthree-minus-core-carrier-defect-ledger} gives
\(\delta(H,5)\le0\).  Since the defect allowance is nonnegative, it
vanishes.  Equations
\eqref{eq:damp-qthree-minus-coordinate-defect-ledger} and
\eqref{eq:damp-qthree-minus-total-defect-ledger} then force
\[
 h=0,\qquad q_S=0\quad(S\in\binom X2),\qquad
 p_x=0\quad(x\in X).
\]
The zero-carrier endpoint charge,
Lemma~\ref{lem:damp-qthree-minus-zero-carrier-endpoint-charge}, instead
gives \(\sum_xp_x\ge2\).  This contradiction proves the claim.
\end{proof}

\begin{lemma}[Three-section endpoint packing on six coordinates]
\label{lem:damp-rank-six-three-section-endpoint-packing}
Let
\[
 P=(p_0,p_1,\ldots,p_6),
 \qquad p_t=\{0,\ldots,t-1\},
\]
be the standard antipodal geodesic in \(Q_{[6]}\), and let
\(F_1,F_2,F_3\) be families containing \(P\).  For each \(i\), let \(G_i\)
be the graph on \([6]=\{0,\ldots,5\}\) whose edges are the two-element
supports shattered by \(F_i\).  Suppose that the three edge sets are
pairwise disjoint and that, for every \(i\),
\[
 \{j_i\}\in F_i\quad\text{for some }1\le j_i\le5,
 \qquad
 [6]\setminus\{k_i\}\in F_i\quad\text{for some }0\le k_i\le4.
\]
Then
\begin{equation}
 \sum_{i=1}^3|E(G_i)|\ge11.
 \label{eq:damp-rank-six-three-section-eleven-pairs}
\end{equation}
Moreover, none of the graphs \(G_i\) contains a triangle.
\end{lemma}

\begin{proof}
We use two elementary consequences of the fixed path order.  First,
\begin{equation}
\begin{aligned}
 L_j&:=\{hj:0\le h<j\}\subseteq E(G_i)
       &&\text{if }\{j\}\in F_i,\\
 R_k&:=\{kh:k<h\le5\}\subseteq E(G_i)
       &&\text{if }[6]\setminus\{k\}\in F_i.
\end{aligned}
\label{eq:damp-rank-six-endpoint-stars}
\end{equation}
Indeed, the displayed word supplies the \(01\)-trace on each listed pair,
while \(P\) supplies the other three traces.

Second, the path order is an umbrella-free order for every \(G_i\): if
\(a<b<c\) and \(ac\in E(G_i)\), then
\begin{equation}
 ab\in E(G_i)\quad\text{or}\quad bc\in E(G_i).
 \label{eq:damp-rank-six-umbrella-free}
\end{equation}
To see this, choose \(w\in F_i\) with \(w_a=0,w_c=1\), which exists
because \(ac\) is shattered and \(P\) already realizes the other three
traces.  If \(w_b=0\), then \(w\) supplies the missing trace on \(bc\);
if \(w_b=1\), it supplies the missing trace on \(ab\).

Pairwise edge-disjointness first implies that the \(j_i\) are distinct
and that the \(k_i\) are distinct.  It also gives
\begin{equation}
 j_i\le k_h\qquad(i\ne h),
 \label{eq:damp-rank-six-cross-endpoint-order}
\end{equation}
because \(k_h<j_i\) would put the edge \(k_hj_i\) in both \(R_{k_h}\)
and \(L_{j_i}\).  Relabel the families so that
\(j_1<j_2<j_3\).  Then \(j_s\ge s\).  The two \(k\)-values belonging to
families other than \(F_3\) are distinct, at least \(j_3\), and at most
four.  Hence \(j_3=3\), and consequently
\[
 (j_1,j_2,j_3)=(1,2,3),
 \qquad \{k_1,k_2,k_3\}=\{2,3,4\},
 \qquad k_3=2.
\]
There are only two possible pairings of the remaining values.  The forced
edges from~\eqref{eq:damp-rank-six-endpoint-stars} are
\begin{equation}
\begin{array}{c|c|c|c}
 &(j_1,k_1)&(j_2,k_2)&(j_3,k_3)\\ \hline
\mathrm{I}&(1,3):(01,34,35)&(2,4):(02,12,45)&
 (3,2):(03,13,23,24,25)\\
\mathrm{II}&(1,4):(01,45)&(2,3):(02,12,34,35)&
 (3,2):(03,13,23,24,25).
\end{array}
\label{eq:damp-rank-six-two-endpoint-packings}
\end{equation}
In either row these are eleven distinct edges, proving
\eqref{eq:damp-rank-six-three-section-eleven-pairs}.  They comprise every
edge of \(K_6\) except
\begin{equation}
 M=\{04,05,14,15\}.
 \label{eq:damp-rank-six-four-unassigned-pairs}
\end{equation}
Thus an additional edge of any \(G_i\) must belong to \(M\).

It remains to exclude a triangle.  In row~I, the only possible triangles
are \(01h\), with \(h\in\{4,5\}\), in \(G_1\), and
\(a45\), with \(a\in\{0,1\}\), in \(G_2\); the forced edges of \(G_3\)
together with \(M\) contain no triangle.  A triangle \(01h\) contradicts
\eqref{eq:damp-rank-six-umbrella-free} on the long edge \(0h\) with
intermediate vertex \(2\): the edge \(02\) belongs to \(G_2\), while
\(2h\) belongs to \(G_3\), so edge-disjointness excludes both from
\(G_1\).  Similarly, a triangle \(a45\) contradicts the umbrella-free
condition on \(a5\) with intermediate vertex \(3\): the edge \(a3\)
belongs to \(G_3\), while \(35\) belongs to \(G_1\).

In row~II, only \(G_1\) can contain a triangle, and it must again be of
the form \(01h\) or \(a45\).  The same two umbrella arguments apply, now
using \(02\in E(G_2)\), \(2h\in E(G_3)\), \(a3\in E(G_3)\), and
\(35\in E(G_2)\).  This excludes every triangle.
\end{proof}

\begin{proposition}[Rank-six punctured-core exclusion]
\label{prop:damp-qthree-minus-rank-six-exclusion}
Let \(H\) have minimum cardinality among all nonempty cornerless ample
families, suppose that \(|H|<64\), and suppose that the maximum core from
Corollary~\ref{cor:damp-least-cornerless-maximum-core} is
\(A=Q_3\setminus\{v\}\) with a repair set \(R\) of order six.  Then \(H\)
has a corner.
\end{proposition}

\begin{proof}
Write \(X\) for the three core coordinates, put
\(\Sigma=\operatorname{Sh}(H)\), and translate the repair coordinates so
that the two prescribed core fibres occur over \(0_R\) and \(1_R\).  The
forced join in the proof of Proposition
\ref{prop:damp-vctwo-core-coordinate-wing-squeeze} contributes exactly
\(3(6+1)=21\) faces through every \(x\in X\).  On the other hand,
\eqref{eq:damp-vctwo-core-coordinate-degree-cap} and \(|H|\le63\) give
\[
 \deg_\Sigma(x)\le21.
\]
Consequently equality holds, and no support through a core coordinate
lies outside the forced join
\begin{equation}
 \binom X{\le2}*
 \left(\{\varnothing\}\cup\binom R1\right).
 \label{eq:damp-rank-six-no-extra-mixed-face}
\end{equation}
In particular, \(H\) has no concept whose core trace is \(v\), since one
such concept together with the two endpoint fibres would shatter \(X\).
Also, no shattered support can contain both a core coordinate and two
repair coordinates.

The three maximal squares of \(A\) have a common repair carrier path.
Indeed, for such a square support \(S\), Lemma
\ref{lem:damp-projected-cube-carrier-support-conservation} and
\eqref{eq:damp-rank-six-no-extra-mixed-face} show that the repair
projection of its full \(S\)-cube carrier is ample and has shadow
\[
 \{\varnothing\}\cup\binom R1.
\]
It contains \(0_R,1_R\), so it has seven vertices and is an antipodal
geodesic.  Two maximal core squares share a core edge.  If their carrier
geodesics ordered some repair pair in opposite ways, their union would
shatter that pair in both sections of the shared core edge.  The union of
the pair with the core-edge coordinate would then be shattered by \(H\),
contrary to~\eqref{eq:damp-rank-six-no-extra-mixed-face}.  Thus the three
carriers are one path \(P\).  Permute \(R\) so that
\begin{equation}
 P=(\varnothing,\{0\},\{0,1\},\ldots,R).
 \label{eq:damp-rank-six-common-repair-path}
\end{equation}
Every core word of \(A\) lies in a maximal core square, and hence every
repair section
\[
 F_a=\{g\in Q_R:(g,a)\in H\}
\]
contains \(P\).  Each \(F_a\) is ample because it is a conditioning of
\(H\).

If \(a\ne b\) and a repair support \(T\) of order at least two were
shattered by both \(F_a\) and \(F_b\), choose a core coordinate \(x\) on
which \(a\) and \(b\) differ.  The two sections would realize every
trace on \(T\) on the two opposite \(x\)-sides, so \(T\cup\{x\}\) would
be shattered.  Hence
\begin{equation}
 \bigl(\operatorname{Sh}(F_a)\setminus\operatorname{Sh}(P)\bigr)
 \cap
 \bigl(\operatorname{Sh}(F_b)\setminus\operatorname{Sh}(P)\bigr)
 =\varnothing
 \qquad(a\ne b).
 \label{eq:damp-rank-six-section-shadow-disjointness}
\end{equation}

Let \(a_1,a_2,a_3\) be the three corners of \(A\).  At the copy of
\(a_i\) over the first endpoint of \(P\), the two incident core directions
and the first repair direction span the product of its unique core square
with the first path edge.  There is no incident core direction leading
outside \(A\).  Cornerlessness therefore supplies another repair
neighbour \(\{j_i\}\), where \(1\le j_i\le5\).  The same argument at the
opposite endpoint supplies \(R\setminus\{k_i\}\), where
\(0\le k_i\le4\).  The pair-shadow graphs of the three corner sections
therefore satisfy Lemma
\ref{lem:damp-rank-six-three-section-endpoint-packing}.  In particular,
their union has at least eleven edges and none of them contains a
triangle.

Suppose first that \(H\) shatters a pure repair support of order at least
three.  By downward closure it shatters a repair triple \(T\), and
ampleness supplies a full \(T\)-cube in one fixed core section \(F_b\).
The trace \(b\) belongs to \(A\), by the first paragraph.  If \(b\) is a
core corner, its pair-shadow graph contains the triangle
\(\binom T2\), contrary to the endpoint-packing lemma.  If \(b\) is not a
core corner, the \(T\)-cube contributes at least four vertices outside
\(P\): an antipodal geodesic meets an affine three-cube in a contiguous
subpath of at most four vertices.  The eleven distinct pair supports in
the corner sections contribute at least eleven further section vertices,
because \(P\subseteq F_a\) and both families are ample.  Since no section
outside \(A\) is nonempty,
\[
 |H|
 =7|P|+\sum_{a\in A}|F_a\setminus P|
 \ge49+11+4=64,
\]
again a contradiction.

It remains that no pure repair triple is shattered.  Then every \(F_a\)
has VC dimension at most two.  Each corner section \(F_{a_i}\) properly
contains \(P\), so Lemma~\ref{lem:damp-relative-peeling-vc-two}, applied
to \(P\subsetneq F_{a_i}\), supplies a corner
\(w\in F_{a_i}\setminus P\).  The word \(w\) is not a prefix in the path
order.  Hence there are \(e<f\) with \(w_e=0,w_f=1\), and
\(P\cup\{w\}\) shatters \(\{e,f\}\).  If \(w\) occurred in any other core
section \(F_b\), that repair pair would be shattered in both sections,
contrary to~\eqref{eq:damp-rank-six-section-shadow-disjointness}.
There is also no section over the missing core trace \(v\).  Thus the
concept \((w,a_i)\) has no incident core edge.  Its unique maximal repair
cube in \(F_{a_i}\) is consequently its unique maximal cube in \(H\), so
\((w,a_i)\) is a corner of \(H\).  This final contradiction proves the
claim.
\end{proof}

\begin{proposition}[VC-two core exclusion through order fifty-four]
\label{prop:damp-vctwo-core-exclusion-through-fifty-four}
Let $H$ have minimum cardinality among all nonempty cornerless ample
families.  If
\[
 48\le |H|\le54,
\]
then $H$ has a corner.
\end{proposition}

\begin{proof}[Computer-assisted proof]
Corollary~\ref{cor:damp-vctwo-core-frontier-below-sixty-four} reduces $H$
to a VC-two maximum core.  For $A=Q_3\setminus\{v\}$, the two-ended
surcharge gives
\[
 |H|\ge7(r+1)+6,
\]
so $r\le5$.  For the pseudogeometric eleven-concept core
$A^{\mathrm{pg}}_4$, it gives
\[
 |H|\ge11(r+1)+6,
\]
so $r=3$.  Thus there are four core--repair profiles, not seven separate
cardinality cases:
\[
 (Q_3\setminus\{v\},r),\quad r\in\{3,4,5\},
 \qquad
 (A^{\mathrm{pg}}_4,3).
\]

For each profile we use the exact static facet-cube model of
Proposition~\ref{prop:damp-facet-cube-reconstruction}.  The variables
choose an antichain of facet supports, one anchored affine cube for each
chosen support, their union $K$, and one order selector.  The clauses
enforce the two prescribed antipodal core fibres, the equality
\[
 |K|=\left|\bigcup_{T\in\mathcal F}2^T\right|,
\]
and incidence multiplicity different from one at every vertex.  These
conditions are equivalent to ampleness and cornerlessness; the formulas
contain neither a recursive minor test nor a peeling search.  Since
$|H|<64$, the strict facet-excess inequality restricts every facet to rank
at most four.

The rank-three punctured-cube formula is already unsatisfiable without
fixing an endpoint-extension orbit.  In every other row, choose a square
core corner.  Proposition
\ref{prop:damp-vctwo-two-ended-facet-orbit-quotient} gives three
unoriented endpoint orbits when $r=3$ and four when $r\ge4$.  The archived
certificate keeps the two endpoint orientations separate and therefore
contains the redundant pair $1/2,2/1$.  One formula for each ordered orbit
fixes the two corresponding maximal cubes.  The resulting formulas are
\[
\begin{array}{c|c|c|r|r}
\text{core}&r&\text{orders}&\text{formulas}&
 (\text{variables},\text{clauses})\\ \hline
Q_3\setminus\{v\}&3&48\text{--}54&1&(3{,}440,32{,}685)\\
Q_3\setminus\{v\}&4&48\text{--}54&5&
 (8{,}024,112{,}468\text{--}112{,}470)\\
Q_3\setminus\{v\}&5&48\text{--}54&5&
 (18{,}524,380{,}182\text{--}380{,}184)\\
A^{\mathrm{pg}}_4&3&50\text{--}54&4&
 (8{,}024,112{,}493\text{--}112{,}495).
\end{array}
\]
Thus fifteen formulas cover every admissible profile and every order in
the interval simultaneously.  Only twelve are distinct modulo the full
endpoint-exchange symmetry; retaining the three duplicates makes the
certificate independent of this final quotient.

All fifteen formulas are unsatisfiable.  An independently compiled
\texttt{drat-trim} at source commit
\texttt{2e3b2dc0ecf938addbd779d42877b6ed69d9a985} verifies every proof;
the checker binary has SHA-256 digest
\[
\texttt{7f24ee5af22745ac2dbe3cf86cc4d61ef684694d3dce46875bb1f2bc6c8a421f}.
\]
For the largest verification core, the formula digest is
\[
\texttt{5b400716b5241d00f5dc4ddc636226bf42c9d3314c1a235c7dd944a1a3cffd05},
\]
and the proof digest is
\[
\texttt{d5c61c2e31887ce4a5d2126c562f20b29714cd2b98609b8448f1b959d42120af}.
\]
The checker retains $59{,}722$ of the $380{,}184$ original clauses and
$1{,}229{,}938$ of the $3{,}249{,}406$ proof lemmas, using
$669{,}133{,}373$ resolution steps and no RAT lemma.  The generators,
manifests, per-profile verification records, and aggregate certificate are
stored under the prefixes \texttt{damp\_vctwo\_core\_o48\_54} and
\texttt{damp\_vctwo\_below55}.  The aggregate mathematical digest is
\[
\texttt{c1903eaf24e3f9cbb979eca762f5ee57c7ea0e338b820001b89eb365832d0839}.
\]
Hence no listed profile exists, proving the claim.
\end{proof}

\begin{proposition}[Rank-four punctured-core exclusion through order fifty-seven]
\label{prop:damp-qthree-minus-rank-four-through-fifty-seven}
Let \(H\) have minimum cardinality among all nonempty cornerless ample
families, suppose that \(|H|\le57\), and suppose that its maximum core is
\(Q_3\setminus\{v\}\) with a repair set of order four.  Then \(H\) has a
corner.
\end{proposition}

\begin{proof}[Computer-assisted proof]
The endpoint-orbit quotient and the stored branch certificates leave only
the \(1/1\) endpoint type.  Translate and permute the repair coordinates so
that its two endpoint supports are the fixed representative
\(\{0\},\{1\}\).  For every core coordinate, the forced join contributes
fifteen shattered supports.  Since \(|H|\le57\), the coordinate-degree
bound gives
\[
 \deg_{\operatorname{Sh}(H)}(x)
 \le\left\lfloor\frac{57-20}{2}\right\rfloor=18.
\]
Thus the core-coordinate extra-face cap is three.  Equivalently, the
defect allowance of
Lemma~\ref{lem:damp-qthree-minus-core-carrier-defect-ledger} satisfies
\(\delta(H,4)\le3\).

We compile the exact static facet-cube model under this cap.  The formula
contains every facet support of rank at most four and every possible
anchor, enforces the two antipodal core fibres and the normalized endpoint
cubes, and imposes ampleness and cornerlessness through exact
facet-cube reconstruction.  The strict facet-excess inequality makes the
rank-four truncation exhaustive below order \(64\).  No cardinality
selector occurs.  The three coordinate inequalities and their weighted
sum
\[
 3h+2\sum_Sq_S+\sum_xp_x\le9
\]
are encoded explicitly.

There are two formulas, according as the missing core word is or is not
realized.  Each has \(7{,}116\) variables and \(98{,}787\) clauses.  Their
CNF SHA-256 digests are
\[
\begin{split}
 &\texttt{1396bc9ad064f8cdce82e77577e82986bfff2b8359a6fe0f16426b23353b1984},\\
 &\texttt{ccf75c091f423c0b8405e714eeb725fe6de757334fad3b5ca6cb9d612e96ca29}.
\end{split}
\]
Kissat returns \textsc{unsat} and emits proof traces with digests
\[
\begin{split}
 &\texttt{27836a554d7c75cdb654e46fb02ca8554b3a4df5140d3564fe6de6d80550fd40},\\
 &\texttt{597b2a593b8532ef850649eaf9cc9831cbbac46c7922b4a1370a3945f02d605e}.
\end{split}
\]
Independently compiled \texttt{drat-trim} verifies both.  For
\(h=0\), it retains \(17{,}540\) original clauses and \(389{,}356\)
lemmas and uses \(52{,}899{,}189\) resolution steps.  For \(h=1\), the
corresponding numbers are \(16{,}881\), \(216{,}279\), and
\(27{,}061{,}423\).  The aggregate verification record is
\path{damp_vctwo_q3minus_r4_deg3_below58_verification.json};
its mathematical digest is
\[
 \texttt{1b133d3afa01d2d0700eb3716b9c9ddb301c43e21eca0baf3597ed3e378e3221}.
\]
The two missing-core states exhaust the normalized \(1/1\) branch, proving
the claim.
\end{proof}

\begin{proposition}[Rank-four punctured-core exclusion through order fifty-nine]
\label{prop:damp-qthree-minus-rank-four-through-fifty-nine}
Let \(H\) have minimum cardinality among all nonempty cornerless ample
families, suppose that \(|H|\le59\), and suppose that its maximum core is
\(Q_3\setminus\{v\}\) with a repair set of order four.  Then \(H\) has a
corner.
\end{proposition}

\begin{proof}[Computer-assisted proof]
As in Proposition
\ref{prop:damp-qthree-minus-rank-four-through-fifty-seven}, the stored
endpoint-orbit certificates reduce the problem to the normalized \(1/1\)
endpoint type.  The forced join contributes fifteen shattered supports
through each core coordinate, while the coordinate-degree bound gives
\[
 \deg_{\operatorname{Sh}(H)}(x)
 \le\left\lfloor\frac{59-20}{2}\right\rfloor=19.
\]
Thus the core-coordinate extra-face cap is four, or equivalently
\(\delta(H,4)\le4\).

We use the same exact facet-cube reconstruction as in the preceding
proposition, now with the three coordinate inequalities at allowance four
and the weighted inequality
\[
 3h+2\sum_Sq_S+\sum_xp_x\le12.
\]
The formula has no cardinality selector.  The strict facet-excess
inequality again makes the rank-four facet truncation exhaustive below
order \(64\).  Splitting only by \(h=0\) and \(h=1\) gives two formulas,
each with \(7{,}116\) variables and \(98{,}787\) clauses.  Their CNF
SHA-256 digests are
\[
\begin{split}
 &\texttt{caf7a41381f2f4cd952f4bfc405af32e1d76465291610067922c467d7920ef80},\\
 &\texttt{0806415174e1ef646c0323b496d70ae75cd48988ddf19aa787571523592876d2}.
\end{split}
\]
Kissat returns \textsc{unsat} and emits proof traces with digests
\[
\begin{split}
 &\texttt{d43895e356786ccad41437149e30d25da9bdd712e2ddc96aacba24fd269dc5f5},\\
 &\texttt{ad897495967cc214d5ee55c597af094dc9aebb99c144412718ffa72aecf3f257}.
\end{split}
\]
Independently compiled \texttt{drat-trim} verifies both traces.  For
\(h=0\), it retains \(21{,}814\) original clauses and
\(2{,}347{,}194\) proof lemmas and uses \(423{,}344{,}899\) resolution
steps.  For \(h=1\), the corresponding numbers are \(18{,}740\),
\(759{,}210\), and \(110{,}604{,}115\).  The aggregate verification
record is
\path{damp_vctwo_q3minus_r4_deg4_below60_verification.json}; its
mathematical digest is
\[
 \texttt{5f68d1cadde08559f6761767bcd6048e324d217161b444d4e0d0cd47151803c3}.
\]
The two values of \(h\) exhaust the normalized branch, proving the claim.
\end{proof}

\begin{proposition}[Rank-five punctured-core exclusion through order fifty-nine]
\label{prop:damp-qthree-minus-rank-five-through-fifty-nine}
Let \(H\) have minimum cardinality among all nonempty cornerless ample
families, suppose that \(|H|\le59\), and suppose that its maximum core is
\(Q_3\setminus\{v\}\) with a repair set of order five.  Then \(H\) has a
corner.
\end{proposition}

\begin{proof}[Computer-assisted proof]
The endpoint-orbit quotient and the stored branch certificates leave only
the \(1/1\) endpoint type.  Its defect allowance satisfies
\(\delta(H,5)\le1\).  Lemma
\ref{lem:damp-qthree-minus-rank-five-allowance-one-normal-form} therefore
gives
\[
 h=0,\qquad q_S=0\quad(S\in\tbinom X2),\qquad
 \sum_xp_x\in\{2,3\}.
\]
After fixing the endpoint representative, Lemma
\ref{lem:damp-qthree-minus-rank-five-allowance-one-carrier-orders}
normalizes one carrier order and leaves thirteen labelled order pairs for
the other two carriers.

We add this order normal form to the exact static facet-cube reconstruction
as redundant clauses.  For each core square and repair word, an auxiliary
literal is equivalent to the conjunction saying that all four concepts of
the positioned core square occur.  Unit clauses fix the normalized first
carrier, and one of thirteen selectors fixes the other two geodesics.  The
resulting formulas split only by singleton-core mass two or three.  They
contain no cardinality selector, and each has \(17{,}757\) variables and
\(333{,}455\) clauses.  Their CNF SHA-256 digests are
\[
\begin{split}
 &\texttt{953b947aeb1f6c24d843948bdb35534b476639145effb7ba92c59d2da45c0f8d},\\
 &\texttt{25d69aab0165358079bbc25cc3ae0963975cc75335832657025e2d01feef7984}.
\end{split}
\]

Kissat returns \textsc{unsat} on both formulas and emits proof traces with
digests
\[
\begin{split}
 &\texttt{a8ebf575f9780163cbeb6015642d082623f12b043459749cbe0a0275a039d984},\\
 &\texttt{9c6c4557797b643ce04bc4c51b65c9c46783f53db63be4fdb87f9744b2e6a1c8}.
\end{split}
\]
Independently compiled \texttt{drat-trim} verifies both.  For singleton
mass two it retains \(44{,}743\) original clauses and \(507{,}751\) proof
lemmas and uses \(96{,}352{,}829\) resolution steps.  For singleton mass
three the corresponding numbers are \(46{,}505\), \(746{,}828\), and
\(143{,}769{,}041\).  The aggregate verification record is
\path{damp_vctwo_q3minus_r5_deg1_ordercuts_below60_verification.json}; its
mathematical digest is
\[
 \texttt{51548309160088e68fe4ed9450cfb60ac1d68b6115fd6737e8e9176df02e37eb}.
\]
The two mass values exhaust the allowance-one normal form, proving the
claim.
\end{proof}

\begin{corollary}[VC-two core exclusion through order fifty-seven]
\label{cor:damp-vctwo-core-exclusion-through-fifty-seven}
Let \(H\) have minimum cardinality among all nonempty cornerless ample
families.  If
\[
 48\le |H|\le57,
\]
then \(H\) has a corner.
\end{corollary}

\begin{proof}
Orders through \(54\) are excluded by
Proposition~\ref{prop:damp-vctwo-core-exclusion-through-fifty-four}.
For orders \(55,56,57\), the core-coordinate wing squeeze and the stored
endpoint certificates leave only the \(1/1\) punctured-cube core at repair
ranks four, five, and six.  Proposition
\ref{prop:damp-qthree-minus-rank-four-through-fifty-seven} excludes rank
four, Corollary
\ref{cor:damp-qthree-minus-rank-five-through-fifty-seven} excludes rank
five, and Proposition~\ref{prop:damp-qthree-minus-rank-six-exclusion}
excludes rank six.
\end{proof}

\begin{corollary}[VC-two core exclusion through order fifty-nine]
\label{cor:damp-vctwo-core-exclusion-through-fifty-nine}
Let \(H\) have minimum cardinality among all nonempty cornerless ample
families.  If
\[
 48\le |H|\le59,
\]
then \(H\) has a corner.
\end{corollary}

\begin{proof}
Orders through \(57\) are excluded by
Corollary~\ref{cor:damp-vctwo-core-exclusion-through-fifty-seven}.  At
orders \(58,59\), the core-coordinate wing squeeze and the stored endpoint
certificates leave only the \(1/1\) punctured-cube core at repair ranks
four, five, and six.  Propositions
\ref{prop:damp-qthree-minus-rank-four-through-fifty-nine} and
\ref{prop:damp-qthree-minus-rank-five-through-fifty-nine} exclude ranks
four and five, while Proposition
\ref{prop:damp-qthree-minus-rank-six-exclusion} excludes rank six.
\end{proof}

\begin{proposition}[Rank-four punctured-core exclusion through order sixty-one]
\label{prop:damp-qthree-minus-rank-four-through-sixty-one}
Let \(H\) have minimum cardinality among all nonempty cornerless ample
families, suppose that \(|H|\le61\), and suppose that its maximum core is
\(Q_3\setminus\{v\}\) with a repair set of order four.  Then \(H\) has a
corner.
\end{proposition}

\begin{proof}[Computer-assisted proof]
The endpoint-orbit certificates again reduce the problem to the normalized
\(1/1\) endpoint type.  Its forced join contributes fifteen shattered
supports through each core coordinate, while the trapped-wing degree bound
gives
\[
 \deg_{\operatorname{Sh}(H)}(x)
 \le\left\lfloor\frac{61-20}{2}\right\rfloor=20.
\]
Thus the core-coordinate extra-face cap is five, or equivalently
\(\delta(H,4)\le5\).

We use the exact facet-cube reconstruction with the three coordinate
inequalities and the weighted defect inequality
\[
 3h+2\sum_Sq_S+\sum_xp_x\le15.
\]
The strict facet-excess bound makes the rank-four facet truncation
exhaustive below order \(64\).  Splitting only by \(h=0\) and \(h=1\)
gives two cardinality-free formulas, each with \(7{,}116\) variables and
\(98{,}787\) clauses.  Their CNF SHA-256 digests are
\[
\begin{split}
 &\texttt{0ccd9f2f256a74751d22e7ef59562826221b7eddb8922b88fb90547dc2cd5acf},\\
 &\texttt{7b2f7513be88239d6a461abf694dffe91be0d112ac98fe1363efd8623c220ac9}.
\end{split}
\]
Kissat returns \textsc{unsat} and emits proof traces with digests
\[
\begin{split}
 &\texttt{00296b53291b81912d77ca423a8a4bafc8d137ecfbe5886b36e7df2c9d276434},\\
 &\texttt{5dff75ce553a2d321d70a7643064d71f48081cf9bb92693166e3c871f8b3ba90}.
\end{split}
\]
Independently compiled \texttt{drat-trim} verifies both traces.  For
\(h=0\), it retains \(24{,}992\) original clauses and
\(12{,}952{,}612\) proof lemmas and uses \(2{,}403{,}191{,}122\)
resolution steps.  For \(h=1\), the corresponding numbers are
\(22{,}906\), \(4{,}277{,}360\), and \(802{,}767{,}916\).  The aggregate
verification record is
\path{damp_vctwo_q3minus_r4_deg5_below62_verification.json}; its
mathematical digest is
\[
 \texttt{1c1dc2f041aaa1f3bb7e576fe5aad72b46384ba7d7dea1ea7e3bf9c308392cf7}.
\]
The two values of \(h\) exhaust the normalized branch.
\end{proof}

\begin{proposition}[Rank-five punctured-core exclusion through order sixty-one]
\label{prop:damp-qthree-minus-rank-five-through-sixty-one}
Let \(H\) have minimum cardinality among all nonempty cornerless ample
families, suppose that \(|H|\le61\), and suppose that its maximum core is
\(Q_3\setminus\{v\}\) with a repair set of order five.  Then \(H\) has a
corner.
\end{proposition}

\begin{proof}[Computer-assisted proof]
The endpoint-orbit certificates reduce the problem to the normalized
\(1/1\) endpoint type.  The forced join contributes eighteen shattered
supports through each core coordinate, while the trapped-wing degree bound
gives
\[
 \deg_{\operatorname{Sh}(H)}(x)
 \le\left\lfloor\frac{61-20}{2}\right\rfloor=20.
\]
Thus the core-coordinate extra-face cap is two, or equivalently
\(\delta(H,5)\le2\).  The cases with smaller allowance are contained in
the allowance-two formula, so it suffices to use
Lemma~\ref{lem:damp-qthree-minus-rank-five-allowance-two-normal-form}.

We compile the same exact static facet-cube reconstruction as above,
with maximum facet rank four and no cardinality selector, and add the
carrier-order normal forms from
Lemma~\ref{lem:damp-qthree-minus-rank-five-allowance-two-carrier-orders}
and Corollary
\ref{cor:damp-rank-five-one-defect-three-carrier-synchronization}.
The structural split is as follows.
\[
\begin{array}{c|l|c}
 \text{state}&\text{further subdivision}&\text{formulas}\\ \hline
 h=0,\ \sum_Sq_S=0
   &\begin{array}{l}
     002,011,012,022,111,112;\\
     122\text{ at distances }000,011,022,112;\\
     222\text{ at distances }000,011,022,112,222
    \end{array}&15\\
 h=1,\ \sum_Sq_S=0&\text{all compatible clean-carrier profiles}&1\\
 h=0,\ \sum_Sq_S=1&\text{all synchronized square-ear profiles}&1
\end{array}
\]
Here a three-digit word in the first row is the sorted singleton-defect
profile, and a word following ``at distances'' is the sorted triple of
pairwise Kendall distances.  The omitted combination \(122/222\) is
empty: three pairwise distances equal to two force every corresponding
singleton defect to be two.

An independent coverage audit enumerates the finite quotients rather than
trusting the formula split.  In the clean \(h=0\) state, the 106 normalized
order pairs and 23 labelled singleton profiles give exactly 422 compatible
order--profile pairs, all covered by the fifteen cells.  In the clean
\(h=1\) state, thirteen order pairs and four labelled singleton profiles
give 28 compatible pairs.  In the remaining state, the four positions of
the square ear give 32 distinct labelled three-carrier patterns: sixteen
with the designated carrier defective and eight for each of the other two
defective-carrier positions.  The audit finds no uncovered case.

The seventeen formulas have between \(17{,}231\) and \(18{,}826\)
variables and between \(330{,}074\) and \(344{,}517\) clauses.  Kissat
returns \textsc{unsat} on every formula.  Independently compiled
\texttt{drat-trim} verifies all seventeen proof traces, using between
\(10{,}017{,}433\) and \(710{,}424{,}249\) resolution steps per trace and
\(4{,}427{,}395{,}811\) steps in total.  The aggregate verification
record is
\path{damp_vctwo_q3minus_r5_deg2_below62_verification.json}; its
mathematical digest is
\[
 \texttt{b821dd601f63ac7e6cacc063be8257a9ed6e8bb0c114364e9e96bba0055dd183}.
\]
Thus every allowance-two normal form is impossible, proving the claim.
\end{proof}

\begin{corollary}[VC-two core exclusion through order sixty-one]
\label{cor:damp-vctwo-core-exclusion-through-sixty-one}
Let \(H\) have minimum cardinality among all nonempty cornerless ample
families.  If
\[
 48\le |H|\le61,
\]
then \(H\) has a corner.
\end{corollary}

\begin{proof}
Orders through \(59\) are excluded by
Corollary~\ref{cor:damp-vctwo-core-exclusion-through-fifty-nine}.  At
orders \(60,61\), the core-coordinate wing squeeze and the stored endpoint
certificates leave only the \(1/1\) punctured-cube core at repair ranks
four, five, and six.  Propositions
\ref{prop:damp-qthree-minus-rank-four-through-sixty-one} and
\ref{prop:damp-qthree-minus-rank-five-through-sixty-one} exclude ranks
four and five, while Proposition
\ref{prop:damp-qthree-minus-rank-six-exclusion} excludes rank six.
\end{proof}

\begin{proposition}[Rank-four punctured-core exclusion through order sixty-three]
\label{prop:damp-qthree-minus-rank-four-through-sixty-three}
Let \(H\) have minimum cardinality among all nonempty cornerless ample
families, suppose that \(|H|\le63\), and suppose that its maximum core is
\(Q_3\setminus\{v\}\) with a repair set of order four.  Then \(H\) has a
corner.
\end{proposition}

\begin{proof}[Computer-assisted proof]
The endpoint-orbit certificates and the core-coordinate wing squeeze leave
the normalized \(1/1\) endpoint type with extra-face allowance six.  The
carrier audit reconstructs all antipodal ample carriers and retains \(50\)
carrier-defect states, \(23{,}511\) carrier patterns, and \(349{,}503\)
compatible carrier--singleton-profile pairs.  Quotienting by exchange of
the two nondesignated carriers and grouping by total defect gives seventeen
certificate cells.  Carrier compatibility itself removes six nominal
cells: totals \(7,8,9\) do not occur when \(h=0\), and totals \(5,6,7\)
do not occur when \(h=1\).

The seventeen formulas have no cardinality selector.  They have between
\(7{,}564\) and \(63{,}744\) variables and between \(100{,}451\) and
\(265{,}984\) clauses.  Kissat returns \textsc{unsat} on every formula,
and independently compiled \texttt{drat-trim} verifies every trace.  The
aggregate checker independently reconstructs the carrier quotient, finds
no uncovered or multiply covered state, and records
\(1{,}728{,}418{,}557\) resolution steps.  Its verification record is
\path{damp_vctwo_q3minus_r4_deg6_below64_verification.json}, with
mathematical digest
\[
 \texttt{03308601e7ce9fa9f6ab858f2ded838b90cdd8ac2c42ebba271fdeb60aed6e4a}.
\]
Thus the allowance-six normal form is impossible.
\end{proof}

\begin{proposition}[Rank-five punctured-core exclusion through order sixty-three]
\label{prop:damp-qthree-minus-rank-five-through-sixty-three}
Let \(H\) have minimum cardinality among all nonempty cornerless ample
families, suppose that \(|H|\le63\), and suppose that its maximum core is
\(Q_3\setminus\{v\}\) with a repair set of order five.  Then \(H\) has a
corner.
\end{proposition}

\begin{proof}[Computer-assisted proof]
The endpoint-orbit certificates reduce again to the normalized \(1/1\)
type, now with extra-face allowance three.  The carrier-order normal form
is split by the missing-core indicator, the three carrier defects, the
sorted singleton-defect profile, and, only for the five resistant profiles,
the sorted triple of pairwise Kendall distances.  This gives sixty-four
certificate cells and no cardinality selector.

An independent reconstruction finds \(1{,}191\) carrier patterns and
\(5{,}810\) compatible carrier--profile pairs.  Every pair lies in exactly
one of the sixty-four cells.  All sixty-four individual formulas are
\textsc{unsat}, and independently compiled \texttt{drat-trim} verifies
all sixty-four traces.  For the final aggregate, one additional unprofiled
\(h=1,q=000\) proof replaces the eight cells in that row.  The resulting
coverage certificate therefore uses fifty-seven proof records while still
covering all sixty-four cells.  The formulas represented in the aggregate
have between \(17{,}310\) and \(24{,}140\) variables and between
\(330{,}104\) and \(375{,}007\) clauses.  The checker records
\(25{,}761{,}684{,}586\) resolution steps in total.

The aggregate verification record is
\path{damp_vctwo_q3minus_r5_deg3_below64_verification.json}.  Its
mathematical digest is
\[
 \texttt{60daf3f6e845b3b858a5013291797fd6016b2c85f3fe56560fe51e8e50caae8e}.
\]
The independently audited coverage map has no uncovered or multiply
covered pair, so every allowance-three normal form is impossible.
\end{proof}

\begin{corollary}[VC-two core exclusion through order sixty-three]
\label{cor:damp-vctwo-core-exclusion-through-sixty-three}
Let \(H\) have minimum cardinality among all nonempty cornerless ample
families.  If
\[
 48\le |H|\le63,
\]
then \(H\) has a corner.
\end{corollary}

\begin{proof}
Corollary~\ref{cor:damp-vctwo-core-exclusion-through-sixty-one} handles
orders through \(61\).  At orders \(62,63\), the endpoint quotient and the
core-coordinate wing squeeze leave only the punctured-cube core at repair
ranks four, five, and six.  Propositions
\ref{prop:damp-qthree-minus-rank-four-through-sixty-three} and
\ref{prop:damp-qthree-minus-rank-five-through-sixty-three} exclude the first
two ranks, while Proposition
\ref{prop:damp-qthree-minus-rank-six-exclusion} excludes rank six.
\end{proof}

\begin{proposition}[Punctured-cube core exclusion below order forty-eight]
\label{prop:damp-qthree-minus-core-below-forty-eight}
Let $X$ be a three-coordinate set, let
$A=Q_X\setminus\{v\}$, and let $R$ be disjoint from $X$ with
$r=|R|\in\{3,4\}$.  Suppose that $H\subseteq Q_{R\cup X}$ is ample and
that two antipodal repair fibres satisfy
\[
 H_f=H_{\overline f}=A.
\]
If $|H|\le47$, then $H$ has a corner.
\end{proposition}

\begin{proof}[Computer-assisted proof]
Translate and permute coordinates so that $f=0_R$ and
$A=Q_X\setminus\{1_X\}$.  Put $\Sigma=\operatorname{Sh}(H)$.  The two
equal antipodal fibres force the join subcomplex
\begin{equation}
 \Sigma_0=
 \binom{X}{\le2}*
 \bigl(\{\varnothing\}\cup\binom R1\bigr)
 \subseteq\Sigma,
 \qquad
 |\Sigma_0|=7(r+1).
 \label{eq:damp-qthree-minus-forced-join-shadow}
\end{equation}
Indeed, both fibres realize all traces on every member of
$\binom X{\le2}$, and their repair bits are opposite in every coordinate.
Thus the available face excess is at most $19$ for $r=3$ and at most $12$
for $r=4$.

We use the static model of
Proposition~\ref{prop:damp-facet-cube-reconstruction}.  Choose an antichain
$\mathcal F$ of facet supports, one anchored affine $T$-cube $C_T$ for
each $T\in\mathcal F$, and put $K=\bigcup_{T\in\mathcal F}C_T$.  The
model enforces
\begin{equation}
 |K|=\left|\bigcup_{T\in\mathcal F}2^T\right|,
 \qquad
 |\{T\in\mathcal F:x\in C_T\}|\ne1
 \quad(x\in Q_{R\cup X}),
 \label{eq:damp-qthree-minus-static-cnf-contract}
\end{equation}
together with $K_f=K_{\overline f}=A$ and $|K|\le47$.  By
Proposition~\ref{prop:damp-facet-cube-reconstruction}, these conditions are
equivalent to ampleness and cornerlessness; no recursive minor or peeling
test occurs in the formula.

The face budget makes the facet-rank truncation exhaustive.  For $r=3$ the
ambient dimension is six, and a rank-six facet already has $64$ vertices,
so it suffices to allow ranks at most five.  For $r=4$, a rank-six facet
has old--repair type $3+3$ or $2+4$ and contributes at least
\[
 \min\left\{
  2^6-\left(1+3+\binom32\right)(1+3),
  2^6-\left(1+2+\binom22\right)(1+4)
 \right\}=36
\]
faces outside $\Sigma_0$, exceeding the twelve-face budget.  Hence ranks
at most five suffice in both cases.  The resulting two global formulas
include every facet support and every anchor of the permitted ranks.  One
order selector ranges over all feasible orders at once.  Their sizes are
\[
\begin{array}{c|r|r}
r&\text{variables}&\text{clauses}\\ \hline
3&34{,}541&215{,}670\\
4&60{,}777&558{,}329.
\end{array}
\]

Lingeling returns \textsc{unsat} on both formulas and emits DRAT proofs.
An independently compiled \texttt{drat-trim} checker returns
\texttt{VERIFIED} on both.  Its exact backward-checking statistics are
\[
\begin{array}{c|r|r|r|r}
r&\text{core clauses}&\text{core lemmas}&\text{proof lemmas}
 &\text{resolution steps}\\ \hline
3&26{,}800&630{,}456&1{,}401{,}951&114{,}807{,}062\\
4&81{,}446&300{,}032&910{,}206&87{,}424{,}450.
\end{array}
\]
Neither core uses a RAT lemma.  The formula hashes are
\[
 \begin{gathered}
 \texttt{462ff85630445e4fa2d9677cf196e1e414ba1357765ab30ad30f89c51e8d00bb},\\
 \texttt{bf850214e54f198586d8bffbc52e2e3951b37410f06dbb0f8ea34927ae38841b},
 \end{gathered}
\]
and the corresponding proof hashes are
\[
 \begin{gathered}
 \texttt{9f359404780c4cbc79d0d7fc7570e3cb5481c2af76b9a76a9d1f93fd9e12a913},\\
 \texttt{3f5107c6a9d47e61add8cb49c8c857e042cf1d41d3010c6c119120ebcc405eec}.
 \end{gathered}
\]
The generator, formulas, proofs, and independent verification record are
stored under the common prefix
\texttt{damp\_q3minus\_core\_below48}.  Unsatisfiability proves the claim.
\end{proof}

\begin{corollary}[Path-only reduction below order forty-eight]
\label{cor:damp-path-only-reduction-below-forty-eight}
Let $H$ have minimum cardinality among all nonempty cornerless ample
families, and suppose that $|H|<48$.  Then the maximum core supplied by
Corollary~\ref{cor:damp-least-cornerless-maximum-core} has VC dimension one.
It is a path on $n+1$ vertices, where $2\le n\le9$; its repair set has
$r\ge3$, with $r\le12$ when $n=2$, and
\[
 (n+1)(r+1)\le40
 \qquad(n\ge3).
\]
\end{corollary}

\begin{proof}
Apply Proposition~\ref{prop:damp-bulk-reduction-below-forty-eight}.  Its
VC-dimension-two alternative has core $Q_3\setminus\{v\}$ and
$r\in\{3,4\}$, contrary to
Proposition~\ref{prop:damp-qthree-minus-core-below-forty-eight}.  The stated
path bounds are the remaining alternative of the bulk reduction.
\end{proof}

\begin{proposition}[Relative-wing first-moment bound]
\label{prop:damp-relative-wing-first-moment-bound}
Let $H$ have minimum cardinality among all nonempty cornerless ample
families, suppose that $|H|<48$, and keep the path-core notation of
Corollary~\ref{cor:damp-path-only-reduction-below-forty-eight}.  Put
\[
 m=|H|,\qquad d=n+r,\qquad
 q=\left\lfloor\frac{m-20}{2}\right\rfloor,
\]
and write $\Sigma=\operatorname{Sh}(H)$.  For $e\in R\mathbin{\dot\cup}X$,
let
\[
 \deg_\Sigma(e)=|\{S\in\Sigma:e\in S\}|.
\]
Then
\begin{align}
 \deg_\Sigma(e)&\le q
       &&(e\in R\mathbin{\dot\cup}X),
 \label{eq:damp-relative-wing-coordinate-degree}\\
 n+1&\le q,\qquad r+1\le q,
 \label{eq:damp-relative-wing-baseline-degrees}\\
 2m-d-2&\le dq.
 \label{eq:damp-relative-wing-first-moment}
\end{align}
Consequently $m\ge35$.  Thus one compatibility calculation for all
coordinate wings excludes every order from $28$ through $34$
simultaneously.
\end{proposition}

\begin{proof}
Fix an active coordinate $e$.  As in Proposition
\ref{prop:damp-two-wing-relative-corner-threshold}, write its two
sections as
\[
 A_b=\rho_e^bH\quad(b\in\{0,1\}),\qquad
 B=A_0\cap A_1,\qquad D_b=A_b\setminus B.
\]
The pairs $B\subsetneq A_b$ are trapped intervals in $\Damp$.
Proposition~\ref{prop:damp-nine-concept-passive-carrier-exclusion} gives
$|D_b|\ge10$ for both $b$.

The two-layer shadow-intersection identity says that the members of
$\Sigma$ containing $e$ are exactly
\[
 \{T\cup\{e\}:T\in\operatorname{Sh}(B)\}.
\]
Since $B$ is ample, their number is $|B|$.  Hence
\[
 m=|D_0|+|D_1|+2|B|\ge20+2\deg_\Sigma(e),
\]
which proves \eqref{eq:damp-relative-wing-coordinate-degree}.

Let $x_1,\ldots,x_n$ be the coordinates of the core path.  Its two
antipodal repair fibres force the following join skeleton in the shattered
complex:
\begin{equation}
 \Sigma_0=\{\varnothing\}
 \cup\{\{e\}:e\in R\mathbin{\dot\cup}X\}
 \cup\{\{t,x_i\}:t\in R,\ 1\le i\le n\}
 \subseteq\Sigma.
 \label{eq:damp-relative-wing-forced-join-skeleton}
\end{equation}
Indeed, the core path shatters every core singleton.  Every adjacent path
carrier contains the two repair antipodes, and therefore shatters every
repair singleton; the two copies of the corresponding core edge then
shatter every pair $\{t,x_i\}$.  In particular, a repair coordinate
belongs to at least $n+1$ members of $\Sigma$, and a core coordinate to
at least $r+1$.  Combining this with
\eqref{eq:damp-relative-wing-coordinate-degree} proves
\eqref{eq:damp-relative-wing-baseline-degrees}.

The family $\Sigma_0$ has $1+d+nr=(n+1)(r+1)$ members and first
support moment
\[
 \sum_{S\in\Sigma_0}|S|=d+2nr.
\]
Every member of $\Sigma\setminus\Sigma_0$ has order at least two,
because $\Sigma_0$ already contains the empty set and every singleton.
Consequently
\begin{align*}
 \sum_{e\in R\mathbin{\dot\cup}X}\deg_\Sigma(e)
 &=\sum_{S\in\Sigma}|S|\\
 &\ge d+2nr+2\bigl(m-(1+d+nr)\bigr)
  =2m-d-2.
\end{align*}
The coordinatewise upper bound
\eqref{eq:damp-relative-wing-coordinate-degree} gives the opposite
estimate $\sum_e\deg_\Sigma(e)\le dq$, proving
\eqref{eq:damp-relative-wing-first-moment}.

It remains only to read these uniform inequalities.  The already proved
exclusion through order $27$ gives $m\ge28$.

If \(28\le m\le29\), then \(q\le4\), so the baseline degree bounds give
\(n,r\le3\) and \(d\le6\).  On the other hand,
\eqref{eq:damp-relative-wing-first-moment} gives
\(2m-2\le5d\), and hence \(d\ge11\), a contradiction.

If \(30\le m\le31\), then \(q\le5\), so \(d\le8\), whereas the first
moment gives \(2m-2\le6d\) and hence \(d\ge10\).  If
\(32\le m\le33\), then \(q\le6\).  The first moment gives \(d\ge9\),
while \(n,r\le5\).  Since \(n\ge2\), \(r\ge3\), and \(n+r\ge9\), one has
\(n\ge4\), and the smallest possible path baseline together with the
seven-unit propagation surcharge is

\[
 (n+1)(r+1)+7\ge5\cdot6+7=37>m,
\]

a contradiction.

Finally suppose \(m=34\).  Then \(q=7\), and the first moment gives

\[
 66\le8(n+r),
\]

so \(n+r\ge9\).  If \(n=2\), the six-unit path surcharge gives
\(34\ge3(r+1)+6\), while \(r+1\le7\); hence \(n+r\le8\), a
contradiction.  If \(n\ge3\), then \(n,r\ge3\), and the seven-unit
surcharge gives \((n+1)(r+1)\le27\).  But \(n+r\ge9\) implies
\((n+1)(r+1)\ge4\cdot7=28\).  Therefore \(m\ge35\).
\end{proof}

\begin{proposition}[Rank-weighted relative-wing inequality]
\label{prop:damp-rank-weighted-relative-wing-inequality}
Let \(H\) have minimum cardinality among all nonempty cornerless ample
families, delete its constant coordinates, and put
\[
 E=\operatorname{idim}(H),\qquad d=|E|,\qquad
 \Sigma=\operatorname{Sh}(H),\qquad m=|H|.
\]
For \(j\in\{1,2\}\), write
\[
 M_j=\sum_{S\in\Sigma}|S|^j.
\]
For every \(e\in E\), write its sections as \(A_e^0,A_e^1\), put
\(B_e=A_e^0\cap A_e^1\), and, for \(b\in\{0,1\}\), choose
\(C_e^b\in\Damp\) inclusion-maximal subject to
\[
 B_e\subseteq C_e^b\subsetneq A_e^b.
\]
Set \(D_e^b=A_e^b\setminus C_e^b\), and define \(\tau(D_e^b)\) by
\eqref{eq:damp-terminal-block-curvature}.  Then
\begin{equation}
 (d+3)M_1-2M_2
 \ge dm+\sum_{e\in E}\sum_{b=0}^1\tau(D_e^b)
 \ge d(m+2).
 \label{eq:damp-rank-weighted-relative-wing-inequality}
\end{equation}
Equivalently, the first coordinatewise wing bound has a second-moment
refinement that remembers the ranks and terminal block curvature of every
coloured acquired-support family.
\end{proposition}

\begin{proof}
Fix now \(e\in E\).  Write the two sections of \(H\) as \(A_e^0,A_e^1\),
put \(B_e=A_e^0\cap A_e^1\), and define the two coloured acquired-support
families
\[
 \mathcal W_e^b=
 \operatorname{Sh}(A_e^b)\setminus\operatorname{Sh}(B_e)
 \qquad(b\in\{0,1\}).
\]
As in Proposition~\ref{prop:damp-two-wing-relative-corner-threshold}, both
pairs \(B_e\subsetneq A_e^b\) are trapped intervals in \(\Damp\).
Applying Proposition~\ref{prop:damp-terminal-block-curvature-tax} to all
\(2d\) pairs gives
\begin{equation}
 \sum_{e\in E}\sum_{b=0}^1
 \sum_{S\in\mathcal W_e^b}|S|
 \ge
 \sum_{e\in E}\sum_{b=0}^1
 \bigl(|A_e^b\setminus B_e|+\tau(D_e^b)\bigr).
 \label{eq:damp-summed-coloured-wing-rank-tax}
\end{equation}

The right-hand wing count is
\begin{equation}
 \sum_{e,b}|A_e^b\setminus B_e|=dm-2M_1.
 \label{eq:damp-coloured-wing-cardinality-identity}
\end{equation}
Indeed, for fixed \(e\) it equals \(m-2|B_e|\), and
\(|B_e|\) is the number of members of \(\Sigma\) containing \(e\).

For the rank-weighted left-hand side, the two-layer shadow identities say
that the union and intersection of the two section shadows are,
respectively, the deletion and the link of \(e\) in \(\Sigma\).  A face
\(S\in\Sigma\) of order \(s\) contributes \(s(d-s)\) over all deletions
and \(s(s-1)\) over all links.  Hence
\begin{equation}
 \sum_{e,b}\sum_{S\in\mathcal W_e^b}|S|
 =\sum_{S\in\Sigma}|S|(d+1-2|S|)
 =(d+1)M_1-2M_2.
 \label{eq:damp-coloured-wing-rank-identity}
\end{equation}
Substituting \eqref{eq:damp-coloured-wing-cardinality-identity} and
\eqref{eq:damp-coloured-wing-rank-identity} into
\eqref{eq:damp-summed-coloured-wing-rank-tax} gives
\[
 (d+1)M_1-2M_2
 \ge dm-2M_1+\sum_{e,b}\tau(D_e^b),
\]
which gives the first inequality in
\eqref{eq:damp-rank-weighted-relative-wing-inequality}.  The second follows
from \(\tau(D_e^b)\ge1\) for all \(2d\) terminal blocks.
\end{proof}

\begin{proposition}[The four order-thirty-two equality shadows]
\label{prop:damp-order-thirty-two-wing-moment-shadows}
Under the hypotheses of Proposition
\ref{prop:damp-relative-wing-first-moment-bound}, suppose that $m=32$.
Then $d=8$ and
\begin{equation}
 (n,r)\in\{(2,6),(3,5),(4,4),(5,3)\}.
 \label{eq:damp-order-thirty-two-path-splits}
\end{equation}
The shattered complex has VC dimension three and contains either one or
two three-element faces.  Let $G$ be its graph of two-element faces and
let $M=K_8\setminus G$.  Thus $M$ has no edge between the core and repair
blocks.  Up to permuting coordinates within the two blocks and, when
appropriate, interchanging the blocks, exactly the following four types
are possible:
\begin{equation}
\begin{array}{c|c|c|l}
|\Sigma\cap\binom E3|&(n,r)&M&\text{three-element faces}\\ \hline
1&(3,5)&P_3\mathbin{\dot\cup}P_5&
 \text{the centre of $P_3$ and two distance-two degree-two vertices of $P_5$}\\
2&(2,6)&K_2\mathbin{\dot\cup}C_6&
 \text{the two alternating triples of $C_6$}\\
2&(4,4)&C_4\mathbin{\dot\cup}P_4&
 \text{the two $C_4$ diagonals, paired with the two internal $P_4$ vertices}\\
2&(4,4)&C_4\mathbin{\dot\cup}K_{1,3}&
 \text{the two $C_4$ diagonals, both joined to the claw centre}.
\end{array}
\label{eq:damp-order-thirty-two-shadow-types}
\end{equation}
The first row also represents the block-swapped split $(5,3)$.
\end{proposition}

\begin{proof}
At $m=32$ one has $q=7$.  Equation
\eqref{eq:damp-relative-wing-first-moment} gives $62\le8d$, so $d\ge8$.
On the other hand,
\eqref{eq:damp-relative-wing-baseline-degrees} gives $n,r\le6$.
If $n=2$, the path surcharge gives
$32\ge3(r+1)+6$, and hence $r\le6$.  If $n\ge3$, it gives
\[
 (n+1)(r+1)\le25.
\]
For $n,r\ge3$, a sum $n+r\ge9$ would make the left-hand side at least
$4\cdot7=28$.  Thus $d\le8$ in every case.  This proves $d=8$ and the
four splits in \eqref{eq:damp-order-thirty-two-path-splits}.

Every ample family of VC dimension at most two is dismantlable, so
cornerlessness gives $\operatorname{VCdim}(H)\ge3$.  Besides the empty set
and the eight singletons, $\Sigma$ has twenty-three faces.  Put
\[
 h=\sum_{\substack{S\in\Sigma\\|S|\ge3}}(|S|-2).
\]
The coordinate-degree bound gives
\[
 54+h=\sum_{S\in\Sigma}|S|
      =\sum_e\deg_\Sigma(e)\le56.
\]
Thus $1\le h\le2$.  A four-element face would bring its four
three-element subsets and contribute at least six to $h$.  Hence there is
no face of order at least four, and $h$ is exactly the number
$t\in\{1,2\}$ of three-element faces.  It follows that $G$ has
$23-t$ edges, so $M$ has $5+t$ edges.

For a coordinate $v$, let $t_v$ be the number of three-element faces
containing it.  Since $\deg_G(v)=7-\deg_M(v)$,
\[
 \deg_\Sigma(v)=1+\deg_G(v)+t_v
               =8-\deg_M(v)+t_v\le7,
\]
and therefore
\begin{equation}
 \deg_M(v)\ge1+t_v.
 \label{eq:damp-order-thirty-two-missing-degree}
\end{equation}
If $t=1$, the two sides of
\eqref{eq:damp-order-thirty-two-missing-degree} have total sums $12$ and
$11$, respectively.  Thus one vertex has one extra missing-pair incidence
and every other vertex has equality.  The triangle is an independent set
in $M$.  Applying the handshaking lemma separately in the two blocks leaves
only block orders three and five: the smaller block contains one triangle
coordinate and has degrees $(2,1,1)$, while the larger contains two and
has degrees $(2,2,2,1,1)$, with the extra incidence on a nontriangle
vertex.  The independence condition makes the two graphs respectively
$P_3$ and $P_5$, in the positions stated in the first row.  For
completeness, the other three block splits fail before anchors are
considered: a triangle vertex in a two-vertex block would require degree
two, three triangle vertices in the six-vertex block would require six
incident missing edges although only five remain there, and the parity
and independence inequalities give the analogous contradiction in the
$4+4$ split.

If $t=2$, both sides of
\eqref{eq:damp-order-thirty-two-missing-degree} have total sum $14$.
Hence equality holds at every vertex.  In the $2+6$ split, the
two-vertex block cannot contain a triangle coordinate.  The two triangles
are disjoint in the other block: if they shared a vertex, that vertex
would require three missing neighbours but would have at most two
coordinates with which it shares neither triangle.  All six degrees in
that block are therefore two, and triangle independence forces the
alternating $C_6$ in the second row.

The $3+5$ split is impossible.  A triangle uses at most one coordinate of
the three-vertex block, because a used coordinate has missing degree two
and hence is adjacent in $M$ to both other coordinates.  Parity forces
exactly one of the two triangles to meet that block.  The other triangle
then lies wholly in the five-vertex block; its three independent vertices
require at least six incident missing edges, although that block has only
five.  The $5+3$ case is symmetric.

Finally consider the $4+4$ split.  Parity and the same independence
argument force total triangle incidences four and two in the two blocks,
up to exchange.  In the first block, every coordinate occurs once; all
missing degrees are two, giving $C_4$, and the two triangle pairs are its
diagonals.  In the second block, either two coordinates occur once, giving
degree sequence $(2,2,1,1)$ and hence $P_4$, or one coordinate occurs
twice, giving $(3,1,1,1)$ and hence $K_{1,3}$.  These are the last two
rows and complete the classification.
\end{proof}

\begin{proposition}[Exclusion of orders thirty-two and thirty-three]
\label{prop:damp-orders-thirty-two-thirty-three-excluded}
No minimum-cardinality nonempty cornerless ample family has order $32$ or
$33$.
\end{proposition}

\begin{proof}
Suppose first that $m=33$.  As in the preceding proof,
$\operatorname{VCdim}(H)\ge3$, so at least one member of
$\Sigma\setminus\Sigma_0$ has order at least three.  The first-moment
estimate therefore improves by one:
\[
 2m-d-1\le dq.
\]
Here $q=7$, so $65\le8d$ and $d\ge9$.  But
\eqref{eq:damp-relative-wing-baseline-degrees} gives $n,r\le6$.
If $n=2$, the path surcharge and $m=33$ give $r\le6$, whence $d\le8$.
If $n\ge3$, the same surcharge gives
$(n+1)(r+1)\le26$, whereas $n,r\ge3$ and $d\ge9$ force this product to
be at least $28$.  Both alternatives are impossible.

It remains to exclude $m=32$.  Proposition
\ref{prop:damp-order-thirty-two-wing-moment-shadows} leaves the four fixed
shattered complexes in
\eqref{eq:damp-order-thirty-two-shadow-types}.  For each one, apply the
anchored-facet reconstruction of Proposition
\ref{prop:damp-facet-cube-reconstruction}: choose one affine cube on every
facet support, require their union to have order $32$, and forbid a covered
vertex from lying in exactly one chosen cube.  These conditions are
equivalent to ampleness and cornerlessness.

The four deterministic Boolean formulas have respectively
\[
\begin{array}{c|r|r|r}
\text{type}&\text{variables}&\text{clauses}&\text{core clauses}\\ \hline
P_3\mathbin{\dot\cup}P_5&10{,}061&139{,}815&49{,}006\\
K_2\mathbin{\dot\cup}C_6&9{,}824&132{,}040&23{,}718\\
C_4\mathbin{\dot\cup}P_4&9{,}782&128{,}655&21{,}243\\
C_4\mathbin{\dot\cup}K_{1,3}&9{,}095&112{,}632&16{,}023.
\end{array}
\]
CaDiCaL~2.1.2 returns \textsc{unsat} on all four formulas.  For each
formula, \texttt{drat-trim} extracts a subformula and independently
verifies its DRAT proof.  The original formula hashes are
\[
\begin{gathered}
\texttt{de88895ed588caeb03d46423ea2018c4a61de90eb63a382f260231e032c39fb1},\\
\texttt{3bbdbdd8be2db7cd7d190c218621b0ace4cee64f4728ce4654ab5492e6a557e9},\\
\texttt{59bb2568d027448bd3f3db3d4eff68c1f439a9ea796f9e3fd8df66d34080497e},\\
\texttt{ab9f9d2438c4733e05ee3943f4e5cc714139ddffdea85349718ffe7ec4bcb863},
\end{gathered}
\]
and the corresponding core-proof hashes are
\[
\begin{gathered}
\texttt{58369ac1d61e09a4bf705186c7ec246d5dd5782d490fa30a93bd20ce5fd8a1d2},\\
\texttt{4f007c122d1980b8389248ee5a4301890d96b162822a3085d3cf338859fe60b0},\\
\texttt{a33913af549b646c44aeb0241ea86779bb4a33cdc476580166e6d817f9783bed},\\
\texttt{603e5c9d789a90c925b1e89a877dcec08c6193bf700895e067ebe34a59862bb7}.
\end{gathered}
\]
The compiler, formulas, extracted cores, proofs, and verifier are stored
under the common prefix
\path{damp_order32_wing_moment}.  The manifest and verification digests
are, respectively,
\[
\begin{gathered}
\texttt{8ee1871622a5af2e8c83a41bd99ac344f2e191acbcc0cc68b0b5abbe45d5505a},\\
\texttt{6ac2a3e3f1ac9f7041c8f177ef19cd21c1e221d76736deb466c4a3a7513fa08d}.
\end{gathered}
\]
Thus none of the four exact equality shadows has a cornerless anchored
realization, excluding order $32$.
\end{proof}

The tip-colouring identity is stronger than the repair-capacity bound in
one essential respect: it records the exact fibres in which projected
corners occur.  The coarse second-moment inequality by itself is too loose
to force a new dimension bound; even Hall's example has substantial slack
in that proxy.  The active problem is therefore to constrain the compatible
system of missing traces $q_A$, or equivalently of tip maps
\(\tau_S:E\setminus S\to C_S\).  Equation
\eqref{eq:damp-cornerless-tip-collision-bound} excludes dimensions five
and six in one argument, and
Corollary~\ref{cor:damp-seven-coordinate-maximum-corner} settles dimension
seven.  The compatibility problem remains in dimensions eight through
eleven.

\subsection{A concrete two-extension obstruction}

The symbolic two-extension theorem also produces examples that are not
cornerless.  For Hall's family, take
\[
 a=183,\qquad b=315.
\]
The two one-concept extensions are corner-peelable, and their new shattered
supports are
\[
 P_a=\{2,3,6,7\},\qquad P_b=\{2,3,8,9\}.
\]
They are distinct, so Corollary
\ref{cor:damp-cornerless-seed-repair-graph} gives an ample forbidden
pc-minor
\[
 X=((C_H\cup\{315\})\times\{0\})
   \cup((C_H\cup\{183\})\times\{1\})
 \subseteq Q_{13}.
\]
It has order \(600\) and exactly the two corners \((315,0)\) and
\((183,1)\).  Deleting either corner leaves a nonpeelable peripheral
expansion containing \(C_H\) as a proper pc-minor, which explains
concretely why direct corner descent fails.

\begin{remark}[Greedy shelling is not the missing invariant]
\label{rem:damp-two-repair-greedy-shelling-failure}
The same example rules out a tempting strengthening of the Yang route.
Its two restrictions in the new coordinate are the corner-peelable
one-concept repairs
\[
 C_H\cup\{315\},\qquad C_H\cup\{183\}.
\]
In the first restriction, \(315\) is a corner whose deletion leaves the
nonpeelable family \(C_H\); the second restriction has the symmetric bad
corner \(183\).  Thus even a proper pc-minor of a forbidden family need
not have the property that every available first peeling step is safe.
Equivalently, its shellable Yang ball need not be greedily shellable in
the sense required by this deletion process.

A worker audit tested every initial corner of all \(39\) elementary
pc-minors of the order-\(600\) family.  Among \(432\) corner--minor pairs,
exactly the two just described are unsafe.  The computation is only a
regression check; the preceding symbolic argument is the reason for the
failure.  The artifact is
\path{damp_two_repair_minor_corner_safety.json}, with mathematical digest
\[
 \texttt{fff90aec268a0b9c85619ae1470e724f2f2ead7965c479431aa9a7f3f421cf5b}.
\]
\end{remark}

\subsection{The minimum-order problem}

The connection with partial shellings also suggests a construction route
that does not enumerate the intervening orders.  Bolognini and Sentinelli
call a proper pure subcomplex of a simplex skeleton \emph{austere} when every
missing facet contains at least two minimal nonfaces
\cite{BologniniSentinelli26}.  Since a possible next shelling facet contains
a unique minimal nonface, austerity is a static certificate that a partial
shelling is stuck.  For a partial shelling of a cross-polytope, the
correspondence of Chalopin, Chepoi, Moran, and Warmuth identifies the missing
facets with an ample cube family; the partial shelling is stuck exactly when
this complementary family has no corner
\cite[Corollary~4.4]{ChalopiChepoiMoran22}.  Thus transversal austerity is
the shelling form of cornerlessness.

This observation does not make the echo construction from
\cite{BologniniSentinelli26} a Yang complex: an echo facet may use both
vertices of one pair and omit another pair, whereas a cross-polytope facet
selects exactly one vertex from every pair.  Their explicit quiet
two-complex nevertheless suggests a natural paired template.  It has eight
vertices, complete one-skeleton, twenty-one triangles, and hence fifty-eight
faces including the empty face.  Let $\Gamma$ denote this complex.  Since
its one-skeleton is complete, every missing triple is a minimal nonface;
quietness ensures that every four-set contains a missing triple, so there
are no larger minimal nonfaces.  Thus the minimal nonfaces are exactly the
thirty-five missing triples.

The signed-trace constraints in fact make this displayed template rigid.
This also serves as a worked application of
Proposition~\ref{prop:damp-coordinate-propagation-rigidity}.
In the facet list below, copied from
\cite{BologniniSentinelli26}, we abbreviate a triple such as
$\{1,2,4\}$ to $124$:
\[
\begin{gathered}
148,348,234,347,368,356,168,357,236,147,167,156,267,257,\\
245,125,458,128,278,123,456.
\end{gathered}
\]
For a coordinate $e$, the following row lists the edges of a spanning tree
of $P_e(\Gamma)$.
\begin{center}
\small
\begin{tabular}{c@{\quad}p{0.82\linewidth}}
\toprule
$e$ & spanning-tree edges\\
\midrule
1 & $124$--$134$, $124$--$145$, $134$--$137$, $134$--$138$,
    $145$--$146$, $145$--$158$, $137$--$135$, $137$--$157$,
    $138$--$136$, $157$--$127$, $136$--$126$, $127$--$178$\\
2 & $124$--$127$, $124$--$247$, $124$--$248$, $127$--$126$,
    $247$--$237$, $248$--$238$, $248$--$258$, $126$--$256$,
    $126$--$268$, $237$--$235$, $256$--$246$\\
3 & $134$--$137$, $134$--$138$, $137$--$136$, $137$--$367$,
    $138$--$238$, $136$--$135$, $367$--$237$, $238$--$378$,
    $135$--$235$, $235$--$345$, $345$--$346$, $345$--$358$\\
4 & $124$--$134$, $124$--$145$, $124$--$248$, $145$--$146$,
    $248$--$247$, $248$--$478$, $146$--$467$, $146$--$468$,
    $247$--$246$, $247$--$457$, $468$--$346$, $457$--$345$\\
5 & $135$--$235$, $235$--$256$, $235$--$345$, $256$--$567$,
    $345$--$358$, $345$--$457$, $567$--$157$, $358$--$568$,
    $457$--$145$, $568$--$158$, $158$--$258$, $258$--$578$\\
6 & $126$--$136$, $126$--$256$, $126$--$268$, $256$--$246$,
    $256$--$567$, $268$--$678$, $246$--$346$, $567$--$367$,
    $346$--$467$, $346$--$468$, $467$--$146$, $468$--$568$\\
7 & $127$--$137$, $127$--$157$, $127$--$178$, $127$--$237$,
    $157$--$567$, $178$--$478$, $178$--$678$, $237$--$247$,
    $237$--$367$, $567$--$457$, $567$--$467$, $478$--$378$,
    $478$--$578$\\
8 & $138$--$238$, $238$--$248$, $238$--$268$, $248$--$258$,
    $268$--$678$, $258$--$158$, $258$--$578$, $678$--$178$,
    $158$--$568$, $578$--$358$, $578$--$378$, $178$--$478$,
    $568$--$468$\\
\bottomrule
\end{tabular}
\end{center}
Each listed pair satisfies one of the two conditions in
Definition~\ref{def:damp-coordinate-propagation-graph}: this can be checked
by inspecting the displayed facet list, since its complement is the set of
missing triples.  The table therefore proves that all eight propagation
graphs are connected.
Proposition~\ref{prop:damp-coordinate-propagation-rigidity} now shows that
every ample $X\subseteq Q_8$ with
$\operatorname{Sh}(X)=\Gamma$ is a cube translate of the canonical face
family of $\Gamma$.  The canonical family has corners: the characteristic
vector of any triangular facet lies only in the corresponding maximal
cube.  Consequently no such $X$ is cornerless.  Thus the displayed quiet
complex does not improve the upper bound.  In fact the same method, followed
by a small residual audit, excludes the entire quiet family from
\cite{BologniniSentinelli26}.

\begin{proposition}[Exclusion of the quiet fifty-eight-face templates]
\label{prop:damp-quiet-58-exclusion}
None of the fifty-nine quiet two-complexes on eight vertices with complete
one-skeleton and twenty-one triangles enumerated in
\cite{BologniniSentinelli26} is the shattered-set complex of a cornerless
ample family.
\end{proposition}

\begin{proof}[Finite audit after coordinate propagation]
We first reproduce the source census.  Fix a Steiner system $S(3,4,8)$.
Its fourteen blocks partition the fifty-six triples.  If $x,y,z$ blocks
contain respectively two, one, and no chosen triangles, then quietness and
the total of twenty-one triangles give
\[
 2x+y=21,\qquad x+y+z=14,
 \]
so $x-z=7$ and at least seven blocks contain two triangles.  The affine
group of the Steiner system has fourteen orbits on seven-subsets of its
blocks.  For one representative of each orbit, we choose two triangles in
the seven marked blocks, impose the total of twenty-one triangles, complete
one-skeleton, and quietness on every four-set, and quotient the surviving
choices by all permutations of the eight vertices.  Exhausting the fourteen
marked problems gives fifty-nine isomorphism classes, agreeing with
\cite{BologniniSentinelli26}.

For forty-seven of these complexes every nonempty graph $P_e(\Sigma)$ is
connected, so Proposition~\ref{prop:damp-coordinate-propagation-rigidity}
excludes them.  Consider one of the remaining twelve.  If
$c_e$ is the number of components of $P_e(\Sigma)$, the
four-coordinate lemma forces the $e$-bit of the omitted traces to be
constant on each component.  Translation in coordinate $e$ complements all
these component bits, so we normalize one of them to zero.  Every admissible
signing therefore occurs among exactly
\[
 2^r,\qquad r=\sum_{e=1}^8(c_e-1),
\]
component signings.  The residual audit is summarized below.
\begin{center}
\begin{tabular}{c@{\qquad}r@{\qquad}r@{\qquad}r@{\qquad}r}
\toprule
$r$ & complexes & signings & ample signings & least corners\\
\midrule
1  & 1 & 2    & 2   & 21\\
2  & 3 & 12   & 10  & 18\\
3  & 4 & 32   & 28  & 15\\
4  & 3 & 48   & 36  & 15\\
11 & 1 & 2048 & 176 & 13\\
\midrule
total & 12 & 2142 & 252 & 13\\
\bottomrule
\end{tabular}
\end{center}
For completeness, the audit evaluates a component signing
$\boldsymbol\sigma$ by forming the common satisfying family
$H(\Sigma,\boldsymbol\sigma)$.  By
Theorem~\ref{thm:damp-fixed-shadow-signed-clutter}, it is ample with
shattered-set complex $\Sigma$ exactly when it has fifty-eight members.
For each such signing we count the members that uniquely realize a trace on
a triangular facet; by
Lemma~\ref{lem:damp-corner-unique-facet-trace}, these are exactly its
corners.  The table exhausts all $2142$ normalized component signings.  All
252 normalized ample signings produce families with at least thirteen
corners, which excludes the remaining twelve complexes.
\end{proof}

The exact basis theorem is a recognition theorem and does not determine the
least order of an ample member of the basis.  Hall's family gives order
\(299\).  We next obtain a conceptual lower bound from the fibers over a
maximal shattered set.  The argument is uniform in the ambient dimension;
the earlier finite exclusions at particular orders remain independent
campaign regressions and are not used here.

The minimum-order question admits a uniform conceptual treatment.
The point is to count how many outside coordinates can fail section
containment in the fibers over one maximal shattered set.  This produces a
bound in arbitrary dimension, rather than a separate argument at each order.

\begin{lemma}[Unmatched coordinates in a rooted fiber]
\label{lem:damp-rooted-fiber-defect}
Let \(C\subseteq Q_U\) be an isometric family containing \(0\), with

\[
 |C|\ge2.
\]

Define

\[
 D(C)=\{e\in U:\text{there is }z\in C\text{ with }
 z_e=1\text{ and }z^{\{e\}}\notin C\}.
\]

Then

\[
 |D(C)|\le |C|-2.
\]

If equality holds and \(s=|C|\), then there are distinct coordinates
\(e_1,\ldots,e_{s-1}\) such that

\[
 C=\{0,\mathbf 1_{\{e_1\}},\mathbf 1_{\{e_1,e_2\}},\ldots,
 \mathbf 1_{\{e_1,\ldots,e_{s-1}\}}\}
\]

and \(D(C)=\{e_1,\ldots,e_{s-2}\}\).
\end{lemma}

\begin{proof}
Root a spanning tree of \(C\) at \(0\) as follows.  For every nonzero
\(z\in C\), choose as its parent a neighbor whose support has cardinality

\[
 |\operatorname{supp}(z)|-1.
\]

Such a neighbor exists because \(C\) is isometric: a shortest path in \(C\)
from \(0\) to \(z\) has length \(|\operatorname{supp}(z)|\).  The chosen
parent edges form a tree,
and each root-to-vertex path in this tree is a hypercube geodesic.

Fix \(e\in D(C)\), and choose a witnessing vertex \(z\).  The tree path from
\(0\) to \(z\) crosses coordinate \(e\) exactly once.  Let \(v_e\) be the
endpoint of that edge farther from the root.  This edge is not the last edge
of the path, since otherwise the preceding vertex would be
\(z^{\{e\}}\in C\), contrary to the choice of \(z\).  Thus \(v_e\) is a
nonroot, nonleaf vertex of the rooted tree.  Distinct coordinates give
distinct vertices \(v_e\), because every nonroot vertex has a unique parent
edge.  A rooted tree with at least two vertices has a nonroot leaf, so it has
at most \(|C|-2\) nonroot, nonleaf vertices.  The injection
\(e\mapsto v_e\) proves the assertion.

Suppose equality holds.  Then the injection uses every nonroot, nonleaf
vertex, and the tree has exactly one nonroot leaf.  A finite rooted tree with
only one nonroot leaf is a path.  Its edge labels
\(e_1,\ldots,e_{s-1}\) are distinct because the path is a hypercube
geodesic.  The vertices of \(C\) are therefore exactly the successive
prefixes displayed in the statement.  Each \(e_i\), for \(i<s-1\), is
unmatched at the final vertex, whereas the only vertex with
\(e_{s-1}\)-value one has its mate at the preceding vertex.  This identifies
\(D(C)\).
\end{proof}

\begin{theorem}[Facet-excess inequality]
\label{thm:damp-facet-excess}
Let \(H\subseteq Q_F\) be ample and cornerless, let \(T\) be a facet of
\(\operatorname{Sh}(H)\), and put \(U=F\setminus T\).  If \(p_T(H)\) is the
number of nonperipheral coordinates in \(U\), then

\begin{equation}
 p_T(H)\le |H|-2^{|T|+1}.
 \label{eq:damp-facet-excess}
\end{equation}

In particular, if \(H\) is periphery-free, then

\begin{equation}
 |H|\ge 2^{|T|+1}+|F|-|T|.
 \label{eq:damp-periphery-free-facet-excess}
\end{equation}
\end{theorem}

\begin{proof}
Since \(H\) is ample, it strongly shatters \(T\).  Thus there is a word
\(y\in Q_U\) such that

\[
 \{(q,y):q\in Q_T\}\subseteq H.
\]

Translate the coordinates in \(U\) so that \(y=0\).  For \(q\in Q_T\), let

\[
 C_q=\{z\in Q_U:(q,z)\in H\}.
\]

Every \(C_q\) contains \(0\).  It is the projection of an intersection of
coordinate semicubes of \(H\), hence is an ample pc-minor and, in
particular, isometric.  Moreover,

\[
 |C_q|\ge2.
\]

Indeed, if \((q,0)\) were the unique member with \(T\)-trace \(q\), then the
facet \(T\) would be uniquely realized by that member, and
Lemma~\ref{lem:damp-corner-unique-facet-trace} would make it a corner.

Let \(e\in U\) be nonperipheral.  In the orientation fixed by the common
base word, the containment

\[
 \rho_e^1H\subseteq\rho_e^0H
\]

fails.  Hence some \(z\in C_q\) has \(z_e=1\) but
\(z^{\{e\}}\notin C_q\), so \(e\in D(C_q)\).  Therefore

\[
 p_T(H)
 \le \left|\bigcup_{q\in Q_T}D(C_q)\right|
 \le \sum_{q\in Q_T}|D(C_q)|
 \le \sum_{q\in Q_T}(|C_q|-2)
 =|H|-2^{|T|+1},
\]

where the third inequality is
Lemma~\ref{lem:damp-rooted-fiber-defect}.  If \(H\) is periphery-free, then
all \(|U|=|F|-|T|\) outside coordinates are counted by \(p_T(H)\), which
gives \eqref{eq:damp-periphery-free-facet-excess}.
\end{proof}

\begin{lemma}[Facets have at least two outside coordinates]
\label{lem:damp-facet-codimension-at-least-two}
Under the hypotheses of Theorem~\ref{thm:damp-facet-excess}, one has
\[
 |F\setminus T|\ge2.
\]
\end{lemma}

\begin{proof}
Put \(U=F\setminus T\), and use the rooted fibers \(C_q\subseteq Q_U\)
from the proof of Theorem~\ref{thm:damp-facet-excess}.  If \(U\) were
empty, the strongly shattered \(T\)-cube would be all of \(Q_F\), so
\(H=Q_F\); this contradicts cornerlessness.  If \(|U|=1\), then every
fiber contains \(0\) and has at least two vertices, and hence every fiber
equals \(Q_U\).  Thus \(H=Q_F\) again.  Equivalently, \(F\) is shattered,
contrary to the maximality of the proper face \(T\) in
\(\operatorname{Sh}(H)\).  Therefore \(|U|\ge2\).
\end{proof}

Equality in \eqref{eq:damp-periphery-free-facet-excess} is itself
incompatible with ampleness.

\begin{corollary}[Strict facet-excess inequality]
\label{cor:damp-strict-facet-excess}
Under the hypotheses of Theorem~\ref{thm:damp-facet-excess}, if \(H\) is
periphery-free, then

\[
 |H|\ge 2^{|T|+1}+|F|-|T|+1.
\]
\end{corollary}

\begin{proof}
Put \(U=F\setminus T\) and \(k=|U|\).  If \(k=0\), the weak inequality would
give \(|H|\ge2^{|T|+1}\), although \(H\subseteq Q_T\) has at most
\(2^{|T|}\) members.  Thus \(k\ge1\).  Suppose equality holds in
\eqref{eq:damp-periphery-free-facet-excess}.  With the notation from the
proof of Theorem~\ref{thm:damp-facet-excess},

\[
 k
 =\left|\bigcup_{q\in Q_T}D(C_q)\right|
 =\sum_{q\in Q_T}|D(C_q)|
 =\sum_{q\in Q_T}(|C_q|-2).
\]

Hence the sets \(D(C_q)\) are pairwise disjoint and partition \(U\), and
every fiber satisfies equality in
Lemma~\ref{lem:damp-rooted-fiber-defect}.  That lemma makes each \(C_q\) a
rooted geodesic path.  Write \(t(q)\) for the label of its last edge.
Then

\[
 \operatorname{supp}(C_q)=D(C_q)\mathbin{\dot\cup}\{t(q)\}.
\]

Because the defect sets partition \(U\), the terminal coordinate \(t(q)\)
belongs to \(D(C_{q'})\) for a unique trace \(q'\), and necessarily
\(q'\ne q\).  Choose \(q\) with \(|C_q|\ge3\), set \(a=t(q)\), and choose
\(b\in D(C_q)\).  Such a fiber exists: if every fiber had two vertices,
then all defect sets would be empty, contrary to \(k\ge1\).  The equality
description in Lemma~\ref{lem:damp-rooted-fiber-defect} then ensures that
\(D(C_q)\ne\varnothing\).

Since \(b\in D(C_q)\) and \(a=t(q)\), the path \(C_q\) contains the zero
trace, the singleton trace on \(b\), and the trace with both \(a\) and \(b\)
equal to one.  The fiber \(C_{q'}\) has \(a\in D(C_{q'})\), so its terminal coordinate is not
\(a\).  Moreover \(b\notin D(C_{q'})\), because the defect sets are
disjoint; if \(b\) occurs on this path, it is therefore its terminal
coordinate and occurs after \(a\).  Thus the path contains a vertex with
\(a\)-value one and \(b\)-value zero.  Together with the common base trace
\(00\), these
vertices realize all four traces on \(\{a,b\}\).  Thus \(\{a,b\}\) is
shattered by \(H\).

For each \(e\in U\), let

\[
 P_e=\{q\in Q_T:\text{some }z\in C_q\text{ has }z_e=1\}.
\]

The sets \(S\cup\{e\}\), with \(S\in\operatorname{Sh}(P_e)\), are shattered
by \(H\): the value-zero \(e\)-section contains the common \(T\)-cube, and
the value-one section projects onto \(P_e\).  Conversely this criterion is
exact.  Since

\[
 |P_e|=\#\{q:e\in\operatorname{supp}(C_q)\},
\]

the fiber identity above gives

\[
 \sum_{e\in U}|P_e|
 =\sum_{q\in Q_T}|\operatorname{supp}(C_q)|
 =\sum_{q\in Q_T}(|D(C_q)|+1)
 =k+2^{|T|}.
\]

The Sandwich inequality for the families \(P_e\subseteq Q_T\) now yields

\[
 \begin{split}
 |\operatorname{Sh}(H)|
 &\ge
 2^{|T|}
 +\sum_{e\in U}|\operatorname{Sh}(P_e)|
 +1\\
 &\ge
 2^{|T|}+\sum_{e\in U}|P_e|+1
 =2^{|T|+1}+k+1>|H|.
 \end{split}
\]

Here the final \(1\) is the shattered pair \(\{a,b\}\), and the counted
families are disjoint because their intersections with \(U\) have sizes
zero, one, and two.  This contradicts the Sandwich equality for \(H\).
\end{proof}

The preceding proof has a stable first layer.  Besides excluding equality,
it excludes one additional vertex uniformly; this is the point at which the
minimum-order argument ceases to be an order-by-order calculation.

\begin{lemma}[First-slack rooted fibers]
\label{lem:damp-first-slack-rooted-fibers}
Let \(C_1,\ldots,C_m\subseteq Q_U\) be ample families containing
\(0\), with \(|C_i|\ge2\).  Put
\[
 D_i=D(C_i),\qquad
 \delta_i=|C_i|-2-|D_i|,
\]
and suppose that
\begin{equation}
 \bigcup_{i=1}^mD_i=U,
 \qquad
 \sum_{i=1}^m(|C_i|-2)=|U|+1.
 \label{eq:damp-first-slack-budget}
\end{equation}
Let \(C=\bigcup_i C_i\), and for any family \(A\subseteq Q_U\) write
\[
 h(A)=|\{R\in\operatorname{Sh}(A):|R|\ge2\}|.
\]
Then
\begin{equation}
 h(C)>\sum_{i=1}^m h(C_i).
 \label{eq:damp-first-slack-new-face}
\end{equation}
\end{lemma}

\begin{proof}
Write
\[
 \mu=\sum_i|D_i|-\left|\bigcup_iD_i\right|.
\]
The defect bound and \eqref{eq:damp-first-slack-budget} give
\[
 \mu+\sum_i\delta_i=1.
\]
Thus there are two cases.

Suppose first that \(\mu=1\).  Then every \(\delta_i=0\), so
Lemma~\ref{lem:damp-rooted-fiber-defect} makes every \(C_i\) a rooted
geodesic path.  Exactly one coordinate, say \(x\), occurs twice among the
defect sets, and every other coordinate occurs once.  There is a coordinate
\(a\ne x\): otherwise two occurrences of the sole coordinate as a defect
would require a path with two distinct edge labels in a one-dimensional
cube.  Let \(P\) be the path with \(a\in D(P)\), and let \(t\) be its
terminal coordinate.  Let \(P'\) be a path with \(t\in D(P')\), which
exists by defect coverage.  The path \(P\) realizes \(00,10,11\) on
\(\{a,t\}\).  Since \(a\) is not the repeated coordinate, it does not
belong to \(D(P')\).  Thus it is either absent from \(P'\) or is the
terminal coordinate of \(P'\), after \(t\).  In both cases \(P'\) realizes
\(01\).  Hence \(\{a,t\}\) is shattered by \(C\).  No individual path
shatters a pair, proving \eqref{eq:damp-first-slack-new-face}.

Now suppose that \(\mu=0\).  The sets \(D_i\) partition \(U\), and there
is a unique index \(j\) with \(\delta_j=1\); all other fibers are rooted
geodesic paths.  Set \(A=C_j\).  Since \(A\) is ample,
\begin{equation}
 |A|=1+|\operatorname{supp}(A)|+h(A).
 \label{eq:damp-first-slack-shadow-count}
\end{equation}
Combining this identity with \(|A|=|D_j|+3\) gives
\[
 |\operatorname{supp}(A)\setminus D_j|+h(A)=2.
\]
The first summand is nonzero.  Otherwise \(h(A)=2\); downward closure says
that the two higher shattered sets are coordinate pairs.  Since \(A\) is
ample, these are precisely the two edges of its crossing graph.  That graph
is a forest, so the crossing-forest theorem
\cite[Theorem~5.5]{klavzar2002partial} makes the one-inclusion graph of \(A\)
a \(C_4\)-cactoid.  Its recursive construction has a peripheral coordinate
whose peripheral side does not contain the root.  This coordinate is not a
member of \(D_j\), a contradiction.

Choose \(a\in\operatorname{supp}(A)\setminus D_j\), and let \(P\ne A\)
be the path fiber with \(a\in D(P)\).  Such a fiber exists because the
defect sets partition \(U\).  Let \(t\) be the terminal coordinate of
\(P\), and let \(B\) be the unique fiber with \(t\in D(B)\).  The path
\(P\) realizes \(00,10,11\) on \(\{a,t\}\).  Since \(a\notin D(B)\), the
fiber \(B\) realizes \(01\): if \(B\) is a path, then \(a\) is absent or
terminal after \(t\); if \(B=A\), take a witness for \(t\in D(A)\), and,
when that witness has \(a\)-value one, use its \(a\)-mate, which belongs to
\(A\) because \(a\notin D(A)\).  Thus \(C\) shatters \(\{a,t\}\).

If \(h(A)=0\), choose distinct coordinates
\(a,b\in\operatorname{supp}(A)\setminus D_j\).  No individual fiber
shatters a pair, whereas the preceding paragraph gives a pair shattered by
\(C\).  This proves \eqref{eq:damp-first-slack-new-face} in this case.

If \(h(A)=1\), let \(\{a,b\}\) be the unique higher shattered set of
\(A\), and let \(a_0\) be the unique member of
\(\operatorname{supp}(A)\setminus D_j\).  The coordinate \(a_0\) is not
one of \(a,b\).  Indeed, the crossing graph of \(A\) has one edge, so the
crossing-forest theorem again makes its one-inclusion graph a
\(C_4\)-cactoid with one square block.  Root its block tree at the block
containing \(0\).  If a leaf block away from the root is a bridge, its
coordinate is nondefect and lies outside \(\{a,b\}\).  Otherwise the unique
square is the leaf farthest from the root.  Both of its directions are then
nondefect: every vertex on its side away from the attachment has its mate
in the square.  This contradicts the uniqueness of \(a_0\).  If the block
tree has only one block, \(A\) is the square and gives the same
contradiction.  Thus the first alternative occurs, and its nondefect bridge
coordinate must be \(a_0\notin\{a,b\}\).
Applying the preceding paragraph with \(a_0\) therefore produces a pair
different from \(\{a,b\}\).  It is a new higher shattered set of \(C\),
which proves \eqref{eq:damp-first-slack-new-face}.
\end{proof}

\begin{lemma}[Exit structure of a low-slack fiber]
\label{lem:damp-low-slack-exit-structure}
Let \(C\subseteq Q_U\) be ample, contain \(0\), and have \(|C|\ge2\).
Write
\[
 E(C)=\operatorname{supp}(C)\setminus D(C),
 \qquad
 h(C)=|\{R\in\operatorname{Sh}(C):|R|\ge2\}|,
\]
and suppose that \(\delta(C)=|C|-2-|D(C)|\le2\).  Then:
\begin{enumerate}
 \item \(E(C)\ne\varnothing\), no set of cardinality at least three is
       shattered by \(C\), and the crossing graph of \(C\) is a forest;
 \item for every \(a\in D(C)\), there is a \(t\in E(C)\) whose value-one
       peripheral side is contained in the value-one \(a\)-semicube.  Thus
       \(C\) realizes \(00,10,11\), but not \(01\), on the ordered pair
       \((a,t)\);
 \item if \(|E(C)|=1\), every pair shattered by \(C\) is contained in
       \(D(C)\).
\end{enumerate}
Moreover,
\begin{equation}
 \delta(C)=|E(C)|+h(C)-1.
 \label{eq:damp-low-slack-potential}
\end{equation}
\end{lemma}

\begin{proof}
Ampleness and the definition of the support give
\[
 |C|=1+|\operatorname{supp}(C)|+h(C).
\]
Since \(D(C)\subseteq\operatorname{supp}(C)\), this proves
\eqref{eq:damp-low-slack-potential}.

We first show that \(E(C)\) is nonempty.  Otherwise
\(h(C)=\delta(C)+1\le3\).  A shattered triple would contribute its three
pairs and the triple itself to \(h(C)\), so all higher shattered sets are
pairs.  These pairs are precisely the edges of the crossing graph.  If that
graph is a forest, the crossing-forest theorem
\cite[Theorem~5.5]{klavzar2002partial} makes the one-inclusion graph of
\(C\) a \(C_4\)-cactoid.  Root its block--cut tree at the block containing
\(0\).  An outer edge of a leaf \(K_2\)-block, or a new direction in a leaf
square of a leaf \(C_4\)-tree, has a peripheral side that avoids the root.
Its direction belongs to \(E(C)\), a contradiction.

If the crossing graph is not a forest, its at most three edges form one
triangle; every other component is an isolated vertex.  Crossing-graph
components are the crossing graphs of the blocks, by
\cite[Theorem~3.4]{klavzar2002partial} applied blockwise.  The triangular
block is one of \(Q_3,Q_3^-\), and \(C_6\), by
\cite[Proposition~5.2]{klavzar2002partial}.  The cube \(Q_3\) has four
higher shattered sets, and \(C_6\) is not ample, so the block is \(Q_3^-\).
Every other block is a \(K_2\).  Root the block--cut tree at the block
containing \(0\) and choose a leaf block away from the root, allowing the
root block itself when it is the only block.  A \(K_2\)-leaf has an outer
peripheral edge direction avoiding the root.  For a \(Q_3^-\)-leaf,
translate its attachment vertex (or the root in the one-block case) to zero
and write \(w\ne0\) for the omitted cube vertex.  For every coordinate
\(t\in\operatorname{supp}(w)\), each present vertex with \(t\)-value one
has its \(t\)-mate.  The corresponding side avoids the attachment vertex,
so it remains peripheral in all of \(C\), and \(t\in E(C)\).  Both cases
contradict \(E(C)=\varnothing\).

Equation~\eqref{eq:damp-low-slack-potential} now gives \(h(C)\le2\).
Thus no triple is shattered, the crossing graph is a forest, and \(C\) is a
\(C_4\)-cactoid by the cited theorem.

We use the peripheral-split lemma for median graphs
\cite[p.~244, paragraph preceding Theorem~5.4]{klavzar2002partial}: inside
either side of any split there is a peripheral split whose peripheral side
is contained in that side.  Let \(a\in D(C)\), and choose an unmatched
vertex \(z\) in the value-one \(a\)-semicube.  That semicube contains an
edge.  Indeed, on a geodesic from \(0\) to \(z\), the first vertex after the
\(a\)-edge cannot be \(z\), since its predecessor would be
\(z^{\{a\}}\in C\).  Apply the lemma to the value-one \(a\)-semicube and
let \(t\) be the resulting peripheral direction.  Its peripheral side is
contained in that semicube and avoids the root.  Thus \(t\ne a\) and
\(t\in E(C)\): the direction cannot be \(a\), since \(a\in D(C)\) says
precisely that the value-one \(a\)-side is not peripheral.  The root and the
two endpoints of a boundary \(t\)-edge realize \(00,10,11\) on \((a,t)\),
whereas containment excludes \(01\).  This proves item~2.

Finally, suppose that the unique member \(t\) of \(E(C)\) crossed a
coordinate \(a\).  Then \(a\in D(C)\), so item~2 supplies an
\(s\in E(C)\) whose peripheral side is contained in the value-one
\(a\)-semicube.  Uniqueness gives \(s=t\).  But crossing of \(a\) and
\(t\) makes the value-one \(t\)-semicube meet both \(a\)-semicubes, so it
cannot be contained in the value-one \(a\)-semicube.  This contradiction
proves item~3.
\end{proof}

\begin{lemma}[Second-slack rooted fibers]
\label{lem:damp-second-slack-rooted-fibers}
Let \(C_1,\ldots,C_m\subseteq Q_U\) be ample families containing \(0\),
with \(|C_i|\ge2\) and \(|U|\ge4\).  Suppose that
\begin{equation}
 \bigcup_{i=1}^mD(C_i)=U,
 \qquad
 \sum_{i=1}^m(|C_i|-2)=|U|+2.
 \label{eq:damp-second-slack-budget}
\end{equation}
Then, with \(C=\bigcup_iC_i\),
\begin{equation}
 h(C)>\sum_{i=1}^m h(C_i).
 \label{eq:damp-second-slack-new-face}
\end{equation}
\end{lemma}

\begin{proof}
Put \(D_i=D(C_i)\), \(E_i=E(C_i)\), \(h_i=h(C_i)\), and
\[
 \mu=\sum_i|D_i|-|U|.
\]
The defect bound and \eqref{eq:damp-second-slack-budget} give
\begin{equation}
 \mu+\sum_i\delta(C_i)=2.
 \label{eq:damp-second-slack-partition}
\end{equation}
In particular every fiber satisfies the hypotheses of
Lemma~\ref{lem:damp-low-slack-exit-structure}, and
\begin{equation}
 \sum_i|E_i|+\sum_i h_i=m+2-\mu.
 \label{eq:damp-second-slack-exit-count}
\end{equation}

We first record the common routing step.  Suppose that a coordinate
\(a\in D_i\) occurs in no other defect set.  By
Lemma~\ref{lem:damp-low-slack-exit-structure}, choose \(t\in E_i\) for
which \(C_i\) realizes \(00,10,11\) on \((a,t)\).  Defect coverage gives
an index \(j\ne i\) with \(t\in D_j\).  Choose a witness
\(z\in C_j\) with \(z_t=1\) and \(z^{\{t\}}\notin C_j\).  If \(z_a=0\),
then \(z\) realizes the missing trace \(01\).  If \(z_a=1\), uniqueness of
the defect occurrence gives \(a\notin D_j\), so
\(z^{\{a\}}\in C_j\) and realizes \(01\).  Thus \(\{a,t\}\) is shattered
by \(C\).  We call it the pair routed from \(a\).

There are three cases in \eqref{eq:damp-second-slack-partition}.  If
\(\mu=2\), all fibers have \(\delta=0\), hence are rooted geodesic paths
and have \(h_i=0\).  The \(|U|+2\) defect occurrences leave a coordinate
that occurs only once, because \(|U|\ge4\).  Its routed pair proves
\eqref{eq:damp-second-slack-new-face}.

Suppose that \(\mu=1\).  Exactly one coordinate is repeated among the
defect sets, and at least \(|U|-1\ge3\) coordinates occur only once.
Moreover \(\sum_i h_i\le\sum_i\delta(C_i)=1\).  If the sum is zero, any
unique defect coordinate gives the required routed pair.  If it is one,
let \(R\) be the unique pair shattered inside a fiber.  Choose a unique
defect coordinate outside the two members of \(R\).  Its routed pair is
different from \(R\), so \(C\) shatters at least two pairs.

It remains that \(\mu=0\).  The defect sets partition \(U\).  If
\(\sum_i h_i\le1\), choose a defect coordinate outside the unique internal
pair when that pair exists; there is such a coordinate because \(|U|\ge4\).
Its routed pair is new.  Finally suppose that \(\sum_i h_i=2\).
Equation~\eqref{eq:damp-second-slack-exit-count} and the nonemptiness of
every \(E_i\) imply \(|E_i|=1\) for all \(i\).  Hence every pair shattered
inside \(C_i\) is contained in \(D_i\), by
Lemma~\ref{lem:damp-low-slack-exit-structure}.  The two internal pairs are
distinct, because the sets \(D_i\) are disjoint.  On the other hand, a
routed pair from \(a\in D_i\) has its other coordinate in some
\(D_j\), \(j\ne i\); it is therefore not internal to any fiber.  Thus the
union shatters the two internal pairs and an additional routed pair.  This
proves \eqref{eq:damp-second-slack-new-face} in every case.
\end{proof}

\begin{lemma}[Structure at fiber slack three]
\label{lem:damp-third-slack-fiber-structure}
Let \(C\subseteq Q_U\) be ample, contain \(0\), have \(|C|\ge2\), and
satisfy \(\delta(C)=3\).  Then \(E(C)\ne\varnothing\), no set of
cardinality at least three is shattered by \(C\), and the crossing graph of
\(C\) is either a forest or consists of one triangle and isolated vertices.
In the forest case, every \(a\in D(C)\) has a \(t\in E(C)\) whose
value-one peripheral side is contained in the value-one \(a\)-semicube.

In the triangular case, \(h(C)=3\), \(|E(C)|=1\), and the unique nontrivial
block is \(Q_3^-\); all other blocks are \(K_2\)'s.  Root the block--cut
tree at the root.  Either its terminal block is a \(K_2\), in which case
some \(a\in D(C)\) and the unique \(t\in E(C)\) have the containment above,
or its terminal block is \(Q_3^-\).  In the latter case, if \(a,b,t\) are
its three directions, then \(t\) is the unique member of \(E(C)\) and the block
realizes every trace on \(\{a,b,t\}\) except the singleton trace on \(t\).
\end{lemma}

\begin{proof}
The potential identity \eqref{eq:damp-low-slack-potential} remains valid and
gives
\begin{equation}
|E(C)|+h(C)=4.
\label{eq:damp-third-slack-potential}
\end{equation}
We claim that \(E(C)\ne\varnothing\).  Otherwise \(h(C)=4\).  A shattered
triple is strongly shattered by ampleness and gives a \(Q_3\)-block; its
block--cut tree has a peripheral side avoiding the root.  If no triple is
shattered, the crossing graph has four edges.  When it is triangle-free,
\cite[Theorem~5.1]{klavzar2002partial} makes \(C\) a cube-free median graph,
and the peripheral-split lemma supplies a peripheral side avoiding the
root.  When it contains a triangle, its fourth edge either forms a separate
crossing component, in which case a terminal block is peripheral, or is
pendant at the triangle.  Any additional isolated crossing-graph vertex is
a bridge block, whose block--cut tree has an outward peripheral leaf.  Thus
only a two-connected block with the four-vertex paw as crossing graph
remains.  Its shattered-set complex has the three minimal nonfaces
\[
 \{a,b,c\},\qquad\{b,t\},\qquad\{c,t\}.
\]
The signed-minimal-nonface equality of
Theorem~\ref{thm:damp-fixed-shadow-signed-clutter} has, modulo cube
translation and interchanging \(b,c\), two solutions.  Their satisfying
families are
\[
 A=\{6,7,9,10,11,12,13,14,15\}
 \quad\hbox{and}\quad
 B=\{2,3,9,10,11,12,13,14,15\}.
\]
For clarity, this two-orbit calculation is just inclusion--exclusion for
the three forbidden subcubes: translate the forbidden trace on
\(\{a,b,c\}\) to \(000\); the equality \(|C|=9\) leaves six pairs of
forbidden traces on \(\{b,t\}\) and \(\{c,t\}\), which form the two orbits
above under flipping \(t\) and interchanging \(b,c\).  The following table
lists, for each possible root, a nonempty set of nondefect directions; the
entries are written as four-bit masks in the order \(a,b,c,t\).
\[
\begin{array}{c|ccccccccc}
C&A\text{-root }6&7&9&10&11&12&13&14&15\\ \hline
E(C)&6&7&9&10&11&12&13&14&15
\end{array}
\]
\[
\begin{array}{c|ccccccccc}
C&B\text{-root }2&3&9&10&11&12&13&14&15\\ \hline
E(C)&2&3&9&10&11&8&9&10&11
\end{array}
\]
Thus the pendant-triangle case also contradicts
\(E(C)=\varnothing\).

Equation~\eqref{eq:damp-third-slack-potential} gives \(h(C)\le3\), so no
triple is shattered.  The crossing graph has at most three edges and is
therefore a forest unless its nontrivial component is a triangle.  In the
forest case \(C\) is a median \(C_4\)-cactoid, and the nested
peripheral-split argument from
Lemma~\ref{lem:damp-low-slack-exit-structure} proves the asserted
containment.

Suppose that the crossing graph contains a triangle.  Then \(h(C)=3\), so
\eqref{eq:damp-third-slack-potential} gives \(E(C)=\{t\}\).  The triangular block is \(Q_3^-\), since
\(Q_3\) shatters a triple and \(C_6\) is not ample; all other blocks are
bridges.  Every leaf block away from the root supplies an outward peripheral
direction, and different blocks use different directions.  Hence the
rooted block--cut tree has a unique terminal block and is a path.

If the terminal block is a bridge of direction \(t\), let \(x\) be its
attachment vertex and \(y\) its outer endpoint.  We have \(x\ne0\): if
\(x=0\), the \(Q_3^-\)-block on the other side has a peripheral direction
avoiding \(0\), distinct from \(t\).  Choose \(a\in\operatorname{supp}(x)\).
The value-one \(t\)-side is \(\{y\}\), lies in the value-one
\(a\)-semicube, and makes \(a\) a defect coordinate.  This gives the first
alternative.

It remains that the terminal block is \(Q_3^-\).  Translate its attachment
vertex, or the root in the one-block case, to zero, and let \(w\ne0\) be
the omitted cube vertex.  For every coordinate in \(\operatorname{supp}(w)\),
the corresponding value-one side is peripheral and avoids the attachment.
It remains peripheral in all of \(C\), so uniqueness of \(E(C)\) forces
\(w=\{t\}\).  The block therefore realizes exactly the other seven traces
on its directions \(a,b,t\), as claimed.
\end{proof}

\begin{lemma}[Third-slack rooted fibers]
\label{lem:damp-third-slack-rooted-fibers}
Let \(C_1,\ldots,C_m\subseteq Q_U\) be ample families containing \(0\),
with \(|C_i|\ge2\) and \(|U|\ge4\).  Suppose that
\begin{equation}
\bigcup_{i=1}^mD(C_i)=U,
 \qquad
 \sum_{i=1}^m(|C_i|-2)=|U|+3.
\label{eq:damp-third-slack-budget}
\end{equation}
Then, with \(C=\bigcup_iC_i\),
\begin{equation}
h(C)>\sum_{i=1}^m h(C_i).
\label{eq:damp-third-slack-new-face}
\end{equation}
\end{lemma}

\begin{proof}
Write \(D_i=D(C_i)\), \(E_i=E(C_i)\), \(h_i=h(C_i)\), and
\[
 \mu=\sum_i|D_i|-|U|.
\]
The defect bound and \eqref{eq:damp-third-slack-budget} give
\begin{equation}
\mu+\sum_i\delta(C_i)=3.
\label{eq:damp-third-slack-distribution}
\end{equation}
Call a coordinate \emph{unique} if it belongs to exactly one defect set.
At most \(\mu\) coordinates are nonunique, so there are at least
\(|U|-\mu\) unique coordinates.  For every unique \(a\in D_i\) in a fiber
of slack at most two, Lemma~\ref{lem:damp-low-slack-exit-structure} gives
an exposed pair \(\{a,t\}\), with \(t\in E_i\), on which \(C_i\) realizes
all traces except \(01\).  Defect coverage supplies a fiber with
\(t\) as a defect; because \(a\) is unique, downward flips in coordinate
\(a\) produce the missing trace exactly as in the common routing step of
Lemma~\ref{lem:damp-second-slack-rooted-fibers}.  Thus the union shatters
this pair.  Routed pairs from distinct unique coordinates are distinct
except possibly when two coordinates route to each other, so \(k\) unique
coordinates produce at least \(\lceil k/2\rceil\) distinct pairs.

If \(\mu=3\), all fibers are paths and one unique coordinate suffices.  If
\(\mu=2\), then \(\sum_i\delta(C_i)=1\) and \(\sum_i h_i\le1\).  When an
internal pair exists, its fiber has a unique exit and the pair is contained
in its defect set by Lemma~\ref{lem:damp-low-slack-exit-structure}.  A
routed pair spans two defect sets and is therefore new.  If \(\mu=1\),
there are at least three unique coordinates and
\(\sum_i h_i\le2\).  When the latter sum is at most one, the two distinct
routed pairs force a new one.  When it equals two, equality in the potential
identities makes every exit set a singleton.  All internal pairs are
contained in individual defect sets, whereas every routed pair spans two of
them.

It remains that \(\mu=0\), so the defect sets partition \(U\) and
\(\sum_i\delta(C_i)=3\).  First suppose that every fiber has slack at most
two, and put \(H=\sum_i h_i\).  The potential identity gives
\begin{equation}
\sum_i(|E_i|-1)=3-H.
\label{eq:damp-third-slack-exit-budget}
\end{equation}
If \(H\le1\), the at least four unique coordinates yield two routed pairs.
If \(H=3\), \eqref{eq:damp-third-slack-exit-budget} makes every \(E_i\) a singleton, so every routed pair is
new.  Finally, if \(H=2\), \eqref{eq:damp-third-slack-exit-budget} gives total exit surplus one.  Thus one
fiber has two exits and the others have one.  The exceptional fiber has at
most one internal pair, since otherwise its slack would be at least three;
all other internal pairs lie within their defect sets.  Of the two distinct
routed pairs, at most one can therefore be internal.

We are left with one fiber \(A\) of slack three and path fibers of slack
zero.  If the crossing graph of \(A\) is a forest and \(D(A)\ne\varnothing\),
Lemma~\ref{lem:damp-third-slack-fiber-structure} gives an exposed pair from
a defect of \(A\), and defect coverage completes it in another fiber.  This
pair is not internal because its peripheral containment omits trace \(01\).
If
\(D(A)=\varnothing\), then \(E(A)=\operatorname{supp}(A)\).  A forest with
\(h(A)\ge2\) has at least \(h(A)+1\) nonisolated vertices, contradicting
\(|E(A)|+h(A)=4\).  Hence \(h(A)\le1\), while the four path defects yield at
least two routed pairs.

Finally suppose that the crossing graph of \(A\) contains a triangle.  In
the bridge-terminal alternative of
Lemma~\ref{lem:damp-third-slack-fiber-structure}, its exposed pair is
completed as before.  In the \(Q_3^-\)-terminal alternative, let
\(a,b,t\) be the three block directions.  The family \(A\) realizes all
traces except the singleton \(t\), and \(a,b\in D(A)\).  The unique exit
\(t\) is a defect of a path fiber \(P\).  Choose a \(t\)-defect witness in
\(P\).  Since the defect sets partition \(U\), neither \(a\) nor \(b\) is
a defect of \(P\); successively flipping either coordinate that has value
one produces a vertex of \(P\) whose trace on \(\{a,b,t\}\) is the
singleton \(t\).  The union therefore shatters this triple.  The three
pairs were the only internal higher faces, so \eqref{eq:damp-third-slack-new-face} follows.  This
exhausts \eqref{eq:damp-third-slack-distribution}.
\end{proof}

\begin{corollary}[Second strict facet-excess inequality]
\label{cor:damp-second-strict-facet-excess}
Under the hypotheses of Theorem~\ref{thm:damp-facet-excess}, if \(H\) is
periphery-free, then
\[
 |H|\ge2^{|T|+1}+|F|-|T|+2.
\]
\end{corollary}

\begin{proof}
Put \(U=F\setminus T\), and suppose that equality held in the first strict
bound, so that
\[
 |H|=2^{|T|+1}+|U|+1.
\]
Use the rooted fibers \(C_q\) from the proof of
Theorem~\ref{thm:damp-facet-excess}.  They satisfy
\eqref{eq:damp-first-slack-budget}.  For \(e\in U\), let
\[
 P_e=\{q\in Q_T:e\in\operatorname{supp}(C_q)\}.
\]
Because each \(C_q\) is ample,
\[
 |C_q|=1+|\operatorname{supp}(C_q)|+h(C_q).
\]
Every set with exactly one outside coordinate that is counted by
\(\operatorname{Sh}(P_e)\), and every higher set shattered by
\(C=\bigcup_qC_q\), is shattered by \(H\).  These collections are
disjoint.  Lemma~\ref{lem:damp-first-slack-rooted-fibers} and the Sandwich
inequality therefore give
\[
 \begin{split}
 |\operatorname{Sh}(H)|
 &\ge 2^{|T|}+\sum_{e\in U}|\operatorname{Sh}(P_e)|+h(C)\\
 &\ge 2^{|T|}+\sum_q|\operatorname{supp}(C_q)|+h(C)\\
 &>2^{|T|}+\sum_q\bigl(|C_q|-1\bigr)=|H|,
 \end{split}
\]
contradicting ampleness.
\end{proof}

\begin{corollary}[Uniform facet-rank bound]
\label{cor:damp-uniform-facet-rank-bound}
Under the hypotheses of Theorem~\ref{thm:damp-facet-excess}, if \(H\) is
periphery-free, then
\[
 |H|\ge2^{|T|+1}+4.
\]
In particular, a facet of rank at least five forces \(|H|\ge68\).
\end{corollary}

\begin{proof}
Lemma~\ref{lem:damp-facet-codimension-at-least-two} and
Corollary~\ref{cor:damp-second-strict-facet-excess} give
\[
 |H|\ge2^{|T|+1}+|F\setminus T|+2
      \ge2^{|T|+1}+4.
\]
\end{proof}

\begin{corollary}[Third strict facet-excess inequality]
\label{cor:damp-third-strict-facet-excess}
Under the hypotheses of Theorem~\ref{thm:damp-facet-excess}, suppose that
\(H\) is periphery-free and \(|F\setminus T|\ge4\).  Then
\[
 |H|\ge2^{|T|+1}+|F|-|T|+3.
\]
\end{corollary}

\begin{proof}
Put \(U=F\setminus T\), and suppose equality held in the second strict
bound.  The rooted fibers \(C_q\) from the proof of
Theorem~\ref{thm:damp-facet-excess} then satisfy
\eqref{eq:damp-second-slack-budget}.  Lemma
\ref{lem:damp-second-slack-rooted-fibers} gives
\[
 h\left(\bigcup_qC_q\right)>\sum_qh(C_q).
\]
Using the one-coordinate incidence families \(P_e\) from the proof of
Corollary~\ref{cor:damp-second-strict-facet-excess}, the same disjoint
Sandwich count gives
\[
 \begin{split}
 |\operatorname{Sh}(H)|
 &\ge2^{|T|}+\sum_{e\in U}|\operatorname{Sh}(P_e)|
       +h\left(\bigcup_qC_q\right)\\
 &>2^{|T|}+\sum_q\bigl(|C_q|-1\bigr)=|H|,
 \end{split}
\]
contradicting ampleness.
\end{proof}

\begin{corollary}[Fourth strict facet-excess inequality]
\label{cor:damp-fourth-strict-facet-excess}
Under the hypotheses of Theorem~\ref{thm:damp-facet-excess}, suppose that
\(H\) is periphery-free and \(|F\setminus T|\ge4\).  Then
\[
 |H|\ge2^{|T|+1}+|F|-|T|+4.
\]
\end{corollary}

\begin{proof}
Put \(U=F\setminus T\), and suppose equality held in the third strict
bound.  The rooted fibers from Theorem~\ref{thm:damp-facet-excess} satisfy
\eqref{eq:damp-third-slack-budget}, so Lemma~\ref{lem:damp-third-slack-rooted-fibers} gives a strict
gain in their union's higher shattered sets.  Adding this gain to the
one-coordinate incidence-family Sandwich count, exactly as in
Corollary~\ref{cor:damp-third-strict-facet-excess}, gives
\(|\operatorname{Sh}(H)|>|H|\), contradicting ampleness.
\end{proof}

The low-dimensional relative-peeling theorem proved in
Appendix~\ref{app:damp-five-coordinate-relative-peeling} also eliminates
two ambient dimensions at once.

\begin{corollary}[Ample families in at most six coordinates]
\label{cor:damp-six-coordinate-peeling}
Every nonempty ample family of isometric dimension at most six can
be corner-peeled to any prescribed concept.
\end{corollary}

\begin{proof}
The assertion in dimensions at most five follows from
Theorem~\ref{thm:damp-five-coordinate-relative-peeling}, applied with
\(S\) the prescribed singleton.  In dimension six, orient an active
coordinate so that the desired endpoint lies in the upper section, and write
\[
 H=(B\times\{0\})\cup(M\times\{1\})
 \subseteq Q_{E\cup\{e\}},
 \qquad |E|=5,
\]
for the family.  If one section is empty, the other uses at most five
coordinates and is already covered.  Otherwise the intersection
\(S=B\cap M\) is nonempty: the one-inclusion graph of an ample family is
connected, so a path between the two sections contains an \(e\)-edge.
Moreover, \(S\) is ample.  Indeed, if a support is shattered by \(S\),
then adjoining \(e\) gives a shattered support of \(H\); ampleness makes
it strongly shattered, and the cube crossing \(e\) makes the original
support strongly shattered in \(S\).  The Sandwich equality follows.

Apply Theorem~\ref{thm:damp-five-coordinate-relative-peeling} to
\(S\subseteq B\), ending at any chosen \(s\in S\).  In \(H\), the same
order deletes the lower copy of \(B\setminus S\): such a vertex has no
\(e\)-mate, so its corner cube is unchanged.  It then deletes the lower
copy of \(S\): each old-coordinate corner cube lies in \(S\) and is
duplicated in \(M\), so adjoining the \(e\)-direction completes the
corner cube in \(H\).  The remaining family is the upper copy of \(M\),
which uses at most five coordinates.  By the first sentence it peels
to the prescribed upper endpoint.
\end{proof}

The next identity records the contribution of indexed fibers to the
shattered sets of a family.  It also supports the finite four-coordinate
calculation used in
Proposition~\ref{prop:damp-order-twenty-seven-core-shadow-budget}.

\begin{lemma}[Indexed fiber shadow decomposition]
\label{lem:damp-indexed-fiber-shadow-decomposition}
Let \(T\) and \(U\) be disjoint coordinate sets, let
\(H\subseteq Q_{T\cup U}\), and write
\[
 C_q=\{z\in Q_U:(q,z)\in H\}
 \qquad(q\in Q_T).
\]
Assume that \(0\in C_q\) for every \(q\in Q_T\).  For
\(R\subseteq U\) and \(\alpha\in Q_R\), put
\[
 A_{R,\alpha}
 =\{q\in Q_T:\alpha\in\pi_R(C_q)\},
 \qquad
 \mathcal K_R
 =\bigcap_{\alpha\in Q_R}\operatorname{Sh}(A_{R,\alpha}).
\]
Then
\begin{equation}
 \operatorname{Sh}(H)
 =\{S\cup R:R\subseteq U,\ S\in\mathcal K_R\}.
 \label{eq:damp-indexed-shadow-decomposition}
\end{equation}

Suppose in addition that every \(C_q\) is ample.  Put
\[
 C=\bigcup_{q\in Q_T}C_q,
 \qquad
 P_e=\{q\in Q_T:e\in\operatorname{supp}(C_q)\}
 \quad(e\in U),
\]
and, for a cube family \(A\), let
\[
 h(A)=|\{R\in\operatorname{Sh}(A):|R|\ge2\}|.
\]
Then the exact shadow gap is
\begin{align}
 |\operatorname{Sh}(H)|-|H|
 ={}&\sum_{e\in U}
   \bigl(|\operatorname{Sh}(P_e)|-|P_e|\bigr)
   +h(C)-\sum_{q\in Q_T}h(C_q) \notag\\
 &+\sum_{\substack{R\subseteq U\\|R|\ge2}}
   \left(|\mathcal K_R|
   -\mathbf 1_{\{R\in\operatorname{Sh}(C)\}}\right).
 \label{eq:damp-indexed-shadow-gap}
\end{align}
The first sum and the last sum in
\eqref{eq:damp-indexed-shadow-gap} are nonnegative.  No sign is asserted for
the middle, unindexed union gap.
\end{lemma}

\begin{proof}
Fix \(S\subseteq T\) and \(R\subseteq U\).  The set \(S\cup R\) is
shattered by \(H\) exactly when, for every \(\alpha\in Q_R\), the fibers
that realize \(\alpha\) project onto all of \(Q_S\).  This is equivalent to
\(S\in\operatorname{Sh}(A_{R,\alpha})\) for every \(\alpha\), and proves
\eqref{eq:damp-indexed-shadow-decomposition}.

The common root gives
\[
 \mathcal K_\varnothing=2^T,
 \qquad
 \mathcal K_{\{e\}}=\operatorname{Sh}(P_e).
\]
Since every fiber is ample,
\[
 |H|
 =\sum_q|C_q|
 =2^{|T|}+\sum_{e\in U}|P_e|+\sum_qh(C_q).
\]
Moreover, \(\mathcal K_R\) is nonempty exactly when every trace on \(R\)
occurs in at least one fiber, which is equivalent to
\(R\in\operatorname{Sh}(C)\).  When it is nonempty, it contains
\(\varnothing\).  Hence
\[
 \sum_{\substack{R\subseteq U\\|R|\ge2}}|\mathcal K_R|
 =h(C)+
 \sum_{\substack{R\subseteq U\\|R|\ge2}}
 \left(|\mathcal K_R|
 -\mathbf 1_{\{R\in\operatorname{Sh}(C)\}}\right).
\]
Substitution in \eqref{eq:damp-indexed-shadow-decomposition} gives
\eqref{eq:damp-indexed-shadow-gap}.  The Sandwich inequality makes every
summand in the first sum nonnegative, and the preceding nonemptiness
criterion makes every summand in the last sum nonnegative.
\end{proof}

\begin{proposition}[Complete indexed total-excess-eleven classification in four coordinates]
\label{prop:damp-complete-indexed-seventh-slack-q4}
Let \(|T|=3\), let \(|U|=4\), and let
\[
 H=\bigcup_{q\in Q_T}\bigl(\{q\}\times C_q\bigr)
 \subseteq Q_{T\cup U},
\]
where every \(C_q\subseteq Q_U\) is ample, contains \(0\), and has at least
two members.  Suppose that
\[
 \bigcup_{q\in Q_T}D(C_q)=U,
 \qquad
 \sum_{q\in Q_T}(|C_q|-2)=11.
\]
If \(H\) is ample, then it has at least four corners.  In particular, no
such \(H\) is cornerless.
\end{proposition}

\begin{proof}[Computer-assisted proof]
Direct enumeration gives all \(2{,}720\) rooted ample subfamilies of
\(Q_4\) having order between two and thirteen.  The exact union-state
program checks all fifty-two partitions of total excess eleven among the
eight fibers.  For a fixed covered-defect mask and union family it retains
the largest internal higher-face count, which is exactly the value that
minimizes the unindexed gap.  Among the full-defect terminal states it finds
\(7{,}293\) states with nonpositive gap, including \(148\) with negative
gap.

One common \(S_4\)-action reduces these to \(423\) union rows.  For every
row, the verifier expands every fiber decomposition whose realized internal
count makes the gap nonpositive; thus the negative-gap rows do not lose
nonextremal decompositions.  This gives \(179{,}937\) fiber multisets in
\(76{,}274\) common-coordinate orbits.  Every orbit has a fiber type of
multiplicity one.  Place the least such type at \(0\in Q_T\), quotient the
remaining placements by the \(S_3\)-stabilizer of \(0\), and give every
representative eight times its stabilizer-orbit weight.  Because the marked
fiber occurs once, this is an exact transversal for the full
forty-eight-element automorphism action on \(Q_T\).

The resulting audit represents all \(1{,}313{,}602{,}248\) labelled
assignments.  Exactly \(217{,}368\) are ample, and every ample assignment
has at least four corners.  For each profile, the minimum indexed shadow is
also checked by a direct seven-coordinate shattered-set count.  On the
\(156\) profiles having at most \(168\) labelled assignments, full-labelled
enumeration independently agrees with the quotient calculation, including
the shadow, ampleness, corner, and peripheral-coordinate data.

The reproducible programs and outputs are
\[
\begin{gathered}
 \texttt{explore\_damp\_q4\_seventh\_slack\_nonpositive.py},\\
 \texttt{damp\_q4\_seventh\_slack\_nonpositive.json},\\
 \texttt{verify\_damp\_q4\_seventh\_slack\_indexed.py},\\
 \texttt{damp\_q4\_seventh\_slack\_indexed\_certificate.json}.
\end{gathered}
\]
The two output digests are respectively
\[
\begin{split}
 &\texttt{c176842bba8798943cdb3d51de47874a204a0d108ba917c972d5f8c91bc9d5b0},\\
 &\texttt{1e4ac11fe91392103e057cfbf490167d059aefb6c96c3f8ef7195edfbcbfeeae}.
\end{split}
\]
Positive unindexed gap is incompatible with ampleness by
\eqref{eq:damp-indexed-shadow-gap}, while the exhaustive indexed audit
covers every nonpositive-gap decomposition.  This proves the claim.
\end{proof}

\begin{theorem}[Order bounds]
\label{prop:damp-minimum-order-profile}
The least order of an ample forbidden pc-minor satisfies
\[
 64\le m_{\mathrm{amp}}\le267.
\]
\end{theorem}

\begin{proof}
Proposition~\ref{prop:damp-order-267-obstruction} gives the upper
bound.  For the lower bound, suppose that an ample forbidden pc-minor
\(O\) has order less than \(64\).  Delete corners from \(O\)
as long as possible.  The ample corner-deletion criterion preserves
ampleness, and the process cannot reach the empty family: such a
sequence would be a corner peeling of \(O\).  It therefore leaves
a nonempty cornerless ample family of order at most \(|O|\).

Choose \(H\) of minimum cardinality among all nonempty cornerless
ample families.  The preceding construction ensures that
\(|H|\le|O|<64\).  There are two exhaustive ranges.

If \(|H|<48\),
Corollary~\ref{cor:damp-path-only-reduction-below-forty-eight}
gives a maximum core of VC dimension one.  The path-core exclusion
of Proposition~\ref{prop:damp-path-core-exclusion-below-sixty-four}
then gives a corner of \(H\), contrary to its choice.  We use that
proposition only in the range \(|H|<48\).

If \(48\le|H|\le63\),
Corollary~\ref{cor:damp-vctwo-core-exclusion-through-sixty-three}
again gives a corner of \(H\), a contradiction.  Both ranges are
impossible, so no ample forbidden pc-minor has order below \(64\).
\end{proof}



\section{Conclusion}
\label{sec:damp-conclusion}

The nonample part of the forbidden pc-minor basis of \(\Damp\) is
the family of antipodally twice-punctured cubes.  The ample part has
exact recognition criteria in cubical, algebraic, topological, and
signed-set language.  A complete intrinsic generative classification
of the ample members remains open.

The canonical three-core reduction and square attachment theorem settle
construction and recovery of every extension over a given forbidden
three-core.  Cut vertices and separating edges give further exact
decompositions, and down-set completion constructs additional examples
from the same rooted factors.  Replicated repairs give a separate
family of constructions; over a cornerless seed, terminal layer data
are precisely partitions of the signed layer facets.  In each case,
the unresolved primitive inputs are explicit: forbidden cores, rooted
factors, resolving additions, or relative peeling properties of pieces.

In the USO formulation, rooted factors obstruct a prescribed sink and
separator pieces obstruct attainable boundary out-maps. The exact gluing
laws in Proposition~\ref{prop:damp-uso-cube-gluing} and
Theorem~\ref{thm:damp-edge-gluing-profile} specify which boundary data
must be compatible. Protected missing cubes give the uniform cycle
certificate of Proposition~\ref{prop:damp-protected-uso-cycle}.
These descriptions keep existence of an acyclic map separate from the
behaviour of one chosen representation map.

Two tasks remain. These primitive properties need to follow from
structural conditions on generative parameters, and the constructions
need to recover all forbidden cores.  The terminal cut-vertex and
two-connected examples outside the common-seed class prevent reducing
coverage to that class.  The infinite family with nonpeelable cube
complements also excludes generation solely by corner deletions from
cubes.  Moreover, a three-core shadow can force corners in every
ample realization (Corollary~\ref{cor:damp-three-core-shadow-forces-corners}),
so changing signs alone cannot normalize all primitive inputs to
cornerless families.  These restrictions leave the intrinsic
classification problem open without requiring a bounded list of
examples or parameters.

The independent extremal question currently has the bounds
\(64\le m_{\mathrm{amp}}\le267\).  Table~\ref{tab:damp-generative-scope}
summarizes the construction and recovery results separately from these
bounds and from the remaining primitive-input problem.

\appendix
\section{Relative peeling in five coordinates}
\label{app:damp-five-coordinate-relative-peeling}

This appendix proves the low-dimensional interval theorem used in
Corollary~\ref{cor:damp-six-coordinate-peeling}.  For nested ample
families $S\subseteq B$, a \emph{relative peeling} means a corner
peeling of $B$ in which every member of $B\setminus S$ is deleted
before every member of $S$.  The proof is symbolic; the bounded
five-coordinate computations retained in the campaign archive are
independent regression checks.

\begin{lemma}[One-concept augmentation as a colon test]
\label{lem:damp-one-concept-colon-test}
Let \(B\subseteq Q_E\) be ample, let \(x\in Q_E\setminus B\), and put
\begin{equation}
m_x=\prod_{e\in E}v_{e,1-x_e},
 \qquad
 \delta(x,y)=\{e\in E:x_e\ne y_e\},
 \qquad
 z_e=v_{e,x_e}.
 \label{eq:damp-relative-augmentation-notation}
\end{equation}
Then the following are equivalent.
\begin{enumerate}
 \item \(B\cup\{x\}\) is ample.
 \item The one-generator colon \((J_B:m_x)\) is generated by variables.
 \item For every \(y\in B\), there is an \(e\in\delta(x,y)\) such that
       \(x^{\{e\}}\in B\).
\end{enumerate}
When they hold, the colon is exactly
\begin{equation}
(J_B:m_x)=(z_e:x^{\{e\}}\in B).
 \label{eq:damp-relative-augmentation-colon}
\end{equation}

More generally, suppose that \(B\subseteq B_0\subseteq Q_E\) are ample
and put \(A=B_0\setminus B\).  If \(x\in A\) cannot be adjoined amply to
\(B\), then there is a \emph{blocking vertex} \(y\in B\) for which
\begin{equation}
\delta(x,y)\cap\{e:x^{\{e\}}\in B\}=\varnothing.
 \label{eq:damp-relative-augmentation-blocker}
\end{equation}
Every \(B_0\)-geodesic from \(x\) to such a \(y\) starts with an edge
from \(x\) to another member of \(A\).  Consequently, if no member of
\(A\) can be adjoined amply to \(B\), then the subgraph induced by \(A\)
has no isolated vertex.
\end{lemma}

\begin{proof}
The monomial-colon formula gives
\begin{equation}
(J_B:m_x)
 =\left(\frac{m_y}{\gcd(m_y,m_x)}:y\in B\right)
 =\left(\prod_{e\in\delta(x,y)}z_e:y\in B\right).
 \label{eq:damp-relative-augmentation-colon-formula}
\end{equation}
A variable \(z_e\) belongs to a minimal monomial generating set on the
right only if \(x^{\{e\}}\in B\).  It follows that this colon is generated
by variables exactly when every displayed monomial is divisible by a
variable coming from a cube neighbor of \(x\) in \(B\).  This is precisely
item~3, and it also proves
\eqref{eq:damp-relative-augmentation-colon}.

If \(B\cup\{x\}\) is ample, it is isometric.  A geodesic from \(x\) to
any \(y\in B\) starts with a neighbor \(x^{\{e\}}\in B\) for some
\(e\in\delta(x,y)\), proving item~3.  Conversely, suppose the colon is
generated by variables.  The exact sequence
\begin{equation}
0\longrightarrow J_B\longrightarrow J_B+(m_x)
 \longrightarrow
 \bigl(S_E/(J_B:m_x)\bigr)(-|E|)\longrightarrow0
 \label{eq:damp-relative-augmentation-exact-sequence}
\end{equation}
has an \(|E|\)-linear module on the left, by ampleness of \(B\), and an
\(|E|\)-linear module on the right, by the Koszul resolution of a
variable-generated ideal.  Its Tor sequence therefore makes
\(J_B+(m_x)=J_{B\cup\{x\}}\) \(|E|\)-linear.  The Yang-ideal criterion
implies that \(B\cup\{x\}\) is ample.

If augmentation fails, the negation of item~3 supplies \(y\) satisfying
\eqref{eq:damp-relative-augmentation-blocker}.  Since \(B_0\) is ample,
it is isometric.  The first edge of a \(B_0\)-geodesic from \(x\) to
\(y\) flips some coordinate in \(\delta(x,y)\).  Its other endpoint
cannot belong to \(B\), by
\eqref{eq:damp-relative-augmentation-blocker}, and hence belongs to
\(A\).  The last assertion follows by applying this argument to every
\(x\in A\).
\end{proof}

\begin{lemma}[A conditioned augmentation lifts or has a transverse blocker]
\label{lem:damp-conditioned-augmentation-transverse-blocker}
Let \(E\) and \(F\) be disjoint coordinate sets, let
\(H\subseteq Q_{E\mathbin{\dot\cup}F}\) be ample, fix
\(\alpha\in Q_F\), and put
\[
 K=\{u\in Q_E:(u,\alpha)\in H\}.
\]
Suppose that \(x\notin K\) and \(K\cup\{x\}\) is ample, and write
\(v=(x,\alpha)\).  Then either \(H\cup\{v\}\) is ample, or there is a set
\(J\subseteq E\mathbin{\dot\cup}F\) such that
\begin{equation}
 |J|\ge2,
 \qquad J\cap F\ne\varnothing,
 \qquad
 (H\cup\{v\})\cap Q_J(v)=\{v,v^J\}.
 \label{eq:damp-conditioned-augmentation-transverse-blocker}
\end{equation}
Here \(Q_J(v)\) denotes the \(J\)-face through \(v\).  Thus a
conditioning-valid one-concept extension can fail globally only through an
antipodal blocker using a coordinate that was fixed in the conditioning.
\end{lemma}

\begin{proof}
If \(H\cup\{v\}\) is nonample, Proposition
\ref{prop:damp-failed-augmentation-antipodal-localization} supplies a set
\(J\) of order at least two for which the last equality in
\eqref{eq:damp-conditioned-augmentation-transverse-blocker} holds.  If
\(J\cap F\) were empty, the whole \(J\)-face would lie in the
\(F\)-conditioning \(\alpha\).  Deleting the fixed coordinates would then
give an antipodal-pair conditioning of \(K\cup\{x\}\), contradicting its
ampleness.  Hence \(J\cap F\ne\varnothing\).
\end{proof}

\begin{lemma}[A blocked puncture forces a complete bipartite two-shadow]
\label{lem:damp-blocked-puncture-bipartite-two-shadow}
Let \(H\subseteq Q_E\) be ample, let \(v\in Q_E\setminus H\), and let
\(T,J\subseteq E\) satisfy \(|T|\ge2\), \(|J|\ge2\),
\begin{equation}
 v^U\in H\quad(\varnothing\ne U\subseteq T),
 \qquad
 (H\cup\{v\})\cap Q_J(v)=\{v,v^J\}.
 \label{eq:damp-blocked-puncture-hypotheses}
\end{equation}
Then \(T\cap J=\varnothing\),
\begin{equation}
 v^{J\cup\{t\}}\in H
 \qquad(t\in T),
 \label{eq:damp-blocked-puncture-opposite-neighbours}
\end{equation}
and \(H\) shatters every pair
\begin{equation}
 \{t,e\},
 \qquad t\in T,\quad e\in J.
 \label{eq:damp-blocked-puncture-bipartite-two-shadow}
\end{equation}
Thus a punctured \(T\)-cube and a transverse antipodal blocker do not
merely give one owner: they force the complete bipartite graph
\(K_{T,J}\) into the two-dimensional shattered shadow.
\end{lemma}

\begin{proof}
For \(t\in T\), the vertex \(v^{\{t\}}\) belongs to \(H\).  If
\(t\in J\), this vertex would be a member of \(Q_J(v)\) different from
the two vertices allowed in
\eqref{eq:damp-blocked-puncture-hypotheses}.  Hence \(T\cap J\) is empty.

Put \(y=v^J\in H\) and fix \(t\in T\).  The vertices \(y\) and
\(v^{\{t\}}\) differ exactly on \(J\mathbin{\dot\cup}\{t\}\).  For every
\(e\in J\), the neighbour
\(y^{\{e\}}=v^{J\setminus\{e\}}\) is excluded by
\eqref{eq:damp-blocked-puncture-hypotheses}.  Since \(H\) is ample, it is
isometric.  An \(H\)-geodesic from \(y\) to \(v^{\{t\}}\) must therefore
start in direction \(t\), proving
\eqref{eq:damp-blocked-puncture-opposite-neighbours}.

Now fix \(e\in J\), and choose \(u\in T\setminus\{t\}\).  Relative to the
trace of \(v\) on \(\{t,e\}\), the four vertices
\[
 v^{\{u\}},\qquad v^{\{t\}},\qquad
 v^J,\qquad v^{J\cup\{t\}}
\]
have traces \(00,10,01,11\), respectively.  They all belong to \(H\), so
\(\{t,e\}\) is shattered.  This proves
\eqref{eq:damp-blocked-puncture-bipartite-two-shadow}.
\end{proof}

\begin{corollary}[Antipodal blockers in a relative-peeling gap]
\label{cor:damp-relative-gap-antipodal-blockers}
Let $S\subsetneq B\subseteq Q_E$ be ample, and suppose that
$S\cup\{x\}$ is nonample for every $x\in B\setminus S$.  For each such
$x$ there are $J_x\subseteq E$, $|J_x|\ge2$, and $y_x=x^{J_x}\in S$
such that the $J_x$-face through $x$ contains no other member of
$S\cup\{x\}$.  Moreover, every $B$-geodesic from $x$ to $y_x$ lies in
that face, has all its internal vertices in $B\setminus S$, and has length
$|J_x|$.
\end{corollary}

\begin{proof}
Apply Proposition~\ref{prop:damp-failed-augmentation-antipodal-localization}
to $S\cup\{x\}$.  Since $B$ is ample, it is isometric.  A shortest
$B$-path between the two resulting antipodes flips exactly the coordinates
of $J_x$ and therefore remains in their face.  No internal vertex belongs to
$S$, by the exact two-point intersection in
\eqref{eq:damp-failed-augmentation-antipodal-localization}.
\end{proof}

\begin{lemma}[Relative peeling and saturated ample intervals]
\label{lem:damp-relative-peeling-saturated-chain}
Let \(\varnothing\ne S\subseteq B\) be ample families.  The following
conditions are equivalent.
\begin{enumerate}
\item The family \(B\) has a relative corner peeling that deletes
      exactly \(B\setminus S\) and retains all of \(S\).
\item There is a chain of ample families
      \[
      S=C_0\subset C_1\subset\cdots\subset C_k=B,
      \qquad |C_i\setminus C_{i-1}|=1
      \quad(1\le i\le k).
      \]
\end{enumerate}
\end{lemma}

\begin{proof}
If the asserted peeling exists, stop it when \(S\) remains and reverse
the deleted prefix.  Every intermediate family is ample because corner
deletion preserves ampleness.  This gives the displayed chain.
Conversely, reverse such a chain.  At each step one vertex is deleted
from an ample family and the remaining family is ample.  The ample
corner-deletion criterion makes that vertex a corner, so the reversed
chain is the required peeling prefix.
\end{proof}

\begin{remark}[The ample-interval conjecture]
\label{rem:damp-relative-peeling-moran-warmuth}
Requiring the equivalent conditions in
Lemma~\ref{lem:damp-relative-peeling-saturated-chain} for every
nested ample pair is equivalent, by finite induction, to the interval
conjecture of Moran and Warmuth
\cite[Conjecture~2]{MoranWarmuth16}: whenever two ample families differ
by at least two concepts, their inclusion interval contains a third ample
family.  The Hall class refutes that conjecture
\cite[Remark~4.6]{ChalopiChepoiMoran22}.  Thus the relative-peeling
theorems below are genuinely low-coordinate results; they cannot be
continued by a dimension-free assertion for arbitrary ample families.
By Proposition~\ref{prop:damp-relative-metric-peeling}, the chain is
equivalently a sequence of isometric one-concept extensions from
\(S\) to \(B\).  No peeling of the retained lower family is part of
this relative assertion.
\end{remark}


\begin{proposition}[Simultaneous shift normalization of ample intervals]
\label{prop:damp-relative-shift-normalization}
Let \(\varnothing\ne S\subsetneq B\subseteq Q_E\) be ample.  For a
coordinate \(e\in E\), define the up-shift \(U_e(X)\) by its two projected
sections:
\[
 U_e(X)_0=X_0\cap X_1,
 \qquad
 U_e(X)_1=X_0\cup X_1.
 \]
Then \(U_e(S)\subsetneq U_e(B)\), both shifted families are ample, and
\[
 \operatorname{Sh}(U_e(S))=\operatorname{Sh}(S),
 \qquad
 \operatorname{Sh}(U_e(B))=\operatorname{Sh}(B).
 \]
After shifting both endpoints once in every coordinate, in any order, the
pair becomes
\begin{equation}
\mathcal H^{\uparrow}_{\operatorname{Sh}(S)}
 \subsetneq
 \mathcal H^{\uparrow}_{\operatorname{Sh}(B)},
 \qquad
 \mathcal H^{\uparrow}_{\Sigma}
 :=\{x\in Q_E:\{e:x_e=0\}\in\Sigma\}.
 \label{eq:damp-relative-standard-shift-pair}
\end{equation}
The terminal pair admits a saturated chain of ample families and hence a
relative corner peeling.  Consequently, if \(S\subsetneq B\) has no
relative peeling, then along every complete simultaneous-shift sequence
there is a first coordinate shift that changes a nonpeelable interval into
a peelable one.  More specifically, if \(S\subsetneq B\) has no ample
one-concept intermediate family, then along every such sequence there is a
first shift after which an ample one-concept intermediate family exists.
\end{proposition}

\begin{proof}
On each fiber over \(E\setminus\{e\}\), the up-shift acts as follows:
\[
 \varnothing\mapsto\varnothing,\qquad
 \{0\}\mapsto\{1\},\qquad
 \{1\}\mapsto\{1\},\qquad
 \{0,1\}\mapsto\{0,1\}.
\]
This map is monotone under inclusion,
so \(U_e(S)\subseteq U_e(B)\); the inclusion remains proper because each
shift preserves cardinality.  The Yang paired-shift theorem
\cite{BousquetYangComplex26} says that an up-shift of an ample family is
ample and preserves its shattered complex.  Its coordinate shifts commute,
and shifting every coordinate gives the standard representatives in
\eqref{eq:damp-relative-standard-shift-pair}.  Since
\(S\subseteq B\), one has
\(\operatorname{Sh}(S)\subseteq\operatorname{Sh}(B)\), so these standard
representatives are nested.

It remains to peel the standard pair.  Order the faces in
\(\operatorname{Sh}(B)\setminus\operatorname{Sh}(S)\) so that every proper
subface occurs first.  Adjoin the corresponding concepts, whose zero sets
are those faces, in that order.  Every intermediate zero-set family is a
simplicial complex and therefore an ample down-set.  Lemma
\ref{lem:damp-relative-peeling-saturated-chain} converts this saturated
ample chain into a relative peeling.  If the original interval is not
peelable, the finite shift sequence starts nonpeelable and ends peelable;
its first peelable term has the asserted predecessor.  The same argument,
using the first member of the terminal saturated chain, proves the final
assertion about one-concept intermediate families.
\end{proof}

\begin{proposition}[One-coordinate inverse-shift criterion]
\label{prop:damp-one-coordinate-inverse-shift}
Let \(X\subseteq Y\subseteq Q_{F\cup\{e\}}\) be ample.  Write \(X_0,X_1\)
for the two \(e\)-sections and put
\begin{equation}
A=X_0\cap X_1,\qquad C=X_0\cup X_1,\qquad \widehat X=U_e(X).
 \label{eq:damp-inverse-shift-sections}
\end{equation}
Suppose that \(p=(z,c)\in U_e(Y)\setminus\widehat X\) and that
\(\widehat X\cup\{p\}\) is ample.

If \(c=0\), then some concept in the \(e\)-fiber over \(z\) belongs to
\(Y\setminus X\) and can be adjoined amply to \(X\).

If \(c=1\), then \(z\notin C\), the family \(C\cup\{z\}\) is ample, and,
for
\begin{equation}
N_C(z)=\{i\in F:z^{\{i\}}\in C\},
 \label{eq:damp-inverse-shift-neighbor-set}
\end{equation}
every available lift \(q_b=(z,b)\in Y\) satisfies
\begin{equation}
X\cup\{q_b\}\text{ is ample}
 \quad\Longleftrightarrow\quad
 z^{\{i\}}\in X_b\quad\text{for every }i\in N_C(z).
 \label{eq:damp-inverse-shift-lift-test}
\end{equation}
More precisely, in the residual literal ring put
\[
 I_b=(J_{X_b}:m_z)\qquad(b\in\{0,1\}).
\]
After extension to the full literal ring,
\begin{equation}
(J_X:m_{q_b})=I_b+v_{e,b}I_{1-b},
 \qquad
 I_0+I_1=(J_C:m_z)
   =(v_{i,z_i}:i\in N_C(z)).
 \label{eq:damp-inverse-shift-colon}
\end{equation}
Consequently, if \(i\in N_C(z)\) but \(z^{\{i\}}\notin X_b\), then
\(z^{\{i\}}\in X_{1-b}\setminus X_b\) and
\(v_{e,b}v_{i,z_i}\) is a minimal quadratic generator of
\((J_X:m_{q_b})\).
\end{proposition}

\begin{proof}
First suppose \(c=0\).  If \(z\in C\), then \(z\notin A\), so exactly one
lift of \(z\) belongs to \(X\).  Since the lower lift belongs to \(U_e(Y)\),
both lifts belong to \(Y\).  The lower section of
\(\widehat X\cup\{p\}\) is \(A\cup\{z\}\), which is ample.  Duplicating
\(z\) into the missing section therefore preserves ampleness, by
Lemma~\ref{lem:damp-expansion-separator-duplication}.

Suppose instead that \(z\notin C\).  Both lifts of \(z\) belong to \(Y\).
For \(w\in C\), apply the one-concept augmentation criterion to \(p\) and
the upper lift \((w,1)\in\widehat X\).  The \(e\)-neighbor of \(p\) is
absent, so there is an \(i\in\delta(z,w)\) with
\((z^{\{i\}},0)\in\widehat X\), or equivalently
\(z^{\{i\}}\in A\).  Both lifts of this neighbor belong to \(X\).
The same criterion now shows that either lift of \(z\) can be adjoined
amply to \(X\).

Now suppose \(c=1\).  Since the upper section of \(\widehat X\) is \(C\),
one has \(z\notin C\).  The upper section of
\(\widehat X\cup\{p\}\) is \(C\cup\{z\}\), so this family is ample.
Assume first that every neighbor in
\eqref{eq:damp-inverse-shift-neighbor-set} belongs to \(X_b\).  Given
\((w,d)\in X\), the augmentation criterion in \(C\cup\{z\}\) supplies
an \(i\in\delta(z,w)\) with \(z^{\{i\}}\in C\).  By hypothesis,
\((z^{\{i\}},b)\in X\), so the same criterion proves that
\(X\cup\{q_b\}\) is ample.

Conversely, if \(z^{\{i\}}\in C\setminus X_b\), then it belongs to
\(X_{1-b}\setminus X_b\).  The concepts \(q_b\) and
\((z^{\{i\}},1-b)\in X\) differ exactly in coordinates \(e,i\), but
neither of the two intermediate cube vertices belongs to \(X\).  This is
a blocking vertex for \(q_b\), so \(X\cup\{q_b\}\) is not ample.  This
proves \eqref{eq:damp-inverse-shift-lift-test}.

Finally, sort the generators in the monomial-colon formula according to
their \(e\)-section.  Generators from section \(b\) give \(I_b\), whereas
those from section \(1-b\) acquire the factor \(v_{e,b}\).  This proves
the first equality in \eqref{eq:damp-inverse-shift-colon}.  Since
\(J_C=J_{X_0}+J_{X_1}\), monomial colons give
\((J_C:m_z)=I_0+I_1\).  The augmentation criterion for
\(C\cup\{z\}\) gives the displayed variable generators.  If the
\(i\)-neighbor is missing from \(X_b\), then \(v_{i,z_i}\notin I_b\)
but \(v_{i,z_i}\in I_{1-b}\).  Neither a unit nor either of the two
variables dividing \(v_{e,b}v_{i,z_i}\) belongs to the full colon, so
this quadratic generator is minimal.
\end{proof}

\begin{corollary}[Minimal repair sets give split corners]
\label{cor:damp-minimal-shift-repair-split-corners}
Let \(S\subsetneq B\subseteq Q_E\) be ample and suppose that no
\(q\in B\setminus S\) makes \(S\cup\{q\}\) ample.  Fix an orientation
of every coordinate and, for \(T\subseteq E\), let \(U_T\) denote the
product of the corresponding commuting up-shifts.  Choose \(T\) inclusion
minimal such that \(U_T(S)\subsetneq U_T(B)\) has an ample one-concept
intermediate family.  Then \(T\ne\varnothing\).

For every \(e\in T\), some first augmentation after the \(e\)-shift has
the upper form \(p_e=(z_e,1)\).  Immediately before that shift, put
\[
 X=U_{T\setminus\{e\}}(S),\qquad
 Y=U_{T\setminus\{e\}}(B).
\]
Then \(C=X_0\cup X_1\) satisfies \(C\cup\{z_e\}\) ample.  Its neighbor
set \(N_e=N_C(z_e)\) is an inclusion-minimal member of
\begin{equation}
\mathcal K=\operatorname{Sh}(B)\setminus\operatorname{Sh}(S),
 \label{eq:damp-minimal-repair-support-difference}
\end{equation}
and \(e\notin N_e\).  For every sign \(b\) for which
\((z_e,b)\in Y\), at least one neighbor \(z_e^{\{i\}}\), \(i\in N_e\),
lies exclusively in the opposite section \(X_{1-b}\setminus X_b\), and
the corresponding colon has the minimal quadratic generator
\(v_{e,b}v_{i,(z_e)_i}\).  Moreover, at least one of
\begin{equation}
(z_e,1-b),\qquad (z_e^{\{i\}},b)
 \label{eq:damp-split-corner-relative-neighbors}
\end{equation}
belongs to \(Y\setminus X\).  Thus every such quadratic generator points
to an adjacent concept in the relative layer, in direction \(e\) or \(i\).
If both signs are available, \(N_e\) has an exclusive neighbor on each
side.  In particular, if \(\mathcal K\) is supported only in ranks two
and three, then
\begin{equation}
|N_e|\in\{2,3\}\qquad(e\in T).
 \label{eq:damp-minimal-repair-support-rank}
\end{equation}
\end{corollary}

\begin{proof}
The fully shifted pair is the canonical nested down-set pair from
Proposition~\ref{prop:damp-relative-shift-normalization}, so it has a
one-concept intermediate family.  The unshifted pair has none, proving
that a nonempty minimal \(T\) exists.  Fix \(e\in T\).  By minimality,
the pair \(X\subsetneq Y\) has no ample one-concept intermediate, whereas
its \(e\)-shift does.  A lower first augmentation would lift to \(X\) by
Proposition~\ref{prop:damp-one-coordinate-inverse-shift}, a contradiction.
Thus it has the asserted upper form, and the same proposition gives the
section split and the minimal quadratic colon generator for every
available sign.  The available lift and its opposite-section blocking
neighbor differ exactly in coordinates \(e,i\).  Since \(Y\) is ample,
it is isometric, so a length-two \(Y\)-geodesic contains one of the two
intermediate vertices in
\eqref{eq:damp-split-corner-relative-neighbors}.  Neither belongs to
\(X\), by construction, proving the relative-layer assertion.

It remains to identify the new support.  The family \(C\cup\{z_e\}\) is
ample and deleting \(z_e\) leaves the ample family \(C\), so \(z_e\) is a
corner.  Its incident directions are exactly \(N_e\), and its unique
maximal cube contains every \(z_e^R\) with
\(\varnothing\ne R\subseteq N_e\).  Hence \(C\cup\{z_e\}\) shatters
\(N_e\).  The family \(C\) does not realize the pattern \(z_e|_{N_e}\):
otherwise an isometric path from \(z_e\) to such a realization would have
to start in a direction outside \(N_e\).  Thus \(N_e\) is the unique new
shattered support.  Coordinate shifts preserve the two endpoint shattered
complexes, so \(N_e\in\mathcal K\).  Since adjoining one concept adds
only \(N_e\), downward closure makes it inclusion-minimal in \(\mathcal K\).
This also proves \(e\notin N_e\) and
\eqref{eq:damp-minimal-repair-support-rank}.
\end{proof}

\begin{remark}[A split corner in three coordinates]
\label{ex:damp-three-coordinate-split-corner}
Let \(F=\{a,b\}\), put \(z=00\), and take
\[
 X_0=\{10,11\},\qquad X_1=\{01,11\}.
\]
The lift \(X\subseteq Q_{F\cup\{e\}}\) is a four-vertex path and is
ample.  Its shifted sections are
\(A=\{11\}\subset C=\{10,01,11\}\).  The upper concept
\(p=(00,1)\) can be adjoined to \(U_e(X)\), because
\(C\cup\{00\}=Q_F\).  The lower lift \(q=(00,0)\), however, cannot be
adjoined to \(X\): the concept \((01,1)\in X\) differs from \(q\) in
coordinates \(b,e\), and neither intermediate vertex belongs to \(X\).
Here \(N_C(00)=\{a,b\}\), but the \(b\)-neighbor \(01\) lies exclusively
in \(X_1\).  Thus \(U_e(X\cup\{q\})=U_e(X)\cup\{p\}\) is ample although
\(X\cup\{q\}\) is not.
\end{remark}

\begin{remark}[What the shift reduction does and does not prove]
\label{rem:damp-first-repair-shift}
Example~\ref{ex:damp-three-coordinate-split-corner} rules out a pointwise
inverse-shift theorem even in three coordinates, while the Hall example
rules out unrestricted lifting of a complete relative peeling.  The
proved gain is the necessary normal form in
Corollary~\ref{cor:damp-minimal-shift-repair-split-corners}: a minimum
five-coordinate gap would carry, after partial standardization, a
constellation of rank-two or rank-three split corners with specified
quadratic colon generators.  Ruling out that constellation using the
complementary signed blockers would replace the relative-ear argument
below.  It is a second sufficient proof route, not an equivalence that has
already been established.
\end{remark}

\begin{lemma}[Relative Cohen--Macaulay form of an ample interval]
\label{lem:damp-relative-yang-cm-flag-h}
Let \(\varnothing\ne S\subseteq B\subseteq Q_E\) be ample, and let
\[
 \Psi(B,S)=(\Delta_B,\Delta_S)
\]
be the corresponding relative Yang complex.  Then \(\Psi(B,S)\) is a
fully Cohen--Macaulay relative complex of dimension \(|E|-1\).  Its
balanced flag \(h\)-numbers are
\begin{equation}
h_R(\Psi(B,S))
 =\begin{cases}
   1,&R\in\operatorname{Sh}(B)\setminus\operatorname{Sh}(S),\\
   0,&\text{otherwise},
  \end{cases}
 \qquad (R\subseteq E).
 \label{eq:damp-relative-yang-flag-h}
\end{equation}
Moreover, a relative shelling of \(\Psi(B,S)\) is exactly an ordering
\(p_1,\ldots,p_k\) of \(B\setminus S\) for which every family
\[
 S_i=S\cup\{p_1,\ldots,p_i\}
 \qquad(0\le i\le k)
\]
is ample.
\end{lemma}

\begin{proof}
The Yang complexes \(\Delta_S\subseteq\Delta_B\) are
Cohen--Macaulay of the same dimension, by the
ampleness--linear-resolution dictionary and Alexander duality.  The
relative face module occurs in the exact sequence
\[
 0\longrightarrow M[\Psi(B,S)]\longrightarrow k[\Delta_B]
 \longrightarrow k[\Delta_S]\longrightarrow0.
\]
The depth lemma therefore makes the relative module Cohen--Macaulay as
well.  This is the fully Cohen--Macaulay presentation in the terminology
of \cite{CodenottiKatthanSanyal19}.

For \(X\in\{S,B\}\), the flag \(f\)-number on a color set \(R\) counts
the faces of \(\Delta_X\) using exactly the colors in \(R\).  Such a face
is a pattern realized by the projection of \(X\) to \(R\), and hence
\begin{equation}
f_R(\Delta_X)=|\pi_R(X)|
 =\bigl|\{T\subseteq R:T\in\operatorname{Sh}(X)\}\bigr|.
 \label{eq:damp-relative-yang-flag-f}
\end{equation}
Here the second equality uses minor closure of ampleness: the projection
is ample, and its shattered supports are precisely the shattered supports
of \(X\) contained in \(R\).  Boolean M\"obius inversion gives
\[
 h_R(\Delta_X)=\mathbf 1_{R\in\operatorname{Sh}(X)}.
\]
Subtracting the two flag \(h\)-vectors proves
\eqref{eq:damp-relative-yang-flag-h}.

Finally, the relative shelling condition for adjoining the facet of
\(p_i\) says precisely that
\((J_{S_{i-1}}:m_{p_i})\) is generated by variables.  By Lemma
\ref{lem:damp-one-concept-colon-test}, this is equivalent to ampleness of
\(S_i\).  Applying this at every step proves the last assertion.
\end{proof}

\begin{lemma}[Balanced Artinian model of an ample interval]
\label{lem:damp-relative-balanced-artinian-reduction}
For \(X\subseteq Q_E\) ample, put
\[
 \theta_e=v_{e,0}+v_{e,1}\qquad(e\in E).
\]
Then \(\boldsymbol\theta=(\theta_e:e\in E)\) is a linear system of
parameters for \(k[\Delta_X]\), and
\begin{equation}
\frac{k[\Delta_X]}{(\boldsymbol\theta)}
 \cong
 \frac{k[x_e:e\in E]}
 {\bigl(x_e^2:e\in E;\ x_R:R\notin\operatorname{Sh}(X)\bigr)},
 \qquad x_R:=\prod_{e\in R}x_e.
 \label{eq:damp-yang-balanced-artinian-reduction}
\end{equation}
In particular, the residue classes \(x_R\),
\(R\in\operatorname{Sh}(X)\), form a vector-space basis.

For a nested ample pair \(S\subseteq B\), quotienting the relative face
module by \(\boldsymbol\theta\) gives the monomial module
\begin{equation}
\frac{M[\Psi(B,S)]}{(\boldsymbol\theta)M[\Psi(B,S)]}
 \cong
 \operatorname{span}_k
 \{x_R:R\in\operatorname{Sh}(B)\setminus
                    \operatorname{Sh}(S)\}.
 \label{eq:damp-relative-artinian-support-module}
\end{equation}
Its nonzero multiplication maps are exactly
\begin{equation}
x_e x_R=x_{R\cup\{e\}}
 \quad\text{when }e\notin R\text{ and }
 R\cup\{e\}\in
 \operatorname{Sh}(B)\setminus\operatorname{Sh}(S).
 \label{eq:damp-relative-artinian-incidence}
\end{equation}
\end{lemma}

\begin{proof}
Every facet of the balanced complex \(\Delta_X\) contains exactly one
vertex of each color.  The coefficient matrix of the color sums on any
facet is therefore the identity matrix, so the color sums form a linear
system of parameters.  Since \(\Delta_X\) is Cohen--Macaulay, they form a
regular sequence.

Modulo the color sums, write \(x_e\) for the image of \(v_{e,0}\); then
the image of \(v_{e,1}\) is \(-x_e\).  The nonface
\(\{v_{e,0},v_{e,1}\}\) gives \(x_e^2=0\).  More generally, a signed face
on a color set \(R\) maps, up to sign, to \(x_R\).  Such a signed face is
absent from \(\Delta_X\) for at least one sign pattern exactly when
\(R\) is not shattered by \(X\).  These are all the Stanley--Reisner
relations after the substitution, proving
\eqref{eq:damp-yang-balanced-artinian-reduction} and its basis statement.

Apply the same regular sequence to the exact sequence
\[
 0\longrightarrow M[\Psi(B,S)]\longrightarrow k[\Delta_B]
 \longrightarrow k[\Delta_S]\longrightarrow0.
\]
Lemma~\ref{lem:damp-relative-yang-cm-flag-h} makes all three modules
Cohen--Macaulay of the same dimension, so successive quotienting by the
regular sequence preserves exactness.  The resulting map sends \(x_R\)
to itself when \(R\in\operatorname{Sh}(S)\) and to zero otherwise.
Its kernel is therefore precisely the span in
\eqref{eq:damp-relative-artinian-support-module}.  Since shattered-set
families are order ideals, multiplication has exactly the form
\eqref{eq:damp-relative-artinian-incidence}.
\end{proof}

\begin{lemma}[The five-color Artinian graph]
\label{lem:damp-q5-artinian-disjointness-graph}
Suppose that the support difference
\(\mathcal K=\operatorname{Sh}(B)\setminus\operatorname{Sh}(S)\)
in Lemma~\ref{lem:damp-relative-balanced-artinian-reduction} is supported
only in ranks two and three, and that \(|E|=5\).  Put
\begin{equation}
G_2=\mathcal K\cap\binom E2,
 \qquad
 G_3=\{E\setminus K:K\in\mathcal K\cap\tbinom E3\}.
 \label{eq:damp-q5-artinian-two-graphs}
\end{equation}
The multiplication graph of the Artinian relative module has bipartition
\(G_2\mathbin{\dot\cup}G_3\), with an edge between
\(R\in G_2\) and \(Q\in G_3\) exactly when \(R\cap Q=\varnothing\).
That edge records multiplication by the unique variable \(x_e\), where
\begin{equation}
\{e\}=E\setminus(R\cup Q).
 \label{eq:damp-q5-artinian-edge-color}
\end{equation}
Thus this multiplication graph is a subgraph of the bipartite double
cover of the Petersen graph \(KG(5,2)\).  It has no four-cycle, and the
complementary ample interval has the same colored graph with its two
parts exchanged.
\end{lemma}

\begin{proof}
By \eqref{eq:damp-relative-artinian-incidence}, a nonzero multiplication
from the basis element \(x_R\) in degree two to the basis element
\(x_K\) in degree three occurs exactly when \(K=R\cup\{e\}\).  Writing
\(Q=E\setminus K\), this is equivalent to \(R\cap Q=\varnothing\), and
the remaining color is the one in
\eqref{eq:damp-q5-artinian-edge-color}.  This is precisely the stated
subgraph of the bipartite double cover of \(KG(5,2)\).

For completeness, a four-cycle would give distinct two-sets
\(R_1,R_2\) and \(Q_1,Q_2\) such that every \(R_i\) is disjoint from
every \(Q_j\).  Then both \(R_i\)'s lie in
\(E\setminus(Q_1\cup Q_2)\), which has size at most two and therefore
cannot contain two distinct two-sets.  Finally, support complementation
interchanges \(G_2\) and \(G_3\), while disjointness and the remaining
edge color are symmetric.
\end{proof}

\begin{lemma}[The support defect of a one-concept augmentation]
\label{lem:damp-relative-singleton-support-defect}
Let \(\varnothing\ne S\subsetneq B\subseteq Q_E\) be ample, put
\[
 \mathcal K=\operatorname{Sh}(B)\setminus\operatorname{Sh}(S),
\]
and, for \(p\in B\setminus S\), put
\begin{equation}
\mathcal H_S(p)
 =\operatorname{Sh}(S\cup\{p\})\setminus
  \operatorname{Sh}(S).
 \label{eq:damp-relative-singleton-defect}
\end{equation}
Then \(\varnothing\ne\mathcal H_S(p)\subseteq\mathcal K\), and the
following conditions are equivalent:
\begin{enumerate}
 \item \(S\cup\{p\}\) is ample;
 \item \(|\mathcal H_S(p)|=1\);
 \item \((J_S:m_p)\) is generated by variables;
 \item the intersection of the facet \(F_p\) with \(\Delta_S\) is pure
       of dimension \(|E|-2\).
\end{enumerate}

For every \(R\subseteq E\),
\begin{equation}
|\pi_R(B)\setminus\pi_R(S)|
 =|\{K\in\mathcal K:K\subseteq R\}|.
 \label{eq:damp-relative-projection-defect-count}
\end{equation}
Consequently, for \(R\in\mathcal K\),
\begin{equation}
R\in\mathcal H_S(p)
 \quad\Longleftrightarrow\quad
 \{K\in\mathcal K:K\subseteq R\}=\{R\}
 \quad\text{and}\quad
 p|_R\notin\pi_R(S).
 \label{eq:damp-relative-singleton-defect-supports}
\end{equation}
\end{lemma}

\begin{proof}
Sauer's inequality and ampleness of \(S\) give
\[
 |\operatorname{Sh}(S\cup\{p\})|
 \ge |S|+1
 =|\operatorname{Sh}(S)|+1,
\]
so \(\mathcal H_S(p)\ne\varnothing\).  Monotonicity of shattering and
\(S\cup\{p\}\subseteq B\) give
\(\mathcal H_S(p)\subseteq\mathcal K\).  Equality in the displayed
Sauer bound is precisely ampleness of \(S\cup\{p\}\), proving the
equivalence of the first two conditions.  The equivalence with the colon
condition is Lemma~\ref{lem:damp-one-concept-colon-test}.  Under the
shellability--linear-quotients dictionary, the same colon condition says
that \(F_p\cap\Delta_S\) is a union of ridges of \(F_p\), which is the
fourth condition.

Subtracting the two flag \(f\)-identities
\eqref{eq:damp-relative-yang-flag-f} gives
\[
 |\pi_R(B)|-|\pi_R(S)|
 =\sum_{K\subseteq R}h_K(\Psi(B,S)).
\]
Lemma~\ref{lem:damp-relative-yang-cm-flag-h} turns the right-hand side
into the count in
\eqref{eq:damp-relative-projection-defect-count}.

If \(R\in\mathcal H_S(p)\), then \(S\cup\{p\}\) shatters \(R\), so
\(B\) shatters \(R\) and \(R\in\mathcal K\).  The projection
\(\pi_R(S)\) has all patterns except \(p|_R\).  The left-hand side of
\eqref{eq:damp-relative-projection-defect-count} is therefore one, which
is equivalent to the first condition on the right of
\eqref{eq:damp-relative-singleton-defect-supports}.  Conversely, that
condition makes \(p|_R\) the unique pattern of \(\pi_R(B)\) absent from
\(\pi_R(S)\).  Adjoining \(p\) fills it and shatters \(R\), proving the
reverse implication.
\end{proof}

\begin{lemma}[Relative peeling in VC dimension two]
\label{lem:damp-relative-peeling-vc-two}
Let \(\varnothing\ne S\subseteq B\) be ample families.  If \(B\)
has VC dimension at most two, then \(B\) has a corner peeling in which
every vertex of \(B\setminus S\) precedes every vertex of \(S\).
\end{lemma}

\begin{proof}
Every support in
\(\operatorname{Sh}(B)\setminus\operatorname{Sh}(S)\) has size
at most two.  Theorem~\ref{thm:damp-relative-small-supports}
therefore reduces \(B\) to \(S\) without deleting any concept
of \(S\).  Choose \(s\in S\).  The same theorem applied to
\(\{s\}\subseteq S\) reduces \(S\) to \(\{s\}\), since
\(S\) also has VC dimension at most two.  Concatenate these two
sequences and finish by deleting \(s\).  This gives the required
full peeling, and even allows any prescribed last concept in \(S\).
\end{proof}

\begin{lemma}[Nested ample families in three coordinates]
\label{lem:damp-nested-ample-relative-peeling-dimension-three}
Let \(\varnothing\ne S\subseteq B\subseteq Q_E\) be ample, and suppose
that \(B\) uses at most three coordinates.  For every \(s\in S\), the
family \(B\) has a corner peeling in which every vertex of \(B\setminus
S\) precedes every vertex of \(S\), and \(s\) is last.
\end{lemma}

\begin{proof}
It is enough to prove that whenever \(S\subsetneq B\), some corner of
\(B\) lies outside \(S\).  Deleting that corner preserves ampleness by
the ample corner-deletion criterion, so iteration first reduces \(B\) to
\(S\).  Applying the same assertion to the nested pair
\(\{s\}\subseteq S\) then peels \(S\) down to \(s\).

For at most two active coordinates the assertion follows by inspecting
a vertex, an edge, a three-vertex tree, and a square.  Suppose that
\(B\) uses all three coordinates.  Up to a symmetry of \(Q_3\), the
possibilities for \(B\) and its set \(\operatorname{Cor}(B)\) of corners
are
\begin{equation}
\begin{array}{c|c|c}
 |B|&B&\operatorname{Cor}(B)\\ \hline
4&\{000,001,010,100\}&\{001,010,100\}\\
4&\{000,001,011,100\}&\{011,100\}\\
5&\{000,001,010,011,100\}&\{001,010,011,100\}\\
6&\{000,001,010,011,100,101\}&\{010,011,100,101\}\\
7&Q_3\setminus\{111\}&\{011,101,110\}\\
8&Q_3&Q_3.
\end{array}
\label{eq:damp-three-coordinate-ample-types}
\end{equation}
Here is a short completeness check.  Full support shatters the empty set
and the three singleton coordinate sets.  If \(|B|=k<8\), ampleness says
that exactly \(k-4\) coordinate pairs are shattered, and each is strongly
shattered, hence represented by a square in \(B\).  For \(k=4\), the
isometric graph of \(B\) is a three-edge tree, giving the star and the
path.  For \(k=5\), the unique square has one pendant vertex in the third
direction.  For \(k=6\), the two squares share their common-coordinate
edge.  The cases \(k=7,8\) are the punctured and full cube.  These are
exactly the six rows of \eqref{eq:damp-three-coordinate-ample-types}.

An ample family is an isometric subgraph of its cube.  In the first row,
an isometric subfamily containing the three leaves also contains their
centre.  In the second row, the two endpoints force the whole path.  In
the third row, the pendant corner forces its attachment vertex.  In the
fourth row, shortest paths between \(010,100\) and between \(011,101\)
force the two vertices of the common edge.  In the fifth row, shortest
paths among \(011,101,110\) force \(001,010,100\).  The only proper
subfamily that could then contain all three corners is the resulting
six-cycle.  That cycle shatters the empty set, all three singletons, and
all three coordinate pairs, so it has at least seven shattered sets but
only six vertices and is not ample.  In the last row every vertex is a
corner.  Thus no proper ample \(S\subsetneq B\) contains every corner of
\(B\), as required.
\end{proof}

\begin{theorem}[Nested ample families in four coordinates]
\label{thm:damp-nested-ample-relative-peeling-dimension-four}
Let \(\varnothing\ne S\subseteq B\subseteq Q_E\) be ample, and suppose
that \(B\) uses at most four coordinates.  For every \(s\in S\), the
family \(B\) has a corner peeling in which every vertex of \(B\setminus
S\) precedes every vertex of \(S\), and \(s\) is last.
\end{theorem}

\begin{proof}
As in Lemma
\ref{lem:damp-nested-ample-relative-peeling-dimension-three}, it is
enough to show that a proper ample subclass \(S\subsetneq B\) omits a
corner of \(B\).  We may assume that \(B\) uses all four coordinates.
If its VC dimension is at most two, this follows from Lemma
\ref{lem:damp-relative-peeling-vc-two}.  If its VC dimension is four,
then \(B=Q_4\) and every vertex is a corner.  It remains to consider VC
dimension three.

Choose a strongly shattered triple and orient the remaining coordinate
so that its cube is in the lower layer.  Then
\begin{equation}
B=(Q_3\times\{0\})\cup(U\times\{1\}),
 \label{eq:damp-four-coordinate-vc-three-form}
\end{equation}
where \(U\) is a nonempty proper ample subfamily of \(Q_3\).  It is
nonempty because \(B\) uses the fourth coordinate, and it is proper
because \(B\) has VC dimension three.  The maximal cubes in
\eqref{eq:damp-four-coordinate-vc-three-form} give
\begin{equation}
\operatorname{Cor}(B)
 =
 ((Q_3\setminus U)\times\{0\})
 \mathbin{\dot\cup}
 (\operatorname{Cor}(U)\times\{1\}).
 \label{eq:damp-four-coordinate-corners}
\end{equation}

The possible sites \(U\), up to cube symmetry, are obtained from the
support-zero, support-one, and support-two cases together with the six
support-three types in
\eqref{eq:damp-three-coordinate-ample-types}.  The corner formula gives
the following short table; the last column is obtained by projecting the
displayed corner set.
\begin{equation}
\begin{array}{c|c|c|c}
 U&|U|&|\operatorname{Cor}(B)|
   &|\operatorname{Sh}(\operatorname{Cor}(B))|\\ \hline
 K_1&1&8&9\\
 K_2&2&8&10\\
 P_3&3&7&10\\
 Q_2&4&8&12\\
 K_{1,3}&4&7&11\\
 P_4&4&6&10\\
 Q_2\text{ with one pendant vertex}&5&7&10\\
 \text{two squares sharing an edge}&6&6&8\\
 Q_3^-&7&4&5
\end{array}
\label{eq:damp-four-coordinate-site-table}
\end{equation}

Put \(C=\operatorname{Cor}(B)\), and suppose first that the last entry of
the table is not in force.  If \(C\subseteq S\), then ampleness and
monotonicity of shattering give
\begin{equation}
|S|=|\operatorname{Sh}(S)|
 \ge |\operatorname{Sh}(C)|\ge8.
 \label{eq:damp-four-coordinate-corner-count}
\end{equation}
Let \(A=Q_4\setminus B\) and \(T=Q_4\setminus S\).  Complements of ample
cube families are ample \cite{Lawrence83,BandeltCDK06}, so
\(A\subsetneq T\) are ample, and
\eqref{eq:damp-four-coordinate-corner-count} gives \(|T|\le8\).
If \(T\) uses at most three coordinates, Lemma
\ref{lem:damp-nested-ample-relative-peeling-dimension-three} applies.
If it uses all four, then it cannot have VC dimension three: a shattered
triple and the remaining active singleton already give at least nine
shattered sets.  Lemma~\ref{lem:damp-relative-peeling-vc-two} therefore
applies instead.

In either case, peel \(T\setminus A\) down to \(A\), and let \(x\) be the
last vertex removed before \(A\) remains.  Then \(A\cup\{x\}\) is ample.
Its complement is \(B\setminus\{x\}\), which is consequently ample.  The
ample corner-deletion criterion makes \(x\) a corner of \(B\).  But
\[
 x\in T\setminus A=B\setminus S,
\]
contradicting \(C\subseteq S\).

It remains to treat \(U=Q_3^-\), for which \(B=Q_4^-\).  Normalize the
missing vertex to \(1111\).  The four corners of \(B\) are its
weight-three vertices.  If an isometric subfamily contains them, then a
shortest path between each pair cannot use the missing vertex and forces
all six weight-two vertices.  Let \(K\) be the resulting ten-vertex
family of all weight-two and weight-three vertices.  It shatters the
empty set, the four singletons, and the six coordinate pairs.

Suppose that an ample \(S\subseteq B\) contains \(K\).  If it also
contains the zero vertex, then every coordinate triple is shattered, so
\(|\operatorname{Sh}(S)|\ge15\); a proper subfamily of \(B\) has order at
most fourteen, a contradiction.  Otherwise let \(k\) be the number of
weight-one vertices in \(S\).  Each such vertex shatters the unique
coordinate triple avoiding its support, and these triples are distinct.
Hence
\[
 |\operatorname{Sh}(S)|\ge11+k>|K|+k=|S|,
\]
again contradicting ampleness.  Thus no proper ample subclass of
\(Q_4^-\) contains all four corners.

We have proved in every case that a proper ample \(S\subsetneq B\)
omits a corner of \(B\).  Delete such a corner and iterate; corner
deletion preserves ampleness.  This peels \(B\setminus S\) first.
Applying the same assertion to \(\{s\}\subseteq S\) completes the
peeling with \(s\) last.
\end{proof}

\begin{lemma}[Augmentation across an expansion separator]
\label{lem:damp-expansion-separator-augmentation}
Let
\[
 H=(C\times\{0\})\cup(Y\times\{1\})
 \subseteq Q_{E\cup\{r\}}
\]
be ample, put \(S=C\cap Y\), and let \(p\in Y\setminus C\).  If
\(S\cup\{p\}\) is ample, then \(C\cup\{p\}\) is ample.
\end{lemma}

\begin{proof}
Put \(Z=C\cup Y\).  This is the contraction of \(H\) in coordinate
\(r\), so it is ample and hence isometric.  Moreover, no edge of \(Z\)
joins \(C\setminus S\) to \(Y\setminus S\).  Indeed, if
\(c\in C\setminus S\) and \(y\in Y\setminus S\) were adjacent, then
\((c,0)\) and \((y,1)\) would have Hamming distance two in \(H\).  The
only two possible intermediate cube vertices are \((y,0)\) and
\((c,1)\), neither of which belongs to \(H\), contradicting isometry.

We verify item~3 of the one-concept augmentation criterion,
Lemma~\ref{lem:damp-one-concept-colon-test}.  Fix \(c\in C\).  If
\(c\in S\), that criterion applied to \(S\cup\{p\}\) gives a coordinate
\(e\in\delta(p,c)\) for which \(p^{\{e\}}\in S\subseteq C\).
Suppose instead that \(c\in C\setminus S\).  Choose a \(Z\)-geodesic
from \(p\) to \(c\).  The preceding separation forces this geodesic to
meet \(S\); let \(s\) be one of its vertices in \(S\).  Since \(s\)
lies on a cube geodesic from \(p\) to \(c\),
\[
 \delta(p,s)\subseteq\delta(p,c).
\]
Applying the augmentation criterion to \(s\in S\) gives
\(e\in\delta(p,s)\) such that
\(p^{\{e\}}\in S\subseteq C\).  Thus every \(c\in C\) satisfies the
criterion, and \(C\cup\{p\}\) is ample.
\end{proof}

\begin{lemma}[Duplicating an augmenting concept across an expansion]
\label{lem:damp-expansion-separator-duplication}
Let
\[
 H=(C\times\{0\})\cup(Y\times\{1\})
 \subseteq Q_{E\cup\{r\}}
\]
be ample, put \(I=C\cap Y\), and let \(p\in Y\setminus C\).  If
\(I\cup\{p\}\) is ample, then
\begin{equation}
H\cup\{(p,0)\}
 =((C\cup\{p\})\times\{0\})\cup(Y\times\{1\})
 \label{eq:damp-expansion-duplicated-augmentation}
\end{equation}
is ample.
\end{lemma}

\begin{proof}
Put \(Z=C\cup Y\).  The two sections \(C,Y\) and the contraction \(Z\)
are ample.  The two-section shattering formula and ampleness of \(H\)
give
\begin{equation}
\operatorname{Sh}(C)\cap\operatorname{Sh}(Y)
 =\operatorname{Sh}(I),
 \label{eq:damp-expansion-section-shattering-intersection}
\end{equation}
and also show that \(I\) is ample.  For completeness, the count is
\[
 |C|+|Y|=|Z|+|I|
 =|\operatorname{Sh}(Z)|
  +|\operatorname{Sh}(C)\cap\operatorname{Sh}(Y)|.
\]
Since \(\operatorname{Sh}(I)\) is contained in the intersection and
Sauer's inequality gives \(|\operatorname{Sh}(I)|\ge|I|\), equality is
forced throughout.

The ample singleton augmentation \(I\cup\{p\}\) adds one shattered
support, say \(R\).  Because it is contained in \(Y\), we have
\(R\in\operatorname{Sh}(Y)\).  Equation
\eqref{eq:damp-expansion-section-shattering-intersection} gives
\(R\notin\operatorname{Sh}(C)\).  Lemma
\ref{lem:damp-expansion-separator-augmentation} shows that
\(C\cup\{p\}\) is ample.  Its unique new shattered support must also be
\(R\), by monotonicity from \(I\cup\{p\}\subseteq C\cup\{p\}\).

In \eqref{eq:damp-expansion-duplicated-augmentation}, the union of the
two sections remains \(Z\), whereas their intersection changes from
\(I\) to \(I\cup\{p\}\).  The two-section formula therefore adds exactly
the support \(R\cup\{r\}\) and no other shattered support.  The family
has gained one concept, so it is ample.
\end{proof}

\begin{lemma}[Section synchronization or a nine-concept cycle]
\label{lem:damp-section-synchronization-or-cycle}
Let
\[
 H=(C\times\{0\})\cup(Y\times\{1\})
 \subseteq Q_{E\cup\{r\}}
\]
be ample, put \(S=C\cap Y\), and suppose that
\[
 \varnothing\ne S\subsetneq Y.
\]
Then at least one of the following holds:
\begin{enumerate}
\item some \(p\in Y\setminus S\) satisfies
\begin{equation}
 S\cup\{p\}\text{ is ample},
 \qquad
 C\cup\{p\}\text{ is ample}.
 \label{eq:damp-synchronized-augmentations}
\end{equation}
For every such \(p\), both
\begin{equation}
 (C\times\{0\})
  \cup((S\cup\{p\})\times\{1\})
 \quad\text{and}\quad
 H\cup\{(p,0)\}
 \label{eq:damp-synchronized-lifts}
\end{equation}
are ample;
\item put \(D=Y\setminus S\), and let \(c(D)\) be the number of connected
      components of the graph induced by \(D\).  Then \(|D|\ge9\), every
      component of that graph contains a cycle, and
      \begin{equation}
       \sum_{T\in\operatorname{Sh}(Y)\setminus\operatorname{Sh}(S)}|T|
       \ge |D|+c(D)\ge10.
       \label{eq:damp-unsynchronized-section-cycle-charge}
      \end{equation}
\end{enumerate}
\end{lemma}

\begin{proof}
The sections \(Y\) and \(S\) are ample by
Lemma~\ref{lem:damp-ample-two-layer-shadow-intersection}.  Suppose first
that some \(p\in Y\setminus S\) makes \(S\cup\{p\}\) ample.  Lemma
\ref{lem:damp-expansion-separator-augmentation} then makes
\(C\cup\{p\}\) ample, proving \eqref{eq:damp-synchronized-augmentations}.

For the first family in
\eqref{eq:damp-synchronized-lifts}, the two-layer trace
formula gives the supports not containing \(r\) as
\(\operatorname{Sh}(C\cup\{p\})\).  The supports containing \(r\) are
indexed by
\[
 \operatorname{Sh}(C)\cap
 \operatorname{Sh}(S\cup\{p\}).
\]
The original ample family \(H\) satisfies
\[
 \operatorname{Sh}(C)\cap\operatorname{Sh}(Y)
 =\operatorname{Sh}(S)
\]
by Lemma~\ref{lem:damp-ample-two-layer-shadow-intersection}.  Since
\(S\subseteq S\cup\{p\}\subseteq Y\), the last intersection is exactly
\(\operatorname{Sh}(S)\).  The first family therefore has
\[
 |\operatorname{Sh}(C\cup\{p\})|+|\operatorname{Sh}(S)|
 =|C|+1+|S|
\]
shattered supports, equal to its number of concepts.  The upper Sandwich
equality proves ampleness.  The second family in
\eqref{eq:damp-synchronized-lifts} is ample by
Lemma~\ref{lem:damp-expansion-separator-duplication}.

It remains to suppose that no \(p\in Y\setminus S\) makes
\(S\cup\{p\}\) ample.  The relative Yang complex of the nested ample pair
\(S\subsetneq Y\) is Cohen--Macaulay by the relative
Cohen--Macaulay criterion of \cite{BousquetYangComplex26}.  The
arbitrary-upper linear-quotient theorem through eight relative facets
therefore gives \(|Y\setminus S|\ge9\).  Put \(D=Y\setminus S\).

The linear presentation of the Alexander-dual relative module splits over
the connected components of the graph induced by \(D\); each component
summand consequently has a linear resolution.  If one component were a
forest, the relative forest theorem of the same paper, applied to that
summand, would supply a concept satisfying the \(S\)-teaching condition.
Adjoining that concept to \(S\) would be ample, contrary to the assumption.
Thus every component contains a cycle, and hence the graph induced by
\(D\) has at least \(|D|\) edges.

Every component of \(D\) has an edge to \(S\), because the ample family
\(Y\) is connected and \(S\ne\varnothing\).  The edge--shadow identity
for the two nested ample families now gives
\[
 \sum_{T\in\operatorname{Sh}(Y)\setminus\operatorname{Sh}(S)}|T|
 =|E(Y)|-|E(S)|
 =|E(D)|+|E(D,S)|
 \ge |D|+c(D).
\]
Together with \(|D|\ge9\), this proves
\eqref{eq:damp-unsynchronized-section-cycle-charge}.
\end{proof}

\begin{theorem}[Complete section synchronization or a terminal cyclic remainder]
\label{thm:damp-complete-section-synchronization-or-cycle}
Keep the notation of Lemma~\ref{lem:damp-section-synchronization-or-cycle}.
There are an ample family \(Z\) with
\[
 S\subseteq Z\subseteq Y
\]
and an ordering \(p_1,\ldots,p_q\) of \(Z\setminus S\) such that, on
putting \(Z_h=S\cup\{p_1,\ldots,p_h\}\), every \(Z_h\) and every family
\begin{equation}
 H_h=((C\cup Z_h)\times\{0\})\cup(Y\times\{1\})
 \label{eq:damp-progressive-section-synchronization}
\end{equation}
is ample.  Moreover, exactly one of the following conclusions holds:
\begin{enumerate}
 \item \(Z=Y\), so all concepts of \(Y\setminus S\) may be duplicated
       across the \(r\)-edge, one at a time, while preserving ampleness;
 \item put \(D=Y\setminus Z\).  Then \(|D|\ge9\), every component of the
       graph induced by \(D\) contains a cycle, and
       \begin{equation}
        \sum_{T\in\operatorname{Sh}(Y)
                    \setminus\operatorname{Sh}(Z)}|T|
        \ge |D|+c(D)\ge10.
        \label{eq:damp-terminal-synchronization-cycle-charge}
       \end{equation}
\end{enumerate}
\end{theorem}

\begin{proof}
Starting with \(Z_0=S\), repeatedly choose
\(p_{h+1}\in Y\setminus Z_h\) for which
\(Z_h\cup\{p_{h+1}\}\) is ample, and stop when no such concept exists.
Every intermediate \(Z_h\) is ample by construction.  The base family in
\eqref{eq:damp-progressive-section-synchronization} is \(H\).  If \(H_h\)
is ample, its two sections are \(C\cup Z_h\) and \(Y\), with intersection
\(Z_h\).  Lemma~\ref{lem:damp-expansion-separator-duplication} shows that
adjoining the lower copy of \(p_{h+1}\) gives the ample family \(H_{h+1}\).
This proves the progressive synchronization assertion.

If the process reaches \(Y\), the first alternative holds.  Otherwise
apply Lemma~\ref{lem:damp-section-synchronization-or-cycle} to the ample
family \(H_q\), whose section intersection is \(Z=Z_q\).  Maximality of
the construction excludes the synchronization alternative of that lemma.
Its cyclic alternative gives all assertions in the second conclusion.
\end{proof}

\begin{corollary}[An adjacent path transition synchronizes or owns a terminal cycle]
\label{cor:damp-adjacent-section-extension-or-cycle}
Keep the path-spine notation of Proposition
\ref{prop:damp-path-spine-euler-identity}.  For every \(1\le i\le n\)
and each \(j\in\{i-1,i\}\), at least one of the following holds:
\begin{enumerate}
 \item writing \(k\) for the other member of \(\{i-1,i\}\), all concepts
       of \(P_j\setminus C_i\) can be ordered so that duplicating their
       missing copies into \(P_k\), in that order, preserves ampleness at
       every step;
 \item after such duplications have synchronized a proper ample
       intermediate family \(Z\), with \(C_i\subseteq Z\subsetneq P_j\),
       the remainder \(D=P_j\setminus Z\) has at least nine concepts,
       every component of its induced graph contains a cycle, and
       \begin{equation}
        \sum_{T\in\operatorname{Sh}(P_j)
                    \setminus\operatorname{Sh}(Z)}|T|
        \ge |D|+c(D)\ge10.
        \label{eq:damp-adjacent-terminal-cycle-charge}
       \end{equation}
\end{enumerate}
\end{corollary}

\begin{proof}
Condition all core coordinates except \(x_i\) to the values shared by
\(a_{i-1}\) and \(a_i\).  The resulting ample two-layer family has
sections \(P_{i-1},P_i\) and overlap \(C_i\).  Apply
Theorem~\ref{thm:damp-complete-section-synchronization-or-cycle} in the
chosen orientation.
\end{proof}

\begin{corollary}[A noncyclic transition exposes an exclusive section corner]
\label{cor:damp-adjacent-transition-corner-or-cycle}
For every oriented incidence \((i,j)\) in the path spine for which
\(P_j\setminus C_i\ne\varnothing\), either the
terminal cyclic alternative of Corollary
\ref{cor:damp-adjacent-section-extension-or-cycle} occurs, or there is a
concept
\begin{equation}
 p_{i,j}\in P_j\setminus C_i
 \label{eq:damp-adjacent-exclusive-section-corner}
\end{equation}
that is a corner of the ample section \(P_j\).
\end{corollary}

\begin{proof}
In the complete-synchronization alternative of Theorem
\ref{thm:damp-complete-section-synchronization-or-cycle}, let \(p_{i,j}\)
be the last concept in the chain from \(C_i\) to \(P_j\).  Deleting it
from \(P_j\) leaves the preceding ample member of the chain.  The ample
corner-deletion criterion therefore makes \(p_{i,j}\) a corner of
\(P_j\), and it lies outside \(C_i\).
\end{proof}

\begin{proposition}[Acquired section shadows obey path-boundary conservation]
\label{prop:damp-section-shadow-boundary-conservation}
Keep the path-spine notation of Proposition
\ref{prop:damp-path-spine-euler-identity}, and put
\[
 \mathcal S=\bigcup_{j=0}^n\operatorname{Sh}(P_j).
\]
For \(T\in\mathcal S\), let \(c_{\rm sh}(T)\) be the number of connected
components of
\[
 I_{\rm sh}(T)=\{j:T\in\operatorname{Sh}(P_j)\}
\]
in the section path.  Then
\begin{align}
 &\sum_{i=1}^n\bigl(
   |P_{i-1}\setminus C_i|+|P_i\setminus C_i|\bigr)
 \notag\\
 &\qquad\le
 2\sum_{T\in\mathcal S}c_{\rm sh}(T)
 =2(|U|+\Phi)
 \le2(|W|+|\Omega|).
 \label{eq:damp-section-shadow-boundary-conservation}
\end{align}
Thus all oriented section wings are controlled at once: repeated acquisition
of one repair support is possible only through a new interval of section
indices, and the path-spine Euler term pays for every additional interval.
\end{proposition}

\begin{proof}
For every path edge, the ample two-layer shadow-intersection identity gives
\[
 \operatorname{Sh}(C_i)
 =\operatorname{Sh}(P_{i-1})\cap\operatorname{Sh}(P_i).
\]
Since all three families are ample, the left side of
\eqref{eq:damp-section-shadow-boundary-conservation} is therefore
\[
 \sum_{i=1}^n\bigl(
  |\operatorname{Sh}(P_{i-1})\setminus\operatorname{Sh}(C_i)|
 +|\operatorname{Sh}(P_i)\setminus\operatorname{Sh}(C_i)|\bigr).
\]
Fix \(T\in\mathcal S\).  Its contribution to this sum is the number of
path edges having exactly one endpoint in \(I_{\rm sh}(T)\).  A union of
\(c_{\rm sh}(T)\) nonempty path intervals has at most
\(2c_{\rm sh}(T)\) boundary edges.  Summing proves the first inequality.

Counting the vertices and internal edges of every set
\(I_{\rm sh}(T)\) gives
\begin{align*}
 \sum_{T\in\mathcal S}c_{\rm sh}(T)
 &=\sum_{j=0}^n|\operatorname{Sh}(P_j)|
   -\sum_{i=1}^n|\operatorname{Sh}(C_i)|\\
 &=\sum_{j=0}^n|P_j|-\sum_{i=1}^n|C_i|
  =|U|+\Phi,
\end{align*}
where the last equality is the word-level interval count in the proof of
Proposition~\ref{prop:damp-path-spine-euler-identity}.  Finally, that
proposition gives the exact cancellation
\[
 |W|+|\Omega|=|U|+\Phi+V,
\]
and \(V\ge0\).  This proves the final inequality.
\end{proof}

\begin{proposition}[Boundary-corner roots are distinct up to fragmentation]
\label{prop:damp-boundary-corner-root-fragmentation}
Let \(\mathcal A\) be any set of oriented incidences for which the corner
alternative in Corollary
\ref{cor:damp-adjacent-transition-corner-or-cycle} holds, and choose one
corner \(p_{i,j}\) for each \((i,j)\in\mathcal A\).  Let \(Z\) be the set
of distinct chosen repair words.  Then
\begin{equation}
 |\mathcal A|\le2(|Z|+\Phi),
 \qquad
 |Z|\ge\frac{|\mathcal A|}{2}-\Phi.
 \label{eq:damp-boundary-corner-root-fragmentation}
\end{equation}
Thus repeated use of one repair word is already paid for by the path-spine
fragmentation term.
\end{proposition}

\begin{proof}
Fix \(p\in Z\).  An incidence selecting \(p\) is an oriented boundary edge
of the subgraph induced by
\[
 I(p)=\{h:p\in P_h\}
\]
in the core path.  A union of \(c(p)\) nonempty path intervals has at most
\(2c(p)\) such boundary incidences.  Hence
\[
 |\mathcal A|
 \le2\sum_{p\in Z}c(p)
 =2|Z|+2\sum_{p\in Z}(c(p)-1)
 \le2(|Z|+\Phi),
\]
which proves both displayed inequalities.
\end{proof}

\begin{proposition}[The only neutral boundary-corner defect is an inward square tip]
\label{prop:damp-boundary-corner-neutral-square-tip}
Fix \(1\le i\le n\), let \(j\in\{i-1,i\}\), and choose any
\[
 p\in\operatorname{Cor}(P_j)\setminus C_i.
\]
Write \(v=(p,a_j)\).  Since \(H\) is cornerless,
Proposition~\ref{prop:damp-column-corner-core-transverse-puncture} gives a
rooted puncture \((v,J)\).  Put
\[
 S=J\cap R,
 \qquad K=J\cap X.
\]
Then \(K\ne\varnothing\) and \(x_i\notin K\).  Moreover, at least one of
the following holds:
\begin{enumerate}
 \item the puncture contains a concept of \(H\) over a core word outside
       \(A\), with repair root determined by \(p\);
 \item \(|K|\ge2\), \(S\ne\varnothing\), and
       \(K\in\Omega_{\ge2}\);
 \item \(K=\{x_\ell\}\), where \(x_\ell\) is the other path direction
       incident with \(a_j\).  For every \(s\in S\), the puncture has a
       provenance-tagged terminal square on \(\{s,x_\ell\}\) whose repair
       root belongs to \(C_\ell\).
\end{enumerate}
In particular, after charging external-core concepts and residual supports,
all uncharged boundary-corner occurrences reduce to inward, spine-aligned
repair--core square tips.  No higher-rank neutral defect remains.
\end{proposition}

\begin{proof}
The transition direction \(x_i\) is not incident with \(v\), because
\(p\notin C_i\).  Since \(J\) consists of directions incident with \(v\),
we have \(x_i\notin K\); core transversality gives \(K\ne\varnothing\).

Suppose first that \(S=\varnothing\).  Then \(|K|=|J|\ge2\), and every
single flip \(v^{\{x\}}\), \(x\in K\), belongs to \(H\).  At most one
direction in \(K\) is the other core-path edge at \(a_j\).  Some such
neighbor therefore lies over a core word outside \(A\), proving the first
alternative.

Suppose next that \(S\ne\varnothing\) and \(|K|\ge2\).  The set \(K\) is
a proper subset of \(J\), so the puncture relations contain the full
\(K\)-cube through \(v\).  Hence \(H\) strongly shatters \(K\).  This
support has at least two core coordinates and therefore belongs to
\(\Omega_{\ge2}\), proving the second alternative.

It remains that \(K=\{x_\ell\}\).  If \(a_j^{\{x_\ell\}}\notin A\), the
present neighbor \(v^{\{x_\ell\}}\) proves the first alternative.  Otherwise
\(x_\ell\) is a path direction incident with \(a_j\).  It differs from the
absent transition direction \(x_i\), so it is the other incident path edge.
In particular, the repair word \(p\) belongs to both path sections across
that edge and hence to \(C_\ell\).

For \(s\in S\), apply the square-tip reduction of Lemma
\ref{lem:damp-shadow-neutral-square-tip-reduction} to
\(L=\{s,x_\ell\}\).  Its repair root is
\(p^{S\setminus\{s\}}\).  Both vertices in the \(x_\ell\)-direction are
proper flips of \(J\), so this repair root also belongs to \(C_\ell\).
The retained parent puncture supplies the asserted provenance tag.
\end{proof}

\begin{proposition}[Closed neutral transfers localize in one repair cube]
\label{prop:damp-neutral-transfer-two-sided-cycle}
For every occurrence of the third alternative in Proposition
\ref{prop:damp-boundary-corner-neutral-square-tip}, retain the parent
puncture and write its support as
\[
 J=S\mathbin{\dot\cup}\{x_\ell\}.
\]
Then \(S\ne\varnothing\), and the opposite repair tip
\begin{equation}
 q=p^{S}
 \label{eq:damp-neutral-transfer-opposite-tip}
\end{equation}
belongs to \(P_j\setminus C_\ell\).  Thus the neutral defect transfers the
chosen boundary occurrence at one side of \(P_j\) to a boundary occurrence
at its other side.

For every interior section \(P_j\), make disjoint copies of its two sets of
boundary corners:
\[
 \mathcal L_j=
 \{(L,p):p\in\operatorname{Cor}(P_j)\cap(P_j\setminus C_j)\},
 \qquad
 \mathcal R_j=
 \{(R,p):p\in\operatorname{Cor}(P_j)\cap(P_j\setminus C_{j+1})\}.
\]
For each retained neutral parent puncture rooted at one of these occurrences,
draw an arrow to the opposite occurrence of the word in
\eqref{eq:damp-neutral-transfer-opposite-tip}, provided that this opposite
word is again a corner with a retained neutral parent puncture.  Every
directed cycle in the resulting neutral-transfer digraph stays in one
\(\mathcal L_j\mathbin{\dot\cup}\mathcal R_j\), has even length, and all
of its repair words lie in one positioned maximal cube \(Q\) of \(P_j\).
If \(T\) is the support of \(Q\) and \(t=|T|\), then \(t\ge2\), and
\begin{equation}
 Q\nsubseteq P_{j-1},
 \qquad
 Q\nsubseteq P_{j+1}.
 \label{eq:damp-neutral-cycle-isolated-cube}
\end{equation}

Choose any arrow in each orientation on such a cycle, and let \(S_+\) be
the repair label of the arrow from \(\mathcal L_j\) to \(\mathcal R_j\)
and \(S_-\) the label of the arrow in the reverse direction.  Put
\(k_+=|S_+|\) and \(k_-=|S_-|\).  Then
\begin{equation}
 \Delta(H;A,R)\ge
 2^t-t-1
 +\max\{0,2^{k_+}-k_+-2\}
 +\max\{0,2^{k_-}-k_--2\}.
 \label{eq:damp-neutral-cycle-shadow-tax}
\end{equation}

A directed two-cycle has a stronger symmetric description.  More precisely,
if \(p\in P_j\setminus C_j\) and
\(q\in P_j\setminus C_{j+1}\) form such a two-cycle, then for one nonempty
\(S\subseteq R\),
\begin{equation}
 q=p^S,\qquad p=q^S,
 \label{eq:damp-neutral-two-cycle-antipodes}
\end{equation}
and
\begin{align}
 &\{p^U:U\subseteq S\}\subseteq P_j,
 \notag\\
 &\{p^U:U\subsetneq S\}\subseteq C_{j+1},
   \qquad p^S\notin C_{j+1},
 \notag\\
 &\{q^U:U\subsetneq S\}\subseteq C_j,
   \qquad q^S\notin C_j.
 \label{eq:damp-neutral-two-cycle-opposite-punctures}
\end{align}
The two chosen concepts \(p,q\) are opposite vertices of one face in their
common positioned maximal cube \(Q\).

If \(k=|S|\), this local cycle gives the path excess bound
\begin{equation}
 \Delta(H;A,R)\ge
 2^t-t-1+2\max\{0,2^k-k-2\}.
 \label{eq:damp-neutral-two-cycle-shadow-tax}
\end{equation}
For \(k\ge2\), this is at least
\[
 3\cdot2^k-3k-5.
\]
Consequently a neutral two-cycle in a family of order below \(46\) has
\(k\le3\).  In the rank-one case the two oppositely oriented punctured
squares on \(\{s,x_j\}\) and \(\{s,x_{j+1}\}\) still pay at least one
non-singleton repair support.
\end{proposition}

\begin{proof}
In the neutral alternative one has
\(J=S\mathbin{\dot\cup}\{x_\ell\}\) and \(|J|\ge2\), so
\(S\ne\varnothing\).  The parent puncture contains the vertex obtained by
the proper flip \(S\), namely \((p^S,a_j)\), but omits its
\(x_\ell\)-neighbour.  This is exactly the assertion that
\(q=p^S\) lies in \(P_j\setminus C_\ell\).

Every arrow stays at the same section and exchanges its two incident path
edges.  Thus the neutral-transfer digraph is the disjoint union of the
bipartite digraphs on
\(\mathcal L_j\mathbin{\dot\cup}\mathcal R_j\), and every directed cycle
is even.

Consider an arrow from a corner \(p\) to a corner \(q=p^S\).  The parent
puncture supplies the full \(S\)-face through \(p\) in \(P_j\).  This face
is contained in the unique maximal cube through \(p\), and that maximal cube
also contains \(q\).  Since \(q\) is a corner, the same cube is its unique
maximal cube.  Following the arrows around a cycle therefore shows that all
of its corners have one common positioned maximal cube \(Q\).  The cycle
contains a vertex of \(\mathcal L_j\) and a vertex of \(\mathcal R_j\), so
\(Q\) contains a word absent from \(P_{j-1}\) and a word absent from
\(P_{j+1}\).  This proves
\eqref{eq:damp-neutral-cycle-isolated-cube}.

If \(t=1\), the two endpoints of \(Q\) have no repair neighbor outside
\(Q\): such a neighbor would lie in a second maximal cube through that
endpoint.  Connectedness would then force \(P_j=Q\), which is impossible
because \(P_j\) contains two repair antipodes and \(r\ge3\).  Hence
\(t\ge2\).

The projection \(W\) contains the full \(T\)-cube \(Q\), so beyond the
empty support and the \(r\) repair singletons it has at least
\(2^t-t-1\) supports contained in \(T\).  An arrow from
\(\mathcal L_j\) to \(\mathcal R_j\), with label \(S_+\), gives a
punctured \(S_+\)-face in \(C_{j+1}\); an arrow in the reverse orientation,
with label \(S_-\), gives a punctured \(S_-\)-face in \(C_j\).  Each such
carrier strongly shatters every proper subset of its arrow label and hence
has at least
\(\max\{0,2^{k_\pm}-k_\pm-2\}\) supports beyond its empty support and
repair singletons.  The three contributions lie in different blocks of the
disjoint path-spine shadow family.  This proves
\eqref{eq:damp-neutral-cycle-shadow-tax}.

Suppose now that the cycle has two arrows.  Write \(S_p\) and \(S_q\) for
their repair labels.  The arrows give \(q=p^{S_p}\) and \(p=q^{S_q}\),
and therefore
\(S_p=S_q=\operatorname{Diff}(p,q)=S\).  The proper-flip relations of the
first parent puncture give the full \(S\)-face in \(P_j\) and the punctured
\(S\)-face rooted at \(p\) in \(C_{j+1}\).  The second parent gives the
oppositely punctured face rooted at \(q\) in \(C_j\).  These are precisely
the inclusions and omissions in
\eqref{eq:damp-neutral-two-cycle-opposite-punctures}.  Substituting
\(S_+=S_-=S\) in \eqref{eq:damp-neutral-cycle-shadow-tax} gives
\eqref{eq:damp-neutral-two-cycle-shadow-tax}.  Since \(t\ge k\), its
right side is at least \(3\cdot2^k-3k-5\) when \(k\ge2\).  If \(k\ge4\), then
\(r\ge k\), \(n\ge2\), and
\[
 |H|\ge(n+1)(r+1)+3\cdot2^k-3k-5\ge46,
\]
with equality in the last lower bound at \(k=4\).  This proves the final
assertions.
\end{proof}

\begin{proposition}[Section excess pays the whole neutral-transfer forest]
\label{prop:damp-neutral-transfer-forest-section-excess}
Fix an interior path section \(P_j\).  Extend its neutral-transfer digraph
as follows.  For every retained neutral boundary-corner occurrence, draw the
arrow to its opposite repair tip \(q=p^S\), whether or not \(q\) is again a
retained neutral corner.  If it is not, adjoin \(q\) as a terminal vertex;
terminal vertices with the same repair word are identified.  Let
\(\kappa_j\) be the number of weak components of this extended digraph.
Then
\begin{equation}
 \kappa_j\le |P_j|-(r+1).
 \label{eq:damp-neutral-transfer-forest-section-excess}
\end{equation}
Consequently, over all interior sections,
\begin{equation}
 \sum_{j=1}^{n-1}\kappa_j
 \le\sum_{j=1}^{n-1}\bigl(|P_j|-(r+1)\bigr)
 \le\Delta(H;A,R).
 \label{eq:damp-neutral-transfer-forest-global-excess}
\end{equation}
Thus one additive budget pays every closed component and every acyclic
terminal component simultaneously.  The bound is independent of the
repair labels, the ranks of terminal maximal cubes, and any reuse of a
terminal support by several arrows in the same component.
\end{proposition}

\begin{proof}
Every nonterminal vertex has out-degree one.  Hence each weak component
contains either one directed cycle or one terminal vertex.  For a weak
component \(D\), let \(W(D)\) be the set of underlying repair words
represented by its occurrence and terminal vertices.  Every component
contains an arrow \(p\longrightarrow p^S\), where \(S\ne\varnothing\),
and therefore
\begin{equation}
 |W(D)|\ge2.
 \label{eq:damp-neutral-component-two-words}
\end{equation}

Choose antipodal repair geodesics
\[
 L\subseteq C_j,
 \qquad
 R'\subseteq C_{j+1},
\]
and let \(\nu\) be the number of vertices exclusive to either one of these
two geodesics.  Both have \(r+1\) vertices, so
\begin{equation}
 |L\cup R'|=r+1+\nu.
 \label{eq:damp-neutral-transfer-two-geodesic-union}
\end{equation}

Consider an arrow whose source lies in \(\mathcal L_j\).  The neutral
parent puncture uses the other path direction \(x_{j+1}\), whereas the
source is absent across \(x_j\).  Its source and target therefore satisfy
\[
 p\in C_{j+1}\setminus C_j,
 \qquad
 q\notin C_{j+1}.
\]
If one of its two words belongs to \(L\cup R'\), that word belongs to the
symmetric difference of the two geodesics.  More precisely, when both
words belong to the union,
\[
 p\in R'\setminus L,
 \qquad
 q\in L\setminus R'.
\]
The reverse orientation gives the symmetric conclusion.  In particular,
no word of \(L\cap R'\) occurs in the neutral digraph.

We count incidences between weak components and their underlying repair
words.  A word outside \(L\cup R'\) belongs to at most two weak
components, one through each of its left and right boundary occurrences.
If it is adjoined as a terminal word, the occurrence on the target side is
absent; all terminal copies of the word are identified, so the same bound
still holds.  A word in \(L\mathbin\triangle R'\) has only the boundary
occurrence on the side of the geodesic which omits it and hence belongs to
at most one weak component.  Therefore
\begin{align*}
 2\kappa_j
 &\le \sum_D|W(D)|\\
 &\le 2|P_j\setminus(L\cup R')|
       +|L\mathbin\triangle R'|\\
 &=2\bigl(|P_j|-|L\cup R'|\bigr)+2\nu\\
 &=2\bigl(|P_j|-(r+1)\bigr),
\end{align*}
where the last equality uses
\eqref{eq:damp-neutral-transfer-two-geodesic-union}.  This proves
\eqref{eq:damp-neutral-transfer-forest-section-excess} while retaining
the possible left--right duplication of one repair word.

Finally, the path-section excess partition
\eqref{eq:damp-path-section-excess-partition} gives
\[
 \Delta(H;A,R)
 =V+\sum_{h=0}^n\bigl(|P_h|-(r+1)\bigr).
\]
All terms are nonnegative, so summing the sectionwise bound over the
interior indices proves
\eqref{eq:damp-neutral-transfer-forest-global-excess}.
\end{proof}

\begin{corollary}[Equality in the neutral-forest budget is two-word]
\label{cor:damp-neutral-forest-equality-two-word}
If equality holds in
\eqref{eq:damp-neutral-transfer-forest-section-excess}, then every weak
component of the extended neutral-transfer digraph uses exactly two
underlying repair words.  Every word of
\(P_j\setminus(L\cup R')\) belongs to exactly two components, and every
word of \(L\mathbin\triangle R'\) belongs to exactly one.  In particular,
every directed cycle in a section attaining equality is a directed
two-cycle.

More precisely, form the bipartite incidence graph \(\mathfrak N_j\)
between the weak components and the underlying repair words which they
use.  Then every component vertex has degree two, every word vertex has
degree one or two, and
\begin{equation}
 \mathfrak N_j\text{ is a disjoint union of paths and cycles}.
 \label{eq:damp-neutral-forest-equality-path-cycle-skeleton}
\end{equation}
Its path endpoints are exactly the words in
\(L\mathbin\triangle R'\); every word on a cyclic component lies outside
\(L\cup R'\).
\end{corollary}

\begin{proof}
Equality at the two ends of the incidence count in the preceding proof
forces equality at every intermediate step.  Thus
\(|W(D)|=2\) for each component, every outside word attains its
multiplicity bound two, and every exclusive geodesic word attains its
multiplicity bound one.  A directed cycle with only two underlying words
alternates between those words and the two occurrence sides, so its
out-degree-one rule closes after the two opposite arrows.  It is a directed
two-cycle.  The asserted degrees in \(\mathfrak N_j\) are the same equality
conditions.  A finite graph of maximum degree two is a disjoint union of
paths and cycles.  Its degree-one word vertices are precisely the exclusive
geodesic words, while all degree-two word vertices lie outside the two
geodesics.  This proves the last assertions.
\end{proof}

\begin{corollary}[Only a two-word component can save a second side demand]
\label{cor:damp-neutral-two-sided-coalescence-star}
In the extended neutral-transfer digraph of one interior section, mark at
most one initial neutral boundary-corner demand in each of
\(\mathcal L_j\) and \(\mathcal R_j\).  Suppose that the two marked
vertices lie in the same weak component \(\mathcal C\).  Then exactly one
of the following holds:
\begin{enumerate}
 \item \(\mathcal C\) uses at least three underlying repair words.  In
       this case
       \begin{equation}
        \kappa_j+1\le |P_j|-(r+1),
        \label{eq:damp-neutral-component-matching-slack}
       \end{equation}
       so this component pays both marked demands;
 \item \(\mathcal C\) uses exactly two underlying repair words and is the
       directed two-cycle whose vertices are the two marked demands.
\end{enumerate}
\end{corollary}

\begin{proof}
If \(\mathcal C\) contains at least three underlying repair words, then
every other weak component contains at least two.  The component--word
incidence count in the proof of Proposition
\ref{prop:damp-neutral-transfer-forest-section-excess} gives
\[
 2\kappa_j+1\le2\bigl(|P_j|-(r+1)\bigr).
\]
Integrality proves the displayed strict-slack alternative.  Suppose
instead that \(\mathcal C\) contains exactly two repair words \(p,q\).
An arrow changes occurrence side and flips the fixed nonempty set
\(\operatorname{Diff}(p,q)\).  Hence the possible nonterminal pairs are
\[
 \{(L,p),(R,q)\},
 \qquad
 \{(R,p),(L,q)\},
\]
and no arrow joins the two pairs.  Since the two marked neutral sources
belong to opposite occurrence sides and lie in the same weak component,
they are the two vertices of one displayed pair.  Each has its prescribed
outgoing arrow to the other, so the component is their directed two-cycle.
\end{proof}

\begin{corollary}[Only a directed two-cycle is a tight coalescence]
\label{cor:damp-neutral-tight-coalescence-two-cycle}
Keep the marked two-sided setting of Corollary
\ref{cor:damp-neutral-two-sided-coalescence-star}.  If the common component
is not the directed two-cycle whose two vertices are the two marked
demands, then
\begin{equation}
 \kappa_j+1\le |P_j|-(r+1).
 \label{eq:damp-neutral-non-two-cycle-coalescence-slack}
\end{equation}
Thus every other common component already pays both initial side demands.
The only coalescence that can attain the
one-component bound is precisely the opposite-puncture two-cycle described
in Proposition~\ref{prop:damp-neutral-transfer-two-sided-cycle}.
\end{corollary}

\begin{proof}
This is the first alternative of Corollary
\ref{cor:damp-neutral-two-sided-coalescence-star}; its second alternative
is precisely the excluded directed two-cycle.
\end{proof}

\begin{proposition}[A neutral transfer to a noncorner exposes a split cube]
\label{prop:damp-neutral-transfer-noncorner-split-cube}
Let \(p\) be a boundary corner of \(P_j\) with a neutral parent puncture,
let \(q=p^S\) be its opposite repair tip as in
\eqref{eq:damp-neutral-transfer-opposite-tip}, and let \(Q\) be the unique
maximal cube of \(P_j\) through \(p\).  Write \(T\) for the support of
\(Q\).  If \(q\) is not a corner of \(P_j\), then \(q\) lies in a second
maximal cube \(Q'\) whose support \(T'\) is incomparable with \(T\).

If \(|T'|\ge2\), then
\begin{equation}
 |P_j|-(r+1)\ge2^{|T|}-|T|.
 \label{eq:damp-neutral-transfer-second-support-tax}
\end{equation}
If no such additional support tax is available, one may choose
\(T'=\{s\}\) with \(s\notin T\): the neutral transfer terminates at a
rank-one cube attached to the maximal repair cube \(Q\).  This is the only
a priori shadow-neutral split-corner configuration; Proposition
\ref{prop:damp-terminal-rank-one-impossible} below shows that the
path-core carrier geometry excludes it.
\end{proposition}

\begin{proof}
The parent puncture contains the full \(S\)-face through \(p\), so
\(q=p^S\) belongs to \(Q\).  Since \(q\) is not a corner, it lies in a
maximal cube \(Q'\ne Q\).  Parallel cubes with the same support are equal
or disjoint, so their common vertex gives \(T'\ne T\).  If
\(T'\subsetneq T\), then \(Q'\) is contained in the full \(T\)-cube
\(Q\), contrary to maximality of \(Q'\).  The reverse containment would
contradict maximality of \(Q\).  Hence \(T\) and \(T'\) are incomparable.

The full cube \(Q\) makes every subset of \(T\) a member of
\(\operatorname{Sh}(P_j)\).  If \(|T'|\ge2\), then \(T'\) is neither a
repair singleton nor a subset of \(T\), so it supplies one further member
of the shadow.  Together with the empty support and all \(r\) repair
singletons, ampleness gives
\[
 |P_j|
 \ge (r+1)+(2^{|T|}-|T|-1)+1,
\]
which is \eqref{eq:damp-neutral-transfer-second-support-tax}.  If this
extra-support alternative is absent, every second maximal cube through
\(q\) has singleton support.  Incomparability then puts its unique direction
outside \(T\), giving the stated attached rank-one cube.
\end{proof}

\begin{proposition}[A neutral transfer cannot terminate in a rank-one attachment]
\label{prop:damp-terminal-rank-one-impossible}
Keep the path-core notation.  Let \(p\) be a boundary corner of \(P_j\)
with a neutral parent puncture of support
\(S\mathbin{\dot\cup}\{x_\ell\}\), and let
\[
 q=p^S\in P_j\setminus C_\ell
\]
be its opposite repair tip.  Suppose that \(q\) is not a corner of
\(P_j\), let \(Q\) be the unique maximal cube through \(p\), and let
\(Q'\ne Q\) be a maximal cube through \(q\).  Then
\[
 |\operatorname{supp}(Q')|\ge2.
\]
Consequently the rank-one alternative in Proposition
\ref{prop:damp-neutral-transfer-noncorner-split-cube} never occurs in a
path-core repair.  Every neutral transfer reaching a noncorner pays the
additional-support bound
\[
 |P_j|-(r+1)\ge2^{|\operatorname{supp}(Q)|}
                 -|\operatorname{supp}(Q)|.
\]
\end{proposition}

\begin{proof}
Suppose for a contradiction that \(Q'\) is the edge
\[
 e=qz,
 \qquad z=q^{\{s\}},
\]
of repair direction \(s\).  Since this edge is an inclusion-maximal cube,
it belongs to no square of \(P_j\).  Lemma
\ref{lem:damp-square-free-edge-bridge} makes \(e\) a bridge.

It is also the only \(s\)-edge of \(P_j\).  Indeed, if \(q'z'\) were a
second \(s\)-edge, choose paths from \(q\) to \(q'\) and from \(z\) to
\(z'\) inside the two \(s\)-sections of \(P_j\).  These sections are
nonempty ample families and hence connected.  The two paths together with
\(e\) and \(q'z'\) would contain a cycle through \(e\), contradicting that
\(e\) is a bridge.

The adjacent carrier \(C_\ell\) contains the two fixed repair antipodes
\(f,\overline f\).  It is ample and therefore connected.  Since the two
antipodes differ in coordinate \(s\), every path between them in
\(C_\ell\) contains an \(s\)-edge.  The only such edge in \(P_j\) is
\(e\), so \(e\subseteq C_\ell\).  In particular \(q\in C_\ell\), contrary
to the choice of the opposite tip.  Hence \(Q'\) cannot have rank one.
The final inequality is now the nonsingleton case of Proposition
\ref{prop:damp-neutral-transfer-noncorner-split-cube}.
\end{proof}
\begin{proposition}[A rank-one closed transfer has length at least six]
\label{prop:damp-rank-one-neutral-cycle-surcharge}
Keep the hypotheses of Proposition
\ref{prop:damp-neutral-transfer-two-sided-cycle}, and suppose that every
arrow of a directed neutral-transfer cycle in \(P_j\) has a singleton
repair label.  If the cycle is not a two-cycle, then its length is at least
six and
\begin{equation}
 |C_j|\ge r+2,
 \qquad
 |C_{j+1}|\ge r+2.
 \label{eq:damp-rank-one-neutral-cycle-carrier-surcharge}
\end{equation}
Consequently, if \(T\) is the support of its common positioned maximal cube
and \(t=|T|\), then
\begin{equation}
 \Delta(H;A,R)\ge2^t-t+1.
 \label{eq:damp-rank-one-neutral-cycle-total-surcharge}
\end{equation}
Thus a four-cycle of rank-one square tips is impossible; every longer
rank-one cycle pays one additional concept in each adjacent carrier, besides
the maximal-cube tax.
\end{proposition}

\begin{proof}
Write the cycle as
\[
 p_0,p_1,\ldots,p_{2m-1},p_0,
\]
where the even vertices belong to \(\mathcal L_j\), the odd vertices belong
to \(\mathcal R_j\), and the arrow from \(p_h\) to \(p_{h+1}\) flips the
repair coordinate \(s_h\).  Indices are read modulo \(2m\).  The repair
words \(p_h\) are distinct.  Indeed, a repetition in the same vertex class
would repeat a vertex of the directed cycle, while a repetition in opposite
classes would give an odd closed walk in the bipartite repair cube.

The carrier \(C_{j+1}\) contains every even vertex and omits every odd
vertex.  Fix consecutive even vertices \(p_{2h}\) and \(p_{2h+2}\).  The
two coordinates \(s_{2h}\) and \(s_{2h+1}\) are distinct, since otherwise
these even vertices would be equal.  They are therefore at Hamming distance
two, with \(p_{2h+1}\) as one of their two common cube neighbours.  Since
\(C_{j+1}\) is ample and hence isometric, and since
\(p_{2h+1}\notin C_{j+1}\), the other common neighbour, say \(z_h\),
belongs to \(C_{j+1}\).

If \(m=2\), the two odd cycle vertices are exactly the two common neighbours
of the two even vertices.  Both are absent from \(C_{j+1}\), contradicting
isometry.  Thus every non-two-cycle has \(m\ge3\).

Suppose for a contradiction that \(|C_{j+1}|=r+1\).  The two fixed repair
antipodes force the empty support and all \(r\) repair singletons into
\(\operatorname{Sh}(C_{j+1})\).  Ampleness shows that these are its entire
shadow.  The edge--shadow identity gives exactly \(r\) edges, while
isometry gives connectedness.  Hence its one-inclusion graph is a tree.
Moreover, the geodesic between the repair antipodes has \(r\) edges, so it
uses the whole tree; the graph is a path.

The length-two paths
\[
 p_{2h},z_h,p_{2h+2}\qquad(0\le h<m)
\]
concatenate to a closed walk in this path.  Every closed walk in a tree has
an immediate backtrack.  No backtrack can occur at a vertex \(z_h\), because
the adjacent even vertices are distinct.  A backtrack at \(p_{2h}\) would
give \(z_{h-1}=z_h\); that common vertex would then be adjacent to the three
distinct vertices \(p_{2h-2},p_{2h},p_{2h+2}\), contradicting that the tree
is a path.  Therefore \(|C_{j+1}|\ge r+2\).  Interchanging the two sides of
the neutral cycle gives \(|C_j|\ge r+2\), proving
\eqref{eq:damp-rank-one-neutral-cycle-carrier-surcharge}.

Finally, the disjoint path-spine decomposition gives the exact excess
identity
\[
 \Delta(H;A,R)
 =|W|-(r+1)
  +\sum_{i=1}^n\bigl(|C_i|-(r+1)\bigr)
  +|\Omega|.
\]
The full \(T\)-cube in \(W\) gives
\(|W|-(r+1)\ge2^t-t-1\), the two displayed carrier surcharges contribute
two more, and every other term is nonnegative.  This proves
\eqref{eq:damp-rank-one-neutral-cycle-total-surcharge}.
\end{proof}

\begin{lemma}[Boolean-chain divergence is bounded by order disagreement]
\label{lem:damp-boolean-chain-divergence-inversions}
Let
\[
 \pi=(a_1,\ldots,a_r),
 \qquad
 \sigma=(b_1,\ldots,b_r)
\]
be two orders of the same coordinate set \(R\), and let
\[
 A_i=\{a_1,\ldots,a_i\},
 \qquad
 B_i=\{b_1,\ldots,b_i\}
 \qquad(0\le i\le r)
\]
be their Boolean chains.  Put
\[
 \nu(\pi,\sigma)
 =|\{i:A_i\ne B_i\}|.
\]
Then
\begin{equation}
 \nu(\pi,\sigma)
 \le \operatorname{inv}(\pi,\sigma),
 \label{eq:damp-boolean-chain-divergence-inversions}
\end{equation}
where \(\operatorname{inv}(\pi,\sigma)\) is the number of coordinate pairs
that occur in opposite orders in \(\pi\) and \(\sigma\).  Consequently,
every matching between the exclusive occurrences of the two chains has at
most \(\operatorname{inv}(\pi,\sigma)\) edges.  The pairs in this matching
need not be cube edges.
\end{lemma}

\begin{proof}
List the common prefix ranks as
\[
 0=i_0<i_1<\cdots<i_m=r,
 \qquad
 A_{i_h}=B_{i_h}.
\]
The coordinates split into consecutive common blocks
\[
 R_h=A_{i_h}\setminus A_{i_{h-1}}
    =B_{i_h}\setminus B_{i_{h-1}},
 \qquad
 \ell_h=|R_h|.
\]
There is no common prefix rank strictly between \(i_{h-1}\) and \(i_h\).
Thus the two orders restricted to \(R_h\) have
\(\ell_h-1\) distinct nonterminal prefix occurrences.  Summing over the
blocks gives
\begin{equation}
 \nu(\pi,\sigma)=\sum_{h=1}^m(\ell_h-1).
 \label{eq:damp-chain-divergence-block-count}
\end{equation}

We claim that the two orders of \(R_h\) have at least \(\ell_h-1\)
inversions.  Transform the first order into the second by a shortest
sequence of adjacent transpositions.  Its length is the inversion number
of the two restricted orders.  If it had fewer than \(\ell_h-1\) steps,
one of the \(\ell_h-1\) boundaries between consecutive positions would
never be used by a transposition.  No coordinate could then cross that
boundary, so the set of coordinates in the positions to its left would be
the same in both orders.  This would give a common prefix rank strictly
between \(i_{h-1}\) and \(i_h\), a contradiction.

Coordinates in different blocks occur in the same relative order under
\(\pi\) and \(\sigma\), so no inversion crosses two blocks.  The total
inversion number is therefore the sum of the block inversion numbers.
Together with \eqref{eq:damp-chain-divergence-block-count}, this proves
\eqref{eq:damp-boolean-chain-divergence-inversions}.

Each chain has exactly one occurrence of every cardinality.  Hence
\(A_i=B_k\) implies \(i=k\), and each chain has exactly
\(\nu(\pi,\sigma)\) exclusive occurrences.  A matching between the two
exclusive sets has at most this many edges, proving the final assertion.
\end{proof}

\begin{proposition}[Carrier excess and order disagreements pay every closed transfer]
\label{prop:damp-neutral-cycle-carrier-disagreement-packing}
Fix an interior path section \(P_j\), and put
\[
 \epsilon_j=|C_j|-(r+1),
 \qquad
 \epsilon_{j+1}=|C_{j+1}|-(r+1).
\]
Choose antipodal repair geodesics \(L\subseteq C_j\) and
\(R'\subseteq C_{j+1}\).  Let \(\nu(L,R')\) be their number of exclusive
vertices on either path, and let \(d(L,R')\) be the number of repair
coordinate pairs that they flip in opposite orders.

For every collection \(\mathcal Z\) of pairwise vertex-disjoint directed
cycles in the neutral-transfer digraph of \(P_j\),
\begin{equation}
 |\mathcal Z|
 \le \epsilon_j+\epsilon_{j+1}+\nu(L,R')
 \le \epsilon_j+\epsilon_{j+1}+d(L,R')
 \le \Delta(H;A,R).
 \label{eq:damp-neutral-cycle-carrier-disagreement-packing}
\end{equation}
In particular, every functional subdigraph has at most
\(\Delta(H;A,R)\) directed cycles.  This pays all same-section closed
transfers simultaneously, independently of their arrow labels, lengths,
common maximal-cube supports, or overlap inside one positioned cube.
\end{proposition}

\begin{proof}
Every ample repair carrier contains the two fixed repair antipodes and is
isometric.  It therefore contains an antipodal repair geodesic and has at
least \(r+1\) concepts, so the displayed excesses are nonnegative.

Every directed neutral cycle alternates between
\(\mathcal L_j\) and \(\mathcal R_j\).  Select from each member of
\(\mathcal Z\) one arrow directed from \(\mathcal L_j\) to
\(\mathcal R_j\).  If this arrow has root \(p\) and target \(q\), the
neutral parent at \(p\) gives
\[
 p\in C_{j+1}\setminus C_j,
 \qquad
 q\in C_j\setminus C_{j+1}.
\]
Indeed, the empty proper repair flip keeps \(p\) in the opposite carrier,
whereas membership in \(\mathcal L_j\) excludes it from \(C_j\); the
corresponding statements at the target follow because \(q\) is a retained
neutral occurrence in \(\mathcal R_j\).  Since the cycles in
\(\mathcal Z\) are vertex-disjoint, the selected pairs \((p,q)\) form a
matching between the two carrier differences.

Charge a selected pair having an endpoint outside \(L\cup R'\) to one such
endpoint.  The matching property makes these charges distinct, and there
are exactly \(\epsilon_j+\epsilon_{j+1}\) carrier vertices outside the two
geodesics.  Every uncharged pair has
\[
 p\in R'\setminus L,
 \qquad
 q\in L\setminus R',
\]
so the uncharged pairs form a matching between the exclusive occurrences
of the two Boolean chains.  Lemma
\ref{lem:damp-boolean-chain-divergence-inversions} bounds their number first
by \(\nu(L,R')\) and then by \(d(L,R')\).  This proves the first two
inequalities in
\eqref{eq:damp-neutral-cycle-carrier-disagreement-packing}.

Every coordinate pair ordered differently on \(L\) and \(R'\) is shattered
by their union and hence by the ample repair projection \(W\).  Thus
\(|W|-(r+1)\ge d(L,R')\).  The path-spine shadow partition gives the exact
identity
\[
 \Delta(H;A,R)
 =|W|-(r+1)
  +\sum_{i=1}^n\bigl(|C_i|-(r+1)\bigr)
  +|\Omega|.
\]
All summands are nonnegative, and the two carrier summands indexed by
\(j,j+1\) are \(\epsilon_j,\epsilon_{j+1}\).  This proves the last
inequality.
\end{proof}

\begin{corollary}[Closed components in distinct sections pay independently]
\label{cor:damp-neutral-cycle-sectionwise-tax}
Let \(\mathcal J\subseteq\{1,\ldots,n-1\}\), and suppose that for every
\(j\in\mathcal J\) the neutral-transfer digraph in \(P_j\) contains a
directed cycle whose common positioned maximal cube has support \(T_j\),
where \(t_j=|T_j|\).  Then
\begin{equation}
 \Delta(H;A,R)
 \ge V+\sum_{j\in\mathcal J}(2^{t_j}-t_j-1)
 \ge V+|\mathcal J|.
 \label{eq:damp-neutral-cycle-sectionwise-tax}
\end{equation}
In particular, overlap among the repair supports \(T_j\) at different
sections causes no double counting: the corresponding concepts lie over
different core words.
\end{corollary}

\begin{proof}
The concepts of \(H\) over the path and outside it form a disjoint
partition, so
\begin{equation}
 \Delta(H;A,R)
 =V+\sum_{h=0}^n\bigl(|P_h|-(r+1)\bigr).
 \label{eq:damp-path-section-excess-partition}
\end{equation}
Every \(P_h\) is ample and contains two repair antipodes; hence it shatters
the empty set and every repair singleton and has order at least \(r+1\).
For \(j\in\mathcal J\), the full \(T_j\)-cube adds all subsets of \(T_j\)
to \(\operatorname{Sh}(P_j)\), and therefore
\[
 |P_j|-(r+1)\ge2^{t_j}-t_j-1.
\]
Summing these inequalities in
\eqref{eq:damp-path-section-excess-partition} proves the first bound.  The
second follows from \(t_j\ge2\), proved in Proposition
\ref{prop:damp-neutral-transfer-two-sided-cycle}.
\end{proof}

\begin{proposition}[Parallel closed-component cubes pay by coordinate degree]
\label{prop:damp-parallel-neutral-cycle-cube-degree}
Fix a path section \(P_j\), a repair support \(T\subseteq R\) of order
\(t\ge2\), and let \(q_{j,T}\ge1\) be the number of distinct positioned
\(T\)-cubes in \(P_j\) that occur as common maximal cubes of directed
neutral-transfer cycles.  Then
\begin{equation}
 |P_j|-(r+1)
 \ge
 2^t-t-1+(q_{j,T}-1)2^{t-1}.
 \label{eq:damp-parallel-neutral-cycle-cube-degree}
\end{equation}
More generally, if the closed components in \(P_j\) use \(d_j\) distinct
maximal-cube supports and \(q_{j,T}\) positioned cubes of support \(T\),
then
\begin{equation}
 |P_j|-(r+1)
 \ge
 \max\left\{
  d_j,
  \max_T\bigl(2^{|T|}-|T|-1
       +(q_{j,T}-1)2^{|T|-1}\bigr)
 \right\}.
 \label{eq:damp-within-section-neutral-cycle-capacity}
\end{equation}
Thus neither different supports nor different positioned cubes with one
support can hide their local concept cost.  Proposition
\ref{prop:damp-neutral-cycle-carrier-disagreement-packing} separately pays
all remaining closed-component multiplicity, including several cycles in
one positioned maximal cube.
\end{proposition}

\begin{proof}
Choose \(e\in T\).  Distinct positioned \(T\)-cubes are disjoint, and
each contains exactly \(2^{t-1}\) edges of direction \(e\).  Hence
\begin{equation}
 |E_e(P_j)|\ge q_{j,T}2^{t-1}.
 \label{eq:damp-parallel-neutral-cycle-edge-count}
\end{equation}
The coordinatewise edge--shadow identity for the ample family \(P_j\)
gives
\[
 |E_e(P_j)|
 =|\{S\in\operatorname{Sh}(P_j):e\in S\}|.
\]
One positioned full \(T\)-cube already supplies the \(2^{t-1}\) subsets
of \(T\) that contain \(e\).  Therefore
\eqref{eq:damp-parallel-neutral-cycle-edge-count} forces at least
\((q_{j,T}-1)2^{t-1}\) further shattered supports outside \(2^T\).
Those supports are distinct from the empty support, all repair singletons,
and the \(2^t-t-1\) nonsingleton subsets of \(T\).  Ampleness now proves
\eqref{eq:damp-parallel-neutral-cycle-cube-degree}.

For the more general assertion, every distinct maximal-cube support has
order at least two and is itself a different member of
\(\operatorname{Sh}(P_j)\) outside the repair baseline.  This gives the
lower bound \(d_j\).  Applying
\eqref{eq:damp-parallel-neutral-cycle-cube-degree} separately to each
support and taking the maximum gives the other term in
\eqref{eq:damp-within-section-neutral-cycle-capacity}.
\end{proof}

\begin{corollary}[Two disjoint cycles in one cube force rank at least three]
\label{cor:damp-repeated-neutral-cycle-common-cube-rank}
Suppose that a functional subdigraph of the neutral-transfer digraph in
\(P_j\) has two vertex-disjoint directed cycles with the same common
positioned maximal cube \(Q\), of support \(T\).  Then \(|T|\ge3\), and
\begin{equation}
 |P_j|-(r+1)\ge4.
 \label{eq:damp-repeated-neutral-cycle-common-cube-tax}
\end{equation}
This rank restriction is a local geometric refinement of the general
closed-transfer packing theorem.
\end{corollary}

\begin{proof}
If \(|T|=2\), the two cycles use at least four distinct vertices and hence
use every vertex of \(Q\).  Every one of these vertices is a corner of
\(P_j\) whose unique maximal cube is \(Q\).  Consequently no vertex of
\(Q\) has a repair neighbour outside \(Q\).  Since \(P_j\) is connected,
this forces \(P_j=Q\).  But \(P_j\) contains two repair antipodes at
distance \(r\ge3\), a contradiction.  Therefore \(|T|\ge3\).  The full
\(T\)-cube then adds at least
\(2^3-3-1=4\) shattered supports beyond the empty support and repair
singletons, proving the displayed bound.
\end{proof}

\begin{proposition}[The bare three-section cube collar is peripheral]
\label{prop:damp-three-section-cube-collar-peripheral}
Let \(T\) be a coordinate set, let \(A,B\subseteq Q_T\) be nonempty ample
families, and introduce two further coordinates \(x,y\).  Then
\begin{equation}
 \mathcal L(A,B)
 =(A\times\{00\})
  \cup(Q_T\times\{10\})
  \cup(B\times\{11\})
 \subseteq Q_{T\mathbin{\dot\cup}\{x,y\}}
 \label{eq:damp-three-section-cube-collar}
\end{equation}
is ample.  If \(A,B\in\Damp\), then
\(\mathcal L(A,B)\in\Damp\).

Consequently, conditioning a same-cube neutral-transfer configuration to
its common repair cube and the three consecutive core sections cannot by
itself produce a smaller forbidden pc-minor: in a least-cardinality
cornerless family, both outer conditionings are smaller dismantlable ample
families, so the bare local collar is dismantlable.  Any contraction theorem
that uses such a local collar must retain parent-puncture provenance or
other attachments outside the three sections.
\end{proposition}

\begin{proof}
First form
\[
 C=(A\times\{0\})\cup(Q_T\times\{1\})
 \subseteq Q_{T\mathbin{\dot\cup}\{x\}}.
\]
After reversing the \(x\)-orientation, this is a peripheral expansion of
the cube \(Q_T\) along its ample subfamily \(A\).  It is therefore ample;
if \(A\in\Damp\), the exact peripheral-expansion criterion puts
\(C\) in \(\Damp\).

Identify \(B\) with the subfamily \(B\times\{1\}\subseteq C\).  With
respect to the new coordinate \(y\), the family in
\eqref{eq:damp-three-section-cube-collar} is
\[
 (C\times\{0\})
 \cup((B\times\{1\})\times\{1\}),
\]
another peripheral expansion, now of \(C\) along the copy of \(B\).
The same exact criterion proves ampleness and, when \(B,C\in\Damp\),
dismantlability.  The final assertion follows because conditionings and
pc-minors of an ample family are ample, while every smaller ample family
belongs to \(\Damp\) by Lemma
\ref{lem:damp-least-cornerless-cardinality-minimality}.
\end{proof}

\begin{proposition}[A fixed repair support has bounded closed-component capacity]
\label{prop:damp-fixed-support-neutral-cycle-capacity}
Keep the hypotheses of Proposition
\ref{prop:damp-neutral-transfer-two-sided-cycle}, and fix a repair support
\(T\subseteq R\) of order \(t\ge2\).  Let \(c_T\) be the number of
distinct pairs \((j,Q)\) such that the neutral-transfer digraph in \(P_j\)
has a directed cycle with common positioned maximal cube \(Q\) of support
\(T\).  Let \(q_T\) be the number of distinct positioned \(T\)-cubes
among these pairs.  If \(c_T>0\), then
\begin{equation}
 |U|\ge2^tq_T,
 \qquad
 V\ge2^{t-1}(c_T-q_T),
 \label{eq:damp-fixed-support-neutral-cycle-capacity}
\end{equation}
and hence
\begin{equation}
 c_T\le
 \frac{|U|+V}{2^{t-1}}-1
 \le
 \frac{|W|+|\Omega|}{2^{t-1}}-1.
 \label{eq:damp-fixed-support-neutral-cycle-count}
\end{equation}
Thus repeated closed components with one repair support cannot be hidden by
reusing the same local picture: a new positioned cube costs \(2^t\) path
repair words, while reappearance of an old positioned cube forces full
\(T\)-cubes over core words outside the path.  Several cycles with the same
pair \((j,Q)\) count only once in this estimate.
\end{proposition}

\begin{proof}
Fix one positioned repair \(T\)-cube \(Q\subseteq Q_R\), and let \(c_Q\)
be the number of section indices \(j\) for which \((j,Q)\) is counted by
\(c_T\).
Consider its core anchor family
\[
 A_Q=\{a\in Q_X:Q\times\{a\}\subseteq H\}.
\]
The family of anchors of all full \(T\)-cubes of \(H\) is ample by the
cube-carrier theorem.  Fixing the repair coordinates in \(R\setminus T\)
to the position of \(Q\) shows that \(A_Q\) is an ample conditioning, and
is therefore connected.

If a cycle using \(Q\) is centered at \(a_j\), then \(Q\subseteq P_j\),
whereas \eqref{eq:damp-neutral-cycle-isolated-cube} gives
\(Q\not\subseteq P_{j-1}\) and \(Q\not\subseteq P_{j+1}\).  Hence \(a_j\)
is an isolated component of \(A_Q\cap A\).  Distinct counted pairs have
distinct section indices, so the latter intersection has at least \(c_Q\)
components.

Put \(B_Q=A_Q\setminus A\).  If \(A_Q\cap A\) has at least two components,
connectivity of \(A_Q\) gives an edge from every component to \(B_Q\).
A core word outside the coordinate-simple path \(A\) has at most two
neighbors on \(A\): two such neighbors have path distance at most two,
distance one is excluded by bipartiteness of the cube, and three path
neighbors are therefore impossible.  Consequently
\[
 |B_Q|\ge\frac{c_Q-1}{2};
\]
the same inequality is trivial when \(c_Q=1\).  Every anchor in \(B_Q\)
supports all \(2^t\) concepts of \(Q\), and all of them are counted by
\(V\).  Thus \(Q\) forces at least \(2^{t-1}(c_Q-1)\) external concepts.

Distinct positioned \(T\)-cubes are disjoint, both over the path and over
external anchors.  Summing over the \(q_T\) cubes gives
\[
 |U|\ge2^tq_T,
 \qquad
 V\ge2^{t-1}\sum_Q(c_Q-1)
   =2^{t-1}(c_T-q_T),
\]
which proves \eqref{eq:damp-fixed-support-neutral-cycle-capacity}.  Adding
the two estimates and using \(q_T\ge1\) gives the first inequality in
\eqref{eq:damp-fixed-support-neutral-cycle-count}.  The second follows
from the path-spine identity
\[
 |W|+|\Omega|=|U|+\Phi+V\ge|U|+V.
\]
\end{proof}

\begin{proposition}[Global conservation of terminal transition charges]
\label{prop:damp-path-terminal-cycle-charge-conservation}
For each of the \(2n\) oriented adjacent-section incidences in Corollary
\ref{cor:damp-adjacent-section-extension-or-cycle}, perform progressive
synchronization and let \(\mathcal T\) be the incidences that stop with a
terminal cyclic remainder.  For \((i,j)\in\mathcal T\), denote that
remainder by \(D_{i,j}\) and put
\[
 \chi_{i,j}=|D_{i,j}|+c(D_{i,j}).
\]
Then
\begin{equation}
 \sum_{(i,j)\in\mathcal T}\chi_{i,j}
 \le
 2\sum_{\ell=0}^n|E(P_\ell)|
 \le
 2\sum_{T\in\operatorname{Sh}(H)}|T\cap R|.
 \label{eq:damp-path-terminal-cycle-global-conservation}
\end{equation}
In particular, if \(b=|\mathcal T|\), then
\begin{equation}
 5b\le\sum_{T\in\operatorname{Sh}(H)}|T\cap R|.
 \label{eq:damp-path-terminal-cycle-count}
\end{equation}
Thus cyclic transition failures cannot be reused without limit along the
path: each failure consumes ten units of section-edge rank, and every
section can pay for at most its two incident transitions.
\end{proposition}

\begin{proof}
For a terminal incidence \((i,j)\), equation
\eqref{eq:damp-adjacent-terminal-cycle-charge} charges \(\chi_{i,j}\) to
the rank sum of a subfamily of \(\operatorname{Sh}(P_j)\).  A fixed
section \(P_j\) is incident with at most two path edges.  Consequently,
after summing all terminal incidences, every member of
\(\operatorname{Sh}(P_j)\) has been charged at most twice.  The
edge--shadow identity gives
\[
 \sum_{(i,j)\in\mathcal T}\chi_{i,j}
 \le2\sum_{j=0}^n
       \sum_{T\in\operatorname{Sh}(P_j)}|T|
 =2\sum_{j=0}^n|E(P_j)|.
\]

The edges counted by \(E(P_j)\) are precisely the repair-coordinate edges
of \(H\) whose core word is \(a_j\).  These sets are disjoint for distinct
\(j\), so their union is contained in the set of all repair-coordinate
edges of \(H\).  A second application of the edge--shadow identity,
resolved by coordinate, counts the latter edges as
\[
 \sum_{T\in\operatorname{Sh}(H)}|T\cap R|.
\]
This proves \eqref{eq:damp-path-terminal-cycle-global-conservation}.  Each
terminal remainder has \(\chi_{i,j}\ge10\), and division by two yields
\eqref{eq:damp-path-terminal-cycle-count}.
\end{proof}

\begin{corollary}[The adjacent-carrier extension-or-owner trichotomy]
\label{cor:damp-adjacent-carrier-extension-or-owner}
Keep the path-spine notation and fix an oriented incidence \((i,j)\), with
\(j\in\{i-1,i\}\).  Run the progressive synchronization in Corollary
\ref{cor:damp-adjacent-section-extension-or-cycle}, and try its proposed
duplications in the full family, in the same order.  Put
\(\widehat H_0=H\), and after \(h\) successful duplications denote the
resulting ample family by \(\widehat H_h\).  One of the following occurs:
\begin{enumerate}
 \item the whole exclusive side is duplicated into the other path section
       by the sequence of ample one-concept extensions \(\widehat H_h\);
 \item a first global extension fails, and it has an antipodal blocker
       \(y=v^J\in\widehat H_h\), where \(v\) is the attempted next
       concept and
       \begin{equation}
        x_i\notin J,
        \qquad J\cap(X\setminus\{x_i\})\ne\varnothing;
        \label{eq:damp-adjacent-carrier-transverse-owner}
       \end{equation}
       in particular, the failure is owned by a core direction other than
       the transition direction \(x_i\);
 \item every proposed global extension succeeds until the local process
       stops at a terminal cyclic remainder \(D\), which satisfies
       \eqref{eq:damp-adjacent-terminal-cycle-charge}.
\end{enumerate}
This trichotomy is independent of the orders of \(H\) and of the core
path.
\end{corollary}

\begin{proof}
At any stage before a terminal remainder is reached, condition all core
coordinates in \(X\setminus\{x_i\}\) to the values common to
\(a_{i-1}\) and \(a_i\).  The conditioned family, with the previously
duplicated concepts included, is exactly the corresponding ample family
in \eqref{eq:damp-progressive-section-synchronization}.  Apply Lemma
\ref{lem:damp-conditioned-augmentation-transverse-blocker} to
\(\widehat H_h\), with
\[
 E=R\mathbin{\dot\cup}\{x_i\},
 \qquad F=X\setminus\{x_i\}.
\]
If an extension fails, its blocker satisfies
\(J\cap(X\setminus\{x_i\})\ne\varnothing\) by Lemma
\ref{lem:damp-conditioned-augmentation-transverse-blocker}.  The concept
being duplicated has its original copy across the \(x_i\)-edge in the
current ample family.  If \(x_i\in J\), that copy would be a third member
of \(Q_J(v)\cap(\widehat H_h\cup\{v\})\), contrary to
\eqref{eq:damp-conditioned-augmentation-transverse-blocker}.  Hence
\(x_i\notin J\), proving
\eqref{eq:damp-adjacent-carrier-transverse-owner}.  If none fails, the
local progressive process either duplicates the whole side or reaches the
terminal cyclic alternative of Theorem
\ref{thm:damp-complete-section-synchronization-or-cycle}.  These are the
other two conclusions.
\end{proof}

\begin{proposition}[Transverse owners move strictly toward an endpoint]
\label{prop:damp-adjacent-carrier-owner-monotonicity}
Suppose that the second alternative of Corollary
\ref{cor:damp-adjacent-carrier-extension-or-owner} occurs while duplicating
a concept into the section over \(a_k\), where
\(k\in\{i-1,i\}\).  Let \(v\) be the missing copy and let
\(y=v^J\in\widehat H_h\) be its antipodal blocker in the current ample
extension state.  Then either the core word of
\(y\) lies outside the path \(A=\{a_0,\ldots,a_n\}\), or it is \(a_h\)
for an index satisfying
\begin{equation}
 h\le i-2\quad\text{if }k=i-1,
 \qquad
 h\ge i+1\quad\text{if }k=i.
 \label{eq:damp-adjacent-carrier-owner-monotonicity}
\end{equation}
Thus an on-spine owner jumps strictly toward the endpoint on the target
side of the transition.  In particular, successive on-spine owner
transfers cannot form a directed cycle and have length at most \(n\).
\end{proposition}

\begin{proof}
Put \(K=J\cap X\).  Equation
\eqref{eq:damp-adjacent-carrier-transverse-owner} gives
\(K\ne\varnothing\) and \(x_i\notin K\).  The core word of \(v\) is
\(a_k\), so the core word of its blocker is \(a_k^K\).  If this word lies
outside \(A\), the first alternative holds.

Otherwise \(a_k^K=a_h\) for some \(h\).  Along the coordinate-simple
path,
\[
 K=\operatorname{Diff}(a_k,a_h)
  =\{x_{\min\{h,k\}+1},\ldots,x_{\max\{h,k\}}\}.
\]
This interval is nonempty and omits \(x_i\).  If \(k=i-1\), it must lie
strictly to the left of \(x_i\), so \(h\le i-2\).  If \(k=i\), it must
lie strictly to the right, so \(h\ge i+1\).  The index therefore moves
strictly toward the corresponding endpoint at every on-spine transfer,
which excludes a directed cycle and bounds the length of a transfer chain.
\end{proof}

\begin{proposition}[Every transverse owner pays the residual shadow]
\label{prop:damp-adjacent-carrier-owner-residual-charge}
Keep the failed-extension notation of Proposition
\ref{prop:damp-adjacent-carrier-owner-monotonicity}.  Write
\(\widehat H\) for the current ample extension state immediately before
the failed augmentation and write the blocker as
\(y=(g,b)\in\widehat H\).  Define the path sections
\(\widehat P_j\), adjacent carriers \(\widehat C_i\), full transverse
carriers \(\widehat D_i\), and residual shadow \(\widehat\Omega\) for
\(\widehat H\) as in Propositions
\ref{prop:damp-path-spine-euler-identity} and
\ref{prop:damp-transverse-path-spine-decomposition}.  Then
\begin{equation}
 y^{\{x_i\}}\in\widehat H,
 \qquad g\in\widehat D_i.
 \label{eq:damp-adjacent-owner-forced-parallel-edge}
\end{equation}
Moreover, at least one of the following residual payments exists:
\begin{enumerate}
 \item if \(g\notin\widehat C_i\), there is a support in
       \(\widehat\Omega\) with exact core part \(\{x_i\}\);
 \item if \(g\in\widehat C_i\), there is an \(\ell\ne i\) such that
       \(\{x_i,x_\ell\}\in\widehat\Omega_{\ge2}\).
\end{enumerate}
Thus every failed lift in the adjacent-carrier trichotomy produces a member
of the residual shadow of its current extension state.  In particular, a
failure on the first proposed duplication pays the original residual shadow
\(\Omega\) of \(H\).
\end{proposition}

\begin{proof}
Let \(v\) be the missing duplicated concept and put
\(w=v^{\{x_i\}}\).  The original copy \(w\) belongs to the current ample
family \(\widehat H\), while \(y=v^J\) is the blocker and
\(x_i\notin J\).  Hence
\[
 \operatorname{Diff}(y,w)=J\mathbin{\dot\cup}\{x_i\}.
\]
For every \(e\in J\), the neighbor \(y^{\{e\}}\) lies in the \(J\)-face
through \(v\) and is different from its two allowed members \(v,y\).
The blocker identity therefore gives
\(y^{\{e\}}\notin\widehat H\).  An isometric
\(\widehat H\)-geodesic from \(y\) to \(w\) must start in a direction
belonging to
\(J\mathbin{\dot\cup}\{x_i\}\).  All directions in \(J\) are absent at
\(y\), so its first edge has direction \(x_i\).  This proves
\eqref{eq:damp-adjacent-owner-forced-parallel-edge}.

If \(g\notin\widehat C_i\), then
\(g\in\widehat D_i\setminus\widehat C_i\).  Both carriers are
ample, so
\(\operatorname{Sh}(\widehat D_i)
 \setminus\operatorname{Sh}(\widehat C_i)\ne\varnothing\).
Every support in this difference, after adjoining \(x_i\), belongs to
\(\widehat\Omega\) by
\eqref{eq:damp-transverse-path-spine-decomposition}; this proves the first
alternative.

It remains to suppose that \(g\in\widehat C_i\).  The fibre
\(\widehat H_g\) then contains the path edge \(a_{i-1}a_i\) as well as the
parallel edge
\(bb^{\{x_i\}}\) from
\eqref{eq:damp-adjacent-owner-forced-parallel-edge}.  The word \(b\) is
neither \(a_{i-1}\) nor \(a_i\): it differs from the target core word of
\(v\) on the nonempty set \(J\cap X\), which omits \(x_i\).  Therefore
some coordinate \(x_\ell\ne x_i\) has a value on \(b\) opposite to its
common value on \(a_{i-1}\) and \(a_i\).  The two parallel \(x_i\)-edges
then realize all four traces on \(\{x_i,x_\ell\}\).  This pair is
shattered by \(\widehat H\), and it lies in
\(\widehat\Omega_{\ge2}\) because the spine family contains no support
with two core coordinates.
\end{proof}

\begin{proposition}[Residual shadows transport backwards through synchronization]
\label{prop:damp-residual-shadow-backward-transport}
Keep the path-spine notation for an ample family \(H\).  Let
\(j,k\in\{i-1,i\}\) be distinct, let
\(p\in P_j\setminus P_k\), and suppose that adjoining the missing copy
\(v=(p,a_k)\) gives an ample family
\[
 H^+=H\cup\{v\}.
\]
Let \(\Omega^+\) be the residual shadow of \(H^+\).  Then there is an
injection
\begin{equation}
 \iota_v:\Omega^+\longrightarrow\Omega
 \label{eq:damp-residual-shadow-backward-transport}
\end{equation}
that fixes every support in \(\Omega^+\cap\Omega\), with at most one
exception.  More precisely, let \(s(v)\in\{0,1\}\) record whether \(a_k\)
has a second path neighbor \(a_h\ne a_j\) for which \(p\in P_h\).  Then
\begin{equation}
 |\Omega^+|=|\Omega|-s(v).
 \label{eq:damp-residual-shadow-successful-duplication-drop}
\end{equation}
If the unique new shattered support of \(H^+\) belongs to \(\Omega^+\),
it is sent by \(\iota_v\) to an old residual support that entered the new
path-spine family; all other members are fixed.
\end{proposition}

\begin{proof}
The repair projection \(W\), its union \(U\), and the external-core count
\(V\) do not change: the repair word \(p\) was already present over
\(a_j\), and the new concept lies over the path core.  The section \(P_k\)
gains \(p\).  Consequently the adjacent carrier \(C_i\) gains \(p\), and
the only other carrier that can change is the one on the second path edge
at \(a_k\); it gains \(p\) exactly when \(s(v)=1\).  All changed carriers
are conditionings of the ample family \(H^+\), hence are ample.  Since each
has gained one concept, its shadow has gained exactly one support.

Let \(\mathcal B\) and \(\mathcal B^+\) be the path-spine families from
\eqref{eq:damp-path-spine-shadow-family}.  The carrier supports are
distinguished by their core coordinates, so the preceding paragraph gives
\begin{equation}
 \mathcal B\subseteq\mathcal B^+,
 \qquad
 |\mathcal B^+\setminus\mathcal B|=1+s(v).
 \label{eq:damp-successful-duplication-spine-growth}
\end{equation}
On the other hand, the ample one-concept extension has a unique new
shattered support; call it \(R_v\).  Thus
\[
 \operatorname{Sh}(H^+)=\operatorname{Sh}(H)\mathbin{\dot\cup}\{R_v\}.
\]
Subtracting \eqref{eq:damp-successful-duplication-spine-growth} from this
identity proves
\eqref{eq:damp-residual-shadow-successful-duplication-drop}.

If \(R_v\notin\Omega^+\), then every member of \(\Omega^+\) already
belongs to \(\Omega\), so inclusion is the desired injection.  Suppose
that \(R_v\in\Omega^+\).  Then none of the \(1+s(v)\) new spine supports
is \(R_v\).  Every one therefore belonged to \(\operatorname{Sh}(H)\),
and, because it did not belong to \(\mathcal B\), it belonged to
\(\Omega\).  Choose one such displaced support \(Q_v\), send
\(R_v\) to \(Q_v\), and fix every common residual support.  The support
\(Q_v\) is absent from \(\Omega^+\), so this map is injective and has the
asserted form.
\end{proof}

\begin{corollary}[All later residual payments return to the original shadow]
\label{cor:damp-residual-payment-amortization}
Let
\[
 H=H_0\subset H_1\subset\cdots\subset H_q
\]
be any sequence of successful adjacent-section duplications of the kind in
Corollary~\ref{cor:damp-adjacent-carrier-extension-or-owner}, with every
\(H_h\) ample.  If \(\Omega_h\) is the residual shadow of \(H_h\), then
the maps in Proposition~\ref{prop:damp-residual-shadow-backward-transport}
compose to an injection
\begin{equation}
 \Omega_q\hookrightarrow\Omega_0.
 \label{eq:damp-residual-payment-amortization}
\end{equation}
Consequently every residual support paid by a later failed lift has a
well-defined representative in the original residual shadow of \(H\).
\end{corollary}

\begin{proof}
Apply Proposition~\ref{prop:damp-residual-shadow-backward-transport} to
each inclusion \(H_{h-1}\subset H_h\) and compose the resulting injections
in reverse chronological order.
\end{proof}

\begin{proposition}[The two progressive synchronizations commute]
\label{prop:damp-two-sided-progressive-synchronization}
Fix an interior path section \(P_j\).  Suppose that the complete-
synchronization alternative of Corollary
\ref{cor:damp-adjacent-section-extension-or-cycle} holds for \(P_j\) at
both adjacent path edges.  Retain orders
\[
 p_1,\ldots,p_a
 \quad\text{of }P_j\setminus C_j,
 \qquad
 q_1,\ldots,q_b
 \quad\text{of }P_j\setminus C_{j+1}
\]
supplied by that corollary.  Duplicate the \(p\)'s into the section over
\(a_{j-1}\) and the \(q\)'s into the section over \(a_{j+1}\), in any
shuffle that preserves the two displayed orders.

At every stage, the next proposed duplication is ample in the corresponding
two-layer conditioning.  Consequently, exactly one of the following occurs:
\begin{enumerate}
 \item a first duplication fails in the full family; it has a transverse
       owner and produces a residual payment which transports injectively
       to the residual shadow of the original family;
 \item every duplication succeeds in the full family.  The resulting ample
       extension \(H^\ast\) has path sections
       \begin{equation}
       P_{j-1}^\ast=P_{j-1}\cup P_j,\qquad
       P_j^\ast=P_j,\qquad
       P_{j+1}^\ast=P_{j+1}\cup P_j,
       \label{eq:damp-two-sided-synchronized-sections}
       \end{equation}
       and its residual shadow injects into the residual shadow of \(H\).
\end{enumerate}
In the second outcome, both adjacent two-layer conditionings are
peripheral toward their outer section.
\end{proposition}

\begin{proof}
Consider a proposed duplication of \(p_h\) across the edge
\(a_{j-1}a_j\).  Condition every core coordinate except \(x_j\) to the
common values of these two core words.  Every copy previously added over
\(a_{j+1}\) is absent from this conditioning, because \(a_{j+1}\) differs
from \(a_j\) in \(x_{j+1}\), which has just been fixed.  The conditioned
family is therefore exactly the \(h\)-th stage of the progressive
synchronization order for the left edge, together with the earlier
\(p\)-prefix.  Corollary
\ref{cor:damp-adjacent-section-extension-or-cycle} says that adjoining the
next copy is ample in this conditioning.  The symmetric argument applies
to every proposed \(q_h\)-duplication.  Hence the two local orders may be
shuffled arbitrarily.

Run the shuffled sequence in the full family.  At a first failed
duplication, apply Lemma
\ref{lem:damp-conditioned-augmentation-transverse-blocker} to the current
ample extension state and the corresponding two-layer conditioning.  The
other-side additions vanish from that conditioning by the preceding
paragraph.  The argument in Corollary
\ref{cor:damp-adjacent-carrier-extension-or-owner} therefore applies
verbatim: the blocker is transverse, and the original copy across the
transition edge excludes that edge from the blocker support.  Proposition
\ref{prop:damp-adjacent-carrier-owner-residual-charge} supplies a residual
payment in the current extension state.  Corollary
\ref{cor:damp-residual-payment-amortization} transports it through all
earlier successful duplications to the original residual shadow.

If no duplication fails, every intermediate family and the final family
\(H^\ast\) are ample.  The left target section has acquired exactly
\(P_j\setminus C_j\); since
\(C_j=P_{j-1}\cap P_j\), it is
\[
 P_{j-1}\cup(P_j\setminus C_j)=P_{j-1}\cup P_j.
\]
The right identity is symmetric, proving
\eqref{eq:damp-two-sided-synchronized-sections}.  Backward residual
transport along the full successful sequence gives the asserted injection.
Finally, \(P_j\) is contained in each outer union in
\eqref{eq:damp-two-sided-synchronized-sections}, which is precisely the
peripheral alternative in the corresponding two-layer conditioning.
\end{proof}

\begin{corollary}[A tight two-cycle hands off to two peripheral transitions]
\label{cor:damp-tight-two-cycle-peripheral-handoff}
Suppose that the two marked initial demands in an interior section form the
tight directed two-cycle of Corollary
\ref{cor:damp-neutral-tight-coalescence-two-cycle}.  Retain the two
complete-synchronization orders for which these marked corners are the last
terms.  Then either a shuffled synchronization produces a transverse
residual payment, or there is an ample extension satisfying
\eqref{eq:damp-two-sided-synchronized-sections}.  Thus a tight two-cycle
creates no further terminal geometry: its only unpaid outcome is a
double-peripheral conditioned handoff.
\end{corollary}

\begin{proof}
Apply Proposition
\ref{prop:damp-two-sided-progressive-synchronization}, using a shuffle in
which all nonfinal terms precede the two marked final terms.  Its two
alternatives are exactly the stated outcomes.
\end{proof}

\begin{lemma}[An adjacent duplication transports one section-shadow boundary]
\label{lem:damp-adjacent-duplication-shadow-transport}
Keep the path-spine notation, let $j,k$ be the two ends of the path edge
$a_{i-1}a_i$, and suppose that
\[
 p\in P_j\setminus P_k,
 \qquad
 H^+=H\cup\{(p,a_k)\}
\]
is ample.  Put $C_i=P_j\cap P_k$.  Then $P_k\cup\{p\}$ and
$C_i\cup\{p\}$ are ample one-concept extensions, and there is a unique
repair support $T$ such that
\begin{align}
 \operatorname{Sh}(P_k\cup\{p\})
   &=\operatorname{Sh}(P_k)\mathbin{\dot\cup}\{T\},\notag\\
 \operatorname{Sh}(C_i\cup\{p\})
   &=\operatorname{Sh}(C_i)\mathbin{\dot\cup}\{T\}.
 \label{eq:damp-adjacent-duplication-common-shadow-token}
\end{align}
Moreover,
\begin{equation}
 T\in\operatorname{Sh}(P_j)\setminus\operatorname{Sh}(P_k).
 \label{eq:damp-adjacent-duplication-shadow-source}
\end{equation}
Thus a successful concept duplication carries exactly one existing
section-shadow support across the same path edge; it creates no new repair
support.
\end{lemma}

\begin{proof}
The two displayed families are conditionings of the ample family $H^+$,
so they are ample.  Each has one more concept than its counterpart before
the duplication and consequently acquires one shattered support.  Call the
new supports $T_k$ and $T_C$, respectively.

Apply the two-layer shadow-intersection identity before and after the
duplication.  Since the source section does not change, it gives
\[
 \operatorname{Sh}(C_i)
  =\operatorname{Sh}(P_j)\cap\operatorname{Sh}(P_k),
 \qquad
 \operatorname{Sh}(C_i\cup\{p\})
  =\operatorname{Sh}(P_j)
       \cap\operatorname{Sh}(P_k\cup\{p\}).
\]
The second intersection is obtained from the first by testing only the one
new support $T_k$.  It must grow, because $C_i\cup\{p\}$ is an ample
one-concept extension.  Hence $T_k\in\operatorname{Sh}(P_j)$ and
$T_C=T_k$.  Taking $T=T_k$ proves both assertions.
\end{proof}

\begin{proposition}[Weighted path boundaries are Lyapunov functions]
\label{prop:damp-adjacent-duplication-boundary-lyapunov}
For a tuple of nonempty ample path sections
$\mathbf P=(P_0,\ldots,P_n)$, define its directed word-boundary potential
by
\begin{equation}
 \Lambda_0(\mathbf P)=
 \sum_{i=1}^n\left(
  i|P_i\setminus P_{i-1}|
  +(n-i+1)|P_{i-1}\setminus P_i|
 \right)
 \label{eq:damp-word-boundary-lyapunov}
\end{equation}
and its rank-weighted shadow-boundary potential by
\begin{equation}
 \begin{split}
 \Lambda_1(\mathbf P)=\sum_{i=1}^n\biggl(&
  i\sum_{T\in\operatorname{Sh}(P_i)
                   \setminus\operatorname{Sh}(P_{i-1})}|T|\\
  &+(n-i+1)\sum_{T\in\operatorname{Sh}(P_{i-1})
                   \setminus\operatorname{Sh}(P_i)}|T|\biggr).
 \end{split}
 \label{eq:damp-shadow-boundary-lyapunov}
\end{equation}
If one concept is duplicated successfully from one path section into its
neighbour, then $\Lambda_0$ decreases by at least one.  If $T$ is the
support in Lemma~\ref{lem:damp-adjacent-duplication-shadow-transport}, then
$\Lambda_1$ decreases by at least $|T|$.  More precisely, each decrease is
either the displayed minimum or the annihilation value obtained when the
transported boundary meets the oppositely oriented boundary at the next
edge.

Put
\[
 \mathcal E(\mathbf P)=\sum_{j=0}^n|E(P_j)|.
\]
Along every sequence of successful adjacent-section duplications,
\begin{equation}
 \mathcal E(\mathbf P)+\Lambda_1(\mathbf P)
 \quad\text{is nonincreasing}.
 \label{eq:damp-section-edge-shadow-lyapunov}
\end{equation}
In particular, every such sequence has length at most
$\Lambda_0(\mathbf P)$ at its initial state.  The initial potentials obey
\begin{align}
 \Lambda_0(\mathbf P)
  &\le2n(|W|+|\Omega|),
 \label{eq:damp-word-boundary-lyapunov-bound}\\
 \Lambda_1(\mathbf P)
  &\le2n\left(
       \sum_{j=0}^n|E(P_j)|-
       \sum_{i=1}^n|E(C_i)|\right).
 \label{eq:damp-shadow-boundary-lyapunov-bound}
\end{align}
Thus no successful handoff can circulate: at concept level and at
rank-weighted shadow level every moved boundary travels strictly toward an
endpoint or annihilates with an opposite boundary.
\end{proposition}

\begin{proof}
Consider first a duplication from $P_i$ into $P_{i-1}$.  It removes the
directed boundary occurrence of weight $i$.  If $i>1$ and the duplicated
word is absent from $P_{i-2}$, it creates the same directed occurrence at
the preceding edge, now of weight $i-1$.  The net decrease is one.  If the
word is present in $P_{i-2}$, the duplication instead removes the
oppositely directed occurrence at that edge; the decrease is larger.  At
$i=1$ the occurrence simply leaves the path.  The argument for a
duplication from $P_{i-1}$ into $P_i$ is symmetric, with weights
$n-i+1$ and $n-i$.  This proves the assertion for $\Lambda_0$.

Lemma~\ref{lem:damp-adjacent-duplication-shadow-transport} says that the
same operation is performed on exactly one section-shadow support $T$.
Since the target section was nonempty, it already shattered the empty set,
so $T\ne\varnothing$.
Multiplying the preceding calculation by $|T|$ proves the assertion for
$\Lambda_1$.  The edge--shadow identity gives
\[
 |E(P_k\cup\{p\})|-|E(P_k)|=|T|.
\]
No other section changes.  Hence the increase of $\mathcal E$ is at most
the simultaneous decrease of $\Lambda_1$, proving
\eqref{eq:damp-section-edge-shadow-lyapunov}.  Since every successful
duplication decreases the nonnegative integer $\Lambda_0$, the length
bound follows.

Every coefficient in \eqref{eq:damp-word-boundary-lyapunov} is at most
$n$.  Proposition~\ref{prop:damp-section-shadow-boundary-conservation}
therefore gives
\[
 \Lambda_0(\mathbf P)
 \le n\sum_{i=1}^n|P_{i-1}\mathbin\triangle P_i|
 \le2n(|W|+|\Omega|),
\]
which is \eqref{eq:damp-word-boundary-lyapunov-bound}.

For the last bound, fix $T\in\bigcup_j\operatorname{Sh}(P_j)$.  If its
section-index set has $c_{\rm sh}(T)$ components, it has at most
$2c_{\rm sh}(T)$ boundary edges.  The rank-weighted version of the interval
count in Proposition
\ref{prop:damp-section-shadow-boundary-conservation} is
\[
 \sum_T |T|c_{\rm sh}(T)
 =\sum_{j=0}^n|E(P_j)|-
   \sum_{i=1}^n|E(C_i)|.
\]
Multiplying the boundary count by the maximum weight $n$ and summing over
$T$ proves \eqref{eq:damp-shadow-boundary-lyapunov-bound}.
\end{proof}

\begin{corollary}[Boundary descent pays every acquired marked copy]
\label{cor:damp-boundary-descent-pays-marked-acquisitions}
Consider any sequence of successful adjacent-section duplications, and
mark any collection of directed boundary demands that disappear during the
sequence.  Charge a marked demand to the first duplication that adjoins its
missing copy.  Then
\begin{equation}
 \#\{\text{disappearing marked demands}\}
 \le
 \Lambda_0(\mathbf P^{\rm in})-
 \Lambda_0(\mathbf P^{\rm out}).
 \label{eq:damp-boundary-descent-marked-demand-payment}
\end{equation}
In particular, two different marked demands can be resolved by the same
duplication only when two oppositely oriented boundaries annihilate there;
that step decreases $\Lambda_0$ by $n+2$, and so pays both demands with
strict slack.
\end{corollary}

\begin{proof}
A duplication always removes the directed boundary across which it is
performed.  It can remove at most one further marked boundary, namely the
oppositely directed boundary at the other edge incident with its target
section.  If it removes no such second boundary, the proof of Proposition
\ref{prop:damp-adjacent-duplication-boundary-lyapunov} gives a decrease of
at least one.  If it removes the second boundary, the two removed weights
sum to
\[
 i+(n-i+2)=n+2
 \quad\text{or}\quad
 (n-i+1)+(i+1)=n+2,
\]
according to the direction of transport.  This is at least two.  Assigning
each marked demand to its first resolving duplication and summing these
stepwise inequalities proves
\eqref{eq:damp-boundary-descent-marked-demand-payment}.
\end{proof}

\begin{proposition}[All original tight two-cycles reduce in two sweeps]
\label{prop:damp-tight-two-cycle-two-sweep-reduction}
In the original path-spine state, mark every interior section whose two
selected initial demands form the tight directed two-cycle of Corollary
\ref{cor:damp-neutral-tight-coalescence-two-cycle}; let $\mathcal J$ be the
set of their indices.  All members of $\mathcal J$ can be treated together
as follows.

First process the leftward complete-synchronization orders in increasing
order of $j\in\mathcal J$.  Then process the unresolved rightward
incidences in decreasing order.  At a rightward incidence use a fresh
progressive synchronization in the current ample two-layer conditioning,
and stop as soon as the marked missing copy is acquired.  For every
$j\in\mathcal J$, one of the following occurs:
\begin{enumerate}
 \item one of the required lifts fails globally and supplies a transverse
       residual payment, which transports to the original residual shadow;
 \item a fresh rightward synchronization stops before the marked copy and
       supplies a terminal cyclic block;
 \item both marked missing copies are acquired by successful ample
       extensions, so the original tight two-cycle disappears.
\end{enumerate}
Every successful duplication in the two sweeps is charged by the strict
descent in Proposition
\ref{prop:damp-adjacent-duplication-boundary-lyapunov}.  Consequently all
successful work in the simultaneous reduction is finite and satisfies the
original-state edge--shadow bound
\eqref{eq:damp-section-edge-shadow-lyapunov}.  No iterated handoff cycle,
owner-repetition alternative, or cardinality-by-cardinality closure is
needed.
\end{proposition}

\begin{proof}
For the leftward sweep, an operation at an earlier centre $k<j$ changes
only the section over $a_{k-1}$; in particular, it changes neither section
over the edge $a_{j-1}a_j$.  After conditioning to that edge, all earlier
additions therefore disappear.  The retained leftward order for the
original tight two-cycle is still locally valid.  Trying its duplications
in the full current family either gives the transverse failure in
Corollary~\ref{cor:damp-adjacent-carrier-extension-or-owner}, with backward
residual transport, or acquires the marked left copy.

After the leftward sweep, consider the indices in decreasing order.  An
already processed larger centre changes only a section strictly to the
right of the current edge.  The leftward sweep may have enlarged the two
current sections, but they remain conditionings of the current ample full
family and their intersection is ample.  Apply Theorem
\ref{thm:damp-complete-section-synchronization-or-cycle} afresh in the
rightward orientation.  If the marked copy has already been acquired, no
operation is needed.  Otherwise, either the progressive process reaches
that copy, a proposed full lift first fails and gives the residual outcome,
or the process stops earlier with the stated terminal cyclic block.  These
are exactly the three alternatives in the proposition.

All full families retained during the construction are ample, and every
successful step is an adjacent-section duplication.  Proposition
\ref{prop:damp-adjacent-duplication-boundary-lyapunov} therefore applies to
their complete concatenation.  Corollary
\ref{cor:damp-boundary-descent-pays-marked-acquisitions} pays all marked
copies acquired in the third alternative, including the case in which one
addition resolves marked demands from two neighbouring centres.  This gives
the final assertions.
\end{proof}

\begin{proposition}[Dynamic terminal blocks obey original-state conservation]
\label{prop:damp-dynamic-terminal-block-conservation}
Start from a path-spine state with section tuple
$\mathbf P^{\rm in}$, perform any sequence of successful adjacent-section
duplications, and process each oriented path incidence at most once.  The
incidences may be processed in different intermediate ample states.  Let
$\mathcal T$ be those that stop with a terminal cyclic remainder, and retain
the notation
\[
 \chi_{i,j}=|D_{i,j}|+c(D_{i,j})
\]
from Proposition~\ref{prop:damp-path-terminal-cycle-charge-conservation}.
If $\mathbf P^{\rm out}$ is the final section tuple, then
\begin{equation}
 \sum_{(i,j)\in\mathcal T}\chi_{i,j}
 \le2\mathcal E(\mathbf P^{\rm out})
 \le2\bigl(
    \mathcal E(\mathbf P^{\rm in})+
    \Lambda_1(\mathbf P^{\rm in})\bigr).
 \label{eq:damp-dynamic-terminal-block-conservation}
\end{equation}
Thus terminal blocks created during the two-sweep reduction require no new
state-dependent budget: their complete rank charge returns to the initial
section-edge rank together with the initial shadow-boundary potential.
\end{proposition}

\begin{proof}
At the state in which $(i,j)$ stops, equation
\eqref{eq:damp-adjacent-terminal-cycle-charge} charges
$\chi_{i,j}$ to the rank sum of a subfamily of
$\operatorname{Sh}(P_j)$ in that state.  Sections only grow during the
duplication sequence, and every intermediate section remains ample.
Consequently its shattered complex is contained in the shattered complex
of the final section $P_j^{\rm out}$.

A fixed section is the source of at most its two oriented path incidences.
Therefore every support in
$\operatorname{Sh}(P_j^{\rm out})$ is counted at most twice after summing
all terminal charges based at that section.  The edge--shadow identity
gives
\[
 \sum_{(i,j)\in\mathcal T}\chi_{i,j}
 \le2\sum_{j=0}^n
       \sum_{T\in\operatorname{Sh}(P_j^{\rm out})}|T|
 =2\mathcal E(\mathbf P^{\rm out}).
\]
Finally apply the Lyapunov inequality
\eqref{eq:damp-section-edge-shadow-lyapunov} to the complete successful
duplication sequence and discard the nonnegative terminal value of
$\Lambda_1$.  This proves the second inequality in
\eqref{eq:damp-dynamic-terminal-block-conservation}.
\end{proof}

\begin{proposition}[Successful acquisitions form a spacetime support forest]
\label{prop:damp-spacetime-support-forest}
Start from an initial path-section tuple
\(\mathbf P^{\rm in}\), perform a sequence of successful adjacent-section
duplications, and let \(\mathbf P^{\rm out}\) be the final tuple.  Fix a
nonbaseline repair support \(T\).  Write
\[
 I_T^{\rm in}
 =\{j:T\in\operatorname{Sh}(P_j^{\rm in})\},
 \qquad
 I_T^{\rm out}
 =\{j:T\in\operatorname{Sh}(P_j^{\rm out})\},
\]
and denote their orders and numbers of path components by
\(m_T^{\rm in},c_T^{\rm in}\) and
\(m_T^{\rm out},c_T^{\rm out}\), respectively.  Let \(q_T\) be the
number of successful duplications whose transported repair support is
\(T\), and let \(\mu_T\) be the number of those steps at which the newly
added section index joins two components of the current \(T\)-index set.
Then
\begin{align}
 m_T^{\rm out}
 &=m_T^{\rm in}+q_T,
 \label{eq:damp-spacetime-support-vertex-conservation}\\
 c_T^{\rm out}
 &=c_T^{\rm in}-\mu_T,
 \label{eq:damp-spacetime-support-component-conservation}\\
 (m_T^{\rm out}-c_T^{\rm out}+1)
  -(m_T^{\rm in}-c_T^{\rm in}+1)
 &=q_T+\mu_T.
 \label{eq:damp-spacetime-support-edge-conservation}
\end{align}

More precisely, start with the rooted support forest from
\eqref{eq:damp-support-interval-boundary-map}.  A nonmerging acquisition
adjoins one occurrence vertex and its processed path edge.  A merging
acquisition adjoins one occurrence vertex and both incident path edges.
The resulting graph is the rooted support forest of \(I_T^{\rm out}\);
in particular, no spacetime support cycle is created.

Consequently, if a demand family \(\mathcal D^{\rm out}\) injects into the
final nonbaseline occurrence set
\[
 \mathfrak O^{\rm out}
 =\{(j,T):T\notin\mathcal L_R,\ j\in I_T^{\rm out}\},
\]
then it and the failed lifts admit a joint injection into the disjoint
union of
\[
 \Xi,\qquad \Omega_0,\qquad
 \mathcal A=\{\text{successful support-acquisition steps}\}.
\]
Equivalently,
\begin{equation}
 |\mathcal D^{\rm out}|+|\mathcal F|
 \le\Delta(H;A,R)+|\mathcal A|.
 \label{eq:damp-spacetime-occurrence-failure-bound}
\end{equation}
Thus the only resource in a final-state terminal payment that is not
already present in the original excess is the successful acquisition
itself.
\end{proposition}

\begin{proof}
At a successful duplication, Lemma
\ref{lem:damp-adjacent-duplication-shadow-transport} supplies one unique
repair support \(T\) that is present in the source-section shadow and
absent from the target-section shadow.  The step adds the target index to
the current set \(I_T\), so every \(T\)-acquisition increases its order by
one.  This proves
\eqref{eq:damp-spacetime-support-vertex-conservation}.

The new index is adjacent to its source index.  If its other path neighbour
is absent from \(I_T\), it extends one component and leaves the component
number unchanged.  If the other neighbour is present, the two neighbours
belong to different components of the path with the target index deleted;
adjoining the target joins those components and decreases their number by
one.  These are the only cases, proving
\eqref{eq:damp-spacetime-support-component-conservation}.  Subtraction
gives \eqref{eq:damp-spacetime-support-edge-conservation}.

The same argument proves the forest description.  In the first case one
adjoins a leaf and one edge.  In the second, one adjoins a degree-two
vertex between two previously different tree components.  Neither
operation creates a cycle.  The final graph has one root edge in a chosen
component and every internal edge of the path induced by
\(I_T^{\rm out}\), exactly as in the proof of Proposition
\ref{prop:damp-path-spine-support-interval-exact-sequence}.

It remains to count the available columns without charging a component
merger twice.  For fixed \(T\),
\begin{align*}
 m_T^{\rm out}
 &=(m_T^{\rm in}-c_T^{\rm in}+1)
   +q_T+(c_T^{\rm in}-1).
\end{align*}
The first term is the number of original rooted and internal-edge supports
in the exact-repair part of \(\Xi\).  Use the \(q_T\) acquisition steps as
formal columns for the second term.  Use the
\(c_T^{\rm in}-1\) original fragmentation generators for the last term.
Proposition~\ref{prop:damp-tight-gap-core-square-reservation} injects all
these fragmentation generators, jointly with all failures, into
\(\Omega_0\).  Therefore any injection of
\(\mathcal D^{\rm out}\) into the final occurrence set lifts to the
asserted joint injection.  Summing over \(T\), and applying
\eqref{eq:damp-support-interval-total-excess}, proves
\eqref{eq:damp-spacetime-occurrence-failure-bound}.
\end{proof}

\begin{proposition}[A dynamic terminal block has an unmarked acquisition
when it needs one]
\label{prop:damp-dynamic-terminal-unmarked-acquisition}
Keep the two-sweep reduction of Proposition
\ref{prop:damp-tight-two-cycle-two-sweep-reduction}.  Suppose that the
original tight two-cycle centred at the interior section \(P_j^{\rm in}\)
ends in the terminal cyclic alternative in a later section
\(P_j^{\rm cur}\).  Put
\[
 e_j^{\rm in}=|P_j^{\rm in}|-(r+1),
\]
and let \(a_j\) be the number of successful duplications before that
terminal state whose target is \(P_j\).  Then
\begin{equation}
 e_j^{\rm in}\ge1,
 \qquad
 e_j^{\rm in}+a_j\ge9.
 \label{eq:damp-dynamic-terminal-original-and-current-excess}
\end{equation}
Consequently exactly one of the following holds:
\begin{enumerate}
 \item \(e_j^{\rm in}\ge2\), so the two sides of the original tight cycle
       can be assigned two distinct original nonbaseline occurrences in
       \(P_j^{\rm in}\);
 \item \(e_j^{\rm in}=1\), and at least one of the acquisitions into
       \(P_j\) is unmarked, meaning that it does not first acquire either
       selected missing copy of an original two-sided demand.
\end{enumerate}
The unmarked acquisition in the second alternative may be chosen
injectively over all dynamic terminal sections.
\end{proposition}

\begin{proof}
The original directed two-cycle has a common positioned \(T\)-cube with
\(|T|\ge2\) by Proposition
\ref{prop:damp-neutral-transfer-two-sided-cycle}.  Hence
Corollary~\ref{cor:damp-neutral-cycle-sectionwise-tax}, applied in the
single section \(j\), gives
\[
 e_j^{\rm in}\ge2^{|T|}-|T|-1\ge1.
\]

At the later terminal state, write \(Z\) for the last synchronized
subfamily and \(D=P_j^{\rm cur}\setminus Z\) for the terminal remainder.
The family \(Z\) contains the current adjacent carrier, which contains the
two fixed repair antipodes.  It is isometric and therefore has at least
\(r+1\) concepts.  The terminal theorem gives \(|D|\ge9\), so
\[
 |P_j^{\rm cur}|-(r+1)
 =|Z|+|D|-(r+1)\ge9.
\]
Every successful duplication into \(P_j\) adds one concept and no operation
deletes a concept.  Thus
\[
 |P_j^{\rm cur}|=|P_j^{\rm in}|+a_j,
\]
which proves
\eqref{eq:damp-dynamic-terminal-original-and-current-excess}.

If \(e_j^{\rm in}\ge2\), ampleness identifies the original nonbaseline
section-shadow supports with exactly \(e_j^{\rm in}\) distinct occurrence
labels, proving the first alternative.  Suppose instead that
\(e_j^{\rm in}=1\).  The displayed inequality gives \(a_j\ge8\).
Only the two oriented incidences directed into \(P_j\) can have a selected
missing copy whose acquisition step targets \(P_j\), and each incidence
has at most one such marked copy.  Therefore at most two of these
\(a_j\) steps are marked; at least six are unmarked.  Choose one.

Finally, an acquisition step has one target section.  Dynamic terminal
cycles with different centres use different target sections \(P_j\), so
the choices made in the second alternative are automatically distinct.
\end{proof}

\begin{corollary}[Terminal remainders add no residual-shadow debt]
\label{cor:damp-terminal-remainders-no-residual-debt}
In the two-sweep assembly, every initial terminal incidence and every
dynamic terminal outcome can be paid by original section-shadow occurrence
columns and by successful-acquisition columns that are not committed to a
marked acquisition.  These choices can be made disjointly from the
fragmentation and failure payments in \(\Omega_0\).
\end{corollary}

\begin{proof}
If a progressive synchronization is terminal already in the initial
state, its terminal remainder \(D\) is disjoint from an ample synchronized
family containing the two repair antipodes.  Hence its base section has
at least \(|D|\ge9\) original nonbaseline occurrence columns.  There are
at most two oriented incidences based at one section, so choose distinct
columns for all such initial terminal demands and for the selected demand
on the other side, if present.

For a dynamic terminal outcome of an original tight two-cycle, the two
selected side demands had coalesced into one neutral component.  This is
the only weak component based at \(P_j\) that contains a selected initial
demand; components containing no selected demand need not be included in
the assembly demand family.  One original occurrence column pays the
selected component by Proposition
\ref{prop:damp-neutral-transfer-forest-section-excess}.  Proposition
\ref{prop:damp-dynamic-terminal-unmarked-acquisition} pays the possible
second-side deficit either with another original occurrence or with an
unmarked acquisition column.  The final sentence of that proposition
makes these latter choices distinct.  Original occurrence columns lift
through Proposition
\ref{prop:damp-support-occurrence-failure-joint-injection}, while
acquisition columns are the additional forest columns in Proposition
\ref{prop:damp-spacetime-support-forest}.  Neither type uses a member of
the residual shadow already assigned to a failure or fragmentation
generator.
\end{proof}

\begin{proposition}[Residual collisions have a bounded genealogy]
\label{prop:damp-transported-residual-bounded-genealogy}
Let
\[
 H_0\subset H_1\subset\cdots\subset H_q
\]
be a sequence of successful adjacent-section duplications, let
$\Omega_h$ be the residual shadow of $H_h$, and use the backward injections
of Proposition
\ref{prop:damp-residual-shadow-backward-transport}.  Interspersed with the
successful steps, process each oriented path incidence at most once.  Let
$\mathcal F$ be a collection of failed lifts, and for a failure across the
path edge $x_i$ choose one of the residual supports supplied by Proposition
\ref{prop:damp-adjacent-carrier-owner-residual-charge}.  Thus its chosen
support $R\in\Omega_h$ satisfies $x_i\in R$.

For a family of supports $\mathcal A$, put
\[
 M_X(\mathcal A)=\sum_{R\in\mathcal A}|R\cap X|.
\]
At a successful step $H_{h-1}\subset H_h$, call the backward injection
\emph{exceptional} if the unique new shattered support $N_h$ belongs to
$\Omega_h$ and is therefore sent to an old residual support displaced by
the enlarged path-spine family.  If $\mathcal X$ is the set of exceptional
steps, then
\begin{equation}
 |\mathcal F|
 \le2M_X(\Omega_0)+
       2\sum_{h\in\mathcal X}|N_h\cap X|.
 \label{eq:damp-transported-residual-genealogy-bound}
\end{equation}
Moreover,
\begin{equation}
 \sum_{h\in\mathcal X}|N_h\cap X|
 \le |E_X(H_q)|-|E_X(H_0)|
 \le \Lambda_0(\mathbf P^{\rm in})-
       \Lambda_0(\mathbf P^{\rm out})+2V,
 \label{eq:damp-exceptional-residual-core-edge-bound}
\end{equation}
where $E_X$ denotes the set of edges whose direction lies in $X$, and $V$
is the number of original concepts over core words outside the path.
Consequently,
\begin{equation}
 |\mathcal F|
 \le2M_X(\Omega_0)+2\Lambda_0(\mathbf P^{\rm in})+4V.
 \label{eq:damp-transported-residual-initial-state-bound}
\end{equation}
Thus failures at different extension times cannot reuse one original
residual support arbitrarily: all extra multiplicity is localized to new
residual-support segments, and their total core rank is paid by boundary
descent and the fixed external-core incidence budget.
\end{proposition}

\begin{proof}
Draw one node $(h,R)$ for every $R\in\Omega_h$ and join it backwards to
$(h-1,\iota_h(R))$.  Every node has one backward edge, and injectivity of
$\iota_h$ gives at most one forward edge at every node.  The resulting
graph is therefore a disjoint union of paths rooted at the members of
$\Omega_0$.  On every nonexceptional backward edge the support label is
unchanged.  An exceptional step introduces at most one new label, namely
$N_h$.  After compressing consecutive nodes with the same label, the paths
therefore consist of the $|\Omega_0|$ initial segments and at most one
further segment for each $h\in\mathcal X$.

Attach every failed lift to the node carrying its chosen residual support
at the state in which the failure occurs.  On a segment with fixed label
$R$, a failure across $x_i$ is possible only when $x_i\in R$.  There are
only two oriented incidences of the path edge $x_i$, and no incidence is
processed twice.  Hence at most $2|R\cap X|$ failures attach to that
segment.  Summing first over the initial segments and then over the
exceptional segments proves
\eqref{eq:damp-transported-residual-genealogy-bound}.

Every successful ample one-concept extension has a different unique new
shattered support.  The coordinatewise edge--shadow identity therefore
gives
\[
 |E_X(H_q)|-|E_X(H_0)|
 =\sum_{h=1}^q|N_h^{\rm new}\cap X|,
\]
where $N_h^{\rm new}$ denotes the unique new support at step $h$.
The exceptional supports form a subfamily of these new supports, proving
the first inequality in
\eqref{eq:damp-exceptional-residual-core-edge-bound}.

It remains to count the new core edges geometrically.  A duplicated path
concept creates one edge to its source section and possibly the edge to its
other path neighbour.  The proof of Proposition
\ref{prop:damp-adjacent-duplication-boundary-lyapunov} shows that the
corresponding decrease of $\Lambda_0$ is at least the number of these new
path edges: it is one when a boundary moves, and $n+2$ when two boundaries
annihilate.  Summing pays every new core edge whose two ends lie over the
path.

Every remaining new core edge joins the duplicated path concept to an
original concept whose core word lies outside the path.  A core word outside
a coordinate-simple cube path has at most two neighbours on that path: any
two such neighbours have path distance two, while bipartiteness excludes
consecutive path vertices, and three path neighbours are impossible.
There are $V$ original external concepts, so at most $2V$ such edges are
created.  This proves the second inequality in
\eqref{eq:damp-exceptional-residual-core-edge-bound}.  Substitute it into
\eqref{eq:damp-transported-residual-genealogy-bound} and discard the
nonnegative final potential to obtain
\eqref{eq:damp-transported-residual-initial-state-bound}.
\end{proof}

\begin{proposition}[Failed lifts have distinct original boundary ancestors]
\label{prop:damp-failed-lift-boundary-ancestor-injection}
Keep the dynamic setting of Proposition
\ref{prop:damp-transported-residual-bounded-genealogy}, and stop processing
an oriented incidence at its first failed global lift.  Then the failures
inject into the directed word-boundary occurrences of the initial section
tuple.  They also inject into the directed section-shadow boundary
occurrences of the initial tuple.  Consequently,
\begin{align}
 |\mathcal F|
 &\le \sum_{i=1}^n
       |P_{i-1}^{\rm in}\mathbin\triangle P_i^{\rm in}|
 \label{eq:damp-failed-lift-word-boundary-injection}\\
 &=\sum_{i=1}^n
       |\operatorname{Sh}(P_{i-1}^{\rm in})
          \mathbin\triangle
         \operatorname{Sh}(P_i^{\rm in})|
 \le2(|W|+|\Omega_0|).
 \label{eq:damp-failed-lift-shadow-boundary-injection}
\end{align}
More precisely, every failure carries two compatible labels: a distinct
initial directed boundary ancestor and a residual-support lineage from
Proposition~\ref{prop:damp-transported-residual-bounded-genealogy}.  Thus
two failures transported to the same original residual support are already
separated by their original boundary labels.
\end{proposition}

\begin{proof}
A directed word-boundary occurrence is a word $p$ present on one end of a
path edge and absent on the other, together with its orientation from the
present end toward the missing end.  Under a successful adjacent
duplication, exactly one of three things happens to the occurrence used by
that duplication: it leaves the path at an endpoint, it annihilates with
the oppositely oriented occurrence at the other edge of the target section,
or it is transported to that other edge.  In the last case the support word
and orientation are unchanged and its endpoint-distance weight decreases
by one.  No operation creates two descendants of one occurrence.  Tracing
backwards therefore injects every directed boundary occurrence in every
intermediate state into a unique initial occurrence.

A failed lift reserves the occurrence across which it was attempted.  That
occurrence cannot later be transported: transport across the same oriented
edge would require continuing the very incidence that was stopped at its
first failure.  An operation from the other edge of the missing section can
only annihilate it.  Hence two failures cannot reserve occurrences with the
same initial ancestor.  This proves
\eqref{eq:damp-failed-lift-word-boundary-injection}.

There is a parallel shadow statement.  Although the full lift fails, the
proposed duplication is ample in its two-layer conditioning.  Its target
section and its section overlap are ample one-concept extensions.  The
proof of Lemma~\ref{lem:damp-adjacent-duplication-shadow-transport}, applied
inside that conditioning, gives one common new support $T$ already present
in the source section.  Thus the failure reserves the directed
section-shadow boundary occurrence labelled by $T$.  Successful global
duplications transport or annihilate these shadow occurrences by Lemma
\ref{lem:damp-adjacent-duplication-shadow-transport}, again without
branching.  After a failure, moving the reserved occurrence across the same
oriented edge would require resuming the stopped incidence.  The same
backward-ancestor argument therefore gives an injection into the initial
shadow-boundary occurrences.

For each edge, the two-layer shadow-intersection identity and ampleness give
\begin{align*}
 |P_{i-1}\mathbin\triangle P_i|
 &=|P_{i-1}|+|P_i|-2|C_i|\\
 &=|\operatorname{Sh}(P_{i-1})|+
   |\operatorname{Sh}(P_i)|-2|\operatorname{Sh}(C_i)|\\
 &=|\operatorname{Sh}(P_{i-1})
       \mathbin\triangle\operatorname{Sh}(P_i)|.
\end{align*}
Finally apply Proposition
\ref{prop:damp-section-shadow-boundary-conservation} to the initial tuple
to obtain the last inequality in
\eqref{eq:damp-failed-lift-shadow-boundary-injection}.  Pairing this
boundary ancestor with the residual genealogy gives the final assertion.
\end{proof}

\begin{proposition}[The residual label of a failed lift is canonical]
\label{prop:damp-failed-lift-canonical-residual}
Keep the notation of a failed adjacent-section lift across the path edge
$a_{i-1}a_i$.  Thus a repair word $p$ is present in the source section,
absent from the target section, and adjoining its missing copy is ample in
the corresponding two-layer conditioning but not in the full current
family $\widehat H$.  Let $T$ be the unique repair support acquired by the
target section in the locally ample extension.  Then
\begin{equation}
 L(p,i):=T\cup\{x_i\}\in\widehat\Omega.
 \label{eq:damp-failed-lift-canonical-residual}
\end{equation}
In particular, every failed lift has a canonical residual payment whose
core part is exactly its transition coordinate.  No choice of a transverse
owner or of a support in a shadow difference is required.
\end{proposition}

\begin{proof}
Write $P_s,P_t$ for the source and target sections and
$C_i=P_s\cap P_t$.  The locally ample extension makes both
$P_t\cup\{p\}$ and $C_i\cup\{p\}$ ample one-concept extensions.  The
proof of Lemma~\ref{lem:damp-adjacent-duplication-shadow-transport}, which
uses only this two-layer conditioning, gives one common new repair support
$T$, with
\[
 T\in\operatorname{Sh}(P_s)\setminus\operatorname{Sh}(P_t),
 \qquad T\notin\operatorname{Sh}(C_i).
\]
In the two-layer extension, the repair projection does not change and the
section-shadow intersection gains exactly $T$.  Hence the unique new
strongly shattered support of that conditioning is $T\cup\{x_i\}$.

Suppose that $T\cup\{x_i\}$ were not shattered by $\widehat H$.  The
family obtained by adjoining the missing copy would then strongly shatter
all $|\widehat H|$ supports of the ample family $\widehat H$, together
with the new support $T\cup\{x_i\}$.  Its number of strongly shattered
supports would equal its number of concepts.  Equality in the lower
Sandwich inequality would make the full extension ample, contrary to the
definition of the failed lift.  Thus
$T\cup\{x_i\}\in\operatorname{Sh}(\widehat H)$.

The support is not in the path-spine family: it has exact core part
$\{x_i\}$, while its repair part $T$ is absent from
$\operatorname{Sh}(C_i)$.  It therefore belongs to
$\widehat\Omega$, proving
\eqref{eq:damp-failed-lift-canonical-residual}.
\end{proof}

\begin{proposition}[Failed lifts inject into the original residual shadow]
\label{prop:damp-canonical-failure-support-birth-injection}
Keep the dynamic two-sweep setting, and let $\mathcal F$ be the failed
lifts.  Give each failure its canonical support $L(p,i)$ from Proposition
\ref{prop:damp-failed-lift-canonical-residual}.  These supports are
pairwise distinct over the complete history, and every one of them belongs
to the original residual shadow $\Omega_0$.  Consequently the canonical
label map is an injection
\begin{equation}
 \mathcal F\hookrightarrow\Omega_0,
 \qquad |\mathcal F|\le|\Omega_0|.
 \label{eq:damp-canonical-failure-initial-bound}
\end{equation}
This removes the residual core moment, the word-boundary potential, and
all multiplicity loss from
\eqref{eq:damp-transported-residual-initial-state-bound}.
\end{proposition}

\begin{proof}
Fix a repair support $T$ and a path edge $x_i$.  At any state in which
$T\cup\{x_i\}$ is the canonical support of a failure, $T$ lies in the
shadow of exactly the source section across that edge.  Section shadows
only grow.  Hence this orientation can never reverse: a later acquisition
may erase the boundary, but cannot recreate it with the opposite
orientation.  The same oriented incidence is stopped at its first failure
and is not processed again.  Therefore $T\cup\{x_i\}$ labels at most one
failure over the whole history.

Suppose first that such a label already belongs to
$\operatorname{Sh}(H_0)$.  At the later failure one has
$T\notin\operatorname{Sh}(C_i)$.  Since carrier shadows only grow, this
was also true initially.  The label has no other core coordinate, so it
does not belong to any other block of the initial path-spine family.
Consequently it lies in $\Omega_0$.

It remains to exclude a canonical label absent from
$\operatorname{Sh}(H_0)$.  Consider the first state in which
$T\cup\{x_i\}$ is shattered.  Every change before a failure is an ample
one-concept extension, so this label is the unique new shattered support of
one earlier successful duplication.  Let $v$ be the concept added at that
birth step, and let $x_k$ be the path coordinate across which it was
duplicated.  The original copy of $v$ is an $x_k$-neighbour already present
before the step.  Lemma~\ref{lem:damp-one-concept-extension-test} says that
the unique new support consists of all directions incident with $v$ before
its addition.  Hence
\[
 x_k\in T\cup\{x_i\}.
\]
The set $T$ contains only repair coordinates, so $k=i$.

In the locally ample target-section extension, the repair directions
incident with $v$ are exactly $T$.  The full
$T\cup\{x_i\}$-cube through $v$ therefore makes the adjacent carrier
$C_i$ acquire $T$ at this birth step.  Thus
$T\cup\{x_i\}$ enters the path-spine family at its first appearance and,
because carrier shadows only grow, remains there forever.  This contradicts
Proposition~\ref{prop:damp-failed-lift-canonical-residual}, which places the
same support in the residual shadow at the later failure.  Every canonical
label must consequently belong to $\operatorname{Sh}(H_0)$, and the first
part of the proof puts it in $\Omega_0$.  Together with pairwise
distinctness, this proves \eqref{eq:damp-canonical-failure-initial-bound}.
\end{proof}

\begin{proposition}[Gap edges reserve the fragmentation payment]
\label{prop:damp-gap-edge-failure-fragmentation-reservation}
Keep the initial path-spine state and the dynamic two-sweep setting of
Proposition
\ref{prop:damp-canonical-failure-support-birth-injection}.  Fix
\(T\in\mathcal T\), and write the components of its initial section-index
set as
\[
 I_T=[\ell_1,r_1]\mathbin{\dot\cup}\cdots
     \mathbin{\dot\cup}[\ell_{c_T},r_{c_T}].
\]
For \(1\le h<c_T\), put
\[
 G_{T,h}=\{r_h+1,\ldots,\ell_{h+1}\},
 \qquad d_{T,h}=|G_{T,h}|=\ell_{h+1}-r_h,
\]
and define
\begin{equation}
 \Gamma_{T,h}
 =\{T\cup\{x_i\}:i\in G_{T,h}\},
 \qquad
 \Gamma_T=\mathop{\dot\bigcup}_{h=1}^{c_T-1}\Gamma_{T,h}.
 \label{eq:damp-gap-edge-residual-reservation}
\end{equation}
Then
\begin{equation}
 \Gamma_T\subseteq\Omega[T],
 \qquad
 |\Gamma_T|=\sum_{h=1}^{c_T-1}d_{T,h}
              =\lambda_T+c_T-1.
 \label{eq:damp-gap-edge-residual-reservation-size}
\end{equation}

Let \(\mathcal F_T\) be the failed lifts whose acquired repair support is
\(T\).  Trace their section-shadow boundaries back to the initial state as
in Proposition
\ref{prop:damp-failed-lift-boundary-ancestor-injection}.  Let
\(f_{T,h}\) be the number of failures descended from the two boundary
occurrences facing the \(h\)-th internal gap, and put
\begin{equation}
 b_T=
 \bigl|\{h:d_{T,h}=2\text{ and }f_{T,h}=2\}\bigr|.
 \label{eq:damp-tight-two-sided-failure-hole-count}
\end{equation}
Then \(0\le f_{T,h}\le2\) and
\begin{equation}
 |\mathcal F_T|+c_T-1
 \le |\Omega[T]|+b_T.
 \label{eq:damp-failure-fragmentation-reservation}
\end{equation}
In particular, canonical failures and the \(c_T-1\) fragmentation
generators admit a joint injection into \(\Omega[T]\) unless some
one-section gap has one failed descendant from each of its two boundary
occurrences.  No longer gap can obstruct this joint payment.
\end{proposition}

\begin{proof}
Fix an internal gap between \(a_{r_h}\) and \(a_{\ell_{h+1}}\).  The ample
projected carrier \(Z_T\) is isometric and contains both endpoints.  Every
shortest path between them uses each of the directions
\[
 x_{r_h+1},\ldots,x_{\ell_{h+1}}
\]
exactly once.  Hence \(Z_T\) has an \(x_i\)-edge for every
\(i\in G_{T,h}\), so \(\{x_i\}\in\operatorname{Sh}(Z_T)\).  The support
conservation identity
\eqref{eq:damp-repair-carrier-core-support-lift} gives
\(T\cup\{x_i\}\in\operatorname{Sh}(H)\).

The two endpoints of the path edge \(a_{i-1}a_i\) cannot both belong to
\(I_T\), because the edge lies between distinct components.  The
two-layer shadow-intersection identity therefore gives
\(T\notin\operatorname{Sh}(C_i)\).  Thus
\(T\cup\{x_i\}\) is absent from the path-spine family and belongs to
\(\Omega[T]\).  Different gap edges have different core coordinates,
which proves the disjoint inclusion in
\eqref{eq:damp-gap-edge-residual-reservation-size}.  Finally,
\[
 d_{T,h}
 =\ell_{h+1}-r_h
 =(\ell_{h+1}-r_h-1)+1;
\]
summing proves the cardinality formula.

Every failure with acquired support \(T\) descends from a distinct initial
directed \(T\)-boundary occurrence.  A boundary occurrence facing an
internal gap can move only through that gap: it either disappears at an
endpoint, annihilates with the occurrence coming from the other side, or
stops at its first failure.  It never branches.  Hence at most the two
boundary occurrences of one gap contribute failures, proving
\(f_{T,h}\le2\).  Failures descended from an outer boundary of \(I_T\)
have canonical labels outside \(\Gamma_T\); these labels are pairwise
distinct members of \(\Omega[T]\) by Proposition
\ref{prop:damp-canonical-failure-support-birth-injection}.

For each internal gap,
\begin{equation}
 f_{T,h}+1
 \le d_{T,h}
   +\mathbf 1_{\{d_{T,h}=2,\ f_{T,h}=2\}}.
 \label{eq:damp-one-gap-failure-fragmentation-reservation}
\end{equation}
Indeed \(d_{T,h}\ge2\), while \(f_{T,h}\le2\); the right-hand correction
is needed only when both inequalities are equalities.  Sum
\eqref{eq:damp-one-gap-failure-fragmentation-reservation} over the
\(c_T-1\) gaps, add the failures descended from outer boundaries, and use
the disjoint reserved family \(\Gamma_T\) together with their canonical
labels.  This gives
\eqref{eq:damp-failure-fragmentation-reservation}.
\end{proof}

\begin{proposition}[Tight support gaps reserve distinct core squares]
\label{prop:damp-tight-gap-core-square-reservation}
Keep the initial path-spine state and the dynamic two-sweep setting of
Proposition
\ref{prop:damp-gap-edge-failure-fragmentation-reservation}.  Let
\(\mathcal H\) be the set of pairs \((T,h)\) counted by \(b_T\).  Thus
the \(h\)-th gap of \(I_T\) has one missing section and both of its
directed boundary descendants end in failure.  If that missing section is
\(a_j\), put
\begin{equation}
 Q_{T,h}=\{x_j,x_{j+1}\}.
 \label{eq:damp-tight-gap-core-square-label}
\end{equation}
Then the map
\begin{equation}
 (T,h)\longmapsto Q_{T,h}
 \label{eq:damp-tight-gap-core-square-injection}
\end{equation}
is an injection into the original residual shadow \(\Omega_0\).  Its
image is disjoint from every exact-repair block
\[
 \Omega_0[T]=\{S\in\Omega_0:S\cap R=T\},
\]
\(T\in\mathcal T\).  Consequently all canonical failures and all support
fragmentation generators fit simultaneously in the original residual
shadow:
\begin{equation}
 |\mathcal F|+\sum_{T\in\mathcal T}(c_T-1)
 \le |\Omega_0|.
 \label{eq:damp-global-failure-fragmentation-complementarity}
\end{equation}
This is independent of the repair positions and antipodal blockers of the
failed lifts.
\end{proposition}

\begin{proof}
Fix \((T,h)\in\mathcal H\).  Write the two components adjacent to its
gap as
\[
 [\ell_h,r_h]\quad\text{and}\quad
 [\ell_{h+1},r_{h+1}].
\]
Since the gap has edge length two,
\(\ell_{h+1}=r_h+2\).  Put \(j=r_h+1\), so that
\[
 a_{j-1},a_{j+1}\in Z_T,
 \qquad a_j\notin Z_T
\]
by \eqref{eq:damp-projected-carrier-path-gap}.  The endpoints have Hamming
distance two.  The two possible middle vertices of a two-edge cube path
between them are \(a_j\) and
\[
 q_j=a_j^{\{x_j,x_{j+1}\}}.
\]
The projected carrier \(Z_T\) is ample and therefore isometric.  Since it
omits \(a_j\), it must contain \(q_j\).

Membership \(q_j\in Z_T\) means that some full positioned \(T\)-cube of
the original family has core word \(q_j\); in particular, some concept of
\(H_0\) projects to \(q_j\).  On the coordinate pair
\(\{x_j,x_{j+1}\}\), the three path vertices
\(a_{j-1},a_j,a_{j+1}\) realize three consecutive traces, while \(q_j\)
realizes the fourth.  The path core \(A\) occurs as a repair fibre of
\(H_0\), so these four traces all occur in \(H_0\).  Hence
\[
 Q_{T,h}=\{x_j,x_{j+1}\}\in\operatorname{Sh}(H_0).
\]
This pure two-core support does not belong to
\(\operatorname{Sh}(W)\), and every other block of the path-spine family
\(\mathcal B\) has exact core part \(\{x_i\}\).  Therefore
\(Q_{T,h}\in\Omega_0\).

It remains to prove distinctness.  A tight gap with missing section
\(a_j\) initially has its two directed \(T\)-boundaries on the incidences
directed into \(P_j\), one across \(x_j\) and one across \(x_{j+1}\).
Transporting either boundary would make \(P_j\) acquire \(T\).  Since both
neighbouring sections already shatter \(T\), that acquisition would
annihilate both boundaries; neither could then have a failed descendant.
Thus, when both descendants fail, they fail on precisely those two
oriented incidences.  The two-sweep protocol processes each oriented
incidence at most once and stops it at its first failure.

It follows that two pairs in \(\mathcal H\) with the same missing index
\(j\) use the same two failed lifts.  A failed lift has one unique acquired
repair support, so the two pairs have the same support \(T\), and then they
are the same gap.  Distinct pairs in \(\mathcal H\) consequently have
distinct missing indices.  Since the edge directions
\(x_1,\ldots,x_n\) of the isometric path \(A\) are pairwise distinct,
their labels in
\eqref{eq:damp-tight-gap-core-square-label} are distinct.  This proves
\eqref{eq:damp-tight-gap-core-square-injection}.

Every acquired section support is nonbaseline and hence belongs to
\(\mathcal T\).  The exact-repair blocks \(\Omega_0[T]\) are pairwise
disjoint, and every member of them has nonempty repair part \(T\), whereas
the labels in \eqref{eq:damp-tight-gap-core-square-injection} have empty
repair part.  Summing
\eqref{eq:damp-failure-fragmentation-reservation} over
\(T\in\mathcal T\), and then paying its total correction
\(\sum_Tb_T=|\mathcal H|\) with these distinct pure core-square labels,
gives \eqref{eq:damp-global-failure-fragmentation-complementarity}.
\end{proof}

\begin{proposition}[The support-interval complex separates occurrence and
failure payments]
\label{prop:damp-support-occurrence-failure-joint-injection}
Keep the initial path-spine state and the two-sweep failure family
\(\mathcal F\).  Put
\[
 \mathfrak O
 =\{(j,T):T\in\mathcal T,\ j\in I_T\}.
\]
Let \(\mathcal D\) be any finite demand family admitting an injection
\(\eta:\mathcal D\hookrightarrow\mathfrak O\).  Then there is an injection
\begin{equation}
 \mathcal D\mathbin{\dot\cup}\mathcal F
 \hookrightarrow \Xi\mathbin{\dot\cup}\Omega_0.
 \label{eq:damp-support-occurrence-failure-joint-injection}
\end{equation}
In particular,
\begin{equation}
 |\mathcal D|+|\mathcal F|
 \le|\Xi|+|\Omega_0|
 =\Delta(H;A,R).
 \label{eq:damp-support-occurrence-failure-joint-bound}
\end{equation}
\end{proposition}

\begin{proof}
For \(T\in\mathcal T\), let \(d_T\) be the number of demands whose
\(\eta\)-label has second coordinate \(T\).  Since \(\eta\) is injective,
\[
 d_T\le m_T.
\]
The exact-repair part of \(\Xi\) has precisely \(m_T-c_T+1\) members:
the pure repair support \(T\), which is the chosen root column of
\(\partial_T\), and the supports \(T\cup\{x_i\}\) on the internal edges of
the components of \(I_T\).  Assign
\[
 \min\{d_T,m_T-c_T+1\}
\]
of the \(T\)-labelled demands injectively to these supports.  The number
left is at most
\[
 \max\{0,d_T-(m_T-c_T+1)\}\le c_T-1.
\]

Proposition~\ref{prop:damp-tight-gap-core-square-reservation} gives one
joint injection of all failures and all
\(\sum_T(c_T-1)\) formal fragmentation generators into \(\Omega_0\).
Use the generators belonging to \(T\) for the remaining \(T\)-labelled
demands.  Different exact-repair parts give disjoint supports in \(\Xi\),
and the fragmentation injection is disjoint from the failure labels.
This proves \eqref{eq:damp-support-occurrence-failure-joint-injection}.
Finally apply
\eqref{eq:damp-support-interval-total-excess} to obtain
\eqref{eq:damp-support-occurrence-failure-joint-bound}.
\end{proof}

\begin{corollary}[Neutral forests and failed lifts share one exact budget]
\label{cor:damp-neutral-forest-failure-joint-injection}
For each interior section \(P_j\), let \(\kappa_j\) be the number of weak
components of its extended neutral-transfer digraph from Proposition
\ref{prop:damp-neutral-transfer-forest-section-excess}.  Then
\begin{equation}
 \sum_{j=1}^{n-1}\kappa_j+|\mathcal F|
 \le\Delta(H;A,R).
 \label{eq:damp-neutral-forest-failure-joint-injection}
\end{equation}
This is a joint injection statement, not the sum of the separate neutral
and failure estimates.
\end{corollary}

\begin{proof}
Because \(P_j\) is ample and contains the two fixed repair antipodes, its
baseline shadow is exactly
\(\mathcal L_R=\{\varnothing\}\cup\binom R1\), of order \(r+1\).
Consequently
\[
 |\{T\in\mathcal T:j\in I_T\}|
 =|P_j|-(r+1).
\]
Proposition
\ref{prop:damp-neutral-transfer-forest-section-excess} gives
\(\kappa_j\le|P_j|-(r+1)\).  For each \(j\), label its neutral weak
components injectively by distinct supports
\(T\in\mathcal T\cap\operatorname{Sh}(P_j)\).  Section indices distinguish
equal supports chosen in different sections, so these labels give an
injection of all neutral components into \(\mathfrak O\).  Apply
Proposition~\ref{prop:damp-support-occurrence-failure-joint-injection}.
\end{proof}

\begin{corollary}[The joint path budget has an anchored normal form]
\label{cor:damp-joint-path-budget-anchored-normal-form}
Keep the hypotheses and notation of Proposition
\ref{prop:damp-support-occurrence-failure-joint-injection}.  For every
demand \(d\in\mathcal D\), write
\(\eta(d)=(j,T)\), and choose a positioned full \(T\)-cube in the section
\(P_j\).  For every failed lift, use its canonical support
\(T\cup\{x_i\}\) from Proposition
\ref{prop:damp-failed-lift-canonical-residual} and choose a positioned
full cube with that support in the original family.  Then these anchored
demands admit one of the following two alternatives:
\begin{enumerate}
 \item they acquire pairwise distinct coface labels in
       \(\operatorname{Sh}(H)\setminus\Sigma_0\);
 \item at least two of them ascend to the same support and the same
       positioned cube of \(H\).
\end{enumerate}
This is a normal form for the chosen occurrence labels.  The positioned
cube chosen for a neutral demand need not contain the puncture from which
that demand arose.  Consequently this statement alone does not localize
the geometry of a neutral demand; the provenance-preserving statement for
the closed neutral components is given below.
If \(|H|<64\), the support \(L\) of such a packet satisfies
\begin{equation}
 2\le |L|\le4.
 \label{eq:damp-anchored-packet-rank-four}
\end{equation}
Thus the internal face flags of every unresolved packet live in one fixed
four-cube, independently of the path length and repair dimension.
\end{corollary}

\begin{proof}
Every \(T\in\mathcal T\) has at least two repair coordinates.  Ampleness
of \(P_j\) makes \(T\) strongly shattered there, so the required
positioned \(T\)-cube exists and lifts to \(H\) over the core word \(a_j\).
The acquired repair support of a failed lift also belongs to
\(\mathcal T\).  Its canonical support is a member of the original
shattered complex by Proposition
\ref{prop:damp-canonical-failure-support-birth-injection}; ampleness of
the original family supplies the chosen positioned cube.

Apply Proposition~\ref{prop:damp-anchored-coface-ascent} simultaneously
to this multiset of anchored faces.  Every initial support contains at
least two repair coordinates, and coface ascent never removes a
coordinate.  Hence every final label lies outside
\(\Sigma_0=\operatorname{Sh}(A)*\mathcal L_R\), whose members have at most
one repair coordinate.  The two alternatives in the anchored-coface
normal form are therefore exactly the alternatives displayed above.

Finally, every final support still contains the initial repair support
\(T\), so its order is at least two.  Extend \(L\) to a facet \(M\) of
\(\operatorname{Sh}(H)\).  The facet-excess inequality,
Theorem~\ref{thm:damp-facet-excess}, gives
\[
 |H|\ge2^{|M|+1}\ge2^{|L|+1}.
\]
If \(|H|<64\), this excludes \(|L|\ge5\) and proves
\eqref{eq:damp-anchored-packet-rank-four}.
\end{proof}

\begin{corollary}[Neutral cycles and canonical failure labels have an
anchored normal form]
\label{cor:damp-neutral-cycle-failure-provenance-normal-form}
Suppose that \(|H|<64\).
For every interior section \(P_j\), let \(\mathcal Z_j\) be the set of
directed cycles in its neutral-transfer digraph.  Keep also the failed-lift
family \(\mathcal F\) from the two-sweep construction.  One can associate
with every member of
\[
 \mathcal Z=\mathop{\dot\bigcup}_{j=1}^{n-1}\mathcal Z_j
 \qquad\text{and with every member of }\mathcal F
\]
a positioned full face of \(H\) so that the associated faces are pairwise
distinct and every one of their supports contains at least two repair
coordinates.  The face of a neutral cycle lies in its actual common
maximal repair cube.  The support of a failure face is its canonical
residual label; the position of that face is an auxiliary realization and
need not lie at the failed path edge.  After anchored coface ascent,
exactly one of the following holds:
\begin{enumerate}
 \item all members of \(\mathcal Z\mathbin{\dot\cup}\mathcal F\) have
       pairwise distinct labels in
       \(\operatorname{Sh}(H)\setminus\Sigma_0\);
 \item at least two members occupy one positioned cube whose support \(L\)
       satisfies
       \[
        3\le |L|\le4.
       \]
\end{enumerate}
Thus a neutral-cycle/canonical-failure packet below order \(64\) has rank
three or four.  In particular, the rank-two packet in the occurrence-label
normal form is an artefact of using arbitrary section occurrences and is
not an active obstruction.
\end{corollary}

\begin{proof}
Fix a positioned maximal repair cube \(Q\) of a section \(P_j\), let \(T\)
be its support, and put \(t=|T|\).  Proposition
\ref{prop:damp-neutral-transfer-two-sided-cycle} gives \(t\ge2\).  We first
assign the cycles based in \(Q\) distinct owner words.  Each cycle contains
at least two distinct repair words, different cycles are disjoint as
boundary occurrences, and a repair word has at most its left and right
occurrence.  Hence every collection of \(k\) cycles collectively contains
at least \(k\) distinct repair words.  Hall's theorem therefore assigns to
each cycle a distinct repair word occurring on that cycle.

We next match these owner words to distinct faces of \(Q\) of dimension at
least two, with every word lying on its assigned face.  For \(t\ge4\), use
the incidence graph between the vertices and the two-faces of \(Q\).  Its
left degree is \(\binom t2\), its right degree is four, and
\(\binom t2\ge4\).  Counting the incidence edges from an arbitrary vertex
set proves Hall's condition.  For \(t=3\), use the six two-faces together
with \(Q\) itself.  A set of at most four vertices has at least four
incident faces, whereas a set of at least five vertices meets every
two-face and hence has all seven faces as neighbours.  There cannot be
eight cycles based in \(Q\): equality would use both boundary occurrences
of every vertex of \(Q\).  All these vertices are corners whose unique
maximal cube is \(Q\), so connectedness would force \(P_j=Q\); but each
adjacent carrier is a nonempty subset of \(P_j\), contradicting the use of
the corresponding boundary occurrence at every vertex of \(Q\).  Thus
Hall's condition again applies to the owner words.
Finally, if \(t=2\), Corollary
\ref{cor:damp-repeated-neutral-cycle-common-cube-rank} gives at most one
cycle, which we assign to \(Q\) itself.  Lift every assigned face to \(H\)
over the core word \(a_j\).

These neutral faces are pairwise distinct.  Within one \(Q\), they are
distinct by the face matching.  Every owner word is a corner of \(P_j\)
whose unique maximal cube is \(Q\).  Hence a face through that owner cannot
also lie in a different positioned maximal cube of the same section.
Faces chosen in different sections have different core anchors.

For a failed lift, use a positioned full cube on its canonical support
\(T\cup\{x_i\}\).  Proposition
\ref{prop:damp-canonical-failure-support-birth-injection} makes these
supports pairwise distinct.  They cannot equal a neutral support, since
they contain one core coordinate, whereas every neutral support chosen
above is contained in \(R\).  This cube realizes the canonical support for
the coface matching; it is not asserted to contain the failed puncture.
Thus all associated faces are pairwise distinct, as claimed.

Apply Proposition~\ref{prop:damp-anchored-coface-ascent} to these faces.
Every final support contains its initial support and therefore has at least
two repair coordinates, so it lies outside
\(\Sigma_0=\operatorname{Sh}(A)*\mathcal L_R\).  Suppose that a terminal
packet had support \(L\) of order two.  The initial support \(K\) of each
member of the packet satisfies \(K\subseteq L\) and \(|K|\ge2\), hence
\(K=L\).  The containing positioned \(L\)-cube is then the initial face
itself.  Two members of the packet would therefore have had the same
initial anchored face, contrary to the preceding paragraph.  Hence every
packet has rank at least three.  The upper bound four follows from the
facet-excess argument in
\eqref{eq:damp-anchored-packet-rank-four}.
\end{proof}

\begin{lemma}[Root-flexible cube-support matching]
\label{lem:damp-root-flexible-cube-support-matching}
Let \(F\in\Damp\), and let \(r\in F\).  There is a bijection
\begin{equation}
 \mu_r:F\longrightarrow\operatorname{Sh}(F)
 \label{eq:damp-root-flexible-cube-support-matching}
\end{equation}
such that
\[
 \mu_r(r)=\varnothing,
 \qquad
 Q_{\mu_r(y)}(y)\subseteq F\quad(y\in F).
\]
The bijection need not be induced by a corner peeling ending at \(r\).
\end{lemma}

\begin{proof}
Take a corner peeling
\[
 F=F_1\supset F_2\supset\cdots\supset F_s=\{y_s\},
 \qquad F_{i+1}=F_i\setminus\{y_i\}.
\]
For \(i<s\), let \(S_i\) be the support of the unique maximal cube
through \(y_i\) in \(F_i\), and put \(S_s=\varnothing\).  Corner deletion
removes exactly the shattered support \(S_i\).  Thus
\[
 M=\{(y_i,S_i):1\le i\le s\}
\]
is a perfect matching in the bipartite incidence graph whose left side is
\(F\), whose right side is \(\operatorname{Sh}(F)\), and where
\(yS\) is an edge precisely when \(Q_S(y)\subseteq F\).

Delete the right vertex \(\varnothing\) and its incident matching edge.
The resulting matching saturates all nonempty shattered supports and
leaves \(y_s\) unmatched.  We claim that every left vertex is reachable
from \(y_s\) by an alternating path.  This follows by descending induction
on the peeling index.  Suppose that every \(y_j\) with \(j>i\) is
reachable.  The family \(F_i\) is connected and has at least two concepts,
so the unique maximal cube
\[
 C_i=Q_{S_i}(y_i)\subseteq F_i
\]
contains a concept \(z_i\ne y_i\).  Its peeling index is larger than
\(i\), so it is reachable by induction.  The edge \(z_iS_i\) is not in
the matching, whereas \(S_iy_i\) is.  Appending these two edges reaches
\(y_i\), proving the claim.

In particular, there is an alternating path from \(y_s\) to the prescribed
concept \(r\).  Switching the matching along this path leaves \(r\)
unmatched and still saturates every nonempty support.  Add the edge
\((r,\varnothing)\).  The resulting perfect matching is the bijection
\(\mu_r\) in~\eqref{eq:damp-root-flexible-cube-support-matching}.
\end{proof}

\begin{proposition}[Carriers ascend strictly to facets]
\label{prop:damp-carrier-strict-facet-ascent}
Let \(H\subseteq Q_E\) have minimum cardinality among all nonempty
cornerless ample families, and put \(\Sigma=\operatorname{Sh}(H)\).
Let \(K\in\Sigma\) be nonempty.  Every proper collection of positioned
\(K\)-cubes can be
matched to pairwise distinct strict cofaces of \(K\), with each cube
contained in a positioned cube on its assigned coface.

Consequently, let \(\mathcal A\) be a finite multiset of anchored faces
whose supports have order at least two.  By repeatedly replacing a face
with a containing positioned coface, one reaches one of the following two
alternatives:

\begin{enumerate}
 \item at least two members occupy the same positioned cube;
 \item all members have pairwise distinct support labels, and every label
       is a facet of \(\Sigma\).
\end{enumerate}
\end{proposition}

\begin{proof}
Fix \(K\).  The full carrier
\[
 \mathcal K_K(H)
 =\{y\in Q_{E\setminus K}:Q_K(y)\subseteq H\}
\]
is ample by Lemma
\ref{lem:damp-projected-cube-carrier-support-conservation}.  Its concepts
are the positioned \(K\)-cubes of \(H\), and the disjoint positioned
cubes give
\[
 2^{|K|}|\mathcal K_K(H)|\le |H|.
\]
Since \(K\ne\varnothing\), the carrier is smaller than \(H\), so
minimum-cardinality puts it in \(\Damp\).

Let \(\mathcal Q\) be a proper collection of its concepts and choose
\(y_\ast\in\mathcal K_K(H)\setminus\mathcal Q\).  Apply
Lemma~\ref{lem:damp-root-flexible-cube-support-matching} to the carrier
with prescribed root \(y_\ast\).  For every \(y\in\mathcal Q\), its
matched support \(S_y\) is nonempty, the supports \(S_y\) are pairwise
distinct, and the full \(S_y\)-cube of carrier anchors through \(y\)
lifts to a full \(K\mathbin{\dot\cup}S_y\)-cube of \(H\) containing
\(Q_K(y)\).  This gives the required matching to strict cofaces.

We now iterate this strict matching.  At a support \(K\), first stop if
two members occupy the same positioned \(K\)-cube.  Otherwise the members
labelled by \(K\) occupy a set \(\mathcal Q\) of distinct positions.  If
\(\mathcal Q\) is proper, move all of them to distinct strict cofaces by
the first part.  If every position is occupied, use the bijection in
Proposition~\ref{prop:damp-anchored-coface-ascent}: at most one member
keeps \(K\), and all the others move strictly upward.  Unless \(K\) is
already a facet, the carrier has at least two concepts, so this operation
strictly increases the sum of the support orders.  A remaining single
member at \(K\) then occupies a proper set of positions and moves strictly
at the next step.

The sum of the support orders is bounded by
\(|E||\mathcal A|\).  Hence the process terminates.  If it does not stop
with a same-position packet, every terminal label is a facet; labels are
pairwise distinct because a repeated terminal label at distinct positions
would still move, whereas repetition at one position is the first
alternative.  This proves the dichotomy.
\end{proof}

\begin{corollary}[Packet-free saturation exposes a corner]
\label{cor:damp-packet-free-saturation-exposes-corner}
Keep the path-core notation, and write
\[
 \Sigma_0=\operatorname{Sh}(A)*
 \bigl(\{\varnothing\}\cup\tbinom R1\bigr),
 \qquad
 \Delta=|\operatorname{Sh}(H)\setminus\Sigma_0|.
\]
Suppose that a saturated anchored demand family consists of exactly
\(\Delta\) faces, every initial support contains at least two repair
coordinates, and strict facet ascent produces no same-position packet.
Then \(H\) has a corner.  In particular, the packet-free equality branch
of the path-wide joint budget is impossible.
\end{corollary}

\begin{proof}
Proposition~\ref{prop:damp-carrier-strict-facet-ascent} assigns the
\(\Delta\) demands pairwise distinct facet labels.  Every label remains
outside \(\Sigma_0\), because coface ascent preserves the two initial
repair coordinates.  There are exactly \(\Delta\) supports outside
\(\Sigma_0\), so the labels exhaust
\(\Sigma\setminus\Sigma_0\).  Consequently every support outside the
baseline is a facet and contains at least two repair coordinates.

Suppose that \(L\in\Sigma\setminus\Sigma_0\) has order at least three.
Choose two repair coordinates \(s,t\in L\cap R\).  Downward closure gives
\(\{s,t\}\in\Sigma\setminus\Sigma_0\).  Both \(\{s,t\}\) and \(L\)
would then be facets, contrary to
\(\{s,t\}\subsetneq L\).  Hence every nonbaseline support is a
two-element subset of \(R\).

The maximum VC-one core satisfies
\(\operatorname{Sh}(A)=\{\varnothing\}\cup\binom X1\).  Thus every
member of \(\Sigma_0\) also has order at most two: besides the empty set
and the singletons, its only members are the mixed pairs
\(\{x,t\}\), \(x\in X\), \(t\in R\).  We have proved that
\(\operatorname{VCdim}(H)\le2\).  Apply Lemma
\ref{lem:damp-relative-peeling-vc-two} to the nested ample pair
\(\{h\}\subseteq H\), for any \(h\in H\).  It supplies a corner peeling
of \(H\), and in particular a corner, contradicting the choice of \(H\).
\end{proof}

\begin{lemma}[Rank-three packets have two local Hall obstructions]
\label{lem:damp-rank-three-packet-local-hall}
Let \(Q_L\) be a positioned full cube of \(H\), where \(|L|=3\), and let
\(\mathcal A\) be a family of pairwise distinct full faces of \(Q_L\),
each with support \(K\subseteq L\) of order at least two.  Put
\[
 c=\begin{cases}
    1,&Q_L\in\mathcal A,\\
    0,&Q_L\notin\mathcal A,
   \end{cases}
\]
and let \(d\) be the number of two-subsets \(K\in\binom L2\) for which
both opposite positioned \(K\)-faces of \(Q_L\) belong to \(\mathcal A\).
The faces in \(\mathcal A\) can be assigned pairwise distinct supports,
using for a \(K\)-face only \(K\) or \(L\), if and only if
\begin{equation}
 c+d\le1.
 \label{eq:damp-rank-three-packet-local-hall}
\end{equation}
Consequently every packet that cannot be separated inside its terminal
three-cube contains one of exactly two minimal configurations:
\begin{enumerate}
 \item the full \(L\)-cube together with the two opposite \(K\)-faces for
       one \(K\in\binom L2\);
 \item the two opposite \(K\)-faces and the two opposite \(K'\)-faces for
       distinct \(K,K'\in\binom L2\).
\end{enumerate}
\end{lemma}

\begin{proof}
For each two-subset \(K\), at most two members of \(\mathcal A\) have
support \(K\), namely the two opposite \(K\)-faces of \(Q_L\).  Assign one
such face, when present, the label \(K\).  Every second face with that
support must use \(L\), and \(Q_L\) itself must also use \(L\).  Thus the
displayed assignment succeeds exactly when the number \(c+d\) of faces
forced to use \(L\) is at most one.

The same count is also necessary.  A full \(L\)-face has only the label
\(L\) among the permitted supports.  Two opposite \(K\)-faces have the
common two-element label set \(\{K,L\}\), so at least one of them must use
\(L\).  The families forced to use \(L\) are disjoint, and only one face
can receive that label.  This proves
\eqref{eq:damp-rank-three-packet-local-hall}.  If the inequality fails,
then either \(c=1\) and some direction is doubled, which contains the first
configuration, or \(c=0\) and at least two directions are doubled, which
contains the second.
\end{proof}

\begin{corollary}[Core profiles of rank-three packets]
\label{cor:damp-rank-three-packet-core-profile}
Apply Lemma~\ref{lem:damp-rank-three-packet-local-hall} to the initial
faces of a packet from Corollary
\ref{cor:damp-neutral-cycle-failure-provenance-normal-form}.
\begin{enumerate}
 \item If \(L\subseteq R\), every member is a neutral-cycle face, and the
       two configurations in Lemma
       \ref{lem:damp-rank-three-packet-local-hall} are the only minimal
       local obstructions.
 \item If \(L=T\mathbin{\dot\cup}\{x_i\}\), where
       \(T\in\binom R2\), every neutral-cycle face has support \(T\), every
       failure face has support \(L\), and at most one failure occurs.
       The only minimal local obstruction is therefore one canonical
       failure representative together with the two opposite
       \(T\)-faces of its positioned \(L\)-cube.
\end{enumerate}
These cases exhaust the rank-three packets.
\end{corollary}

\begin{proof}
Every initial face contains at least two repair coordinates.  A neutral
face has support contained in \(R\), while a failure face has canonical
support consisting of at least two repair coordinates and exactly one core
coordinate.  If \(L\subseteq R\), no failure support can be contained in
\(L\), proving the first assertion.

Otherwise the three-element support \(L\) has exactly two repair
coordinates and one core coordinate, so write it as displayed.  A neutral
initial support contained in \(L\) must then equal \(T\), and a failure
initial support contained in \(L\) must equal \(L\).  Canonical failure
supports are pairwise distinct, so there is at most one such failure.  The
second minimal configuration of Lemma
\ref{lem:damp-rank-three-packet-local-hall} needs two different doubled
two-supports and is impossible.  Its first configuration has precisely the
form asserted.  No third core profile is possible because every initial
support has at least two repair coordinates.
\end{proof}

\begin{corollary}[Mixed rank-three packets separate]
\label{cor:damp-mixed-rank-three-packet-separates}
In the anchored normal form of Corollary
\ref{cor:damp-neutral-cycle-failure-provenance-normal-form}, every
rank-three packet whose support contains a core coordinate can be assigned
pairwise distinct labels using supports carried by its terminal cube.
Consequently every locally deficient rank-three packet has pure repair
support and contains one of the two neutral-cycle configurations in
Lemma~\ref{lem:damp-rank-three-packet-local-hall}.
\end{corollary}

\begin{proof}
By Corollary~\ref{cor:damp-rank-three-packet-core-profile}, a locally
deficient mixed packet would have support
\[
 L=T\mathbin{\dot\cup}\{x_i\},
 \qquad T\in\binom R2,
\]
and would consist of one failure representative of support \(L\) together
with the two opposite \(T\)-faces of its positioned \(L\)-cube.  The
failure occurred across the path edge of direction \(x_i\), and
Proposition~\ref{prop:damp-failed-lift-canonical-residual} gives
\begin{equation}
 T\notin\operatorname{Sh}(C_i).
 \label{eq:damp-mixed-rank-three-failure-carrier-absence}
\end{equation}

Each of the two \(T\)-faces is a neutral-cycle face and hence, by the
construction in Corollary
\ref{cor:damp-neutral-cycle-failure-provenance-normal-form}, lies over a
path-core word.  Write those words as \(a_h\) and \(a_k\).  The faces are
opposite in one positioned \(L\)-cube, so their core anchors differ exactly
in \(x_i\).  A coordinate-simple cube path flips every path direction
once; two of its vertices differ in only \(x_i\) precisely when they are
the two endpoints of its \(x_i\)-edge.  Therefore
\(\{a_h,a_k\}=\{a_{i-1},a_i\}\).

The two opposite faces of a positioned \(L\)-cube have the same repair
anchor outside \(T\).  Thus one positioned full \(T\)-cube belongs to both
sections \(P_{i-1}\) and \(P_i\), and hence to their intersection \(C_i\).
This gives \(T\in\operatorname{Sh}(C_i)\), contradicting
\eqref{eq:damp-mixed-rank-three-failure-carrier-absence}.  The mixed
minimal obstruction is impossible.  Lemma
\ref{lem:damp-rank-three-packet-local-hall} now assigns distinct local
labels to every mixed packet, while Corollary
\ref{cor:damp-rank-three-packet-core-profile} leaves only the two stated
pure-repair configurations.
\end{proof}

\begin{corollary}[A pure rank-three packet ascends or owns a maximal
three-cube]
\label{cor:damp-pure-rank-three-packet-maximal-owner}
Let a locally deficient rank-three packet from Corollary
\ref{cor:damp-mixed-rank-three-packet-separates} have pure repair support
\(L\).  All of its initial faces lie in one path section \(P_j\), and its
terminal positioned \(L\)-cube \(Q_L\) is contained in one common
positioned maximal repair cube \(Q\) of \(P_j\).  If the support of \(Q\)
properly contains \(L\), the packet has a common strict coface of rank four.
Consequently a deficient alternating component can stop at rank three only
when \(Q=Q_L\) is itself a maximal repair three-cube of \(P_j\).
\end{corollary}

\begin{proof}
Every member of a pure-support packet is a neutral-cycle face.  Such a face
was lifted over the path word of its owner section.  Since all faces belong
to the same positioned repair cube \(Q_L\), their core anchors agree, and
they therefore lie in one section \(P_j\).

Each chosen neutral face contains its assigned owner word \(p\), which is a
corner of \(P_j\).  Its unique positioned maximal cube \(Q(p)\) contains
that face.  It also contains the terminal cube \(Q_L\): both are cubes
through \(p\), and every cube through a corner is contained in its unique
maximal cube.  If \(p\) and \(p'\) are owners of two packet faces, then
\(Q_L\subseteq Q(p)\cap Q(p')\).  In particular \(p\in Q(p')\), so
uniqueness of the maximal cube through \(p\) gives
\(Q(p)=Q(p')\).  This proves the common-owner assertion.

Write \(M\) for the support of \(Q\).  If \(s\in M\setminus L\), the
positioned \(L\cup\{s\}\)-face of \(Q\) containing \(Q_L\) is a full
rank-four cube and contains every member of the packet.  It supplies the
claimed common strict coface.  If no rank-four ascent is available, then
\(M=L\); maximality of \(Q\) gives \(Q=Q_L\), as required.
\end{proof}

\begin{proposition}[An overloaded maximal repair three-cube pays an
exterior tax]
\label{prop:damp-maximal-repair-three-cube-exterior-tax}
Let a positioned maximal repair three-cube \(Q\) of an interior section
\(P_j\) be the common owner of \(c\) directed neutral cycles.  Then
\begin{equation}
 \Delta(H;A,R)\ge4+\max\{0,2c-8\}.
 \label{eq:damp-maximal-repair-three-cube-exterior-tax}
\end{equation}
In the rank-three motif branch \(c\ge3\), and
\begin{align*}
 c=3&\ \Longrightarrow\ \Delta(H;A,R)\ge c+1,\\
 c=4&\ \Longrightarrow\ \Delta(H;A,R)\ge c+1,\\
 c\ge5&\ \Longrightarrow\ \Delta(H;A,R)\ge2c-4\ge c+1.
\end{align*}
Thus every maximal-three-cube rank-three motif supplies at least one local
support unit beyond its cycle population.
\end{proposition}

\begin{proof}
The full three-cube \(Q\) contributes its three two-subsets and its
three-element support beyond the empty support and the repair singletons.
Hence
\[
 |P_j|-(r+1)\ge4.
\]

The \(c\) cycles are disjoint as vertices of the left--right occurrence
digraph and have length at least two, so they use at least \(2c\) boundary
occurrences.  The cube \(Q\) has eight repair words, and each word has at
most its left and right occurrence.  Consequently at least
\(\max\{0,2c-8\}\) repair words have both occurrences on these cycles.
Every such word \(p\) is a corner of \(P_j\) whose unique maximal repair
cube is \(Q\), and
\[
 p\notin C_j\cup C_{j+1}.
\]

Consider the concept \((p,a_j)\) of \(H\).  It has no core neighbour in
either adjacent path section.  If it had no core neighbour over a word
outside the path, then every cube of \(H\) through it would lie in the
section \(P_j\) and hence in \(Q\); it would be a corner of \(H\), contrary
to the standing hypothesis.  Thus every doubled boundary word forces a
concept of \(H\) over a core word outside the path.  Distinct repair words
force distinct concepts, and therefore
\[
 V\ge\max\{0,2c-8\}.
\]
There is one extra point when \(c=4\).  If none of the eight minimum cycle
occurrences shared an underlying repair word, they would use all eight
vertices of \(Q\).  Every vertex of \(Q\) would then be a corner of \(P_j\)
whose unique maximal cube is \(Q\).  No repair edge would leave \(Q\), and
connectedness would force \(P_j=Q\).  Since \(P_j\) contains the two fixed
repair antipodes, this would force \(r=3\) and make \(Q\) the full repair
cube.  Both antipodes belong to each adjacent carrier, so neither has a
left or right boundary occurrence, contradicting the use of all eight
words.  Therefore some word occurs on both sides, and the preceding
ambient-corner argument gives
\begin{equation}
 c=4\quad\Longrightarrow\quad V\ge1.
 \label{eq:damp-four-cycle-maximal-three-cube-exterior-unit}
\end{equation}

Substitute this and the section bound into the exact path-section
partition
\[
 \Delta(H;A,R)
 =V+\sum_{h=0}^n\bigl(|P_h|-(r+1)\bigr)
\]
to obtain
\eqref{eq:damp-maximal-repair-three-cube-exterior-tax}.  If \(c=3\), the
initial cube contribution four is already \(c+1\).  If \(c=4\), combine
the four cube supports with
\eqref{eq:damp-four-cycle-maximal-three-cube-exterior-unit}.  If \(c\ge5\),
the remaining inequalities are immediate.  Either local Hall motif uses
at least three cycles, so all possible cycle populations have at least one
unit of strict local tax.
\end{proof}

\begin{corollary}[A saturated maximal-three-cube motif has three
two-cycles]
\label{cor:damp-saturated-maximal-three-cube-three-two-cycles}
Suppose that the neutral-forest bound is tight in the section \(P_j\):
\[
 \kappa_j=|P_j|-(r+1).
\]
If a positioned maximal repair three-cube \(Q\) of \(P_j\) owns a
locally deficient rank-three packet, then it owns exactly three directed
neutral cycles.  They are directed two-cycles and use six distinct repair
words of \(Q\).
\end{corollary}

\begin{proof}
Corollary~\ref{cor:damp-neutral-forest-equality-two-word} makes every
directed neutral cycle in \(P_j\) a two-cycle.  In the notation of
\eqref{eq:damp-neutral-two-cycle-opposite-punctures}, one endpoint belongs
to \(C_j\setminus C_{j+1}\) and the other to
\(C_{j+1}\setminus C_j\).  Each endpoint consequently has only one
boundary occurrence, so two distinct two-cycles cannot share an underlying
repair word.  Since \(Q\) has eight vertices, it owns at most four such
cycles.

A locally deficient rank-three motif uses at least three cycles.  If there
were four, their eight distinct words would be all vertices of \(Q\).
Every vertex would then be a corner of \(P_j\) with unique maximal cube
\(Q\), so connectedness would force \(P_j=Q\).  As in the proof of
Proposition~\ref{prop:damp-maximal-repair-three-cube-exterior-tax}, the two
fixed repair antipodes would then lie in both adjacent carriers and could
not be boundary occurrences.  This is a contradiction.  Hence exactly
three cycles occur, and they use six distinct words.
\end{proof}

\begin{lemma}[Three opposite punctures force a short common pair]
\label{lem:damp-complementary-five-carrier-punctures}
Let \(Q\) be a three-cube and let \(D_-,D_+\subseteq Q\) be ample families
such that
\begin{equation}
 |D_-|=|D_+|=5,
 \qquad
 D_-\cup D_+=Q,
 \qquad
 |D_-\cap D_+|=2.
 \label{eq:damp-complementary-five-carrier-punctures}
\end{equation}
Suppose that the three words of \(D_-\setminus D_+\) can be paired with
the three words of \(D_+\setminus D_-\) so that every pair \((x,y)\), with
\(S=\operatorname{Diff}(x,y)\), satisfies
\begin{equation}
 \{x^U:U\subsetneq S\}\subseteq D_-,
 \qquad
 \{y^U:U\subsetneq S\}\subseteq D_+.
 \label{eq:damp-complementary-five-carrier-puncture-test}
\end{equation}
Then, after a cube automorphism and possibly interchanging the two sides,
there are directions \(a,b,e\) such that
\begin{equation}
 D_-=\{0,a,b,ab,e\},
 \qquad
 D_+=\{e,ae,be,abe,z\},
 \qquad
 z\in\{0,a,b\}.
 \label{eq:damp-complementary-five-carrier-short-normal-form}
\end{equation}
Here a word denotes the set of directions in which it differs from \(0\).
In particular, the two common words are \(e,z\), at distance one or two.
Consequently, if the two common words are opposite in \(Q\), at most two
disjoint pairs can satisfy
\eqref{eq:damp-complementary-five-carrier-puncture-test}.
\end{lemma}

\begin{proof}
An ample five-vertex subfamily of a three-cube uses all three coordinates:
five vertices cannot lie in one facet.  It therefore shatters the empty
set, all three singletons, and exactly one pair.  Thus it is a full square
with one vertex attached to a vertex of that square.  Normalize \(D_-\) as
\[
 D_-=\{0,a,b,ab,e\},
\]
where \(\{0,a,b,ab\}\) is its square and \(e\) is attached at \(0\).
Since the two families cover \(Q\), the three vertices outside \(D_-\),
namely \(ae,be,abe\), belong to \(D_+\).  The remaining two vertices of
\(D_+\) form the intersection \(I=D_-\cap D_+\).

There are only ten possible pairs \(I\in\binom{D_-}{2}\), and the ample
condition on \(D_+\) has a direct square--pendant interpretation.  The
complete list is
\[
\begin{array}{c|c|c}
I& D_+\text{ is ample}&
 \text{a three-pair matching satisfies
 \eqref{eq:damp-complementary-five-carrier-puncture-test}}\\ \hline
\{0,a\},\{0,b\},\{0,ab\},\{a,b\}&\text{no}&\text{no}\\
\{0,e\}&\text{yes}&\text{yes}\\
\{a,ab\},\{b,ab\}&\text{yes}&\text{no}\\
\{a,e\},\{b,e\}&\text{yes}&\text{yes}\\
\{ab,e\}&\text{yes}&\text{no}.
\end{array}
\]
For completeness, the three successful matchings are respectively
\begin{align*}
 I=\{0,e\}:&\quad
  a\leftrightarrow ae,\quad
  b\leftrightarrow be,\quad
  ab\leftrightarrow abe,\\
 I=\{a,e\}:&\quad
  0\leftrightarrow ae,\quad
  b\leftrightarrow be,\quad
  ab\leftrightarrow abe,\\
 I=\{b,e\}:&\quad
  0\leftrightarrow be,\quad
  a\leftrightarrow ae,\quad
  ab\leftrightarrow abe.
\end{align*}
Every failed ample row is checked by one proper flip: in each of
\(I=\{a,ab\},\{b,ab\},\{ab,e\}\), the exclusive word \(0\) has no
admissible partner in \(\{ae,be,abe\}\).  For example, pairing it with
\(ae\) requires both proper intermediate words \(a,e\) on the two
prescribed sides, and the analogous obstruction holds for \(be\) and
\(abe\).  In each of the four nonample rows, the displayed five-set
shatters at least two of the three coordinate pairs, in addition to the
empty set and the three singletons, so its shadow has order at least six.
This proves both the list and
\eqref{eq:damp-complementary-five-carrier-short-normal-form}.  In the last
row the common words \(ab,e\) are opposite, which proves the final
assertion.
\end{proof}

\begin{proposition}[Three neutral two-cycles force six section units]
\label{prop:damp-three-two-cycles-six-section-units}
Let \(Q\) be a positioned maximal repair three-cube in an interior
section \(P_j\).  If \(Q\) owns three directed neutral two-cycles whose
six underlying repair words are distinct, then
\begin{equation}
 |P_j|-(r+1)\ge6.
 \label{eq:damp-three-two-cycles-six-section-units}
\end{equation}
In particular, neither the four nonbaseline supports of \(Q\) nor those
four supports together with one additional section support can realize
this three-cycle profile.
\end{proposition}

\begin{proof}
Write \(T\) for the support of \(Q\), so \(|T|=3\), and suppose first for a
contradiction that
\begin{equation}
 |P_j|-(r+1)=4.
 \label{eq:damp-three-two-cycles-four-unit-assumption}
\end{equation}
The section \(P_j\) contains the two fixed repair antipodes.  It therefore
shatters the empty set and all \(r\) repair singletons.  The full cube
\(Q\) supplies the three pairs in \(\binom T2\) and \(T\) itself.  These
are four further supports, so ampleness and
\eqref{eq:damp-three-two-cycles-four-unit-assumption} give the exact
identity
\begin{equation}
 \operatorname{Sh}(P_j)
 =\{\varnothing\}\cup\binom R1\cup\binom T2\cup\{T\}.
 \label{eq:damp-three-two-cycles-four-unit-shadow}
\end{equation}

We first extract the graph forced by this equality.  For every
\(s\in R\setminus T\), no square of \(P_j\) uses \(s\).  Hence every
\(s\)-edge is a bridge by Lemma~\ref{lem:damp-square-free-edge-bridge}.
There is exactly one such edge.  Indeed, there is at least one because the
repair antipodes differ in \(s\).  If there were two, paths between their
same-side endpoints inside the two convex \(s\)-semicubes would form a
cycle containing both edges, contradicting that they are bridges.

Thus \(P_j\) has exactly \(r-3\) edges whose directions lie outside \(T\).
On the other hand,
\eqref{eq:damp-three-two-cycles-four-unit-assumption} gives
\[
 |P_j\setminus Q|=(r+5)-8=r-3.
\]
Deleting the \(r-3\) bridge edges produces \(r-2\) connected components.
There are only the cube \(Q\) and \(r-3\) vertices outside it, so equality
forces every outside component to be a single vertex.  Consequently the
one-inclusion graph of \(P_j\) is obtained from \(Q\) by attaching a tree
of \(r-3\) bridge edges and no other edges.

Every one of these bridges separates the two repair antipodes, since they
have opposite values in its direction.  After contracting \(Q\), every
edge of the attachment tree therefore lies on the path between the two
antipodes.  The tree is that path.  In \(P_j\), it consists of two
possibly empty pendant paths attached to vertices \(u,v\in Q\).  A
geodesic between the repair antipodes uses all \(r-3\) bridge directions
and has total length \(r\); hence its segment in \(Q\) has length three.
Thus \(u\) and \(v\) are opposite vertices of \(Q\).

Both adjacent carriers \(C_j\) and \(C_{j+1}\) are connected and contain
the repair antipodes.  They must therefore contain both pendant paths and
the vertices \(u,v\).  The endpoints of each neutral two-cycle belong to
opposite exclusive carrier sides by
\eqref{eq:damp-neutral-two-cycle-opposite-punctures}.  In particular, none
of the six cycle words is \(u\) or \(v\).  They are distinct by hypothesis,
so they are precisely the other six vertices of \(Q\), three on each
exclusive side.  Put
\[
 D_-=C_j\cap Q,
 \qquad
 D_+=C_{j+1}\cap Q.
\]
Repeated coordinate restrictions preserve ampleness, and therefore
\(D_-\) and \(D_+\) are ample five-vertex subfamilies of \(Q\) satisfying
\begin{equation}
 D_-\cup D_+=Q,
 \qquad
 D_-\cap D_+=\{u,v\}.
 \label{eq:damp-three-two-cycles-complementary-carriers}
\end{equation}

The two words of every neutral two-cycle satisfy
\eqref{eq:damp-complementary-five-carrier-puncture-test} by
Proposition~\ref{prop:damp-neutral-transfer-two-sided-cycle}.  Since
\(u,v\) are opposite, the final assertion of Lemma
\ref{lem:damp-complementary-five-carrier-punctures} permits at most two
such disjoint pairs, contradicting the three cycles.  Thus
\eqref{eq:damp-three-two-cycles-four-unit-assumption} is impossible.

It remains to exclude section excess five.  In that case the shadow in
\eqref{eq:damp-three-two-cycles-four-unit-shadow} has exactly one further
member \(S\).  Downward closure forces \(S\) to be a two-set not contained
in \(T\).  Moreover \(P_j\) contains exactly one positioned \(S\)-square:
the full \(S\)-cube carrier is ample, and its shadow is the link of \(S\),
which here consists only of the empty set.

Every direction outside \(T\cup S\) belongs to no square.  Its unique edge
is a bridge, by the argument above, and both adjacent carriers contain
that bridge because they contain the repair antipodes.  Contract all these
bridge directions.  Contraction preserves ampleness, the full cube \(Q\),
the two adjacent carriers, and the six neutral-cycle words.  None of those
words is incident with a contracted bridge, since it is a corner whose
unique maximal cube is \(Q\).  We obtain one of two reduced configurations.

Suppose first that \(S=\{t,e\}\), where \(t\in T\) and \(e\notin T\).
The reduced family has ten vertices: the eight vertices of \(Q\) and two
others.  Its unique \(S\)-square must therefore share a \(t\)-edge with
\(Q\), and its other two vertices are precisely the vertices outside
\(Q\).  Exactly one repair antipode \(u\) belongs to \(Q\).  The other
antipode belongs to the opposite edge of this square, so the shared
\(t\)-edge does not contain \(u\): its two fixed \(T\setminus\{t\}\)
coordinates are opposite to those of \(u\).  The two vertices of the
shared edge lie in both \(Q\) and the \(S\)-square and hence are not
corners.  The word \(u\) lies in both adjacent carriers.  Thus at most
\(8-2-1=5\) vertices of \(Q\) can be distinct exclusive carrier corners,
whereas the three cycles require six.  This is impossible.

Suppose next that \(S=\{e,f\}\subseteq R\setminus T\).  The reduced
family has eleven vertices.  Its unique \(S\)-square shares one vertex
\(v\) with \(Q\) and contributes the other three; otherwise the order
would be at least twelve.  The two repair antipodes show that the vertex
\(u\in Q\) belonging to one antipode and the attachment vertex \(v\) are
opposite in \(Q\).  The square is the only connection from \(Q\) to the
other repair antipode.  Hence both adjacent carriers contain \(u\) and
\(v\).  The six distinct cycle words are consequently exactly the other
six vertices of \(Q\), three exclusive to each carrier.  Their
intersections with \(Q\) are ample five-vertex families satisfying
\eqref{eq:damp-complementary-five-carrier-punctures}.  Since their two
common words are opposite, Lemma
\ref{lem:damp-complementary-five-carrier-punctures} permits only two
opposite-puncture pairs, a contradiction.

Both possible forms of \(S\) are excluded.  The section excess is neither
four nor five, and \eqref{eq:damp-three-two-cycles-six-section-units}
follows.
\end{proof}

\begin{proposition}[A three-cycle attachment is short or costs seven]
\label{prop:damp-three-cycle-short-or-seven}
Let \(Q\) be a positioned maximal repair three-cube in an interior
section \(P_j\), and suppose that it owns three directed neutral
two-cycles on six distinct repair words.  Let \(u,v\) be the two vertices
of \(Q\) not used by those cycles.  Then at least one of the following
alternatives holds:
\begin{enumerate}
 \item
 \begin{equation}
  |P_j|-(r+1)\ge7;
  \label{eq:damp-three-cycle-seven-unit-alternative}
 \end{equation}
 \item both adjacent carriers contain \(u,v\), and their restrictions to
       \(Q\) have the short common-pair normal form
       \eqref{eq:damp-complementary-five-carrier-short-normal-form}.
       In particular, \(\operatorname{dist}_Q(u,v)\in\{1,2\}\).
\end{enumerate}
Consequently, if the section excess is six, the second alternative is the
only remaining maximal-three-cube profile.
\end{proposition}

\begin{proof}
The six cycle words are corners of \(P_j\) whose unique maximal cube is
\(Q\).  Hence no edge of \(P_j\) leaves \(Q\) at any of those words.

Suppose first that one adjacent carrier, say \(C_j\), omits one of \(u,v\).
The carrier is ample, contains the two fixed repair antipodes, and is
therefore isometric.  Choose an antipodal repair geodesic
\(G\subseteq C_j\), so \(|G|=r+1\).  Its intersection with \(Q\) has at
most one vertex.  Indeed, a coordinate-simple cube geodesic meets a fixed
subcube in a contiguous subpath: for every coordinate fixed by \(Q\), the
times at which its value agrees with the fixed value form an initial or a
terminal interval, and their intersection is an interval.  If this
subpath were nontrivial, each of its endpoints would be \(u\) or \(v\).
For an internal endpoint of \(G\), this follows because its next edge
leaves \(Q\); for an endpoint of \(G\), it follows because a repair
antipode belongs to both adjacent carriers and hence is not one of the six
exclusive cycle words.  The two endpoints of a nontrivial subpath are
distinct, so \(C_j\) would contain both \(u,v\), contrary to the choice of
\(C_j\).  Therefore
\[
 |P_j|
 \ge |Q\cup G|
 \ge8+(r+1)-1=r+8,
\]
which is \eqref{eq:damp-three-cycle-seven-unit-alternative}.

It remains that both carriers contain \(u,v\).  Each cycle contributes one
word exclusively to \(C_j\) and its mate exclusively to \(C_{j+1}\).
Thus
\[
 D_-=C_j\cap Q,
 \qquad
 D_+=C_{j+1}\cap Q
\]
are ample five-vertex restrictions, cover \(Q\), and intersect exactly in
\(\{u,v\}\).  Their three exclusive pairs satisfy
\eqref{eq:damp-complementary-five-carrier-puncture-test} by
Proposition~\ref{prop:damp-neutral-transfer-two-sided-cycle}.  Lemma
\ref{lem:damp-complementary-five-carrier-punctures} gives the asserted
short common-pair normal form and the distance conclusion.
\end{proof}

\begin{lemma}[Two-square attachment antipode dichotomy]
\label{lem:damp-two-square-attachment-antipode-dichotomy}
Let \(F\subseteq Q_R\) be ample and contain two antipodal words
\(\alpha,\beta\).  Suppose that \(T\subseteq R\) has order three, that
\(Q\subseteq F\) is a positioned full \(T\)-cube, and that
\begin{equation}
 \operatorname{Sh}(F)
 =
 \{\varnothing\}\cup\binom R1\cup\binom T2\cup\{T,S_1,S_2\},
 \label{eq:damp-two-square-attachment-shadow}
\end{equation}
where \(S_1,S_2\) are distinct pairs not contained in \(T\) and
\begin{equation}
 R=T\cup S_1\cup S_2.
 \label{eq:damp-two-square-attachment-active-coordinates}
\end{equation}
If at least six vertices of \(Q\) are corners of \(F\) whose unique
maximal cube is \(Q\), put
\[
 N=Q\setminus\operatorname{Cor}(F).
\]
Then \(1\le|N|\le2\), and the following alternatives hold.
\begin{enumerate}
 \item If exactly one of \(\alpha,\beta\), say \(\alpha\), belongs to
       \(Q\), then \(\alpha\notin N\).  If moreover \(|N|=1\), its
       unique word is opposite to \(\alpha\) in \(Q\).
 \item If neither antipode belongs to \(Q\), then \(N\) consists of two
       opposite words of \(Q\).
\end{enumerate}
\end{lemma}

\begin{proof}
For \(i\in\{1,2\}\), the full \(S_i\)-cube carrier is ample and its
shadow is the link of \(S_i\) in
\(\operatorname{Sh}(F)\).  That link contains only the empty set.
Consequently there is a unique positioned \(S_i\)-square; denote it by
\(A_i\).  The same carrier argument shows that the positioned
\(T\)-cube is unique.  For \(U\in\binom T2\), the link of \(U\) consists
of the empty set and the singleton \(T\setminus U\).  Its carrier
therefore has two vertices, supplied by the two opposite \(U\)-faces of
\(Q\).  Thus those six faces and \(A_1,A_2\) are all the squares of
\(F\).

Every coordinate in \(R\) belongs to a shattered pair.  Its full
edge-carrier is therefore a connected ample family with at least two
vertices.  No carrier vertex is isolated, so every edge of \(F\) lies in
a square.  The only squares are the six faces of \(Q\) and the two
squares \(A_1,A_2\).  Since \(F\) is connected,
\begin{equation}
 F=Q\cup A_1\cup A_2,
 \label{eq:damp-two-square-attachment-block-cover}
\end{equation}
and the intersection graph of these three cube blocks is connected.
Moreover,
\begin{equation}
 N=Q\cap(A_1\cup A_2),
 \qquad 1\le|N|\le2.
 \label{eq:damp-two-square-attachment-footprint}
\end{equation}
The set \(R\setminus T\) is nonempty, so \(Q\) contains at most one of
the two antipodes.

Suppose first that \(\alpha\in Q\).  We claim that
\(\alpha\notin A_1\cup A_2\).  Otherwise, after interchanging the
squares, assume \(\alpha\in A_1\).  The other antipode cannot lie in
\(Q\) or \(A_1\), and hence belongs to \(A_2\).  The two squares are
disjoint: if \(z\in A_1\cap A_2\), then
\[
 R=\operatorname{Diff}(\alpha,\beta)
 \subseteq S_1\cup S_2,
\]
whereas two pairs not contained in the three-set \(T\) cover at most two
coordinates of \(T\).  Connectivity of the block-intersection graph now
forces both squares to meet \(Q\).  Their intersections with \(Q\) are
disjoint and together have at most two vertices by
\eqref{eq:damp-two-square-attachment-footprint}.  Thus each intersection
is a singleton, so \(S_1,S_2\) are disjoint from \(T\).  If
\(y\in A_2\cap Q\), then
\[
 \operatorname{Diff}(\alpha,\beta)\subseteq T\cup S_2.
\]
The distinct pure pairs \(S_1,S_2\) leave a coordinate of
\(S_1\) outside \(T\cup S_2\), another contradiction.  This proves
\(\alpha\notin N\).

Assume in addition that \(N=\{w\}\).  The two squares cannot both meet
\(Q\): both intersections would be the singleton \(w\), so both supports
would be pure, and whichever square contains \(\beta\) would give
\(\operatorname{Diff}(\alpha,\beta)\subseteq T\cup S_i\), missing a
coordinate of the other pure pair.  Hence, after relabelling,
\[
 Q\cap A_1=\{w\},\qquad Q\cap A_2=\varnothing,
 \qquad A_1\cap A_2\ne\varnothing.
 \label{eq:damp-two-square-attachment-one-footprint-chain}
\]
In particular, \(S_1\cap T=\varnothing\).  Put
\[
 k=|S_2\cap T|,
 \qquad h=|S_1\cap S_2|.
\]
Here \(k,h\in\{0,1\}\), and
\[
 |R|=7-k-h.
\]
Since a nonempty intersection of two positioned cubes is a cube on the
common support, \(|A_1\cap A_2|=2^h\).  On the other hand, ampleness and
\eqref{eq:damp-two-square-attachment-shadow} give
\(|F|=|R|+7\), while inclusion--exclusion in
\eqref{eq:damp-two-square-attachment-one-footprint-chain} gives
\[
 |F|=8+4+4-1-2^h.
\]
Therefore
\[
 2^h=1+k+h,
\]
which forces \(k=0\).  Both square supports are pure.  The antipode
\(\beta\) belongs to \(A_2\), since membership in \(A_1\) would miss a
coordinate of \(S_2\setminus S_1\).  Choose
\(z\in A_1\cap A_2\).  The words \(w,z,\beta\) agree on \(T\), whereas
\(\alpha\) and \(\beta\) differ on all three coordinates of \(T\).
Thus \(w\) is opposite to \(\alpha\) in \(Q\).

Finally suppose that neither antipode belongs to \(Q\).  They cannot
belong to the same square, so, after relabelling,
\(\alpha\in A_1\) and \(\beta\in A_2\).  The squares are disjoint, for a
word in their intersection would again imply
\(R\subseteq S_1\cup S_2\).  Their block-intersection graph can therefore
be connected only if both squares meet \(Q\).  The two intersections are
nonempty and disjoint, and their union is \(N\), which has order at most
two.  Hence
\[
 Q\cap A_1=\{x\},\qquad Q\cap A_2=\{y\},
\]
and both square supports are disjoint from \(T\).  The pairs
\(\alpha,x\) and \(\beta,y\) agree on \(T\).  Since the antipodes differ
on all three \(T\)-coordinates, \(x,y\) are opposite in \(Q\).  This
proves the second alternative.
\end{proof}

\begin{proposition}[Three neutral two-cycles force seven section units]
\label{prop:damp-three-two-cycles-seven-section-units}
Let \(Q\) be a positioned maximal repair three-cube in an interior
section \(P_j\).  If \(Q\) owns three directed neutral two-cycles whose
six underlying repair words are distinct, then
\begin{equation}
 |P_j|-(r+1)\ge7.
 \label{eq:damp-three-two-cycles-seven-section-units}
\end{equation}
\end{proposition}

\begin{proof}
Proposition~\ref{prop:damp-three-two-cycles-six-section-units} gives a
lower bound of six.  Suppose for a contradiction that equality holds.
Besides the empty set and the repair singletons, the four supports
\(\binom T2\cup\{T\}\) of \(Q\) leave room for exactly two further
shattered supports.  Downward closure makes them distinct pairs
\(S_1,S_2\) not contained in \(T\).

If a repair coordinate lies outside \(T\cup S_1\cup S_2\), it belongs to
no shattered pair.  Its edge-carrier is a singleton, so there is one
edge in that direction and the edge is a bridge.  Contract all such
directions.  Each contraction lowers both the order and the dimension by
one and preserves the section excess, the owner cube, the adjacent
carriers, and the six cycle words.  None of those six words is incident
with a contracted edge, because its unique maximal cube is \(Q\).
We obtain an ample family \(F\) satisfying
\eqref{eq:damp-two-square-attachment-shadow} and
\eqref{eq:damp-two-square-attachment-active-coordinates}.  The six cycle
words remain corners of \(F\) with unique maximal cube \(Q\).

Let \(u,v\) be the two unused words of \(Q\).  Equality excludes the
first alternative of
Proposition~\ref{prop:damp-three-cycle-short-or-seven}.  Hence both
adjacent carriers contain \(u,v\), and
\begin{equation}
 \operatorname{dist}_Q(u,v)\in\{1,2\}.
 \label{eq:damp-six-unit-short-common-pair}
\end{equation}
Every repair antipode belongs to both carriers, so an antipode in \(Q\)
must be one of \(u,v\).

Apply Lemma~\ref{lem:damp-two-square-attachment-antipode-dichotomy} and
write \(N=Q\setminus\operatorname{Cor}(F)\).  If one antipode
\(\alpha\) lies in \(Q\), then \(\alpha\) is a corner.  When
\(|N|=2\), the cube has only six corners, one of which is \(\alpha\);
there are not six further corners available for the cycle words.  When
\(|N|=1\), the cycle words are exactly the six vertices outside
\(\{\alpha\}\cup N\).  Thus \(\{u,v\}=\{\alpha\}\cup N\), and the lemma
makes \(u,v\) opposite, contradicting
\eqref{eq:damp-six-unit-short-common-pair}.

If neither antipode lies in \(Q\), the lemma makes \(N\) a pair of
opposite words.  The six cycle words are precisely \(Q\setminus N\), so
\(\{u,v\}=N\), again contradicting
\eqref{eq:damp-six-unit-short-common-pair}.  These cases are exhaustive
because \(Q\) contains at most one repair antipode.  Section excess six
is impossible, proving
\eqref{eq:damp-three-two-cycles-seven-section-units}.
\end{proof}

\begin{corollary}[A saturated maximal three-cube has four further
components]
\label{cor:damp-saturated-maximal-three-cube-four-further-components}
Under the hypotheses of Corollary
\ref{cor:damp-saturated-maximal-three-cube-three-two-cycles}, one has
\begin{equation}
 |P_j|-(r+1)=\kappa_j\ge7.
 \label{eq:damp-saturated-maximal-three-cube-four-further-components}
\end{equation}
Thus besides the three two-cycles owned by \(Q\), the section contains at
least four further weak components, all using two underlying repair
words.  The six-unit short-common-pair profile does not belong to the
global equality branch.
\end{corollary}

\begin{proof}
The three cycles are directed two-cycles on six distinct words by
Corollary~\ref{cor:damp-saturated-maximal-three-cube-three-two-cycles}.
Proposition~\ref{prop:damp-three-two-cycles-seven-section-units} gives the
lower bound in the display, while section saturation gives its equality
on the left.  Three of the \(\kappa_j\) components are the cycles owned by
\(Q\), so at least four others remain.  Corollary
\ref{cor:damp-neutral-forest-equality-two-word} makes every component in
a saturated section a two-word component.
\end{proof}

\begin{lemma}[Coface matching is controlled by upper ideals]
\label{lem:damp-coface-matching-upper-ideals}
Let \(Q_L\) be a positioned cube, let
\(\mathcal P\subseteq2^L\) be an upper ideal, and let
\(\mathcal A\) be a family of pairwise distinct full faces of \(Q_L\)
whose supports belong to \(\mathcal P\).  For \(K\in\mathcal P\), write
\[
 m(K)=|\{A\in\mathcal A:\operatorname{supp}(A)=K\}|.
\]
Each face of support \(K\) may use any label
\(J\in\mathcal P\) with \(K\subseteq J\).  The faces admit pairwise
distinct labels if and only if
\begin{equation}
 \sum_{K\in\mathcal U}m(K)\le|\mathcal U|
 \qquad
 \text{for every upper ideal \(\mathcal U\subseteq\mathcal P\)}.
 \label{eq:damp-coface-matching-upper-ideal-criterion}
\end{equation}
\end{lemma}

\begin{proof}
Necessity follows by considering all faces whose supports belong to
\(\mathcal U\).  Every allowed label for such a face remains in
\(\mathcal U\), because \(\mathcal U\) is an upper ideal.

Conversely, let \(\mathcal B\subseteq\mathcal A\), and let
\(\mathcal U\) be the upward closure in \(\mathcal P\) of the supports
of the faces in \(\mathcal B\).  This is exactly the set of labels
available to at least one member of \(\mathcal B\).  Hence
\[
 |\mathcal B|
 \le\sum_{K\in\mathcal U}m(K)
 \le|\mathcal U|.
\]
Hall's theorem now gives a matching of all faces to distinct labels.
\end{proof}

\begin{corollary}[Weighted-graph criterion inside a four-cube]
\label{cor:damp-rank-four-weighted-graph-hall}
Let \(|L|=4\) and
\(\mathcal P=\{K\subseteq L:|K|\ge2\}\).  In the notation of
Lemma~\ref{lem:damp-coface-matching-upper-ideals}, put
\[
 q=m(L),\qquad
 t_v=m(L\setminus\{v\}),\qquad
 p_{uv}=m(L\setminus\{u,v\})
\]
for \(v\in L\) and \(\{u,v\}\in\binom L2\).  Then the faces can be
assigned distinct coface labels if and only if
\begin{equation}
 q-1+\sum_{v\in W}(t_v-1)
 +\sum_{\{u,v\}\in\binom W2}(p_{uv}-1)_+
 \le0
 \qquad(W\subseteq L).
 \label{eq:damp-rank-four-weighted-graph-hall}
\end{equation}
Thus local rank-four deficiency is a positive induced-subgraph weight,
not a list of positioned face configurations.
\end{corollary}

\begin{proof}
Index the four triple supports by their missing elements.  Index a pair
support \(L\setminus\{u,v\}\) by the edge \(\{u,v\}\); its two triple
cofaces are indexed by \(u,v\).  Every nonempty upper ideal contains
\(L\).  It is therefore specified by a vertex set \(W\subseteq L\) of
triple supports and an arbitrary edge set
\(E\subseteq\binom W2\) of pair supports.  Its demand minus its order is
\[
 q-1+\sum_{v\in W}(t_v-1)
       +\sum_{\{u,v\}\in E}(p_{uv}-1).
\]
For fixed \(W\), this expression is maximized by taking precisely the
edges with \(p_{uv}>1\).  Lemma
\ref{lem:damp-coface-matching-upper-ideals} now gives
\eqref{eq:damp-rank-four-weighted-graph-hall}.
\end{proof}

\begin{proposition}[An overloaded maximal repair four-cube pays an
exterior tax]
\label{prop:damp-maximal-repair-four-cube-exterior-tax}
Let a positioned maximal repair four-cube \(Q\) of an interior section
\(P_j\) be the common owner of \(c\) directed neutral cycles.  Then
\begin{equation}
 \Delta(H;A,R)\ge11+\max\{0,2c-16\}.
 \label{eq:damp-maximal-repair-four-cube-exterior-tax}
\end{equation}
In particular, if \(c\ge3\), then
\begin{equation}
 \Delta(H;A,R)\ge c+1.
 \label{eq:damp-maximal-repair-four-cube-strict-tax}
\end{equation}
\end{proposition}

\begin{proof}
The four-cube contributes all six repair pairs, all four repair triples,
and its four-element support beyond the empty support and the repair
singletons.  Thus
\[
 |P_j|-(r+1)\ge11.
\]

The \(c\) cycles are disjoint as boundary occurrences and use at least
\(2c\) occurrences.  A repair word of \(Q\) has at most its left and
right occurrence, so at least
\(\max\{0,2c-16\}\) words occur on both boundary sides.  Every such word
is a corner of \(P_j\) whose unique maximal repair cube is \(Q\), and it
belongs to neither adjacent carrier.  As in the proof of
Proposition~\ref{prop:damp-maximal-repair-three-cube-exterior-tax},
ambient cornerlessness forces a distinct concept over a core word outside
the path for every doubled word.  Hence the external-concept term in the
path-section partition satisfies
\[
 V\ge\max\{0,2c-16\}.
\]
Together with the eleven section supports, this proves
\eqref{eq:damp-maximal-repair-four-cube-exterior-tax}.

For \(3\le c\le10\), the initial eleven units already give
\(\Delta(H;A,R)\ge c+1\).  For \(c\ge11\), the right-hand side of
\eqref{eq:damp-maximal-repair-four-cube-exterior-tax} is
\(2c-5\ge c+1\).  This proves the strict-tax assertion.
\end{proof}

\begin{corollary}[Pure rank-four packets have strict excess]
\label{cor:damp-pure-rank-four-packet-strict-excess}
Let \(|H|<64\), and let a locally deficient packet from Corollary
\ref{cor:damp-neutral-cycle-failure-provenance-normal-form} occupy a
positioned four-cube whose support lies in \(R\).  Then all packet members
are neutral cycles in one section and have one common positioned maximal
repair four-cube owner.  This owner supplies at least one unit of strict
excess beyond its cycle population.  Consequently a pure rank-four packet
does not belong to the global equality branch.
\end{corollary}

\begin{proof}
Every packet member is a neutral-cycle face, because a canonical failure
support contains a core coordinate.  Since the terminal cube has pure
repair support, all its faces have the same core anchor and lie in one
section \(P_j\).

Each initial face contains its assigned owner word, which is a corner of
\(P_j\).  Its unique maximal repair cube contains both that face and the
terminal four-cube.  The argument in the proof of Corollary
\ref{cor:damp-pure-rank-three-packet-maximal-owner} shows that all packet
members have the same maximal owner \(Q\).  If its support had order at
least five, the facet-excess inequality would give
\[
 |H|\ge2^6=64,
\]
contrary to \(|H|<64\).  Hence \(Q\) is the terminal four-cube itself.

Any family of at most two distinct faces has distinct coface labels:
use their own supports when these differ, and use the common support and
the full four-support when they agree.  A deficient packet therefore
contains at least three cycles.  Apply
Proposition~\ref{prop:damp-maximal-repair-four-cube-exterior-tax} to all
cycles owned by \(Q\).  Equation
\eqref{eq:damp-maximal-repair-four-cube-strict-tax} gives the claimed
strict unit.
\end{proof}

\begin{corollary}[Two-repair two-core rank-four packets separate]
\label{cor:damp-two-repair-two-core-rank-four-separates}
Let a rank-four packet from Corollary
\ref{cor:damp-neutral-cycle-failure-provenance-normal-form} have support
\begin{equation}
 L=T\mathbin{\dot\cup}\{x_i,x_k\},
 \qquad T\in\binom R2,\quad x_i,x_k\in X.
 \label{eq:damp-two-repair-two-core-rank-four-support}
\end{equation}
Then its members can be assigned pairwise distinct coface labels inside
the terminal four-cube.
\end{corollary}

\begin{proof}
Every neutral face has support \(T\).  A failure face has support
\(T\cup\{x_i\}\) or \(T\cup\{x_k\}\), and canonical failure supports are
distinct, so at most one of each kind occurs.  The neutral \(T\)-faces
are indexed by core anchors in the square on
\(\{x_i,x_k\}\).  They arise over vertices of the coordinate-simple path
\(A\), with every other core coordinate fixed by the terminal cube.
The intersection of \(A\) with this positioned core square is a path of
order at most three.  Thus at most three neutral faces occur.

The four available labels are
\[
 T,\qquad T\cup\{x_i\},\qquad T\cup\{x_k\},\qquad L.
\]
An explicit matching fails only if all three possible neutral faces and
both failure faces occur, giving five demands for these four labels.
Assume this profile.

After interchanging \(i,k\), the three core anchors are consecutive path
vertices
\[
 a_{i-1},a_i,a_{i+1},
 \qquad k=i+1.
\]
Indeed, they agree outside \(\{x_i,x_k\}\), so no other path direction
can occur between them.  The three neutral faces belong to one positioned
terminal cube and therefore have the same repair anchor outside \(T\).
The corresponding positioned \(T\)-cube lies in all three sections
\(P_{i-1},P_i,P_{i+1}\), and hence
\[
 T\in\operatorname{Sh}(C_i)\cap
       \operatorname{Sh}(C_{i+1}).
\]
On the other hand, the two failure supports are
\(T\cup\{x_i\}\) and \(T\cup\{x_{i+1}\}\).
Proposition~\ref{prop:damp-failed-lift-canonical-residual} says that their
acquired repair support \(T\) is absent from \(C_i\) and \(C_{i+1}\),
respectively.  This contradiction eliminates the only deficient profile.
\end{proof}

\begin{corollary}[The remaining rank-four packet is a three-arm
potential]
\label{cor:damp-rank-four-three-arm-potential}
Let a rank-four packet have support
\begin{equation}
 L=B\mathbin{\dot\cup}\{x_i\},
 \qquad B\in\binom R3.
 \label{eq:damp-three-repair-one-core-rank-four-support}
\end{equation}
For \(s\in B\), put \(K_s=B\setminus\{s\}\), and define
\begin{align*}
 p_s&=\#\{\text{neutral faces of support }K_s\},\\
 b&=\#\{\text{neutral faces of support }B\},\\
 f_s&=\#\{\text{failure faces of support }K_s\cup\{x_i\}\},\\
 g&=\#\{\text{failure faces of support }L\}.
\end{align*}
Then
\begin{equation}
 0\le p_s\le4,\qquad 0\le b\le2,
 \qquad f_s,g\in\{0,1\}.
 \label{eq:damp-rank-four-three-arm-basic-bounds}
\end{equation}
Canonical carrier absence also gives
\begin{equation}
 f_s=1\Longrightarrow p_s\le2,
 \qquad
 g=1\Longrightarrow b\le1.
 \label{eq:damp-rank-four-three-arm-carrier-bounds}
\end{equation}
The packet can be assigned distinct coface labels inside its terminal
cube if and only if
\begin{equation}
 \Phi_4
 :=
 g+b-2+\sum_{s\in B}(p_s+f_s-2)_+
 \le0.
 \label{eq:damp-rank-four-three-arm-potential}
\end{equation}
Consequently every locally deficient rank-four packet below order \(64\)
has the support profile
\eqref{eq:damp-three-repair-one-core-rank-four-support} and satisfies
\(\Phi_4>0\).  This is the sole remaining rank-four profile.
\end{corollary}

\begin{proof}
The support types follow from provenance.  A neutral face has pure repair
support of order at least two, hence support \(K_s\) or \(B\).  A failure
support has exactly one core coordinate, hence is
\(K_s\cup\{x_i\}\) or \(L\).  The numbers in
\eqref{eq:damp-rank-four-three-arm-basic-bounds} are the numbers of
parallel faces of the relevant supports; failure multiplicities are at
most one because canonical failure supports are distinct.

If \(f_s=1\), then
\(K_s\notin\operatorname{Sh}(C_i)\) by
Proposition~\ref{prop:damp-failed-lift-canonical-residual}.  For each of
the two values of the remaining repair coordinate \(s\), at most one of
the two core-side \(K_s\)-faces can therefore occur.  Hence \(p_s\le2\).
The same argument for the two opposite \(B\)-faces gives \(b\le1\) when
\(g=1\).  This proves
\eqref{eq:damp-rank-four-three-arm-carrier-bounds}.

It remains to apply Lemma
\ref{lem:damp-coface-matching-upper-ideals}.  The admissible support poset
has bottom elements \(K_s\), middle elements \(B\) and
\(K_s\cup\{x_i\}\), and top element \(L\).  An upper ideal containing a
nonempty set \(D\subseteq B\) of bottom arms must contain \(B,L\) and,
for every \(s\in D\), the mixed triple
\(K_s\cup\{x_i\}\).  Its maximum demand minus order is
\begin{equation}
 g+b-2+\sum_{s\in D}(p_s+f_s-2).
 \label{eq:damp-rank-four-three-arm-upper-ideal-excess}
\end{equation}
Optional mixed triples have multiplicity at most one and cannot increase
this excess.  An ideal containing no bottom arm has nonpositive excess:
the only potentially positive middle contribution is \(b-1\), and
\(g=1\) forces \(b\le1\).

The carrier bounds give \(g+b-2\le0\).  Every nonempty arm set \(D\)
has excess at most \(\Phi_4\).  Hence \(\Phi_4\le0\) makes all
upper-ideal excesses nonpositive.  If \(\Phi_4>0\), at least one arm has
positive value; taking exactly all positive arms gives a nonempty upper
ideal whose excess is \(\Phi_4\).  The upper-ideal criterion proves
\eqref{eq:damp-rank-four-three-arm-potential}.  Pure rank-four packets
have strict excess by Corollary
\ref{cor:damp-pure-rank-four-packet-strict-excess}, and the two-repair
two-core profile separates by Corollary
\ref{cor:damp-two-repair-two-core-rank-four-separates}.  Since every
initial support has at least two repair coordinates, these cases and
\eqref{eq:damp-three-repair-one-core-rank-four-support} exhaust rank four.
\end{proof}

\begin{lemma}[Two occupied core sides synchronize the repair cube]
\label{lem:damp-rank-four-two-core-sides-synchronize}
Let \(Q_L\) be a positioned full cube of \(H\), where

\[
 L=B\mathbin{\dot\cup}\{x_i\},
 \qquad B\subseteq R.
\]

Suppose that \(Q_L\) contains a neutral-cycle face on each of its two
\(x_i\)-sides.  Then those two faces lie over the consecutive path words
\(a_{i-1},a_i\), and the positioned full \(B\)-cube given by either
\(x_i\)-side belongs to the carrier \(C_i=P_{i-1}\cap P_i\).  In
particular,

\begin{equation}
 B\in\operatorname{Sh}(C_i).
 \label{eq:damp-rank-four-two-core-sides-synchronize}
\end{equation}
\end{lemma}

\begin{proof}
A neutral-cycle face lies over the path word of its owner section.  Let
the two path words supplied by the two \(x_i\)-sides of \(Q_L\) be
\(a_h,a_k\).  The positioned cube fixes every core coordinate other than
\(x_i\), so \(a_h\) and \(a_k\) differ exactly in \(x_i\).  A
coordinate-simple path flips \(x_i\) only on the edge
\(a_{i-1}a_i\).  Hence

\[
 \{a_h,a_k\}=\{a_{i-1},a_i\}.
\]

The two \(x_i\)-sides of \(Q_L\) are full positioned \(B\)-cubes with
the same repair anchor.  They therefore give one positioned full
\(B\)-cube in both \(P_{i-1}\) and \(P_i\), hence in \(C_i\).  This
proves \eqref{eq:damp-rank-four-two-core-sides-synchronize}.
\end{proof}

\begin{proposition}[Rank-four three-arm deficiency descends to rank
three]
\label{prop:damp-rank-four-three-arm-descends-rank-three}
Keep the notation of Corollary
\ref{cor:damp-rank-four-three-arm-potential}, and suppose that
\(\Phi_4>0\).  Then the initial neutral faces contain, on one
\(x_i\)-side of the terminal cube, a locally deficient pure-repair
rank-three packet in the positioned \(B\)-cube.  Thus a
three-repair one-core packet has no rank-four-irreducible Hall
obstruction.
\end{proposition}

\begin{proof}
Orient the two \(x_i\)-sides of the terminal cube by
\(\epsilon\in\{0,1\}\).  Let \(p_{s,\epsilon}\) be the number of
neutral \(K_s\)-faces on side \(\epsilon\), and let
\(b_\epsilon\in\{0,1\}\) record whether the neutral \(B\)-face on that
side occurs.  Thus

\[
 p_s=p_{s,0}+p_{s,1},
 \qquad b=b_0+b_1.
\]

For \(s\in B\), call side \(\epsilon\) a \emph{loaded \(s\)-row} if

\begin{equation}
 p_{s,\epsilon}=2
 \quad\text{and}\quad
 \bigl(f_s=1\ \text{or}\ p_{s,1-\epsilon}>0\bigr).
 \label{eq:damp-rank-four-loaded-row}
\end{equation}

Let \(\rho_{s,\epsilon}\) be its indicator, and put

\[
 q_\epsilon=b_\epsilon+\sum_{s\in B}\rho_{s,\epsilon}.
\]

Finally, let \(d\) count the coordinates \(s\in B\) for which
\(f_s=1\) and

\[
 p_{s,0}=p_{s,1}=1.
\]

For such an \(s\), canonical carrier absence prevents the two neutral
faces from having the same fixed value of \(s\); they are the two
diagonal \(K_s\)-faces, one on each core side.

The four possible \(K_s\)-faces form a \(2\)-by-\(2\) array, indexed by
the \(x_i\)-side and the fixed value of \(s\).  Reading this array gives
the exact identity

\begin{equation}
 (p_s+f_s-2)_+
 =\rho_{s,0}+\rho_{s,1}
  +\mathbf 1_{\{f_s=1,\ p_{s,0}=p_{s,1}=1\}}.
 \label{eq:damp-rank-four-arm-row-decomposition}
\end{equation}

Indeed, if \(f_s=0\), three selected faces give one loaded row and four
give two, while at most two give neither.  If \(f_s=1\), the carrier
bound gives \(p_s\le2\); two selected faces form either one row or the
diagonal pair just described.  Summing
\eqref{eq:damp-rank-four-arm-row-decomposition} yields

\begin{equation}
 \Phi_4=g+q_0+q_1+d-2.
 \label{eq:damp-rank-four-side-potential-decomposition}
\end{equation}

If \(q_\epsilon\ge2\) for some side \(\epsilon\), choose any two of the
events counted by \(q_\epsilon\).  A \(b_\epsilon\)-event is the full
\(B\)-face, while a loaded-row event consists of the two opposite
\(K_s\)-faces for one \(s\).  Two events therefore contain either the
full \(B\)-cube together with one opposite pair, or two opposite pairs
with distinct supports.  These are exactly the two minimal
rank-three Hall obstructions in Lemma
\ref{lem:damp-rank-three-packet-local-hall}.  All their faces are
neutral and lie on one core side, so this is the required pure-repair
rank-three subpacket.

It remains to rule out \(q_0,q_1\le1\).  Suppose first that \(g=1\).
Positivity in
\eqref{eq:damp-rank-four-side-potential-decomposition} gives
\(q_0+q_1+d\ge2\).  If \(d>0\), a diagonal arm supplies a neutral face
on each core side.  If \(d=0\), then \(q_0=q_1=1\), and again each core
side contains a neutral face.  Lemma
\ref{lem:damp-rank-four-two-core-sides-synchronize} gives
\(B\in\operatorname{Sh}(C_i)\).  But \(g=1\) is a canonical failure of
support \(B\cup\{x_i\}\), so Proposition
\ref{prop:damp-failed-lift-canonical-residual} gives
\(B\notin\operatorname{Sh}(C_i)\), a contradiction.

Finally suppose that \(g=0\).  Positivity gives
\(q_0+q_1+d\ge3\), and hence \(d>0\).  Choose a diagonal arm \(s\).
The same synchronization lemma gives
\(B\in\operatorname{Sh}(C_i)\), and downward closure gives
\(K_s\in\operatorname{Sh}(C_i)\).  On the other hand, the diagonal arm
has \(f_s=1\), whose canonical failure identity gives
\(K_s\notin\operatorname{Sh}(C_i)\).  This final contradiction proves
that some \(q_\epsilon\) is at least two, and the first part of the proof
then supplies the rank-three subpacket.
\end{proof}

\begin{corollary}[Rank four has no new packet obstruction]
\label{cor:damp-rank-four-packets-closed}
Below order \(64\), every locally deficient rank-four packet either is
strictly paid by the pure-repair four-cube tax or contains one of the
already classified pure-repair rank-three packets.  The two-repair
two-core profile and every rank-four-irreducible three-repair one-core
profile separate.  Consequently rank four contributes no new equality
profile beyond the rank-three row.
\end{corollary}

\begin{proof}
Pure-repair packets have strict excess by Corollary
\ref{cor:damp-pure-rank-four-packet-strict-excess}.  Two-repair
two-core packets separate by Corollary
\ref{cor:damp-two-repair-two-core-rank-four-separates}.  The only other
support profile is
\(B\mathbin{\dot\cup}\{x_i\}\).  If it is locally deficient, then
\(\Phi_4>0\) by Corollary
\ref{cor:damp-rank-four-three-arm-potential}, and Proposition
\ref{prop:damp-rank-four-three-arm-descends-rank-three} finds a proper
rank-three deficient subpacket unless the configuration is already
contradictory.  This exhausts the rank-four support profiles.
\end{proof}

\begin{proposition}[Repair-position carriers conserve strict cofaces]
\label{prop:damp-repair-position-carrier-strict-cofaces}
Keep the notation of Proposition
\ref{prop:damp-projected-carrier-doubled-support-fragmentation}.  Besides
the core projection \(Z_T\), project the full \(T\)-cube carrier to the
coordinates outside \(T\) in the repair block:
\[
 Y_T=\mathcal K_T(H)|_{R\setminus T}.
\]
Then \(Y_T\) is ample and
\begin{equation}
 \operatorname{Sh}(Y_T)
 =\{S\subseteq R\setminus T:
       T\mathbin{\dot\cup}S\in\operatorname{Sh}(H)\}.
 \label{eq:damp-repair-position-carrier-strict-cofaces}
\end{equation}
Consequently, if two positioned full \(T\)-cubes of \(H\) have distinct
repair anchors outside \(T\), then for every coordinate \(s\) on which
those anchors differ,
\begin{equation}
 T\cup\{s\}\in\operatorname{Sh}(W)\setminus\mathcal L_R
 \subseteq\Xi.
 \label{eq:damp-distinct-repair-positions-strict-cofaces}
\end{equation}
\end{proposition}

\begin{proof}
Apply Lemma
\ref{lem:damp-projected-cube-carrier-support-conservation} with \(T\) as
the cube support and \(R\setminus T\) as the projection block.  This gives
ampleness of \(Y_T\) and
\eqref{eq:damp-repair-position-carrier-strict-cofaces}.  If two anchors
differ on \(s\), an isometric \(Y_T\)-geodesic between them has an
\(s\)-edge.  Hence \(\{s\}\in\operatorname{Sh}(Y_T)\), and the displayed
identity gives \(T\cup\{s\}\in\operatorname{Sh}(H)\).  This support is
purely in \(R\), so it lies in \(\operatorname{Sh}(W)\); it has order at
least three because \(|T|\ge2\), and therefore is outside
\(\mathcal L_R\) and belongs to \(\Xi\).
\end{proof}

\begin{corollary}[A tight two-sided failure is a coface or one word hole]
\label{cor:damp-tight-two-sided-failure-anchor-dichotomy}
Fix a one-section gap of \(I_T\) and suppose that one failed lift descends
from each of its two initial boundary occurrences.  For a failure \(F\),
let \(p_F\) be its attempted repair word and put
\[
 u_F=p_F|_{R\setminus T}.
\]
Then \(u_F\in Y_T\).  For the two failures \(F_-,F_+\), exactly one of the
following holds:
\begin{enumerate}
 \item the anchors \(u_{F_-}\) and \(u_{F_+}\) differ, and every coordinate
       in their difference supplies the strict pure-repair coface
       in \eqref{eq:damp-distinct-repair-positions-strict-cofaces};
 \item \(p_{F_-}=p_{F_+}\).  Thus the two failure lineages terminate at
       the same missing concept in the middle section: one repair word is
       present on both sides of the gap and absent in its centre.
\end{enumerate}
Hence the position-sensitive part of a tight hole has only one residual
profile: a common missing repair word.  Distinct positions already produce
a strict-coface candidate in \(\Xi\).
\end{corollary}

\begin{proof}
At a failed lift, adjoining \(p_F\) to the target section and to the
adjacent carrier is locally ample with unique new repair support \(T\).
The one-concept extension test gives
\[
 p_F^U\in C_i\qquad(\varnothing\ne U\subseteq T).
\]
The source section also contains \(p_F\).  It therefore contains the full
positioned \(T\)-cube through \(p_F\), proving \(u_F\in Y_T\).

If the two anchors differ, Proposition
\ref{prop:damp-repair-position-carrier-strict-cofaces} gives the first
alternative.  Suppose that they agree.  The two repair words then belong
to the same positioned \(T\)-face.  Both failed lifts target the unique
middle section of the tight gap.  At the time of either attempt, that
section contains every other vertex of the corresponding \(T\)-face.
Section families grow throughout the two sweeps.  If the two attempted
words were different, the earlier punctured face would already contain the
word attempted later, contradicting its absence at the later attempt.
Thus the attempted words are equal, proving the second alternative.
\end{proof}

\begin{corollary}[Transverse blockers supply pure-repair cross-supports]
\label{cor:damp-failed-lift-pure-repair-cross-supports}
Consider a failed adjacent-section lift with acquired repair support
\(T\), attempted concept \(v\), and antipodal blocker \(v^J\).  Put
\[
 S=J\cap R,
 \qquad K=J\cap X.
\]
Then \(|T|\ge2\), \(T\cap S=\varnothing\), and
\begin{equation}
 \{t,s\}\in\operatorname{Sh}(H)
 \qquad(t\in T,\ s\in S).
 \label{eq:damp-failed-lift-pure-repair-cross-supports}
\end{equation}
If the blocker has exactly one core direction, then \(S\ne\varnothing\),
so the failure forces at least \(|T|\) nonbaseline pure-repair supports.
In particular, a tight two-sided hole whose two blockers lie over the two
opposite neighbouring path sections exposes such a support family on each
side.
\end{corollary}

\begin{proof}
Every repair section contains the two fixed repair antipodes, so it already
shatters the empty set and every repair singleton.  The support newly
acquired by a locally ample section extension is therefore nonbaseline,
and \(|T|\ge2\).  The local one-concept extension test gives
\[
 v^U\in H\qquad(\varnothing\ne U\subseteq T).
\]
Apply Lemma~\ref{lem:damp-blocked-puncture-bipartite-two-shadow}.  It gives
\(T\cap J=\varnothing\) and shatters \(\{t,e\}\) for every
\(t\in T\) and \(e\in J\).  Restricting \(e\) to \(S\) proves
\eqref{eq:damp-failed-lift-pure-repair-cross-supports}.  Finally the
blocker has order at least two and has nonempty core part by
Corollary~\ref{cor:damp-adjacent-carrier-extension-or-owner}; if that core
part is a singleton, its repair part \(S\) must be nonempty.  Distinct
pairs \(\{t,s\}\) have order two and hence lie outside the common repair
baseline.
\end{proof}
\begin{remark}[The remaining path-wide proof contract]
\label{rem:damp-adjacent-carrier-remaining-contract}
Corollary~\ref{cor:damp-adjacent-carrier-extension-or-owner} separates the
local problem from the global one.  A transition that stops locally pays a
cyclic block governed by Proposition
\ref{prop:damp-path-terminal-cycle-charge-conservation}.  A transition
that fails only after lifting to \(H\) has a transverse core owner, and
Proposition~\ref{prop:damp-adjacent-carrier-owner-monotonicity} makes every
on-spine transfer move strictly toward an endpoint.  Proposition
\ref{prop:damp-adjacent-carrier-owner-residual-charge} already assigns a
current-state residual support to every failed lift.  Thus the remaining
theorem need not construct local peeling states, exclude owner cycles, or
compare different extension states: Corollary
\ref{cor:damp-residual-payment-amortization} transports every later payment
back to the original residual shadow.  Proposition
\ref{prop:damp-section-shadow-boundary-conservation} controls repeated
section acquisitions by path intervals.

The neutral part now has one uniform accounting theorem.  Proposition
\ref{prop:damp-neutral-transfer-forest-section-excess} adjoins every
opposite tip as a terminal vertex and counts incidences between weak
components and their underlying repair words.  Every component uses at
least two words.  A word outside two chosen antipodal carrier geodesics can
belong to at most its left and right components, while an exclusive
geodesic word belongs to at most one.  The two multiplicity bounds cancel
exactly to give
\[
 \kappa_j\le |P_j|-(r+1).
\]
Summing over the interior sections pays every cyclic and acyclic neutral
component directly from the additive path-section excess.  In particular,
terminal rank, terminal-support reuse, arrow labels, and repeated cycles in
one positioned cube create no separate multiplicity problem.  The
closed-cycle packing, rank-one surcharge, sectionwise cube tax,
positioned-cube degree bound, fixed-support capacity, and common-cube rank
restriction above remain valid quantitative refinements, but they are no
longer active proof-debt rows.

Corollary~\ref{cor:damp-neutral-forest-equality-two-word} also classifies
the only way this bound can be tight.  Every weak component then uses two
repair words, and the component--word incidence graph is a disjoint union
of paths and cycles.  Its path endpoints are exactly the exclusive words
of the two carrier geodesics; its cyclic word vertices lie outside both.
In particular, every closed neutral component in the global equality
branch is already a directed two-cycle.  Longer neutral cycles and
higher-word acyclic components create strict section slack and are not
part of the remaining saturation problem.

There is also no hidden singleton exception.  Lemma
\ref{lem:damp-square-free-edge-bridge} makes a maximal singleton attachment
a bridge and the only edge in its repair direction.  Since the target
carrier contains the two repair antipodes, it must use that direction and
therefore must contain the bridge, contradicting the definition of the
opposite tip.  This is Proposition
\ref{prop:damp-terminal-rank-one-impossible}.  It is a structural
strengthening, although the forest estimate already makes the rank of an
acyclic terminal irrelevant.

The path-wide accounting can consequently be assembled without a terminal
classification.  For the \(2n\) oriented adjacent-section
incidences, compare the selected initial boundary-corner demands with the
terminal cyclic charges, transverse-owner payments, and neutral weak
components, all in the original state.  By Corollaries
\ref{cor:damp-neutral-two-sided-coalescence-star} and
\ref{cor:damp-neutral-tight-coalescence-two-cycle}, a component containing
the two selected demands from opposite sides of one section is already paid
twice unless those demands form a directed two-cycle.  Proposition
\ref{prop:damp-neutral-transfer-two-sided-cycle} identifies its two parent
punctures exactly: they have one common repair label and opposite missing
tips in the two adjacent carriers.

The directed two-cycle is no longer an unclassified terminal.  Proposition
\ref{prop:damp-two-sided-progressive-synchronization} and Corollary
\ref{cor:damp-tight-two-cycle-peripheral-handoff} show that it either
produces a residual support which transports to the original shadow, or
has an ample extension in which its middle section is contained in both
neighbouring sections.

The handoff itself is now globally controlled, rather than left as an
iteration problem.  Lemma
\ref{lem:damp-adjacent-duplication-shadow-transport} shows that each
successful concept duplication carries exactly one existing repair support
across the same section edge.  Proposition
\ref{prop:damp-adjacent-duplication-boundary-lyapunov} weights a boundary by
its remaining distance to the endpoint toward which it moves.  Both the
word potential and the rank-weighted shadow potential decrease strictly;
the latter decrease pays the simultaneous increase of the section-edge
rank.  Corollary
\ref{cor:damp-boundary-descent-pays-marked-acquisitions} also handles the
only possible collision: if one duplication resolves marked demands from
both neighbouring centres, two opposite boundaries annihilate and the
potential has more, not less, slack.

Proposition~\ref{prop:damp-tight-two-cycle-two-sweep-reduction} therefore
treats all tight two-cycles visible in the original state in two directed
sweeps.  Every one is either erased with an explicit boundary-descent
payment, replaced by a terminal cyclic block, or replaced by a transverse
residual payment transported to the original shadow.  There is no remaining
infinite closure, owner-repetition, or endpoint-confluence branch.

The remaining assembly row is narrower still.  Proposition
\ref{prop:damp-dynamic-terminal-block-conservation} combines every terminal
block created before or during the sweeps into the original section-edge
rank plus the original shadow-boundary potential.  Successful acquisitions
are already paid by Corollary
\ref{cor:damp-boundary-descent-pays-marked-acquisitions}.  Only the
transverse failures remain.  Proposition
\ref{prop:damp-failed-lift-canonical-residual} identifies their residual
payments exactly: a failure across $x_i$ carries the support
$T\cup\{x_i\}$, where $T$ is its transported section-shadow boundary.
Proposition
\ref{prop:damp-canonical-failure-support-birth-injection} shows that these
labels are pairwise distinct over the complete history and all belong to
$\Omega_0$.  Thus transverse failures inject directly into the original
residual shadow.  The residual core moment, $\Lambda_0$, and every
multiplicity factor in Proposition
\ref{prop:damp-transported-residual-bounded-genealogy} are no longer active
proof debt.

The dynamic collision problem is therefore closed.  Proposition
\ref{prop:damp-path-spine-support-interval-exact-sequence} also removes a
second apparent overlap.  For each repair support \(T\), successful
acquisitions are columns of the forest boundary map \(\partial_T\).  They
identify section occurrences along one interval and are not an additional
cardinality charge.  The original excess has the exact static partition
\[
 \Delta(H;A,R)=|\Xi|+|\Omega_0|,
\]
where \(\Xi\) consists of one pure-support root, the internal carrier edges
of every section-support interval, and the repair supports seen only in
\(W\).  The only cokernel has one generator for every interval component
after the distinguished one; equation
\eqref{eq:damp-support-interval-residual-balance} records exactly how those
generators, external concepts, and repair supports seen only in \(W\)
balance the residual shadow.

Thus boundary descent is no longer a separate active row.  Proposition
\ref{prop:damp-projected-carrier-doubled-support-fragmentation} assembles
all positioned \(T\)-cubes in one ample projected carrier \(Z_T\).  If the
section-index set of \(T\) has \(c_T\) components, then the residual shadow
with exact repair part \(T\) has at least \(2(c_T-1)\) members.  One unit
comes from the cokernel component in
\eqref{eq:damp-support-interval-exact-sequence}; the other is forced by an
external core vertex on an isometric carrier path between consecutive
components.

Proposition
\ref{prop:damp-gap-edge-failure-fragmentation-reservation} makes those two
units explicit.  A gap of edge length \(d\) reserves all \(d\) labels
\[
 T\cup\{x_i\}
\]
on its path edges.  At most two failure lineages can enter one gap.  Hence
the failure labels and one fragmentation generator fit in the reserved
family whenever \(d\ge3\), and also when \(d=2\) unless both lineages fail.
The former supportwise complementarity problem is therefore reduced to one
normal form: a single missing section between two \(T\)-intervals, with a
failed descendant from each side.

The exceptional tight hole has a simpler payment that forgets the repair
positions altogether.  Its projected carrier \(Z_T\) contains the two
path vertices on either side of the missing section and is isometric.  It
must therefore use the fourth vertex of their core square.  Together with
the three consecutive path vertices, this external carrier vertex shatters
the pure core pair \(\{x_j,x_{j+1}\}\).  Two different exceptional holes
cannot have the same missing section, because that would reuse the same two
failed oriented incidences with two different unique acquired supports.
Proposition~\ref{prop:damp-tight-gap-core-square-reservation} consequently
pays all correction terms \(b_T\) by distinct original residual supports
and proves the global complementarity inequality
\[
 |\mathcal F|+\sum_{T\in\mathcal T}(c_T-1)\le|\Omega_0|.
\]
The repair-position and blocker results remain useful local refinements,
but they are no longer required for this payment.

Thus the failure--fragmentation row is closed without a support-incidence
Hall theorem.  The occurrence row is closed as well.  Proposition
\ref{prop:damp-support-occurrence-failure-joint-injection} embeds every
chosen original-state occurrence demand jointly with all failures in
\(\Xi\mathbin{\dot\cup}\Omega_0\).  Proposition
\ref{prop:damp-spacetime-support-forest} adds exactly one formal column per
successful acquisition, and Corollary
\ref{cor:damp-terminal-remainders-no-residual-debt} shows that dynamic
terminal remainders use only unmarked columns of this kind.  Hence no local
collision between neutral components, terminal remainders, successful
acquisitions, fragmentation generators, or failures remains.

This conclusion is a capacity theorem, not yet a contradiction.  An
injection of the selected demands into the excess budget proves that the
payments can coexist; it does not prove that cornerlessness creates more
payments than the budget contains.  The independent fixed-spine argument
now supplies the lower endpoint \(40\), while the path-core target remains
\(46\).

The arbitrary occurrence-label normal form in Corollary
\ref{cor:damp-joint-path-budget-anchored-normal-form} is only a capacity
statement: its chosen section cube need not contain the puncture that
created the demand.  Corollary
\ref{cor:damp-neutral-cycle-failure-provenance-normal-form} supplies the
geometric replacement for the closed neutral components and retains the
canonical residual labels of the failures.  It anchors each cycle in its
actual owner cube and represents each failure by an auxiliary cube on its
canonical support.  The full positioned-cube carrier then moves every
collision between different positions strictly upward in the support
poset.  Below order \(48\), every terminal packet consequently has rank
three or four; rank two occurs only in the occurrence proxy.  Acyclic
neutral components require no such packet analysis, because the forest
accounting charges them at their terminal outcomes.

The final contract is therefore an exact matching statement for these
anchored faces.  Lemma~\ref{lem:damp-rank-three-packet-local-hall} reduces
the rank-three branch to two minimal local motifs: a full cube with one
opposite pair of square faces, or two opposite pairs in different
directions.  Corollary
\ref{cor:damp-mixed-rank-three-packet-separates} excludes the mixed
core--repair instance of the first motif, because it would put the
failure's repair support back into the path-edge carrier from which it is
canonically absent.  Thus both surviving motifs are pure-repair
neutral-cycle configurations.  Corollary
\ref{cor:damp-pure-rank-three-packet-maximal-owner} then shows that either
the alternating component already reaches rank four or the terminal cube
is a maximal repair three-cube in its common owner section.  It remains to
show that the two motifs in this maximal-cube normal form supply strict
surplus or expose a corner.  Proposition
\ref{prop:damp-maximal-repair-three-cube-exterior-tax} already gives strict
tax for every possible cycle population, including the apparent
four-cycle equality case.  Thus no isolated rank-three equality profile
remains; in a globally saturated component, its spare unit must be consumed
by another alternating demand.  The equality skeleton makes this handoff
considerably more rigid: Corollary
\ref{cor:damp-saturated-maximal-three-cube-three-two-cycles} shows that a
maximal-three-cube motif in a saturated section consists of exactly three
directed two-cycles on six distinct cube words.  All other cycle
populations create strict section slack.  Proposition
\ref{prop:damp-three-two-cycles-six-section-units} now excludes both the
four-unit cube profile and its one-support extension.  In a saturated
section the three cycles therefore coexist with at least three further
two-word components; the former one- and two-attachment profiles do not
occur.  Proposition~\ref{prop:damp-three-cycle-short-or-seven} reduces
section excess six to two unused cube words that belong to both adjacent
carriers and have distance one or two.

The two-square attachment lemma closes that residue.  At exact excess six,
the shadow has only two pair supports beyond the four supports of the
owner three-cube.  After contracting singleton bridge directions,
edge-carrier connectedness makes the section the union of exactly three
maximal blocks: the owner cube and the two positioned squares.  The
intersection nerve and the two repair antipodes then force the following
dichotomy.  If one antipode lies in the owner cube, it is a cube corner;
when only one attachment word remains, that word is opposite to the
antipode.  If neither antipode lies in the cube, the two attachment words
are opposite.  With six distinct cycle corners, either there are too few
corners outside the common antipode or the two unused words are therefore
opposite.  Both conclusions contradict the short common-pair normal form.
Proposition~\ref{prop:damp-three-two-cycles-seven-section-units} consequently
gives
\[
 |P_j|-(r+1)\ge7.
\]
In a saturated section the three cycles coexist with at least four further
two-word components.

Thus the last rank-three equality profile is closed by one cube-block
nerve argument, not by the fixed-rank realization list that suggested it.
Every deficient anchored packet below order \(48\) now either reaches rank
four or already creates strict excess.

The rank-four row also has an exact order-free compression.  Coface
matching is equivalent to the upper-ideal inequalities
\eqref{eq:damp-coface-matching-upper-ideal-criterion}; in a pure
four-cube these are the induced-subgraph weights in
\eqref{eq:damp-rank-four-weighted-graph-hall}.  A pure-repair packet has
a common maximal repair four-cube owner, whose eleven nonbaseline supports
give strict excess.  A packet with two repair and two core coordinates
separates: its only Hall-deficient profile would put the same repair
two-cube in two path-edge carriers from which the two canonical failure
labels require it to be absent.

The final mixed row initially appears as one three-arm potential.  Its
terminal support is \(B\mathbin{\dot\cup}\{x_i\}\), with three repair
coordinates and one core coordinate, and it is deficient exactly when
\[
 \Phi_4
 =g+b-2+\sum_{s\in B}(p_s+f_s-2)_+>0.
\]
This potential has a second, geometric decomposition.  Its positive
units are loaded rows on the two core sides together with diagonal
failure arms.  Two loaded events on one side are exactly a deficient
pure-repair rank-three subpacket.  If neither side contains two events,
positivity forces neutral faces on both sides of the \(x_i\)-edge.  The
terminal cube then puts the full \(B\)-cube in \(C_i\), contradicting
the canonical carrier absence of either the top failure or a diagonal
arm.  Proposition
\ref{prop:damp-rank-four-three-arm-descends-rank-three} therefore shows
that every positive \(\Phi_4\) descends to the already closed rank-three
row.  Corollary~\ref{cor:damp-rank-four-packets-closed} closes rank four
without a face-placement list.

The packet-free equality case is closed by strict facet ascent.  Every
proper carrier has a corner peeling, whose cube-support matching can be
rerooted by alternating-path exchange at any prescribed carrier position.
This matching flexibility holds even when no peeling ends at that position.
A proper set of positioned faces can consequently be moved entirely to
strict cofaces, with no carrier-order threshold.  In the absence of a
packet, iteration ends with distinct facet labels.  Equality exhausts the
nonbaseline shadow with these facets.
Downward closure then forbids a support of rank at least three: any two of
its repair coordinates would form a proper nonbaseline subface that is
also a facet.  The whole shattered complex has VC dimension at most two,
and relative VC-two peeling exposes a corner.  Thus packet-free equality
is impossible, simultaneously for all orders \(40,\ldots,47\).

Together with the rank-three and rank-four closures, this completes the
path-core repair theorem below order \(48\).  No cardinality or path-profile
census enters the final step.

The shift-criticality theorem and the two-sided shield-capacity lemma remain
independent structural checks on a least obstruction; they are not needed
for the completed lower bound.
\end{remark}

\begin{lemma}[Four-coordinate section synchronization]
\label{lem:damp-four-coordinate-section-synchronization}
Let \(|E|\le4\), let \(C,Y\subseteq Q_E\), and suppose that
\begin{equation}
H=(C\times\{0\})\cup(Y\times\{1\})
 \subseteq Q_{E\cup\{r\}}
 \label{eq:damp-section-synchronization-lift}
\end{equation}
is ample.  Put \(S=C\cap Y\).  If \(S\ne\varnothing\) and
\(Y\setminus C\ne\varnothing\), then some \(p\in Y\setminus C\)
satisfies
\begin{equation}
S\cup\{p\}\text{ is ample},
 \qquad
 C\cup\{p\}\text{ is ample}.
 \label{eq:damp-section-synchronized-augmentations}
\end{equation}
For every such \(p\), the subfamily
\begin{equation}
(C\times\{0\})
 \cup((S\cup\{p\})\times\{1\})
 \label{eq:damp-section-synchronized-sublift}
\end{equation}
is ample.
\end{lemma}

\begin{proof}
The coordinate sections and their intersection are ample.  Theorem
\ref{thm:damp-nested-ample-relative-peeling-dimension-four}, applied to
the proper inclusion \(S\subsetneq Y\), gives a concept
\(p\in Y\setminus S=Y\setminus C\) such that
\(S\cup\{p\}\) is ample.  Lemma
\ref{lem:damp-expansion-separator-augmentation} then shows that
\(C\cup\{p\}\) is ample as well.  This proves
\eqref{eq:damp-section-synchronized-augmentations} without a finite
classification of the section pairs.

It remains to justify the final assertion without computation.  Both
singleton augmentations in
\eqref{eq:damp-section-synchronized-augmentations} add exactly one new
shattered support; call them \(R_S\) and \(R_C\), respectively.  Since
\(S\cup\{p\}\subseteq Y\), we have \(R_S\in\operatorname{Sh}(Y)\).
Ampleness of \(H\), applied to supports containing \(r\), gives
\begin{equation}
\operatorname{Sh}(C)\cap\operatorname{Sh}(Y)
 =\operatorname{Sh}(S).
 \label{eq:damp-section-synchronization-no-synergy}
\end{equation}
Thus \(R_S\notin\operatorname{Sh}(C)\).  On the other hand,
\(S\cup\{p\}\subseteq C\cup\{p\}\), so monotonicity of shattering
forces \(R_S=R_C\).  The two-section formula now gives both required
equalities for \eqref{eq:damp-section-synchronized-sublift}: supports
not containing \(r\) add the common support \(R_S\), while supports
containing \(r\) remain exactly \(\operatorname{Sh}(S)\).  Hence the
sublift is ample.
\end{proof}

\begin{lemma}[A peripheral upper endpoint lowers the interval dimension]
\label{lem:damp-peripheral-upper-endpoint-reduction}
Suppose that every proper nested pair of nonempty ample families using at
most \(d-1\) coordinates has an ample one-concept intermediate family.
Let \(\varnothing\ne S\subsetneq B\subseteq Q_E\) be ample, with
\(|E|=d\).  If \(B\) has a peripheral coordinate, then some
\(p\in B\setminus S\) makes \(S\cup\{p\}\) ample.
\end{lemma}

\begin{proof}
Orient a peripheral coordinate \(r\) so that
\[
 B=(C\times\{0\})\cup(Y\times\{1\}),
 \qquad Y\subseteq C,
\]
and write
\[
 S=(A\times\{0\})\cup(D\times\{1\}),
 \qquad A\subseteq C,\quad D\subseteq Y.
\]
If \(A=\varnothing\), choose any \(p\in D\).  Adding its lower copy is
a peripheral expansion of \(D\) along the singleton \(\{p\}\), and is
ample.  Thus assume \(A\ne\varnothing\).

First suppose that \(D\setminus A\ne\varnothing\).  Connectivity of
\(S\) gives \(I=A\cap D\ne\varnothing\).  The lower-dimensional
hypothesis applied to \(I\subsetneq D\) supplies
\(p\in D\setminus I=D\setminus A\) for which \(I\cup\{p\}\) is ample.
Lemma~\ref{lem:damp-expansion-separator-duplication}, applied to \(S\),
then shows that adjoining the lower copy of \(p\) to \(S\) is ample.
Since \(p\in D\subseteq Y\subseteq C\), that copy belongs to \(B\).

It remains to suppose that \(D\subseteq A\).  If \(A\subsetneq C\),
the lower-dimensional hypothesis supplies \(p\in C\setminus A\) with
\(A\cup\{p\}\) ample.  The sections
\(D\subseteq A\cup\{p\}\) are nested, so adding the lower copy of \(p\)
preserves ampleness.  Finally, if \(A=C\), properness of
\(S\subset B\) gives \(D\subsetneq Y\).  Apply the lower-dimensional
hypothesis to \(D\subsetneq Y\), or choose any \(p\in Y\) when
\(D=\varnothing\), and add the upper copy.  Again the resulting sections
are nested ample families.  This exhausts the cases.
\end{proof}

\begin{lemma}[Sparse crossing graphs force a peripheral coordinate]
\label{lem:damp-sparse-q5-crossing-peripheral}
Let \(A\subseteq Q_5\) be ample and use all five coordinates.  If the
crossing graph of its five \(\Theta\)-classes has at most four edges,
then \(A\) has a peripheral coordinate.
\end{lemma}

\begin{proof}
Let \(G\) be the one-inclusion graph of \(A\), and write \(G^\sharp\)
for its crossing graph.  By hypothesis, \(|E(G^\sharp)|\le4\).

If \(G^\sharp\) is a forest, then \(G\) is a \(C_4\)-cactoid by
\cite[Theorem~5.5]{klavzar2002partial}.  A leaf block is a \(K_2\) or a
\(C_4\)-tree.  In the former case its outer edge class is peripheral.
In the latter, root the tree of edge-glued squares at the unique cut
vertex through which the rest of \(G\) is attached.  A farthest leaf
square has a new \(\Theta\)-class whose outer two-vertex side avoids the
cut vertex and is peripheral in all of \(G\).  The same conclusion is
immediate if \(G\) itself is a single block.

Suppose that \(G^\sharp\) contains a cycle.  A connected graph on five
vertices with a cycle has at least five edges, so \(G^\sharp\) is
disconnected.  It has only one cyclic component: two cyclic components
would use at least six vertices.  The components of the crossing graph
are the crossing graphs of the blocks of \(G\), with a \(K_1\) component
corresponding to a bridge block; this follows from
\cite[Theorem~3.4]{klavzar2002partial} applied blockwise.  The block--cut
tree has at least two leaf blocks, and at most one of them has cyclic
crossing graph.  Choose a leaf block whose crossing graph is a tree.
By \cite[Theorem~5.4]{klavzar2002partial}, it is a \(K_2\) or a
\(C_4\)-tree.  The preceding leaf-block argument again supplies a
peripheral \(\Theta\)-class of \(G\).  Thus the full-support family
always has a peripheral coordinate.
\end{proof}

\begin{corollary}[Small five-coordinate endpoints are peripheral]
\label{lem:damp-small-q5-endpoint-peripheral}
If \(A\subseteq Q_5\) is ample and \(|A|\le10\), then either \(A\) is
contained in a coordinate hyperplane or \(A\) has a peripheral
coordinate.
\end{corollary}

\begin{proof}
Suppose that \(A\) uses all five coordinates.  A crossing pair is a
shattered coordinate pair.  Write \(G^\sharp\) for the crossing graph
of the one-inclusion graph of \(A\).  Then
\[
 |A|=|\operatorname{Sh}(A)|
 \ge1+5+|E(G^\sharp)|.
\]
Thus \(|E(G^\sharp)|\le4\), and Lemma
\ref{lem:damp-sparse-q5-crossing-peripheral} applies.
\end{proof}

\begin{lemma}[Two-sided normal form of a minimal ample-interval gap]
\label{lem:damp-minimal-interval-gap-two-sided}
Fix a coordinate set \(E\), and suppose that some proper nested ample
pair has no ample one-concept intermediate family.  Choose such a pair
\(S\subsetneq B\subseteq Q_E\) minimizing
\(|U|\), where \(U=B\setminus S\), and put
\begin{equation}
A=Q_E\setminus B,
 \qquad T=Q_E\setminus S.
 \label{eq:damp-minimal-gap-four-endpoints}
\end{equation}
Then there is no ample family strictly between \(S\) and \(B\).  The
complementary pair \(A\subsetneq T\) has the same property.  In
particular, for every \(p\in U\), none of
\begin{equation}
S\cup\{p\},\qquad B\setminus\{p\},\qquad
 A\cup\{p\},\qquad T\setminus\{p\}
 \label{eq:damp-minimal-gap-four-failures}
\end{equation}
is ample.
Equivalently, every corner of \(B\) belongs to \(S\), and every corner
of \(T\) belongs to \(A\).

Moreover, every \(p\in U\) has two blocking vertices
\(s_p\in S\) and \(a_p\in A\).  A \(B\)-geodesic from \(p\) to
\(s_p\) and a \(T\)-geodesic from \(p\) to \(a_p\) both start with an
edge whose other endpoint lies in \(U\).  Finally,
\(\Psi(B,S)=(\Delta_B,\Delta_S)\) is fully Cohen--Macaulay but loses
this property when any one of its relative facets is filled into
\(\Delta_S\).
\end{lemma}

\begin{proof}
If an ample \(C\) satisfied \(S\subsetneq C\subsetneq B\), then the
pair \(S\subset C\) has smaller difference.  Minimality would give an
ample augmentation of \(S\) by a member of \(C\setminus S\), contrary
to the choice of \(S\subset B\).  Thus no such \(C\) exists, and the
first two failures in \eqref{eq:damp-minimal-gap-four-failures} follow.
Set complementation in \(Q_E\) reverses the interval and preserves
ampleness, proving the corresponding assertions for \(A\subset T\)
and the other two failures.
The ample corner-deletion criterion says that deleting a vertex preserves
ampleness exactly when that vertex is a corner.  Applying it to \(B\) and
\(T\) gives the two corner containments.

Apply the blocking-vertex assertion of Lemma
\ref{lem:damp-one-concept-colon-test} first to \(S\subset B\) and then
to \(A\subset T\).  It supplies \(s_p,a_p\) and the two asserted first
edges.  Lemma~\ref{lem:damp-relative-yang-cm-flag-h} makes
\(\Psi(B,S)\) fully Cohen--Macaulay.  Filling the relative facet of
\(p\) replaces \(S\) by \(S\cup\{p\}\), which is not ample and hence
does not have a Cohen--Macaulay Yang complex.  This proves the final
claim.
\end{proof}

\begin{proposition}[Central relative-CM profile in five coordinates]
\label{prop:damp-q5-central-relative-cm-profile}
If the five-coordinate ample-interval theorem fails, then it has a
minimal counterexample \(S\subsetneq B\subseteq Q_5\) with the following
properties.  With \(A,T,U\) as in
\eqref{eq:damp-minimal-gap-four-endpoints}, all four endpoint families
\(S,B,A,T\) use all five coordinates and have no peripheral coordinate.
Writing
\[
 m_X=|\operatorname{Sh}(X)\cap\tbinom{E}{2}|,
 \qquad
 t_X=|\operatorname{Sh}(X)\cap\tbinom{E}{3}|,
\]
one has
\begin{equation}
5\le m_B\le10,qquad 1\le t_B\le5,qquad
 12\le|B|=6+m_B+t_B\le21,
 \label{eq:damp-q5-central-endpoint-profile}
\end{equation}
and the same bounds hold for \(T\).  Consequently
\begin{equation}
11\le|S|,|A|\le20.
 \label{eq:damp-q5-central-complement-orders}
\end{equation}

The new shattered supports
\[
 \mathcal K=\operatorname{Sh}(B)\setminus\operatorname{Sh}(S)
\]
have only sizes two and three.  If
\[
 a=|\mathcal K\cap\tbinom E2|,
 \qquad b=|\mathcal K\cap\tbinom E3|,
\]
then
\begin{equation}
|U|=a+b,qquad
 h(\Psi(B,S))=(0,0,a,b,0,0).
 \label{eq:damp-q5-central-relative-h-vector}
\end{equation}
Every nonzero balanced flag \(h\)-entry is equal to one.  For the
complementary interval \(A\subset T\), the new two-supports are the
complements of the members of \(\mathcal K\cap\binom E3\), and the new
three-supports are the complements of the members of
\(\mathcal K\cap\binom E2\).
\end{proposition}

\begin{proof}
Choose the counterexample by Lemma
\ref{lem:damp-minimal-interval-gap-two-sided}.  The four-coordinate
theorem makes every coordinate active in the two upper endpoints
\(B\) and \(T\).  The upper endpoint \(B\) cannot have a peripheral coordinate, by Lemma
\ref{lem:damp-peripheral-upper-endpoint-reduction}; applying the same
argument to the complementary counterexample \(A\subset T\) shows that
\(T\) has none either.  Set complementation reverses the inclusions of
the two coordinate sections, so it preserves the existence of a
peripheral coordinate.  Hence \(S\) and \(A\) are periphery-free as
well; if either omitted a coordinate, its complement would have nested
sections in that coordinate.  Thus all four endpoints have full support.

For an ample \(X\subseteq Q_E\), complement duality gives
\begin{equation}
R\in\operatorname{Sh}(Q_E\setminus X)
 \quad\Longleftrightarrow\quad
 E\setminus R\notin\operatorname{Sh}(X).
 \label{eq:damp-shattering-complement-duality}
\end{equation}
Since \(A\) uses every coordinate, all singleton supports are shattered
by \(A\); therefore no four-subset is shattered by \(B\).  Neither
\(B\) nor \(A\) is the full cube, so the five-set is not shattered
either.  Thus
\[
 |B|=|\operatorname{Sh}(B)|=1+5+m_B+t_B.
\]

If \(m_B\le4\), Lemma
\ref{lem:damp-sparse-q5-crossing-peripheral} makes \(B\) peripheral.
Hence \(m_B\ge5\).  Applying the same lemma to \(A\), and using
\eqref{eq:damp-shattering-complement-duality}, gives
\[
 5\le m_A=10-t_B,
\]
so \(t_B\le5\).  If \(t_B=0\), then \(B\) has VC dimension at most
two, and Lemma~\ref{lem:damp-relative-peeling-vc-two} supplies the
forbidden augmentation.  Hence \(t_B\ge1\), proving
\eqref{eq:damp-q5-central-endpoint-profile}.  Repeating the argument for
\(A\subset T\) gives the same bounds for \(T\), and complementation
then gives \eqref{eq:damp-q5-central-complement-orders}.

All four endpoints have full support, so their shattered-set systems
agree in ranks zero and one.  The preceding argument excludes ranks
four and five for \(B\), and therefore \(\mathcal K\) is supported in
ranks two and three.  Ampleness gives
\[
 |U|=|B|-|S|
 =|\operatorname{Sh}(B)|-|\operatorname{Sh}(S)|=a+b.
\]
Equation~\eqref{eq:damp-q5-central-relative-h-vector} and the flag
statement now follow from Lemma
\ref{lem:damp-relative-yang-cm-flag-h}.  Applying
\eqref{eq:damp-shattering-complement-duality} to both ends of the
interval interchanges new supports with their set complements, proving
the last assertion.
\end{proof}

\begin{lemma}[Paired defect cylinders force two neighbours in the intermediate layer]
\label{lem:damp-q5-paired-defect-cylinders}
Let \(S\subsetneq B\subseteq Q_E\) be a minimum five-coordinate gap in
the normal form of Proposition
\ref{prop:damp-q5-central-relative-cm-profile}.  Put
\[
 U=B\setminus S,\qquad A=Q_E\setminus B,
 \qquad T=Q_E\setminus S,
\]
and write
\[
 \mathcal K_i=\mathcal K\cap\binom{E}{i},
 \qquad
 \mathcal K^*=\{E\setminus R:R\in\mathcal K\},
 \qquad
 \mathcal H_A(p)=\operatorname{Sh}(A\cup\{p\})
                  \setminus\operatorname{Sh}(A).
\]
Define the eligible supports
\begin{equation}
\mathcal E(\mathcal K)
 =\mathcal K_2\mathbin{\dot\cup}
  \{R\in\mathcal K_3:
       \text{no member of \(\mathcal K_2\) is contained in \(R\)}\}.
 \label{eq:damp-q5-eligible-defect-supports}
\end{equation}
Then, for every \(p\in U\),
\begin{equation}
\mathcal H_S(p)\subseteq\mathcal E(\mathcal K),
 \qquad
 \mathcal H_A(p)\subseteq\mathcal E(\mathcal K^*),
 \qquad
 |\mathcal H_S(p)|,|\mathcal H_A(p)|\ge2.
 \label{eq:damp-q5-two-sided-defect-degree}
\end{equation}

Each \(R\in\mathcal E(\mathcal K)\) has a unique pattern
\(\sigma_R\in Q_R\setminus\pi_R(S)\).  The cylinder
\begin{equation}
C_R=\{x\in Q_E:x|_R=\sigma_R\}
 \label{eq:damp-q5-lower-defect-cylinder}
\end{equation}
is contained in \(T\), and
\(R\in\mathcal H_S(p)\) exactly when \(p\in C_R\).  Dually, every
\(Q\in\mathcal E(\mathcal K^*)\) defines a cylinder \(D_Q\subseteq B\),
and \(Q\in\mathcal H_A(p)\) exactly when \(p\in D_Q\).  Hence
\begin{equation}
p^{\{e\}}\in U
 \quad\text{for every }e\in E\setminus(R\cup Q),
 \quad R\in\mathcal H_S(p),\quad Q\in\mathcal H_A(p).
 \label{eq:damp-q5-cross-cylinder-u-neighbor}
\end{equation}

Let \(N_U(p)=\{e\in E:p^{\{e\}}\in U\}\).  It is nonempty.  If
\(N_U(p)=\{e\}\), then the other four coordinates can be labelled
\(a,b,c,d\) so that
\begin{equation}
\begin{aligned}
  \mathcal H_S(p)&=\{\{a,b,c\},\{a,b,d\}\},\\
  \mathcal H_A(p)&=\{\{a,c,d\},\{b,c,d\}\}.
 \end{aligned}
 \label{eq:damp-q5-split-star-defect-supports}
\end{equation}
In particular,
\begin{equation}
\{e,a\},\{e,b\}\in\mathcal K_2,
 \qquad
 \{e,c\},\{e,d\}
 \in\{E\setminus R:R\in\mathcal K_3\}.
 \label{eq:damp-q5-split-star-support-graphs}
\end{equation}
Moreover,
\begin{equation}
N_S(p)=\{a,b\},\qquad N_A(p)=\{c,d\},\qquad N_U(p)=\{e\},
 \label{eq:damp-q5-split-star-neighbor-partition}
\end{equation}
and \(B\) contains the two squares through \(p\) on
\(\{e,a\}\) and \(\{e,b\}\), while \(T\) contains the two squares
through \(p\) on \(\{e,c\}\) and \(\{e,d\}\).
This singleton alternative is impossible.  Consequently,
\begin{equation}
|N_U(p)|\ge2
 \qquad\text{for every }p\in U.
 \label{eq:damp-q5-two-middle-neighbours}
\end{equation}
\end{lemma}

\begin{proof}
We first convert the two one-concept failures into intersecting signed
cylinders.  We then show that a unique neighbor in \(U\) forces the four
triple supports displayed in
\eqref{eq:damp-q5-split-star-defect-supports}.  Finally, a minimal blocking
vertex and an isometric geodesic in \(S\) rule out that configuration.

The central profile supports \(\mathcal K\) only in ranks two and
three.  Equation~\eqref{eq:damp-relative-singleton-defect-supports}
therefore says exactly that a support contributing to
\(\mathcal H_S(p)\) is either a member of \(\mathcal K_2\), or a member
of \(\mathcal K_3\) containing no member of \(\mathcal K_2\).  This is
the first inclusion in
\eqref{eq:damp-q5-two-sided-defect-degree}; the complementary interval
gives the second.  None of the four one-concept families in
\eqref{eq:damp-minimal-gap-four-failures} is ample, so Lemma
\ref{lem:damp-relative-singleton-support-defect} gives both lower bounds.

For \(R\in\mathcal E(\mathcal K)\), the projection-defect count
\eqref{eq:damp-relative-projection-defect-count} is one.  Since \(B\)
shatters \(R\), the sole pattern absent from \(\pi_R(S)\) is
\(\sigma_R\).  This proves the cylinder statement, and the dual argument
gives \(D_Q\subseteq B\).  If both cylinders contain \(p\), their
intersection is the subcube through \(p\) whose free coordinates are
\(E\setminus(R\cup Q)\).  It lies in
\(B\cap T=U\), proving
\eqref{eq:damp-q5-cross-cylinder-u-neighbor}.

We use the following elementary five-set observation.  If
\(R\in\mathcal E(\mathcal K)\),
\(Q\in\mathcal E(\mathcal K^*)\), and \(R\cup Q=E\), then
\begin{equation}
Q=E\setminus R.
 \label{eq:damp-q5-eligible-complement-only}
\end{equation}
Indeed, two two-sets cannot cover \(E\).  If their sizes are two and
three, equality forces complementation.  If both have size three, then
\(E\setminus Q\in\mathcal K_2\) and
\(E\setminus Q\subseteq R\), contradicting the eligibility of \(R\).

Choose two distinct members of each defect family in
\eqref{eq:damp-q5-two-sided-defect-degree}.  A fixed \(R\) has only one
set complement, so one of the four cross pairs is not complementary.
By \eqref{eq:damp-q5-eligible-complement-only}, its union is not \(E\),
and \eqref{eq:damp-q5-cross-cylinder-u-neighbor} proves
\(N_U(p)\ne\varnothing\).

Suppose now that \(N_U(p)=\{e\}\), and put \(H=E\setminus\{e\}\).
For any two chosen lower supports \(R_1,R_2\) and upper supports
\(Q_1,Q_2\), equation
\eqref{eq:damp-q5-cross-cylinder-u-neighbor} gives
\begin{equation}
H\subseteq R_i\cup Q_j\qquad(1\le i,j\le2).
 \label{eq:damp-q5-split-star-covering-law}
\end{equation}
None of the four supports can be a two-set.  Suppose, for example, that
\(R_1\) is a two-set and put \(C=H\setminus R_1\).  A two-set \(Q_j\)
satisfying \eqref{eq:damp-q5-split-star-covering-law} must equal \(C\).
Then \(E\setminus Q_j=R_1\cup\{e\}\in\mathcal K_3\), while the covering
law forces \(R_2\) to contain \(R_1\); it can neither be the same edge nor
an eligible triple.  Thus both \(Q_j\)'s are triples.  A triple through
\(e\) that covers \(H\) with \(R_1\) is uniquely \(E\setminus R_1\);
if one \(Q_j\) has this form, the covering law again forces \(R_2\) to
contain \(R_1\).  The only remaining two triples are
\(C\cup\{r\}\), one for each \(r\in R_1\).  Their complements are the
two pairs \(\{e,r\}\), and covering both with \(R_2\) once more forces
\(R_1\subseteq R_2\), contradicting distinctness or eligibility.  The
dual argument excludes a two-set among the \(Q_j\)'s.

Nor can one of the four three-sets contain \(e\).  If, say, \(e\in R_i\),
put \(C=H\setminus R_i\), a two-set.  A two-set \(Q_j\) covering \(H\)
with \(R_i\) must equal \(C=E\setminus R_i\), so at most one of the two
distinct \(Q_j\)'s has this form.  Every covering triple contains \(C\),
and its complementary pair lies inside \(R_i\).  Since such a triple is
eligible for \(\mathcal K^*\), that complementary pair belongs to
\(\mathcal K_2\), contradicting the eligibility of \(R_i\).  Hence a
second upper support cannot exist.  The dual argument handles the
\(Q_j\)'s.

The four supports are therefore three-subsets of the four-set \(H\).
Write the two lower supports as \(H\setminus\{c\}\) and
\(H\setminus\{d\}\), and denote the remaining elements by \(a,b\).
The covering law forces each upper support to contain \(c,d\); since the
two are distinct, they are \(H\setminus\{a\}\) and
\(H\setminus\{b\}\).  This is
\eqref{eq:damp-q5-split-star-defect-supports}.  No further lower support
is possible: it would have to contain \(a,b\); the edge \(\{a,b\}\) is
excluded by eligibility of the two displayed lower triples, a triple
through \(e\) contains one of the forced edges \(\{e,a\},\{e,b\}\), and
the two remaining triples are already listed.  The dual argument excludes
further upper supports.

Taking complements of the two upper triples gives
\(\{e,a\},\{e,b\}\in\mathcal K_2\).  Taking complements of the two
lower triples gives the two asserted edges of the complementary graph,
proving \eqref{eq:damp-q5-split-star-support-graphs}.

Finally, the cylinder of \(H\setminus\{c\}\) is the square through
\(p\) with free coordinates \(c,e\), and it lies in \(T\); the other
lower cylinder gives the \(d,e\)-square.  The two upper cylinders give
the \(a,e\)- and \(b,e\)-squares in \(B\).  Since the only
\(U\)-neighbor of \(p\) is its \(e\)-neighbor, its \(a,b\)-neighbors
belong to \(S\), and its \(c,d\)-neighbors belong to \(A\).  This proves
\eqref{eq:damp-q5-split-star-neighbor-partition} and the final assertion.

It remains to rule out this normal form.  By
Lemma~\ref{lem:damp-one-concept-colon-test}, there is a blocking vertex
\(s\in S\) for \(p\).  Choose one for which
\[
 D=\delta(p,s)
\]
is inclusion-minimal.  The blocking condition and
\eqref{eq:damp-q5-split-star-neighbor-partition} give
\(D\cap\{a,b\}=\varnothing\).  Because \(B\) is isometric and
\(p,s\in B\), some first step of a \(B\)-geodesic from \(p\) to \(s\)
flips a coordinate of \(D\).  The \(c\)- and \(d\)-neighbours of \(p\)
belong to \(A\), whereas its \(e\)-neighbour belongs to \(U\).
Thus \(e\in D\).

The two \(T\)-squares in the conclusion contain
\(p^{\{e,c\}}\) and \(p^{\{e,d\}}\), so neither vertex belongs to
\(S\).  Also \(p^{\{e\}}\in U\), not in \(S\).  Hence
\begin{equation}
D=\{c,d,e\},
 \qquad s=p^{\{c,d,e\}}.
 \label{eq:damp-q5-split-star-blocker}
\end{equation}
Minimality of \(D\) says that no vertex \(p^F\), with
\(F\subsetneq D\), belongs to \(S\): such a vertex would itself be a
blocking vertex with a smaller difference set, because
\(F\cap N_S(p)=\varnothing\).

Now \(p^{\{a\}}\in S\).  Consider an \(S\)-geodesic from this vertex
to \(s\), which exists because \(S\) is ample and hence isometric.  The
coordinate \(a\) cannot be removed before all of \(c,d,e\) have been
added: otherwise the geodesic would contain a vertex \(p^F\in S\) for
some proper subset \(F\subsetneq D\).  It follows that
\(p^{\{a,c,d,e\}}\in S\).  The four vertices
\[
 p^{\{b\}},\quad p^{\{a\}},\quad
 p^{\{c,d,e\}},\quad p^{\{a,c,d,e\}}
\]
realize all four traces on \(\{e,a\}\).  Thus
\(\{e,a\}\in\operatorname{Sh}(S)\), contradicting
\(\{e,a\}\in\mathcal K_2
=\operatorname{Sh}(B)\setminus\operatorname{Sh}(S)\) from
\eqref{eq:damp-q5-split-star-support-graphs}.  The singleton alternative
is impossible, proving \eqref{eq:damp-q5-two-middle-neighbours}.
\end{proof}

\begin{corollary}[Minimum degree of the intermediate layer]
\label{cor:damp-q5-minimum-gap-cycle-core}
In the setting of Lemma~\ref{lem:damp-q5-paired-defect-cylinders}, the
subgraph of \(Q_E\) induced by \(U=B\setminus S\) has minimum degree at
least two.  Consequently, each of its connected components contains an
even cycle.  Moreover,
\[
 4\le |U|\le10,
\]
so a shortest such cycle has length \(4\), \(6\), \(8\), or \(10\).
\end{corollary}

\begin{proof}
The degree assertion is \eqref{eq:damp-q5-two-middle-neighbours}.  Every
finite graph of minimum degree at least two contains a cycle, and every
subgraph of a hypercube is bipartite.  Proposition
\ref{prop:damp-q5-central-relative-cm-profile} gives
\(|S|\ge11\) and \(|B|\le21\), so \(|U|\le10\).  This proves the
remaining assertions.
\end{proof}

\begin{lemma}[Persistence of defect labels along an intermediate edge]
\label{lem:damp-q5-defect-edge-transport}
Retain the setting of
Lemma~\ref{lem:damp-q5-paired-defect-cylinders}, and let
\(p,q=p^{\{e\}}\in U\).  For \(X=S\) and for \(X=A\), one has
\begin{equation}
\{R\in\mathcal H_X(p):e\notin R\}
 =\{R\in\mathcal H_X(q):e\notin R\}
 =\mathcal H_X(p)\cap\mathcal H_X(q).
 \label{eq:damp-q5-defect-edge-transport}
\end{equation}
Consequently, along a cycle of the subgraph induced by \(U\), membership of a fixed
support \(R\) in the defect family can change only across an edge whose
color belongs to \(R\).  Across such an edge, \(R\) cannot occur at both
endpoints.  In particular, the number of changes encountered while
traversing the cycle is even.
\end{lemma}

\begin{proof}
If \(e\notin R\), then \(p|_R=q|_R\).  The characterization
\eqref{eq:damp-relative-singleton-defect-supports} therefore gives
\(R\in\mathcal H_X(p)\) if and only if
\(R\in\mathcal H_X(q)\).  If \(e\in R\), the two traces on \(R\) are
different.  An eligible support has a unique pattern missing from the
projection of \(X\), so it cannot belong to both defect families.  This
proves \eqref{eq:damp-q5-defect-edge-transport}.  The cycle assertions
follow by reading these two alternatives around the closed walk.
\end{proof}


\begin{lemma}[Convex signed defect carriers]
\label{lem:damp-q5-defect-carrier-convexity}
Retain the setting and notation of
Lemma~\ref{lem:damp-q5-paired-defect-cylinders}.  For
\(R\in\mathcal E(\mathcal K)\) and
\(Q\in\mathcal E(\mathcal K^*)\), put
\begin{equation}
\Gamma_R^-=B\cap C_R=U\cap C_R,
 \qquad
 \Gamma_Q^+=T\cap D_Q=U\cap D_Q.
 \label{eq:damp-q5-defect-carriers}
\end{equation}
After the coordinates fixed by the corresponding cylinder are suppressed,
each nonempty carrier is ample.  Moreover, its one-inclusion graph is a
convex subgraph of the graph induced by \(U\).  Finally,
\begin{equation}
R\in\mathcal H_S(p)\Longleftrightarrow p\in\Gamma_R^-,
 \qquad
 Q\in\mathcal H_A(p)\Longleftrightarrow p\in\Gamma_Q^+.
 \label{eq:damp-q5-carrier-label-equivalence}
\end{equation}
\end{lemma}

\begin{proof}
The cylinder \(C_R\) is disjoint from \(S\), so
\(B\cap C_R=(S\mathbin{\dot\cup}U)\cap C_R=U\cap C_R\).
Fixing the pattern \(\sigma_R\) on \(R\) is a sequence of restrictions
of the ample family \(B\).  Restrictions preserve ampleness, proving the
first assertion for \(\Gamma_R^-\).  The upper assertion follows in the
same way from \(D_Q\cap A=\varnothing\) and the ample family \(T\).

Let \(x,y\in\Gamma_R^-\).  Since \(B\) is isometric, every \(B\)-geodesic
from \(x\) to \(y\) has their Hamming distance.  Such a geodesic cannot
flip a coordinate of \(R\), since the two endpoints agree on \(R\);
hence it stays in \(C_R\), and therefore in \(\Gamma_R^-\).
In particular, a geodesic of Hamming length lies in \(U\).  Conversely,
every \(U\)-geodesic between \(x\) and \(y\) has at least their Hamming
distance, and equality is attained inside \(\Gamma_R^-\).  A shortest
one cannot flip a coordinate of \(R\), so it also stays in the carrier.
Thus the carrier is convex in the graph induced by \(U\).  The proof for
\(\Gamma_Q^+\) is identical, using \(T\).  The equivalences in
\eqref{eq:damp-q5-carrier-label-equivalence} are the two cylinder
characterizations from Lemma~\ref{lem:damp-q5-paired-defect-cylinders}.
\end{proof}

\begin{corollary}[Interval law on a shortest relative cycle]
\label{cor:damp-q5-shortest-cycle-carrier-intervals}
Let \(C\) be a shortest cycle in the graph induced by \(U\), of length
\(2k\).  After a cyclic shift and a relabelling of its active coordinates,
the edge-color word of \(C\) is
\begin{equation}
e_1,e_2,\ldots,e_k,e_1,e_2,\ldots,e_k,
 \qquad 2\le k\le5.
 \label{eq:damp-q5-isometric-cycle-word}
\end{equation}
The intersection of \(C\) with any signed defect carrier is either the
whole cycle or a single cyclic interval.  If the carrier has support
\(R\), a proper interval has at most \(5-|R|\) edges, and the colors of
its two boundary edges belong to \(R\).

If \(k\ge3\), the whole-cycle alternative is impossible.  Consequently,
on a shortest ten-cycle every pair carrier contains at most four
consecutive cycle vertices, every triple carrier contains at most three,
and the two boundary-edge colors of each carrier are distinct.
\end{corollary}

\begin{proof}
A shortest cycle in any graph is isometric: a path shorter than the
shorter of its two cycle arcs would create a shorter cycle.  Hence every
arc of \(C\) with at most \(k\) edges is a hypercube geodesic and uses
each coordinate at most once.  The first \(k\) edge colors are therefore
distinct.  The other half is another geodesic between the same two
vertices and uses the same \(k\) colors.  Applying the distinctness
statement to every cyclic block of \(k\) edges gives
\(e_{i+k}=e_i\), proving
\eqref{eq:damp-q5-isometric-cycle-word}.

By Lemma~\ref{lem:damp-q5-defect-carrier-convexity}, a carrier is convex
in the graph induced by \(U\).  Its intersection with the isometric cycle
is therefore geodesically convex in that cycle, hence is the whole cycle
or one interval.  The endpoints of a proper interval agree on the
\(|R|\) fixed coordinates, so their Hamming distance, which equals the
length of the interval, is at most \(5-|R|\).  If a boundary edge had a
color outside \(R\), its two endpoints would have the same trace on
\(R\), contrary to maximality of the interval.

Suppose that a carrier contains \(C\) and \(k\ge3\).  The canonical cycle
in \eqref{eq:damp-q5-isometric-cycle-word} shatters any pair of its active
coordinates.  The carrier is ample, so it strongly shatters that pair
and contains a square.  This square lies in \(U\), contradicting the
choice of a cycle of length greater than four as shortest.  Thus every
carrier interval is proper when \(k\ge3\).  On a ten-cycle, two
occurrences of the same edge color are separated by five edges.  An
interval bounded by them has more than \(5-|R|\) edges for
\(|R|\in\{2,3\}\), so the boundary colors are distinct.
\end{proof}

\begin{corollary}[Endpoint profile forced by a shortest ten-cycle]
\label{cor:damp-q5-ten-cycle-endpoint-profile}
If a shortest cycle in \(U\) has length ten, then it contains every
vertex of \(U\), and
\begin{equation}
\begin{array}{c|cccc}
 X&S&A&B&T\\ \hline
 |X|&11&11&21&21\\
 m_X&5&5&10&10\\
 t_X&0&0&5&5.
 \end{array}
 \label{eq:damp-q5-ten-cycle-endpoint-table}
\end{equation}
In particular,
\begin{equation}
|\mathcal K_2|=|\mathcal K_3|=5.
 \label{eq:damp-q5-ten-cycle-support-split}
\end{equation}

The lower and upper defect carriers give two covers of this ten-cycle,
each with multiplicity at least two at every vertex.  Their traces on the
cycle are the intervals described in
Corollary~\ref{cor:damp-q5-shortest-cycle-carrier-intervals}.  If a lower
carrier of support \(R\) and an upper carrier of support \(Q\) meet, then
\begin{equation}
|R\cup Q|\ge4.
 \label{eq:damp-q5-ten-cycle-cross-cover}
\end{equation}
When equality holds, the sole color outside \(R\cup Q\) is the color of
one of the two cycle edges incident with every common vertex.
\end{corollary}

\begin{proof}
The cycle has ten distinct vertices, while
Corollary~\ref{cor:damp-q5-minimum-gap-cycle-core} gives \(|U|\le10\);
hence \(U\) is exactly the cycle.  Its canonical word uses all five
coordinates and realizes all four patterns on every coordinate pair.
Thus \(m_B=m_T=10\).  Write \(t_B,t_T\) for the numbers of shattered
triples.  The central profile gives
\[
 |B|=16+t_B,\qquad |T|=16+t_T,\qquad t_B,t_T\le5.
\]
Since \(S=Q_E\setminus T\) and \(|B|-|S|=|U|=10\), one obtains
\(t_B+t_T=10\).  Therefore \(t_B=t_T=5\), and complement duality gives
\[
 m_S=m_A=5,\qquad t_S=t_A=0.
\]
This proves \eqref{eq:damp-q5-ten-cycle-endpoint-table}.
Taking the differences of the two shattered-set systems gives
\eqref{eq:damp-q5-ten-cycle-support-split}.

The multiplicity-two assertion is
\eqref{eq:damp-q5-two-sided-defect-degree}, and the interval description
is Corollary~\ref{cor:damp-q5-shortest-cycle-carrier-intervals}.  If two
opposite carriers meet at \(p\), their defining cylinders have intersection
equal to the full subcube through \(p\) on the coordinates outside
\(R\cup Q\), and that subcube lies in \(U\).  The ten-cycle contains no
square, proving \eqref{eq:damp-q5-ten-cycle-cross-cover}.  If the union
has size four, the intersection is an edge of \(U\); since \(U\) is the
cycle, it must be one of the two cycle edges at each of its vertices.
\end{proof}

\begin{theorem}[Five-color circular-carrier exclusion]
\label{thm:damp-q5-circular-carrier-exclusion}
Let \(C\) be the canonical edge-colored ten-cycle in
\eqref{eq:damp-q5-isometric-cycle-word}.  There is no support system
\(\mathcal K\subseteq\binom E2\mathbin{\dot\cup}\binom E3\), with
\(|\mathcal K_2|=|\mathcal K_3|=5\), having the following data.
For every support in \(\mathcal E(\mathcal K)\), choose one trace pattern
on \(C\); do the same for every support in
\(\mathcal E(\mathcal K^*)\).  Require every chosen trace fiber to be a
nonempty convex interval, every cycle vertex to lie in at least two
chosen fibers on each side, and every intersecting lower--upper pair of
fibers to have supports whose union has size at least four.
\end{theorem}

\begin{proof}
Write
\[
 X=\mathcal K_2,
 \qquad
 Z=\{E\setminus R:R\in\mathcal K_3\}.
\]
Thus \(X\) and \(Z\) are five-element families of two-subsets of
\(E\).  Join \(x\in X\) to \(z\in Z\) when \(x\cap z=\varnothing\),
and let \(X_0\) and \(Z_0\) be the isolated vertices in the two parts.
By \eqref{eq:damp-q5-eligible-defect-supports}, the lower eligible
supports are
\[
 X\mathbin{\dot\cup}\{E\setminus z:z\in Z_0\},
\]
and the upper eligible supports are
\[
 Z\mathbin{\dot\cup}\{E\setminus x:x\in X_0\}.
\]
We call the pair supports mandatory and the displayed triple supports
extra.

Every cycle vertex lies in an extra carrier on at least one side.
Indeed, otherwise two mandatory lower carriers and two mandatory upper
carriers contain that vertex.  The cross-union condition says that each
of the two lower pairs is disjoint from each of the two upper pairs,
which gives a four-cycle in the bipartite disjointness graph.  This is
impossible by Lemma~\ref{lem:damp-q5-artinian-disjointness-graph}.
Every triple trace fiber on the canonical ten-cycle has at most three
vertices, by
Corollary~\ref{cor:damp-q5-shortest-cycle-carrier-intervals}.  Hence
\begin{equation}
10\le 3\bigl(|X_0|+|Z_0|\bigr).
 \label{eq:damp-q5-ten-cycle-isolate-cover}
\end{equation}

We next note that \(|X_0|,|Z_0|\le2\).  To see this, three distinct
edges of \(K_5\) have at most four common edge transversals, where an
edge transversal means an edge meeting all three.  If two of the three
edges are disjoint, their four cross edges are the only possibilities.
Otherwise choose two that meet at \(c\).  If the third also contains
\(c\), only the four edges through \(c\) can meet all three; if it does
not, at most two of those four edges and the edge joining the two other
endpoints of the chosen pair can do so.  If \(X_0\) contained three
members, every member of the five-element family \(Z\) would be such a
transversal, a contradiction.  The argument with \(X\) and \(Z\)
interchanged is identical.  Equation
\eqref{eq:damp-q5-ten-cycle-isolate-cover} now gives
\begin{equation}
|X_0|=|Z_0|=2.
 \label{eq:damp-q5-ten-cycle-two-isolates}
\end{equation}

For \(x\in X_0\), consider its nonempty mandatory lower pair fiber.
It cannot meet any mandatory upper pair fiber, because \(x\) meets
every member of \(Z\).  Both upper labels at each of its vertices must
therefore be the two extra carriers.  In particular, the pair fiber is
contained in the upper triple fiber with support \(E\setminus x\).
These two trace patterns together fix every coordinate, so their fibers
meet in at most one cube vertex.  The mandatory pair fiber is therefore
a singleton.  Dually, the mandatory upper pair fiber of every
\(z\in Z_0\) is a singleton.

Index the colors in their cyclic order in
\eqref{eq:damp-q5-isometric-cycle-word}.  A pair trace fiber on \(C\)
is a singleton exactly when its two colors are consecutive in this
order.  Thus \(X_0\) and \(Z_0\) are two-edge subsets of the color
pentagon.  Every edge in \(X_0\) meets every edge in \(Z_0\).  Two
nonadjacent pentagon edges have only one common pentagon-edge
transversal, so the two members of \(X_0\) are adjacent.  The only
pentagon edges meeting both are those two edges themselves.  Therefore
\(X_0=Z_0\).  After a cyclic relabelling, and writing
\(ij\) for \(\{i,j\}\),
\begin{equation}
X_0=Z_0=\{01,04\}.
 \label{eq:damp-q5-ten-cycle-isolate-normal-form}
\end{equation}
The five edges of \(K_5\) meeting both \(01\) and \(04\) are
\(01,02,03,04,14\).  Since every member of \(X\) and of \(Z\) must
meet both isolated edges, it follows that
\begin{equation}
X=Z=\{01,02,03,04,14\}.
 \label{eq:damp-q5-ten-cycle-support-normal-form}
\end{equation}

It remains only to read the forced trace intervals.  Write the vertices
of \(C\) as \(v_0,\ldots,v_9\), with the edge
\(v_iv_{i+1}\) colored \(i\pmod5\).  The singleton fibers of \(01\)
are \(\{v_1\}\) and \(\{v_6\}\), while those of \(04\) are
\(\{v_0\}\) and \(\{v_5\}\).  By the antipodal symmetry, suppose
the lower \(01\)-fiber is \(\{v_1\}\).  It is contained in both upper
extra fibers, so the latter are forced to be
\begin{equation}
\{v_0,v_1,v_2\}\quad\text{on }234,
 \qquad
 \{v_9,v_0,v_1\}\quad\text{on }123.
 \label{eq:damp-q5-ten-cycle-upper-extra-fibers}
\end{equation}
Their common vertices force the lower \(04\)-fiber to be
\(\{v_0\}\).

The upper \(01\)-fiber cannot be \(\{v_1\}\): the two corresponding
lower extra fibers would then equal the intervals in
\eqref{eq:damp-q5-ten-cycle-upper-extra-fibers}, whereas lower and upper
fibers with the same triple support may not meet.  Hence the upper
\(01\)-fiber is \(\{v_6\}\), and the same containment argument forces
the lower extra fibers to be
\begin{equation}
\{v_5,v_6,v_7\}\quad\text{on }234,
 \qquad
 \{v_4,v_5,v_6\}\quad\text{on }123,
 \label{eq:damp-q5-ten-cycle-lower-extra-fibers}
\end{equation}
and the upper \(04\)-fiber to be \(\{v_5\}\).

The vertex \(v_3\) lies in none of the four extra fibers and none of
the four isolated mandatory fibers.  Its two lower and two upper labels
must therefore be chosen from the pair supports
\(02,03,14\).  Every chosen lower pair must be disjoint from every
chosen upper pair, again producing a four-cycle in the bipartite
disjointness graph.  This final contradiction proves the theorem.
\end{proof}

\begin{corollary}[Exclusion of the ten-cycle branch]
\label{cor:damp-q5-no-ten-cycle-branch}
In a minimum five-coordinate gap, a shortest cycle in the intermediate
layer has length \(4\), \(6\), or \(8\).
\end{corollary}

\begin{proof}
A shortest ten-cycle would have the exact endpoint and carrier profile
of Corollary~\ref{cor:damp-q5-ten-cycle-endpoint-profile}, contradicting
Theorem~\ref{thm:damp-q5-circular-carrier-exclusion}.
\end{proof}

\begin{lemma}[Exhaustion by a shortest eight-cycle]
\label{lem:damp-q5-eight-cycle-exhaustion}
If a minimum five-coordinate gap has a shortest intermediate-layer cycle
\(C\) of length eight, then \(U=V(C)\).  Moreover, the vertices of
\(U\) agree on the unique coordinate not used by the edges of \(C\).
\end{lemma}

\begin{proof}
Corollary~\ref{cor:damp-q5-minimum-gap-cycle-core} gives \(|U|\le10\),
so at most two vertices of \(U\) lie outside \(C\).  No such vertex can
have two neighbours on \(C\).  Indeed, those two cycle vertices would
be at Hamming distance two.  Since \(C\) is isometric by
Corollary~\ref{cor:damp-q5-shortest-cycle-carrier-intervals}, their
two-edge path on \(C\), together with the outside vertex, would give a
four-cycle in \(U\), contrary to the choice of \(C\).

One outside vertex would consequently have degree at most one in \(U\),
contradicting \eqref{eq:damp-q5-two-middle-neighbours}.  If there were
two, each would have to be adjacent to the other and to a distinct
cycle vertex.  They would give a three-edge path between two vertices of
\(C\).  The distance of those endpoints along the isometric cycle is
odd and at most three, so the outside path and a shortest cycle arc
would form a four- or six-cycle in \(U\), again a contradiction.  Thus
\(U=V(C)\).

The repeated color word
\eqref{eq:damp-q5-isometric-cycle-word} uses exactly four coordinates
when \(|C|=8\).  No edge of \(C=U\) flips the remaining coordinate,
so all vertices of \(U\) agree there.
\end{proof}

\begin{lemma}[Elimination of an unused coordinate]
\label{lem:damp-q5-unused-coordinate-elimination}
Let \(S\subsetneq B\subseteq Q_E\) be a minimum five-coordinate gap,
and suppose that all vertices of \(U=B\setminus S\) agree in coordinate
\(r\).  Orient \(r\) so that \(U\) lies in its zero-section, and write
the two sections of the upper and lower endpoints as
\begin{equation}
B=(M\times\{0\})\cup(Y\times\{1\}),
 \qquad
 S=(L\times\{0\})\cup(Y\times\{1\}),
 \qquad M\setminus L=U.
 \label{eq:damp-q5-unused-coordinate-sections}
\end{equation}
Then
\begin{equation}
U\cap Y=\varnothing.
 \label{eq:damp-q5-unused-coordinate-unique-fibers}
\end{equation}
Moreover, no member of
\(\mathcal K=\operatorname{Sh}(B)\setminus\operatorname{Sh}(S)\)
contains \(r\).
\end{lemma}

\begin{proof}
Put \(I_M=M\cap Y\) and \(I_L=L\cap Y\).  The two-section shattering
calculation used in Lemma
\ref{lem:damp-expansion-separator-duplication} shows that both
intersections are ample.  Suppose that
\(I_L\subsetneq I_M\).  If \(I_L=\varnothing\), choose any
\(p\in I_M\); otherwise Theorem
\ref{thm:damp-nested-ample-relative-peeling-dimension-four} gives
\(p\in I_M\setminus I_L\) for which
\(I_L\cup\{p\}\) is ample.  In either case,
Lemma~\ref{lem:damp-expansion-separator-duplication}, applied to the
lower endpoint \(S\), shows that adjoining the missing zero-lift of
\(p\) preserves ampleness.  This contradicts the definition of a
minimum gap.  Thus \(I_L=I_M\), which is
\eqref{eq:damp-q5-unused-coordinate-unique-fibers}.

For an ample two-section family, the supports containing \(r\) are
obtained from the intersection of the shattered sets of its two
sections.  Equation
\eqref{eq:damp-expansion-section-shattering-intersection} applied to
the two endpoints, together with \(I_L=I_M\), gives
\[
 \operatorname{Sh}(M)\cap\operatorname{Sh}(Y)
 =\operatorname{Sh}(I_M)
 =\operatorname{Sh}(I_L)
 =\operatorname{Sh}(L)\cap\operatorname{Sh}(Y).
\]
Thus the supports containing \(r\) are the same at the two endpoints.
\end{proof}

\begin{theorem}[Exclusion of the eight-cycle branch]
\label{thm:damp-q5-no-eight-cycle-branch}
In a minimum five-coordinate gap, a shortest cycle in the intermediate
layer has length \(4\) or \(6\).
\end{theorem}

\begin{proof}
Suppose that a shortest cycle has length eight, and use the section
notation of \eqref{eq:damp-q5-unused-coordinate-sections}.  Put
\(\mathcal K=\operatorname{Sh}(B)\setminus\operatorname{Sh}(S)\).
Lemma~\ref{lem:damp-q5-unused-coordinate-elimination} says that no
member of \(\mathcal K\) contains \(r\).

Write \(a=|\mathcal K_2|\) and \(b=|\mathcal K_3|\).  The central
profile and \(|U|=8\) give \(a+b=8\).  All these supports use the four
active cycle colors, so \(a\le6\) and \(b\le4\); hence
\begin{equation}
a\ge4,
 \qquad
 b\ge2.
 \label{eq:damp-q5-eight-cycle-support-bounds}
\end{equation}

Let
\[
 X=\mathcal K_2,
 \qquad
 Z=\{E\setminus R:R\in\mathcal K_3\}.
\]
Each member of \(Z\) has the form \(\{r,i\}\), where \(i\) is an
active color.  As in the proof of
Theorem~\ref{thm:damp-q5-circular-carrier-exclusion}, join
\(x\in X\) to \(z\in Z\) when they are disjoint.  A lower eligible
triple would correspond to an isolated \(z=\{r,i\}\), which would mean
that every pair in \(X\) contains \(i\).  Only three active pairs
contain a fixed color, whereas \(a\ge4\).  Thus every lower eligible
support is a pair.

An upper eligible triple corresponds to an isolated \(x\in X\).  Such
an \(x\) must contain the active color of every member of \(Z\).  Since
\(b\ge2\), there is at most one such \(x\), and hence at most one upper
triple carrier.  Consider a cycle vertex outside that carrier, or any
cycle vertex if the carrier does not exist.  Its two lower labels and
two upper labels are pair supports.  The cross-cylinder law makes each
chosen lower pair disjoint from each chosen upper pair, producing a
four-cycle in the bipartite disjointness graph.  This contradicts
Lemma~\ref{lem:damp-q5-artinian-disjointness-graph}.  Therefore the sole
upper triple carrier would have to contain the entire eight-cycle.  This
is impossible by
Corollary~\ref{cor:damp-q5-shortest-cycle-carrier-intervals}.  The
eight-cycle branch is excluded.
\end{proof}

\begin{lemma}[Small-ear normal form around a shortest six-cycle]
\label{lem:damp-q5-six-cycle-small-ear-normal-form}
Suppose that a minimum five-coordinate gap has a shortest
intermediate-layer cycle \(C\) of length six, and put
\(W=U\setminus V(C)\).  Then the graph induced by \(U\) is connected.
If \(W\ne\varnothing\), its induced graph is one of
\begin{equation}
P_3,
 \qquad
 P_4,
 \qquad
 K_{1,3}.
 \label{eq:damp-q5-six-cycle-ear-trees}
\end{equation}
No vertex of \(W\) has two neighbours on \(C\), and every leaf of the
tree induced by \(W\) has exactly one neighbour on \(C\).

More generally, suppose that \(w,w'\in W\) have respective cycle
neighbours \(u,u'\), and let \(t\) be their distance in the tree induced
by \(W\).  Then
\begin{equation}
d_C(u,u')+t\ge4,
 \qquad
 d_C(u,u')\equiv t\pmod2.
 \label{eq:damp-q5-six-cycle-ear-attachment-law}
\end{equation}
Consequently, if \(W=P_3\), only its endpoints attach to \(C\), and
their cycle neighbours are at distance two.  If \(W=K_{1,3}\), only
its leaves attach to \(C\), at the three alternating cycle vertices.
\end{lemma}

\begin{proof}
By \eqref{eq:damp-q5-two-middle-neighbours}, every component of the
graph induced by \(U\) has minimum degree at least two.  Hence each
contains a cycle.  A component disjoint from \(C\) would contain at
least six vertices, giving \(|U|\ge12\), contrary to
Corollary~\ref{cor:damp-q5-minimum-gap-cycle-core}.  Thus \(U\) is
connected and \(|W|\le4\).

No vertex outside \(C\) has two cycle neighbours: as in the proof of
Lemma~\ref{lem:damp-q5-eight-cycle-exhaustion}, isometry of \(C\) would
turn those two edges into a four-cycle.  The graph induced by \(W\) is
a forest, because it has at most four vertices and \(U\) has no
four-cycle.  Each of its leaves must therefore use its remaining
incident edge to attach to \(C\), and the attachment is unique.

A component of this forest cannot have one vertex, by minimum degree.
Nor can it have two vertices.  In the latter case both vertices attach
to \(C\), giving a three-edge path between two cycle vertices.  Parity
and the absence of a four-cycle force those endpoints to be antipodal on
\(C\).  The path is then a cube geodesic of length three and uses the
three colors of \(C\).  It lies in the same three-cube as \(C\).  But a
canonical six-cycle in a three-cube omits an antipodal pair, whereas the
two internal vertices of the path are adjacent.  This is impossible.
Thus every forest component has at least three vertices.  Since
\(|W|\le4\), there is only one, and the three possible trees are exactly
those in \eqref{eq:damp-q5-six-cycle-ear-trees}.

For the attachment law, the path through \(W\) and its two attachment
edges has length \(t+2\).  Together with a shortest \(u,u'\)-arc of
\(C\), it forms a cycle in \(U\).  Its length is even and at least six,
which gives both assertions in
\eqref{eq:damp-q5-six-cycle-ear-attachment-law}.

If \(W=P_3\), an attachment at its middle vertex would have to be
antipodal on \(C\) to each of the two leaf attachments.  The two leaves
would then attach to the same cycle vertex and form a four-cycle through
the middle vertex.  Hence only the leaves attach, and
\eqref{eq:damp-q5-six-cycle-ear-attachment-law} puts their cycle
neighbours at distance two.  The same argument excludes an attachment
at the center of \(K_{1,3}\).  Its three leaf attachments are pairwise
at cycle distance two, so they are the three alternating vertices of
\(C\).
\end{proof}

\begin{corollary}[Elimination of the bare and three-vertex-ear six-cycles]
\label{cor:damp-q5-six-cycle-no-small-ear}
In the setting of
Lemma~\ref{lem:damp-q5-six-cycle-small-ear-normal-form}, the outside
graph \(W\) is either \(P_4\) or \(K_{1,3}\).  In particular,
\begin{equation}
|U|=10.
 \label{eq:damp-q5-six-cycle-ten-vertex-layer}
\end{equation}
\end{corollary}

\begin{proof}
Put
\(\mathcal K=\operatorname{Sh}(B)\setminus\operatorname{Sh}(S)\).
Suppose first that \(W=\varnothing\).  Then \(U=C\) uses three
coordinates and is constant in the other two.  Applying
Lemma~\ref{lem:damp-q5-unused-coordinate-elimination} to both unused
coordinates shows that no support in \(\mathcal K\) contains either of
them.  The central profile supports \(\mathcal K\) only in ranks two
and three, so there are at most
\(\binom32+\binom33=4\) possible supports.  This contradicts
\(|\mathcal K|=|U|=6\).

Now suppose that \(W=P_3\).  Its two leaf attachments are separated by
a two-edge arc of \(C\).  The path through \(W\), together with that
arc, is another shortest six-cycle.  If the two arc colors are \(e_1\)
and \(e_2\), its canonical repeated color word makes the four-edge ear
word
\begin{equation}
f,e_1,e_2,f
 \label{eq:damp-q5-six-cycle-three-ear-word}
\end{equation}
up to reversal.  The color \(f\) is not the third color of \(C\): if
it were, all three vertices of \(W\) would lie in the three-cube of
\(C\), which has only the two vertices omitted by its canonical
six-cycle.  Thus \(U\) uses exactly four coordinates and is constant in
the fifth, say \(r\).

Lemma~\ref{lem:damp-q5-unused-coordinate-elimination} implies that
every member of \(\mathcal K\) avoids \(r\).  Write
\(a=|\mathcal K_2|\) and \(b=|\mathcal K_3|\).  Since
\(a+b=|U|=9\), while there are only six pairs and four triples on the
four active colors,
\begin{equation}
a\ge5,
 \qquad
 b\ge3.
 \label{eq:damp-q5-six-cycle-three-ear-support-bounds}
\end{equation}
As before, let \(X=\mathcal K_2\) and
\(Z=\{E\setminus R:R\in\mathcal K_3\}\).  Each member of \(Z\) is
\(\{r,i\}\) for an active color \(i\).  No member of \(Z\) is isolated
in the bipartite disjointness graph: isolation would require all at
least five pairs in \(X\) to contain the same color, but only three
active pairs do.  No member of \(X\) is isolated either, because it
would have to contain the at least three distinct active colors indexed
by \(Z\).  Hence both eligible-support families contain only pairs.
At any vertex of \(U\), two lower and two upper defect labels therefore
give a four-cycle in the bipartite disjointness graph, contradicting
Lemma~\ref{lem:damp-q5-artinian-disjointness-graph}.  Thus \(W\ne P_3\),
and the small-ear normal form proves the assertion.
\end{proof}

\begin{lemma}[A ten-vertex intermediate layer uses all five colors]
\label{lem:damp-q5-ten-layer-five-colors}
In a minimum five-coordinate gap with \(|U|=10\), the graph induced by
\(U\) uses all five coordinates.
\end{lemma}

\begin{proof}
Suppose that every vertex of \(U\) agrees in coordinate \(r\).  By
Lemma~\ref{lem:damp-q5-unused-coordinate-elimination}, every support in
\(\mathcal K=\operatorname{Sh}(B)\setminus\operatorname{Sh}(S)\)
avoids \(r\).  There are exactly
\(\binom42+\binom43=10\) rank-two and rank-three supports on the other
four colors.  Since the central profile gives
\(|\mathcal K|=|U|=10\), all of them belong to \(\mathcal K\).

Every lower eligible support is consequently a pair: each active triple
contains an active pair in \(\mathcal K\).  After complementation, every
upper eligible support is also a pair.  Two lower and two upper labels at
any vertex of \(U\) would give a four-cycle in the bipartite disjointness
graph, contrary to Lemma
\ref{lem:damp-q5-artinian-disjointness-graph}.  Thus no coordinate is
unused by \(U\).
\end{proof}

\begin{proposition}[Five-color bowtie normal form for the six-cycle branch]
\label{prop:damp-q5-six-cycle-bowtie-normal-form}
If a minimum five-coordinate gap has a shortest cycle of length six,
then, after a cube symmetry, its intermediate layer is
\begin{equation}
\begin{split}
 U={}&\{\varnothing,0,01,012,12,2\}\\
    &{}\cup\{\varnothing,0,04,034,34,3\}.
 \end{split}
 \label{eq:damp-q5-six-cycle-bowtie}
\end{equation}
Thus \(U\) is the union of a canonical six-cycle on colors
\(\{0,1,2\}\) and a canonical six-cycle on colors \(\{0,3,4\}\),
sharing exactly the edge \(\varnothing 0\).  Moreover,
\begin{equation}
|\mathcal K_2|=|\mathcal K_3|=5.
 \label{eq:damp-q5-six-cycle-bowtie-support-split}
\end{equation}
\end{proposition}

\begin{proof}
By Corollary~\ref{cor:damp-q5-six-cycle-no-small-ear}, the outside tree
is \(P_4\) or \(K_{1,3}\), and \(|U|=10\).

Suppose first that it is a claw.  Its three leaves attach to alternating
vertices of the base cycle.  For each pair of leaves, the path through
the center and the two-edge arc between their attachments form a
canonical six-cycle.  Reading its repeated color word shows that the
three leaf-to-cycle edges have one common color \(f\), while the three
center-to-leaf edges have the three base-cycle colors.  The common color
\(f\) is different from every base color, because each of the three
two-edge arcs omits a different base color.  Hence \(U\) uses only four
coordinates, contradicting
Lemma~\ref{lem:damp-q5-ten-layer-five-colors}.

Now let the outside tree be the path \(P_4\), and let \(u,v\) be the
cycle neighbours of its leaves.  The attachment law gives
\(d_C(u,v)\in\{1,3\}\).  If the distance is three, the five-edge path
between the antipodal vertices \(u,v\) flips the three base colors an
odd number of times and all other colors an even number of times.  Its
two excess flips therefore use one color twice.  That color cannot be a
base color, for then all four internal vertices would lie in the
three-cube of \(C\), outside a six-cycle that omits only two vertices.
Thus the path, and hence \(U\), uses only four colors, again contradicting
Lemma~\ref{lem:damp-q5-ten-layer-five-colors}.

It remains that \(u,v\) are adjacent, say by an edge of color \(0\).
Together with that edge, the five-edge path is a canonical six-cycle.
Its path word is
\begin{equation}
f,g,0,f,g
 \label{eq:damp-q5-six-cycle-four-ear-word}
\end{equation}
up to reversal and interchange of \(f,g\).  If both \(f,g\) were base
colors, the four path vertices would again have to lie among the two
vertices of the base three-cube outside \(C\).  If exactly one were a
base color, \(U\) would use only four coordinates.  Therefore \(f,g\)
are the two colors outside the base three-cube.  Relabel them as
\(3,4\); the resulting vertex set is exactly
\eqref{eq:damp-q5-six-cycle-bowtie}.  A direct Hamming-distance check
also shows that the internal path vertices have no further neighbours
on the base cycle.

Finally, write \(m_B,t_B,m_T,t_T\) for the numbers of shattered pairs
and triples at the two upper endpoints.  The central profile gives
\[
 |B|=6+m_B+t_B,
 \qquad
 |T|=6+m_T+t_T.
\]
Since \(|B|-|S|=10\) and \(T=Q_E\setminus S\), their sum is \(42\).
The four support counts therefore sum to \(30\), their maximum possible
value.  Hence
\(m_B=m_T=10\) and \(t_B=t_T=5\).  Complement duality gives five new
pairs and five new triples, proving
\eqref{eq:damp-q5-six-cycle-bowtie-support-split}.
\end{proof}

\begin{lemma}[Trace fibers on the five-color bowtie]
\label{lem:damp-q5-bowtie-trace-table}
Let \(U\) be the bowtie in
\eqref{eq:damp-q5-six-cycle-bowtie}.  Every trace fiber on a
three-element support has at most three vertices.

Now let \(x_1,x_2\in\binom E2\) be distinct and intersecting.  For each
\(i\), choose a singleton trace fiber \(F_i\) on \(x_i\) and a trace
fiber \(H_i\) on \(E\setminus x_i\).  If
\[
 F_1\cup F_2\subseteq H_1\cap H_2,
\]
then, up to interchanging \(x_1,x_2\), the complete list of possibilities
is the following.  The last column records \(H_1\cup H_2\).
\begin{equation}
\begin{array}{c|c|c}
 \{x_1,x_2\}&F_1\cup F_2&H_1\cup H_2\\ \hline
 \{01,02\}&\{12,012\}&\{01,012,12,2\}\\
 \{01,12\}&\{12,2\}&\{\varnothing,012,12,2\}\\
 \{02,12\}&\{012,01\}&\{0,01,012,12\}\\ \hline
 \{03,04\}&\{034,34\}&\{04,034,34,3\}\\
 \{03,34\}&\{034,04\}&\{0,04,034,34\}\\
 \{04,34\}&\{34,3\}&\{\varnothing,3,34,034\}.
\end{array}
\label{eq:damp-q5-bowtie-trace-table}
\end{equation}
In particular, \(x_1,x_2\) are two edges of one of the color triangles
\(\{0,1,2\}\) and \(\{0,3,4\}\).  Both singleton fibers and both
complementary triple fibers lie in the corresponding six-cycle.
\end{lemma}

\begin{proof}
The bowtie contains no square.  A trace fiber on a three-element support
lies in a two-dimensional subcube, so four vertices in such a fiber
would form a square.  This proves the first assertion.

For the second assertion, the pair supports admitting a singleton trace
fiber are
\[
 01,\ 02,\ 12,\ 03,\ 04,\ 34.
\]
Their singleton fibers are, respectively,
\[
 \{12\},\quad \{012\},\quad \{01\}\text{ or }\{2\},
 \quad \{034\},\quad \{34\},\quad
 \{3\}\text{ or }\{04\}.
\]
Two complementary triple fibers \(H_1,H_2\) containing both singleton
vertices force those vertices to agree outside
\(x_1\cap x_2\).  Thus they are the endpoints of an edge of that common
color.  Comparing the displayed singleton fibers leaves exactly the six
rows of \eqref{eq:damp-q5-bowtie-trace-table}.  Reading the two
complementary projections gives the last column.
\end{proof}

\begin{theorem}[Five-color bowtie-carrier exclusion]
\label{thm:damp-q5-bowtie-carrier-exclusion}
Let \(U\) be the bowtie in
\eqref{eq:damp-q5-six-cycle-bowtie}.  There is no support system
\(\mathcal K\subseteq\binom E2\mathbin{\dot\cup}\binom E3\), with
\[
 |\mathcal K_2|=|\mathcal K_3|=5,
\]
having the following data.  For every support in
\(\mathcal E(\mathcal K)\), choose one nonempty trace fiber on \(U\);
do the same for every support in \(\mathcal E(\mathcal K^*)\).  Require
every vertex of \(U\) to lie in at least two chosen fibers on each side,
and require every intersecting lower--upper pair of fibers to have
supports whose union has size at least four.
\end{theorem}

\begin{proof}
Put
\[
 X=\mathcal K_2,\qquad
 Z=\{E\setminus R:R\in\mathcal K_3\},
\]
and join \(x\in X\) to \(z\in Z\) when \(x\cap z=\varnothing\).
Let \(X_0,Z_0\) be the isolated vertices in the two parts.  As in the
proof of Theorem~\ref{thm:damp-q5-circular-carrier-exclusion}, the
mandatory supports are the pairs in \(X\) and \(Z\), while the extra
supports are
\[
 \{E\setminus z:z\in Z_0\}
 \quad\text{and}\quad
 \{E\setminus x:x\in X_0\}.
\]

Every vertex of \(U\) lies in an extra fiber on at least one side.
Otherwise its two lower and two upper labels are pair supports, and the
cross-union condition makes each lower pair disjoint from each upper
pair.  This is a four-cycle in the bipartite disjointness graph,
contrary to Lemma~\ref{lem:damp-q5-artinian-disjointness-graph}.
Lemma~\ref{lem:damp-q5-bowtie-trace-table} bounds every extra fiber by
three vertices, whence
\[
 10\le3\bigl(|X_0|+|Z_0|\bigr).
\]

The common-transversal argument in the proof of
Theorem~\ref{thm:damp-q5-circular-carrier-exclusion} gives
\(|X_0|,|Z_0|\le2\): five edges cannot all meet three prescribed
distinct edges of \(K_5\).  Consequently
\[
 |X_0|=|Z_0|=2.
\]
Write \(X_0=\{x_1,x_2\}\).  The nonempty lower pair fiber on \(x_i\)
cannot meet a mandatory upper pair fiber, since \(x_i\) meets every
member of \(Z\).  At each of its vertices, the two required upper
labels must therefore be the two extra upper fibers.  In particular,
it is contained in the extra fiber on \(E\setminus x_i\), so fixing the
two complementary supports makes it a singleton.  Both singleton
fibers are contained in both extra upper fibers.

There are five members of \(Z\), all meeting both \(x_1,x_2\).
The two edges \(x_1,x_2\) cannot be disjoint, because only their four
cross edges meet both.  They are therefore intersecting, and
Lemma~\ref{lem:damp-q5-bowtie-trace-table} says that \(X_0\) is a
two-edge wedge in one of the two color triangles.  Its two isolated
pair fibers and the two upper extra fibers lie in the associated
six-cycle.  The dual argument gives the same conclusion for \(Z_0\),
its upper pair fibers, and the two lower extra fibers.

Every member of \(X_0\) meets every member of \(Z_0\).  If the two
wedges lie in the same color triangle, all four extra fibers and all
four isolated mandatory fibers lie in the same six-cycle.  Any vertex
of the other cycle outside the shared edge lies in none of them.  If
the wedges lie in different triangles, the cross-intersection condition
forces
\[
 X_0=\{01,02\},\qquad Z_0=\{03,04\},
\]
or the reverse.  The first and fourth rows of
\eqref{eq:damp-q5-bowtie-trace-table} then show that neither endpoint
of the shared edge \(\varnothing 0\) lies in an extra or isolated
mandatory fiber.

In either case choose a vertex \(p\) outside all extra fibers and all
isolated mandatory fibers.  Its two lower labels are distinct pairs in
\(X\setminus X_0\), and its two upper labels are distinct pairs in
\(Z\setminus Z_0\).  The cross-union condition makes every chosen lower
pair disjoint from every chosen upper pair, producing a four-cycle in
the bipartite disjointness graph.  This final contradiction proves the
theorem.
\end{proof}

\begin{corollary}[Exclusion of the six-cycle branch]
\label{cor:damp-q5-no-six-cycle-branch}
In a minimum five-coordinate gap, a shortest cycle in the intermediate
layer cannot have length six.
\end{corollary}

\begin{proof}
Proposition~\ref{prop:damp-q5-six-cycle-bowtie-normal-form} gives the
bowtie and its five-pair/five-triple support split.  Its lower and upper
defect carriers satisfy the hypotheses of
Theorem~\ref{thm:damp-q5-bowtie-carrier-exclusion} by
Lemma~\ref{lem:damp-q5-paired-defect-cylinders} and the
cross-cylinder law \eqref{eq:damp-q5-cross-cylinder-u-neighbor}.
\end{proof}

\begin{corollary}[Four-cycle reduction]
\label{cor:damp-q5-four-cycle-reduction}
Every minimum five-coordinate gap has a four-cycle in its intermediate
layer.
\end{corollary}

\begin{proof}
Corollary~\ref{cor:damp-q5-minimum-gap-cycle-core} supplies a shortest
cycle of length \(4,6,8\), or \(10\).  The other three lengths are
excluded by Corollaries~\ref{cor:damp-q5-no-six-cycle-branch},
\ref{thm:damp-q5-no-eight-cycle-branch}, and
\ref{cor:damp-q5-no-ten-cycle-branch}.
\end{proof}

\begin{lemma}[The intermediate layer contains no three-cube]
\label{lem:damp-q5-no-three-cube-layer}
In a minimum five-coordinate gap, the graph induced by \(U\) contains
no three-dimensional subcube.
\end{lemma}

\begin{proof}
Suppose that a three-face \(F\) of \(Q_E\) is contained in \(U\).
Since \(8\le |U|\le10\), there are at most two vertices outside \(F\).
If \(U=F\), then \(U\) uses only three coordinates.  Applying
Lemma~\ref{lem:damp-q5-unused-coordinate-elimination} to the other two
coordinates leaves at most
\(\binom32+\binom33=4\) possible supports in \(\mathcal K\), contrary
to \(|\mathcal K|=|U|=8\).

If there is exactly one vertex outside \(F\), it has at most one
neighbour in \(F\), contradicting the minimum-degree conclusion of
Lemma~\ref{lem:damp-q5-paired-defect-cylinders}.  Suppose finally that
there are two outside vertices.  Each has at most one neighbour in
\(F\), so minimum degree forces the two outside vertices to be adjacent
and each to have a neighbour in \(F\).  A vertex adjacent to \(F\)
differs from it in one of the two coordinates fixed on \(F\).  The two
outside vertices must differ in the same fixed coordinate; otherwise
they differ in both fixed coordinates and cannot be adjacent.  Hence
\(U\) uses only the three coordinates of \(F\) and one further
coordinate.  This contradicts
Lemma~\ref{lem:damp-q5-ten-layer-five-colors}.
\end{proof}

\begin{lemma}[A degree-two vertex in the square branch]
\label{lem:damp-q5-square-branch-degree-two-vertex}
In a minimum five-coordinate gap, the graph induced by \(U\) has a
vertex of degree exactly two.
\end{lemma}

\begin{proof}
Fix a square \(C\subseteq U\), with active colors \(a,b\), and translate
so that its three inactive coordinates are zero.  Suppose for a
contradiction that every vertex of \(U\) has degree at least three.
For each inactive color \(i\), let \(S_i\subseteq C\) be the set of
square vertices whose \(i\)-neighbour belongs to \(U\).  Put
\[
 m=\sum_i|S_i|,
 \qquad
 r=|\{i:S_i\ne\varnothing\}|.
\]
Every vertex of \(C\) already has its two square neighbours and
therefore needs an inactive neighbour.  Hence
\[
 \bigcup_iS_i=C,
 \qquad m\ge4.
\]
No \(S_i\) equals \(C\), since that would give a three-cube in \(U\),
contrary to Lemma~\ref{lem:damp-q5-no-three-cube-layer}.  Thus every
nonempty \(S_i\) has size at most three, and in particular \(r\ge2\).

Let \(W_1\) be the set of the \(m\) recorded inactive neighbours, and
put \(W_2=U\setminus(C\cup W_1)\), with \(t=|W_2|\).  Since
\(|U|\le10\),
\[
 t\le6-m.
\]
Write \(e_{11}\) for the number of edges within \(W_1\) and \(e_{12}\)
for the number between \(W_1\) and \(W_2\).  Edges within \(W_1\)
join lifts in the same inactive direction.  A proper nonempty induced
subgraph of the square on \(S_i\) has at most \(|S_i|-1\) edges, so
\[
 e_{11}\le m-r.
\]
A vertex of \(W_2\) has at most two neighbours in \(W_1\): to have any,
its inactive support must have size two, and deleting either inactive
color gives the only two possibilities.  Hence \(e_{12}\le2t\).

Every vertex of \(W_1\) has exactly one neighbour in \(C\) and, by the
degree assumption, at least two neighbours in \(W_1\cup W_2\).
Counting its incident edges gives
\[
 2m\le2e_{11}+e_{12}
 \le2(m-r)+2t,
\]
and therefore \(r\le t\).  Combining the inequalities yields
\[
 2\le r\le t\le6-m\le2.
\]
Thus \(r=t=2\) and \(m=4\), with equality throughout the edge count.
Let \(i,j\) be the two inactive colors used in \(W_1\).  Each of the two
vertices of \(W_2\) then has two neighbours in \(W_1\), so its inactive
support is \(\{i,j\}\), and its square trace belongs to
\(S_i\cap S_j\).  The two vertices have distinct square traces; hence
\(|S_i\cap S_j|\ge2\).  Since
\(|S_i|+|S_j|=4\), this forces
\[
 S_i=S_j
\]
with both sets of size two.  Their union cannot be all four vertices of
\(C\), a final contradiction.  The minimum-degree lemma already gives
degree at least two, so some vertex has degree exactly two.
\end{proof}

\begin{corollary}[Opposite carriers meet in at most a square]
\label{cor:damp-q5-cross-carrier-square-or-edge}
If a lower carrier of support \(R\) and an upper carrier of support
\(Q\) meet, then
\[
 |R\cup Q|\ge3.
\]
If equality holds, their defining cylinders meet in a square contained
in \(U\); if \(|R\cup Q|=4\), they meet in an edge contained in \(U\).
\end{corollary}

\begin{proof}
The intersection of the two defining cylinders is the subcube through
any common vertex on the coordinates in \(E\setminus(R\cup Q)\).
As in the proof of
\eqref{eq:damp-q5-cross-cylinder-u-neighbor}, the lower cylinder lies in
\(T\), the upper cylinder lies in \(B\), and their intersection is
therefore contained in \(B\cap T=U\).  Lemma
\ref{lem:damp-q5-no-three-cube-layer} bounds its dimension by two.
\end{proof}

\begin{lemma}[Degree-two vertices require a triple carrier]
\label{lem:damp-q5-degree-two-needs-triple}
Let \(p\in U\), and suppose that its two neighbours in \(U\) have edge
colors \(a,b\).  Then \(p\) lies in an eligible triple carrier on at
least one of the lower and upper sides.
\end{lemma}

\begin{proof}
Suppose instead that every lower and upper defect label at \(p\) is a
pair support.  Choose distinct lower labels \(x_1,x_2\) and distinct
upper labels \(z_1,z_2\), which is possible by
\eqref{eq:damp-q5-two-sided-defect-degree}.  For every \(i,j\), the
intersection of their cylinders contains the subcube through \(p\) on
\(E\setminus(x_i\cup z_j)\).  Every free coordinate is therefore the
color of a neighbour of \(p\).  On putting
\(I=E\setminus\{a,b\}\), we obtain
\[
 I\subseteq x_i\cup z_j
 \qquad(1\le i,j\le2).
\]

Fix \(x_1,x_2\).  Each \(z_j\) must contain
\[
 (I\setminus x_1)\cup(I\setminus x_2)
 =I\setminus(x_1\cap x_2).
\]
Two distinct two-sets \(z_1,z_2\) can both contain this set only if it
has size at most one.  Thus \(x_1\cap x_2\) contains at least two
elements of the three-set \(I\).  Since \(x_1,x_2\) are themselves
two-sets, this forces \(x_1=x_2\), a contradiction.
\end{proof}

\begin{lemma}[Signed split normal form at a degree-two vertex]
\label{lem:damp-q5-degree-two-signed-split}
Let \(p\in U\) have two neighbours in \(U\).  After possibly exchanging
the complementary gaps \(S\subseteq B\) and \(A\subseteq T\), the five
colors can be labelled \(a,b,c,d,e\) so that
\begin{equation}
N_U(p)=\{a,b\},
 \qquad
 N_S(p)=\{c\},
 \qquad
 N_A(p)=\{d,e\}.
\label{eq:damp-q5-degree-two-neighbour-split}
\end{equation}
Every support in \(\mathcal H_S(p)\) contains \(c\), while
\begin{equation}
\mathcal H_A(p)
 \subseteq
 \bigl\{\{a,d,e\},\{b,d,e\},\{c,d,e\}\bigr\},
 \qquad
 |\mathcal H_A(p)|\ge2.
\label{eq:damp-q5-degree-two-upper-triples}
\end{equation}

Put \(X=\mathcal K_2\), and let \(X_0\) be the set of members of \(X\)
isolated from
\(\{E\setminus R:R\in\mathcal K_3\}\) in the bipartite disjointness
graph.  Then at least two of
\begin{equation}
\{a,b\},\qquad \{a,c\},\qquad \{b,c\}
\label{eq:damp-q5-degree-two-isolated-triangle}
\end{equation}
belong to \(X_0\).
\end{lemma}

\begin{proof}
For \(R\in\mathcal H_S(p)\), every coordinate in \(N_S(p)\) belongs to
\(R\).  Indeed, if \(f\in N_S(p)\setminus R\), then \(p\) and
\(p^{\{f\}}\) have the same trace on \(R\), whereas the former realizes
the unique pattern absent from \(S\) and the latter belongs to \(S\).
This is impossible.  Dually, every member of \(\mathcal H_A(p)\)
contains \(N_A(p)\).

The three colors outside \(N_U(p)\) partition into \(N_S(p)\) and
\(N_A(p)\).  Neither part has size three.  For if, say,
\(|N_S(p)|=3\), then a rank-two or rank-three support containing it must
be that unique three-set, contradicting
\(|\mathcal H_S(p)|\ge2\).  The dual argument applies to \(N_A(p)\).
Their sizes are consequently one and two, giving
\eqref{eq:damp-q5-degree-two-neighbour-split} after exchanging the two
sides if necessary.

Every upper label contains \(\{d,e\}\).  If the pair \(\{d,e\}\) were
an upper label, it would be an upper eligible pair.  No upper triple
containing it would then be eligible, by
\eqref{eq:damp-q5-eligible-defect-supports}; since there is no second
pair containing both \(d,e\), this would contradict the two-label lower
bound.  Thus all upper labels are triples, and the only possibilities
are those in \eqref{eq:damp-q5-degree-two-upper-triples}.

An upper eligible triple \(Q\) has the form \(E\setminus x\) for an
isolated pair \(x\in X_0\).  Complementing the three triples in
\eqref{eq:damp-q5-degree-two-upper-triples} gives, respectively,
\(\{b,c\},\{a,c\},\{a,b\}\).  At least two upper labels occur, proving
\eqref{eq:damp-q5-degree-two-isolated-triangle}.
\end{proof}

\begin{lemma}[Marked complementary vertices]
\label{lem:damp-q5-marked-complementary-carriers}
Retain the notation of
Lemma~\ref{lem:damp-q5-degree-two-signed-split}.  If
\(x\in X_0\), then the mandatory lower pair carrier and its
complementary upper triple carrier meet in exactly one vertex:
\begin{equation}
\Gamma_x^-\cap\Gamma_{E\setminus x}^+
 =\{v_x\}\subseteq U.
 \label{eq:damp-q5-marked-complementary-vertex}
\end{equation}
Thus every isolated pair support has a canonically marked relative
vertex, determined by the two missing trace patterns.
\end{lemma}

\begin{proof}
Every pair in \(\mathcal K_2\) is eligible.  The assumption
\(x\in X_0\) says exactly that the complementary triple
\(E\setminus x\) is eligible for \(\mathcal K^*\).  Hence the two
carriers in \eqref{eq:damp-q5-marked-complementary-vertex} are defined.
Their cylinders fix complementary sets of coordinates, so their two
trace patterns have a unique common extension \(v_x\in Q_E\).  The
lower cylinder is contained in \(T\), while the upper cylinder is
contained in \(B\).  Consequently their common vertex belongs to
\(B\cap T=U\), and the carrier definitions give the displayed equality.
\end{proof}

\begin{lemma}[Blocker face forced by the local upper triples]
\label{lem:damp-q5-degree-two-blocker-face}
Use the signed split normal form
\eqref{eq:damp-q5-degree-two-neighbour-split}, and put
\begin{equation}
I=\{i\in\{a,b,c\}:\{i,d,e\}\in\mathcal H_A(p)\},
 \qquad J=\{a,b,c\}\setminus I.
 \label{eq:damp-q5-degree-two-upper-index-set}
\end{equation}
Then \(2\le |I|\le3\), and the restriction of \(A\) to the
three-face through \(p\) on \(\{a,b,c\}\) is a nonempty ample family
satisfying
\begin{equation}
\varnothing\ne
 A\cap\{p^F:F\subseteq\{a,b,c\}\}
 \subseteq
 \{p^I,p^{I\cup J}\}.
 \label{eq:damp-q5-degree-two-blocker-edge}
\end{equation}
When \(|I|=2\), the right-hand side is an edge; when \(|I|=3\), it
is the singleton \(\{p^{\{a,b,c\}}\}\).

More precisely, set \(z_F=\prod_{f\in F}v_{f,p_f}\).  There is an
\(M\in\{I,I\cup J\}\) such that
\begin{equation}
(J_A:m_p)=(v_{d,p_d},v_{e,p_e},z_M).
 \label{eq:damp-q5-degree-two-blocker-colon}
\end{equation}
If \(|I|=3\), then necessarily \(M=\{a,b,c\}\), the vertex
\(p^{\{a,b,c\}}\) is the unique blocker for adjoining \(p\) to
\(A\), and \(A\) contains the full \(d,e\)-square based at that
vertex.  In this case the four vertices
\begin{equation}
p,\quad p^{\{a\}},\quad p^{\{b\}},\quad p^{\{a,b\}}
 \label{eq:damp-q5-degree-two-forced-u-square}
\end{equation}
all belong to \(U\).
\end{lemma}

\begin{proof}
The bounds on \(|I|\) are
\eqref{eq:damp-q5-degree-two-upper-triples}.  For every \(i\in I\),
the upper defect cylinder with support \(\{i,d,e\}\) contains \(p\),
is disjoint from \(A\), and is free on the other two coordinates.
Thus a vertex of \(A\) that agrees with \(p\) in coordinates \(d,e\)
must differ from \(p\) in every coordinate of \(I\).  This proves the
containment in \eqref{eq:damp-q5-degree-two-blocker-edge}.

The family on the left of
\eqref{eq:damp-q5-degree-two-blocker-edge} is a restriction of the ample
family \(A\), and is therefore ample.  It is nonempty because a blocking
vertex for adjoining \(p\) to \(A\) differs from \(p\) in none of the
coordinates in \(N_A(p)=\{d,e\}\).  This also proves the point-versus-edge
description.

The monomial-colon formula
\eqref{eq:damp-relative-augmentation-colon-formula} now gives
\eqref{eq:damp-q5-degree-two-blocker-colon}.  Indeed, the two neighbours
\(p^{\{d\}},p^{\{e\}}\in A\) supply the two variables.  Every other
vertex outside the displayed three-face contributes a monomial divisible
by one of them.  Inside the face there are only the one or two vertices
in \eqref{eq:damp-q5-degree-two-blocker-edge}; their unique
inclusion-minimal difference support is some
\(M\in\{I,I\cup J\}\).

Suppose finally that \(|I|=3\).  The blocker face is then the singleton
\(p^{\{a,b,c\}}\).  Since \(p^{\{d\}},p^{\{e\}}\in A\) and \(A\) is
isometric, the other intermediate vertex on a geodesic between them is
\(p^{\{d,e\}}\in A\).  These three vertices and
\(p^{\{a,b,c\}}\) realize all four patterns on \(\{d,e\}\), so \(A\)
shatters that pair.  Ampleness makes shattering strong.  The zero-pattern
vertex of a full \(d,e\)-square must lie in the blocker face, and hence
is \(p^{\{a,b,c\}}\).  This proves the asserted square in \(A\).

The upper cylinder with support \(\{c,d,e\}\) contains the full
\(a,b\)-square through \(p\) in \(B\).  Its first three displayed
vertices belong to \(U\) by
\eqref{eq:damp-q5-degree-two-neighbour-split}.  If
\(p^{\{a,b\}}\) belonged to \(S\), then the restriction of \(S\) to
the face where \(d,e\) agree with \(p\) would contain both
\(p^{\{c\}}\) and \(p^{\{a,b\}}\).  Within that face, every neighbour
of the latter vertex is one of
\(p^{\{a\}},p^{\{b\}},p^{\{a,b,c\}}\), none of which belongs to
\(S\).  This would disconnect a nonempty restriction of the ample, hence
isometric, family \(S\).  Therefore \(p^{\{a,b\}}\notin S\); it lies
in \(B\), so it belongs to \(U\), proving
\eqref{eq:damp-q5-degree-two-forced-u-square}.
\end{proof}

\begin{lemma}[Two upper labels force a peripheral contraction]
\label{lem:damp-q5-two-label-peripheral-contraction}
In the setting of
Lemma~\ref{lem:damp-q5-degree-two-blocker-face}, suppose that
\(|I|=2\).  Write \(I=\{r,s\}\) and \(J=\{t\}\).  Then
\begin{equation}
p^{\{r,s\}}\in A
 \qquad\text{and}\qquad
 (J_A:m_p)=(v_{d,p_d},v_{e,p_e},z_{\{r,s\}}).
 \label{eq:damp-q5-two-label-quadratic-colon}
\end{equation}

Translate \(p\) to the empty set, identify vertices with their difference
supports from \(p\), and contract coordinate \(t\) in \(A\).  The resulting
ample family \(C=\pi_tA\subseteq Q_{\{r,s,d,e\}}\) satisfies
\begin{equation}
\begin{split}
 \{rs,d,e,de,rsd,rse,rsde\}
 &\subseteq C\\
 &\subseteq
 \{rs,d,e,de,rsd,rse,rsde,rd,sd,re,se,rde,sde\}.
 \end{split}
 \label{eq:damp-q5-two-label-contraction-profile}
\end{equation}
Moreover, the \(t=p_t\) section of \(A\), after suppressing \(t\),
contains the seven vertices on the left-hand side.
At least one of \(d,e\) is a peripheral coordinate of \(C\).
\end{lemma}

\begin{proof}
By Lemma~\ref{lem:damp-q5-degree-two-blocker-face}, the part of \(A\)
agreeing with \(p\) in coordinates \(d,e\) is a nonempty subset of
\(\{p^{\{r,s\}},p^{\{r,s,t\}}\}\).  Suppose that it is the singleton
\(\{p^{\{r,s,t\}}\}\).  The vertices \(p^{\{d\}}\) and
\(p^{\{e\}}\) belong to \(A\), and isometry of \(A\), applied to this
pair, also gives \(p^{\{d,e\}}\in A\).  These three vertices together
with \(p^{\{r,s,t\}}\) show that \(A\) shatters \(\{d,e\}\).
Ampleness makes this shattering strong.  The base of the resulting full
\(d,e\)-square must agree with \(p\) in \(d,e\), so the singleton
description forces that base to be \(p^{\{r,s,t\}}\).

That square realizes the four traces on \(\{t,d,e\}\) with \(t\ne p_t\).
The vertices \(p^{\{d\}},p^{\{e\}},p^{\{d,e\}}\) realize the three
nonconstant traces with \(t=p_t\), while no vertex of \(A\) realizes the
trace of \(p\): any such vertex would lie in the displayed singleton and
would instead differ in \(t\).  Thus adjoining \(p\) to \(A\) newly
shatters \(\{t,d,e\}\).  This would be a third member of
\(\mathcal H_A(p)\), contrary to \(I=\{r,s\}\).  Hence
\(p^{\{r,s\}}\in A\), and the colon formula in
Lemma~\ref{lem:damp-q5-degree-two-blocker-face} has its
inclusion-minimal third generator at \(\{r,s\}\).  This proves
\eqref{eq:damp-q5-two-label-quadratic-colon}.

We now use the translated set notation in the statement.  The two upper
defect cylinders have supports \(rde\) and \(sde\).  They exclude both
lifts of each of \(\varnothing,r,s\) from \(A\), so
\(\varnothing,r,s\notin C\).  The vertices \(d,e,de,rs\) belong to
\(C\).  Let \(X=\rho_t^{p_t}A\).  This ample section contains
\(d,e,de,rs\), so it shatters \(\{d,e\}\).  Strong shattering in \(X\)
forces a full \(d,e\)-square.  Its zero-pattern vertex is \(rs\), because
the blocker-face description leaves no other vertex of \(X\) agreeing
with \(p\) in \(d,e\).  Hence \(rsd,rse,rsde\in X\), proving the
section assertion and the corresponding containment in \(C\).
These seven forced vertices and the three excluded vertices leave exactly
the six optional vertices displayed in
\eqref{eq:damp-q5-two-label-contraction-profile}.

It remains to prove the peripheral assertion.  Put
\[
 \alpha=[rd\in C],\quad \beta=[sd\in C],\quad
 \gamma=[re\in C],\quad \delta=[se\in C],\quad
 \epsilon=[rde\in C],\quad \phi=[sde\in C].
\]
Here brackets denote truth values.  Since the upper defect on \(rde\)
makes \(C\) shatter \(\{r,d\}\), ampleness of \(C\) supplies a full
\(r,d\)-square.  The missing vertices \(\varnothing,s\) rule out the
two squares with \(e=0\); the two remaining bases give
\begin{equation}
\begin{gathered}
  \alpha\vee\beta,\qquad \gamma\vee\delta,\\
  (\gamma\wedge\epsilon)\vee(\delta\wedge\phi),\qquad
  (\alpha\wedge\epsilon)\vee(\beta\wedge\phi).
 \end{gathered}
 \label{eq:damp-q5-two-label-ladder-laws}
\end{equation}
Here the first line follows from isometry between \(d\) or \(e\) and
\(rs\).  The final disjunction is obtained in the same way as the
preceding one, using strong shattering of \(\{r,e\}\).
Isometry of \(C\), applied respectively to the pairs \(rd,re\) and
\(sd,se\), gives
\begin{equation}
\alpha\wedge\gamma\Longrightarrow\epsilon,
 \qquad
 \beta\wedge\delta\Longrightarrow\phi.
 \label{eq:damp-q5-two-label-isometry-implications}
\end{equation}
Reading the two sections from
\eqref{eq:damp-q5-two-label-contraction-profile}, coordinate \(d\) is
peripheral precisely when
\begin{equation}
\gamma\Longrightarrow\epsilon
 \quad\text{and}\quad
 \delta\Longrightarrow\phi,
 \label{eq:damp-q5-two-label-d-peripheral}
\end{equation}
and coordinate \(e\) is peripheral precisely when
\begin{equation}
\alpha\Longrightarrow\epsilon
 \quad\text{and}\quad
 \beta\Longrightarrow\phi.
 \label{eq:damp-q5-two-label-e-peripheral}
\end{equation}
The same section check gives
\begin{equation}
\begin{split}
 r\text{ is peripheral}
 &\Longleftrightarrow
 \alpha\wedge\gamma\wedge\epsilon,\\
 s\text{ is peripheral}
 &\Longleftrightarrow
 \beta\wedge\delta\wedge\phi.
 \end{split}
 \label{eq:damp-q5-two-label-rs-peripheral}
\end{equation}

Suppose that \(d\) is not peripheral.  If
\(\gamma\wedge\neg\epsilon\) holds, then
 \eqref{eq:damp-q5-two-label-ladder-laws} forces
\(\delta\wedge\phi\).  Failure of peripheralness for \(e\) would give
either \(\alpha\wedge\neg\epsilon\), contradicting the first implication
in \eqref{eq:damp-q5-two-label-isometry-implications}, or
\(\beta\wedge\neg\phi\), contradicting \(\phi\).  If instead
\(\delta\wedge\neg\phi\) holds, then
 \eqref{eq:damp-q5-two-label-ladder-laws} forces
\(\gamma\wedge\epsilon\).  Nonperipheralness of \(e\) would now give
either \(\alpha\wedge\neg\epsilon\), immediately impossible, or
\(\beta\wedge\neg\phi\), contrary to the second implication in
\eqref{eq:damp-q5-two-label-isometry-implications}.  Therefore \(d\) or
\(e\) is peripheral.
\end{proof}

\begin{lemma}[Carrier compatibility in the two-label core]
\label{lem:damp-q5-two-label-carrier-compatibility}
Let \(C\subseteq Q_{\{r,s,d,e\}}\) be an ample family satisfying
\eqref{eq:damp-q5-two-label-contraction-profile}, and put
\[
 K=\{rs,d,e,de,rsd,rse,rsde\}.
\]
Let \(D\subseteq C\) be a nonempty isometric family.  Orient a peripheral
coordinate \(q\) of \(C\) so that its \(q=0\) section is contained in its
\(q=1\) section.  Call this orientation \emph{\(D\)-compatible} if
\begin{equation}
x\in D,\quad q\notin x
 \quad\Longrightarrow\quad
 x\cup\{q\}\in D.
 \label{eq:damp-q5-two-label-carrier-compatible}
\end{equation}
If some connected component of \(C\setminus D\) is disjoint from \(K\),
then \(C\) has a \(D\)-compatible peripheral coordinate.
\end{lemma}

\begin{proof}
Use the six optional-vertex indicators
\(\alpha,\beta,\gamma,\delta,\epsilon,\phi\) from the proof of
Lemma~\ref{lem:damp-q5-two-label-peripheral-contraction}.  The optional
vertices induce two disjoint three-vertex paths,
\begin{equation}
rd-rde-re
 \qquad\text{and}\qquad
 sd-sde-se.
 \label{eq:damp-q5-two-label-optional-paths}
\end{equation}
Let \(P\) be a component of \(C\setminus D\) disjoint from \(K\).
After interchanging \(r,s\), assume that
\(P\subseteq\{sd,sde,se\}\).  For
\(x\in\{d,e,de\}\), if \(sx\in P\), then both \(x\) and \(rsx\)
belong to \(D\): they are core vertices adjacent to \(sx\), and otherwise
would belong to the same component of \(C\setminus D\).  The two vertices
\(x,rsx\) have precisely \(rx,sx\) as their common cube neighbours.
Isometry of \(D\) and \(sx\notin D\) therefore force
\begin{equation}
sx\in P
 \quad\Longrightarrow\quad
 x,rx,rsx\in D.
 \label{eq:damp-q5-two-label-opposite-path}
\end{equation}

Up to interchanging \(d,e\), the connected subpath \(P\) has four
possible shapes.  We treat them in decreasing order.

\smallskip
\noindent\emph{The whole path \(P=\{sd,sde,se\}\).}
Formula \eqref{eq:damp-q5-two-label-opposite-path} puts
\[
 d,rd,rsd,\quad de,rde,rsde,\quad e,re,rse
\]
in \(D\).  Thus \(r\) is peripheral by
\eqref{eq:damp-q5-two-label-rs-peripheral}.  Every \(r=0\) vertex of
\(D\) has its \(r\)-mate in the displayed list; the three optional
\(s\)-vertices lie in \(P\), not in \(D\).  Hence \(r\) is
\(D\)-compatible.

\smallskip
\noindent\emph{The edge \(P=\{sd,sde\}\).}
The opposite-path rule gives
\begin{equation}
d,rd,rsd,de,rde,rsde\in D.
 \label{eq:damp-q5-two-label-edge-forced-carrier}
\end{equation}
Consequently \(e\) is peripheral by
\eqref{eq:damp-q5-two-label-e-peripheral}.  All its compatibility
conditions are already supplied by
\eqref{eq:damp-q5-two-label-edge-forced-carrier}, except possibly
\begin{equation}
rs\in D\quad\Longrightarrow\quad rse\in D.
 \label{eq:damp-q5-two-label-edge-only-failure}
\end{equation}
Suppose this implication fails.  If \(se\in C\), then
\(se\in D\), because it is adjacent to the component \(P\).
But \(rs,se\in D\) have Hamming distance two, and their two common
neighbours are \(s\notin C\) and \(rse\notin D\), contradicting
isometry.  Hence \(se\notin C\).  The ladder laws
 \eqref{eq:damp-q5-two-label-ladder-laws} now force
\(re,rde\in C\), so \(d\) is peripheral.  Its \(D\)-compatibility follows
from \eqref{eq:damp-q5-two-label-edge-forced-carrier}.  Explicitly, the
possible lower-side pairs are
\[
 rs\to rsd,\quad e\to de,\quad rse\to rsde,\quad
 re\to rde,\quad se\to sde;
\]
the first four required mates are in \(D\) whenever their sources are,
and \(se\notin C\).

\smallskip
\noindent\emph{The middle vertex \(P=\{sde\}\).}
We have \(de,rsde\in D\), and their isometry forces \(rde\in D\).
If \(sd\) or \(se\) is present in \(C\), then it belongs to \(D\), since
it is adjacent to \(P\).  Because \(rde,sde\in C\), both \(d\) and \(e\)
are peripheral by
\eqref{eq:damp-q5-two-label-d-peripheral} and
\eqref{eq:damp-q5-two-label-e-peripheral}.

The only possible compatibility failures for \(d\) are
\begin{equation}
rs\in D,\ rsd\notin D,
 \qquad\text{or}\qquad se\in D;
 \label{eq:damp-q5-two-label-middle-d-failures}
\end{equation}
for \(e\), they are
\begin{equation}
rs\in D,\ rse\notin D,
 \qquad\text{or}\qquad sd\in D.
 \label{eq:damp-q5-two-label-middle-e-failures}
\end{equation}
All other required mates are \(de,rde,rsde\).  The vertices \(sd,se\)
cannot both lie in \(D\), because their common neighbours are
\(s\notin C\) and \(sde\notin D\).
If \(sd\in D\), isometry between \(sd\) and \(rsde\) forces
\(rsd\in D\), so \(d\) is compatible.  Symmetrically, if \(se\in D\),
then \(rse\in D\), so \(e\) is compatible.  If neither lies in \(D\),
the only possible failures in
\eqref{eq:damp-q5-two-label-middle-d-failures}--\eqref{eq:damp-q5-two-label-middle-e-failures}
require \(rs\in D\).  Isometry between \(rs\) and \(rsde\) then puts
\(rsd\) or \(rse\) in \(D\), making \(d\) or \(e\), respectively,
compatible.

\smallskip
\noindent\emph{The endpoint \(P=\{sd\}\).}
The opposite-path rule gives \(d,rd,rsd\in D\).  If \(sde\in C\), then
\(sde\in D\), because it is adjacent to \(P\).  The ladder laws give
\begin{equation}
re\vee se,\qquad
 re\Longrightarrow rde,\qquad
 se\Longrightarrow sde.
 \label{eq:damp-q5-two-label-endpoint-laws}
\end{equation}
Thus \(d\) is peripheral.  Its compatibility cannot fail at
\(e\to de\): the vertices \(d,e\in D\) would force \(de\in D\).
It cannot fail at \(re\to rde\): the vertices \(rd,re\in D\) would force
\(rde\in D\).  Nor can it fail at \(se\to sde\), by
\eqref{eq:damp-q5-two-label-endpoint-laws} and the component boundary.
The only possible failure is therefore
\begin{equation}
rse\in D,\qquad rsde\notin D.
 \label{eq:damp-q5-two-label-endpoint-only-failure}
\end{equation}

Assume \eqref{eq:damp-q5-two-label-endpoint-only-failure}.
Isometry between \(rsd,rse\in D\) forces \(rs\in D\).  If \(sde\in C\),
then \(sde\in D\), but \(rsd,sde\) would have neither of their common
neighbours \(sd,rsde\) in \(D\).  Hence \(sde\notin C\), and
\eqref{eq:damp-q5-two-label-endpoint-laws} gives
\[
 se\notin C,\qquad re,rde\in C.
\]
Coordinate \(r\) is consequently peripheral.  It is also compatible.
Indeed, \(d\to rd\) is already contained in \(D\).  If \(e\in D\),
isometry between \(e,rse\) forces \(re\in D\), since \(se\notin C\).
Finally suppose \(de\in D\).  A \(D\)-geodesic from \(de\) to \(rse\)
must start at \(rse\) through \(re\), because \(se\notin C\) and
\(rsde\notin D\).  Its next vertex is \(e\) or \(rde\).  In the former
case, isometry between \(rd,re\in D\) forces \(rde\in D\); in the latter
case this is immediate.  Thus every \(r=0\) carrier vertex has its
\(r\)-mate in \(D\).  This completes all four path shapes.
\end{proof}

\begin{corollary}[The two-label branch is impossible in a minimum gap]
\label{cor:damp-q5-no-two-label-branch}
In the minimum-gap setting of
Proposition~\ref{prop:damp-q5-central-relative-cm-profile} and
Lemma~\ref{lem:damp-q5-degree-two-blocker-face}, one has \(|I|=3\).
Equivalently, after the signed normalization,
\begin{equation}
\mathcal H_A(p)=\{ade,bde,cde\}.
 \label{eq:damp-q5-three-upper-labels-forced}
\end{equation}
\end{corollary}

\begin{proof}
Suppose that \(|I|=2\), use \(I=\{r,s\}\), \(J=\{t\}\), and translate
\(p\) to the empty set as in
Lemma~\ref{lem:damp-q5-two-label-peripheral-contraction}.  Let
\begin{equation}
X=\rho_t^{p_t}A,\qquad
 Y=\rho_t^{1-p_t}A,\qquad
 C=X\cup Y=\pi_tA,\qquad D=X\cap Y.
 \label{eq:damp-q5-two-label-expansion-sections}
\end{equation}
 The section assertion of that lemma gives \(K\subseteq X\).
The two-section shattering formula for an ample expansion, used in
Lemma~\ref{lem:damp-expansion-separator-duplication}, shows that \(D\)
is ample, hence isometric.  It is nonempty whenever both signed sections
are nonempty, because otherwise the two sections would be disconnected in
the one-inclusion graph of \(A\).

All four endpoints of a minimum gap are periphery-free by
Proposition~\ref{prop:damp-q5-central-relative-cm-profile}.  In
particular, \(t\) is not peripheral in \(A\), so \(X\setminus Y\) and
\(Y\setminus X\) are both nonempty.  Moreover, no edge of \(C\) joins
these two sets: otherwise its two signed lifts in \(A\) would be at
Hamming distance two with neither possible intermediate vertex present,
contradicting isometry of \(A\).  Hence each connected component of
\(C\setminus D\) lies wholly in one of the two differences.  A component
meeting \(Y\setminus X\) is disjoint from \(K\subseteq X\).

Lemma~\ref{lem:damp-q5-two-label-carrier-compatibility} now supplies a
\(D\)-compatible peripheral coordinate \(q\) of \(C\).  This coordinate
is peripheral in \(A\).  Indeed, if a vertex on its smaller side belongs
to \(D\), compatibility puts its \(q\)-mate in \(D\).  If it lies in
\(C\setminus D\), its mate either lies in \(D\), or the edge between them
stays in one component of \(C\setminus D\) and therefore in the same
signed section.  In every case each lift of the smaller-side vertex has
its \(q\)-mate in the same \(t\)-section.  This contradicts
periphery-freeness of \(A\).  Therefore \(|I|=3\).
\end{proof}

\begin{corollary}[Marked-square profile with three upper labels]
\label{cor:damp-q5-three-label-marked-square}
In the setting of Lemma~\ref{lem:damp-q5-degree-two-blocker-face},
suppose that \(|I|=3\).  Then
\begin{equation}
\begin{aligned}
  \Gamma_{\{c,d,e\}}^+
   &=\{p,p^{\{a\}},p^{\{b\}},p^{\{a,b\}}\},\\
  \Gamma_{\{a,d,e\}}^+\cap\Gamma_{\{c,d,e\}}^+
   &=\{p,p^{\{b\}}\},\\
  \Gamma_{\{b,d,e\}}^+\cap\Gamma_{\{c,d,e\}}^+
   &=\{p,p^{\{a\}}\}.
 \end{aligned}
 \label{eq:damp-q5-three-label-upper-square-traces}
\end{equation}
The marked vertices of the three complementary isolated pairs satisfy
\begin{equation}
v_{\{a,b\}}\in
  \{p^{\{a\}},p^{\{b\}},p^{\{a,b\}}\},
 \quad
 v_{\{a,c\}}\in\{p,p^{\{a\}},p^{\{a,c\}}\},
 \quad
 v_{\{b,c\}}\in\{p,p^{\{b\}},p^{\{b,c\}}\}.
 \label{eq:damp-q5-three-label-marked-locations}
\end{equation}
In particular, the mandatory \(a,b\)-pair carrier meets the forced
intermediate square in the single vertex \(v_{\{a,b\}}\ne p\).
\end{corollary}

\begin{proof}
The three upper cylinders through \(p\) are free, respectively, on
\(\{b,c\}\), \(\{a,c\}\), and \(\{a,b\}\).  The last cylinder is
the square in \eqref{eq:damp-q5-degree-two-forced-u-square}, so its
upper carrier is that entire square.  Intersecting it with either of the
other two cylinders fixes, respectively, the \(a\)- or the \(b\)-coordinate
at its value in \(p\), giving the two displayed edges.  All their vertices
belong to \(U\), proving
\eqref{eq:damp-q5-three-label-upper-square-traces}.

For an isolated pair \(x\), its marked vertex lies in the upper cylinder
on \(E\setminus x\), by
Lemma~\ref{lem:damp-q5-marked-complementary-carriers}.  Reading the three
corresponding squares and deleting the vertex \(p^{\{c\}}\in S\) gives
the three containments in
\eqref{eq:damp-q5-three-label-marked-locations}.  Finally,
\(\{a,b\}\notin\mathcal H_S(p)\), because every lower defect support at
\(p\) contains \(c\).  The carrier-label equivalence
\eqref{eq:damp-q5-carrier-label-equivalence} therefore excludes
\(v_{\{a,b\}}=p\), while
\eqref{eq:damp-q5-marked-complementary-vertex} makes its intersection
with the square unique.
\end{proof}

\begin{lemma}[Antipodal four-section rigidity]
\label{lem:damp-q5-antipodal-four-section-rigidity}
Let \(F=\{a,b,c\}\), write \(abc\) for the all-one vertex of
\(Q_F\), and let \(X,Y,Z\subseteq Q_F\).  In the coordinate order
\((d,e)\), suppose that
\begin{equation}
\mathcal A
 =\bigl(\{abc\}\times\{00\}\bigr)
  \mathbin{\dot\cup}(X\times\{10\})
  \mathbin{\dot\cup}(Y\times\{01\})
  \mathbin{\dot\cup}(Z\times\{11\})
 \subseteq Q_{F\cup\{d,e\}}
 \label{eq:damp-q5-antipodal-four-sections}
\end{equation}
is ample and
\begin{equation}
\varnothing,abc\in X\cap Y\cap Z.
 \label{eq:damp-q5-antipodal-section-endpoints}
\end{equation}
If \(\mathcal A\) has no peripheral coordinate, then, after permuting
\(a,b,c\) and possibly interchanging \(d,e\),
\begin{equation}
\begin{aligned}
  X&=Q_F\setminus\{b,c,bc\},\\
  Y&=Q_F\setminus\{a,c,ac\},\\
  Z&=Q_F\setminus\{c,ac,bc\}.
 \end{aligned}
 \label{eq:damp-q5-pentagonal-four-sections}
\end{equation}
Consequently, \(\mathcal A\) is a maximum class of VC dimension two.
In particular, it has order sixteen.
\end{lemma}

\begin{proof}
We first encode each of the three nontrivial sections by one interval in
a six-cycle.  The six vertices of \(Q_F\) other than
\(\varnothing,abc\) induce the cycle
\begin{equation}
a-ab-b-bc-c-ac-a.
 \label{eq:damp-q5-middle-six-cycle}
\end{equation}
Each of \(X,Y,Z\) is ample, because coordinate face restrictions preserve
ampleness.  We shall use the following elementary description of such a
section.
If an ample \(P\subseteq Q_F\) contains \(\varnothing\) and \(abc\),
then
\begin{equation}
H_P:=Q_F\setminus P
 \label{eq:damp-q5-antipodal-holes}
\end{equation}
is either empty or a cyclic interval of
\eqref{eq:damp-q5-middle-six-cycle} having at most four vertices.
Indeed, isometry supplies a three-edge geodesic from
\(\varnothing\) to \(abc\).  Its two middle vertices are adjacent on
the cycle, so \(H_P\) is contained in the complementary four-vertex
path.  Since \(P\) contains the two antipodes, it shatters all three
singletons.  If \(H_P\ne\varnothing\), it does not shatter \(F\).
Moreover, a pair is not shattered precisely when \(H_P\) contains the
two endpoints of an edge of the corresponding color in the middle cycle.
The three edges in any four-vertex subpath of the cycle have distinct
colors, so the number of unshattered pairs is exactly \(|E(H_P)|\).
Ampleness therefore gives
\[
 8-|H_P|=|\operatorname{Sh}(P)|
 =4+3-|E(H_P)|,
\]
and hence
\[
 |E(H_P)|=|H_P|-1.
\]
A subset of a path has this many induced
edges exactly when it is an interval.  Conversely, the same count shows
that the complement of each listed interval is ample.  Thus the interval
description is exact.

Put \(H_X=Q_F\setminus X\), and define \(H_Y,H_Z\) similarly.  For an
arbitrary subset \(H\) of the middle cycle, let
\[
 \mathcal N(H)
 =\{R\subseteq F:R\notin
        \operatorname{Sh}(Q_F\setminus H)\}.
\]
This set is read directly from the cycle: it is empty for \(H=\varnothing\);
otherwise it consists of \(F\), together with the pair \(F\setminus\{i\}\)
for every cycle edge of color \(i\) whose two endpoints lie in \(H\).

We next record the exact ampleness equation for
\eqref{eq:damp-q5-antipodal-four-sections}.  A support \(R\subseteq F\)
is shattered by \(\mathcal A\) precisely when it is shattered by
\(X\cup Y\cup Z\).  The support \(R\cup\{d\}\) is shattered precisely
when \(R\) is shattered by both \(Y\) and \(X\cup Z\), and the analogous
condition for \(R\cup\{e\}\) uses \(X\) and \(Y\cup Z\).  Finally, among
the supports containing both \(d,e\), only \(\{d,e\}\) is shattered,
because the \(00\)-section is the singleton \(\{abc\}\).  Since
\(X\cup Z=Q_F\setminus(H_X\cap H_Z)\), and similarly for the other
unions, the equality
\[
 |\operatorname{Sh}(\mathcal A)|=|\mathcal A|
\]
is equivalent to
\begin{equation}
\begin{split}
 |\mathcal N(H_X\cap H_Y\cap H_Z)|
 &+|\mathcal N(H_Y)\cup\mathcal N(H_X\cap H_Z)|\\
 &+|\mathcal N(H_X)\cup\mathcal N(H_Y\cap H_Z)|
  =|H_X|+|H_Y|+|H_Z|.
 \end{split}
 \label{eq:damp-q5-antipodal-hole-budget}
\end{equation}

Nonperipherality of \(d\) and \(e\) gives two simple interval
noncontainments.  If \(H_Z\subseteq H_Y\), then \(Y\subseteq Z\), so the
\(d=0\) section \(Y\) is contained in the \(d=1\) section \(X\cup Z\).
Thus \(d\) would be peripheral.  The analogous argument for \(e\) shows
that nonperipherality forces
\begin{equation}
H_Z\nsubseteq H_Y,
 \qquad H_Z\nsubseteq H_X.
 \label{eq:damp-q5-antipodal-hole-noncontainment}
\end{equation}

For completeness, the short interval calculation from
\eqref{eq:damp-q5-antipodal-hole-budget} and
\eqref{eq:damp-q5-antipodal-hole-noncontainment} is displayed below.
We write, for example, \(bc\) for the vertex \(\{b,c\}\) of \(Q_F\).
Up to a permutation of \(F\) and interchange of \(X,Y\), these are all
the possibilities.  The last column lists an old coordinate whose
zero-section is contained in its one-section in \(\mathcal A\).
\begin{equation}
\begin{array}{c|c|c|c}
H_X&H_Y&H_Z&\text{peripheral old coordinate}\\ \hline
\{c,bc\}&\{b,bc\}&\{b,c,bc\}&a\\
\{c,bc\}&\{c,ac\}&\{c,ac,bc\}&a,b\\
\{c,ac,bc\}&\{b,bc\}&\{b,c,bc\}&a\\
\{b,c,bc\}&\{c,ac\}&\{c,ac,bc\}&a\\
\{c,ac,bc\}&\{b,bc\}&\{b,c,ac,bc\}&a\\
\{b,c,bc\}&\{c,ac\}&\{b,c,ac,bc\}&a\\
\{c,ac,bc\}&\{b,ab,bc\}&\{b,c,bc\}&a\\
\{b,c,bc\}&\{a,c,ac\}&\{c,ac,bc\}&\text{none}\\
\{c,ac,bc\}&\{b,c,bc\}&\{b,c,ac,bc\}&a
\end{array}
\label{eq:damp-q5-antipodal-hole-table}
\end{equation}
Here is a direct way to check exhaustiveness.  Each \(H\) is one of the
empty interval and the six intervals of each length \(1,2,3,4\) in
\eqref{eq:damp-q5-middle-six-cycle}.  For each possible length of
\(H_Z\), place one vertex of \(H_Z\setminus H_X\) and one of
\(H_Z\setminus H_Y\), as required by
\eqref{eq:damp-q5-antipodal-hole-noncontainment}.  The rule defining
\(\mathcal N\) then determines the three summands on the left of
\eqref{eq:damp-q5-antipodal-hole-budget}.  Equality leaves exactly the
nine rows above.  Thus the table checks only interval endpoints on the
displayed six-cycle; no endpoint family is being classified.

For \(f\in F\), the section of \(\mathcal A\) with \(f=0\) is contained
in the section with \(f=1\) exactly when, in each of \(X,Y,Z\), every
vertex with \(f=0\) has its \(f\)-mate.  This verifies the last column of
\eqref{eq:damp-q5-antipodal-hole-table} directly.  Since \(\mathcal A\)
is periphery-free, only the eighth row remains.  Taking complements gives
exactly \eqref{eq:damp-q5-pentagonal-four-sections}.

Each of \(X,Y,Z\) has five vertices, so
\[
 |\mathcal A|=1+5+5+5=16.
\]
The section \(X\) shatters \(\{b,c\}\),
\(Y\) shatters \(\{a,c\}\), and \(Z\) shatters \(\{a,b\}\).
The endpoints \(\varnothing,abc\) in all three sections also show that
every pair containing one of \(d,e\) is shattered, and the vertices in
the four \(d,e\)-sections shatter \(\{d,e\}\).  Hence \(\mathcal A\)
shatters all ten coordinate pairs.  Together with the empty support and
the five singletons, these are sixteen shattered supports.  Ampleness and
\(|\mathcal A|=16\) leave no shattered triple.  Thus \(\mathcal A\) is
maximum of VC dimension two.
\end{proof}

\begin{corollary}[The three-label branch is impossible]
\label{cor:damp-q5-no-three-label-branch}
The three-label profile in
Corollary~\ref{cor:damp-q5-three-label-marked-square} cannot occur in a
minimum five-coordinate gap.
\end{corollary}

\begin{proof}
Translate the selected vertex \(p\) to \(\varnothing\), and use the
coordinates \(a,b,c,d,e\) from
\eqref{eq:damp-q5-degree-two-neighbour-split}.  The blocker-face
conclusion of Lemma~\ref{lem:damp-q5-degree-two-blocker-face} says that
the four \(d,e\)-sections of \(A\), with \(F=\{a,b,c\}\), have the form
\eqref{eq:damp-q5-antipodal-four-sections}: the \(00\)-section is the
singleton \(\{abc\}\), while the other three sections all contain
\(\varnothing\) and \(abc\).  The endpoint \(A\) is ample and
periphery-free by
Proposition~\ref{prop:damp-q5-central-relative-cm-profile}.  Lemma
\ref{lem:damp-q5-antipodal-four-section-rigidity} therefore makes \(A\)
maximum of VC dimension two.

Complement duality
\eqref{eq:damp-shattering-complement-duality} now makes
\(B=Q_E\setminus A\) maximum of VC dimension two as well: a support
\(R\) is shattered by \(B\) exactly when \(E\setminus R\) is not
shattered by \(A\), which here is equivalent to
\(|R|\le2\).  Lemma~\ref{lem:damp-relative-peeling-vc-two} supplies an
ample one-concept intermediate family between \(S\) and \(B\), contrary
to the definition of a minimum gap.
\end{proof}

\begin{remark}
Theorem~\ref{thm:damp-q5-circular-carrier-exclusion} is a purely
signed set-system statement on five colors.  By
Corollary~\ref{cor:damp-q5-ten-cycle-endpoint-profile}, it rules out the
ten-cycle branch of a minimum gap without reconstructing either endpoint.
The bowtie-carrier argument in
Theorem~\ref{thm:damp-q5-bowtie-carrier-exclusion} similarly excludes
the six-cycle branch.  Corollary~\ref{cor:damp-q5-four-cycle-reduction}
therefore leaves only the four-cycle branch.  A carrier may contain that
whole cycle, so the interval argument requires a different replacement.
Lemmas~\ref{lem:damp-q5-marked-complementary-carriers} and
\ref{lem:damp-q5-degree-two-blocker-face} initially reduce that
replacement to two signed local profiles.  The distinction is the number
of upper triple cylinders through the selected vertex, not the total
number of isolated pair supports: a globally isolated pair need not have
its complementary cylinder pass through that vertex.  The two-label
profile is now excluded by
Corollary~\ref{cor:damp-q5-no-two-label-branch}.  Its contraction consists
of a seven-vertex core and two optional three-vertex paths.  Any
nonconstant signed expansion has a core-free excess component in one of
those paths; isometry of the duplicated carrier then forces a compatible
peripheral semicube, which remains peripheral after lifting.  Thus every
minimum gap would have all three local upper triples.  Lemma
\ref{lem:damp-q5-antipodal-four-section-rigidity} then identifies its
periphery-free complementary endpoint with the pentagonal maximum
VC-two class, and Corollary~\ref{cor:damp-q5-no-three-label-branch}
excludes it by relative VC-two peeling.  This closes the four-cycle
branch without choosing between the adjacent and opposite positions of
the marked pair carrier.
\end{remark}

\begin{remark}[The exact symbolic replacement for the bounded audit]
\label{rem:damp-q5-relative-ear-target}
The replacement is now entirely symbolic.  The two-sided defect
cylinders give minimum degree two in the intermediate layer.  Convex
carrier intervals exclude cycles of lengths ten, eight, and six, leaving
a square and a degree-two vertex.  The signed blocker at that vertex has
two or three upper labels.  Carrier compatibility excludes the two-label
case.  In the three-label case, the complementary endpoint splits into
four antipodal \(Q_3\)-sections.  The holes of each section are short
intervals of the middle six-cycle
\eqref{eq:damp-q5-middle-six-cycle}; the single ampleness budget
\eqref{eq:damp-q5-antipodal-hole-budget} and periphery-freeness force the
pentagonal maximum VC-two class.  Relative VC-two peeling gives the final
contradiction.

The balanced Artinian graph and the relative Yang-ball formulation remain
useful interpretations of the reduction, but neither is needed as an
additional hypothesis.  In particular, the earlier order-through-twenty-one
files are now regression checks only; they are not part of the proof of the
five-coordinate theorem.
\end{remark}

\begin{theorem}[Relative peeling in five coordinates]
\label{thm:damp-five-coordinate-relative-peeling}
Let \(\varnothing\ne S_0\subseteq B\subseteq Q_E\) be ample, and
suppose that \(B\) uses at most five coordinates.  For every
\(s\in S_0\), the family \(B\) has a corner peeling in which every
vertex of \(B\setminus S_0\) precedes every vertex of \(S_0\), and
\(s\) is last.
\end{theorem}

\begin{proof}
It is enough to prove that whenever \(S_0\subsetneq B\), some
\(p\in B\setminus S_0\) makes \(S_0\cup\{p\}\) ample.  Otherwise choose
a minimum counterexample as in Lemma
\ref{lem:damp-minimal-interval-gap-two-sided}.  Proposition
\ref{prop:damp-q5-central-relative-cm-profile} puts it in the central
five-coordinate profile.  Corollary
\ref{cor:damp-q5-four-cycle-reduction} gives a square in its intermediate
layer, and Lemma~\ref{lem:damp-q5-square-branch-degree-two-vertex}
supplies a degree-two vertex on the square.  Lemmas
\ref{lem:damp-q5-degree-two-signed-split} and
\ref{lem:damp-q5-degree-two-blocker-face} give the signed blocker normal
form at that vertex.  Corollary
\ref{cor:damp-q5-no-two-label-branch} excludes two upper labels, while
Corollary~\ref{cor:damp-q5-no-three-label-branch} excludes three.
These are the only possibilities, so the minimum counterexample does not
exist.

We may consequently adjoin one ample concept repeatedly until \(B\) is
reached.  Reversing this chain peels \(B\setminus S_0\) first.  Applying
the same assertion to \(\{s\}\subseteq S_0\) then completes the peeling
with \(s\) last.
\end{proof}

\begin{proposition}[A six-coordinate failure of relative peeling]
\label{prop:damp-six-coordinate-relative-failure}
There is a maximum class $P\subseteq Q_6$ of VC dimension three
containing a maximum class $O$ of VC dimension two such that $P$
cannot be corner-peeled to $O$ while retaining every vertex of $O$.
Thus the five-coordinate relative-peeling theorem cannot be extended
to six coordinates, even for maximum classes of consecutive ranks.
Equivalently, $P$ admits no acyclic CCMW USO with inward boundary at $O$.
Moreover, every representation map of $O$ extends literally to a
representation map of $P$, but the induced CCMW USO of every such extension
has a directed cycle in $P\setminus O$. Thus even universal literal
extendibility does not imply relative peeling.
\end{proposition}
\begin{proof}
Number the coordinates $0,\ldots,5$ and identify a word with the
integer whose binary digits are its coordinate values. Take
\[
\begin{aligned}
P={}&\{0,1,2,4,5,6,7,8,9,10,11,12,13,14,15,\\
 &16,17,20,21,22,23,25,27,29,31,32,33,40,41,43,\\
 &47,48,49,52,53,54,55,56,57,59,61,63\},\\
O={}&\{0,2,4,6,7,15,16,20,22,23,31,47,48,49,\\
 &52,53,54,55,56,57,61,63\}.
\end{aligned}
\]
Both inclusions $\varnothing\ne O\subsetneq P\subsetneq Q_6$ are
immediate. Direct finite verification gives the following numbers of
shattered supports by cardinality; the strongly shattered counts are
identical:
\[
\begin{array}{c|rrrrrrr|r}
 &0&1&2&3&4&5&6&\text{total}\\\hline
 P&1&6&15&20&0&0&0&42\\
 O&1&6&15&0&0&0&0&22.
\end{array}
\]
Thus $P$ and $O$ are maximum and ample. The complete corner list of
$P$, with the incident coordinate sets, is
\[
\begin{array}{c|c}
\text{corner}&\text{incident directions}\\\hline
2&\{1,2,3\}\\
47&\{2,4,5\}\\
54&\{0,1,5\}\\
56&\{0,3,4\}.
\end{array}
\]
Each of these four vertices belongs to $O$. Consequently no corner
can be deleted while retaining $O$, so the relative peeling cannot start.
The class $O$ is corner-peelable by
Corollary~\ref{cor:damp-six-coordinate-peeling}; the USO conclusion
follows from Proposition~\ref{prop:damp-relative-uso-extension}.

The relative shattered supports are precisely the twenty triples, an
antichain. The antichain relative-support theorem of
\cite[Theorem ``Antichain relative supports'']{BousquetPartitionability26}
therefore extends every prescribed representation map of $O$ literally
to $P$. This is an application of that theorem, not a new extension proof.
Every such extension has inward boundary at $O$: if a crossing edge of
direction $i$ were directed out of $v\in O$, its unchanged label $r(v)$
would contain $i$. Condition (C1) for the prescribed USO on $O$ would then
supply the same $i$-neighbour within $O$, contradicting that the edge crosses
the boundary. If its induced orientation on $P\setminus O$ were acyclic,
Proposition~\ref{prop:damp-relative-uso-extension} would give the impossible
relative peeling. Hence a directed cycle occurs outside $O$ in every
literal extension, even when the prescribed map on $O$ is acyclic.

The finite checks use only the displayed sets, not the SAT solver
that found them. For each triple, a full cube in $P$ certifies strong
shattering, and for each four-set a missing trace certifies nonshattering.
The corresponding witnesses for $O$ use pairs and triples. The direct
certificate also supplies one missing vertex of the incident cube for
every noncorner of $P$. All witnesses are retained in the companion project file
\path{../ample-corners/evidence/relative_pair_q6/direct_certificate.json}
and checked by
\path{../ample-corners/scripts/verify_relative_pair_q6.py}.
\end{proof}

The corner transfer in
Theorem~\ref{thm:damp-two-coordinate-low-vc-core-corner} assumed equal
opposite sections. The following extension permits nested opposite
sections and also records how the exclusive components are assigned.

\begin{theorem}[Rectangular boundaries give forbidden pc-minors]
\label{thm:damp-rectangular-boundary-forbidden}
For $i=1,2$, let $\varnothing\ne K_i\subsetneq A_i$ be ample
classes on disjoint coordinate sets $E_i$, with every coordinate active
in $A_i$. Assume the following:
\begin{enumerate}
 \item $A_i$ and $K_i$ admit complete corner peelings;
 \item every corner $v$ of $A_i$ belongs to $K_i$, but its unique
       maximal cube $M_i(v)$ is not contained in $K_i$;
 \item for each elementary restriction or contraction $\mu$ on $E_i$
       with $\mu K_i\ne\varnothing$, the class $\mu A_i$ peels
       relatively to $\mu K_i$.
\end{enumerate}
Then
\[
 H=(A_1\times K_2)\cup(K_1\times A_2)
\]
is a cornerless ample forbidden pc-minor for $\Damp$. Its order is
\[
 |H|=|A_1||K_2|+|K_1||A_2|-|K_1||K_2|.
\]
\end{theorem}
\begin{proof}
We use the rectangular-boundary shattering formula
\cite[Theorem \texttt{relative-rectangular-boundary-tensor-factorization}]
{BousquetPartitionability26}:
\[
 \operatorname{Sh}(H)=
 (\operatorname{Sh}(A_1)\times\operatorname{Sh}(K_2))
 \cup
 (\operatorname{Sh}(K_1)\times\operatorname{Sh}(A_2)),
\]
where a pair of supports is identified with their disjoint union.
For completeness, if a support in each block is shattered by the larger
class but by neither smaller class, choose one missing trace in each
smaller class. Their combined trace is absent from $H$. In all other
cases the appropriate product slab supplies the witnessing product cube.
Thus the formula is exact and all its supports are strongly shattered,
so $H$ is ample. Inclusion--exclusion gives its order.

Consider a possible corner $(x,y)$ of $H$. If $x\notin K_1$, then
$y\in K_2$ and the section at $y$ is $A_1$. A corner of the whole
class is a corner in this section, so $x$ would be a corner of $A_1$
outside $K_1$, contrary to (2). The symmetric argument handles
$y\notin K_2$. It remains that $x\in K_1$ and $y\in K_2$.
The two coordinate sections are now $A_1,A_2$, so $x,y$ must both
be corners of these classes. All their incident directions occur in
$H$, and the required incident cube is $M_1(x)\times M_2(y)$.
By (2) choose $x'\in M_1(x)\setminus K_1$ and
$y'\in M_2(y)\setminus K_2$. The vertex $(x',y')$ is absent from
$H$, a contradiction. Hence $H$ is cornerless and is not in $\Damp$.

We show directly why a relative peeling in either factor makes the
boundary peelable. Suppose $A_1$ peels to $K_1$. In that relative
order, delete the slice $\{x\}\times K_2$ for each
$x\in A_1\setminus K_1$, using a fixed peeling of $K_2$ inside the
slice. At an intermediate deletion, the incident cube is the product
of the current corner cube at $x$ in the remaining $A_1$ and the
current corner cube in the slice. Later outside slices still contain
all of $K_2$, while slices over $K_1$ contain $A_2\supseteq K_2$,
so this product cube is present. Thus every deletion is valid.
The remainder $K_1\times A_2$ peels by the product of their peelings.

An elementary operation $\mu$ on $E_1$ sends $H$ exactly to
\[
 (\mu A_1\times K_2)\cup(\mu K_1\times A_2).
\]
If $\mu K_1$ is empty, this is the product of two peelable classes,
since peelability is closed under pc-minors. Otherwise (3) and the
preceding slice argument give a peeling; $\mu K_1$ is peelable by
(1) and minor closure. The same argument applies to $E_2$.
Every coordinate is active in $H$, because the smaller classes are
nonempty, so these are proper elementary minors. All proper pc-minors
therefore belong to $\Damp$, proving forbiddenness.
\end{proof}

\begin{corollary}[An explicit rectangular-boundary obstruction]
\label{cor:damp-rectangular-boundary-1364}
Let $O\subset P\subset Q_6$ be the displayed $22/42$ pair in
Proposition~\ref{prop:damp-six-coordinate-relative-failure}.
On two disjoint copies of its six coordinates, the class
\[
 H=(P\times O)\cup(O\times P)
\]
is an ample forbidden pc-minor of order $1364$, dimension $12$,
and VC dimension $5$.
\end{corollary}
\begin{proof}
Both factors peel by the prescribed-endpoint theorem through seven
coordinates. Every elementary image of the nested pair uses at most
five coordinates, so Theorem~\ref{thm:damp-five-coordinate-relative-peeling} supplies condition (3).
The verified corner list of $P$ lies in $O$. Each listed corner has
three incident directions and its full three-cube, whereas $O$ has
VC dimension two, so condition (2) holds. The theorem applies.
Its order is $2\cdot42\cdot22-22^2=1364$. Both coordinate blocks
are active. The shattering formula has maximal support sizes $3+2$
and $2+3$, giving VC dimension five.
\end{proof}

More generally, choose independently, in each six-coordinate block,
a maximum rank-three class $P_i$ and a maximum rank-two subclass
$O_i$ containing all corners of $P_i$. Then
$(P_1\times O_2)\cup(O_1\times P_2)$ is forbidden of the same order,
dimension, and VC dimension. Indeed, every corner of a maximum
rank-three class has a maximal three-cube, so no such cube lies in
$O_i$; the remaining hypotheses follow from the same small-dimension
theorems. Thus this finite input rule needs no separate peeling or
minor tests. The displayed pair proves that its input set is nonempty.

The construction uses no rooted negative input. It converts two
explicit relative failures, with their proper-image peeling conditions,
into an ordinary forbidden pc-minor. The rectangular-boundary operation
and its ampleness are reused from the cited result; the cornerlessness
and minor-minimality criterion above are the added conclusions. This
is a sound generative subfamily, not an exhaustive description of all
forbidden pc-minors. Its $1364$-vertex instance does not improve the
$64$--$267$ minimum-order interval. The independent verifier
\texttt{verify\_rectangular\_obstruction.py} checks its full shattered
family, the absence of every corner, and explicit corner orders for
all $36$ elementary pc-minors.

\begin{proposition}[Intrinsic recovery of a rectangular boundary]
\label{prop:damp-rectangular-boundary-recovery}
Fix a partition of the coordinates of a nonempty class $H$ into two
nonempty blocks $E_1,E_2$. Let $A_1,A_2$ be its projections, and put
\[
 K_1=\{x\in A_1:\{x\}\times A_2\subseteq H\},\qquad
 K_2=\{y\in A_2:A_1\times\{y\}\subseteq H\}.
\]
Suppose $\varnothing\ne K_i\subsetneq A_i$ for $i=1,2$.
Then $H$ has a rectangular-boundary presentation on this split if
and only if
\[
 |H|=|A_1||K_2|+|K_1||A_2|-|K_1||K_2|.
\]
When it exists, this presentation is unique and equals
$(A_1\times K_2)\cup(K_1\times A_2)$.
If $H$ is ample, all four recovered classes are ample. If $H$ is
forbidden for $\Damp$, all four belong to $\Damp$.

Whenever all four recovered classes belong to $\Damp$ and the
presentation exists, $H$ is cornerless if and only if every corner
of each $A_i$ belongs to $K_i$ and its unique maximal cube is not
contained in $K_i$.
\end{proposition}
\begin{proof}
By definition both slabs are contained in $H$. Their union has the
cardinality in the display, so equality of cardinalities is equivalent
to equality of classes. Conversely, in any presentation with nonempty
proper smaller factors, the projections are the larger factors. A
row indexed outside the first smaller factor is exactly the second
smaller factor and hence is not full. Thus the full rows recover the
first smaller factor; full columns recover the second. This proves
uniqueness. The larger factors are projections of $H$. Fixing a word
in $A_2\setminus K_2$ gives the section $K_1$, and symmetrically for
$K_2$. Ampleness follows from minor closure. For a forbidden $H$
these projections and sections are proper pc-minors, so they belong
to $\Damp$.

For necessity of the corner conditions, suppose a corner $x$ of
$A_1$ lies outside $K_1$. Choose a corner $y$ of $K_2$, which exists
because $K_2$ is nonempty and peelable. At $(x,y)$ the incident cube
is the product of the corner cubes in $A_1$ and $K_2$ and lies wholly
in the first slab. This is a corner of $H$, a contradiction. Thus
all corners of $A_1$ lie in $K_1$, and symmetrically for $A_2$.
If a corner cube $M_1(x)$ were wholly in $K_1$, choose any corner
$y$ of $A_2$. It belongs to $K_2$ by what was just proved.
The full incident cube $M_1(x)\times M_2(y)$ lies in the second
slab, again giving a corner. This proves the necessary cube condition
in both blocks. Conversely, the cornerlessness argument in
Theorem~\ref{thm:damp-rectangular-boundary-forbidden} uses exactly these
conditions: outside the two smaller factors a corner cannot project
to an outer corner, and inside them its product cube has a missing
outside--outside vertex. The proper-image hypotheses are not needed
for that argument.
\end{proof}

This is marked recovery and an exact cornerless criterion, not a
converse to the entire forbiddenness theorem: necessity of its relative
peeling conditions on proper images has not been proved. For two
six-coordinate maximum rank-three/rank-two input pairs those conditions
hold automatically, so the stated input rule is exhaustive within
that cornerless rectangular subfamily. There is no assertion that
every cornerless forbidden class has such a split.

For the $1364$-vertex class above, all $2047$ unordered nontrivial
coordinate bipartitions were checked by the recovery criterion.
Exactly one works: the original two six-coordinate blocks, recovering
orders $42,22,42,22$. The script
\texttt{recognize\_rectangular\_boundaries.py} also agrees with direct
enumeration of all possible nonempty proper lower factors on every
nonempty three-coordinate class. The recognition certificate is
retained with the explicit obstruction.

The same construction also supplies a direct small-order target without
maximum assumptions. Any nonempty proper ample pair $K\subset A\subseteq Q_6$
satisfying the corner condition above satisfies the peeling hypotheses
of the construction theorem, by the seven-coordinate endpoint and
five-coordinate relative theorems. Thus its self-pair gives a forbidden
object of order $k(2a-k)$, where $a=|A|$ and $k=|K|$.
Moreover mixing two different admissible pairs cannot improve the minimum
over all self-pairs. If their self-pair orders are
$m_i=k_i(2a_i-k_i)$ and their mixed order is
$m_{12}=a_1k_2+k_1a_2-k_1k_2$, then
\[
 m_{12}^2-m_1m_2=(a_1k_2-a_2k_1)^2\ge0.
\]
All orders are positive, so $m_{12}\ge\min(m_1,m_2)$.
Consequently the entire six-coordinate-input rectangular construction
has an output below $267$ if and only if one admissible pair satisfies
$k(2a-k)\le266$. This is a smaller-obstruction search criterion;
existence of such a pair is not asserted.

\begin{proposition}[Small rectangular inputs have three size profiles]
\label{prop:damp-rectangular-small-input-profiles}
Suppose a nonempty proper ample pair $K\subset A\subseteq Q_6$
retains every corner of $A$, and $k(2a-k)\le266$, where
$a=|A|$ and $k=|K|$. Then $K$ uses all six coordinates and
\[
 (k,a)\in
 \{(7,a):16\le a\le22\}\cup
 \{(8,a):17\le a\le20\}\cup
 \{(9,a):18\le a\le19\}.
\]
In these three cases $K$ has VC dimension at most two and its
shattered family consists of the empty set, all six singletons,
and exactly $k-7$ pairs. The conclusion does not require the
additional corner-cube-escape condition of the construction theorem.
\end{proposition}
\begin{proof}
First, $A$ must use all six coordinates. Otherwise the five-coordinate
relative theorem supplies a first corner outside $K$, a contradiction.
We claim that each side of every active coordinate of $A$ contains
a corner of $A$. Write its two nonempty sections as $A_0,A_1$ and
put $D=A_0\cap A_1$. Connectedness of $A$ implies $D\ne\varnothing$,
and $D$ is ample as a coordinate reduction. If $A_1\setminus D$ is
nonempty, the five-coordinate relative theorem gives a corner of
$A_1$ outside $D$. It has no neighbour across the split coordinate,
so its incident cube is unchanged and it is a corner of $A$.
If $A_1=D$, take any corner of this nonempty five-coordinate class.
Its incident cube doubles across the split coordinate inside $A$,
so its copy in $A_1$ is a corner of $A$. The same argument applies
to $A_0$, proving the claim. Since $K$ retains all corners, it meets
both sides of all six coordinates. Consequently $k\ge7$ by the
Sandwich equality and the six shattered singletons.

The relative peeling theorem for at most eight added concepts
proved in the argument of Corollary~\ref{cor:damp-small-component-wing-inverse} implies $a-k\ge9$: otherwise this nested
ample pair would admit a first relative deletion. (It is also the
special case of the small graph-core relative theorem where the
whole difference has at most eight vertices.) Therefore
\[
 k(2a-k)\ge k(k+18).
\]
The size bound forces $k\le9$. Substituting $k=7,8,9$ in
$a\ge k+9$ and $a\le\lfloor(266+k^2)/(2k)\rfloor$ gives the stated
ranges. Finally, any shattered triple would add its three pairs
and itself to the seven already shattered empty/singleton supports,
forcing $|K|\ge11$. Thus no triple is shattered, and the Sandwich
equality leaves exactly $k-7$ pair supports.
\end{proof}

\begin{theorem}[Lower bounds for rectangular forbidden classes]
\label{thm:damp-rectangular-input-order-exclusion}
Every class generated by the six-coordinate rectangular-input rule
has at least $272$ vertices. More generally, let a cornerless ample
forbidden pc-minor admit a presentation
\[
 H=(A_1\times K_2)\cup(K_1\times A_2),\qquad
 \varnothing\ne K_i\subsetneq A_i,
\]
on disjoint active coordinate blocks. Then it uses at least twelve
coordinates; if it uses exactly twelve, it has at least $272$ vertices.
The dimension bound is attained by the $1364$-vertex example.
In particular the known twelve-coordinate obstruction of order $267$
has no such rectangular presentation.
\end{theorem}
\begin{proof}[Computer-assisted proof]
We first exclude every strict corner-retaining pair $K\subset A\subseteq Q_6$
whose self-boundary order $k(2a-k)$ is at most $271$.
The proof of Proposition~\ref{prop:damp-rectangular-small-input-profiles}
continues to give exactly the same three profiles when $266$ is
replaced by $271$: $k(k+18)\le271$ forces $7\le k\le9$, and the
upper bounds for $a$ remain $22,20,19$, respectively.

The pre-existing rooted ample inventory on six coordinates was audited
independently through order nine. Starting with $\{0\}$, for every
representative at each order through eight, every possible one-vertex
augmentation was tested by direct shattered-set counting. The resulting
coordinate-permutation orbits coincide exactly with the next inventory
level. Every nonempty ample class on six coordinates peels to any
prescribed vertex, so reversing a peeling ending at zero proves that
this augmentation induction covers every rooted class. Each listed
representative was also checked directly for ampleness. Restricting
to six active coordinates gives $48,162,347$ rooted representatives
at orders $7,8,9$. Taking every translation by a member and every
coordinate permutation removes the root marking and leaves the
following complete cover:
\[
\begin{array}{c|c|c}
 |K|&\text{unrooted lower classes}&\text{allowed }|A|\\\hline
 7&11&16\text{--}22\\
 8&28&17\text{--}20\\
 9&47&18\text{--}19.
\end{array}
\]

For each of these $86$ fixed lower classes, the formula imposes:
$K\subseteq A$; ampleness of $A$ by total asymmetry; the displayed
order bounds; absence of any corner in $A\setminus K$; and an
$A$-neighbour outside $K$ at every corner of $K$. The last condition
is exactly corner-cube escape. Indeed, if a corner of $A$ has its
incident cube in $K$, its neighbours in $A$ and $K$ coincide, making
it a corner of $K$ without an outside neighbour. Conversely such a
corner of $K$ remains a corner of $A$ with the same cube. Thus no
additional restriction is being imposed on the admissible pairs.

All $86$ formulas are unsatisfiable, each with a separately verified
DRAT proof. There are no unresolved cases or unchecked solver answers
in this cover. The exact formulas, proofs, checker logs, source snapshots,
and input inventory are retained under
\path{evidence/rectangular_fixed_lower/} in the ample-corners project.
The program \texttt{prepare\_rectangular\_lower\_cases.py} audits the
inventory; \texttt{certify\_small\_rectangular\_inputs.py} produces and
checks the proofs. A separate coverage audit reproduces all $86$ formula
hashes and checks all $258$ retained artifact hashes. The local escape
equivalence also agrees on every one of the $6050$ nonempty proper
nested pairs on three coordinates. These finite checks supplement the
symbolic coverage and encoding arguments above.

Consequently every admissible self-boundary has order at least $272$.
The mixed-order inequality proved after
Proposition~\ref{prop:damp-rectangular-boundary-recovery} gives the same
lower bound when the two inputs differ.

Finally, for an arbitrary cornerless forbidden $H$ with such a
presentation, Proposition~\ref{prop:damp-rectangular-boundary-recovery}
makes the four input classes peelable and forces the strict
corner-retention conditions. Neither upper class can use at most
five coordinates: relative peeling to its nonempty proper lower
class would supply a corner outside that lower class. Each block
therefore has at least six active coordinates. If their total is
twelve, both have exactly six and the proved order bound applies.
The explicit $1364$ example proves sharpness of the dimension bound.
\end{proof}

This closes the search for a sub-$267$ obstruction within this
rectangular family. It does not raise the global lower bound $64$,
nor does it show that order $272$ is attained in the family. As a
separate direct check, the rectangular recognizer finds no presentation
of the known $267$-vertex class among all $2047$ coordinate splits.
The full generative characterization must therefore include other
structures even at twelve coordinates.

\begin{corollary}[A size criterion for relative peeling on six coordinates]
\label{cor:damp-relative-six-weighted-size}
Let $\varnothing\ne K\subsetneq A\subseteq Q_6$ be ample, with
$k=|K|$ and $a=|A|$. If
\[
 k(2a-k)\le271,
\]
then $A$ peels relatively to $K$. No corner-retention, maximum-class,
or corner-cube-escape hypothesis is required.
\end{corollary}
\begin{proof}
Suppose no such relative peeling exists. Delete corners outside $K$
as long as possible. The process stops at an ample class $B$ properly
containing $K$, with every corner of $B$ in $K$; reaching $K$ would
itself give a relative peeling of $A$.

Whenever a corner $v$ of $B$ has its entire maximal cube contained
in $K$, delete $v$ from both classes. It is also a corner of $K$:
all neighbours of $v$ in $B$ lie in that cube and hence in $K$, so
the incident cubes in the two classes coincide. Both endpoints
therefore remain ample. This is the equal-support case of
common-corner cancellation
\cite[Lemma \texttt{relative-common-corner-cancellation}]
{BousquetPartitionability26}; the relative face set is unchanged.
There is also a direct check that the deletion remains stalled.
Every vertex outside $K$ is nonadjacent to $v$, so its incident
directions do not change on deleting $v$. A missing vertex of its
incident cube cannot become present after a deletion. It therefore
remains a noncorner. The difference of the two classes is unchanged
and nonempty. The smaller class cannot become empty: its last vertex
would have an incident cube consisting only of itself in the larger
class, contradicting connectedness of that larger ample class.

After finitely many such simultaneous deletions, we obtain
$\varnothing\ne K'\subsetneq B'\subseteq Q_6$ such that every
corner of $B'$ lies in $K'$ and every such corner cube escapes $K'$.
The six-coordinate rectangular construction applies to this pair.
Theorem~\ref{thm:damp-rectangular-input-order-exclusion} therefore gives
\[
 |K'|\bigl(2|B'|-|K'|\bigr)\ge272.
\]
But $|B'|\le a$ and $|K'|\le k$, and $y(2x-y)$ is increasing in
each variable on $0<y<x$. Thus the left side is at most
$k(2a-k)\le271$, a contradiction.
\end{proof}

The added conclusion here is a relative-peeling criterion for arbitrary
nested ample pairs of the indicated size. The cancellation operation is
reused from the cited result; it is essential that the two corner cubes
coincide. Deleting an arbitrary common corner with different incident
cubes need not preserve the relative face set. This corollary does not
assert that every pair outside the displayed size range fails to peel.

\begin{lemma}[Nested opposite sections and their cubes]
\label{lem:damp-nested-diagonal-cube-cover}
Let an ample class $X\subseteq Q_n\times Q_2$ have sections
$A,P,Q,C$ at $00,10,01,11$, with $\varnothing\ne C\subseteq A$.
Put $R=C\cap P\cap Q$. Then $R$ is ample, every cube of $C$ is
contained in $P$ or in $Q$, and every connected component of
$C\setminus R$ lies wholly in $P\setminus Q$ or in $Q\setminus P$.
Every corner of $C$ outside $R$ lifts to a corner of $X$.

If $C$ has a corner and $X$ is cornerless, then
$\varnothing\ne R\subsetneq C$, every corner of $C$ lies in $R$,
and $C$ cannot peel relatively to $R$. In particular $X$ has a corner
whenever $C$ has VC dimension at most two or uses at most five active
coordinates.
\end{lemma}
\begin{proof}
Let $K$ be any cube in $C$. Restrict $X$ to the face generated by $K$
and the two last coordinates. Its $00$ and $11$ sections both contain
all of $K$, because $C\subseteq A$. If neither middle section contained
$K$, the nonempty cube complement of this face restriction would be
the union of two nonempty parts at $10$ and $01$, with no edge between
them. This contradicts connectedness of a nonempty ample complement.
Thus $K\subseteq P$ or $K\subseteq Q$. Applying this first to vertices
and then to edges shows $C\subseteq P\cup Q$ and that no edge of
$C\setminus R$ joins its two exclusive sides. This proves the component
assertion. The set $R=A\cap P\cap Q\cap C$ is the double coordinate
reduction of $X$, so is ample.

If $v$ is a corner of $C$ outside $R$, it belongs to exactly one middle
section. Its unique maximal cube $K$ is contained in that section by
the preceding cube-cover argument. Doubling $K$ across the corresponding
last-coordinate edge gives its entire incident cube in $X$, so its
$11$ copy is a corner. If $R=C$, every corner cube of $C$ occurs in
all four sections and extends across both last coordinates, again
giving a corner of $X$.

Consequently, if $X$ is cornerless and $C$ has a corner, then $R$ is
nonempty and proper and contains every corner of $C$. No first relative
deletion is possible. Finally, relative peeling holds for every nested
ample pair with upper VC dimension at most two, and for every pair
using at most five coordinates (Theorems~\ref{thm:damp-relative-small-supports}
and~\ref{thm:damp-five-coordinate-relative-peeling}); these classes
also have corners. These results exclude the required stalled interval
in either stated case.
\end{proof}

\begin{proposition}[Exact compatibility over a full base]
\label{prop:damp-full-base-shadow-defect}
Let $P,Q,C\subseteq Q_n$ be arbitrary, including possibly empty
families, and let $X\subseteq Q_n\times Q_2$ have sections
$Q_n,P,Q,C$ at $00,10,01,11$. For any family $H$, put
$\delta(H)=|\operatorname{Sh}(H)|-|H|$, and define
\[
\begin{split}
 E_P&=\operatorname{Sh}(P\cup C)\setminus
       (\operatorname{Sh}(P)\cup\operatorname{Sh}(C)),\\
 E_Q&=\operatorname{Sh}(Q\cup C)\setminus
       (\operatorname{Sh}(Q)\cup\operatorname{Sh}(C)),\\
 E_C&=\operatorname{Sh}(C)\setminus
       (\operatorname{Sh}(P)\cup\operatorname{Sh}(Q)).
\end{split}
\]
Then the following identity is a sum of nonnegative terms:
\begin{equation}
 \delta(X)=\delta(P)+\delta(Q)+\delta(C)
                  +|E_P|+|E_Q|+|E_C|.
 \label{eq:damp-full-base-shadow-defect}
\end{equation}
Consequently $X$ is ample if and only if $P,Q,C$ are ample and
$E_P=E_Q=E_C=\varnothing$.

Suppose these conditions hold. For a corner $v$ of an ample family
$H$, write $K_H(v)$ for its unique maximal cube, viewed as a set of
words in $Q_n$, and put $R=P\cap Q\cap C$. Then $X$ is cornerless
if and only if all the following conditions hold:
\begin{enumerate}
 \item $P\cup Q=Q_n$, with $P\ne Q_n$ and $Q\ne Q_n$;
 \item every corner $v$ of $P$ lies in $C$ and satisfies
       $K_P(v)\not\subseteq Q\cap C$;
 \item every corner $v$ of $Q$ lies in $C$ and satisfies
       $K_Q(v)\not\subseteq P\cap C$;
 \item every corner $v$ of $C$ lies in $R$ and satisfies
       $K_C(v)\not\subseteq R$.
\end{enumerate}
\end{proposition}
\begin{proof}
Write $\mathcal A=\operatorname{Sh}(P)$,
$\mathcal B=\operatorname{Sh}(Q)$, and
$\mathcal D=\operatorname{Sh}(C)$. Splitting shattered supports
according to their intersection with the two last coordinates gives
\[
 |\operatorname{Sh}(X)|=2^n+
 |\operatorname{Sh}(P\cup C)|+
 |\operatorname{Sh}(Q\cup C)|+
 |\mathcal A\cap\mathcal B\cap\mathcal D|.
\]
Indeed, a support using just the first of the two last coordinates
requires every old trace in $P\cup C$; its other section is full.
The second coordinate is symmetric. A support using both requires
every old trace separately in $P,Q,C$.
For any three finite sets one has
\[
 |\mathcal A\cup\mathcal D|+|\mathcal B\cup\mathcal D|
 +|\mathcal A\cap\mathcal B\cap\mathcal D|
 =|\mathcal A|+|\mathcal B|+|\mathcal D|
   +|\mathcal D\setminus(\mathcal A\cup\mathcal B)|.
\]
Substitute this equality and $|X|=2^n+|P|+|Q|+|C|$ to obtain
\eqref{eq:damp-full-base-shadow-defect}. The Sandwich inequality makes
each $\delta$ nonnegative, with equality exactly for ample families.
This proves the ampleness criterion, including empty sections.

For corners, inspect the cube spanned by all incident directions.
At a base vertex, the old directions span the full $Q_n$. If neither
middle neighbour is present, this vertex is a corner. If exactly one
is present, it is a corner precisely when that middle section is
$Q_n$. If both are present, it is a corner precisely when all three
other sections are full. Moreover a cornerless class cannot have
$P=Q_n$ or $Q=Q_n$: it would contain a full codimension-one cube,
and avoiding a corner at each vertex of that cube forces every
opposite vertex to be present, giving the full cube and a contradiction.
These observations give exactly condition (1).

At $(v,10)$ the base neighbour is always present. A corner there
requires $v$ to be a corner of $P$. If $v\notin C$, its incident cube
is automatically complete; if $v\in C$, that cube is complete exactly
when $K_P(v)\subseteq Q\cap C$. This gives (2), and interchanging
the last coordinates gives (3).

At $(v,11)$ a corner requires $v$ to be a corner of $C$. Lemma~\ref{lem:damp-nested-diagonal-cube-cover} shows that a corner outside $R$ lifts
to a corner of $X$. For a corner in $R$, both middle neighbours are
present, and its incident cube is complete exactly when
$K_C(v)\subseteq P\cap Q$; the base is full. Since $K_C(v)\subseteq C$,
this is equivalent to containment in $R$, proving (4).
Every vertex lies in one of these four sections, so the conditions
are also sufficient.
\end{proof}

This is an exact reduction of the full-base search to three families
on two fewer coordinates, with no maximum-class assumption. It
specializes the general counting of shattered supports by section
indices: the full base makes every remaining defect nonnegative
and gives the explicit compatibility conditions above. It does not
classify the triples satisfying those conditions. In particular it
is not an exhaustive generator for all forbidden pc-minors. The
identity and the local corner formulas are independently checked
on all $4096$ labelled triples of subsets of $Q_2$, including the
$819$ ample full-base families, by
\texttt{verify\_full\_base\_defect.py}.

\begin{corollary}[A single edge-fiber normalization]
\label{cor:damp-full-base-edge-normalization}
Let a cornerless ample class have sections $Q_n,P,Q,C$.
Then both $P\setminus(Q\cup C)$ and $Q\setminus(P\cup C)$ are
nonempty. In particular, translating only the old coordinates puts
its fiber over zero in the form $\{00,10\}$. This one fiber type
suffices to cover all cornerless full-base classes up to translation.
\end{corollary}
\begin{proof}
The exact corner criterion makes $P,Q$ proper and $P\cup Q=Q_n$.
Also $C\ne Q_n$: otherwise Lemma~\ref{lem:damp-nested-diagonal-cube-cover},
applied to the full cube $C$, would put $Q_n$ in $P$ or $Q$.
The ampleness criterion gives
\[
 \operatorname{Sh}(Q\cup C)
 =\operatorname{Sh}(Q)\cup\operatorname{Sh}(C).
\]
If $Q\cup C=Q_n$, the full coordinate support would be shattered
by $Q$ or $C$, forcing that family to be $Q_n$, a contradiction.
Hence some $x$ lies outside $Q\cup C$. Since $P\cup Q=Q_n$, it lies
in $P$, so its four-section fiber is exactly $\{00,10\}$.
The symmetric argument proves $Q\setminus(P\cup C)\ne\varnothing$.
Translation by $x$ preserves the full base and gives the normalization.
\end{proof}

Thus the earlier three-case fiber cover can be replaced by a single
case. This does not assert that the other fiber types cannot occur at
other old vertices. It chooses one guaranteed edge fiber, without
assuming invariance of the family under translations or permutations.

\begin{proposition}[A certified degree restriction at an exclusive edge fiber]
\label{prop:damp-full-base-exclusive-degree}
Let $X\subseteq Q_6\times Q_2$ be cornerless and ample, with sections
$Q_6,P,Q,C$ at $00,10,01,11$. Every vertex of
$P\setminus(Q\cup C)$ has at least four neighbors in $P$.
Symmetrically, every vertex of $Q\setminus(P\cup C)$ has at least
four neighbors in $Q$.
\end{proposition}
\begin{proof}
Choose $x\in P\setminus(Q\cup C)$ and translate the six old
coordinates so that $x=0$. The exact corner criterion implies that
every corner of $P$ belongs to $C$. Thus $0$ is not a corner of $P$;
in particular its degree in $P$ is at least two, since a vertex of
degree zero or one belongs to a unique maximal cube.
Permuting the six old coordinates preserves the full base and the
fiber over zero. If the degree is $k$, we may therefore require
that its neighbors are exactly the first $k$ unit vectors.

For $k=2,3$, the exact full-base shadow encoding with
$0\in P\setminus(Q\cup C)$ and these six membership literals is
unsatisfiable. Each formula has $9629$ variables and $81747$ clauses.
The formulas were reconstructed byte-for-byte, and their DRAT
refutations independently checked locally after the remote checks.
The formulas, compressed proofs, hashes and replay records are in
\texttt{papers/ample-corners/evidence/q8\_full\_base\_remote};
\texttt{verify\_q8\_full\_base\_remote.py} reproduces the checks.
This rules out degrees two and three. Exchanging the two new
coordinates proves the symmetric statement.
\end{proof}

The degree-four, degree-five and degree-six cases remain unresolved:
the remote solver reached its $600$-second limit in each case, as did
two unsplit encodings. The proposition applies only to the indicated
exclusive fibers. It neither excludes all six-cube classes nor changes
the interval $[64,267]$ for the smallest ample obstruction.

\begin{proposition}[Triple shattering and a small overlap force a corner]
\label{prop:damp-maximum-relative-pair-no-completion}
Suppose an ample class $X\subseteq Q_n\times Q_2$ has sections
$Q_n,P,Q,C$ at $00,10,01,11$, with $C\ne\varnothing$.
If $P$ shatters every three-coordinate support and
$\operatorname{VCdim}(P\cap C)\le2$, then $X$ has a corner.
In particular this excludes every maximum rank-three side with a
maximum rank-two overlap, without requiring either maximum hypothesis
in the general statement.
\end{proposition}
\begin{proof}
Ampleness of the face with sections $P,C$ gives
$\operatorname{Sh}(P)\cap\operatorname{Sh}(C)
=\operatorname{Sh}(P\cap C)$: a support shattered in both sections,
with the separating coordinate adjoined, is strongly shattered in the
face. Since every triple belongs to $\operatorname{Sh}(P)$ and no
triple belongs to the right side, $C$ has VC dimension at most two.
Lemma~\ref{lem:damp-nested-diagonal-cube-cover} now supplies a corner of $X$.
\end{proof}

\begin{corollary}[A connected difference excludes a maximum opposite section]
\label{cor:damp-connected-maximum-opposite-exclusion}
Let a cornerless ample class have sections $Q_n,P,Q,C$, and put
$R=C\cap P\cap Q$. Suppose $C$ is maximum of VC dimension $d\ge1$
and $R$ is maximum of VC dimension $d-1$. Then the induced graph
on $C\setminus R$ is disconnected.
\end{corollary}
\begin{proof}
If this graph were connected, the component assertion of
Lemma~\ref{lem:damp-nested-diagonal-cube-cover} would place it entirely
in one exclusive side. After interchanging the middle coordinates,
$C\subseteq P$ and $C\cap Q=R$. The section-shadow identity for the
face with sections $Q,C$ gives
$\operatorname{Sh}(Q)\cap\operatorname{Sh}(C)=\operatorname{Sh}(R)$.
All $d$-sets lie in $\operatorname{Sh}(C)$ and none in
$\operatorname{Sh}(R)$, so $\operatorname{VCdim}(Q)\le d-1$.
Since $Q$ contains the maximum class $R$, the Sauer bound forces $Q=R$.

Cornerlessness at the full-base vertices forces $P\cup Q=Q_n$:
a base vertex with neither outside neighbour is a corner. As
$Q=R\subseteq C\subseteq P$, this gives $P=Q_n$.
The class then contains a full $(n+1)$-cube. Every one-coordinate
fiber above this cube must be nontrivial to avoid a corner at its
base, forcing the full $(n+2)$-cube, which has corners. This is the
required contradiction.
\end{proof}

The $42/22$ pair above is excluded in both marked positions. The
proposition excludes placing its larger class in a middle section
with the smaller class as its overlap with $C$. The corollary also
excludes placing the larger class at $11$ and the smaller class at
$R$: their $20$-vertex difference is connected. A spanning tree is
included in the direct finite certificate for the pair. These statements
concern the indicated full-base completions, not every use of the pair
in a larger construction.

\begin{remark}[Why the five-coordinate proof is conceptual]
\label{rem:damp-q5-extendable-shellability}
Proposition~\ref{prop:damp-q5-central-relative-cm-profile} still explains
why ranks two and three are the only relevant relative \(h\)-supports.
The proof does not classify the endpoint families in that order range.
Instead it follows one shortest cycle in the intermediate layer, uses
signed carrier convexity to eliminate the three longer cycle lengths,
and reduces the remaining square to one degree-two vertex.  The
two-label profile is controlled by two three-vertex paths; the
three-label profile is controlled by three intervals in the fixed
six-cycle \eqref{eq:damp-q5-middle-six-cycle}.  Thus both finite-looking
steps are consequences of constant-size structural lemmas, and the
order-through-twenty-one verification is only an independent regression
check.
\end{remark}
